%% -*- coding: utf-8 -*-
%% Time-stamp: <Chen Wang: 2024-04-02 10:48:01>

\documentclass[zihao=-4,AutoFakeBold=true]{ctexbook}

\ctexset {
  part = {
    numbering = false
  },
  chapter = {
    numbering = false
  },
  section = {
    numbering = false
  },
  subsection = {
    numbering = false,
    format = \centering\bfseries,
  }
}

% \usepackage[margin=0.5cm]{geometry}
% For Meizu Hornor V10
\usepackage[
% showframe,
% includefoot,
% paperwidth=2.68in, % 2.68in,
paperwidth=3.79in,              %If for A4 paper size.
paperheight=5.36in,
left=0.1in,
right=0.1in,
top=0.1in,
bottom=0.20in,
footskip=10pt
]
{geometry}
% For Kindle 6'
% \usepackage[
% % showframe,
% paperwidth=3.6in,
% paperheight=4.8in,
% left=0.15in,
% right=0.15in,
% top=0.18in,
% bottom=0.18in,
% % footskip=10pt
% ]
% {geometry}

\usepackage{eso-pic}
\newcommand\BackgroundPic{%
  \put(0,0){%
    \parbox[b][\paperheight]{\paperwidth}{%
      \vfill
      \centering
      \includegraphics[width=\paperwidth,height=\paperheight,%
      keepaspectratio]{../bg.jpg}%
      \vfill
    }}}

\usepackage{xcolor}
\definecolor{BlueGrey}{HTML}{DBE1F1}
\definecolor{GreenRead}{HTML}{A8F29A}
\definecolor{GreyRead}{HTML}{D8D3D6}
\definecolor{YellowRead}{HTML}{F8FD89}
\definecolor{BlueRead}{HTML}{96ADFC}
\definecolor{Turquoise}{HTML}{A5F7E1}

\usepackage{xeCJKfntef}
\usepackage{xeCJK}
\xeCJKsetup{
  AutoFallBack = true,
  PunctStyle = kaiming,
  CheckFullRight = true,
  CheckSingle = false,
  AllowBreakBetweenPuncts = true,
  CJKmath = true,
}
% \setCJKmainfont[BoldFont={FZYouHei_GBK 513B},FallBack={BabelStone Han}]{FZYaSong-R-GBK}
\setCJKmainfont[BoldFont={FZYouHei_GBK 513B},FallBack={SimSun-ExtB}]{FZYaSong-R-GBK}% {FZNewKai-Z03}
% \setCJKmathfont{FZSongKeBenXiuKai-R-GBK}
% \setCJKmainfont[BoldFont={FZYouHei_GBK}]{FZYaSong-R-GBK}
\setCJKsansfont{Microsoft YaHei}
% \usepackage{CJKnumb}
\usepackage{zhnumber}
% 中文字体命令
\newCJKfontfamily[simsunb]\simsunb{SimSun-ExtB}
\newCJKfontfamily[hanaa]\hanaa{HanaMinA}
\newCJKfontfamily[hanab]\hanab{HanaMinB}
\newCJKfontfamily[fzbyks]\fzbyks{FZBoYaKanSongK}
\newCJKfontfamily[fzkt]\fzkt{FZNewKai-Z03}
\newCJKfontfamily[fzbys]\fzbys{FZYaSong-R-GBK}
\newCJKfontfamily[fzkeb]\fzkeb{FZKaiS-Extended(SIP)}
\newCJKfontfamily[jigmo]\jigmo{Jigmo} % For basic and ext A
\newCJKfontfamily[jigmotwo]\jigmotwo{Jigmo2} % For ext B.C.D.E.F
\newCJKfontfamily[jigmothree]\jigmothree{Jigmo3} % For ext G.H
\newfontfamily{\iosev}{IosevkaTerm NFM}


\newCJKfontfamily{\dejavusanszh}{DejaVu Sans} % this font has the required glyphs
\newfontface{\dejavusans}{DejaVu Sans} % this font has the required glyphs

\newcommand{\iching}[1]{{\dejavusanszh\char\numexpr"4DC0+#1}}
\newcommand{\trigram}[1]{{\dejavusans\char\numexpr"2630+#1}}

% 设置段落不缩进
% \setlength{\parindent}{0pt}

\usepackage{setspace}
\setstretch{1.3}

\usepackage{lastpage}
\usepackage{fancyhdr}
\fancyhead{}
\pagestyle{fancy}
\fancyfoot[C]{\scriptsize 第 \thepage\ 页（共 \pageref{LastPage} 页）}
% \fancyfoot[C]{\scriptsize 第 \thepage\ 页}
\renewcommand{\footrulewidth}{-10pt}
% \fancyhf{} % clear all header and footers
% \renewcommand{\headrulewidth}{0pt} % remove the header rule
% \rfoot{\thepage}
\fancypagestyle{plain}{%
  \renewcommand{\headrulewidth}{0pt}%
  \fancyhf{}%
}

\usepackage{fancyvrb}
\usepackage{fvextra}

\usepackage{framed}
\newenvironment{poem}[1][\hsize]{%
  \def\FrameCommand{{\color{red}\vrule width 0pt}\hspace{0pt}\fboxsep=\FrameSep}%
  \MakeFramed{\hsize#1\advance\hsize-\width\FrameRestore}}
{\endMakeFramed}

\usepackage{changepage}

% \usepackage{amssymb}
% \usepackage{mathabx}
\usepackage{latexsym}
\usepackage{amsmath}
\usepackage{txfonts}
\usepackage{amsfonts}
\usepackage{enumitem}

\usepackage{emptypage}

%% Need to fix, bookmarksdepth=4 not work
\usepackage[hidelinks,pdfpagelayout=SinglePage]{hyperref}
% \usepackage[hidelinks,bookmarksdepth=4]{hyperref}
\hypersetup{
  colorlinks=true,
  linkcolor=blue,
  filecolor=magenta,
  urlcolor=cyan,
}

%% 使每页的脚注的数字为黑色圆圈数字
\usepackage{pifont}
\makeatletter
\newcommand*{\bcircnum}[1]{%
  \expandafter\@bcircnum\csname c@#1\endcsname
}

\newcommand*{\@bcircnum}[1]{%
  \ifnum#1<1 %
  \@ctrerr
  \else
  \ifnum#1>30 %
  \@ctrerr
  \else
  \ding{\the\numexpr 181+(#1)\relax}%
  \fi
  \fi
}
\makeatother

\renewcommand*{\thefootnote}{\bcircnum{footnote}}

% \usepackage{booktabs}

\usepackage{fontspec}
\usepackage{xunicode}
\usepackage{xltxtra}
% \setmainfont{FZYaSongS-R-GB}
\XeTeXlinebreaklocale "zh"
\XeTeXlinebreakskip=0pt plus 1pt minus .1pt

% \newfontfamily\fzkeb[Path=./fonts/]{FZKaiSExtB.TTF}
% \newCJKfontfamily[kai]\fzkeb[Path=../../fonts/]{FZKaiSExtB.TTF}
% \newCJKfontfamily[kai]\hanaeb[Path=../../fonts/]{HanaMinB.ttf}
% \newCJKfontfamily[song]\hanaa[Path=~/.fonts/]{HanaMinA.ttf}
% \newCJKfontfamily[song]\hanab[Path=~/.fonts/]{HanaMinB.ttf}

%% Has problems when use \input or \include
% \usepackage{underscore}

\usepackage{xpinyin}
% \xpinyinsetup{
%   multiple={\color{red}}
%   % Sep={.1em plus .1em},
%   % font={\Large \bfseries}
% }

\usepackage[hang,flushmargin,perpage,stable]{footmisc}
% \renewcommand*{\footnotelayout}{\fontsize{6pt}{8pt}\selectfont\fzkt}
\renewcommand*{\footnotelayout}{\zihao{6}\fzkt}

% \usepackage{tcolorbox}
% \renewcommand{\footnotesize}{\fontsize{6pt}{8pt}\selectfont}

% \usepackage{etoolbox}
% \makeatletter
% \patchcmd{\@footnotetext}{\footnotesize}{\footnotesize\KaiTi}{}{}
% \makeatother

\newcommand{\sy}{%
  \begingroup\normalfont
  \includegraphics[height=1.3\fontcharht\font`\B]{sy.png}%
  \endgroup
}
\newcommand{\jj}{%
  \begingroup\normalfont
  \includegraphics[height=1.3\fontcharht\font`\B]{jj.png}%
  \endgroup
}
\newcommand{\zy}{%
  \begingroup\normalfont
  \includegraphics[height=1.3\fontcharht\font`\B]{zy.png}%
  \endgroup
}
\newcommand{\jinjie}{%
  \begingroup\normalfont
  \includegraphics[height=1.3\fontcharht\font`\B]{jinjie.png}%
  \endgroup
}
\newcommand{\jinyi}{%
  \begingroup\normalfont
  \includegraphics[height=1.3\fontcharht\font`\B]{jinyi.png}%
  \endgroup
}
\newcommand{\shu}{%
  \begingroup\normalfont
  \includegraphics[height=1.1\fontcharht\font`\B]{shu.png}%
  \endgroup
}
\newcommand{\zhu}{%
  \begingroup\normalfont
  \includegraphics[height=1.1\fontcharht\font`\B]{zhu.png}%
  \endgroup
}
\newcommand{\jian}{%
  \begingroup\normalfont
  \includegraphics[height=1.3\fontcharht\font`\B]{jian.png}%
  \endgroup
}
\newcommand{\zhuan}{%
  \begingroup\normalfont
  \includegraphics[height=1.3\fontcharht\font`\B]{zhuan.png}%
  \endgroup
}

\usepackage{makeidx}
\makeindex
% \DeclareRobustCommand{\myidx}[1]{{\Unifont #1}\index{#1}}
\DeclareRobustCommand{\myidx}[1]{{\fzzcys #1}\index{#1}}
% \DeclareRobustCommand{\myidxp}[1]{{\fzzcys \CJKunderline-{#1}}\index{#1}}
% \DeclareRobustCommand{\myidxb}[1]{{\fzzcys \CJKunderwave-{#1}}\index{#1}}

\usepackage{xparse}
\NewDocumentCommand{\myidxp}{mO{#1}}{{\fzzcys \CJKunderline-{#1}}\index{#2}}
\NewDocumentCommand{\myidxb}{mO{#1}}{{\fzzcys \CJKunderwave-{#1}}\index{#2}}

\usepackage[totoc,font=small,rule=0.2pt]{idxlayout}
\usepackage{parskip}

%% title & author
\title{\bfseries 尚\ 書\ 正\  義}
\author{漢·孔安國（注） \\ 唐·孔穎達（疏）}
\date{}

\begin{document}

%% Looks has problems when use \input or \include.
\AddToShipoutPicture*{\BackgroundPic}

\clearpage\maketitle
\clearpage\tableofcontents\thispagestyle{empty}
\newpage

\pagenumbering{arabic}
% \pagecolor{BlueGrey}
\zihao{4}

%% 正文

%% -*- coding: utf-8 -*-
%% Time-stamp: <Chen Wang: 2024-07-23 00:25:13>

% {\noindent \zhu \zihao{5} \fzbyks } -> 注(△ ○)
% {\noindent \shu \zihao{5} \fzkt } -> 疏

\part{序}


\chapter{卷一}

\section{尚書序}

 {\noindent \zhu\zihao{6}\fzbyks \CJKunderwave{釋文}：“此\CJKunderline{孔氏}所作，述\CJKunderwave{尚書}起之時代，並敘為注之由。故相承講之，今依舊為音。”\par}

 {\noindent \shu\zihao{5}\fzkt 正義曰：道本衝寂，非有名言。既形以道生，物由名舉，則凡諸經史，因物立名。物有本形，形從事著，聖賢闡教，事顯於言，言愜群心，書而示法，既書有法，因號曰“書”。後人見其久遠，自於上世，“尚”者，上也。言此上代以來之書，故曰“尚書”。且言者意之聲，書者言之記，是故存言以聲意，立書以記言。故易曰：“書不盡言，言不盡意。”是言者意之筌蹄，書言相生者也。書者，舒也。\CJKunderwave{書緯·璿璣鈐}云：“書者，如也。”則書者，寫其言，如其意，情得展舒也。又\CJKunderline{劉熙}\CJKunderwave{釋名}云：“書者，庶也，以記庶物。又為著。”言事得彰著。五經六籍皆是筆書，此獨稱“書”者，以彼五經者非是君口出言，即書為法，所書之事，各有云為，遂以所為別立其稱。稱以事立，故不名“書”。至於此書者，本書君事，事雖有別，正是君言，言而見書，因而立號，以此之故，名異諸部。但諸部之書，隨事立名，名以事舉，要名立之後，亦是筆書，故百氏六經總曰“書”也。\CJKunderwave{論讖}所謂“題意別名，各自載耳”。昭二年\CJKunderwave{左傳}曰，晉\CJKunderline{韓起}適魯，“觀書於太史氏，見易象與\CJKunderwave{魯春秋}”。此總名“書”也。“序”者，言序述\CJKunderwave{尚書}起、存亡注說之由，序為\CJKunderwave{尚書}而作，故曰“尚書序”。\CJKunderwave{周頌}曰：“繼序思不忘。”\CJKunderwave{毛傳}云：“序者，緒也。”則緒述其事，使理相胤續，若繭之抽緒。但易有\CJKunderwave{序卦}，\CJKunderline{子夏}作\CJKunderwave{詩序}，\CJKunderline{孔子}亦作\CJKunderwave{尚書序}，故\CJKunderline{孔君}因此作序名也。\CJKunderline{鄭玄}謂之“贊”者，以序不分散，避其序名，故謂之“贊”。贊者，明也，佐也。佐成序義，明以註解故也。安國以\CJKunderline{孔子}之序分附篇端，故已之總述亦謂之“序”。事不煩重，義無所嫌故也。\par}

\textcolor{red}{古者}伏犧氏之王天下也，始畫八卦，造書契，以代結繩之政，由是文籍\textcolor{red}{生焉}。\footnote{伏犧氏，伏古作虙，犧本又作羲，亦作戲，辭皮反。\CJKunderwave{說文}云，賈侍中說此犧非古字。張揖\CJKunderwave{字詁}云：“羲古字，戲今字。”一號包羲氏。三皇之最先，風姓，母曰華胥，以木德王，即太皞也。王，於況反。畫，乎麥反。卦，俱賣反。契，苦計反。書者，文字；契者，刻木而書其側：故曰“書契”也。一雲以書契約其事也。\CJKunderline{鄭玄}云：“以書書木邊，言其事，刻其木，謂之書契也。”結繩，\CJKunderwave{易·繫辭}上：“上古結繩以治，後世聖人易之以書契。”文，文字也。籍，籍書。}

{\noindent\shu\zihao{5}\fzkt “古者”至“生焉”。正義曰：“代結繩”者，言前世之政用結繩，今有書契以代之。則伏犧時始有文字以書事，故曰“由是文籍生焉”。自今本昔曰“古”。古者以聖德伏物教人取犧牲，故曰“伏犧”。字或作宓犧，音亦同。\CJKunderwave{律曆志}曰：結作網罟，以取犧牲，故曰“伏犧”。或曰“包犧”，言取犧而包之。顧氏讀包為庖，取其犧牲以供庖廚。顧氏又引\CJKunderwave{帝王世紀}云：“伏犧母曰華胥，有巨人跡，出於雷澤，華胥以足履之，有娠，生伏犧於成紀，蛇身人首。”\CJKunderwave{月令}云：“其帝太昊。”\CJKunderwave{繫辭}云：“古者包犧氏之王天下也。”是直變“包”言“伏”耳。則伏犧是皇，言“王天下”者，以皇與帝、王據跡為優劣，通亦為王。故\CJKunderwave{禮運}雲“昔者先王”，亦謂上代為王。但自下言之，則以上身為王，據王身於下，謂之“王天下”也。知伏犧“始畫八卦”者，以\CJKunderwave{繫辭}雲“包犧氏之王天下也”，後乃云：“始畫八卦以通神明之德，以類萬物之情”，故知之也。知時“造書契以代結繩之政”者，亦以\CJKunderwave{繫辭}雲“上古結繩而治，後世聖人易之以書契，蓋取諸夬”，是造書契可以代結繩也。彼直言“後世聖人”，知是伏犧者，以理比況而知。何則？八卦畫萬物之象，文字書百事之名，故\CJKunderwave{繫辭}曰：“仰則觀象於天，俯則觀法於地，觀鳥獸之文與地之宜，近取諸身，遠取諸物，始畫八卦。”是萬象見於卦。然畫亦書也，與卦相類，故知書契亦伏犧時也。由此孔意正欲須言伏犧時有書契，本不取於八卦。今雲“八卦”者，明書、卦相類，據\CJKunderwave{繫辭}有畫八卦之成文而言，明伏犧造書契也。言“結繩”者，當如鄭注云：“為約，事大大其繩，事小小其繩。王肅亦曰“結繩，識其政事”是也。言“書契”者，鄭云：“書之於木，刻其側為契，各持其一，後以相考合，若結繩之為治。”孔無明說，義或當然。\CJKunderwave{說文}云：“文者，物象之本也。”“籍”者，借也，藉此簡書以記錄政事，故曰“籍”。“蓋取諸夬”，“夬”者，決也，言文籍所以決斷，宣揚王政，是以夬。\CJKunderwave{繇}曰：“揚於王庭。”\CJKunderwave{繫辭}雲“包犧氏之王天下”，又云“作結繩而為罔罟，蓋取諸離”。彼謂結罔罟之繩，與結為政之繩異也。若然，\CJKunderwave{尚書緯}及\CJKunderwave{孝經讖}皆雲三皇無文字，又班固、馬融、\CJKunderline{鄭玄}、王肅諸儒皆以為文籍初自五帝，亦云三皇未有文字，與此說不同。何也？又蒼頡造書出於\CJKunderwave{世本}，蒼頡豈伏犧時乎？且\CJKunderwave{繫辭}雲黃帝、堯、舜為九事之目，未乃云：“上古結繩以治，後世聖人易之以書契。”是後世聖人即黃帝、堯、舜，何得為伏犧哉？孔何所據而更與\CJKunderwave{繫辭}相反，如此不同者？\CJKunderwave{藝文志}曰：“仲尼沒而微言絕，七十子喪而大義乖。”況遭秦焚書之後，群言競出，其緯文鄙近，不出聖人，前賢共疑，有所不取。通人考正，偽起哀、平，則\CJKunderline{孔君}之時，未有此\CJKunderwave{緯}，何可引以為難乎？其馬、鄭諸儒，以據文立說，見後世聖人在九事之科，便謂書起五帝，自所見有異，亦不可難孔也。而\CJKunderwave{繫辭}雲後世聖人在九事之下者，有以而然。案彼文先歷說伏犧、神農“蓋取”，下乃云：“黃帝、堯、舜垂衣裳而天下治，蓋取諸乾坤。”是黃帝、堯、舜之事也。又舟楫取渙，服牛取隨，重門取豫，臼杵取小過，弧矢取睽，此五者時無所繫，在黃帝、堯、舜時以否皆可以通也。至於宮室、葬與書契，皆先言“上古”、“古者”，乃言“後世聖人易之”，則別起事之端，不指黃帝、堯、舜時。以此葬事雲“古者”不雲“上古”，而云“易之以棺槨”。棺槨自殷湯而然，非是彼時之驗，則上古結繩何廢伏犧前也？其蒼頡，則說者不同，故\CJKunderwave{世本}云：“蒼頡作書。”司馬遷、班固、韋誕、宋忠、傅玄皆云：“蒼頡，黃帝之史官也。”崔瑗，曹植，蔡邕，索靖皆直云：“古之王也。”徐整云：“在神農、黃帝之間。”譙周云：“在炎帝之世。”衛氏云：“當在庖犧、蒼帝之世。”慎到云：“在庖犧之前。”張揖云：“蒼頡為帝王，生於禪通之紀。”\CJKunderwave{廣雅}曰：自開闢至獲麟二百七十六萬歲，分為十紀。則大率一紀二十七萬六千年。十紀者，九頭一也，五龍二也，攝提三也，合雒四也，連通五也，序命六也，循飛七也，因提八也，禪通九也，流訖十也。如揖此言，則蒼頡在獲麟前二十七萬六千餘年。是說蒼頡其年代莫能有定，亦不可以難孔也。然紀自燧人而下，揖以為自開闢而設，又伏犧前六紀後三紀，亦為據張揖、慎到、徐整等說，亦不可以年斷。其涗訖之紀，似自黃帝為始耳。又依\CJKunderwave{易緯·通卦驗}，燧人在伏犧前，“表計置其刻曰：蒼牙通靈，昌之成，孔演命，明道經”。\CJKunderline{鄭玄}注云：“刻謂刻石而記識之。”據此，伏犧前己有文字矣。又\CJKunderwave{陰陽書}稱天老對黃帝云：“鳳皇之象，首戴德，揹負仁，頸荷義，膺抱信，足履政，尾系武。”又\CJKunderwave{山海經}云：“鳳皇首文曰德，背文曰義，翼文曰順，膺文曰仁，腹文曰信。”又\CJKunderwave{易·繫辭}云：“河出圖，洛出書，聖人則之。是文字與天地並興焉。又\CJKunderwave{韓詩外傳}稱古封太山、禪梁甫者萬餘人，仲尼觀焉不能盡識。又\CJKunderwave{管子}書稱管仲對齊桓公曰：“古之封太山者七十二家，夷吾所識十二而已。”首有“無懷氏封太山，禪云云”。其登封者皆刻石紀號，但遠者字有雕毀，故不可識，則夷吾所不識者六十家，又在無懷氏前，\CJKunderline{孔子}睹而不識，又多於夷吾，是文字在伏犧之前己自久遠，何怪伏犧而有書契乎？如此者，蓋文字在三皇之前未用之教世，至伏犧乃用造書契以代結繩之政，是教世之用，猶燧人有火，中古用以燔黍捭豚，後聖乃修其利相似。文字理本有之，用否隨世而漸也。若然，惟\CJKunderwave{繫辭}至神農始有噬嗑與益，則伏犧時其卦未重，當無雜卦，而得有取諸夬者，此自\CJKunderline{鄭玄}等說耳。案\CJKunderwave{說卦}曰：“昔者聖人幽贊於神明而生蓍。”\CJKunderwave{繫辭}曰：“天生神物，聖人則之。”則伏犧用蓍而筮矣。故鄭注\CJKunderwave{說卦}亦曰：“昔者聖人，謂伏犧\CJKunderline{文王}也。”\CJKunderwave{繫辭}又曰：“十有八變而成卦。”是言爻皆三歸奇為三變，十八變則六爻明矣。則筮皆六爻，伏犧有筮，則有六爻，何為不重而怪有夬卦乎？ \par}

\textcolor{red}{伏犧}、神農、黃帝之書，謂之\CJKunderwave{三墳}，言大道也。少昊、顓頊、高辛、唐、虞之書，謂之\CJKunderwave{五典}，言\textcolor{red}{常道}也。\footnote{少，詩照反。昊，胡老反。少昊。金天氏，名摯，字青陽，一曰玄器，已姓。黃帝之子，母曰女節。以金德王，五帝之最先。顓音專。頊，許玉反。顓頊，高陽氏，姬姓。黃帝之孫，昌意之子，母曰景僕，謂之女樞。以水德王，五帝之二也。高辛，帝嚳也，姬姓。嚳，口毒反。母曰不見。以木德王，五帝之三也。唐，\CJKunderline{帝堯}也，姓伊耆氏。堯初為唐侯，後為天子，都陶，故號\CJKunderline{陶唐氏}。帝嚳之子，帝摯之弟，母曰慶都。以火德王，五帝之四也。虞，\CJKunderline{帝舜}也，姓姚氏，國號有虞。顓頊六世孫，瞽瞍之子，母曰握登。以土德王，五帝之五也。先儒解三皇五帝與\CJKunderline{孔子}同，並見發題。}

{\noindent\shu\zihao{5}\fzkt “伏犧”至“常道”也。正義曰：墳”，大也。以所論三皇之事，其道至大，故曰“言大道也”。以“典”者，常也，言五帝之道，可以百代常行，故曰“言常道也”。此三皇五帝，或舉德號，或舉地名，或直指其人，言及便稱，不為義例。顧氏引\CJKunderwave{帝王世紀}云：農母曰女登，有神龍首感女登而生炎帝，人身牛首。黃帝母曰附寶，見大電光繞北斗樞星，附寶感而懷孕，二十四月而生黃帝，日角龍顏。少昊金天氏母曰女節，有星如虹下流，意感而生少昊。顓頊母曰景僕，昌意正妃，謂之女樞，有星貫月如虹，感女樞於幽房之宮而生顓頊。堯母曰慶都，觀河遇赤龍，晻然陰風，感而有孕，十四月而生堯。又云舜母曰握登，見大虹感而生舜。此言“謂之三墳”，“謂之五典”者，因\CJKunderwave{左傳}有“三墳五典”之文，故指而謂之。然五帝之書皆謂之典，則\CJKunderwave{虞書·皋陶謨}、\CJKunderwave{益稷}之屬，亦應稱典。所以別立名者，若主論帝德，則以典為名，其臣下所為，隨義立稱。其\CJKunderwave{三墳}直雲“言大道也”，\CJKunderwave{五典}直雲“言常道也”，不訓“墳”、“典”之名者，以墳大典常，常訓可知，故略之也。“常道”所以與“大道”為異者，以帝者公平天下，其道可以常行，故以“典”言之。而皇優於帝，其道不但可常行而已，又更大於常，故言“墳”也。此為對例耳。雖少有優劣，皆是大道，並可常行。故\CJKunderwave{禮運}云：“以大道之行”為五帝時也。然帝號同天，名所莫加，優而稱“皇”者，以“皇”是美大之名，言大於帝也。故後代措廟立主，尊之曰“皇”，生者莫敢稱焉。而士庶祖父稱曰“皇”者，以取美名，可以通稱故也。案\CJKunderwave{左傳}上有“三墳五典”，不言墳是三皇之書，典是五帝之書。孔知然者，案今\CJKunderwave{堯典}、\CJKunderwave{舜典}是二帝二典，推此二典而上，則五帝當五典，是五典為五帝之書。今三墳之書在五典之上，數與三皇相當，墳又大名，與皇義相類，故云三皇之書為三墳。\CJKunderline{孔君}必知三皇有書者，案\CJKunderwave{周禮·小史職}“掌三皇五帝之書”，是其明文也。\CJKunderline{鄭玄}亦云其書即三墳五典。但\CJKunderline{鄭玄}以三皇無文，或據後錄定。\CJKunderline{孔君}以為書者記當時之事，不可以在後追錄，若當時無書，後代何以得知其道也？此亦\CJKunderline{孔君}所據三皇有文字之驗耳。\CJKunderline{鄭玄}注\CJKunderwave{中候}，依\CJKunderwave{運鬥樞}以伏犧、女媧、神農為三皇，又云五帝坐，帝鴻、金天、高陽、高辛、唐、虞氏。知不爾者，\CJKunderline{孔君}既不依緯，不可以緯難之。又易興作之條，不見有女媧，何以輒數？又\CJKunderline{鄭玄}云：“女媧修伏犧之道，無改作則。”已上修舊者眾，豈皆為皇乎？既不數女媧，不可不取黃帝以充三皇耳。又\CJKunderline{鄭玄}數五帝，何以六人？或為之說云，德協五帝，座不限多少，故六人亦名五帝。若六帝何有五座？而皇指大帝，所謂“耀魄寶”，止一而已，本自無三皇，何雲三皇？豈可三皇數人，五帝數座，二文舛互，自相乖阻也。其諸儒說三皇，或數燧人，或數祝融以配犧、農者，其五帝皆自軒轅，不數少昊，斯亦非矣。何燧人說者以為伏犧之前，據易曰“帝出於震”，震，東方，其帝太昊。又云“古者包犧氏之王天下也”，言古者製作莫先於伏犧，何以燧人側在前乎？又祝融及顓頊以下火官之號，金天已上百官之號，以其徵五經，無雲祝融為皇者，縱有，不過如\CJKunderline{共工氏}。\CJKunderline{共工}有水瑞，乃與犧、農、軒、摯相類，尚雲霸其九州，祝融本無此瑞，何可數之乎？\CJKunderwave{左傳}曰：“少昊之立，鳳鳥適至。”於\CJKunderwave{月令}又在秋享食，所謂白帝之室者也，何為獨非帝乎？故\CJKunderline{孔君}以黃帝上數為皇，少昊為五帝之首耳。若然，案今\CJKunderwave{世本·帝系}及\CJKunderwave{大戴禮·五帝德}並\CJKunderwave{家語}宰我問、太史公\CJKunderwave{五帝本紀}皆以黃帝為五帝，此乃史籍明文，而\CJKunderline{孔君}不從之者。孟軻曰：“信\CJKunderwave{書}不如其無\CJKunderwave{書}，吾於\CJKunderwave{武成}取二三策而已。”言書以漸染之濫也。孟軻已然，況後之說者乎？又\CJKunderwave{帝系}、\CJKunderwave{本紀}、\CJKunderwave{家語}、\CJKunderwave{五帝德}皆云：“少昊即黃帝子青陽是也，顓頊黃帝孫、昌意子，帝嚳高辛氏為黃帝曾孫、玄囂孫、僑極子，堯為帝嚳子，舜為顓頊七世孫。此等之書，說五帝而以黃帝為首者，原由\CJKunderwave{世本}。經於暴秦，為儒者所亂。\CJKunderwave{家語}則王肅多私定，\CJKunderwave{大戴禮}、\CJKunderwave{本紀}出於\CJKunderwave{世本}，以此而同。蓋以少昊而下皆出黃帝，故不得不先說黃帝，因此謬為五帝耳。亦由\CJKunderwave{繫辭}以黃帝與堯、舜同事，故儒者共數之焉。\CJKunderline{孔君}今者意以\CJKunderwave{月令}春曰太昊，夏曰炎帝，中央曰黃帝，依次以為三皇。又依\CJKunderwave{繫辭}，先包犧氏王。沒，神農氏作。又沒，黃帝氏作。亦文相次。皆著作見於易，此三皇之明文也。\CJKunderwave{月令}秋曰少昊，冬曰顓頊，自此為五帝。然黃帝是皇，今言“帝”不雲“皇”者，以皇亦帝也，別其美名耳。太昊為皇，\CJKunderwave{月令}亦曰“其帝太昊”，易曰“帝出於震”是也。又軒轅之稱黃帝，猶神農之雲炎帝，神農於\CJKunderwave{月令}為炎帝，不怪炎帝為皇，何怪軒轅稱帝？而梁主云：“書起軒轅，同以燧人為皇，其五帝自黃帝至堯而止。知帝不可以過五，故曰舜非三王，亦非五帝，與三王為四代而已。”其言與詩之為體，不雅則風，除皇已下不王則帝，何有非王非帝，以為何人乎？典，謨皆雲“帝曰”，非帝如何！ \par}

\textcolor{red}{至於}夏、商、周之書，雖設教不倫，雅誥奧義，其歸\textcolor{red}{一\xpinyin*{揆}}。\footnote{夏，\CJKunderline{禹}天下號也。以金德王，三王之最先。商，湯天下號，亦號殷，以水德王，三王之二也。周，\CJKunderline{文王}、\CJKunderline{武王}有天下號也。以木德王，三王之三也。誥，故報反，告也，示也。奧，烏報反，深也。揆，葵癸反，度也，道理也。}

{\noindent\shu\zihao{5}\fzkt “至於”至“一揆”。正義曰：既皇書稱“墳”，帝書稱“典”，除皇與帝墳、典之外，以次累陳，故言“至於”。夏、商、周三代之書，雖復當時所設之教，與皇及帝墳、典之等不相倫類，要其言皆是雅正辭誥，有深奧之義，其所歸趣與墳、典一揆。明雖事異墳、典而理趣終同，故所以同入\CJKunderwave{尚書}，共為世教也。\CJKunderline{孔君}之意，以墳、典亦是\CJKunderwave{尚書}，故此因墳、典而及三代。下雲“討論墳、典，斷自唐、虞以下”，是墳、典亦是\CJKunderwave{尚書}之內，而小史偏掌之者，以其遠代故也。此既言墳、典，不依外文連類，解八索、九丘，而言三代之書廁於其間者，孔意以墳、典是\CJKunderwave{尚書}，丘、索是\CJKunderwave{尚書}外物，欲先說\CJKunderwave{尚書}事訖，然後及其外物，故先言之也。夏、商、周之書，皆訓誥誓命之事，言“設教”者，以此訓誥誓命即為教而設，故云“設教”也。言“不倫”者，倫，類也。三代戰爭不與皇帝等類，若然，五帝稱“典”，三王劣而不倫，不得稱“典”。則三代非典，不可常行，何以垂法乎？然三王世澆，不如上代，故隨事立名，雖篇不目典，理實是典，故曰“雅誥奧義，其歸一揆”，即為典之謂也。然三王之書，惟無典謨，以外訓、誥、誓、命、歌、貢、徵、範，類猶有八，獨言“誥”者，以別而言之。其類有八，文從要約，一“誥”兼焉。何者？以此八事皆有言以誥示，故總謂之“誥”。又言“奧義”者，指其言謂之誥，論其理謂之義，故以義配焉。言“其歸一揆”，見三代自歸於一，亦與墳、典為一揆者，況喻之義。假譬人射，莫不皆發，志揆度於的，猶如聖人立教，亦同揆度於至理，故云“一揆”。 \par}

是故歷代寶之，以為大訓。

{\noindent\shu\zihao{5}\fzkt 正義曰：\CJKunderwave{顧命}云：“越玉五重，陳寶。”即以赤刀、大訓在西序，是“寶之，以為大訓”之文。彼注以典謨為之，與此相當。要六藝皆是，此直為\CJKunderwave{書}者，指而言之，故彼注亦然也。彼直周時寶之，此知歷代者，以墳、典久遠，周尚寶之，前代可知，故言“歷代”耳。 \par}

\textcolor{red}{八卦}之說，謂之\CJKunderwave{八索}，求其義也。九州之志，謂之\CJKunderwave{九丘}。言九州所有，土地所生，風氣所宜，皆聚\textcolor{red}{此書}也。\footnote{丘，聚也。八索，所白反，下同，求也。徐音素，本或作素。}

{\noindent\shu\zihao{5}\fzkt “八卦”至“此書”也。正義曰：以墳、典因外文而知其丘、索與墳、典文連，故連而說之，故總引傳文以充足己意，且為於下見與墳、典俱被黜削，故說而以為首引。言為論八卦事義之說者，其書謂之\CJKunderwave{八索}。其論九州之事所有志記者，其書謂之九丘，所以名“丘”者，以丘，聚也，言於九州當有土地所生之物，風氣所宜之事，莫不皆聚見於此書，故謂之\CJKunderwave{九丘}焉。然八卦言之“說”，九州言之“志”，不同者，以八卦交互相說其理，九州當州有所志識，以此而不同。此“索”謂求索，亦為搜索，以易八卦為主，故易曰：“八卦成列，象在其中矣。因而重之，爻在其中矣。”又曰：“八卦相蕩。”是六十四卦，三百八十四爻，皆出於八卦。就八卦而求其理，則萬有一千五百二十策，天下之事得，故謂之“索”，非一索再索而已。此“索”於\CJKunderwave{左傳}亦或謂之“索”，說有不同，皆後人失其真理，妄穿鑿耳。其\CJKunderwave{九丘}取名於聚，義多如山丘，故為聚。\CJKunderwave{左傳}或謂之“九區”，得為說當九州之區域，義亦通也。又言“九州所有”，此一句與下為總，即“土地所生，風氣所宜”是所有也。言“土地所生”，即其動物、植物，大率土之所生不出此二者。又云“風氣所宜”者，亦與土地所生大同。何者？以九州各有土地，有生與不生，由風氣所宜與不宜。此亦\CJKunderwave{職方}、\CJKunderwave{禹貢}之類。別而言之，“土地所生”若\CJKunderwave{禹貢}之“厥貢”、“厥篚”也，“風氣所宜”若\CJKunderwave{職方}其畜宜若干、其民若干男、若干女是也。上“墳”、“典”及“索”不別訓之，以可知，故略之。“丘”訓既難，又須別言“九州所宜”已下，故先訓之，於下結義，故云“皆聚此書”也。 \par}

\CJKunderwave{\textcolor{red}{春秋}左氏傳}曰，楚左史倚相“能讀三墳、五典、八索、九丘”，即謂上世帝王\textcolor{red}{遺書}也。\footnote{左史，史官左右。倚，於綺反，劉琴綺反。相，息亮反。倚相，楚靈王時史官也。}

{\noindent\shu\zihao{5}\fzkt “春秋”至“遺書”也。正義曰：以上因有外文言墳、典、丘、索而謂之，故引成文以證結之。此昭十二年\CJKunderwave{左傳}楚靈王見倚相趨過，告右尹子革以此辭。知“倚相”是其名字，蓋為太史，而主記左動之事，謂之“左史”。不然，或楚俗與諸國不同，官多以左右為名，或別有此左史乎？彼子革答王云：“倚相，臣問\CJKunderwave{祈招}之詩而不知，若問遠焉，其焉能知之？”彼以為倚相不能讀之。此雲“能”者，以此據\CJKunderwave{左傳}成文，因王言而引之。假不能讀，事亦無妨，況子革欲開諫王之路，倚相未必不能讀也。言此墳、典、丘，索即此書是謂上世帝王遺餘之書也。以楚王論時已在三王之末，故云“遺書”。其丘、索知是前事，亦不知在何代，故直總言“帝王”耳。 \par}

\textcolor{red}{先君}\CJKunderline{孔子}，生於周末，睹史籍之煩文，懼覽之者不一，遂乃定\CJKunderwave{禮}、\CJKunderwave{樂}，明舊章，刪\CJKunderwave{詩}為三百篇，約史記而修\CJKunderwave{春秋}，贊\CJKunderwave{易}道以黜\CJKunderwave{八索}，述\CJKunderwave{職方}以除\CJKunderwave{\textcolor{red}{九丘}}。

{\noindent\shu\zihao{5}\fzkt “先君”至“九丘”。正義曰：既結申帝王遺書，欲言\CJKunderline{孔子}就而刊定。\CJKunderwave{孔子世家}云，安國是\CJKunderline{孔子}十一世孫，而上尊先祖，故曰“先君”。\CJKunderwave{穀梁}以為魯襄公二十一年冬十一月庚子\CJKunderline{孔子}生，\CJKunderwave{左傳}哀公十六年夏四月已醜\CJKunderline{孔子}卒，計以周靈王時生，敬王時卒，故為“周末”。上雲“文籍”，下雲“滅先代典籍”，此言“史籍”。“籍”者，古書之大名。由文而有籍。謂之“文籍”；因史所書謂之“史籍”；可以為常，故曰“典籍”，義亦相通也。但上因書契而言“文”，下傷秦滅道以稱“典”，於此言“史”者，不但義通上下，又以此“史籍”不必是先王正史，是後代好事者作，以此懼其不一，故曰：“蓋有不知而作之者，我無是也。”先言“定\CJKunderwave{禮}、\CJKunderwave{樂}”者，欲明\CJKunderline{孔子}欲反於聖道以歸於一，故先言其舊行可從者。修而不改曰“定”，就而減削曰“刪”，準依其事曰“約”，因而佐成曰“贊”，顯而明之曰“述”，各從義理而言。獨\CJKunderwave{禮}、\CJKunderwave{樂}不改者，以\CJKunderwave{禮}、\CJKunderwave{樂}聖人制作，已無貴位，故因而定之。又云“明舊章”者，即\CJKunderwave{禮}、\CJKunderwave{樂}、\CJKunderwave{詩}、易、\CJKunderwave{春秋}是也。以“易道”、“職方”與“黜八索”、“除九丘”相對，其約史記以刪\CJKunderwave{詩}、\CJKunderwave{書}為偶，其定\CJKunderwave{禮}、\CJKunderwave{樂}文孤，故以“明舊章”配之，作文之體也。易亦是聖人所作，不言“定”者，以易非如\CJKunderwave{禮}、\CJKunderwave{樂}，人之行事，不須雲“定”。又因而為作\CJKunderwave{十翼}，故云“贊”耳。易文在下者，亦為“黜八索”與“除九丘”相近故也。為文之便，不為義例。\CJKunderline{孔子}之修六藝，年月孔無明說。\CJKunderwave{論語}曰：“吾自衛反魯，然後樂正，\CJKunderwave{雅}、\CJKunderwave{頌}各得其所。”則\CJKunderline{孔子}以魯哀公十一年反魯為大夫，十二年\CJKunderwave{孟子}卒，\CJKunderline{孔子}吊，則致仕時年七十以後。“修”，述也。\CJKunderwave{詩}有序三百一十一篇，全者三百五篇，雲“三百”者，亦舉全數計。\CJKunderwave{職方}在\CJKunderwave{周禮·夏官}，亦武帝時出於山岩屋壁，即藏秘府，世人莫見。以\CJKunderline{孔君}為武帝博士，於秘府而見為。知必“黜八索”、“除九丘”者，以三墳、五典本有八，今序只有二典而已，其三典、三墳今乃寂寞，明其除去，既墳、典書內之正尚有去者，況書外乎？故知丘、索亦黜除也。“黜”與“除”其義一也，黜退不用而除去之。必雲“贊易道以黜”者，以不有所興，孰有所廢故也。\CJKunderwave{職方}即\CJKunderwave{周禮}也，上已雲“定\CJKunderwave{禮}、\CJKunderwave{樂}”，即\CJKunderwave{職方}在其內。別雲述之，以為“除九丘”，舉其類者以言之。則雲“述”者，以定而不改即是遵述，更有書以述之。 \par}

\textcolor{red}{討論}\CJKunderwave{墳}、\CJKunderwave{典}，斷自唐虞以下，訖於周。芟夷煩亂，翦截浮辭，舉其宏綱，撮其機要，足以垂世立教，典、謨、訓、誥、誓、命之文凡\textcolor{red}{百篇}。\footnote{斷，丁亂反。訖，居乙反，又許乙反。芟，色鹹反。翦，諮淺反。撮，七活反。機，本又作幾。典凡十五篇，正典二，攝十三，十一篇亡。謨，莫胡反。凡三篇，正二，攝一。訓凡十六篇，正二篇亡，攝十四，三篇亡。誥凡三十八篇，正八，攝三十，十八篇亡。誓，市制反。凡十篇，正八，攝二，十篇亡。命凡十八篇，正十二，三篇亡，攝六，四篇亡。}

{\noindent\shu\zihao{5}\fzkt “討論”至“百篇”。正義曰：言\CJKunderline{孔子}既懼覽之者不一，不但刪\CJKunderwave{詩}、約史、定\CJKunderwave{禮}、贊易，有所黜除而已，又討整論理此三墳、五典並三代之書也。\CJKunderwave{論語}曰：“世叔討論之。”鄭以“討論”為整理，\CJKunderline{孔君}既取彼文，義亦當然。以書是亂物，故就而整理之。若然，墳、典\CJKunderline{周公}制禮，使小史掌之；而\CJKunderline{孔子}除之者，蓋隨世不同亦可，\CJKunderline{孔子}之時，墳、典已雜亂，故因去之。\CJKunderwave{左傳}曰“芟夷蘊崇之”，又曰“俘翦惟命”，\CJKunderwave{詩}曰“海外有截”，此\CJKunderline{孔君}所取之文也。“芟夷”者，據全代、全篇似草隨次皆芟，使平夷。若自帝嚳己上三典、三墳是芟夷之文，自夏至周雖有所留，全篇去之而多者，即“芟夷”也。“翦截”者，就代就篇辭有浮者翦截而去之，去而少者為“翦截”也。“舉其宏綱”即上“芟夷煩亂”也，“撮其機要”即上“翦截浮辭”也。且“宏綱”雲“舉”，是據篇、代大者言之；“機要”雲“撮”，為就篇、代之內而撮出之耳。“宏”，大也；“綱”者，網之索，舉大綱則眾目隨之。“機”者，機關，撮取其機關之要者，“斷自唐虞以下”者，孔無明說。\CJKunderwave{書緯}以為帝嚳以上，樸略難傳，唐虞已來，煥炳可法。又禪讓之首，至周五代一意故耳，孔義或然。“典”即\CJKunderwave{堯典}、\CJKunderwave{舜典}，“謨”即\CJKunderwave{大禹謨}、\CJKunderwave{皋陶謨}，“訓”即\CJKunderwave{伊訓}、\CJKunderwave{高宗之訓}，“誥”即\CJKunderwave{湯誥}、\CJKunderwave{大誥}，“誓”即\CJKunderwave{甘誓}、\CJKunderwave{湯誓}，“命”即\CJKunderwave{畢命}、\CJKunderwave{顧命}之等是也。說者以\CJKunderwave{書}體例有十，此六者之外尚有徵、貢、歌、範四者，並之則十矣。若\CJKunderwave{益稷}、\CJKunderwave{盤庚}，單言附於十事之例。今孔不言者，不但舉其機約，亦自徵、貢、歌、範非君出言之名，六者可以兼之。此雲“凡百篇”，據序而數故耳。或雲百二篇者，誤有所由。以前漢之時，有東萊張霸偽造\CJKunderwave{尚書}百兩篇，而為緯者附之。因此鄭云：“異者其在大司徒、大僕正乎？此事為不經也。”鄭作\CJKunderwave{書論}，依\CJKunderwave{尚書緯}云：“\CJKunderline{孔子}求書，得黃帝玄孫帝魁之書，迄於秦穆公，凡三千二百四十篇。斷遠取近，定可以為世法者百二十篇，以百二篇為\CJKunderwave{尚書}，十八篇為\CJKunderwave{中候}。”以為去三千一百二十篇，以上取黃帝玄孫，以為不可依用。今所考核\CJKunderwave{尚書}，首自舜之末年以禪於\CJKunderline{禹}，上錄舜之得用之事，由堯以為\CJKunderwave{堯典}，下取舜禪之後，以為舜讓得人，故史體例別，而不必君言。若\CJKunderwave{禹貢}全非君言，而\CJKunderline{禹}身事受禪之後無入\CJKunderwave{夏書}之言。是舜史自錄成一法，後代因之耳。 \par}

\textcolor{red}{所以}恢弘至道，示人主以軌範也。帝王之制，坦然明白，可舉而行，三千之徒並受\textcolor{red}{其義}。\footnote{恢，苦回反，大也。坦，土管反。}

{\noindent\shu\zihao{5}\fzkt “所以”至“其義”。正義曰：此論\CJKunderline{孔子}正理群經已畢，總而結之，故為此言。\CJKunderwave{家語}及\CJKunderwave{史記}皆雲“\CJKunderline{孔子}弟子三千人”，故云“三千之徒”也。 \par}

\textcolor{red}{及秦}始皇滅先代典籍，焚書坑儒，天下學士，逃難解散，我先人用藏其家書於\textcolor{red}{屋壁}。\footnote{始皇名政，二十六年初並六國，自號始皇帝。焚\CJKunderwave{書}、\CJKunderwave{詩}在始皇之三十四年，坑儒在三十五年。坑，苦庚反。難，乃旦反。解音蟹。}

{\noindent\shu\zihao{5}\fzkt “及秦”至“屋壁”。正義曰：言\CJKunderline{孔子}既定此書後，雖曰明白，反遭秦始皇滅除之。依\CJKunderwave{秦本紀}云，秦王正二十六年平定天下，尊為皇帝，不復立諡，以為初並天下，故號始皇。為滅先代典籍，故云“坑儒焚書”。以即位三十四年，因置酒於咸陽宮，丞相李斯奏請“天下敢有藏\CJKunderwave{詩}、\CJKunderwave{書}、百家語者，悉詣守、尉親燒之。有敢偶語\CJKunderwave{詩}、\CJKunderwave{書}者棄市。令下三十日不燒，黥為城旦”。制曰：“可。”是“焚書”也。三十五年，始皇以方士盧生求仙藥不得，以為誹謗，諸生連相告引，四百六十餘人皆坑之咸陽。是“坑儒”也。又衛宏\CJKunderwave{古文奇字序}云：“秦改古文以為篆隸，國人多誹謗。秦患天下不從，而召諸生，至者皆拜為郎，凡七百人。又密令冬月種瓜於驪山硎谷之中溫處，瓜實，乃使人上書曰：‘瓜冬有實。’有詔天下博士諸生說之，人人各異，則皆使往視之。而為伏機，諸生方相論難，因發機，從上填之以土，皆終命也。”“我先人用藏其家書於屋壁”者，\CJKunderwave{史記·孔子世家}云，\CJKunderline{孔子}生鯉，字伯魚。魚生伋，字子思。思生白，字子上。上生求，字子家。家生箕，字子京。京生穿，字子高。高生慎，慎為魏相。慎生鮒，鮒為陳涉博士。鮒弟子襄，為惠帝博士，長沙太守。襄生中。中生武。武生延陵及安國，為武帝博士，臨淮太守。\CJKunderwave{家語序}云：“子襄以秦法峻急，壁中藏其家書。”是安國祖藏之。 \par}

\textcolor{red}{漢室}龍興，開設學校，旁求儒雅，以闡大猷。濟南伏生，年過九十，失其本經，口以傳授。裁二十餘篇。以其上古之書，謂之\CJKunderwave{尚書}。百篇之義，世莫\textcolor{red}{得聞}。\footnote{校，戶教反。\CJKunderwave{詩}箋云：“鄭國謂學為校。”闡，尺善反，大也，明也。濟，子禮反，郡名也。伏生，名勝。過，古臥反，後同。傳，直專反，下“傳之”同。“二十餘篇”即馬、鄭所注二十九篇也。}

{\noindent\shu\zihao{5}\fzkt “漢室”至“得聞”。正義曰：將言所藏之書得之所由，故本之也。言“龍興”者，以易龍能變化，故比之聖人。九五“飛龍在天”，猶聖人在天子之位，故謂之“龍興”也。言“學校”者，校，學之一名也。故\CJKunderwave{鄭詩序}雲“子衿，刺學校廢”，\CJKunderwave{左傳}雲“然明請毀鄉校”是也。\CJKunderwave{漢書}云：“惠帝除挾書之律，立學興教，招聘名士。文景以後儒者更眾，至武帝尤甚。”故云“旁求儒雅”。\CJKunderwave{詩·小雅}曰：“匪先民是程，匪大猷是經。”彼注云：“猷，道也。”大道即先王六籍是也。伏生名勝，為秦二世博士，\CJKunderwave{儒林傳}云：“孝文帝時，求能治\CJKunderwave{尚書}者，天下無有，聞伏生治之，欲召。時伏生年已九十有餘，老不能行，於是詔太常，使掌故臣晁錯往受之。得二十九篇，即以教於齊魯之間。”是“年過九十”也。案\CJKunderwave{史記}：“秦時焚書，伏生壁藏之。其後兵火起，流。漢定天下，伏生求其書，亡數十篇，獨得二十九篇，以教於齊魯之閒。”則伏生壁內得二十九篇。而云“失其本經，口以傳授”者，蓋伏生初實壁內得之以教齊魯，傳教既久，誦文則熟，至其末年，因其習誦，或亦目暗，至年九十晁錯往受之時，不執經而口授之故也。又言“裁二十餘篇”者，意在傷亡，為少之文勢。何者？以數法隨所近而言之，若欲多之，當雲得三十篇，今“裁二十餘篇”，言“裁”亦意以為少之辭。又二十九篇自是計卷，若計篇則三十四，去\CJKunderwave{泰誓}猶有三十一。案\CJKunderwave{史記}及\CJKunderwave{儒林傳}皆雲“伏生獨得二十九篇，以教齊魯”，則今之\CJKunderwave{泰誓}，非初伏生所得。案馬融雲“\CJKunderwave{泰誓}後得”，\CJKunderline{鄭玄}\CJKunderwave{書論}亦云“民間得\CJKunderwave{泰誓}”。\CJKunderwave{別錄}曰：“武帝末，民有得\CJKunderwave{泰誓}書於壁內者，獻之。與博士使讀說之，數月皆起，傳以教人。”則\CJKunderwave{泰誓}非伏生所傳。而言二十九篇者，以司馬遷在武帝之世見\CJKunderwave{泰誓}出而得行，入於伏生所傳內，故為史總之，並雲伏生所出，不復曲別分析。雲民間所得，其實得時不與伏生所傳同也。但伏生雖無此一篇，而\CJKunderwave{書}傳有八百諸侯俱至孟津，白魚入舟之事，與\CJKunderwave{泰誓}事同，不知為伏生先為此說？不知為是\CJKunderwave{泰誓}出後。後人加增此語？案王充\CJKunderwave{論衡}及\CJKunderwave{後漢史}獻帝建安十四年黃門侍郎房宏等說云，宣帝本始元年，河內女子有壞老子屋，得古文\CJKunderwave{泰誓}三篇。\CJKunderwave{論衡}又云：“以掘地所得者。”今\CJKunderwave{史}、\CJKunderwave{漢}書皆雲伏生傳二十九篇，則司馬遷時已得\CJKunderwave{泰誓}，以並歸於伏生，不得雲宣帝時始出也。則雲宣帝時女子所得，亦不可信。或者爾時重得之，故於後亦據而言之。\CJKunderwave{史記}雲伏生得二十九篇，\CJKunderwave{武帝記}載今文\CJKunderwave{泰誓}末篇，由此劉向之作\CJKunderwave{別錄}，班固為\CJKunderwave{儒林傳}，不分明，因同於\CJKunderwave{史記}。而劉向雲武帝末得之\CJKunderwave{泰誓}，理當是一。而古今文不同者，即馬融所云：“吾見書傳多矣，凡諸所引，今之\CJKunderwave{泰誓}皆無此言，而古文皆有。”則古文為真，亦復何疑？但於先有張霸之徒偽造\CJKunderwave{泰誓}，以藏壁中，故後得而惑世也。亦可今之\CJKunderwave{泰誓}百篇之外，若\CJKunderwave{周書}之例，以於時實有觀兵之誓，但不錄入\CJKunderwave{尚書}。故古文\CJKunderwave{泰誓}曰“皇天震怒，命我文考，肅將天威，大勳未集。肆予小子發，以爾友邦冢君，觀政於商”是也。又云：“以其上古之書，謂之\CJKunderwave{尚書}”者，此文繼在“伏生”之下，則言“以其上古之書，謂之\CJKunderwave{尚書}”，此伏生意也。若以伏生指解\CJKunderwave{尚書}之名，名已先有，有則當雲名之\CJKunderwave{尚書}。既言“以其上古之書”，今先雲“以其”，則伏生意之所加，則知“尚”字乃伏生所加也。以“尚”解上，則“尚”訓為上。上者，下所慕尚，故義得為通也。\CJKunderline{孔君}既陳伏生此義，於下更無是非，明即用伏生之說，故書此而論之。馬融雖不見\CJKunderline{孔君}此說，理自然同，故曰“上古有虞氏之書，故曰\CJKunderwave{尚書}”是也。王肅曰：“上所言，史所書，故曰\CJKunderwave{尚書}。”鄭氏云：“尚者上也，尊而重之，若天書然，故曰\CJKunderwave{尚書}。”二家以“尚”與“書”相將，則上名不正出於伏生？\CJKunderline{鄭玄}依\CJKunderwave{書緯}，以“尚”字是\CJKunderline{孔子}所加，故\CJKunderwave{書贊}曰：“\CJKunderline{孔子}乃尊而命之曰\CJKunderwave{尚書}。”\CJKunderwave{璿璣鈐}云：“因而謂之\CJKunderwave{書}，加尚以尊之。”又曰：“\CJKunderwave{書}務以天言之。”\CJKunderline{鄭玄}溺於\CJKunderwave{書緯}之說，何有人言而須系之於天乎？且\CJKunderline{孔君}親見伏生，不容不悉，自雲伏生“以其上古之書，謂之\CJKunderwave{尚書}”，何雲\CJKunderline{孔子}加也？王肅雲“上所言，史所書”，則“尚”字與“書”俱有，無先後。既直雲“尚“，何以明上之所言？書者，以筆畫記之辭，群書皆是，何知\CJKunderwave{書}要責史所為也？此其不若前儒之說密耳。雲“上古”者亦無指定之目，自伏生言之，則於漢世，仰遵前代，自周已上皆是。馬融云：“有虞氏為\CJKunderwave{書}之初耳。”若易歷三世，則伏犧為上古，\CJKunderline{文王}為中古，\CJKunderline{孔子}為下古。\CJKunderwave{禮運}\CJKunderline{鄭玄}以先王食腥與易“上古結繩”同時，為上古，神農為中古，五帝為下古。其不相對，則無例耳。且太之與上為義不異，\CJKunderwave{禮}以唐虞為太古，以下有三代，冠而推之為然，是為不定則。但今世已上，仰之已古，便為上古耳。以\CJKunderwave{書}是本名，“尚”是伏生所加，故諸引\CJKunderwave{書}直雲“\CJKunderwave{書}曰”，若有配代而言，則曰“\CJKunderwave{夏書}”無言“\CJKunderwave{尚書}”者。 \par}

\textcolor{red}{至魯}共王好治宮室，壞\CJKunderline{孔子}舊宅，以廣其居，於壁中得先人所藏古文虞、夏、商、周之書及傳\CJKunderwave{論語}、\CJKunderwave{孝經}，皆科斗文字。王又升\CJKunderline{孔子}堂，聞金石絲竹之音，乃不\textcolor{red}{壞宅}。\footnote{共音恭，亦作龔，又作恭。共王，漢景帝之子，名餘。好，呼報反，下“好古”同。壞音怪，下同。\CJKunderwave{字林}作𣀤，雲公壞反，毀也。傳謂\CJKunderwave{春秋}也，一雲周易十翼，非經謂之傳。論如字，又音倫。科，苦禾反。科斗，蟲名，蝦蟆子。書形似之。}

{\noindent\shu\zihao{5}\fzkt “至魯”至“壞宅”。正義曰：欲雲得百篇之由，故序其事。漢景帝之子名餘，封於魯，為王，死諡曰共。存日以居於魯，近\CJKunderline{孔子}宅。好治宮室，故欲裒益，乃壞\CJKunderline{孔子}舊宅，以增廣其居。於所壞壁內得安國先人所藏古文虞夏商周之書及傳\CJKunderwave{論語}、\CJKunderwave{孝經}，皆是科斗文字。王雖得此書，猶壞不止。又升\CJKunderline{孔子}廟堂，聞金鐘石磬絲琴竹管之音，以懼其神異乃止，不復敢壞宅也。上言藏家書於屋壁，此亦屋壁內得書也，亦得及傳\CJKunderwave{論語}、\CJKunderwave{孝經}等。不從約雲“得\CJKunderwave{尚書}”，而煩文言“虞夏商周之書”者，以壁內所得，上有題目“虞夏商周書”，其序直雲“書序”，皆無“尚”字，故其目錄亦然，故不雲“尚書”而言“虞夏商周之書”。安國亦以此知“尚”字是伏生所加。惟此壁內所無，則書本無“尚”字明矣。凡書，非經則謂之傳。言“及傳\CJKunderwave{論語}、\CJKunderwave{孝經}”，正謂\CJKunderwave{論語}、\CJKunderwave{孝經}是傳也。漢武帝謂東方朔云：“傳曰：‘時然後言，人不厭其言。’”又漢東平王劉雲與其太師策書云：“傳曰：‘陳力就列，不能者止。’”又成帝賜翟方進策書云：“傳曰：‘高而不危，所以長守貴也。’”是漢世通謂\CJKunderwave{論語}、\CJKunderwave{孝經}為傳也。以\CJKunderwave{論語}、\CJKunderwave{孝經}非先王之書，是\CJKunderline{孔子}所傳說，故謂之傳，所以異於先王之書也。上已雲“壞\CJKunderline{孔子}舊宅”，又云“乃不壞宅”者，初王意欲壞之，已壞其屋壁，聞八音之聲，乃止，餘者不壞，明知已壞者亦不敢居，故云“乃不壞宅”耳。 \par}

\textcolor{red}{悉以}書還\CJKunderline{孔氏}。科斗書廢已久，時人無能知者，以所聞伏生之書考論文義，定其可知者，為隸古定，更以竹簡寫之，增多伏生二十五篇。伏生又以\CJKunderwave{舜典}合於\CJKunderwave{堯典}，\CJKunderwave{益稷}合於\CJKunderwave{皋陶謨}，\CJKunderwave{盤庚}三篇合為一，\CJKunderwave{康王之誥}合於\CJKunderwave{顧命}，復出此篇，並序，凡五十九篇，為四十六卷。其餘錯亂摩滅，弗可復知，悉上送官，藏之書府，以待\textcolor{red}{能者}。\footnote{隸音麗，謂用隸書寫古文。二十五篇謂\CJKunderwave{虞書·大禹謨}，\CJKunderwave{夏書·五子之歌}、\CJKunderwave{胤徵}，\CJKunderwave{商書·仲虺之誥}、\CJKunderwave{湯誥}、\CJKunderwave{伊訓}、\CJKunderwave{太甲}三篇、\CJKunderwave{鹹有一德}、\CJKunderwave{說命}三篇，\CJKunderwave{周書·泰誓}三篇、\CJKunderwave{武成}、\CJKunderwave{旅獒}、\CJKunderwave{微子之命}、\CJKunderwave{蔡仲之命}、\CJKunderwave{周官}、\CJKunderwave{君陳}、\CJKunderwave{畢命}、\CJKunderwave{君牙}、\CJKunderwave{冏命}。合舊音合，又如字，下同。皋音高，本又作咎。陶音遙，本又作繇。盤，步幹反，本又作般。復，扶又反，下同。五十九篇即今所行五十八篇，其一是百篇之序。謂\CJKunderwave{虞書·汩作}、\CJKunderwave{九共}九篇、\CJKunderwave{膏飫}，\CJKunderwave{夏書·帝告}、\CJKunderwave{釐沃}、\CJKunderwave{湯徵}、\CJKunderwave{汝鳩}、\CJKunderwave{汝方}，\CJKunderwave{商書·夏社}、\CJKunderwave{疑至}、\CJKunderwave{臣扈}、\CJKunderwave{典寶}、\CJKunderwave{明居}、\CJKunderwave{肆命}、\CJKunderwave{徂後}、\CJKunderwave{沃丁}、\CJKunderwave{鹹乂}四篇、\CJKunderwave{伊陟}、\CJKunderwave{原命}、\CJKunderwave{仲丁}、\CJKunderwave{河亶甲}、\CJKunderwave{祖乙}、\CJKunderwave{高宗之訓}，\CJKunderwave{周書·分器}、\CJKunderwave{旅巢命}、\CJKunderwave{歸禾}、\CJKunderwave{嘉禾}、\CJKunderwave{\CJKunderline{成王}政}、\CJKunderwave{將蒲姑}、\CJKunderwave{賄肅慎之命}、\CJKunderwave{亳姑}，凡四十二篇，亡。上，時掌反。}

{\noindent\shu\zihao{5}\fzkt “悉以”至“能者”。正義曰：既雲王不壞宅，以懼神靈，因還其書。已前所得，言“悉以書還\CJKunderline{孔氏}”，則上“傳\CJKunderwave{論語}、\CJKunderwave{孝經}”等皆還之，故言“悉”也。科斗書，古文也，所謂蒼頡本體，周所用之，以今所不識，是古人所為，故名“古文”。形多頭粗尾細，狀腹團圓，似水蟲之科斗，故曰“科斗”也。以古文經秦不用，故云廢已久矣，時人無能知識者。\CJKunderline{孔君}以人無能知識之故，己欲傳之，故以所聞伏生之書，比校起發，考論古文之義。考文而云“義”者，以上下事義推考其文，故云“義”也。“定其可知”者，就古文內定可知識者為隸古定。不言就伏生之書而云以其所聞者，明用伏生書外亦考之。故云“可知者”，謂並伏生書外有可知，不徒伏生書內而已。言“隸古”者，正謂就古文體而從隸定之。存古為可慕，以隸為可識，故曰“隸古”，以雖隸而猶古。由此故謂\CJKunderline{孔君}所傳為古文也。古文者，蒼頡舊體，周世所用之文字。案班固\CJKunderwave{漢志}及許氏\CJKunderwave{說文}，書本有六體：一曰指事，上下；二曰象形，日月；三曰形聲，江河；四曰會意，武信；五曰轉註，考老；六曰假借，令長。此造字之本也。自蒼頡以至今，字體雖變，此本皆同，古今不易也。自蒼頡以至周宣，皆蒼頡之體，未聞其異。宣王紀其史籀始有大篆十五篇，號曰篆籀，惟篆與蒼頡二體而已。衛恆曰：“蒼頡造書，觀於鳥跡，因而遂滋，則謂之字。字有六義，其文至於三代不改。及秦用篆書，焚燒先代典籍，古文絕矣。”許慎\CJKunderwave{說文}言自秦有八體：一曰大篆，二曰小篆，三曰刻符，四曰蟲書，五曰摹印，六曰署書，七曰殳書，八曰隸書。亡新居攝，以應制作，改定古文，使甄豐校定，時有六書：一曰古文，\CJKunderline{孔子}壁內書也；二曰奇字，即古字有異者；三曰篆書，即小篆，下杜人程邈所作也；四曰佐書，秦隸書也；五曰繆篆，所以摹印也；六曰鳥蟲書，所以書幡信也。由此而論，即秦罷古文而有八體，非古文矣。以至亡新六書並八體，亦用書之六體以造其字。其亡新六書於秦八體，用其小篆、蟲書、摹印、隸書，去其大篆、刻符、殳書、署書，而加以古文與奇字，其刻符及署書蓋同摹印，殳書同於繆篆，大篆正古文之別，以慕古故乃用古文與奇字而不用大篆也。是\CJKunderline{孔子}壁內古文即蒼頡之體，故\CJKunderline{鄭玄}云：“書初出屋壁，皆周時象形文字，今所謂科斗書。”以形言之為科斗，指體即周之古文。\CJKunderline{鄭玄}知者，若於周時秦世所有，至漢猶當識之，不得雲無能知者。又亡新古文亦云即\CJKunderline{孔氏}壁內古文，是其證也。或以古文即大篆，非也。何者？八體六書自大篆，與古文不同；又秦有大篆，若大篆是古文，不得雲“古文遂絕”，以此知大篆非古文也。六書古文與蟲書本別，則蟲書非科斗書也。\CJKunderline{鄭玄}雲周之象形文字者，總指六書象科斗之形，不謂六書之內“一曰象形”也。又云“更以竹簡寫之”，明留其壁內之本也。顧氏云：“策長二尺四寸，簡長一尺二寸。”“曾多伏生二十五篇”者，以壁內古文篇題殊別，故知“以\CJKunderwave{舜典}合於\CJKunderwave{堯典}，\CJKunderwave{益稷}合於\CJKunderwave{皋陶謨}”。伏生之本亦壁內古文而合之者，蓋以老而口授之時，因誦而連之，故殊耳。其\CJKunderwave{盤庚}本當同卷，故有並也。\CJKunderwave{康王之誥}以一時之事，連誦而同卷，當以“王出在應門之內”為篇首，及以“王若曰，庶邦”亦誤矣。以伏生本二十八篇，\CJKunderwave{盤庚}出二篇，加\CJKunderwave{舜典}、\CJKunderwave{益稷}、\CJKunderwave{康王之誥}凡五篇為三十三篇，加所增二十五篇為五十八，加序一篇為五十九篇，雲“復出此篇，並序，凡五十九篇”。此雲“為四十六卷”者，謂除序也。下雲“定五十八篇，既畢”，不更雲卷數，明四十六卷故爾。又伏生二十九卷而序在外，故知然矣。此雲“四十六卷”者，不見安國明說，蓋以同序者同卷，異序者異卷，故五十八篇為四十六卷。何者？五十八篇內有\CJKunderwave{太甲}、\CJKunderwave{盤庚}、\CJKunderwave{說命}、\CJKunderwave{泰誓}皆三篇共卷，減其八，又\CJKunderwave{大禹謨}、\CJKunderwave{皋陶謨}、\CJKunderwave{益稷}又三篇同序共卷，其\CJKunderwave{康誥}、\CJKunderwave{酒誥}、\CJKunderwave{梓材}亦三篇同序共卷，則又減四，通前十二，以五十八減十二，非四十六卷而何？其\CJKunderwave{康王之誥}乃與\CJKunderwave{顧命}別卷，以別序故也。“其餘錯亂摩滅”，五十八篇外四十二篇也，以不可復知，亦上送官。其可知者己用竹簡寫得其本，亦俱送入府，故在秘府得有古文也。以後生可畏，或賢聖間出，故須藏之以待能整理讀之者。 \par}

\textcolor{red}{承詔}為五十九篇作傳，於是遂研精覃思，博考經籍，採\xpinyin*{摭}群言，以立訓傳。約文申義，敷暢厥旨，庶幾有補於\textcolor{red}{將來}。\footnote{為，於偽反。覃，徒南反，深也。思，息嗣反。採，本又作採。摭，之石反，一音之若反。敷，芳夫反。暢，醜亮反。}

{\noindent\shu\zihao{5}\fzkt “承詔”至“將來”。正義曰：安國時為武帝博士，\CJKunderline{孔君}考正古文之日，帝之所知，亦既定訖，當以聞於帝，帝令註解，故云“承詔為五十九篇作傳”。以注者多言曰“傳”，“傳”者，傳通故也。以“傳”名出自丘明。賓牟賈對\CJKunderline{孔子}曰“史失其傳”。又\CJKunderwave{喪服}儒者皆雲\CJKunderline{子夏}作傳，是“傳”名久矣。但大率秦漢之際，多名為“傳”；於後儒者以其傳多，或有改之別雲“註解”者；仍有同者，以當時之意耳。說者為例院“前漢稱傳，於後皆稱注”，誤矣。何者？馬融、王肅亦稱註名為“傳”，傳何有例乎？以聖道弘深，當須詳悉，於是研核精審，覃靜思慮以求其理，冀免乖違，既顧察經文，又取證於外，故須廣博推考群經六籍，又捃拾採摭群書之言，以此文證造立訓解，為之作傳。明不率爾。雖復廣證，亦不煩多，為傳直約省文，令得申盡其義。明文要義通，不假煩多也。以此得申，故能遍佈通暢\CJKunderwave{書}之旨意，是辭達而已，不求於煩。既義暢而文要，則觀者曉悟，故云庶幾有所補益於將來，讀之者得悟而有益也。敷，布也。厥，其也。庶，幸也。幾，冀也。\CJKunderwave{爾雅}有訓。既雲“經籍”，又稱“群言”者，“經籍”，五經是也；“群言”，子史是也。以\CJKunderwave{書}與經籍理相因通，故云“博考”；子史時有所須，故云“採摭”耳。案\CJKunderline{孔君}此傳辭旨不多，是“約文”也。要文無不解，是“申義”也。其義既申，故云敷暢其義之旨趣耳。考其此注，不但言少，\CJKunderwave{書}之為言多須詁訓，而\CJKunderline{孔君}為例，一訓之後，重訓者少，此亦約文也。 \par}

\CJKunderwave{\textcolor{red}{書序}}，序所以為作者之意。昭然義見，宜相附近，故引之各冠其篇首，定五十八篇。既畢，會國有巫蠱事，經籍道息，用不復以聞，傳之子孫，以貽後代。若好古博雅君子，與我同志，亦所不\textcolor{red}{隱也}。\footnote{為，於偽反，又如字。見，賢遍反。冠，工亂反。巫蠱，漢武帝末徵和中，江充造蠱敗戾太子，故經籍道息焉。巫音無。蠱音古。貽，以之反。}

{\noindent\shu\zihao{5}\fzkt “書序”至“隱也”。正義曰：\CJKunderline{孔君}既言己立傳之意，又當斟酌所宜。而\CJKunderwave{書序}雖名為序，不是總陳書意泛論，乃篇篇各序作意，但作序者不敢廁於正經，故謙而聚於下。而注述者不可代作者之謙，須從利益而欲分之，從便雲序，序所以當篇為作此書之意，則是當篇作意觀序而昭然，意義顯見。既義見由序，此序宜各與其本篇相從附近，不宜聚於一處。故每篇引而分之，各冠加於篇首，令意昭見。序既分散，損其一篇，故定五十八篇。然此本承詔而作，作畢當以上奏聞知，但會值國家有巫蠱之事，好愛經籍之道滅息，假奏亦不能行用，為此之故，不復以此傳奏聞。亦以既傳成不得聞上，惟自傳於己之子孫，以遺與後世之人使行之。亦不敢望後世必行，故云若後世有好愛古道、廣博學問、志懷雅正如此之君子，冀能與我同於慕古之志，以行我道。我道得此人流行，亦所以傳不隱蔽。是弘道由人也。言“巫蠱”者，\CJKunderwave{王制}曰：“執左道以亂政者殺。”\CJKunderline{鄭玄}注云：“左道謂巫蠱之屬。”以非正道，故謂之左道。以蠱皆巫之所行，故云巫蠱。蠱者總名。\CJKunderwave{左傳}云：“惑蠱其君。”則蠱者怪惑之名。指體則藥毒害人者是，若行符厭俗之為魅，令人蠱惑夭年傷性皆是也。依\CJKunderwave{漢書}，此時武帝末年，上已年老，淫惑鬼神，崇信巫術。由此奸人江充因而行詐，先於太子宮埋桐人，告上云：“太子宮有蠱氣。”上信之，使江充治之，於太子宮果得桐人。太子知己不為此，以江充故為陷己，因而殺之。而帝不知太子實心，謂江充言為實，即詔丞相劉屈犛發三輔兵討之。太子看長安因與鬥，不勝而出走，奔湖關自殺。此即巫蠱事也。言“不隱”者，不謂恐隱藏己道，以己道人所不知，懼其幽隱，人能行之使顯，為不隱蔽耳。易曰：“謙謙君子。”仁者好謙，而\CJKunderline{孔君}自作揄揚。雲君子知己者亦意在教世，欲令人睹此言，知己傳是深遠，因而有所曉寤，令之有益，故不可以苟謙也。亦猶\CJKunderline{孔子}曰：“何有於我哉”。 \par}

%%% Local Variables:
%%% mode: latex
%%% TeX-engine: xetex
%%% TeX-master: "../Main"
%%% End:

%% -*- coding: utf-8 -*-
%% Time-stamp: <Chen Wang: 2024-04-02 11:42:41>

% {\noindent \zhu \zihao{5} \fzbyks } -> 注 (△ ○)
% {\noindent \shu \zihao{5} \fzkt } -> 疏

\part{虞書}

\chapter{卷二}

\section{堯典第一（堯典上）}

 {\noindent\zhu\zihao{6}\fzbyks \CJKunderwave{釋文}：“凡十六篇，十一篇亡，五篇見存。” \par}

 {\noindent\shu\zihao{5}\fzkt “古文尚書堯典第一”。正義曰：檢古本並石經，直言“\CJKunderwave{堯典}第一”，無“古文尚書”。以\CJKunderline{孔君}從隸古，仍號“古文”，故後人因而題於此，以別伏生所出、大小夏侯及歐陽所傳為今文故也。“堯典第一”，篇之名，當與眾篇相次。“第”訓為次也，於次第之內而處一，故曰“堯典第一”。以此第一者，以五帝之末接三王之初，典策既備，因機成務，交代揖讓，以垂無為，故為第一也。然\CJKunderwave{書}者理由舜史，勒成一家，可以為法，上取堯事，下終禪禹，以至舜終，皆為舜史所錄。其堯、舜之典，多陳行事之狀，其言寡矣，\CJKunderwave{禹貢}即全非君言，準之後代，不應入\CJKunderwave{書}，此其一體之異，以此禹之身事於禪後，無入\CJKunderwave{夏書}之理。自\CJKunderwave{甘誓}已下，皆多言辭，則古史所書於是乎始。知\CJKunderwave{五子之歌}亦非上言，典書草創，以義而錄，但致言有本，名隨其事。檢其此體，為例有十。一曰典，二曰謨，三曰貢，四曰歌，五曰誓，六曰誥，七曰訓，八曰命，九曰徵，十曰範。\CJKunderwave{堯典}、\CJKunderwave{舜典}二篇，典也。\CJKunderwave{大禹謨}、\CJKunderwave{皋陶謨}二篇，謨也。\CJKunderwave{禹貢}一篇，貢也。\CJKunderwave{五子之歌}一篇，歌也。\CJKunderwave{甘誓}、\CJKunderwave{泰誓}三篇、\CJKunderwave{湯誓}、\CJKunderwave{牧誓}、\CJKunderwave{費誓}、\CJKunderwave{泰誓}八篇，誓也。\CJKunderwave{仲虺之誥}、\CJKunderwave{湯誥}、\CJKunderwave{大誥}、\CJKunderwave{康誥}、\CJKunderwave{酒誥}、\CJKunderwave{召誥}、\CJKunderwave{洛誥}、\CJKunderwave{康王之誥}八篇，誥也。\CJKunderwave{伊訓}一篇，訓也。\CJKunderwave{說命}三篇，\CJKunderwave{微子之命}、\CJKunderwave{蔡仲之命}、\CJKunderwave{顧命}、\CJKunderwave{畢命}、\CJKunderwave{冏命}、\CJKunderwave{文侯之命}九篇，命也。\CJKunderwave{胤徵}一篇，徵也。\CJKunderwave{洪範}一篇，範也。此各隨事而言。\CJKunderwave{益稷}亦謨也，因其人稱言以別之。其\CJKunderwave{太甲}、\CJKunderwave{咸有一德}，伊尹訓道王，亦訓之類。\CJKunderwave{盤庚}亦誥也，故王肅云：“不言誥，何也？取其徒而立功，非但錄其誥。”\CJKunderwave{高宗肜日}與訓序連文，亦訓辭可知也。\CJKunderwave{西伯戡黎}云，“祖伊恐，奔告於受”，亦誥也。\CJKunderwave{武成}云，“識其政事”，亦誥也。\CJKunderwave{旅獒}戒王，亦訓也。\CJKunderwave{金縢}自為一體，祝亦誥辭也。\CJKunderwave{梓材}，\CJKunderwave{酒誥}分出，亦誥也。\CJKunderwave{多士}以王命誥，自然誥也。\CJKunderwave{無逸}戒王，亦訓也。\CJKunderwave{君奭}周公誥召公，亦誥也。\CJKunderwave{多方}、\CJKunderwave{周官}上誥於下，亦誥也。\CJKunderwave{君陳}、\CJKunderwave{君牙}與\CJKunderwave{畢命}之類，亦命也。\CJKunderwave{呂刑}陳刑告王，亦誥也。\CJKunderwave{書}篇之名，因事而立，既無體例，隨便為文。其百篇次第，於序孔、鄭不同。孔以\CJKunderwave{湯誓}在\CJKunderwave{夏社}前，於百篇為第二十六；鄭以為在\CJKunderwave{臣扈}後，第二十九。孔以\CJKunderwave{咸有一德}次\CJKunderwave{太甲}後，第四十；鄭以為在\CJKunderwave{湯誥}後，第三十二。孔以\CJKunderwave{蔡仲之命}次\CJKunderwave{君奭}後，第八十三；鄭以為在\CJKunderwave{費誓}前，第九十六。孔以\CJKunderwave{周官}在\CJKunderwave{立政}後，第八十八；鄭以為在\CJKunderwave{立政}前，第八十六。孔以\CJKunderwave{費誓}在\CJKunderwave{文侯之命}後，第九十九；鄭以為在\CJKunderwave{呂刑}前，第九十七。不同者孔依壁內篇次及序為文，鄭依賈氏所奏\CJKunderwave{別錄}為次，孔未入學官。以此不同，考論次第，孔義是也。 \par}

 {\noindent\shu\zihao{5}\fzkt 正義曰：\CJKunderwave{堯典}雖曰唐事，本以虞史所錄，末言舜登庸由堯，故追堯作典，非唐史所錄，故謂之\CJKunderwave{虞書}也。\CJKunderline{鄭玄}雲“舜之美事，在於堯時”是也。案馬融、\CJKunderline{鄭玄}、王肅、\CJKunderwave{別錄}題皆曰\CJKunderwave{虞夏書}，以虞、夏同科，雖虞事亦連夏。此直言\CJKunderwave{虞書}，本無\CJKunderwave{夏書}之題也。案鄭序以為\CJKunderwave{虞夏書}二十篇，\CJKunderwave{商書}四十篇，\CJKunderwave{周書}四十篇，\CJKunderwave{贊}雲“三科之條，五家之教”，是虞、夏同科也。其孔於\CJKunderwave{禹貢}注云“禹之王以是功，故為\CJKunderwave{夏書}之首”，則虞、夏別題也，以上為\CJKunderwave{虞書}則十六篇。又\CJKunderwave{帝告}、\CJKunderwave{釐沃}、\CJKunderwave{湯徵}、\CJKunderwave{汝鳩}、\CJKunderwave{汝方}於\CJKunderline{鄭玄}為\CJKunderwave{商書}，而孔並於\CJKunderwave{胤徵}之下，或以為夏事，猶\CJKunderwave{西伯戡黎}，則\CJKunderwave{夏書}九篇，\CJKunderwave{商書}三十五篇，此與鄭異也。或孔因\CJKunderwave{帝告}以下五篇亡，並注於\CJKunderwave{夏書}不廢，猶\CJKunderwave{商書}乎？別文所引皆雲“\CJKunderwave{虞書}曰”、“\CJKunderwave{夏書}曰”，無並言\CJKunderwave{虞夏書}者。又伏生雖有一\CJKunderwave{虞夏傳}，以外亦有\CJKunderwave{虞傳}、\CJKunderwave{夏傳}，此其所以宜別也，此孔依虞、夏各別而存之。莊八年\CJKunderwave{左傳}引“\CJKunderwave{夏書}曰‘皋陶邁種德’”，僖二十四年\CJKunderwave{左傳}引夏書曰“地平天成”，二十七年引\CJKunderwave{夏書}“賦納以言”，襄二十六年引\CJKunderwave{夏書}曰“與其殺不辜，寧失不經”，皆在\CJKunderwave{大禹謨}、\CJKunderwave{皋陶謨}。當雲\CJKunderwave{虞書}而云\CJKunderwave{夏書}者，以事關禹，故引為\CJKunderwave{夏書}。若\CJKunderwave{洪範}以為\CJKunderwave{周書}，以箕子至周，商人所陳而傳引之，即曰\CJKunderwave{商書}也。案壁內所得，孔為傳者凡五十八篇，為四十六卷。三十三篇與鄭注同，二十五篇增多鄭注也。其二十五篇者，\CJKunderwave{大禹謨}一，\CJKunderwave{五子之歌}二，\CJKunderwave{胤徵}三，\CJKunderwave{仲虺之誥}四，\CJKunderwave{湯誥}五，\CJKunderwave{伊訓}六，\CJKunderwave{太甲}三篇九，\CJKunderwave{咸有一德}十，\CJKunderwave{說命}三篇十三，\CJKunderwave{泰誓}三篇十六，\CJKunderwave{武成}十七，\CJKunderwave{旅獒}十八，\CJKunderwave{微子之命}十九，\CJKunderwave{蔡仲之命}二十，\CJKunderwave{周官}二十一，\CJKunderwave{君陳}二十二，\CJKunderwave{畢命}二十三，\CJKunderwave{君牙}二十四，\CJKunderwave{冏命}二十五。但\CJKunderline{孔君}所傳，值巫蠱不行以終。前漢諸儒知孔本有五十八篇，不見孔傳，遂有張霸之徒於鄭注之外偽造\CJKunderwave{尚書}凡二十四篇，以足鄭注三十四篇為五十八篇。其數雖與孔同，其篇有異。孔則於伏生所傳二十九篇內無古文\CJKunderwave{泰誓}，除\CJKunderwave{序}尚二十八篇，分出\CJKunderwave{舜典}、\CJKunderwave{益稷}、\CJKunderwave{盤庚}二篇、\CJKunderwave{康王之誥}為三十三，增二十五篇為五十八篇。\CJKunderline{鄭玄}則於伏生二十九篇之內分出\CJKunderwave{盤庚}二篇，\CJKunderwave{康王之誥}、又\CJKunderwave{泰誓}三篇，為三十四篇更增益偽書二十四篇為五十八。所增益二十四篇者，則鄭注\CJKunderwave{書序}，\CJKunderwave{舜典}一，\CJKunderwave{汩作}二，\CJKunderwave{九共}九篇十一，\CJKunderwave{大禹謨}十二，\CJKunderwave{益稷}十三，\CJKunderwave{五子之歌}十四，\CJKunderwave{胤徵}十五，\CJKunderwave{湯誥}十六，\CJKunderwave{咸有一德}十七，\CJKunderwave{典寶}十八，\CJKunderwave{伊訓}十九，\CJKunderwave{肆命}二十，\CJKunderwave{原命}二十一，\CJKunderwave{武成}二十二，\CJKunderwave{旅獒}二十三，\CJKunderwave{冏命}二十四。以此二十四為十六卷，以\CJKunderwave{九共}九篇共卷，除八篇，故為十六。故\CJKunderwave{藝文志}、劉向\CJKunderwave{別錄}雲“五十八篇”。\CJKunderwave{藝文志}又云：“孔安國者，\CJKunderline{孔子}後也。悉得其書，以古文又多十六篇。”篇即卷也。即是偽書二十四篇也。劉向作\CJKunderwave{別錄}，班固作\CJKunderwave{藝文志}並雲此言，不見孔傳也。劉歆作\CJKunderwave{三統曆}，論武王伐紂，引今文\CJKunderwave{泰誓}雲“丙午逮師”，又引\CJKunderwave{武成}“越若來三月五日甲子，咸劉商王受”，並不與孔同，亦不見孔傳也。後漢初賈逵\CJKunderwave{奏尚書疏}雲“流為烏”，是與孔亦異也。馬融\CJKunderwave{書序}云：“經傳所引\CJKunderwave{泰誓}，\CJKunderwave{泰誓}並無此文。”又云：“逸十六篇，絕無師說。”是融亦不見也。服虔、杜預注\CJKunderwave{左傳}“亂其紀綱”，並雲夏桀時，服虔、杜預皆不見也。\CJKunderline{鄭玄}亦不見之，故注\CJKunderwave{書序}、\CJKunderwave{舜典}雲“入麓伐木”，注\CJKunderwave{五子之歌}雲“避亂於洛汭”，注\CJKunderwave{胤徵}雲“胤徵，臣名”，又注\CJKunderwave{禹貢}引\CJKunderwave{胤徵}雲“厥匪玄黃，昭我周王”，又注\CJKunderwave{咸有一德}雲“伊陟臣扈曰”，又注\CJKunderwave{典寶}引\CJKunderwave{伊訓}雲“載孚在毫”，又曰“徵是三朡”，又注\CJKunderwave{旅獒}雲“獒讀曰豪，謂是遒豪之長”，又古文有\CJKunderwave{仲虺之誥}、\CJKunderwave{太甲}、\CJKunderwave{說命}等見在而云亡，其\CJKunderwave{汩典}、\CJKunderwave{典寶}之等一十三篇見亡而云已，逸是不見古文也。案伏生所傳三十四篇者謂之今文，則夏侯勝、夏侯建、歐陽和伯等三家所傳及後漢末蔡邕所勒石經是也。孔所傳者，膠東庸生、劉歆、賈逵、馬融等所傳是也。\CJKunderline{鄭玄}\CJKunderwave{書贊}云：“我先師棘子下生安國，亦好此學，衛、賈、馬二三君子之業，則雅才好博，既宣之矣。”又云：“歐陽氏失其本義，今疾此蔽冒，猶復疑惑未悛。”是鄭意師祖孔學，傳授膠東庸生、劉歆、賈逵、馬融等學，而賤夏侯、歐陽等；何意鄭注\CJKunderwave{尚書}，亡逸並與孔異，篇數並與三家同？又劉歆、賈逵、馬融之等並傳孔學，雲十六篇逸，與安國不同者，良由孔注之後，其書散逸，傳注不行。以庸生、賈、馬之等惟傳孔學經文三十三篇，故鄭與三家同，以為古文。而鄭承其後，所注皆同賈逵、馬融之學，題曰\CJKunderwave{古文尚書}，篇與夏侯等同，而經字多異。夏侯等書“宅嵎夷”為“宅嵎鐵”，“昧谷”，曰“柳谷”，“心腹腎腸”曰“憂腎陽”，“劓刵劅剠”雲“臏宮劓割頭庶剠”，是鄭注不同也。三家之學傳孔業者，\CJKunderwave{漢書·儒林傳}云，安國傳都尉朝子俊，俊傳膠東庸生，生傳清河胡常，常傳徐敖，敖傳王璜及塗惲，惲傳河南桑欽。至後漢初衛、賈、馬亦傳孔學，故\CJKunderwave{書贊}云：“自世祖興後漢，衛、賈、馬二三君子之業是也，所得傳者三十三篇古經，亦無其五十八篇，及傳說絕無傳者。”至晉世王肅注\CJKunderwave{書}，始似竊見孔傳，故注“亂其紀綱”為夏太康時。又\CJKunderwave{晉書·皇甫謐傳}云：“姑子外弟梁柳邊得\CJKunderwave{古文尚書}，故作\CJKunderwave{帝王世紀}，往往載孔傳五十八篇之書。”\CJKunderwave{晉書}又云：“晉太保公鄭衝以古文授扶風蘇愉，愉字休預。預授天水梁柳，字洪季，即謐之外弟也。季授城陽臧曹，字彥始。始授郡守子汝南梅賾，字仲真，又為豫章內史，遂於前晉奏上其書而施行焉。”時已亡失\CJKunderwave{舜典}一篇，晉末範寧為解時已不得焉。至齊蕭鸞建武四年姚方興於大航頭得而獻之，議者以為孔安國之所注也。值方興有罪，事亦隨寢。至隋開皇二年購慕遺典，乃得其篇焉。然孔注之後歷及後漢之末，無人傳說。至晉之初猶得存者，雖不列學官，散在民間，事雖久遠，故得猶存。 \par}

\CJKunderline{孔氏}傳\footnote{傳即注也，以傳述為義，舊說漢已前稱傳。}

{\noindent\shu\zihao{5}\fzkt 正義曰：以注者多門，故云其氏以別眾家。或當時自題\CJKunderline{孔氏}，亦可以後人辨之。 \par}

\textcolor{red}{昔在}帝堯，聰明文思，光宅天下。\footnote{\textcolor{red}{言聖德之遠著}。昔，古也。堯，唐帝名。馬融云：“諡也，翼善傳聖曰堯。”宅，拥有。聰，千公反。思，息嗣反，又如字，下同。著，張慮反。}將遜於位，讓於虞舜，\footnote{遜，遁也。\textcolor{red}{老使}攝，遂\textcolor{red}{禪之}。遁本作遯，徒遜反，退也，避也。“遂禪”音時戰反，讓也，授也。}作\CJKunderwave{\textcolor{red}{堯典}}。

{\noindent\zhuan\zihao{6}\fzbyks 傳“言聖德之遠著”。正義曰：“聖德”解“聰明文思”，“遠著”解“光宅天下”。 \par}

{\noindent\zhuan\zihao{6}\fzbyks 傳“老使”至“禪之”。正義曰：“老使攝”者解“將遜於位”，雲“遂禪之”者，解“讓於虞舜”也。以己年老，故遜之。使攝之，後功成而禪。禪即讓也。言“攝”者，“納於大麓”是也。禪者，“汝陟帝位”是也。雖舜受而攝之，而堯以為禪。或雲“汝陟帝位”為攝，因即直言為讓，故云“遂”也。\CJKunderline{鄭玄}云：“堯尊如故，舜攝其事”是也。 \par}

{\noindent\shu\zihao{5}\fzkt “昔在”至“堯典”。正義曰：此序\CJKunderline{鄭玄}、馬融、王肅並雲\CJKunderline{孔子}所作，孔義或然。\CJKunderwave{詩}、\CJKunderwave{書}理不應異，夫子為\CJKunderwave{書}作序，不作\CJKunderwave{詩}序者，此自或作或否，無義例也。鄭知\CJKunderline{孔子}作者，依緯文而知也。安國既以同序為卷，檢此百篇，凡有六十三序，序其九十六篇。\CJKunderwave{明居}、\CJKunderwave{咸有一德}、\CJKunderwave{立政}、\CJKunderwave{無逸}不序所由，直雲“咎單作\CJKunderwave{明居}”、“伊尹作\CJKunderwave{咸有一德}”、“周公作\CJKunderwave{立政}”、“周公作\CJKunderwave{無逸}”。六十三序者，若\CJKunderwave{汩作}、\CJKunderwave{九共}九篇，\CJKunderwave{槁飫}，十一篇共序；其\CJKunderwave{咸乂}四篇同序；其\CJKunderwave{大禹謨}、\CJKunderwave{皋陶謨}、\CJKunderwave{益稷}、\CJKunderwave{夏社}、\CJKunderwave{疑至}、\CJKunderwave{臣扈}、\CJKunderwave{伊訓}、\CJKunderwave{肆命}、\CJKunderwave{徂後}、\CJKunderwave{太甲}三篇、\CJKunderwave{盤庚}三篇、\CJKunderwave{說命}三篇、\CJKunderwave{泰誓}三篇、\CJKunderwave{康誥}、\CJKunderwave{酒誥}、\CJKunderwave{梓材}，二十四篇，皆三篇同序；其\CJKunderwave{帝告}、\CJKunderwave{釐沃}、\CJKunderwave{汝鳩}、\CJKunderwave{汝方}、\CJKunderwave{伊陟}、\CJKunderwave{原命}、\CJKunderwave{高宗肜日}、\CJKunderwave{高宗之訓}八篇皆共卷，類同，故同序。同序而別篇者三十三篇，通\CJKunderwave{明居}、\CJKunderwave{無逸}等四篇為三十七篇，加六十三即百篇也。序者以序別行辭為形勢。言昔日在於帝號堯之時也，此堯身智無不知聰也，神無不見明也。以此聰明之神智足可以經緯天地，即“文”也；又神智之運，深敏於機謀，即“思”也。“聰明文思”即其聖性行之於外，無不備知，故此德充滿居止於天下而遠著。德既如此，政化有成，天道衝盈，功成者退，以此故將遜遁避於帝位，以禪其有聖德之虞舜。史序其事，而作\CJKunderwave{堯典}之篇。言“昔在”者，\CJKunderline{鄭玄}云：“\CJKunderwave{書}以堯為始，獨雲“昔在”，使若無先之典然也。”\CJKunderwave{詩}云：“自古在昔。”言“在昔”者，自下本上之辭。言“昔在”者，從上自下為稱，故曰“使若無先之”者。據代有先之，而書無所先，故云“昔”也。言“帝”者，天之一名，所以名“帝”。帝者，諦也。言天蕩然無心，忘於物我，言公平通遠，舉事審諦，故謂之“帝”也。五帝道同於此，亦能審諦，故取其名。若然，聖人皆能同天，故曰“大人”。大人者與天地合其德，即三王亦大人。不得稱帝者，以三王雖實聖人，內德同天，而外隨時運，不得盡其聖，用逐跡為名，故謂之為王。\CJKunderwave{禮運}曰，“大道之行，天下為公”，即帝也。“大道既隱，各親其親”，即王也。則聖德無大於天，三皇優於帝，豈過乎天哉！然則三皇亦不能過天，但遂同天之名，以為優劣。五帝有為而同天，三皇無為而同天，立名以為優劣耳。但有為無為亦逐多少以為分，三王亦順帝之則而不盡，故不得名帝。然天之與帝，義為一也。人主可得稱帝，不可得稱天者，以天隨體而立名，人主不可同天之體也。無由稱天者，以天德立號，王者可以同其德焉，所以可稱於帝。故繼天則謂之“天子”，其號謂之“帝”，不得雲“帝子”也。言“堯”者，孔無明解。案下傳云：“虞，氏。舜，名。”然堯、舜相配為義，既舜為名，則堯亦名也。以此而言，禹、湯亦名。於下都無所解，而放勳、重華、文命注隨其事而解其文以為義，不為堯、舜及禹之名。據此，似堯、舜及禹與湯相類，名則俱名，不應殊異。案鄭以下亦云：“虞，氏。舜，名。”與孔傳不殊。及鄭注\CJKunderwave{中侯}，云：“重華，舜名。”則舜不得有二名。鄭注\CJKunderwave{禮記}云：“舜之言充。”是以舜為號諡之名，則下注雲“舜，名，亦號諡之名”也。推此則\CJKunderline{孔君}亦然。何以知之？既湯類堯、舜當為名，而孔注\CJKunderwave{論語}“曰予小子履”云，“履是殷湯名堯舜，是湯名履，而湯非名也。又此不雲堯、舜是名，則堯及舜、禹非名，於是明矣。既非名，而放勳、重華、文命蓋以為三王之名，同於\CJKunderline{鄭玄}矣。鄭知名者，以\CJKunderwave{帝系}雲“禹名文命”，以上類之亦名。若然，名本題情記意，必有義者，蓋運命相符，名與運接，所以異於凡平。或說以其有義，皆以為字。古代尚質，若名之不顯，何以著字？必不獲已，以為非名非字可也。譙周以堯為號，皇甫謐以放勳、重華、文命為名。案\CJKunderwave{諡法}“翼善傳聖曰堯，仁義盛明曰舜”，是堯、舜諡也。故馬融亦云諡也。又曰“淵源流通曰禹，雲行雨施曰湯”，則禹、湯亦是諡法。而馬融云：“禹湯不在\CJKunderwave{諡法}。”故疑之。將由\CJKunderwave{諡法}或本不同，故有致異。亦可本無禹、湯為諡，後來所加，故或本曰“除虐去殘曰湯”，是以異也。\CJKunderwave{檀弓}曰：“死諡，周道也。”\CJKunderwave{周書}諡法周公所作，而得有堯、舜、禹、湯者，以周法死後乃追，故謂之為諡。諡者，累也，累其行而號也。隨其行以名之，則死諡猶生號。因上世之生號陳之為死諡，明上代生死同稱。上世質，非至善至惡無號，故與周異。以此堯、舜或雲號，或雲諡也。若然，湯名履而王侯，\CJKunderwave{世本}“湯名天乙”者，安國意蓋以湯受命之王，依殷法以乙日生，名天乙。至將為王，又改名為履，故二名也。亦可。安國不信\CJKunderwave{世本}，無天乙之名。皇甫謐巧欲傅會，雲“以乙日生，故名履，字天乙”。又云祖乙，亦云“乙日生，複名乙”，引\CJKunderwave{易緯}“\CJKunderline{孔子}所謂天之錫命，故可同名”。既以天乙為字，何雲同名乎？斯又妄矣，號之曰堯者，\CJKunderwave{釋名}以為“其尊高堯堯然，物莫之先，故謂之堯也”。\CJKunderwave{諡法}云：“翼善傳聖曰堯。”堯者以天下之生善，因善欲禪之，故二八顯升，所謂為翼。能傳位於聖人，天下為公，此所以出眾而高也。言“聰明者，據人近驗，則聽遠為聰，見微為明，若離婁之視明也，師曠之聽聰也；以耳目之聞見，喻聖人之智慧，兼知天下之事，故在於聞見而已，故以“聰明”言之。智之所用，用於天地，經緯天地謂之“文”，故以聰明之用為文。須當其理，故又云“思”而會理也。經雲“欽明”，此為聰明者，彼方陳行事，故美其敬，此序其聖性，故稱其“聰”，隨事而變文。下\CJKunderwave{舜典}直雲“堯聞之聰明”，不雲“文思”者，此將言堯用，故云“文思”，彼要雲舜德，故直雲“聰明”，亦自此而可知也。言“光宅”者，經傳云：“光，充也。”不訓“宅”者，可知也。不於此訓“光”者，從經為正也。下“將遜於位”傳雲“遜，遁”者，以經無“遜”字，故在序訓之。 \par}

\textcolor{red}{堯典}\footnote{言堯可為百代常行之道。}

{\noindent\shu\zihao{5}\fzkt “堯典”。正義曰：序已雲“作堯典”而重言此者，此是經之篇目，不可因序有名，略其舊題，故諸篇皆重言本目而就目解之。稱“典”者，以道可百代常行。若堯、舜禪讓聖賢，禹、湯傳授子孫，即是堯、舜之道不可常行，但惟德是與，非賢不授。授賢之事，道可常行，但后王德劣不能及古耳。然經之與典俱訓為常，名典不名經者，以經是總名，包殷、周以上，皆可為後代常法，故以經為名。典者，經中之別，特指堯、舜之德，於常行之內道最為優，故名典不名經也。其太宰六典及司寇三典者，自由當代常行，與此別矣。 \par}

\textcolor{red}{曰若}稽古，帝堯\footnote{\textcolor{red}{若，順}。稽，考也。能順考古道而行之者\textcolor{red}{帝堯}。}曰放勳，欽明文思安安\footnote{\textcolor{red}{勳，功}。欽，敬也。言堯放上世之功，化而以敬明文思之四德，安天下之當\textcolor{red}{安者}。放，方往反，注同。徐云：“鄭、王如字。”勳，許雲反，功也。馬云：“放勳，堯名。”皇甫謐同。一云：“放勳，堯字。”“欽明文思”，馬云：“威儀表備之欽，照臨四方謂之明，經緯天地謂之文，道德純備謂之思。”}，允恭克讓，光被四表，格於\textcolor{red}{上下}。\footnote{\textcolor{red}{允，信}。克，能。光，充。格，至也。既有四德，又信恭能讓，故其名聞充溢四外，至於\textcolor{red}{天地}。被，皮寄反，徐扶義反。聞音問，本亦作問。溢音逸。}

{\noindent\zhuan\zihao{6}\fzbyks 傳“若順”至“帝堯”。正義曰：“若、順”，\CJKunderwave{釋言}文。\CJKunderwave{詩}稱“考卜惟王”，\CJKunderwave{洪範}考卜之事謂之“稽疑”，是“稽”為考，經傳常訓也。\CJKunderwave{爾雅}一訓一也，孔所以約文，故數字俱訓，其末以一“也”結之。又已經訓者，後傳多不重訓。顯見可知，則徑言其義，皆務在省文故也。言“順考古道”者，古人之道非無得失，施之當時又有可否，考其事之是非，知其宜於今世，乃順而行之。言其行可否，順是不順非也。考“古”者自己之前，無遠近之限，但事有可取，皆考而順之。今古既異時，政必殊古，事雖不得盡行，又不可頓除古法，故\CJKunderwave{說命}曰：“事不師古，以克永世，匪說攸聞。”是後世為治當師古法，雖則聖人，必須順古。若空欲追遠，不知考擇，居今行古，更致禍災。若宋襄慕義，師敗身傷，徐偃行仁，國亡家滅，斯乃不考之失。故美其能順考也。\CJKunderline{鄭玄}信緯，訓“稽”為同，訓“古”為天，言“能順天而行之，與之同功”。\CJKunderwave{論語}稱惟堯則天，\CJKunderwave{詩}美文王“順帝之則”，然則聖人之道莫不同天合德，豈待同天之語，然後得同之哉？\CJKunderwave{書}為世教，當因之人事，以人系天，於義無取，且“古”之為天，經無此訓。高貴鄉公皆以鄭為長，非篤論也。 \par}

{\noindent\zhuan\zihao{6}\fzbyks 傳“勳功”至“安者”。正義曰：“勳、功”、“欽、敬”，\CJKunderwave{釋詁}文。此經述上稽古之事，放效上世之功，即是考於古道也。經言“放勳”，放其功而已。傳兼言“化”者，據其勳業謂之功，指其教人則為化，功之與化所從言之異耳。\CJKunderline{鄭玄}云：“敬事節用謂之欽，照臨四方謂之明，經緯天地謂之文，慮深通敏謂之思。”孔無明說，當與之同。四者皆在身之德，故謂之“四德”。凡是臣人王者皆須安之，故廣言“安天下之當安者”。所安者則下文“九族、百姓、萬邦”是也。其“敬明文思”為此次者，顧氏云：“隨便而言，無義例也。”知者此先“聰”後“明”，\CJKunderwave{舜典}雲“明四目，達四聰”，先明後“聰”，故知無例也。今考\CJKunderwave{舜典}云：“浚哲文明”，又先“文”後“明”，與此不類，知顧氏為得也。 \par}

{\noindent\zhuan\zihao{6}\fzbyks 傳“允信”至“天地”。正義曰：“允、信”、“格、至”，\CJKunderwave{釋詁}文。“克，能”、“光，充”，\CJKunderwave{釋言}文。在身為德，施之曰行。\CJKunderline{鄭玄}云：“不懈於位曰恭，推賢尚善曰讓。”恭讓是施行之名。上言堯德，此言堯行，故傳以文次言之。言堯既有敬明文思之四德，又信實、恭勤、善能、推讓，下人愛其恭讓，傳其德音，故其名遠聞，旁行則充溢四方，上下則至於天地。持身能恭，與人能讓，自己及物，故先恭後讓。恭言信，讓言克，交互其文耳。皆言信實能為也。傳以“溢”解“被”，言其饒多盈溢，故被及之也。表裡內外相對之言，故以表為外，向不向上至有所限，旁行四方無復限極，故四表言“被”，上下言“至”。“四外”者，以其無限，自內言之，言其至於遠處，正謂四方之外畔者，當如\CJKunderwave{爾雅}所謂“四海”、“四荒”之地也。先“四表”後“上下”者，人之聲名，宜先及於人，後被四表，是人先知之，故先言至人。後言至於上下，言至於天地，喻其聲聞遠耳。\CJKunderwave{禮運}稱聖人為政，能使“天降膏露，地出醴泉”，是名聞遠達，使天地效靈，是亦格於上下之事。 \par}

{\noindent\shu\zihao{5}\fzkt “曰若”至“上下”。正義曰：史將述堯之美，故為題目之辭曰，能順考校古道而行之者，是帝堯也。又申其順考古道之事曰，此帝堯能放效上世之功而施其教化，心意恆敬，智慧甚明，發舉則有文謀，思慮則能通敏，以此四德安天下之當安者。在於己身則有此四德，其於外接物又能信實、恭勤、善能、謙讓。恭則人不敢侮，讓則人莫與爭，由此為下所服，名譽著聞，聖德美名充滿被溢於四方之外，又至於上天下地。言其日月所照，霜露所墜，莫不聞其聲名，被其恩澤。此即稽古之事也。 \par}


\textcolor{red}{克明}俊德，以親九族\footnote{\textcolor{red}{能明}俊德之士任用之，以睦高祖玄孫\textcolor{red}{之親}。九族，上自高祖，下至玄孫，凡九族。馬、鄭同。}。九族既睦，平章百姓\footnote{\textcolor{red}{既，已}也。百姓，百官。言化九族而平和\textcolor{red}{章明}。}。百姓昭明，協和萬邦。黎民於變\textcolor{red}{時雍}\footnote{\textcolor{red}{昭亦}明也。協，合。黎，眾。時，是。雍，和也。言天下眾民皆變化化上，是以風俗\textcolor{red}{大和}。黎，力兮反。}。

{\noindent\zhuan\zihao{6}\fzbyks 傳“能明”至“之親”。正義曰：\CJKunderline{鄭玄}云：“‘俊德’，賢才兼人者。”然則“俊德”謂有德。人能明俊德之士者，謂命為大官，賜之厚祿，用其才智，使之高顯也。以其有德，故任用之。以此賢臣之化，親睦高祖玄孫之親。上至高祖，下及玄孫，是為九族。同出高曾，皆當親之，故言之“親”也。\CJKunderwave{禮記·喪服小記}云：“親親以三為五，以五為九。”又\CJKunderwave{異義}、夏侯、歐陽等以為九族者，父族四、母族三、妻族二，皆據“異姓有服”。\CJKunderline{鄭玄}駁云：“異姓之服不過緦麻，言不廢昏。又\CJKunderwave{昏禮}請期雲‘惟是三族之不虞’，恐其廢昏，明非外族也。”是鄭與孔同。“九族”謂帝之九族，“百姓”謂百官族姓，“萬邦”謂天下眾民，自內及外，從高至卑，以為遠近之次也。知“九族”非民之九族者，以先親九族，次及百姓，百姓是群臣弟子，不宜越百姓而先下民。若是民之九族，則“九族既睦”，民已和矣，下句不當復言“協和萬邦”，以此知帝之九族也。堯不自親九族，而待臣使之親者，此言用臣法耳，豈有聖人在上，疏其骨肉者乎？若以堯自能親，不待臣化，則化萬邦百姓，堯豈不能化之，而待臣化之也？且言“親九族”者，非徒使帝親之，亦使臣親之，帝亦令其自相親愛，故須臣子之化也。 \par}

{\noindent\zhuan\zihao{6}\fzbyks 傳“既，已”至“章明”。正義曰：“既”、“已”義同，故訓“既”為已，經傳之言。“百姓”或指天下百姓，此下句乃有“黎民”，故知“百姓”即百官也。百官謂之百姓者，隱八年\CJKunderwave{左傳}云：“天子建德，因生以賜姓。”謂建立有德以為公卿，因其所生之地而賜之以為其姓，令其收斂族親，自為宗主。明王者任賢不任親，故以“百姓”言之。\CJKunderwave{周官}篇云：“唐、虞稽古，建官惟百。”\CJKunderwave{大禹謨}云：“率百官若帝之初。”是唐、虞之世經文皆稱“百官”。而\CJKunderwave{禮記·明堂位}雲“有虞氏之官五十”，後世所記不合經也。“平章”與“百姓”其文非九族之事，傳以此經之事文勢相因，先化九族，乃化百官，故云“化九族而平和章明”。謂九族與百官皆須導之以德義，平理之使之協和，教之以禮法，章顯之使之明著。 \par}

{\noindent\zhuan\zihao{6}\fzbyks 傳“昭亦”至“大和”。正義曰：\CJKunderwave{釋詁}以“昭”為光，光、明義同，經已有“明”，故云“昭亦明也”。\CJKunderwave{釋詁}以“協”為和，和、合義同，故訓“協”為合也。“黎，眾”、“時，是”，\CJKunderwave{釋詁}文。“雍，和”，\CJKunderwave{釋訓}文。堯民之變，明其變惡從善，人之所和，惟風俗耳。故知謂“天下眾人皆變化化上，是以風俗大和”，人俗大和，即是太平之事也。此上經三事相類，古史交互立文。以“親”言“既睦”，“平章”言“昭明”，“協和”言“時雍”。“睦”即“親”也，“章”即“明”也，“雍”即“和”也，各自變文以類相對。平九族使之親，平百姓使之明，正謂使從順禮義，恩情和合，故於萬邦變言“協和”，明“以親九族”、“平章百姓”亦是協和之也。但九族宜相親睦，百姓宜明禮義，萬邦宜盡和協，各因所宜為文，其實相通也。民言“於變”，謂從上化，則“九族既睦”、“百姓昭明”亦是變上，故得睦得明也。 \par}

{\noindent\shu\zihao{5}\fzkt “克明”至“時雍”。正義曰：言堯能名聞廣遠，由其委任賢哲，故復陳之。言堯之為君也，能尊明俊德之士，使之助己施化。以此賢臣之化，先令親其九族之親。九族蒙化已親睦矣，又使之和協顯明於百官之族姓。百姓蒙化皆有禮儀，昭然而明顯矣，又使之合會調和天下之萬國。其萬國之眾人於是變化從上，是以風俗大和，能使九族敦睦，百姓顯明，萬邦和睦，是“安天下之當安”者也。 \par}

\textcolor{red}{乃命}羲、和，欽若昊天，曆象日月星辰，敬授人時\footnote{\textcolor{red}{重黎}之後羲氏、和氏世掌天地四時之官，故堯命之，使敬順昊天。昊天言元氣廣大。星，四方中星。辰，日月所會。曆象其分節。敬記天時以授人也。此舉其目，下別\textcolor{red}{序之}。羲和，馬云：“羲氏掌天官，和氏掌地官，四子掌四時。”昊，胡老反。重，直龍反，少昊之後。黎，高陽之後。“日月所會”謂日月交會於十二次也。寅曰析木，卯曰大火，辰曰壽星，巳曰鶉尾，午曰鶉火，未曰鶉首，申曰實沈，酉曰大梁，戌曰降婁，亥曰娵訾，子曰玄枵，醜曰星紀。}。

{\noindent\zhuan\zihao{6}\fzbyks 傳“重黎”至“序之”。正義曰：\CJKunderwave{楚語}云：“少昊氏之衰，九黎亂德，人神雜擾，不可方物。顓頊受之，乃命南正重司天以屬神，火正黎司地以屬民，使復舊常，無相侵瀆。其後三苗復九黎之惡，堯覆育重黎之後，不忘舊者，使復典之。以至於夏商。”據此文則自堯及商無他姓也。堯育重黎之後，是此羲和可知。是羲和為重黎之後，世掌天地之官文所出也。\CJKunderwave{呂刑}先“重”後“黎”，此文先“羲”後“和”，楊子\CJKunderwave{法言}云：“羲近重，和近黎。”是“羲”承“重”而“和”承“黎”矣。\CJKunderwave{呂刑}稱“乃命重黎”與此“命羲和”為一事也。故\CJKunderwave{呂刑}傳云：“重即羲也，黎即和也。”羲和雖別為氏族而出自重黎，故\CJKunderwave{呂刑}以“重黎”言之。\CJKunderwave{鄭語}云：“為高辛氏火正。”則高辛亦命重黎。故\CJKunderline{鄭玄}於此注云：“高辛氏世，命重為南正司天，黎為火正司地。”據“世掌”之文，用\CJKunderwave{楚語}為說也。\CJKunderwave{楚世家}云：“重黎為帝嚳火正，能光融天下，啻嚳命曰祝融。共工氏作亂，帝嚳使重黎誅之而不盡。帝乃以庚寅日誅重黎，而以其弟吳回為重黎，復居火正，為祝融。”案昭二十九年\CJKunderwave{左傳}稱少昊氏有子曰重，顓頊氏有子曰黎。則重黎二人，各出一帝。而\CJKunderwave{史記}並以重黎為楚國之祖，吳回為重黎，以重黎為官號，此乃\CJKunderwave{史記}之謬。故束晳譏馬遷並兩人以為一，謂此是也。\CJKunderwave{左傳}稱重為句芒，黎為祝融，不言何帝使為此官。但黎是顓頊之子，其為祝融，必在顓頊之世。重雖少昊之胤，而與黎同命，明使重為句芒亦是顓頊時也。祝融火官可得稱為火正，句芒木官不應號為南正，且木不主天，火不主地，而\CJKunderwave{外傳}稱顓頊命南正司天，火正司地者，蓋使木官兼掌天，火官兼掌地。南為陽位，故掌天謂之南正。黎稱本官，故掌地猶為火正。鄭答趙商云：“先師以來，皆雲火掌為地，當雲黎為北正。”孔無明說，未必然也。昭十七年\CJKunderwave{左傳}郯子稱少昊氏以鳥名官，自顓頊已來乃命以民事。句芒、祝融皆以人事名官，明此當顓頊之時也。傳言少昊氏有四叔，當為後代子孫，非親子也。何則？傳稱共工氏有子曰句龍，共工氏在顓頊之前多歷年代，豈復共工氏親子至顓頊時乎？明知少昊四叔亦非親子，高辛所命重黎或是重黎子孫，未必一人能歷二代。又高辛前命後誅，當是異人。何有罪而誅，不容列在祀典。明是重黎之後，世以重黎為號，所誅重黎是有功重黎之子孫也。\CJKunderwave{呂刑}說羲和之事，猶尚謂之重黎，況彼尚近重黎，何故不得稱之？以此知異世重黎號同人別。顓頊命重司天，黎司地，羲氏掌天，和氏掌地，其實重、黎、羲、和通掌之也。此雲“乃命羲和，欽若昊天”，是羲和二氏共掌天地之事。以乾坤相配，天地相成，運立施化者天，資生成物者地，天之功成其見在地，故下言“日中，星鳥”之類是天事也，“平秩東作”之類是地事也，各分掌其時，非別職矣。案\CJKunderwave{楚語}云，重司天以屬神，黎言地以屬人。天地既別，人神又殊，而云通掌之者，外傳之文說\CJKunderwave{呂刑}之義，以為少昊之衰，天地相通，人神雜擾，顓頊乃命重黎分而異之，以解絕地天通之言，故云各有所掌。天地相通，人神雜擾，見其能離絕天地，變異人神耳，非即別掌之。下文別序所掌，則羲主春夏，和主秋冬，俱掌天時，明其共職。彼又言：“至於夏商，世掌天地。”\CJKunderwave{胤徵}云：“羲和湎淫，廢時亂日。”不知日食，羲和同罪，明其世掌天地共職。可知顓頊命掌天地，惟重黎二人，堯命羲和則仲叔四人者，以羲和二氏賢者既多，且後代稍文，故分掌其職事，四人各職一時，兼職方岳，以有四嶽，故用四人。顓頊之命重黎，惟司天地，主嶽以否不可得知。設令亦主方岳，蓋重黎二人分主東西也。馬融、\CJKunderline{鄭玄}皆以此“命羲和”者，命為天地之官。下雲“分命”，申命為四時之職。天地之與四時於周則冢宰司徒之屬，六卿是也。孔言“此舉其目，下別序之”，則惟命四人，無六官也。下傳雲四嶽即羲和四子，\CJKunderwave{舜典}傳稱禹、益六人新命有職，與四嶽十二牧凡為二十二人。然新命之六人，禹命為百揆，契作司徒，伯夷為秩宗，皋陶為士，垂作共工，亦禹、契之輩即是卿官。卿官之外別有四嶽，四嶽非卿官也。孔意以羲和非是卿官，別掌天地，但天地行於四時，四時位在四方，平秩四時之人因主方岳之事，猶自別有卿官分掌諸職。\CJKunderwave{左傳}稱少昊氏以鳥名官，五鳩氏即周世之卿官也。五鳩之外別有鳳鳥氏，歷正也，班在五鳩之上。是上代以來皆重歷數，故知堯於卿官之外別命羲和掌天地也。於時羲和似尊於諸卿，後世以來稍益卑賤。\CJKunderwave{周禮}“太史掌正歲年以序事”，即古羲和之任也。桓十七年\CJKunderwave{左傳}雲“日官居卿以底日”，猶尚尊其所掌。周之卿官明是堯時重之，故特言“乃命羲和”。此“乃命羲和”重述“克明俊德”之事，得致雍和所由。已上論堯聖性，此說堯之任賢，據堯身而言用臣，故云“乃命”，非“時雍”之後方始命之。“使敬順昊天”，昊天者混元之氣，昊然廣大，故謂之“昊天”也。\CJKunderwave{釋天}云：“春為蒼天，夏為昊天，秋為旻天，冬為上天。”\CJKunderwave{毛詩}傳云：“尊而君之則稱皇天，元氣廣大則稱昊天，仁覆閔下則稱旻天，自上降監則稱上天，據遠視之蒼蒼然則稱蒼天。”\CJKunderwave{爾雅}四時異名，\CJKunderwave{詩}傳即隨事立稱。\CJKunderline{鄭玄}讀\CJKunderwave{爾雅}雲“春為昊天，夏為蒼天”，故駁\CJKunderwave{異義}云：“春氣博施，故以廣大言之。夏氣高明，故以遠言之。秋氣或生或殺，故以閔下言之。冬氣閉藏而清察，故以監下言之。皇天者尊而號之也。”六籍之中，諸稱天者以情所求言之耳，非必於其時稱之。然此言堯敬大四天，故以“廣大”言之。“星，四方中星”者，二十八宿，布在四方，隨天轉運，更互在南方，每月各有中者。\CJKunderwave{月令}每月昏旦，惟舉一星之中，若使每日視之，即諸宿每日昏旦莫不常中，中則人皆見之，故以中星表宿，“四方中星”總謂二十八宿也。或以\CJKunderwave{書傳}雲“主春者張，昏中，可以種穀。主夏者火，昏中，可以種黍。主秋者虛，昏中，可以種麥。主冬者昴，昏中，可以收斂。皆雲上告天子，下賦臣人。天子南面而視四方星之中，知人緩急，故曰敬授人時”，謂此“四方中星”如\CJKunderwave{書}傳之說。孔於虛昴諸星本無取中之事，用\CJKunderwave{書傳}為孔說非其旨矣。“辰，日月所會”者，昭七年\CJKunderwave{左傳}士文伯對晉侯之辭也。日行遲，月行疾，每月之朔月行及日而與之會，其必在宿。分二十八宿，是日月所會之處。辰，時也，集會有時，故謂之辰。“日月所會”與“四方中星”俱是二十八宿。舉其人目所見，以星言之。論其日月所會，以辰言之，其實一物，故星、辰共文。\CJKunderwave{益稷}稱古人之象，日月星辰共為一象，由其實同故也。日月與星，天之三光。四時變化，以此為政。故命羲和，令以算術推步，累歷其所行，法象其所在，具有分數節候，參差不等，敬記此天時以為歷而授人。此言星辰共為一物。\CJKunderwave{周禮·大宗伯}云：“實柴祀日月星辰。”\CJKunderline{鄭玄}雲“星謂五緯，辰謂日月所會十二次”者，以星、辰為二者。五緯與二十八宿俱是天星，天之神祇，禮無不祭，故\CJKunderline{鄭玄}隨事而分之。以此“敬授人時”無取五緯之義，故\CJKunderline{鄭玄}於此注亦以星、辰為一，觀文為說也。然則五星與日月皆別行，不與二十八宿同為不動也。 \par}

分命羲仲，宅嵎夷，曰暘谷\footnote{\textcolor{red}{宅，居}也。東表之地稱嵎夷。暘，明也。日出於谷而天下明，故稱暘谷。暘谷、嵎夷一也。羲仲居治東方\textcolor{red}{之官}。嵎音隅，馬曰：“嵎，海嵎也。”夷，萊夷也。\CJKunderwave{尚書、·考靈耀}及\CJKunderwave{史記}作“禺銕”。暘音陽。谷，工木反，又音欲，下同。馬云：“暘谷，海嵎夷之地名。”“日出於谷”本或作“日出於陽穀”，“陽”衍字。}。寅賓出日，平秩東作\footnote{\textcolor{red}{寅，敬}。賓，導。秩，序也。歲起於東而始就耕，謂之東作。東方之官敬導出日，平均次序東作之事，以\textcolor{red}{務農}也。寅，徐以真反，又音夷，下同。賓如字，徐音殯，馬云：“從也。”出，尺遂反，又如字，注同。平如字，馬作蘋，普庚反，云：“使也。”下皆放此。秩如字。}。日中，星鳥，以殷仲春\footnote{\textcolor{red}{日中}謂春分之日。鳥，南方朱鳥七宿。殷，正也。春分之昏，鳥星畢見，以正仲春之氣節，轉以推季孟則\textcolor{red}{可知}。中，貞仲反，又如字。殷，於勤反，馬、鄭云：“中也。”宿音秀，下同。見，賢遍反，下同。}。厥民析，鳥獸孳尾\footnote{\textcolor{red}{冬寒}無事，併入室處。春事既起，丁壯就功。厥，其也。言其民老壯分析。乳化曰孳，交接\textcolor{red}{曰尾}。析，星曆反。孳音字。乳，儒付反。\CJKunderwave{說文}云：“人及鳥生子曰乳，獸曰產。”}。

{\noindent\zhuan\zihao{6}\fzbyks 傳“宅居”至“之官”。正義曰：“宅，居”，\CJKunderwave{釋言}文。\CJKunderwave{禹貢}青州云：“嵎夷既略。”青州在東界外之畔為表，故云“東表之地稱嵎夷”也。陰陽相對，陰暗而陽明也，故以“暘”為明。谷無陰陽之異，以日出於谷而天下皆明，故謂日出之處為“暘谷”。冬南夏北不常厥處，但日由空道，似行自谷，故以“谷”言之，非實有深谷而日從谷之出也。據日所出謂之“暘谷”，指其地名即稱“嵎夷”，故云“暘谷、嵎夷一也”。又解“居”者，居其官不居其地，故云“羲仲居治東方之官”。此言“分命”者，上雲“乃命羲和”，總舉其目，就“乃命”之內分其職掌，使羲主春夏，和主秋冬，分一歲而別掌之，故言“分命”。就羲和之內又重分之，故於夏變言“申命”。既命仲而覆命叔，是其重命之也。所命無伯、季者，蓋時無伯、季，或有而不賢，則\CJKunderwave{外傳}稱“堯育重黎之後，不忘舊者，使復典之”，明仲叔能守舊業，故命之也。此羲和掌序天地，兼知人事，因主四時而分主四方，故舉東表之地，以明所舉之域。地東舉嵎夷之名，明分三方皆宜有地名，此為其始，故特詳舉其文。羲仲居治東方之官，居在帝都而遙統領之。王肅雲“皆居京師而統之，亦有時述職”，是其事也。以春位在東，因治於東方，其實本主四方春政，故於和仲之下云：“此居治西方之官，掌秋天之政。”明此，掌春天之政，孔以經事詳，故就下文而互發之。 \par}

{\noindent\zhuan\zihao{6}\fzbyks 傳“寅敬”至“務農”。正義曰：“寅，敬”，\CJKunderwave{釋詁}文。賓者主行導引，故“賓”為導也。\CJKunderwave{釋詁}以“秩”為常，常即次第有序，故“秩”為序也。一歲之事，在東則耕作，在南則化育，在西則成熟，在北則改易，故以方名配歲事為文，言順天時氣以勸課人務也。春則生物，秋則成物。日之出也，物始生長，人當順其生長，致力耕耘。日之入也，物皆成熟，人當順其成熟，致力收斂。東方之官當恭敬導引日出，平秩東作之事，使人耕耘。西方之官當恭敬從送日入，平秩西成之事，使人收斂。日之出入，自是其常，但由日出入，故物有生成。雖氣能生物，而非人不就。勤於耕稼，是導引之。勤於收藏，是從送之。冬夏之文無此類者，南北二方非日所出入，“平秩南訛”亦是導日之事，“平在朔易”亦是送日之事。依此春秋而共為賓餞，故冬夏二時無此一句。勸課下民，皆使致力，是敬導之。平均次序，即是授人，田裡各有疆埸，是平均之也。耕種收斂使不失其次序，王者以農為重，經主於農事。“寅賓出日”為“平秩”設文，故並解之也。言“敬導出日”者，正謂平秩次序東作之事以務農也。鄭以“作”為生，計秋言西成，春宜言東生。但四時之功皆須作力，不可不言力作，直說生成，明此以歲事初起，時言“東作”，以見四時亦當力作，故孔以耕作解之。\CJKunderline{鄭玄}云：“寅賓出日，謂春分朝日。”又以“寅餞納日，謂秋分夕日”也。 \par}

{\noindent\zhuan\zihao{6}\fzbyks 傳“日中”至“可知”。正義曰：其仲春、仲秋、冬至、夏至，馬融云：“古制刻漏晝夜百刻。晝長六十刻，夜短四十刻。晝短四十刻，夜長六十刻。晝中五十刻，夜亦五十刻。”融之此言據日出見為說。天之晝夜以日出入為分，人之晝夜以昏明為限。日未出前二刻半為明，日入後二刻半為昏，損夜五刻以裨於晝，則晝多於夜，覆校五刻。古今歷術與太史所候皆云，夏至之晝六十五刻，夜三十五刻。冬至之晝四十五刻，夜五十五刻。春分秋分之晝五十五刻，夜四十五刻。此其不易之法也。然今太史細候之法，則校常法半刻也。從春分至於夏至，晝暫長，增九刻半。夏至至於秋分，所減亦如之。從秋分至於冬至，晝暫短，減十刻半。從冬至至於春分，其增亦如之。又於每氣之間增減刻數，有多有少，不可通而為率。漢初未能審知，率九日增減一刻，和帝時待詔霍融始請改之。鄭注\CJKunderwave{書緯·考靈曜}仍雲“九日增減一刻”，猶尚未覺悟也。鄭注此云：“日長者日見之漏五十五刻，日短者日見之漏四十五刻。”與歷不同。故王肅難云：“知日見之漏減晝漏五刻，不意馬融為傳已減之矣。因馬融所減而又減之，故日長為五十五刻，因以冬至反之，取其夏至夜刻，以為冬至晝短，此其所以誤耳。”“鳥，南方朱鳥七宿”者，在天成象，星作鳥形。\CJKunderwave{曲禮}說軍陳象天之行，“前朱雀，後玄武，左青龍，右白虎”。“雀”即鳥也。“武”謂龜甲捍禦，故變文“玄武”焉。是天星有龍虎鳥龜之形也。四方皆有七宿，各成一形。東方成龍形，西方成虎形，皆南首而北尾。南方成鳥形，北方成龜形，皆西首而東尾。以南方之宿象鳥，故言鳥謂朱鳥七宿也。此經舉宿，為文不類。春言“星鳥”，總舉七宿。夏言“星火”，獨指房、心。虛、昴惟舉一宿。文不同者，互相通也。\CJKunderwave{釋言}以“殷”為中，中、正義同，故“殷”為正也。此經冬夏言“正”，春秋言“殷”者，其義同。春分之昏，觀鳥星畢見，以正仲春之氣節，計仲春日在奎、婁而入於酉地，則初昏之時井、鬼在午，柳、星、張在巳，軫、翼在辰，是朱鳥七宿皆得見也。春有三月，此經直雲“仲春”，故傳辨之云，既正仲春，轉以推季孟之月，則事亦可知也。天道左旋，日體右行，故星見之方與四時相逆。春則南方見，夏則東方見，秋則北方見，冬則西方見，此則勢自當然。而\CJKunderwave{書緯}為文生說，言“春夏相與交，秋冬相與互，謂之母成子，子助母”，斯假妄之談耳。馬融、\CJKunderline{鄭玄}以為“星鳥、星火謂正在南方。春分之昏七星中，仲夏之昏心星中，秋分之昏虛星中，冬至之昏昴星中，皆舉正中之星，不為一方盡”，見此其與孔異也。至於舉仲月以統一時，亦與孔同。王肅亦以星鳥之屬為昏中之星，其要異者以所宅為孟月，日中、日永為仲月，星鳥、星火為季月，“以殷”、“以正”皆總三時之月，讀“仲”為中，言各正三月之中氣也。以馬融、\CJKunderline{鄭玄}之言，不合天象，星火之屬仲月未中，故為每時皆歷陳三月，言日以正仲春，以正春之三月中氣。若正春之三月中，當言“以正春中”，不應言“以正仲春”。王氏之說非文勢也。\CJKunderline{孔氏}直取“畢見”，稍為迂闊，比諸王、馬，於理最優。 \par}

{\noindent\zhuan\zihao{6}\fzbyks 傳“冬寒”至“曰尾”。正義曰：“厥，其”，\CJKunderwave{釋言}文。其人老弱在室，丁壯適野，是老壯分析也。孳、字，古今同耳。字訓愛也，產生為乳，胎孕為化，孕產必愛之，故乳化曰“孳”。鳥獸皆以尾交接，故交接曰“尾”。計當先尾後孳，隨便言之。 \par}

申命羲叔，宅南交\footnote{\textcolor{red}{申，重}也。南交言夏與春交。舉一隅以見之。此居治南方\textcolor{red}{之官}。重，直用反。}。平秩南訛，敬致\footnote{\textcolor{red}{訛，化}也。掌夏之官平敘南方化育之事，敬行其教，以致其功。四時同之，亦舉\textcolor{red}{一隅}。訛，五和反。}。日永，星火，以正仲夏\footnote{\textcolor{red}{永，長}也，謂夏至之日。火，蒼龍之中星，舉中則七星見\textcolor{red}{可知}。以正仲夏之氣節，季孟亦可知。}。厥民因，鳥獸希革\footnote{\textcolor{red}{因，謂}老弱因就在田之丁壯以助農也。夏時鳥獸毛羽希少改易。\textcolor{red}{革，改}也。}。

{\noindent\zhuan\zihao{6}\fzbyks 傳“申重”至“之官”。正義曰：“申、重”，\CJKunderwave{釋詁}文。此官既主四時，亦主方面，經言“南交”，謂南方與東方交，傳言“夏與春交”，見其時、方皆掌之。春盡之日與立夏之初，時相交也，東方之南，南方之東，位相交也，言羲叔所掌與羲仲相交際也。四時皆舉仲月之候，嫌其不統季孟，於此言“交”，明四時皆然，故傳言“舉一隅以見之”。春上無冬，不得見其交接，至是夏與春交，故此言之。 \par}

{\noindent\zhuan\zihao{6}\fzbyks 傳“訛化”至“一隅”。正義曰：“訛，化”，\CJKunderwave{釋言}文。禾苗秀穗，化成子實，亦胎生乳化之類，故“掌夏之官平序南方化育之事”，謂勸課民耘耨，使苗得秀實。“敬行其教，以致其功”，謂敬行平秩之教，以致化育之功。農功歲終乃畢，敬行四時皆同，於此言之，見四時皆然，故云“亦舉一隅”也。夏日農功尤急，故就此言之。 \par}

{\noindent\zhuan\zihao{6}\fzbyks 傳“永長”至“可知”。正義曰：“永，長”，\CJKunderwave{釋詁}文。夏至之日日最長，故知謂夏至之日。計七宿房在其中，但房、心連體，心統其名。\CJKunderwave{左傳}言“火中”、“火見”，\CJKunderwave{詩}稱“七月流火”，皆指房、心為火，故曰“火，蒼龍之中星”。特舉一星，與鳥不類，故云“舉中則七星見可知”。計仲夏日在東井而入於酉地，即初昏之時角、亢在午，氐、房、心在巳，尾、箕在辰，是東方七宿皆得見也。 \par}

{\noindent\zhuan\zihao{6}\fzbyks 傳“因謂”至“革改”。正義曰：春既分析在外，今日因往就之，故言“因，謂老弱因就在田之丁壯以務農”也。鳥獸冬毛最多，春猶未脫，故至夏始毛羽希少，改易往前。“革”謂變革，故為改也。傳之訓字，或先或後，無義例也。 \par}

分命和仲，宅西，曰昧谷\footnote{\textcolor{red}{昧，冥}也。日入於谷而天下冥，故曰昧谷。昧谷曰西，則嵎夷東可知。此居治西方之官，掌秋天\textcolor{red}{之政}也。昧，武內反。冥，莫定反。}。寅餞納日，平秩西成\footnote{\textcolor{red}{餞，送}也。日出言導，日入言送，因事之宜。秋，西方，萬物成。平序其政，助\textcolor{red}{成物}。餞，賤衍反。馬云：“滅也，滅猶沒也。”}。宵中，星虛，以殷仲秋\footnote{\textcolor{red}{宵，夜}也。春言日，秋言夜，互相備。虛，玄武之中星，亦言七星皆以秋分日見，以正\textcolor{red}{三秋}。}。厥民夷，鳥獸毛毨\footnote{\textcolor{red}{夷，平}也，老壯在田與夏平也。毨，理也，毛更生\textcolor{red}{整理}。“毛毨”，下先典反。\CJKunderwave{說文}云：“仲秋鳥獸毛盛，可選取以為器用也。”}。

{\noindent\zhuan\zihao{6}\fzbyks 傳“昧冥”至“之政”。正義曰：\CJKunderwave{釋言}云：“晦，冥也。”冥是暗，故“昧”為冥也。“谷”者日所行之道，日入於谷而天下皆冥，故謂日入之處為“昧谷”，非實有谷而日入也。此經春秋相對，春不言“東”，但舉昧谷曰“西”，則嵎夷東可知。然則東言“嵎夷”，則西亦有地明矣，闕其文所以互見之。傳於春言“東方之官”，不言“掌春”，夏言“掌夏之官”，不言“南方”，此言“居治西方之官，掌秋天之政”，互文明四時皆同。 \par}

{\noindent\zhuan\zihao{6}\fzbyks 傳“餞送”至“成物”。正義曰：送行飲酒謂之餞，故“餞”為送也。導者引前之言，送者從後之稱，因其欲出，導而引之，因其欲入，從而送之，是其因事之宜而立此文也。秋位在西，於時萬物成熟，平序其秋天之政，未成則耘耨，既熟則收斂，助天成物，以此而從送入日也。納、入義同，故傳以入解“納”。 \par}

{\noindent\zhuan\zihao{6}\fzbyks 傳“宵夜”至“三秋”。正義曰：“宵，夜”，\CJKunderwave{釋言}文。舍人曰：“宵，陽氣消也”。三時皆言日，惟秋言夜，故傳辨之云：“春言日，秋言夜，互相備”也，互著明也。明日中宵亦中，宵中日亦中，因此而推之，足知日永則宵短，日短則宵長，皆以此而備知也。正於此時變文者，以春之與秋日夜皆等，春言“出日”即以“日”言之，秋雲“納日”即以“夜”言之，亦事之宜也。北方七宿則虛為中，故虛為玄武之中星。計仲秋日在角、亢而入於酉地，初昏之時鬥、牛在午，女、虛、危在巳，室、壁在辰，舉虛中星言之，亦言七星皆以秋分之日昏時並見，以正秋之三月。 \par}

{\noindent\zhuan\zihao{6}\fzbyks 傳“夷平”至“整理”。正義曰：\CJKunderwave{釋詁}云：“夷、平，易也。”俱訓為易，是“夷”得為平。秋禾未熟，農事猶煩，故“老壯在田與夏平”也。“毨”者，毛羽美悅之狀，故為理也。夏時毛羽希少，今則毛羽復生，夏改而少，秋更生多，故言“更生整理”。 \par}

申命和叔，宅朔方，曰幽都。平在朔易\footnote{\textcolor{red}{北稱}朔，亦稱方，言一方則三方見矣。北稱幽都，南稱明從可知也。都謂所聚也。易謂歲改易北方，平均在察其政，以順天常。上總言羲和敬順昊天，此分別仲叔，各有\textcolor{red}{所掌}。別音彼列反，下同。}。日短，星昴，以正仲冬\footnote{日短，冬至之日。昴，白虎之中星，亦以七星並見，以正冬之三節。}。厥民隩，鳥獸氄毛\footnote{\textcolor{red}{隩，室}也。民改歲入此室處，以闢風寒。鳥獸皆生而毳細毛以自\textcolor{red}{溫焉}。隩，於六反，馬云：“暖也。”氄，如勇反，徐又音而充反，馬云：“溫柔貌。”闢音避。耎，如兗反，本或作濡，音儒。毳，尺銳反。}。

{\noindent\zhuan\zihao{6}\fzbyks 傳“北稱”至“所掌”。正義曰：\CJKunderwave{釋訓}云：“朔，北方也。”舍人曰：“朔，盡也。北方萬物盡，故言朔也。”李巡曰：“萬物盡於北方，蘇而復生，故言北方。”是“北稱朔”也。羲和主四方之官，四時皆應言“方”，於此言“方”者，即三方皆見矣。春為歲首，故舉地名；夏與春交，故言“南交”；秋言“西”以見嵎夷當為東，冬言“方”以見三時皆有方。古史要約，其文互相發見也。“幽”之與“明”文恆相對，北既稱“幽”，則南當稱“明”，從此可知，故於夏無文。經冬言“幽都”，夏當雲“明都”，傳不言“都”者，從可知也。鄭云：“夏不言‘曰明都’三字，摩滅也。”伏生所誦與壁中舊本並無此字，非摩滅也。王肅以“夏無‘明都’，避‘敬致’，然即‘幽’足見‘明’，闕文相避”，如肅之言，義可通矣。“都謂所聚”者，總言此方是萬物所聚之處，非指都邑聚居也。“易謂歲改易於北方”者，人則三時在野，冬入隩室，物則三時生長，冬入囷倉，是人之與物皆改易也。王肅云：“改易者，謹約蓋藏，循行積聚。”引\CJKunderwave{詩}“嗟我婦子，曰為改歲，入此室處”。王肅言人物皆易，孔意亦當然也。\CJKunderwave{釋詁}云：“在，察也。舍人曰：“在，見物之察。”是“在”為察義，故言“平均在察其政，以順天常”。以“在察”須與“平均”連言，不復訓“在”為察，故\CJKunderwave{舜典}之傳別更訓之。三時皆言“平秩”，此獨言“平在”者，以三時乃役力田野，當次序之，冬則物皆藏入，須省察之，故異其文。秋日物成就，故傳言“助成物”，冬日蓋藏，天之常道，故言“順天常”，因明“東作”、“南訛”亦是助生物，類常道也。上總言羲和敬順昊天，此分別仲叔各有所掌，明此四時之節，即順天之政，實恐人以“敬順昊天”直是曆象日月，嫌仲叔所掌非順天之事，故重明之。 \par}

{\noindent\zhuan\zihao{6}\fzbyks 傳“隩室”至“溫焉”。正義曰：\CJKunderwave{釋宮}云：“西南隅謂之隩。”孫炎云：“室中隱隩之處也。”隩是室內之名，故以“隩”為室也。物生皆盡，野功咸畢，是歲改矣。以天氣改歲，故入此室處，以避風寒。天氣既至，故鳥獸皆生耎毳細毛以自溫焉。經言“氄毛”，謂附肉細毛，故以“耎毳”解之。 \par}

帝曰：“諮！汝羲暨和。期三百有六旬有六日，以閏月定四時，成歲\footnote{\textcolor{red}{諮，嗟}。暨，與也。匝四時曰期。一歲十二月，月三十日，正三百六十日；除小月六，為六日，是為一歲有餘十二日；未盈三歲足得一月，則置閏焉，以定四時之氣節，成一歲之\textcolor{red}{曆象}。暨，其器反。期，居其反，下同。旬，似遵反，十日為旬。匝，子合反。}。允釐百工，庶績咸熙。”\footnote{\textcolor{red}{允，信}。釐，治。工，官。績，功。咸，皆。熙，廣也。言定四時成歲歷，以告時授事，則能信治百官，眾功皆廣，嘆\textcolor{red}{其善}。釐，力之反。熙，許其反，興也。}

{\noindent\zhuan\zihao{6}\fzbyks 傳“諮嗟”至“曆象”。正義曰：“諮，嗟”、“暨，與”，皆\CJKunderwave{釋詁}文也。“迎四時曰期”，“期”即“迎”也。故王肅云：“期，四時是也。”然古時真歷遭戰國及秦而亡，漢存六歷雖詳於五紀之論，皆秦漢之際假託為之，實不得正要有梗概之言。周天三百六十五度四分度之一。而日日行一度，則一期三百六十五日四分日之一。今\CJKunderwave{考靈曜}、\CJKunderwave{幹鑿度}諸緯皆然。此言三百六十六日者，王肅云：“四分日之一又入六日之內，舉全數以言之，故云三百六十六日也。”傳又解所以須置閏之意，皆據大率以言之，云：“一歲十二月，月三十日，正三百六十日也；除小月六，又為六日。”今經雲三百六十六日，故云“餘十二日”，不成期。以一月不整三十日，今一年餘十二日，故未至盈滿三歲足得一月，則置閏也。以時分於歲，故云“氣節”，謂二十四氣，時月之節。歲總於時，故云“曆象日月星辰，敬授人時”，以相配成也。六歷、諸緯與\CJKunderwave{周髀}皆云，日行一度，月行十三度十九分度之七，為每月二十九日過半。日之於法，分為日九百四十分日之四百九十九，即月有二十九日半強。為十二月，六大之外有日分三百四十八，是除小月無六日，又大歲三百六十六日，小歲三百五十五日，則一歲所餘無十二日。今言“十二日”者，皆以大率據整而計之，其實一歲所餘正十一日弱也。以為十九年七閏，十九年年十一日則二百九日，其七月四大三小猶二百七日，況無四大乎？為每年十一日弱分明矣。所以弱者，以四分日之一於九百四十分，則一分為二百三十五分，少於小月餘分三百四十八。以二百三十五減三百四十八，不盡一百一十三，是四分日之一餘矣。皆以五日為率，其小月雖為歲日殘分所減，猶餘一百一十三，則實餘尚無六日。就六日抽一日為九百四十分減其一百一十三分，不盡八百二十七分。以不抽者五日並三百六十日外之五日為十日，其餘九百四十分日之八百二十七，為每歲之實餘。今十九年年，十日得整日一百九十。又以十九乘八百二十七分，得一萬五千七百一十三。以日法九百四十除之，得十六日。以並一百九十日為二百六日，不盡六百七十三分為日餘。今為閏月得七，每月二十九日，七月為二百三日。又每四百九十九分以七乘之得三千四百九十三，以日法九百四十分除之得三日。以二百三日亦為二百六日，不盡亦六百七十三為日餘，亦相當矣。所以無閏時不定，歲不成者，若以閏無，三年差一月，則以正月為二月，每月皆差；九年差三月，即以春為夏；若十七年差六月，即四時相反；時何由定，歲何得成乎？故須置閏以定四時。故\CJKunderwave{左傳}雲“履端於始，序則不愆；舉正於中，民則不惑；歸餘於終，事則不悖”是也，先王以重閏焉。王肅云：“鬥之所建，是為中氣，日月所在。鬥指兩辰之間，無中氣，故以為閏也。” \par}

{\noindent\zhuan\zihao{6}\fzbyks 傳“允信”至“其善”。正義曰：\CJKunderwave{釋訓}云：“鬼之為言歸也。”\CJKunderwave{鄉飲酒義}云：“春之為言蠢也。”然則\CJKunderwave{釋訓}之例有以聲相近而訓其義者，“釐，治”，“工，官”，皆以聲近為訓，他皆放此類也。“績，功”、“咸、皆”，\CJKunderwave{釋詁}文。“熙、廣”，\CJKunderwave{周語}文。此經文義承“成歲”之下，傳以文勢次之，言定歷授事能使眾功皆廣。“嘆其善”謂帝嘆羲和之功也。 \par}

{\noindent\shu\zihao{5}\fzkt “乃命”至“咸熙”。正義曰：上言能明俊德，又述能明之事，堯之聖德美政如上所陳。但聖不必獨理，必須賢輔。堯以須臣之故，乃命有俊明之人羲氏、和氏敬順昊天之命，歷此法象。其日之甲乙，月之大小，昏明遞中之星，日月所會之辰，定其所行之數，以為一歲之歷。乃依此歷，敬授下人以天時之早晚。其總為一歲之歷，其分有四時之異，既舉總目，更別序之。堯於羲和之內，乃分別命其羲氏而字仲者，令居治東方嵎夷之地也。日所出處名曰暘明之谷，於此處所主之職，使羲仲主治之。既主東方之事，而日出於東方，令此羲仲恭敬導引將出之日，平均次序東方耕作之事，使彼下民務勤種植。於日晝夜中分，刻漏正等，天星朱鳥南方七宿合昏畢見，以此天之時候調正仲春之氣節。此時農事已起，不居室內，其時之民宜分析適野。老弱居室，丁壯就功。於時鳥獸皆孕胎卵，孳尾匹合。又就所分羲氏之內，重命其羲氏而字叔者，使之居治南方之職，又於天分南方與東交，立夏以至立秋時之事，皆主之。均平次序南方化育之事，敬行其教，以致其功，於日正長，晝漏最多，天星大火東方七宿合昏畢見，以此天時之候調正仲夏之氣節。於時苗稼已殖，農事尤煩，其時之民，老弱因一丁壯就在田野。於時鳥獸羽毛希少，變改寒時。又分命和氏而字仲者，居治西方日所入處，名曰昧冥之谷。於此處所主之職，使和仲主治之。既主西方之事，而日入在於西方，令此和仲恭敬從送既入之日，平均次序西方成物之事，使彼下民務勤收斂。於晝夜中分，漏刻正等，天星之虛北方七宿合昏畢見，以此天時之候調正仲秋之氣節。於時禾苗秀實，農事未閒，其時之民與夏齊平，盡在田野。於時鳥獸毛羽更生，已稍整治。又重命和氏而字叔者，令居治北方名曰幽都之地，於此處所主之職，使和叔主治之。平均視察北方歲改之事。於日正短，晝漏最少，天星之昴西方七宿合昏畢見，以此天時之候調正仲冬之氣節。於時禾稼已入，農事閒暇，其時之人皆處深隩之室，鳥獸皆生耎毳細毛以自溫暖。此是羲和敬天授人之實事也。羲和所掌如是，故帝堯乃述而嘆之曰：“諮嗟！汝羲仲、羲叔與和仲、和叔。一期之間三百有六旬有六日，分為十二月，則餘日不盡，令氣朔參差，若以閏月補闕，令氣朔得正定四時之氣節，成一歲之曆象，是汝之美可嘆也。又以此歲歷告時授事，信能和治百官，使之眾功皆廣也。”嘆美羲和能敬天之節，眾功皆廣，則是風俗大和。 \par}

\textcolor{red}{帝曰}：“疇諮若時？登庸。”\footnote{\textcolor{red}{疇，誰}。庸，用也。誰能咸熙庶績，順是事者，將登\textcolor{red}{用之}。疇，直由反。}

{\noindent\zhuan\zihao{6}\fzbyks 傳“疇誰”至“用之”。正義曰：“疇，誰”，\CJKunderwave{釋詁}文。“庸”聲近“用”，故為用也。馬融以“羲和為卿官堯之末年，皆以老死，庶績多闕。故求賢順四時之職，欲用以代羲和”。孔於下傳云：“四嶽，即上羲和之四子。”帝就羲和求賢，則所求者別代他官，不代羲氏、和氏。孔以羲和掌天地之官，正在敬順昊天，告時授事而已，其施政者乃是百官之事，非復羲和之職。但羲和告時授事，流行百官，使百官庶績咸熙，今雲“咸熙庶績，順是事者”，指謂求代百官之闕，非求代羲和也。此經文承“庶績”之下而言“順是事者”，故孔以文勢次之，此言“誰能咸熙庶績，順是事者，將登用之”，蓋求卿士用任也。計堯即位至洪水之時六十餘年，百官有闕，皆應求代。求得賢者，則史亦不錄。不當帝意，乃始錄之，為求舜張本。故惟帝求一人，放齊以一人對之，非六十餘年止求一人也。堯以聖德在位，庶績咸熙，蓋應久矣。此繼“咸熙”之下，非知早晚求之，史自歷序其事，不必與治水同時也。計四嶽職掌天地，當是朝臣之首。下文求治水者，帝“諮四嶽”，此不言“諮四嶽”者，帝求賢者固當博訪朝臣，但史以有嶽對者言“諮！四嶽”，此不言“諮”者，但此無嶽對，故不言耳。 \par}

放齊曰：“胤子朱啟明。”帝曰：“籲！嚚訟，可乎？”\footnote{\textcolor{red}{放齊}，臣名。胤，國。子，爵。朱，名。啟，開也。籲，疑怪之辭。言不忠信為嚚，又好爭訟，可乎！言\textcolor{red}{不可}。放，方往反，注同。胤，引信反，馬云：“嗣也。”籲，況於反，徐往付反，一音於。嚚，魚巾反。訟，才用反，馬本作庸。好，呼報反，下注同。爭，鬥也。}

{\noindent\zhuan\zihao{6}\fzbyks 傳“放齊”至“不可”。正義曰：以放齊舉人對帝，故知臣名，為名為字，不可得知。傳言“名”者，辯此是為臣之名號耳，未必是臣之名也。夏王仲康之時，胤侯命掌六師，\CJKunderwave{顧命}陳寶有胤之舞衣，故知古有胤國。“胤”既是國，自然“子”為爵，“朱”為名也。馬融、\CJKunderline{鄭玄}以為“帝之胤子曰朱也”。求官而薦太子，太子下愚以為啟明，揆之人情，必不然矣。“啟”之為開，書傳通訓，言此人心志開解而明達。“籲”者必有所嫌而為此聲，故以為“疑怪之辭”。僖二十四年\CJKunderwave{左傳}曰：“口不道忠信之言為嚚。”是“言不忠信為嚚”也。其人心既頑嚚，又好爭訟，此實不可，而帝雲“可乎”，故籲聲而反之。“可乎”，言不可也。唐堯聖明之主，應任賢哲，放齊聖朝之臣，當非庸品，人有善惡，無容不知，稱“嚚訟”以為“啟明”，舉愚臣以對聖帝，何哉？將以知人不易，人不易知，密意深心，固難照察。胤子矯飾容貌，但以惑人，放齊內少鑑明，未能圓備，謂其實可任用，故承意舉之。以帝堯之聖，乃知其嚚訟之事，放齊所不知也。驩兜薦舉共工，以為比周之惡，謂之四凶，投之遠裔。放齊舉胤子，不為兇人者，胤子雖有嚚訟之失，不至滔天之罪，放齊謂之實賢，非是苟為阿比。驩兜則志不在公，私相朋黨，共工行背其言，心反於貌，其罪並深，俱被流放，其意異於放齊舉胤子故也。 \par}

帝曰：“疇諮若予採？”\footnote{\textcolor{red}{採，事}也。復求誰能順我\textcolor{red}{事者}。予音餘，又羊汝反。採，七在反，馬云：“官也。”復，扶又反。}

{\noindent\zhuan\zihao{6}\fzbyks 傳“採事”至“事者”。正義曰：“採，事”，\CJKunderwave{釋詁}文。上已求順時，不得其人，故復求順我事者。順時順事其義一也。史以上承“庶績”之下，故言順時，謂順是庶績之事，此不可復同前文，故變言順我帝事，其意亦如前經，當求卿士之任也。順我事之下亦宜有“登用”之言，上文已具，故於此略之。 \par}

驩兜曰：“都！共工方鳩僝功。”\footnote{\textcolor{red}{驩兜}，臣名。都，於，嘆美之辭。共工，官稱。鳩，聚。僝，見也。嘆共工能方方聚見\textcolor{red}{其功}。驩，呼端反。兜，丁侯反。共音恭，注同。僝，仕簡反，徐音撰，馬云：“具也。”於音烏。稱，尺證反。}

{\noindent\zhuan\zihao{6}\fzbyks 傳“驩兜”至“其功”。正義曰：驩兜亦舉人對帝，故知臣名。“都，於”，\CJKunderwave{釋詁}文。“於”即“嗚”字，嘆之辭也。將言共工之善，故先嘆美之。\CJKunderwave{舜典}命垂作共工，知共工是官稱。鄭以為“其人名氏未聞。先祖居此官。故以官氏也”。計稱人對帝。不應舉先世官名．孔直雲“官稱”，則其人於時居此官也。時見居官，則是已被任用，復舉之者，帝求順事之人，欲置之上位，以為大臣，所欲尊於共工，故舉之也。“鳩，聚”，\CJKunderwave{釋詁}文。僝然，見之狀，故為見。“嘆共工能方方聚見其功”，謂每於所在之方，皆能聚集善事，以見其功，言可用也。若能共工實有見功，則是可任用之人，帝言其庸違滔天不可任者，共工言是行非，貌恭心很，取人之功以為己功，其人非無見功，但功非己有。\CJKunderwave{左傳}說驩兜雲“醜類惡物”，是與比周；“天下之人謂之渾敦”，言驩兜以共工比周，妄相薦舉，知所言見功非其實功也。 \par}

帝曰：“籲！靜言庸違，象恭滔天。”\footnote{\textcolor{red}{靜，謀}。滔，漫也。言共工自為謀言，起用行事而違背之，貌象恭敬而心傲很，若漫天。言不\textcolor{red}{可用}。滔，吐刀反。漫，末旦反，下同，又末寒反。背音佩。傲，五報反，下同。很，恨懇反。}

{\noindent\zhuan\zihao{6}\fzbyks 傳“靜謀”至“可用”。正義曰：“靜，謀”，\CJKunderwave{釋詁}文。滔者，漫浸之名，浸必漫其上，故“滔”為漫也。共工險偽之人，自為謀慮之言皆合於道，及起用行事而背違之，言其語是而行非也。貌象恭敬而心傲很，其侮上陵下，若水漫天，言貌恭而心很也。行與言違，貌恭心反，乃是大佞之人，不可任用也。明君聖主莫先於堯，求賢審官王政所急，乃有放齊之不識是非，驩兜之朋黨惡物，共工之巧言令色，崇伯之敗善亂常，聖人之朝不才總萃，雖曰難之，何其甚也！此等諸人，才實中品，亦雖行有不善，未為大惡，故能仕於聖代，致位大官。以帝堯之末，洪水為災，欲責非常之功，非復常人所及，自非聖舜登庸，大禹致力，則滔天之害未或可平。以舜禹之成功，見此徒之多罪。勳業既謝，愆釁自生，為聖所誅，其咎益大。且虞史欲盛彰舜德，歸過前人，\CJKunderwave{春秋}史克以宣公比堯，辭頗增甚，知此等並非下愚，未有大惡。其為不善，惟帝所知，將言求舜，以見帝之知人耳。 \par}

帝曰：“諮！四嶽，\footnote{四嶽，即上羲和之四子，分掌四嶽之諸侯，故稱焉。}湯湯洪水方割，\footnote{湯湯，流貌。洪，大。割，害也。言大水方方為害。湯音傷。洪音戶工反。}蕩蕩懷山襄陵，浩浩滔天。\footnote{蕩蕩，言之奔突有所滌除。懷，包。襄，上也。包山上陵，浩浩盛大，若漫天。浩，胡老反。滌，大曆反。上，時掌反。}下民其諮，有能俾乂？”\footnote{俾，使。乂，治也。言民諮嗟憂愁，病水困苦，故問四嶽，有能治者將使之。俾，必爾反。}

{\noindent\zhuan\zihao{6}\fzbyks 傳“四嶽”至“稱焉”。正義曰：上列羲和所掌雲宅嵎夷、朔方，言四子居治四方，主於外事。嶽者，四方之大山。今王朝大臣皆號稱“四嶽”，是與羲和所掌其事為一，以此知“四嶽，即上羲和之四子”也。又解謂之嶽者，以其“分掌四嶽之諸侯，故稱焉”。\CJKunderwave{舜典}稱“巡守至於岱宗，肆覲東後”，\CJKunderwave{周官}說巡守之禮云，諸侯各朝於方岳之下，是四方諸侯分屬四嶽也。計堯在位六十餘年，乃命羲和蓋應早矣。若使成人見命，至此近將百歲，故馬、鄭以為羲和皆死。孔以為四嶽即是羲和至今仍得在者。以羲和世掌天地，自當父子相承，不必仲叔之身皆悉在也。\CJKunderwave{書傳}雖出自伏生，其常聞諸先達，虞傳雖說\CJKunderwave{舜典}之四嶽尚有羲伯、和伯，是仲叔子孫世掌嶽事也。 \par}

{\noindent\zhuan\zihao{6}\fzbyks 傳“湯湯”至“為害”。正義曰：湯湯，波動之狀，故為“流貌”。“洪，大”，\CJKunderwave{釋詁}文。刀害為割，故“割”為害也。“言大水方方為害”謂其遍害四方也。 \par}

{\noindent\zhuan\zihao{6}\fzbyks 傳“蕩蕩”至“漫天”。正義曰：蕩蕩，廣平之貌，“言水勢奔突有所滌除”，謂平地之水，除地上之物，為水漂流，無所復見，蕩然惟有水耳。懷，藏，包裹之義，故“懷”為包也。\CJKunderwave{釋言}以“襄”為駕，駕乘牛馬皆車在其上，故“襄”為上也。“包山”謂繞其傍，“上陵”謂乘其上，平地已皆蕩蕩，又復繞山上陵，故為盛大之勢，總言浩浩盛大若漫天然也。天者無上之物，漫者加陵之辭，甚其盛大，故云“若漫天”也。 \par}

{\noindent\zhuan\zihao{6}\fzbyks 傳“俾，使”、“乂，治也”。正義曰：“俾，使”、“乂，治”，\CJKunderwave{釋詁}文。 \par}

僉曰：“於，鯀哉！”\footnote{僉，皆也。鯀，崇伯之名。朝臣舉之。僉，七廉反，又七劍反。於音烏。鯀，故本反，馬云：“禹父也。朝，直遙反。}

{\noindent\zhuan\zihao{6}\fzbyks 傳“僉皆”至“舉之”。正義曰：“僉，皆”，\CJKunderwave{釋詁}文。\CJKunderwave{周語}雲“有崇伯鯀”，即鯀是崇君；伯，爵；故云“鯀，崇伯之名”。帝以嶽為朝臣之首，故特言四嶽，其實求能治者，普問朝臣，不言嶽對而云皆曰，乃眾人舉之，非獨四嶽，故言“朝臣舉之”。 \par}

帝曰：“籲，咈哉！方命圮族。”\footnote{凡言“籲”者皆非帝意。咈，戾。圮，毀。族，類也。言鯀性很戾，好比方名，命而行事，輒毀敗善類。咈，扶弗反，忿戾也。方如字，馬云：“方，放也。”徐云：“鄭、王音放。”圮音皮美反。戾音力計反。}

{\noindent\zhuan\zihao{6}\fzbyks 傳“凡言”至“善類”。正義曰：自上以來三經求人，所舉者帝言其惡，而辭皆稱“籲”，故知凡言“籲”者皆非帝之所當意也。“咈”者相乖詭之意，故為戾也。“圮，毀”，\CJKunderwave{釋詁}文。\CJKunderwave{左氏}稱“非我族類，其心必異”，族、類義同，故“族”為類也。言鯀性很戾，多乖異眾人，好此方直之名，內有奸回之志，命而行事輒毀敗善類。何則？心性很戾，違眾用己，知善不從，故云“毀敗善類”。\CJKunderwave{詩}稱“貪人敗類”，與此同。鄭、王以“方為放，謂放棄教命”。\CJKunderwave{易·坤卦}六二“直、方、大”，是直方之事為人之美名。此經雲“方”，故依經為說。 \par}

嶽曰：“異哉！試可乃已。”\footnote{異，已也。已，退也。言餘人盡已，唯鯀可試，無成乃退。異，徐云：“鄭音異，孔、王音怡。”}

{\noindent\zhuan\zihao{6}\fzbyks 傳“異，已。已，退也”。正義曰：“異”聲近已，故為已也。已訓為止，是停住之意，故為退也。 \par}

帝曰：“往，欽哉！”\footnote{敕鯀往治水，命使敬其事。堯知其性很戾圮族，未明其所能，而據總言可試，故遂用之。}九載，績用\textcolor{red}{弗成}。\footnote{載，年也。三考九年，功用不成，則放退之。}

{\noindent\zhuan\zihao{6}\fzbyks 傳“敕鯀”至“用之”。正義曰：傳解鯀非帝所意而命使之者，堯知其性很戾圮族，未明其所能。夫管氏之好奢尚僣，翼贊霸圖；陳平之盜嫂受金，弼諧帝業，然則人有性雖不善，才堪立功者。而眾皆據之言鯀可試，冀或有益，故遂用之。孔之此說，據跡立言，必其盡理而論，未是聖人之實。何則？禹稱“帝德廣運，乃聖乃神”，夫以聖神之資，聰明之鑑，既知鯀性很戾，何故使之治水者？馬融云：“堯以大聖，知時運當然，人力所不能治，下民其諮，亦當憂勞。屈己之是，從人之非，遂用於鯀。”李顒云：“堯雖獨明於上，眾多不達於下，故不得不副倒懸之望，以供一切之求耳。” \par}

{\noindent\zhuan\zihao{6}\fzbyks 傳“載年”至“退之”。正義曰：\CJKunderwave{釋天}云：“載，歲也。夏曰歲，商曰祀，周曰年，唐虞曰載。”李巡云：“各自紀事，示不相襲也。”孫炎曰：“歲，取歲星行一次也。祀，取四時祭祀一訖也。年，取禾穀一熟也。載，取萬物終而更始；是載者年之別名，故以載為年也。”\CJKunderwave{舜典}云：“三載考績，三考，黜陟幽明。”是“三考，九年”也。功用不成，水害不息，故放退之，謂退使不復治水。至明年得舜，乃殛之羽山。\CJKunderwave{周禮·太宰職}云：“歲終則令百官各正其治，而詔王廢置。三年則大計群吏之治而誅賞。”然則考課功績必在歲終，此言“功用不成”，是九年歲終三考也。下雲“朕在位七十載”，而求得虞舜歷試三載，即數登用之年，至七十二年為三載，即知“七十載”者與此異年，此時堯在位六十九年。鯀初治水之時，堯在位六十一年。若然，鯀既無功，早應黜廢。而待九年無成始退之者，水為大災，天之常運；而百官不悟，謂鯀能治水，及遣往治，非無小益，下人見其有益，謂鯀實能治之。日復一日，以終三考，三考無成，眾人乃服，然後退之，故至九年。\CJKunderwave{祭法}云：“鯀障洪水而殛死，禹能修鯀之功。”然則禹之大功，顧亦因鯀，是治水有益之驗。但不能成功，故誅殛之耳。若然，災以運來，時不可距，假使興禹，未必能治。何以治水之功不成而便殛鯀者？以鯀性傲很，帝所素知，又治水無功，法須貶黜，先有很戾之惡，復加無功之罪，所以殛之羽山，以示其罪。若然，禹既聖人，當知洪水時未可治，何以不諫父者？梁主以為“舜之怨慕，由己之私；鯀之治水，乃為國事。上令必行，非禹能止。時又年小，不可干政也”。 \par}

{\noindent\shu\zihao{5}\fzkt “帝曰疇諮若予”至“九載績用弗成”。正義曰：史又序堯事。堯任羲和，眾功已廣，及其末年，群官有闕，復求賢人，欲任用之。帝曰：“誰乎？諮嗟。”嗟人之難得也。“有人能順此咸熙庶績之事者，我將登而用之”。有臣放齊者對帝：“有胤國子爵之君，其名曰朱，其人心志開達，性識明悟。”言此人可登用也。帝疑怪嘆之曰：“籲！此人既頑且嚚，又好爭訟，豈可用乎？”總言不可也。史又記堯復求人。帝曰：“誰乎？諮嗟。”嗟人之難得也。“今有人能順我事者否乎？”言有即欲用之也。有臣驩兜者對帝曰：“嗚乎！”嘆有人之大賢也。“帝臣共工之官者，此人於所在之方能立事業，聚見其功”。言此人可用也。帝亦疑怪之曰：“籲！此人自作謀計之言，及起用行事而背違之，貌象恭敬而心傲很，若漫天。”言此人不可用也。頻頻求人，無當帝意。於是洪水為災，求人治之。帝曰：“諮嗟！”嗟水災之大也，呼掌嶽之官而告以須人之意。“汝四嶽等，今湯湯流行之水，所在方方為害。又其勢奔突蕩蕩然，滌除在地之物，包裹高山，乘上丘陵，浩浩盛大，勢若漫天。在下之人其皆諮嗟，困病其水矣。有能治者將使治之”。群臣皆曰：“嗚呼！”嘆其有人之能。“惟鯀堪能治之”。帝又疑怪之曰：“籲！其人心很戾哉！好此方直之名，命而行事，輒毀敗善類。”言其不可使也。朝臣已共薦舉，四嶽又復然之。嶽曰：“帝若謂鯀為不可，餘人悉皆已哉。”言不及鯀也。“惟鯀一人試之可也。試若無功，乃黜退之”。言洪水必須速治，餘人不復及鯀，故勸帝用之。帝以群臣固請，不得已而用之。乃告敕鯀曰：“汝往治水，當敬其事哉！”鯀治水九載，已經三考而功用不成。言帝實知人，而朝無賢臣，致使水害未除，待舜乃治。此經三言求人，未必一時之事，但歷言朝臣不賢，為求舜張本故也。 \par}


\textcolor{red}{帝曰}：“諮！四嶽，朕在位七十載，\footnote{堯年十六以唐侯升為天子，在位七十年，則時年八十六，老將求代。朕，直錦反，馬云：“我也。”}汝能庸命，巽朕位？”\footnote{巽，順也。言四嶽能用帝命，故欲使順行帝位之事。巽音遜，馬云：“讓也。”}

{\noindent\zhuan\zihao{6}\fzbyks 傳“堯年”至“求代”。正義曰：遍檢今之書傳，無堯即位之年。\CJKunderline{孔氏}博考群書，作為此傳，言“堯年十六以唐侯升為天子”，必當有所案據，未知出何書。計十六為天子，其歲稱元年，在位七十載，應年八十五。孔雲“八十六”者，\CJKunderwave{史記}諸書皆言，堯帝嚳之子，帝摯之弟。嚳崩，摯立。摯崩，乃傳位於堯。然則堯以弟代兄，蓋逾年改元，據其改元年則七十載，數其立年故八十六。下句求人巽位是“老將求代”也。此經文承“績用不成”之下，計治水之事於時最急，不求治水之人而先求代己者，堯以身既年老，臣無可任治水之事，非己所能，故求人代己，令代者自治。是虞史盛美舜功，言堯不能治水，以大事付舜，美舜能消大災，成堯美也。 \par}

{\noindent\zhuan\zihao{6}\fzbyks 傳“巽順”至“之事”。正義曰：“巽，順”，\CJKunderwave{易·說卦}文。帝呼四嶽，言“汝能庸命”，四嶽自謙，言“己否德”，故知“汝”，四嶽。言四嶽能用帝命，故帝欲使之順行帝位之事，將使攝也。在位之臣，四嶽為長，故讓位於四嶽也。 \par}

嶽曰：“否德，忝帝位。”\footnote{否，不。忝，辱也。辭不堪。否，方久反，又音鄙。忝音他簟反。}

{\noindent\zhuan\zihao{6}\fzbyks 傳“否不”至“不堪”。正義曰：否，古今不字。“忝，辱”，\CJKunderwave{釋言}文。己身不德，恐辱帝位，自辭不堪。嶽為群臣之首，自度既不堪，意以為在位之臣皆亦不堪，由是自辭而已，不薦餘人。故帝使之明舉側陋之處。 \par}

曰：“明明揚側陋。”\footnote{堯知子不肖，有禪位之志，故明舉明人在側陋者。廣求賢也。肖音笑，\CJKunderwave{說文}云：“肖，骨肉相似也，不似其先，故曰不肖。”}

{\noindent\zhuan\zihao{6}\fzbyks 傳“堯知”至“求賢”。正義曰：此經“曰”上無“帝”，以可知而省文也。傳解四嶽既辭，而復言此者，堯知子不肖，不堪為主，有禪位與人之志，故令四嶽明舉明人令其在側陋者，欲使廣求賢也。鄭注\CJKunderwave{雜記}云：“肖，似也，言不如人也。”\CJKunderwave{史記·五帝本紀}云：“堯知子丹朱之不肖，不足授天下，於是權授舜。授舜則天下得其利而丹朱病，授丹朱則天下病而丹朱得其利。堯曰‘終不以天下之病而利一人’，而卒授舜以天下。”是堯知子不肖而禪舜之意也。\CJKunderwave{文王世子}論舉賢之法云：“或以事舉，或以言揚。”揚亦舉也，故以舉解“揚”。經之“揚”字在於二“明”之下，傳進“舉”字於兩“明”之中，經於“明”中宜有“揚”字，言明舉明人於側陋之處。“明”下有“揚”，故上闕“揚”文。傳進“舉”於“明”上，互文以足之也。“側陋”者，僻側淺陋之處。意言不問貴賤，有人則舉，是令朝臣廣求賢人也。堯知有舜而朝臣不舉，故令廣求賢以啟之。臣亦以堯知側陋有人，故不得不舉舜耳。此言堯知子不肖，有志禪位，然則自有賢子，必不禪人。授賢爰自上代，堯舜而已。非堯舜獨可，彼皆不然。將以子不肖，時無聖者，乃運值汙隆，非聖有優劣。而\CJKunderwave{緯侯}之書附會其事，乃云：“河洛之符，名字之錄。”何其妄且俗也！ \par}

師錫帝曰：“有鰥在下，曰虞舜。”\footnote{師，眾。錫，與也。無妻曰鰥。虞，氏。舜，名。在下民之中。眾臣知舜聖賢，恥己不若，故不舉。乃不獲已而言之。錫，星曆反。鰥，故頑反。虞舜，虞，氏；舜，名也。馬云：“舜，諡也。舜死後賢臣錄之，臣子為諱，故變名言諡。”}

{\noindent\zhuan\zihao{6}\fzbyks 傳“師眾”至“言之”。正義曰：“師，眾”、“錫，與”，\CJKunderwave{釋詁}文。“無妻曰鰥”，\CJKunderwave{釋名}云：“愁悒不寐，目恆鰥鰥然，故鰥字從魚，魚目恆不閉。”\CJKunderwave{王制}云：“老而無妻曰鰥。”舜於時年未三十而謂之“鰥”者，\CJKunderwave{書傳}稱\CJKunderline{孔子}對子張曰：“舜父頑，母嚚，無室家之端，故謂之鰥。”“鰥”者無妻之名，不拘老少。少者無妻可以更娶，老者即不復更娶，謂之天民之窮，故\CJKunderwave{禮}舉老者耳。\CJKunderwave{詩}云：“何草不玄，何人不鰥。”暫離室家尚謂之鰥，不獨老而無妻始稱鰥矣。\CJKunderwave{書傳}以舜年尚少為之說耳。“虞，氏。舜，名”者，舜之為虞，猶禹之為夏，外傳稱禹氏曰“有夏”，則此舜氏曰“有虞”。顓頊已來，地為國號，而舜有天下，號曰“有虞氏”，是地名也。王肅云：“虞，地名也。”皇甫謐云：“堯以二女妻舜，封之於虞，今河東太陽山西虞地是也。”然則舜居虞地，以虞為氏，堯封之虞為諸侯，及王天下，遂為天子之號，故從微至著，常稱虞氏。舜為生號之名，前已具釋。傳又解眾人以舜與帝，則眾人盡知有舜。但舜在下人之中，未有官位，眾臣德不及之，而位居其上，雖知舜實聖賢，而恥己不若，故不舉之。以帝令舉及側陋，意謂帝知有舜，乃不獲己而言之耳。知然者，正以初不薦舉，至此始言，明是恥己不若，故不早舉。舜實聖人，而連言“賢”者，對則事有優劣，散即語亦相通。舜謂禹曰“惟汝賢”，是言聖德稱“賢”也。傳以“師”為眾臣，為朝臣之眾，或亦通及吏人。王肅云：“古者將舉大事，訊群吏，訊萬人。堯將讓位，諮四嶽，使問群臣。眾舉側陋，眾皆願與舜。堯計事之大者莫過禪讓，必應博詢吏人，非獨在位。”王氏之言得其實矣。鄭以“師為諸侯之師”，帝諮四嶽，遍訪群臣，安得諸侯之師獨對帝也。 \par}

帝曰：“俞，予聞，如何？”\footnote{俞，然也。然其所舉，言我亦聞之，其德行如何？俞，羊朱反。行，不孟反，下“其行”同。}

{\noindent\zhuan\zihao{6}\fzbyks 傳“俞然”至“如何”。正義曰：“俞，然”，\CJKunderwave{釋言}文。“然其所舉，言我亦聞也，其德行如何？”恐所聞不審，故詳問之。堯知有舜，不召取禪之而訪四嶽，令眾舉薦者，以舜在卑賤，未有名聞，率暴禪之，則下人不服。故\CJKunderline{鄭玄}\CJKunderwave{六藝論}云：“若堯知命在舜，舜知命在禹，猶求於群臣，舉於側陋，上下交讓，務在服人。\CJKunderline{孔子}曰：‘人可使由之，不可使知之。’此之謂也。”是解堯使人舉舜之意也。 \par}

嶽曰：“瞽子，父頑，母嚚，象傲，\footnote{無目曰瞽。舜父有目，不能分別好惡，故時人謂之瞽，配字曰瞍。瞍無目之稱。心不則德義之經為頑。象，舜弟之字，傲慢不友。言並惡。瞽音古。傲，五報反。瞍，素後反。稱，尺證反，又如字。}克諧以孝，烝烝乂，不格奸。”\footnote{諧，和。烝，進也。言能以至孝和諧頑嚚昏傲，使進進以善自治，不至於奸惡。諧，戶皆反。烝，之承反。奸，古顏反。}

{\noindent\zhuan\zihao{6}\fzbyks 傳“無目”至“並惡”。正義曰：\CJKunderwave{周禮}樂官有瞽蒙之職，以其無目使悤了相之，是“無目曰瞽”。又解稱瞽之意，舜父有目但不能識別好惡，與無目者同，“故時人謂之瞽”。“配字曰瞍”，“瞍”亦無目之稱，故或謂之為“瞽瞍”。\CJKunderwave{詩}雲“矇瞍奏公”，是“瞍”為瞽類。\CJKunderwave{大禹謨}雲“祗載見瞽瞍”，是相配之文。\CJKunderwave{史紀}云：“舜父瞽瞍盲。”以為“瞽瞍”是名，身實無目也。孔不然者，以經說舜德行，美其能養惡人，父自名瞍，何須言之？若實無目，即是身有固疾，非善惡之事，輒言舜是盲人之子，意欲何所見乎？\CJKunderwave{論語}云：“未見顏色而言謂之瞽。”則言瞽者非謂無目。\CJKunderwave{史記}又說瞽瞍使舜上廩，從下縱火焚廩；使舜穿井，下土實井。若其身自能然，不得謂之無目，明以不識善惡故稱瞽耳。“心不則德義之經為頑”，僖二十四年\CJKunderwave{左傳}文。“象，舜弟之字”，以字表象是人之名號，其為名字未可詳也。\CJKunderwave{釋訓}云：“善兄弟為友。”\CJKunderwave{孟子}說象與父母共謀殺舜，是“傲慢不友”。言舜父母與弟並皆惡也。此經先指舜身，因言瞽子，又稱父頑者，欲極其惡，故文重也。 \par}

{\noindent\zhuan\zihao{6}\fzbyks 傳“諧和”至“奸惡”。正義曰：“諧，和”、“烝，進”，\CJKunderwave{釋詁}文。上歷言三惡，此美舜能養之，言舜能和之以至孝之行，和頑嚚昏傲，使皆進進於善道，以善自治，不至於奸惡。以下愚難變化，令慕善是舜之美行，故以此對堯。案\CJKunderwave{孟子}及\CJKunderwave{史紀}稱瞽瞍縱火焚廩，舜以兩笠自扞而下；以土實井，舜從旁空井出；象與父母其分財物。舜之大孝升聞天朝，堯妻之二女，三惡尚謀殺舜，為奸之大莫甚於此。而言“不至奸”者，此三人性實下愚，動罹刑網，非舜養之，久被刑戮，猶尚有心殺舜，餘事何所不為？舜以權謀自免厄難，使瞽無殺子之愆，象無害兄之罪，“不至於奸惡於此益驗。終令瞽亦允若，象封有鼻，是“不至於奸惡”也。 \par}

帝曰：“我其試哉！\footnote{○言欲試舜，觀其行跡。}女於時，觀厥刑于二女。”\footnote{女，妻。刑，法也。堯於是以二女妻舜，觀其法度接二女，以治家觀治國。女音而據反。妻音千計反。}釐降二女於媯汭，嬪於虞。\footnote{降，下。嬪，婦也。舜為匹夫，能以義理下帝女之心於所居媯水之汭，使行婦道於虞氏。媯音居危反。汭音如銳反，水之內也。杜預注\CJKunderwave{左傳}云：“水之隈曲曰汭。”嬪音毗人反。}

{\noindent\zhuan\zihao{6}\fzbyks 傳“言欲”至“行跡”。正義曰：下言妻舜以女，觀其治家，是試舜觀其行跡也。馬、鄭、王本說“此經皆無‘帝曰’”，當時庸生之徒漏之也。\CJKunderline{鄭玄}云：“試以為臣之事。”王肅云：“試之以官。”鄭、王皆以\CJKunderwave{舜典}合於此篇，故指歷試之事充此“試哉”之言。孔據古今別卷，此言“試哉”正謂以女試之，既善於治家，別更試以難事，與此異也。 \par}

{\noindent\zhuan\zihao{6}\fzbyks 傳“女妻”至“治國”。正義曰：\CJKunderwave{左傳}稱“宋雍氏女於鄭莊公”，“晉伐驪戎，驪戎男女以驪姬”，以女妻人謂之女，故云“女，妻”也。“刑，法”，\CJKunderwave{釋詁}文。此已下皆史述堯事，非復堯語。言“女於時”，謂妻舜於是，故傳倒文以曉，民“堯於是以二女妻舜”。必妻之者，舜家有三惡，身為匹夫，忽納帝女，難以和協，觀其施法度於二女，以法治家觀治國。將使治國，故先使治家。敵夫曰妻，不得有二女，言“女於時”者，總言之耳。二女之中當有貴賤長幼，劉向\CJKunderwave{列女傳}云：“二女長曰娥皇，次曰女英。舜既升為天子，娥皇為後，女英為妃。”然則初適舜時，即娥皇為妻。鄭“不言妻者，不告其父，不序其正”。又注\CJKunderwave{禮記}云：“舜不告而娶，不立正妃。”此則鄭自所說，未有書傳云然。案\CJKunderwave{世本}“堯是黃帝玄孫，舜是黃帝八代之孫”，計堯女於舜之曾祖為四從姊妹，以之為妻，於義不可。\CJKunderwave{世本}之言未可據信，或者古道質故也。 \par}

{\noindent\zhuan\zihao{6}\fzbyks 傳“降下”至“虞氏”。正義曰：“降，下”，\CJKunderwave{釋詁}文。\CJKunderwave{周禮}九嬪之職“掌婦學之法”，嬪是婦之別名，故以“嬪”為婦。“釐降”，未能以義理下之，則女意初時不下，故傳解之，言舜為匹夫，帝女下嫁，以貴適賤，必自驕矜，故美舜能以義理下帝女尊亢之心於所居媯水之汭，使之服行婦道於虞氏。“虞”與“媯汭”為一地，見其心下，乃行婦道，故分為二文。言“匹夫”者，士大夫已上則有妾媵，庶人無妾媵，惟夫妻相匹，其名既定，雖單亦通，謂之匹夫匹婦。媯水在河東虞鄉縣歷山西，西流至蒲阪縣，南入於河，舜居其旁。周武王賜陳胡公之姓為媯，為舜居媯水故也。舜仕堯朝，不家在於京師，而令二女歸虞者，蓋舜以大孝示法，使妻歸事於其親，以帝之賢女事頑嚚舅姑，美其能行婦道，故云“嬪於虞”。 \par}

帝曰：“\textcolor{red}{欽哉}！”\footnote{嘆舜能修己行敬以安人，則其所能者大矣。}

{\noindent\zhuan\zihao{6}\fzbyks 傳“嘆舜”至“大矣”。正義曰：二女行婦道，乃由舜之敬，故帝言“欽哉”。嘆能修己行敬以安民也。能修己及安人，則是所能者大，故嘆之。\CJKunderwave{論語}云：“修己以安百姓，堯舜其猶病諸。”傳意出於彼也。 \par}

{\noindent\shu\zihao{5}\fzkt “帝曰諮四”至“欽哉”。正義曰：帝以鯀功不成，又已年老，求得授位明聖，代御天災，故諮嗟：“汝四嶽等，我在天子之位七十載矣。”言己年老，不堪在位。“汝等四嶽之內，有能用我之命，使之順我帝位之事。”言欲讓位與之也。四嶽對帝曰：“我等四嶽皆不有用命之德，若使順行帝事，即辱於帝位。”言己不堪也。帝又言曰：“汝當明白舉其明德之人於僻隱鄙陋之處，何必在位之臣乃舉之也。”於是朝廷眾臣乃與帝之明人曰：“有無妻之鰥夫，在下民之內，其名曰虞舜。”言側陋之處有此賢人。帝曰：“然，我亦聞之，其德行如何？”四嶽又對帝曰：“其人愚瞽之子，其父頑，母嚚，其弟字象，性又傲慢。家有三惡，其人能諧和以至孝之行，使此頑嚚傲慢者皆進進於善以自治，不至於奸惡。”言能調和惡人，是為賢也。帝曰：“其行如此，當可任用，我其召而試之哉！欲配女與試之也。”即以女妻舜，於是欲觀其居家治否也。舜能以義理下二女之心於媯水之汭，使行婦道於虞氏。帝嘆曰：“此舜能敬其事哉！”嘆其善治家，知其可以治國，故下篇言其授以官位而歷試諸難。 \par}


%%% Local Variables:
%%% mode: latex
%%% TeX-engine: xetex
%%% TeX-master: "../Main"
%%% End:

%% -*- coding: utf-8 -*-
%% Time-stamp: <Chen Wang: 2024-04-02 11:42:41>

% {\noindent \zhu \zihao{5} \fzbyks } -> 注 (△ ○)
% {\noindent \shu \zihao{5} \fzkt } -> 疏

\chapter{卷三}


\section{舜典第二（堯典下）}


 {\noindent\zihao{6}\fzbyks \CJKunderwave{釋文}：“王氏注，相承云：‘梅頤上\CJKunderline{孔氏}傳\CJKunderwave{古文尚書}，雲\CJKunderwave{舜典}一篇。’時以王肅注頗類\CJKunderline{孔氏}，故取王注從‘謹徽五典’以下為\CJKunderwave{舜典}，以續孔傳徐仙民亦音此本，今依舊音之。” \par}

\textcolor{red}{\CJKunderline{虞舜}}側微\footnote{為庶人，故微賤。}，堯聞之聰明，將使嗣位，歷試諸難\footnote{嗣，繼也。試以治民之難事。難，乃丹反。}，作\CJKunderwave{\textcolor{red}{舜典}}。

{\noindent\zhuan\zihao{6}\fzbyks 傳“為庶人，故微賤”。正義曰：此雲“側微”即\CJKunderwave{堯典}“側陋”也。不在朝廷謂之“側”，其人貧賤謂之“微”，居處褊隘故言“陋”，此指解“微”，故云“為庶人，故微賤”也。\CJKunderwave{帝系}云：“顓頊生窮蟬，窮蟬生敬康，敬康生句芒，句芒生蟜牛，蟜牛生瞽瞍，瞽瞍生舜。”昭八年\CJKunderwave{左傳}云：“自幕至於瞽瞍無違命。”似其繼世相傳，當有國土。孔言“為庶人”者，\CJKunderwave{堯典}雲“有鰥在下”，此雲“\CJKunderline{虞舜}側微”，必是為庶人矣，蓋至瞽瞍始失國也。 \par}

{\noindent\zhuan\zihao{6}\fzbyks 傳“嗣繼”至“難事”。正義曰：“嗣，繼”，\CJKunderwave{釋詁}文。經所云“慎徽五典”、“納於百揆”、“賓於四門”皆是試以治民之難事也。 \par}

舜典\footnote{“典”之義與堯同。}

{\noindent\shu\zihao{5}\fzkt “\CJKunderline{虞舜}”至“舜典”。正義曰：\CJKunderline{虞舜}所居側陋，身又微賤，堯聞之有聰明聖德，將使之繼己帝位，歷試於諸所難為之事，史述其事，故作\CJKunderwave{舜典} \par}

\textcolor{red}{曰若}稽古，\CJKunderline{帝舜}\footnote{亦言其順考古道而行之。}，曰\CJKunderline{重華}，協於帝\footnote{華謂文德，言其光文重合於堯，俱聖明。“曰若稽古，\CJKunderline{帝舜}，曰\CJKunderline{重華}，協於帝”，此十二字是姚方興所上，\CJKunderline{孔氏}傳本無。阮孝緒\CJKunderwave{七錄}亦云然。方興本或此下更有“濬哲文明，溫恭允塞，玄德升聞，乃命以位”。此二十八字異，聊出之，於王注無施也。}。濬哲文明，溫恭允塞\footnote{濬，深。哲，智也。舜有深智文明溫恭之德，信允塞上下。}，玄德升聞，乃命\textcolor{red}{以位}\footnote{玄謂幽潛，潛行道德，升聞天朝，遂見徵用。}。

{\noindent\zhuan\zihao{6}\fzbyks 傳“濬深”至“上下”。正義曰：“濬，深”、“哲，智”皆\CJKunderwave{釋言}文。舍人曰：“濬，下之深也。哲，大智也。”舜有深智，言其智之深，所知不淺近也。經緯天地曰“文”，照臨四方曰“明”。\CJKunderwave{詩}云：“溫溫恭人。”言其色溫而貌恭也。舜既有深遠之智，又有文明溫恭之德，信能充實上下也。\CJKunderwave{詩}毛傳訓“塞”為實，言能充滿天地之間，\CJKunderwave{堯典}所謂“格於上下”是也。不言四表者，以四表外無限極，非可實滿，故不言之。堯舜道同，德亦如一，史官錯互為文，故與上篇相類，是其所“合於堯”也。 \par}

{\noindent\zhuan\zihao{6}\fzbyks 傳“玄謂”至“徵用”。正義曰：老子云：“玄之又玄，眾妙之門。”則“玄”者微妙之名，故云“玄謂幽潛”也。舜在畎畝之間，潛行道德，顯彰於外，升聞天朝。“天朝”者，天子之朝也。從下而上謂之為“升”。天子聞之，故遂見徵用。 \par}

{\noindent\shu\zihao{5}\fzkt “曰若”至“以位”。正義曰：昔東晉之初，豫章內史梅賾上\CJKunderline{孔氏}傳，猶闕\CJKunderwave{舜典}。自此“乃命以位”已上二十八字，世所不傳。多用王、範之注補之，而皆以“慎徽”已下為\CJKunderwave{舜典}之初。至齊蕭鸞建武四年，吳興姚方興於大航頭得\CJKunderline{孔氏}傳古文\CJKunderwave{舜典}，亦類太康中書，乃表上之。事未施行，方興以罪致戮。至隋開皇初購求遺典，始得之。史將錄舜之美，故為題目之辭曰，能順而考案古道而行之者，是為\CJKunderline{帝舜}也。又申其順考古道之事曰，此舜能繼堯，重其文德之光華，用此德合於\CJKunderline{帝堯}，與堯俱聖明也。此舜性有深沈智慧，文章明鑑，溫和之色，恭遜之容，由名聞遠達，信能充實上下，潛行道德，升聞天朝，堯乃徵用，命之以位而試之也。 \par}

\textcolor{red}{慎徽}五典，五典克從。\footnote{徽，美也。五典，五常之教，父義、母慈、兄友、弟恭、子孝。舜慎美篤行斯道，舉八元使布之於四方，五教能從，無違命。徽，許韋反，王雲美，馬云：“善也。”從，才容反。八元，\CJKunderwave{左傳}“高辛氏有才子八人，伯奮、仲堪、叔獻、季仲、伯虎、仲熊、叔豹”。}納於百揆，百揆時敘。\footnote{揆，度也。度百事，總百官，納舜於此官。舜舉八凱，使揆度百事，百事時敘，無廢事業。揆音葵癸反。凱音開在反。\CJKunderwave{左傳}“高陽氏有才子八人，蒼舒、隤敱、檮戭、大臨、尨降、庭堅、仲容、叔達，齊聖廣淵，明允篤誠，天下之民謂之八凱”。}

{\noindent\zhuan\zihao{6}\fzbyks 傳“徽美”至“違命”。正義曰：\CJKunderwave{釋詁}雲“徽，善也”，善亦美也。此\CJKunderwave{五典}與下文“五品”、“五教”其事一也。一家之內品有五，謂父母兄弟子也。教此五者各以一事，教父以義，教母以慈，教兄以友，教弟以恭，教子以孝，是為五教也。五者皆可常行，謂之\CJKunderwave{五典}，是五者同為一事，所從言之異耳。文十八年\CJKunderwave{左傳}曰：“昔高辛氏有才子八人，伯奮、仲堪、叔獻、季仲、伯虎、仲熊、叔豹、季貍，忠肅恭懿，宣慈惠和，天下之民謂之八元。舜臣堯，舉八元，使布五教於四方，父義、母慈、兄友、弟恭、子孝。”以此知\CJKunderwave{五典}是五常之教，謂此父義之等五事也。\CJKunderwave{皋陶謨}云：“天敘有典，敕我五典五惇哉！”惇，厚也，行此五典須厚行之，篤亦厚也。言舜謹慎美善，篤行斯道，舉八元使布之於四方，命教天下之民。以此五教能使天下皆順從之，無違逆舜之命也。\CJKunderwave{左傳}又云：“故\CJKunderwave{虞書}數舜之功，曰‘慎徽五典，五典克從’，無違教也。”父母於子並宜為慈，今分之者，以父主教訓，母主撫養；撫養在於恩愛，故以慈為名；教訓愛而加嚴，故以義為稱。義者宜也，理也，教之以義，方使得事理之宜，故為義也。\CJKunderwave{釋訓}雲“善兄弟為友”，則兄弟之恩俱名為友，今雲“兄友、弟恭”者，以其同志曰友，友是相愛之名；但兄弟相愛，乃有長幼，故分其弟使之為恭，恭敬於兄，而兄友愛之。 \par}

{\noindent\zhuan\zihao{6}\fzbyks 傳“揆度”至“事業”。正義曰：“揆，度”，\CJKunderwave{釋言}文。“百揆”者，言百事皆度之。國事散在諸官，故度百事為“總百官”也。\CJKunderwave{周官}云：“唐虞稽古，建官惟百，內有百揆四嶽。”則“百揆”為官名，故云“納舜於此官”也。文十八年\CJKunderwave{左傳}云：“昔高陽氏有才子八人，蒼舒、隤敱、檮戭、大臨、尨降、庭堅、仲容、叔達，齊聖廣淵，明允篤誠，天下之民謂之八凱。舜臣堯，舉八凱，使主后土，以揆百事，莫不時敘，地平天成。”又云：“\CJKunderwave{虞書}數舜之功，曰‘納於百揆，百揆時敘’，無廢事業也。”是言百官於是得其次敘，皆無廢事業。舜既臣堯，乃舉元、凱，主后土，布五教，同時為之。史官立文，自以人事外內為次，故孔先言“八元”。若\CJKunderwave{左傳}據所出代之先後，故先舉“八凱”。堯既得舜，庶事委之。舜既臣堯，任無不統。非“五典克從”之後方始“納於百揆”，“百揆時敘”之後方始“賓於四門”，“四門穆穆”謂流四凶，流放四凶最在於前矣。\CJKunderwave{洪範}云：“\CJKunderline{鯀}則殛死，\CJKunderline{禹}乃嗣興。”是先誅\CJKunderline{鯀}而後用\CJKunderline{禹}明。此言三事皆同時為之，但言“百揆時敘”，故言“納於百揆”，其實“納於百揆”初得即然，由舜既居百揆，故得舉用二八，若偏居一職，不得分使元、凱。 \par}

賓於四門，四門穆穆。\footnote{穆穆，美也。四門，四方之門。舜流四凶族，四方諸侯來朝者，舜賓迎之，皆有美德，無兇人。朝，直遙反。}納於大麓，烈風雷雨弗迷。\footnote{麓，錄也。納舜使大錄萬機之政，陰陽和，風雨時，各以其節，不有迷錯愆伏。明舜之德合於天。麓音鹿，王云：“錄也。”馬、鄭云：“山足也。”愆音起虔反。}

{\noindent\zhuan\zihao{6}\fzbyks 傳“穆穆美”至“兇人”。正義曰：“穆穆，美也”，\CJKunderwave{釋詁}文。“四門，四方之門”，謂四方諸侯來朝者從四門而入。文十八年\CJKunderwave{左傳}歷言四凶之行乃云：“舜臣堯，流四凶族，渾敦、窮奇、檮杌、饕餮投諸四裔，以御螭魅。”又曰：“\CJKunderwave{虞書}數舜之功，曰‘賓於四門，四門穆穆’，無兇人也。”是言“皆有美德，無兇人”也。案驗四凶之族，皆是王朝之臣，舜流王朝之臣，而言諸侯無兇人者，以外見內，諸侯無兇人，則於朝必無矣。\CJKunderline{鄭玄}以“賓”為擯，謂“舜為上擯以迎諸侯”。今孔不為擯者，則謂舜既錄攝事，無不統以諸侯為賓，舜主其禮，迎而待之，非謂身為擯也。 \par}

{\noindent\zhuan\zihao{6}\fzbyks 傳“麓錄”至“於天”。正義曰：“麓”聲近錄，故為錄也。\CJKunderwave{皋陶謨}云：“一日二日萬幾。”言天下之事，事之微者有萬，喻其多無數也。納舜使大錄萬機之政，還是納於百揆，揆度百事，大錄萬機，總是一事，不為異也。但此言“德合於天”，故以“大錄”言耳。\CJKunderwave{論語}稱\CJKunderline{孔子}曰：“迅雷風烈必變。”\CJKunderwave{書傳}稱：“越常之使久矣，天之無烈風淫雨。”則烈風是猛疾之風，非善風也。經言“烈風雷雨弗迷”，言舜居大錄之時，陰陽和，風雨時，無此猛烈之風，又雷雨各以其節，不有迷錯愆伏也。“迷錯”者，應有而無，應無而有也。昭四年\CJKunderwave{左傳}雲“冬無愆陽，夏無伏陰”。無愆伏者，無冬溫夏寒也。舜錄大政，天時如此，明舜之德合於天也。此文與上三事亦同時也。上為變人，此為動天，故最後言之，以為功成之驗。王肅云：“堯得舜任之，事無不統，自‘慎徽五典’以下是也。”其言合孔意。 \par}

帝曰：“格，汝\CJKunderline{舜}。詢事考言，乃言\xpinyin*{厎}可績，三載。汝陟帝位。”\footnote{格，來。詢，謀。乃，汝。厎，致。陟，升也。堯呼舜曰：“來，汝所謀事，我考汝言。汝言致，可以立功，三年矣。三載考績，故命使升帝位。”將禪之。詢音荀。厎音之履反，王云：“致也。”馬云：“定也。”}\CJKunderline{舜}讓於德，\textcolor{red}{弗嗣}。\footnote{辭讓於德不堪，不能嗣成帝位。}

{\noindent\zhuan\zihao{6}\fzbyks 傳“格來”至“禪之”。正義曰：“格，來”，\CJKunderwave{釋言}文。“詢，謀”、“陟，升”\CJKunderwave{釋詁}文。“厎”聲近致，故為致也。經傳言“汝”多呼為“乃”，知“乃”、“汝”義同。凡事之始必先謀之，後為之，堯呼舜曰：“來，汝舜。”呼使前而與之言也。“汝所謀事，我考汝言。汝所為之事，皆副汝所謀，致可以立功，於今三年矣”。從徵得至此為三年也。君之馭臣，必三年考績。考既有功，故使升帝位，將禪之也。\CJKunderline{鯀}三考乃退，此一考使升者，\CJKunderline{鯀}待三考，冀其有成，無成功乃黜，為緩刑之義；舜既有成，更無所待，故一考即升之。且大聖之事，不可以常法論也。若然，\CJKunderwave{禹貢}“兗州作十有三載乃同”，是\CJKunderline{禹}治兗州之水，乃積十有三年。此始三年已言“地平天成”者，\CJKunderwave{祭法}云：“\CJKunderline{鯀}障洪水而殛死，\CJKunderline{禹}能修\CJKunderline{鯀}之功。”先儒馬融等皆以為“\CJKunderline{鯀}既九年，又加此三年為十二年，惟兗州未得盡平，至明年乃畢”。八州已平，一州未畢，足以為成功也。 \par}

{\noindent\shu\zihao{5}\fzkt “慎徽”至“弗嗣”。正義曰：此承“乃命以位”之下，言命之以位，試之以事也。堯使舜慎美篤行五常之教，而五常之教皆能順從而行之，無違命也。又納於百官之事，命揆度行之，而百事所揆度者於是皆得次序，無廢事也。又命使賓迎諸侯於四門，而來入者穆穆然皆有美德，無兇人也。又納於大官，總錄萬機之政，而陰陽和，風雨時，烈風雷雨不有迷惑錯謬。明舜之德合於天，天人和協，其功成矣。\CJKunderline{帝堯}乃謂之曰：“來，汝舜，有所謀之事，我考驗汝舜之所言。汝言致可以立功，於今三年，汝功已成，汝可升處帝位。”告以此言，欲禪之也。舜辭讓於德，言已德不堪嗣成帝也。 \par}

\textcolor{red}{正月}上日，受終於文祖。\footnote{上日，朔日也。終謂堯終帝位之事。文祖者堯文德之祖廟。正音政，又音徵。王云：“文祖，廟名。”馬云：“文祖，天也。天為文萬物之祖，故曰文祖。”}在\xpinyin*{璿}璣玉衡，以齊七政。\footnote{在，察也。璿，美玉。璣、衡，王者正天文之器，可運轉者。七政，日月五星各異政。舜察天文，齊七政，以審己當天心與否。璿音旋。}

{\noindent\zhuan\zihao{6}\fzbyks 傳“上日”至“祖廟”。正義曰：月之始日謂之朔日，每月皆有朔日，此是正月之朔，故云“上日”，言一歲日之上也。下雲“元日”亦然。\CJKunderline{鄭玄}以為“帝王易代，莫不改正。堯正建醜，舜正建子。此時未改堯正，故云‘正月上日’。即位，乃改堯正，故云‘月正元日’”，故以異文。先儒王肅等以為“惟殷周改正，易民視聽。自夏已上，皆以建寅為正。此篇二文不同，史異辭耳”。孔意亦然。下雲“歲二月”，傳雲“既班瑞之明月”，以此為建寅之月也。“受終”者，堯為天子，於此事終而授與舜，故知“終謂堯終帝位之事”，“終”言堯終舜始也。禮有大事，行之於廟，況此是事之大者，知“文祖者，堯文德之祖廟”也。且下云：“歸，格於藝祖。”“藝”、“文”義同。知“文祖”是廟者，\CJKunderwave{咸有一德}云：“七世之廟，可以觀德。”則天子七廟，其來自遠。堯之文祖，蓋是堯始祖之廟，不知為誰也。\CJKunderwave{帝系}及\CJKunderwave{世本}皆雲“黃帝生玄囂，玄囂生僑極，僑極生帝嚳，帝嚳生堯”。即如彼言，黃帝為堯之高祖，黃帝以上不知復祭何人，充此七數，況彼二書未必可信，堯之文祖不可強言。 \par}

{\noindent\zhuan\zihao{6}\fzbyks 傳“在察”至“與否”。正義曰：“在，察”，\CJKunderwave{釋詁}文。\CJKunderwave{說文}云：“璿，美玉也。”玉是大名，璿是玉之別稱。璣衡俱以玉飾，但史之立文，不可以玉璣、玉衡一指玉體，一指玉名，猶\CJKunderwave{左傳}雲“瓊弁玉纓”。所以變其文，傳以總言玉名，故云“美玉”，其實玉衡亦美玉也。\CJKunderwave{易·賁卦}彖云：“觀乎天文以察時變。”日月星宿運行於天，是為天之文也。璣衡者，璣為轉運，衡為橫簫，運璣使動，於下以衡望之，是“王者正天文之器”。漢世以來，謂之渾天儀者是也。馬融云：“渾天儀可旋轉，故曰璣。衡，其橫簫，所以視星宿也。以璿為璣，以玉為衡，蓋貴天象也。”蔡邕雲“玉衡長八尺，孔徑一寸，下端望之以視星辰。蓋懸璣以象天而衡望之，轉璣窺衡以知星宿”，是其說也。“七政”，其政有七，於璣衡察之，必在天者，知“七政”謂日月與五星也。木曰歲星，火曰熒惑星，土曰鎮星，金曰太白星，水曰辰星。\CJKunderwave{易·繫辭}云：“天垂象，見吉凶，聖人象之。”此日月五星有吉凶之象，因其變動為佔，七者各自異政，故為七政。得失由政，故稱政也。舜既受終，乃察璣衡，是“舜察天文，齊七政，以審已之受禪當天心與否”也。馬融云：“日月星皆以璿璣玉衡度知其盈縮、進退、失政所在。聖人謙讓，猶不自安，視璿璣玉衡以驗齊日月五星行度，知其政是與否，重審已之事也。”上天之體，不可得知，測天之事，見於經者唯有此“璿璣玉衡”一事而已。蔡邕\CJKunderwave{天文志}云：“言天體者有三家，一曰周髀，二曰宣夜，三曰渾天。宣夜絕無師說。周髀術數具在，考驗天象，多所違失，故史官不用。惟渾天者近得其情，今史所用候臺銅儀，則其法也。”虞喜云：“宣，明也。夜，幽也。幽明之數，其術兼之，故曰宣夜。”但絕無師說，不知其狀如何。周髀之術以為天似覆盆，蓋以斗極為中，中高而四邊下，日月旁行繞之。日近而見之為晝，日遠而不見為夜。渾天者以為地在其中，天周其外，日月初登於天，後入於地。晝則日在地上，夜則日入地下。王蕃\CJKunderwave{渾天說}曰：“天之形狀似鳥卵，天包地外，猶卵之裡黃，圓如彈丸，故曰渾天。”言其形體渾渾然也。其術以為天半覆地上，半在地下。其天居地上，見有一百八十二度半強，地下亦然。北極出地上三十六度，南極入地下亦三十六度，而嵩高正當天之中極，南五十五度當嵩高之上。又其南十二度為夏至之日道，又其南二十四度為春秋分之日道，又其南二十四度為冬至之日道，南下去地三十一度而已。是夏至日北去極六十七度，春秋分去極九十一度，冬至去極一百一十五度，此其大率也。其南北極持其兩端，其天與日月星宿斜而回轉，此必古有其法，遭秦而滅。楊子\CJKunderwave{法言}云：“或問渾天，曰，落下閎營之，鮮于妄人度之，耿中丞象之，幾乎，幾乎，莫之能違也！”是楊雄之意，以唬骸半天而問之也。閎與妄人，武帝時人。宣帝時司農中丞耿壽昌始鑄銅為之象，史官施用焉。後漢張衡作\CJKunderwave{靈憲}以說其狀。蔡邕、\CJKunderline{鄭玄}、陸績、吳時王藩、晉世姜岌、張衡、葛洪皆論渾天之義，並以渾說為長。江南宋元嘉年皮延宗又作是\CJKunderwave{渾天論}，太史丞錢樂鑄銅作渾天儀，傳於齊梁，周平、江陵遷其器於長安，今在太史書矣。衡長八尺，璣徑八尺，圓周二丈五尺強，轉而望之，有其法也。 \par}

肆類於上帝，\footnote{堯不聽舜讓，使之攝位。舜察天文，考齊七政而當天心，故行其事。肆，遂也。類謂攝位事類。遂以攝告天及五帝。王云：“上帝，天也。”馬云：“上帝，太一神，在紫微宮，天之最尊者。”}禋於六宗，\footnote{精意以享謂之禋。宗，尊也。所尊祭者，其祀有六，謂四時也、寒暑也、日也、月也、星也、水旱也。祭亦以攝告。禋音因，王云：“絜祀也。”馬云：“精意以享也。”六宗，王云：“四時、寒暑、日、月、星、水旱也。”馬云：“天地四時也。”}望于山川，遍於群神。\footnote{九州名山大川、五嶽四瀆之屬，皆一時望祭之。群神謂丘陵、墳衍、古之聖賢，皆祭之。墳，扶雲反。衍音演。}

{\noindent\zhuan\zihao{6}\fzbyks 傳“堯不”至“五帝”。正義曰：傳以既受終事，又察璣衡，方始祭於群神，是舜察天文，考齊七政，知己攝位而當於天心，故行其天子之事也。\CJKunderwave{祭法}云：“有而下者祭百神”遍祭群神是天子事也。“肆”是縱緩之言，此因前事而行後事，故以“肆”為遂也。“類”謂攝位事類，既知攝當天心，遂以攝位事類告天帝也。此“類”與下“禋”、“望”相次，當為祭名。\CJKunderwave{詩}雲“是類是禡”，\CJKunderwave{周禮·肆師}雲“類造上帝”，\CJKunderwave{王制}雲“天子將出類乎上帝”，所言“類”者皆是祭天之事，言以事類而祭也。\CJKunderwave{周禮·小宗伯}云：“天地之大災，類社稷，則為位。”是類之為祭，所及者廣。而傳雲“類謂攝位事類”者，以攝位而告祭，故類為祭名。\CJKunderwave{周禮·司服}云：“王祀昊天上帝，則服大裘而冕，祀五帝亦如之。”是昊天外更有五帝，上帝可以兼之，故以“告天及五帝”也。\CJKunderline{鄭玄}篤信讖緯，以為“昊天上帝謂天皇大帝，北辰之星也。五帝謂靈威仰等，太微宮中有五帝座星是也”。如鄭之言，天神有六也。\CJKunderwave{家語}云：“季康子問五帝之名，\CJKunderline{孔子}曰：‘天有五行：金木水火土。分時化育，以成萬物，其神謂之五帝。’”王肅云：“五行之神，助天理物者也。”孔意亦當然矣。此經惟有祭天，不言祭地及社稷，必皆祭之，但史略文耳。 \par}

{\noindent\zhuan\zihao{6}\fzbyks 傳“精意”至“攝告”。正義曰：\CJKunderwave{國語}云：“精意以享禋也。”\CJKunderwave{釋詁}云：“禋，祭也。”孫炎曰：“禋，絜敬之祭也。”\CJKunderwave{周禮·大宗伯}云：“以禋祀祀昊天上帝，以實柴祀日月星辰，以槱燎祀司中、司命、風師、雨師。”鄭云：“禋之言禋，周人尚臭，煙氣之臭聞者也。”鄭以“禋祀”之文在“燎”、“柴”之上，故以“禋”為此解耳。而\CJKunderwave{洛誥}雲“秬鬯二卣，曰‘明禋’”，又曰“禋於\CJKunderline{文王}、\CJKunderline{武王}”，又曰“王賓殺禋咸格”。經傳之文，此類多矣，非燔柴祭之也，知“禋”是精誠絜敬之名耳。“宗”之為尊，常訓也。名曰“六宗”，明是所尊祭者有六，但不知六者為何神耳。\CJKunderwave{祭法}云：“埋少牢於太昭，祭時。相近於坎壇，祭寒暑。王宮，祭日。夜明，祭月。幽禜，祭星。雩禜，祭水旱也。”據此言六宗，彼祭六神，故傳以彼六神謂此六宗。必謂彼之所祭是此六宗者，彼文上有祭天、祭地，下有山谷、丘陵，此“六宗”之文在上帝之下，山川之上，二者次第相類，故知是此六宗。王肅亦引彼文乃云：“禋於六宗，此之謂矣。”\CJKunderline{鄭玄}注彼云：“四時謂陰陽之神也。”然則陰陽寒暑水旱各自有神，此言“禋於六宗”，則六宗常禮也。禮無此文，不知以何時祀之。鄭以彼皆為祈禱之祭，則不可用\CJKunderline{鄭玄}注以解此傳也。漢世以來，說六宗者多矣。歐陽及大小夏侯說\CJKunderwave{尚書}皆云：“所祭者六，上不謂天，下不謂地，旁不謂四方，在六者之間，助陰陽變化，實一而名六宗矣。”孔光、劉歆以“六宗謂乾坤六子：水火雷風山澤也”。賈逵以為：“六宗者，天宗三，日月星也；地宗三，河海岱也。”馬融云：“萬物非天不覆，非地不載，非春不生，非夏不長，非秋不收，非冬不藏，此其謂六也。”\CJKunderline{鄭玄}以六宗言“禋”，與祭天同名，則六者皆是天之神祗，謂“星、辰、司中、司命、風師、雨師。星謂五緯星，辰謂日月所會十二次也。司中、司命文昌第五第四星也。風師，箕也。雨師，畢也”。晉初幽州秀才張髦上表云：“臣謂禋於六宗，祀祖考所尊者六，三昭三穆是也。”司馬彪又上表云，歷難諸家及自言己意“天宗者，日月星辰寒暑之屬也；地宗，社稷五祀之屬也；四方之宗，四時五帝之屬”。惟王肅據\CJKunderwave{家語}六宗與孔同。各言其志，未知孰是。司馬彪\CJKunderwave{續漢書}云：“安帝元初六年，立六宗祠於洛陽城西北亥地，祀比大社，魏亦因之。晉初荀顗定新祀，以六宗之神諸說不同廢之。摯虞駁之，謂：‘宜依舊，近代以來皆不立六宗之祠也。’” \par}

{\noindent\zhuan\zihao{6}\fzbyks 傳“九州”至“祭之”。正義曰：“望于山川”，大總之語，故知九州之內所有名山大川、五嶽四瀆之屬皆一時望祭之也。\CJKunderwave{王制}云：“名山大川不以封。”山川大，乃有名，是“名”、“大”互言之耳。\CJKunderwave{釋山}云：“泰山為東嶽，華山為西嶽，霍山為南嶽，恆山為北嶽，嵩高山為中嶽。”\CJKunderwave{白虎通}云：“嶽者何？捔也，捔考功德也。”應劭\CJKunderwave{風俗通}云：“嶽者，捔考功德黜陟也。”然則四萬方有一大山，天子巡守至其下，捔考諸侯功德而黜陟之，故謂之“嶽”。\CJKunderwave{釋水}云：“江河淮濟為四瀆。四瀆者，發源注海者也。”\CJKunderwave{釋名}云：“瀆，獨也，各獨出其水而入海也。”嶽是名山，瀆是大川，故先言名山大川，又舉嶽瀆以見之。嶽瀆之外猶有名山大川，故言“之屬”以包之。\CJKunderwave{周禮·大司樂}云：“四鎮五嶽崩，令去樂。”鄭云：“四鎮，山之重大者，謂揚州之會稽山，青州之沂山，幽州醫無閭山，冀州之霍山。”是五嶽之外名山也。\CJKunderwave{周禮·職方氏}每州雲“其川”、“其浸”，若雍州雲“其川涇、汭，其浸渭、洛”，如此之類，是四瀆之外大川也。言“遍於群神”，則神無不遍，故“群神謂丘陵、墳衍、古之聖賢，皆祭之”。\CJKunderwave{周禮·大司樂}云：“凡六樂者，一變而致川澤之示，再變而致山林之示，三變而致丘陵之示，四變而致墳衍之示。”\CJKunderline{鄭玄}\CJKunderwave{大司徒}注云：“積石曰山，竹木曰林，注瀆曰川，水鍾曰澤，土高曰丘，大阜曰陵，水崖曰墳，下平曰衍。”此傳舉“丘陵墳衍”則林澤亦包之矣。“古之聖賢”謂\CJKunderwave{祭法}所云“在祀典”者，黃帝、顓頊、句龍之類皆祭之也。 \par}

輯五瑞，既月乃日，覲四嶽群牧，班瑞於\textcolor{red}{群后}。\footnote{輯，斂。既，盡。覲，見。班，還。後，君也。舜斂公侯伯子男之瑞圭璧，盡以正月中，乃日日見四嶽及九州牧監，還五瑞於諸侯，與之正始。輯，徐音集，王雲合，馬云：“斂也。”瑞，垂偽反，信也。牧，牧養之牧，徐音目。}

{\noindent\zhuan\zihao{6}\fzbyks 傳“輯斂”至“正始”。正義曰：“覲，見”，“後，君”，\CJKunderwave{釋詁}文。\CJKunderwave{釋言}云：“輯，合也。”“輯”是合聚之義，故為斂也。日月食盡謂之既，是“既”為盡也。\CJKunderwave{釋言}云：“班，賦也。”孫炎曰：“謂布與也。”“輯”是斂聚，“班”為散佈，故為還也。下雲“班瑞於群后”，則知“輯”者從群后而斂之，故云“舜斂公侯伯子男之瑞圭璧”也。\CJKunderwave{周禮·典瑞}云：“公執桓圭，侯執信圭，伯執躬圭，子執谷璧，男執蒲璧。”是圭璧為五等之瑞。諸侯執之以為王者瑞信，故稱“瑞”也。舜以朔日受終於文祖，又遍祭群神及斂五瑞，則入月以多日矣。“盡以正月中”謂從斂瑞以後至月末也。“乃日日見四嶽及九州牧監”，舜初攝位，當發號出令，日日見之，與之言也。州牧各監一州諸侯，故言“監”也。更復還五瑞於諸侯者，此瑞本受於堯，斂而又還之，若言舜新付之，改為舜臣，與之正新君之始也。 \par}

{\noindent\shu\zihao{5}\fzkt “正月”至“群后”。正義曰：舜既讓而不許，乃以堯禪之。明年正月上日，受堯終帝位之事於堯文祖之廟。雖受堯命，猶不自安。又以璿為璣、以玉為衡者，是為王者正天文之器也，乃復察此璿璣玉衡，以齊整天之日月五星七曜之政。觀其齊與不齊，齊則受之是也，不齊則受之非也。見七政皆齊，知己受為是，遂行為帝之事，而以告攝事類祭於上帝，祭昊天及五帝也。又禋祭於六宗等尊卑之神，望祭於名山大川、五嶽四瀆，而又遍祭于山川、丘陵、墳衍、古之聖賢之群神，以告已之受禪也。告祭既畢，乃斂公侯伯子男五等之瑞玉。其圭與璧悉斂取之盡，以正月之中，乃日月見四嶽及群牧，既而更班所斂五瑞於五等之群后，而與之更始，見己受堯之禪，行天子之事也。 \par}

\textcolor{red}{歲二月}，東巡守，至於岱宗，柴，\footnote{諸侯為天子守土，故稱守，巡行之。既班瑞之明月，乃順春東巡。岱宗，泰山，為四嶽所宗。燔柴祭天告至。○巡，似遵反，徐養純反。守，時救反，本或作狩。岱音代，泰山也。柴，士皆反。\CJKunderwave{爾雅}：“祭天曰燔柴。”馬曰：“祭時積柴，加牲其上而燔之。”行，下孟反。燔，扶袁反，又扶雲反。}望秩于山川，\footnote{東嶽諸侯竟內名山大川如其秩次望祭之。謂五嶽牲禮視三公，四瀆視諸侯，其餘視伯子男。○瀆，徒木反。}肆覲東後。\footnote{遂見東方之國君。}

{\noindent\zhuan\zihao{6}\fzbyks 傳“諸侯”至“告至”。正義曰：王者所為巡守者，以諸侯自專一國，威福在己，恐其擁遏上命，澤不下流，故時自巡行，問民疾苦。\CJKunderwave{孟子}稱晏子對齊景公云：“天子適諸侯曰巡守。巡守者，巡所守也。”是言天子巡守主謂巡行諸侯，故言諸侯為天子守土，故稱守，而往巡行之。定四年\CJKunderwave{左傳}祝鮀言衛國“取相土之東都，以會王之東搜”，“搜”是獵之名也。王者因巡諸侯，或亦獵以教戰，其守皆作“狩”。\CJKunderwave{白虎通}云：“王者所以巡狩者何？巡者循也，狩者收也，為天子循收養人。”彼因名以附說，不如晏子之言得其本也。正月班瑞，二月即行，故云“既班瑞之明月，乃順春東巡”。春位在東，故“順春”也。\CJKunderwave{爾雅}：“泰山為東嶽。”此巡守至於岱，岱之與泰，其山有二名也。\CJKunderwave{風俗通}云：“泰山，山之尊者，一曰岱宗。岱，始也。宗，長也。萬物之始，陰陽交代，故為五嶽之長。”是解岱即泰山，為四嶽之宗，稱岱宗也。\CJKunderwave{郊特牲}云：“天子適四方，先柴。”是燔柴為祭天告至也。 \par}

{\noindent\zhuan\zihao{6}\fzbyks 傳“東嶽”至“子男”。正義曰：四時各至其方岳，望祭其方岳山川，故云“東嶽諸侯境內名山大川如其秩次望祭之”也。言秩次而祭，知遍於群神，故云“五嶽牲禮視三公，四瀆視諸侯，其餘視伯子男”也。其尊卑所視，\CJKunderwave{王制}及\CJKunderwave{書傳}之文，“牲禮”二字孔增之也。諸侯五等，三公為上等，諸侯為中等，伯子男為下等，則所言諸侯，惟謂侯爵者耳。其言所視，蓋視其祭祀。祭五嶽如祭三公之禮，祭四瀆如祭諸侯之禮，祭山川如祭伯子男之禮。公侯伯子男尊卑既有等級，其祭禮必不同，但古典亡滅，不可復知。\CJKunderline{鄭玄}注\CJKunderwave{書傳}云：“所視者，謂其牲幣粢盛籩豆爵獻之數。”案五等諸侯適天子皆膳用太牢，禮諸侯祭皆用太牢，無上下之別。又\CJKunderwave{大行人}云，上公九獻，侯伯七獻，子男五獻。\CJKunderwave{掌客}上公饔餼九牢、飧五牢，侯伯饔餼七牢、飧四牢，子男饔餼五牢、飧三牢。又上公豆四十，侯伯三十二，子男二十四。並\CJKunderline{伯與}侯同。又鄭注\CJKunderwave{禮器}“四望”、“五獻”據此諸文。與孔傳\CJKunderwave{王制}不同者，\CJKunderwave{掌客}、\CJKunderwave{行人}自是周法，孔與\CJKunderwave{王制}先代之禮。必知然者，以\CJKunderwave{周禮}侯與伯同，\CJKunderwave{公羊}及\CJKunderwave{左氏傳}皆以公為上，伯子男為下，是其異也。 \par}

協時月正日，同律度量衡。\footnote{合四時之氣節，月之大小，日之甲乙，使齊一也。律法制及尺丈、斛鬥、斤兩，皆均同。○同律，王云：“同，齊也。律，六律也。”馬云：“律，法也。”鄭云：“陰呂陽律也。”度如字，丈尺也。量，力尚反，鬥斛也。衡，稱也。}修五禮、五玉、\footnote{修吉、兇、賓、軍、嘉之禮。五等諸侯執其玉。}三帛、二生、一死贄，\footnote{三帛，諸侯世子執纁，公之孤執玄，附庸之君執黃。二生，卿執羔，大夫執雁。一死，士執雉。玉、帛、生、死，所以為贄以見之。贄音至，本又作摯。纁，許雲反。}如五器，卒乃復。\footnote{卒，終。復，還也。器謂圭璧。如五器，禮終則還之。三帛、生、死則否。復，扶又反，下同。還音旋。}

{\noindent\zhuan\zihao{6}\fzbyks 傳“合四”至“均同”。正義曰：上篇已訓“協”為合，故注即以合言之也，他皆仿此。\CJKunderwave{周禮·太史}云：“正歲年，頒告朔於邦國。”則節氣晦朔皆天子頒之。猶恐諸侯國異或不齊同，故因巡守而合和之。節是月初，氣是月半也。\CJKunderwave{世本}云：“容成作歷。大撓作甲子。”二人皆黃帝之臣，蓋自黃帝已來始用甲子紀日，每六十日而甲子一周。\CJKunderwave{史記}稱紂為長夜之飲，忘其日辰。恐諸侯或有此之類，故須合日之甲乙也。時也，月也，日也，三者皆當勘檢諸國使齊一也。“律”者候氣之管，而度量衡三者法制皆出於律，故云“律法制”也。度有丈尺，量有斛鬥，衡有斤兩，皆取法於律，故孔解“律”為法制，即雲“及尺丈、斛鬥、斤兩皆均同之”。\CJKunderwave{漢書日·律曆志}云，度量衡出於黃鐘之律也。度者，分、寸、尺、丈、引，所以度長短也。本起於黃鐘之管長。以子谷秬黍中者，以一黍之廣度之，千二百黍為一分。十分為寸，十寸為尺，十尺為丈，十丈為引，而五度審矣。量謂龠、合、升、鬥、斛，所以量多少也。本起於黃鐘之龠，以子谷秬黍中者千有二百實為一龠。十龠為合，十合為升，十升為鬥，十鬥為斛，而五量嘉矣。權者，銖、兩、斤、鈞、石，所以稱物知輕重也。本起於黃鐘之龠，一龠容千二百黍，重十二銖，兩之為兩，十六兩為斤，三十斤為鈞，四鈞為石，而五權謹矣。權、衡一物。衡，平也；權，重也；稱上謂之衡，稱錘謂之權，所從言之異耳。如彼\CJKunderwave{志}文，是度量衡本起於律也。時月言“協”，日言“正”，度量衡言“同”者，以時月須與他月和合，故言“協”；日有正與不正，故言“正”；度量衡俱是明之，所用恐不齊同，故言“同”；因事宜而變名耳。 \par}

{\noindent\zhuan\zihao{6}\fzbyks 傳“修吉”至“其玉”。正義曰：\CJKunderwave{周禮·大宗伯}云：“以吉禮事邦國之鬼神示，以凶禮哀邦國之憂，以賓禮親邦國，以軍禮同邦國，以嘉禮親萬民之昏姻。”知“五禮”謂此也。帝王之名既異，古今之禮或殊，而以周之五禮為此“五禮”者，以帝王相承，事有損益，後代之禮亦當是前代禮也。且歷驗此經，亦有五事：此篇“類於上帝”，吉也；“如喪考妣”，兇也；“群后四朝”，賓也；\CJKunderwave{大禹謨}雲“汝徂徵”，軍也；\CJKunderwave{堯典}雲“女於時”，嘉也。五禮之事，並見於經，知與後世不異也。此雲“五玉”，即上文“五瑞”，故知“五等諸侯執其玉”也。\CJKunderline{鄭玄}云：“執之曰瑞，陳列曰玉。” \par}

{\noindent\zhuan\zihao{6}\fzbyks 傳“諸侯”至“執黃”。正義曰：\CJKunderwave{周禮·典命}云：“凡諸侯之適子，誓於天子，攝其君，則下其君之禮一等。未誓，則以皮帛繼子男之下。公之孤四命，以皮帛視小國之君。”是諸侯世子、公之孤執帛也。附庸雖則無文，而為南面之君，是一國之主，春秋時附庸之君適魯皆稱“來朝”，未有爵命，不得執玉，則亦繼小國之君同執帛也。經言“三帛”，必有三色，所云纁、玄、黃者，孔時或有所據，末知出何書也。王肅云：“三帛，纁、玄、黃也。附庸與諸侯之適子、公之孤執皮帛，其執之色未詳。聞或曰孤執玄，諸侯之適子執纁，附庸執黃。”王肅之注\CJKunderwave{尚書}，其言多同孔傳。\CJKunderwave{周禮}孤與世子皆執皮帛，\CJKunderline{鄭玄}云：“皮帛者，束帛而表，之以皮為之飾。皮，虎豹皮也。”此三帛不言皮，蓋於時未以皮為飾。 \par}

{\noindent\zhuan\zihao{6}\fzbyks 傳“卿執”至“執雉”。正義曰：此皆\CJKunderwave{大宗伯}文也。\CJKunderline{鄭玄}曰：“羔，小羊，取其群而不失其類也。雁，取其候時而行也。雉，取其守介，死不失節也。\CJKunderwave{曲禮}雲‘飾羔雁者以繢’，謂衣之以布而又畫之。雉執之無飾。\CJKunderwave{士相見之禮}，卿大夫飾贄以布。不言繢，此諸侯之臣與天子之臣異也。”鄭之此言，論周之禮耳，虞時每事猶質，羔雁不必有飾。 \par}

{\noindent\zhuan\zihao{6}\fzbyks 傳“玉帛”至“見之”。正義曰：\CJKunderwave{曲禮}云：“贄，諸侯圭，卿羔，大夫雁，士雉。”雉不可生，知“一死”是雉，“二生”是羔、雁也。\CJKunderline{鄭玄}云：“贄之言至，所執以自至也。”自“五玉”以下，蒙上“修”文者，執之使有常也。若不言“贄”，則不知所用，故言“贄”以結上，又見玉、帛、生、死皆所以為贄，以見君與自相見，其贄同也。 \par}

{\noindent\zhuan\zihao{6}\fzbyks 傳“卒終”至“則否”。正義曰：“卒，終”，\CJKunderwave{釋詁}文。\CJKunderwave{釋言}云：“還、復，返也。”是還、復同義，故為還也。“五器”文在“贄”下，則是贄內之物。\CJKunderwave{周禮·大宗伯}云：“以玉作六器。”知“器謂圭璧”，即五玉是也。如，若也。言諸侯贄之內，若是五器，禮終乃還之，如三帛、生、死則不還也。\CJKunderwave{聘義}云：“以圭璋聘，重禮也。已聘而還圭璋，此輕財而重禮之義也。”\CJKunderwave{聘義}主於說聘，其朝禮亦然。\CJKunderwave{周禮·司儀}云：“諸公相見為賓，還圭，如將幣之儀。”是圭璧皆還之也。\CJKunderwave{士相見禮}言大夫以下見國君之禮云：“若他邦之人，則使擯者還其贄。”己臣皆不還其贄，是“三帛、生、死則否”。 \par}

五月南巡守，至於南嶽，如岱禮。\footnote{南嶽，衡山。自東嶽南巡，五月至。}八月西巡守，至於西嶽，如初。\footnote{西嶽，華山。初謂岱宗。華，戶化反。華山在弘農。}十有一月朔巡守，至於北嶽，如西禮。\footnote{北嶽，恆山。有，如字，徐於救反。如西禮，方興本同，馬本作“如初”。}歸，格於藝祖，用特。\footnote{巡守四嶽然後歸，告至文祖之廟。藝，文也。言祖則考著。特，一牛。藝，魚世反，馬、王云：“禰也。”}五載一巡守，群后四朝。\footnote{各會朝於方岳之下，凡四處，故曰“四朝”。將說“敷奏”之事，故申言之。堯舜同道，舜攝則然，堯又可知。四朝，馬、王皆云：“四面朝於方岳之下。”鄭云：“四朝，四季朝京師也。”朝音直遙反，注同。}敷奏以言，明試以功，車服\textcolor{red}{以庸}。\footnote{敷，陳。奏，進也。諸侯四朝，各使陳進治禮之言。明試其言，以要其功，功成則賜車服以表顯其能用。敷音孚。}

{\noindent\zhuan\zihao{6}\fzbyks 傳“南嶽”至“月至”。正義曰：\CJKunderwave{釋山}云：“河南華，河東岱，河北恆，江南衡。”李巡云：“華，西嶽華山也。岱，東嶽泰山也。恆，北嶽恆山也。衡，南嶽衡山也。”郭璞云：“恆山一名常山，避漢文帝諱。”\CJKunderwave{釋山}又云：“泰山為東嶽，華山為西嶽，霍山為南嶽，恆山為北嶽。”岱之與泰，衡之與霍，皆一山而有兩名也。張揖云：“天柱謂之霍山。”\CJKunderwave{漢書·地理志}云，天柱在廬江灊縣，則霍山在江北。而與江南衡為一者，郭璞\CJKunderwave{爾雅}注云：“霍山今在廬江灊縣，潛水出焉，別名天柱山。漢武帝以衡山遼曠，故移其神於此。今其彼土俗人皆呼之為南嶽。南嶽本自以兩山為名，非從近來也。而學者多以霍山不得為南嶽，又云漢武帝來始乃名之。即如此言，謂武帝在\CJKunderwave{爾雅}前乎？斯不然矣。”是解衡、霍二名之由也。書傳多雲“五嶽”，以嵩高為中嶽，此雲“四嶽”者，明巡守至於四嶽故也。\CJKunderwave{風俗通}云：“泰山，山之尊者，一曰岱宗。岱，始也。宗，長也。萬物之始，陰陽交代，故為五嶽之長。王者受命恆封禪之。衡山一名霍山，言萬物霍然大也。華，變也，萬物變由西方也。恆，常也，萬物伏北方有常也。”二月至於岱宗，不指“嶽”名者，巡守之始，故詳其文，三時言嶽名，明岱亦是嶽，因事宜而互相見也。四巡之後乃雲“歸，格”，則是一出而周四嶽。故知自東嶽而即南行，以五月至也。王者順天道以行人事，故四時之月各當其時之中，故以仲月至其嶽。上雲“歲二月，東巡守”，以二月始發者，此四時巡守之月，皆以至嶽為文，東巡以二月至，非發時也，但舜以正月有事，二月即發行耳。\CJKunderline{鄭玄}以為“每嶽禮畢而歸，仲月乃復更去”。若如鄭言，當於東巡之下即言“歸，格”，後以“如初”包之，何當北巡之後始言歸乎？且若來而復去，計程不得周遍，此事不必然也。其經南雲“如岱禮”，西雲“如初”，北雲“如西禮”者，見四時之禮皆同，互文以明耳。不巡中嶽者，蓋近京師，有事必聞，不慮枉滯，且諸侯分配四方，無屬中嶽，故不須巡之也。“朔巡守”。正義曰：\CJKunderwave{釋訓}云：“朔，北方也。”故\CJKunderwave{堯典}及此與\CJKunderwave{禹貢}皆以“朔”言北，史變文耳。 \par}

{\noindent\zhuan\zihao{6}\fzbyks 傳“巡守”至“一牛”。正義曰：此承四巡之下，是巡守既遍，然後歸也。以上受終在文祖之廟，知此以“告至文祖之廟”。才、藝、文、德其義相通，故“藝”為文也。“文祖”、“藝祖”史變文耳。\CJKunderwave{王制}說巡守之禮云：“歸，格於祖禰，用特。”此不言“禰”，故傳推之。“言祖則考著”，考近於祖，舉尊以及卑也。“特”者獨也，故為“一牛”。此惟言“文祖”，故云“一牛”。遍告諸廟，廟用一牛，故鄭注：“彼雲祖下及禰皆一牛也。”此時舜始攝位，未自立廟，故知告堯之文祖也。 \par}

{\noindent\zhuan\zihao{6}\fzbyks 傳“各會”至“可知”。正義曰：此總說巡守之事，而言“群后四朝”，是言四方諸侯各自會朝於方岳之下。凡四處別朝，故云“四朝”。上文“肆覲東後”是為一朝，四嶽禮同，四朝見矣。計此不宜須重言之。為將說“敷奏”之事，“敷奏”因朝而為，故申言之。申，重也。此是巡守大法，文在舜攝位之時，嫌堯本不然，故云“堯舜同道，舜攝則然，堯又可知”也。堯法已然，舜無增改，而言此以美舜者，道同於堯，足以為美，故史錄之。 \par}

{\noindent\zhuan\zihao{6}\fzbyks 傳“敷陳”至“能用”。正義曰：“敷”者佈散之言，與陳設義同，故為陳也。“奏”是進上之語，故為進也。諸侯四處來朝，每朝之處，舜各使陳進其治理之言，令自說巳之治政。既得其言，乃依其言明試之，以要其功。以如其言，即功實成，則賜之車服，以表顯其人有才能可用也。人以車服為榮，故天子之賞諸侯，皆以車服賜之。\CJKunderwave{覲禮}雲“天子賜侯氏以車服”是也。 \par}

{\noindent\shu\zihao{5}\fzkt “歲二月”至“以庸”。正義曰：舜既班瑞群后，即以其歲二月東行巡省守土之諸侯，至於岱宗之嶽，燔柴告至，又望而以秩次祭於其方岳山川。柴望既畢，遂以禮見東方諸侯諸國之君，於此諸國協其四時氣節、月之大小，正其日之甲乙，使之齊一。均同其國之法制、度之丈尺、量之斛鬥、衡之斤兩，皆使齊同，無輕重大小。又修五禮：吉、兇、賓、軍、嘉之禮。修五玉：公侯伯子男所執之圭璧也。又修三帛：諸侯世子、公之孤、附庸之君所執玄、纁、黃之帛也。又修二生：卿所執羔、大夫所執雁也。又修一死：士所執雉也。自“五玉”至於“一死”，皆蒙上“修”文，總言所用。玉、帛、生、死皆為贄以見天子也。其贄之內，如五玉之器，禮終乃復還之。其帛與生、死則不還也。東嶽禮畢，即向衡山，五月南巡守，至於南嶽之下，柴望以下一如岱宗之禮。南嶽禮畢，即向華山，八月西巡守，至於西嶽之下，其禮如初時，如岱宗所行。西嶽禮畢，即向恆山。朔，北也。十有一月北巡守，至於北嶽之下，一如西嶽之禮。巡守既周，乃歸京師。藝，文也。至於文祖之廟，用特牛之牲設祭以告巡守歸至也。從是以後每五載一巡守，其巡守之年，諸侯群后四方各朝天子於方岳之下。其朝之時，各使自陳進其所以治化之言。天子明試其言，以考其功，功成有驗，則賜之車服，以表顯其有功，能用事。 \par}

\textcolor{red}{肇十}有二州，\footnote{肇，始也。\CJKunderline{禹}治水之後，舜分冀州為幽州、幷州，分青州為營州，始置十二州。肇音兆。十有二州，謂冀、兗、青、徐、荊、楊、豫、梁、雍、並、幽、營也。}封十有二山，濬川。\footnote{封，大也。每州之名山殊大之，以為其州之鎮。有流川則深之，使通利。濬，荀俊反。}

{\noindent\zhuan\zihao{6}\fzbyks 傳“肇始”至“二州”。正義曰：“肇，始”，\CJKunderwave{釋詁}文。\CJKunderwave{禹貢}治水之時，猶為九州，今始為十二州，知“\CJKunderline{禹}治水之後”也。\CJKunderline{禹}之治水，通\CJKunderline{鯀}九載，為作十有三載，則舜攝位元年，九州始畢。當是二年之後，以境界太遠，始別置之。知“分冀州為幽州、幷州”者，以王者廢置，理必相沿。\CJKunderwave{周禮·職方氏}九州之名有幽、並，無徐、梁。周立州名必因於古，知舜時當有幽、並。\CJKunderwave{職方}幽、並山川於\CJKunderwave{禹貢}皆冀州之域，知分冀州之域為之也。\CJKunderwave{爾雅·釋地}九州之名於\CJKunderwave{禹貢}無樑、青，而有幽、營，雲“燕曰幽州，齊曰營州”。孫炎以\CJKunderwave{爾雅}之文與\CJKunderwave{職方}、\CJKunderwave{禹貢}並皆不同，疑是殷制。則營州亦有所因，知舜時亦有營州。齊即青州之地，知分青州為之。於此居攝之時，始置十有二州，蓋終舜之世常然。宣三年\CJKunderwave{左傳}云：“昔夏之方有德也，貢金九牧。”則\CJKunderline{禹}登王位，還置九州，其名蓋如\CJKunderwave{禹貢}，其境界不可知也。 \par}

{\noindent\zhuan\zihao{6}\fzbyks 傳“封大”至“通利”。正義曰：\CJKunderwave{釋詁}云：“冢，大也。”舍人曰：“冢，封之大也。”定四年\CJKunderwave{左傳}雲“封豕長蛇”，相對是“封”為大也。\CJKunderwave{周禮·職方氏}每州皆雲“其山鎮曰某山”，揚州會稽，荊州衡山，豫州華山，雍州吳山，冀州霍山，幷州恆山，幽州醫無閭，青州沂山，兗州岱山，是周時九州之內最大之山。舜時十有二山，事亦然也。州內雖有多山，取其最高大者，以為其州之鎮，特舉其名，是殊大之也。其有川，無大無小，皆當深之，故云“濬川”，有“流川則深之，使通利也。”\CJKunderwave{職方氏}每州皆雲其川、其浸，亦舉其州內大川，但令小大俱通，不復舉其大者，故直雲濬之而已。 \par}

象以典刑，\footnote{象，法也。法用常刑，用不越法。}流宥五刑，\footnote{宥，寬也。以流放之法寬五刑。宥音又，馬云：“宥，二宥也。”}鞭作官刑，\footnote{鞭以作為治官事之刑。}撲作教刑，\footnote{撲，榎楚也。不勤道業則撻之。撲，普卜反，徐敷卜反。榎，皆雅反。}金作贖刑。\footnote{金，黃金。誤而入刑，出金以贖罪。○贖，石欲反，徐音樹。}

{\noindent\zhuan\zihao{6}\fzbyks 傳“象法”至“越法”。正義曰：\CJKunderwave{易·繫辭}云：“象也者，象此者也。”又曰：“天垂象，聖人則之。”是“象”為仿法，故為法也。五刑雖有常法，所犯未必當條，皆須原其本情，然後斷決。或情有差降，俱被重科；或意有不同，失出失入，皆是違其常法。故令依法用其常刑，用之使不越法也。 \par}

{\noindent\zhuan\zihao{6}\fzbyks 傳“宥寬”至“五刑”。正義曰：“宥，寬”，\CJKunderwave{周語}文，“流”謂徙之遠方；“放”，使生活；以流放之法寬縱五刑也。此惟解以流寬之刑，而不解宥寬之意。\CJKunderline{鄭玄}云：“其輕者或流放之，四罪是也。”王肅云：“謂君不忍刑殺，宥之以遠方。”然則知此是據狀合刑，而情差可恕，全赦則太輕，致刑即太重，不忍依例刑殺，故完全其體，宥之遠方。應刑不刑，是寬縱之也。上言“典刑”，此言“五刑”者，其法是常，其數則五，“象以典刑”謂其刑之也，“流宥五刑”謂其遠縱之也。“流”言“五刑”，則“典刑”亦五，其文互以相見。王肅云：“言宥五刑，則正五刑見矣。”是言二文相通之意也。“典刑”是其身，“流宥”離其鄉，流放致罪為輕，比鞭為重，故次“典刑”之下。先言“流宥”，鞭撲雖輕，猶虧其體，比於出金贖罪又為輕，且\CJKunderwave{呂刑}五罰雖主贖五刑，其鞭撲之罪亦容輸贖，故後言之。此正刑五與流宥鞭撲俱有常法，“典”字可以統之，故發首言“典刑”也。 \par}

{\noindent\zhuan\zihao{6}\fzbyks 傳“以鞭”至“之刑”。正義曰：此有鞭刑，則用鞭久矣。\CJKunderwave{周禮·滌狼氏}：“誓大夫曰，敢不關，鞭五百。”\CJKunderwave{左傳}有鞭徒人費、圉人犖是也，子玉使鞭七人，衛侯鞭師曹三百，日來亦皆施用。大隨造律，方使廢之。“治官事之刑”者，言若於官事不治則鞭之，蓋量狀加之，未必有定數也。 \par}

{\noindent\zhuan\zihao{6}\fzbyks 傳“撲榎”至“撻之”。正義曰：\CJKunderwave{學記}云：“榎楚二物，以收其威。”\CJKunderline{鄭玄}云：“榎，槄也。楚，荊也。二物可以撲撻犯禮者。”知“撲”是榎楚也。既言“以收其威”，知“不勤道業則撻之”。\CJKunderwave{益稷}云：“撻以記之。”又\CJKunderwave{大射}、\CJKunderwave{鄉射}皆雲司馬搢撲。則撲亦官刑，惟言“作教刑”者，官刑鞭撲俱用，教刑惟撲而已，故屬撲於教。其實官刑亦當用撲，蓋重者鞭之，輕者撻之。 \par}

{\noindent\zhuan\zihao{6}\fzbyks 傳“金黃”至“贖罪”。正義曰：此以“金”為黃金。\CJKunderwave{呂刑}“其罰百鍰”，傳為“黃鐵”。俱是贖罪而金鐵不同者，古之金銀銅鐵總號為金，別之四名耳。\CJKunderwave{釋器}云：“黃金謂之蕩，白金謂之銀。”是黃金白銀俱名金也。\CJKunderwave{周禮·考工記}攻金之工築氏為削，冶氏為殺矢，鳧氏為鍾，㮚氏為量，段氏為鎛，桃氏為劍，其所為者有銅，有鐵，是銅鐵俱名為金，則鐵名亦包銅矣。此傳“黃金”、\CJKunderwave{呂刑}“黃鐵”，皆是今之銅也。古之贖罪者皆用銅，漢始改用黃金，但少其斤兩，令與銅相敵。故\CJKunderline{鄭玄}駁\CJKunderwave{異義}言：“贖死罪千鍰，鍰六兩大半兩，為四百一十六斤十兩六半兩銅，與今贖死罪金三斤為價相依附。”是古贖罪皆用銅也。實謂銅而謂之金、鐵，知傳之所言謂銅為金、鐵耳。漢及後魏贖罪皆用黃金，後魏以金難得，合金一兩收絹十匹。今律乃復依古，死罪贖銅一百二十斤，於古稱為三百六十斤。孔以鍰為六兩，計千鍰為三百七十五斤，今贖輕於古也。誤而入罪，出金以贖，即律“過失殺傷人，各依其狀以贖論”是也。\CJKunderwave{呂刑}所言“疑赦”乃罰者，即今律“疑罪各從其實，以贖論”是也。疑謂虛實之證等，是非之理均，或事涉疑似，旁無證見，或雖有證見，事非疑似，如此之類，言皆為疑罪。疑而罰贖，\CJKunderwave{呂刑}已明言。誤而輸贖，於文不顯，故此傳指言誤而入罪以解此“贖”。鞭撲加於人身，可雲“撲作教刑”，金非加人之物，而言“金作贖刑”，出金之與受撲俱是人之所患，故得指其所出，以為刑名。 \par}

眚災肆赦，怙終賊刑。\footnote{眚，過。災，害。肆，緩。賊，殺也。過而有害，當緩赦之。怙奸自終，當刑殺之。眚，所景反。怙音戶。}欽哉，欽哉，惟刑之恤哉！\footnote{舜陳典刑之義，敕天下使敬之，憂欲得中。恤，峻律反，憂也。}

{\noindent\zhuan\zihao{6}\fzbyks 傳“眚過”至“殺之”。正義曰：\CJKunderwave{春秋}言“肆眚”者，皆謂緩縱過失之人，是“肆”爰緩也，“眚”爰過也。\CJKunderwave{公羊傳}雲“害物曰災”，是為害也。宣二年\CJKunderwave{左傳}晉侯殺趙盾，“使鉏麑賊之”，是“賊”為殺也。此經二句承上“典刑”之下，總言用刑之要。過而有害，雖據狀合罪，而原心非故，如此者當緩赦之。小則恕之，大則宥之，上言“流宥”、“贖刑”是也。怙恃奸詐，欺罔時人，以此自終，無心改悔，如此者當刑殺之。小者刑之，大者殺之，上言“典刑”及“鞭”、“撲”皆是也。經言“賊刑”，傳雲“刑殺”，不順經文者，隨便言之。 \par}

{\noindent\zhuan\zihao{6}\fzbyks 傳“舜陳”至“得中”。正義曰：此經二句，舜之言也。不言“舜曰”，以可知而略之。舜既制此典刑，又陳典刑之義，以敕天下百官，使敬之哉，敬之哉，惟刑之憂哉。憂念此刑，恐有濫失，欲使得中也。 \par}

流\CJKunderline{共工}於幽洲，\footnote{象恭滔天，足以惑世，故流放之。幽洲，北裔。水中可居者曰州。共音恭。\CJKunderwave{左傳}：“少皞氏有不才子，毀信廢忠，崇飾惡言，靖僣庸回，服讒搜慝，以誣盛德，天下之民謂之窮奇。”杜預云：“即\CJKunderline{共工}。”裔，以制反。}放\CJKunderline{驩兜}於崇山，\footnote{黨於\CJKunderline{共工}，罪惡同。崇山，南裔。驩，呼端反。兜，丁侯反。\CJKunderwave{左傳}：“帝鴻氏有不才子，掩義隱賊，好行兇德，醜類惡物，頑嚚不友，是與比周，天下之民謂之渾敦。”杜預云：“即\CJKunderline{驩兜}也。帝鴻，黃帝也。”}竄三苗於三危，\footnote{三苗，國名。縉雲氏之後，為諸侯，號饕餮。三危，西裔。竄，七亂反。三苗，馬、王云：“國名也。縉雲氏之後，為諸侯，蓋饕餮也。”\CJKunderwave{左傳}：“縉雲氏有不才子，貪於飲食，冒於貨賄，侵欲崇侈，不可盈厭，聚斂積實，不知紀極，不念孤寡，不恤窮匱，天下之民以比三兇，謂之饕餮。”杜預云：“縉云，黃帝時官名，非帝子孫，故以比三兇也。貪財曰饕，貪食曰餮。”縉音晉。饕，土刀反。餮，他節反。}殛\CJKunderline{鯀}於羽山，\footnote{方命圮族，績用不成，殛竄放流，皆誅也。異其文，述作之體。羽山，東裔，在海中。殛，紀力反。\CJKunderline{鯀}，故本反。\CJKunderwave{左傳}：“顓頊氏有不才子，不可教訓，不知話言，告之則頑，舍之則嚚，傲狼明德，以亂天常，天下之民謂之檮杌。”杜預云：“即\CJKunderline{鯀}也。檮杌，兇頑無儔匹之貌。”}四罪而天下\textcolor{red}{咸服}。\footnote{皆服舜用刑當其罪，故作者先敘典刑而連引四罪，明皆徵用所行，於此總見之。}

{\noindent\zhuan\zihao{6}\fzbyks 傳“象恭”至“曰洲”。正義曰：\CJKunderwave{堯典}言\CJKunderline{共工}之行云：“靜言庸違，象恭滔天。”言貌象恭敬，傲狠漫天，足以疑惑世人，故流放也。\CJKunderwave{左傳}說此事言“投諸四裔”。\CJKunderwave{釋地}雲“燕曰幽州”，知“北裔”也。“水中可居者曰洲”，\CJKunderwave{釋水}文。李巡曰：“四方有水，中央高，獨可居，故曰洲。”天地之勢，四邊有水，鄒衍書說“九州之外有瀛海環之”，是九州居水內，故以州為名，共在一洲之上，分之為九耳。州取水內為名，故引\CJKunderwave{爾雅}解“州”也。“投之四裔”，“裔”訓遠也，當在九州之外，而言“於幽州”者，在州境之北邊也。\CJKunderwave{禹貢}羽山在徐州，三危在雍州，故知北裔在幽州。下三者所居皆言山名，此\CJKunderline{共工}所處不近大山，故舉州言之。此流四凶在治水前，於時未作十有二州，則無幽州之名，而云“幽州”者，史據後定言之。 \par}

{\noindent\zhuan\zihao{6}\fzbyks 傳“黨於”至“南裔”。正義曰：\CJKunderline{共工}象恭滔天而\CJKunderline{驩兜}薦之，是“黨於\CJKunderline{共工}，罪惡同”，故放之也。\CJKunderwave{左傳}說此事云：“流四凶族，投諸四裔。”則四方方各有一人，幽州在北裔，雍州三危在西裔，徐州羽山在東裔，三方既明，知崇山在南裔也。\CJKunderwave{禹貢}無崇山，不知其處，蓋在衡嶺之南也。 \par}

{\noindent\zhuan\zihao{6}\fzbyks 傳“三苗”至“西裔”。正義曰：昭元年\CJKunderwave{左傳}說自古諸侯不用王命者，“虞有三苗，夏有觀扈”，知三苗是國，其國以三苗為名，非三國也。杜預言“三苗地闕，不知其處”。三兇皆是王臣，則三苗亦應是諸夏之國入仕王朝者也。文十八年\CJKunderwave{左傳}言：“縉雲氏有不才子，貪於飲食，冒於貨賄，侵欲崇侈，不可盈厭，聚斂積實，不知紀極，不分孤寡，不恤窮匱天下之民以比三兇，謂之饕餮。”即此三苗是也。知其然者，以\CJKunderwave{左傳}說此事言：“舜臣堯，流四凶族渾敦、窮奇、檮杌、饕餮，投諸四裔，以御螭魅。”謂此\CJKunderline{驩兜}、\CJKunderline{共工}、三苗與\CJKunderline{鯀}也。雖知彼言四凶，此等四人，但名不同，莫知孰是，惟當驗其行跡，以別其人。\CJKunderwave{左傳}說窮奇之行雲“靖譖庸回”，\CJKunderwave{堯典}言\CJKunderline{共工}之行雲“靜言庸違”，其事既同，知窮奇是\CJKunderline{共工}也。\CJKunderwave{左傳}說渾敦之行雲“醜類惡物，是與比周”，\CJKunderwave{堯典}言\CJKunderline{驩兜}薦舉\CJKunderline{共工}，與惡比周，知渾敦是\CJKunderline{驩兜}也。\CJKunderwave{左傳}說檮杌之行言“不可教訓，不知話言，傲狠明德，以亂天常”，\CJKunderwave{堯典}言\CJKunderline{鯀}之行雲“咈哉，方命圮族”，其事既同，知檮杌是\CJKunderline{鯀}也。惟三苗之行\CJKunderwave{堯典}無文，\CJKunderline{鄭玄}具引\CJKunderwave{左傳}之文乃云：“命\CJKunderline{驩兜}舉\CJKunderline{共工}，則\CJKunderline{驩兜}為渾敦也，\CJKunderline{共工}為窮奇也，\CJKunderline{鯀}為檮杌也，而三苗為饕餮亦可知。”是先儒以書傳相考，知三苗是饕餮也。\CJKunderwave{禹貢}雍州言“三危既宅，三苗丕敘”，知三危是西裔也。 \par}

{\noindent\zhuan\zihao{6}\fzbyks 傳“方命”至“海中”。正義曰：“方命圮族”，是其本性；“績用不成”，試而無功；二者俱是其罪，故並言之。\CJKunderwave{釋言}云：“殛，誅也。”傳稱流四凶族者，皆是流而謂之“殛竄放流，皆誅”者，流者移其居處，若水流然，罪之正名，故先言也。放者使之自活，竄者投棄之名，殛者誅責之稱，俱是流徙，異其文，述作之體也。四者之次，蓋以罪重者先。\CJKunderline{共工}滔天，為罪之最大。\CJKunderline{驩兜}與之同惡，故以次之。\CJKunderwave{祭法}以\CJKunderline{鯀}障洪水，故列諸祀典，功雖不就，為罪最輕，故後言之。\CJKunderwave{禹貢}徐州雲“蒙羽其藝”，是羽山為東裔也。\CJKunderwave{漢書·地理志}羽山在東海郡祝其縣西南，海水漸及，故言“在海中”也。 \par}

{\noindent\zhuan\zihao{6}\fzbyks 傳“皆服”至“見之”。正義曰：此四罪者徵用之初即流之也，舜以微賤超升上宰，初來之時，天下未服，既行四罪，故天下皆服舜用刑得當其罪也。自“象以典刑”以下，徵用而即行之，於此居攝之後，追論成功之狀。故作者先敘典刑，言舜重刑之事，而連引四罪，述其刑當之驗，明此諸事皆是徵用之時所行，於此總見之也。知此等諸事皆“徵用所行”者，\CJKunderwave{洪範}雲“\CJKunderline{鯀}則殛死，\CJKunderline{禹}乃嗣興”，僖三十三年\CJKunderwave{左傳}雲“舜之罪也殛\CJKunderline{鯀}，其舉也興\CJKunderline{禹}”，襄二十一年\CJKunderwave{左傳}雲“\CJKunderline{鯀}殛而\CJKunderline{禹}興”，此三者皆言殛\CJKunderline{鯀}而後用\CJKunderline{禹}，為治水是徵用時事，四罪在治水之前，明是“徵用所行”也。又下雲\CJKunderline{禹}讓稷、契、\CJKunderline{皋陶}，帝因追美三人之功，所言稷播百穀、契敷五教、\CJKunderline{皋陶}作士皆是徵用時事，\CJKunderline{皋陶}所行“五刑有服”、“五流有宅”即是“象以典刑”、“流宥五刑”，此為徵用時事足可明矣。而\CJKunderline{鄭玄}以為“\CJKunderline{禹}治水事畢，乃流四凶”。故王肅難鄭言：“若待\CJKunderline{禹}治水功成，而後以\CJKunderline{鯀}為無功殛之，是為舜用人子之功，而流放其父，則\CJKunderline{禹}之勤勞適足使父致殛，為舜失五典克從之義，\CJKunderline{禹}陷三千莫大之罪，進退無據，亦甚迂哉！” \par}

{\noindent\shu\zihao{5}\fzkt “肇十”至“咸服”。正義曰：史言舜既攝位，出行巡守，復分置州域，重慎刑罰。於\CJKunderline{禹}治水後，始分置十有二州，每州以一大山為鎮，殊大者十有二山。深其州內之川，使水通利。又留意於民，詳其罪罰，依法用其常刑，使罪各當，刑不越法。用流放之法寬宥五刑。五刑雖有犯者，或以恩減降，不使身服其罪，所以流放宥之。五刑之外，更有鞭作治官事之刑；有撲作師儒教訓之刑；其有意善功惡，則令出金贖罪之刑；若過誤為害，原情非故者，則緩縱而赦放之；若怙恃奸詐，終行不改者，則賊殺而刑罪之。舜慎刑如此，又設言以誡百官曰：“敬之哉！敬之哉！推此刑罰之事最須憂念之哉！”令勤念刑罰，不使枉濫也。又言舜非於攝位之後，方始重慎刑罪，初於登用之日即用刑當其罪，流徙\CJKunderline{共工}於北裔之幽州，放逐\CJKunderline{驩兜}於南裔之崇山，竄三苗於西裔之三危，誅殛伯\CJKunderline{鯀}於東裔之羽山。行此四罪，各得其實，而天下皆服從之。 \par}

\textcolor{red}{二十}有八載，帝乃殂落。\footnote{殂落，死也。堯年十六即位，七十載求禪，試舜三載，自正月上日至崩二十八載，堯死壽一百一十七歲。殂，才枯反。}百姓如喪考妣，\footnote{考妣，父母。言百官感德思慕。喪如字，又息浪反。妣，必履反，父曰考，母曰妣。}三載，四海遏密\textcolor{red}{八音}。\footnote{遏，絕。密，靜也。八音，金、石、絲、竹、匏、土、革、木。四夷絕音三年，則華夏可知。言盛德恩化所及者遠。遏，安葛反，或音謁。八音謂金，鍾也；石，磬也；絲，琴瑟也；竹，篪笛也；匏，笙也；土，壎也；革，鼓也；木，柷敔也。匏，白交反。}

{\noindent\zhuan\zihao{6}\fzbyks 傳“殂落”至“七歲”。正義曰：“殂落，死也”，\CJKunderwave{釋詁}文。李巡曰：“殂落，堯死之稱。”郭璞曰：“古死尊卑同稱。故\CJKunderwave{書}堯曰‘殂落’，舜曰‘陟方乃死’。”謂之“殂落”者，蓋“殂”為往也，言人命盡而往；“落”者若草木葉落也。堯以十六即位，明年乃為元年。七十載求禪，求禪之時八十六也。試舜三年，自正月上日至崩二十八載，總計其數，凡壽一百一十七歲。案\CJKunderwave{堯典}求禪之年即得舜而試之，求禪試舜共在一年也。更得二年，即為歷試三年，故下傳雲“歷試二年”。與攝位二十八年，合得為“三十在位”。故王肅云：“徵用三載，其一在徵用之年，其餘二載，與攝位二十八年凡三十歲也。”故孔傳云：“歷試二年。”明其一年在徵用之限。以此計之，淮有一百一十六歲，不得有七，蓋誤為七也。 \par}

{\noindent\zhuan\zihao{6}\fzbyks 傳“考妣”至“思慕”。正義曰：\CJKunderwave{曲禮}云：“生曰父母，死曰考妣。”\CJKunderline{鄭玄}云：“考，成也，言其德行之成也。妣之言媲也，媲於考也。”\CJKunderwave{喪服}為父為君，同服斬衰。\CJKunderwave{檀弓}說事君之禮云：“服勤至死，方喪三年。”\CJKunderline{鄭玄}云：“方喪資於事父，凡此以義為制。”義重則恩輕，其情異於父。“如喪考妣”言百官感德，情同父母，思慕深也。諸經傳言“百姓”或為百官，或為萬民，知此“百姓”是百官者，以\CJKunderwave{喪服}庶民為天子齊衰三月，畿外之民無服，不得如考妣，故知百官也。 \par}

{\noindent\zhuan\zihao{6}\fzbyks 傳“遏絕”至“者遠”。正義曰：“密，靜”，\CJKunderwave{釋詁}文。“遏”，止絕之義，故為絕也。\CJKunderwave{周禮·太師}云：“播之以八音：金、石、土、革、絲、木、匏、竹。”鄭云：“金，鍾鎛也；石，磬也；土，壎也；革，鼓鼗也；絲，琴瑟也；木，柷敔也；匏，笙也；竹，管簫也。”傳言“八音”與彼次不同者，隨便言耳。\CJKunderwave{釋地}云：“九夷八狄七戎六蠻謂之四海。”夷狄尚絕音三年，則華夏內國可知也。\CJKunderwave{喪服}諸侯之大夫為天子正服穗衰，既葬除之。今能使四夷三載絕音，言堯有盛德，恩化所及遠也。 \par}

{\noindent\shu\zihao{5}\fzkt “二十”至“八音”。正義曰：舜受終之後，攝天子之事二十有八載，\CJKunderline{帝堯}乃死。百官感德思慕，如喪考妣。三載之內，四海之人，蠻夷戎狄皆絕靜八音而不復作樂。是堯盛德恩化所及者遠也。 \par}

\textcolor{red}{月正}元日，\CJKunderline{舜}格於文祖，\footnote{月正，正月。元日，上日也。舜服堯喪三年畢，將即政，故復至文祖廟告。復，扶又反。}詢於四嶽，闢四門，\footnote{詢，謀也。謀政治於四嶽，開闢四方之門未開者，廣致眾賢。闢，婢亦反，徐甫亦反。}明四目，達四聰。\footnote{廣視聽於四方，使天下無壅塞。}

{\noindent\zhuan\zihao{6}\fzbyks 傳“月正”至“廟告”。正義曰：“正”訓長也，“月正”言月之最長，正月長於諸月，“月正”還是正月也。上日，日之最上；元日，日之最長，“元日”還是上日。王肅云：“月正元日，猶言正月上日。”變文耳。\CJKunderwave{禮}雲“令月吉日”，又變文言“吉月令辰”，此之類也。知“舜服堯喪三年畢，將即政”者，以堯存且攝其位，堯崩謙而不居。\CJKunderwave{孟子}云：“堯崩，三年喪畢，舜避\CJKunderline{丹朱}於南河之南。天下諸侯朝覲者不之堯子而之舜，獄訟者不之堯子而之舜，謳歌者不之堯子而謳歌舜。曰‘天也。’然後之中國踐天子位。”既言然矣，此文又承“三載”之下，故知舜服堯喪三年畢，將欲即政。“復至文祖廟告”，前以攝位告，今以即政告也。此猶是堯之文祖，自此以後舜當自立文祖之廟，堯之文祖當遷於\CJKunderline{丹朱}之國也。 \par}

{\noindent\zhuan\zihao{6}\fzbyks 傳“詢謀”至“眾賢”。正義曰：“詢，謀”，\CJKunderwave{釋詁}文。“闢”訓開，開四方之門，謂開仕路，引賢人也。\CJKunderwave{論語}云：“從我於陳蔡者，皆不及門也。”門者行之所由，故以門言仕路。以堯舜之聖，求賢久矣，今更言開門，是開“其未開”者，謂多設取士之科，以此廣致眾賢也。 \par}

{\noindent\zhuan\zihao{6}\fzbyks 傳“廣視”至“壅塞”。正義曰：“聰”謂耳聞之也。既雲“明四目”不雲“聰四耳”者，目視苦其不明，耳聰貴其及遠，“明”謂所見博，“達”謂聽至遠，二者互以相見。故傳總申其意“廣視聽於四方，使天下無壅塞”。天子之聞見在下，必由近臣四嶽親近之官，故與謀此事也。 \par}

諮十有二牧，曰：“食哉，惟時！\footnote{諮亦謀也。所重在於民食，惟當敬授民時。}柔遠能邇，惇德允元，\footnote{柔，安。邇，近。敦，厚也。元，善之長。言當安遠，乃能安近。厚行德信，使足長善。惇音敦。長，張丈反，下同。}而難任人，蠻夷\textcolor{red}{率服}。”\footnote{任，佞。難，拒也。佞人斥遠之，則忠信昭於四夷，皆相率而來服。難音乃旦反。任音壬，又音而鴆反。}

{\noindent\zhuan\zihao{6}\fzbyks 傳“諮亦”至“民時”。正義曰：“諮，謀”，\CJKunderwave{釋詁}文。以上“帝曰，諮”，上連“帝曰”，故為諮嗟，此則上有“詢於四嶽”，言“諮十有二牧”，故為謀也。立君所以牧民，民生在於粒食，是君之所重。\CJKunderwave{論語}雲“所重民食”，謂年穀也。種殖收斂，及時乃獲，故“惟當敬授民時”。 \par}

{\noindent\zhuan\zihao{6}\fzbyks 傳“柔安”至“長善”。正義曰：“柔，安”、“邇，近”、“惇，厚”，皆\CJKunderwave{釋詁}文。“元，善之長”，\CJKunderwave{易}文言也。安近不能安遠，遠人或來擾亂，雖欲安近，近亦不安。人君為政，若其不能安近，但戒使之柔遠，故能安近。言當安彼遠人，乃能安近。欲令遠言皆安也。王肅云：“能安遠者先能安近。”知不然者，以牧在遠方，故據遠言之。“惇德”者，令人君厚行德也。“允元”者，信使足為長善也。言人君厚行德之與信，使足為善長，民必效之為善而行也。 \par}

{\noindent\zhuan\zihao{6}\fzbyks 傳“任佞”至“來服”。正義曰：“任，佞”，\CJKunderwave{釋詁}文。孫炎云：“似可任之佞也。”\CJKunderwave{論語}說為邦之法雲“遠佞人”，“佞人殆”，故以難距佞人為“斥遠之”，令不幹朝政。朝無佞人，則“忠信昭於四夷，皆相率而來服”也。舉蠻夷而戎狄亦見矣。 \par}

{\noindent\shu\zihao{5}\fzkt “月正”至“率服”。正義曰：自此已下言舜真為天子，命百官授職之事。舜既除堯喪，以明年之月正元日，舜至於文祖之廟，告己將即正位為天子也。告廟既訖，乃謀政治於四嶽之官。所謀開四方之門，大為仕路致眾賢也。明四方之目，使為己遠視四方也。達四方之聰，使為己遠聽聞四方也。恐遠方有所擁塞，令為己悉聞見之。既謀於四嶽，又別敕州牧，諮十有二牧曰：“人君最所重者，在於民之食哉！惟當敬授民之天時，無失其農要。為政務在安民，當安彼遠人，則能安近人耳。遠人不安，則近亦不安。”欲令遠近皆安之也。“又當厚行德信，而使足為善長”。欲令諸侯皆厚行其德，為民之師長。“而難拒佞人，斥遠之，使不幹朝政，如是則誠信昭於四夷，自然蠻夷皆相率而來服也”。 \par}

\textcolor{red}{\CJKunderline{舜}曰}：“諮！四嶽，有能奮庸熙帝之載，\footnote{奮，起。庸，功。載，事也。訪群臣有能起發其功，廣堯之事者。言“舜曰”以別堯。奮，弗運反。}使宅百揆，亮採惠疇？”\footnote{亮，信。惠，順也。求其人使居百揆之官，信立其功，順其事者誰乎？}僉曰：“\CJKunderline{伯禹}作司空。”\footnote{四嶽同辭而對，\CJKunderline{禹}代\CJKunderline{鯀}為宗伯，入為天子司空。治洪水有成功，言可用之。}

{\noindent\zhuan\zihao{6}\fzbyks 傳“奮起”至“別堯”。正義曰：“奮”是起動之意，故為起也。\CJKunderwave{釋詁}云：“庸，勞也。”勞亦功也。\CJKunderline{鄭玄}云：“載，行也。”王肅云：“載，成也。”孔以“載”為事也，各自以意訓耳。舜受堯禪，當繼行其道。行之在於任臣，百揆臣之最貴，求能起發其功，廣大\CJKunderline{帝堯}之事者，欲任之。舜既即位，可以稱帝，而言“舜曰”者，承堯事下，言“舜曰”以別堯，於此一別，以下稱帝也。 \par}

{\noindent\zhuan\zihao{6}\fzbyks 傳“亮信”至“誰乎”。正義曰：“亮，信”，\CJKunderwave{釋詁}文。“惠，順”，\CJKunderwave{釋言}文。上雲舜“納於百揆”，“百揆”是官名，故求其人，使居百揆之官。居官則當信立其功，能順其事者誰乎？此官任重，當統群職繼堯之功，故歷言所順而後始問誰乎，異於餘官先言“疇”也。 \par}

{\noindent\zhuan\zihao{6}\fzbyks 傳“四嶽”至“用之”。正義曰：“僉”訓為皆，故云“四嶽皆同辭而對”也。\CJKunderwave{國語}云：“有崇伯\CJKunderline{鯀}，堯殛之於羽山。”賈逵云：“崇，國名。伯，爵也。”\CJKunderline{禹}代\CJKunderline{鯀}為崇伯，入為天子司空，以其伯爵，故稱“\CJKunderline{伯禹}”。言人之賢而舉其為官，知\CJKunderline{禹}“治洪水自成功，言可用”也。 \par}

帝曰：“俞，諮！\CJKunderline{禹}，汝平水土，惟時懋哉！”\footnote{然其所舉，稱\CJKunderline{禹}前功以命之。懋，勉也。惟居是百揆，勉行之。俞，以朱反。懋音茂，王云：“勉也。”馬云：“美也。”}\CJKunderline{禹}拜稽首，讓於\CJKunderline{稷}、\CJKunderline{契}暨\CJKunderline{皋陶}。\footnote{居稷官者棄也。契、\CJKunderline{皋陶}，二臣名。稽首，首至地。稽音啟。稽首，首至地，臣事君之禮。契，息列反。陶音遙。}帝曰：“俞，汝\textcolor{red}{往哉}！”\footnote{然其所推之賢，不許其讓，敕使往宅百揆。}

{\noindent\zhuan\zihao{6}\fzbyks 傳“然其”至“行之”。正義曰：\CJKunderline{禹}平水土，往前之事，嫌其今覆命之令平水土，故云“稱\CJKunderline{禹}前功以命之”。“懋，勉”，\CJKunderwave{釋詁}文。 \par}

{\noindent\zhuan\zihao{6}\fzbyks 傳“居稷”至“首至地”。正義曰：下文帝述三人，遂變“稷”為“棄”，故解之“居稷官者棄也”。獨稱“官”者，出自\CJKunderline{禹}意耳，不必著義。鄭云：“時天下賴后稷之功，故以官名通稱。”或當然也。經因“稷、契”名單，共文言“暨\CJKunderline{皋陶}”，為文勢耳。三人為此次者，蓋以官尊卑為先後也。\CJKunderwave{周禮·太祝}：“辨九拜，一曰稽首。”稽首為敬之極，故為“首至地”。稽首是拜內之別名，\CJKunderline{禹}拜乃稽首，故云“拜稽首”也。 \par}

{\noindent\shu\zihao{5}\fzkt “舜曰”至“往哉”。正義曰：舜本以百揆攝位，今既即政，故求置其官。曰：“諮嗟！四嶽等，汝於群臣之內，有能起發其功，廣大\CJKunderline{帝堯}之事者，我欲使之居百揆之官。在官而信立其功，於事能順者，其是唯乎？”四嶽皆曰：“\CJKunderline{伯禹}作司空，有成功，惟此人可用。”帝曰：“然。”然其所舉得人也。乃諮嗟敕\CJKunderline{禹}：“汝本平水土，實有成功，惟當居是百揆而勉力行哉！”\CJKunderline{禹}拜稽首，讓於稷、契與\CJKunderline{皋陶}。帝曰：“然。”然其所讓實賢也。“汝但往居此職”。不許其讓也。 \par}

\textcolor{red}{帝曰}：“\CJKunderline{棄}，黎民阻飢，汝后稷，播時\textcolor{red}{百穀}。”\footnote{阻，難。播，布也。眾人之難在於飢，汝后稷，布種是百穀以濟之。美其前功以勉之。阻，莊呂反，王云：“難也。”播，波左反。}

{\noindent\zhuan\zihao{6}\fzbyks 傳“阻難”至“勉之”。正義曰：“阻，難”，\CJKunderwave{釋詁}文。“播”是分散之義，故為布也。王肅云：“播，敷也。”堯遭洪水，民不粒食，故眾民之難在於飢也。“稷”是五穀之長，立官主此稷事。“後”訓君也。帝言：“汝君此稷官，布種是百穀以濟救之。”追美其功以勸勉之。上文“讓於稷、契”，\CJKunderwave{益稷}雲“暨稷”\CJKunderwave{呂刑}雲“稷降播種”，\CJKunderwave{國語}雲“稷為天官”。單名為“稷”，尊而君之稱為“后稷”，故\CJKunderwave{詩傳}、\CJKunderwave{孝經}皆以“后稷”為言，非官稱“後”也。 \par}

{\noindent\shu\zihao{5}\fzkt “帝曰棄”至“百穀”。正義曰：帝因\CJKunderline{禹}讓三人而官不轉，各述其功以勸之。帝呼稷曰：“棄，往者洪水之時，眾民之難難在於飢，汝君為此稷之官，教民布種是百穀以濟活之。”言我佑汝功，當勉之。 \par}

\textcolor{red}{帝曰}：“\CJKunderline{契}，百姓不親，五品不遜，\footnote{五品謂五常。遜，順也。}汝作司徒，敬敷五教，\textcolor{red}{在寬}。”\footnote{布五常之教務在寬，所以得人心，亦美其前功。}

{\noindent\zhuan\zihao{6}\fzbyks 傳“五品”至“順也”。正義曰：“品”謂品秩，一家之內尊卑之差，即父母兄弟子是也。教之義慈友恭孝，此事可常行，乃為“五常”耳。傳上雲“五典克從”，即此五品能順。上傳以解\CJKunderwave{五典}為五常，又解此以同之，故云“五品謂五常”。其實五常據教為言，不據品也。“遜，順”，常訓也。不順謂不義、不慈、不友、不恭、不孝也。 \par}

{\noindent\zhuan\zihao{6}\fzbyks 傳“布五”至“前功”。正義曰：文十八年\CJKunderwave{左傳}雲“布五教於四方，父義、母慈、兄友、弟恭、子孝”，是“布五常之教”也。\CJKunderwave{論語}雲“寬則得眾”，故“務在寬，所以得民心”也。治不遜之罪，宜峻法以繩之，而貴其務在寬者，此五品不遜，直是禮教不行，風俗未淳耳，未有殺害之罪，故教之務在於寬。若其不孝不恭，其人至於逆亂而後治之，於事不得寬也。 \par}

{\noindent\shu\zihao{5}\fzkt “帝曰契”至“在寬”。正義曰：帝又呼契曰：“往者天下百姓不相親睦，家內尊卑五品不能和順。汝作司徒之官，謹敬布其五常之教，務在於寬，故使五典克從。是汝之功，宜當勉之。” \par}

\textcolor{red}{帝曰}：“\CJKunderline{皋陶}，蠻夷猾夏，寇賊奸宄，\footnote{猾，亂也。夏，華夏。群行攻劫曰寇，殺人曰賊。在外曰奸，在內曰宄。言無教之致。猾，戶八反。寇，苦豆反。宄音軌。}汝作士。五刑有服，\footnote{士，理官也。五刑，墨、劓、剕、宮、大辟。服，從也。言得輕重之中正。劓，魚器反，截鼻也。剕，扶味反，刖足也。大辟，婢亦反，死刑也。}五服三就。\footnote{既從五刑，謂服罪也。行刑當就三處，大罪於原野，大夫於朝，士於市。處，昌慮反。朝，直遙反。}五流有宅，五宅三居。\footnote{謂不忍加刑，則流放之，若四凶者。五刑之流，各有所居。五居之差，有三等之居，大罪四裔，次九州之外，次千里之外。}惟明\textcolor{red}{克允}。”\footnote{言\CJKunderline{皋陶}能明信五刑，施之遠近，蠻夷猾夏，使咸信服，無敢犯者。因\CJKunderline{禹}讓三臣，故歷述之。}

{\noindent\zhuan\zihao{6}\fzbyks 傳“猾亂”至“之致”。正義曰：“猾”者狡猾相亂，故“猾”為亂也。“夏”訓大也，中國有文章光華，禮義之大。定十年\CJKunderwave{左傳}雲“裔不謀夏，夷不亂華”，是中國為華夏也。“寇”者眾聚為之，賊者殺害之稱，故“群行攻劫曰寇，殺人曰賊”。成十七年\CJKunderwave{左傳}云：“亂在外為奸，在內為宄。”是“在外曰奸，在內曰宄”也。“寇賊奸宄”皆是作亂害物之名也。“蠻夷猾夏”，興兵犯邊，害大，故先言之。“寇賊奸宄”，皆國內之害，小，故後言之。\CJKunderwave{管子}曰：“倉廩實知禮節，衣食足知榮辱，讓生於有餘，爭生於不足。”往者洪水為災，下民飢困，內有寇賊為害，外則四夷犯邊，皆言無教之致也。唐堯之聖，協和萬邦，不應末年，頓至於此，蓋少有其事，辭頗增甚，歸功於人，作與奪之勢耳。 \par}

{\noindent\zhuan\zihao{6}\fzbyks 傳“士理”至“中正”。正義曰：“士”即\CJKunderwave{周禮}司寇之屬，有士師、卿士等，皆以“士”為官名。\CJKunderline{鄭玄}云：“士，察也，主察獄訟之事。”\CJKunderwave{月令}云：“命大理。”昭十四年\CJKunderwave{左傳}云：“叔魚攝理。”是謂獄官為理官也。準\CJKunderwave{呂刑}文，知五刑謂墨、劓、剕、宮、大辟也。人心服罪是順從之義，故為從也。所以服者，言得輕重之中正也。\CJKunderwave{呂刑}雲“咸庶中正”是也。 \par}

{\noindent\zhuan\zihao{6}\fzbyks 傳“既從”至“於市”。正義曰：經言“五服”，謂\CJKunderline{皋陶}所斷五刑皆服其罪，傳既訓“服”為從，故云“既從五刑謂服罪也”。“行刑當就三處”，惟謂大辟罪耳。\CJKunderwave{魯語}云：“刑五而已，無有隱者。大刑用甲兵，次刑斧鉞，中刑刀鋸，其次鑽筰，薄刑鞭撲，以威民。故大者陳之原野，小者致之市朝，五刑三次，是無隱也。”孔用彼為說，故以“三就”為原野與朝、市也。\CJKunderwave{國語}賈逵注云：“用兵甲者，諸侯逆命，征討之刑也。大夫已上於朝，士已下於市。”傳雖不言“已上、已下”，為義亦當然也。\CJKunderwave{國語}雲五刑者，謂甲兵也，斧鉞也，刀鋸也，鑽筰也，鞭撲也，與\CJKunderwave{呂刑}之五刑異也。所言“三次”即此“三就”是也。惟死罪當分就處所，其墨、劓、剕、官無常處可就也。馬、鄭、王三家皆以“三就為原野也、市朝也、甸師氏也”。案刑于甸師氏者，王之同族，刑于隱者，不與國人，慮兄弟耳，非所刑之正處。此言正刑，不當數甸師也。又市、朝異所，不得合以為一，且皆\CJKunderwave{國語}之文，其義不可通也。 \par}

{\noindent\zhuan\zihao{6}\fzbyks 傳“謂不”至“之外”。正義曰：此“五流有宅”即“流宥五刑”也。當在五刑而流放之，故知謂“不忍加刑，則流放之，若四凶”也。\CJKunderline{鄭玄}云：“舜不刑此四人者，以為堯臣，不忍刑之。”王肅云：“謂在八議之闢，君不忍殺，宥之以遠。”八議者，\CJKunderwave{周禮·小司寇}所云議親、議故、議賢、議能、議貴、議賓、議勤是也。以君恩不忍殺，罪重不可全赦，故流之也。“五刑之流，各有所居”，謂徙置有處也。“五居之差，有三等之居”，量其罪狀為遠近之差也。四裔最遠，在四海之表，故“大罪四裔”，謂不犯死罪也。故\CJKunderwave{周禮·調人職}雲“父之仇，闢諸海外”，即與“四裔”為一也。“次九州之外”，即\CJKunderwave{王制}雲入學不率教者，“屏之遠方，西方曰僰，東方曰寄”，注云“逼寄於夷狄也”，與此“九州之外”同也。“次千里之外”者，即\CJKunderwave{調人職}雲“兄弟之仇，闢諸千里之外”也。\CJKunderwave{立政}雲“中國之外”，不同者，言“中國”者，據罪人所居之國定千里也。據其遠近，其實一也。\CJKunderwave{周禮}與\CJKunderwave{王制}既有三處之別，故約以為言。\CJKunderline{鄭玄}云：“三處者，自九州之外至於四海，三分其地，遠近若周之夷、鎮、蕃也。”然罪有輕重不同，豈五百里之校乎？不可從也。 \par}

{\noindent\zhuan\zihao{6}\fzbyks 傳“言皋”至“述之”。正義曰：“惟明”謂\CJKunderline{皋陶}之明，“克允”謂受罪者信服。故王肅云：“惟明其罪，能使之信服。”是信施於彼也。但彼人信服，由\CJKunderline{皋陶}有信，故傳言：“\CJKunderline{皋陶}能明信五刑，施之遠近蠻夷，使咸信服。”主言信者，見其\CJKunderline{皋陶}有信，故彼信之也。 \par}

{\noindent\shu\zihao{5}\fzkt “帝曰\CJKunderline{皋陶}”至“克允”。正義曰：帝呼\CJKunderline{皋陶}曰：“往者蠻夷戎狄猾亂華夏，又有強寇劫賊外奸內宄者，為害甚大。汝作士官治之，皆能審得其情，致之五刑之罪，受罪者皆有服從之心。”言輕重得中，悉無怨恨也。“五刑有服從者，於三處就而殺之。其有不忍刑其身者，則斷為五刑而流放之。五刑之流，各有所居處。五刑所居，於三處居之。所以輕重罪得其宜，受罪無怨者，惟汝識見之明，能使之信服，故奸邪之人無敢更犯。是汝之功，宜當勉之”。因\CJKunderline{禹}之讓，以次誡之。 \par}

帝曰：“疇若予工？”僉曰：“\CJKunderline{垂}哉！”\footnote{問：“誰能順我百工事者？”朝臣舉垂。垂，臣名。垂如字，徐音睡。}帝曰：“俞，諮！\CJKunderline{垂}，汝共工。”\footnote{共謂供其職事。共音恭。}

{\noindent\zhuan\zihao{6}\fzbyks 傳“問誰”至“臣名”。正義曰：\CJKunderwave{考工記}云：“國有六職，百工與居一焉。”“工”即百工，故云“問誰能順我百工事者”。直言“帝曰”，無所偏諮，故知“僉曰”是朝臣共舉垂也。 \par}

{\noindent\zhuan\zihao{6}\fzbyks 傳“共謂供其職事”。正義曰：\CJKunderwave{堯典}傳云：“共工，官稱。”即彼以“共工”二字為官名。上雲“疇若予工”，單舉“工”名，今命此人云“汝作共工”，明是帝謂此人堪供此職，非是呼此官名為“共工”也。其官或以“共工”為名，要帝意言“共”謂供此職也。 \par}

\CJKunderline{垂}拜稽首，讓於\CJKunderline{殳斨}暨\CJKunderline{伯與}\footnote{\CJKunderline{殳斨}、\CJKunderline{伯與}二臣名。\xpinyin*{斨}，七良反。與音餘。}帝曰：“俞，往哉！汝諧。”\footnote{汝能諧和此官。}帝曰：“疇若予上下草木鳥獸？”僉曰：“\CJKunderline{益}哉！”\footnote{上謂山，下謂澤，順謂施其政教，取之有時，用之有節。言伯益能之。益，\CJKunderline{皋陶}子也。}帝曰：“俞，諮！\CJKunderline{益}，汝作朕虞。”\footnote{虞，掌山澤之官。}

{\noindent\zhuan\zihao{6}\fzbyks 傳“上謂”至“能之”。正義曰：言“上下草木鳥獸”，則上之與下各有草木鳥獸，即\CJKunderwave{周禮}山虞、澤虞之官各掌其教，知“上謂山，下謂澤”也。順其草木鳥獸之宜，明是“施其政教，取之有時，用之有節”也。馬、鄭、王本皆為“\CJKunderline{禹}曰：‘益哉！’”是字相近而彼誤耳。 \par}

{\noindent\zhuan\zihao{6}\fzbyks “作朕虞”。正義曰：此官以“虞”為名，帝言作我虞耳，“朕”非官名也。\CJKunderline{鄭玄}云：“言朕虞，重鳥獸草木。”\CJKunderwave{漢書}王莽自稱為予，立予虞之官。則莽謂此官名為“朕虞”，其義必不然也。 \par}

\CJKunderline{益}拜稽首，讓於\CJKunderline{朱虎}、\CJKunderline{熊羆}。帝曰：“俞，往哉！汝諧。”\footnote{\CJKunderline{朱虎}、\CJKunderline{熊羆}，二臣名。垂、益所讓四人皆在元凱之中。羆，彼皮反。}帝曰：“諮！四嶽，有能典朕三禮？”僉曰：“\CJKunderline{伯夷}。”\footnote{三禮，天地人之禮。\CJKunderline{伯夷}，臣名，姜姓。}

{\noindent\zhuan\zihao{6}\fzbyks 傳“\CJKunderline{朱虎}”至“之中”。正義曰：知“垂、益所讓四人皆在元凱之中”者，以文十八年\CJKunderwave{左傳}八元之內有伯虎、仲熊，即此“\CJKunderline{朱虎}、\CJKunderline{熊羆}”是也。虎、熊在元凱之內，明\CJKunderline{殳斨}、\CJKunderline{伯與}亦在其內，但不知彼誰當之耳。益是\CJKunderline{皋陶}之子，\CJKunderline{皋陶}即庭堅也。益在八凱之內，垂則不可知也。傳不在\CJKunderline{伯夷}、夔龍之下為此言者，以\CJKunderline{伯夷}姜姓，不在元凱之內，夔龍亦不可知，惟言此四人耳。傳雖言\CJKunderline{殳斨}、\CJKunderline{伯與}，亦難知也。 \par}

{\noindent\zhuan\zihao{6}\fzbyks 傳“三禮”至“姜姓”。正義曰：此時“秩宗”，即\CJKunderwave{周禮}之宗伯也，其職雲“掌天神、人鬼、地祇之禮”，雖三者併為吉禮，要言三禮者是天地人之事，故知三禮是“天地人之禮”。上文舜之巡守言“修五禮”，此雲“典朕三禮”，各有其事，則五禮皆據其所施於三處。五禮所施於天地人耳。言“三”足以包五，故舉“三”以言之。\CJKunderwave{鄭語}云：“姜，\CJKunderline{伯夷}之後也。\CJKunderline{伯夷}能禮於神以佐堯。”是\CJKunderline{伯夷}為姜姓也。此經不言“疇”者，訪其有能，是問誰可知，上文已具，此略之也。 \par}

帝曰：“俞，諮！\CJKunderline{伯}，汝作秩宗。\footnote{秩序宗尊也，主郊廟之官。}夙夜惟寅，直哉惟清。”\footnote{夙，早也。言早夜敬思其職，典禮施政教，使正直而清明”寅如字，徐音夷。}

{\noindent\zhuan\zihao{6}\fzbyks 傳“秩序”至“之官”。正義曰：\CJKunderwave{堯典}傳已訓“秩”為序，此複訓者，此為官名，須辨官名之義，故詳之也。“宗”之為尊，常訓也。主郊廟之官，掌序鬼神尊卑，故以“秩宗”為名。“郊”謂祭天南郊，祭地北郊；“廟”謂祭先祖，即\CJKunderwave{周禮}所謂“天神、人鬼、地祇之禮”是也。 \par}

{\noindent\zhuan\zihao{6}\fzbyks 傳“夙，早也。言早夜敬思其職，典禮施政教，使正直而清明”。正義曰：“夙，早”，\CJKunderwave{釋詁}文。“早夜敬服其職”，謂侵早已起，深夜乃臥，謹敬其職事也。典禮之官施行教化，使正直而清明。正直，不枉曲也。清明，不闇昧也。 \par}

伯拜稽首，讓於\CJKunderline{夔}、\CJKunderline{龍}。\footnote{夔、龍，二臣名。夔音求龜反。}帝曰：“俞，往，欽哉！”\footnote{然其賢，不許讓。}\textcolor{red}{帝曰}：“\CJKunderline{夔}，命汝典樂，教胄子，\footnote{胄，長也，謂元子以下至卿大夫子弟。以歌詩蹈之舞之，教長國子中、和、祇、庸、孝、友。胄，直又反，王云：“胄子，國子也。”馬云：“胄，長也，教長天下之子弟。”}直而溫，寬而慄，\footnote{教之正直而溫和，寬弘而能莊慄。莊慄，戰慄也。}剛而無虐，簡而無傲。\footnote{剛失入虐，簡失入傲，教之以防其失。}

{\noindent\zhuan\zihao{6}\fzbyks 傳“胄長”至“孝友”。正義曰：\CJKunderwave{說文}云：“胄，胤也。”\CJKunderwave{釋詁}云：“胤，繼也。”繼父世者惟長子耳，故以“胄”為長也。“謂元子已下至卿大夫子弟”者，\CJKunderwave{王制}云：“樂正崇四術，立四教。王太子、王子、群后之太子、卿大夫元士之適子皆造焉。”是“下至卿大夫”也。不言“元士”，士卑，故略之。彼鄭注云：“王子，王之庶子也。”此傳兼言“弟”者，蓋指太子之弟耳。或孔意公卿大夫之弟亦教之，國子以適為主，故言“胄子”也。命典樂之官，使教胄子。下句又言詩歌之事，是令夔以歌詩蹈之舞之，教此適長國子也。\CJKunderwave{周禮·大司樂}云：“以樂德教國子中、和、祇、庸、孝、友。”鄭云：“中猶忠也。和，剛柔適也。祇，敬也。庸，有常也。善父母曰孝。善兄弟曰友。”是言樂官用樂教之，使成此六德也。\CJKunderwave{樂記}又云：“樂在宗廟之中，君臣上下同聽之，則莫不和敬。在族黨鄉里之中，長幼同聽之，則莫不和順。在閨門之內，父子兄弟同聽之，則莫不和親。”是樂之感人，能成忠、和、祇、庸、孝、友之六德也。 \par}

{\noindent\zhuan\zihao{6}\fzbyks 傳“教之”至“莊慄”。正義曰：此“直而溫”與下三句皆使夔教胄子，令性行當然，故傳發首言“教之”也。正直者失於太嚴，故令“正直而溫和”；寬弘者失於緩慢，故令“寬弘而莊慄”；謂矜莊嚴慄。慄者，謹敬也。 \par}

{\noindent\zhuan\zihao{6}\fzbyks 傳“剛失”至“其失”。正義曰：剛彊之失入於苛虐，故令人剛而無虐。簡易之失入於傲慢，故令簡而無傲。剛、簡是其本性，教之使無虐、傲，是言教之以防其失也。由此而言之，上二句亦直、寬是其本性，直失於不溫，寬失於不慄，故教之使溫、慄也。直、寬、剛、簡即\CJKunderline{皋陶}所謀之九德也。九德而獨舉此四事者，人之大體，故特言之。 \par}

詩言志，歌永言，\footnote{謂詩言志以導之，歌詠其義以長其言。永徐音詠，又如字。}聲依永，律和聲。\footnote{聲謂五聲：宮、商、角、徵、羽。律謂六律、六呂，十二月之音氣。言當依聲律以和樂。}八音克諧，無相奪倫，神人以和。”\footnote{倫，理也。八音能諧，理不錯奪，則神人咸和。命夔使勉之。}

{\noindent\zhuan\zihao{6}\fzbyks 傳“謂詩”至“其言”。正義曰：作詩者自言己志，則詩是言志之書，習之可以生長志意，故教其詩言志以導胄子之志，使開悟也。作詩者直言不足以申意，故長歌之，教令歌詠其詩之義以長其言，謂聲長續之。定本經作“永”字，明訓“永”為長也。 \par}

{\noindent\zhuan\zihao{6}\fzbyks 傳“聲謂”至“和樂”。正義曰：\CJKunderwave{周禮·太師}云：“文之以五聲：宮、商、角、徵、羽。”言五聲之清濁有五品，分之為五聲也。又“太師掌六律、六呂以合陰陽之聲。陽聲黃鐘、太簇、姑洗、蕤賓、夷則、無射。陰聲大呂、應鐘、南呂、林鐘、仲呂、夾鍾”。是六律、六呂之名也。\CJKunderwave{漢書·律曆志}云：“律有十二，陽六為律，陰六為呂。”是陰律名同，亦名呂也。\CJKunderline{鄭玄}云：“律述氣也，同助陰宣氣，與之同也。”又云：“呂，旅也，言旅助陽宣氣也。”\CJKunderwave{志}又云：“律黃帝之所作也，黃帝使伶倫氏自大夏之西、崑崙之陰，取竹於嶰谷之中各生、其竅厚薄均者，斷兩節之間吹之，以為黃鐘之宮。制十二籥，以聽鳳皇之鳴，其雄聲為六，雌鳴亦六，以比黃鐘之宮，是為律之本。”言律之所作如此。聖人之作律也，既以出音，又以候氣，布十二律於十二月之位，氣至則律應，是六律、六呂述十二月之音氣也。“聲依永”者，謂五聲依附長言而為之，其聲未和，乃用此律呂調和其五聲，使應於節奏也。 \par}

{\noindent\zhuan\zihao{6}\fzbyks 傳“倫理”至“勉之”。正義曰：“倫”之為理，常訓也。八音能諧，相應和也。各自守分，不相奪道理，是言理不錯亂相奪也。如此則神人咸和矣。帝言此者，命夔使勉之也。\CJKunderwave{大司樂}云：“大合樂以致鬼神示，以和邦國，以諧萬民，以安賓客，以說遠人。”是神人和也。 \par}

\CJKunderline{夔}曰：“於！予擊石拊石，百獸\textcolor{red}{率舞}。”\footnote{石，磬也。磬，音之清者。拊亦擊也。舉清者和則其餘皆從矣。樂感百獸，使相率而舞，則神人和可知。於如字，或音烏而絕句者，非。拊音撫，徐音府。}

{\noindent\zhuan\zihao{6}\fzbyks 傳“石磬”至“可知”。正義曰：樂器惟磬以石為之，故云：“石，磬也。”八音之音石磬最清，故知磬是音之聲清者。磬必擊以鳴之，故云拊亦擊之。重其文者，擊有大小，“擊”是大擊，“拊”是小擊。音聲濁者粗，清者精，精則難和，舉清者和，則其餘皆從矣。\CJKunderwave{商頌}云：“依我磬聲。”是言磬聲清，諸音來依之。“百獸率舞”即\CJKunderwave{大司樂}雲“以作動物”、\CJKunderwave{益稷}雲“鳥獸蹌蹌”是也。人神易感，鳥獸難感，百獸相率而舞，則神人和可知也。夔言此者，以帝戒之雲“神人以和”，欲使勉力感神人也。乃答帝雲“百獸率舞”，則神人以和，言帝德及鳥獸也。 \par}

{\noindent\shu\zihao{5}\fzkt “帝曰夔”至“率舞”。正義曰：帝因\CJKunderline{伯夷}所讓，隨才而任用之。帝呼\CJKunderline{夔}曰：“我念命女典掌樂事，當以詩樂教訓世適長子，使此長子正直而溫和，寬弘而莊慄，剛毅而不苛虐，簡易而不傲慢。教之詩樂，所以然者，詩言人之志意，歌詠其義以長其言。樂聲依此長歌為節，律呂和此長歌為聲。八音皆能和諧，無令相奪道理，如此則神人以此和矣。”夔答舜曰：“嗚呼！我擊其石磬，拊其石磬，諸音莫不和諧，百獸相率而舞。”樂之所感如此，是人神既已和矣。 \par}

\textcolor{red}{帝曰}：“\CJKunderline{龍}，朕\xpinyin*{堲}讒說\xpinyin*{殄}行，震驚朕師。\footnote{堲，疾。殄，絕。震，動也。言我疾讒說絕君子之行而動驚我眾，欲遏絕之。堲，徐在力反。讒，\CJKunderwave{切韻}仕咸反。說如字，注同，徐失銳反。殄，\CJKunderwave{切韻}徒典反。行，下孟反，注同。}命汝作納言，夙夜出納朕命，\textcolor{red}{惟允}。”\footnote{納言，喉舌之官。聽下言納於上，受上言宣於下，必以信。喉音侯。}

{\noindent\zhuan\zihao{6}\fzbyks 傳“堲疾”至“絕之”。正義曰：“堲”聲近疾，故為疾也。“殄，絕”、“震，動”皆\CJKunderwave{釋詁}文。讒人以善為惡，以惡為善，故言“我疾讒說絕君子之行”。眾人畏其讒口，故為讒也，“動驚我眾”，欲遏止之。 \par}

{\noindent\zhuan\zihao{6}\fzbyks 傳“納言”至“以信”。正義曰：\CJKunderwave{詩}美仲山甫為王之喉舌。喉舌者，宣出王命，如王咽喉口舌，故納言為“喉舌之官”也。此官主聽下言納於上，故以“納言”為名。亦主受上言宣於下，故言出朕命。“納言”不納於下，“朕命”有出無入，官名“納言”，雲“出納朕命”，互相見也。“必以信”者，不妄傳下言，不妄宣帝命，出納皆以信也。 \par}

{\noindent\shu\zihao{5}\fzkt “帝曰龍”至“惟允”。正義曰：帝呼龍曰：“龍，我憎疾人為讒佞之說，絕君子之行，而動驚我眾人，欲遏之。故命汝作納言之官。從早至夜出納我之教命，惟以誠信。”每事皆信則讒言自絕，命龍使勉之。 \par}

\textcolor{red}{帝曰}：“諮！汝二十有二人，\footnote{禹、垂、益、\CJKunderline{伯夷}、夔、龍六人新命有職，四嶽、十二牧凡二十二人，特敕命之。}欽哉！惟時亮\textcolor{red}{天功}。”\footnote{各敬其職，惟是乃能信立天下之功。}

{\noindent\zhuan\zihao{6}\fzbyks 傳“\CJKunderline{禹}垂”至“命之”。正義曰：傳以此文總結上事，據上文“詢於四嶽”，“諮，十有二牧”，及新命六官等，適滿二十二人，謂此也。其稷、契、\CJKunderline{皋陶}、\CJKunderline{殳斨}、\CJKunderline{伯與}、\CJKunderline{朱虎}、\CJKunderline{熊羆}七人仍舊，故不須敕命之。嶽、牧亦應是舊而敕命之者，嶽牧外內之官，常所諮詢，故亦敕之。\CJKunderline{鄭玄}云：“自‘諮，十有二牧’至‘帝曰龍’，皆月正元日格於文祖所敕命也。”案經“格於文祖”之後方始詢於四嶽，諮十二州牧，未必一日之內即得行此諸事，傳既不說，或歷日命授，乃總敕之，未必即是元日之事也。鄭以為“二十二人數\CJKunderline{殳斨}、\CJKunderline{伯與}、\CJKunderline{朱虎}、\CJKunderline{熊羆}，不數四嶽”。彼四人者直被讓而已，不言居官，何故敕使敬之也？嶽、牧俱是帝所諮詢，何以敕牧不敕嶽也？必非經旨，故孔說不然。 \par}

{\noindent\shu\zihao{5}\fzkt “帝曰諮”至“天功”。正義曰：帝既命用眾官，乃總戒敕之曰：“諮嗟！汝新命六人，及四嶽、十二牧凡二十有二人，汝各當敬其職事哉！惟是汝等敬事，則信實能立天下之功。天下之功，成王在於汝，可得不敬之哉！” \par}

\textcolor{red}{三載}考績，三考，黜陟幽明，\footnote{三年有成，故以考功。九歲則能否幽明有別，黜退其幽者，升進其明者。黜，醜律反。}庶績咸熙。分北\textcolor{red}{三苗}。\footnote{考績法明，眾功皆廣。三苗幽暗，君臣善否，分北流之，不令相從。善惡明。北如字，又音佩。令，力呈反。}

{\noindent\zhuan\zihao{6}\fzbyks 傳“三年”至“明者”。正義曰：三年一閏，天道成，人亦可以成功，故以三年考校其功之成否也。九年三考，則人之能否可知，幽明有別。“黜退其幽者”，或奪其官爵，或徙之遠方。“升進其明者”，或益其土地，或進其爵位也。 \par}

{\noindent\zhuan\zihao{6}\fzbyks 傳“考績”至“惡明”。正義曰：考績法明，人皆自勵，故得“眾功皆廣”也。“分北三苗”即是黜幽之事，故於“考績”之下言其流之。“分”謂別之。雲“北”者，言相背，必善惡不同。故知三苗幽暗，宜黜其君臣，乃有善否，分背流之，不令相從。俱徙之則善從惡，俱不徙則惡從善，言善惡不使相從，言舜之黜陟善惡明也。\CJKunderline{鄭玄}以為“流四凶者，卿為伯子，大夫為男，降其位耳，猶為國君”，故以“三苗為西裔諸侯，猶為惡，乃復分北流之”，謂分北西裔之三苗也。孔傳“竄三苗”為誅也，其身無復官爵，必非黜陟之限，其所“分北”，非彼“竄”者。王肅云：“三苗之民有赦宥者，復不從化，不令相從，分北流之。”王肅意彼赦宥者復繼為國君，至不復從化，故分北流之。\CJKunderline{禹}繼\CJKunderline{鯀}為崇伯，三苗未必絕後，傳意或如肅言。 \par}

{\noindent\shu\zihao{5}\fzkt “三載”至“三苗”。正義曰：自此以下史述舜事，非帝語也。言帝命群官之後，經三載乃考其功績，經三考則九載。“黜陟幽明”，明者升之，暗者退之。群官懼黜思升，各敬其事，故得“眾功皆廣”。前流四凶時，三苗之君竄之西裔，更紹其嗣，不滅其國。舜即政之後，三苗復不從化，是暗當黜之。其君臣有善有惡，舜復分北流其三苗。北，背也。善留惡去，使分背也。 \par}

\CJKunderline{舜}生三十，徵庸\footnote{言其始見試用。}三十，在位\footnote{歷試二年，攝位二十八年。}五十載，陟方乃死。\footnote{方，道也。舜即位五十年，升道南方巡守，死於蒼梧之野而葬焉。陟方，猶巡狩，天子外出巡視。三十徵庸，三十在位，服喪三年，其一在三十之數，為天子五十年，凡壽百一十二歲。}

{\noindent\zhuan\zihao{6}\fzbyks 傳“歷試”至“八年”。正義曰：上云：“乃言厎可績，三載”，則歷試當三年。雲“二年”者，其一即是徵用之年，已在上句三十之數，故惟有二年耳。受終居攝，尚在臣位，故歷試併為三十。“在位”謂在臣位也。 \par}

{\noindent\zhuan\zihao{6}\fzbyks 傳“方道”至“十二歲”。正義曰：\CJKunderwave{論語}云：“可謂仁之方也已。”孔注亦以“方”為道，常訓也。“舜即位五十年”，從格於文祖之後數之。“升道”謂乘道而行也。天子之行必是巡其所守之國，故通以“巡守”為名，未必以仲夏之月巡守南嶽也。\CJKunderwave{檀弓}雲“舜葬蒼梧之野”，是舜死蒼梧之野因而葬焉。孔以“月正元日”在“三載”、“遏密”之下，又\CJKunderwave{孟子}云，舜服堯三年喪畢，避堯之子，故“服喪三年”。三年之喪，二十五月而畢，其一年即在三十在位之數，惟有二年。是舜年六十二，為天子五十年，是舜“凡壽百一十二歲”也。\CJKunderwave{大禹謨}雲“帝曰‘朕宅帝位三十有三載’”，乃求禪\CJKunderline{禹}。\CJKunderwave{孟子}云：“舜薦\CJKunderline{禹}於天子，十有七年。”是在位五十年，其文明矣。\CJKunderline{鄭玄}讀此經云：“‘舜生三十’，謂生三十年也。‘登庸二十’，謂歷試二十年。‘在位五十載，陟方乃死’，謂攝位至死為五十年。舜年一百歲也。”\CJKunderwave{史記}云：“舜年三十堯舉用之，年五十攝行天子事，年五十八堯崩，年六十一而踐天子位，三十九年崩。”皆謬耳。 \par}

\textcolor{red}{帝釐}下土，方設居方，\footnote{言舜理四方諸侯，各設其官居其方。釐，力之反，馬云：“賜也，理也。”下土，絕句；一讀至“方”字絕句。}別生分類。\footnote{生，姓也。別其姓族，分其類，使相從。別，彼列反。分，方雲反，徐扶問反。}作\CJKunderwave{汩作}、\footnote{汩，治。作，興也。言其治民之功興，故為\CJKunderwave{汩作}之篇。亡。汩音骨。}\CJKunderwave{九共}九篇、\CJKunderwave{\textcolor{red}{槁飫}}。\footnote{槁，勞也。飫，賜也。凡十一篇，皆亡。共音恭，王己勇反，法也，馬同。槁，苦報反。飫，於據反。\CJKunderwave{槁飫}亦\CJKunderwave{書}篇名也。\CJKunderwave{汩作}等十一篇同此序。其文皆亡，而序與百篇之序同編，故存。今馬、鄭之徒百篇之序總為一卷，孔以各冠其篇首，而亡篇之序即隨其次篇居見存者之間。眾家經文並盡此，唯王注本下更有“\CJKunderwave{汩作}、\CJKunderwave{九共}故逸。故亦作古”。}

{\noindent\zhuan\zihao{6}\fzbyks 傳“言舜”至“其方”。正義曰：在\CJKunderwave{虞書}，知“帝”是舜也。“下土”對天子之辭，故云“理四方諸侯，各為其官居其方”。不知若為設之。凡此三篇之序，亦既不見其經，暗射無以可中。\CJKunderline{孔氏}為傳，復順其文為其傳耳，是非不可知也。他皆仿此。 \par}

{\noindent\zhuan\zihao{6}\fzbyks 傳“汩治”至“篇亡”。正義曰：“汩”之為治，無正訓也。“作”是起義，故為興也。“言其治民之功興”，以意言之耳。 \par}

{\noindent\zhuan\zihao{6}\fzbyks 傳“槁，勞。飫，賜也”。正義曰：\CJKunderwave{左傳}言“槁師”者，以師枯槁，用酒食勞之，是“槁”得為勞也。襄二十六年\CJKunderwave{左傳}云：“將賞，為之加膳，加膳則飫賜。”是“飫”得為賜也。亦不知勞賜之何所謂也。 \par}

{\noindent\shu\zihao{5}\fzkt “帝釐”至“槁飫”。正義曰：此序也，孔以\CJKunderwave{書序}序所以為作者之意，宜相附近，故引之各冠其篇首。其經亡者，以序附於本篇次而為之傳，故此序在此也。\CJKunderline{帝舜}治理下土諸侯之事，為各於其方置設其官，居其所在之方而統治之。又為民別其姓族之生，分別異類，各使相從作\CJKunderwave{汩作}篇，又作\CJKunderwave{九共}九篇，又作\CJKunderwave{槁飫}之篇，凡十一篇，皆亡。 \par}

%%% Local Variables:
%%% mode: latex
%%% TeX-engine: xetex
%%% TeX-master: "../Main"
%%% End:

%% -*- coding: utf-8 -*-
%% Time-stamp: <Chen Wang: 2024-04-02 11:42:41>

% {\noindent \zhu \zihao{5} \fzbyks } -> 注 (△ ○)
% {\noindent \shu \zihao{5} \fzkt } -> 疏

\chapter{卷四}


\section{大禹謨第三【偽】}


 {\noindent\zihao{6}\fzbyks \CJKunderwave{釋文}：“徐云：‘本\CJKunderwave{虞書}裛為一卷，凡十二卷，今依\CJKunderwave{七志}、\CJKunderwave{七錄}為十三卷。’” \par}

\textcolor{red}{\CJKunderline{皋陶}}矢厥謨，\footnote{矢，陳也。皋音高。陶音遙。}\CJKunderline{禹}成厥功，\footnote{陳其成功。}\CJKunderline{帝舜}申之。\footnote{申，重也，重美二子之言。○重，直用反。}作\CJKunderwave{大禹}、\CJKunderwave{皋陶謨}、\footnote{大禹謀九功，\CJKunderline{皋陶}謀九德。謨亦作謩。}\CJKunderwave{\textcolor{red}{益稷}}。\footnote{凡三篇。}

{\noindent\zhuan\zihao{6}\fzbyks 傳“矢，陳也”。正義曰：“矢，陳”，\CJKunderwave{釋詁}文。 \par}

{\noindent\zhuan\zihao{6}\fzbyks 傳“陳其成功”。正義曰：此是謨篇，\CJKunderline{禹}成其功，陳其言耳。蒙上“矢”文，故傳明之，言“陳其成功”也。序“成”在“厥”上，傳“成”在下者，序順上句，傳從便文，故倒也。 \par}

{\noindent\zhuan\zihao{6}\fzbyks 傳“申重”至“之言”。正義曰：“申，重”，\CJKunderwave{釋詁}文。\CJKunderwave{大禹謨}云：“帝曰：‘俞！地平天成，時乃功。’”又“帝曰：‘\CJKunderline{皋陶}，惟茲臣庶，罔或於予正。時乃功，懋哉！”\CJKunderwave{益稷}云：“迪朕德，時乃功。”皆是重美二子之言也。 \par}

{\noindent\zhuan\zihao{6}\fzbyks 傳“\CJKunderline{大禹}”至“九德”。正義曰：二篇皆是謨也，序以一“謨”總二篇，故傳明之。\CJKunderline{大禹}治水能致九功而言“謨”，以其序有“謨”文，故云“謨”也。 \par}

{\noindent\zhuan\zihao{6}\fzbyks 傳“凡三篇”。正義曰：\CJKunderwave{益稷}亦\CJKunderline{大禹}所謀，不言“謨”者，\CJKunderline{禹}謀言及益稷，非是益稷為謀，不得言\CJKunderwave{益稷謨}也。其篇雖有“夔曰”，夔言樂和，本非謀慮，不得謂之“夔謨”。 \par}

{\noindent\shu\zihao{5}\fzkt “\CJKunderline{皋陶}”至“益稷”。正義曰：\CJKunderline{皋陶}為\CJKunderline{帝舜}陳其謀，\CJKunderline{禹}為\CJKunderline{帝舜}陳已成所治水之功，\CJKunderline{帝舜}因其所陳從而重美之。史錄其辭，作\CJKunderwave{大禹}、\CJKunderwave{皋陶}二篇之謨，又作\CJKunderwave{益稷}之篇，凡三篇也。篇先\CJKunderwave{大禹}，序先言\CJKunderline{皋陶}者，\CJKunderwave{皋陶}之篇，\CJKunderline{皋陶}自先發端，\CJKunderline{禹}乃然而問之，\CJKunderline{皋陶}言在\CJKunderline{禹}先，故序先言\CJKunderline{皋陶}。其此篇以功大為先，故先\CJKunderwave{禹}也。\CJKunderwave{益稷}之篇亦是\CJKunderline{禹}之所陳，因\CJKunderline{皋陶}之言。而\CJKunderline{禹}論益稷，在\CJKunderwave{皋陶謨}後，故後其篇。 \par}

大禹謨\footnote{禹稱大，大其功。謨，謀也。}

{\noindent\zhuan\zihao{6}\fzbyks 傳“\CJKunderline{禹}稱”至“謀也”。正義曰：余文單稱“\CJKunderline{禹}”，而此獨加“大”者，故解之：\CJKunderline{禹}與\CJKunderline{皋陶}同為舜謀，而\CJKunderline{禹}功實大，\CJKunderline{禹}與\CJKunderline{皋陶}不等，史加大其功，使異於\CJKunderline{皋陶}，於此獨加“大”字與\CJKunderline{皋陶}總言故也。“謨，謀”，\CJKunderwave{釋詁}文。此三篇皆是舜史所錄，上取堯事，下錄\CJKunderline{禹}功，善於堯之知己，又美所禪得人，故包括上下以為\CJKunderwave{虞書}。其事以類相從，非由事之先後。若其不然，上篇已言舜死，於此豈死後言乎？此篇已言禪\CJKunderline{禹}，下篇豈受禪後乎？聰明史以類聚為文。計此三篇，\CJKunderline{禹}謨最在後，以\CJKunderline{禹}功大，故進之於先。\CJKunderwave{孟子}稱“舜薦\CJKunderline{禹}於天，十有七年”，則\CJKunderline{禹}攝一十七年，舜陟方乃死。不知\CJKunderline{禹}徵有苗，在攝幾年。史述\CJKunderline{禹}之行事，不必以攝位之年即徵苗民也。 \par}

曰若稽古，\CJKunderline{大禹}\footnote{順考古道而言之。}曰：“文命敷於四海，祗承於帝。\footnote{言其外布文德教命，內則敬承堯舜。文命，孔云：“文德教命也。”先儒云：“文命，\CJKunderline{禹}名。”}

{\noindent\zhuan\zihao{6}\fzbyks 傳“順考”至“言之”。正義曰：典是常行，謨是言語，故傳於典雲“行之”，於謨雲“言之”，皆是順考古道也。 \par}

{\noindent\zhuan\zihao{6}\fzbyks 傳“言其”至“堯舜”。正義曰：“敷於四海”即敷此文命，故言“外布文德教命”也。“四海”舉其遠地，故傳以“外”、“內”言之。“祗”訓敬也，\CJKunderline{禹}承堯舜二帝，故云“敬承堯舜”。傳不訓“祗”而直言“敬”，以易知而略之。 \par}

{\noindent\shu\zihao{5}\fzkt “曰若”至“於帝”。正義曰：史將錄\CJKunderline{禹}之事，故為題目之辭曰，能順而考案古道而言之者，是大功之\CJKunderline{禹}也。此\CJKunderline{禹}能以文德教命布陳於四海，又能敬承堯舜。外布四海，內承二帝，言其道周備。 \par}

\textcolor{red}{曰}：“後克艱厥後，臣克艱厥臣，政乃乂，黎民敏德。”\footnote{敏，疾也。能知為君難，為臣不易，則其政治，而眾民皆疾修德。○易，以豉反。治，直吏反。}帝曰：“俞！允若茲，嘉言罔攸伏，野無遺賢，萬邦咸寧。\footnote{攸，所也。善言無所伏，言必用。如此則賢才在位，天下安寧。俞，羊朱反。攸音由，徐以帚反。}稽於眾，捨己從人，不虐無告，不廢困窮，惟帝\textcolor{red}{時克}。”\footnote{帝謂堯也，舜因嘉言無所伏，遂稱堯德以成其義。考眾從人，矜孤愍窮，凡人所輕，聖人所重。舍音舍。告，故毒反。矜，居陵反。}

{\noindent\zhuan\zihao{6}\fzbyks 傳“敏疾”至“修德”。正義曰：許慎\CJKunderwave{說文}云：“敏，疾也。”是相傳為訓。“為君難，為臣不易”，\CJKunderwave{論語}文。能知為君難，為臣不易，則當謹慎恪勤，求賢自輔，故其政自然治矣。見善則用，知賢必進，眾民各自舉，則皆疾修德矣。此經上不言\CJKunderline{禹}者，承上\CJKunderline{禹}事，以可知而略之。 \par}

{\noindent\zhuan\zihao{6}\fzbyks 傳“攸所”至“下安”。正義曰：“攸，所”，\CJKunderwave{釋言}文。“善言無所伏”者，言其必用之也。言之善者必出賢人之口，但言之易，行之難，或有人不賢而言可用也，故“嘉言”與“賢”異其文也。如此用善言，任賢才在位，則天下安。 \par}

{\noindent\zhuan\zihao{6}\fzbyks 傳“帝謂”至“所重”。正義曰：舜稱為“帝”，故知“帝謂堯”也。舜因嘉言無所伏，以為堯乃能然，故遂稱堯德以成其義。此\CJKunderline{禹}言之義，以堯之聖智，無所不能，惟言其“考眾從人，矜孤愍窮”，以為堯之美者，此是“凡人所輕，聖人所重”。“不虐”、“不廢”，皆謂矜撫愍念之，互相通也。\CJKunderwave{王制}云：“少而無父謂之孤，老而無子謂之獨，老而無妻謂之鰥，老而無夫謂之寡。此四者天民之窮而無告者。”故此“無告”是彼四者。彼四者而此惟言孤者，四者皆孤也，言“孤”足以總之。言“困窮”，謂貧無資財也。 \par}

{\noindent\shu\zihao{5}\fzkt “曰後”至“時克”。正義曰：\CJKunderline{禹}為\CJKunderline{帝舜}謀曰：“君能重難其為君之事，臣能重難其為臣之職，則上之政教乃治，則下之眾民皆化而疾修其德。而帝曰：“然。信能如此，君臣皆能自難，並願善以輔己，則下之善言無所隱伏，在野無遺逸之賢，賢人盡用，則萬國皆安寧也。為人上者考於眾言，觀其是非，舍已之非，從人之是。不苛虐鰥寡孤獨無所告者，必哀矜之；不廢棄困苦貧窮無所依者，必愍念之。惟\CJKunderline{帝堯}於是能為此行，餘人所不能。”言“克艱”之不易也。 \par}

\CJKunderline{益}曰：“都！帝德廣運，乃聖乃神，乃武乃文。\footnote{益因舜言又美堯也。“廣”謂所覆者大，“運”謂所及者遠。聖無所不通，神妙無方，文經天地，武定禍亂。}皇天眷命，奄有四海，為天下君。”\footnote{眷，視。奄，同也。言堯有此德，故為天所命，所以勉舜也。眷，居倦反。奄，於檢反。}

{\noindent\zhuan\zihao{6}\fzbyks 傳“益因”至“禍亂”。正義曰：“廣”者闊之義，故為“所覆者大”。“運”者動之言，故為“所及者遠”。\CJKunderwave{洪範}雲“睿作聖”，言通知眾事，故為“無所不通”。案\CJKunderwave{易}曰“神者妙萬物而為言也”，又曰“神妙無方”，此言神道微妙，無可比方，不知其所以然。\CJKunderwave{易}又云：“陰陽不測之謂神。”\CJKunderwave{諡法}云：“經緯天地曰文，克定禍亂曰武。”經傳“文”、“武”倒者，經取韻句，傳以文重故也。 \par}

{\noindent\zhuan\zihao{6}\fzbyks 傳“眷視”至“勉舜”。正義曰：\CJKunderwave{詩}云：“乃眷西顧。”謂視而回首。\CJKunderwave{說文}亦以“眷”為視。“奄，同”，\CJKunderwave{釋言}文。益因帝言盛稱堯善者，亦勸勉舜，冀之必及堯也。 \par}

{\noindent\shu\zihao{5}\fzkt “\CJKunderline{益}曰”至“下君”。正義曰：益承帝言，嘆美堯德曰：“嗚呼！\CJKunderline{帝堯}之德，廣夫運行。乃聖而無所不通，乃神而微妙無方，乃武能克定禍亂，乃文能經緯天地。以此為大天顧視而命之，使同有四海之內，為天下之君。” \par}

\CJKunderline{禹}曰：“惠迪吉，從逆兇，惟影響。”\footnote{迪，道也。順道吉從逆兇。吉凶之報，若影之隨形，響之應聲。言不虛。迪，徒力反。響，許丈反。}\CJKunderline{益}曰：“籲！戒哉！儆戒無虞，罔失法度。\footnote{先籲後戒，欲使聽者精其言。虞，度也。無億度，謂無形。戒於無形，備慎深。秉法守度，言有恆。籲，況俱反。度，徒布反。虞度，徙洛反。}罔遊於逸，罔淫於樂。\footnote{淫，過也。遊逸過樂，敗德之原。富貴所忽，故特以為戒。樂音洛。}任賢勿貳，去邪勿疑。疑謀勿成，百志惟熙。\footnote{一意任賢，果於去邪，疑則勿行，道義所存於心，日以廣矣。去，起呂反。熙，火其反。罔違道以幹百姓之譽，幹，求也。失道求名，古人賤之。}罔咈百姓以從己之慾。\footnote{咈，戾也。專欲難成，犯眾興禍，故戒之。咈，扶弗反。戾，連弟反。}無怠無荒，四夷來王。”\footnote{言天子常戒慎，無怠惰荒廢，則四夷歸往之。怠音待。惰，徒臥反。}

{\noindent\zhuan\zihao{6}\fzbyks 傳“迪，道也”。正義曰：\CJKunderwave{釋詁}文。 \par}

{\noindent\zhuan\zihao{6}\fzbyks 傳“先籲”至“有恆”。正義曰：\CJKunderwave{堯典}傳云：“籲，疑怪之辭。”此無可怪，聞善驚而為聲耳。“先籲後戒”者，驚其言之美，然後設戒辭，欲使聽者精審其言。“虞，度”，\CJKunderwave{釋詁}文。“無億度”者，謂不有此事，無心億度之。\CJKunderwave{曲禮}云：“凡為人子者，聽於無聲，視於無形。”戒於無形見之事，言備慎深也。安不忘危，治不忘亂，是其慎無形也。法度當執守之，故以“秉法守度”解不失，言有恆也。 \par}

{\noindent\zhuan\zihao{6}\fzbyks 傳“淫過”至“為戒”。正義曰：“淫”者過度之意，故為過也。“逸”謂縱體，“樂”謂適心，縱體在於逸遊，適心在於淫恣，故以“遊逸過樂”為文。二者敗德之源，富貴所忽，故特以為戒。 \par}

{\noindent\zhuan\zihao{6}\fzbyks 傳“幹求”至“賤之”。正義曰：“幹，求”，\CJKunderwave{釋言}文。“失道求名”謂曲取人情，苟悅眾意，古人賤之。 \par}

{\noindent\zhuan\zihao{6}\fzbyks 傳“咈戾”至“戒之”。正義曰：\CJKunderwave{堯典}已訓“咈”悉戾。彼謂戾朋儕，此謂戾在下，故詳其文耳。“專欲難成，犯眾興禍”，襄十年\CJKunderwave{左傳}文。 \par}

{\noindent\shu\zihao{5}\fzkt “\CJKunderline{禹}曰”至“來王”。正義曰：\CJKunderline{禹}曰：“益言謀及世事，言人順道則吉，從逆則兇。吉凶之報，惟若影之隨形，響之應聲。”言其無不報也。益聞\CJKunderline{禹}語，驚懼而言曰：“籲！誠如此言，宜誡慎之哉！所誡者，當儆誡其心無億度之事。”謂忽然而有，當誡慎之：“無失其守法度，使行必有恆，無違常也。無遊縱於逸豫，無過耽於戲樂，當誡慎之以保己也。任用賢人勿有二心，逐去回邪勿有疑惑。所疑之謀勿成用之，如是則百種志意惟益廣也。無違越正道以求百姓之譽，無反戾百姓以從己心之慾。常行此事，無怠惰荒廢，則四夷之國皆來歸往之。”此亦所以勸勉舜也。 \par}

\CJKunderline{禹}曰：“於！帝念哉！德惟善政，政在養民。\footnote{嘆而言“念”，重其言。為政以德，則民懷之。}水火金木土谷惟修，\footnote{言養民之本在先修六府。}正德、利用、厚生惟和，\footnote{正德以率下，利用以阜財，厚生以養民，三者和，所謂善政。}九功惟叙，九叙惟歌。\footnote{言六府三事之功有次敘，皆可歌樂，乃德政之致。樂音洛。}

{\noindent\zhuan\zihao{6}\fzbyks 傳“嘆而”至“懷之”。正義曰：“於”，嘆辭。嘆而言“念”，自重其言，欲使帝念之。此史以類相從，共為篇耳。非是一時之事，不使念益言也。\CJKunderline{禹}謀以九功為重，知“重其言”者，九功之言也。 \par}

{\noindent\zhuan\zihao{6}\fzbyks 傳“言養”至“六府”。正義曰：下文帝言“六府”即此經六物也。六者民之所資，民非此不生，故言“養民之本在先修六府”也。“府”者藏財之處，六者貨財所聚，故稱“六府”。襄二十七年\CJKunderwave{左傳}云：“天生五材，民並用之。”即是水火金木土，民用此自資也。彼惟五材，此兼以谷為六府者，谷之於民尢急，谷是土之所生，故於土下言之也。此言五行，與\CJKunderwave{洪範}之次不同者，\CJKunderwave{洪範}以生數為次，此以相刻為次，便文耳。六府是民之急，先有六府乃可施教，故先言“六府”，後言“三事”也。 \par}

{\noindent\zhuan\zihao{6}\fzbyks 傳“止德”至“善政”。正義曰：“正德”者，自正其德，居上位者正己以治民，故所以率下人。“利用”者，謂在上節儉，不為縻費，以利而用，使財物殷阜，利民之用，為民興利除害，使不匱乏，故所以阜財。“阜財”謂財豐大也。“厚生”謂薄徵徭，輕賦稅，不奪農時，令民生計溫厚，衣食豐足，故所以養民也。“三者和”謂德行正、財用利、生資厚。立君所以養民，人君若能如此，則為君之道備矣。故謂“善政”，結上“德惟善政”之言。此三者之次，人君自正乃能正下，故以“正德”為先；利用然後厚生，故後言“厚生”。“厚生”謂財用足，禮讓行也。 \par}

{\noindent\zhuan\zihao{6}\fzbyks 傳“言六”至“之致”。正義曰：上六下三，即是“六府三事”，此總雲“九功”，知六府三事之功為九功。“惟敘”者，即上“惟修”、“惟和”為次序。事皆有敘，民必歌樂君德，故九敘皆可歌樂，乃人君德政之致也。言下民必有歌樂，乃為善政之驗，所謂和樂興而頌聲作也。 \par}

戒之用休，董之用威，勸之以九歌，俾勿壞。”\footnote{休，美。董，督也。言善政之道，美以戒之，威以督之，歌以勸之。使政勿壞，在此三者而已。○俾，必爾反。壞，乎怪反。}帝曰：“俞！地平天成，六府三事允治，萬世永賴，時乃功。”\footnote{水土治曰“平”，五行敘曰“成”。因\CJKunderline{禹}陳九功而嘆美之，言是汝之功，明眾臣不及。}

{\noindent\zhuan\zihao{6}\fzbyks 傳“休美”至“而已”。正義曰：“休，美”，\CJKunderwave{釋詁}文。又云：“董，督，正也。”是“董”為督也。此“戒之”、“董之”、“勸之”皆謂人君自戒勸，欲使善政勿壞，在此三事而已。文七年\CJKunderwave{左傳}云，晉卻缺言於趙宣子，引此一經，乃言：“九功之德皆可歌也，謂之九歌。若吾子之德莫可歌也，其誰來之？盍使睦者歌吾子乎？”言“九功之德皆可歌”者，若水能灌溉，火能烹飪，金能斷割，木能興作，土能生殖，谷能養育，古之歌詠各述其功，猶如漢魏已來樂府之歌事，歌其功用，是舊有成辭。人君修治六府以自勸勉，使民歌詠之，三事亦然。 \par}

{\noindent\zhuan\zihao{6}\fzbyks 傳“水土”至“不及”。正義曰：\CJKunderwave{釋詁}云：“平，成也。”是“平”、“成”義同，天、地文異而分之耳。天之不成，由地之不平，故先言“地平”，本之於地以及天也。\CJKunderline{禹}平水土，故“水土治曰平”。五行之神，佐天治物，系之於天，故“五行敘曰成”。\CJKunderwave{洪範}雲“\CJKunderline{鯀}堙洪水，汩陳其五行，彝倫攸斁”，\CJKunderline{禹}治洪水，“彝倫攸敘”，是\CJKunderline{禹}命五行敘也。帝因\CJKunderline{禹}陳九功而嘆美之，指言是汝之功，明眾臣不及。 \par}

{\noindent\shu\zihao{5}\fzkt “\CJKunderline{禹}曰”至“乃功”。正義曰：\CJKunderline{禹}因益言，又獻謀於帝曰“嗚呼！帝當念之哉！”言：“所謂德者惟是善於政也。政之所為，在於養民。養民者，使水火金木土谷此六事惟當修治之。正身之德，利民之用，厚民之生，此三事惟當諧和之。修和六府三事，九者皆就有功，九功惟使皆有次敘，九事次敘惟使皆可歌樂，此乃德之所致。是德能為善政之道，終當不得怠惰。但人雖為善，或寡令終，故當戒敕之念用美道，使民慕美道行善。又督察之用威罰。”言其不善當獲罪。“勸勉之以九歌之辭。但人君善政，先致九歌成辭自勸勉也。用此事，使此善政勿有敗壞之時”。勸帝使長為善也。帝答\CJKunderline{禹}曰：“汝之所言為然。汝治水土，使地平天成，六府三事信皆治理，萬代長所恃賴，是汝之功也。”歸功於\CJKunderline{禹}，明群臣不及。 \par}

帝曰：“格，汝\CJKunderline{禹}。朕宅帝位三十有三載，耄期倦於勤。汝惟不怠，總朕師。”\footnote{八十、九十曰耄，百年曰期頤。言己年老，厭倦萬機，汝不懈怠於位，稱總我眾，欲使攝。格，庚白反。朕，直錦反。耄，莫報反。倦，其卷反。頤，以之反。厭，於豔反。解，於賣反。}\CJKunderline{禹}曰：“朕德罔克，民不依。\CJKunderline{皋陶}邁種德，德乃降，黎民懷之。\footnote{邁，行。種，布。降，下。懷，歸也。言己無德，民所不能依。\CJKunderline{皋陶}布行其德，下治於民，民歸服之。種，章用反。降，江巷反。}帝念哉！念茲在茲，釋茲在茲，\footnote{茲，此。釋，廢也。念此人在此功，廢此人在此罪。言不可誣。}名言茲在茲，允出茲在茲，惟帝念功。”\footnote{名言此事，必在此義；信出此心，亦在此義。言\CJKunderline{皋陶}之德以義為主，所宜念之。}

{\noindent\zhuan\zihao{6}\fzbyks 傳“八十”至“使攝”。正義曰：“八十、九十曰耄，百年曰期頤”，\CJKunderwave{曲禮}文也。如\CJKunderwave{舜典}之傳，計舜年六十三即政，至今九十五矣。年在耄、期之間，故並言之。鄭云：“期，要也。頤，養也。不知衣服食味，孝子要盡養之道而已。”孔意當然。 \par}

{\noindent\zhuan\zihao{6}\fzbyks 傳“邁行”至“服之”。正義曰：“邁，行”、“降，下”，\CJKunderwave{釋言}文。又云：“懷，來也。”來亦歸也。種物必佈於地，故為布也。 \par}

{\noindent\zhuan\zihao{6}\fzbyks 傳“茲此”至“可誣”。正義曰：“茲，此”，\CJKunderwave{釋詁}文。“釋”為舍義，故為廢也。\CJKunderline{禹}之此意，欲令帝念\CJKunderline{皋陶}。下雲“惟帝念功”，“念”是念功，知“廢”是廢罪，言念、廢心依其實，不可誣罔也。 \par}

{\noindent\zhuan\zihao{6}\fzbyks 傳“名言”至“念之”。正義曰：“名言”謂己發於口，“信出”謂始發於心，皆據欲舉\CJKunderline{皋陶}，必先念慮於心，而後宣之於口。先言“名言”者，己對帝讓\CJKunderline{皋陶}，即是名言之事，故先言其意。然後本其心，故後言“信出”。“以義為主”者，言己讓\CJKunderline{皋陶}，事非虛妄，以義為尚。 \par}

{\noindent\shu\zihao{5}\fzkt “帝曰格”至“念功”。正義曰：此舜言。將禪\CJKunderline{禹}帝，呼\CJKunderline{禹}曰：“來，汝\CJKunderline{禹}。我居帝位已三十有三載，在耄、期之間，厭倦於勤勞。汝惟在官不懈怠，可代我居帝位，總領我眾。”\CJKunderline{禹}讓之曰：“我德實無所能，民必不依就我也。”言己不堪總眾也。“\CJKunderline{皋陶}行佈於德，德乃下洽於民，眾皆歸服之，可令\CJKunderline{皋陶}攝也。我所言者，帝當念之哉！凡念愛此人，在此功勞，知有功乃用之。釋廢此人，在此罪釁，知有罪乃廢之”。言進人退人不可誣也。“名目言談此事，必在此事之義而名言之。若信實出見此心，必在此心之義而出見之”。言己名言其口，出見其心，以舉\CJKunderline{皋陶}，皆在此義，不有虛妄。“帝當念錄其功以禪之”。言\CJKunderline{皋陶}堪攝位也。 \par}

帝曰：“\CJKunderline{皋陶}，惟茲臣庶，罔或干予正。\footnote{或，有也。無有干我正。言順命。}汝作士，明於五刑，以弼五教，期於予治。\footnote{弼，輔。期，當也。嘆其能以刑輔教，當於治體。治音稚。當，丁浪反，又如字。}刑期於無刑，民協於中，時乃功，懋哉！”\footnote{雖或行刑，以殺止殺，終無犯者。刑期於無所刑，民皆命於大中之道，是汝之功，勉之。懋音茂。}\CJKunderline{皋陶}曰：“帝德罔愆，臨下以簡，御眾以寬。\footnote{愆，過也。善則歸君，人臣之義。愆音牽。}罰弗及嗣，賞延於世。\footnote{嗣亦世，俱謂子。延，及也。父子罪不相及，而及其賞。道德之政。}宥過無大，刑故無小。\footnote{過誤所犯，雖大必宥。不忌故犯，雖小必刑。宥音又。}罪疑惟輕，功疑惟重。\footnote{刑疑附輕，賞疑從重，忠厚之至。}與其殺不辜，寧失不經。好生之德，洽於民心，茲用不犯於有司。”\footnote{辜，罪。經，常。司，主也。\CJKunderline{皋陶}因帝勉己，遂稱帝之德，所以明民不犯上也。寧失不常之罪，不枉不辜之善，仁愛之道。辜音孤。好，呼報反。}帝曰：“俾予從欲以治，四方風動，惟乃之休。”\footnote{使我從心所欲而政以治，民動順上命，若草應風，是汝能明刑之美。}

{\noindent\zhuan\zihao{6}\fzbyks 傳“弼輔”至“治體”。正義曰：\CJKunderwave{書傳}稱“左輔右弼”，是“弼”亦輔也。期要是相當之言，故為當也。傳言“當於治體”，言\CJKunderline{皋陶}用刑，輕重得中，於治體與正相當也。 \par}

{\noindent\zhuan\zihao{6}\fzbyks 傳“雖或”至“勉之”。正義曰：言\CJKunderline{皋陶}或行刑，乃是以殺止殺。為罪必將被刑，民終無犯者。要使人無犯法，是期於無所用刑，刑無所用。此“期”為限，與前經“期”義別，而\CJKunderwave{論語}所謂“勝殘去殺”矣。“民皆合於大中”，言舉動每事得中，不犯法憲，是“合大中”即\CJKunderwave{洪範}所謂“皇極”是也。 \par}

{\noindent\zhuan\zihao{6}\fzbyks 傳“愆過”至“之義”。正義曰：“愆，過”，\CJKunderwave{釋言}文。\CJKunderwave{坊記}云：“善則稱君，過則稱己，則民作忠。”是善則稱君，人君之義也。“臨下”據其在上，“御眾”斥其治民，簡易、寬大，亦不異也。\CJKunderwave{論語}云：“居敬而行簡，以臨其民，不亦可乎？”是臨下宜以簡也。又曰：“寬則得眾。”“居上不寬，吾何以觀之哉？”是御眾宜以寬也。 \par}

{\noindent\zhuan\zihao{6}\fzbyks 傳“嗣亦”至“及也”。正義曰：“嗣”謂繼父，“世”謂後胤，故“俱謂子”也。“延”訓長，以長及物，故“延”為及也。 \par}

{\noindent\zhuan\zihao{6}\fzbyks 傳“辜罪”至“之道”。正義曰：“辜，罪”，\CJKunderwave{釋詁}文。“經，常”，“司，主”，常訓也。\CJKunderline{皋陶}因帝勉己，遂稱帝之德。所以明民不犯上者，自由帝化使然，非己力也。“不常之罪”者，謂罪大，非尋常小罪也。枉殺無罪，妄免有罪，二者皆失，必不得民心。寧妄免大罪，不枉殺無罪，以好生之心故也。大罪尚赦，小罪可知。欲極言不可枉殺不辜，寧放有罪故也，故言非常大罪以對之耳。“寧失不經”與“殺不辜”相對，故為放赦罪人，原帝之意，等殺無罪，寧放有罪。傳言帝德之善，寧失有罪，不枉殺無罪，是仁愛之道。各為文勢，故經傳倒也。“治”謂沾漬優渥，洽於民心，言潤澤多也。 \par}

{\noindent\shu\zihao{5}\fzkt “帝曰\CJKunderline{皋陶}”至“之休”。正義曰：帝以\CJKunderline{禹}讓\CJKunderline{皋陶}，故述而美之。帝呼之曰：“\CJKunderline{皋陶}，惟此群臣眾庶，皆無敢有干犯我正道者。由汝作士官，明曉於五刑，以輔成五教，當於我之治體。用刑期於無刑，以殺止殺，使民合於中正之道，令人每事得中，是汝之功，當勉之哉！”\CJKunderline{皋陶}以帝美己，歸美於君曰：“民合於中者，由帝德純善，無有過失，臨臣下以簡易，御眾庶以優寬。罰人不及後嗣，賞人延於來世。宥過失者無大，雖大亦有之。刑其故犯者無小，雖小必刑之。罪有疑者，雖重，從輕罪之。功有疑者，雖輕，從重賞之。與其殺不辜非罪之人，寧失不經不常之罪。以等枉殺無罪，寧妄免有罪也。由是故帝之好生之德，下洽於民心，民服帝德如此，故用是不犯於有司。”言民之無刑非已力也。帝又述之曰：“使我從心所欲而為政，以大治四方之民，從我化如風之動草，惟汝用刑之美。”言己知其有功也。 \par}

帝曰：“來，\CJKunderline{禹}。降水儆予，成允成功，惟汝賢。\footnote{水性流下，故曰下水。儆，戒也。能成聲教之信，成治水之功，言\CJKunderline{禹}最賢，重美之。儆，居領反。重，直用反。}克勤於邦，克儉於家，不自滿假，惟汝賢。\footnote{滿謂盈實。假，大也。言\CJKunderline{禹}惡衣薄食，卑其宮室，而盡力為民，執心謙沖，不自盈大。假，工雅反。盡，津忍反。為，於偽反。}汝惟不矜，天下莫與汝爭能。汝惟不伐，天下莫與汝爭功。\footnote{自賢曰矜，自功曰伐。言\CJKunderline{禹}推善讓人而不失其能，不有其勞而不失其功，所以能絕眾人。}

{\noindent\zhuan\zihao{6}\fzbyks 傳“水性”至“美之”。正義曰：“降水”，洪水也。水性下流，故曰下水。\CJKunderline{禹}以治水之事儆戒於予。\CJKunderwave{益稷}云：“予創若時，娶於塗山，辛壬癸甲，啟呱呱而泣，予弗子，惟荒度土功之事。”雖文在下篇，實是欲禪前事，故帝述而言之。\CJKunderwave{禹貢}言治水功成云：“朔南暨聲教。”故知“成允”是“成聲教之信”，“成功”是“成治水之功”也。前已言地平天成是汝功，今復說治水之事，“言\CJKunderline{禹}最賢，重美之”也。\CJKunderline{禹}實聖人，美其賢者，其性為聖，其功為賢，猶\CJKunderwave{易·繫辭}雲“可久則賢人之德，可大則賢人之業”，亦是聖人之事。 \par}

{\noindent\zhuan\zihao{6}\fzbyks 傳“滿謂”至“盈大”。正義曰：“滿”以器喻，故為盈實也。“假，大”，\CJKunderwave{釋詁}文。言己無所不知，是為自滿。言己無所不能，是為自大。\CJKunderline{禹}實不自滿大，故為賢也。\CJKunderwave{論語}美\CJKunderline{禹}之功德云：“惡衣服，菲飲食，卑宮室，而盡力乎溝洫。”故傳引彼。惡衣、薄食、卑其宮室是“儉於家”，盡力為民是“勤於邦”。上言其功，此言其德，故再雲“惟汝賢”。 \par}

{\noindent\zhuan\zihao{6}\fzbyks 傳“自賢”至“眾人”。正義曰：自言己賢曰矜，自言己功曰伐。\CJKunderwave{論語}云：“願無伐善。”\CJKunderwave{詩}云：“矜其車甲。”“矜”與“伐”俱是誇義，以經有“爭能”、“爭功”，故別解之耳。弗矜莫與汝爭能，即矜者矜其能也。賢、能大同小異，故“自賢”解“矜”。\CJKunderwave{老子}云：“夫惟不爭，故天下莫能與之爭。”是故不矜伐而不失其功能，此所以能絕異於眾人也。 \par}

予懋乃德，嘉乃丕績，天之歷數在汝躬，汝終陟元后。\footnote{丕，大也。歷數謂天道。元，大也；大君，天子。舜善\CJKunderline{禹}有治水之大功，言天道在汝身，汝終當升為天子。丕，普悲反。}人心惟危，道心惟微，惟精惟一，允執厥中。\footnote{危則難安，微則難明，故戒以精一，信執其中。}無稽之言勿聽，弗詢之謀勿庸。\footnote{無考無信驗，不詢專獨，終必無成，故戒勿聽用。聽，徐天定反。}

{\noindent\zhuan\zihao{6}\fzbyks 傳“丕大”至“天子”。正義曰：“丕，大”，\CJKunderwave{釋詁}文。“歷數”謂天曆運之數，帝王易姓而興，故言“歷數謂天道”。\CJKunderline{鄭玄}以“歷數在汝身謂有圖籙之名”，孔無讖緯之說，義必不然。當以大功既立，眾望歸之，即是天道在身。\CJKunderwave{釋詁}“元”訓為首，首是體之大也。\CJKunderwave{易}曰“大君有命”，是“大君”謂天子也。 \par}

{\noindent\zhuan\zihao{6}\fzbyks 傳“危則”至“其中”。正義曰：居位則治民，治民必須明道，故戒之以“人心惟危，道心惟微”。道者經也，物所從之路也。因言“人心”，遂雲“道心”。人心惟萬慮之主，道心為眾道之本。立君所以安人，人心危則難安。安民必須明道，道心微則難明。將欲明道，必須精心。將欲安民，必須一意。故以戒精心一意。又當信執其中，然後可得明道以安民耳。 \par}

{\noindent\zhuan\zihao{6}\fzbyks 傳“無考”至“聽用”。正義曰：為人之君不當妄用人言，故又戒之：“無可考校之言謂無信驗，不詢於眾人之謀謂專獨用意。”言無信驗是虛妄之言，獨為謀慮是偏見之說，二者終必無成，故戒令勿聽用也。“言”謂率意為語，“謀”謂豫計前事，故互文也。 \par}

可愛非君？可畏非民？非元后何戴？後非眾罔與守邦？\footnote{民以君為命，故可愛。君失道，民叛之，故可畏。言眾戴君以自存，君恃眾以守國，相須而立。}欽哉！慎乃有位，敬修其可願，四海困窮，天祿永終。\footnote{有位，天子位。可原謂道德之美。困窮謂天民之無告者。言為天子勤此三者，則天之祿籍長終汝身。}惟口出好興戎，朕言不再。”\footnote{好謂賞善，戎謂伐惡。言口榮辱之主，慮而宣之，成於一也。出如字，徐尺遂反。好如字，徐許到反。}

{\noindent\zhuan\zihao{6}\fzbyks 傳“民以”至“而立”。正義曰：百人無主，不散則亂，故“民以君為命”。君尊，民畏之，嫌其不愛，故言“愛”也。民賤，君忽之，嫌其不畏，故言“畏”也。 \par}

{\noindent\zhuan\zihao{6}\fzbyks 傳“有位”至“汝身”。正義曰：上雲“汝終陟元后”，命昇天位，知其慎汝有位，慎天子位也。道德人之可願，知“可願”者，是道德之美也。惟言“四海困窮”，不結言民之意，必謂四海之內困窮之民，令天子撫育之。故知如\CJKunderwave{王制}所云，孤獨鰥寡“此四者，天民之窮而無告者”，此是困窮者也。言為天子，當慎天位，修道德，養窮民，勤此三者，則天之祿籍長終汝身。祿謂福祿，籍謂名籍，言享大福，保大名也。 \par}

{\noindent\zhuan\zihao{6}\fzbyks 傳“好謂”至“於一”。正義曰：昭二十八年\CJKunderwave{左傳}云：“慶賞刑威曰君。”君出言有賞有刑，“出好”謂愛人而出好言，故為賞善。“興戎”謂疾人而動甲兵，故謂伐惡。\CJKunderwave{易·繫辭}曰：“言語者，君子之樞機。樞機之發，榮辱之主。”必當慮之於心，然後宣之於口，故成之於一而不可再。帝言“我命汝昇天位”者，是慮而宣之，此言故不可再。 \par}

{\noindent\shu\zihao{5}\fzkt “帝曰來”至“不再”。正義曰：帝不許\CJKunderline{禹}讓，呼之曰：“來，\CJKunderline{禹}。下流之水儆戒於我，我恐不能治之。汝能成聲教之信，能成治水之功，惟汝之賢。汝能勤勞於國。”謂盡力於溝洫。“能節儉於家”。謂薄飲食，卑宮室。“常執謙沖，不自滿溢誇大，惟汝之賢也”。又申美之：“汝惟不自矜誇，故天下莫敢與汝爭能。汝惟不自稱伐，故天下莫敢與汝爭功。”美功之大也。“我今勉汝之德，善汝大功，天之歷運之數帝位當在汝身，汝終當升此大君之位，宜代我為天子”。因戒以為君之法：“民心惟甚危險，道心惟甚幽微。危則難安，微則難明，汝當精心，惟當一意，信執其中正之道，乃得人安而道明耳。又為人君，不當妄受用人語。無可考驗之言，勿聽受之。不是詢眾之謀，勿信用之。”言“民所愛者，豈非人君乎？民以君為命，故愛君也”。言“君可畏者，豈非民乎？君失道則民叛之，故畏民也。眾非大君而何所奉戴？無君則民亂，故愛君也。君非眾人無以守國，無人則國亡，故畏民也。君民相須如此，當宜敬之哉！謹慎汝所有之位，守天子之位，勿使失也。敬修其可原之事”。謂道德之美，人所原也。“養彼四海困窮之民，使皆得存立，則天之祿籍長終汝身矣”。又告\CJKunderline{禹}：“惟口之所言，出好事，興戎兵，非善思慮無以出口，我言不可再發。”令\CJKunderline{禹}受其言也。 \par}

\CJKunderline{禹}曰：“枚卜功臣，惟吉之從。”\footnote{枚謂歷卜之而從其吉。此\CJKunderline{禹}讓之志。枚音梅。}帝曰：“\CJKunderline{禹}，官佔，惟先蔽志，昆命於元龜。\footnote{帝王立卜佔之官，故曰官佔。蔽，斷。昆，後也。官佔之法，先斷人志，後命於元龜，言志定然後卜。蔽，必世反，徐甫世反。斷，丁亂反。}朕志先定，詢謀僉同，鬼神其依，龜筮協從，卜不習吉。”\footnote{習，因也。然已謀之於心，謀及卜筮，四者合從，卜不因吉，無所枚卜。僉，七潛反。}\CJKunderline{禹}拜稽首，固辭。\footnote{再辭曰固。}帝曰：“毋！惟汝諧。”\footnote{言毋，所以禁其辭。\CJKunderline{禹}有大功德，故能諧和元后之任。禁，今鴆反，又音金。}

{\noindent\zhuan\zihao{6}\fzbyks 傳“枚謂”至“之志”。正義曰：\CJKunderwave{周禮}有銜枚氏，所銜之物狀如箸。今人數物雲一枚、兩枚，則“枚”是籌之名也。“枚卜”謂人人以次歷申卜之，似老枚數然。然請卜不請筮者，舉動也。 \par}

{\noindent\zhuan\zihao{6}\fzbyks 傳“帝王”至“後卜”。正義曰：佔是卜人之佔，而云“官佔”者，帝王立卜筮之官，故曰“官佔”。\CJKunderwave{洪範}“稽疑”云：“擇建立卜筮人。”是帝王立卜筮之官。\CJKunderwave{周禮}司寇斷獄為蔽獄，是“蔽”為斷也。“昆，後”，\CJKunderwave{釋言}文。官佔之法，先斷人志，後命元龜，言志定然後卜也。\CJKunderwave{洪範}雲“汝則有大疑，謀及乃心，謀乃卿士，謀及庶人”，是先斷人志；乃雲“謀及卜筮”，是後命元龜。“元龜”謂大龜也。 \par}

{\noindent\zhuan\zihao{6}\fzbyks 傳“習因”至“枚上”。正義曰：\CJKunderwave{表記}雲“卜筮不相襲”。鄭云：“襲，因也。”然則“習”與“襲”同。重衣謂之襲，習是後因前，故為因也。“朕志先定”言已謀之於心。“龜筮協從”是謀及卜筮。經言“詢謀僉同”，謀及卿士、庶人，謀皆同心。“鬼神其依”，即是龜筮之事，卜筮通鬼神之意，故言“鬼神其依”。“龜筮協從”，謂卜得吉，是依從也。志先定也，謀僉同也，鬼神依也，龜筮從也，四者合從，然後命汝。卜法不得因吉，無所復枚卜也。如帝此言，既謀既卜，方始命\CJKunderline{禹}，仍請枚卜者，帝與朝臣私謀私卜，將欲命\CJKunderline{禹}，\CJKunderline{禹}不預謀，故不在，更請卜也。 \par}

{\noindent\zhuan\zihao{6}\fzbyks 傳“言毋”至“之任”。正義曰：\CJKunderwave{說文}云：“毋，止之也。”其字從女，內有一畫，象有奸之者，禁止令勿奸也。古人言毋，猶今人言莫，是“言毋者，所以禁其辭”，令勿辭。 \par}

{\noindent\shu\zihao{5}\fzkt “\CJKunderline{禹}曰”至“汝諧”。正義曰：\CJKunderline{禹}以讓而不許，更請帝曰：“每以一枚歷卜功臣，惟吉之人，從而受之。”帝曰：“\CJKunderline{禹}，卜官之佔，惟能先斷人志，後乃命其大龜。我授汝之志先以定矣，又詢於眾人，其謀又皆同美矣。我後謀及鬼神，加之卜筮，鬼神其依我矣，龜筮複合從矣。卜法不得因前之吉更復卜之，不須復卜也。”\CJKunderline{禹}猶拜而後稽首，固辭。帝曰：“毋！”毋者，禁止其辭也。“惟汝能諧和此元后之任，汝宜受之”。 \par}

正月朔旦，受命於神宗，\footnote{受舜終事之命。神宗，文祖之宗廟，言“神”尊之。○正音政，徐音徵。}率百官，若帝之初。\footnote{順舜初攝帝位故事奉行之。}

{\noindent\zhuan\zihao{6}\fzbyks 傳“受舜”至“尊之”。正義曰：\CJKunderwave{舜典}說舜之初“受終於文祖”，此言“若帝之初”，知“受命”即是“舜終事之命”也。神宗猶彼文祖，故云“文祖之宗廟”。“文祖”言祖有文德，“神宗”言神而尊之，名異而實同。神宗當舜之始祖。案\CJKunderwave{帝嚳}云，黃帝生昌意，昌意生顓頊，顓頊生窮蟬，窮蟬生敬康，敬康生句芒，句芒生蟜牛，蟜牛生瞽瞍，瞽瞍生舜。即是舜有七廟，黃帝為始祖，其顓頊與窮蟬為二祧，敬康、句芒、蟜牛、瞽瞍為親廟，則文祖為黃帝顓頊之等也。 \par}

{\noindent\zhuan\zihao{6}\fzbyks 傳“順舜”至“行之”。正義曰：“若”不得為“如”也。\CJKunderwave{舜典}巡守之事，言如初者皆言“如”，不言“若”，知此“若”為順也。順舜初攝帝位故事而盡行之，其奉行者當如\CJKunderwave{舜典}在“璿璣”以下，“班瑞群后”以上也。其巡守非率百官之事，舜尚自為陟方，\CJKunderline{禹}攝帝位，未得巡守，此是舜史所錄，以為\CJKunderwave{虞書}，故言順帝之初，奉行帝之事故，自美禪之得人也。 \par}

{\noindent\shu\zihao{5}\fzkt “正月”至“之初”。正義曰：舜即政三十三年，命\CJKunderline{禹}代己，\CJKunderline{禹}辭不獲免。乃以明年正月朔旦，受終事之命於舜神靈之宗廟，總率百官。順帝之初攝故事，言與舜受禪之初，其事悉皆同也。此年舜即政三十四年，九十六也。 \par}

帝曰：“諮！\CJKunderline{禹}，惟時有苗弗率，汝徂徵。”\footnote{三苗之民數幹王法。率，循。徂。往也。不循常道，言亂逆。命\CJKunderline{禹}討之。數音朔。}\CJKunderline{禹}乃會群后，誓於師曰：“濟濟有眾，咸聽朕命。\footnote{會諸侯共伐有苗。軍旅曰誓。濟濟，眾盛之貌。濟，子禮反。}蠢茲有苗，昏迷不恭，\footnote{蠢，動。昏，暗也。言其所以宜討之。蠢，春允反。}

{\noindent\zhuan\zihao{6}\fzbyks 傳“三苗”至“討之”。正義曰：\CJKunderwave{呂刑}稱苗民“作五虐之刑”，皇帝“遏絕苗民，無世在下”，謂堯初誅三苗。\CJKunderwave{舜典}雲“竄三苗於三危”，謂舜居攝之時，投竄之也。\CJKunderwave{舜典}又云“庶績咸熙，分北三苗”，謂舜即位之後，往徙三苗也。今復不率命，命\CJKunderline{禹}徂徵，是三苗之民數千王誅之事，\CJKunderline{禹}率眾徵之，猶尚逆命。即三苗是諸侯之君，而謂之“民”者，以其頑愚，號之為“民”。\CJKunderwave{呂刑}雲“苗民弗用靈”，是謂為“民”也。\CJKunderwave{呂刑}稱堯誅三苗雲“無世在下”，而得有苗國曆代常存者，“無世在下”謂誅叛者，絕後世耳，蓋不滅其國，又立其近親紹其先祖。\CJKunderline{鯀}既殛死於羽山，\CJKunderline{禹}乃代為崇伯，三苗亦竄其身而存其國。故舜時有被宥者，復不從化，更分北流之。下傳雲“三苗之國，左洞庭，右彭蠡”，其國在南方。蓋分北之時，使為南國君，今復不率帝道。“率，循”、“徂，往”皆\CJKunderwave{釋詁}文。不循帝道，言其亂逆，以其亂逆，故命\CJKunderline{禹}討之。案\CJKunderwave{舜典}皆言舜受終之後，萬事皆舜主之。舜自巡守，不稟堯命。此言“若帝之初”，其事亦應同矣。而此言命\CJKunderline{禹}徵苗，舜復陟方乃死，與舜受堯禪事不同者，以題曰\CJKunderwave{虞書}，即舜史所錄，明其詳於舜，事略於堯、\CJKunderline{禹}也。 \par}

{\noindent\zhuan\zihao{6}\fzbyks 傳“會諸”至“之貌”。正義曰：軍眾曰誓”，\CJKunderwave{曲禮}文也。隱八年\CJKunderwave{穀梁傳}曰：“誥誓不及五帝，盟詛不及三土，交質不及二伯。”“二伯”謂齊桓公、晉文公也。“不及者”，言於時未有也。據此文，五帝之世有誓。\CJKunderwave{周禮}立司盟之官，三王之世有盟也。\CJKunderwave{左傳}雲平王與鄭交質，二伯之前有質也。\CJKunderwave{穀梁傳}漢初始作，不見經文，妄言之耳。美軍眾而言“濟濟”，知是“眾盛之貌”。 \par}

{\noindent\zhuan\zihao{6}\fzbyks 傳“蠢動”至“討之”。正義曰：蠢，動”，\CJKunderwave{釋詁}文。\CJKunderwave{釋訓}云：“蠢，不遜也。”郭璞云：“蠢動為惡，不謙遜也。”日入為“昏”，是為暗也。動為惡而暗於事，言其所以宜討之。 \par}

侮慢自賢，反道敗德，\footnote{狎侮先王，輕慢典教，反正道，敗德義。侮，亡甫反。慢，亡諫反。}君子在野，小人在位，\footnote{廢仁賢，任奸佞。}民棄不保，天降之咎，\footnote{言民叛，天災之。咎，其九反。}肆予以爾眾士，奉辭罰罪。\footnote{肆，故也。辭謂“不恭”，罪謂“侮慢”以下事。}爾尚一乃心力，其克有勳。”\footnote{尚，庶幾。一汝心力，以從我命。}

{\noindent\zhuan\zihao{6}\fzbyks 傳“狎侮”至“德義”。正義曰：“侮”謂輕人身，“慢”謂忽言語，故為“狎侮先王，輕慢典教”。“侮”、“慢”義同，因有二字而分釋之。\CJKunderwave{論語}雲“狎大人，侮聖人之言”，則“狎”、“侮”為異。\CJKunderwave{旅獒}雲“狎侮君子”，則“狎”、“侮”意亦同。\CJKunderline{鄭玄}云：“狎，慣忽也。”慣見而忽之，是“侮”之義。傳取“狎”、“侮”連言之。慢先王典，教自謂己賢，不知先王訓教。“道”者，物所由之路，“德”謂自得於心。反正道，從邪徑，敗德義，毀正行也。 \par}

{\noindent\zhuan\zihao{6}\fzbyks 傳“廢仁賢，任奸佞”。正義曰：雖則下愚之君，皆雲好賢疾佞，非知賢而廢之，知佞而任之。但愚人所好，必同於民，賢求其心，佞從其欲，以賢為惡，謂佞為善，故仁賢見廢，奸佞被任，此則“昏迷”之狀也。 \par}

{\noindent\zhuan\zihao{6}\fzbyks 傳“肆故”至“下事”。正義曰：“肆，故”，\CJKunderwave{釋詁}文。所奉之辭即所伐之罪，但天子責其不恭，數其身罪，因其文異而分之。 \par}

{\noindent\zhuan\zihao{6}\fzbyks 傳“尚庶”至“我命”。正義曰：\CJKunderwave{釋言}云：“庶幾，尚也。”反以相解，故“尚”為庶幾。 \par}

{\noindent\shu\zihao{5}\fzkt “帝曰諮”至“有勳”。正義曰：史言\CJKunderline{禹}雖攝位，帝尊如故，時有苗國不順，帝曰：“諮嗟！汝\CJKunderline{禹}，惟時有苗之國不循帝道，汝往徵之。”\CJKunderline{禹}得帝命，乃會群臣諸侯，告誓於眾曰：“濟濟美盛之有眾，皆聽從我命。今蠢蠢然動而不遜者，是此有苗之君。昏暗迷惑，不恭敬王命。侮慢眾常，自以為賢。反戾正道，敗壞德義。君子在野，小人在位。由此民棄叛之，不保其有眾，上天降之殃咎。故我以爾眾士，奉此譴責之辭，伐彼有罪之國。汝等庶幾同心盡力，以從我命，其必能有大功勳，不可懈惰。” \par}

三旬，苗民逆命。\footnote{旬，十日也。以師臨之，一月不服，責舜不先有文誥之命、威讓之辭，而便憚之以威，脅之以兵，所以生辭。誥，古報反。憚，徒旦反，一音丹末反。脅，許業反。}\CJKunderline{益}贊於\CJKunderline{禹}曰：“惟德動天，無遠弗\xpinyin*{屆}。\footnote{贊，佐。屆，至也。益以此義佐\CJKunderline{禹}，欲其修德致遠。屆音戒。}滿招損，謙受益，時乃天道。\footnote{自滿者人損之，自謙者人益之，是天之常道。}


{\noindent\zhuan\zihao{6}\fzbyks 傳“旬十”至“生辭”。正義曰：\CJKunderwave{堯典}雲“三百有六旬”，是知“旬，十日”也。以師臨之，一月不服者，責舜不先有文告之命、威讓之辭，而便憚之以威，脅之以兵，所以有苗得生辭也。傳知然者，昭十三年\CJKunderwave{左傳}論征伐之事云：“告之以文辭，董之以武師。”是用兵者先告，不服然後伐之。今經無先告之文，而有逆命之事，故知責舜不先有文告之命，而即脅之以兵。其“文告之命、威讓之辭”，\CJKunderwave{國語}亦有其事。夫以大舜足達用兵之道，而不為文告之命，使之得生辭者，有苗數幹王誅，逆者難以言服，故憚之以威武，任其生辭。待其有辭，為之振旅，彼若師退而服，我復更有何求？為退而又不降，復往必無辭說。不恭而徵之，有辭而舍之，正是柔服之道也。若先告以辭，未必即得從命。不從而後行師，必將大加殺戮。不以文誥，感德自來，固是大聖之遠謀也。 \par}

{\noindent\zhuan\zihao{6}\fzbyks 傳“贊佐”至“致遠”。正義曰：\CJKunderwave{禮}有贊佐，是助祭之人，故“贊”為佐也。“屆，至也”，\CJKunderwave{釋詁}文。經云：“惟德動天。”天遠而難動，德能動遠。又言：“無遠不屆。”乃據人言德動遠人，無不至也。益以此義佐\CJKunderline{禹}，欲修德致遠，使有苗自來也。德之動天，經傳多矣。\CJKunderwave{禮運}云，聖人順民，“天不愛其道”，地不愛其寶，故天降膏露，地出醴泉”。如此之類，皆德動之也。 \par}

{\noindent\zhuan\zihao{6}\fzbyks 傳“自滿”至“常道”。正義曰：自以為滿，人必損之；自謙受物，人必益之。\CJKunderwave{易寶·謙卦}彖曰：“天道虧盈而益謙，地道變盈而流謙，鬼神害盈而福謙，人道惡盈而好謙。”是滿招損，謙受益為天道之常也。益言此者，欲令\CJKunderline{禹}修德息師，持謙以待有苗。 \par}

帝初於歷山，往于田，日號泣於旻天、於父母，\footnote{仁覆愍下謂之旻天。言舜初耕於歷山之時，為父母所疾，日號泣於旻天及父母，克己自責，不責於人。田，本或作畋。號，戶高反。旻，武巾反。}負罪引慝，祗載見瞽瞍，夔夔齋慄，瞽亦允若。\footnote{慝，惡。載，事也。夔夔，悚懼之貌。言舜負罪引惡，敬以事見於父，悚懼齋莊，父亦信順之。言能以至誠感頑父。慝，他則反。見，賢遍反。瞽音古。瞍，素後反。夔，求龜反。齋音側皆反。}至諴感神，矧茲有苗。”\footnote{諴，和。矧，況也。至和感神，況有苗乎！言易感。諴音咸。矧，失忍反。易，以豉反。}

{\noindent\zhuan\zihao{6}\fzbyks 傳“仁覆”至“責於人”。正義曰：“仁覆愍下謂之旻天”，\CJKunderwave{詩}毛傳文也。旻，愍也。求天愍己，故呼曰“旻天”。\CJKunderwave{書傳}言：“舜耕於歷山。”\CJKunderline{鄭玄}云：“歷山在河東。”是耕於歷山之時，為父母所疾，故往于田，日號泣於旻天。何為其然也！\CJKunderwave{孟子}曰：“怨慕也。長息問於公明高曰：‘舜往于田，則予既聞命矣，號泣於旻天及父母，即吾不知矣。’公明高曰：‘非爾所知也。’我竭力耕田，供為子職而已，父母不愛我，何哉？大孝終身慕父母，五十而慕者，予於大舜見之矣。”言舜之號泣怨慕者，克己自責，不責於人也。 \par}

{\noindent\zhuan\zihao{6}\fzbyks 傳“慝惡”至“頑父”。正義曰：“慝”之為惡，常訓耳。\CJKunderwave{舜典}已訓“載”為事，以非常訓，故詳其文。“夔夔”與“齋慄”共文，故為“悚懼之貌”。自負其罪，引惡歸己，事瞽同耳，丁寧深言之。“敬以事見於父”者，謂恭敬，自因事務須見父，恭敬以見。夔夔然悚懼齊慄，是見時之貌。“父亦信順之”者，謂當以事見之時，順帝意不悖怒也。言“能以至誠感頑父”者，言感使當時暫以順耳，不能使每事信順，變為善人。故\CJKunderwave{孟子}說舜既被堯徵用，堯妻之以二女，瞽瞍猶與象欲謀殺舜而分其財物，是下愚之性，終不可改。但舜善養之，使不至於奸惡而已。 \par}

{\noindent\zhuan\zihao{6}\fzbyks 傳“諴和”至“易感”。正義曰：“諴”亦咸也，咸訓為皆，皆能相從，亦和之義也。“矧，況”，\CJKunderwave{釋言}文。上言德能動天，次言帝能感瞽。天以玄遠難感，瞽以頑愚難感，言苗民近於天而智於瞽，故言感天感瞽以況之。天是神也，覆動上天，言至和尚能感天神，而況於有苗乎！言有苗易感。神覆動天而不覆言瞽者，以瞽雖愚，猶是人類，天神事與人隔，感天難於感瞽，故舉難者以見之。其實天與瞽俱言難感，以況有苗易於彼二者。 \par}

禹拜昌言曰：“俞！”班師振旅。\footnote{昌，當也。以益言為當，故拜受而然之，遂還師。兵入曰振旅。言整眾。當，丁浪反，下同。還，經典皆音旋。}帝乃誕敷文德，\footnote{遠人不服，大布文德以來之。誕音但。}舞幹羽於兩階，\footnote{幹，楯。羽，翳也。皆舞者所執。修闡文教，舞文舞於賓主階間，抑武事。階徐音皆。楯，食允反。翳，於計反。闡，尺善反。}七旬，有苗格。\footnote{討而不服，不討自來，明御之者必有道。三苗之國，左洞庭，右彭蠡，在荒服之例，去京師二千五百里也。洞，徒弄反。蠡音禮。}

{\noindent\zhuan\zihao{6}\fzbyks 傳“昌當”至“整眾”。正義曰：“昌，當也”，\CJKunderwave{釋詁}文。\CJKunderline{禹}以益言為當，拜舜而已即還。還不請者，\CJKunderwave{春秋}襄十九年“晉士匄帥師侵齊，聞齊侯卒，乃還”，\CJKunderwave{公羊傳}曰“大夫以君命出，進退在大夫”，是言進退由將，不須請也。或可當時請帝乃還，文不具耳。“兵入曰振旅”，\CJKunderwave{釋天}文。與\CJKunderwave{春秋}二傳皆有此文。振，整也。言整眾而還。 \par}

{\noindent\zhuan\zihao{6}\fzbyks 傳“遠人”至“來之”。正義曰：“遠人不服，文德以來之”，\CJKunderwave{論語}文也。益贊於\CJKunderline{禹}，使修德，而帝自誕敷者，言君臣同心．“大布”者，多設文德之教，君臣共行之也。 \par}

{\noindent\zhuan\zihao{6}\fzbyks 傳“幹楯”至“武事”。正義曰：\CJKunderwave{釋言}云：“幹，扞也。”孫炎曰：“幹楯，自蔽扞也。”以楯為人扞。通以“幹”為楯名，故“幹”為楯。\CJKunderwave{釋言}又云：“纛，翳也。“郭璞云：“舞者持以自蔽翳也。”故\CJKunderwave{明堂位}云，朱干玉鏚以舞大武。戚，斧也，是武舞執斧執楯。\CJKunderwave{詩}云：“左手執籥，右手秉翟。”是文舞執籥，故幹羽皆舞者所執。修闡文教，不復征伐，故舞文德之舞於賓主階間，言帝抑武事也。經雲“舞幹羽”，即亦舞武也。傳惟言舞文者，以據器言之，則有武有文，俱用以為舞，而不用於敵，故教為文也。 \par}

{\noindent\zhuan\zihao{6}\fzbyks 傳“討而”至“百里”。正義曰：“御之必有道”者，不恭而往徵，得辭而振旅，而御之以道。\CJKunderwave{史記}吳起對魏武侯云：“昔三苗氏左洞庭，右彭蠡，德義不修，而\CJKunderline{禹}滅之。”此言來服，則是不滅。吳起言滅者，以武侯恃險，言滅以懼之。辯士之說，不必皆依實也。知“在荒服之例”者，以其地驗之為然。\CJKunderwave{禹貢}五服，甸、侯、綏、要、荒。荒最在外，王畿面五百里，其外四服又每服五百里，是去京師為二千五百里。 \par}

{\noindent\shu\zihao{5}\fzkt “三旬”至“苗格”。正義曰：\CJKunderline{禹}既誓於眾，而以師臨苗。經三旬，苗民逆帝命，不肯服罪。益乃進謀以佐於\CJKunderline{禹}曰：“惟是有德，能動上天。苟能修德，無有遠而不至。”因言行德之事：“自滿者招其損，謙虛者受其益，是乃天之常道。”欲\CJKunderline{禹}修德，謙虛以來苗。既說其理，又言其驗：“帝乃初耕於歷山之時，為父母所疾。往至於田，日號泣於旻天。於父母乃自負其罪，自引其惡，恭敬以事見父瞽瞍，夔夔然悚懼，齋莊戰慄，不敢言己無罪。舜謙如此，雖瞽瞍之頑愚，亦能信順。帝至和之德尚能感於冥神，況此有苗乎！”言其苗易感於瞽瞍。\CJKunderline{禹}拜益受之當言，曰：“然。”然益語也。遂還師整眾而歸。\CJKunderline{帝舜}乃大布文德，舞幹羽於兩階之間，七旬而有苗自服來至。言主聖臣賢，御之有道也。 \par}

\section{皋陶謨第四（皋陶謨上）}


皋陶謨\footnote{謨，謀也。\CJKunderline{皋陶}為\CJKunderline{帝舜}謀。為，於偽反。}

{\noindent\zhuan\zihao{6}\fzbyks 傳“謨謀”至“舜謀”。正義曰：孔以此篇惟與\CJKunderline{禹}言，嫌其不對\CJKunderline{帝舜}，故言“為\CJKunderline{帝舜}謀”。將言“為\CJKunderline{帝舜}謀”故又訓“謨”為謀，以詳其文。 \par}

曰若稽古，\CJKunderline{皋陶}\footnote{亦順考古道以言之。夫典、謨，聖帝所以立治之本，皆師法古道以成不易之則。夫音扶。治，直吏反，下同。}曰：“允迪厥德，謨明弼諧。”\footnote{迪，蹈。厥，其也。其，古人也。言人君當信蹈行古人之德，謀廣聰明以輔諧其政。蹈，徒報反。}\CJKunderline{禹}曰：“俞，如何？”\footnote{然其言，問所以行。}

{\noindent\zhuan\zihao{6}\fzbyks 傳“亦順”至“之則”。正義曰：二謨其目正同，故云“亦順考古道以言”也。堯舜考古以行，謂之為“典”，\CJKunderline{大禹}\CJKunderline{皋陶}考古以言謂之為“謨”。“典”、“謨”之文不同，其目皆雲考古，故傳明言其意：“夫典、謨，聖帝所以立治之本。”雖言行有異，皆是考法古道以成不易之則，故史皆以“稽古”為端目。但君則行之，臣則言之，以尊卑不同，故“典”、“謨”名異。\CJKunderline{禹}亦為君而云“謨者，\CJKunderline{禹}在舜時未為君也。顧氏亦同此解。\CJKunderline{皋陶}德劣於\CJKunderline{禹}，皆是考古以言，故得同其題目。但\CJKunderline{禹}能“敷於四海，祇承於帝”，\CJKunderline{皋陶}不能然，故此下更無別辭耳。 \par}

{\noindent\zhuan\zihao{6}\fzbyks 傳“迪蹈”至“其政”。正義曰：\CJKunderwave{釋詁}云：“迪，道也。”聲借為導，導音與蹈同，故“迪”又為蹈也。其德即其上“稽古”，故曰“其，古人也”。而臣為君謀，故云“言人君當信蹈行古人之德”，謂蹈履依行之也。“謀廣聰明”，“聰明”者自是己性，又當受納人言，使多所聞見，以博大此聰明，以輔弼和諧其政。經惟言“明”，傳亦有“聰”者，以耳目同是所用，故以“聰明”言之。此“曰”上不言“\CJKunderline{皋陶}”，猶\CJKunderline{大禹}為謀，“曰”上不言“\CJKunderline{禹}”。\CJKunderline{鄭玄}云：“以\CJKunderline{皋陶}下屬為句。”則“稽古”之下無人名，與上三篇不類甚矣。 \par}

\CJKunderline{皋陶}曰：“都！慎厥身修，思永。\footnote{嘆美之重也。慎修其身，思為長久之道。“身修”絕句。}惇敘九族，庶明勵翼，邇可遠，在茲。”\footnote{言慎修其身，厚次敘九族，則眾庶皆明其教，而自勉勵翼戴上命，近可推而遠者，在此道。惇，\CJKunderwave{切韻}都昆反。}\CJKunderline{禹}拜昌言曰：“俞。”\footnote{以\CJKunderline{皋陶}言為當，故拜受而然之。當，丁浪反，下同。}

{\noindent\zhuan\zihao{6}\fzbyks 傳“嘆美”至“之道”。正義曰：案傳之言以“修”為上讀，顧氏亦同也。 \par}

{\noindent\zhuan\zihao{6}\fzbyks 傳“言慎”至“此道”。正義曰：自身以外，九族為近，故慎修其身，又厚次敘九族，猶堯之為政，先以親九族也。人君既能如此，則眾庶皆明其教，而各自勉勵翼戴上命。昭九年\CJKunderwave{左傳}說晉叔向言“翼戴天子”，故以為“翼戴上命”，言如鳥之羽翼而奉戴之。王者率己以化物，親親以及遠，故從近可推而至於遠者，在修己身、親九族之道。王肅云：“以眾賢明為砥礪，為羽翼。”鄭云：“厲，作也，以眾賢明作輔翼之臣。”與孔不同。 \par}

{\noindent\shu\zihao{5}\fzkt “曰若”至“曰俞”。正義曰：史將言\CJKunderline{皋陶}之能謀，故為題目之辭曰，能順而考案古道而言之者，是\CJKunderline{皋陶}也。其為帝謀曰：“為人君者當信實蹈行古人之德，而謀廣其聰明之性，以輔諧己之政事，則善矣。”\CJKunderline{禹}曰：“然。”然其謀是也。“此當如何行之？”\CJKunderline{皋陶}曰：“嗚呼！重其事而嘆美之。“行上謀者，當謹慎其己身，而修治人之事，思為久長之道。又厚次敘九族之親而不遺棄，則眾人皆明曉上意，而各自勉勵翼戴上命，行之於近，而可推而至遠者，在此道也。”\CJKunderline{禹}乃拜受其當理之言，曰：“然。”美其言而拜受之。 \par}

\CJKunderline{皋陶}曰：“都！在知人，在安民。”\footnote{嘆修身親親之道在知人所信任，在能安民。}\CJKunderline{禹}曰：“籲！咸若時，惟帝其難之。\footnote{言\CJKunderline{帝堯}亦以知人安民為難，故曰：“籲！”}知人則哲，能官人。安民則惠，黎民懷之。\footnote{哲，智也。無所不知，故能官人。惠，愛也。愛則民歸之。}能哲而惠，何憂乎\CJKunderline{驩兜}？\footnote{佞人亂德，堯憂其敗政，故流放之。}何遷乎有苗？何畏乎巧言令色孔壬？”\footnote{孔，甚也。巧言，靜言庸違。令色，象恭滔天。\CJKunderline{禹}言有苗、\CJKunderline{驩兜}之徒甚佞如此，堯畏其亂政，故遷放之。}


{\noindent\zhuan\zihao{6}\fzbyks 傳“哲智”至“歸之”。正義曰：“哲，智”，\CJKunderwave{釋言}文。舍人曰：“哲，大智也。”無所不知，知人之善惡，是能官人。“惠，愛”，\CJKunderwave{釋詁}文。君愛民則民歸之。 \par}

{\noindent\zhuan\zihao{6}\fzbyks 傳“孔甚”至“放之”。正義曰：“孔，甚”，\CJKunderwave{釋詁}文。上句既言“\CJKunderline{驩兜}”、“有苗”，則此“巧言令色”\CJKunderline{共工}之行也，故以\CJKunderwave{堯典}\CJKunderline{共工}之事解之：“巧言，靜言庸違”也，“令色，象恭滔天”也。“孔壬”之文在三人之下，總上三人皆甚佞也。“苗”言其名，“巧言令色”言其行，令其文首尾互相見，故傳通言之：“\CJKunderline{禹}言有苗、\CJKunderline{驩兜}之徒甚佞如此，堯畏其亂政，故遷放之。”傳不言\CJKunderline{共工}，故云“之徒”以包之。“遷”與“憂”、“畏”亦互相承言，畏之而憂，乃遷之也。四凶惟言三者，馬融云：“\CJKunderline{禹}為父隱，故不言\CJKunderline{鯀}也。” \par}

{\noindent\shu\zihao{5}\fzkt “\CJKunderline{皋陶}曰都在”至“孔壬”。正義曰：\CJKunderline{皋陶}以\CJKunderline{禹}然其言，更述修身親親之道，嘆而言曰：“人君行此道者，在於知人善惡，擇善而信任之；在於能安下民，為政以安定之也。”\CJKunderline{禹}聞此言，乃驚而言曰：“籲！人君皆如是。能知人，能安民，惟\CJKunderline{帝堯}猶其難之，況餘人乎？知人善惡，則為大智。能用官，得其人矣。能安下民，則為惠政，眾民皆歸之矣。此甚不易也，若\CJKunderline{帝堯}能智而惠，則當朝無奸佞，何憂懼於\CJKunderline{驩兜}之佞而流放之？何須遷徙於有苗之君？何所畏懼於彼巧言令色為甚佞之人？”三兇見惡，\CJKunderline{帝堯}方始去之，是知人之難。 \par}

\CJKunderline{皋陶}曰：“都！亦行有九德。\footnote{言人性行有九德，以考察真偽則可知。○行，下孟反，注“性行”、“行正直”之“行”同。}亦言其人有德，乃言曰，載采采。”\footnote{載，行。採，事也。稱其人有德，必言其所行某事某事以為驗。}


{\noindent\zhuan\zihao{6}\fzbyks 傳“言人”至“可知”。正義曰：“言人性行有九德”，下文所云是也。如此九者考察其真偽，則人之善惡皆可知矣。然則\CJKunderline{皋陶}之賢不及\CJKunderline{帝堯}遠矣，\CJKunderline{皋陶}知有此術，\CJKunderline{帝堯}無容不知；而有四凶在朝，\CJKunderline{禹}言帝難之者，堯朝之有四凶，晦跡以顯舜爾。\CJKunderline{禹}言惟帝難之，說彼甚佞，因其成敗以示教法，欲開\CJKunderline{皋陶}之志，故舉大事以為戒；非是此實甚佞，堯不能知也。顧氏亦云：“堯實不以此為難。今雲難者，俯同流俗之稱也。” \par}

{\noindent\zhuan\zihao{6}\fzbyks 傳“載行”至“為驗”。正義曰：“載”者，運行之義，故為行也。此謂薦舉人者稱其人有德，欲使在上用之，必須言其所行之事，雲見此人常行其某事某事，由此所行之事以為有德之驗。\CJKunderwave{論語}云：“如有所譽者，其有所試矣。”是言試之於事，乃可知其德。 \par}

{\noindent\shu\zihao{5}\fzkt “\CJKunderline{皋陶}”至“采采”。正義曰：\CJKunderline{禹}既言知人為難，\CJKunderline{皋陶}又言行之有術，故言曰：“嗚呼！人性雖則難知，亦當考察其所行有九種之德。人慾稱薦人者，不直言可用而已，亦當言其人有德。問其德之狀，乃言曰其德之所行某事某事。以所行之事為九德之驗，如此則可知也。” \par}

\CJKunderline{禹}曰：“何？”\footnote{問九德品例。}\CJKunderline{皋陶}曰：“寬而慄，\footnote{性寬弘而能莊慄。}柔而立，\footnote{和柔而能立事。}愿而恭，\footnote{愨愿而恭恪。願音願。愨，\CJKunderwave{切韻}苦角反。恪，苦各反。}亂而敬，\footnote{亂，治也。有治而能謹敬。}擾而毅，\footnote{擾，順也。致果為毅。擾，而小反，徐音饒。毅，五既反。}直而溫，\footnote{行正直而氣溫和。}簡而廉，\footnote{性簡大而有廉隅。}剛而塞，\footnote{剛斷而實塞。斷，丁亂反。}彊而義。\footnote{無所屈撓，動必合義。撓，女孝反。}彰厥有常，吉哉！\footnote{彰，明。吉，善也。明九德之常，以擇人而官之，則政之善。}

{\noindent\zhuan\zihao{6}\fzbyks 傳“性寬”至“莊慄”。正義曰：此九德之文，\CJKunderwave{舜典}雲“寬而慄，直而溫”，與此正同。彼雲“剛而無虐，簡而無傲”，與此小異。彼言“剛失入虐”，此言“剛斷而能實塞”，“實塞”亦是不為虐。彼言“簡失入傲”，此言“簡大而有廉隅”，“廉隅”亦是不為傲也。九德皆人性也。\CJKunderline{鄭玄}云：“凡人之性有異，有其上者，不必有下；有其下者，不必有上。上下相協，乃成其德。”是言上下以相對，各令以相對兼而有之，乃為一德。此二者雖是本性，亦可以長短自矯。寬弘者失於緩慢，故性寬弘而能矜莊嚴慄，乃成一德。九者皆然也。 \par}

{\noindent\zhuan\zihao{6}\fzbyks 傳“愨愿而恭恪”。正義曰：“願”者愨謹良善之名。謹願者失於遲鈍，貌或不恭，故愨願而能恭恪乃為德。 \par}

{\noindent\zhuan\zihao{6}\fzbyks 傳“亂治”至“謹敬”。正義曰：“亂，治”，\CJKunderwave{釋詁}文。有能治者，謂才高於人也，堪撥煩理劇者也。負才輕物，人之常性，故有治而能謹敬乃為德也。“願”言“恭”，“治”雲“敬”者，恭在貌，敬在心；願者遲鈍，外失於儀，故言“恭”以表貌；治者輕物，內失於心，故稱“敬”以顯情。“恭”與“敬”其事亦通，“願”其貌恭而心敬也。 \par}

{\noindent\zhuan\zihao{6}\fzbyks 傳“擾順”至“為毅”。正義曰：\CJKunderwave{周禮·大宰}云：“以擾萬民。”\CJKunderline{鄭玄}云：“擾猶馴也。”\CJKunderwave{司徒}云：“安擾邦國。”鄭云：“擾亦安也。”“擾”是安馴之義，故為順也。“致果為毅”，宣二年\CJKunderwave{左傳}文。彼文以“殺敵為果，致果為毅”，謂能致果敢殺敵之心，是為強毅也。和順者失於不斷，故順而能決乃為德也。 \par}

{\noindent\zhuan\zihao{6}\fzbyks 傳“性簡”至“廉隅”。正義曰：“簡”者，寬大率略之名。志遠者遺近，務大者輕細，弘大者失於不謹，細行者不修廉隅，故簡大而有廉隅乃為德也。 \par}

{\noindent\zhuan\zihao{6}\fzbyks 傳“剛斷”至“實塞”。正義曰：“塞”訓實也。剛而能斷失於空疏，必性剛正而內充實乃為德也。 \par}

{\noindent\zhuan\zihao{6}\fzbyks 傳“無所”至“合義”。正義曰：強直自立，無所屈撓，或任情違理，失於事宜，動合道義乃為德也。鄭注\CJKunderwave{論語}云：“剛謂強，志不屈撓。”即“剛”、“強”義同。此“剛”、“強”異者，“剛”是性也，“強”是志也。當官而行，無所避忌，剛也。執己所是，不為眾撓，強也。“剛”、“強”相近，鄭連言之。“寬”謂度量寬弘，“柔”謂性行和柔，“擾”謂事理擾順，三者相類，即\CJKunderwave{洪範}雲“柔克”也。“願”謂容貌恭正，“亂”謂剛柔治理，“直”謂身行正直，三者相類，即\CJKunderwave{洪範}雲“正直”也。“簡”謂器量凝簡，“剛”謂事理剛斷，“強”謂性行堅強，三者相類，即\CJKunderwave{洪範}雲“剛克”也。而九德之次，從“柔”而至“剛”也，惟“擾而毅”在“願”、“亂”之下耳。其\CJKunderwave{洪範}三德，先人事而後天地，與此不同。 \par}

{\noindent\zhuan\zihao{6}\fzbyks 傳“彰明”至“之善”。正義曰：“彰，明”、“吉，善”，常訓也。此句言用人之義。所言九德，謂彼人常能然者。若暫能為之，未成為德。故人君取士，必明其九德之常，知其人常能行之，然後以此九者之法擇人而官之，則為政之善也。“明”謂人君明知之。王肅云：“明其有常則善也，言有德當有恆也。”其意亦言彼能有常，人君能明之也。鄭云：“人能明其德，所行使有常，則成善人矣。”其意謂彼人自明之，與孔異也。 \par}

{\noindent\shu\zihao{5}\fzkt “\CJKunderline{禹}曰”至“吉哉”。正義曰：\CJKunderline{皋陶}既言其九德，\CJKunderline{禹}乃問其品例曰：“何謂也？”\CJKunderline{皋陶}曰：“人性有寬弘而能莊慄也，和柔而能立事也，愨願而能恭恪也，治理而能謹敬也，和順而能果毅也，正直而能溫和也，簡大而有廉隅也，剛斷而能實塞也，強勁而合道義也。人性不同，有此九德。人君明其九德所有之常，以此擇人而官之，則為政之善哉！” \par}

日宣三德，夙夜浚明有家。\footnote{三德，九德之中有其三。宣，布。夙，早。浚，須也。卿大夫稱家。言能日日布行三德，早夜思之，須明行之，可以為卿大夫。浚，息俊反，馬云：“大也。”}日嚴祗敬六德，亮採有邦。\footnote{有國，諸侯。日日嚴敬其身，敬行六德，以信治政事，則可以為諸侯。嚴如字，馬、徐魚檢反。}


{\noindent\zhuan\zihao{6}\fzbyks 傳“三德”至“大夫”。正義曰：此文承“九德”之下，故知“三德”是九德之內課有其三也。\CJKunderwave{周語}云：“宣佈哲人之令德。”“宣”以布義，故為布也。“夙，早”，\CJKunderwave{釋詁}文。又云：“須，待也。”此經之意，謂夜思之，明旦行之，“須”為待之意，故“浚”為須也。大夫受采邑，賜氏族，立宗廟，世不絕祀，故稱“家”。位不虛受，非賢臣不可，言能日日布行三德，早夜思之，待明行之，如此念德不懈怠者，乃可以為大夫也。以士卑，故言不及也。計有一德二德，即可以為士也。鄭以“‘三德’、‘六德’，皆‘亂而敬’以下之文”，經無此意也。 \par}

{\noindent\zhuan\zihao{6}\fzbyks 傳“有國”至“諸侯”。正義曰：天子分地建國，諸侯專為己有，故“有國”謂諸侯也。“祗”亦為敬，“敬”有二文，上謂敬身，下謂敬德，“嚴”則敬之狀也。故言“日日嚴敬其身，敬行六德，以信治政事，則可以為諸侯”也。諸侯大夫皆言“日日”者，言人之行德，不可暫時舍也。臣當行君之令，故早夜思之。君是出令者，故言敬身行德。此文以小至大，總以天子之事，故先大夫而後諸侯。 \par}

翕受敷施，九德咸事，俊乂在官。\footnote{翕，和也。能合受三六之德而用之，以佈施政教，使九德之人皆用事。謂天子如此，則俊德治能之士並在官。翕，許及反。俊乂，馬曰：“千人曰俊，百人曰乂。”}百僚師師，百工惟時，\footnote{僚、工皆官也。師師，相師法。百官皆是，言政無非。僚本又作寮。}撫於五辰，庶績其凝。”\footnote{凝，成也。言百官皆撫順五行之時，眾功皆成。撫，方武反。凝，魚陵反，馬云：“定也。”}

{\noindent\zhuan\zihao{6}\fzbyks 傳“翕合”至“在官”。正義曰：“翕，合”，\CJKunderwave{釋詁}文。以文承“三德”、“六德”之下，故言合受三六之德而用之。以此人為官，令其佈施政教，使此九德之人皆居官用事。謂天子各任其所能。大夫所行三德，或在諸侯六德之內，但並此三六之德，既充九數，故言九德皆用事，謂用為大夫，用為諸侯，使之治民事也。大夫諸侯當身自行之，故言“日宣”、“日嚴”。天子當任人使行之，故言“合受而用之”。其實天子亦備九德，故能任用三德、六德也，則俊德治能之士並在官矣。“乂”訓為治，故云“治能”。馬、王、鄭皆云：“才德過千人為俊，百人為乂。” \par}

{\noindent\zhuan\zihao{6}\fzbyks 傳“僚工”至“無非”。正義曰：“僚，官”，\CJKunderwave{釋詁}文。“工，官”，常訓也。“師師”謂相師法也。 \par}

{\noindent\zhuan\zihao{6}\fzbyks 傳“凝成”至“皆成”。正義曰：\CJKunderline{鄭玄}亦云：“凝，成也。”王肅云：“凝猶定也。”皆以意訓耳。文承“百工”之下，“撫於五辰”還是百工撫之，故云“百官皆撫順五行之時，則眾功皆成”也。“五行之時”即四時也。\CJKunderwave{禮運}曰“播五行於四時”，土寄王四季，故為“五行之時”也。所撫順者，即\CJKunderwave{堯典}“敬授民時”，“平秩東作”之類是也。 \par}

{\noindent\shu\zihao{5}\fzkt “日宣”至“其凝”。正義曰：\CJKunderline{皋陶}既陳人有九德，宜擇而官之，此又言官之所宜：“若人能日日宣佈三德，早夜思念而須明行之，此人可以為卿大夫，使有家也。若日日嚴敬其身，又能敬行六德，信能治理其事，此人可以為諸侯，使有國也。然後撫以天子之任，合受有家有國三六之德而用之，佈施政教，使九德之人皆得用事，事各盡其能，無所遺棄，則天下俊德治能之士並在官矣。皆隨賢才任職，百官各師其師，轉相教誨，則百官惟皆是矣，無有非者。以此撫順五行之時，以化隘天下之民，則眾功其皆成矣。”結上知人安民之意。 \par}

無教逸欲，有邦\footnote{不為逸豫貪慾之教，是有國者之常。}兢兢業業，一日二日萬幾。\footnote{兢兢，戒慎。業業，危懼。幾，微也。言當戒懼萬事之微。兢，居凌反。業如字，徐五荅反。幾，徐音機。}無曠庶官，天工人其代之。\footnote{曠，空也。位非其人為空官。言人代天理官，不可以天官私非其才。}


{\noindent\zhuan\zihao{6}\fzbyks 傳“不為”至“之常”。正義曰：“毋”者禁戒之辭，人君身為逸欲，下則效之，是以禁人君使不自為耳。不為逸豫貪慾之教，是有國者之常也。此文主於天子，天子謂天下為國。\CJKunderwave{詩}雲“生此王國”之類是也。 \par}

{\noindent\zhuan\zihao{6}\fzbyks 傳“兢兢”至“之微”。正義曰：\CJKunderwave{釋訓}云：“兢兢，戒也。業業，危也。”戒必慎，危必懼，傳言“慎”、“懼”以足之。\CJKunderwave{易·繫辭}云：“幾者動之微。”故“幾”為微也。一日二日之間，微者乃有萬事，言當戒慎萬事之微。微者尚有萬，則大事必多矣。且微者難察，察則勞神，以言不可逸耳。馬、王皆云：“一日、二日，猶日日也。” \par}

{\noindent\zhuan\zihao{6}\fzbyks 傳“曠空”至“其才”。正義曰：“曠”之為空，常訓也。位非其人，所職不治，是為空官。天不自治，立君乃治之。君不獨治，為臣以佐之。下典、禮、德、刑，無非天意者。天意既然，人君當順天，是言人當代天治官。官則天之官，居天之官，代天為治，苟非其人，不堪此任，人不可以天之官而私非其才。王肅云：“天不自下治之，故人代天居之，不可不得其人也。” \par}

天敘有典，敕我五典五惇哉！\footnote{天次敘人之常性，各有分義，當敕正我五常之敘，使合於五厚，厚天下。有典，馬本作\CJKunderwave{五典}。分，扶問反。}天秩有禮，自我五禮有庸哉！\footnote{庸，常。自，用也。天次秩有禮，當用我公、侯、伯、子、男五等之禮以接之，使有常。有庸，馬本作“五庸”。}同寅協恭，和衷哉！\footnote{衷，善也。以五禮正諸侯，使同敬合恭而和善。衷音中。}


{\noindent\zhuan\zihao{6}\fzbyks 傳“天次”至“天下”。正義曰：天敘有典，有此五典，即父義、母慈、兄友、弟恭、子孝是也。五者人之常性，自然而有，但人性有多少耳。天次敘人之常性，使之各有分義。義，宜也。今此義、慈、友、恭、孝各有定分，合於事宜。此皆出天然，是為天次敘之。天意既然，人君當順天之意，敕正我五常之教，使合於五者皆厚，以教天下之民也。五常之教，人君為之，故言“我”也。五教遍於海內，故以“天下”言之。 \par}

{\noindent\zhuan\zihao{6}\fzbyks 傳“庸常”至“有常”。正義曰：“庸，常”，\CJKunderwave{釋詁}文。又云：“由，自也。”“由”是用，故“自”為用也。“天次敘有禮”，謂使賤事貴，卑承尊，是天道使之然也。天意既然，人君當順天意，用我公、侯、伯、子、男五等之禮以接之，使之貴賤有常也。此文主於天子，天子至於諸侯，車旗衣服、國家禮儀、饗食燕好、饔餼飧牢，禮各有次秩以接之。上言“天敘”，此雲“天秩”者，“敘”謂定其倫次，“秩”謂制其差等，義亦相通。上雲“敕我”，此言“自我”者，五等以教下民，須敕戒之；五禮以接諸侯，當用我意；故文不同也。上言“五惇”，此言“五庸”者，五典施於近親，欲其恩厚；五禮施於臣下，欲其有常；故文異也。王肅云：“五禮謂王、公、卿、大夫、士。”\CJKunderline{鄭玄}云：“五禮，天子也，諸侯也，卿大夫也，士也，庶民也。”此無文可據，各以意說耳。 \par}

{\noindent\zhuan\zihao{6}\fzbyks 傳“衷善”至“和善”。正義曰：“衷”之為善，常訓也。故\CJKunderwave{左傳}雲“天誘其衷”，說者皆以衷為善。此文承“五禮”之下，禮尚恭敬，故“以五禮正諸侯，使同敬合恭而和善”也。\CJKunderline{鄭玄}以為“並上之禮共有此事”。五典室家之內，務在相親，非復言以恭敬，恭敬惟為五禮而已，孔言是也。 \par}

天命有德，五服五章哉！\footnote{五服，天子、諸侯、卿、大夫、士之服也。尊卑彩章各異，所以命有德。}天討有罪，五刑五用哉！\footnote{言天以五刑討五罪，用五刑宜必當。}政事懋哉！懋哉！\footnote{言敘典秩禮，命德討罰無非天意者，故人君居天官，聽政治事，不可以不自勉。}

{\noindent\zhuan\zihao{6}\fzbyks 傳“五服”至“有德”。正義曰：\CJKunderwave{益稷}云：“以五采彰施於五色，作服，汝明。”是天子、諸侯、卿、大夫、士之服也。其“尊卑採章各異”，於彼傳具之。天命有德，使之居位，命有貴賤之倫，位有上下之異，不得不立名，以此等之，象物以彰之。先王制為五服，所以表貴賤也。服有等差，所以別尊卑也。 \par}

{\noindent\shu\zihao{5}\fzkt “無教”至“懋哉”。正義曰：\CJKunderline{皋陶}既言用人之法，又戒以居官之事：“上之所為，下必效之。無教在下為逸豫貪慾之事，是有國之常道也。為人君當兢兢然戒慎，業業然危懼。”言當戒慎。“一日二日之間而有萬種幾微之事，皆須親自知之，不得自為逸豫也。萬幾事多，不可獨治，當立官以佐己，無得空廢眾官，使才非其任。此官乃是天官，人其代天治之，不可以天之官而用非其人”。又言：“典禮德刑皆從天出，天次敘人倫，使有常性，故人君為政，當敕正我父、母、兄、弟、子五常之教教之，使五者皆惇厚哉！天又次敘爵命，使有禮法，故人君為政，當奉用我公、侯、伯、子、男五等之禮接之，使五者皆有常哉！接以常禮，當使同敬合恭而和善哉！天又命用有九德，使之居官，當承天意為五等之服，使五者尊卑彰明哉！天又討治有罪，使之絕惡，當承天意為五等之刑，使五者輕重用法哉！典禮德刑，無非天意，人君居天官，聽治政事，當須勉之哉！” \par}

天聰明，自我民聰明。\footnote{言天因民而降之福，民所歸者天命之。天視聽人君之行，用民為聰明。}天明畏，自我民明威。\footnote{天明可畏，亦用民成其威。民所叛者天討之，是天明可畏之效。畏如字，徐音威，馬本作威。}達於上下，敬哉有土！”\footnote{言天所賞罰，惟善惡所在，不避貴賤。有土之君，不可不敬懼。}\CJKunderline{皋陶}曰：“朕言惠，可厎行。”\footnote{其所陳“九德”以下之言，順於古道，可致行。}\CJKunderline{禹}曰：“俞，乃言厎可績。”\footnote{然其所陳，從而美之曰：“用汝言，致可以立功。”}\CJKunderline{皋陶}曰：“予未有知，思曰贊贊襄哉！”\footnote{言我未有所知，未能思致於善，徒亦贊奏上古行事而言之。因\CJKunderline{禹}美之，承以謙辭，言之序。知如字，徐音智。思如字，徐音息吏反。襄，息羊反，上也。馬云：“因也。”案\CJKunderwave{爾雅}作儴，因也，如羊反。}

{\noindent\zhuan\zihao{6}\fzbyks 傳“言天”至“聰明”。正義曰：皇天無心，以百姓之心為心。此經大意言民之所欲，天必從之。“聰明”謂聞見也，天之所聞見，用民之所聞見也。然則“聰明”直是見聞之義，其言未有善惡；以下言“明威”，是天降之禍，知此“聰明”是天降之福。此即\CJKunderwave{泰誓}所云“天聽自我民聽，天視自我民視”。故“民所歸者，天命之。”大而言之，民所歸就，天命之為天子也。小而言之，雖公卿大夫之任，亦為民所歸向，乃得居之。此文主於天子，故言“天視聽人君之行，用民為聰明”，戒天子使順民心，受天之福也。 \par}

{\noindent\zhuan\zihao{6}\fzbyks 傳“言天”至“敬懼”。正義曰：上句有賞罰，故言“天所賞罰，不避貴賤”。此之“達於上下”，言天子亦不免也。\CJKunderwave{喪服}\CJKunderline{鄭玄}注云：“天子諸侯及卿大夫有地者皆曰君。”即此“有土”可兼大夫以上。但此文本意實主於天子，戒天子不可不敬懼也。 \par}

{\noindent\zhuan\zihao{6}\fzbyks 傳“言我”至“之序”。正義曰：\CJKunderline{皋陶}目言“可致行”，\CJKunderline{禹}言“致可績”，此承而為謙，知其目言未有所知，未能思致於善也。“思”字屬上句，王肅云：“贊贊猶贊奏也。”顯氏云：“襄，上也。謂贊奏上古行事而言之也。”經雲“曰”者，謂我上之所言也。傳不訓“襄”為上，已從“襄陵”而釋之。故二劉並以“襄”為因，若必為因，孔傳無容不訓其意。言進習上古行事，因贊成其辭而言之也。傳雖不訓“襄”字，其義當如王說。\CJKunderline{皋陶}慮忽之，自雲“言順可行”。因\CJKunderline{禹}美之，即承謙辭。一揚一抑，言之次序也。\CJKunderline{鄭玄}云：“贊，明也。襄之言暢，言我未有所知，所思徒贊明帝德，暢我忠言而已。謙也。” \par}

{\noindent\shu\zihao{5}\fzkt “天聰”至“襄哉。正義曰：此承上“懋哉”以下言所勉之者。以天之聰明視聽，觀人有德。用我民以為耳目之聰明，察人言善者，天意歸賞之。又天之明德可畏，天威者，用我民言惡而叛之，因討而伐之，成其明威。天所賞罰，達於上下，不避貴賤，故須敬哉，有土之君！\CJKunderline{皋陶}既陳此戒，欲其言入之，故曰：“我之此言，順於古道，可致行，不可忽也。”\CJKunderline{禹}即受之曰：“然，汝言用而致可以立功。”重其言以深戒帝。\CJKunderline{皋陶}乃承之以謙曰：“我未有所知，未能思致於善，我所言曰，徒贊奏上古所行而言之哉！非己知思而所自能。”是其謙也。 \par}

%%% Local Variables:
%%% mode: latex
%%% TeX-engine: xetex
%%% TeX-master: "../Main"
%%% End:

%% -*- coding: utf-8 -*-
%% Time-stamp: <Chen Wang: 2024-04-02 11:42:41>

% {\noindent \zhu \zihao{5} \fzbyks } -> 注 (△ ○)
% {\noindent \shu \zihao{5} \fzkt } -> 疏

\chapter{卷五}


\section{益稷第五（皋陶謨中、下）}


益稷\footnote{禹稱其人，因以名篇。}

{\noindent\zhuan\zihao{6}\fzbyks 傳“禹稱”至“名篇”。正義曰：禹言“暨益”、“暨稷”，是禹稱其二人，二人佐禹有功，因以此二人名篇。既美大禹，亦所以彰此二人之功也。禹先言“暨益”，故“益”在“稷”上。馬、鄭、王所據\CJKunderwave{書序}此篇名為\CJKunderwave{棄稷}。“棄”“稷”一人，不宜言名又言官，是彼誤耳。又合此篇於\CJKunderwave{皋陶謀}，謂其別有\CJKunderwave{棄稷}之篇，皆由不見古文，妄為說耳。 \par}

帝曰：“來，禹，汝亦昌言。”\footnote{因皋陶謨九德，故呼禹使亦陳當言。當，丁浪反，本亦作讜，當湯反。李登\CJKunderwave{聲類}云：“讜言，善言也。”}禹拜曰：“都！帝，予何言？予思日孜孜。”\footnote{拜而嘆，辭不言，欲使帝重皋陶所陳。言己思日孜孜不怠，奉承臣功而已。思，徐如字，又息吏反。孜音茲。}皋陶曰：“籲！如何？”\footnote{問所以孜孜之事。}禹曰：“洪水滔天，浩浩懷山襄陵，下民昏墊。\footnote{言天下民昏瞀墊溺，皆困水災。○浩，戶老反。墊，丁念反。瞀音務，一音茂，本或作務。溺，乃歷反。}


{\noindent\zhuan\zihao{6}\fzbyks 傳“因皋”至“當言”。正義曰：上篇皋陶謀九德，此帝呼禹，令亦陳當言。“亦”者，亦皋陶也。明上篇皋陶雖與益相應，其言亦對帝也。上傳雲“皋陶為帝舜謀”者，以此而知也。 \par}

{\noindent\zhuan\zihao{6}\fzbyks 傳“拜而”至“而已”。正義曰：既已拜而嘆，必有所美，復辭而不言，是知欲使帝重皋陶所陳，言己無以加也。王肅雲“帝在上，皋陶陳謀於下，已備矣，我復何所言乎”是也。既無所言，故言已思惟日孜孜不敢怠惰，奉成臣職而已。孜孜者，勉功不怠之意。 \par}

{\noindent\zhuan\zihao{6}\fzbyks 傳“言天下”至“水災”。正義曰：“瞀”者眩惑之意，故言“昏瞀”。“墊”是下溼之名，故為溺也。言天下之人遭此大水，精神昏瞀迷惑，無有所知，又苦沉溺，皆困此水災也。鄭云：“昏，沒也。墊，陷也。禹言洪水之時，人有沒陷之害。” \par}

予乘四載，隨山刊木，\footnote{所載者四，謂水乘舟，陸乘車，泥乘輴，山乘樏。隨行九州之山林，刊槎其木，開通道路以治水也。○乘音繩。刊，苦安反。輴，醜倫反。\CJKunderwave{漢書}作橇，如淳音蕝，以板置泥上。服虔云：“木橇，形如木箕，擿行泥上。”\CJKunderwave{屍子}云：“澤行乘蕝。”蕝音子絕反，樏，力追反。\CJKunderwave{史記}作橋，徐音丘遙反。\CJKunderwave{漢書}作梮，九足反。行，下孟反。槎，士雅反，下同。\CJKunderwave{說文}云：“袤斫。”又莊下反。}暨益奏庶鮮食。\footnote{奏謂進於民。鳥獸新殺曰鮮。與益槎木，獲鳥獸，民以進食。○塈，其器反。鮮，徐音仙，馬云：“鮮，生也。”}予決九川，距四海，濬畎澮距川。\footnote{距，至也。決九州名川，通之至海。一畎之間，廣尺、深尺曰畝。方百里之間，廣二尋、深二仞曰澮。澮畎深之至川，亦入海。○畎，工犬反。澮，故外反。廣，光浪反。深，屍鴆反。}


{\noindent\zhuan\zihao{6}\fzbyks 傳“所載”至“治水”。正義曰：\CJKunderwave{史記·河渠書}云：“\CJKunderwave{夏書}曰：‘禹湮洪水十三年，三過家不入門。陸行載車，水行載舟，泥行蹈橇。橇音蕝，山行即橋。橋，丘遙反。’”徐廣曰：“橋一作輂。輂，幾玉反，直轅車也。”\CJKunderwave{屍子}云：“山行乘樏，泥行乘蕝。蕝，子絕反。”\CJKunderwave{漢書·溝洫志}云：“泥行乘毳，山行則梮。梮，居足反。”毳行如箕，擿行泥上。如淳云：“毳謂以板置泥上，以通行路也。”\CJKunderwave{慎子}云：“為毳者，患塗之泥也。”應劭云：“梮或作樏，為人所牽引也。”如淳云：“梮謂以鐵如錐，頭長半寸，施之履下以上山，不蹉跌也。”韋昭云：“梮，木器也。如今輿床，人輿以行也。”此經惟言“四載”，傳言所載者四，同彼\CJKunderwave{史記}之說。古書屍子、慎子之徒有此言也。輴與毳為一，樏與梮、輿為一。古篆變形，字體改易，說者不同，未知孰是。禹之施功，本為治水，此經乃雲“隨山刊木”，刊木為治水，治水遍於九州，故云“隨行九州之山林”。襄二十五年\CJKunderwave{左傳}云：“井堙木刊。”“刊”是除木之義也。\CJKunderwave{毛傳}云：“除木曰槎。”故曰“刊槎其木，開通道路以治水”。\par}

{\noindent\zhuan\zihao{6}\fzbyks 傳“奏謂”至“進食”。正義曰：黎民阻飢，為人治水，故知“奏謂進食於人”也。\CJKunderwave{禮}有鮮魚臘，以其新殺鮮淨，故名為鮮，是鳥獸新殺曰鮮，魚鱉新殺亦曰鮮也。此承“山”下，故為鳥獸，下承水後，故為魚鱉，其新殺之意同也。既言刊木乃進鮮食，食是除木所得，故言“與益槎木，獲鳥獸，人以進食”。 \par}

{\noindent\zhuan\zihao{6}\fzbyks 傳“距至”至“入海”。正義曰：“距”者相抵之名，故為至也。非是名川不能至海，故“決九州之名川通之至海”也。\CJKunderwave{考工記}云：“匠人為溝洫，耜廣五寸，二耜為耦。一耦之伐，廣尺、深尺謂之畎。田首倍之，廣二尺、深二尺謂之遂。九夫為井，井間廣四尺、深四尺謂之溝。方十里為成，成間廣八尺、深八尺謂之洫。方百里為同，同間廣二尋、深二仞謂之澮。”是畎、遂、溝、洫、澮皆通水之道也。以小注大，故從畎、遂、溝、洫乃以入澮，澮入於川，川入於海，是畎內之水亦入海也。惟言“畎、澮”，舉大小而略其餘也。先言決川至海，後言濬畎至川者，川既入海，然後澮得入川，故先言川也。 \par}

暨稷播，奏庶艱食鮮食。\footnote{艱，難也。眾難得食處，則與稷教民播種之，決川有魚鱉，使民鮮食之。○艱，工間反，馬本作根，云：“根生之食，謂百穀。”處，昌慮反。鱉，必滅反。}懋遷有無化居。\footnote{化，易也。居謂所宜居積者。勉勸天下，徙有之無，魚鹽徙山，林木徙川澤，交易其所居積。懋音茂。鹽，餘廉反。}烝民乃粒，萬邦作乂。”\footnote{米食曰粒。言天下由此為治本。○烝，之承反。粒音立。治，直吏反，下同。}皋陶曰：“俞！師汝昌言。”\footnote{言禹功甚當，可師法。○當，丁浪反。}

{\noindent\zhuan\zihao{6}\fzbyks 傳“艱難”至“鮮食之”。正義曰：“艱，難也”，\CJKunderwave{釋詁}文。禹主治水，稷主教播種，水害漸除，則有可耕之地，難得食處，先須教導以救之，故云“眾難得食處，則與稷教人播種之”。易得食處，人必自能得之，意在救人艱危之厄，故舉難得食處以言之。於時雖漸播種，得谷猶少，人食未足，故決川有魚鱉，使人鮮食之。言食魚以助谷也。鄭玄云：“與稷教人種澤物菜蔬艱厄之食。”稷功在於種穀，不主種菜蔬也。言后稷種菜蔬艱厄之食，傳記未有此言也。 \par}

{\noindent\zhuan\zihao{6}\fzbyks 傳“化易”至“居積”。正義曰：變化是改易之義，故“化”為易也。“居謂所宜居積者”，近水者居魚鹽，近山者居林木也。“勉勸天下徙有之無”者，謂徙我所有，往彼無鄉；取彼所有，以濟我之所無。“魚鹽徙山，林木徙川澤，交易其所宜居積”，言此“遷”者，謂將物去，不得空取彼物也。王肅云：“易居者不得空去，當滿而去，當滿而來也。” \par}

{\noindent\zhuan\zihao{6}\fzbyks 傳“米食”至“治本”。正義曰：\CJKunderwave{說文}云：“粒，糂也。”今人謂飯為米糂，遺餘之飯謂之一粒、兩粒，是米食曰粒，言是用米為食之名也。人非谷不生，政由谷而就，言天下由此谷為治政之本也。君子之道以謙虛為德，禹盛言己功者，為臣之法當孜孜不怠，自言己之勤苦，所以勉勸人臣，非自伐也。 \par}

{\noindent\shu\zihao{5}\fzkt “帝曰來”至“汝昌言”。正義曰：皋陶既為帝謀，帝又呼禹進之，曰：“來，禹，汝亦宜陳其當言。”禹拜曰：“嗚乎！帝，皋陶之言既已美矣，我更何所言？我之所思者，每日孜孜勤於臣職而已。”皋陶怪禹不言，故謂之曰：“籲！”問其所以孜孜之事如何。禹曰：“往者洪水漫天，浩浩然盛大，包山上陵，下民昏惑沉溺，皆困水災。我乘舟車輴樏等四種之載，隨其所往之山，槎木通道而治之。與益所進於人者，惟有槎木所獲眾鳥獸鮮肉為食也。我又通決九州名川，通之至於四海。深其畎澮，以至於川，水漸除矣。與稷播種五穀，進於眾人難得食處，乃決水所得魚鱉鮮肉為食也，人既皆得食矣。又勸勉天下徙有之無，交易其所居積。於是天下眾人乃皆得米粒之食，萬國由此為治理之政。我所言孜孜者在此也。”皋陶曰：“然。可以為師法者，是汝之當言。” \par}

禹曰：“都！帝，慎乃在位。”帝曰：“俞。”\footnote{然禹言，受其戒。}禹曰：“安汝止，惟幾惟康，其弼直，\footnote{言慎在位，當先安好惡所止，念慮幾微，以保其安，其輔臣必用直人。好惡，上呼報反，下烏路反，又並如字。}惟動丕應徯志。\footnote{徯，待也。帝先安所止，動則天下大應之，順命以待帝志。應，應對之應。徯，胡啟反。}


{\noindent\zhuan\zihao{6}\fzbyks 傳“言慎”至“直人”。正義曰：此禹重戒帝，覆上“慎乃在位”。“當先安好惡所止”，謂心之所止，當止好不止惡，言惡以刑好也。\CJKunderwave{大學}云：“為人君止於仁，為人臣止於敬。”“好惡所止”謂此類也。傳意以上“惟”為念，下“惟”為辭，故云念慮幾微，然後以保其好惡所安寧耳。 \par}

{\noindent\zhuan\zihao{6}\fzbyks 傳“徯待”至“帝志”。正義曰：“徯，待”，\CJKunderwave{釋詁}文。帝先能自安所止，心之所止，止於好事，其有舉動，發號出令，則天下大應之，順命以待帝志。謂靜以待命，有命則從也。 \par}

以昭受上帝，天其申命用休。”\footnote{昭，明也。非但人應之，又乃明受天之報施，天又重命用美。施，始豉反。重，直用反。}帝曰：“籲！臣哉鄰哉！鄰哉臣哉！”禹曰：“俞。”\footnote{鄰，近也。言君臣道近，相須而成。}

{\noindent\zhuan\zihao{6}\fzbyks 傳“昭明”至“用美”。正義曰：\CJKunderwave{堯典}已訓“昭”為明，此重訓，詳之。皇天無親，惟德是輔，人之所欲，天必從之。帝若能安所止，非但人歸之，又乃明受天之報施。天下太平，祚胤長遠，是天之報施也。“天又重命用美”，謂四時和祥瑞臻之類也。或當前後非一，故傳言“又”也。 \par}

{\noindent\zhuan\zihao{6}\fzbyks 傳“鄰近”至“而成”。正義曰：\CJKunderwave{周禮}“五家為鄰”，取相近之義，故“鄰”為近也。禹言君當好善，帝言須得臣力，再言鄰哉，言君臣之道當相須而成，鄭玄云：“臣哉，汝當為我鄰哉！鄰哉，汝當為我臣哉！反覆言此，欲其志心入禹。” \par}

{\noindent\shu\zihao{5}\fzkt “禹曰都”至“曰俞”。正義曰：禹以皋陶然已，因嘆而戒帝曰：“嗚呼！帝當謹慎汝所在之位。”帝受其戒，曰：“然。”禹又戒帝曰：“若欲慎汝在位，當須先安定汝心好惡所止，念慮事之微細，以保安其身，其輔弼之臣必用正直之人。若能如此，惟帝所動，則天下大應之，以待帝志。以明受天之佈施，於天其重命帝用美道也。”帝以禹言已重，乃驚而言曰：“籲！臣哉近哉，臣當親近君也！近哉臣哉，君當親近臣也！”言君臣當相親近，共與成政道也。禹應帝曰：“然。”言君臣宜相親近也。 \par}

帝曰：“臣作朕股肱耳目。\footnote{言大體若身。股音古。肱，古弘反。}予欲左右有民，汝翼。\footnote{左右，助也。助我所有之民，富而教之，汝翼成我。}予欲宣力四方，汝為。\footnote{布力立治之功，汝群臣當為之。}予欲觀古人之象，\footnote{欲觀示法象之服制。觀，舊音官，又官喚反。}


{\noindent\zhuan\zihao{6}\fzbyks 傳“言大體若身”。正義曰：君為元首，臣為股肱耳目，大體如一身也。足行手取，耳聽目視，身雖百體，四者為大，故舉以為言。鄭玄云：“動作視聽皆由臣也。” \par}

{\noindent\zhuan\zihao{6}\fzbyks 傳“左右”至“成我”。正義曰：\CJKunderwave{釋詁}云：“左、右、助，慮也”，同訓為慮，是“左右”得為助也。立君所以牧人，人之自營生產，人君當助救之。\CJKunderwave{論語}稱\CJKunderline{孔子}適衛，欲先富民而後教之，故云“助我所有之民，欲富而教之”也。君子施教，本為養人，故先雲助人，舉其重者。以其為人事重，當須翼成，故言“汝翼”。次顯君施教化，須臣為之，故言“汝為”。次明衣服上下，標顯尊卑，故云“汝明”。次雲六律、五聲，故云“汝聽”。各隨事立文，其實不異。 \par}

{\noindent\zhuan\zihao{6}\fzbyks 傳“布力”至“為之”。正義曰：\CJKunderwave{詩}雲“四方於宣”，\CJKunderwave{論語}雲“陳力就列”，是佈政用力，故言“布力立治之功，汝群臣當為之”。 \par}

{\noindent\zhuan\zihao{6}\fzbyks 傳“欲觀”至“服制”。正義曰：“觀示法象之服制”者，謂欲申明古人法象之衣服，垂示在下使觀之也。\CJKunderwave{易·繫辭}云：“黃帝、堯、舜垂衣裳而天下治。”象物制服，蓋因黃帝以還，未知何代而具彩章。舜言己欲觀古，知在舜之前耳。 \par}

日、月、星辰、山、龍、華蟲，\footnote{日月星為三辰。華象草華蟲雉也。畫三辰、山、龍、華蟲於衣服旌旗。蟲，直弓反。}作會，宗彝，\footnote{會，五采也，以五采成此畫焉。宗廟彝樽亦以山、龍、華蟲為飾。會，胡對反，馬、鄭作繪。彝音夷，馬同，鄭雲“宗彝，虎也。”}


{\noindent\zhuan\zihao{6}\fzbyks 傳“日月”至“旌旗”。正義曰：桓二年\CJKunderwave{左傳}云：“三辰旗旗，昭其明也。”三辰謂此日月星也，故“日月星為三辰”。辰即時也，三者皆是示人時節，故並稱辰焉。傳言此者，以“辰”在“星”下，總上三事為辰，辰非別為物也。\CJKunderwave{周禮·大宗伯}云：“實柴祀日月星辰。”鄭玄云：“星謂五緯也，辰謂日月所會十二次也。”“星”、“辰”異者，彼鄭以遍祭天之諸神十二次也，次亦當祭之，故令“辰”與“星”別。此雲畫之於衣，日月合宿之辰，非有形容可畫，且\CJKunderwave{左傳}雲三辰即日月星也。\CJKunderwave{周禮}“司常掌九旗之物”，惟日月為常，不言畫星，蓋太常之上又畫星也。\CJKunderwave{穆天子傳}稱天子葬盛姬，畫日月七星，蓋畫北斗也。草木雖皆有華，而草華為美，故云“華象草華蟲雉”也。\CJKunderwave{周禮·司服}有“鷩冕”，鷩則雉焉，雉五色，象草華也。\CJKunderwave{月令}五時皆雲其蟲，“蟲”是鳥獸之總名也。下雲“作服，汝明”，知“畫三辰、山、龍、華蟲於衣服”也。又言“旌旗”者，\CJKunderwave{左傳}言“三辰旗旗”，\CJKunderwave{周禮·司常}雲“日月為常”，王者禮有沿革，後因於前，故知舜時三辰亦畫之於旌旗也。下傳雲“天子服日月而下”，則三辰畫之於衣服，又畫於旌旗也。\CJKunderwave{周禮·司服}云：“享先王則袞冕。”袞者，卷也，言龍首卷然。以袞為名，則所畫自龍已下，無日月星也。\CJKunderwave{郊特牲}云：“祭之日，王被袞冕以象天也。”又曰：“龍章而設日月，以象天也。”鄭玄云：“謂有日月星辰之章”，“設日月畫於衣服旌旗也”。據此記文，袞冕之服亦畫日月。鄭注\CJKunderwave{禮記}言\CJKunderwave{郊特牲}所云“謂魯禮也”。要其文稱王被服袞冕，非魯事也。或當二代天子衣上亦畫三辰，自龍章為首，而使袞統名耳。\CJKunderwave{禮}文殘缺，不可得詳，但如孔解，舜時天子之衣畫日月耳。鄭玄亦以為然。王肅以為“舜時三辰即畫於旌旗，不在衣也，天子山、龍、華蟲耳。 \par}

{\noindent\zhuan\zihao{6}\fzbyks 傳“會五”至“為飾”。正義曰：“會”者合聚之名，下雲“以五彩彰施於五色，作服”，知“會”謂五色也。\CJKunderwave{禮}衣畫而裳繡，“五色備謂之繡”，知畫亦備五色，故云“以五彩成此畫焉”，謂畫之於衣、宗彝。文承“作會”之下，故云“宗廟、彝樽亦以山、龍、華蟲為飾”。知不以日月星為飾者，孔以三辰之尊不宜施於器物也。\CJKunderwave{周禮}有山罍、龍勺、雞彝、鳥彝，以類言之，知彝樽以山、龍、華蟲為飾，亦畫之以為飾也。\CJKunderwave{周禮}彝器所云犧、象、雞、鳥者，鄭玄皆為畫飾，與孔意同也。\CJKunderwave{周禮}彝器無山、龍、華蟲為飾者，帝王革易，所尚不同，故有異也。 \par}

藻、火、粉、米、黼、黻、絺、繡，\footnote{藻，水草有文者。火為火字，粉若粟冰，米若聚米，黼若斧形，黻為兩己相背，葛之精者曰絺，五色備曰繡。藻音早，本又作薻。粉米，\CJKunderwave{說文}作黺𪓋，徐米作䋛，音米。黼音甫，白與黑謂之黼。黻音弗，黑與青謂之黻。絺，徐敕私反，又敕其反，馬同，鄭陟裡反，刺也。繡音秀。背音佩。}以五采彰施於五色，作服，汝明。\footnote{天子服日月而下，諸侯自龍袞而下至黼黻，士服藻火，大夫加粉米。上得兼下，下不得僣上。以五采明施於五色，作尊卑之服，汝明制之。袞，工本反。僣，子念反。}


{\noindent\zhuan\zihao{6}\fzbyks 傳“藻水”至“曰繡”。正義曰：\CJKunderwave{詩}雲“魚在在藻”，是“藻”為水草。草類多矣，獨取此草者，謂此草有文故也。“火為火字”，謂刺繡為“火”字也。\CJKunderwave{考工記}云：“火以圜。”鄭司農云：“謂圜形似火也。鄭玄云：“形如半環。”然\CJKunderwave{記}是後人所作，何必能得其真？今之服章繡為“火”字者，如孔所說也。“粉若粟冰”者，粉之在粟，其狀如冰。“米若聚米”者，刺繡為文，類聚米形也。“黼若斧形”，\CJKunderwave{考工記}云：“白與黑謂之黼。”\CJKunderwave{釋器}云：“斧謂之黼。”孫炎云：“黼文如斧形”，蓋半白半黑，似斧刃白而身黑。“黻謂兩己相背”，謂刺繡為“己”字，兩“己”字相背也。\CJKunderwave{考工記}云：“黑與青謂之黻。”刺繡為兩\CJKunderwave{己}字，以青黑線繡也。\CJKunderwave{詩·葛覃}雲“為絺為綌”，是絺用葛也。\CJKunderwave{玉藻}云：“浴用二巾，上絺下綌。”\CJKunderwave{曲禮}云：“為天子削瓜者副之，巾以絺。為國君者華之，巾以綌。”皆以絺貴而綌賤，是絺精而綌粗，故“葛之精者曰絺”。“五色備謂之繡”，\CJKunderwave{考工記}文也。計此所陳，皆述祭服。祭服玄纁為之，後代無用絺者，蓋於時仍質，暑月染絺為纁而繡之以為祭服。孔以“華象草華蟲雉”，則合華蟲為一，\CJKunderwave{周禮}鄭玄注亦然，則以日、月、星辰、山、龍、華蟲六章畫於衣也。藻、火、粉、米、黼、黻六章繡於裳也。天之大數不過十二，故王者製作皆以十二象天也。顧氏取先儒等說，以為“日月星取其照臨，山取能興雲雨，龍取變化無方，華取文章，雉取耿介”。顧氏雖以唬骸藹蟲為二，其取象則同。又云：“藻取有文，火取炎上，粉取絜白，米取能養，黼取能斷，黻取善惡相背。”鄭玄云：“會讀為繪。宗彝謂宗廟之鬱鬯樽也。故虞夏以上，蓋取虎彝蜼彝而已。粉米，白米也。絺讀為黹。黹，紩也。自日月至黼黻凡十二章，天子以飾祭服。凡畫者為繪，刺者為繡。此繡與繪各有六，衣用繪，裳用繡。至周而變之，以三辰為旗旗，謂龍為袞，宗彝為毳，或損益上下，更其等差。”鄭意以“華蟲”為一，“粉米”為一，加“宗彝”謂虎蜼也。\CJKunderwave{周禮}宗廟彝器有虎彝、蜼彝，故以“宗彝”為虎蜼也。此經所云凡十二章，日也，月也，星也，山也，龍也，華蟲也，六者畫以作繪，施於衣也；宗彝也，藻也，火也，粉米也，黼也，黻也，此六者紩以為繡，施之於裳也。鄭玄雲“至周而變易之，損益上下，更其等差”，\CJKunderwave{周禮·司服}之注具引此文，乃云：“此古天子冕服十二章也。王者相變，至周而以日月星畫於旌旗。冕服九章，登龍于山，登火於宗彝，尊其神明也。九章，初一曰龍，次二曰山，次三曰華蟲，次四曰火，次五曰宗彝，皆畫以為繢；次六曰藻，次七曰粉米，次八曰黼，次九曰黻，以絺為繡。則袞之衣五章，裳四章，凡九也。鷩畫以雉，謂華蟲也。其衣三章，裳四章，凡七也。毳畫虎蜼，謂宗彝也。其衣三章，裳二章，凡五也。”是鄭以冕服之名皆取章首為義，袞冕九章，以龍為首，龍首卷然，故以袞為名。鷩冕七章，華蟲為首，華蟲即鷩雉也。毳冕五章，虎蜼為首，虎蜼毛淺，毳是亂毛，故以毳為名。如鄭此解，配文甚便，於絺繡之義，總為消帖。但解“宗彝”為虎蜼，取理太回，未知所說誰得經旨。 \par}

{\noindent\zhuan\zihao{6}\fzbyks 傳“天子”至“制之”。正義曰：此言“作服，汝明”，故傳辨其等差。天子服日月而下十二章，諸侯自龍袞而下至黼黻八章，再言“而下”，明天子諸侯皆至黼黻也。士服藻火二章，大夫加粉米四章。孔註上篇“五服”，謂“天子、諸侯、卿、大夫、士”，則卿與大夫不同，當加之以黼黻為六章。孔略而不言，孔意蓋以\CJKunderwave{周禮}制諸侯有三等之服，此諸侯同八章者，上古樸質，諸侯俱南面之尊，故合三為一等。且\CJKunderwave{禮}諸侯多同為一等，故\CJKunderwave{雜記}雲“天子九虞，諸侯七虞”，\CJKunderwave{左傳}雲“天子七月而葬，諸侯五月而葬”，是也。孔以此經上句“日、月、星辰、山、龍、華蟲”尊者在上，下句“藻、火、粉、米、黼、黻”尊者在下，黼黻尊於粉米，粉米尊於藻火，故從上以尊卑差之，士服藻火，大夫加以粉米，並藻火為四章。馬融不見孔傳，其注亦以為然，以古有此言，相傳為說也。蓋以衣在上為陽，陽統於上，故所尊在先。裳在下為陰，陰統於下，故所重在後。\CJKunderwave{詩}稱“玄袞及黼”\CJKunderwave{顧命}雲“麻冕黼裳”，當以黼為裳，故首舉黼以言其事如孔說也。天子諸侯下至黼黻，大夫粉米兼服藻火，是“上得兼下”也。士不得服粉米，大夫不得服黼黻，是“下不得僣上”也。訓“彰”為明，以五種之彩明施於五色，作尊卑之服，汝當分明制之，令其勿使僣濫也。鄭玄云：“性曰採，施曰色。”以本性施於繒帛，故云“以五采施於五色”也。鄭云：“作服者，此十二章為五服，天子備有焉，公自山龍而下，侯伯自華蟲而下，子男自藻火而下，卿大夫自粉米而下。”亦是以意說也。此雲“作服”，推據衣服，所以經有“宗彝”，及孔雲旌旗亦以山、龍、華蟲為飾者，但此雖以服為主，上既雲“古人之象”，則法象分在器物，皆悉明之，非止衣服而已。旌旗器物皆是彩飾，被服以明尊卑，故總雲“作服”以結之。 \par}

予欲聞六律、五聲、八音，在治忽，以出納五言，汝聽。\footnote{言欲以六律和聲音，在察天下治理及忽怠者，又以出納仁義禮智信五德之言，施於民以成化，汝當聽審之。出如字，又尺遂反，注同。納如字，又音內。}予違，汝弼。汝無面從，退有後言。\footnote{我違道，汝當以義輔正我。無得面從我違，而退後有言我不可弼。}欽四鄰。庶頑讒說，若不在時，\footnote{四近前後左右之臣，敕使敬其職。眾頑愚讒說之人，若所行不在於是而為非者，當察之。}


{\noindent\zhuan\zihao{6}\fzbyks 傳“言欲”至“審之”。正義曰：此經大意，令臣審聽樂音，察世之治否以報君也。金、石、絲、竹、匏、土、革、木，八物各出其音，謂之“八音”。八音之聲皆有清濁，聖人差之以為五品，宮、商、角、徵、羽，謂之“五聲”。五聲高下各有所準則，聖人制為六律，與五聲相均，作樂者以律均聲，聲從器出。帝言我欲以六律和彼五聲八音，以此樂之音聲，察世之治否。\CJKunderwave{詩序}云：“治世之音安以樂，其政和；亂世之音怨以怒，其政乖。”此則聽聲知政之道也。言今聽作樂，若其音安樂和平，則時政辨治而修理也；若其音怨怒乖離，則時政忽慢而怠惰也；是用樂之聲音察天下治理及忽怠者也。知其治理，則保以修之；知其忽怠，則改以修之；此治理忽怠，人君所願聞也。又樂之感人，使和易調暢，若樂音合度，則言必得理。以此樂音出納仁義禮智信五德之言，乃君之發言，合彼五德，施之於人，可以成其教化，是出五言也。人之五言，合彼五德，歸之於君，可以成諷諫，是納五言也。君言可以利民，民言可以益君，是言之善惡由樂音而知也。此言之善惡，亦人君之所原聞也。政之理忽，言之善惡，皆是上所願聞，欲令察知以告己，得守善而改惡，故帝令臣，汝當為我聽審之也。六律、六呂，當有十二，惟言“六律”者，鄭玄云：“舉陽，陰從可知也。”傳以“五言”為“五德之言”者，\CJKunderwave{漢書·律曆志}稱五聲播於五常，則角為仁，商為義，徵為禮，羽為智，宮為信，\CJKunderwave{志}之所稱必有舊說也。言五聲與五德相協，此論樂事而云“出納五言”，知是出納五德之言也。樂音和，則五德之言得其理；音不和，則五德之言違其度；故亦以樂音察五言也。帝之此言，自說臣之大法。於舜所聽，使聽韶樂也。襄二十九年\CJKunderwave{左傳}吳季札見舞韶樂而嘆曰：“德至矣哉，大矣！如天之無不幬也，如地之無不載也。”然則韶樂盡善盡美，有理無忽，而並言“忽”者，韶樂自美，取樂採人歌為曲，若其怠忽，則音辭亦有焉，故常使聽察之也。 \par}

{\noindent\zhuan\zihao{6}\fzbyks 傳“四近”至“察之”。正義曰：\CJKunderwave{冏命}云：“惟予一人無良，實賴左右前後有位之士，匡其不及。”知“四近”謂前後左右四者近君之臣，敕使敬其職也。更欲告以此下之辭，故敕之。眾頑愚讒說之人，若有所行不在於是而為非者，當察之。知其非，乃撻之書之。此與以下發端也。“庶頑讒說”謂朝廷之臣，“格則承之”乃謂天下之人。舜之朝廷當無讒說之人，故設為大法，戒慎之耳。四近之臣，普謂近君之臣耳，無常人也。鄭玄以“四近為左輔右弼，前疑後承”，惟伏生\CJKunderwave{書傳}有此言，\CJKunderwave{文王世子}雲“有師保，有疑承”，以外經傳無此官也。 \par}

侯以明之，撻以記之，\footnote{當行射侯之禮，以明善惡之教。笞撻不是者，使記識其過。撻，他末反，又他達反。笞，敕疑反。}書用識哉，欲並生哉！\footnote{書識其非，欲使改悔，與共並生。}工以納言，時而颺之，\footnote{工，樂官，掌誦詩以納諫，當是正其義而颺道之。颺音揚。}格則承之庸之，否則威之。”\footnote{天下人能至於道則承用之，任以官。不從教則以刑威之。否，方有反，徐音鄙。任，汝鴆反。}

{\noindent\zhuan\zihao{6}\fzbyks 傳“當行”至“其過”。正義曰：\CJKunderwave{禮}射皆張侯射之，知“侯以明之”，“當行射侯之禮，以明善惡之教”。射禮有序賓以賢，詢眾擇善之義，是可以明善惡也。“笞撻不是者，使記識其過”，謂過輕者也，大罪刑殺之矣。古之射侯之士，無以言之。案\CJKunderwave{周禮·司裘}云：“王大射則供虎侯、熊侯、豹侯、設其鵠。諸侯則供熊侯、豹侯，卿大夫則供麋侯，皆設其鵠。”鄭玄注云：“虎九十弓，即方一丈八尺。熊七十弓，方一丈四尺。豹、麋五十弓，方一丈。”鄭又引\CJKunderwave{梓人}“為侯，廣與崇方，三分其廣，而鵠居一焉”。則丈八之侯，鵠方六尺。丈四之侯，鵠方四尺六寸大半寸。一丈之侯，鵠方三尺三寸少半寸，此皆大射之侯也。\CJKunderwave{射人}云：“王以六耦射三侯，五正。諸侯以四耦射二侯，三正。孤卿大夫以三耦射一侯，二正。士以三耦射豹侯，二正。”鄭玄注云：“五正者，五采。中朱，次白，次蒼，次黃，玄居外。三正者，去玄、黃。二正者，去白、蒼而畫以朱、綠。”此賓射之侯也。鄭以賓射三侯步數高廣，與大射侯同，正大如鵠。\CJKunderwave{司裘}及\CJKunderwave{射人}所云諸侯者，謂圻內諸侯。若圻外諸侯，則\CJKunderwave{儀禮·大射}云，大侯九十弓，熊侯七十弓，豹侯五十弓，皆以三耦；其賓射則無文。若天子已下之燕射，案\CJKunderwave{鄉射記}云：“天子熊侯，白質。諸侯麋侯，赤質。大夫布侯，畫以虎豹。士布侯，畫以鹿豕。”熊侯已下同五十弓，即侯身高一丈，君臣共射之。 \par}

{\noindent\zhuan\zihao{6}\fzbyks 傳“書識”至“並生”。正義曰：書識其非，亦是小過者也，欲並生哉。總上三者，“侯以明之，撻以記之，書用識哉”，皆是欲其改悔，與無過之人共並生也。 \par}

{\noindent\zhuan\zihao{6}\fzbyks 傳“工樂”至“道之”。正義曰：\CJKunderwave{禮}通謂樂官為工，知\CJKunderwave{工}是樂官，則\CJKunderwave{周禮}大師、瞽蒙之類也。樂官掌頌詩言以納諫，以詩之義理或微，人君聽之，若有不悟，當正其義而揚道之。揚，舉也，舉而道向君也。 \par}

{\noindent\zhuan\zihao{6}\fzbyks 傳“天下”至“威之”。正義曰：言“承之用之”，則此人未在官也，故言謂天下民必也。能至於道即賢者，故承用之而任以官也。“否”謂不從教者，則以刑威之而罪其身也。臣過必小，故撻之書之；人罪或大，故以刑威之。 \par}

{\noindent\shu\zihao{5}\fzkt “帝曰臣”至“威之”。正義曰：帝以禹然己言，又說須臣之事：“作我股肱耳目。”言已動作視聽皆由臣也。“我欲助我所有之人，使之家給人足，汝當翼讚我也。我欲布陳智力於天下四方，為立治之功，汝等當與我為之。我欲觀示君臣上下以古人衣服之法象，其日、月、星辰、山、龍、華蟲作會，合五采而畫之。又畫山、龍、華蟲於宗廟彝樽。其藻、火、粉、米、黼、黻於絺葛而刺繡，以五種之彩明施於五色，製作衣服，汝當為我明其差等而制度之。我欲聞知六律，和五聲，播之於八音，以此音樂察其政治與忽怠者，其樂音又以出納五德之言，汝當為我聽審之。我有違道，汝當以義輔成我。汝無得知我違非而對面從我，退而後更有言，雲我不可輔也”。既言其須臣之力，乃總敕之：“敬其職事哉，汝在我前後左右四旁鄰近之臣也。其眾類頑愚讒說之人，若有所行不在於是而為非者，汝當察之以法，行射侯之禮，知其善惡以明別之。行有不是者，又撻其身以記之。書其過者以以識哉。所以撻之書之者，冀其改悔，欲與並生活哉！工樂之官以納諫言於上，當是正其義而顯揚之，使我自知得失也。”又總言御下之法：“天下之人有能至於道者，則當承受而進用之，當任以官也。不從教者，則以刑罰威之，當罪其身也。此等皆汝臣之所為。” \par}

禹曰：“俞哉！帝光天之下，至於海隅蒼生，\footnote{光天之下，至於海隅，蒼蒼然生草木。言所及廣遠。}萬邦黎獻，共惟帝臣，惟帝時舉。敷納以言，明庶以功，車服以庸。\footnote{獻，賢也。萬國眾賢，共為帝臣。帝舉是而用之，使陳布其言，明之皆以功大小為差，以車服旌其能用之。}誰敢不讓？敢不敬應。\footnote{上惟賢是用，則下皆敬應上命而讓善。應，應對之應。}帝不時，敷同日奏罔功。\footnote{帝用臣不是，則遠近布同而日進於無功，以賢愚並位，優劣共流故。}


{\noindent\zhuan\zihao{6}\fzbyks 傳“光天”至“廣遠”。正義曰：\CJKunderwave{堯典}之序，訓“光”為充，即此亦為充，言充滿大天之下也。據其方面即四隅，為遠至於海隅，舉極遠之處，言帝境所及廣遠，其內多賢人也。 \par}

{\noindent\zhuan\zihao{6}\fzbyks 傳“獻賢”至“用之”。正義曰：\CJKunderwave{釋言}云：“獻，聖也。賢是聖之次，臣德不宜言聖，故為賢也。“萬國眾賢，共為帝臣”，言求臣之處多也。帝舉是眾賢而用之，使陳布其言，令其自說己之所能，聽其言而納受之，依其言而考試之，顯明眾臣，皆以功大小為差，然後賜車服以旌別其人功能事用，是舉賢用人之法也。\CJKunderwave{舜典}云：“敷奏以言，明試以功。”“奏”、“試”二字與此異者，彼言施於諸侯，其人見為國君，故令奏言試功；此謂方始擢用，故言“納”、“庶”。“納”謂受取之，“庶”謂在群眾。 \par}

{\noindent\zhuan\zihao{6}\fzbyks 傳“帝用”至“流故”。正義曰：帝用臣不是，不以言考功，在下知帝不分別善惡，則無遠近遍佈同心，日日進於無功之人，由其賢愚並位，優劣共流故也。“敷”是布之義，故言“遠近布同”，同心妄舉也。 \par}

無若丹朱傲，惟慢遊是好。\footnote{丹朱，堯子。舉以戒之。傲，五報反，字又作奡。好，呼報反。}傲虐是作，罔晝夜頟頟。\footnote{傲戲而為虐，無晝夜，常崑崙肆惡無休息。傲，五羔反，徐五報反，注同。頟，五客反。}罔水行舟，朋淫於家，用殄厥世。\footnote{朋，群也。丹朱習於無水陸地行舟，言無度。群淫於家，妻妾亂。用是絕其世，不得嗣。殄，徒見反。}


{\noindent\zhuan\zihao{6}\fzbyks 傳“丹朱，堯子”。正義曰：\CJKunderwave{漢書·律曆志}云：“堯讓舜，使子朱處於丹淵，為諸侯。”則“朱”是名，“丹”是國也。 \par}

{\noindent\zhuan\zihao{6}\fzbyks 傳“傲戲”至“休息”。正義曰：\CJKunderwave{詩}美衛武公云：“善戲謔兮，不為虐兮。”丹朱反之，故“傲戲而為虐”也。“崑崙”是不休息之意，“肆”謂縱恣也，晝夜常崑崙然縱恣為惡，無休息時也。 \par}

{\noindent\zhuan\zihao{6}\fzbyks 傳“朋群”至“得嗣”。正義曰：朋輩與群聚義同，故“朋”為群也。聖人作車以行陸，作舟以行水，丹朱乃習於無水而陸地行舟，言其所為惡事無節度也。此乃稟受惡性，習惡事也。鄭玄云：“丹朱見洪水時人乘舟，今水已治，猶居舟中，崑崙使人推行之。”案下句云，予創若時，乃勤治水，則丹朱行舟之時，水尚未除，非效洪水之時人乘舟也。“群淫於家”，言群聚妻妾，恣意淫之，無男女之別，故言“妻妾亂”也。用是之惡，故絕其世位，不得嗣父也。此用“殄厥世”一句，禹既見世絕，今始言之，以明行惡之驗。此句非禹所創，創之者，創其行之惡耳。 \par}

予創若時，娶於塗山，辛、壬、癸、甲。\footnote{創，懲也。塗山，國名。懲丹朱之惡，辛日娶妻，至於甲日，復往治水，不以私害公。娶，促住反。復，扶又反。}啟呱呱而泣，予弗子，惟荒度土功。\footnote{啟，禹子也。禹治水，過門不入，聞啟泣聲，不暇子名之，以大治度水土之功故。呱音孤。予如字，鄭將吏反。度，徒洛反。}


{\noindent\zhuan\zihao{6}\fzbyks 傳“創懲”至“害公”。正義曰：“創”與“懲”皆是見惡自止之意，故云“創，懲也”。哀七年\CJKunderwave{左傳}云：“禹會諸侯於塗山。”杜預云：“塗山在壽春縣東北。”“塗山，國名”，蓋近彼山也。“娶於塗山”，言其所娶之國耳，非就妻家見妻也。懲丹朱之惡，故不可不勤，故辛日娶妻，至於甲日復往治水。孔雲“復往”，則已嘗治水，而輟事成昏也。鄭玄云：“登用之年，始娶於塗山氏，三宿而為帝所命治水。”鄭意娶後始受帝命，娶前未治水也。然娶後始受帝命，當雲聞命即行，不須計辛之與甲日數多少，當如孔說，輟事成昏也。此時禹父新殛，而得為昏者，鯀放而未死，不妨禹娶。且治水四年，兗州始畢，禹娶不必在殛鯀之年也。 \par}

{\noindent\zhuan\zihao{6}\fzbyks 傳“啟禹”至“功故”。正義曰：“啟，禹子”，\CJKunderwave{世本}文也。\CJKunderwave{孟子}稱禹治水，三過其門而不入，是至門而聞啟泣聲，不暇如人父子，名為己子而愛念之，以其為大治度水土之功故也。訓“荒”為“大治”，謂去其水。“度”謂量其功，故“治度”連言之。 \par}

弼成五服，至於五千，州十有二師。\footnote{五服，侯、甸、綏、要、荒服也。服五百里，四方相距為方五千裡，治洪水輔成之。一州用三萬人功，九州二十七萬庸。至於五千，馬云：“面五千裡，為方萬里。”鄭云：“五服已五千，又弼成為萬里。”州十有二師，二千五百人為師，鄭云：“師，長也。”要，一遙反。}


{\noindent\zhuan\zihao{6}\fzbyks 傳“五服”至“萬庸”。正義曰：據\CJKunderwave{禹貢}所云五服之名數，知五服即甸、侯、綏、要、荒服也。彼五服每服五百里，四面相距為方五千裡也。王肅云：“五千裡者，直方之數。若其迥邪委曲，動有倍加之較。”是直路五千裡也。“治洪水輔成之”者，謂每服之內，為其小數，定其差品，各有所掌，是禹輔成之也。\CJKunderwave{周禮}大司馬法，二千五百人為師。每州十有二師，通計之一州用三萬人功，總計九州用二十七萬庸。“庸”亦功也。州境既有闊狹，用功必有多少，例言“三萬人”者，大都通率為然，惟言用“三萬人”者，不知用功日數多少，治水四年乃畢，用功蓋多矣，不知用幾日也。鄭玄云：“輔五服而成之，至於面方，各五千裡，四面相距為方萬里。九州州立十二人為諸侯師，以佐牧。堯初制五服，服各五百里。要服之內方四千裡，曰九州。其外荒服，曰四海。此禹所受，\CJKunderwave{地記書}曰‘崑崙山東南，地方五千裡，名曰神州’者。禹弼五服之殘數，亦每服者合五百里，故有萬里之界、萬國之封焉。猶用要服之內為九州，州更方七千裡。七七四十九，得方千里者四十九。其一以為圻內，餘四十八，八州分而各有六。\CJKunderwave{春秋}傳曰：‘禹朝群臣於會稽，執玉帛者萬國。’言執玉帛者，則九州之內諸侯也。其制特置牧，以諸侯賢者為之師。蓋百國一師，州十有二師，則州千二百國也。八州凡九千六百國，其餘四百國在圻內。與\CJKunderwave{王制}之法準之，八州通率封公侯百里之國者一，伯七十里之國二，子男五十里之國四，方百里者三，封國七有畸，至於圻內，則子男而已。”鄭云：“禹弼成五服，面各五千裡。”王肅\CJKunderwave{禹貢}之注已難之矣。傳稱“萬”，盈數也，“萬國”舉盈數而言，非謂其數滿萬也。\CJKunderwave{詩}桓曰“綏萬邦烝民”，曰“揉此萬邦”，豈周之建國復有萬乎？天地之勢，平原者甚少，山川所在不啻居半，豈以不食之地，亦封建國乎？王圻千里，封五十里之國四百，則圻內盡以封人，王城宮室無建立之處，言不顧實，何至此也？百國一師，不出典記，自造此語，何以可從？“禹朝群臣於會稽”，\CJKunderwave{魯語}文也。“執玉帛者萬國”，\CJKunderwave{左傳}文也，採合二事，亦為謬矣。 \par}

外薄四海，咸建五長。\footnote{薄，迫也。言至海。諸侯五國立賢者一人為方伯，謂之五長，以相統治，以獎帝室。薄，蒲各反，徐扶各反。長，丁丈反。五長，眾官之長。}各迪有功，苗頑弗即工，帝其念哉！\footnote{九州五長各蹈為有功，惟三苗頑兇，不得就官。善惡分別。別，彼列反。}帝曰：“迪朕德，時乃功惟敘。”\footnote{言天下蹈行我德，是汝治水之功有次序，敢不念乎！}

{\noindent\zhuan\zihao{6}\fzbyks 傳“薄迫”至“帝室”。正義曰：\CJKunderwave{釋言}云：“逼，迫也。”“薄”者逼近之義，故云迫也。外迫四海，言從京師而至於四海也。\CJKunderwave{釋地}云：“九夷、八狄、七戎、六蠻謂之四海。”謂九州之外也。\CJKunderwave{王制}云：“五國以為屬，屬有長。”此“建五長”亦如彼文，故云“諸侯五國立賢者一人為方伯，謂之五長，以相統治，欲以共獎帝室故”也。僖元年\CJKunderwave{公羊傳}曰：“上無天子，下無方伯。”“方伯”謂\CJKunderwave{周禮}“九命作伯”者也。\CJKunderwave{王制}云：“千里之外設方伯。”“方伯”一為之長，謂\CJKunderwave{周禮}“八命作牧”者也。傳言五國立一人為方伯，直是五國之長耳，與彼異也，以其是當方之長，故傳以“方伯”言之。 \par}

{\noindent\zhuan\zihao{6}\fzbyks 傳“九州”至“分別”。正義曰：蹈為有功之長，言蹈履典法，持之有功。惟三苗頑兇，不得就官，謂舜分北三苗之時，苗君有罪，不得就其諸侯國君之官，而被流於遠方也。言“九州五長各蹈為有功”，則海內諸侯皆有功矣。惟有三苗不得就官，以見天下大治，而惡者少耳。頑則不得就官，言善惡分別也。 \par}

{\noindent\shu\zihao{5}\fzkt “禹曰”至“惟敘”。正義曰：禹既得帝言，乃答帝曰：“然。既帝之任臣，又言當擇人，充滿大天之下，旁至四海之隅，蒼蒼然生草木之處，皆是帝德所及。其內有萬國眾賢，皆共為帝臣。”言其可用者甚眾也。“帝當就是眾賢之內，舉而用之。其舉用之法，各使陳布其言，納受之，以其言之所能，從其所能而驗試之。明顯眾人所能，當以功之大小。既知有功，乃賜之以車服，以表其功有能用。帝以此法用人，即在下之人，知官不妄授，必用度才能而使之。如此，誰敢不讓有德？敢不敬應帝命而推先善人也？若帝用臣不是，不宜試驗，不知臧否，則群臣遠近，遍佈同心，而日進無功之人”。既戒帝擇人，又勸帝自勤。“無若丹朱之傲，惟慢褻之遊是其所好。傲戲而為虐，是其所為。為此惡事，不問晝夜，而頟頟然恆為之無休息。又無水而陸地行舟，群朋淫泆於室家之內。用此之故，絕其世嗣，不得居位。我本創丹朱之惡若是也，故娶於塗山之國，歷辛、壬、癸、甲四日而即往治水。其後過門不入，聞啟呱呱而泣，我不暇入而子名之，惟以大治度水土之功故也。水土既平，乃輔成五服，四面相距，至於五千裡。州十有二師，其治水之時，所役人功，每州用十有二師，各用三萬人也。自京師外迫及四海，其間諸侯五國皆立一長，迆相統領。以此諸侯各蹈行所職，併為有功，惟有三苗頑兇，不能就官，我以供勤之故，得使災消沒。帝念此事哉！不可不自勤也”。帝答禹曰：“天下之人皆蹈行我德，是汝治水之功，惟有次敘故也。”受其戒而美其功也。 \par}

皋陶方祗厥敘，方施象刑惟明。\footnote{方，四方。禹五服既成，故皋陶敬行其九德考績之次序於四方，又施其法刑，皆明白。史因禹功重美之。重，直用反。}

{\noindent\zhuan\zihao{6}\fzbyks 傳“方四”至“美之”。正義曰：皋陶為帝所任，遍及天下，故“方”為四方也。天下蹈行帝德，水土既治，亦由刑法彰明，若使水害不息，皋陶法無所施，若無皋陶以刑，人亦未能奉法。天下蹈行帝德，二臣共有其功，故史因帝歸功於禹，兼記皋陶之功。\CJKunderwave{舜典}與\CJKunderwave{大禹謨}已美皋陶，故言“重美之”也。傳言“考績之次敘”者，皋陶所言九德，依德以考其功績，亦是刑法之事，故兼言也。鄭雲“歸美於二臣”，則以此經為帝語。此文上無所由，下無所結，形勢非語辭也，故傳以為史因記之 \par}

{\noindent\shu\zihao{5}\fzkt “皋陶”至“惟明”。正義曰：此經史述為文，非帝言也。史以禹成五服，帝念禹功，故因美皋陶。言禹既弼成五服，故皋陶於其四方敬行九德考績之法，有次敘也。又於四方施其刑法，惟明白也。由禹有此大功，故史重美之也。 \par}

夔曰：“戛擊鳴球，搏拊琴瑟以詠。祖考來格。\footnote{戛擊，柷敔，所以作止樂。搏拊以韋為之，實之以糠，所以節樂。球，玉磬。此舜廟堂之樂，民悅其化，神歆其祀，禮備樂和，故以祖考來至明之。夔，求龜反。戛，居八反，徐古八反，馬云：“櫟也。”球音求。搏音博。拊音撫。柷，尺叔反，所以作樂。敔，魚呂反，所以止樂。糠音康。歆，許金反。}虞賓在位，群后德讓。\footnote{丹朱為王者後，故稱賓。言與諸侯助祭，班爵同，推先有德。}


{\noindent\zhuan\zihao{6}\fzbyks 傳“戛擊”至“明之”。正義曰：“戛擊”是作用之名，非樂器也，故以“戛擊”為柷敔。柷敔之狀，經典無文，漢初已來學者相傳，皆雲柷如漆桶，中有椎柄，動而擊其旁也。敔狀如伏虎，背上有刻，戛之以為聲也。樂之初，擊柷以作之；樂之將末，戛敔以止之，故云“所以作止樂”雙解之。\CJKunderwave{釋樂}云：“所以鼓柷謂之止，所以鼓敔謂之籈。”郭璞云：“柷如漆桶，方二尺四寸，深一尺八寸，中有椎，柄連氐，挏之令左右擊。止者，其椎名也。敔如伏虎，背上有二十七鉏鋙刻，以木長一尺櫟之。籈者，其名也。”是言擊柷之椎名為“止”，戛敔之木名為“籈”，“戛”即櫟也。漢禮器制度及\CJKunderwave{白虎通}、馬融、鄭玄、李巡其說皆為然也。惟郭璞為詳，據見作樂器而言之。搏拊形如鼓，以韋為之，實之以糠，擊之以節樂，漢初相傳為然也。\CJKunderwave{釋器}云：“球，玉也。”“鳴球”謂擊球使鳴，樂器惟磬用玉，故球為玉磬。\CJKunderwave{商頌}雲“依我磬聲”，磬亦玉磬也。鄭玄云：“磬，懸也，而以合堂上之樂。玉磬和，尊之也。”然則鄭以球玉之磬懸於堂下，尊之，故進之使在上耳。此“舜廟堂之樂”，謂廟內堂上之樂，言“祖考來格”，知在廟內，下雲“下管”，知此在堂上也。馬融見其言“祖考”，遂言“此是舜除瞽瞍之喪，祭宗廟之樂”，亦不知舜父之喪在何時也。但此論韶樂，必在即政後耳。此說樂音之和，而云“祖考來格”者，聖王先成於人，然後致力於神。言“人悅其化，神歆其祀，禮備樂和，所以“祖考來至“明矣。以祖考來”至“明樂之和諧也。詩稱神之格思不可度思而云：“祖考來至”者，王肅云：“祖考來至者，見其光輝也。”蓋如\CJKunderwave{漢書·郊祀志}稱武帝郊祭天祠，上有美光也。此經文次，以“柷敔”是樂之始終，故先言“戛擊”。其“球”與“搏拊琴瑟”，皆當彈擊，故使“鳴”冠於“球”上，使下共蒙之也。鄭玄以“戛擊鳴”球三者，皆總下樂，櫟擊此四器也”。樂器惟敔當櫟耳，四器不櫟，鄭言非也。 \par}

{\noindent\zhuan\zihao{6}\fzbyks 傳“丹朱”至“有德”。正義曰：\CJKunderwave{微子之命}雲“作賓於王家”，\CJKunderwave{詩}頌微子之來，謂之“有客”，是王者之後，為時王所賓也。故知“虞賓”謂丹朱，為王者後，故稱賓也。王者立二代之後，而獨言丹朱者，蓋高辛氏之後，無文而言，故惟指丹朱也。王者之後，尊於群后，故殊言“在位”。群后亦在位也，後言德讓，丹朱亦以德讓矣，故言“與諸侯助祭，年爵同者，推先有德”也。二王之後併為上公，亦有與丹朱爵同，故丹朱亦讓也。丹朱之性下愚，堯不能化，此言“有德”者，猶上雲“瞽亦允若”，暫能然也。 \par}

下管鼗鼓，合止柷敔，\footnote{堂下樂也。上下合止樂，各有柷敔，明球、弦、鍾、籥，各自互見。鼗音桃。合如字，徐音合。籥，餘若反。互音護。見，賢遍反，下“見細器”同。}笙鏞以間，鳥獸蹌蹌。\footnote{鏞，大鐘。跡葴迭也。吹笙擊鐘，鳥獸化德，相率而舞，蹌蹌然。鏞音庸。間，間側之間。鳥獸，孔以為自舞也。馬云：“鳥獸，筍篪也。”蹌，七羊反，舞貌。\CJKunderwave{說文}作蹌，云：“鳥獸求食聲。”迭，直結反。}


{\noindent\zhuan\zihao{6}\fzbyks 傳“堂下”至“互見”。正義曰：經言“下管”，知是“堂下樂”也。敔當戛之，柷當擊之，上言“戛擊”，此言“柷敔”，其事是一，故云“上下合止樂，各有柷敔”也。言堂下堂上合樂各以柷，止樂各以敔也。上言作用，此言器名，兩相備也。上下皆有“柷敔”，兩見其文，明球、弦、鍾、籥，上下樂器不同，各自更互見也。弦謂琴瑟。鍾，鏞也。籥，管也。琴瑟在堂，鍾籥在庭，上下之器各別，不得兩見其名，各自更互見之。依\CJKunderwave{大射}禮，鐘磬在庭，今鳴球於廟堂之上者，案\CJKunderwave{郊特牲}雲“歌者在上”，貴人聲也。\CJKunderwave{左傳}雲“歌鐘二肆”，則堂上有鍾，明磬亦在堂上，故漢魏已來登歌皆有鐘磬。\CJKunderwave{燕禮}、\CJKunderwave{大射}堂上無鐘磬者，諸侯樂不備也。 \par}

{\noindent\zhuan\zihao{6}\fzbyks 傳“鏞大”至“蹌蹌然”。正義曰：\CJKunderwave{釋樂}云：“大鐘謂之鏞。”李巡曰：“大鐘音聲大。鏞，大也。”孫炎曰：“鏞，深長之聲。”\CJKunderwave{釋詁}云：“間，代也。”孫炎曰：“間廁之代也。”\CJKunderwave{釋言}云：“遞，迭也。”李巡曰：“遞者，更迭間廁，相代之義。”故“間”為迭也。吹笙擊鐘，更迭而作，鳥獸化德，相率而舞，蹌蹌然。下雲“百獸率舞”，知此“蹌蹌然”亦是舞也。\CJKunderwave{禮}雲“凡行容愓愓”，“大夫濟濟，士蹌蹌”，是為行動之貌，故為舞也。 \par}

簫韶九成，鳳皇來儀。”\footnote{韶，舜樂名。言簫，見細器之備。雄曰鳳，雌曰皇，靈鳥也。儀，有容儀。備樂九奏而致鳳皇，則餘鳥獸不待九而率舞。韶，時昭反。}夔曰：“於！予擊石拊石，百獸率舞，庶尹允諧。”\footnote{尹，正也，眾正官之長。信皆和諧，言神人治。始於任賢，立政以禮，治成以樂，所以太平。於、予並如字。}

{\noindent\zhuan\zihao{6}\fzbyks 傳“韶舜”至“率舞”。正義曰：“韶”是舜樂，經傳多矣，但余文不言“簫”。“簫”乃樂器，非樂名，簫是樂器之小者。“言簫，見細器之備”，謂作樂之時，小大之器皆備也。\CJKunderwave{釋鳥}云：“鶠，鳳。其雌皇。”是此鳥“雄曰鳳，雌曰皇”。\CJKunderwave{禮運}云：“麟、鳳、龜、龍謂之四靈。”是鳳皇為神靈之鳥也。\CJKunderwave{易·漸卦}上九：“鴻漸於陸，其羽可用為儀。”是儀為“有容儀”也。“成”謂樂曲成也。鄭云：“成猶終也。”每曲一終，必變更奏，故經言“九成”，傳言“九奏”，\CJKunderwave{周禮}謂之“九變”，其實一也。言“簫”見細器之備，備樂九奏而致鳳皇，則其餘鳥獸不待九而率舞也。尊者體盤，靈瑞難致，故“九成”之下始言“鳳皇來儀”。“鳥獸蹌蹌”乃在上句，傳據此文言鳥獸易來，鳳皇難致，故云“鳥獸不待九”也。樂之作也，依上下遞奏間合而後曲成，神物之來，上下共致，非堂上堂下別有忻感。以祖考尊神，配堂上之樂；鳥獸賤物，故配堂下之樂。總上下之樂，言九成致鳳。尊異靈瑞，故別言爾，非堂上之樂獨致神來，堂下之樂偏令獸舞也。鄭玄注\CJKunderwave{周禮}具引此文，乃雲“此其在於宗廟九奏效應也”。是言祖考來格、百獸率舞皆是九奏之事也。\CJKunderwave{大司樂}云：“凡六樂者，六變而致象物及天神。”鄭玄云：“象物，有象在天，所謂四靈者。”彼謂大蜡之祭，作樂以致其神。此謂鳳皇身至，故九奏也。 \par}

{\noindent\zhuan\zihao{6}\fzbyks 傳“尹正”至“太平”。正義曰：“尹，正”，\CJKunderwave{釋言}文。“眾正官之長”，謂每職之首，\CJKunderwave{周官}所謂“唐虞稽古，建官惟百”是也。“信皆和諧”，言職事修理也。上雲“祖考來格”，此言眾正官治，言神人洽，樂音和也。此篇初說用臣之法，末言樂音之和，言其“始用任賢，立政以禮，治成以樂，所以得致太平”，解史錄夔言之意。 \par}

{\noindent\shu\zihao{5}\fzkt “夔曰”至“允諧”。正義曰：皋陶、大禹為帝設謀，大聖納其昌言，天下以之致治，功成道洽，禮備樂和，史述夔言，繼之於後。夔曰：“在舜廟堂之上，戛敔擊柷，鳴球玉之磬，擊搏拊，鼓琴瑟，以歌詠詩章，樂音和協，感致幽冥，祖考之神來至矣。虞之賓客丹朱者在於臣位，與群君諸侯以德相讓。此堂上之樂，所感深矣。又於堂下吹竹管，擊鼗鼓，合樂用柷，止樂用敔，吹笙擊鐘，以次迭作，鳥獸相率而舞，其容蹌蹌然。堂下之樂，感亦深矣。簫韶之樂，作之九成，以致鳳皇來而有容儀也。”夔又曰：“嗚呼！”嘆舜樂之美。“我大擊其石磬，小拊其石磬，百獸相率而舞，鳥獸感德如此，眾正官長信皆和諧矣。”言舜致教平而樂音和，君聖臣賢，謀為成功所致也。 \par}

帝庸作歌曰：“敕天之命，惟時惟幾。”\footnote{用庶尹允諧之政，故作歌以戒，安不忘危。敕，正也。奉正天命以臨民，惟在順時，惟在慎微。}乃歌曰：“股肱喜哉！元首起哉！百工熙哉！”\footnote{元首，君也。股肱之臣喜樂盡忠，君之治功乃起，百官之業乃廣。樂者洛。盡，津忍反。}皋陶拜手稽首，颺言曰：“念哉！\footnote{大言而疾曰颺。承歌以戒帝。颺音揚。}率作興事，慎乃憲，欽哉！\footnote{憲，法也。天子率臣下為起治之事，當慎汝法度，敬其職。}


{\noindent\zhuan\zihao{6}\fzbyks 傳“用庶”至“慎微”。正義曰：此承夔言之下，既得總言而歌，故知“帝庸作歌”者，用“庶尹允諧之政，故作歌以自戒之，安不忘危”也。“敕”是正齊之意，故為正也。言人君奉正天命，以臨下民，惟在順時，不妨農務也，惟在慎微，不忽細事也。鄭玄以為戒臣，孔以為自戒者，以正天之命是人君之事故也。 \par}

{\noindent\zhuan\zihao{6}\fzbyks 傳“元首”至“乃廣”。正義曰：\CJKunderwave{釋詁}云：“元、良，首也。”僖三十三年\CJKunderwave{左傳}稱狄人歸先軫之元，則“元”與“首”各為頭之別名，此以“元首”共為頭也。君臣大體猶如一身，故“元首，君也”。“股肱之臣喜樂盡忠”，謂樂行君之化。“君之治功乃起”言無變事業，事業在於百官，故眾功皆起，百官之業乃廣也。 \par}

{\noindent\zhuan\zihao{6}\fzbyks 傳“憲法”至“其識”。正義曰：“憲，法”，\CJKunderwave{釋詁}文。此言“興事”，對上“起哉”。“天子率臣下為起治之事”，言臣不能獨使起也。 \par}

屢省乃成，欽哉！\footnote{屢，數也。當數顧省汝成功，敬終以善，無懈怠。屢，力具反。省，悉井反。數，色角反。懈，隹賣反。}乃賡載歌曰：“元首明哉！股肱良哉！庶事康哉！”\footnote{賡，續。載，成也。帝歌歸美股肱，義未足，故續歌。先君後臣，眾事乃安，以成其義。賡，加孟反，劉皆行反，\CJKunderwave{說文}以為古續字。}又歌曰：“元首叢脞哉！股肱惰哉！萬事墮哉！\footnote{叢脞，細碎無大略。君如此，則臣懈惰，萬事墮廢，其功不成。歌以申戒。叢，太公反。脞，倉果反。徐音鎖，馬云：“叢，總也。脞，小也。”惰，徒臥反。墮，許規反。}帝拜曰：“俞，往欽哉！”\footnote{拜受其歌，戒群臣自今以往，敬其職事哉。}

{\noindent\zhuan\zihao{6}\fzbyks 傳“屢數”至“懈怠”。正義曰：\CJKunderwave{釋詁}云：“屢、數，疾也。”俱訓為疾，故“屢”為數也。“顧省汝成功”，謂已有成功，令數顧省之，敬終以善，無懈怠也。恐其惰於已成功，故以此為戒。 \par}

{\noindent\zhuan\zihao{6}\fzbyks 傳“賡續”至“其義”。正義曰：\CJKunderwave{詩}云：“西有長賡。”\CJKunderwave{毛傳}亦以“賡”為續，是相傳有此訓也。鄭玄以“載”為始，孔以“載”為成，各以意訓耳。“帝歌歸美股肱，義未足”者，非君之明，為臣不能盡力，空責臣功是其義未足。以此續成帝歌，必先君後臣，眾事乃安，故以此言成其義也。 \par}

{\noindent\zhuan\zihao{6}\fzbyks 傳“叢脞”至“申戒”。正義曰：孔以“叢脞”為細碎無大略，鄭以“叢脞，總聚小小之事以亂大政”，皆是以意言耳。君無大略，則不能任賢，功不見知，則臣皆懈惰，萬事墮廢，其功不成，故又歌以重戒也。“庶事”、“萬事”，為義同而文變耳。 \par}

{\noindent\shu\zihao{5}\fzkt “帝庸”至“往欽哉”。正義曰：帝既得夔言，用此庶尹允諧之政，故乃作歌自戒。將歌而先為言曰：“人君奉正天命，以臨下民，惟當在於順時，惟當在於慎微。”既為此言，乃歌曰：“股肱之臣喜樂其事哉！元首之君政化乃起哉！百官事業乃得廣大哉！”言君之善政由臣也。皋陶拜手稽首，颺聲大言曰：“帝當念是言哉！率領臣下，為起政治之事，慎汝天子法度而敬其職事哉！又當數自顧省已之成功而敬終之哉！”乃續載帝歌曰：“會是元首之君能明哉！則股肱之臣乃善哉！？事皆得安寧哉！”既言其美，又戒其惡：“元首之君叢脞細碎哉！則股肱之臣懈怠緩慢哉！眾事悉皆墮廢哉！”言政之得失由君也。帝拜而受之曰：“然。”然其所歌顯是也。“汝群臣自今已往，各敬其職事哉！” \par}

%%% Local Variables:
%%% mode: latex
%%% TeX-engine: xetex
%%% TeX-master: "../Main"
%%% End:

%% -*- coding: utf-8 -*-
%% Time-stamp: <Chen Wang: 2024-04-02 11:42:41>

% {\noindent \zhu \zihao{5} \fzbyks } -> 注 (△ ○)
% {\noindent \shu \zihao{5} \fzkt } -> 疏

\part{夏書}


\chapter{卷六}


\section{禹貢第一（禹貢上、中、下）}

\CJKunderline{禹}別九州，\footnote{分其圻界。別，彼列反。九州，\CJKunderwave{周公職錄}云：“黃帝受命，風后受圖，割地布九州。”\CJKunderwave{鄹子}云：“中國為赤縣，內有九州。”\CJKunderwave{春秋說題辭}云：“州之言殊也。”圻，其依反。}隨山濬川，\footnote{刊其木，深其流。濬，思俊反。刊，苦安反。}任土作貢。\footnote{任其土地所有，定其貢賦之差。此堯時事，而在\CJKunderwave{夏書}之首，\CJKunderline{禹}之王以是功。任，而鴆反。貢，字或作贛。王，於況反。}

{\noindent\zhuan\zihao{6}\fzbyks 傳“分其圻界”。正義曰：\CJKunderwave{詩}傳云：“圻，疆也。”分其疆界，使有分限。計九州之境，當應舊定，而云“\CJKunderline{禹}別”者，以堯遭洪水，萬事改新，此為作貢生文，故言“\CJKunderline{禹}別”耳。 \par}

{\noindent\zhuan\zihao{6}\fzbyks 傳“刊其木，深其流”。正義曰：經言“隨山刊木”，序以較略為文，直雲“隨山”，不雲隨山為何事，故傳明之隨山刊其木也。“濬川”，深其流也。“隨山”本為“濬川”，故連言之。 \par}

{\noindent\zhuan\zihao{6}\fzbyks 傳“任其”至“是功”。正義曰：九州之土，物產各異，任其土地所有，以定貢賦之差，既任其所有，亦因其肥瘠多少不同，製為差品。\CJKunderline{鄭玄}云：“任土謂定其肥磽之所生。”是言用肥瘠多少為差也。“賦”者，自上稅下之名，謂治田出谷，故經定其差等，謂之“厥賦”。“貢”者，從下獻上之稱，謂以所出之谷，市其土地所生異物，獻其所有，謂之“厥貢”。雖以所賦之物為貢用，賦物不盡有也，亦有全不用賦物，直隨地所有，採取以為貢者，此之所貢，即與\CJKunderwave{周禮·太宰}“九貢”不殊，但\CJKunderwave{周禮}分之為九耳。其賦與\CJKunderwave{周禮}“九賦”全異，彼賦謂口率出錢。不言“作賦”而言“作貢”者，取下供上之義也。諸序皆言作某篇，此序不言“作\CJKunderwave{禹貢}”者，以發首言“\CJKunderline{禹}”，句末言“貢”，篇名足以顯矣。百篇之序，此類有三。“微子作誥父師、少師”，不言“作\CJKunderwave{微子}”，“仲虺作誥”，不言“作\CJKunderwave{仲虺之誥}”，與此篇皆為理足而略之也。又解篇在此之意，此治水是堯末時事，而在\CJKunderwave{夏書}之首，\CJKunderline{禹}之得王天下，以是治水之功，故以為\CJKunderwave{夏書}之首。此篇史述時事，非是應對言語，當是水土既治，史即錄此篇，其初必在\CJKunderwave{虞書}之內，蓋夏史抽入\CJKunderwave{夏書}，或仲尼始退其第，事不可知也。 \par}

{\noindent\shu\zihao{5}\fzkt “\CJKunderline{禹}別”至“作貢”。正義曰：\CJKunderline{禹}分別九州之界，隨其所至之山，刊除其木，深大其川，使得注海。水害既除，地複本性，任其土地所有，定其貢賦之差，史錄其事，以為\CJKunderwave{禹貢}之篇。 \par}

禹貢\footnote{禹制九州貢法。}

{\noindent\zhuan\zihao{6}\fzbyks 傳“\CJKunderline{禹}制九州貢法”。正義曰：\CJKunderline{禹}制貢法，故以“\CJKunderline{禹}貢”名篇。貢賦之法其來久矣，治水之後更復改新，言此篇貢法是\CJKunderline{禹}所制，非\CJKunderline{禹}始為貢也。 \par}

{\noindent\shu\zihao{5}\fzkt “\CJKunderline{禹}貢”。正義曰：此篇史述為文，發首“奠高山大川”，言\CJKunderline{禹}治九州之水，水害既除，定山川次秩，與諸州為引序。自“導岍”至“嶓冢”，條說所治之山，言其首尾相及也。自“導弱水”至“導洛”，條說所治之水，言其發源注海也。自“九州攸同”至“成賦中邦”，總言水土既平，貢賦得常之事也。“錫土姓”三句，論天子於土地布行德教之事也。自“五百里甸服”至“二百里流”，總言四海之內，量其遠近，分為五服之事也。自“東漸於海”以下，總結\CJKunderline{禹}功成受錫之事也。 \par}

\CJKunderline{禹}敷土，隨山刊木，\footnote{洪水泛溢，\CJKunderline{禹}布治九州之土，隨行山林，斬木通道。○敷，芳無反，馬云：“分也。”泛，敷劍反。行，下孟反。}奠高山大川。\footnote{奠，定也。高山，五嶽。大川，四瀆。定其差秩，祀禮所視。奠，田遍反。瀆音獨，下同。}

{\noindent\zhuan\zihao{6}\fzbyks 傳“洪水”至“通道”。正義曰：\CJKunderwave{詩}傳云：“泛，泛流也。”泛是水流之貌，洪水流而泛溢，浸壞民居，故\CJKunderline{禹}分佈治之。知者，文十八年\CJKunderwave{左傳}雲“舉八凱使主后土”，則伯益之輩佐\CJKunderline{禹}多矣，\CJKunderline{禹}必身行九州，規謀設法，乃使佐巳之人分佈治之。於時平地盡為流潦，鮮有陸行之路，故將欲治水，隨行山林，斬木通道。鄭云：“必隨州中之山而登之，除木為道，以望觀所當治者，則規其形而度其功焉。”是言\CJKunderline{禹}登山之意也。\CJKunderwave{孟子}曰，\CJKunderline{禹}三過門不入。其家門猶三過之，則其餘所歷多矣。來而復往，非止一處，故言分佈治之之。 \par}

{\noindent\zhuan\zihao{6}\fzbyks 傳“奠定”至“所視”。正義曰：\CJKunderwave{禮}定器於地，通名為“奠”，是“奠”為定也。山之高者，莫高於嶽；川之大者，莫大於瀆；故言“高山，五嶽”，謂、嵩、岱、衡、華、恆也；“大川，四瀆”，謂江、河、淮、濟也。此舉高大為言，卑小亦定之矣。\CJKunderwave{舜典}云：“望秩于山川。”故言“定其差秩”，定其大小次敘也。定其“祀禮所視”，謂\CJKunderwave{王制}所云“五嶽視三公，四瀆視諸侯，其餘視伯子男”。往者洪水滔天，山則為水所包，川則水皆泛溢，祭祀禮廢，今始定之，以見水土平，復舊制也。經雲“荊岐既旅”，“蔡蒙旅平”，“九山刊旅”，是次秩既定，故旅祭之。 \par}

{\noindent\shu\zihao{5}\fzkt “\CJKunderline{禹}敷”至“大川”。正義曰：言\CJKunderline{禹}分佈治此九州之土，其治之也，隨行所至之山，除木通道，決流其水，水土既平，乃定其高山大川。謂定其次秩尊卑，使知祀禮所視。言\CJKunderline{禹}治其山川，使復常也。 \par}

\textcolor{red}{冀州}既載，\footnote{堯所都也。先施貢賦役，載於書。冀，居器反。州，九州名義見\CJKunderwave{爾雅音}。載如字。載，載於書也；馬同，鄭、韋昭云：“載，事也。”}

{\noindent\zhuan\zihao{6}\fzbyks 傳“堯所”至“於書”。正義曰：史傳皆雲堯都平陽，\CJKunderwave{五子之歌}曰：“惟彼陶唐，有此冀方。”是冀州堯所都也。諸州冀為其先，治水先從冀起，為諸州之首。記其役功之法，“既載”者，言先施貢賦役，載於書也。謂計人多少，賦功配役，載於書籍，然後徵而用之，以治水也。冀州如此，則餘州亦然，故於此特記之也。王肅云：“言巳賦功屬役，載於書籍。”傳意當然，鄭云：“載之言事，事謂作徒役也。\CJKunderline{禹}知所當治水，又知用徒之數，則書於策以告帝，徵役而治之。”惟解“載”字為異，其意亦同孔也。 \par}

{\noindent\shu\zihao{5}\fzkt “冀州”。正義曰：九州之次，以治為先後。以水性下流，當從下而洩，故治水皆從下為始。冀州，帝都，於九州近北，故首從冀起。而東南次兗，而東南次青，而南次徐，而南次揚，從揚而西次荊，從荊而北次豫，從豫而西次梁，從梁而北次雍，雍地最高，故在後也。自兗已下，皆準地之形勢，從下向高，從東向西。青、徐、揚三州併為東偏，雍州高於豫州，豫州高於青、徐，雍、豫之水從青、徐而入海也。梁高於荊，荊高於揚，梁、荊之水從揚而入海也。兗州在冀州東南，冀、兗二州之水，各自東北入海也。冀州之水不經兗州，以冀是帝都，河為大患，故先從冀起，而次治兗。若使冀州之水東入兗州，水無去處，治之無益，雖是帝都，不得先也。此經大體每州之始先言山川，後言平地。青州、梁州先山後川，徐州、雍州先川后山，兗、揚、荊、豫有川無山，揚、豫不言平地，冀州田賦之下始言“恆、衛既從”，史以大略為文，不為例也。每州之下言水路相通，通向帝都之道，言\CJKunderline{禹}每州事了，入朝以白帝也。 \par}

壺口治梁及岐。\footnote{壺口在冀州，梁、岐在雍州，從東循山治水而西。壺音胡，馬云：“壺口，山名。”治如字。岐，其宜反。雍，於用反，後州名同。案，此句或斷為：“冀州，既載壺口，治梁及岐。”載，事也，施工也。}

{\noindent\zhuan\zihao{6}\fzbyks 傳“壺口”至“而西”。正義曰：\CJKunderwave{史記}稱高祖入咸陽，蕭何先收圖籍，則秦焚詩書，圖籍皆在。\CJKunderline{孔君}去漢初七八十年耳，身為武帝博士，必當具見圖籍，其山川所在，必是驗實而知。“壺口在冀州，梁、岐在雍州”，當時疆界為然也。此於冀州之分，言及雍州之山者，“從東循山治水而西”故也。鄭云：“於此言‘治梁及岐’者，蓋治水從下起，以襄水害易也。”班固作\CJKunderwave{漢書·地理志}，據前漢郡縣言山川所在。\CJKunderwave{志}雲壺口在河東北屈縣東南。應劭云：“已有南屈，故稱北屈。”梁山在左馮翊夏陽縣西北，岐山在右扶風美陽縣西北，然則壺口西至梁山，梁山西至岐山，從東而向西言之也。經於“壺口”之下言“治”者，孔意蓋雲欲見上下皆治也。 \par}

既修太原，至於岳陽。\footnote{高平曰太原，今以為郡名。嶽，太嶽，在太原西南。山南曰陽。嶽，字又作嶽。太嶽，山名。陽，山南曰陽，水北亦曰陽。}

{\noindent\zhuan\zihao{6}\fzbyks 傳“高平”至“曰陽”。正義曰：“太原”，原之大者，\CJKunderwave{漢書}以為郡名，傳欲省文，故云“高平曰太原，今以為郡名”，即晉陽縣是也。\CJKunderwave{釋地}云：“廣平曰原，高平曰陸。”孔以太原地高，故言“高平”，其地高而廣也。下文導山雲“壺口、雷首至於太嶽”，知此“嶽”即太嶽也。屬河東郡，在太原西南也。\CJKunderwave{地理志}河東彘縣東有霍太山，此彘縣周厲王所奔，順帝改為永安縣，\CJKunderwave{周禮·職方氏}冀州其山鎮曰霍山，即此太嶽是也。山南見日，故“山南曰陽”。此說循理平地，言從太原至嶽山之南，故云“岳陽”也。 \par}

覃懷厎績，至於衡漳。\footnote{覃懷，近河地名。漳水橫流入河，從覃懷致功至橫漳。覃，徒南反。厎，之履反。衡如字，橫也；馬云：“水名。”漳音章。近，附近之近。}

{\noindent\zhuan\zihao{6}\fzbyks 傳“覃懷”至“衡漳”。正義曰：\CJKunderwave{地理志}河內郡有懷縣，在河之北，蓋“覃懷”二字共為一地，故云“近河地名”。“衡”即古“橫”字，漳水橫流入河，故云“橫漳”。漳在懷北五百餘里，從覃懷致功而北至橫漳也。\CJKunderwave{地理志}云，清漳水出上黨沾縣大黽谷，東北至渤海阜城縣入河，過郡五，行千六百八十里，此沾縣因水為名。\CJKunderwave{志}又云：“沾水出壺關。\CJKunderwave{志}又云，濁漳水出長子縣，東至鄴縣入清漳。\CJKunderline{鄭玄}亦云：“橫漳，漳水橫流。”王肅云：“衡、漳，二水名。” \par}

厥土惟白壤，\footnote{無塊曰壤，水去土復，其性色白而壤。壤，若丈反。馬云：“天性和美也。”塊，苦對反。}

{\noindent\zhuan\zihao{6}\fzbyks 傳“無塊”至“而壤”。正義曰：\CJKunderwave{九章算術}“穿地四，為壤五。壤為息土”，則“壤”是土和緩之名，故云“無塊曰壤”。此土本色為然，水去土復其性，色白而壤。雍州色黃而壤，豫州直言“壤”，不言其色，蓋州內之土不純一色，故不得言色也。 \par}

厥賦惟上上，錯，\footnote{賦謂土地所生，以供天子。上上，第一。錯，雜，雜出第二之賦。上如字，賦第一。錯，倉各反，馬云：“地有上下相錯，通率第一。”供音恭。}

{\noindent\zhuan\zihao{6}\fzbyks 傳“賦謂”至“之賦”。正義曰：以文承“厥土”之下，序雲“任土作貢”，又“賦”者稅斂之名，往者洪水為災，民皆墊溺，九州賦稅蓋亦不行，水災既除，土複本性，以作貢賦之差，故云“賦謂土地所生，以供天子”。謂稅谷以供天子，\CJKunderline{鄭玄}雲“此州入谷不貢”是也。因九州差為九等，“上上”是第一也。交錯是間雜之義，故“錯”為雜也。顧氏雲“上上之下即次上中”，故云“雜出第二之賦”也。\CJKunderwave{孟子}稱稅什一為正；輕之於堯舜，為大貊小貊；重之於堯舜，為大桀小桀；則此時亦什一。稅俱什一，而得為九等差者，人功有強弱，收穫有多少。傳以荊州“田第八，賦第三”，為“人功修”也，雍州“田第一，賦第六”，為“人功少”也，是據人功多少總計以定差。此州以上上為正，而雜為次等，言出上上時多，而上中時少也。多者為正，少者為雜，故云“第一”。此州言“上上錯”者，少在正下，故先言“上上”，而後言“錯”。豫州言“錯上中”者，少在正上，故先言“錯”，而後言“上中”。揚州雲“下上上錯”，不言“錯下上”者，以本設九等，分三品為之上中下，下上本是異品，故變文言“下上上錯”也。梁州雲“下中三錯”者，梁州之賦凡有三等，其出下中時多，故以“下中”為正，上有下上，下有下下，三等雜出，故言“三錯”，是明雜有下上、下下可知也。此九等所較無多，諸州相準為等級耳。此計大率所得，非上科定也。但治水據田責其什一，隨上豐瘠，是上之任土，而下所獻自有差降，即以差等為上之定賦也。然一升一降，不可常同。冀州自出第二，與豫州同，時則無第一之賦。豫州與冀州等一同，時則無第二之賦。或容如此，事不可恆。\CJKunderline{鄭玄}云：“賦之差，一井，上上出九夫稅，下下出一夫稅，通率九州一井稅五夫。”如鄭此言，上上出稅，九倍多於下下。鄭\CJKunderwave{詩箋}云：“井稅一夫，其田百畝。”若上上一井稅一夫，則下下九井乃出一夫，稅太少矣；若下下井稅一夫，則上上全入官矣，豈容輕重頓至是乎？ \par}

厥田惟中中。\footnote{田之高下肥瘠，九州之中為第五。中，丁仲反，又如字。中，馬云：“土地有高下。”肥，符非反。瘠，在亦反。}

{\noindent\zhuan\zihao{6}\fzbyks 傳“田之”至“第五”。正義曰：\CJKunderline{鄭玄}雲“田著高下之等者，當為水害備也”，則鄭謂地形高下為九等也。王肅雲“言其土地各有肥瘠”，則肅定其肥瘠以為九等也。如鄭之義，高處地瘠，出物既少，不得為上。如肅之義，肥處地下，水害所傷，出物既少，不得為上。故孔雲“高下肥瘠”，共相參對，以為九等。上言“敷土”，此言“厥田”，“田”、“土”異者，\CJKunderline{鄭玄}云：“地當陰陽之中，能吐生萬物者曰土。據人功作力競得而田之，則謂之田。”“田”、“土”異名，義當然也。 \par}

恆、衛既從，大陸既作。\footnote{二水已治，從其故道，大陸之地巳可耕作。從，才容反。}

{\noindent\zhuan\zihao{6}\fzbyks 傳“二水”至“耕作”。正義曰：二水泛溢漫流已治，從其故道故，今已可耕作也。青州“濰、淄其道”，與此“恆、衛既從”同，是從故道也。荊州“雲土、夢作乂”，與此“大陸既作”同，是水治可耕作也。其文不同，史異辭耳，無義例也。“壺口”與雍州之山連文，故傳言“壺口在冀州”。此無所嫌，故不言在冀州，以下皆如此也。\CJKunderwave{地理志}云，恆水出常山上曲陽縣，東入滱水。衛水出常山靈壽縣，東北入滹。大陸在鉅鹿縣北。\CJKunderwave{釋地}“十藪”云：“晉有大陸。”孫炎等皆云：“今鉅鹿縣北廣阿澤也。”郭璞云：“廣阿，猶大陸，以地名言之。”近為是也。\CJKunderwave{春秋}：“魏獻子畋於大陸，焚焉，還，卒於寧。”杜氏\CJKunderwave{春秋}說云，嫌鉅鹿絕遠，以為汲郡修武縣吳澤也。寧即修武也。然此二澤相去其遠，所以得為“大陸”者，以\CJKunderwave{爾雅}“廣平曰陸”，但廣而平者則名大陸，故異所而同名焉。然此二澤地形卑下，得以廣平為陸者，澤雖卑下，旁帶廣平之地，故統名焉。故大陸，澤名；廣河，以旁近大陸故也。 \par}

島夷皮服，\footnote{海曲謂之島。居島之夷還服其皮，明水害除。○島，當老反，馬云：“島夷，北夷國。”}

{\noindent\zhuan\zihao{6}\fzbyks 傳“海曲”至“害除”。正義曰：孔讀“鳥”為“島”。島是海中之山，\CJKunderwave{九章算術}所云“海島邈絕不可踐量”是也。傳雲“海曲謂之島”，謂其海曲有山。夷居其上，此居島之夷，常衣鳥獸之皮，為遭洪水，衣食不足，今還得衣其皮服，以明水害除也。\CJKunderline{鄭玄}云：“鳥夷，東方之民，搏食鳥獸者也。”王肅云：“鳥夷，東北夷國名也。”與孔不同。 \par}

夾右碣石，入於河。\footnote{碣石，海畔山。\CJKunderline{禹}夾行此山之右，而入河逆上。此州帝都，不說境界，以餘州所至則可知。先賦後田，亦殊於餘州。不言貢篚，亦差於餘州。夾音協，注同帶也。碣，其列反，韋昭其逝反。上，時掌反。篚，方尾反。}

{\noindent\zhuan\zihao{6}\fzbyks 傳“碣石”至“餘州”。正義曰：\CJKunderwave{地理志}碣石山在北平驪城縣西南，是碣石為海畔山也。鄭云：“\CJKunderwave{戰國策}碣石在九門縣，今屬常山郡，蓋別有碣石與此名同。今驗九門無此山也。”下文“導河入於海”，傳云：“入於渤海。”渤海之郡當以此海為名。計渤海北距碣石五百餘里，河入海處遠在碣石之南，\CJKunderline{禹}行碣石不得入於河也。蓋遠行通水之處，北盡冀州之境，然後南迴入河而逆上也。“夾右”者，孔雲“夾行此山之右”，則行碣石山西，南行入河，在碣石之右，故云“夾右”也。顧氏亦云：“山西曰右。”\CJKunderline{鄭玄}云：“\CJKunderline{禹}由碣石山西北行，盡冀州之境，還從山東南行，入河。”鄭以北行則東為右，南行西為右，故夾山兩旁，山常居右，與孔異也。“梁州”傳云：“浮東渡河而還帝都，白所治也。”則入河逆上，為還都白所治也。\CJKunderline{禹}之治水，必每州巡行，度其形勢，計其人功，施設規模，指授方略，令人分佈並作，還都白帝所治。於時帝都近河，故於每州之下皆言浮水達河，記\CJKunderline{禹}還都之道也。冀、兗、徐、荊、豫、梁、雍州各自言河，惟青、揚二州不言河耳。兗州雲“浮於濟、漯，達於河”，故青州直雲“達於濟”；徐州雲“浮於淮、泗，達於河”，故揚州雲“達於淮、泗”；皆記\CJKunderline{禹}入河之道也。王肅云：“凡每州之下說諸治水者，\CJKunderline{禹}功主於治水，故詳記其所治之州往還所乘涉之水名。”肅雖不言“還都白帝”，亦謂為治水，故浮水也。\CJKunderline{鄭玄}以為“治水既畢，更復行之，觀地肥瘠，定貢賦上下”．其意與孔異也。八州皆言境界，而此獨無，故解之，“此州帝都，不說境界，以餘州所至則可知”也。兗州雲“濟、河”，自東河以東也。豫州雲“荊、河”，自南河以南也。雍州雲“西河”，自西河以西也。明東河之西，西河之東，南河之北是冀州之境也。馬、鄭皆雲“冀州不書其界者，時帝都之，使若廣大然”。文既局以州名，復何以見其廣大？是妄說也。又解餘州先田後賦，此州先賦後田，亦如境界，殊於餘州也。言“殊”者，當為田賦以收穫為差，田以肥瘠為等。若田在賦上，則賦宜從田，田美則宜賦重，無以見人功修否，故令賦先於田也。以見賦由人功，此州既見此理，餘州從而可知，皆令賦在田下，欲見賦從田出，為此故殊於餘州也。\CJKunderline{鄭玄}云：“此州入谷不貢。”下雲“五百里甸服”，傳雲“為天子服治田”，是田入谷，故不獻貢篚，差異於餘州也。甸服止方千里，冀之北土境界甚遙，遠都之國，必有貢篚，舉大略而言也。 \par}

濟、河惟\textcolor{red}{兗州}。\footnote{東南據濟，西北距河。濟，子禮反，下同。兗，悅轉反。}

{\noindent\zhuan\zihao{6}\fzbyks 傳“東南”至“距河”。正義曰：此下八州，發首言山川者，皆謂境界所及也。“據”謂跨之，“距”，至也。濟、河之間相去路近，兗州之境，跨濟而過，東南越濟水，西北至東河也。李巡注\CJKunderwave{爾雅}解州名云：“兩河間其氣清，性相近，故曰冀。冀，近也。濟、河間其氣專，體性信謙，故云兗。兗，信也。淮海間其氣寬舒，稟性安徐，故曰徐。徐，舒也。江南其氣燥勁，厥性輕揚，故曰揚。揚，輕也。荊州其氣燥剛，稟性彊梁，故曰荊。荊，強也。河南其性安舒，厥性寬豫，故曰豫。豫，舒也。河西其氣蔽壅，受性急兇，故云雍。雍，壅也。”\CJKunderwave{爾雅}“九州”無樑、清，故李巡不釋，所言未必得其本也。 \par}

九河既道，\footnote{河水分為九道，在此州界，平原以北是。九河，徒駭一，太史二，馬頰三，覆釜四，胡蘇五，簡六，絜七，鉤盤八，鬲津九；出\CJKunderwave{爾雅}。}

{\noindent\zhuan\zihao{6}\fzbyks 傳“河水”至“北是”。正義曰：河自大陸之北敷為九河，謂大陸在冀州，嫌九河亦在冀州，故云“在此州界”也。河從大陸東畔北行，而東北入海。冀州之東境，至河之西畔。水分大河，東為九道，故知在兗州界平原以北是也。\CJKunderwave{釋水}載“九河”之名云：“徒駭、太史、馬頰、覆釜、胡蘇、簡、絜、鉤盤、鬲津。”李巡曰：“徒駭，\CJKunderline{禹}疏九河，以徒眾起，故云徒駭。太史，\CJKunderline{禹}大使徒眾通其水道，故曰太史。馬頰，河勢上廣下狹，狀如馬頰也。覆釜，水中多渚，往往而處，形如覆釜。胡蘇，其水下流，故曰胡蘇。胡，下也。蘇，流也。簡，大也，河水深而大也。絜，言河水多山石，治之苦絜。絜，苫也。鉤盤，言河水曲如鉤，屈折如盤也。鬲津，河水狹小，可鬲以為津也。”孫炎曰：“徒駭，\CJKunderline{禹}疏九河，用功雖廣，眾懼不成，故曰徒駭。胡蘇，水流多散胡蘇然。”其餘同李巡。郭璞云：“徒駭今在成平。東光縣今有胡蘇亭。”覆釜之名同李巡，餘名皆雲其義未詳。計\CJKunderline{禹}陳九河，雲復其故道，則名應先有，不宜“徒駭”、“太史”因\CJKunderline{禹}立名，此郭氏所以未詳也。或九河雖舊有名，至\CJKunderline{禹}治水，更別立名，即\CJKunderwave{爾雅}所云是也。\CJKunderwave{漢書·溝洫志}成帝時，河堤都尉許商上書曰：“古記九河之名，有徒駭、胡蘇、鬲津，今見在成平、東光、鬲縣界中。自鬲津以北至徒駭，其間相去二百餘里。”是知九河所在，徒駭最北，鬲津最南。蓋徒駭是河之本道，東出分為八枝也。許商上言三河，下言三縣，則徒駭在成平，胡蘇在東光，鬲津在鬲縣，其餘不復知也。\CJKunderwave{爾雅}“九河”之次，從北而南。既知三河之處，則其餘六者，太史、馬頰、覆釜在東光之北，成平之南；簡、絜、鉤盤在東光之南，鬲縣之北也。其河填塞，時有故道。\CJKunderline{鄭玄}云：“周時齊桓公塞之，同為一河。今河間弓高以東，至平原鬲津，往往有其遺處。”\CJKunderwave{春秋緯·寶乾圖}云：“移河為界在齊呂，填閼八流以自廣。”\CJKunderline{鄭玄}蓋據此文為“齊桓公塞之”也。言閼八流拓境，則塞其東流八枝，並使歸於徒駭也。 \par}

雷夏既澤，灉沮會同。\footnote{雷夏，澤名。灉、沮，二水，會同此澤。灉，徐音邕，王於用反。沮，七餘反。}

{\noindent\zhuan\zihao{6}\fzbyks 傳“雷夏”至“此澤”。正義曰：洪水之時，高原亦水，澤不為澤。“雷夏既澤”，高地水盡，此復為澤也。於“澤”之下言“灉、沮會同”，謂二水會合而同入此澤也。\CJKunderwave{地理志}云，雷澤在濟陰城陽縣西北。 \par}

桑土既蠶，是降丘宅土。\footnote{地高曰丘。大水去，民下丘，居平土，就桑蠶。蠶，在南反。}

{\noindent\zhuan\zihao{6}\fzbyks 傳“地高”至“桑蠶”。正義曰：\CJKunderwave{釋丘}云：“非人為之丘。”孫炎曰：“地性自然也。”是“地高曰丘”也。“降丘宅土”與“既蠶”連文，知“下丘，居平土，就桑蠶”也。計下丘居土，諸處皆然，獨於此州言之者，\CJKunderline{鄭玄}云：“此州寡于山，而夾川兩大流之間，遭洪水，其民尤困。水害既除，於是下丘居土，以其免於厄，尤喜，故記之。 \par}

{\noindent\shu\zihao{5}\fzkt “桑土”至“宅土”。正義曰：宜桑之土既得桑養蠶矣。洪水之時，民居丘上，於是得下丘陵，居平土矣。 \par}

厥土黑墳，\footnote{色黑而墳起。墳，扶粉反，後同；韋昭音勃僨反，起也；馬云：“有膏肥也。”}厥草惟繇，厥木惟條。\footnote{繇，茂。條，長也。繇音遙，馬云：“抽也。”}

{\noindent\zhuan\zihao{6}\fzbyks 傳“繇，茂。條，長也”。正義曰：“繇”是茂之貌，“條”是長之體，言草茂而木長也。九州惟此州與徐揚三州言草木者，三州偏宜之也。宜草木，則地美矣。而田非上者，為土下溼故也。 \par}

厥田惟中下，\footnote{田第六。}厥賦貞，\footnote{貞，正也。州第九，賦正與九相當。}

{\noindent\zhuan\zihao{6}\fzbyks 傳“貞正”至“相當”。正義曰：\CJKunderwave{周易}彖、象皆以“貞”為正也。諸州賦無下下，“貞”即下下，為第九也。此州治水最在後畢，州為第九成功，其賦亦為第九，列賦於九州之差，與第九州相當，故變文為“貞”，見此意也。 \par}

作十有三載，乃同。\footnote{治水十三年，乃有賦法，與他州同。載，馬、鄭本作年。}

{\noindent\zhuan\zihao{6}\fzbyks 傳“治水”至“州同”。正義曰：“作”者，役功作務，謂治水也。治水十三年，乃有賦法，始得貢賦，與他州同也。他州十二年，此州十三年，比於他州最在後也。\CJKunderwave{堯典}言\CJKunderline{鯀}治水九載，績用不成，然後堯命得舜，舜乃舉\CJKunderline{禹}治水，三載功成，堯即禪舜。此言“十三載”者，並\CJKunderline{鯀}九載數之。\CJKunderwave{祭法}雲“\CJKunderline{禹}能修\CJKunderline{鯀}之功”，明\CJKunderline{鯀}已加功，而\CJKunderline{禹}因之也。此言“十三載”者，記其治水之年，言其水害除耳，非言十三年內皆是\CJKunderline{禹}之治水施功也。馬融曰：“\CJKunderline{禹}治水三年，八州平，故堯以為功而禪舜。”是十二年而八州平，十三年而兗州平，兗州平在舜受終之年也。 \par}

厥貢漆絲，厥篚織文。\footnote{地宜漆林，又宜桑蠶。織文，錦綺之屬。盛之篚篚而貢焉。漆音七。盛音成。}

{\noindent\zhuan\zihao{6}\fzbyks 傳“地宜”至“貢焉”。正義曰：任土作貢，此州貢漆，知“地宜漆林”也。\CJKunderwave{周禮·載師}雲“漆林之徵”，故以“漆林”言之。“綺”是織繒之有文者，是綾錦之別名，故云“錦綺之屬”，皆是織而有文者也。“篚”是入貢之時盛在於篚，故云“盛之篚篚而貢焉”。\CJKunderline{鄭玄}云：“貢者百功之府受而藏之。其實於篚者，入於女功，故以貢篚別之。”歷檢篚之所盛，皆供衣服之用，入於女功，如鄭言矣。“檿絲”中琴瑟之弦，亦是女功所為也。“織貝”，\CJKunderline{鄭玄}以為織如貝文，傳謂“織為細紵。貝為水物”，則貝非服飾所須，蓋恐其損缺，故以筐篚盛之也。諸州無“厥篚”者，其諸州無入篚之物，故不貢也。漢世陳留襄邑縣置服官，使製作衣服，是兗州綾錦美也。 \par}

浮於濟、漯，達於河。\footnote{順流曰浮。濟、漯兩水名。因水入水曰達。漯，天答反，\CJKunderwave{篇韻}作他合反。}

{\noindent\zhuan\zihao{6}\fzbyks 傳“順流”至“曰達”。正義曰：\CJKunderwave{地理志}云，漯水出東郡東武陽縣，至樂安千乘縣入海，過郡三，行千二十里。其濟則下文具矣，是“濟、漯”為二水名也。言“因水入水曰達當”，謂從水入水，不須舍舟而陸行也。揚州雲“沿於江海，達於淮泗”，傳雲“沿江入海，自海入淮，自淮入泗”，是言水路相通，得乘舟經達也。案青州雲“浮於汶，達於濟”，經言濟會於汶，浮汶則達濟也。此雲“浮於濟、漯，達於河”，從漯入濟，自濟入河徐州雲“浮於淮泗，達於河”，蓋以徐州北接青州，既浮淮泗，當浮汶入濟，以達於河也。 \par}

海、岱惟\textcolor{red}{青州}。\footnote{東北據海，西南距岱。岱音代，泰山也。}

{\noindent\zhuan\zihao{6}\fzbyks 傳“東北”至“距岱”。正義曰：海非可越而言“據”者，東萊東境之縣，浮海入海曲之間，青州之境，非至海畔而已，故言“據”也。漢末有公孫度者，竊據遼東，自號青州刺史，越海收東萊諸郡。堯時青州當越海而有遼東也。舜為十二州，分青州為營州，營州即遼東也。 \par}

嵎夷既略，濰、淄其道。\footnote{嵎夷，地名。用功少曰略。濰、淄二水復其故道。嵎音隅。濰音惟，本亦作惟，又作維。淄，側其反。}

{\noindent\zhuan\zihao{6}\fzbyks 傳“嵎夷”至“故道”。正義曰：“嵎夷，地名”，即\CJKunderwave{堯典}“宅嵎夷”是也。“嵎夷”、“萊夷”、“和夷”為地名，“淮夷”為水名，“島夷”為狄名，皆觀文為說也。“略”是簡易之義，故“用功少為略”也。\CJKunderwave{地理志}云，濰水出琅邪箕屋山，北至都昌縣入海，過郡三，行五百二十里。淄水出泰山萊蕪縣原山，東北至千乘博昌縣入海。 \par}

厥土白墳，海濱廣斥。\footnote{濱，涯也。言復其斥鹵。濱，必人反。斥，徐音尺。\CJKunderwave{說文}云：“東方謂之斥，西方謂之滷。”鄭云：“斥謂地咸鹵。”涯，魚佳反。}

{\noindent\zhuan\zihao{6}\fzbyks 傳“濱涯”至“斥鹵”。正義曰：“濱，涯”，常訓也。\CJKunderwave{說文}云：“滷，咸地也。東方謂之斥，西方謂之滷。”海畔迥闊，地皆斥鹵，故云“廣斥”，言水害除，復舊性也。 \par}

厥田惟上下，厥賦中上。\footnote{田第三，賦第四。}厥貢鹽、絺，海物惟錯。\footnote{絺，細葛。錯，雜，非一種。鹽，餘佔反。絺，敕其反。種，章勇反。}岱畎絲、枲、鉛、松、怪石。\footnote{畎，谷也。怪，異；好石似玉者。岱山之谷，出此五物，皆貢之。畎，工犬反，徐本作畎谷。枲，思似反。鉛，寅專反，字從㕣；㕣音以選反。怪如字。怪石，碔砆之屬。}

{\noindent\zhuan\zihao{6}\fzbyks 傳“畎谷”至“貢之”。正義曰：\CJKunderwave{釋水}云：“水注川曰谿，注谿曰谷。”谷是兩山之閒流水之道，“畎”言畎去水，故言“谷”也。“怪石”，奇怪之石，故云“好石似玉”也。“枲”，麻也。“鉛”，錫也。岱山之谷有此五物，美於他方所有，故貢之也。 \par}

萊夷作牧，\footnote{萊夷，地名。可以放牧。萊音來。牧，牧養之牧，徐音目，一音茂，注同。}厥篚檿絲。\footnote{檿桑蠶絲，中琴瑟弦。檿，烏簟反，山桑也。}

{\noindent\zhuan\zihao{6}\fzbyks 傳“檿桑”至“瑟弦”。正義曰：\CJKunderwave{釋木}云：“檿桑，山桑。”郭璞曰：“柘屬也。”“檿絲”，是蠶食檿桑，所得絲韌，中琴瑟弦也。 \par}

浮於汶，達於濟。

{\noindent\shu\zihao{5}\fzkt “浮於汶”。正義曰：\CJKunderwave{地理志}云，汶水出泰山萊蕪縣原山，西南入濟也。 \par}

海、岱及淮惟\textcolor{red}{徐州}。\footnote{東至海，北至岱，南及淮。汶音問。}淮、沂其乂，蒙、羽其藝。\footnote{二水已治，二山已可種藝。沂，魚依反，水名。藝，魚世反。}

{\noindent\zhuan\zihao{6}\fzbyks 傳“二水”至“種藝”。正義曰：“乂”訓治也，故云“二水已治”。\CJKunderwave{地理志}云，沂水出泰山，蓋縣臨樂子山南至下邳入泗，過郡五，行六百里。淮出桐柏山，發源遠矣，於此州言之者，淮水至此而大，為害尤甚，喜得其治，故於此記之。\CJKunderwave{地理志}云，蒙山在泰山蒙陰縣西南。羽山在東海祝其縣南。\CJKunderwave{詩}雲“藝之荏菽”，故“藝”為種也。 \par}

大野既豬，東原底平。\footnote{大野，澤名。水所停曰豬。東原致功而平，言可耕。豬，張魚反，馬云：“水所停止深者曰豬。”劉東胡反。}

{\noindent\zhuan\zihao{6}\fzbyks 傳“大野”至“可耕”。正義曰：\CJKunderwave{地理志}云，大野澤在山陽鉅野縣北。“巨”即大也。\CJKunderwave{檀弓}云：“汙其宮而豬焉。”又澤名孟豬，停水處也，故云“水所停曰豬”。往前漫溢，今得豬水為澤也。“東原”即今之東平郡也。致功而地平，言其可耕也。 \par}

厥土赤埴墳，草木漸包。\footnote{土黏曰埴。漸，進長。包，叢生。埴，市力反，鄭作戠，徐、鄭、王皆讀曰熾，韋昭音試。漸如字，本又作蔪，\CJKunderwave{字林}才冉反，草之相包裹也。包，必茅反，字或作苞，非叢生也，馬云：“相包裹也。”黏，女佔反。長，丁丈反。叢，才公反。}

{\noindent\zhuan\zihao{6}\fzbyks 傳“土黏”至“叢生”。正義曰：“戠”“埴”、音義同。\CJKunderwave{考工記}用土為瓦，謂之“摶埴之工”，是“埴”謂黏土，故“土黏曰埴”。\CJKunderwave{易·漸卦}彖云：“漸，進也。”\CJKunderwave{釋言}云：“苞，稹也。”孫炎曰“物叢生曰苞，齊人名曰稹”。郭璞曰：“今人呼叢致者為稹。”“漸苞”謂長進叢生，言其美也。 \par}

厥田惟上中，厥賦中中。\footnote{田第二，賦第五。}厥貢惟土五色，\footnote{王者封五色土為社，建諸侯則各割其方色土與之，使立社。燾以黃土，苴以白茅，茅取其潔，黃取王者覆四方。燾，徒報反，覆也。苴，子餘反，包裹也。}

{\noindent\zhuan\zihao{6}\fzbyks 傳“王者”至“四方”。正義曰：傳解貢土之意，王者封五色土以為社，若封建諸侯，則各割其方色土與之，使歸國立社。其土燾以黃土。燾，覆也。四方各依其方色，皆以黃土覆之。其割土與之時，苴以白茅，用白茅裹土與之。必用白茅者，取其絜清也。\CJKunderwave{易}稱“藉用白茅”，茅色白而絜美。\CJKunderwave{韓詩外傳}云：“天子社廣五丈，東方青，南方赤，西方白，北方黑，上冒以黃土。將封諸侯，各取其方色土，苴以白茅，以為社。明有土謹敬絜清也。”蔡邕\CJKunderwave{獨斷}云：“天子大社，以五色土為壇。皇子封為王者，授之大社之土，以所封之方色苴以白茅，使之歸國以立社，謂之茅社。”是必古書有此說，故先儒之言皆同也。 \par}

羽畎夏翟，嶧陽孤桐，\footnote{夏翟，翟，雉名。羽中旌旄，羽山之谷有之。孤，特也。嶧山之陽特生桐，中琴瑟。夏，行雅反。翟，徒歷反。嶧音亦，一音夕。}

{\noindent\zhuan\zihao{6}\fzbyks 傳“夏翟”至“琴瑟”。正義曰：\CJKunderwave{釋鳥}云：“翟，山雉。”此言“夏翟”，則“夏翟”共為雉名。\CJKunderwave{周禮}立夏採之官，取此名也。\CJKunderwave{周禮·司常}云：“全羽為旞，析羽為旌。”用此羽為之，故云“羽中旌旄”也。\CJKunderwave{地理志}云，東海下邳縣西有葛嶧山，即此山也。 \par}

泗濱浮磬，淮夷蠙珠暨魚。\footnote{泗，水涯。水中見石，可以為磬。蠙珠，珠名。淮、夷二水出蠙珠及美魚。泗音四，水名。淮夷，鄭云：“淮水之夷民也。”馬云：“淮、夷，二水名。”孔傳雲“淮夷之水”，本亦有作“淮夷二水”也。蠙音蒲邊反，徐扶堅反，字又作蚍，韋昭薄迷反，蚌也。暨，其器反。見，賢遍反。}

{\noindent\zhuan\zihao{6}\fzbyks 傳“泗水”至“美魚”。正義曰：泗水旁山而過，石為泗水之涯，石在水旁，水中見石，似若水中浮然，此石可以為磬，故謂之“浮磬”也。貢石而言磬者，此石宜為磬，猶如“砥礪”然也。蠙是蚌之別名，此蚌出珠，遂以蠙為珠名。蠙之與魚皆是水物，而以“淮夷”冠之，知“淮夷”是二水之名。“淮”即四瀆之淮也。“夷”蓋小水，後來竭涸，不復有其處耳。王肅亦以“淮夷”為水名。\CJKunderline{鄭玄}以為“淮水之上，夷民獻此珠與魚”也。\CJKunderwave{地理志}泗水出濟陰乘氏縣，東南至臨淮雎陵縣入淮，行千一百一十里也。 \par}

厥篚玄纖縞。\footnote{玄，黑繒。縞，白繒。纖，細也。纖在中，明二物皆當細。纖，息廉反。縞，古老反，徐古到反。繒，似陵反。}

{\noindent\zhuan\zihao{6}\fzbyks 傳“玄黑”至“當細”。正義曰：篚之所盛，例是衣服之用，此單言“玄”，玄必有質，玄是黑色之別名，故知玄是黑繒也。\CJKunderwave{史記}稱高祖為義帝發喪，諸侯皆縞素，是縞為白繒也。 \par}

浮於淮、泗，達於河。\footnote{河如字，\CJKunderwave{說文}作菏，工可反，云：“水出山陽湖陵南。”}

淮、海惟\textcolor{red}{揚州}。\footnote{北據淮，南距海。}彭蠡既豬，陽鳥攸居。\footnote{彭蠡，澤名。隨陽之鳥，鴻雁之屬，冬月所居於此澤。蠡音禮，張勃\CJKunderwave{吳錄}云：“今名洞庭湖。”案今在九江郡界。}

{\noindent\zhuan\zihao{6}\fzbyks 傳“彭蠡”至“此澤”。正義曰：“彭蠡”是江漢合處，下雲“導漾水，南入於江，東匯為彭蠡”是也。日之行也，夏至漸南，冬至漸北，鴻雁之屬，九月而南，正月而北，左思\CJKunderwave{蜀都賦}所云“木落南翔，冰泮北徂”是也。日，陽也，此鳥南北與日進退，隨陽之鳥，故稱“陽鳥”，冬月所居於此彭蠡之澤也。 \par}

三江既入，震澤厎定。\footnote{震澤，吳南大湖名。言三江已入，致定為震澤。三江，韋昭云：“謂吳松江、錢唐江、浦陽江也。”\CJKunderwave{吳地記}云：“松江東北行七十里，得三江口，東北入海為婁江，東南入海為東江，並松江為三江。震澤，吳都太湖。厎，之履反，致也，\CJKunderwave{史記}音致。大湖，音太湖。}

{\noindent\zhuan\zihao{6}\fzbyks 傳“震澤”至“震澤”。正義曰：\CJKunderwave{地理志}云，會稽吳縣，故周泰伯所封國也。具區在西，古文以為震澤，是“吳南大湖名”。蓋縣治居澤之東北，故孔傳言“南”，\CJKunderwave{志}言“西”。大澤畜水，南方名之曰湖，三江既入此湖也。治水致功，今江入此澤，故“致定為震澤”也。下傳云：“自彭蠡江分為三，入震澤，遂為北江而入海。”是孔意江從彭蠡而分為三，又共入震澤，從震澤復分為三，乃入海。鄭云：“三江分於彭蠡，為三孔東入海。”其意言“三江既入”，入海耳，不入震澤也。又案\CJKunderwave{周禮·職方}揚州藪曰具區，浸曰五湖。五湖即震澤。若如\CJKunderwave{志}雲具區即震澤，則浸、藪為一。案餘州浸、藪皆異而揚州同者，蓋揚州浸、藪同處，論其水謂之浸，指其澤謂之藪。 \par}

篠簜既敷，\footnote{篠，竹箭。簜，大竹。水去巳布生。篠，西了反。簜，徒黨反，或作𥯕，他莽反。}

{\noindent\zhuan\zihao{6}\fzbyks 傳“篠，竹箭。簜，大竹”。正義曰：\CJKunderwave{釋草}云：“篠，竹箭。”郭璞云：“別二名也。”又云“簜，竹”。李巡曰：“竹節相去一丈曰簜。”孫炎曰：“竹闊節者曰簜。”郭璞云：“竹別名。”是篠為小竹，簜為大竹。 \par}

厥草惟夭，厥木惟喬。\footnote{少長曰夭。喬，高也。夭，於嬌反，馬云：“長也。”喬，其嬌反，徐音驕。少，詩照反。長，丁丈反。}

{\noindent\zhuan\zihao{6}\fzbyks 傳“少長曰夭。喬，高也”。正義曰：夭是少長之貌，\CJKunderwave{詩}曰：“桃之夭夭”是也。“喬，高”，\CJKunderwave{釋詁}文。\CJKunderwave{詩}曰“南有喬木”是也。 \par}

厥土惟塗泥，\footnote{地泉溼。}厥田惟下下，厥賦下上錯。\footnote{田第九，賦第七，雜出第六。}厥貢惟金三品，\footnote{金、銀、銅也。}

{\noindent\zhuan\zihao{6}\fzbyks 傳“金、銀、銅也”。正義曰：“金”既總名，而云“三品”，黃金以下惟有白銀與銅耳，故為“金、銀、銅也”。\CJKunderwave{釋器}云：“黃金謂之璗，其美者謂之鏐。白金謂之銀，其美者謂之鐐。”郭璞曰：“此皆道金銀之別名及其美者也。鏐即紫磨金也。”\CJKunderline{鄭玄}以為“金三品者，銅三色也。” \par}

瑤、琨、篠簜，\footnote{瑤、琨皆美玉。瑤音遙。琨音昆，美石也，馬本作瑻，韋昭音貫。}

{\noindent\zhuan\zihao{6}\fzbyks 傳“瑤、琨，皆美玉”。正義曰：美石似玉者也。玉、石其質相類，美惡別名也。王肅云：“瑤、琨，美石次玉者也。” \par}

齒、革、羽、毛、惟木。\footnote{齒，象牙。革，犀皮。羽，鳥羽。毛，旄牛尾。木，楩、軺、豫章。犀，細兮反。旄音毛。楩音緶，又婢善反。}

{\noindent\zhuan\zihao{6}\fzbyks 傳“齒象”至“豫章”。正義曰：\CJKunderwave{詩}雲“元龜象齒”，知“齒”是象牙也。\CJKunderwave{說文}云：“齒，口齗骨也。牙，牡齒也。”隱五年\CJKunderwave{左傳}雲“齒牙骨角”，“牙”、“齒”小別，統而名之，“齒”亦“牙”也。\CJKunderwave{考工記}：“犀甲七屬，兕甲六屬。”宣二年\CJKunderwave{左傳}云：“犀兕尚多，棄甲則那？”是甲之所用，犀革為上；革之所美，莫過於犀；知“革”是犀皮也。\CJKunderwave{說文}云：“獸皮治去其毛為革。”革與皮去毛為異耳。\CJKunderwave{說文}云：“羽，鳥長毛也。”知“羽”是鳥羽。南方之鳥，孔雀、翡翠之屬，其羽可以為飾，故貢之也。\CJKunderwave{說文}云：“犛，西南夷長旄牛也。”此犛牛之尾可為旌旗之飾，經傳通謂之“旄”。\CJKunderwave{牧誓}雲“右秉白旄”，\CJKunderwave{詩}雲“建旐設旄”，皆謂此牛之尾，故知“毛”是旄牛尾也。直雲“惟木”，不言木者，故言“楩、軺、豫章”，此三者是揚州美木，故傳舉以言之，所貢之木不止於此。 \par}

島夷卉服，\footnote{南海島夷草服葛越。卉，徐許貴反。}

{\noindent\zhuan\zihao{6}\fzbyks 傳“南海”至“葛越”。正義曰：上傳“海曲謂之島”，知此“島夷”是南海島上之夷也。\CJKunderwave{釋草}云：“卉，草。”舍人曰：“凡百草一名卉。”知“卉服”是“草服葛越”也。葛越，南方布名，用葛為之。左思\CJKunderwave{吳都賦}雲“蕉葛升越，弱於羅紈”是也。冀州雲“島夷皮服”，是夷自服皮，皮非所貢也。此言“島夷卉服”，亦非所貢也。此與“萊夷作牧”並在貢篚之間，古史立文不次也。\CJKunderline{鄭玄}云：“此州下溼，故衣草服。貢其服者，以給天子之官。”與孔異也。 \par}

厥篚織、貝。\footnote{織，細繒。貝，水物。}

{\noindent\zhuan\zihao{6}\fzbyks 傳“織細”至“水物”。正義曰：傳以“貝”非織物，而云“織貝”，則貝、織異物，織是織而為之，揚州紵之所出，此物又以篚盛之，為衣服之用，知是“細紵”，謂細紵布也。\CJKunderwave{釋魚}之篇貝有居陸居水，此州下溼，故云“水物”。\CJKunderwave{釋魚}有“玄貝，貽貝。餘貾，黃白文。餘泉，白黃文”，當貢此有文之貝以為器物之飾也。\CJKunderline{鄭玄}云：“貝，錦名。\CJKunderwave{詩}云：‘萋兮斐兮，成是貝錦。’凡為織者先染其絲，乃織之則文成矣。\CJKunderwave{禮記}曰：‘土不衣織。’”與孔異也。 \par}

厥包橘柚錫貢。\footnote{小曰橘，大曰柚。其所包裹而致者，錫命乃貢。言不常。橘，均必反。柚，由究反。裹音果。}

{\noindent\zhuan\zihao{6}\fzbyks 傳“小曰”至“不常”。正義曰：橘、柚二果，其種本別，以實相比，則柚大橘小，故云“小曰橘，大曰柚”。猶\CJKunderwave{詩傳}雲“大曰鴻，小曰雁”，亦別種也。此物必須裹送，故云其所包裹而送之。以須之有時，故待錫命乃貢，言不常也。文在“篚”下，以不常故耳。荊州“納錫大龜”，豫州“錫貢磬錯”，皆為非常，並在“篚”下。荊州言“包”，傳雲“橘柚”也，文在“篚”上者，荊州橘柚為善，以其常貢。此州則不常也。王肅云：“橘與柚錫其命而後貢之，不常入，當繼荊州乏無也。”鄭云：“有錫則貢之，此州有錫而貢之，或時無，則不貢。錫，所以柔金也。\CJKunderwave{周禮·考工記}云，攻金之工掌執金錫之齊故也。” \par}

沿於江、海，達於淮、泗。\footnote{順流而下曰沿。沿江入海，自海入淮，自淮入泗。沿，悅專反；鄭本作松，松當為沿；馬本作均，云：“均，平。”}

{\noindent\zhuan\zihao{6}\fzbyks 傳“順流”至“入泗”。正義曰：文十年\CJKunderwave{左傳}云：“沿漢溯江。”溯是逆，沿是順，故“順流而下曰沿”。“沿江入海”，順也。“自海入淮，自淮入泗”，逆也。 \par}

荊及衡陽惟\textcolor{red}{荊州}。\footnote{北據荊山，南及衡山之陽。}

{\noindent\zhuan\zihao{6}\fzbyks 傳“北據”至“之陽”。正義曰：此州北界至荊山之北，故言“據”也。“南及衡山之陽”，其境過衡山也。以衡是大山，其南無復有名山大川可以為記，故言“陽”見其南至山南也。 \par}

江、漢朝宗於海，\footnote{二水經此州而入海，有似於朝，百川以海為宗。宗，尊也。朝，直遙反。}

{\noindent\zhuan\zihao{6}\fzbyks 傳“二水”至“宗尊也”。正義曰：\CJKunderwave{周禮·大宗伯}諸侯見天子之禮，“春見曰朝，夏見曰宗”。鄭云：“朝猶朝也，欲其來之早也。宗，尊也，欲其尊王也。”“朝宗”是人事之名，水無性識，非有此義。以海水大而江、漢小，以小就大，似諸侯歸於天子，假人事而言之也。\CJKunderwave{詩}云：“沔彼流水，朝宗於海。”\CJKunderwave{毛傳}云：“水猶有所朝宗。”“朝宗”是假人事而言水也。\CJKunderwave{老子}云：“滄海所以能為百谷王者，以其下之。”是百川以海為宗。鄭云：“江水、漢水其流遄疾，又合為一，共赴海也。猶諸侯之同心尊天子而朝事之。荊楚之域，國有道則後服，國無道則先強，故記其水之義，以著人臣之禮。 \par}

九江孔殷，\footnote{江於此州界分為九道，甚得地勢之中。九江，\CJKunderwave{潯陽地記}云：“一曰烏白江，二曰蚌江，三曰烏江，四曰嘉靡江，五曰畎江，六曰源江，七曰累江，八曰提江，九曰箘江。”張須元\CJKunderwave{緣江圖}云：“一曰三里江，二曰五州江，三曰嘉靡江，四曰烏土江，五曰白蚌江，六曰白烏江，七曰箘江，八曰沙提江，九曰廩江。參差隨水長短，或百里，或五十里。始於鄂陵，終於江口，會於桑落洲。\CJKunderwave{太康地記}曰：“九江，劉歆以為湖漢九水，入彭蠡澤也。”}

{\noindent\zhuan\zihao{6}\fzbyks 傳“江於”至“之中”。正義曰：傳以“江”是此水大名，“九江”謂大江分而為九，猶大河分為九河，故言“江於此州之界分為九道”。訓“孔”為甚，“殷”為中，言“甚得地勢之中”也。鄭云：“殷猶多也。九江從山谿所出，其孔眾多，言治之難也。\CJKunderwave{地理志}九江在今廬江潯陽縣南，皆東合為大江。”如鄭此意，九江各自別源，其源非大江也，下流合於大江耳。然則江以南水無大小，俗人皆呼為江，或從江分出，或從外合來，故孔、鄭各為別解。應劭注\CJKunderwave{地理志}雲“江自潯陽分為九道”，符於孔說，\CJKunderwave{潯陽記}有九江之名：“一曰烏江，二曰蚌江，三曰烏白江，四曰嘉靡江，五曰畎江，六曰源江，七曰廩江，八曰提江，九曰箘江。”雖名起近代，義或當然。 \par}

沱、潛既道，\footnote{沱，江別名。潛，水名。皆復其故道。沱，徒河反。潛，捷廉反，馬云：“沱，湖也。其中泉出而不流者謂之潛。”}

{\noindent\zhuan\zihao{6}\fzbyks 傳“沱江”至“故道”。正義曰：下文“泯山導江，東別為沱”，是“沱”為江之別名也。經無“潛”之本源，直雲“水名”。\CJKunderwave{釋水}云，水自江出為沱，漢為潛。鄭注此，既引\CJKunderwave{爾雅}，乃云：“今南郡枝江縣有沱水，其尾入江耳，首不於江出也。華容有夏水，首出江，尾入沔蓋，此所謂沱也。潛則未聞象類。”此解荊州之沱、潛發源此州。若如鄭言，此水南流，不入荊州界，非此潛也。此下樑州注云：“二水亦謂自江漢出者。\CJKunderwave{地理志}在今蜀郡鄆縣江沱及漢中安陽皆有沱水、潛水，其尾入江漢耳，首不於此出。江源有𨞪江，首出江，南至犍為武陽又入江，豈沱之類與？蓋漢西漢，出嶓冢，東南至巴郡江州入江，行二千七百六十里。”此解梁州之沱、潛也。郭璞\CJKunderwave{爾雅音義}云：“沱水自蜀郡都水縣揃山與江別而更流。”璞又云：“有水從漢中沔陽縣南流，至梓潼漢壽入大穴中，通峒山下西南潛出，一名沔水，舊俗雲即\CJKunderwave{禹貢}潛也。”郭璞此言，亦解梁州沱、潛，與鄭又異。然\CJKunderwave{地理志}及鄭皆以荊、梁二州各有沱、潛，又郭氏所解沱、潛惟據梁州，不言荊州之沱、潛，而孔梁州注云“沱、潛發源此州，入荊州”，以二州沱、潛為一者。然彼州山水古今不可移易，孔為武帝博士，\CJKunderwave{地理志}無容不知，蓋以水從江漢出者皆曰“沱潛”，但地勢西高東下，雖於梁州合流，還從荊州分出，猶如濟水入河，還從河出，故孔舉大略為發源梁州耳。 \par}

雲土、夢作乂。\footnote{雲夢之澤在江南，其中有平土丘，水去可為耕作畎畝之治。云，徐本作雲。夢，亡弄反，一音武仲反，徐莫公反。治，直吏反。}

{\noindent\zhuan\zihao{6}\fzbyks 傳“雲夢”至“之治”。正義曰：昭三年\CJKunderwave{左傳}楚子與鄭伯田於江南之夢，是“雲夢之澤在江南”也。\CJKunderwave{地理志}南郡華容縣南有云夢澤，杜預雲“南郡枝江縣西有云夢城”，江夏安陸縣亦有云夢，或曰南郡華容縣東南有巴丘湖。江南之夢，雲夢一澤，而每處有名者，司馬相如\CJKunderwave{子虛賦}雲“雲夢者方八九百里”，則此澤跨江南北，每處名存焉。定四年\CJKunderwave{左傳}稱楚昭王寢於雲中，則此澤亦得單稱“雲”，單稱“夢”。經之“土”字在二字之間，蓋史文兼上下也。此澤既大，其內有平土，有高丘，水去可為耕作畎畝之治。 \par}

厥土惟塗泥，厥田惟下中，厥賦上下。\footnote{田第八，賦第三，人功修。}厥貢羽、毛、齒、革，惟金三品，\footnote{土所出與揚州同。}

{\noindent\zhuan\zihao{6}\fzbyks 傳“土所”至“州同”。正義曰：與揚州同，而揚州先“齒、革”，此州先“羽、毛”者，蓋以善者為先。由此而言之，諸州貢物多種，其次第皆以當州貴者為先也。 \par}

\xpinyin*{杶}、幹、\xpinyin*{栝}、柏，\footnote{幹，柘也。柏葉松身曰栝。杶，敕倫反，徐敕荀反，木名，又作櫄。幹，本又作幹，故旦反。栝，古活反。馬云：“白栝也。”柘，章夜反。}

{\noindent\zhuan\zihao{6}\fzbyks 傳“幹柘”至“曰栝”。正義曰：“幹”為弓幹，\CJKunderwave{考工記}云，弓人取幹之道也，以柘為上，知此“幹”是柘也。\CJKunderwave{釋木}云：“栝，柏葉松身。”陸機\CJKunderwave{毛詩義疏}雲“杶、㯉、栲、漆相似如一”，則杶似㯉漆也。杶、栝、柏皆木名也，以其所施多矣，柘木惟用為弓幹，弓幹莫若柘木，故舉其用也。 \par}

礪、砥、砮、丹，\footnote{砥細於礪，皆磨石也。砮，石，中矢鏃。丹，朱類。礪，力世反。砥音脂，徐之履反，韋昭音旨。砮音奴，韋昭乃固反。磨，末佐反。鏃，子木反，一音七木反。}

{\noindent\zhuan\zihao{6}\fzbyks 傳“砥細”至“朱類”。正義曰：“砥”以細密為名，“礪”以粗糲為稱，故“砥細於礪，皆磨石也”。鄭云：“礪，磨刀刃石也。精者曰砥。”\CJKunderwave{魯語}曰：“肅慎氏貢楛矢石砮。”賈逵云：“砮，矢鏃之石也。”故曰“砮，石，中矢鏃”。“丹”者，丹砂，故云“朱類”。王肅云：“丹可以為採。” \par}

惟\xpinyin*{菌}、\xpinyin*{簵}、\xpinyin*{楛}，三邦厎貢厥名。\footnote{箘、簵，美竹。楛中矢幹。三物皆出雲夢之澤，近澤三國常致貢之，其名天下稱善。箘，求隕反，韋昭一名聆風。簵音路。楛音戶，馬云：“木名，可以為箭。”\CJKunderwave{毛詩草木疏}云：“葉如荊而赤，莖似蓍。”近，附近之近。}

{\noindent\zhuan\zihao{6}\fzbyks 傳“箘簵”至“稱善”。正義曰：“箘、簵，美竹”，當時之名猶然。鄭云：“箘簵，䉁風也。”竹有二名，或大小異也，箘、簵是兩種竹也。“肅慎氏貢楛矢”，知“楛中矢幹”。“三物皆出雲夢之澤”，當時驗之猶然。經言“三邦厎貢”，知近澤三國致此貢也。文續“”，則其物特有美名，故云“其名天下稱善”。\CJKunderline{鄭玄}以“厥名”下屬“包匭菁茅”。 \par}

包\footnote{橘柚。}匭菁茅，\footnote{匭，匣也。菁以為菹，茅以縮酒。匭音軌。菁，子丁反，徐音精，馬同，鄭云：“茅有毛刺曰菁茅。”匣，胡甲反，又音甲。菹，\CJKunderwave{切韻}側魚反。縮，所六反。}

{\noindent\zhuan\zihao{6}\fzbyks 傳“橘柚”。正義曰：“包”下言“匭菁茅”，\CJKunderwave{說文}云：“匚，受物之器。象形也。凡匚之屬皆從匚。”“匱”、“匣”之字皆從匚，“匭”亦從匚，故“匭”是匣也。“菁茅”既以匭盛，非所包之物，明包必有裹也。此州所出與揚州同，揚州“厥包橘柚”，知此“包”是橘柚也。王肅云：“揚州‘厥包橘柚’，從省而可知也。” \par}

{\noindent\zhuan\zihao{6}\fzbyks 傳“匭，匣也。菁以為菹，茅以縮酒”。正義曰：“匭”是匱之別名，匱之小者。菁茅所盛，不須大匱，故用匣也。\CJKunderwave{周禮·醢人}有“菁菹”、“鹿臡”，故知“菁以為菹”。鄭云：“菁，蔓菁也。”蔓菁處處皆有，而令此州貢者，蓋以其味善也。僖四年\CJKunderwave{左傳}齊桓公責楚雲“爾貢包茅不入，王祭不供，無以縮酒”，是“茅以縮酒”也。\CJKunderwave{郊特牲}云：“縮酒用茅，明酌也。”鄭注云：“以茅縮酒也。”，\CJKunderwave{周禮·甸師}云：“祭祀供蕭茅。”鄭興云：“蕭字或為莤，莤讀為縮，束茅立之祭前，酒沃其上，酒滲下，若神飲之，故謂之縮。”杜預解\CJKunderwave{左傳}，用鄭興之說，未知誰同孔旨。特令此州貢茅，茅當異於諸處。杜預云：“茅之為異，未審也。”或雲茅有三脊，案\CJKunderwave{史記}齊桓公欲封禪，管仲睹其不可窮以辭，因設以無然之事云：“古之封禪，江淮之間，三脊茅以為藉。”此乃懼桓公耳，非荊州所有也。\CJKunderline{鄭玄}以“菁茅”為一物，“匭猶纏結也。菁茅之有毛刺者重之，故既包裹而又纏結也”。 \par}

厥篚玄纁、璣組，\footnote{此州染玄纁色善，故貢之。璣，珠類，生於水。組，綬類。纁，許雲反。璣，其依反，又音機，馬同，\CJKunderwave{說文}云：“珠不圜也。”\CJKunderwave{字書}云：“小珠也。”\CJKunderwave{玉篇}渠依、居沂二反。組音祖，馬云：“組，文也。”}

{\noindent\zhuan\zihao{6}\fzbyks 傳“此州”至“綬類”。正義曰：\CJKunderwave{釋器}云：“三染謂之纁。”李巡云：“三染其色已成為絳，纁、絳一名也。”\CJKunderwave{考工記}云：“三入為纁，五入為緅，七入為緇。”鄭云：“纁者三入而成，又再染以黑則為緅，又再染以黑則為緇。玄色在緅、緇之間，其六人者是染玄纁之法也。”此州染玄纁色善，故令貢之。\CJKunderwave{說文}雲“璣，珠不圓者”，故為“珠類”。\CJKunderwave{玉藻}說佩玉所懸者皆雲“組綬”，是組、綬相類之物也。 \par}

九江納錫大龜。\footnote{尺二寸曰大龜，出於九江水中。龜不常用，錫命而納之。馬云：“納，入也。”}

{\noindent\zhuan\zihao{6}\fzbyks 傳“尺二”至“納之”。正義曰：\CJKunderwave{史記·龜策傳}云：“龜千歲滿尺二寸”，\CJKunderwave{漢書·食貨志}雲“元龜距髯長尺二寸”，故以“尺二寸為大龜”。冠以“九江”，知“出九江水中”也。文在“篚”下而言“納錫”，是言“龜不常用，故錫命乃納之”，言此大龜錫命乃貢之也。 \par}

浮於江、沱、潛、漢，逾於洛，至於南河。\footnote{逾，越也。河在冀州，南東流，故越洛而至南河。江、沱、潛、漢，四水名。本或作“潛於漢”，非。逾，羊朱反。}

{\noindent\shu\zihao{5}\fzkt “浮於江、沱、潛、漢”。正義曰：浮此四水乃得至洛。本或“潛”下有“於”，誤耳。 \par}

荊、河惟\textcolor{red}{豫州}。\footnote{西南至荊山，北距河水。}伊、洛、瀍、澗既入於河，\footnote{伊出陸渾山，洛出上洛山，澗出沔池山，瀍出河南北山，四水合流而入河。瀍，直然反。澗，故晏反。渾音魂，又胡囷、胡昆二反。沔，亡淺反，又亡忍反，下同。陸渾、沔池二縣屬河南郡。}

{\noindent\zhuan\zihao{6}\fzbyks 傳“伊出”至“入河”。正義曰：\CJKunderwave{地理志}云，伊水出弘農盧氏縣東熊耳山，東北入洛。洛水出弘農上洛縣冢領山，東北至鞏縣入河。瀍水出河南谷城縣潛亭北，東南入洛。澗水出弘農新安縣，東南入洛。\CJKunderwave{志}與傳異者，熊耳山在陸渾縣西，冢領山在上洛縣境之內，沔池在新安縣西、谷城潛亭北，此即是河南境內之北山也。\CJKunderwave{志}詳而傳略，所據小異耳。伊、瀍、澗三水入洛合流而入河，言其不復為害也。 \par}

滎波既豬，\footnote{滎澤波水已成遏豬。zhuan，戶扃反；6澤也。波如字，馬本又播；滎播，澤名。遏，byks反。}

{\noindent\shu\zihao{5}\fzkt 傳“滎澤”至“遏豬”。正義曰：沇水入河而溢為滎，“滎”是澤名。洪水之時，此澤水大，動成波浪。此澤其時波水已成遏豬，言壅遏而為豬，畜水而成澤，不濫溢也。鄭云：“今塞為平地，滎陽民猶謂其處為滎澤，在其縣東。”言在滎澤縣之東也。馬、鄭王本皆作“滎播”，謂此澤名“滎播”。\CJKunderwave{春秋}閔二年衛侯“及狄人戰於滎澤”，不名“播”也。\CJKunderline{鄭玄}謂衛狄戰在此地，杜預雲“此滎澤當在河北，以衛敗方始渡河，戰處必在河北”。蓋此澤跨河南北，多而得名耳。 \par}

導菏澤，被孟豬。\footnote{菏澤在胡陵。孟豬，澤名，在菏東北，水流溢覆被之。導音道，下同。菏，徐音柯，又土可反，注同，韋胡阿反。被，皮寄反，徐扶義反，注同。豬，張魚反，又音諸；\CJKunderwave{左傳}及\CJKunderwave{爾雅}皆作“孟諸”，宋藪澤也。}

{\noindent\zhuan\zihao{6}\fzbyks 傳“菏澤”至“被之”。正義曰：\CJKunderwave{地理志}山陽郡有胡陵縣，不言其縣有菏澤也。又云，菏澤在濟陰定陶縣東。孟豬在梁國雎陽縣東北。以今地驗之，則胡陵在雎陽之東，定陶在雎陽之北，其水皆不流溢東北被孟豬也。然郡縣之名，隨代變易，古之胡陵當在雎陽之西北，故得東出被孟豬也。於此作“孟豬”，\CJKunderwave{左傳}、\CJKunderwave{爾雅}作“孟諸”，\CJKunderwave{周禮}作“望諸”，聲轉字異，正是一地也。 \par}

厥土惟壤，下土墳壚。\footnote{高者壤，下者壚，壚疏。壚音盧，\CJKunderwave{說文}：“黑剛土也。”}厥田惟中上，厥賦錯上中。\footnote{田第四，賦第二，又雜出第一。}厥貢漆、枲、絺、紵，厥篚纖纊，\footnote{纊，細綿。絺，敕其反。紵，直呂反。纊音曠。綿，\CJKunderwave{切韻}武延反。}

{\noindent\zhuan\zihao{6}\fzbyks 傳“纊，細綿”。正義曰：\CJKunderwave{禮·喪大記}候死者“屬纊以俟絕氣”，即“纊”是新綿耳。“纖”是細，故言“細綿”。 \par}

錫貢磬錯。\footnote{治玉石曰錯。治磬錯。}

{\noindent\zhuan\zihao{6}\fzbyks 傳“治玉”至“磬錯”。正義曰：\CJKunderwave{詩}云：“佗山之石，可以攻玉。”又曰：“可以為錯。”磬有以玉為之者，故云“治玉石曰錯”，謂“治磬錯”也。 \par}

浮於洛，達於河。

華陽、黑水惟\textcolor{red}{梁州}。\footnote{東據華山之南，西距黑水。華，胡化反，又胡瓜反。}

{\noindent\zhuan\zihao{6}\fzbyks 傳“東據”至“黑水”。正義曰：\CJKunderwave{周禮·職方氏}豫州其山鎮曰華山，在豫州界內。此梁州之境東據華山之南，不得其山，故言“陽”也。此山之西，雍州之境也。 \par}

岷、\xpinyin*{嶓}既藝，沱、潛既道。\footnote{岷山、嶓冢皆山名。水去已可種藝。沱、潛發源此州，入荊州。岷，武巾反。嶓音波，徐甫河反，韋音播。}

{\noindent\zhuan\zihao{6}\fzbyks 傳“岷山”至“荊州”。正義曰：漢制，縣有羌夷曰道。\CJKunderwave{地理志}云，蜀郡有湔道，岷山在西徼外，江水所出也。隴西郡西縣冢山，西漢水所出。是二者皆山名也。沱出於江，潛出於漢，二水發源此州而入荊州，故荊州亦云“沱、潛既道”。 \par}

蔡、蒙旅平，和夷厎績。\footnote{蔡，蒙二山名，今之峨眉山、蒙頂山。祭山曰旅。平言治功畢。和夷之地，致功可藝。旅如字，韋音盧。和如字，又作龢，鄭云：“和讀曰洹。”治，直吏反，下同。}

{\noindent\zhuan\zihao{6}\fzbyks 傳“蔡蒙”至“可藝”。正義曰：\CJKunderwave{地理志}云，蒙山在蜀郡青衣縣。應劭云：“順帝改曰漢嘉縣。”蔡山不知所在。\CJKunderwave{論語}云：“季氏旅於泰山。”是“祭山曰旅”也。“平”者言其治水畢，猶上“既藝”也。“和夷”平地之名，致功可藝。“藝”與“平”互言耳。 \par}

厥土青黎，\footnote{色青黑而沃壤。黎，鄭力兮反，徐力私反，馬云：“小疏也。”}

{\noindent\zhuan\zihao{6}\fzbyks 傳“色青黑而沃壤”。正義曰：孔以“黎”為黑，故云“色青黑”。其地“沃壤”，言其美也。王肅曰：“青，黑色。黎，小疏也。” \par}

厥田惟下上，厥賦下中三錯。\footnote{田第七，賦第八，雜出第七、第九三等。}

{\noindent\zhuan\zihao{6}\fzbyks 傳“田第”至“三等”。正義曰：傳以既言“下中”，復雲“三錯”，舉下中第八為正，上下取一，故雜出第七、第九與第八為三也。鄭云：“三錯者，此州之地有當出下之賦者少耳，又有當出下上、中下者差復益小。”與孔異也。 \par}

厥貢璆、鐵、銀、鏤、砮、磬，\footnote{璆，玉名。鏤，剛鐵。璆音蚪，徐又居蚪反，又閭幼反，馬同，韋昭、郭璞云：“紫磨金。”案郭注\CJKunderwave{爾雅}璆即紫磨金。鐵，天結反。鏤，婁豆反。}

{\noindent\zhuan\zihao{6}\fzbyks 傳“璆玉”至“剛鐵”。正義曰：\CJKunderwave{釋器}云：“璆、琳，玉也。”郭璞云：“璆、琳，美玉之別名。”“鏤”者，可以刻鏤，故為“剛鐵”也。 \par}

熊、羆、狐、貍織皮。\footnote{貢四獸之皮，織金罽。熊音雄。羆，彼宜反，如熊而黃。貍，力疑反。罽，紀例反。}

{\noindent\zhuan\zihao{6}\fzbyks 傳“貢四”至“金罽”。正義曰：與“織皮”連文，必不貢生獸，故云“貢四獸之皮”。\CJKunderwave{釋言}云：“犛，罽也。”舍人曰：“犛謂毛罽也。胡人續羊毛作衣。”孫炎曰：“毛犛為罽。”織毛而言“皮”者，毛附於皮，故以“皮”表毛耳。 \par}

西傾因桓是來，浮於潛，逾於沔，\footnote{西傾，山名。桓水自西傾山南行，因桓水是來，浮於潛。漢上曰沔。傾，窺井反。}

{\noindent\zhuan\zihao{6}\fzbyks 傳“西傾”至“曰沔”。正義曰：下文導山有“西傾”，知是山名也。\CJKunderwave{地理志}云，西傾在隴西臨洮縣西南。西傾在雍州，自西傾山南行，因桓水是來，浮於潛水也。\CJKunderwave{地理志}云，桓水出蜀郡蜀山，西南行羌中，入南海，則初發西傾未有水也，不知南行幾里得桓水也。下傳雲“泉始出山為漾水，東南流為沔水，至漢中東行為漢水”，是“漢上曰沔”。 \par}

入於渭，亂於河。\footnote{越沔而北入渭，浮東渡河而還帝都，白所治。正絕流曰亂。渭音謂。}

{\noindent\zhuan\zihao{6}\fzbyks 傳“越沔”至“曰亂”。正義曰：計沔在渭南五百餘里，故越沔陸行而北入渭。渭水入河，故浮渭而東。帝都在河之東，故渡河陸行而還帝都也。以每州之下言入河之事，河近帝都，知是還都白所治也。“正絕流曰亂”，\CJKunderwave{釋水}文。孫炎曰：“橫渡也。” \par}

黑水、西河惟\textcolor{red}{雍州}。\footnote{西距黑水，東據河。龍門之河在冀州西。雍，於用反。}

{\noindent\zhuan\zihao{6}\fzbyks 傳“西距”至“州西”。正義曰：\CJKunderline{禹}治豫州，乃次梁州，自東向西，故言梁州之境，先以華陽而後黑水。從梁適雍，自南向北，故先黑水而後西河。計雍州之境，被荒服之外，東不越河，而西逾黑水，王肅雲“西據黑水，東距西河”，所言得其實也。遍檢孔本，皆雲“西距黑水，東據河”，必是誤也。又河在雍州之東，而謂之“西河”者，龍門之河在冀州西界，故謂之“西河”。\CJKunderwave{王制}云：“自東河至於西河，千里而近。”是河相對而為東西也。 \par}

弱水既西，\footnote{導之西流，至於合黎。}

{\noindent\zhuan\zihao{6}\fzbyks 傳“導之”至“合黎”。正義曰：諸水言“既導”，此言“既西”，由地勢不同，導之使西流也。鄭云：“眾水皆東，此水獨西，故記其西下也。” \par}

涇屬渭汭。\footnote{屬，逮也。水北曰汭。言治涇水入於渭。涇音經。屬，之蜀反。汭，本又作內，同；如銳反，馬云：“入也。”逮音代。}

{\noindent\zhuan\zihao{6}\fzbyks 傳“屬逮”至“於渭”。正義曰：“屬”謂相連屬，故訓為逮。逮，及也，言水相及。\CJKunderwave{詩}毛傳云：“汭，水涯也。”鄭云：“汭之言內也。”蓋以人皆南面望水，則北為汭也。且涇水南入渭，而名為“渭汭”，知“水北曰汭”。言治涇水使之入渭，亦是從故道也。\CJKunderwave{地理志}云，涇水出安定涇陽縣西岍頭山，東南至馮翊陽陵縣入渭，行千六百里。 \par}

漆沮既從，\xpinyin*{灃}水攸同。\footnote{漆沮之水，已從入渭。灃水所同，同於渭也。沮，七徐反。灃，芳弓反。}

{\noindent\zhuan\zihao{6}\fzbyks 傳“漆沮”至“於渭”。正義曰：\CJKunderwave{詩}云：“自土沮漆。”\CJKunderwave{毛傳}云：“沮水、漆水也。”則“漆沮”本為二水。\CJKunderwave{地理志}云，漆水出扶風漆縣西。闞駰\CJKunderwave{十三州志}云：“漆水出漆縣西北岐山，東入渭。”沮則不知所出，蓋東入渭時已與漆合。渭發源遠，以渭為主，上雲“涇屬渭”是矣。故此言“漆沮既從”，已從於渭；灃水所同，亦同於渭；以渭為主故也。\CJKunderwave{地理志}灃水出扶風鄠縣東南，北過上林苑入渭也。 \par}

荊、岐既旅，\footnote{已旅祭，言治功畢。此荊在岐東，非荊州之荊。治，直吏反。}

{\noindent\zhuan\zihao{6}\fzbyks 傳“已旅”至“之荊”。正義曰：洪水之時，祭祀禮廢，已旅祭而言治功畢。治水從下，自東而西，先荊後岐，荊在岐東，嫌與上荊為一，故云“非荊州之荊”也。\CJKunderwave{地理志}云，\CJKunderwave{禹貢}北條荊山在馮翊懷德縣南，南條荊山在南郡臨沮縣北。彼是荊州之荊也。 \par}

終南、惇物，至於鳥鼠。\footnote{三山名，言相望。終南，山名。\CJKunderwave{漢書·地理志}一名太一。\CJKunderwave{山秦記}云：“又名地肺。”惇物，山名，\CJKunderwave{漢書}云：“垂山也。”}

{\noindent\zhuan\zihao{6}\fzbyks 傳“三山”至“相望”。正義曰：以“荊”、“岐”單名，此山複名，故辯之雲“三山名”也。至於為首尾之辭。故“言相望”也。三山空舉山名，不言治意，蒙上“既旅”之文也。\CJKunderwave{地理志}云，扶風武功縣有太一山，古文以為終南。垂山，古文以為惇物。皆在縣東。 \par}

原隰厎績，至於豬野。\footnote{下溼曰隰。豬野，地名。言皆致功。}

{\noindent\zhuan\zihao{6}\fzbyks 傳“下溼”至“致功”。正義曰：“下溼曰隰”，\CJKunderwave{釋地}文。\CJKunderwave{地理志}云，豬野澤，在武威縣東北有休屠澤，古文以為豬野澤。\CJKunderline{鄭玄}以為“\CJKunderwave{詩}雲‘度其隰原’，即此‘原隰’是也。原隰，豳地。從此致功，西至豬野之澤也”。 \par}

三危既宅，三苗丕敘。\footnote{西裔之山已可居，三苗之族大有次敘。美\CJKunderline{禹}之功。丕，普悲反。}

{\noindent\zhuan\zihao{6}\fzbyks 傳“西裔”至“之功”。正義曰：\CJKunderwave{左傳}稱舜去四凶，投之四裔，\CJKunderwave{舜典}雲“竄三苗於三危”，是“三危”為西裔之山也。其山必是西裔，未知山之所在。\CJKunderwave{地理志}杜林以為敦煌郡，即古瓜州也。昭九年\CJKunderwave{左傳}云：“先王居檮杌於四裔，故允姓之奸居於瓜州。”杜預云：“允姓之祖與三苗俱放於三危。瓜州，今敦煌也。”\CJKunderline{鄭玄}引\CJKunderwave{地記書}云：“三危之山在鳥鼠之西，南當岷山，則在積石之西南。”\CJKunderwave{地記}乃妄書，其言未必可信。要知三危之山必在河之南也。\CJKunderline{禹}治水未已竄三苗，水災既除，彼得安定，故云三危之山已可居，三苗之族大有次敘，記此事以美\CJKunderline{禹}治之功也。 \par}

厥土惟黃壤，厥田惟上上，厥賦中下。\footnote{田第一，賦第六，人功少。}

{\noindent\zhuan\zihao{6}\fzbyks 傳“田第一”至“功少”。正義曰：此與荊州賦田升降皆較六等，荊州升之極，故云“人功修”；此州降之極，故云“人功少”。其餘相較少者，從此可知也。\CJKunderwave{王制}云：“凡居民，量地以致邑，度地以居民，地、邑、民居必參相得也。”則民當相準，而得有人功修、人功少者，\CJKunderwave{記}言初置邑者，可以量之，而州境闊遠，民居先定，新遭洪水，存亡不同，故地勢有美，惡人功有多少。治水之後即為此差，在後隨人少多必得更立其等，此非永定也。 \par}

厥貢惟球、琳、琅玕。\footnote{球、琳皆玉名。琅玕，石而似珠。球音求。琳，韋音來金反。琅音郎。玕音幹，\CJKunderwave{山海經}云：“崑崙山有琅玕樹。”}

{\noindent\zhuan\zihao{6}\fzbyks 傳“球琳”至“似珠”。正義曰：\CJKunderwave{釋地}云：“西北之美者，有崑崙虛之璆琳琅玕焉。”說者皆云：“球、琳，美玉名。琅玕，石而似珠者。”必相傳驗，實有此言也。 \par}

浮於積石，至於龍門西河，\footnote{積石山在金城西南，河所經也。沿河順流而北，千里而東，千里而南。龍門山在河東之西界。}

{\noindent\zhuan\zihao{6}\fzbyks 傳“積石”至“西界”。正義曰：\CJKunderwave{地理志}云，積石山在金城河關縣西南羌中，河行塞外東，北入塞內。積石非河之源，故云“河所經也”。河從西來，至此北流，故\CJKunderline{禹}“沿河順流而北”。\CJKunderwave{釋水}云：“河千里一曲一直。”故千里而東，千里而南，至於龍門西河也。\CJKunderwave{地理志}云，龍門山在馮翊夏陽縣北。此山當河之道，\CJKunderline{禹}鑿以通河東郡之西界也。\CJKunderline{禹}至此渡河而還都白帝也。“沿”或誤為“治”，此說\CJKunderline{禹}行，不說治水也。 \par}

會於渭汭。\footnote{逆流曰會。自渭北涯逆水西上。上，時掌反。}

{\noindent\zhuan\zihao{6}\fzbyks 傳“逆流”至“西上”。正義曰：“會”，合也。人行逆流而水相向，故“逆流曰會”。從河入渭，自渭北涯逆水西上，言\CJKunderline{禹}白帝訖，從此而西上，更入雍州界也。諸州之末，惟言還都之道，此州事終，言發都更去，明諸州皆然也。 \par}

織皮崑崙、析支、渠搜、西戎即敘。\footnote{織皮，毛布。有此四國，在荒服之外，流沙之內，羌髳之屬皆就次敘。美\CJKunderline{禹}之功及戎狄也。侖，曾門反，馬云：“崑崙在臨羌西。”析，星曆反，馬云：“析支在河關西。”搜，所由反。\CJKunderwave{漢書·志}朔方郡有渠搜縣，\CJKunderwave{武紀}雲“北發渠搜”是也。髳音謀，又音毛。西戎，國名。}

{\noindent\zhuan\zihao{6}\fzbyks 傳“織皮”至“戎狄也”。正義曰：四國皆衣皮毛，故以“織皮”冠之。傳言“織皮，毛布。有此四國”，崑崙也，析支也，渠也，搜也，四國皆是戎狄也。末以“西戎”總之。此戎在荒服之外，流沙之內。\CJKunderwave{牧誓}云，武王伐紂，有羌髳從之。此是羌髳之屬，\CJKunderline{禹}皆就次敘。美\CJKunderline{禹}之功遠及戎狄，故記之也。\CJKunderline{鄭玄}云：“衣皮之民，居此崑崙、析支、渠搜三山之野者，皆西戎也。”王肅云：“崑崙在臨羌西，析支在河關西。西戎，西域也。”王肅不言“渠搜”，鄭並“渠搜”為一，孔傳不明。或亦以“渠搜”為一，通“西戎”為四也。鄭以“崑崙”為山，謂別有崑崙之山，非河所出者也。所以孔意或是地名國號，不必為山也。 \par}

導\xpinyin*{岍}及岐，至於荊山，\footnote{更理說所治山川首尾所在，治山通水，故以山名之。三山皆在雍州。導音道，從首起也。岍音牽，字又作汧，山名，一名吳嶽，馬本作開。}

{\noindent\zhuan\zihao{6}\fzbyks 傳“更理”至“雍州”。正義曰：“荊”、“岐”上已具矣，而此復言之，以山勢相連而州境隔絕，更從上理說所治山川首尾所在，總解此下導山水之意也。其實通水而文稱導山者，導山本為治水，故以導山名之。\CJKunderwave{地理志}云，吳嶽在扶風岍縣西，古文以為岍山，岐山在美陽縣西北，荊山在懷德縣。三山皆在雍州。 \par}

{\noindent\shu\zihao{5}\fzkt “導岍及岐”。正義曰：上文每州說其治水登山，從下而上，州境隔絕，未得徑通。今更從上而下，條說所治之山，本以通水，舉其山相連屬，言此山之傍所有水害皆治訖也。因冀州在北，故自北為始。從此“導岍”至“敷淺原”，舊說以為三條。\CJKunderwave{地理志}云，\CJKunderwave{禹貢}北條荊山，在馮翊懷德縣南，南條荊山，在南郡臨沮縣東北。是舊有三條之說也。故馬融、王肅皆為三條，“導岍”北條，“西傾”中條，“嶓冢”南條。\CJKunderline{鄭玄}以為四列，“導岍”為陰列，“西傾”為次陰列，“嶓冢”為次陽列，“岷山”為正陽列。\CJKunderline{鄭玄}創為此說，孔亦當為三條也。“岍”與“嶓冢”言“導”，“西傾”不言“導”者，史文有詳略，以可知，故省文也。 \par}

逾於河。\footnote{此謂梁山龍門西河。}

{\noindent\zhuan\zihao{6}\fzbyks 傳“此謂”至“西河”。正義曰：“逾於河”謂山逾之也。此處山勢相望，越河而東，故云此謂龍門西河，言此處山不絕，從此而渡河也。 \par}

壺口、雷首，至於太嶽。\footnote{三山在冀州。太嶽，上黨西。}

{\noindent\zhuan\zihao{6}\fzbyks 傳“三山”至“黨西”。正義曰：\CJKunderwave{地理志}云，壺口在河東北屈縣東南，雷首在河東蒲阪縣南，太嶽在河東彘縣東。是“三山在冀州”。以太嶽東近上黨，故云“在上黨西”也。 \par}

厎柱、析城，至於王屋。\footnote{此三山在冀州南河之北東行。厎，之履反。柱如字，韋知父反，又知女反。厎柱，山名，在河水中。}

{\noindent\zhuan\zihao{6}\fzbyks 傳“此三”至“東行”。正義曰：\CJKunderwave{地理志}云，析城在河東濩澤縣西，王屋在河東垣縣東北。\CJKunderwave{地理志}不載厎柱。厎柱在太陽關東，析城之西。從厎柱至王屋，在冀州南河之北東行也。 \par}

大行、恆山，至於碣石，入於海。\footnote{此二山連延東北，接碣石而入滄海。百川經此眾山，\CJKunderline{禹}皆治之，不可勝名，故以山言之。行，戶剛反，又如字。滄音倉。勝音升。}

{\noindent\zhuan\zihao{6}\fzbyks 傳“此二”至“言之”。正義曰：\CJKunderwave{地理志}云，大行山在河內山陽縣西北，恆山在常山上曲陽縣西北。大行去恆山太遠，恆山去碣石又遠，故云“此二山連延東北，接碣石而入滄海”，言山傍之水皆入海，山不入海也。又解治水言山之意，“百川經此眾山，\CJKunderline{禹}皆治之，川多不可勝名，故以山言之”也。謂漳、潞、汾、涑在壺口、雷首、太行，經厎柱、析城，濟出王屋，淇近大行，恆、衛、滹沲、滱、易近恆山、碣石之等也。 \par}

西傾、朱圉、鳥鼠，\footnote{西傾、朱圉在積石以東。鳥鼠，渭水所出，在隴西之西。三者雍州之南山。傾，窺井反。圉，魚呂反。}

{\noindent\zhuan\zihao{6}\fzbyks 傳“西傾”至“南山”。正義曰：\CJKunderwave{地理志}云，西傾在隴西臨洮縣西南，朱圉在天水冀縣南。言“在積石以東”，見河所經也。\CJKunderwave{地理志}云，鳥鼠同穴山，在隴西首陽縣西南，渭水所出，在隴西郡之西。是三者皆雍州之南山也。 \par}

至於太華。\footnote{相首尾而東。華如字，又戶化反。}

{\noindent\zhuan\zihao{6}\fzbyks 傳“相首尾而東”。正義曰：\CJKunderwave{地理志}云，太華在京兆華陰縣南。鳥鼠東望太華太遠，故云“相首尾而東”也。 \par}

熊耳、外方、桐柏，至於陪尾。\footnote{四山相連，東南在豫州界。洛經熊耳，伊經外方，淮出桐柏，經陪尾。凡此皆先舉所施功之山於上，而後條列所治水於下，互相備。陪音裴；陪尾，山名；\CJKunderwave{漢書}作“橫尾”。列如字，本或作別，彼列反。}

{\noindent\zhuan\zihao{6}\fzbyks 傳“四山”至“相備”。正義曰：\CJKunderwave{地理志}云，熊耳山在弘農盧氏縣東，伊水所出。嵩高山在穎川嵩高縣，古文以為外方山。桐柏山在南陽平氏縣東南。橫尾山在江夏安陸縣東北，古文以為陪尾山。是四山接華山而相連，東南皆在豫州界也。凡舉山名，皆為治水，故言水之所經，“洛出熊耳，伊經外方，淮出桐柏，經陪尾”。導山本為治水，故云“皆先舉所施功之山於上，而後條列所治水於下，互相備”也。 \par}

導嶓冢，至於荊山。\footnote{漾水出嶓冢，在梁州，經荊山。荊山在荊州。漾，羊尚反。}

{\noindent\zhuan\zihao{6}\fzbyks 傳“漾水”至“荊州”。正義曰：下云：“嶓冢導漾”，梁州雲“岷、嶓既藝”，是嶓冢在梁州也。荊州以荊山為名，知“荊山在荊州”也。 \par}

內方，至於大別。\footnote{內方、大別，二山名。在荊州，漢所經。}

{\noindent\zhuan\zihao{6}\fzbyks 傳“內方”至“所經”。正義曰：\CJKunderwave{地理志}云，章山在江夏竟陵縣東北，古文以為內方山。\CJKunderwave{地理志}無大別。\CJKunderline{鄭玄}云：“大別在廬江安豐縣。”杜預解\CJKunderwave{春秋}云：“大別闕，不知何處。”或曰大別在安豐縣西南，\CJKunderwave{左傳}云，吳既與楚夾漢，然後楚“乃濟漢而陳，自小別至於大別”。然則二別近漢之名，無緣得在安豐縣。如預所言，雖不知其處，要與內方相接，漢水所經，必在荊州界也。 \par}

岷山之陽，至於衡山。\footnote{岷山，江所出，在梁州。衡山，江所經，在荊州。}

{\noindent\zhuan\zihao{6}\fzbyks 傳“岷山”至“荊州”。正義曰：其下雲“岷山導江”，梁州“岷、嶓既藝”，是岷山在梁州也。\CJKunderwave{地理志}云，衡山在長沙湘南縣東南，上言“衡陽惟荊州”，是“江所經，在荊州”也。 \par}

過九江，至於敷淺原。\footnote{言衡山連延過九江，接敷淺原。言“導”從首起，言“陽”從南。敷淺原，一名博陽山，在揚州豫章界。}

{\noindent\zhuan\zihao{6}\fzbyks 傳“言衡”至“章界”。正義曰：“衡”即橫也，東西長，今之人謂之為嶺。東行連延過九江之水，而東接於敷淺原之山也。經於岍及嶓冢言“導”，岷山言“陽”，故解之“言‘導’從首起，言‘陽’從南”，言岷山之南至敷淺原，別以岷山為首，不與大別相接。由江所經，別記之耳，以見岷非三條也。\CJKunderwave{地理志}豫章歷陵縣南有博陽山，古文以為敷淺原。 \par}

導弱水，至於合黎，\footnote{合黎，水名，在流沙東。弱，本或作溺。合如字。黎，力兮反，馬云：“地名。”}

{\noindent\zhuan\zihao{6}\fzbyks 傳“合黎”至“沙東”。正義曰：弱水得入合黎，知“合黎”是水名。顧氏云：“\CJKunderwave{地說書}合黎，山名。”但此水出合黎，因山為名。\CJKunderline{鄭玄}亦以為山名。\CJKunderwave{地理志}張掖郡刪丹縣，桑欽以為導弱水自此，西至酒泉、合黎。張掖郡又有居延澤，在縣東北，古文以為流沙。如\CJKunderwave{志}之言，酒泉郡在張掖郡西，居延屬張掖，合黎在酒泉，則流沙在合黎之東，與此傳不合。案經弱水西流，水既至於合黎，餘波入於流沙，當如傳文合黎在流沙之東，不得在其西也。 \par}

{\noindent\shu\zihao{5}\fzkt “導弱水”。正義曰：此下所導，凡有九水，大意亦自北為始。以弱水最在西北，水又西流，故先言之。黑水雖在河南，水從雍、梁西界南入南海，與諸水不相參涉，故又次之。四瀆江、河為大，河在北，故先言河也。漢入於江，故先漢後江。其濟發源河北，越河而南，與淮俱為四瀆，故次濟，次淮。其渭與洛俱入於河，故後言之。計流水多矣，此舉大者言耳。凡此九水，立文不同，弱水、黑水、沇水不出于山，文單，故以“水”配。其餘六水，文與山連，既繫於山，不須言“水”。積石山非河上源，記施功之處，故云“導河積石”，言發首積石起也。漾、江先山後水，淮、渭洛先水後山，皆是史文詳略，無義例也。又淮、渭、洛言“自某山”者，皆是發源此山，欲使異於導河，故加“自”耳。\CJKunderline{鄭玄}云：“凡言‘導’者，發源於上，未成流。凡言‘自’者，亦發源於上，未成流。”必其俱未成流，何須別“導”與“自”？河出崑崙，發源甚遠，豈至積石，猶未成流而云“導河”也？ \par}

餘波入於流沙。\footnote{弱水餘波西溢入流沙。溢音逸。}導黑水，至於三危，入於南海。\footnote{黑水自北而南，經三危，過樑州，入南海。}

{\noindent\zhuan\zihao{6}\fzbyks 傳“黑水”至“南海”。正義曰：\CJKunderwave{地理志}益州郡計在蜀郡西南三千餘里，故滇王國也。武帝元封二年始開為郡。郡內有滇池縣，縣有黑水祠，止言有其祠，不知水之所在。鄭云：“今中國無也。”傳之此言，順經文耳。案酈元\CJKunderwave{水經}：“黑水出張掖雞山，南流至敦煌，過三危山，南流入於南海。”然張掖、敦煌並在河北，所以黑水得越河入南海者，河自積石以西皆多伏流，故黑水得越而南也。 \par}

導河積石，至於龍門；\footnote{施功發於積石，至於龍門，或鑿山，或穿地以通流。}

{\noindent\zhuan\zihao{6}\fzbyks 傳“施功”至“通流”。正義曰：河源不始於此，記其施功處耳，故言“施功發於積石”。\CJKunderwave{釋水}雲“河千里一曲一直”，則河從積石北行，又東，乃南行至於龍門，計應三千餘里。龍門厎柱，鑿山也。其餘平地，穿地也。“或鑿山、或穿地以通流”，言自積石至海皆然也。\CJKunderwave{釋水}云：“河出崑崙虛，色白。”李巡曰：“崑崙，山名。虛，山下地也。”郭璞云：“發源高處激湊，故水色白。潛流地中，受渠眾多，渾濁，故水色黃。”\CJKunderwave{漢書·西域傳}云：“河有兩源，一出蔥嶺，一出於闐。于闐在南山下，其河北流，與蔥嶺河合，東注蒲昌海。蒲昌海，一名鹽澤者，去玉門、陽關三百餘里，廣袤三四百里。其水停居，冬夏不增減，皆以為潛行地下，南出於積石，為中國河。”郭璞云：“其去崑崙，裡數遠近未得詳也。” \par}

南至於華陰，\footnote{河自龍門南流至華山，北而東行。}東至於厎柱；\footnote{厎柱，山名。河水分流，包山而過，山見水中若柱然，在西虢之界。見，賢遍反。虢，寡白反。}又東至於孟津，\footnote{孟津，地名。在洛北，都道所湊，古今以為津。孟津，如字。洛北，地名。湊，七豆反。}

{\noindent\zhuan\zihao{6}\fzbyks 傳“孟津”至“為津”。正義曰：“孟”是地名，“津”是渡處，在孟地致津，謂之“孟津”。傳雲“地名”，謂“孟”為地名耳。杜預云：“孟津，河內河陽縣南孟津也。在洛陽城北，都道所湊，古今常以為津。武王渡之，近世以來呼為武濟。” \par}

東過洛汭，至於大伾；\footnote{洛汭，洛入河處。山再成曰伾。至於大伾而北行。伾，本或作岯，音丕，又皮鄙反；徐扶眉反，又敷眉反；韋音嚭；郭撫梅反，字或作 。}

{\noindent\zhuan\zihao{6}\fzbyks 傳“洛汭”至“北行”。正義曰：“洛汭，洛入河處”，河南鞏縣東也。\CJKunderwave{釋山}云：“再成英，一成岯。”李巡曰：“山再重曰英，一重曰岯。”傳雲“再成曰岯”，與\CJKunderwave{爾雅}不同，蓋所見異也。\CJKunderline{鄭玄}云：“大岯在修武武德之界。”張揖云：“成皋縣山也。”\CJKunderwave{漢書音義}有臣瓚者，以為：“修武武德無此山也。成皋縣山，又不一成，今黎陽縣山臨河。豈不是大岯乎？”瓚言當然。 \par}

北過降水，至於大陸；\footnote{降水，水名，入河。大陸，澤名。降如字，鄭戶江反。}

{\noindent\zhuan\zihao{6}\fzbyks 傳“降水”至“澤名”。正義曰：\CJKunderwave{地理志}云，降水在信都縣。案班固\CJKunderwave{漢書}以襄國為信都，在大陸之內。或降水發源在此，下尾至今之信都，故得先過降水，乃至大陸。若其不爾，則降水不可知也。鄭以“降讀為降，下江反，聲轉為共，河內共縣，淇水出焉，東至魏郡黎陽縣入河。北近降水也，周時國於此地者惡言降水，改謂之共”。此鄭胸臆，不可從也。\par}

又北播為九河，\footnote{北分為九河，以殺其溢，在兗州界。殺，所界反。溢，字又作隘，於賣反。}同為逆河，入於海。\footnote{同合為一大河，名逆河，而入於渤海。皆\CJKunderline{禹}所加功，故敘之。渤，蒲兀反。}

{\noindent\zhuan\zihao{6}\fzbyks 傳“同合”至“敘之”。正義曰：傳言九河將欲至海，更同合為一大河，名為逆河，而入於渤海也。\CJKunderline{鄭玄}云：“下尾合，名為逆河，言相向迎受。”王肅云：“同逆一大河，納之於海。”其意與孔同。 \par}

嶓冢導漾，東流為漢；\footnote{泉始出山為漾水，東南流為沔水，至漢中東流為漢水。}

{\noindent\zhuan\zihao{6}\fzbyks 傳“泉始”至“漢水”。正義曰：傳之此言，當據時人之名為說也。\CJKunderwave{地理志}云，漾水出隴西氐道縣，至武都為漢水。不言中為沔水。孔知嶓冢之東、漢水之西而得為沔水者，以\CJKunderline{禹}治梁州，入帝都白所治，雲“逾於沔，入於渭”，是沔近於渭，當梁州向冀州之路也。應劭云：“沔水自江別，至南郡華容縣為夏水，過江夏郡入江。”既雲“江別”，明與此沔別也。依\CJKunderwave{地理志}，漢水之尾變為夏水，是應劭所云沔水下尾亦與漢合，乃入於江也。 \par}

又東為滄浪之水；\footnote{別流在荊州。浪音郎。}

{\noindent\zhuan\zihao{6}\fzbyks 傳“別流在荊州”。正義曰：傳言“別流”，似分為異水。案經首尾相連，不是分別，當以名稱別流也。又上在梁州，故此雲“在荊州”。 \par}

過三澨，至於大別，\footnote{三澨，水名，入漢。大別，山名。澨，市制反。}南入於江。\footnote{觸山回南，入江。觸，\CJKunderwave{切韻}尺玉反。}東匯澤為彭蠡，\footnote{匯，回也。水東回為彭蠡大澤。匯，徐胡罪反，韋空為反。}東為北江，入於海。\footnote{自彭蠡江分為三，入震澤，遂為北江而入海。}

{\noindent\zhuan\zihao{6}\fzbyks 傳“自彭”至“入海”。正義曰：揚州云：“三江既入，震澤厎定”，孔為“三江既入”，入震澤也，故言江自彭蠡分而為三江，復共入震澤。出澤又分為三，此水遂為北江而入於海。\CJKunderline{鄭玄}以為：“‘三江既入’，入於海，不入震澤也。”孔必知入震澤者，以震澤屬揚州，彭蠡在揚州之西界，今從彭蠡有三江，則震澤之西三江具矣。今雲“三江既入”，繼以“震澤厎定”，故知三江入震澤矣。今南人以大江不入震澤，震澤之東別有松江等三江。案\CJKunderwave{職方}揚州“其川曰三江”，宜舉州內大川，其松江等雖出震澤，入海既近，\CJKunderwave{周禮}不應舍岷山大江之名，而記松江等小江之說。山水古今變易，故鄭雲既知今，亦當知古，是古今同之驗也。 \par}

岷山導江，東別為沱；\footnote{江東南流，沱東行。沱，唐河反。}

{\noindent\zhuan\zihao{6}\fzbyks 傳“江東”至“東行”。正義曰：以上雲“浮於江、沱、潛、漢”，其次自南而北，江在沱南，知“江東南流，而沱東行”。 \par}

又東至於澧，\footnote{澧，水名。澧音禮。}

{\noindent\zhuan\zihao{6}\fzbyks 傳“澧，水名”。正義曰：\CJKunderline{鄭玄}以此經自“導弱水”已下，言“過”言“會”者，皆是水名，言“至於”者或山或澤，皆非水名，故以“合黎”為山名，“澧”為陵名。\CJKunderline{鄭玄}云：“今長沙郡有澧陵縣，其以陵名為縣乎？”孔以“合黎”與“澧”皆為水名，弱水餘波入於流沙，則本源入合黎矣。合黎得容弱水，知是水名。\CJKunderwave{楚辭}曰：“濯餘佩兮澧浦”是，“澧”亦為水名。 \par}

過九江，至於東陵；\footnote{江分為九道，在荊州。東陵，地名。}

{\noindent\zhuan\zihao{6}\fzbyks 傳“江分”至“地名”。正義曰：九江之水，\CJKunderline{禹}前先有其處。\CJKunderline{禹}今導江，過歷九江之處，非是別有九江之水。 \par}

東迆北會於匯；\footnote{迆，溢也。東溢分流，都共北會為彭蠡。迆，以爾反，馬雲靡也。}

{\noindent\zhuan\zihao{6}\fzbyks 傳“迆溢”至“彭蠡”。正義曰：“迆”言靡迆，邪出之言，故為溢也。東溢分流，又都共聚合，北會彭蠡，言散流而複合也。鄭雲“東迆者為南江”，孔意或然。“至”之與“會”，史異文耳。 \par}

東為中江，入於海。\footnote{有北，有中，南可知。}

{\noindent\zhuan\zihao{6}\fzbyks 傳“有北，有中，南可知”。正義曰：\CJKunderwave{地理志}云，南江從會稽吳縣南東入海，中江從丹陽無湖縣西東至會稽陽羨縣東入海，北江從會稽毗陵縣北入於海。 \par}

導沇水，東流為濟，\footnote{泉源為沇，流去為濟，在溫西北平地。沇音兗，又以轉反。}

{\noindent\zhuan\zihao{6}\fzbyks 傳“泉源”至“平地”。正義曰：\CJKunderwave{地理志}云，濟水出河東垣縣王屋山，東南至河內武德縣入河。傳言“在溫西北平地”者，濟水近在河內，孔必驗而知之。見今濟水所出，在溫之西北七十餘里，溫是古之舊縣，故計溫言之。 \par}

入於河，溢為滎；\footnote{濟水入河，並流十數里，而南截河。又並流數里，溢為滎澤，在敖倉東南。數，色主反，下同，一本作十所。}

{\noindent\zhuan\zihao{6}\fzbyks 傳“濟水”至“東南”。正義曰：此皆目驗為說也。濟水既入於河，與河相亂，而知截河過者，以河濁濟清，南出還清，故可知也。 \par}

東出於陶丘北，\footnote{陶丘，丘再成。陶音桃。}

{\noindent\zhuan\zihao{6}\fzbyks 傳“陶丘，丘再成”。正義曰：\CJKunderwave{釋丘}云：“再成為陶丘。”李巡曰：“再成，其形再重也。”郭璞云：“今濟陰定陶城中有陶丘。”\CJKunderwave{地理志}云，定陶縣西南有陶丘亭。 \par}

又東至於菏；\footnote{菏澤之水。}又東北，會於汶；\footnote{濟與汶合。}又北東入於海。\footnote{北折而東。折，之設反。}導淮自桐柏，\footnote{桐柏山，在南陽之東。}

{\noindent\zhuan\zihao{6}\fzbyks 傳“桐柏”至“之東”。正義曰：\CJKunderwave{地理志}云，桐柏山在南陽平氏縣東南，淮水所出。\CJKunderwave{水經}云：“出胎簪山，東北過桐柏山。”胎簪蓋桐柏之傍小山，傳言南陽郡之東也。 \par}

東會於泗、沂，東入於海。\footnote{與泗、沂二水合，入海。}

{\noindent\zhuan\zihao{6}\fzbyks 傳“與泗”至“入海”。正義曰：\CJKunderwave{地理志}云，沂水出泰山蓋縣，南至下邳，入泗。泗水出濟陰乘氏縣，至臨淮雎陵縣入淮。乃沂水先入泗，泗入淮耳。以沂水入泗處去淮已近，故連言之。 \par}

導渭自鳥鼠同穴，\footnote{鳥鼠共為雌雄，同穴處此山，遂名山曰鳥鼠，渭水出焉。}

{\noindent\zhuan\zihao{6}\fzbyks 傳“鳥鼠”至“出焉”。正義曰：\CJKunderwave{釋鳥}云：“鳥鼠同穴，其鳥為鵭，其鼠為鼵。”李巡曰：“鵭鼵鳥鼠之名，共處一穴，天性然也。”郭璞曰：“鼵如人家鼠而短尾，鵭似鵽而小，黃黑色。穴入地三四尺，鼠在內，鳥在外，今在隴西首陽縣有鳥鼠同穴山。\CJKunderwave{尚書}孔傳雲‘共為雄雌’，張氏\CJKunderwave{地理記}雲‘不為牝牡’。”璞並載此言，未知誰得實也。\CJKunderwave{地理志}云，隴西首陽西南有鳥鼠同穴山，渭水所出，至京兆北船司空縣入河，過郡四，行千八百七十里。 \par}

東會於灃，又東會於涇，\footnote{灃水自南，涇水自北而合。灃音豐。}又東過漆沮，入於河。\footnote{漆沮，一水名，亦曰洛水，出馮翊北。翊，與職反。}

{\noindent\zhuan\zihao{6}\fzbyks 傳“漆沮”至“翊北”。正義曰：\CJKunderwave{地理志}云，漆水出扶風漆縣。依\CJKunderwave{十三州記}，漆水在岐山東入渭，則與漆沮不同矣。此雲“會於涇”，又“東過漆沮”，是漆沮在涇水之東，故孔以為洛水一名漆沮。\CJKunderwave{水經}沮水出北池直路縣，東入洛水。又云鄭渠在太上皇陵東南，濯水入焉，俗謂之漆水，又謂之漆沮，其水東流注於洛水。\CJKunderwave{志}云，出馮翊懷德縣，東南入渭。以水土驗之，與\CJKunderwave{毛詩}古公“自土沮漆”者別也。彼漆即扶風漆水也，彼沮則未聞。 \par}

導洛自熊耳，\footnote{在宜陽之西。}東北會於澗、瀍，\footnote{會於河南城南。}又東會於伊，\footnote{合於洛陽之南。}又東北入於河。\footnote{合於鞏之東。鞏，恭勇反，縣名，屬河南郡。}九州攸同，\footnote{所同事在下。}四隩既宅，\footnote{四方之宅巳可居。隩，於六反，\CJKunderwave{玉篇}於報反。}九州刊旅，九川滌源，九澤既陂，\footnote{九州名山與槎木通道而旅祭矣，九州之川已滌除泉源無壅塞矣，九州之澤已陂障無決溢矣。滌，待歷反。陂，彼宜反。槎，仕雅反。障，章尚反。}四海會同，六府孔修。\footnote{四海之內會同京師，九州同風，萬國共貫，水、火、金、木、土、谷甚修治。言政化和。貫，工喚反。}庶土交正，厎慎財賦，\footnote{交，俱也。眾土俱得其正，謂壤、墳、壚。致所慎者，財貨貢賦。言取之有節，不過度。}咸則三壤成賦，（中邦。）\footnote{皆法壤田上中下大較三品，成九州之賦，明水害除。較音角。}

{\noindent\zhuan\zihao{6}\fzbyks 傳“所同事在下”。正義曰：九州所同，與下為目，故言“所同事在下”。“四隩既宅”已下皆是也，其言“九山”、“九川”、“九澤”，最是同之事矣。 \par}

{\noindent\zhuan\zihao{6}\fzbyks 傳“四方”至“可居”。正義曰：室隅為“隩”，“隩”是內也。人之造宅為居，至其隩內，遂以“隩”表宅，故傳以“隩”為宅，以宅內可居，言四方舊可居之處皆可居也。 \par}

{\noindent\zhuan\zihao{6}\fzbyks 傳“九州”至“溢矣”。正義曰：上文諸州有言山川澤者，皆舉大言之。所言不盡，故於此復更總之。“九山”、“九川”、“九澤”，言九州之內所有山川澤，無大無小，皆刊槎決除已訖，其皆旅祭。惟據名山大川言“旅”者，往前大水，旅祭禮廢，已旅見已治也。山非水體，故以“旅”見治。其實水亦旅矣，發首雲“奠高山大川”，但是定位，皆已旅祭也。川言“滌除泉源”，從其所出，至其所入，皆蕩除之，無壅塞也。澤言“既陂”，往前濫溢，今時水定，或作陂以障之，使無決溢。\CJKunderwave{詩}云：“彼澤之陂。”\CJKunderwave{毛傳}云：“陂，澤障也。” \par}

{\noindent\zhuan\zihao{6}\fzbyks 傳“四海”至“化和”。正義曰：\CJKunderwave{禮}諸侯之見天子，“時見曰會，殷見曰同”。此言“四海會同”，乃謂官之與民皆得聚會京師，非據諸侯之身朝天子也。夷狄戎蠻謂之四海，但天子之於夷狄，不與華夏同風，故知“四海”謂“四海之內”，即是九州之中，乃有萬國。萬國同其風化，若物在繩索之貫，故云“九州同風，萬國共貫”。\CJKunderwave{大禹謨}云，水、火、金、木、土、谷謂之六府。皆修治者，言政化和也。由政化和平，民不失業，各得殖其資產，故六府修治也。 \par}

{\noindent\zhuan\zihao{6}\fzbyks 傳“交俱”至“過度”。正義曰：交錯、更互，“俱”之義，故“交”為俱也。洪水之時，高下皆水，土失本性。今水災既除，“眾土俱得其正，謂壤、墳、壚”，還復其壤、墳、壚之性也。諸州之土，“青黎”是色，“塗泥”是溼，土性之異，惟有“壤、墳、壚”耳，故舉三者以言也。致所慎者，財貨貢賦，謹慎其事，不使害人，言取民有節，什一而稅，不過度也。 \par}

{\noindent\zhuan\zihao{6}\fzbyks 傳“皆法”至“害除”。正義曰：土壤各有肥瘠，貢賦從地而出，故分其土壤為上中下，計其肥瘠等級甚多，但舉其大較，定為三品，法則地之善惡，以為貢賦之差。雖細分三品，以為九等，人功修少，當時小異，要民之常稅必準其土，故皆法三壤成九州之賦。言得施賦法，以明水害除也。“九州”即是“中邦”，故傳以“九州”言之。 \par}

{\noindent\shu\zihao{5}\fzkt “九州”至“中邦”。正義曰：昔堯遭洪水，道路阻絕，今水土既治，天下大同，故總敘之，今九州所共同矣。所同者，四方之宅已儘可居矣，九州之山刊槎其木旅祭之矣，九州之川滌除泉源無壅塞矣，九州之澤已皆陂障無決溢矣，四海之內皆得會同京師無乖異矣，六材之府甚修治矣。言海內之人皆豐足矣。水災已除，天下眾土墳壤之屬俱得其正，複本性故也。民既豐足，取之有藝，致所重慎者惟財貨賦稅也。慎之者，皆法則其三品上壤，準其地之肥瘠，為上中下三等，以成其貢賦之法於中國。美\CJKunderline{禹}能治水土，安海內，於此裛結之。 \par}

（中邦）錫土姓，祗臺德先，不距朕行。\footnote{臺，我也。天子建德，因生以賜姓。謂有德之人生此地，以此地名賜之姓以顯之。王者常自以敬我德為先，則天下無距違我行者。}

{\noindent\zhuan\zihao{6}\fzbyks 傳“臺我”至“行者”。正義曰：“臺，我”，\CJKunderwave{釋詁}文。“天子建德，因生以賜姓”，隱八年\CJKunderwave{左傳}文。既引其文，又解其義：土，地也，謂有德之人生於此地，天子以地名賜之姓以尊顯之。\CJKunderwave{周語}稱帝嘉\CJKunderline{禹}德，賜姓曰姒；祚四嶽，賜姓曰姜；\CJKunderwave{左傳}稱周賜陳胡公之姓為媯，皆是因生賜姓之事也。臣蒙賜姓，其人少矣，此事是用賢大者，故舉以為言。王者既能用賢，又能謹敬，其立意也常自以敬我德為先，則天下無有距違我天子之行者。\CJKunderwave{論語}云：“上好禮，則民莫敢不敬。上好義，則民莫敢不服。上好信，則民莫敢不用情。”王者自敬其德，則民豈敢不敬之？人皆敬之，誰敢距違者？聖人行而天下皆悅，動而天下皆應，用此道也。 \par}

{\noindent\shu\zihao{5}\fzkt “錫土”至“朕行”。正義曰：此一經皆史美\CJKunderline{禹}功，言九州風俗既同，可以施其教化，天子惟當擇任其賢者，相與共治之。選有德之人，賜與所生之土為姓，既能尊賢如是，又天子立意，常自以敬我德為先，則天下之民無有距違我天子所行者。皆\CJKunderline{禹}之使然，故敘而美之。 \par}

五百里甸服。\footnote{規方千里之內謂之甸服。為天子服治田，去王城面五百里。甸，田遍反。為，於偽反。}

{\noindent\zhuan\zihao{6}\fzbyks 傳“規方”至“百里”。正義曰：“先王規方千里，以為甸服”，\CJKunderwave{周語}文。\CJKunderwave{王制}亦云：“千里之內曰甸。”\CJKunderline{鄭玄}云：“服治田，出谷稅也。言甸者，主治田，故服名甸也。” \par}

{\noindent\shu\zihao{5}\fzkt “五百里甸服”。正義曰：既言九州同風，法壤成賦，而四海之內路有遠近，更敘弼成五服之事。甸、侯、綏、要、荒五服之名，堯之舊制。洪水既平之後，\CJKunderline{禹}乃為之節文，使賦役有恆，職掌分定。甸服去京師最近，賦稅尤多，故每於百里即為一節。侯服稍遠，近者供役，故二百里內各為一節，三百里外共為一節。綏、要、荒三服，去京師益遠，每服分而為二，內三百里為一節，外二百里為一節。以遠近有較，故其任不等。甸服入谷，故發首言賦稅也。賦令自送入官，故三百里內每皆言“納”。四百里、五百里不言“納”者，從上省文也。於三百里言“服”者，舉中以明上下，皆是服王事也。侯服以外貢不入谷，侯主為斥候。二百里內徭役差多，故各為一名。三百里外同是斥候，故共為一名。自下皆先言三百里，而後二百里，舉大率為差等也。 \par}

百里賦納緫，\footnote{甸服內之百里近王城者。禾稿曰緫，入之供飼國馬。納如字，本又作內，音同，下如字。緫音揔。近，附近之近。稿，故老反。供音恭。飼音嗣。}

{\noindent\zhuan\zihao{6}\fzbyks 傳“甸服”至“國馬”。正義曰：去王城五百里總名甸服，就其甸服內又細分之。從內而出，此為其首，故云“甸服之內近王城者”，“緫”者，總下“銍”、“秸”，禾穗與稿，緫皆送之，故云“禾稿曰緫，入之供飼國馬”。\CJKunderwave{周禮}掌客待諸侯之禮有芻，有禾，此緫是也。 \par}

二百里納銍，\footnote{銍，刈，謂禾穗。銍，珍慄反。穗亦作穟，音遂。}

{\noindent\zhuan\zihao{6}\fzbyks 傳“銍，刈，謂禾穗”。正義曰：\CJKunderline{劉熙}\CJKunderwave{釋名}云：“銍，獲禾鐵也。”\CJKunderwave{說文}云：“銍，獲禾短鐮也。”\CJKunderwave{詩}雲“奄觀銍刈”，用銍刈者，謂禾穗也。禾穗用銍以刈，故以“銍”表禾穗也。 \par}

三百里納秸服，\footnote{秸，稿也，服瑽役。秸，本或作稭，工八反，馬云：“去其穎，音鞂。”}

{\noindent\zhuan\zihao{6}\fzbyks 傳“秸，稿也，服稿役”。正義曰：\CJKunderwave{郊特牲}云：“莞簟之安，而稿秸之設。”“秸”亦“稿”也，雙言之耳。去穗送稿，易於送穗，故為遠彌輕也。然計什一而得，稿粟皆送，則秸服重於納銍，則乖近重遠輕之義。蓋納粟之外，斟酌納稿。“服稿役”者，解經“服”字，於此言“服”，明上下服皆並有所納之役也。四百里猶尚納粟，此當稿、粟別納，非是徒納稿也。 \par}

四百里粟，五百里米。\footnote{所納精者少，粗者多。}

{\noindent\zhuan\zihao{6}\fzbyks 傳“所納”至“者多”。正義曰：直納粟米為少，禾稿俱送為多。其於稅也。皆當什一，但所納有精粗，遠輕而近重耳。 \par}

五百里侯服。\footnote{甸服外之五百里。侯，候也。斥候而服事。}

{\noindent\zhuan\zihao{6}\fzbyks 傳“甸服”至“服事”。正義曰：“侯”聲近候，故為候也。襄十八年\CJKunderwave{左傳}稱晉人伐齊，使司馬斥山澤之險。“斥”謂檢行之也。“斥候”謂檢行險阻，伺候盜賊。此五百里主為斥候而服事天子，故名“侯服”。因見諸言“服”者，皆是服事也。 \par}

百里採，\footnote{侯服內之百里，供王事而已，不主一。}

{\noindent\zhuan\zihao{6}\fzbyks 傳“侯服”至“主一”。正義曰：“採”訓為事，此百里之內主供王事而已。“事”謂役也，有役則供，不主於一，故但言“採”。 \par}

二百里男邦，\footnote{男，任也，任王者事。任，而針反，又而鴆反，下同。}

{\noindent\zhuan\zihao{6}\fzbyks 傳“男，任也，任王者事”。正義曰：“男”聲近任，故訓為任。“任王者事”，任受其役，此任有常，殊於“不主一”也。言“邦”者，見上下皆是諸侯之國也。 \par}

三百里諸侯。\footnote{三百里同為王者斥候，故合三為一名。為，於偽反。}

{\noindent\zhuan\zihao{6}\fzbyks 傳“三百”至“一名”。正義曰：經言“諸侯”者，三百里內同為王者斥候，在此內所主事同，故合三百、四百、五百共為一名，言“諸侯”以示義耳。 \par}

五百里\xpinyin*{綏}服。\footnote{綏，安也。侯服外之五百里，安服王者之政教。綏，息遺反。}

{\noindent\zhuan\zihao{6}\fzbyks 傳“綏安”至“政教”。正義曰：“綏，安”，\CJKunderwave{釋詁}文。要服去京師已遠，王者以文教要束使服。此綏服路近，言“王者政教”，以示不待要束言安服自服也。\CJKunderwave{周語}云：“先王之制，邦內甸服，邦外侯服，侯衛賓服，夷蠻要服，戎狄荒服。”彼“賓服”當此“綏服”。韋昭云：“以文武侯衛為安，王賓之，因以名服。”然則“綏”者據諸侯安王為名，“賓”者據王敬諸侯為名，故云“先王之制”，則此服舊有二名。 \par}

三百里揆文教，\footnote{揆，度也。度王者文教而行之，三百里皆同。揆，葵癸反。度，待洛反。}

{\noindent\zhuan\zihao{6}\fzbyks 傳“揆度”至“皆同”。正義曰：\CJKunderwave{釋詁}訓“揆”為度，故雙言之。以王者有文教，此服諸侯揆度王者政教而行之，必自揆度，恐其不合上耳。即是安服王者之義。 \par}

二百里奮武衛。\footnote{文教外之二百里奮武衛，天子所以安。奮，方問反。}

{\noindent\zhuan\zihao{6}\fzbyks 傳“文教”至“以安”。正義曰：既言“三百”，又言“二百”，嫌是“三百”之內，以下二服文與此同，故於此解之，此是“文教外之二百里”也。由其心安王化，奮武以衛天子，所以名此服為安也。內文而外武，故先“揆文教”，後言“奮武衛”，所從言之異，與安之義同。奮武衛天子，是其安之驗也。言服內諸侯，心安天子，非言天子賴諸侯以安也。 \par}

五百里要服。\footnote{綏服外之五百里，要束以文教。要，一遙反。束如字，一音來。}

{\noindent\zhuan\zihao{6}\fzbyks 傳“綏服”至“文教”。正義曰：“要”者約束之義。上言“揆文教”，知“要”者，“要束以文教”也。綏服自揆天子文教，恐其不稱上旨。此要服差遠，已慢王化，天子恐其不服，乃以文教要服之。名為“要”，見其疏遠之義也。 \par}

三百里夷，\footnote{守平常之教，事王者而已。馬云：“夷，易也。”}二百里蔡。\footnote{蔡，法也。法三百里而差簡。差，初佳反，又初賣反。}

{\noindent\zhuan\zihao{6}\fzbyks 傳“蔡法”至“差簡”。正義曰：“蔡”之為法，無正訓也。上言“三百里夷”，“夷”訓平也，言守平常教耳。此名為“蔡”，義簡於“夷”，故訓“蔡”為法。法則三百里者，去京師彌遠，差復簡易，言其不能守平常也。 \par}

五百里荒服。\footnote{要服外之五百里。言荒又簡略。}

{\noindent\zhuan\zihao{6}\fzbyks 傳“要服”至“簡略”。正義曰：服名“荒”者，王肅云：“政教荒忽，因其故俗而治之。”傳言“荒又簡略”，亦當以為荒忽，又簡略於要服之蔡也。 \par}

三百里蠻，\footnote{以文德蠻來之，不制以法。}

{\noindent\zhuan\zihao{6}\fzbyks 傳“以文”至“以法”。正義曰：鄭云：“蠻者聽從其俗，羈縻其人耳。故云蠻，蠻之言緡也。”其意言“蠻”是緡也，緡是繩也，言“蠻”者以繩束物之名。揆度文教，\CJKunderwave{論語}稱“遠人不服，則修文德以來之”，故傳言“以文德蠻來之”，不制以國內之法強逼之。王肅云：“蠻，慢也，禮儀簡慢。”與孔異。然甸、侯、綏、要四服，俱有三日之役，什一而稅，但二百里蔡者，稅微差簡，其荒服力役田稅並無，故鄭注云：“蔡之言殺，減殺其賦。”荒服既不役作其人，又不賦其田事也。其侯綏等所出稅賦，各入本國，則亦有納緫、納銍之差，但此據天子立文耳。要服之內，皆有文教，故孔於要服傳雲“要束以文教”，則知已上皆有文教可知。獨於綏服三百里雲“揆文教”者，以去京師既遠，更無別供，又不近外邊，不為武衛。其要服又要束始行文教，無事而能揆度文教而行者，惟有此三百里耳。“奮武衛”者，在國習學兵武，有事則征討夷狄。不於要服內奮武衛者，以要服逼近夷狄，要束始來，不可委以兵武。 \par}

二百里流。\footnote{流，移也。言政教隨其俗。凡五服相距為方五千裡。}

{\noindent\zhuan\zihao{6}\fzbyks 傳“流移”至“千里”。正義曰：流，如水流，故云“移也”。其俗流移無常，故政教隨其俗，任其去來，不服蠻來之也。凡五服之別，各五百里，是王城四面，面別二千五百里，四面相距為方五千裡也。賈逵、馬融以為“甸服之外百里至“五百里米特有此數，去王城千里；其侯、綏、要、荒服各五百里，是面三千里，相距為方六千里”。\CJKunderline{鄭玄}以為“五服服別五百里，是堯之舊制。及\CJKunderline{禹}弼之，每服之間更增五百里，面別至於五千裡，相距為方萬里”。司馬遷與孔意同，王肅亦以為然，故肅注此云：“賈、馬既失其實，\CJKunderline{鄭玄}尤不然矣。\CJKunderline{禹}之功在平治山川，不在拓境廣土。土地之廣三倍於堯，而書傳無稱也，則\CJKunderline{鄭玄}創造，難可據信。漢之孝武，疲弊中國，甘心夷狄，天下戶口至減太半，然後僅開緣邊之郡而已。\CJKunderline{禹}方憂洪水，三過其門不入，未暇以征伐為事，且其所以為服之名，輕重顛倒，遠近失所，難得而通矣。先王規方千里，以為甸服，其餘均分之公、侯、伯、子、男，使各有寰宇，而使甸服之外諸侯入禾稿，非其義也。”史遷之旨蓋得之矣，是同於孔也。若然，\CJKunderwave{周禮}王畿之外別有九服，服別五百里，是為方萬里，復以何故三倍於堯？又\CJKunderwave{地理志}言漢之土境東西九千三百二里，南北萬三千三百六十八里。驗其所言山川，不出\CJKunderwave{禹貢}之域，山川戴地，古今必同，而得裡數異者，堯與周漢其地一也，\CJKunderwave{尚書}所言，據其虛空鳥路方直而計之，\CJKunderwave{漢書}所言，乃謂著地人跡屈曲而量之，所以數不同也。故王肅上篇注云：“方五千裡者，直方之數，若其回邪委曲，動有倍加之較。”是言經指直方之數，漢據回邪之道。有九服、五服，其地雖同，王者革易，自相變改其法，不改其地也。\CJKunderline{鄭玄}不言\CJKunderline{禹}變堯法，乃雲地倍於堯，故王肅所以難之。\CJKunderwave{王制}云：“西不盡流沙，東不盡東海，南不盡衡山，北不盡恆山。”凡四海之內，斷長補短，方三千里者，彼自言“不盡”，明未至遠界，且\CJKunderwave{王制}漢世為之，不可與經合也。 \par}

東漸於海，西被於流沙，朔南暨聲教，\footnote{漸，入也。被，及也。此言五服之外皆與王者聲教而朝見。漸，子廉反。被，皮寄反。朔，朔北也。與音預。朝，直遙反。見，賢遍反。}訖於四海。\CJKunderline{禹}錫玄圭，告厥成功。\footnote{玄，天色。\CJKunderline{禹}功盡加於四海，故堯賜玄圭以彰顯之。言天功成。訖，斤密反。}

{\noindent\zhuan\zihao{6}\fzbyks 傳“漸入”至“朝見”。正義曰：“漸”是沾溼，故為入，謂入海也。覆被是遠及之辭，故為及也。海多邪曲，故言“漸入”；流沙長遠，故言“被及”，皆是過之急也。五服之下乃說此事，故言“此五服之外皆與王者聲教而朝見”，言其聞風感德而來朝也。\CJKunderline{鄭玄}云：“南北不言所至，容逾之。”此言“西被於流沙”，流沙當是西境最遠者也。而\CJKunderwave{地理志}以流沙為張掖居延澤是也，計三危在居延之西，太遠矣，\CJKunderwave{志}言非也。 \par}

{\noindent\zhuan\zihao{6}\fzbyks 傳“玄天”至“功成”。正義曰：\CJKunderwave{考工記}“天謂之玄”，是“玄”為天色。\CJKunderline{禹}之蒙賜，必是堯賜，故史敘其事，“\CJKunderline{禹}功盡加於四海，故堯賜玄圭以玄顯之”。必以天色圭者，“言天功成”也。\CJKunderwave{大禹謨}舜美\CJKunderline{禹}功雲“地平天成”，是天功成也。 \par}

{\noindent\shu\zihao{5}\fzkt “東漸”至“成功”。正義曰：言五服之外，又東漸入於海，西被及於流沙，其北與南雖在服外，皆與聞天子威聲文教，時來朝見，是\CJKunderline{禹}治水之功盡加於四海。以\CJKunderline{禹}功如是，故帝賜以玄色之圭，告其能成天之功也。 \par}

%%% Local Variables:
%%% mode: latex
%%% TeX-engine: xetex
%%% TeX-master: "../Main"
%%% End:

%% -*- coding: utf-8 -*-
%% Time-stamp: <Chen Wang: 2024-04-02 11:42:41>

% {\noindent \zhu \zihao{5} \fzbyks } -> 注 (△ ○)
% {\noindent \shu \zihao{5} \fzkt } -> 疏

\chapter{卷七}

\section{甘誓第二}

啟與有扈戰於甘之野，作\CJKunderwave{甘誓}。\footnote{夏啟嗣\CJKunderline{禹}位，伐有扈之罪。○啟，\CJKunderline{禹}子，嗣\CJKunderline{禹}為天子也。扈音戶。有扈，國名，與夏同姓。馬云：“姒姓之國，為無道者。”案京兆鄠縣即有扈之國也。甘，有扈郊地名，馬云：“南郊地也。”甘，水名，今在鄠縣西。誓，馬云：“軍旅曰誓，會同曰誥。”}

{\noindent\zhuan\zihao{6}\fzbyks 傳“夏啟”至“之罪”。正義曰：\CJKunderwave{孟子}稱，\CJKunderline{禹}薦益於天，七年，\CJKunderline{禹}崩之後，益避啟於箕山之陰，天下諸侯不歸益而歸啟，曰：“吾君之子也。”啟遂即天子位。\CJKunderwave{史記·夏本紀}稱，啟立，有扈氏不服，故伐之。蓋由自堯舜受禪相承，啟獨見繼父，以此不服，故云“夏啟嗣\CJKunderline{禹}立，伐有扈之罪”，言繼立者，見其由嗣立，故不服也。 \par}

{\noindent\shu\zihao{5}\fzkt “啟與”至“甘誓”。正義曰：夏王啟之時，諸侯有扈氏叛，王命率眾親征之。有扈氏發兵拒啟，啟與戰於甘地之野。將戰，集將士而誓戒之。史敘其事，作\CJKunderwave{甘誓}。 \par}

甘誓\footnote{甘，有扈郊地名。將戰先誓。}

{\noindent\zhuan\zihao{6}\fzbyks 傳“甘有”至“先誓”。正義曰：\CJKunderwave{地理志}扶風鄠縣，古扈國，夏啟所伐者也。“鄠”、“扈”音同，未知何故改也。啟伐有扈，必將至其國，乃出兵與啟戰，故以“甘”為有扈之郊地名。馬融云：“甘，有扈南郊地名。”計啟西行伐之，當在東郊。融則扶風人，或當知其處也。“將戰先誓”，誓是臨戰時也。\CJKunderwave{甘誓}、\CJKunderwave{牧誓}、\CJKunderwave{費誓}皆取誓地為名，\CJKunderwave{湯誓}舉其王號，\CJKunderwave{泰誓}不言“武誓”者，皆史官不同，故立名有異耳。\CJKunderwave{泰誓}未戰而誓，故別為之名。\CJKunderwave{秦誓}自悔而誓，非為戰誓，自約其心，故舉其國名。 \par}

{\noindent\shu\zihao{5}\fzkt “甘誓”。正義曰：發首二句敘其誓之由，其“王曰”已下皆是誓之辭也。\CJKunderwave{曲禮}云：“約信曰誓。”將與敵戰，恐其損敗，與將士設約，示賞罰之信也。將戰而誓，是誓之大者。\CJKunderwave{禮}將祭而號令齊百官，亦謂之誓。\CJKunderwave{周禮·大宰}云：“祀五帝則掌百官之誓戒。”\CJKunderline{鄭玄}云：“誓戒，要之以刑，重失禮也。”\CJKunderwave{明堂位}所謂“各揚其職，百官廢職，服大刑”，是誓辭之略也。彼亦是約信，但小於戰之誓。馬融云：“軍旅曰誓，會同曰誥。”“誥”、“誓”俱是號令之辭，意小異耳。 \par}

大戰於甘，乃召六卿。\footnote{天子六軍，其將皆命卿。將，子匠反。}王曰：“嗟！六事之人，\footnote{各有軍事，故曰六事。}予誓告汝：有扈氏威侮五行，怠棄三正，\footnote{五行之德，王者相承所取法。有扈與夏同姓，恃親而不恭，是則威虐侮慢五行，怠惰棄廢天地人之正道。言亂常。侮，亡甫反。正如字，徐音徵，馬云：“建子、建醜、建寅，三正也。”惰，徒臥反。}天用剿絕其命，\footnote{用其失道故。剿，截也。截絕，謂滅之。剿，子六反，\CJKunderwave{玉篇}子小反，馬本作巢，與\CJKunderwave{玉篇}、\CJKunderwave{切韻}同。}今予惟恭行天之罰。\footnote{恭，奉也，言欲截絕之。罰音伐。}


{\noindent\zhuan\zihao{6}\fzbyks 傳“天子”至“命卿”。正義曰：將戰而召六卿，明是卿為軍將。“天子六軍，其將皆命卿”，\CJKunderwave{周禮·夏官序}文也。\CJKunderline{鄭玄}云：“夏亦然，則三王同也。”經言“大戰”者，\CJKunderline{鄭玄}云：“天子之兵，故曰大。”孔無明說，蓋以六軍並行，威震多大，故稱“大戰”。 \par}

{\noindent\zhuan\zihao{6}\fzbyks 傳“各有”至“六事”。正義曰：卿為軍將，故云“乃召六卿”，及其誓之，非六卿而已。\CJKunderline{鄭玄}云：“變六卿言六事之人者，言軍吏下及士卒也。”下文戒左右與御，是遍敕在軍之士，步卒亦在其間。六卿之身及所部之人，各有軍事，故“六事之人”為總呼之辭。 \par}

{\noindent\zhuan\zihao{6}\fzbyks 傳“五行”至“亂帝”。正義曰：“五行”，水、火、金、木、土也。分行四時，各有其德。\CJKunderwave{月令}孟春三日，太史謁於天子，曰：“某日立春，盛德在木。”夏雲“盛德在火”，秋雲“盛德在金”，冬雲“盛德在水”。此五行之德，王者雖易姓，相承其所取法同也。言王者共所取法，而有扈氏獨侮慢之，所以為大罪也。且五行在人為仁、義、禮、智、信，威侮五行，亦為侮慢此五常而不行也。有扈與夏同姓，恃親而不恭天子，廢君臣之義，失相親之恩，五常之道盡矣，是“威侮五行”也。無所畏忌，作威虐而侮慢之，故云“威虐侮慢”。\CJKunderwave{易·說卦}云：“立天之道曰陰與陽，立地之道曰柔與剛，立人之道曰仁與義。”物之為大，無大於此者，\CJKunderwave{周易}謂之“三才”。人生天地之間，莫不法天地而行事，以此知“怠惰棄廢天地人之正道”。棄廢此道，“言亂常也”。孔、馬、鄭、王與皇甫謐等皆言有扈與夏同姓，並依\CJKunderwave{世本}之文。\CJKunderwave{楚語}云，昭王使觀射父傅太子，射父辭之曰：“堯有\CJKunderline{丹朱}，舜有商均，夏有觀扈，周有管蔡。”是其“恃親而不恭”也。\CJKunderwave{周語}云，帝嘉\CJKunderline{禹}德，賜姓曰姒，\CJKunderline{禹}始得姓。有扈與夏同姓，則為啟之兄弟。知此者，蓋\CJKunderline{禹}未賜姓之前，以姒為姓，故\CJKunderline{禹}之親屬舊已姓姒，帝嘉其德，又以姒姓顯揚之。猶若\CJKunderline{\CJKunderline{伯夷}}\CJKunderwave{國語}稱賜姓曰姜，然\CJKunderline{\CJKunderline{伯夷}}是炎帝之後，未賜姓之前先為姜姓，與此同也。故有扈以為夏之同姓。 \par}

{\noindent\zhuan\zihao{6}\fzbyks 傳“用其”至“滅之”。正義曰：天子用兵，稱“恭行天罰”，諸侯討有罪，稱“肅將王誅”，皆示有所稟承，不敢專也。有扈既有大罪，宜其絕滅，故原天之意，言天用其失道之故，欲截絕其命，謂滅之也。“剿”是斬斷之義，故為截也。 \par}

左不攻於左，汝不恭命。\footnote{左，車左，左方主射。絕之也，治其職。}右不攻於右，汝不恭命。\footnote{右，車右，勇力之士，執戈矛以退敵。}御非其馬之正，汝不恭命。\footnote{御以正馬為政。三者有失，皆不奉我命。御，魚慮反。}


{\noindent\zhuan\zihao{6}\fzbyks 傳“左車”至“其職”。正義曰：歷言“左”、“右”及“御”，此三人人在一車之上也，故“左”為車左，則“右”為車右明矣宣十二年\CJKunderwave{左傳}云：“楚許伯御樂伯，攝叔為右，以致晉師。樂伯曰：‘吾聞致師者，左射以菆。’攝叔曰：‘吾聞致師者，右入壘，折馘執俘而還。’”是左方主射，右主擊刺，而御居中也。御言“正馬”，而左右不言所職者，以戰主殺敵，左右用兵是戰之常事，故略而不言；御惟主馬，故特言之，互相明也。此謂凡常兵車，甲士三人，所主皆如此耳。若將之兵車，則御者在左，勇力之士在右，將居鼓下，在中央，主擊鼓，與軍人為節度。成二年\CJKunderwave{左傳}說晉伐齊云：“晉解張御郤克，\CJKunderline{鄭玄}緩為右。卻克傷於矢，未絕鼓音，曰：‘餘病矣。’張侯曰：‘自始合，而矢貫餘手及肘，餘折以御，左輪朱殷，豈敢言病？’”郤克傷於矢而鼓音未絕，張侯為御而血染左輪，是御在左而將居中也。“攻”之為治，常訓也。“治其職”者，左當射人，右當擊刺，是其所掌職事也。 \par}

{\noindent\zhuan\zihao{6}\fzbyks 傳“御以”至“我命”。正義曰：“御以正馬為政”，言御之政事，事在正馬，故馬不正則罪之。\CJKunderwave{詩}云：“兩驂如手。”傳云：“進止如御者之手。”是為馬之正也。左、右與御三者有失，言“皆不奉我命”，以御在後，故總解之。 \par}

用命，賞於祖。\footnote{天子親征，必載遷廟之祖主行，有功則賞祖主前，示不專。}弗用命，戮於社，\footnote{天子親征，又載社主，謂之社事，不用命奔北者，則戮之於社主前。社主陰，陰主殺，親祖嚴社之義。戮音六。北如字，又音佩，軍走曰北。}予則孥戮汝。”\footnote{孥，子也。非但止汝身，辱及汝子。言恥累也。孥音奴，累，劣偽反。}

{\noindent\zhuan\zihao{6}\fzbyks 傳“天子”至“不專”。正義曰：\CJKunderwave{曾子問}云：“\CJKunderline{孔子}曰：‘天子巡守，以遷廟之主行，載於齊車，言必有尊也。’”巡守尚然，征伐必也。故云“天子親征，必載遷廟之祖主行，有功則賞祖主前，示不專”也。\CJKunderwave{周禮·大司馬}云：“若師不功，則厭而奉主車。”\CJKunderline{鄭玄}云：“厭，伏冠也。奉，猶送也。”送主歸於廟與社，亦是征伐載主之事也。 \par}

{\noindent\zhuan\zihao{6}\fzbyks 傳“天子”至“之義”。正義曰：定四年\CJKunderwave{左傳}云：“君以軍行，祓社釁鼓，祝奉以從。”是天子親征，又載社主行也。\CJKunderwave{郊特牲}云：“惟為社事，單出裡。”故以“社事”言之。“不用命奔北者，則戮之於社主之前”，“奔北”謂背陳走也。所以刑賞異處者，社主陰，陰主殺，則祖主陽，陽主生。\CJKunderwave{禮}左宗廟，右社稷，是祖陽而社陰。就祖賞，就社殺，親祖嚴社之義也。大功大罪則在軍賞罰，其遍敘諸勳，乃至太祖賞耳。 \par}

{\noindent\zhuan\zihao{6}\fzbyks 傳“孥子”至“恥累也”。正義曰：\CJKunderwave{詩}雲“樂爾妻孥”，對“妻”別文，是“孥”為子也。非但止辱汝身，並及汝子亦殺，言以恥惡累之。\CJKunderwave{湯誓}云：“予則孥戮汝。”傳曰：“古之用刑，父子兄弟罪不相及，今雲‘孥戮汝’，權以脅之，使勿犯。”此亦然也。 \par}

{\noindent\shu\zihao{5}\fzkt “大戰”至“戮汝”。正義曰：史官自先敘其事，啟與有扈大戰於甘之野，將欲交戰，乃召六卿，令與眾士俱集。王乃言曰：“嗟！”重其事，故嗟嘆而呼之：“汝六卿者，各有軍事之人。我設要誓之言以敕告汝：今有扈氏威虐侮慢五行之盛德，怠惰棄廢三才之正道，上天用失道之故，今欲截絕其命。天既如此，故我今惟奉行天之威罰，不敢違天也。我既奉天，汝當奉我。汝諸士眾在車左者，不治理於車左之事，是汝不奉我命。在車右者，不治理於車右之事，是汝不奉我命。御車者非其馬之正，令馬進退違戾，是汝不奉我命。汝等若用我命，我則賞之於祖主之前。若不用我命，則戮之於社主之前。所戮者非但止汝身而已，我則並殺汝子以戮辱汝。汝等不可不用我命以求殺敵。”戒之使齊力戰也。 \par}

\section{五子之歌第三【偽】}


太康失邦，\footnote{啟子也。盤於遊田，不恤民事，為羿所逐，不得反國。}昆弟五人須於洛汭，作\CJKunderwave{五子之歌}。\footnote{太康五弟與其母待太康於洛水之北，怨其不反，故作歌。五子名字，書傳無聞，仲康蓋其一也。須，馬云：“止也。”汭，如銳反，本又作內，音同。}


{\noindent\zhuan\zihao{6}\fzbyks 傳“太康”至“作歌”。正義曰：“昆弟五人”，自有長幼，故稱“昆弟”，嫌是太康之昆，故云“太康之五弟”。 \par}

{\noindent\shu\zihao{5}\fzkt “太康”至“之歌”。正義曰：啟子太康，以遊畋棄民，為羿所逐，失其邦國。其未失國之前，畋於洛水之表，太康之弟，更有昆弟五人，從太康畋獵，與其母待太康於洛水之北。太康為羿所距，不得反國，其弟五人，即啟之五子，並怨太康，各自作歌。史敘其事，作\CJKunderwave{五子之歌}。 \par}

五子之歌\footnote{啟之五子，因以名篇。}


{\noindent\zhuan\zihao{6}\fzbyks 傳“啟之”至“名篇”。正義曰：直言“五子”，不知謂誰，故言“啟之五子”。太康之弟，敘怨作歌，不言“五弟”而言“五子”者，以其迷祖之訓，故系父以言之。 \par}

{\noindent\shu\zihao{5}\fzkt “五子之歌”。正義曰：史述作歌之由，先敘失國之事，“其一曰”以下乃是歌辭。此五子作歌五章，每章各是一人之作，辭相連接，自為終始。初言“皇祖有訓”，未必則指怨太康。必是五子之歌相顧，從輕至甚。“其一”、“其二”蓋是昆弟之次，或是作歌之次，不可知也。 \par}

太康尸位以逸豫，\footnote{屍，主也。主以尊位，為逸豫不勤。逸，本又作佾。豫，本又作忬，音同。}滅厥德，黎民咸貳。\footnote{君喪其德，則眾民皆二心矣。黎，力兮反。喪，息浪反。}乃盤遊無度，\footnote{盤樂遊逸無法度。盤，步幹反，本或作槃。度如字。樂如字。}畋於有洛之表，十旬弗反。\footnote{洛水之表，水之南。十日曰旬。田獵過百日不還。畋音田。}有窮后羿，因民弗忍，距於河。\footnote{有窮，國名。羿，諸侯名。距太康於河，不得入國，遂廢之。羿，五計反，徐胡細反。距音巨。}厥弟五人御其母以從，\footnote{御，待也，言從畋。從如字，或作才用反，非。}徯於洛之汭。五子咸怨，\footnote{待太康，怨其久畋失國。徯，胡啟反。}述\CJKunderline{大禹}之戒以作歌。\footnote{述，循也。歌以敘怨。}

{\noindent\zhuan\zihao{6}\fzbyks 傳“屍，主也”。正義曰：\CJKunderwave{釋詁}文。 \par}

{\noindent\zhuan\zihao{6}\fzbyks 傳“有窮”至“廢之”。正義曰：襄四年\CJKunderwave{左傳}曰：“夏之方衰也，后羿自鉏遷於窮石。”然則羿居窮石，故曰“有窮，國名”。“窮”是諸侯之國，“羿”是其君之名也。\CJKunderwave{說文}云：“羿，帝嚳射官也。”賈逵云：“羿之先祖，世為先王射官，故帝賜羿弓矢使司射。”\CJKunderwave{淮南子}云：“堯時十日並生，堯使羿射九日而落之。”\CJKunderwave{楚辭·天問}云：“羿焉彃日烏解羽？”\CJKunderwave{歸藏易}亦云：“羿彃十日。”\CJKunderwave{說文}云：“彃者，射也。”此三者言雖不經以取信，要言帝嚳時有羿，堯時亦有羿，則羿是善射之號，非復人之名字。信如彼言，則不知羿名為何也。夏都河北，洛在河南，距太康於河北，不得入國，遂廢太康耳。羿猶立仲康，不自立也。 \par}

{\noindent\zhuan\zihao{6}\fzbyks 傳“述循”至“敘怨”。正義曰：“述，循”，\CJKunderwave{釋詁}文。循其所戒，用作歌以敘怨也。其一曰“皇祖有訓”，其二曰“訓有之”，是“述\CJKunderline{大禹}之戒”也。其三恨亡國都，其四恨絕宗祀，其五言追悔無及，直是指怨太康，非為述祖戒也。本述戒作歌，因即言及時事，故言祖戒以總之。 \par}

{\noindent\shu\zihao{5}\fzkt “太康”至“作歌”。正義曰：天子之在天位，職當牧養兆民。太康主以尊位，用為逸豫，滅其人君之德，眾人皆有二心。太康乃復愛樂遊，逸無有法度，畋獵於洛水之表，一出十旬不反。有窮國君其名曰羿，因民不能堪忍太康之惡，率眾距之於河，不得反國。太康初去之時，其弟五人侍其母以從太康。太康畋於洛南，五弟待於洛北，太康久而不反，致使羿距於河。五子皆怨太康，追述\CJKunderline{大禹}之戒以作歌，而各敘已怨之志也。其弟待母以從太康，太康初去即然。待於洛水之北，以冀太康速反。羿既距之，五子乃怨。史述太康之惡既盡，然後言其作歌，故令“羿距”之文乃在“母從”之上，作文之勢當然也。 \par}

其一曰：“皇祖有訓，民可近，不可下，\footnote{皇，君也。君祖\CJKunderline{禹}有訓戒。近謂親之。下謂失分。近，附近之近。分，扶問反。}民惟邦本，本固邦寧。\footnote{言人君當固民以安國。}予視天下，愚夫愚婦，一能勝予，\footnote{言能畏敬小民，所以得眾心。}一人三失，怨豈在明？不見是圖。\footnote{三失，過非一也。不見是謀，備其微。三如字，又息暫反。見，賢遍反。}予臨兆民，懍乎若朽索之馭六馬，\footnote{十萬曰億，十億曰兆，言多。懍，危貌。朽，腐也。腐索馭六馬，言危懼甚。懍，力甚反。朽，許久反。馭音御。腐，扶甫反。}為人上者，奈何不敬？”\footnote{能敬則不驕，在上不驕，則高而不危。}


{\noindent\zhuan\zihao{6}\fzbyks 傳“皇君”至“失分”。正義曰：“皇，君”，\CJKunderwave{釋詁}文。述\CJKunderline{禹}之戒，知“君祖”是\CJKunderline{禹}，\CJKunderline{禹}有訓也。“民可近”者，據“君”為文。“近”謂親近之也，“下”謂卑下輕忽之，失本分也。奪其農時，勞以橫役，是失分也，故下雲“予視天下，愚夫愚婦，一能勝予”，是畏敬下民也。 \par}

{\noindent\zhuan\zihao{6}\fzbyks 傳“言能”至“眾心”。正義曰：我視愚夫愚婦，當能勝我身，是畏敬小民也。由能畏敬小民，故以小民從命，是“得眾心”也。 \par}

{\noindent\zhuan\zihao{6}\fzbyks 傳“三失”至“其微”。正義曰：顧氏云：“怨豈在明？未必皆在明著之時，必於未形之日思善道以自防衛之。”是備慎其微也。 \par}

{\noindent\zhuan\zihao{6}\fzbyks 傳“十萬”至“懼甚”。正義曰：古數十萬曰億，十億曰兆，言多也。懍懍，心懼之意，故為危貌。“朽，腐”，常訓也。腐索馭六馬，索絕馬驚，馬驚則逸，言危懼甚也。經傳之文，惟此言“六馬”，漢世此經不傳，餘書多言駕四者，\CJKunderwave{春秋·公羊}說天子駕六，\CJKunderwave{毛詩}說天子至大夫皆駕四，許慎案\CJKunderwave{王度記}雲天子駕六，\CJKunderline{鄭玄}以\CJKunderwave{周禮}校人養馬，“乘馬一師四圉”，四馬曰乘，\CJKunderwave{康王之誥}雲“皆布乘黃朱”，以為天子駕四。漢世天子駕六，非常法也。然則此言馬多懼深，故舉六以言之。 \par}

{\noindent\shu\zihao{5}\fzkt “其一”至“不敬”。正義曰：我君祖\CJKunderline{大禹}有訓戒之事，言民可親近，不可卑賤輕下。令其失分，則人懷怨，則事上之心不固矣。民惟邦國之本，本固則邦寧。言在上不可使人怨也。我視天下之民，愚夫愚婦，一能過勝我，安得不敬畏之也？所以畏其怨者，一人之身，三度有失；凡所過失，為人所怨，豈在明著？大過皆由小事而起。言小事不防，易致大過，故於不見細微之時，當於是豫圖謀之，使人不怨也。我臨兆民之上，常畏人怨，懍懍乎危懼，若腐索之馭六馬。索絕則馬逸，言危懼之甚。人之可畏如是，為民上者奈何不敬慎乎？怨太康之不恤下民也。 \par}

其二曰：“訓有之，內作色荒，外作禽荒。\footnote{作，為也。迷亂曰荒。色，女色。禽，鳥獸。}

{\noindent\zhuan\zihao{6}\fzbyks 傳“作為”至“鳥獸”。正義曰：“作，為”，\CJKunderwave{釋言}文。昭元年\CJKunderwave{左傳}晉平公近女色過度，惑以喪志。\CJKunderwave{老子}云：“馳騁田獵，令人心發狂。”好色好田則精神迷亂，故“迷亂曰荒”。女有美色，男子悅之，經傳通謂女人為“色”。獵則鳥獸並取，故以“禽”為鳥獸也。 \par}

甘酒嗜音，峻宇雕牆。\footnote{甘嗜無厭足。峻，高大。雕，飾畫。甘，一音戶甘反。嗜，市志反。峻，思俊反。牆，慈羊反。厭，於鹽反，又於豔反。有一於此，未或不亡。”此六者，棄德之君必有其一。}有一必亡，況兼有乎！

其三曰：“惟彼陶唐，有此冀方。\footnote{陶唐，\CJKunderline{帝堯}氏，都冀州，統天下四方。}

{\noindent\zhuan\zihao{6}\fzbyks 傳“陶唐”至“四方”。正義曰：\CJKunderwave{世本}云：“\CJKunderline{帝堯}為陶唐氏。”韋昭云：“陶、唐皆國名，猶湯稱殷商也。”案書傳皆言堯以唐侯升為天子，不言封於陶唐，“陶唐”二字或共為地名，未必如昭言也。以天子王有天下，非獨冀州一方，故以“冀方”為“都冀州，統天下四方”。堯都平陽，舜都蒲阪，\CJKunderline{禹}都安邑，相去不盈二百，皆在冀州，自堯以來其都不出此地，故舉陶唐以言之。 \par}

今失厥道，亂其紀綱，乃厎滅亡。”\footnote{言失堯之道，亂其法制，自致滅亡。厎，之履反。}

其四曰：“明明我祖，萬邦之君。有典有則，貽厥子孫。\footnote{君萬國為天子。典謂經籍。則，法。貽，遺也。言仁及後世。貽，以之反。遺，唯季反。}關石和鈞，王府則有。荒墜厥緒，覆宗絕祀。”\footnote{金鐵曰石，供民器用，通之使和平，則官民足。言古制存，而太康失其業，以取亡。覆，芳服反。供音恭。}


{\noindent\zhuan\zihao{6}\fzbyks 傳“君萬”至“後世”。正義曰：“萬邦之君”，謂君統萬國為天子也。“典”謂先王之典，可憑據而行之，故為經籍。“則，法”，\CJKunderwave{釋詁}文。“典”謂先王舊典，“法”謂當時所制，其事不為大異，重言以備文耳。“貽，遺”，\CJKunderwave{釋言}文。以典法遺子孫，言仁恩及後世。 \par}

{\noindent\zhuan\zihao{6}\fzbyks 傳“金鐵”至“取亡”。正義曰：“關”者，通也。名“石”而可通者，惟衡量之器耳。\CJKunderwave{律曆志}云：“二十四銖為兩，十六兩為斤，三十斤為鈞，四鈞為石。”是石為稱之最重，以石而稱則為重物，故“金鐵曰石”。言絲綿止於斤兩，金鐵乃至於石，舉“石”而言之，則所稱之物皆通之也。傳取金鐵重物以解言“石”之意，非謂所關通者惟金鐵耳。米粟則鬥斛以量之，布帛則丈尺以度之，惟言關通權衡，則度量之物，懋遷有無，亦關通矣，舉一以言之耳。衡石所稱之物，以供民之器用，其土或有或無，通使和平也。\CJKunderwave{論語}云：“百姓足，君孰與不足？”民既足用，則官亦富饒，故“通之使和平，則官民皆足”。有典有法可依而行，官民足可坐而守，言古制存，而太康失其業，所以亡也。訓“緒”為業，費氏、顧氏等意云，通金鐵於人，官不禁障，民得取之以供器用。器既具，所以上下充足。以金鐵皆從石而生，則金鐵亦石之類也。故\CJKunderwave{漢書·五行志}雲石為怪異，入金不從革之條。費、顧之義，亦得通也。 \par}

{\noindent\shu\zihao{5}\fzkt “其四”至“絕祀”。正義曰：有明明之德，我祖\CJKunderline{大禹}也。以有明德為萬邦之君，謂為天子也。有治國之典，有為君之法，遺其後世之子孫，使法則之。又關通衡石之用，使之和平。人既足用，王之府藏則皆有矣。典存國富，宜以為政，今太康荒廢墜失其業，覆滅宗族，斷絕祭祀。言太康棄典法，所以滅宗祀也。 \par}

其五曰：“嗚呼！曷歸？予懷之悲。\footnote{曷，何也。言思而悲。曷，戶割反。}萬姓仇予，予將疇依？\footnote{仇，怨也。言當依誰以復國乎？}鬱陶乎予心，顏厚有忸怩。\footnote{鬱陶，言哀思也。顏厚，色愧。忸怩，心慚，慚愧於仁人賢士。鬱音蔚。陶音桃。鬱陶，憂思也。忸，女六反。怩，女姬反，徐乃私反。思，息嗣反。}弗慎厥德，雖悔可追？”\footnote{言人君行己不慎其德，以速滅敗，雖欲改悔，其可追及乎？言無益。雖如字，或作雎。}


{\noindent\zhuan\zihao{6}\fzbyks 傳“仇怨”至“國乎”。正義曰：桓二年\CJKunderwave{左傳}雲“怨耦曰仇”，故為怨也。羿距於河，不得復反，乃思太康，欲歸依之，言當依誰以復國乎？ \par}

{\noindent\zhuan\zihao{6}\fzbyks 傳“鬱陶”至“賢士”。正義曰：\CJKunderwave{孟子}稱舜弟象見舜云：“思君正鬱陶。”“鬱陶”，精神憤結積聚之意，故為哀思也。\CJKunderwave{詩}云：“顏之厚矣。”羞愧之情見於面貌，似如面皮厚然，故以“顏厚”為色愧忸怩，羞不能言，心慚之狀。小人不足以知得失，故“慚愧於仁人賢士”。 \par}

{\noindent\shu\zihao{5}\fzkt “其五”至“可追”。正義曰：嗚呼！太康已覆滅矣，我將何所依歸？我以此故，思之而悲。太康為惡，毒遍天下，萬姓皆共仇我，我將誰依就乎？鬱陶而哀思乎，我之心也！我以此故，外貌顏厚而內情忸怩羞慚。由太康不慎其德，以致此見距，雖欲改悔，其可追及之乎？事已往矣，不可如何。從首漸怨，至此為深，皆是羿距時事也。 \par}

\section{胤徵第四【偽】}


羲和湎淫，廢時亂日，\footnote{羲氏、和氏，世掌天地四時之官，自唐虞至三代，世職不絕。承太康之後，沈湎於酒，過差非度，廢天時，亂甲乙。湎，徐音緬，面善反。差，初賣反，又初佳反。}胤往徵之，作\CJKunderwave{胤徵}。\footnote{胤國之君受王命往徵之。胤，國名。}

胤徵\footnote{奉辭罰罪曰徵。}


{\noindent\zhuan\zihao{6}\fzbyks 傳“羲氏”至“甲乙”。正義曰：“羲氏、和氏世掌天地四時之官”，\CJKunderwave{堯典}所言是其事也。羲和是重黎之後，\CJKunderwave{楚語}稱堯育重黎之後，使典天地，以至於夏商，是“自唐虞至三代，世職不絕”，故此時羲和仍掌時日。以太康逸豫，臣亦縱弛。此承太康之後，於今仍亦懈惰，沈湎於酒，過差非度，廢天時，亂甲乙，是其罪也。經雲“酒荒於厥邑”，惟言荒酒，不言好色，故訓“淫”為過，言耽酒為過差也。聖人作歷數以紀天時，不存歷數，是“廢天時”也。日以甲乙為紀，不知日食，是“亂甲乙”也。 \par}

{\noindent\zhuan\zihao{6}\fzbyks 傳“奉辭罰罪”。正義曰：奉責讓之辭，伐不恭之罪，名之曰“徵”。徵者，正也。伐之以正其罪。 \par}

{\noindent\shu\zihao{5}\fzkt “羲和”至“胤徵”。正義曰：羲氏、和氏，世掌天地四時之官，今乃沈湎於酒，過差非度，廢天時，亂甲乙，不以所掌為意，胤國之侯受王命往徵之。史敘其事，作\CJKunderwave{胤徵}。 \par}

惟仲康肇位四海，\footnote{羿廢太康，而立其弟仲康為天子。○肇音兆。}胤侯命掌六師。\footnote{仲康命胤侯掌王六師，為大司馬。}羲和廢厥職，酒荒於厥邑，\footnote{舍其職官，還其私邑，以酒迷亂，不修其業。舍音舍。}胤後承王命徂徵。\footnote{徂，往也，就其私邑往討之。}


{\noindent\zhuan\zihao{6}\fzbyks 傳“羿廢”至“天子”。正義曰：以羿距太康於河，於時必廢之也。\CJKunderwave{夏本紀}云：“太康崩，弟仲康立。”襄四年\CJKunderwave{左傳}云：“羿因夏民以伐夏政。”則羿於其後篡天子之位，仲康不能殺羿，必是羿握其權，知仲康之立，是羿立之矣，故云“羿廢太康，而立其弟仲康為天子”。計\CJKunderwave{五子之歌}，仲康當是其一，仲康必賢於太康，但形勢既衰，政由羿耳。羿在夏世為一代大賊，\CJKunderwave{左傳}稱羿既篡位，寒浞殺之。羿滅夏後相，相子少康始滅浞復夏政。計羿、浞相承，向有百載，為夏亂甚矣。而\CJKunderwave{夏本紀}云：“太康崩，其弟仲康立。仲康崩，子相立。相崩，子少康立。”都不言羿、浞之事，是馬遷之說疏矣。 \par}

{\noindent\shu\zihao{5}\fzkt “惟仲康”至“徂徵”。正義曰：惟仲康始即王位，臨四海，胤國之侯受王命為大司馬，掌六師。於是有羲氏、和氏廢其所掌之職，縱酒荒迷，亂於私邑，胤國之君承王命往徵之。 \par}

告於眾曰：“嗟予有眾，\footnote{誓敕之。}聖有謨訓，明徵定保。\footnote{徵，證。保，安也。聖人所謀之教訓，為世明證，所以定國安家。}先王克謹天戒，臣人克有常憲，\footnote{言君能慎戒，臣能奉有常法。}百官修輔，厥後惟明明。\footnote{修職輔君，君臣俱明。}每歲孟春，遒人以木鐸徇於路，\footnote{遒人，宣令之官。木鐸，金鈴木舌，所以振文教。遒，在由反。鐸，待洛反。鈴音令。}官師相規，工執藝事以諫。\footnote{官眾，從官。更相規闕，百工各執其所治技藝以諫，諫失常。藝，本又作埶。更音庚。技，其綺反。}其或不恭，邦有常刑。\footnote{言百官廢職，服大刑。}

{\noindent\zhuan\zihao{6}\fzbyks 傳“徵證”至“安家”。正義曰：成八年\CJKunderwave{左傳}稱晉殺趙括，欒、郤為徵。“徵”是證驗之義，故為證也。能自保守是安定之義，故為安也。聖人將為教訓，必謀而後行，故言“所謀之教訓”。聖人之言，必有其驗，故為“世之明證”。用聖人之謨訓，必有成功，故“所以定國安家”。 \par}

{\noindent\zhuan\zihao{6}\fzbyks 傳“言君”至“常法”。正義曰：王者代天理官，故稱“天戒”。臣人奉主法令，故言“常憲”。君當奉天，臣當奉君，言君能戒慎天戒也，臣能奉有常法，奉行君法也。此謂大臣，下雲“百官修輔”，謂眾臣。 \par}

{\noindent\zhuan\zihao{6}\fzbyks 傳“遒人”至“文教”。正義曰：以執木鐸徇於路，是宣令之事，故言“宣令之官”。\CJKunderwave{周禮}無此官，惟\CJKunderwave{小宰}云：“正歲，帥理官之屬而觀治象之法，徇以木鐸曰：‘不用法者，國有常刑。’”宣令之事，略與此同。此似別置其官，非如周之小宰。名曰“遒人”，不知其意，蓋訓“遒”為聚，聚人而令之，故以為名也。\CJKunderwave{禮}有“金鐸”、“木鐸”，“鐸”是鈴也，其體以金為之，明舌有金木之異，知木鐸是木舌也。\CJKunderwave{周禮}教鼓人“以金鐸通鼓”，大司馬“教振旅，兩司馬執鐸”，\CJKunderwave{明堂位}雲“振木鐸於朝”，是武事振金鐸，文事振木鐸。今雲“木鐸”，故云“所以振文教”也。 \par}

{\noindent\zhuan\zihao{6}\fzbyks 傳“官眾”至“失常”。正義曰：“相規”，相平等之辭，故“官眾”謂“眾官”，“相規”謂“更相規闕”。平等有闕，已尚相規，見上之過，諫之必矣。“百工各執其所治技藝以諫”，謂被遣作器，工有奢儉，若\CJKunderwave{月令}雲“無作淫巧，以蕩上心”，見其淫巧不正，當執之以諫，諫失常也。百工之賤，猶令進諫，則百工以上，不得不諫矣。 \par}

{\noindent\zhuan\zihao{6}\fzbyks 傳“言百”至“大刑”。正義曰：“百官廢職，服大刑”，\CJKunderwave{明堂位}文也。顧氏云：“百官群臣其有廢職懈怠不恭謹者，國家當有常刑。” \par}

{\noindent\shu\zihao{5}\fzkt “告於”至“常刑”。正義曰：胤侯將徵羲和，告於所部之眾曰：“嗟乎！我所有之眾人，聖人有謨之訓，所以為世之明證，可以定國安家。其所謀者，言先王能謹慎敬畏天戒，臣人者能奉先王常法，百官修常職輔其君，君臣相與如是，則君臣俱明，惟為明君明臣。”言君當謹慎以畏天，臣當守職以輔君也。“先王恐其不然，大開諫爭之路。每歲孟春，遒人之官以木鐸徇於道路，以號令臣下，使在官之眾更相規闕；百工雖賤，令執其藝能之事以諫上之失常。其有違諫不恭謹者，國家則有常刑”。 \par}

“惟時羲和，顛覆厥德，\footnote{顛覆言反倒。將陳羲和所犯，故先舉孟春之令，犯令之誅。覆，芳服反。倒，丁老反。}沈亂於酒，畔宮離次，\footnote{沈謂醉冥。失次位也。離如字，又力智反。冥，莫定反，又亡丁反。}俶擾天紀，遐棄厥司。\footnote{俶，始。擾，亂。遐，遠也。紀謂時日。司，所主也。俶，本又作㑐，亦作叔，同尺六反。擾，而小反。}乃季秋月朔，辰弗集於房，\footnote{辰，日月所會。房，所舍之次。集，合也。不合即日食可知。}

{\noindent\zhuan\zihao{6}\fzbyks 傳“顛覆”至“之誅”。正義曰：“顛覆言反倒”，謂人反倒也。人當豎立，今乃反倒，猶臣當事君，今乃廢職，似人之反倒然。言臣以事君為德，故言“顛覆厥德”。胤侯將陳羲和之罪，故先舉孟春之令，犯令之誅，舉輕以見重，小事犯令猶有常刑，況叛官離次為大罪乎！ \par}

{\noindent\zhuan\zihao{6}\fzbyks 傳“沈謂”至“次位”。正義曰：沒水謂之沈，大醉冥然，無所復知，猶瀋水然，故謂醉為“沈”。 \par}

{\noindent\zhuan\zihao{6}\fzbyks 傳“俶始”至“所主”。正義曰：“俶，始”、“遐，遠”皆\CJKunderwave{釋詁}文。“擾”謂煩亂，故為亂也。\CJKunderwave{洪範}五紀，“五曰歷數”，歷數所以紀天時。此言“天紀”，謂時日。此時日之事是羲和所司，言棄其所主。 \par}

{\noindent\zhuan\zihao{6}\fzbyks 傳“辰日”至“可知”。正義曰：昭七年\CJKunderwave{左傳}曰，晉侯問於士文伯曰：“何謂辰？”對曰：“日月之會是謂辰。”是“辰”為日月之會。日月俱右行於天，日行遲，月行疾，日每日行一度，月日行十三度十九分度之七，計二十九日過半，月已行天一周，又逐及日而與日聚會，謂此聚會為“辰”。一歲十二會，故為十二辰，即子、醜、寅、卯之屬是也。“房”謂室之房也，故為“所舍之次”。計九月之朔，日月當會於大火之次。\CJKunderwave{釋言}云：“集，會也。”會即是合，故為“合”也。日月當聚會共舍，今言日月不合於舍，則是日食可知也。日食者，月掩之也，月體掩日，日被月映，即不成共處，故以不集言日食也。或以為“房”謂房星，九月日月會於大火之次。房、心共為大火，言辰在房星，事有似矣。知不然者，以集是止舍之處，言其不集於舍，故得以表日食；若言不集於房星，似太遲太疾，惟可見歷錯，不得以表日食也。且日之所在，星宿不見，正可推算以知之，非能舉目見之。君子慎疑，寧當以日在之宿為文，以此知其必非房星也。 \par}

瞽奏鼓、嗇夫馳，庶人走。\footnote{凡日食，天子伐鼓於社，責上公。瞽，樂官，樂官進鼓則伐之。嗇夫，主幣之官，馳取幣禮天神。眾人走，供救日食之百役也。嗇音色。馳，車馬曰馳。走，步曰走。供音恭。}羲和屍厥官，罔聞知，\footnote{主其官而無聞知於日食之變異，所以罪重。}昏迷於天象，以幹先王之誅。\footnote{暗錯天象，言昏亂之甚。幹，犯也。}政典曰：‘先時者殺無赦，\footnote{政典，夏後為政之典籍。若\CJKunderwave{周官}六卿之治典。先時，謂曆象之法，四時節氣，弦望晦朔。先天時則罪死無赦。先，悉薦反，又如字，注“先時”、“先天”同。赦亦作亦攵。治，直吏反。}不及時者殺無赦。’\footnote{不及謂曆象後天時。雖治其官，苟有先後之差，則無赦，況廢官乎！後，胡豆反。}

{\noindent\zhuan\zihao{6}\fzbyks 傳“凡日”至“百役”。正義曰：文十五年\CJKunderwave{左傳}云：“日有食之，天子不舉，伐鼓於社。諸侯用幣於社，伐鼓於朝。”杜預以為“伐鼓於社”，“責群陰”也。此傳言“責上公”者，\CJKunderwave{郊特牲}云：“社祭土而主陰氣也，君南向北墉下，答陰之義也。”是言社主陰也。日食陰侵陽，故杜預以為“責群陰”也。昭二十九年\CJKunderwave{左傳}云：“封為上公，祀為貴神。社稷五祀，是尊是奉。”是社祭句龍為上公之神也。日食臣侵君之象，故傳以為“責上公”，亦當群陰上公並責之也。\CJKunderwave{周禮}瞽蒙之官掌作樂，瞽為樂官。樂官用無目之人，以其無目，於音聲審也。\CJKunderwave{詩}雲“奏鼓簡簡”，謂伐鼓為奏鼓，知“樂官進鼓則伐之”。\CJKunderwave{周禮·太僕}：“軍旅、田役，贊王鼓。救日月，亦如之。”\CJKunderline{鄭玄}云：“王通鼓，佐聲其餘面。”則救日之時，王或親鼓。莊二十五年\CJKunderwave{穀梁傳}曰：“天子救日，置五麾，陳五兵、五鼓。”陳既多，皆樂人伐之。\CJKunderwave{周禮}無嗇夫之官，\CJKunderwave{禮}云：“嗇夫承命，告於天子。”\CJKunderline{鄭玄}云：“嗇夫蓋司空之屬也。”嗇夫主幣，\CJKunderwave{禮}無其文，此雲“嗇夫馳”，必馳走有所取也。\CJKunderwave{左傳}云：“諸侯用幣。”則天子亦當有用幣之處，嗇夫必是主幣之官，馳取幣也。社神尊於諸侯，故諸侯用幣於社以請救。天子伐鼓於社，必不用，知嗇夫“馳取幣禮天神”。“庶人走”，蓋是庶人在官者，謂諸侯胥徒也。其走必有事，知為“供救日食之百役”也。\CJKunderwave{曾子問}云：“諸侯從天子救日食，各以方色與其兵。”\CJKunderwave{周禮·庭氏}云：“救日之弓矢。”是救日必有多役，庶人走供之。鄭注\CJKunderwave{庭氏}云，以救日為太陽之弓，救月為太陰之弓，救日以枉矢，救月以恆矢。其鼓則蓋用祭天之雷鼓也。昭十七年“夏六月甲戌朔，日有食之”。\CJKunderwave{左傳}云：“季平子曰：‘惟正月朔，慝未作，日有食之，於是乎有伐鼓用幣，禮也，其餘則否。’太史曰：‘在此月也，當夏四月，是謂孟夏。’”如彼傳文，惟夏四月有伐鼓用幣禮，余月則不然，此以九月日食亦奏鼓用幣者，顧氏云：“夏禮異於周禮也。” \par}

{\noindent\zhuan\zihao{6}\fzbyks 傳“政典”至“無赦”。正義曰：胤侯，夏之卿士，引政典而不言古典，則當時之書，知是“夏後為政之典籍”也。\CJKunderwave{周禮}：“太宰掌建邦之六典，以佐王治邦國。一曰治典，二曰教典，三曰禮典，四曰政典，五曰刑典，六曰事典。”若周官六卿之治典，謂此也。“先時”、“不及”者，謂此曆象之法，四時節氣，弦望晦朔，不得先天時，不得後天時。四時時各九十日有餘，分為八節，節各四十五日有餘也。節氣者，周天三百六十五日四分日之一，四時分之，均分為十二月，則月各得三十日十六分日之七，以初為節氣，半為中氣，故一歲有二十四氣也。計十二月，每月二十九日強半也。以月初為朔，月盡為晦，當月之中，日月相望，故以月半為望。望去晦、朔，皆不滿十五日也。又半此望去晦、朔之數，名之曰弦。弦者，言其月光正半如弓弦也。晦者，月盡無月，言其暗也。朔者，蘇也，言月死而更蘇也。先天時者，所名之日，在天時之先。假令天之正時，當以甲子為朔，今歷乃以癸亥為朔，是造歷先天時也；若以乙丑為朔，是造歷後天時也。後即是“不及時”也。其氣、望等皆亦如此。 \par}

{\noindent\shu\zihao{5}\fzkt “惟時”至“無赦”。正義曰：言不諫尚有刑，廢職懈怠，是為大罪。惟是羲和，顛倒其奉上之德，而沉沒昏亂於酒，違叛其所掌之官，離其所居位次，始亂天之紀綱，遠棄所主之事。乃季秋九月之朔，日月當合於辰。其日之辰，日月不合於舍；不得合辰，謂日被月食，日有食之。\CJKunderwave{禮}有救日之法，於是瞽人樂官進鼓而擊之，嗇夫馳騁而取幣以禮天神，庶人奔走供救日食之百役。此為災異之大，群官促遽若此，羲和主其官而不聞知日食，是大罪也。此羲和昏暗迷錯於天象，以犯先王之誅，此罪不可赦也。故先王為政之典曰：“主歷之官，為歷之法，節氣先天時者殺無赦，不及時者殺無赦。”失前失後尚猶合殺，況乎不知日食；其罪不可赦也，況彼罪之大。言已所以徵也。 \par}

“今予以爾有眾，奉將天罰。\footnote{將，行也。奉王命行王誅，謂殺湎淫之身，立其賢子弟。}

爾眾士同力王室，尚弼予欽承天子威命。\footnote{以天子威命督其士眾，使用命。}火炎崐岡，玉石俱焚。\footnote{山脊曰岡。崐山出玉，言火逸而害玉。崐音昆。}天吏逸德，烈於猛火。\footnote{逸，過也。天王之吏為過惡之德，其傷害天下，甚於火之害玉。猛火烈矣，又烈於火。}殲厥渠魁，脅從罔治。\footnote{殲，滅。渠，大。魁，帥也。指謂羲和罪人之身，其脅從距王師者，皆無治。殲，子廉反。魁，苦回反。脅，虛業反。帥，色類反。}舊染汙俗，咸與惟新。\footnote{言其餘人久染汙俗，本無噁心，皆與更新，一無所問。汙，烏故反，汙辱之汙；又音烏，涴泥著物也，一音烏臥反。}嗚呼！威克厥愛，允濟。\footnote{嘆能以威勝所愛，則必有成功。}愛克厥威，允罔功。\footnote{以愛勝威，無以濟眾，信無功。}其爾眾士，懋戒哉！”\footnote{言當勉以用命，戒以闢戮。懋音茂。闢音避。}


{\noindent\zhuan\zihao{6}\fzbyks 傳“將行”至“子弟”。正義曰：“將”之為行，常訓也。天欲加罪，王者順天之罰，則王誅也。“奉王命行王誅，謂殺淫湎之身”，羲和之罪，不及其嗣，故知殺其身，立其賢子弟。\CJKunderwave{楚語}云，重黎之後，世掌天地四時之官，至於夏商。則此不滅其族，故傳言此也。 \par}

{\noindent\zhuan\zihao{6}\fzbyks 傳“山脊”至“害玉”。正義曰：\CJKunderwave{釋山}云：“山脊，岡。”孫炎曰：“長山之脊也。”以崐山出玉，言火逸害玉，喻誅惡害善也。 \par}

{\noindent\zhuan\zihao{6}\fzbyks 傳“逸過”至“於火”。正義曰：“逸”即佚也，佚是淫縱之名，故為過也。“天王之吏”，言位貴而威高，乘貴勢而逞毒心，或睚眥而害良善，故為“過惡之德，其傷害天下，甚於火之害玉”。猛火為烈甚矣，又復烈之於火，言其害之深也。 \par}

{\noindent\zhuan\zihao{6}\fzbyks 傳“殲滅”至“無治”。正義曰：“殲，盡也”，\CJKunderwave{釋詁}文。舍人曰：“殲，眾之盡也。”眾皆死盡為殲也。“渠，大”、“魁，帥”，無正訓。以上“殲厥渠魁”謂滅其元首，故以“渠”為大，“魁”為帥。史傳因此謂賊之首領為“渠帥”，本原出於此。 \par}

{\noindent\shu\zihao{5}\fzkt “今予”至“戒哉”。正義曰：“羲和所犯如上，故今我用汝所有之眾，奉王命，行天罰。汝等眾士，當同心盡力於王室，庶幾輔我敬承天子之命，使我伐必克之。”又恐兵威所及，濫殺無辜，故假喻以戒之：“火炎崐山之岡，玉石俱被焚燒。天王之吏為過惡之德，則酷烈甚於猛火。宜誅惡存善，不得濫殺。滅其為惡大帥，罪止羲和之身，其被迫脅而從距王師者，皆無治責其罪。久染汙穢之俗，本無噁心，皆與惟德更新，一無所問。”又言將軍之法，必有殺戮。“嗚呼！”重其事，故嘆而言之。“將軍威嚴能勝其愛心，有罪者雖愛必誅，信有成功。若愛心勝其威嚴，親愛者有罪不殺，信無功矣”。言我雖愛汝，有罪必殺。其汝眾士宜勉力以戒慎哉！勿違我命以取殺也。 \par}

自契至於成湯，八遷。\footnote{十四世，凡八徙國都。契，息列反，殷之始祖。八遷之書，史唯見四。}湯始居亳，從先王居，\footnote{契父帝嚳都亳，湯自商丘遷焉，故曰“從先王居”。亳，旁各反，徐扶各反。嚳，苦毒反。}作\CJKunderwave{帝告}、\CJKunderwave{釐沃}。\footnote{告來居，治沃土，二篇皆亡。告，工毒反。釐，力之反。沃，徐烏酷反。此五亡篇，舊解是\CJKunderwave{夏書}，馬、鄭之徒以為\CJKunderwave{商書}，兩義並通。}


{\noindent\zhuan\zihao{6}\fzbyks 傳“十四”至“國都”。正義曰：\CJKunderwave{周語}曰：“玄王勤商，十四世而興。”“玄王”謂契也，勤殖功業十四世，至湯而興，為天子也。\CJKunderwave{殷本紀}云，契生昭明。“昭明卒，子相土立。相土卒，子昌若立。昌若卒，子曹圉立。曹圉卒，子冥立。冥卒，子振立。振卒，子微立。微卒，子報丁立。報丁卒，子報乙立。報乙卒，子報丙立。報丙卒，子主壬立。主壬卒，子主癸立。主癸卒，子天乙立。天乙是為成湯”是也。契至成湯十四世，凡八遷國都者，\CJKunderwave{商頌}雲“帝立子生商”，是契居商也。\CJKunderwave{世本}云：“昭明居砥石。”\CJKunderwave{左傳}稱相土居商丘，及今湯居亳，事見經傳者有此四遷，其餘四遷未詳聞也。\CJKunderline{鄭玄}云：“契本封商，國在太華之陽。”皇甫謐云：“今上洛商是也。”襄九年\CJKunderwave{左傳}云：“陶唐氏之火正閼伯居商丘，相土因之。”杜預云：“今梁國睢陽宋都是也。”其“砥石”先儒無言，不知所在。自契至湯，諸侯之國而得數遷都者，蓋以時王命之使遷。至湯乃以商為天下號，則都雖數遷，商名不改。今湯遷亳，乃作此篇，若是諸侯遷都，則不得史錄其事以為\CJKunderwave{商書}之首。文在“湯徵諸侯”、“伊尹去亳”之上，是湯將欲為王時事。史以商有天下，乃追錄初興，並\CJKunderwave{湯徵}與\CJKunderwave{汝鳩}、\CJKunderwave{汝方}，皆是伐桀前事，後追錄之也。 \par}

{\noindent\zhuan\zihao{6}\fzbyks 傳“契父”至“王居”。正義曰：“先王”，天子也。自契已下，皆是諸侯，且文稱契至湯，今雲“從先王居”者，必從契之先世天子所居也。\CJKunderwave{世本}、\CJKunderwave{本紀}皆雲契是帝嚳子，知先王是契父帝嚳。帝嚳本居亳，今湯往從之。嚳實帝也，言“先王”者，對文論優劣，則有皇與帝及王之別，散文則雖皇與帝皆得言王也。故\CJKunderwave{禮運}云：“昔者先王未有宮室。”乃謂上皇為王，是其類也。孔言“湯自商丘遷焉”，以相土之居商丘，其文見於\CJKunderwave{左傳}，因之言自商丘徙耳。此言不必然也，何則？相土，契之孫也，自契至湯凡八遷，若相土至湯都遂不改，豈契至相土三世而七遷也？相土至湯必更遷都，但不知湯從何地而遷亳耳，必不從商丘遷也。\CJKunderline{鄭玄}云：“亳，今河南偃師縣有湯亭。”\CJKunderwave{漢書音義}臣瓚者云：“湯居亳，今濟陰亳縣是也，今亳有湯冢。己氏有伊尹冢。”杜預云：“梁國蒙縣北有亳城，城中有成湯冢，其西又有伊尹冢。”皇甫謐云：“\CJKunderwave{孟子}稱湯居亳，與葛為鄰，葛伯不祀，湯使亳眾為之耕。葛即今梁國寧陵之葛鄉也。若湯居偃師，去寧陵八百餘里，豈當使民為之耕乎？亳今梁國谷熟縣是也。”諸說不同，未知孰是。 \par}

{\noindent\zhuan\zihao{6}\fzbyks 傳“告來”至“皆亡”。正義曰：經文既亡，其義難明，孔以意言耳。所言“帝告”，不知告誰，序言“從先王居”，或當告帝嚳也。 \par}

{\noindent\shu\zihao{5}\fzkt “自契”至“釐沃”。正義曰：自此已下皆\CJKunderwave{商書}也。序本別卷，與經不連，孔以經序宜相附近，引之各冠其篇首。此篇經亡序存，文無所託，不可以無經之序為卷之首，本書在此，故附此卷之末。契是商之始祖，故遠本之。自契至於成湯，凡八遷都。至湯始往居亳，從其先王帝嚳舊居。當時湯有言告，史序其事，作\CJKunderwave{帝告}、\CJKunderwave{釐沃}二篇。 \par}

湯徵諸侯，\footnote{為夏方伯，得專征伐。}葛伯不祀，湯始徵之，\footnote{葛，國。伯，爵也。廢其土地山川及宗廟神祗，皆不祀，湯始伐之。伐始於葛。祗，巨支反。}

{\noindent\shu\zihao{5}\fzkt 傳“葛國”至“於葛”。正義曰：序言“湯徵諸侯”，知其人是葛國之君，伯爵。直雲“不祀”，文無指斥。\CJKunderwave{王制}云：“山川神祗有不舉者為不敬，不敬者君削以地。宗廟有不順者為不孝，不孝者君黜以爵。”是言不祀必廢其土地山川之神祗及宗廟，皆不祀，故湯始徵之。湯伐諸侯，伐始於葛，\CJKunderwave{仲虺之誥}雲“初征自葛”是也。\CJKunderwave{孟子}云：“湯居亳，與葛為鄰。葛伯不祀，湯使人問之曰：‘何為不祀？’曰：‘無以供犧牲也。’湯使遺之牛羊。葛伯食之，又不祀。湯又使人問之曰：‘何為不祀？’曰：‘無以供粢盛也。’湯使亳往為之耕，老弱饋食。葛伯率其人，要其酒食黍稻者，劫而奪之，不授者殺之。有童子以黍肉餉，殺而奪之。\CJKunderwave{書}曰：‘葛伯仇餉。’此之謂也。”是說伐始於葛之事也。 \par}

作\CJKunderwave{湯徵}。\footnote{述始徵之義也，亡。}伊尹去亳適夏，\footnote{伊尹，字剩䂍湯進於桀。}

{\noindent\zhuan\zihao{6}\fzbyks 傳“伊尹”至“於桀”。正義曰：伊，氏；尹，字，故云“字氏”，倒文以曉人也。伊尹不得叛湯，知湯貢之於桀。必貢之者，湯欲以誠輔桀，冀其用賢以治；不可匡輔，乃始伐之，此時未有伐桀之意，故貢伊尹使輔之。\CJKunderwave{孫武兵書·反間篇}曰：“商之興也，伊尹在夏。周之興也，呂牙在殷。”言使之為反間也，與此說殊。 \par}

既醜有夏，復歸於亳。\footnote{醜惡其政。不能用賢，故退還。復，扶又反。}入自北門，乃遇汝鳩、汝方，\footnote{鳩、方，二人湯之賢臣。不期而會曰遇。}作\CJKunderwave{汝鳩}、\CJKunderwave{汝方}。\footnote{言所以醜夏而還之意，二篇皆亡。}

{\noindent\zhuan\zihao{6}\fzbyks 傳“鳩方”至“曰遇”。正義曰：伊尹與之言，知是賢臣也。“不期而會曰遇”，隱八年\CJKunderwave{穀梁傳}文也。 \par}

%%% Local Variables:
%%% mode: latex
%%% TeX-engine: xetex
%%% TeX-master: "../Main"
%%% End:

%% -*- coding: utf-8 -*-
%% Time-stamp: <Chen Wang: 2024-04-02 11:42:41>

% {\noindent \zhu \zihao{5} \fzbyks } -> 注 (△ ○)
% {\noindent \shu \zihao{5} \fzkt } -> 疏


\part{商書}


\chapter{卷八}


\section{湯誓第一}


 {\noindent\zhuan\zihao{6}\fzbyks \CJKunderwave{釋文}：凡三十四篇，十七篇存。\par}

\CJKunderline{伊尹}相湯，伐\CJKunderline{桀}，升自陑，\footnote{桀都安邑，湯升道從陑，出其不意。陑在河曲之南。相，息亮反。湯如字。馬云：“俗儒以湯為諡，或為號。號者似非其意，言諡近之。然不在\CJKunderwave{諡法}，故無聞焉。及\CJKunderline{禹}，俗儒以為名，\CJKunderwave{帝系}\CJKunderline{禹}名文命，\CJKunderwave{王侯世本}湯名天乙，推此言之，\CJKunderline{禹}豈復非諡乎？亦不在\CJKunderwave{諡法}，故疑焉。”桀，其列反，夏之末天子。升音升。陑音而。}遂與\CJKunderline{桀}戰于鳴條之野，\footnote{地在安邑之西，桀逆拒湯。}作\CJKunderwave{湯誓}。


{\noindent\zhuan\zihao{6}\fzbyks 傳“桀都”至“之南”。正義曰：此序湯自伐桀，必言“\CJKunderline{伊尹}相湯”者，序其篇次，自為首尾，以上雲\CJKunderline{伊尹}醜夏，遂相\CJKunderline{成湯}伐之，故文次言“\CJKunderline{伊尹}”也。計太公之相\CJKunderline{武王}，猶如\CJKunderline{伊尹}之相\CJKunderline{成湯}，\CJKunderwave{泰誓}不言“太公相”者，彼文無其次也。且\CJKunderline{武王}之時，有周、召之倫，聖賢多矣。湯稱\CJKunderline{伊尹}云：“聿求元聖，與之戮力。”\CJKunderline{伊尹}稱：“惟尹躬暨湯，咸有一德。”則\CJKunderline{伊尹}相湯，其功多于太公，故特言“\CJKunderline{伊尹}相湯”也。“桀都安邑”，相傳為然，即漢之河東郡安邑縣是也。\CJKunderwave{史記}吳起對魏武侯云：“\CJKunderline{夏桀}之居，左河濟，右太華，伊闕在其南，羊腸在其北，修政不仁，湯放之也。”\CJKunderwave{地理志}云：上黨郡壺關縣有羊腸阪，在安邑之北。是桀都安邑必當然矣。將明陑之所在，故先言“桀都安邑”。桀都在亳西，當從東而往，今乃升道從陑。“升”者，從下向上之名。言陑當是山阜之地，歷險迂路，為出不意故也。陑在河曲之南，蓋今潼關左右。河曲在安邑西南，從陑向北，渡河乃東向安邑。鳴條在安邑之西，桀西出拒湯，故“戰于鳴條之野”。陑在河曲之南，鳴條在安邑之西，皆彼有其跡，相傳云然。湯以至聖伐暴，當顯行用師，而出其不意，掩其不備者，湯承禪代之後，嘗為桀臣，慚而且懼，故出其不意。\CJKunderline{武王}則三分天下有其二，久不事紂，紂有浮桀之罪，地無險要之勢，故顯然致罰，以明天誅。又殷勤誓眾，與湯有異，所以湯惟一誓，\CJKunderline{武王}有三。 \par}

{\noindent\zhuan\zihao{6}\fzbyks 傳“地在”至“拒湯”。正義曰：\CJKunderline{鄭玄}云：“鳴條，南夷地名。\CJKunderwave{孟子}雲舜卒于鳴條，東夷之地，或雲陳留平邱縣今有鳴條亭是也。”皇甫謐云：“\CJKunderwave{伊訓}曰：‘造攻自鳴條，朕哉自亳。’又曰：‘夏師敗績，乃伐三朡。’\CJKunderwave{湯誥}曰：‘王歸自克夏，至于亳。’三朡在定陶，于義不得在陳留與東夷也。今安邑見有鳴條陌、昆吾亭，\CJKunderwave{左氏}以為昆吾與桀同以乙卯日亡，韋顧亦爾。故\CJKunderwave{詩}曰：‘韋顧既伐，昆吾\CJKunderline{夏桀}。’于\CJKunderwave{左氏}昆吾在衛，乃在濮陽，不得與桀異處同日而亡，明昆吾亦來安邑，欲以衛桀，故同日亡，而安邑有其亭也。且吳起言險以指安邑，安邑于此而言，何得在南夷乎？”謐言是也。 \par}

{\noindent\shu\zihao{5}\fzkt “\CJKunderline{伊尹}”至“湯誓”。正義曰：\CJKunderline{伊尹}以夏政醜惡，去而歸湯。輔相\CJKunderline{成湯}，與之伐桀，升道從陑，出其不意，遂與桀戰于鳴條之野。將戰而誓戒士眾，史敘其事，作\CJKunderwave{湯誓}。 \par}

湯誓\footnote{戒誓湯士眾。}

{\noindent\shu\zihao{5}\fzkt “湯誓”。正義曰：此經皆誓之辭也。\CJKunderwave{甘誓}、\CJKunderwave{泰誓}、\CJKunderwave{牧誓}發首皆有序引，別言其誓意，記其誓處。此與\CJKunderwave{費誓}惟記誓辭、不言誓處者，史非一人，辭有詳略，序以經文不具，故備言之也。 \par}

王曰：“格爾眾庶，悉聽朕言。\footnote{契始封商，湯遂以為天下號。湯稱王，則比桀于一夫。格，庚白反。}非台小子，敢行稱亂。有夏多罪，天命殛之。\footnote{稱，舉也。舉亂，以諸侯伐天子。非我小子敢行此事，桀有昏德，天命誅之，今順天。臺，以之反，下同。殛，居力反。}今爾有眾，汝曰：‘我后不恤我眾，舍我穡事，而割正夏。’\footnote{汝，汝有眾。我后，桀也。正，政也。言奪民農功而為割剝之政。恤，荀律反。舍音舍，廢也。}予惟聞汝眾言，\footnote{不憂我眾之言。}\CJKunderline{夏氏}有罪，予畏上帝，不敢不正。\footnote{不敢不正桀罪誅之。}今汝其曰：‘夏罪，其如台？’\footnote{今汝其復言桀惡，其亦如我所聞之言。復，扶又反。}\CJKunderline{夏王}率遏眾力，率割夏邑。\footnote{言桀君臣相率為勞役之事以絕眾力，謂廢農功。相率割剝夏之邑居，謂徵賦重。遏，于葛反，徐音謁，馬云：“止也。”}


{\noindent\zhuan\zihao{6}\fzbyks 傳“契始”至“一夫”。正義曰：以湯于此稱“王”，故本其號商之意，“契始封商，湯遂以商為天下之號”。\CJKunderline{鄭玄}之說亦然。惟\CJKunderline{王肅}云：“相士居商丘，湯取商為號。”若取商丘為號，何以不名“商丘”，而單名“商”也？若八遷，國名商不改，則此商猶是契商。非相士之商也。若八遷。遷即改名。則相士至湯改名多矣。相士既非始祖，又非受命，何故用其所居之地以為天下號名？\CJKunderline{成湯}之意，復何取乎？知其必不然也。湯取契封商，以商為天下之號，周不取后稷封邰為天下之號者，契後八遷，商名不改，\CJKunderline{成湯}以商受命，故宜以商為號；后稷之後，隨遷易名，\CJKunderline{公劉}為豳，\CJKunderline{大王}為周，\CJKunderline{文王}以周受命，故當以周為號。二代不同，理則然矣。\CJKunderwave{泰誓}云：“獨夫受。”此湯稱為王，則比桀于一夫；桀既同于一夫，故湯可稱王矣。是言湯于伐桀之時始稱王也。\CJKunderwave{周書·泰誓}稱王，則亦伐紂之時始稱王也。\CJKunderline{鄭玄}以\CJKunderline{文王}生稱王，亦謬也。 \par}

{\noindent\zhuan\zihao{6}\fzbyks 傳“稱舉”至“順天”。正義曰：“稱，舉”，\CJKunderwave{釋言}文。常法以臣伐君，則為亂逆，故舉亂謂“以諸侯伐天子”。“桀有昏德”，宣三年\CJKunderwave{左傳}文。以有昏德，天命誅之，今乃順天行誅，非復臣伐君也。以此解眾人守常之意也。 \par}

{\noindent\zhuan\zihao{6}\fzbyks 傳“今汝”至“之言”。正義曰：如我者，謂湯之自稱我也。湯謂其眾云：“汝言桀之罪，如我誓言所述也。” \par}

{\noindent\zhuan\zihao{6}\fzbyks 傳“言桀”至“賦重”。正義曰：此經與上“舍我穡事，而割正夏”其意一也。上言夏王之身，此言“君臣相率”，再言所以積桀之非也。力施于農，財供上賦，故以止絕眾力謂廢農功，割剝夏邑謂徵賦重。言以農時勞役，又重斂其財，致使民困而怨深，賦斂重則民不安矣。 \par}

有眾率怠弗協，曰：‘時日曷喪？予及汝皆亡！’\footnote{眾下相率為怠惰，不與上和合。比桀于日，曰：“是日何時喪？我與汝俱亡！”欲殺身以喪桀。喪，息浪反。惰，徒臥反。}夏德若茲，今朕必往。\footnote{凶德如此，我必往誅之。}爾尚輔予一人，致天之罰，予其大賚汝。\footnote{賚，與也。汝庶幾輔成我，我大與汝爵賞。罰音伐。賚，力代反，徐音來。}爾無不信，朕不食言。\footnote{食盡其言，偽不實。}爾不從誓言，\footnote{不用命。}予則孥戮汝，罔有攸赦。”\footnote{古之用刑，父子兄弟罪不相及，今雲孥戮汝，無有所赦，權以脅之，使勿犯。}

{\noindent\zhuan\zihao{6}\fzbyks 傳“眾下”至“喪桀”。正義曰：上既馭之非道，下亦不供其命，故“眾下相率為怠惰，不與上和合”，不肯每事順從也。比桀于日，曰：“是日何時喪亡？”欲令早喪桀命也。“我與汝俱亡”者，民相謂之辭，言並欲殺身以喪桀也。所以比桀于日者，以日無喪之理，猶雲桀不可喪，言喪之難也。不避其難，與汝俱亡，欲殺身以喪桀，疾之甚也。鄭云：“桀見民欲叛，乃自比于日，曰：‘是日何嘗喪乎？日若喪亡，我與汝亦皆喪亡。’引不亡之徵以脅恐下民也。” \par}

{\noindent\zhuan\zihao{6}\fzbyks 傳“食盡”至“不實”。正義曰：\CJKunderwave{釋詁}云：“食，偽也。”孫炎曰：“食言之偽也。”哀二十五年\CJKunderwave{左傳}云，孟武伯惡郭重，曰：“何肥也？”公曰：“是食言多矣，能無肥乎？”然則言而不行如食之消盡，後終不行前言為偽，故通謂偽言為“食言”，故\CJKunderwave{爾雅}訓“食”為偽也。 \par}

{\noindent\zhuan\zihao{6}\fzbyks 傳“古之”至“勿犯”。正義曰：昭二十年\CJKunderwave{左傳}引\CJKunderwave{康誥}曰：“父子兄弟，罪不相及。”是古之用刑如是也。既刑不相及，必不殺其子，權時以迫脅之，使勿犯刑法耳。不于\CJKunderwave{甘誓}解之者，以\CJKunderline{夏啟}承舜、\CJKunderline{禹}之後，刑罰尚寬，殷、周以後，其罪或相緣坐，恐其實有孥戮，故于此解之。\CJKunderline{鄭玄}云：“大罪不止其身，又孥戮其子孫。\CJKunderwave{周禮}云：‘其奴男子入于罪隸，女子入于舂稿。’”鄭意以為實戮其子，故\CJKunderwave{周禮}注云：“奴謂從坐而沒入縣官者也。”孔以孥戮為權脅之辭，則\CJKunderwave{周禮}所云非從坐也。\CJKunderline{鄭玄}云：“謂坐為盜賊而為奴者，輸于罪隸。”舂人、稿人之官引此“孥戮汝”。又引\CJKunderwave{論語}云：“\CJKunderline{箕子}為之奴。”或如眾言，別有沒入，非緣坐者也。 \par}

{\noindent\shu\zihao{5}\fzkt “王曰”至“攸赦”。正義曰：商王\CJKunderline{成湯}將與桀戰，呼其將士曰：“來，汝在軍之眾庶，悉聽我之誓言。我伐夏者，非我小子輒敢行此以臣伐君，舉為亂事，乃由有夏君桀多有大罪，上天命我誅之。桀既失君道，我非復桀臣，是以順天誅之，由其多罪故也。桀之罪狀，汝盡知之。今汝桀之所有之眾，即汝輩是也。汝等言曰：‘我君\CJKunderline{夏桀}，不憂念我等眾人，舍廢我稼穡之事，奪我農功之業，而為割剝之政于夏邑，斂我貨財。’我惟聞汝眾言，夏氏既有此罪，上天命我誅，桀我畏上天之命，不敢不正桀罪而誅之。又質而審之，今汝眾人其必言曰：‘夏王之罪其實如我所言。’夏王非徒如此，又與臣下相率遏絕眾力，使不得事農。又相率為割剝之政于此夏邑，使不得安居。上下同惡，民困益甚，由是汝等相率怠，惰不與在上和協。比桀于日，曰：‘是日何時能喪？君其可喪，我與汝皆亡身殺之。’寧殺身以亡桀，是其惡之甚。夏王惡德如此，今我必往誅之。汝庶幾輔成我一人，致行天之威罰，我其大賞賜汝。汝無得不信我語，我終不食盡其言，為虛偽不實。汝若不從我之誓言，我則並殺汝子，以戮汝身，必無有所赦。”勸使勉力，勿犯法也。“庶”亦“眾”也，古人有此重言，猶云“艱難”也。 \par}

湯既勝夏，欲遷其社，不可。\footnote{湯承堯、舜禪代之後，順天應人，逆取順守而有慚德，故革命創制，改正易服，變置社稷，而後世無及句龍者，故不可而止。社，后土之神。禪，時戰反。應，應對之應。創，初亮反。正音徵，又音政。句音鉤。句龍，\CJKunderline{共工}之子，為后土。}作\CJKunderwave{夏社}、\CJKunderwave{疑至}、\CJKunderwave{臣扈}。\footnote{言夏社不可遷之義，\CJKunderwave{疑至}及\CJKunderwave{臣扈}，三篇皆亡。扈音戶。}

{\noindent\zhuan\zihao{6}\fzbyks 傳“湯承”至“而止”。正義曰：傳解湯遷社之意，湯承堯、舜禪代之後，已獨伐而取之，雖復應天順人，乃是逆取順守，而有慚愧之德，自恨不及古人，故革命創制改正易服，因變置社稷也。\CJKunderwave{易·革卦·彖}曰：“湯武革命，順乎天而應乎人。”下篇言湯有慚德。\CJKunderwave{大傳}云：“改正朔，易服色，此其所得與民變革者也。”所以變革此事，欲易人之視聽，與之更新，故于是之時變置社稷。昭二十九年\CJKunderwave{左傳}云：“\CJKunderline{共工氏}有子曰句龍，為后土，后土為社。有烈山氏之子曰柱，為稷，自夏已上祀之。周棄亦為稷，自商已來祀之。”\CJKunderwave{祭法}云：“厲山氏之有天下也，其子曰農，能殖百穀。夏之衰也，周棄繼之，故祀以為稷。\CJKunderline{共工氏}之霸九州也，其子曰后土，能平九州，故祀以為社。”是言變置之事也。\CJKunderwave{魯語}文與\CJKunderwave{祭法}正同，而云“夏之興也，周棄繼之”，“興”當為“衰”字之誤耳。湯于初時，社稷俱欲改之。周棄功多于柱，即令廢柱祀棄。而上世治水土之臣，其功無及句龍者，故不可遷而止。此序之次在\CJKunderwave{湯誓}之下，云“湯既勝夏”，下云“夏師敗績，湯遂從之”，是未及逐桀，已為此謀。\CJKunderline{鄭玄}等注此序，乃在\CJKunderwave{湯誓}之上，若在作誓之前，不得云“既勝夏”也。\CJKunderwave{孟子}曰：“犧牲既成，粢盛既絜，祭祀以時，然而旱乾水益，則變置社稷。”\CJKunderline{鄭玄}因此乃云：“湯伐桀之時，大旱，既置其禮祀，明德以薦而猶旱，至七年，故更致社稷。”乃謂湯即位之後，七年大旱，方始變之。若實七年乃變，何當系之“勝夏”？“勝夏”猶尚不可，況在\CJKunderwave{湯誓}前乎？且\CJKunderwave{禮記}云：“夏之衰也，周棄繼之。”商興七年乃變，安得以“夏衰”為言也？若商革夏命，猶七年祀柱，\CJKunderwave{左傳}亦不得斷為自夏已上祀柱，自商以來祀棄也。由此而言，孔稱改正朔而變置社稷，所言得其旨也。漢世儒者說社稷有二，\CJKunderwave{左傳}說社祭句龍，稷祭柱、棄，惟祭人神而已。\CJKunderwave{孝經}說社為土神，稷為穀神，句龍、柱、棄是配食者也。孔無明說，而此經云“遷社”，孔傳云“無及句龍”，即同賈逵、\CJKunderline{馬融}等說，以社為句龍也。 \par}

{\noindent\zhuan\zihao{6}\fzbyks 傳“言夏”至“皆亡”。正義曰：“疑至”與“\CJKunderline{臣扈}”相類，當是二臣名也。蓋亦言其不可遷之意。\CJKunderline{馬融}云：“聖人不可自專，複用二臣自明也。” \par}

{\noindent\shu\zihao{5}\fzkt “湯既”至“\CJKunderline{臣扈}”。正義曰：湯既伐而勝夏，革命創制，變置社稷，欲遷其社，無人可代句龍，故不可而止。于時有言議論其事，故史敘之，為\CJKunderwave{夏社}、\CJKunderwave{疑至}、\CJKunderwave{臣扈}三篇，皆亡。 \par}

夏師敗績，湯遂從之，\footnote{大崩曰敗績。從謂遂討之。績，子寂反。從，才容反。}遂伐三\xpinyin*{朡}，俘厥寶玉。\footnote{三朡，國名，桀走保之，今定陶也。桀自安邑東入山，出太行，東南涉河。湯緩追之，不迫，遂奔南巢。俘，取也。玉以禮神，使無水旱之災，故取而寶之。朡，子公反。俘音孚。行，戶剛反，一音如字。}\CJKunderline{誼伯}、\CJKunderline{仲伯}作\CJKunderwave{典寶}。\footnote{二臣作\CJKunderwave{典寶}一篇，言國之常寶也，亡。誼，本或作義。}

{\noindent\shu\zihao{5}\fzkt 傳“三朡”至“寶之”。正義曰：湯伐三朡，知是國名。逐桀而伐其國，知“桀走保之”也。“今定陶”者，相傳為然。安邑在洛陽西北，定陶在洛陽東南，孔跡其所往之路，桀自安邑東入山，出太行，乃東南涉河，往奔三朡；湯緩追之，不迫，遂奔南巢。“俘、取”，\CJKunderwave{釋詁}文。桀必載寶而行，棄于三朡。取其寶玉，取其所棄者也。楚語云：“玉足以庇廕嘉，谷使無水旱之災，則寶之。”韋昭云：“玉，禮神之玉也。”言用玉禮神，神享其德，使風雨調和，可以庇廕嘉穀，故取而寶之。 \par}

\section{仲虺之誥第二【偽】}

湯歸自夏，至于大坰，\footnote{自三朡而還。大坰，地名。夏，亥雅反。坰，故螢反，徐，欽螢反，又古螢反。}\CJKunderline{仲虺}作誥。\footnote{為湯左相，奚仲之後。虺，許鬼反。誥，故報反。相，息亮反。奚，弦雞反。}


{\noindent\zhuan\zihao{6}\fzbyks 傳“為湯”至“之後”。正義曰：定元年\CJKunderwave{左傳}云：“薛之皇祖奚仲居薛，以為夏車正。\CJKunderline{仲虺}居薛，以為湯左相。”是其事也。 \par}

{\noindent\shu\zihao{5}\fzkt “湯歸”至“作誥”。正義曰：湯歸自伐夏，至于大坰之地，其臣\CJKunderline{仲虺}作誥以誥湯，使錄其言，作\CJKunderwave{仲虺之誥}。上言“遂伐三朡”，故傳言“自三朡而還”。不言“歸自三朡”，而言“歸自夏”者，伐夏而遂逐桀，于今方始旋歸，以自夏告廟，故序言“自夏”。傳本其來處，故云“自三朡”耳。“大坰，地名”，未知所在，當是定陶向亳之路所經。湯在道而言“予恐來世以臺為口實”，故\CJKunderline{仲虺}至此地而作誥也。序不言“作\CJKunderline{仲虺}之誥”，以理足文便，故略之。 \par}

仲虺之誥\footnote{仲虺，臣名，以諸侯相天子。會同曰誥。}


{\noindent\zhuan\zihao{6}\fzbyks 傳“\CJKunderline{仲虺}”至“曰誥”。正義曰：伯仲叔季，人字之常，“\CJKunderline{仲虺}”必是其名，或字仲而名虺。古人名或不可審知，縱使是字，亦得謂之為名，言是人之名號也。\CJKunderwave{左傳}稱居薛，為湯左相，是“以諸侯相天子”也。\CJKunderwave{周禮·士師}云：“以五戒先後刑罰，一曰誓，用之于軍旅。二曰誥，用之于會同。”是“會同曰誥”。“誥”謂于會之所，設言以誥眾，此惟誥湯一人而言“會同”者，因解諸篇“誥”義，且\CJKunderline{仲虺}必對眾誥湯，亦是“會同曰誥”。 \par}

{\noindent\shu\zihao{5}\fzkt “\CJKunderline{仲虺}之誥”。正義曰：發首二句，史述\CJKunderline{成湯}之心。次二句，湯言己慚之意，\CJKunderline{仲虺}乃作誥。以下皆勸湯之辭。自“曰嗚呼”至“用爽厥師”，言天以桀有罪，命伐夏之事。自“簡賢輔勢”至“言足聽聞”，說湯在桀時怖懼之事。自“惟王弗邇聲色”至“厥惟舊哉”，言湯有德行加民，民歸之事。自“佑賢輔德”以下說天子之法，當擢用賢良，屏黜昏暴，勸湯奉行此事，不須以放桀為惡。\CJKunderwave{康誥}、\CJKunderwave{召誥}之類，二字足以為文，“\CJKunderline{仲虺}誥”三字不得成文，以“之”字足成其句。\CJKunderwave{畢命}、\CJKunderwave{冏命}不言“之”，\CJKunderwave{微子之命}、\CJKunderwave{文侯之命}言“之”，與此同，猶\CJKunderwave{周禮·司服}言“大裘而冕”，亦足句也。 \par}

\CJKunderline{成湯}放\CJKunderline{桀}于南巢，惟有慚德，\footnote{湯伐桀，武功成，故以為號。南巢，地名。有慚德，慚德不及古。湯伐桀，武功成，故號\CJKunderline{成湯}。一云：“成，諡也。”}曰：“予恐來世以台為口實。”\footnote{恐來世論道我放天子，常不去口。}\CJKunderline{仲虺}乃作誥，\footnote{陳義誥湯，可無慚。}曰：“嗚呼！惟天生民有欲，無主乃亂，\footnote{民無君主則恣情慾，必致禍亂。}惟天生聰明時乂。\footnote{言天生聰明，是治民亂。}有夏昏德，民墜塗炭。\footnote{\CJKunderline{夏桀}昏亂，不恤下民，民之危險，若陷泥墜火，無救之者。}天乃錫王勇智，表正萬邦，纘\CJKunderline{禹}舊服，\footnote{言天與王勇智，應為民主，儀表天下，法正萬國，繼\CJKunderline{禹}之功，統其故服。纘，子管反。應，應對之應。}茲率厥典，奉若天命。\footnote{天意如此，但當循其典法，奉順天命而已，無所慚。}

{\noindent\shu\zihao{5}\fzkt “\CJKunderline{成湯}放桀于南巢”。正義曰：桀奔南巢，湯縱而不迫，故稱“放”也。傳言“南巢，地名”，不知地之所在。\CJKunderwave{周書}序有“\CJKunderline{巢伯}來朝”，傳云：“南方遠國。”\CJKunderline{鄭玄}云：“巢，南方之國。世一見者，桀之所奔，蓋彼國也。以其國在南，故稱南耳。”傳並以“南巢”為地名，不能委知其處，故未明言之。 \par}

“夏王有罪，矯誣上天，以布命于下。\footnote{言託天以行虐于民，乃桀之大罪。矯，居表反。誣音無。}帝用不臧，式商受命，用爽厥師。\footnote{天用桀無道，故不善之。式，用。爽，明也。用商受王命，用明其眾，言為主也。臧，作郎反。}簡賢附勢，實繁有徒。\footnote{簡，略也。賢而無勢則略之，不賢有勢則附之。若是者繁多有徒眾，無道之世所常。繁音煩。}肇我邦于有夏，若苗之有莠，若粟之有秕。\footnote{始我商家，國于夏世，欲見翦除，若莠生苗，若秕在粟，恐被鋤治簸颺。莠，羊九反。秕，悲裡反，徐，甫裡反，又必履反。鋤，仕魚反。簸，彼我反。颺音揚。}小大戰戰，罔不懼于非辜。矧予之德，言足聽聞。\footnote{言商家小大憂危，恐其非罪見滅。矧，況也。況我之道德善言足聽聞乎！無道之惡有道，自然理。惡，烏路反。}惟王不邇聲色，不殖貨利。\footnote{邇，近也。不近聲樂，言清簡。不近女色，言貞固。殖，生也。不生資貨財利，言不貪也。既有聖德，兼有此行。近，附近之近。行，下孟反。}


{\noindent\zhuan\zihao{6}\fzbyks 傳“式，用。爽，明也”。正義曰：“式，用”，\CJKunderwave{釋言}文。昭七年\CJKunderwave{左傳}云：“是以有精爽至于神明。”從爽以至于明，則“爽”是明之始，故“爽”為明也。經稱“昧爽”，謂未大明也。 \par}

{\noindent\shu\zihao{5}\fzkt “夏王”至“厥師”。正義曰：“矯”，詐也。“誣”，加也。夏王自有所欲，詐加上天，言天道須然，不可不爾，假此以布苛虐之命于天下，以困苦下民。上天用桀無道之故，故不善之，用使商家受此為王之命，以王天下。用命商王，明其所有之眾，謂湯教之使修德行善以自安樂，是明之也。 \par}

德懋懋官，功懋懋賞。用人惟己，改過不吝。\footnote{勉于德者，則勉之以官。勉于功者，則勉之以賞。用人之言，若自己出；有過則改，無所吝惜，所以能\CJKunderline{成王}業。}克寬克仁，彰信兆民。\footnote{言湯寬仁之德明信于天下。}

{\noindent\shu\zihao{5}\fzkt “德懋”至“不吝”。正義曰：于德能勉力行之者，王則勸勉之以官。于功能勉力為之者，王則勸勉之以賞。用人之言，惟如己之所出；改悔過失，無所吝惜。美湯之行如此。凡庸之主，得人之言，恥非己智，雖知其善，不肯遂從。己有愆失，恥于改過，舉事雖覺其非，不肯更悔，是惜過不改。故以此美湯也。\CJKunderline{成湯}之為此行，尚為\CJKunderline{仲虺}所稱歎，凡人能勉者鮮矣。 \par}

乃\CJKunderline{葛伯}仇餉，初征自葛，東征西夷怨，南征北狄怨，\footnote{\CJKunderline{葛伯}遊行，見農民之餉于田者，殺其人，奪其餉，故謂之仇餉。仇，怨也。湯為是以不祀之罪伐之，從此後遂徵無道。西夷、北狄，舉遠以言，則近者著矣。仇音求。餉，式亮反。}曰：‘奚獨後予？’\footnote{怨者辭也。}攸徂之民，室家相慶，曰：‘徯予後，後來其蘇。’\footnote{湯所往之民，皆喜曰：“待我君來，其可蘇息。”徯，胡啟反。蘇，字亦作穌。}民之戴商，厥惟舊哉！\footnote{舊，謂初征自葛時。}佑賢輔德，顯忠遂良。\footnote{賢則助之，德則輔之，忠則顯之，良則進之。明王之道。}兼弱攻昧，取亂侮亡。\footnote{弱則兼之，暗則攻之，亂則取之，有亡形則侮之。言正義。}推亡固存，邦乃其昌。\footnote{有亡道，則推而亡之；有存道，則輔而固之。王者如此，國乃昌盛。推，土雷反。}

{\noindent\zhuan\zihao{6}\fzbyks 傳“賢則”至“之道”。正義曰：\CJKunderwave{周禮·鄉大夫}云：“三年則大比，考其德行道藝，而興賢者。”\CJKunderline{鄭玄}云：“賢者謂有德行者。”\CJKunderwave{詩序}云：“忠臣良士皆是善也。”然則“賢”是德盛之名，“德”是資賢之實，“忠”是盡心之事，“良”是為善之稱，俱是可用之人，所從言之異耳。“佑之”與“輔、顯之”與“遂”，隨便而言之。 \par}

{\noindent\zhuan\zihao{6}\fzbyks 傳“弱則”至“正義”。正義曰：力少為“弱”，不明為“昧”，政荒為“亂”，國滅為“亡”，“兼”謂包之，“攻”謂擊之，“取”謂取為己有，“侮”謂侮慢其人。“弱”、“昧”、“亂”、“亡”，俱是彼國衰微之狀。“兼”、“攻”、“取”、“侮”，是此欲吞併之意。“弱”、“昧”是始衰之事，來服則製為己屬，不服則以兵攻之。此二者始欲服其人，末是滅其國。“亂”是已亂，“亡”謂將亡，二者衰甚，已將滅其國。亡形已著，無可忌憚，故陵侮其人。既侮其人，必滅其國，故以“侮”言之。此是人君之正義。\CJKunderline{仲虺}陳此者，意亦言桀亂亡，取之不足為愧。下言“推亡”及“覆昏暴”，其意亦在桀也。 \par}

{\noindent\shu\zihao{5}\fzkt “乃\CJKunderline{葛伯}仇餉”。正義曰：此言“乃”者，卻說已過之事。\CJKunderwave{胤徵}云“乃季秋月朔”，其義亦然。\CJKunderwave{左傳}稱“怨耦曰仇”，謂彼人有負于我，我心怨之，是名為“仇”也。餉田之人不負\CJKunderline{葛伯}，\CJKunderline{葛伯}奪其餉而殺之，是\CJKunderline{葛伯}以餉田之人為己之仇。言非所怨而妄殺，故湯為之報也。\CJKunderwave{孟子}稱湯使亳眾往為之耕，有童子以黍肉餉，\CJKunderline{葛伯}奪而殺之。則\CJKunderline{葛伯}所殺，殺亳人也。傳言“\CJKunderline{葛伯}遊行，見農人之餉于田者，殺其人而奪其餉，故謂之仇餉”，乃似\CJKunderline{葛伯}自殺己人，與\CJKunderwave{孟子}違者，湯之徵葛，以人之枉死而為之報耳，不為亳人乃報之，非亳人則赦之，故傳指言殺餉，不辨死者何人。亳人、葛人，義無以異，故不復言“亳”，非是故違\CJKunderwave{孟子}。 \par}

德日新，萬邦惟懷。志自滿，九族乃離。\footnote{日新，不懈怠。自滿，志盈溢。懈，工債反。}

{\noindent\shu\zihao{5}\fzkt “德日”至“乃離”。正義曰：\CJKunderwave{易·繫辭}云：“日新之謂盛德。”修德不怠，日日益新，德加于人，無遠不屆，故萬邦之眾惟盡歸之。志意自滿則陵人，人既被陵，情必不附，雖九族之親，乃亦離之。“萬邦”，舉遠以明近；“九族”，舉親以明疏也。漢代儒者說九族有二，案\CJKunderwave{禮戴}及\CJKunderwave{尚書緯}、歐陽說九族，乃異姓有屬者，父族四，母族三，妻族二。\CJKunderwave{古尚書}說九族，從高祖至玄孫凡九族。\CJKunderwave{堯典}云“以親九族”，傳云“以睦高祖玄孫之親”，則此言“九族”，亦謂高祖玄孫之親也。謂“萬邦惟懷”，實歸之。“九族乃離”，實離之。聖賢設言為戒，容辭頗甚，父子之間，便以志滿相棄。此言“九族”，以為外姓九族有屬，文便也。 \par}

王懋昭大德，建中于民，以義制事，以禮制心，垂裕後昆。\footnote{欲王自勉，明大德，立大中之道于民，率義奉禮，垂優足之道示後世。中如字；本或作忠，非。裕，徐以樹反。}予聞曰：‘能自得師者王，\footnote{求賢聖而事之。王，徐于況反，又如字。}謂人莫已若者亡。\footnote{自多足，人莫之益，亡之道。}好問則裕，自用則小。’\footnote{問則有得，所以足，不問專固，所以小。好，呼報反。}嗚呼！慎厥終，惟其始。\footnote{靡不有初，鮮克有終，故戒慎終如其始。鮮，息淺反。}殖有禮，覆昏暴。\footnote{有禮者封殖之，昏暴者覆亡之。覆，芳服反。暴，蒲報反，字或作虣。}欽崇天道，永保天命。”\footnote{王者如此上事，則敬天安命之道。}

\section{湯誥第三【偽】}

\CJKunderline{湯}既黜夏命，\footnote{黜，退也，退其王命。}復歸于亳，作\CJKunderwave{湯誥}。

湯誥\footnote{以伐桀大義告天下。}

{\noindent\shu\zihao{5}\fzkt “湯既”至“湯誥”。正義曰：湯既黜夏王之命，復歸于亳，以伐桀大義誥示天下。史錄其事，作\CJKunderwave{湯誥}。\CJKunderline{仲虺}在路作誥，此至亳乃作，故次\CJKunderwave{仲虺}之下。 \par}

王歸自克夏，至于亳，誕告萬方。\footnote{誕，大也。以天命大義告萬方之眾人。誕音但。告，工毒反。}

{\noindent\shu\zihao{5}\fzkt “王歸自克夏”。正義曰：湯之伐桀，當有諸侯從之，不從行者必應多矣。既已克夏，改正名號，還至于亳，海內盡來，猶如\CJKunderwave{武成}篇所云“庶邦冢君暨百工，受命于周”也。湯于此時大誥諸侯以伐桀之義，故云“誕告萬方”。“誕，大”，\CJKunderwave{釋詁}文。“萬”者舉盈數。下云“凡我造邦”，是誥諸侯也。 \par}

王曰：“嗟！爾萬方有眾，明聽予一人誥。\footnote{天子自稱曰予一人，古今同義。}惟皇上帝，降衷于下民。\footnote{皇，大。上帝，天也。衷，善也。}

{\noindent\shu\zihao{5}\fzkt “降衷于下民”。正義曰：天生烝民，與之五常之性，使有仁義禮智信，是天降善于下民也。天既與善于民，君當順之，故下傳云，順人有常之性，則是為君之道。 \par}

若有恆性，克綏厥猷惟後。\footnote{順人有常之性，能安立其道教，則惟為君之道。}夏王滅德作威，以敷虐于爾萬方百姓。\footnote{\CJKunderline{夏桀}滅道德，作威刑以布行虐政于天下百官。言殘酷。}爾萬方百姓，罹其兇害，弗忍荼毒，\footnote{罹，被。荼毒，苦也。不能堪忍，虐之甚。罹，力之反；本亦作羅，洛河反。荼音徒。}

{\noindent\shu\zihao{5}\fzkt “弗忍荼毒”。正義曰：\CJKunderwave{釋草}云：“荼，苦菜。”此菜味苦，故假之以言人苦。“毒”謂螫人之蟲，蛇虺之類。實是人之所苦，故並言“荼毒”以喻苦也。 \par}

並告無辜于上下神祇。\footnote{言百姓兆民並告無罪，稱冤訴天地。冤，紆元反。}天道福善禍淫，降災于夏，以彰厥罪。\footnote{政善天福之，淫過天禍之，故下災異以明桀罪惡，譴寤之而桀不改。譴，遣戰反。寤，五故反。}

“肆台小子，將天命明威，不敢赦。\footnote{行天威，謂誅之。臺音怡。}敢用玄牡，敢昭告于上天神後，請罪有夏。\footnote{明告天，問桀百姓有何罪而加虐乎？牡，茂後反。}

{\noindent\shu\zihao{5}\fzkt “敢用玄牡”。正義曰：\CJKunderwave{檀弓}云：“殷人尚白，牲用白。”今云“玄牡”，夏家尚黑，于時未變夏禮，故不用白也。故安國注\CJKunderwave{論語}“敢用玄牡”之文云，“殷家尚白，未變夏禮，故用玄牡”，是其義也。\CJKunderline{鄭玄}說：“天神有六，周家冬至祭皇天大帝于圜丘，牲用蒼。夏至祭靈威仰于南郊，則牲用騂。”孔注\CJKunderwave{孝經}，圜丘與郊共為一事，則孔之所說無六天之事，\CJKunderwave{論語·堯曰}之篇所言“敢用玄牡”即此事是也。孔注\CJKunderwave{論語}以為“堯曰”之章“有二帝三王之事，錄者採合以成章。檢\CJKunderwave{大禹謨}及此篇與\CJKunderwave{泰誓}、\CJKunderwave{武成}，則‘堯曰’之章其文略矣。”。\CJKunderline{鄭玄}解\CJKunderwave{論語}云：“‘用玄牡’者，為舜命\CJKunderline{禹}事，于時裛告五方之帝，莫適用，用皇天大帝之牲。”其意與孔異。 \par}

聿求元聖，與之戮力，以與爾有眾請命。\footnote{聿，遂也。大聖陳力，謂\CJKunderline{伊尹}。放桀除民之穢，是請命。聿，允橘反，述也。戮舊音六，又力雕反，\CJKunderwave{說文}力周反，\CJKunderwave{史記}音力消反。穢，于廢反。}

{\noindent\shu\zihao{5}\fzkt 傳“聿遂也”至“請命”。正義曰：“聿”訓述也，述前所以申遂，故“聿”為遂也。“戮力”猶勉力也。\CJKunderwave{論語}云：“陳力就列。”湯臣大賢惟有\CJKunderline{伊尹}，故知“大聖陳力，謂\CJKunderline{伊尹}”也。\CJKunderline{伊尹}賢人而謂之“聖”者，相對則“聖”極而“賢”，次散文則“賢”、“聖”相通。舜謂\CJKunderline{禹}曰：“惟汝賢。”是“聖”得謂之“賢”，則“賢”亦可言“聖”。\CJKunderline{鄭玄}\CJKunderwave{周禮注}云：“聖通而先識也。”解先識則為聖名，故\CJKunderline{伊尹}可為聖也。\CJKunderwave{孟子}云：“\CJKunderline{伯夷}聖人之清者也，\CJKunderline{伊尹}聖人之任者也，柳下惠聖人之和者也，\CJKunderline{孔子}聖人之時者也。”是謂\CJKunderline{伊尹}為聖人者也。桀為殘虐，人不自保，故伐桀除人之穢，是為請命。 \par}

上天孚佑下民，罪人黜伏。\footnote{孚，信也。天信佑助下民，桀知其罪，退伏遠屏。}天命弗僣，賁若草木，兆民允殖。\footnote{僣，差。賁，飾也。言福善禍淫之道不差，天下惡除，煥然咸飾，若草木同華，民信樂生。僣，子念反，忒也；劉創林反。賁，彼義反，徐扶雲反。煥，呼亂反。樂音洛。}

{\noindent\shu\zihao{5}\fzkt “天命”至“允殖”。正義曰：桀以大罪，身即黜伏，是天之福善禍淫之命信而不僣差也。既除大惡，天下煥然修飾，若草木同生華，兆民信樂生也。昔日不保性命，今日樂生活矣。僣差，不齊之意，故傳以“僣”為差。“賁，飾”，\CJKunderwave{易·序卦}文也。 \par}

俾予一人，輯寧爾邦家。\footnote{言天使我輯安汝國家。國，諸侯。家，卿大夫。}茲朕未知獲戾于上下，\footnote{此伐桀未知得罪于天地。謙以求眾心。戾，力計反。}

{\noindent\shu\zihao{5}\fzkt 傳“此伐”至“眾心”。正義曰：經言“茲”者，謂此伐桀也。顧氏云：“‘未知得罪于天地’，言伐桀之事，未知得罪于天地以否。”湯之伐桀，上應天心，下符人事，本實無罪，而云未知得罪以否者，謙以求眾心。 \par}

慄慄危懼，若將隕于深淵。\footnote{慄慄危心，若墜深淵。危懼之甚。慄音慄。隕，于敏反。}凡我造邦，無從匪彝，無即慆淫，\footnote{戒諸侯與之更始。彝，常。慆，慢也。無從非常，無就慢過，禁之。彝，徐音夷。慆，他刀反。}各守爾典，以承天休。\footnote{守其常法，承大美道。}爾有善，朕弗敢蔽。罪當朕躬，弗敢自赦，惟簡在上帝之心。\footnote{所以不蔽善人，不赦己罪，以其簡在天心故也。}

{\noindent\shu\zihao{5}\fzkt “惟簡在上帝之心”。正義曰：\CJKunderline{鄭玄}注\CJKunderwave{論語}云：“簡閱在天心，言天簡閱其善惡也。” \par}

其爾萬方有罪，在予一人。\footnote{自責化不至。}予一人有罪，無以爾萬方。\footnote{無用爾萬方，言非所及。}嗚呼！尚克時忱，乃亦有終。”\footnote{忱，誠也。庶幾能是誠道，乃亦有終世之美。忱，市林反。}

\CJKunderline{咎單}作\CJKunderwave{明居}。\footnote{\CJKunderline{咎單}，臣名，主土地之官。作明居民法一篇，亡。單音善，卷末同。}

{\noindent\shu\zihao{5}\fzkt “\CJKunderline{咎單}作明居”。正義曰：百篇之序此類有四，\CJKunderline{伊尹}作\CJKunderwave{咸有一德}、\CJKunderline{周公}作\CJKunderwave{無逸}、作\CJKunderwave{立政}，與此篇。直言其所作之人，不言其作者之意，蓋以經文分明，故略之。\CJKunderline{馬融}云：“\CJKunderline{咎單}為湯司空。”傳言“主土地之官”，蓋亦為司空也。 \par}

\section{伊訓第四【偽】}


\CJKunderline{成湯}既沒，\CJKunderline{太甲}元年，\footnote{太甲，太丁子，湯孫也。太丁未立而卒，及湯沒而\CJKunderline{太甲}立，稱元年。}\CJKunderline{伊尹}作\CJKunderwave{伊訓}、\CJKunderwave{肆命}、\CJKunderwave{徂後}。\footnote{凡三篇，其二亡。}


{\noindent\zhuan\zihao{6}\fzbyks 傳“\CJKunderline{太甲}”至“元年”。正義曰：“\CJKunderline{太甲}，太丁子”，\CJKunderwave{世本}文也。此序以“\CJKunderline{太甲}元年”繼“湯沒”之下，明是太丁未立而卒，\CJKunderline{太甲}以孫繼祖，故湯沒而\CJKunderline{太甲}代立，即以其年稱為元年也。周法以逾年即位，知此即以其年稱元年者，此經云“元祀十有二月，\CJKunderline{伊尹}祠于先王。奉嗣王祗見厥祖”，\CJKunderwave{太甲}中篇云“惟三祀十有二月朔，\CJKunderline{伊尹}以冕服奉嗣王歸于亳”，二者皆云“十有二月”，若是逾年即位，二者皆當以正月行事，何以用十二月也？明此經“十二月”是湯崩之逾月，\CJKunderwave{太甲}中篇“三祀十有二月”是服闋之逾月，以此知湯崩之年，\CJKunderline{太甲}即稱元年也。舜\CJKunderline{禹}以受帝終事，自取歲首，遭喪嗣位，經無其文，夏後之世或亦不逾年也。顧氏云：“殷家猶質，逾月即改元年，以明世異，不待正月以為首也。”商謂年為祀，序稱“年”者，序以周世言之故也。據此經序及\CJKunderwave{太甲}之篇，\CJKunderline{太甲}必繼湯後，而\CJKunderwave{殷本紀}云：“湯崩，太子太丁未立而卒，于是乃立太丁之弟外丙。三年崩，別立外丙之弟仲壬。四年崩，\CJKunderline{伊尹}乃立太丁之子\CJKunderline{太甲}。”與經不同，彼必妄也。劉歆、班固不見古文，謬從\CJKunderwave{史記}。皇甫謐既得此經，作\CJKunderwave{帝王世紀}，乃述馬遷之語，是其疏也。顧氏亦云：“止可依經誥大典，不可用傳記小說。” \par}

{\noindent\shu\zihao{5}\fzkt “\CJKunderline{成湯}”至“徂後”。正義曰：\CJKunderline{成湯}既沒，其歲即\CJKunderline{太甲}元年。\CJKunderline{伊尹}以\CJKunderline{太甲}承湯之後，恐其不能纂修祖業，作書以戒之。史敘其事，作\CJKunderwave{伊訓}、\CJKunderwave{肆命}、\CJKunderwave{徂後}三篇。 \par}

伊訓\footnote{作訓以教道\CJKunderline{太甲}。}

惟元祀十有二月乙丑，\CJKunderline{伊尹}祠于先王。\footnote{此湯崩逾月，\CJKunderline{太甲}即位，奠殯而告。祀，年也。夏曰歲，商曰祀，周曰年，\CJKunderline{唐}、\CJKunderline{虞}曰載。尹祠音辭，祭也。}


{\noindent\zhuan\zihao{6}\fzbyks 傳“此湯”至“而告”。正義曰：\CJKunderwave{太甲}中篇云：“三祀十有二月，\CJKunderline{伊尹}以冕服奉嗣王。”則是除喪即吉，明十二月服終。\CJKunderwave{禮記}稱：“三年之喪，二十五月而畢。”知此年十一月湯崩，此祠先王是“湯崩逾月，\CJKunderline{太甲}即位，奠殯而告”也。此“奠殯而告”，亦如周\CJKunderline{康王}受顧命屍于天子。春秋之世既有奠殯即位、逾年即位，此逾月即位當奠殯即位也。此言“\CJKunderline{伊尹}祠于先王”，是特設祀也，“嗣王祗見厥祖”是始見祖也。特設祀禮而王始見祖，明是初即王位，告殯為喪主也。 \par}

{\noindent\shu\zihao{5}\fzkt “惟元祀”。正義曰：“\CJKunderline{伊尹}祠于先王”，謂祭湯也。“奉嗣王祗見厥祖”，謂見湯也。故傳解“祠先王”為“奠殯而告”，“見厥祖”為“居位主喪”，“群后咸在”為“在位次”，皆述在喪之事，是言“祠”是奠也。祠喪于殯，斂、祭皆名為奠，虞祔卒哭始名為祭。知“祠”非宗廟者，“元祀”即是初喪之時，未得祠廟，且湯之父祖不追為王，所言“先王”惟有湯耳，故知“祠”實是奠，非祠宗廟也。祠之與奠有大小耳，祠則有主有屍，其禮大；奠則奠器而已，其禮小。奠、祠俱是享神，故可以“祠”言奠，亦由于時猶質，未有節文，周時則祠、奠有異，故傳解“祠”為奠耳。 \par}

奉嗣王祗見厥祖，\footnote{居位主喪。見，賢遍反。}侯甸群后咸在，\footnote{在位次。甸，徒遍反。}百官緫己以聽冢宰。\footnote{\CJKunderline{伊尹}制百官，以三公攝冢宰。緫音揔。}\CJKunderline{伊尹}乃明言烈祖之成德，以訓于王。\footnote{湯有功烈之祖，故稱焉。}

{\noindent\zhuan\zihao{6}\fzbyks 傳“湯有”至“稱焉”。正義曰：“湯有功烈之祖”，\CJKunderwave{毛詩}傳文也。“烈”訓業也，湯有定天下之功業，為商家一代之大祖，故以“烈祖”稱焉。 \par}

曰：“嗚呼！古有夏先後，方懋厥德，罔有天災。\footnote{先君謂\CJKunderline{禹}以下、少康以上賢王。言能以德禳災。少，詩照反。上，時掌反。禳，如羊反。}

{\noindent\zhuan\zihao{6}\fzbyks 傳“先君”至“禳災”。正義曰：有夏先君，總指桀之上世，有德之王皆是也。傳舉聖賢者言“\CJKunderline{禹}已下、少康已上”，惟當\CJKunderline{禹}與啟及少康耳。\CJKunderwave{魯語}云：“杼能師\CJKunderline{禹}者也。”杼少康之子，傳蓋以其德衰薄，故斷自少康已上耳。由勉行其德，故無有天災。言能以德禳災也。 \par}

山川鬼神，亦莫不寧，\footnote{莫，無也。言皆安之。}暨鳥獸魚鱉咸若。\footnote{雖微物皆順之，明其餘無不順。暨，其器反。鱉，必滅反。}

{\noindent\shu\zihao{5}\fzkt “山川”至“咸若”。正義曰：“山川鬼神”，謂山川之鬼神也。“亦莫不寧”者，謂鬼神安人君之政。政善則神安之，神安之則降福人君，無妖孽也。“鳥獸魚鱉咸若”者，謂人君順禽魚，君政善而順彼性，取之有時，不夭殺也。鳥獸在陸，魚鱉在水，水陸所生微細之物，人君為政皆順之，明其餘無不順也。 \par}

于其子孫弗率，皇天降災，假手于我有命，\footnote{言桀不循其祖道，故天下禍災，藉手于我有命商王誅討之。}造攻自鳴條，朕哉自亳。\footnote{造、哉，皆始也。始攻桀伐無道，由我始修德于亳。亳，旁各反，徐扶各反。}

{\noindent\shu\zihao{5}\fzkt “于其”至“自亳”。正義曰：“于其子孫”，于有夏先君之子孫，謂桀也。“不循其祖之道，天下禍災”，謂滅其國而誅其身也。天不能自誅于桀，故藉手于我有命之人，謂\CJKunderline{成湯}也。言湯有天命，將為天子，就湯藉手使誅桀也。既受天命誅桀，始攻從鳴條之地而敗之。天所以命我者，由湯始自修德于亳故也。 \par}

惟我商王，布昭聖武，代虐以寬，兆民允懷。\footnote{言湯布明武德，以寬政代桀虐政，兆民以此皆信懷我商王之德。}今王嗣厥德，罔不在初。\footnote{言善惡之由無不在初，欲其慎始。}立愛惟親，立敬惟長，始于家邦，終于四海。\footnote{言立愛敬之道，始于親長，則家國並化，終洽四海。長，丁丈反。}

{\noindent\shu\zihao{5}\fzkt “立愛”至“四海”。正義曰：王者之馭天下，撫兆人，惟愛敬二事而已。\CJKunderwave{孝經·天子之章}盛論愛敬之事，言天子當用愛敬以接物也。行之所立，自近為始。立愛惟親，先愛其親，推之以及疏。立敬惟長，先敬其長，推之以及幼。即\CJKunderwave{孝經}所云“愛親者不敢惡于人，敬親者不敢慢于人”。是推親以及物，始則行于家國，終乃治于四海，即\CJKunderwave{孝經}所云“德教加于百姓，刑于四海”是也。所異者\CJKunderwave{孝經}論愛敬並始于親，令緣親以及疏，此分敬屬長，言從長以及幼耳。 \par}

嗚呼！先王肇修人紀，從諫弗\xpinyin*{咈}，先民時若。\footnote{言湯始修為人綱紀，有過則改，從諫如流，必先民之言是順。咈，扶弗反。}

{\noindent\shu\zihao{5}\fzkt “先民時若”。正義曰：賈逵注\CJKunderwave{周語}云：“先民，古賢人也。”\CJKunderwave{魯語}云“古曰在昔，昔曰先民”，然則先民在古昔之前，遠言之也。遠古賢人亦是民內之一人，故以“民”言之。先民之言于是順從，言其動皆法古賢也。 \par}

居上克明，\footnote{言理恕。}

{\noindent\shu\zihao{5}\fzkt “居上克明”。正義曰：見下之謂“明”，言其以理恕物，照察下情，是能明也。 \par}

為下克忠，\footnote{事上竭誠。}與人不求備，檢身若不及，\footnote{使人必器之。常如不及，恐有過。}

{\noindent\shu\zihao{5}\fzkt “檢身若不及”。正義曰：“檢”謂自攝斂也，檢敕其身，常如不及，不自大以卑人，不恃長以陵物也。 \par}

以至于有萬邦，茲惟艱哉！\footnote{言湯操心常危懼，動而無過，以至為天子，此自立之難。操，七曹反，又七報反。}敷求哲人，俾輔于爾後嗣，\footnote{布求賢智，使師輔于爾嗣王。言仁及後世。哲，本又作喆。俾，必爾反。}制官刑，儆于有位。\footnote{言湯制治官刑法，以儆戒百官。儆，居領反。}曰：‘敢有恆舞于宮，酣歌于室，時謂巫風。\footnote{常舞則荒淫。樂酒曰酣，酣歌則廢德。事鬼神曰巫。言無政。酣，戶甘反。巫音無。樂音洛。}敢有殉于貨色，恆于遊畋，時謂淫風。\footnote{殉，求也。昧求財貨美色，常遊戲畋獵，是淫過之風俗。殉，辭俊反，徐辭荀反。畋音田。}


{\noindent\zhuan\zihao{6}\fzbyks 傳“常舞”至“無政”。正義曰：酣歌常舞併為耽樂無度，荒淫廢德，俱是敗亂政事，其為愆過不甚異也。恆舞酣歌乃為愆耳，若不恆舞、不酣歌非為過也。“樂酒曰酣”，言耽酒以自樂也。\CJKunderwave{說文}亦云：“酣，樂酒也。”\CJKunderwave{楚語}云：“民之精爽不攜貳者，則明神降之，在男曰覡，在女曰巫。”又\CJKunderwave{周禮}有男巫女巫之官，皆掌接神，故“事鬼神曰巫”也。廢棄德義，專為歌舞，似巫事鬼神然，言其無政也。 \par}

{\noindent\zhuan\zihao{6}\fzbyks 傳“殉求”至“風俗”。正義曰：“殉”者心循其事，是貪求之意，故為求也。志在得之，不顧禮義，“昧求”謂貪昧以求之。\CJKunderwave{無逸}云“于遊、于畋”，是“遊”與“畋”別，故為遊戲與畋獵為之無度，是淫過之風俗也。 \par}

敢有侮聖言，逆忠直，遠耆德，比頑童，時謂亂風。\footnote{狎侮聖人之言而不行，拒逆忠直之規而不納，耆年有德疏遠之，童稚頑嚚親比之，是荒亂之風俗。遠，于萬反，注同。耆，巨夷反。比，毗志反，徐扶至反。稚，直利反。嚚，魚巾反。}惟茲三風十愆，卿士有一于身，家必喪；\footnote{有一過則德義廢，失位亡家之道。愆，去幹反。喪如字，又息浪。}邦君有一于身，國必亡。\footnote{諸侯犯此，國亡之道。}臣下不匡，其刑墨，具訓于蒙士。’\footnote{邦君卿士則以爭臣自匡正。臣不正君，服墨刑，鑿其頟，涅以墨。蒙士，例謂下士，士以爭友僕隸自匡正。爭，諫爭之爭。鑿，在洛反。頟，魚白反。涅，乃結反。未，郎計反。}

{\noindent\zhuan\zihao{6}\fzbyks 傳“狎侮”至“風俗”。正義曰：“侮”謂輕慢，“狎”謂慣忽，故傳以“狎”配“侮”而言之。\CJKunderwave{旅獒}云“德盛不狎侮”，是“狎”、“侮”意相類也。 \par}

{\noindent\zhuan\zihao{6}\fzbyks 傳“邦君”至“匡正”。正義曰：言十愆有一，則亡國喪家，邦君卿士慮其喪亡之故，則宜以爭臣自匡正。犯顏而諫，臣之所難，故設不諫之刑以勵臣下，故言“臣不正君，則服墨刑”。墨刑，五刑之輕者。謂“鑿其頟，涅以墨”，\CJKunderwave{司刑}所謂“墨罪五百”者也。“蒙”謂蒙稚，卑小之稱，故“蒙士例謂下士”也。顧氏亦以為“蒙”謂蒙暗之士，“例”字宜從下讀，言此等流例謂下士也。 \par}

{\noindent\shu\zihao{5}\fzkt “曰敢有”至“蒙士”。正義曰：此皆湯所制治官之刑，以儆戒百官之言也。“三風十愆”，謂巫風二，舞也，歌也；淫風四，貨也，色也，遊也，畋也；與亂風四為十愆也。舞及遊、畋，得有時為之，而不可常然，故三事特言“恆”也。歌則可矣，不可樂酒而歌，故以“酣”配之。巫以歌舞事神，故歌舞為巫覡之風俗也。貨色人所貪慾，宜其以義自節，而不可專心殉求，故言“殉于貨色”。心殉貨色，常為遊畋，是謂淫過之風俗也。侮慢聖人之言，拒逆忠直之諫，疏遠耆年有德，親比頑愚幼童，愛惡憎善，國必荒亂，故為“荒亂之風俗”也。此“三風十愆”，雖惡有大小，但有一于身，皆喪國亡家，故各從其類，相配為風俗。“臣下不匡，其刑墨”，言臣無貴賤，皆當匡正君也。“具訓于蒙士”者，謂湯制官刑，非直教訓邦君卿大夫等，使之受諫，亦備具教訓下士，使受諫也。 \par}

“嗚呼！嗣王祇厥身，念哉！\footnote{言當敬身，念祖德。}聖謨洋洋，嘉言孔彰。\footnote{洋洋，美善。言甚明可法。洋音羊，徐音翔。}

{\noindent\shu\zihao{5}\fzkt “聖謨”至“孔彰”。正義曰：此嘆聖人之謨洋洋美善者，謂上湯作官刑，所言三風十愆，令受下之諫，是善言甚明可法也。 \par}

惟上帝不常，作善降之百祥，作不善降之百殃。\footnote{祥，善也。天之禍福，惟善惡所在，不常在一家。}爾惟德罔小，萬邦惟慶。\footnote{修德無小，則天下賚慶。賚，力代反。}爾惟不德罔大，墜厥宗。”\footnote{苟為不德無大，言惡有類，以類相致，必墜失宗廟。此\CJKunderline{伊尹}至忠之訓。}


{\noindent\zhuan\zihao{6}\fzbyks 傳“苟為”至“之訓”。正義曰：“爾惟德”，謂修德以善也。“爾惟不德”，謂不修德為惡也。\CJKunderwave{易·繫辭}曰：“善不積不足以成名，惡不積不足以滅身。”乃謂大善始為福，大惡乃成禍。此訓作勸誘之辭，言為善無小，小善萬邦猶慶，況大善乎？而為惡無，大言小惡猶墜厥宗，況大惡乎？此經二事辭反而意同也。傳“言惡有類”者，解小惡墜宗之意。初為小惡，小惡有族類，以類相致，至于大惡，若致于大惡，必墜失宗廟。言至于大惡乃墜，非小惡即能墜也。\CJKunderwave{晉語}云：“趙文子冠，見韓獻子，曰：‘戒之，此謂成人。成人在始，始與善，善進，不善蔑由至矣。始與不善，不善進，善亦蔑由至矣。’”言惡有類，以類相致也。今\CJKunderline{太甲}初立，恐其親近惡人，以惡類相致禍害，故以言戒之。此是\CJKunderline{伊尹}至忠之訓也。 \par}

{\noindent\shu\zihao{5}\fzkt “爾惟”至“厥宗”。正義曰：又戒王，爾惟修德而為善。德無小，德雖小猶萬邦賴慶，況大善乎？爾惟不德而為惡，惡無大，惡雖小猶墜失其宗廟，況大惡乎？。 \par}

肆命\footnote{陳天命以戒\CJKunderline{太甲}，亡。}

徂後\footnote{陳往古明君以戒，亡。}

\section{太甲上第五【偽】}


\CJKunderline{太甲}既立，不明，\footnote{不用\CJKunderline{伊尹}之訓，不明居喪之禮。}\CJKunderline{伊尹}放諸桐。\footnote{湯葬地也。不知朝政，故曰放。朝，直遙反。}三年復歸于亳，思庸，\footnote{念常道。}\CJKunderline{伊尹}作\CJKunderwave{太甲}三篇。


{\noindent\zhuan\zihao{6}\fzbyks 傳“不用”至“之禮”。正義曰：此篇承\CJKunderwave{伊訓}之下，經稱“不惠于阿衡”，知“不明”者，“不用\CJKunderline{伊尹}之訓”也。“王徂桐宮”，始云“居憂”，是未放已前不明居喪之禮也。 \par}

{\noindent\zhuan\zihao{6}\fzbyks 傳“湯葬”至“曰放”。正義曰：經稱“營于桐宮，密邇先王”，知桐是“湯葬地”也。舜放四凶，徙之遠裔；春秋放其大夫，流之他境；嫌此亦然，故辨之云“不知朝政，故曰放”。使之遠離國都，往居墓側，與彼放逐事同，故亦稱“放”也。古者天子居喪三年，政事聽于冢宰，法當不知朝政，而云“不知朝政，曰放”者，彼正法三年之內，君雖不親政事，冢宰猶尚諮稟，此則全不知政，故為放也。 \par}

{\noindent\shu\zihao{5}\fzkt “\CJKunderline{太甲}”至“三篇”。正義曰：\CJKunderline{太甲}既立為君，不明居喪之禮，\CJKunderline{伊尹}放諸桐宮，使之思過，三年復歸于亳都，以其能改前過，思念常道故也。自初立至放而復歸，\CJKunderline{伊尹}每進言以戒之，史敘其事作\CJKunderwave{太甲}三篇。案經上篇是放桐宮之事，中下二篇是歸亳之事，此序歷言其事以總三篇也。 \par}

太甲\footnote{戒\CJKunderline{太甲}，故以名篇。}

{\noindent\zhuan\zihao{6}\fzbyks 傳“戒\CJKunderline{太甲}，故以名篇”。正義曰：\CJKunderwave{盤庚}、\CJKunderwave{仲丁}、\CJKunderwave{祖乙}等皆是發言之人名篇，此\CJKunderwave{太甲}及\CJKunderwave{沃丁}、\CJKunderwave{君奭}以被告之人名篇，史官不同，故以為名有異。且\CJKunderwave{伊訓}、\CJKunderwave{肆命}、\CJKunderwave{徂後}與此三篇及\CJKunderwave{咸有一德}皆是\CJKunderline{伊尹}戒\CJKunderline{太甲}，不可同名\CJKunderwave{伊訓}，故隨事立稱，以\CJKunderwave{太甲}名篇也。 \par}

惟嗣王不惠于阿衡，\footnote{阿，倚。衡，平。言不順\CJKunderline{伊尹}之訓。倚，于綺反。}


{\noindent\zhuan\zihao{6}\fzbyks 傳“阿倚”至“之訓”。正義曰：古人所讀“阿”、“倚”同音，故“阿”亦倚也。稱上謂之“衡”，故“衡”為平也。\CJKunderwave{詩}毛傳云：“阿衡，\CJKunderline{伊尹}也。”\CJKunderline{鄭玄}亦云：“阿，倚。衡，平也。\CJKunderline{伊尹}，湯倚而取平，故以為官名。” \par}

{\noindent\shu\zihao{5}\fzkt “惟嗣”至“阿衡”。正義曰：\CJKunderline{太甲}以元年十二月即位，比至放桐之時，未知凡經幾月。必是\CJKunderline{伊尹}數諫，久而不順，方始放之，蓋以三五月矣，必是二年放之。序言“三年復歸”者，謂即位三年，非在桐宮三年也。史錄其\CJKunderline{伊尹}訓王，有\CJKunderwave{伊訓}、\CJKunderwave{肆命}、\CJKunderwave{徂後}，其餘忠規切諫，固應多矣。\CJKunderline{太甲}終不從之，故言“不惠于阿衡”。史為作書發端，故言此為目也。 \par}

\CJKunderline{伊尹}作書曰：“先王顧諟天之明命，以承上下神祇。\footnote{顧謂常目在之。諟，是也。言敬奉天命以承順天地。顧音故。諟音是，\CJKunderwave{說文}：“理也。”祇，巨支反。}

{\noindent\zhuan\zihao{6}\fzbyks 傳“顧謂”至“天地”。正義曰：\CJKunderwave{說文}云：“顧，還視也。”“諟”與“是”，古今之字異，故變文為“是”也。言先王每有所行，必還回視是天之明命，謂常目在之。言其想象如目前，終常敬奉天命，以承上天下地之神祇也。 \par}

社稷宗廟，罔不祇肅。\footnote{肅，嚴也。言能嚴敬鬼神而遠之。遠，于萬反。}天監厥德，用集大命，撫綏萬方。\footnote{監，視也。天視湯德，集王命于其身，撫安天下。監，工暫反。}惟\CJKunderline{尹}躬克左右厥辟宅師，\footnote{\CJKunderline{伊尹}言能助其君居業天下之眾。闢，必亦反，徐甫亦反。}

{\noindent\shu\zihao{5}\fzkt “惟尹躬”。正義曰：\CJKunderwave{孫武兵書}及\CJKunderwave{呂氏春秋}皆雲\CJKunderline{伊尹}名摯，則“尹”非名也。今自稱“尹”者，蓋湯得之，使尹正天下，故號曰“\CJKunderline{伊尹}”；人既呼之為“尹”，故亦以“尹”自稱。禮法君前臣名，不稱名者，古人質直，不可以後代之禮約之。 \par}

肆嗣王丕承基緒。\footnote{肆，故也。言先祖勤德，致有天下，故子孫得大承基業，宜念祖修德。丕，普悲反，徐甫眉反。}惟\CJKunderline{尹}躬先見于西邑夏，自周有終，相亦惟終。\footnote{周，忠信也。言身先見夏君臣用忠信有終。夏都在亳西。先見，並如字，注同。}其後嗣王，罔克有終，相亦罔終。\footnote{言桀君臣滅先人之道德，不能終其業，以取亡。相，悉亮反。}嗣王戒哉！祗爾厥辟，闢不闢，忝厥祖。”\footnote{以不終為戒慎之至，敬其君道，則能終。忝，辱也。為君不君，則辱其祖。}王惟庸，罔念聞。\footnote{言\CJKunderline{太甲}守常不改，無念聞\CJKunderline{伊尹}之戒。}\CJKunderline{伊尹}乃言曰：“先王昧爽丕顯，坐以待旦。\footnote{爽，顯皆明也。言先王昧明思大明其德，坐以待旦而行之。昧音妹。}旁求俊彥，啟迪後人，\footnote{旁非一方。美士曰彥。開道後人。言訓戒。俊，本亦作畯。迪，大歷反。}無越厥命以自覆。\footnote{越，墜失也。無失亡祖命而不勤德，以自顛覆。越，于月反，本又作粵。覆，芳服反，注同。}慎乃儉德，惟懷永圖。\footnote{言當以儉為德，思長世之謀。}若虞機張，往省括于度，則釋。\footnote{機，弩牙也。虞，度也。度機，機有度以準望，言修德夙夜思之，明旦行之，如射先省矢括于度，釋則中。省，息井反。括，故活反。度如字。虞度，待洛反。中，丁仲反。}欽厥止，率乃祖攸行，\footnote{止謂行所安止，君止于仁，子止于孝。}惟朕以懌，萬世有辭。”\footnote{言能循汝祖所行，則我喜悅，王亦見嘆美無窮。懌音亦。}

{\noindent\zhuan\zihao{6}\fzbyks 傳“爽顯”至“行之”。正義曰：昭七年\CJKunderwave{左傳}云：“是以有精爽至于神明。”從“爽”以至于“明”，是“爽”謂未大明也。“昧”是晦冥，“爽”是未明，謂夜向晨也。\CJKunderwave{釋詁}云：“丕，大也。顯，光也。”光亦明也。于夜昧冥之時，思欲大明其德，既思得之，坐以待旦而行之。言先王身之勤也。 \par}

{\noindent\zhuan\zihao{6}\fzbyks 傳“旁非”至“訓戒”。正義曰：“旁”謂四方求之，故言“非一方”也。“美士曰彥”，\CJKunderwave{釋訓}文。舍人曰：“國有美士，為人所言道也。” \par}

{\noindent\zhuan\zihao{6}\fzbyks 傳“機弩”至“則中”。正義曰：“括”謂矢末，“機張”、“省括”，則是以射喻也。“機”是轉關，故為弩牙。“虞”訓度也。度機者，機有法度，以準望所射之物，“準望”則解經“虞”也。如射者弩以張訖機關，先省矢括與所射之物，三者于法度相當，乃後釋弦發矢，則射必中矣。言為政亦如是也。 \par}

{\noindent\shu\zihao{5}\fzkt “\CJKunderline{伊尹}”至“有辭”。正義曰：\CJKunderline{伊尹}作書以告，\CJKunderline{太甲}不念聞之。\CJKunderline{伊尹}乃又言曰：“先王以昧爽之時，思大明其德，既思得其事，則坐以待旦，明則行之。其身既勤于政，又乃旁求俊彥之人，置之于位，令以開導後人。先王之念子孫，其憂勤若是，嗣王今承其後，無得墜失其先祖之命，以自覆敗。王當慎汝儉約之德，令其以儉為德而謹慎守之，惟思其長世之謀。謀為政之事，譬若以弩射也。可準度之機已張之，又當以意往省視矢括，當于所度，則釋而放之。如是而射，則無不中矣。猶若人君所修政教，欲發命也，當以意夙夜思之，使當于民心，明旦行之，則無不當矣。王又當敬其身所安止，循汝祖之所行。若能如此，惟我以此喜悅，王于萬世常有善辭，言有聲譽，亦見嘆美無窮也。” \par}

王未克變。\footnote{未能變，不用訓。\CJKunderline{太甲}性輕脫，\CJKunderline{伊尹}至忠，所以不已。輕，遣政反。}

{\noindent\zhuan\zihao{6}\fzbyks 傳“未能”至“不已”。正義曰：“未能變”者，據在後能變，故當時為未能也。時既未變，是不用\CJKunderline{伊尹}之訓也。\CJKunderline{太甲}終為人主，非是全不可移，但體性輕脫，與物推遷，雖有心向善，而為之不固。\CJKunderline{伊尹}至忠，所以進言不已。是\CJKunderline{伊尹}知其可移，故誨之不止，冀其終從己也。 \par}

\CJKunderline{伊尹}曰：“茲乃不義，習與性成。\footnote{言習行不義，將成其性。義，本亦作誼。}予弗狎于弗順，營于桐宮，密邇先王其訓，無俾世迷。\footnote{狎，近也。經營桐墓立宮，令\CJKunderline{太甲}居之。近先王，則訓于義，無成其過，不使世人迷惑怪之。俾，必爾反，後篇同。近，附近之近。令，力呈反。}


{\noindent\zhuan\zihao{6}\fzbyks 傳“狎近”至“怪之”。正義曰：狎習是相近之義，故訓為近也。不順即是近不順也。習為不義，近于不順，則當日日益惡，必至滅亡，故\CJKunderline{伊尹}言已不得使王近于不順，故經營桐墓，立宮墓旁，令\CJKunderline{太甲}居之，不使復知朝政，身見廢退，必當改悔為善也。 \par}

{\noindent\shu\zihao{5}\fzkt “\CJKunderline{伊尹}”至“世迷”。正義曰：\CJKunderline{伊尹}以王未變，乃告于朝廷群臣曰：“此嗣王所行，乃是不義之事。習行此事，乃與性成。”言為之不已，將以不義為性也。“我不得令王近于不順之事，當營于桐墓立宮，使此近先王，當受人教訓之，無得成其過失，使後世人迷惑怪之”。 \par}

王徂桐宮居憂，\footnote{往入桐宮，居憂位。}克終允德。”\footnote{言能思念其祖，終其信德。}

{\noindent\zhuan\zihao{6}\fzbyks 傳“往入”至“憂位”。正義曰：亦既不知朝政之事，惟行居喪之禮。“居憂位”謂服治喪禮也。\CJKunderline{伊尹}亦使兵士衛之，選賢俊教之，故\CJKunderline{太甲}能終信德也。 \par}

\section{太甲中第六【偽】}


惟三祀十有二月朔，\footnote{湯以元年十一月崩，至此二十六月，三年服闋。闋，苦穴反。}\CJKunderline{伊尹}以冕服奉嗣王歸于亳。\footnote{冕，冠也。逾月即吉服。冕音免。}

{\noindent\shu\zihao{5}\fzkt “惟三”至“于亳”。正義曰：周制，君薨之年屬前君，明年始為新君之元年。此殷法，君薨之年而新君即位，即以其年為新君之元年。“惟三祀”者，\CJKunderline{太甲}即位之三年也。湯以元年十一月崩，至此年十一月為再期，除喪服也。至十二月服闋。闋，息也。如喪服息即吉服。舉事貴初始，故于十二月朔以冕服奉嗣王歸于亳。冕是在首之服，冠內之別名，冠是首服之大名，故傳以“冕”為冠。案\CJKunderwave{王制}云：“殷人冔而祭。”\CJKunderwave{大雅}云：“常服黼冔。”冔是殷之祭冠，今云“冕”者，蓋“冕”為通名。\CJKunderwave{王制}又云：“有虞氏皇而祭，夏后氏收而祭，殷人冔而祭，周人冕而祭。”並是當代別名。殷禮不知天子幾冕，\CJKunderwave{周禮}天子六冕，大裘之冕，祭天尚質。弁師惟掌五冕，備物盡文，惟袞冕耳。此以“冕服”，蓋以袞冕之服也。顧氏云：“祥禫之制，前儒不同。”案\CJKunderwave{士虞禮}云“期而小祥”，又“期而大祥”，“中月而禫”。\CJKunderline{王肅}云：“祥月之內又禫祭，服彌寬而變彌數也。”\CJKunderwave{禮記·檀弓}云“祥而縞，是月禫，徙月樂”。\CJKunderline{王肅}云：“是祥之月而禫，禫之明月可以樂矣。”案此孔傳云“二十六月，服闋”，則與\CJKunderline{王肅}同。\CJKunderline{鄭玄}以中月為間一月，云“祥後復更有一月而禫”，則三年之喪凡二十七月，與孔為異。 \par}

作書曰：“民非後，罔克胥匡以生。\footnote{無能相匡，故須君以生。胥，息餘反。}後非民，罔以闢四方。\footnote{須民以君四方。}皇天眷佑有商，俾嗣王克終厥德，實萬世無疆之休。”\footnote{言王能終其德，乃天之顧佑商家，是商家萬世無窮之美。疆，居良反。}王拜手稽首，曰：“予小子不明于德，自底不類。\footnote{君而稽首于臣，謝前過。類，善也。暗于德，故自致不善。底，之履反。}欲敗度，縱敗禮，以速戾于厥躬。\footnote{速，召也。言己放縱情慾，毀敗禮儀法度，以召罪于其身。敗，必邁反，徐甫邁反。縱，子用反。戾，郎計反。}

{\noindent\zhuan\zihao{6}\fzbyks 傳“速召”至“其身”。正義曰：\CJKunderwave{釋言}云：“速，徵也。徵，召也。”轉以相訓，故“速”為召也。“欲”者本之于情，“縱”者放之于外，有欲而縱之，“縱”、“欲”為一也。準法謂之“度”，體見謂之“禮”，“禮”、“度”一也。故傳並釋之，“言己放縱情Q欲Y，毀敗禮儀法度，以召罪于其身”也。 \par}

天作孽，猶可違。自作孽，不可\xpinyin*{逭}。\footnote{孽，災。逭，逃也。言天災可避，自作災不可逃。孽，魚列反。逭，胡亂反。}

{\noindent\zhuan\zihao{6}\fzbyks 傳“孽災”至“可逃”。正義曰：\CJKunderwave{洪範五行傳}有“妖、孽、眚、祥”，\CJKunderwave{漢書·五行志}說云：“凡草物之類謂之妖，妖猶夭胎，言尚微也。蟲豸之類謂之孽，孽則牙孽矣。甚則異物生，謂之眚。自外來謂之祥。”是“孽”為災初生之名，故為災也。“逭，逃也”，\CJKunderwave{釋言}文。樊光云：“行相避逃謂之逭，亦行不相逢也。”天作災者，謂若\CJKunderline{太戊}桑谷生朝，\CJKunderline{高宗}雊雉升鼎耳。可修德以禳之，是“可避”也。“自作災”者，謂若桀放鳴條，紂死宣室，是“不可逃”也。據其將來，修德可去；及其已至，改亦無益。天災自作，逃否亦同。且天災亦由人行而至，非是橫加災也。此\CJKunderline{太甲}自悔之深，故言自作甚于天災耳。 \par}

既往背師保之訓，弗克于厥初，尚賴匡救之德，圖惟厥終。”\footnote{言己已往之前，不能修德于其初，今庶幾賴教訓之德，謀終于善。悔過之辭。背音佩，徐扶代反。}\CJKunderline{伊尹}拜手稽首，\footnote{拜手，首至手。}

{\noindent\zhuan\zihao{6}\fzbyks 傳“拜手，首至手”。正義曰：\CJKunderwave{周禮·太祝}：“辨九拜，一曰稽首，二曰頓首，三曰空首。”\CJKunderline{鄭玄}云：“稽首，拜頭至地也。頓首，拜頭叩地也。空首，拜頭至手，所謂拜手也。”鄭惟解此三者拜之形容，所以為異也。稽首頭至地，頭下至地也。頓首頭下至地，暫一叩之而已。此言“拜手稽首”者，初為拜頭至手，乃復申頭以至于地，至手是為“拜手”，至地乃為“稽首”。然則凡為稽首者，皆先為拜手，乃後為稽首。故“拜手稽首”連言之，諸言“拜手稽首”，義皆同也。\CJKunderwave{太祝}又云：“四曰振動，五曰吉拜，六曰兇拜，七曰奇拜，八曰褒拜，九曰肅拜。”鄭注云，振動者，戰慄變動而拜。吉拜者，拜而後稽顙，謂齊衰不杖以下者之拜。兇拜者，稽顙而後拜，即三年喪拜也。奇拜者，謂君答臣一拜也。褒拜者，謂再拜拜神與屍也。肅拜者，謂揖拜也，禮介者不拜，及婦人之拜也。\CJKunderwave{左傳}云：“天子在，寡君無所稽首。”則諸侯于天子稽首也，諸侯相于則頓首也，君于臣則空首也。 \par}

曰：“修厥身，允德協于下，惟明後。\footnote{言修其身，使信德合于群下，惟乃明君。}先王子惠困窮，民服厥命，罔有不悅。\footnote{言湯子愛睏窮之人，使皆得其所，故民心服其教令，無有不忻喜。}並其有邦厥鄰，乃曰：‘徯我后，後來無罰。’\footnote{湯俱與鄰並有國，鄰國人乃曰：“待我君來。”言忻戴。“君來無罰”，言仁惠。徯，胡啟反。}

{\noindent\shu\zihao{5}\fzkt “並其”至“無罰”。正義曰：言湯昔為諸侯之時，與湯並居其有邦國，謂諸侯之國也。此諸侯國人其與湯鄰近者，皆原以湯為君。乃言曰：“待我后，後來無罰于我。”言羨慕湯德，忻戴之也。 \par}

王懋乃德，視乃厥祖，無時豫怠。\footnote{言當勉修其德，法視其祖而行之，無為是逸豫怠惰。懋音茂。}奉先思孝，接下思恭。\footnote{以念祖德為孝，以不驕慢為恭。}視遠惟明，聽德惟聰。\footnote{言當以明視遠，以聰聽德。}朕承王之休無斁。”\footnote{王所行如此，則我承王之美無厭。斁音亦。厭，于豔反。}

{\noindent\zhuan\zihao{6}\fzbyks 傳“言當”至“聽德”。正義曰：人之心識所知在于聞見，聞見所得在于耳目，故欲言人之聰明，以視聽為主。視若不見，故言“惟明”，“明”謂監察是非也。聽若不聞，故言“惟聰”，“聰”謂識知善惡也。視戒見近迷遠，故言“視遠”。聽戒背正從邪，故言“聽德”。各準其事，相配為文。 \par}

\section{太甲下第七【偽】}


\CJKunderline{伊尹}申誥于王曰：“嗚呼！惟天無親，克敬惟親。\footnote{言天于人無有親疏，惟親能敬身者。}

{\noindent\shu\zihao{5}\fzkt “\CJKunderline{伊尹}申誥于王”。正義曰：\CJKunderline{伊尹}以至忠之心喜王改悔，重告于王，冀王大善，一篇皆誥辭也。天親克敬，民歸有仁，神享克誠，言天民與神皆歸于善也。奉天宜其敬謹，養民宜用仁恩，事神當以誠信，亦準事相配而為文也。 \par}

民罔常懷，懷于有仁。\footnote{民所歸無常，以仁政為常。}鬼神無常享，享于克誠。\footnote{言鬼神不保一人，能誠信者則享其祀。}天位艱哉！\footnote{言居天子之位難，以此三者。}德惟治，否德亂。\footnote{為政以德則治，不以德則亂。治，直吏反，注及下同。}與治同道，罔不興。與亂同事，罔不亡。\footnote{言安危在所任，治亂在所法。}

{\noindent\zhuan\zihao{6}\fzbyks 傳“言安”至“所法”。正義曰：任賢則興，鹽佞則亡，故“安危在所任”。于善則治，于惡則亂，故“治亂在所法”。總言治國則稱“道”，單指所行則言“事”。興難而亡易，道大而事小，故大言“興”而小言“亡”也。此所云“惟言治亂在所法”耳。下句云“終始慎厥與”，言當與賢不與佞，治亂在于用臣，故傳于此言“安危在所任”也。 \par}

“終始慎厥與，惟明明後。\footnote{明慎其所與治亂之機，則為明王明君。}

{\noindent\shu\zihao{5}\fzkt “惟明明後”。正義曰：重言“明明”，言其為大明耳。傳因文重，故言“明王明君”，君、王猶是一也。 \par}

“先王惟時懋敬厥德，克配上帝。\footnote{言湯推是終始所與之難，勉修其德，能配天而行之。}今王嗣有令緒，尚監茲哉！\footnote{令，善也。繼祖善業，當夙夜庶幾視祖此配天之德而法之。}若升高，必自下。若陟遐，必自邇。\footnote{言善政有漸，如登高升遠必用下近為始，然後終致高遠。}無輕民事，惟難。\footnote{無輕為力役之事，必重難之乃可。}無安厥位，惟危。\footnote{言當常自危懼，以保其位。}慎終于始。\footnote{于始慮終，于終思始。}

{\noindent\shu\zihao{5}\fzkt “慎終于始”。正義曰：欲慎其終，于始即須慎之，故傳云“于始慮終”。傳以將終戒惰，故又云“于終思始”，言終始皆當慎也。 \par}

有言逆于汝心，必求諸道。\footnote{人以言咈違汝心，必以道義求其意，勿拒逆之。}有言遜于汝志，必求諸非道。\footnote{遜，順也。言順汝心，必以非道察之，勿以自臧。}嗚呼！弗慮胡獲？弗為胡成？一人元良，萬邦以貞。\footnote{胡，何。貞，正也。言常念慮道德，則得道德；念為善政，則成善政。一人，天子。天子有大善，則天下得其正。}

{\noindent\zhuan\zihao{6}\fzbyks 傳“胡何”至“其正”。正義曰：“胡”之與“何”，方言之異耳。\CJKunderwave{易}彖、象皆以“貞”為正也。\CJKunderline{伊尹}此言，勸王為善，“弗慮”、“弗為”，必是善事。人君善事，惟有道德政教。言不慮何獲，是念慮有所得，知心所念慮是道德也。不為何成，則為之有所成，則知心所念是為善政也。謂天子為“一人”者，其義有二。一則天子自稱“一人”，是為謙辭，言己是人中之一耳。一則臣下謂天子為“一人”，是為尊稱，言天下惟一人而已 \par}

君罔以辯言亂舊政，\footnote{利口覆國家，故特慎焉。}臣罔以寵利居成功\footnote{成功不退，其志無限，故為之極以安之。}，邦其永孚于休。”\footnote{言君臣各以其道，則國長信保于美。}

{\noindent\zhuan\zihao{6}\fzbyks 傳“成功”至“安之”。正義曰：四時之序，成功者退。臣既成功，不知退謝，其志貪慾無限，其君不堪所求，或有怨恨之心，君懼其謀，必生誅殺之計，自古以來，人臣有功不退者皆喪家滅族者眾矣。經稱臣無以寵利居成功者，為之限極以安之也。\CJKunderline{伊尹}告君而言及臣事者，雖復泛說大理，亦見已有退心也。 \par}

\section{咸有一德第八【偽】}


\CJKunderline{伊尹}作\CJKunderwave{咸有一德}。\footnote{言君臣皆有純一之德，以戒\CJKunderline{太甲}。}

{\noindent\shu\zihao{5}\fzkt “\CJKunderline{伊尹}作\CJKunderwave{咸有一德}”。正義曰：\CJKunderline{太甲}既歸于亳，\CJKunderline{伊尹}致仕而退，恐\CJKunderline{太甲}德不純一，故作此篇以戒之。經稱尹躬及湯咸有一德，言己君臣皆有純一之德，戒\CJKunderline{太甲}使君臣亦然。此主戒\CJKunderline{太甲}而言臣有一德者，欲令\CJKunderline{太甲}亦任一德之臣。經云“任官惟賢材，左右惟其人”，是戒\CJKunderline{太甲}使善用臣也。\CJKunderline{伊尹}既放\CJKunderline{太甲}，又迎而復之，是\CJKunderline{伊尹}有純一之德，己為\CJKunderline{太甲}所信，是己君臣純一，欲令\CJKunderline{太甲}法之。 \par}

咸有一德\footnote{即政之後恐其不一，故以戒之。}

{\noindent\shu\zihao{5}\fzkt “咸有一德”。正義曰：此篇終始皆言一德之事，發首至“陳戒于德”敘其作戒之由，已下皆戒辭也。“德”者，得也，內得于心，行得其理，既得其理，執之必固，不為邪見更致差貳，是之謂“一德”也。而凡庸之主，監不周物，志既少決，性復多疑，與智者謀之，與愚者敗之，則是二三其德，不為一也。經云：“德惟一，動罔不吉。德二三，動罔不兇。”是不二三則為一德也。又曰：“終始惟一，時乃日新。”言守一必須固也。\CJKunderline{太甲}新始即政，\CJKunderline{伊尹}恐其二三，故專以一德為戒。 \par}

\CJKunderline{伊尹}既復政厥辟，\footnote{還政\CJKunderline{太甲}。}將告歸，乃陳戒于德。\footnote{告老歸邑，陳德以戒。}


{\noindent\zhuan\zihao{6}\fzbyks 傳“告老”至“以戒”。正義曰：\CJKunderline{伊尹}湯之上相，位為三公，必封為國君。又受邑于畿內，告老致政事于君，欲歸私邑以自安。將離王朝，故陳戒以德也。\CJKunderwave{無逸}云“肆\CJKunderline{祖甲}之享國三十三年”，傳稱\CJKunderline{祖甲}即\CJKunderline{太甲}也。\CJKunderwave{殷本紀}云：“\CJKunderline{太甲}崩，子\CJKunderline{沃丁}立。”\CJKunderwave{沃丁}序云：“\CJKunderline{沃丁}既葬\CJKunderline{伊尹}于亳。”則\CJKunderline{伊尹}卒在\CJKunderline{沃丁}之世。湯為諸侯之時已得\CJKunderline{伊尹}，此至\CJKunderline{沃丁}始卒，\CJKunderline{伊尹}壽年百有餘歲。此告歸之時，已應七十左右也。\CJKunderwave{殷本紀}云：“\CJKunderline{太甲}既立三年，\CJKunderline{伊尹}放之于桐宮。居桐宮三年，悔過反善，\CJKunderline{伊尹}乃迎而受之政。”謂\CJKunderline{太甲}歸亳之歲已為即位六年，與此經相違，馬遷之說妄也。\CJKunderwave{紀年}云，殷仲壬“即位，居亳，其卿士\CJKunderline{伊尹}”。仲壬崩，\CJKunderline{伊尹}乃放\CJKunderline{太甲}于桐而自立也。\CJKunderline{伊尹}即位于\CJKunderline{太甲}七年。\CJKunderline{太甲}潛出自桐，殺\CJKunderline{伊尹}，乃立其子\CJKunderline{伊陟}、伊奮，命復其父之田宅而中分之。案此經序\CJKunderline{伊尹}奉\CJKunderline{太甲}歸于亳，其文甚明。\CJKunderwave{左傳}又稱“\CJKunderline{伊尹}放\CJKunderline{太甲}而相之”，\CJKunderwave{孟子}云“有\CJKunderline{伊尹}之志則可，無\CJKunderline{伊尹}之志則篡”，\CJKunderline{伊尹}不肯自立，\CJKunderline{太甲}不殺\CJKunderline{伊尹}也。必若\CJKunderline{伊尹}放君自立，\CJKunderline{太甲}起而殺之，則\CJKunderline{伊尹}死有餘罪，義當汙宮滅族，\CJKunderline{太甲}何所感德而復立其子，還其田宅乎？\CJKunderwave{紀年}之書，晉太康八年汲郡民發魏安僖王冢得之，蓋當時流俗有此妄說，故其書因記之耳。 \par}

{\noindent\shu\zihao{5}\fzkt “\CJKunderline{伊尹}”至“于德”。正義曰：自\CJKunderline{太甲}居桐，而\CJKunderline{伊尹}秉政。\CJKunderline{太甲}既歸于亳，\CJKunderline{伊尹}還政其君，將欲告老歸其私邑，乃陳言戒王于德，以一德戒王也。\CJKunderline{太甲}既得復歸，\CJKunderline{伊尹}即應還政，其告歸陳戒，未知在何年也。下云“今嗣王新服厥命”，則是初始即政，蓋\CJKunderline{太甲}居亳之後即告老也。\CJKunderwave{君奭}云：“在\CJKunderline{太甲}，時則有若保衡。”保衡，\CJKunderline{伊尹}也。襄二十一年\CJKunderwave{左傳}云：“\CJKunderline{伊尹}放\CJKunderline{太甲}而相之，卒無怨色。”則\CJKunderline{伊尹}又相\CJKunderline{太甲}。蓋\CJKunderline{伊尹}此時將欲告歸，\CJKunderline{太甲}又留之為相，如\CJKunderline{成王}之留\CJKunderline{周公}，不得歸也。 \par}

曰：“嗚呼！天難\xpinyin*{諶}，命靡常。\footnote{以其無常，故難信。諶，徐巿林反。}常厥德，保厥位。厥德匪常，九有以亡。\footnote{人能常其德，則安其位。九有，諸侯。桀不能常其德，湯伐而兼之。}

{\noindent\shu\zihao{5}\fzkt “九有以亡”。正義曰：\CJKunderwave{毛詩}傳云：“九有，九州也。”此傳云“九有，諸侯”，謂九州所有之諸侯。\CJKunderline{伊尹}此言，泛說大理，未指\CJKunderline{夏桀}，但傳顧下文比桀，為此言之驗，故云“桀不能常其德，湯伐而兼之”。 \par}

夏王弗克庸德，慢神虐民。\footnote{言桀不能常其德，不敬神明，不恤下民。}皇天弗保，監于萬方，啟迪有命，\footnote{言天不安桀所為，廣視萬方，有天命者開道之。}眷求一德，俾作神主。\footnote{天求一德，使伐桀為天地神祇之主。}惟\CJKunderline{尹}躬暨\CJKunderline{湯}，咸有一德，克享天心，受天明命，\footnote{享，當也。所徵無敵，謂之受天命。}

{\noindent\zhuan\zihao{6}\fzbyks 傳“享當”至“天命”。正義曰：德當神意，神乃享之，故以“享”為當也。天道遠而人道近，天之命人，非有言辭文誥，正以神明祐之，使之所徵無敵，謂之受天命也。緯候之書乃稱有黃龍白龜白魚赤雀負圖銜書以授聖人，正典無其事也。漢自哀平之間，緯候始起，假託鬼神，妄稱祥瑞，孔時未有其說，縱使時已有之，亦非孔所信也。 \par}

以有九有之師，爰革夏正。\footnote{爰，于也。于得九有之眾，遂伐夏勝之，改其正。}非天私我有商，惟天佑于一德。\footnote{非天私商而王之，佑助一德，所以王。王，于況反，下同；或如字。}非商求于下民，惟民歸于一德。\footnote{非商以力求民，民自歸于一德。}德惟一，動罔不吉。德二三，動罔不兇。\footnote{二三，言不一。}惟吉凶不僣在人，惟天降災祥在德。\footnote{行善則吉，行惡則兇，是不差。德一，天降之善；不一，天降之災；是在德。僣，子念反。}

{\noindent\shu\zihao{5}\fzkt “惟吉”至“在德”。正義曰：指其已然，則為“吉凶”；言其徵兆，則曰“災祥”；其事不甚異也。吉凶已成之事，指人言之，故曰“在人”。災祥未至之徵，行之所招，故言“在德”。“在德”謂為德有一與不一，“在人”謂人行有善與不善也。吉凶已在其身，故不言來處；災祥自外而至，故言“天降”；其實吉凶亦天降也。 \par}

“今嗣王新服厥命，惟新厥德。\footnote{其命，王命。新其德，戒勿怠。}終始惟一，時乃日新。\footnote{言德行終始不衰殺，是乃日新之義。行，下孟反。殺，色界反，衰微也，殺害也，言小小害也。}任官惟賢材，左右惟其人。\footnote{官賢才而任之，非賢材不可任。選左右，必忠良。不忠良，非其人。}


{\noindent\zhuan\zihao{6}\fzbyks 傳“其命”至“勿怠”。正義曰：\CJKunderwave{說命}云：“王言惟作命。”成十八年\CJKunderwave{左傳}云：“人之求君，使出命也。”是言人君職在發命。“新服厥命”，新始服行王命，故云“其命，王命”也。“新其德”者，勤行其事，日日益新，戒王勿懈怠也。 \par}

{\noindent\zhuan\zihao{6}\fzbyks 傳“言德”至“之義”。正義曰：“日新”者，日日益新也。若今日勤而明日惰，昨日是而今日非，自旁觀之，則有新有舊。言王德行終始皆同，不有衰殺，從旁觀之，每日益新，是乃“日新”之義也。 \par}

{\noindent\zhuan\zihao{6}\fzbyks 傳“官賢”至“其人”。正義曰：“任官”謂任人以官，故云“官賢才而任之”，言官用賢才而委任之。\CJKunderwave{詩序}云“任賢使能”，非賢才不可任也。\CJKunderwave{冏命}云：“小大之臣，咸懷忠良。”故言“選左右，必忠良”，不忠良，即是非其人。“任官”是用人為官，“左右”亦是任而用之，故言“選左右”也。直言其人，“人”字不見，故據\CJKunderwave{冏命}之文，以“忠良”充之。 \par}

臣為上為德，為下為民。\footnote{言臣奉上佈德，順下訓民，不可官所私，任非其人。為上，于偽反，下“為民”同。為德，如字，下“為下”同，徐皆于偽反。}其難其慎，惟和惟一。\footnote{其難無以為易，其慎無以輕之，群臣當和一心以事君，政乃善。易，以豉反。}

{\noindent\zhuan\zihao{6}\fzbyks 傳“言臣”至“其人”。正義曰：“言臣奉上佈德”者，“奉上”謂奉為在上，解經“為上”也；“佈德”者謂布為道德，解經“為德”也。“順下訓民”者，“順下”謂卑順以為臣下，解經“為下”也；“訓民”者，謂以善道訓助下民，解經“為民”也。顧氏亦同此解。 \par}

{\noindent\zhuan\zihao{6}\fzbyks 傳“其難”至“乃善”。正義曰：此經申上臣事既所為如此，其難無以為易，其慎無以輕忽之，戒臣無得輕易臣之職也。既事不可輕，宜和協奉上，群臣當一心以事君，如此政乃善耳。一心即一德，言臣亦當一德也。 \par}

{\noindent\shu\zihao{5}\fzkt “今嗣王”至“惟一”。正義曰：上既言“在德”，此指戒嗣王，今新始服其王命，惟當新其所行之德。所云“新”者，終始所行，惟常如一，無有衰殺之時，是乃“日新”也。王既身行一德，臣亦當然。任人為官，惟用其賢材。輔弼左右，惟當用其忠良之人，乃可為左右耳。此“任官”、“左右”即王之臣也。臣之為用，所施多矣。何者？言臣之助為在上，當施為道德；身為臣下，當須助為于民也。臣之既當為君，又須為民，故不可任非其才，用非其人。此臣之所職，其事甚難，無得以為易。其事須慎，無得輕忽。為臣之難如此，惟當群臣和順，惟當共秉一心，以此事君，然後政乃善耳。言君臣宜皆有一德。 \par}

德無常師，主善為師。\footnote{德非一方，以善為主，乃可師。}善無常主，協于克一。\footnote{言以合于能一為常德。}俾萬姓咸曰：‘大哉！王言。’\footnote{一德之言，故曰大。}又曰：‘一哉！王心。’\footnote{能一德，則一心。}克綏先王之祿，永厎烝民之生。\footnote{言為王而令萬姓如此，則能保安先王之寵祿，長致眾民所以自生之道，是明王之事。烝，之丞反。}嗚呼！七世之廟，可以觀德。\footnote{天子立七廟，有德之王則為祖宗，其廟不毀，故可觀德。}萬夫之長，可以觀政。\footnote{能整齊萬夫，其政可知。長，丁丈反。}


{\noindent\zhuan\zihao{6}\fzbyks 傳“天子”至“觀德”。正義曰：天子立七廟，是其常事，其有德之王則列為祖宗，雖七廟親盡，而其廟不毀，故于七廟之外可以觀德矣。下云“萬夫之長，可以觀政”，謂觀其萬夫之長。此“七世之廟，可以觀德”，謂觀七世之外。文雖同而義小異耳，所謂辭不害意。漢氏以來，論七廟諸多矣，其文見于記傳，\CJKunderwave{禮器}、\CJKunderwave{家語}、\CJKunderwave{荀卿書}、\CJKunderwave{穀梁傳}皆曰天子立七廟，以為天子常法，不辨其廟之名。\CJKunderwave{王制}云：“天子七廟，三昭三穆，與太祖之廟而七。”\CJKunderwave{祭法}云：“王立七廟，曰考廟，曰王考廟，曰皇考廟，曰顯考廟，曰祖考廟，皆月祭之。遠廟為祧，有二祧，享嘗乃止。”\CJKunderwave{漢書}韋玄成議曰：“周之所以七廟者，后稷始封，\CJKunderline{文王}\CJKunderline{武王}受命而王，是以三廟不毀，與親廟四而七也。”\CJKunderline{鄭玄}用此為說。惟周有七廟，二祧為\CJKunderline{文王}\CJKunderline{武王}廟也，故\CJKunderline{鄭玄}\CJKunderwave{王制}注云：“此周制。七者，太祖及\CJKunderline{文王}\CJKunderline{武王}二祧，與親廟四。太祖，后稷也。殷則六廟，契及湯與二昭二穆。夏則五廟，無太祖，\CJKunderline{禹}與二昭二穆而已。”良由不見古文，故為此謬說。此篇乃是\CJKunderwave{商書}，已云“七世之廟”，則天子立七廟，王者常禮，非獨周人始有七廟也。文武則為祖宗，不在昭穆之數，\CJKunderwave{王制}之文不得云“三昭三穆”也。劉歆、\CJKunderline{馬融}、\CJKunderline{王肅}雖則不見古文，皆以七廟為天子常禮。所言二祧者，\CJKunderline{王肅}以為高祖之父及祖也，並高祖已下共為三昭三穆耳。\CJKunderwave{喪服小記}云：“王者禘其祖之所自出，以其祖配之而立四廟。庶子王亦如之。”所以不同者，\CJKunderline{王肅}等以為受命之王是初基之王，故立四廟。“庶子王”者，謂庶子之後自外繼立，雖承正統之後，自更別立己之高祖已下之廟，猶若漢宣帝別立戾太子悼皇考廟之類也。或可庶子初基為王，亦得與嫡子同，正立四廟也。 \par}

{\noindent\shu\zihao{5}\fzkt “嗚呼”至“觀政”。正義曰：此又勸王修德以立後世之名。禮王者祖有功、宗有德，雖七世之外其廟不毀。嗚呼！七世之廟其外則猶有不毀者，可以觀知其有明德也。立德在于為政，萬夫之長能使其整齊，可以觀知其善政也。萬夫之長尚爾，況天子乎？觀王使為善政也。 \par}

後非民罔使，民非後罔事。\footnote{君以使民自尊，民以事君自生。}無自廣以狹人，匹夫匹婦，不獲自盡，民主罔與成厥功。”\footnote{上有狹人之心，則下無所自盡矣。言先盡其心，然後乃能盡其力，人君所以成功。狹，戶夾反。盡，徐子忍反，注同。}

{\noindent\shu\zihao{5}\fzkt “無自”至“厥功”。正義曰：既言君民相須，又戒王虛心待物。凡為人主，無得自為廣大，以狹小前人，勿自以所知為大，謂彼所知為小。若謂彼狹小，必待之輕薄。彼知遇薄，則意不自盡。匹夫匹婦不得自盡其意，則在下不肯親上，在上不得下情，如是則人主無與成其功也。 \par}

沃丁既葬\CJKunderline{伊尹}于亳，\footnote{沃丁，\CJKunderline{太甲}子。\CJKunderline{伊尹}既致仕老終，以三公禮葬。沃，烏毒反，徐于毒反。}\CJKunderline{咎單}遂訓\CJKunderline{伊尹}事，\footnote{訓暢其所行功德之事。作\CJKunderwave{沃丁}。\CJKunderline{咎單}，忠臣名。作此篇以戒也，亡。}

{\noindent\zhuan\zihao{6}\fzbyks 傳“\CJKunderline{沃丁}”至“禮葬”。正義曰：\CJKunderwave{世本}、\CJKunderwave{本紀}皆云“\CJKunderline{太甲}崩，子\CJKunderline{沃丁}立”，是為\CJKunderline{太甲}子也。\CJKunderline{伊尹}本是三公，上篇言其告歸，知“致仕老終，以三公禮葬”。皇甫謐云：“\CJKunderline{沃丁}八年，\CJKunderline{伊尹}卒，卒年百有餘歲。大霧三日。\CJKunderline{沃丁}葬之以天子禮，葬祀以太牢，親臨喪，以報大德。”晉文請遂，襄王不許，\CJKunderline{沃丁}不當以天子之禮葬\CJKunderline{伊尹}也。孔言三公禮葬，未必有文，要情事當然也。 \par}

{\noindent\shu\zihao{5}\fzkt “\CJKunderline{沃丁}”至“作\CJKunderline{沃丁}”。正義曰：\CJKunderline{沃丁}，殷王名也。“\CJKunderline{沃丁}既葬\CJKunderline{伊尹}”，言重其賢德，備禮而葬之。\CJKunderline{咎單}以\CJKunderline{沃丁}愛慕\CJKunderline{伊尹}，遂訓暢\CJKunderline{伊尹}之事以告\CJKunderline{沃丁}。史錄其事，作\CJKunderwave{沃丁}之篇。 \par}

\CJKunderline{伊陟}相\CJKunderline{大戊}，\footnote{伊陟，\CJKunderline{伊尹}子。\CJKunderline{太戊}，\CJKunderline{沃丁}弟之子。陟，張力反。相，息亮反。\CJKunderline{太戊}，馬云：“\CJKunderline{太甲}子。”}亳有祥，桑谷共生于朝。\footnote{祥，妖怪。二木合生，七日大拱，不恭之罰。桑，蘇臧反。谷，工木反，楮也。朝，直遙反。}\CJKunderline{伊陟}贊于\CJKunderline{巫咸}，作\CJKunderwave{咸乂}四篇。\footnote{贊，告也。\CJKunderline{巫咸}，臣名。皆亡。\CJKunderline{巫咸}，馬云：“巫，男巫也，名咸，殷之巫也。”}

{\noindent\zhuan\zihao{6}\fzbyks 傳“\CJKunderline{伊陟}”至“之子”。正義曰：“\CJKunderline{伊陟}，\CJKunderline{伊尹}子”，相傳為然。\CJKunderwave{殷本紀}云：“\CJKunderline{沃丁}崩，弟太庚立。崩，子小甲立。崩，弟雍己立。崩，弟\CJKunderline{太戊}立。”是\CJKunderline{太戊}為小甲弟，太庚之子。 \par}

{\noindent\zhuan\zihao{6}\fzbyks 傳“祥妖”至“之罰”。正義曰：\CJKunderwave{漢書·五行志}云：“凡草物之類謂之妖，自外來謂之祥。”“祥”是惡事先見之徵，故為“妖怪”也。“二木合生”謂共處生也。“七日大拱”，\CJKunderline{伏生}\CJKunderwave{書傳}有其文，或當別出餘書，則孔用之也。\CJKunderline{鄭玄}注\CJKunderwave{書傳}云：“兩手扼之曰拱。”生七日而見其大滿兩手也。\CJKunderwave{殷本紀}云“一暮大拱”，言一夜即滿拱，所聞不同，故說異也。\CJKunderwave{五行傳}曰：“貌之不恭，是謂不肅。時則有青眚之祥。”\CJKunderwave{漢書·五行志}夏侯始昌、劉向等說云：“肅，敬也。內曰恭，外曰敬。人君行己，體貌不恭，怠慢驕蹇，則不能敬。木色青，故有青眚之祥。”是言木之變怪，是貌不恭之罰。人君貌不恭，天將罰之，木怪見其徵也。皇甫謐云：“\CJKunderline{太戊}問于\CJKunderline{伊陟}，\CJKunderline{伊陟}曰：‘臣聞妖不勝德，帝之政事有闕。’白帝修德。\CJKunderline{太戊}退而佔之曰：‘桑谷野木而不合生于朝，意者朝亡乎？’\CJKunderline{太戊}懼，修先王之政，明養老之禮，三年而遠方重譯而至七十六國。”是言妖不勝德也。 \par}

{\noindent\zhuan\zihao{6}\fzbyks 傳“贊告”至“臣名”。正義曰：禮有贊者，皆以言告人，故“贊”為告也。\CJKunderwave{君奭}傳曰：“巫，氏也。”當以“巫”為氏，名“咸”。此言“臣名”者，言是臣之名號也。\CJKunderline{鄭玄}云“\CJKunderline{巫咸}謂之巫官”者，案\CJKunderwave{君奭}咸子又稱，賢父子併為大臣，必不世作巫官，故孔言“巫，氏”是也。 \par}

{\noindent\shu\zihao{5}\fzkt “\CJKunderline{伊陟}”至“四篇”。正義曰：\CJKunderline{伊陟}輔相\CJKunderline{太戊}，于亳都之內，有不善之祥，桑谷二木共生于朝。朝非生木之處，是為不善之徵，\CJKunderline{伊陟}以此桑谷之事告于\CJKunderline{巫咸}。史錄其事，作\CJKunderwave{咸乂}四篇。“乂”訓治也，言所以致妖，須治理之，故名篇為\CJKunderwave{咸乂}也。\CJKunderline{伊陟}不先告\CJKunderline{太戊}而告\CJKunderline{巫咸}者，\CJKunderwave{君奭}云：“在\CJKunderline{太戊}，時則有若\CJKunderline{巫咸}乂王家。”則咸是賢臣，能治王事，大臣見怪而懼，先共議論，而後以告君。下篇序云：“\CJKunderline{太戊}贊于\CJKunderline{伊陟}。”明先告于\CJKunderline{巫咸}，而後告\CJKunderline{太戊}。 \par}

\CJKunderline{太戊}贊于\CJKunderline{伊陟}，\footnote{告以改過自新。}作\CJKunderwave{伊陟}、\CJKunderwave{原命}。\footnote{原，臣名。\CJKunderwave{原命}、\CJKunderwave{伊陟}二篇皆亡。}

{\noindent\shu\zihao{5}\fzkt “\CJKunderline{太戊}”至“原命”。正義曰：言\CJKunderline{太戊}贊于\CJKunderline{伊陟}，惟告\CJKunderline{伊陟}，不告原也。史錄其事，而作\CJKunderwave{伊陟}、\CJKunderwave{原命}二篇，則\CJKunderline{太戊}告\CJKunderline{伊陟}，亦告原，俱以桑谷事告，故序裛以為文也。“原”是臣名而云“原命”，謂以言命原，故以“原命”名篇，猶如\CJKunderwave{冏命}、\CJKunderwave{畢命}也。 \par}

\CJKunderline{仲丁}遷于囂，\footnote{\CJKunderline{太戊}子。去亳。囂，地名。囂，五羔反。}

{\noindent\zhuan\zihao{6}\fzbyks 傳“\CJKunderline{太戊}”至“地名”。正義曰：此及下傳言\CJKunderline{仲丁}是\CJKunderline{太戊}之子，\CJKunderline{河亶甲}，\CJKunderline{仲丁}弟也。\CJKunderline{祖乙}，\CJKunderline{河亶甲}子，皆\CJKunderwave{世本}文也。\CJKunderline{仲丁}是\CJKunderline{太戊}之子，\CJKunderline{太戊}之時仍雲亳有祥，知\CJKunderline{仲丁}遷于囂去亳也。 \par}

{\noindent\shu\zihao{5}\fzkt “\CJKunderline{仲丁}遷于囂”。正義曰：此三篇皆是遷都之事，俱以君名名篇，並陳遷都之義，如\CJKunderwave{盤庚}之誥民也。發其舊都謂之“遷”，到彼新邑謂之“居”，“遷于囂”與“居相”亦事同也。以“\CJKunderline{河亶甲}”三字句長，不言“于”，其實亦是俱于相也。“圮于耿”者，孔意以為毀于相地，乃遷于耿地，其篇蓋言毀意，故序特言“圮”也。李顒云：“囂在陳留浚儀縣。”皇甫謐云：“\CJKunderline{仲丁}自亳徙囂，在河北也。或曰今河南敖倉，二說未知孰是也。”相地孔云“在河北”，蓋有文而知也。謐又以耿在河東，皮氏縣耿鄉是也。 \par}

作\CJKunderwave{仲丁}。\footnote{陳遷都之義，亡。}\CJKunderline{河亶甲}居相，\footnote{\CJKunderline{仲丁}弟。相，地名，在河北。亶，丁但反。相，息亮反，在河北，今魏郡有相縣。}作\CJKunderwave{河亶甲}。\footnote{亡。}\CJKunderline{祖乙}圮于耿，\footnote{\CJKunderline{亶甲}子。圮于相，遷于耿。河水所毀曰圮。圮，備美反，徐扶鄙反。馬云：“毀也。”}作\CJKunderwave{祖乙}。\footnote{亡。}

{\noindent\zhuan\zihao{6}\fzbyks 傳“亶甲”至“曰圮”。正義曰：孔以\CJKunderline{河亶甲}居相，\CJKunderline{祖乙}即亶甲之子，故以為圮于相地，乃遷都于耿。\CJKunderwave{釋詁}云：“圮，毀也。”故云“河水所毀曰圮”。據文“圮于耿”也，知非圮毀于耿，更遷餘處，必雲圮于相地，遷于耿者，明與其上文連。上云“遷于囂”，謂遷來向囂；“居于相”，謂居于相地，故知“圮于耿”謂遷來于耿，以文相類，故孔為此解。謂古人之言，雖尚要約，皆使言足其文，令人曉解，若圮于相，遷居于耿，經言“圮于耿”，大不辭乎？且亶甲居于相，\CJKunderline{祖乙}居耿，今為水所毀，更遷他處，故言毀于耿耳，非既毀乃遷耿也。\CJKunderwave{盤庚}云：“不常厥邑，于今五邦。”及其數之，惟有亳、囂、相、耿四處而已。知此既毀于耿，更遷一處，\CJKunderline{盤庚}又自彼處而遷于殷耳。\CJKunderwave{殷本紀}云：“\CJKunderline{祖乙}遷于邢。”馬遷所為說耳。\CJKunderline{鄭玄}云：“\CJKunderline{祖乙}又去相居耿，而國為水所毀，于是修德以御之，不復徙也。錄此篇者善其國圮毀，改政而不徙。”如鄭所言，稍為文便。但上有\CJKunderwave{仲丁}、\CJKunderwave{亶甲}，下有\CJKunderwave{盤庚}，皆為遷事，作書述其遷意，匆陽毀而不遷，序當改文見義，不應文類遷居，更以不遷為義。汲冢古文云“\CJKunderline{盤庚}自奄遷于殷”者，蓋\CJKunderline{祖乙}圮于耿，遷于奄，\CJKunderline{盤庚}自奄遷于殷，亳、囂、相、耿與此奄五邦者。此蓋不經之書，未可依信也。 \par}

%%% Local Variables:
%%% mode: latex
%%% TeX-engine: xetex
%%% TeX-master: "../Main"
%%% End:

%% -*- coding: utf-8 -*-
%% Time-stamp: <Chen Wang: 2024-04-02 11:42:41>

% {\noindent \zhu \zihao{5} \fzbyks } -> 注 (△ ○)
% {\noindent \shu \zihao{5} \fzkt } -> 疏

\chapter{卷九}


\section{盤庚上第九（盤庚）}


盤庚五遷，將治亳殷，\footnote{自湯至盤庚凡五遷都，盤庚治亳殷。盤，本又作般，步幹反。治，直吏反。}民諮胥怨，\footnote{胥，相也。民不欲徙，乃諮嗟憂愁，相與怨上。胥，徐思餘反。怨，紆萬反。}作\CJKunderwave{盤庚}三篇。

{\noindent\zhuan\zihao{6}\fzbyks 傳“自湯”至“亳怨”。正義曰：經言“不常厥邑，於今五邦”，故序言“盤庚五遷”。傳嫌一身五遷，故辨之雲“自湯至盤庚凡五遷都”也。上文言“自契至於成湯八遷”，並數湯為八，此言“盤庚五遷”，又並數湯為五，湯一人再數，故班固云：“殷人屢遷，前八後五，其實正十二也。”此序雲盤庚“將治亳殷”，下傳雲“殷，亳之別名”，則“亳殷”即是一都，湯遷還從先王居也。汲冢古文云：“盤庚自奄遷於殷，殷在鄴南三十里。”束晳云：“\CJKunderwave{尚書序}‘盤庚五遷，將治亳殷’，舊說以為居亳，亳殷在河南。\CJKunderline{孔子}壁中\CJKunderwave{尚書}雲‘將始宅殷’，是與古文不同也。\CJKunderwave{漢書·項羽傳}雲‘洹水南殷墟上’，今安陽西有殷。”束晳以殷在河北，與亳異也。然\CJKunderline{孔子}壁內之書，安國先得其本，此“將治亳殷”不可作“將始宅殷”。“亳”字摩滅，容或為“宅”。壁內之書，安國先得，“始”皆作“亂”，其字與“始”不類，無緣誤作“始”字，知束晳不見壁內之書，妄為說耳。若洹水南有殷墟，或當餘王居之，非盤庚也。盤庚治於亳殷，紂滅在於朝歌，則盤庚以後遷於河北，蓋盤庚后王有從河有亳地遷於洹水之南，後又遷於朝歌。 \par}

{\noindent\zhuan\zihao{6}\fzbyks 傳“胥相”至“怨上”。正義曰：\CJKunderwave{釋詁}云：“胥，皆也。”“相”亦是皆義，故通訓“胥”為相也。民不欲徙，乃諮嗟憂愁，相與怨上，經雲“民不適有居”，是怨上之事也。仲丁、祖乙亦是遷都，序無民怨之言，此獨有怨者，盤庚，祖乙之曾孫也，祖乙遷都於此，至今多歷年世，民居已久，戀舊情深；前王三徙，誥令則行，曉喻之易，故無此言；此則民怨之深，故序獨有此事。彼各一篇，而此獨三篇者，謂民怨上，故勸誘之難也。民不欲遷，而盤庚必遷者，\CJKunderline{鄭玄}云：“祖乙居耿後，奢侈逾禮，土地迫近山川，嘗圮焉。至陽甲立，盤庚為之臣，乃謀徙居湯舊都。”又序注云：“民居耿久，奢淫成俗，故不樂徙。”王肅云：“自祖乙五世至盤庚，元兄陽甲，宮室奢侈，下民邑居墊隘，水泉瀉滷，不可以行政化，故徙都於殷。”皇甫謐云：“耿在河北，迫近山川，自祖辛已來，民皆奢侈，故盤庚遷於殷。”此三者之說皆言奢侈，\CJKunderline{鄭玄}既言君奢，又言民奢，王肅專謂君奢，皇甫謐專謂民奢。言君奢者以天子宮室奢侈，侵奪下民；言民奢者以豪民室宇過度，逼迫貧乏；皆為細民弱劣無所容居，欲遷都改制以寬之。富民戀舊，故違上意，不欲遷也。案檢孔傳無奢侈之語，惟下篇雲“今我民用蕩析離居，罔有定極”，傳云：“水泉沈溺，故蕩析離居，無安定之極，徙以為之極。”孔意蓋以地勢洿下，又久居水變，水泉瀉滷，不可行化，故欲遷都，不必為奢侈也。此以君名名篇，必是為君時事，而\CJKunderline{鄭玄}以為上篇是盤庚為臣時事，何得專輒謬妄也！ \par}

{\noindent\shu\zihao{5}\fzkt “盤庚”至“三篇”。正義曰：商自成湯以來屢遷都邑，仲丁、河亶甲、祖乙皆有言誥，歷載於篇。盤庚最在其後，故序總之，“自湯至盤庚凡五遷都”。今盤庚將欲遷居，而治於亳之殷治，民皆戀其故居，不欲移徙，諮嗟憂愁，相與怨上，盤庚以言辭誥之。史敘其事，作\CJKunderwave{盤庚}三篇。 \par}

盤庚\footnote{盤庚，殷王名。殷質，以名篇。盤庚，馬云：“祖乙曾孫，祖丁之子。不言‘盤庚誥’何？非但錄其誥也，取其徙而立功，故以“盤庚”名篇。}


{\noindent\zhuan\zihao{6}\fzbyks 傳“殷質，以名篇”。正義曰：\CJKunderwave{周書}諡法成王時作，故桓六年\CJKunderwave{左傳}云：“周人以諱事神。”殷時質，未諱君名，故以王名名篇也。上\CJKunderwave{仲丁}、\CJKunderwave{祖乙}亦是王名，於此始作傳者，以上篇經亡，此經稱\CJKunderwave{盤庚}，故就此解之。\CJKunderwave{史記·殷本紀}云：“盤庚崩，弟小辛立。殷復衰，百姓思盤庚，乃作\CJKunderwave{盤庚}三篇。”與此序違，非也。\CJKunderline{鄭玄}云：“盤庚，湯十世孫，祖乙之曾孫，以五遷繼湯，篇次\CJKunderwave{祖乙}，故繼之。於上累之，祖乙為湯玄孫，七世也，又加祖乙，復其祖父，通盤庚，故十世。”\CJKunderwave{本紀}云：“祖乙崩，子祖辛立。崩，子開甲立。崩，弟祖丁立。崩，門甲之子南庚立。崩，祖丁子陽甲立。崩，弟盤庚立。”是祖乙生祖辛，祖辛生祖丁，祖丁生盤庚，故為曾孫。 \par}

{\noindent\shu\zihao{5}\fzkt “盤庚”。正義曰：此三篇皆以民不樂遷，開解民意，告以不遷之害，遷都之善也。中上二篇，未遷時事，下篇既遷後事。上篇人皆怨上，初啟民心，故其辭尤切。中篇民已少悟，故其辭稍緩。下篇民既從遷，故辭復益緩。哀十一年\CJKunderwave{左傳}引此篇雲“盤庚之誥”，則此篇皆誥辭也。題篇不曰“盤庚誥”者，王肅云：“取其徙而立功，故但以‘盤庚’名篇。”然\CJKunderwave{仲丁}、\CJKunderwave{祖乙}、\CJKunderwave{河亶甲}等皆以王名篇，則是史意異耳，未必見他義。 \par}

盤庚遷於殷，\footnote{亳之別名。}民不適有居。\footnote{適，之也，不欲之殷有邑居。}率籲眾戚，出矢言，\footnote{籲，和也。率和眾憂之人，出正直之言。籲音喻。戚，千歷反。}曰：“我王來，既爰宅於茲，\footnote{我王祖乙居耿。爰，於也。言祖乙已居於此。}重我民，無盡劉。\footnote{劉，殺也。所以遷此，重我民，無慾盡殺故。盡，子忍反。}



{\noindent\zhuan\zihao{6}\fzbyks 傳“亳之別名”。正義曰：此序先“亳”後“殷”，“亳”是大名，“殷”是亳內之別名。\CJKunderline{鄭玄}云：“商家自徙此而號曰殷。”鄭以此前未有殷名也。中篇云：“殷降大虐。”將遷於殷，先正其號，明知於此號為殷也。雖兼號為殷，而商名不改，或稱商，或稱殷，又有兼稱殷商。\CJKunderwave{商頌}雲“商邑翼翼”、“撻彼殷武”是單稱之也。又\CJKunderwave{大雅}雲“殷商之旅”、“諮汝殷商”，是兼稱之也。亳是殷地大名，故殷社謂之亳社，其亳\CJKunderline{鄭玄}以為偃師，皇甫謐以為梁國谷熟縣，或雲濟陰亳縣。說既不同，未知誰是。 \par}

{\noindent\zhuan\zihao{6}\fzbyks 傳“適之”至“邑居”。正義曰：\CJKunderwave{釋詁}云：“適、之，往也。”俱訓為往，故“適”得為之，不欲往彼殷地，別有新邑居也。 \par}

{\noindent\zhuan\zihao{6}\fzbyks 傳“籲和”至“之言”。正義曰：“籲”即裕也，是寬裕，故為和也。憂則不和，“戚”訓憂也，故“率和眾憂之人，出正直之言”。\CJKunderwave{詩}雲“其直如矢”，故以“矢言”為“正直之言”。 \par}

{\noindent\zhuan\zihao{6}\fzbyks 傳“我王”至“於此”。正義曰：孔以祖乙圮於相地，遷都於耿，今盤庚自耿遷於殷，以“我王”為祖乙，此謂耿也。 \par}

{\noindent\zhuan\zihao{6}\fzbyks 傳“劉殺”至“殺故”。正義曰：“劉，殺”，\CJKunderwave{釋詁}文。水泉咸鹵，不可行化，王化不行，殺民之道。先王所以去彼遷此者，重我民，無慾盡殺故也。 \par}

不能胥匡以生，卜稽曰：‘其如臺。’\footnote{言民不能相匡以生，則當卜稽於龜以徙，曰：“其如我所行。”稽，工兮反。臺音怡。}先王有服，恪謹天命，茲猶不常寧，\footnote{先王有所服行，敬謹天命，如此尚不常安，有可遷輒遷。恪，苦各反。}不常厥邑，於今五邦。\footnote{湯遷亳，仲丁遷囂，河亶甲居相，祖乙居耿，我往居亳，凡五徙國都。馬云：“五邦謂商丘、亳、囂、相、耿也。”}今不承於古，罔知天之斷命，\footnote{今不承古而徙，是無知天將斷絕汝命。斷又音短。}矧曰其克從先王之烈？\footnote{天將絕命，尚無知之，況能從先王之業乎？從，才容反。}若顛木之有由櫱，\footnote{言今往遷都，更求昌盛，如顛仆之木，有用生櫱哉。櫱，五達反，本又作枿，馬云：“顛木而肄生曰枿。”僕音赴，又步北反。}天其永我命於茲新邑，\footnote{言天其長我命於此新邑，不可不徙。}紹復先王之大業，厎綏四方。\footnote{言我徙欲如此。厎之履反。}

{\noindent\zhuan\zihao{6}\fzbyks 傳“言民”至“所行”。正義曰：不徙所以不能相匡以生者，謂水泉沉溺，人民困苦，不能以義相匡正以生。又考卜於龜以徙，\CJKunderwave{周禮·大卜}：“大遷考貞龜。”是遷必卜也。 \par}

{\noindent\zhuan\zihao{6}\fzbyks 傳“先王”至“輒遷”。正義曰：下雲“於今五邦”，自湯以來數之，則此言“先王”總謂成湯至祖乙也。“先王有所服行”，謂行有典法，言能敬順天命，即是“有所服行”也。盤庚言先王敬順天命，如此尚不常安，有可遷輒遷；況我不能敬順天命，不遷民必死矣，故不可不遷也。 \par}

{\noindent\zhuan\zihao{6}\fzbyks 傳“湯遷”至“國都”。正義曰：孔以盤庚意在必遷，故通數“我往居亳”為“五邦”。鄭、王皆云，湯自商徙亳，數商、亳、囂、相、耿為五。計湯既遷亳，始建王業，此言先王遷都，不得遠數居亳之前充此數也。 \par}

{\noindent\zhuan\zihao{6}\fzbyks 傳“言今”至“櫱哉”。正義曰：\CJKunderwave{釋詁}云：“枿，餘也。”李巡曰：“枿，槁木之餘也。”郭璞云：“晉衛之間曰枿。”是言木死顛仆，其根更生櫱哉。此都毀壞，若枯死之木，若棄去毀壞之邑，更得昌盛，猶顛仆枯死之木用生櫱哉。 \par}

{\noindent\shu\zihao{5}\fzkt “盤庚”至“四方”。正義曰：盤庚欲遷於亳之殷地，其民不欲適彼殷地別有邑居，莫不憂愁，相與怨上。盤庚率領和諧其眾憂之人，出正直之言以曉告曰：“我先王初居此者，從舊都來，於是宅於此地。所以遷於此者，為重我民，無慾盡殺故。先王以久居墊隘，不遷則死，見下民不能相匡正以生，故謀而來徙。以徙為善，未敢專決，又考卜於龜以徙。既獲吉兆，乃曰：‘其如我所行欲徙之吉。’先王成湯以來，凡有所服行，敬順天命，如此尚不常安，可徙則徙，不常其邑，於今五邦矣。今若不承於古，徙以避害，則是無知天將斷絕汝命矣。天將絕命，尚不能知，況曰其能從先王之基業乎？今我往遷都，更求昌盛，若顛仆之木，有用生櫱哉。人衰更求盛，猶木死生櫱哉。我今遷向新都，上天其必長我殷之王命於此新邑，繼復先王之大業，致行其道，以安四方之人。我徙欲如此耳，汝等何以不原徙乎？”前雲若不徙以避害，則天將絕汝命，謂絕臣民之命，明亦絕我殷王之命。復雲若遷往新都，天其長我殷之王命，明亦長臣民之命，互文也。 \par}

盤庚斆於民，由乃在位，以常舊服，正法度。\footnote{斆，教也。教人使用汝在位之命，用常故事，正其法度。斆，戶教反，下如字。度如字。}曰：“無或敢伏小人之攸箴。”\footnote{言無有故伏絕小人之所欲箴規上者。戒朝臣。箴，之林反，馬云：“諫也。”朝，直遙反。}


{\noindent\zhuan\zihao{6}\fzbyks 傳“斆教”至“朝臣”。正義曰：\CJKunderwave{文王世子}云：“小樂正斆幹，大胥贊之。籥師斆戈，籥師丞贊之。”彼並是教舞干戈，知“斆”為教也。小民等患水泉沉溺，欲箴規上而徙，汝臣下勿抑塞伏絕之。\CJKunderline{鄭玄}云：“奢侈之俗，小民咸苦之，欲言於王。今將屬民而詢焉，故敕以無伏之。” \par}

{\noindent\shu\zihao{5}\fzkt “盤庚”至“攸箴”。正義曰：前既略言遷意，今復並戒臣民。盤庚先教於民云：“汝等當用汝在位之命，用舊常故事，正其法度。”欲令民徙，從其臣言也。民從上命，即是常事法度也。又戒臣曰：“汝等無有敢伏絕小人之所欲箴規上者。” \par}

王命眾悉至於庭。\footnote{眾，群臣以下。}

{\noindent\zhuan\zihao{6}\fzbyks 傳“眾，群臣以下”。正義曰：\CJKunderwave{周禮}：“小司寇掌外朝之政，以致萬民而詢焉，一曰詢國危，二曰詢國遷，三曰詢立君。”是國將大遷，必詢及於萬民。故知眾悉至王庭是“群臣以下”，謂及下民也。民不欲徙，由臣不助王勸民，故以下多是責臣之辭。 \par}

王若曰：“格汝眾，予告汝訓，\footnote{告汝以法教。}汝猷黜乃心，無傲從康。\footnote{謀退汝違上之心，無傲慢，從心所安。傲，五報反。}古我先王，亦惟圖任舊人共政。\footnote{先王謀任久老成人共治其政。任，而鴆反。}

{\noindent\zhuan\zihao{6}\fzbyks 傳“先王”。正義曰：此篇所言“先王”，其文無指斥者，皆謂成湯已來諸賢王也。下言“神後”、“高後”者，指謂湯耳。下篇言“古我先王，適於山”者，乃謂遷都之王仲丁、祖乙之等也。此言“先王”謂先世賢王。此既言“先王”，下句“王播告之”、“王用丕欽”蒙上之“先”，不言“先”，省文也。 \par}

王播告之修，不匿厥指，\footnote{王佈告人以所修之政，不匿其指。播，波餓反。匿，女力反。}

{\noindent\zhuan\zihao{6}\fzbyks 傳“王布”至“其指”。正義曰：上句言先王用舊人共政，下雲“王播告之修”，當謂告臣耳。傳言“佈告人”者，以下雲“民用丕變”，是必告臣，亦又告民。 \par}

王用丕欽，罔有逸言，民用丕變。\footnote{王用大敬其政教，無有逸豫之言，民用大變從化。}今汝聒聒，起信險膚，予弗知乃所訟。\footnote{聒聒，無知之貌。起信險偽膚受之言，我不知汝所訟言何謂。聒，古活反，馬雲\CJKunderwave{說文}皆云：“拒善自用之意。”}

{\noindent\zhuan\zihao{6}\fzbyks 傳“聒聒”至“何謂”。正義曰：\CJKunderline{鄭玄}云：“聒讀如‘聒耳’之聒，聒聒，難告之貌。”王肅云：“聒聒，善自用之意也。”此傳以“聒聒”為“無知之貌”，以“聒聒”是多言亂人之意也。“起信險膚”者，言發起所行，專信此險偽膚受淺近之言。信此浮言，妄有爭訟，我不知汝所訟言何謂。言無理也。 \par}

非予自荒茲德，惟汝含德，不惕予一人。予若觀火。\footnote{我之慾徙，非廢此德。汝不從我命，所含惡德，但不畏懼我耳。我視汝情如視火。惕，他歷反。}

{\noindent\shu\zihao{5}\fzkt “非予”至“觀火”。正義曰：言先王敬其教，民用大變。我命教汝，汝不肯徙。非我自廢此丕欽之德，惟汝之所含德甚惡，不畏懼我一人故耳。汝含藏此意，謂我不知。我見汝情若觀火。言見之分明如見火也。 \par}

予亦拙謀，作乃逸。\footnote{逸，過也。我不威脅汝徙，是我拙謀成汝過。拙，之劣反。}

{\noindent\zhuan\zihao{6}\fzbyks 傳“逸過”至“汝過”。正義曰：“逸，過”，\CJKunderwave{釋言}文。我若以威加汝，汝自不敢不遷，則無違上之過也。我不威脅汝徙，乃是我亦拙謀，作成汝過也。恨民以恩導之而不從己也。 \par}

若網在綱，有條而不紊。若農服田力穡，乃亦有秋。\footnote{紊，亂也。穡，耕稼也。下之順上，當如網在綱，各有條理而不亂也。農勤穡則有秋，下承上則有福。紊音問，徐音文。}

{\noindent\zhuan\zihao{6}\fzbyks 傳“紊亂”至“有福”。正義曰：“紊”是絲亂，故為亂也。“稼”、“穡”相對，則種之曰“稼”，斂之曰“穡”。“穡”是秋收之名，得為耕穫總稱，故云“穡，耕稼”。“下承上則有福”，“福”謂祿賞。 \par}

汝克黜乃心，施實德於民，至於婚友，丕乃敢大言，汝有積德。\footnote{汝群臣能退去傲上之心，施實德於民，至於婚姻僚友，則我大乃敢言汝有積德之臣。}乃不畏戎毒於遠邇，惰農自安，不昬作勞，不服田畝，越其罔有黍稷。\footnote{戎，大。昬，強。越，於也。言不欲徙，則是不畏大毒於遠近。如怠惰之農，苟自安逸，不強作勞于田畝，則黍稷無所有。昬，馬同；本或作暋，音敏。\CJKunderwave{爾雅}昬、暋皆訓強，故兩存。越，本又作粵，音曰，於也。強，其丈反。}

{\noindent\zhuan\zihao{6}\fzbyks 傳“戎大”至“所有”。正義曰：“戎，大”、“昬，強”、“越，於”皆\CJKunderwave{釋詁}文。孫炎曰：“昬，夙夜之強也。\CJKunderwave{書}曰：‘不昬作勞。’”引此解彼，是亦讀此為昬也。\CJKunderline{鄭玄}讀“昬”為暋，訓為勉也，與孔不同。傳雲“言不欲徙，則是不畏大毒於遠近”，其意言不徙則有毒，“毒”為禍患也；“遠近”謂賒促，言害至有早晚也。不強於作勞，則黍稷無所獲，以喻不遷於新邑，則福祿無所有也。此經惰農弗昬無黍稷，對上“服田力穡，乃亦有秋”，但其文有詳略耳。 \par}

“汝不和吉言於百姓，惟汝自生毒，\footnote{責公卿不能和喻百官，是自生毒害。}

{\noindent\zhuan\zihao{6}\fzbyks 傳“責公”至“毒害”。正義曰：此篇上下皆言“民”，此獨雲“百姓”，則知百姓是百官也。百姓既是百官，和吉言者又在百官之上，知此經是責公卿不能和喻善言於百官，使之樂遷也。不和百官，必將遇禍，是公卿自生毒害。 \par}

乃敗禍奸宄，以自災於厥身。\footnote{言汝不相率共徙，是為敗禍奸宄以自災之道。宄音軌。}乃既先惡於民，乃奉其恫，汝悔身何及？\footnote{群臣不欲徙，是先惡於民。恫，痛也。不徙則禍毒在汝身，徙奉持所痛而悔之，則於身無所及。奉，孚勇反，注同。恫，敕動反，又音通，痛也。}

{\noindent\zhuan\zihao{6}\fzbyks 傳“群臣”至“所及”。正義曰：群臣是民之師長，當倡民為善，群臣亦不欲徙，是乃先惡於民也。“恫，痛”，\CJKunderwave{釋言}文。 \par}

相時憸民，猶胥顧於箴言，其發有逸口，矧予制乃短長之命？\footnote{言憸利小民，尚相顧於箴誨，恐其發動有過口之患，況我制汝死生之命，而汝不相教從我，是不若小民。相時，相，息亮反，馬云：“視也。”徐息羊反。憸，息廉反，馬云：“憸利，小小見事之人也。”徐七漸反。}汝曷弗告朕，而胥動以浮言，恐沈於眾？\footnote{曷，何也。責其不以情告上，而相恐欲以浮言，不徙，恐汝沉溺於眾，有禍害。曷，何末反。}若火之燎於原，不可鄉邇，其猶可撲滅。\footnote{火炎不可向近，尚可撲滅。浮言不可信用，尚可得遏絕之。燎，力召反，又力鳥反，又力紹反。向，許亮反。撲，普卜反。近，附近之近。}則惟汝眾自作弗靖，非予有咎。\footnote{我刑戮汝，非我咎也。靖，謀也。是汝自為非謀所致。}

{\noindent\zhuan\zihao{6}\fzbyks 傳“曷何”至“禍害”。正義曰：“曷”、“何”同音，故“曷”為何也。顧氏云：“汝以浮雲恐動不徙，更是無益。我恐汝自取沉溺於眾人，不免禍害也。” \par}

{\noindent\zhuan\zihao{6}\fzbyks 傳“我刑”至“所致”。正義曰：我刑戮汝，汝自招之，非我咎也。“靖，謀”，\CJKunderwave{釋詁}文。告民不徙者，非善謀也。由此而被刑戮，是汝自為非謀所致也。 \par}

{\noindent\shu\zihao{5}\fzkt “相時”至“有咎”。正義曰：又責大臣不相教遷徙，是不如小民。我視彼憸利小民，猶尚相顧於箴規之言，恐其發舉有過口之患，故以言相規。患之小者尚知畏避，況我為天子制汝短長之命？威恩甚大，汝不相教從我，乃是汝不如小民。汝若不欲徙，何不以情告我，而輒相恐動以浮華之言？乃語民云：“國不可徙，我恐汝自取沉溺於眾人，而身被刑戮之禍害。”此浮言流行，若似火之燎於原野，炎熾不可向近，其猶可撲之使滅，以喻浮言不可止息，尚可刑戮使絕也。若以刑戮加汝，則是汝眾自為非謀所致此耳，非我有咎過也。 \par}

“遟任有言曰：‘人惟求舊，器非求舊，惟新。’\footnote{遟任，古賢。言人貴舊，器貴新，汝不徙，是不貴舊。遟，直疑反，徐持夷反。任，而今反，馬云：“古老成人。”}古我先王，暨乃祖乃父，胥及逸勤，予敢動用非罰？\footnote{言古之君臣相與同勞逸，子孫所宜法之，我豈敢動用非常之罰脅汝乎？}世選爾勞，予不掩爾善。\footnote{選，數也。言我世世數汝功勤，不掩蔽汝善，是我忠於汝。選，息轉反，又蘇管反。掩，本又作弇。數，色主反。}茲予大享於先王，爾祖其從與享之。\footnote{古者天子錄功臣配食於廟。大享，烝嘗也。所以不掩汝善。與音預。烝，之丞反。}作福作災，予亦不敢動用非德。\footnote{善自作福，惡自作災，我不敢動用非罰加汝，非德賞汝乎？從汝善惡而報之。}

{\noindent\zhuan\zihao{6}\fzbyks 傳“遟任”至“貴舊”。正義曰：其人既沒，其言立於後世，知是古賢人也。\CJKunderline{鄭玄}云：“古之賢史。”王肅云：“古老成人。”皆謂賢也。 \par}

{\noindent\zhuan\zihao{6}\fzbyks 傳“選數”至“於汝”。正義曰：\CJKunderwave{釋詁}云：“算，數也。”舍人曰：“釋數之曰算。”“選”即算也，故訓為數。經言世世數汝功勞，是從先王至己常行此事，故云“是我忠於汝”也。言己之忠，責臣之不忠也。 \par}

{\noindent\zhuan\zihao{6}\fzbyks 傳“古者”至“汝善”。正義曰：\CJKunderwave{周禮·大宗伯}祭祀之名，天神曰“祀”，地祇曰“祭”，人鬼曰“享”。此“大享於先王”，謂天子祭宗廟也。傳解天子祭廟，得有臣祖與享之意，言“古者天子錄功臣配食於廟”，故臣之先祖得與享之也。“古者”\CJKunderline{孔氏}據已而道前世也，此殷時已然矣。“大享，烝嘗”者，烝嘗是秋冬祭名，謂之“大享”者，以事各有對。若烝嘗對禘祫，則禘祫為大，烝嘗為小。若四時自相對，則烝嘗為大，礿祠為小。以秋冬物成，可薦者眾，故烝學為大；春夏物未成，可薦者少，故禘祫為小也。知烝嘗有功臣與祭者，案\CJKunderwave{周禮·司勳}雲“凡有功者，銘書於王之太常，祭於大烝，司勳詔之”是也。嘗是烝之類，而傳以嘗配之，\CJKunderwave{魯頌}曰“秋而載嘗”是也。\CJKunderwave{祭統}雲“內祭則大嘗禘是也，外祭則郊社是也”。然彼以祫為大嘗，知此不以烝嘗時為禘祫，而直據時祭者，以殷祫於三時，非獨烝嘗也。秋冬之祭，尚及功臣，則禘祫可知。惟春夏不可耳，以物末成故也。近代已來，惟禘祫乃祭功臣配食，時祭不及之也。近代已來，功臣配食各配所事之君，若所事之君其廟已毀，時祭不祭毀廟，其君尚不時祭，其臣固當止矣。禘祫則毀廟之主亦在焉，其時功臣亦當在也。\CJKunderwave{王制}云：“犆礿，祫禘，祫嘗，祫烝，諸侯礿犆，禘，一犆一祫，嘗祫，烝祫。”此\CJKunderwave{王制}之文，夏殷之制，天子春惟時祭，其夏秋冬既為祫，又為時祭。諸侯亦春為時祭，夏惟作祫，不作祭，秋冬先作時祭，而後祫。周則春曰祠，夏曰礿，三年一祫在秋，五年一禘在夏，故\CJKunderwave{公羊傳}云：“五年再殷祭。”\CJKunderwave{禮緯}云：“三年一祫，五年一禘。”此是鄭氏之義，未知孔意如何。 \par}

{\noindent\shu\zihao{5}\fzkt “遟任”至“非德”。正義曰：可遷則遷，是先王舊法。古之賢人遟任有言曰：“人惟求舊，器非求舊，惟新。”言人貴舊，器貴新，汝不欲徙，是不貴舊，反遟任也。古者我之先王及汝祖汝父相與同逸豫，同勤勞，汝為人子孫，宜法汝父祖，當與我同其勞逸。我豈敢動用非常之罰脅汝乎？自我先王以至於我，世世數汝功勞，我不掩蔽汝善，是我忠於汝也。以此故我大享祭於先王，汝祖其從我先王與在宗廟而歆享之，是我不掩汝善也。汝有善自作福，汝有惡自作災，我亦不敢動用非德之賞妄賞汝，各從汝善惡而報之耳。其意告臣言從上必有賞，違命必有罰也。 \par}

予告汝於難，若射之有志。\footnote{告汝行事之難，當如射之有所準志，必中所志乃善。射，食夜反。準音準。中，丁仲反。}

{\noindent\zhuan\zihao{6}\fzbyks 傳“告汝”至“乃善”。正義曰：此傳惟順經文，不言喻意。\CJKunderline{鄭玄}云：“我告汝，於我心至難矣。夫射者，張弓屬矢而志在所射，必中然後發之。為政之道亦如是也，以己心度之，可施於彼，然後出之。” \par}

{\noindent\shu\zihao{5}\fzkt “予告”至“有志”。正義曰：既言作福作災由人行有善惡，故復教臣行善：“我告汝於行事之難，猶如射之有所準志。志之所主，欲得中也，必中所志，乃為善耳。”以喻人將有行，豫思念之，行得其道為善耳。其意言遷都是善道，當念從我言也。 \par}

汝無侮老成人，無弱孤有幼。\footnote{不用老成人之言，是侮老之。不徙則孤幼受害，是弱易之。侮，亡甫反。易，以豉反。}

{\noindent\zhuan\zihao{6}\fzbyks 傳“不用”至“易之”。正義曰：“老”謂見其年老，謂其無所復知。“弱”謂見其幼弱，謂其未有所識。鄭云：“老弱皆輕忽之意也。”老成人之言云可徙，不用其言，是侮老之也。不徙則水泉咸鹵，孤幼受害，不念其害，則是卑弱輕易之也。 \par}

各長於厥居，勉出乃力，聽予一人之作猷。\footnote{盤庚群臣下各思長於其居，勉盡心出力，聽從遷徙之謀。長，丁丈反。}

{\noindent\zhuan\zihao{6}\fzbyks 傳“盤庚”至“之謀”。正義曰：於時群臣難毀其居宅，惟見目前之利，不思長久之計。其臣非一，共為此心。盤庚群臣下各思長久於其居處，勉強盡心出力，聽從我遷徙之謀。自此已下皆是也。 \par}

無有遠邇，用罪伐厥死，用德彰厥善。\footnote{言遠近待之如一，罪以懲之，使勿犯，伐去其死道。德以明之，使勸慕，競為善。去，起呂反。}

{\noindent\shu\zihao{5}\fzkt “無有”至“厥善”。正義曰：此即遷徙之謀也。言我至新都，撫養在下，無有遠之與近，必當待之如一。用刑殺之罪伐去其死道，用照察之德彰明其行善。有過，罪以懲之，使民不犯非法。死刑不用，是“伐去其死道”。“伐”若伐樹然，言止而不復行用也。有善者，人主以照察之德加賞祿以明之，使競慕為善，是彰其善也。此二句相對，上言“用罪伐厥死”，下宜言“用賞彰厥生”，不然者，上言用刑，下言賞善，死是刑之重者，舉重故言“死”；有善乃可賞，故言“彰厥善”；行賞是德，故以“德”言賞；人生是常，無善亦生，不得言“彰厥生”，故文互。 \par}

邦之臧，惟汝眾。\footnote{有善則群臣之功。}\footnote{臧，徐子郎反。}邦之不臧，惟予一人有佚罰。\footnote{佚，失也。是己失政之罰。罪己之義。佚音逸。}凡爾眾，其惟致告：\footnote{致我誠，告汝眾。}自今至於後日，各恭爾事，齊乃位，度乃口。\footnote{奉其職事，正齊其位，以法度居汝口，勿浮言。度，徐如字，亦作渡。}罰及爾身，弗可悔。”\footnote{不從我謀，罰及汝身，雖悔可及乎？}

{\noindent\shu\zihao{5}\fzkt “度乃口”。正義曰：“度”，法度也，故傳言“以法度居汝口”也。 \par}

\section{盤庚中第十（盤庚）}

盤庚作，惟涉河以民遷。\footnote{為此南渡河之法，用民徙。}乃話民之弗率，誕告用亶其有眾。\footnote{話，善言。民不循教，發善言大告用誠於眾。話，胡快反，馬云：“告也，言也。”誕，徐音但。亶，丁但反，馬本作單，音同，誠也。}咸造勿褻在王庭，\footnote{造，至也。眾皆至王庭，無褻慢。造，七報反，注同；馬在早反，云：“為也。”褻，息列反。}盤庚乃登進厥民。\footnote{升進，命使前。}

{\noindent\zhuan\zihao{6}\fzbyks 傳“為此”至“民徙”。正義曰：\CJKunderline{鄭玄}雲“作渡河之具”，王肅雲“為此思南渡河之事”，此傳言“南渡河之法”，皆謂造舟舡渡河之具，是濟水先後之次，思其事而為之法也。 \par}

{\noindent\zhuan\zihao{6}\fzbyks 傳“話善”至“於眾”。正義曰：\CJKunderwave{釋詁}云：“話，言也。”孫炎曰：“話，善人之言也。”王苦民不從教，必發善言告之，故以“話”為善言。\CJKunderline{鄭玄}\CJKunderwave{詩箋}亦云：“話，善言也。”曰：“明聽朕言，無荒失朕命。荒，廢。 \par}

{\noindent\shu\zihao{5}\fzkt “盤庚”至“厥民”。正義曰：盤庚於時見都河北，欲遷向河南，作惟南渡河之法，欲用民徙，乃出善言以告曉民之不循教者，大為教告，用誠心於其所有之眾人。於時眾人皆至，無有褻慢之人，盡在於王庭。盤庚乃升進其民，延之使前而教告之。史敘其事，以為盤庚發誥之目。 \par}

嗚呼！古我前後，罔不惟民之承。\footnote{言我先世賢君，無不承安民而恤之。}保後胥戚，鮮以不浮於天時。\footnote{民亦安君之政，相與憂行君令。浮，行也。少以不行於天時者，言皆行天時。鮮，息淺反。}

{\noindent\zhuan\zihao{6}\fzbyks 傳“民亦”至“天時”。正義曰：以君承安民而憂之，故民亦安君之政，相與憂行君令，使君令必行。責時群臣不憂行君令也。舟舡浮水而行，故以“浮”為行也。行天時也，順時佈政，若\CJKunderwave{月令}之為也。 \par}

殷降大虐，先王不懷。\footnote{我殷家於天降大災，則先王不思故居而行徙。}

{\noindent\zhuan\zihao{6}\fzbyks 傳“我殷”至“行徙”。正義曰：遷徙者止為邑居墊隘，水泉咸鹵，非為避天災也。此傳以“虐”為災，“懷”為思，言“殷家於天降大災，則先王不思故居而行徙”者，以天時人事終是相將，邑居不可行化，必將天降之災。上雲“不能相匡以生”、“罔知天之斷命”，即是天降災也。 \par}

厥攸作視民利，用遷。\footnote{其所為視民有利，則用徙。}汝曷弗念我古後之聞？\footnote{古後先王之聞，謂遷事。曷，何末反，下同。}承汝俾汝，惟喜康共，非汝有咎，比於罰。\footnote{今我法先王惟民之承，故承汝使汝徙，惟與汝共喜安，非謂汝有惡徙汝，令比近於殃罰。俾，必爾反。咎，其九反。比，毗志反，徐扶志反，注及下同。共，群用反。令，力呈反。近，附近之近。}

{\noindent\shu\zihao{5}\fzkt “承汝”至“於罰”。正義曰：先王為政，惟民之承。今我亦法先王，故承安汝使汝徙。惟歡喜安樂皆與汝共之，非謂汝有咎惡而徙汝，令比近於殃罰也。 \par}

予若籥懷茲新邑，亦惟汝故，以丕從厥志。\footnote{言我順和懷此新邑，欲利汝眾，故大從其志而徙之。籲，羊戍反。}

{\noindent\shu\zihao{5}\fzkt “予若”至“厥志”。正義曰：盤庚言，我順於道理，和協汝眾，歸懷此新邑者，非直為我王家，亦惟利汝眾，故為此大從我本志而遷徙，不有疑也。 \par}

“今予將試以汝遷，安定厥邦。\footnote{試，用。}汝不憂朕心之攸困，\footnote{所困，不順上命。}乃咸大不宣乃心，欽念以忱，動予一人。\footnote{汝皆大不布腹心，敬念以誠感我，是汝不盡忠。忱，市林反。}爾惟自鞠自苦，\footnote{鞠，窮也。言汝為臣不忠，自取窮苦。鞠，居六反。}若乘舟，汝弗濟，臭厥載。\footnote{言不徙之害，如舟在水中流不渡，臭敗其所載物。臭，徐尺售反。載如字，又在代反。}

{\noindent\shu\zihao{5}\fzkt “臭厥載”。正義曰：“臭”是氣之別名，古者香氣穢氣皆名為“臭”。\CJKunderwave{易}雲“其臭如蘭”，謂香氣為“臭”也。\CJKunderwave{晉語}雲“惠公改葬申生，臭徹於外”，謂穢氣為“臭”也。下文覆述此意雲“無起穢以自臭”，則此“臭”謂穢氣也。肉敗則臭，故以“臭”為敗。船不渡水，則敗其所載物也。 \par}

爾忱不屬，惟胥以沈。不其或稽，自怒曷瘳？\footnote{汝忠誠不屬逮古，苟不欲徙，相與沉溺，不考之先王，禍至自怒，何瘳差乎？屬音燭，注同，馬云：“獨也。”沈，直林反。瘳，敕留反。}

{\noindent\shu\zihao{5}\fzkt “爾忱”至“曷瘳”。正義曰：盤庚責其臣民，汝等不用徙者，由汝忠誠不能屬逮於古賢。苟不欲徙，惟相與沉溺於眾。不欲徙之言，不其有考驗於先王遷徙之事。汝既不考於古，及其禍至，乃自忿怒，何所瘳差也？ \par}

汝不謀長，以思乃災，汝誕勸憂。\footnote{汝不謀長久之計，思汝不徙之災，苟不欲徙，是大勸憂之道。}

{\noindent\shu\zihao{5}\fzkt “汝誕勸憂”。正義曰：凡人以善自勸，則善事多。若以憂自勸，則憂來眾。今不徙則憂來眾，是自勸勵以憂愁之道。 \par}

今其有今罔後，汝何生在上？\footnote{言不徙無後計，汝何得久生在人上，禍將及汝。}

{\noindent\shu\zihao{5}\fzkt “今其”至“在上”。正義曰：顧氏云：“責群臣：汝今日其且有今目前之小利，無後日久長之計，患禍將至，汝何得久生在民上也？” \par}

今予命汝一，無起穢以自臭。\footnote{我一心命汝，汝違我是自臭敗。穢，於廢反。}

{\noindent\shu\zihao{5}\fzkt “今予”至“自臭”。正義曰：今我命汝，是我之一心也。汝當從我，無得起為穢惡，以自臭敗。汝違我命，是起穢以自臭也。 \par}

恐人倚乃身，迂乃心。\footnote{言汝既不欲徙，又為他人所誤。倚，曲。迂，僻。倚，於綺反，徐於奇反。迂音於。僻，匹亦反。}

{\noindent\zhuan\zihao{6}\fzbyks 傳“言汝”至“迂僻”。正義曰：人心不能自決，則好用非理之謀。言汝既不欲遷徙，又為他人所誤。盤庚疑其被誤，故言此也。以物倚物者必曲，故“倚”為曲也。“迂”是回也，回行必僻，故“迂”為僻也。 \par}

{\noindent\shu\zihao{5}\fzkt “恐人”至“乃心”。正義曰：言汝心既不欲徙，旁人或更誤汝。我有恐他人倚曲汝身，迂僻汝心，使汝益不用徙也。 \par}

予迓續乃命於天，予豈汝威？用奉畜汝眾。\footnote{迓，迎也。言我徙，欲迎續汝命於天，豈以威脅汝乎？用奉畜養汝眾。迓，五駕反。畜，許竹反，下同。脅，虛業反。}

{\noindent\zhuan\zihao{6}\fzbyks 傳“迓迎”至“汝眾”。正義曰：“迓，迎”，\CJKunderwave{釋詁}文。不遷必將死矣，天欲遷以延命。天意向汝，我欲迎之。天斷汝命，我欲續之。我今徙者，欲迎續汝命於天，豈以威脅汝乎？遷都惟用奉養汝群臣民耳。 \par}

“予念我先神後之勞爾先，予丕克羞爾，用懷爾然。\footnote{言我亦法湯大能進勞汝，以義懷汝心，而汝違我，是汝反先人。勞，力報反，又如字，注同。}

{\noindent\zhuan\zihao{6}\fzbyks 傳“言我”至“先人”。正義曰：\CJKunderwave{易}稱：“神者，妙萬物而為言也。”殷之先世，神明之君惟有湯耳，故知“神後”謂湯也。下“高後”、“先後”與此“神後”一也。“神”者，言其通聖。“高”者，言其德尊。此“神後”言“先”，於“高後”略而不言“先”，其下直言“先後”，又略而不言“高”，從上省文也。“勞爾先”，謂愛之也。“勞”者，勤也，閔其勤勞而慰勞之，“勞”亦愛之義，故\CJKunderwave{論語}云：“愛之，能勿勞乎？”是“勞”為愛也。此言湯勞汝先，則此所責之臣，其祖於成湯之世已在朝廷。世仕王朝而不用己命，故責之深也。 \par}

{\noindent\shu\zihao{5}\fzkt “予念”至“爾然”。正義曰：我念我先世神後之君成湯，愛勞汝之先人，故我大能進用汝，與汝爵位，用以道義懷安汝心耳。然汝乃違我命，是汝反先人也。 \par}

失於政，陳於茲，高後丕乃崇降罪疾，曰：‘曷虐朕民？’\footnote{崇，重也。今既失政，而陳久於此而不徙，湯必大重下罪疾於我，曰：“何為虐我民而不徙乎？”重，直勇反，又直恭反。}汝萬民乃不生生，暨予一人猷同心，\footnote{不進進謀同心徙。}先後丕降與汝罪疾，曰：‘曷不暨朕幼孫有比？’\footnote{言非但罪我，亦將罪汝。幼孫，盤庚自謂。比，同心。}故有爽德，自上其罰汝，汝罔能迪。\footnote{湯有明德在天，見汝情，下罰汝，汝無能道。言無辭。}

{\noindent\zhuan\zihao{6}\fzbyks 傳“崇重”至“徙乎”。正義曰：“崇，重”，\CJKunderwave{釋詁}文。又云：“塵，久也。”孫炎曰：“陳居之久，久則生塵矣。”古者“陳”、“塵”同也，故“陳”為久之義。 \par}

{\noindent\zhuan\zihao{6}\fzbyks 傳“不進”至“心徙”。正義曰：物之生長，則必漸進，故以“生生”為進進。王肅亦然。“進進”是同心原樂之意也。此實責群臣而言“汝萬民”者，民心亦然，因博及之。 \par}

{\noindent\zhuan\zihao{6}\fzbyks 傳“湯有”至“無辭”。正義曰：訓“爽”為明，言其見下，故稱“明德”。\CJKunderwave{詩}稱“三後在天”，死者精神在天，故言下見汝。 \par}

{\noindent\shu\zihao{5}\fzkt “失於”至“能迪”。正義曰：盤庚以民不願遷，言神將罪汝，欲懼之使從己也。我所以必須徙者，我今失於政教，陳久於此，民將有害，高德之君成湯必忿我不徙，大乃重下罪疾於我，曰：“何為殘虐我民而不徙乎？”我既欲徙，而汝與萬民，乃不進進與我一人謀計同心，則我先君成湯大下與汝罪疾，曰：“何故不與我幼孫盤庚有相親比同心徙乎？”汝不與我同心，故湯有明德，從上見汝之情，其下罪罰於汝。汝實有罪，無所能道。言無辭以有解說也。 \par}

古我先後，既勞乃祖乃父，\footnote{勞之共治人。}汝共作我畜民。汝有戕，則在乃心。\footnote{戕，殘也。汝共我治民，有殘人之心而不欲徙，是反父祖之行。戕，在良反，又士良反。行，下孟反。}我先後綏乃祖乃父，乃祖乃父乃斷棄汝，不救乃死。\footnote{言我先王安汝父祖之忠，今汝不忠汝父祖，必斷絕棄汝命，不救汝死。斷，丁緩反。}

{\noindent\zhuan\zihao{6}\fzbyks 傳“勞之共治人”。正義曰：下句責臣之身雲“汝共作我畜民”，明先後勞其祖父，是勞之共治民也。 \par}

{\noindent\zhuan\zihao{6}\fzbyks 傳“戕殘”至“之行”。正義曰：\CJKunderwave{春秋}宣十八年“邾人戕鄫子”，\CJKunderwave{左傳}云：“凡自虐其君曰弒，自外曰戕。”“戕”為殘害之義，故為殘也。先後愛勞汝祖汝父，與共治民，汝祖父必有愛人之心。“作”訓為也。汝今共為我養民之官，而有殘民之心，而不用徙以避害，是汝反祖父之行。盤庚距湯，年世多矣，臣父不及湯世而云“父”者，與“祖”連言之耳。 \par}

{\noindent\shu\zihao{5}\fzkt “古我”至“乃死”。正義曰：又責群臣：“古我先君成湯，既愛勞汝祖汝父，與之共治民矣。汝今共為我養民之官，是我於汝與先君同也。而汝有殘虐民之心，非我令汝如此，則在汝心自為此惡，是汝反祖父之行。雖汝祖父，亦不祐汝。我先君安汝祖汝父之忠，汝祖汝父忠於先君，必忿汝違我，乃斷絕棄汝命，不救汝死。”言汝違我命，故汝祖父亦忿見湯罪汝，不救汝死也。 \par}

茲予有亂政同位，具乃貝玉。\footnote{亂，治也。此我有治政之臣，同位於父祖，不念盡忠，但念貝玉而已。言其貪。治，直吏反。盡，子忍反。}乃祖先父丕乃告我高後曰：‘作丕刑于朕孫。’\footnote{言汝父祖見汝貪而不忠，必大乃告湯曰：“作大刑于我子孫，求討不忠之罪。”告，工號反。我高後，本又作“乃祖乃父”。}迪高後，丕乃崇降弗祥。\footnote{言汝父祖開道湯，大重下不善以罰汝。陳忠孝之義以督之。}


{\noindent\zhuan\zihao{6}\fzbyks 傳“言汝”至“之罪”。正義曰：上句言成湯罪此諸臣，其祖父不救子孫之死，此句言臣之祖父請成湯討其子孫，以不從已，故責之益深。先祖請討，非盤庚所知，原神之意而為之辭，以懼其子孫耳。 \par}

{\noindent\zhuan\zihao{6}\fzbyks 傳“言汝”至“督之”。正義曰：訓“迪”為道，言汝父祖開道湯也。不從君為不忠，違父祖為不孝，父祖開道湯下罰，欲使從君順祖，陳忠孝之義以督勵之。 \par}

{\noindent\shu\zihao{5}\fzkt “茲予”至“弗祥”。正義曰：又責臣云：“汝祖父非徒不救汝死，乃更請與汝罪。於此我有治政之臣，同位於其父祖。其位與父祖同，心與父祖異。不念忠誠，但念具汝貝玉而已。”言其貪而不忠也。“汝先祖先父以汝如此，大乃告我高後曰：‘為大刑于我子孫。’以此言開道我高後，故我高後大乃下不善之殃以罰汝。成湯與汝祖父皆欲罪汝，汝何以不從我徙乎？”亂治”至“其貪”。正義曰：“亂，治”，\CJKunderwave{釋詁}文。舍人曰：“亂，義之治也。”孫炎曰：“亂，治之理也。”大臣理國之政，此者所責之人，故言於此我有治政大臣。言其同位於父祖，責其位同而心異也。貝者，水蟲。古紉選其甲以為貨，如今之用錢然。\CJKunderwave{漢書·食貨志}具有其事。貝是行用之貨也，貝玉是物之最貴者，責其貪財，故舉二物以言之。當時之臣不念盡忠於君，但念具貝玉而已，言其貪也。 \par}

“嗚呼！今予告汝不易。\footnote{凡所言皆不易之事。易，以豉反，注同。}永敬大恤，無胥絕遠。\footnote{長敬我言，大憂行之，無相與絕遠棄廢之。遠，於萬反，又如字，注同。}汝分猷念以相從，各設中於乃心。\footnote{群臣當分明相與謀念，和以相從，各設中正於汝心。分，扶問反，又如字，注同。}乃有不吉不迪，\footnote{不善不道，謂兇人。}顛越不恭，暫遇奸宄，\footnote{顛，隕。越，墜也。不恭，不奉上命。暫遇人而劫奪之，為奸於外，為宄於內。暫，才淡反。隕，于敏反。}我乃劓殄滅之，無遺育，無俾易種於茲新邑。\footnote{劓，割。育，長也。言不吉之人當割絕滅之，無遺長其類，無使易種於此新邑。劓，魚器反，徐吾氣反。殄，徒典反。易如字，又以豉反，注同。長，丁丈反，下“遺長”同。}往哉生生！今予將試以汝遷，永建乃家。”\footnote{自今以往，進進於善。我乃以汝徙，長立汝家。卿大夫稱家。}

{\noindent\zhuan\zihao{6}\fzbyks 傳“不易之事”。正義曰：此“易”讀為“難易”之“易”，“不易”言其難也。王肅云：“告汝以命之不易。”亦以“不易”為難。\CJKunderline{鄭玄}云：“我所以告汝者不變易，言必行之。”謂盤庚自道己言必不改易，與孔異。 \par}

{\noindent\zhuan\zihao{6}\fzbyks 傳“顛隕”至“於內”。正義曰：\CJKunderwave{釋詁}云：“隕，落也。隕，墜也。”“顛越”是從上倒下之言，故以“顛”為隕，“越”是遺落，為墜也。\CJKunderwave{左傳}僖九年齊桓公云：“恐隕越於下。”文十八年史克云：“弗敢失墜。”“隕”、“越”是遺落廢失之意，故以隕墜不恭為“不奉上命”也。“暫遇人而劫奪之”謂逄人即劫，為之無已。成十七年\CJKunderwave{左傳}曰“亂在外為奸，在內為宄”，是劫奪之事，故以劫奪解其“奸宄”也。 \par}

{\noindent\zhuan\zihao{6}\fzbyks 傳“劓割”至“新邑”。正義曰：五刑截鼻為劓，故“劓”為割也。“育，長”，\CJKunderwave{釋詁}文。“不吉之人當割絕滅之，無遺長其類”，謂早殺其人，不使得子孫，有此惡類也。“易種”者，即今俗語云“相染易”也。惡種在善人之中，則善人亦變易為惡，故絕其類，無使易種於此新邑也。滅去惡種，乃是常法，而言“於此新邑”，言己若至新都，當整齊使絜清。 \par}

{\noindent\zhuan\zihao{6}\fzbyks 傳“自今”至“稱家”。正義曰：“長立汝家”謂賜之以族，使子孫不絕，\CJKunderwave{左傳}所謂“諸侯命氏”是也。王朝大夫，天子亦命之氏，故云“立汝家”也。 \par}

{\noindent\shu\zihao{5}\fzkt “嗚呼”至“乃家”。正義曰：盤庚以言事將畢，欲戒使入之，故“嗚呼”而嘆之。今我告汝皆不易之事，言其難也。事既不易，當長敬我言，大憂行之，無相絕遠棄廢之，必須存心奉行。汝群臣臣分輩相與計謀念，和協以相從，各設中正於汝心，勿為殘害之事。汝群臣若有不善不道，隕墜禮法，不恭上命，暫逄遇人，即為奸宄而劫奪之，我乃割絕滅之，無有遺餘生長。所以然者，欲無使易其種類於此新邑故耳。自今已往哉，汝當進進於善。今我將用以汝遷，長立汝家，使汝在位，傳諸子孫。勿得違我言也。 \par}

\section{盤庚下第十一（盤庚）}

盤庚既遷，奠厥攸居，乃正厥位，\footnote{定其所居，正郊廟朝社之位。奠，田薦反。朝，直遙反。}綏爰有眾，曰：“無戲怠，懋建大命。\footnote{安於有眾，戒無戲怠，勉立大教。}今予其敷心腹腎腸，歷告爾百姓於朕志。\footnote{布心腹，言輸誠於百官以告志。腎，時忍反。腸，徐待良反。}罔罪爾眾，爾無共怒，協比讒言予一人。\footnote{群臣前有此過，故禁其後。今我不罪汝，汝勿共怒我，合比兇人而妄言。比，毗志反。讒，仕咸反。}

{\noindent\zhuan\zihao{6}\fzbyks 傳“定其”至“之位”。正義曰：訓“攸”為所，“定其所居”，總謂都城之內官府萬民之居處也。\CJKunderline{鄭玄}云：“徙主於民，故先定其裡宅所處，次乃正宗廟朝廷之位。”如鄭之意，“奠厥攸居”者，止謂定民之居，豈先令民居使足，待其餘剩之處，然後建王宮乎？若留地以擬王宮，即是先定王居，不得為先定民矣。孔惟言“定其所居”，知是官民之居並定之也。\CJKunderwave{禮}郊在國外，左祖右社，面朝後市，“正厥位”謂正此郊廟朝社之位也。 \par}

{\noindent\zhuan\zihao{6}\fzbyks 傳“安於”至“大教”。正義曰：\CJKunderline{鄭玄}云：“勉立我大命，使心識教令，常行之。”王肅云：“勉立大教，建性命，致之五福。”又案下句爾然共怒予一人，是恐其不從已命，此句宜言我有教命，汝當勉力立之。鄭說如孔旨也。 \par}

{\noindent\zhuan\zihao{6}\fzbyks 傳“布心”至“告志”。正義曰：此論心所欲言，腹內之事耳。以心為五臟之主，腹為六腑之總，腸在腹內，腎在心下，舉“腎腸”以配“腹心”，\CJKunderwave{詩}曰：“公侯腹心。”宣十二年\CJKunderwave{左傳}云：“敢布腹心。”是“腹心”足以表內，“腎腸”配言之也。 \par}

{\noindent\shu\zihao{5}\fzkt “盤庚”至“一人”。正義曰：盤庚既遷至殷地，定其國都處所，乃正其郊廟朝社之位。又屬民而聚之，安慰於其所有之眾，曰：“汝等自今以後，無得遊戲怠惰，勉力立行教命。今我其布心腹腎腸，輸寫誠信，歷遍告汝百姓於我心志者。”欲遷之日，民臣共怒盤庚盤庚，恐其怖懼，故開解之。“今我無復罪汝眾人。我既不罪汝，汝無得如前共為忿怒，協比讒言毀惡我一人”。恕其前愆，與之更始也。 \par}

古我先王，將多於前功，\footnote{言以遷徙多大前人之功美。}適於山，用降我凶德，嘉績於朕邦。\footnote{徙必依山之險，無城郭之勞。下去兇惡之德，立善功於我國。降，工巷反。去，羌呂反。}今我民用蕩析離居，罔有定極。\footnote{水泉沉溺，故蕩析離居，無安定之極，徙以為之極。}

{\noindent\zhuan\zihao{6}\fzbyks 傳“言以”至“功美”。正義曰：“古我先王”，謂遷都者。“前人”謂未遷者。前人久居舊邑，民不能相匡以生，則是居無功矣。盤庚言先王以此遷，徙故多大前人之功美，故我今遷，亦欲多前功矣。 \par}

{\noindent\zhuan\zihao{6}\fzbyks 傳“徙必”至“我國”。正義曰：先王至此五邦，不能盡知其地，所都皆近山，故總稱“適於山”也。\CJKunderwave{易·坎卦}彖云：“王公設險以守其國。”徙必依山之險，欲使下民無城郭之勞。雖則近山，不可全無城郭，言其防守易耳。徙必近山，則舊處新居皆有山矣。而云“適於山”者，言其徙必依山，不適平地，不謂舊處無山，故徙就山也。水泉咸鹵，民居墊隘，時君不為之徙，即是兇惡之德。其徙者是下去兇惡之德，立善功於我新遷之國也。言“下”者，凶德在身，下而墜去之。 \par}

{\noindent\zhuan\zihao{6}\fzbyks 傳“水泉”至“之極”。正義曰：民居積世，穿掘處多，則水泉盈溢，令人沈深而陷溺。其處不可安居，播蕩分析，離其居宅，無安定之極。“極”訓中也。\CJKunderwave{詩}云：“立我烝民，莫匪爾極。”言民賴后稷之功，莫不得其中。今為民失中，故徙以為之中也。 \par}

{\noindent\shu\zihao{5}\fzkt “古我”至“定極”。正義曰：言古者我之先王，將欲多大於前人之功，是故徙都而適於山險之處，用下去我兇惡之德，立善功於我新國。但徙來已久，水泉沉溺，今我在此之民，用播蕩分析，離其居宅，無有安定之極，我今徙而使之得其中也。說其遷都之意，亦欲多大前人之功，定民極也。 \par}

爾謂朕：‘曷震動萬民以遷？’\footnote{言皆不明己本心。}肆上帝將復我高祖之德，亂越我家。\footnote{以徙故，天將復湯德，治理於我家。治，直吏反。}朕及篤敬，恭承民命，用永地於新邑。\footnote{言我當與厚敬之臣，奉承民命，用長居新邑。}肆予沖人，非廢厥謀，吊由靈。\footnote{衝，童。童人，謙也。吊，至。靈，善也。非廢，謂動謀於眾，至用其善。吊音的，或如字。}各非敢違卜，用宏茲賁。\footnote{宏、賁皆大也。君臣用謀，不敢違卜，用大此遷都大業。賁，扶雲反。}

{\noindent\zhuan\zihao{6}\fzbyks 傳“以徙”至“我家”。正義曰：民害不徙，違失湯德。以徙之故，天必祐我，將使復奉湯德，令得治理於我家。言由徙故天福之也。 \par}

{\noindent\zhuan\zihao{6}\fzbyks 傳“衝童”至“其善”。正義曰：“衝”、“童”聲相近，皆是幼小之名。自稱“童人”，言己幼小無知，故為“謙也”。“吊，至”、“靈，善”皆\CJKunderwave{釋詁}文。\CJKunderwave{禮}將有大事，必謀於眾。謀眾乃是常理，故言“非廢，謂動謀於眾”，言己不自專也。眾謀必有異見，故至極用其善者。 \par}

{\noindent\zhuan\zihao{6}\fzbyks 傳“宏賁”至“大業”。正義曰：“宏、賁皆大也”，\CJKunderwave{釋詁}文。樊光曰：“\CJKunderwave{周禮}雲‘其聲大而宏’，\CJKunderwave{詩}雲‘有賁其首’，是宏、賁皆為大之義也。”“各”者非一之辭，故為“君臣用謀，不敢違卜”。\CJKunderwave{洪範}云：“汝則有大疑，謀及卿士，謀及卜筮。”言“非敢違卜”，是既謀及於眾，又決於蓍龜也。“用大此遷都”，“大”謂立嘉績以大之也。 \par}

{\noindent\shu\zihao{5}\fzkt “爾謂”至“茲賁”。正義曰：言我徙以為民立中，汝等不明我心，乃謂我何故震動萬民以為此遷。我以此遷之故，上天將復我高祖成湯之德，治理於我家。我當與厚敬之臣，奉承民命，用是長居於此新邑。以此須遷之故，我童蒙之人，非敢廢其詢謀。謀於眾人，眾謀不同，至用其善者。言善謀者，皆欲遷都也。又決之於龜卜而得吉，我與汝群臣各非敢違卜，用是必遷，光大此遷都之大業。我徙本意如此耳。 \par}

“嗚呼！邦伯師長，百執事之人，尚皆隱哉！\footnote{國伯，二伯及州牧也。眾長，公卿也。言當庶幾相隱括共為善政。長，丁丈反，注同。}予其懋簡相爾，念敬我眾。\footnote{簡，大。相，助也。勉大助汝，念敬我眾民。相，息亮反。}朕不肩好貨，敢恭生生。鞠人謀人之保居，敘欽。\footnote{肩，任也。我不任貪貨之人，敢奉用進進於善者。人之窮困能謀安其居者，則我式序而敬之。好，呼報反。任，而林反。}

{\noindent\zhuan\zihao{6}\fzbyks 傳“國伯”至“善政”。正義曰：“邦伯”，邦國之伯，諸侯師長，故為東西二伯及九州之牧也。\CJKunderline{鄭玄}注\CJKunderwave{禮記}云：“殷之州長曰伯，虞夏及周皆曰牧。”此殷時而言“牧”者，此乃鄭之所約，孔意不然，故總稱“牧”也。“師”訓為眾，“眾長”，眾官之長，故為三公六卿也。其“百執事”，謂大夫以下，諸有職事之官皆是也。此部敕眾臣，故二伯已下及執事之人皆戒之也。\CJKunderwave{釋言}云：“庶幾，尚也。”反覆相訓，故“尚”為庶幾。“庶”，幸也。“幾”，冀也。“隱”謂隱審也。幸冀相與隱審檢括，共為善政，欲其同心共為善也。“隱括”必是舊語，不知本出何書。何休\CJKunderwave{公羊序}云：“隱括使就繩墨焉。” \par}

{\noindent\zhuan\zihao{6}\fzbyks 傳“簡大”至“眾民”。正義曰：“簡，大”，\CJKunderwave{釋詁}文。又云：“相、助，慮也。”俱訓為“慮”，是“相”得為助也。盤庚欲使群臣同心為善，欲勉力大佐助之，使皆念敬我眾民也。 \par}

{\noindent\zhuan\zihao{6}\fzbyks 傳“肩任”至“敬之”。正義曰：\CJKunderwave{釋詁}云：“肩，勝也。”舍人曰：“肩，強之勝也。”強能勝重，是堪任之義，故為任也。我今不委任貪貨之人。以“恭”為奉。人有向善而心不決志，故美其人能果敢奉用進進於善者，言其人好善不倦也。“鞠”訓為窮，“鞠人”謂窮困之人。“謀人之保居”，謂謀此窮人之安居，若見人之窮困，能謀安其居。愛人而樂安存之者，則我式序而敬之。\CJKunderwave{詩}云：“式序在位。”言其用次序在官位也。鄭、王皆以“鞠”為養，言“能謀養人安其居者，我則次序而敬之”，與孔不同。 \par}

{\noindent\shu\zihao{5}\fzkt “嗚呼”至“敘欽”。正義曰：言遷事已訖，故嘆而敕之：“嗚呼！國之長伯，及眾官之長與百執事之人，庶幾皆相與隱括共為善政哉！我其勉力大助汝等為善，汝當思念愛敬我之眾民。我不任用好貨之人，有人果敢奉用進進於善，見窮困之人能謀此窮困之人安居者，我乃次序而敬用之。” \par}

今我既羞告爾於朕志，若否，罔有弗欽。\footnote{已進告汝之後，順於汝心與否，當以情告我，無敢有不敬。告，故報反。}無總於貨寶，生生自庸。\footnote{無總貨寶以求位，當進進皆自用功德。}式敷民德，永肩一心。”\footnote{用布示民，必以德義，長任一心以事君。}

{\noindent\shu\zihao{5}\fzkt “今我”至“一心”。正義曰：今我既進而告汝於我心志矣，其我所告，順合於汝心以否，當以情告我，無得有不敬者。汝等無得總於貨寶以求官位，當進進自用功德，不當用富也。用此布示於民，必以德義，長任一心以事君，不得懷二意。以遷都既定，故殷勤以戒之。 \par}

%%% Local Variables:
%%% mode: latex
%%% TeX-engine: xetex
%%% TeX-master: "../Main"
%%% End:

%% -*- coding: utf-8 -*-
%% Time-stamp: <Chen Wang: 2024-04-02 11:42:41>

% {\noindent \zhu \zihao{5} \fzbyks } -> 注 (△ ○)
% {\noindent \shu \zihao{5} \fzkt } -> 疏

\chapter{卷十}


\section{說命上第十二【偽】}


高宗夢得說，\footnote{盤庚弟小乙子，名武丁，德高可尊，故號高宗。夢得賢相，其名曰說。說，本又作兌，音悅，注及下篇同。相，息亮反。下同。}使百工營求諸野，得諸傅巖，\footnote{使百官以所夢之形象經求之於野，得之於傅巖之谿。}作\CJKunderwave{說命}三篇。\footnote{命說為相，使攝政。}


{\noindent\zhuan\zihao{6}\fzbyks 傳“盤庚”至“曰說”。正義曰：\CJKunderwave{世本}云：“盤庚崩，弟小辛立。崩，弟小乙立。崩，子武丁立。”是武丁為“盤庚弟小乙子”也。\CJKunderwave{喪服四制}云：“高宗者，武丁。武丁者，殷之賢王也。當此之時，殷衰而復興，禮廢而復起，中而高之，故謂之高宗。”是“德高可尊，故號高宗”也。經雲“爰立作相”，王呼之曰“說”，知其“名曰說”。 \par}

{\noindent\zhuan\zihao{6}\fzbyks 傳“使百”至“之谿”。正義曰：以“工”為官，見其求者眾多，故舉“百官”言之。使百官以所夢之形象經營求於外野。皇甫謐雲“使百工寫其形象”，則謂“工”為工巧之人，與孔異也。\CJKunderwave{釋水}云：“水注川曰谿。”李巡曰：“水出於山，入於川曰谿。”然則“谿”是水流之處，“巖”是山崖之名。序稱“得諸傅巖”，傳雲“得之於傅巖之谿”，以“巖”是總名，故序言之耳。 \par}

{\noindent\zhuan\zihao{6}\fzbyks 傳“命說”至“攝政”。正義曰：經稱“爰立作相”，是命為相也。“惟說命總百官”，是“使攝位”也。 \par}

{\noindent\shu\zihao{5}\fzkt “高宗”至“三篇”。正義曰：殷之賢王有高宗者，夢得賢相，其名曰“說”。群臣之內既無其人，使百官以所夢之形象經營求之於野外，得之於傅氏之巖，遂命以為相。史敘其事，作\CJKunderwave{說命}三篇。 \par}

說命\footnote{始求得而命之。}

{\noindent\shu\zihao{5}\fzkt “說命”。正義曰：此三篇上篇言夢說，始求得而命之；中篇說既總百官，戒王為政；下篇王欲師說而學，說報王為學之有益，王又厲說以伊尹之功。相對以成章，史分序以為三篇也。 \par}

王宅憂，亮陰三祀。\footnote{陰，默也。居憂，信默三年不言。亮，本又作諒，如字，又力章反。}


{\noindent\zhuan\zihao{6}\fzbyks 傳“陰默”至“不言”。正義曰：“陰”者，幽暗之義，“默”亦暗義，故為默也。\CJKunderwave{易}稱“君子之道，或默或語”，則“默”者，不言之謂也。\CJKunderwave{無逸}傳雲“乃有信默，三年不言”，有此“信默”，則“信”謂信任冢宰也。 \par}

{\noindent\shu\zihao{5}\fzkt “王宅憂，亮陰三祀”。正義曰：言王居父憂，信任冢宰，默而不言已三年矣。三年不言，自是常事，史錄此句於首者，謂既免喪事，可以言而猶不言，故述此以發端也。 \par}

既免喪，其惟弗言，\footnote{除喪，猶不言政。}群臣咸諫於王曰：“嗚呼！知之曰明哲，明晢實作則。\footnote{知事則為明智，明智則能製作法則。哲，本又作喆。}天子惟君萬邦，百官承式，\footnote{天下待令，百官仰法。}王言惟作命，不言臣下罔攸稟令。”\footnote{稟，受。令亦命也。}王庸作書以誥曰：“以臺正於四方，臺恐德弗類，茲故弗言。\footnote{用臣下怪之，故作誥。類，善也。我正四方，恐德不善，此故不言。誥，故報反。臺音怡。}恭默思道，夢帝賚予良弼，其代予言。”\footnote{夢天與我輔弼良佐，將代我言政教。賚，力代反，徐音來。}乃審厥象，俾以形旁求於天下。\footnote{審所夢之人，刻其形象，以四方旁求之於民間。俾，必爾反。}說築傅巖之野，惟肖。\footnote{傅氏之巖，在虞虢之界，通道所經，有澗水壞道，常使胥靡刑人築護此道。說賢而隱，代胥靡築之以供食。肖，似。似所夢之形。肖音笑。虢，寡白反。壞音怪。供音恭。}

{\noindent\zhuan\zihao{6}\fzbyks 傳“傅氏”至“之形”。正義曰：傳以“傅”為氏，此巖以“傅”為名，明巖傍有姓傅之民，故云“傅氏之巖”也。\CJKunderwave{屍子}云：“傅巖在北海之洲。”傳言“虞虢之界”，孔必有所案據而言之也。\CJKunderwave{史記·殷本紀}云：“是時說為胥靡，築於傅險。”晉灼\CJKunderwave{漢書音義}云：“胥，相也。靡，隨也。古者相隨坐輕刑之名。”言於時築傅險，則以杵築地，傅說賢人，必身不犯罪，言其說為胥靡，當是時代胥靡也。傳云：“通道所經，有澗水壞道，常使胥靡刑人築護此道。說賢而隱，代胥靡築之以供食。”或亦有成文也。\CJKunderwave{殷本紀}又云，武丁得說，“舉以為相，遂以傅險姓之，號曰傅說”。鄭云：“得諸傅巖，高宗因以傅命說為氏。”案序直言“夢得說”，不言“傅”，或如馬鄭之言。如高宗始命為傅氏，不知舊何氏也。皇甫謐云：“高宗夢天賜賢人，胥靡之衣蒙之而來，且云：‘我徒也，姓傅名說，天下得我者豈徒也哉！’武丁悟而推之曰：‘傅者，相也。說者，歡悅也。天下當有傅我而說民者哉！’明以夢視百官，百官皆非也。乃使百工寫其形象，求諸天下，果見築者胥靡衣褐帶索，執役於虞虢之間、傅巖之野，名說。以其得之傅巖，謂之傅說。”案謐言初夢即雲“姓傅名說”，又言“得之傅巖，謂之傅說”，其言自不相副，謐惟見此書傅，會為近世之語，其言非實事也。 \par}

爰立作相，王置諸其左右。\footnote{於是禮命立以為相，使在左右。}命之曰：“朝夕納誨，以輔臺德。\footnote{言當納諫誨直辭，以輔我德。朝，張遙反。}若金，用汝作礪。\footnote{鐵須礪以成利器。礪，力世反。}若濟巨川，用汝作舟楫。\footnote{渡大水待舟楫。楫音接，徐音集。}若歲大旱，用汝作霖雨。\footnote{霖，三日雨。霖以救旱。}

{\noindent\zhuan\zihao{6}\fzbyks 傳“霖，三日雨”。正義曰：隱九年\CJKunderwave{左傳}云：“凡雨自三日已往為霖。” \par}

啟乃心，沃朕心。若藥弗瞑眩，厥疾弗瘳。\footnote{開汝心，以沃我心。如服藥必瞑眩極，其病乃除。欲其出切言以自警。}


{\noindent\zhuan\zihao{6}\fzbyks 傳“開汝”至“自警”。正義曰：“瞑眩”者，令人憤悶之意也。\CJKunderwave{方言}云：“凡飲藥而毒東齊海岱間或謂之瞑，或謂之眩。”郭璞云：“瞑眩亦通語也。”然則藥之攻病，先使人瞑眩憤亂，病乃得瘳。傳言“瞑眩極”者，言悶極藥乃行也。\CJKunderwave{楚語}稱衛武公作懿以自警，“懿”即\CJKunderwave{大雅·抑}詩也。切言出於傅說，據王以為自警也。 \par}

{\noindent\shu\zihao{5}\fzkt “啟乃”至“弗瘳”。正義曰：當開汝心所有，以灌沃我心。欲令以彼所見，教己未知故也。其沃我心，須切至，若服藥不使人瞑眩憤亂，則其疾不得瘳愈。言藥毒乃得除病，言切乃得去惑也。 \par}

若跣弗視地，厥足用傷。\footnote{跣必視地，足乃無害。言欲使為已視聽。跣，先典反，徐七顯反。為，於偽反。}惟暨乃僚，罔不同心，以匡乃闢。\footnote{與汝並官，皆當倡率，無不同心以匡正汝君。闢，必亦反。}俾率先王，迪我高後，以康兆民。\footnote{言匡正汝君，使循先王之道，蹈成湯之蹤，以安天下。}嗚呼！欽予時命，其惟有終。”\footnote{敬我是命，修其職，使有終。}說復於王曰：“惟木從繩則正，後從諫則聖。\footnote{言木以繩直，君以諫明。}後克聖，臣不命其承，\footnote{君能受諫，則臣不待命，其承意而諫之。}疇敢不祗若王之休命？”\footnote{言王如此，誰敢不敬順王之美命而諫者乎？}

\section{說命中第十三【偽】}

惟說命總百官，\footnote{在冢宰之任。總音揔。}

{\noindent\shu\zihao{5}\fzkt “惟說命總百官”。正義曰：惟此傅說，受王命總百官之職，謂在“冢宰之任”也。說以官高任重，乃進言於王，故史特標此句為發言之端也。 \par}

乃進於王曰：“嗚呼！明王奉若天道，建邦設都，\footnote{天有日月北斗五星二十八宿，皆有尊卑相正之法，言明王奉順此道，以立國設都。宿音秀。}

{\noindent\zhuan\zihao{6}\fzbyks 傳“天有”至“設都”。正義曰：\CJKunderwave{晉語}云：“大者天地，其次君臣。”\CJKunderwave{易·繫辭}云：“天垂象，見吉凶，聖人象之。”皆言人君法天以設官，順天以致治也。天有日月照臨晝夜，猶王官之伯率領諸侯也。北斗環繞北極，猶卿士之周衛天子也。五星行於列宿，猶州牧之省察諸侯也。二十八宿佈於四方，猶諸侯為天子守土也。天象皆有尊卑相正之法，言明王奉順天道以立國設都也。“立國”謂立王國及邦國，“設都”謂設帝都及諸侯國都，總言建國立家之事。 \par}

樹后王君公，承以大夫師長，\footnote{言立君臣上下，將陳為治之本，故先舉其始。王，於方反。長，丁丈反。治，直吏反，下同。}

{\noindent\shu\zihao{5}\fzkt “樹後”至“師長”。正義曰：此又總言設官分職之事也。“樹”，立也。“后王”謂天子也。“君公”謂諸侯也。“承”者奉上之名。“后王君公”，人主也。“大夫師長”，人臣也。臣當奉行君命，故以“承”言之。\CJKunderwave{周禮}立官多以“師”為名，“師”者眾所法，亦是長之義也。大夫已下，分職不同，每官各有其長，故以“師長”言之。三公則“君公”之內包之，卿則“大夫”之文兼之，“師長”之言亦通有士。將陳為治之本，故先舉其始，略言設官，故辭不詳備。為治之本，“惟天聰明”已下皆是也。 \par}

不惟逸豫，惟以亂民。\footnote{不使有位者逸豫民上，言立之主使治民。豫，羊慮反。}惟天聰明，惟聖時憲，惟臣欽若，惟民從乂。\footnote{憲，法也。言聖王法天以立教，臣敬順而奉之，民以從上為治。從，才容反。}

{\noindent\zhuan\zihao{6}\fzbyks 傳“憲法”至“為治”。正義曰：“憲，法”，\CJKunderwave{釋詁}文。人之聞見於耳目，天無形體，假人事以言也。“聰”謂無所不聞，“明”謂無所不見。惟聖人於是法天，言法天以立教，於下無不聞見，除其所惡，納之於善。雖復運有推移，道有升降，其所施為未嘗不法天也。“臣敬順而奉之”，“奉”即上文“承”也，奉承君命而布之於民。“民以從上為治”，不從上命則亂，故“從乂”也。 \par}

惟口起羞，惟甲冑起戎，\footnote{甲，鎧。胄，兜鍪也。言不可輕教令，易用兵。胄，直又反。鎧，苦代反。兜，丁侯反。鍪，莫侯反。易，以豉反。}惟衣裳在笥，惟干戈省厥躬。\footnote{言服不可加非其人，兵不可任非其才。○笥，息嗣反。省，息井反，一本作眚。}


{\noindent\zhuan\zihao{6}\fzbyks 傳“甲鎧”至“用兵”。正義曰：經傳之文無“鎧”與“兜鍪”，蓋秦漢已來始有此名，傳以今曉古也。古之甲冑皆用犀兕，未有用鐵者，而“鍪”、“鎧”之字皆從金，蓋後世始用鐵耳。口之出言為教令，甲冑興師乃用之，言不可輕教令，易用兵也。“易”亦輕也。安危在出令，令之不善，則人違背之，是“起羞”也。靜亂在用兵，伐之無罪，則人叛違之，是“起戎”也。 \par}

{\noindent\zhuan\zihao{6}\fzbyks 傳“言服”至“其才”。正義曰：“非其人”、“非其才”，義同而互文也。\CJKunderwave{周禮·大宗伯}：“以九儀之命正邦國之位，一命受職，再命受服，三命受位，四命受器，五命賜則，六命賜官，七命賜國，八命作牧，九命作伯。”鄭云：“一命始見，命為正吏。受職，治職事也。列國之士一命，王之下士亦一命。再命受服，受玄冕之服。列國之大夫再命，王之中士亦再命。”然則“再命”已上始受衣服，未賜之時在官之篋笥也。甲冑干戈俱是軍器，上言不可輕用兵，此言不可妄委人，雖文重而意異也。 \par}

{\noindent\shu\zihao{5}\fzkt “惟口”至“厥躬”。正義曰：言王者法天施化，其舉止不可不慎。惟口出令不善，以起羞辱；惟甲冑伐非其罪，以起戎兵；言不可輕教令，易用兵也。惟衣裳在篋笥，不可加非其人，觀其能足稱職，然後賜之。惟干戈在府庫，不可任非其才，省其身堪將帥，然後授之。上二句事相類，下二句文不同者，衣裳言在篋笥，干戈不言所在，干戈雲“省厥躬”，衣裳不言視其人，令其互相足也。 \par}

王惟戒茲，允茲克明，乃罔不休。\footnote{言王戒慎此四“惟”之事，信能明，政乃無不美。}惟治亂在庶官。\footnote{言所官得人則治，失人則亂。}官不及私暱，惟其能。\footnote{不加私暱，惟能是官。}\footnote{暱，女乙反。}爵罔及惡德，惟其賢。\footnote{言非賢不爵。}

{\noindent\shu\zihao{5}\fzkt “官不”至“其賢”。正義曰：\CJKunderwave{王制}云：“論定然後官之，任官然後爵之。”鄭云：“官之，使之試守也。爵之，命之也。”然則治其事謂之“官”，受其位謂之“爵”，“官”、“爵”一也，所從言之異耳。“賢”謂德行，“能”謂才用；治事必用能，故“官”雲“惟其能”；受位宜得賢，故“爵”雲“惟其賢”。\CJKunderwave{詩序}云：“任賢使能。”\CJKunderwave{周禮·鄉大夫}：“三年則大比，考其德行道藝，而與賢者能者。”鄭云：“賢者，有德行者。能者，有道藝者。”是“賢”、“能”為異耳。“私暱”謂知其不可而用之，“惡德”謂不知其非而任之，戒王使審求人，絕私好也。 \par}

慮善以動，動惟厥時。\footnote{非善非時不可動。}有其善，喪厥善。矜其能，喪厥功。\footnote{雖天子亦必讓以得之。喪，息浪反。}

{\noindent\shu\zihao{5}\fzkt “有其”至“厥功”。正義曰：人生尚謙讓而憎自取，自有其善，則人不以為善，故實善而喪其善。自誇其能，則人不以為能，故實能而喪其能。由其自取，故人不與之。“有其善”則伐善也。舜美禹云：“汝惟不矜，天下莫與汝爭能。汝惟不伐，天下莫與汝爭功。”是言推而不有，故名反歸之也。 \par}

惟事事乃其有備，有備無患。\footnote{事事，非一事。}無啟寵納侮，\footnote{開寵非其人，則納侮之道。}

{\noindent\shu\zihao{5}\fzkt “無啟寵納侮”。正義曰：君子位高益恭，小人得寵則慢。若寵小人，則必恃寵慢主，無得開小人以寵，自納此輕侮也。“開”謂君出恩以寵臣，“納”謂臣入慢以輕王，據君而言“開”、“納”，以出、入為文也。 \par}

無恥過作非。\footnote{恥過誤而文之，遂成大非。}

{\noindent\zhuan\zihao{6}\fzbyks 傳“恥過”至“大非”。正義曰：仲虺之美成湯云：“改過不吝。”明小人有過，皆惜而不改。\CJKunderwave{論語}云：“小人之過也必文。”恥有過誤而更以言辭文飾之，望人不覺，其非彌甚，故“遂成大非”也。 \par}

惟厥攸居，政事惟醇。\footnote{其所居行，皆如所言，則王之政事醇粹。醇音純。粹，雖遂反。}黷於祭祀，時謂弗欽。禮煩則亂，事神則難。”\footnote{祭不欲數，數則黷，黷則不敬。事神禮煩，則亂而難行。高宗之祀特豐數近廟，故說因以戒之。黷，徒木反。數，色角反。}

{\noindent\zhuan\zihao{6}\fzbyks 傳“祭不”至“戒之”。正義曰：“祭不欲數，數則黷，黷則不敬”，\CJKunderwave{禮記·祭義}文也。此一經皆言祭祀之事，“禮煩”亦謂祭祀之煩，故傳總云：“事神禮煩，則亂而難行。”孔以\CJKunderwave{高宗肜日}祖已訓諸王“祀無豐於暱”，謂傅說此言為彼事而發，故云高宗之祀特豐數於近廟，故說因而戒之。 \par}

王曰：“旨哉！說，乃言惟服。\footnote{旨，美也。美其所言皆可服行。}乃不良於言，予罔聞於行。”\footnote{汝若不善於所言，則我無聞於所行之事。}說拜稽首，曰：“非知之艱，行之惟艱。\footnote{言知之易，行之難。以勉高宗。}王忱不艱，允協於先王成德，\footnote{王心誠不以行之為難，則信合於先王成德。忱，市林反。}惟說不言有厥咎。”\footnote{王能行善，而說不言，則有其咎罪。}

\section{說命下第十四【偽】}

王曰：“來，汝說。臺小子舊學於甘盤，\footnote{學先王之道。甘盤，殷賢臣有道德者。臺音怡。}

{\noindent\shu\zihao{5}\fzkt “王曰”至“甘盤”。正義曰：“舊學於甘盤”謂為王子時也。\CJKunderwave{君奭}篇周公仰陳殷之賢臣云：“在武丁，時則有若甘盤。”然則甘盤於高宗之時有大功也。上篇高宗免喪不言，即求傅說，似得說時無賢臣矣。蓋甘盤於小乙之世以為大臣，小乙將崩，受遺輔政，高宗之初得有大功。及高宗免喪，甘盤已死，故\CJKunderwave{君奭}傳曰：“高宗即位，甘盤佐之，後有傅說。”是言傅說之前有甘盤也。但下句言“既乃遯於荒野”，是學訖乃遁，非即位之初從甘盤學也。 \par}

既乃遁於荒野，入宅於河。\footnote{既學而中廢業，遁居田野。河，洲也。其父欲使高宗知民之艱苦，故使居民間。遁，徒頓反。}

{\noindent\zhuan\zihao{6}\fzbyks 傳“既學”至“民間”。正義曰：“河”是水名，水不可居，而云“入宅於河”，知在河之洲也。\CJKunderwave{釋水}云：“水中可居者曰洲。”初遁田野，後入河洲，言其徙居無常也。\CJKunderwave{無逸}云：“其在高宗，時舊勞於外，爰暨小人。”言“其父欲使高宗知民之艱苦，故使居民間”也。於時蓋未為太子，殷道雖質，不可既為太子，更得與民雜居。 \par}

自河徂亳，暨厥終罔顯。\footnote{自河往居亳，與今其終，故遂無顯明之德。}爾惟訓於朕志，\footnote{言汝當教訓於我，使我志通達。}若作酒醴，爾惟曲糵；\footnote{酒醴須曲糵以成，亦言我須汝以成。曲，起六反。糵，魚列反。}若作和羹，爾惟鹽梅。\footnote{鹽，咸。梅，醋。羹須咸醋以和之。羹音庚，一音衡。鹽，餘廉反。梅亦作楳。醋，七故反。和如字，又胡臥反。}爾交修予，罔予棄，予惟克邁乃訓。”\footnote{交，非一之義。邁，行也。言我能行汝教。}

{\noindent\zhuan\zihao{6}\fzbyks 傳“交非”至“汝教”。正義曰：“爾交修予”，令其交更修治己也。故以“交”為“非一之義”，言交互教之，非一事之義。“邁，行”，\CJKunderwave{釋詁}文。 \par}

說曰：“王，人求多聞，時惟建事，學於古訓，乃有獲。\footnote{王者求多聞以立事，學於古訓，乃有所得。}事不師古，以克永世，匪說攸聞。\footnote{事不法古訓而以能長世，非說所聞。言無是道。}惟學遜志，務時敏，厥修乃來。\footnote{學以順志，務是敏疾，其德之修乃來。}

{\noindent\shu\zihao{5}\fzkt “惟學”至“乃來”。正義曰：人志本欲求善，欲學順人本志，學能務是敏疾，則其德之修乃自來。言務之既疾，則德自來歸己也。 \par}

允懷於茲，道積於厥躬。\footnote{信懷此學志，則道積於其身。}惟斆學半，念終始典於學，厥德修罔覺。\footnote{斆，教也。教然後知所困，是學之半。終始常念學，則其德之修，無能自覺。斆，戶孝反。}

{\noindent\shu\zihao{5}\fzkt “惟斆”至“罔覺”。正義曰：教人然後知困，知困必將自強，惟教人乃是學之半，言其功半於學也。於學之法，念終念始，常在於學，則其德之修漸漸進益，無能自學其進。言日有所益，不能自知也。 \par}

監於先王成憲，其永無愆。\footnote{愆，過也。視先王成法，其長無過，其惟學乎！愆，起虔反。}惟說式克欽承，旁招俊乂，列於庶位。”\footnote{言王能志學，說亦用能敬承王志，廣招俊乂，使列眾官。俊，本又作畯。}王曰：“嗚呼！說，四海之內，咸仰朕德，時乃風。\footnote{風，教也。使天下皆仰我德，是汝教。仰如字，徐五亮反。}股肱惟人，良臣惟聖。\footnote{手足具，乃成人。有良臣，乃成聖。}昔先正保衡，作我先王，\footnote{保衡，伊尹也。作，起。正，長也。言先世長官之臣。長，丁丈反，下同。}

{\noindent\zhuan\zihao{6}\fzbyks 傳“保衡”至“之臣”。正義曰：保衡、阿衡俱伊尹也。\CJKunderwave{君奭}傳曰：“伊尹為保衡，言天下所取安所取平也。”鄭箋云：“阿，倚。衡，平也。伊尹湯所依倚而取平也，故以為官名。”又云：“太甲時曰保衡。”鄭不見古文\CJKunderwave{太甲}雲“不惠於阿衡”，故此為解，孔所不用。計此阿衡、保衡非常人之官名，蓋當時特以此名號伊尹也。“作”訓為起，言起而助湯也。“正，長”，\CJKunderwave{釋詁}文。 \par}

乃曰：‘予弗克俾厥後惟堯舜，其心愧恥，若撻於市。”\footnote{言伊尹不能使其君如堯舜，則恥之，若見撻於市，故成其能。俾，必爾反。撻，他達反。}一夫不獲，則曰時予之辜。\footnote{伊尹見一夫不得其所，則以為己罪。}佑我烈祖，格於皇天。\footnote{言以此道左右成湯，功至大天，無能及者。}爾尚明保予，罔俾阿衡，專美有商。\footnote{汝庶幾明安我事，則與伊尹同美。阿，烏何反。}惟後非賢不乂，惟賢非後不食。\footnote{言君須賢治，賢須君食。治，直吏反。}其爾克紹乃闢於先王，永綏民。”\footnote{能繼汝君於先王，長安民，則汝亦有保衡之功。闢，必亦反。}說拜稽首，曰：“敢對揚天子之休命。”\footnote{對，答也。答受美命而稱揚之。}

\section{高宗肜日第十五}


高宗祭成湯，有飛雉升鼎耳而雊，\footnote{耳不聰之異。雊，鳴。雊，工豆反。}祖已訓諸王，\footnote{賢臣也，以訓道諫王。己音紀。}作\CJKunderwave{高宗肜日}、\CJKunderwave{高宗之訓}。\footnote{所以訓也，亡。肜音融。}


{\noindent\zhuan\zihao{6}\fzbyks 傳“耳不”至“雊鳴”。正義曰：經言“肜日，有雊雉”，不知祭何廟，鳴何處，故序言“祭成湯”、“升鼎耳”以足之。禘祫與四時之祭，祭之明日皆為肜祭，不知此肜是何祭之肜也。\CJKunderwave{洪範}“五事”有貌、言、視、聽、思，若貌不恭、言不從、視不明、聽不聰、思不睿，各有妖異興焉。雉乃野鳥，不應入室，今乃入宗廟之內，升鼎耳而鳴，孔以雉鳴在鼎耳，故以為“耳不聰之異”也。\CJKunderwave{洪範·五行傳}云：“視之不明，時則有羽蟲之孽。聽之不聰，時則有介蟲之孽。”言之不從，時則有毛蟲之孽。貌之不恭，時則有鱗蟲之孽。思之不睿，時則有倮蟲之孽。”先儒多以此為羽蟲之孽，非為“耳不聰”也。\CJKunderwave{漢書·五行志}：“劉歆以為鼎三足，三公象也，而以耳行。野鳥居鼎耳，是小人將居公位，敗宗廟之祀也。”鄭云：“鼎，三公象也，又用耳行，雉升鼎耳而鳴，象視不明，天意若雲當任三公之謀以為政。”劉、鄭雖小異，其為羽蟲之孽則同，與孔意異。\CJKunderwave{詩}云：“雉之朝雊，尚求其雌。”\CJKunderwave{說文}云：“雊，雄雉鳴也。雷始動，雉乃鳴而雊其頸。” \par}

{\noindent\zhuan\zihao{6}\fzbyks 傳“所以訓也，亡”。正義曰：名\CJKunderwave{高宗之訓}，所以訓高宗也。此二篇俱是祖己之言，並是訓王之事，經雲“乃訓於王”，此篇亦是訓也。但所訓事異，分為二篇，標此為發言之端，故以“肜日”為名。下篇總諫王之事，故名之“訓”，終始互相明也。\CJKunderwave{肆命}、\CJKunderwave{徂後}，孔歷其名於\CJKunderwave{伊訓}之下，別為之傳。此\CJKunderwave{高宗之訓}因序為傳，不重出名者，此以訓王事同，因解文便作傳，不為例也。 \par}

{\noindent\shu\zihao{5}\fzkt “高宗”至“之訓”。正義曰：高宗祭其太祖成湯於肜祭之日，有飛雉來升祭之鼎耳而雊鳴，其臣祖已以為王有失德而致此祥，遂以道義訓王，勸王改修德政。史敘其事，作\CJKunderwave{高宗肜日}、\CJKunderwave{高宗之訓}二篇。 \par}

高宗肜日\footnote{祭之明日又祭。殷曰肜，周曰繹。○繹音亦，字書作釋。\CJKunderwave{爾雅}云：“又祭也。周曰繹，商曰肜，夏曰復胙。”}

{\noindent\zhuan\zihao{6}\fzbyks 傳“祭之”至“曰繹”。正義曰：\CJKunderwave{釋天}云：“繹，又祭也。周曰繹，商曰肜。”孫炎曰：“祭之明日尋繹復祭也。”“肜”者，相尋不絕之意。\CJKunderwave{春秋}宣八年六月“辛巳，有事於太廟。壬午，猶繹”。\CJKunderwave{穀梁傳}曰：“繹者，祭之旦日之享賓也。”是肜者，“祭之明日又祭”也。\CJKunderwave{爾雅}因繹祭而本之上世，故先周後商，此以上代先後，故與\CJKunderwave{爾雅}倒也。\CJKunderwave{釋天}又云“夏曰復胙”，郭璞雲“未見所出”，或無此一句。孔傳不言“夏曰復胙”，於義非所須，或本無此事也。\CJKunderwave{儀禮·有司徹}上大夫曰“儐屍”，與正祭同日。鄭康成注\CJKunderwave{詩·鳧鷖}云，祭天地社稷山川，五祀皆有繹祭。 \par}

高宗肜日，越有雊雉。\footnote{於肜日有雉異。}祖己曰：“惟先格王，正厥事。”\footnote{言至道之王遭變異，正其事而異自消。}


{\noindent\zhuan\zihao{6}\fzbyks 傳“言至”至“自消”。正義曰：“格”訓至也。“至道之王”謂用心至極，行合於道。遭遇變異，改修德教，正其事而異自消。大戊拱木，武丁雊雉，皆感變而懼，殷道復興，是異自消之驗也。至道之王，當無災異，而云遭變消災者，天或有譴告，使之至道，未必為道不至而致此異。且匆尋戒之辭，不可執文以害意也。此經直雲“祖己曰”，不知與誰語，鄭雲“謂其黨”，王肅雲“言於王”。下句始言“乃訓於王”，此句未是告王之辭，私言告人，鄭說是也。 \par}

{\noindent\shu\zihao{5}\fzkt “高宗”至“厥事”。正義曰：高宗既祭成湯，肜祭之日，於是有雊鳴之雉在於鼎耳，此乃怪異之事。賢臣祖已見其事而私自言曰：“惟先世至道之王遭遇變異，則正其事而異自消也。”既作此言，乃進言訓王。史錄其事，以為訓王之端也。 \par}

乃訓於王。曰：“惟天監下民，典厥義。\footnote{祖己既言，遂以道訓諫王，言天視下民，以義為常。}降年有永有不永，非天夭民，民中絕命。\footnote{言天之下年與民，有義者長，無義者不長，非天欲夭民，民自不修義以致絕命。中，丁仲反，又如字。}民有不若德，不聽罪。天既孚命正厥德。\footnote{不順德，言無義。不服罪，不改修。天已信命正其德，謂有永有不永。}


{\noindent\zhuan\zihao{6}\fzbyks 傳“言天”至“絕命”。正義曰：經惟言“有永有不永”，安和由義者？以上句雲“惟天監下民，典厥義”，天既以義為常，知命之長短莫不由義，故云“天之下年與民，有義者長。無義者不長”也。民有五常之性，謂仁義禮智信也，此獨以“義”為言者，五常指體則別，理亦相通；“義”者，宜也，得其事宜；五常之名，皆以適宜為用，故稱“義”可以總之也。民有貴賤貧富愚智好醜，不同多矣，獨以夭壽為言者，鄭玄云：“年命者，蠢愚之人尤愒焉，故引以諫王也。”愒，貪也。\CJKunderwave{洪範}“五福”以壽為首，“六極”以短折為先，是年壽者最是人之所貪，故祖己引此以諫王也。 \par}

{\noindent\zhuan\zihao{6}\fzbyks 傳“不順”至“不永”。正義曰：傳亦顧上經，故“不順德，言無義”也。“聽”謂聽從，故以“不聽”為“不服罪”。言既為罪，過而不肯改修也。“天已信命正其德”，言天自信命，賞有義，罰無義，此事必信也。天自正其德，福善禍淫，其德必不差也。謂民有永有不永，天隨其善惡而報之。勸王改過修德以求永也。 \par}

{\noindent\shu\zihao{5}\fzkt “乃訓”至“厥德”。正義曰：祖己既私言其事，乃以道訓諫於王曰：“惟天視此下民，常用其義。”言以義視下，觀其為義以否。“其下年與民，有長者，有不長者”。言與為義者長，不義者短。“短命者非是天欲夭民，民自不修義，使中道絕其性命。但人有為行不順德義，有過不服聽罪，過而不改，乃致天罰，非天欲夭之也。天既信行賞罰之命，正其馭民之德，欲使有義者長，不義者短，王安得不行義事，求長命也？” \par}

乃曰：‘其如臺。’\footnote{祖己恐王未受其言，故乃復曰，天道其如其所言。臺音怡。復，扶又反。}嗚呼！王司敬民，罔非天胤，典祀無豐於暱。”\footnote{胤，嗣。暱，近也。嘆以感王入其言，王者主民，當敬民事。民事無非天所嗣常也，祭祀有常，不當特豐於近廟。欲王因異服罪改修之。豐，芳弓反。暱，女乙反。\CJKunderwave{屍子}云：“不避遠暱。”暱，近也。又乃禮反。馬云：“暱，考也，謂禰廟也。”}


{\noindent\zhuan\zihao{6}\fzbyks 傳“胤嗣”至“改修之”。正義曰：\CJKunderwave{釋詁}云：“胤、嗣，繼也。”俱訓為繼，是“胤”德為嗣，嗣亦繼之義也。\CJKunderwave{釋詁}云：“即，尼也。”孫炎曰：“即猶今也，尼者近也。”郭璞引\CJKunderwave{屍子}曰“悅尼而來遠”，是“尼”為近也。“尼”與“暱”音義同。烝民不能自治，自立君以主之，是“王者主民”也。既與民為主，當敬慎民事。民事無大小，無非天所嗣常也，言天意欲令繼嗣行之，所以為常道也。“祭祀有常”，謂犧牲粢盛尊彝俎豆之數禮有常法。“不當特豐於近廟”，謂犧牲禮物多也。祖己知高宗豐於近廟，欲王因此雊雉之異，服罪改修以從禮耳，其異不必由豐近而致之也。王肅亦云：“高宗豐於禰，故有雊雉升遠祖成湯廟鼎之異。” \par}

{\noindent\shu\zihao{5}\fzkt “嗚呼”至“於暱”。正義曰：祖己恐其言不入王意，又嘆而戒之：“嗚呼！王者主民，當謹敬民事。民事無非天所繼嗣以為常道者也。天以其事為常，王當繼天行之。祀禮亦有常，無得豐厚於近廟，若特豐於近廟，是失於常道。”高宗豐於近廟，欲王服罪改修也。 \par}

\section{西伯戡黎第十六}


殷始咎周，\footnote{咎，惡。○咎，其九反，馬云：“咎周者，為周所咎。}周人乘黎。\footnote{乘，勝也。所以見惡。黎，力兮反，國名，\CJKunderwave{尚書大傳}作耆。}祖伊恐，\footnote{祖己後賢臣。}告於受，\footnote{受，紂也，音相亂。帝乙之子，嗣立，暴虐無道。○受如字，傳云：“受，紂也。音相亂。”馬云：“受讀曰紂。或曰受婦人之言，故號曰受也。”}作\CJKunderwave{西伯戡黎}。\footnote{戡亦勝也。○伯亦作柏。戡音堪，\CJKunderwave{說文}作𢦟，雲“殺也”。以此戡訓刺，音竹甚反。勝，詩證反。}

{\noindent\zhuan\zihao{6}\fzbyks 傳“咎惡”，又云“乘勝”至“見惡”。正義曰：\CJKunderwave{易·繫辭}雲“無咎者善補過也”，則“咎”是過之別名，以彼過而憎惡之，故“咎”為惡也。以其勝黎，所以見惡，釋其見惡之由，是周人勝黎之後始惡之。\CJKunderwave{詩毛傳}云：“乘，陵也。”乘駕是加陵之意，故“乘”為勝也。鄭玄云：“紂聞文王斷虞芮之訟，又三伐皆勝，而始畏惡之。”所言據\CJKunderwave{書傳}為說，伏生\CJKunderwave{書傳}雲“文王受命，一年斷虞芮之質，二年伐邘，三年伐密須，四年伐犬夷，五年伐耆，六年伐崇，七年而崩”。耆即黎也。乘黎之前始言惡周，故鄭以伐邘、伐密須、伐犬夷三伐皆勝，始畏惡之。\CJKunderwave{武成}篇文王“誕膺天命”，九年乃崩，則伐國之年不得如\CJKunderwave{書傳}所說，未必見三伐皆勝始畏之。 \par}

{\noindent\zhuan\zihao{6}\fzbyks 傳“祖己後賢臣”。正義曰：此無所出，正以同為祖氏，知是其後，明能先覺，故知賢臣。 \par}

{\noindent\zhuan\zihao{6}\fzbyks 傳“受紂”至“無道”。正義曰：經雲“奔告於王”，王無諡號，故序言“受”以明之。此及\CJKunderwave{泰誓}、\CJKunderwave{武成}皆呼此君為“受”，自外書傳皆呼為“紂”。“受”即“紂”也，音相亂，故字改易耳。\CJKunderwave{殷本紀}云：“帝乙崩，子辛立，是為帝辛，天下謂之紂。”鄭玄云：“紂，帝乙之少子，名辛。帝乙愛而欲立焉，號曰受德，時人傳聲轉作紂也。”史掌書，知其本，故曰“受”，與孔大同。\CJKunderwave{諡法}云：“殘義損善曰紂。”殷時未有諡法，後人見其惡，為作惡義耳。 \par}

{\noindent\zhuan\zihao{6}\fzbyks 傳“戡亦勝也”。正義曰：“戡，勝”，\CJKunderwave{釋詁}文。孫炎曰：“戡，強之勝也。” \par}

{\noindent\shu\zihao{5}\fzkt “殷始”至“戡黎”。正義曰：文王功業稍高，王兆漸著，殷之朝廷之臣始畏惡周家。所以畏惡之者，以周人伐而勝黎邑故也。殷臣祖伊見周克黎國之易，恐其終必伐殷，奔走告受，言殷將滅。史敘其事，作\CJKunderwave{西伯戡黎}。 \par}

西伯戡黎

西伯既戡黎，\footnote{近王圻之諸侯，在上黨東北。近，附近之近。圻，巨依反。}


{\noindent\zhuan\zihao{6}\fzbyks 傳“近王”至“東北”。正義曰：黎國，漢之上黨郡壺關所治黎亭是也。紂都朝歌，王圻千里，黎在朝歌之西，故為“近王圻之諸侯”也。鄭云：“入紂圻內。”文王猶尚事紂，不可伐其圻內。所言“圻內”，亦無文也。 \par}

{\noindent\shu\zihao{5}\fzkt “西伯戡黎”。正義曰：鄭玄云：“西伯，周文王也。時國於岐，封為雍州伯也。國在西，故曰西伯。”王肅云：“王者中分天下，為二公總治之，謂之二伯，得專行征伐，文王為西伯。黎侯無道，文王伐而勝之。”兩說不同，孔無明解。下傳雲“文王率諸侯以事紂”，非獨率一州之諸侯也。\CJKunderwave{論語}稱“三分天下有其二，以服事殷”，謂文王也。終乃三分有二，豈獨一州牧乎？且言“西伯”對東為名，不得以國在西而稱“西伯”也，蓋同王肅之說。 \par}

祖伊恐，奔告於王。曰：“天子，天既訖我殷命，\footnote{文王率諸侯以事紂，內秉王心，紂不能制，今又克有黎國，迫近王圻，故知天已畢訖殷之王命。言將化為周。王心，於況反，下注“宜王者”同。}

{\noindent\zhuan\zihao{6}\fzbyks 傳“文王”至“為周”。正義曰：襄四年\CJKunderwave{左傳}云：“文王率殷之叛國以事紂。”是率諸侯共事紂也。貌雖事紂，內秉王心，佈德行威，有將王之意。而紂不能制，日益強大，今復克有黎國，迫近王圻，似有天助之力，故云“天已畢訖殷之王命”，言殷祚至此而畢，將欲化為周也。 \par}

格人元龜，罔敢知吉。\footnote{至人以人事觀殷，大龜以神靈考之，皆無知吉。}

{\noindent\zhuan\zihao{6}\fzbyks 傳“至人”至“知吉”。正義曰：“格”訓為至，“至人”謂至道之人，有所識解者也。至人以人事觀殷，大龜有神靈逆知來物，故“大龜以神靈考之”。二者皆無知殷有吉者，言必兇也。祖伊未必問至人，親灼龜，但假之以為言耳。 \par}

非先王不相我後人，惟王淫戲用自絕。\footnote{非先祖不助子孫，以王淫過戲怠，用自絕於先王。相，息亮反。}故天棄我，不有康食。不虞天性，不迪率典。\footnote{以紂自絕於先王，故天亦棄之，宗廟不有安食於天下。而王不度知天性命所在，而所行不蹈循常法。言多罪。度，待洛反。}

{\noindent\zhuan\zihao{6}\fzbyks 傳“以紂”至“多罪”。正義曰：\CJKunderwave{禮記}稱“萬物本於天，人本於祖”，則天與先王俱是人君之本。紂既自絕於先王，亦自絕於天。上經言紂自絕先王，此言天棄紂，互明紂自絕，然後天與先王棄絕之。故傳申通其意，“以紂自絕先王，故天亦棄之”。“亦”者，亦先王，言先王與天俱棄之也。\CJKunderwave{孝經}言天子得萬國之歡心，以事其先王，然後祭則鬼享之。今紂既自絕於先王，先王不有安食於天下，言紂雖以天子之尊事宗廟，宗廟之神不得安食也。而王不度知天命所在，不知已之性命當盡也，而所行不蹈循常法，動悉違法，言多罪。 \par}

今我民罔弗欲喪，曰：‘天曷不降威？大命不摯？’今王其如臺。”\footnote{摯，至也。民無不欲王之亡，言：“天何不下罪誅之？有大命宜王者，何以不至？”王之兇害，其如我所言。摯音至，本又作𡠗。}

{\noindent\zhuan\zihao{6}\fzbyks 傳“摯至也”至“所言”。正義曰：摯”、“至”同音，故“摯”為至也。“言天何不下罪誅之”，恨其久行虐政，欲得早殺之也。“有大命宜王者，何以不至”，向望大聖之君，欲令早伐紂也。“王之兇禍，其如我之所言”，以王不信，故審告之也。 \par}

王曰：“嗚呼！我生不有命在天？”\footnote{言我生有壽命在天，民之所言，豈能害我。遂惡之辭。}祖伊反曰：“嗚呼！乃罪多參在上，乃能責命於天？”\footnote{反，報紂也，言汝罪惡眾多，參列於上天，天誅罰汝，汝能責命於天，拒天誅乎？參，七南反，馬云：“參字累在上。”}殷之即喪，指乃功，不無戮於爾邦。”\footnote{言殷之就亡，指汝功事所致，汝不得無死戮於殷國，必將滅亡，立可待。}

\section{微子第十七}


殷既錯天命，\footnote{錯，亂也。錯，七各反，馬云：“廢也。”}微子作誥父師、少師。\footnote{告二師而去紂。少，詩照反。}


{\noindent\zhuan\zihao{6}\fzbyks 傳“錯，亂也”。正義曰：交錯是渾亂之義，故為亂也。不指言紂惡而言“錯亂天命”者，天生烝民，立君以牧之，為君而無君道，是錯亂天命，為惡之大，故舉此以見惡之極耳。 \par}

{\noindent\shu\zihao{5}\fzkt “殷既”至“少師”。正義曰：殷紂既暴虐無道，錯亂天命，其兄微子知紂必亡，以作言誥告父師箕子、少師比干。史敘其事而作此篇也。名曰\CJKunderwave{微子}而不言“作\CJKunderwave{微子}”者，已言“微子作誥”，以可知而省文也。 \par}

微子\footnote{微，圻內國名。子，爵。為紂卿士，去無道。}

{\noindent\zhuan\zihao{6}\fzbyks 傳“微圻”至“無道”。正義曰：微國在圻內，先儒相傳為然。鄭玄以為微與箕俱在圻內，孔雖不言箕，亦當在圻內也。王肅云：“微，國名。子，爵。入為王卿士。”肅意蓋以微為圻外，故言“入”也。微子名啟，\CJKunderwave{世家}作開，避漢景帝諱也。啟與其弟仲衍，皆是紂之同母庶兄，\CJKunderwave{史記}稱“微仲衍”。衍亦稱“微”者，微子封微，以微為氏，故弟亦稱微，猶如春秋之世虞公之弟稱虞叔，祭公之弟稱祭叔。微子若非大臣，則無假憂紂，亦不必須去，以此知其為卿士也。傳雲“去無道”者，以“去”見其為卿士也。 \par}

微子若曰：“父師、少師，\footnote{父師，太師，三公，箕子也。少師，孤卿，比干。微子以紂距諫，知其必亡，順其事而言之。}殷其弗或亂正四方。\footnote{或，有也。言殷其不有治正四方之事，將必亡。治，直吏反。}我祖底遂陳於上，\footnote{言湯致遂其功，陳列於上世。}我用沈酗於酒，用亂敗厥德於下。\footnote{我，紂也。沈湎酗醟，敗亂湯德於後世。沈，徐直金反。酗，況具反，以酒為兇曰酗，\CJKunderwave{說文}作䣱，云：“酒醟。”湎，面善反。醟音詠，\CJKunderwave{說文}於命反，䣱酒也。}


{\noindent\zhuan\zihao{6}\fzbyks 傳“父師”至“言之”。正義曰：以\CJKunderwave{畢命}之篇王呼畢公為“父師”，畢公時為太師也。\CJKunderwave{周官}云：“太師、太傅、太保，茲惟三公。少師、少傅、少保曰三孤。”\CJKunderwave{家語}云：“比干官則少師。”少師是比干，知太師是箕子也。遍檢書傳，不見箕子之名，惟司馬彪注\CJKunderwave{莊子}云：“箕子名胥餘。”不知出何書也。\CJKunderwave{周官}以少師為孤，此傳言“孤卿”者，孤亦卿也，\CJKunderwave{考工記}曰“外有九室，九卿朝焉”，是三孤六卿共為九卿也。比干不言封爵，或本無爵，或有而不言也。\CJKunderwave{家語}云：“比干是紂之親，則諸父。”知比干是紂之諸父耳。箕子則無文。\CJKunderwave{宋世家}云：“箕子者，紂親戚也。”止言親戚，不知為父為兄也。鄭玄、王肅皆以箕子為紂之諸父，服虔、杜預以為紂之庶兄，既無正文，各以意言之耳。微子以紂距諫，知其必亡，心欲去之，故順其去事而言，呼二師以告之。 \par}

{\noindent\zhuan\zihao{6}\fzbyks 傳“或有”至“必亡”。正義曰：“或”者不定之辭，其事欲當然，則是有此事，故以“或”為有也。鄭玄\CJKunderwave{論語注}亦云：“或之言有也，不有言無也。”天子，天下之主，所以治正四方，“言殷其不有治正四方之事”，言將必亡。 \par}

{\noindent\zhuan\zihao{6}\fzbyks 傳“我紂”至“後世”。正義曰：嗜酒亂德，是紂之行，故知“我”，我紂也。人以酒亂，若沈於水，故以耽酒為“沈”也。湎然是齊同之意，\CJKunderwave{詩}云：“天不湎爾以酒。”鄭云：“天不同汝顏色以酒。”是“湎”謂酒變面色，湎然齊同，無復平時之容也。\CJKunderwave{說文}云：“酗，醟也。”然則“酗”、“醟”一物，謂飲酒醉而發怒。經言亂敗其德，必有所屬，上言“我祖”指謂成湯，知言“敗亂湯德於後世”也。上謂前世，故下為後世也。 \par}

殷罔不小大，好草竊奸宄。\footnote{草野竊盜，又為奸宄於內外。好，呼報反。宄音軌。}卿士師師非度，凡有辜罪，乃罔恆獲。\footnote{六卿典士相師效，為非法度，皆有辜罪，無秉常得中者。度如字。}小民方興，相為敵讎。\footnote{卿士既亂，而小人各起一方，共為敵讎。言不和同。讎，常周反。}今殷其淪喪，若涉大水，其無津涯。\footnote{淪，沒也。言殷將沒亡，如涉大水，無涯際，無所依就。淪音倫，徐力允反。喪，息浪反。涯，五皆反，又宜佳反。}殷遂喪，越至於今。”\footnote{言遂喪亡於是，至於今，到不待久。}

{\noindent\zhuan\zihao{6}\fzbyks 傳“六卿”至“中者”。正義曰：“士”訓事也，故“卿士”為“六卿典事”。“師師”言相師效為非法度之事也。止言“卿士”，以貴者尚爾，見賤者皆然。故王肅云：“卿士以下，轉相師效為非法度之事也。”鄭云：“凡猶皆也。”傳意亦然，以“凡”為皆，言卿士以下在朝之臣，其所舉動皆有辜罪，無人能秉常行得中正者。 \par}

{\noindent\shu\zihao{5}\fzkt “微子”至“於今”。正義曰：微子將欲去殷，順其去事而言，曰“父師”、“少師”，呼二師與之言也。今殷國其將不復有治正四方之事，言其必滅亡也。昔我祖成湯，致行其道，遂其功業，陳列於上世矣。今我紂惟用沈湎酗Ｅ１於酒，用是亂敗其祖之德於下。由紂亂敗之故，今日殷人無不小大皆好草竊奸宄。雖在朝卿士，相師師為非法度之事。朝廷之臣皆有辜罪，乃無有一人能秉常得中者。在外小人，方方各起，相與共為敵讎。荒亂如此，今殷其沒，亡若涉大水，其無津濟涯岸。殷遂喪亡，言不復久也。“此喪亡於是，至於今，到必不得更久也”。 \par}

曰：“父師、少師，我其發出狂，吾家耄遜於荒。\footnote{我念殷亡，發疾生狂，在家耄亂，故欲遯出於荒野。言愁悶。出，尺遂反。耄，字又作旄，莫報反，注同。遯，徒困反，徐徒頓反，一音都困反。}今爾無指，告予顛隮，若之何其？”\footnote{汝無指意告我殷邦顛隕隮墜，如之何其救之？隮，子細反，\CJKunderwave{玉篇}子兮反，\CJKunderwave{切韻}祖稽反。隕，于敏反。}


{\noindent\zhuan\zihao{6}\fzbyks 傳“我念”至“愁悶”。正義曰：狂生於心而出於外，故傳以“出狂”為“生狂”。應璩詩云“積念發狂痴”，此其事也。在家思念之深，精神益以耄亂。鄭玄云：“耄，昏亂也。”在家不堪耄亂，故欲遯出於荒野，言愁悶之至。\CJKunderwave{詩}云：“駕言出遊，以寫我憂。”亦此意也。 \par}

{\noindent\zhuan\zihao{6}\fzbyks 傳“汝無”至“救之”。正義曰：“無指意告我者”，謂無指殷亡之事告我，言殷將隕墜，欲留我救之。“顛”謂從上而隕，“隮”謂墜於溝壑，皆滅亡之意也。昭十三年\CJKunderwave{左傳}曰：“小人老而無子，知隮於溝壑矣。”王肅云：“隮，隮溝壑。”言此“隮”之義如\CJKunderwave{左傳}也。 \par}

{\noindent\shu\zihao{5}\fzkt “曰父師”至“何其”。正義曰：微子既言紂亂，乃問身之所宜，止而復言，故別加一“曰父師少師”，更呼而告之也。“我念殷亡之故，其心發疾生狂，吾在家心內耄亂，欲遜遯出於荒野。今汝父師少師無指滅亡之意告我云，殷邦其隕墜，則當如之何其救之乎？”恐其留己共救之也。 \par}

父師若曰：“王子，\footnote{比干不見，明心同，省文。微子帝乙元子，故曰王子。見，賢遍反。省，所景反。}天毒降災荒殷邦，方興沈酗於酒，\footnote{天生紂為亂，是天毒下災，四方化紂沈湎，不可如何。}乃罔畏畏，咈其耇長舊有位人。\footnote{言起沈湎，上不畏天災，下不畏賢人。違戾耇老之長致仕之賢，不用其教，法紂故。咈，扶勿反。耇，工口反。長，丁丈反，注同。}


{\noindent\zhuan\zihao{6}\fzbyks 傳“比干”至“王子”。正義曰：諮二人而一人答，“明心同，省文”也。鄭云：“少師不答，志在必死。”然則箕子本意豈必求生乎？身若求生，何以不去？既“不顧行遁”，明期於必死，但紂自不殺之耳。若比干意異，箕子則別有答，安得默而不言？孔解“心同”是也。“微子帝乙元子”，\CJKunderwave{微子之命}有其文也。父師言微子為“王子”，則父師非王子矣，鄭、王等以為紂之諸父當是實也。 \par}

{\noindent\zhuan\zihao{6}\fzbyks 傳“天生”至“如何”。正義曰：“荒殷邦”者，乃是紂也，而云“天毒降災”，故言“天生紂為亂”，本之於天，天毒下災也。以微子云“若之何”，此答彼意，故言“四方化紂沈湎，不可如何”。 \par}

{\noindent\zhuan\zihao{6}\fzbyks 傳“言起”至“紂故”。正義曰：文在“方興沈酗”之下，則此無所畏畏者，謂當時四方之民也。民所當畏，惟畏天與人耳，故知二畏者，上不畏天，下不畏賢人。違戾耇長與舊有位人，即是不畏賢人，故不用其教，紂無所畏，此民無所畏，謂法紂故也。 \par}

今殷民乃攘竊神祇之犧牷牲用，以容將食，無災。\footnote{自來而取曰攘。色純曰犧。體完曰牷。牛羊豕曰牲。器實曰用。盜天地宗廟牲用，相容行食之，無災罪之者。言政亂。攘，如羊反，因來而取曰攘。竊，馬云：“往盜曰竊。”神祇，天曰神，地曰祇。犧，許宜反。牷音全。}降監殷民，用乂讎斂，召敵讎不怠。\footnote{下視殷民，所用治者，皆重賦傷民、斂聚怨讎之道，而又亟行暴虐，自召敵讎不解怠。讎如字，下同。徐云：“鄭音疇。”馬本作稠，云：“數也。”豔，力檢反；馬、鄭力豔反，謂賦斂也；徐云：“鄭力劍反。”治，直吏反。亟，欺忌反，數也；又紀力反；本又作極，如字，至也。解，佳賣反。}罪合於一，多瘠罔詔。\footnote{言殷民上下有罪，皆合於一法紂，故使民多瘠病，而無詔救之者。瘠，在益反。}商今其有災，我興受其敗。\footnote{災滅在近，我起受其敗，言宗室大臣義不忍去。}


{\noindent\zhuan\zihao{6}\fzbyks 傳“自來”至“政亂”。正義曰：“攘”、“竊”同文，則“攘”是竊類。\CJKunderwave{釋詁}云：“攘，因也。”是因其自來而取之名“攘”也。\CJKunderwave{說文}云：“犧，宗廟牲也。”\CJKunderwave{曲禮}云：“天子以犧牛。”天子祭牲必用純色，故知“色純曰犧”也。\CJKunderwave{周禮}：“牧人掌牧六牲，以供祭祀之牲牷。”以“牷”為言，必是體全具也，故“體完曰牷”。經傳多言“三牲”，知“牲”是牛羊豕也。以“犧”、“牷”、“牲”三者既為俎實，則“用”者簠簋之實，謂黍稷稻粱，故云“器實曰用”，謂粢盛也。\CJKunderwave{禮}“天曰神，地曰祗”，舉天地則人鬼在其間矣，故總雲“盜天地宗廟牲用”也。訓“將”為行，“相容行食之”謂所司相通容，使盜者得行盜而食之。大祭祀之物，物之重者，盜而無罪，言政亂甚也。漢魏以來著律皆云：“敢盜郊祀宗廟之物，無多少皆死。”為特重故也。 \par}

{\noindent\zhuan\zihao{6}\fzbyks 傳“下視”至“懈怠”。正義曰：箕子身為三公，下觀世俗，故云“下視殷民”。“所用治者”謂卿士已下是治民之官也。以紂暴虐，務稱上旨，“皆重賦傷民”。民既傷矣，則以上為讎，\CJKunderwave{泰誓}所謂“虐我則讎”是也。重斂民財，乃是“聚斂怨讎之道”。既為重斂，而又亟行暴虐。亟，急也。急行暴虐，欲以威民，乃是“自召敵讎”。勤行虐政，是“不懈怠”也。 \par}

商其淪喪，我罔為臣僕。詔王子出迪。\footnote{商其沒亡，我二人無所為臣僕，欲以死諫紂。我教王子出，合於道。臣僕，一本無臣字。}我舊雲刻子，王子弗出，我乃顛隮。\footnote{刻，病也。我久知子賢，言於帝乙。病立子，帝乙不肯。病子不得立，則宜為殷後者子。今若不出逃難，我殷家宗廟乃隕墜無主。舊云，馬云：“言也。”刻音克，馬云：“侵刻也。”難，乃旦反。}自靖。人自獻於先王，\footnote{各自謀行其志，人人自獻達於先王，以不失道。靖，馬本作清，謂潔也。}我不顧行遁。”\footnote{言將與紂俱死，所執各異，皆歸於仁，明君子之道，出處語默非一途。顧音故，徐音鼓。}

{\noindent\zhuan\zihao{6}\fzbyks 傳“商其”至“於道”。正義曰：“有災”與“淪喪”一事，而重出文者，上言“商今其有災，我興受其敗”，逆言災雖未至，至則己必受禍；此言“商其淪喪，我罔為臣僕”，豫言殷滅之後，言己不事異姓，辭有二意，故重出其文。我無所為臣僕，言不能與人為臣僕，必欲以死諫紂。但箕子之諫，值紂怒不甚，故得不死耳。“我教王子出，合於道”，保全身命，終為殷後，使宗廟有主，享祀不絕，是合其道也。 \par}

{\noindent\zhuan\zihao{6}\fzbyks 傳“刻病”至“無主”。正義曰：“刻”者，傷害之義，故為病也。\CJKunderwave{呂氏春秋·仲冬紀}云：“紂之母生微子啟與仲衍，其時猶尚為妾，改而為妻後生紂。紂之父欲立微子啟為太子，太史據法而爭，曰：‘有妻之子，不可立妾之子。’故立紂為後。”於時箕子蓋謂請立啟而帝乙不聽，今追恨其事，我久知子賢，言於帝乙，欲立子為太子，而帝乙不肯，我病子不得立，則宜為殷後。 \par}

{\noindent\zhuan\zihao{6}\fzbyks 傳“言將”至“一途”。正義曰：不肯遁以求生，“言將與紂俱死”也。或去或留，所執各異，皆歸於仁。\CJKunderline{孔子}稱“殷有三仁焉”，是“皆歸於仁”也。\CJKunderwave{易·繫辭}曰：“君子之道，或出或處，或默或語。”是“非一途”也。何晏云：“仁者愛人，三人行異而同稱仁者，以其俱在憂亂寧民。” \par}

{\noindent\shu\zihao{5}\fzkt “父師”至“行遁”。正義曰：父師亦順其事而報微子曰：“王子，今天酷毒下災，生此昏虐之君，以荒亂殷之邦國。紂既沈湎，四方化之，皆起而沈湎酗醟於酒，不可如何。小人皆自放恣，乃無所。上不畏天災，下不畏賢人，違戾其耇老之長與舊有爵位致仕之賢人。今殷民乃攘竊祭祀神祗之犧牷牲用，以相通容，行取食之，無災罪之者。”盜天地大祀之物用而不得罪，言政亂甚也。“我又下視殷民，所用為治者，皆讎怨斂聚之道”也。言重賦傷民，民以在上為讎，重賦乃是斂讎也。“既為重賦，又急行暴虐，此所以益招民怨，是乃自召敵讎不懈怠也。上下各有罪，合於一紂之身”。言紂化之使然也。“故使民多瘠病，而無詔救之者。商今其有滅亡之災，我起而受其敗。商其沒亡喪滅，我無所為人臣僕”。言不可別事他人，必欲諫取死也。“我教王子出奔於外，是道也。我久雲子賢，言於帝乙，欲立子，不肯。我乃病傷子不得立為王，則宜終為殷後。若王子不出，則我殷家宗廟乃隕墜無主”。既勸之出，即與之別云：“各自謀行其志，人人各自獻達於先王，我不顧念行遁之事。”明期與紂俱死。 \par}

%%% Local Variables:
%%% mode: latex
%%% TeX-engine: xetex
%%% TeX-master: "../Main"
%%% End:

%% -*- coding: utf-8 -*-
%% Time-stamp: <Chen Wang: 2024-04-02 11:42:41>

% {\noindent \zhu \zihao{5} \fzbyks } -> 注 (△ ○)
% {\noindent \shu \zihao{5} \fzkt } -> 疏


\part{周書}


\chapter{卷十一}


\section{泰誓上第一【偽】}

 {\noindent\zhuan\zihao{6}\fzbyks 案：孫星衍之\CJKunderwave{尚書今古文註疏}中\CJKunderwave{泰誓}與此文不同。 \par}

惟十有一年，武王伐殷。\footnote{周自虞芮質厥成，諸侯並附，以為受命之年。至九年而文王卒，武王三年服畢，觀兵孟津，以卜諸侯伐紂之心。諸侯僉同，乃退以示弱。芮，如銳反。虞、芮，二國名。僉，七廉反。}一月戊午，師渡孟津，\footnote{十三年正月二十八日，更與諸侯期而共伐紂。孟津，地名也。}作\CJKunderwave{泰誓}三篇。\footnote{渡津乃作。}


{\noindent\zhuan\zihao{6}\fzbyks 傳“周自”至“示弱”。正義曰：\CJKunderwave{武成}篇云：“我文考文王，誕膺天命，以撫方夏。惟九年，大統未集。”則文王以九年而卒也。\CJKunderwave{無逸}稱文王“享國五十年”至嗣位至卒非徒九年而已。知此十一年者，文王改稱元年，至九年而卒，至此年為十一年也。\CJKunderwave{詩}云：“虞芮質厥成。”\CJKunderwave{毛傳}稱“天下聞虞芮之訟息，歸周者四十餘國”，故知“周自虞芮質厥成，諸侯並附，以為受命之年。至九年而文王卒”，至此十一年，武王居父之喪“三年服畢”也。案\CJKunderwave{周書}云：“文王受命九年，惟暮春在鎬，召太子發作文傳。”其時猶在，但未知崩月。就如暮春即崩，武王服喪至十一年三月大祥，至四月觀兵，故今文\CJKunderwave{泰誓}亦云“四月觀兵”也。知此十一年非武王即位之年者，\CJKunderwave{大戴禮}雲“文王十五而生武王”，則武王少文王十四歲也。\CJKunderwave{禮記·文王世子}云：“文王九十七而終，武王九十三而終。”計其終年，文王崩時武王已八十三矣。八十四即位，至九十三而崩，適滿十年，不得以十三年伐紂。知此十一年者，據文王受命而數之。必繼文王年者，為其卒父業故也。緯候之書言受命者，謂有黃龍玄龜白魚赤雀負圖銜書以命人主，其言起於漢哀平之世，經典無文焉，孔時未有此說。\CJKunderwave{咸有一德}傳云：“所徵無敵謂之受天命。”此傳云：“諸侯並附，以為受命之年。”是孔解受命皆以人事為言，無瑞應也。\CJKunderwave{史記}亦以斷虞芮之訟為受命元年，但彼以文王受命七年而崩，不得與孔同耳。三年之喪，二十五月而畢，故九年文王卒，至此一年服畢。此經武王追陳前事，云：“肆予小子發，以爾友邦冢君，觀政於商。”是十一年伐殷者，止為觀兵孟津，以卜諸侯伐紂之心，言“於商”，知亦至孟津也。 \par}

{\noindent\zhuan\zihao{6}\fzbyks 傳“十三年正月”至“伐紂”。正義曰：以“一月戊午”，乃是作誓月日。經言“十三年春，大會於孟津”，又云“戊午，次於河朔”，知此“一月戊午”是十三年正月戊午日，非是十一年正月也。序不別言“十三年”，而以“一月”接“十一年”下者，序以觀兵至而即還，略而不言月日，誓則經有“年”有“春”，故略而不言“年春”，止言“一月”，使其互相足也。戊午是二十八日，以歷推而知之，據經亦有其驗。\CJKunderwave{漢書·律曆志}載舊說云：“死魄，朔也。生魄，望也。”\CJKunderwave{武成}篇說此伐紂之事云：“惟一月壬辰，旁死魄。”則壬辰近朔而非朔，是為月二日也。二日壬辰，則此月辛卯朔矣。以次數之，知戊午是二十八日也。不言“正月”而言“一月”者，以\CJKunderwave{武成}經言“一月”，故此序同之。\CJKunderwave{武成}所以稱“一月”者，\CJKunderwave{易·革卦}彖曰：“湯武革命，順乎天而應乎人。”象曰：“君子以治歷明時。”然則改正治歷，必自武王始矣。武王以殷之十二月發行，正月四日殺紂，既入商郊，始改正朔，以殷之正月為周之二月。其初發時猶是殷之十二月，未為周之正月，改正在後，不可追名為“正月”，以其實是周之一月，故史以“一月”名之。顧氏以為“古史質，或雲正月，或雲一月，不與\CJKunderwave{春秋}正月同”，義或然也。\CJKunderwave{易緯}稱“文王受命，改正朔，布王號於天下”。\CJKunderline{鄭玄}依而用之，言文王生稱王，已改正。然天無二日，民無二王，豈得殷紂尚在而稱周王哉？若文王身自稱王，已改正朔，則是功業成矣，武王何得雲“大勳未集”，欲卒父業也？\CJKunderwave{禮記大傳}云：“牧之野，武王之大事也。既事而退，追王大王亶父、王季歷、文王昌。”是追為王，何以得為文王身稱王，已改正朔也？\CJKunderwave{春秋}“王正月”謂周正月也，\CJKunderwave{公羊傳}曰：“王者孰謂？謂文王。”其意以正為文王所改。\CJKunderwave{公羊傳}漢初俗儒之言，不足以取正也。\CJKunderwave{春秋}之“王”，自是當時之王，非改正之王。晉世有王愆期者，知其不可，注\CJKunderwave{公羊}以為春秋制，文王指\CJKunderline{孔子}耳，非周昌也。\CJKunderwave{文王世子}稱武王對文王云：“西方有九國焉，群王其終撫諸。”呼文王為“王”，是後人追為之辭，其言未必可信，亦非實也。 \par}

{\noindent\zhuan\zihao{6}\fzbyks 傳“渡津乃作”。正義曰：“孟”者，河北地名，\CJKunderwave{春秋}所謂“向盟”是也。於孟地置津，謂之“孟津”，言師渡孟津，乃作\CJKunderwave{泰誓}，知三篇皆“渡津乃作”也。然則中篇獨言“戊午，次於河朔”者，三篇皆河北乃作，分為三篇耳。上篇未次時作，故言“十三年春”。中篇既次乃作，故言“戊午”之日。下篇則明日乃作，言“時厥明”。各為首引，故文不同耳。\CJKunderwave{尚書}遭秦而亡，漢初不知篇數，武帝時有太常蓼侯孔臧者，安國之從兄也，與安國書云：“時人惟聞\CJKunderwave{尚書}二十八篇，取象二十八宿，謂為信然，不知其有百篇也。”然則漢初惟有二十八篇，無\CJKunderwave{泰誓}矣。後得偽\CJKunderwave{泰誓}三篇，諸儒多疑之。馬融\CJKunderwave{書序}曰：“\CJKunderwave{泰誓}後得，案其文似若淺露。又云：‘八百諸侯，不召自來，不期同時，不謀同辭。’及‘火復於上，至於王屋，流為雕，至五，以谷俱來。’舉火神怪，得無在子所不語中乎？又\CJKunderwave{春秋}引\CJKunderwave{泰誓}曰：‘民之所欲，天必從之。’\CJKunderwave{國語}引\CJKunderwave{泰誓}曰：‘朕夢協朕卜，襲於休祥，戎商必克。’\CJKunderwave{孟子}引\CJKunderwave{泰誓}曰：‘我武惟揚，侵於之疆，取彼兇殘，我伐用張，於湯有光。’\CJKunderwave{孫卿}引\CJKunderwave{泰誓}曰：‘獨夫受。’\CJKunderwave{禮記}引\CJKunderwave{泰誓}曰：‘予克受，非予武，惟朕文考無罪。受克予，非朕文考有罪，惟予小子無良。’今文\CJKunderwave{泰誓}，皆無此語。吾見書傳多矣，所引\CJKunderwave{泰誓}而不在\CJKunderwave{泰誓}者甚多，弗復悉記，略舉五事以明之亦可知矣。”王肅亦云：“\CJKunderwave{泰誓}近得，非其本經。”馬融惟言後得，不知何時得之。\CJKunderwave{漢書}婁敬說高祖云：“武王伐紂，不期而會盟津之上者八百諸侯。”偽\CJKunderwave{泰誓}有此文，不知其本出何書也。武帝時董仲舒對策云：“\CJKunderwave{書}曰：‘白魚入於王舟，有火入於王屋，流為烏。周公曰：“復哉！復哉！”’”今引其文，是武帝之時已得之矣。李顒集註\CJKunderwave{尚書}，於偽\CJKunderwave{泰誓}篇每引“孔安國曰”，計安國必不為彼偽書作傳，不知顒何由為此言。梁王兼而存之，言“本有兩\CJKunderwave{泰誓}，古文\CJKunderwave{泰誓}伐紂事，聖人選為\CJKunderwave{尚書}。今文\CJKunderwave{泰誓}觀兵時事，別錄之以為\CJKunderwave{周書}”，此非辭也。彼偽書三篇，上篇觀兵時事，中下二篇亦伐紂時事，非盡觀兵時事也。且觀兵示弱即退，復何誓之有？設有其誓，不得同以\CJKunderwave{泰誓}為篇名也。 \par}

{\noindent\shu\zihao{5}\fzkt “惟十”至“三篇”。正義曰：惟文王受命十有一年，武王服喪既畢，舉兵伐殷，以卜諸侯伐紂之心。雖諸侯僉同，乃退以示弱。至十三年紂惡既盈，乃復往伐之。其年一月戊午之日，師渡孟津，王誓以戒眾。史敘其事，作\CJKunderwave{泰誓}三篇。 \par}

泰誓\footnote{大會以誓眾。}

{\noindent\shu\zihao{5}\fzkt 傳“大會以示眾”。正義曰：經云：“大會於孟津”，知名曰\CJKunderwave{泰誓}者，其“大會以示眾”也。王肅云：“武王以大道誓眾。”肅解彼偽文，故說謬耳。\CJKunderwave{湯誓}指湯為名，此不言“武誓”而別立名者，以武誓非一，故史推義作名\CJKunderwave{泰誓}，見大會也。\CJKunderwave{牧誓}舉戰地，時史意也。顧氏以為：“泰者，大之極也。猶如天子諸侯之子曰太子，天子之卿曰太宰，此會中之大，故稱\CJKunderwave{泰誓}也。” \par}

惟十有三年春，大會於孟津。\footnote{三分二諸侯，及諸戎狄。此周之孟春。“惟十有三年春”或作“十有一年”，後人妄看序文輒改之。}

{\noindent\zhuan\zihao{6}\fzbyks 傳“三分”至“孟春”。正義曰：\CJKunderwave{論語}稱“三分天下有其二”，中篇言“群后以師畢會”，則周之所有諸國皆集。\CJKunderwave{牧誓}所呼有“庸、蜀、羌、髳、微、盧、彭、濮人”，知此大會，謂三分有二之諸侯及諸戎狄皆會也。序言“一月”，知此春是“周之孟春”，謂建子之月也。知者案\CJKunderwave{三統曆}以殷之十二月武王發師，至二月甲子咸劉商王紂，彼十二月即周之正月建子之月也。 \par}

{\noindent\shu\zihao{5}\fzkt “惟十”至“孟津”。正義曰：此三篇俱是孟津之上大告諸國之君，而發首異者，此見大會誓眾，故言“大會於孟津”；中篇徇師而誓，故言“以師畢會”；下篇王更徇師，故言“大巡六師”，皆史官觀事而為作端緒耳。 \par}

王曰：“嗟！我友邦冢君，越我御事庶士，明聽誓。\footnote{冢，大。御，治也。友諸侯，親之。稱大君，尊之。下及我治事眾士，大小無不皆明聽誓。}

{\noindent\zhuan\zihao{6}\fzbyks 傳“冢大”至“聽誓”。正義曰：“冢，大”，\CJKunderwave{釋詁}文。侍御是治理之事，故通訓“御”為治也。同志為“友”，天子友諸侯，親之也。\CJKunderwave{牧誓}傳曰：“言志同滅紂。”令總呼國君皆為大君，尊之也。“下及治事眾士”，謂國君以外卿大夫及士諸掌事者。“大小無不皆明聽誓”，自士以上皆總戒之也。 \par}

惟天地萬物父母，惟人萬物之靈。\footnote{生之謂父母。靈，神也。天地所生，惟人為貴。}

{\noindent\zhuan\zihao{6}\fzbyks 傳“生之”至“為貴”。正義曰：萬物皆天地生之，故謂天地為父母也。\CJKunderwave{老子}云：“神得一以靈。”“靈”、“神”是一，故“靈”為神也。\CJKunderwave{禮運}云：“人者天地之心，五行之端也，食味別聲被色而生者也。”言人能兼此氣性，餘物則不能然。故\CJKunderwave{孝經}云：“天地之性人為貴。”此經之意，天地是萬物之父母，言天地之意，欲養萬物也。人是萬物之最靈，言其尤宜長養也。紂違天地之心而殘害人物，故言此以數之，與下句為首引也。 \par}

亶聰明，作元后，元后作民父母。\footnote{人誠聰明，則為大君，而為眾民父母。亶，丁但反。}今商王受，弗敬上天，降災下民。沈湎冒色，敢行暴虐，\footnote{沈湎嗜酒，冒亂女色，敢行酷暴，虐殺無辜。湎，面善反。冒，莫報反，注下同。嗜，市志反，\CJKunderwave{切韻}常利反。酷，苦毒反。}

{\noindent\shu\zihao{5}\fzkt 傳“沈湎”至“無辜”。正義曰：人被酒困，若沈於水，酒變其色，湎然齊同，故“沈湎”為嗜酒之狀。“冒”訓貪也，亂女色，荒也。“酷”解經之“暴”，“殺”辭經之“虐”，皆果敢為之。案\CJKunderwave{說文}云：“酷，酒厚味也。”酒味之厚必嚴烈，人之暴虐與酒嚴烈同，故謂之“酷”。 \par}

罪人以族，官人以世，\footnote{一人有罪，刑及父母兄弟妻子，言淫濫。官人不以賢才，而以父兄，所以政亂。}

{\noindent\zhuan\zihao{6}\fzbyks 傳“一人”至“政亂”。正義曰：秦政酷虐，有三族之刑，謂非止犯者之身，乃更上及其父，下及其子。經言“罪人以族”，故以三族解之。父母，前世也；兄弟及妻，當世也；子孫，後世也。一人有罪，刑及三族，言淫濫也。古者臣有大功，乃得繼世在位。而紂之官人，不以賢才，而以父兄，已濫受寵，子弟頑愚亦用，不堪其職，所以政亂。“官人以世”，惟當用其子耳，而傳兼言“兄”者，以紂為惡，或當因兄用弟，故以“兄”協句耳。 \par}

惟宮室、臺榭、陂池、侈服，以殘害於爾萬姓。\footnote{土高曰臺，有木曰榭，澤障曰陂，停水曰池，侈謂服飾過制。言匱民財力為奢麗。榭，\CJKunderwave{爾雅}云：“有木曰榭。”本又作謝。陂，彼皮反。障，之亮反。匱，其愧反。}

{\noindent\zhuan\zihao{6}\fzbyks 傳“土高”至“奢麗”。正義曰：\CJKunderwave{釋宮}云：“宮謂之室，室謂之宮。”李巡曰：“所以古今通語，明實同而兩名。”此傳不解“宮室”，義當然也。\CJKunderwave{釋宮}又云：“闍謂之臺。有木者謂之榭。”李巡曰：“臺積土為之，所以觀望也。臺上有屋謂之榭。”又云：“無室曰榭，四方而高曰臺。”孫炎曰：“榭但有堂也。”郭璞曰：“榭即今之堂堭也。”然則榭是臺上之屋，歇前無室，今之廳是也。\CJKunderwave{詩}云：“彼澤之陂。”\CJKunderwave{毛傳}云：“陂，澤障也。”障澤之水，使不流溢謂之“陂”，停水不流謂之“池”。“侈”亦奢也，謂依服採飾過於制度，言匱竭民之財力為奢麗也。顧氏亦云：“華侈服飾。”二劉以為宮室之上而加侈服。據孔傳雲“服飾過制”，即謂人之服飾，二劉之說非也。\CJKunderwave{殷本紀}云：“紂厚賦稅以實鹿臺之錢，而盈巨橋之粟。益收狗馬奇物，充牣宮室。益廣沙丘苑臺，多聚野獸飛鳥置其中。大聚樂戲於沙丘，以酒為池，懸肉為林，使男女倮相逐其間。”說紂奢侈之事，書傳多矣。 \par}

焚炙忠良，刳剔孕婦。\footnote{忠良無罪焚炙之，懷子之婦刳剔視之。言暴虐。刳，口胡反。剔，他歷反。孕，以證反，徐養證反。}

{\noindent\zhuan\zihao{6}\fzbyks 傳“忠良”至“暴虐”。正義曰：“焚炙”俱燒也，“刳剔”謂割剝也。\CJKunderwave{說文}云：“刳，刲也。”今人去肉至骨謂之“剔去”，是則亦\CJKunderwave{剔}之義也。武王以此數紂之惡，必有忠良被炙，孕婦被刳，不知其姓名為誰也。\CJKunderwave{殷本紀}云，紂為長夜之飲。時諸侯或叛，妲己以為罰輕，紂欲重刑，乃為熨斗，以火燒之然，使人舉輒爛其手，不能勝。紂怒，乃更為銅柱，以膏塗之，亦加於炭火之上，使有罪者緣之，足滑跌墜入中。紂與妲己以為大樂，名曰炮烙之刑。是紂焚炙之事也。後文王獻洛西之地，赤壤之田方千里，請紂除炮烙之刑，紂許之。皇甫謐作\CJKunderwave{帝王世紀}亦云然。謐又云：“紂剖比干妻，以視其胎。”即引此為“刳剔孕婦”也。 \par}

皇天震怒，命我文考，肅將天威，大勳未集。\footnote{言天怒紂之惡，命文王敬行天罰，功業未成而崩。}肆予小子發，以爾友邦冢君，觀政於商。\footnote{父業未就之故，故我與諸侯觀紂政之善惡。謂十一年自孟津還時。}惟受罔有悛心，乃夷居弗事上帝神祗，遺厥先宗廟弗祀。\footnote{悛，改也，言紂縱惡無改心，平居無故廢天地百神宗廟之祀。慢之甚。悛，七全反。}

{\noindent\zhuan\zihao{6}\fzbyks 傳“悛改”至“之甚”。正義曰：\CJKunderwave{左傳}稱“長惡不悛”，“悛”是退前創改之義，故為改也。觀政於商，計當恐怖，言紂縱惡無改悔之心，平居無故不事神祗，是紂之大惡。“上帝”，舉其尊者，謂諸神悉皆不事，故傳言“百神”以該之。“不事”亦是“不祀”，別言“遺厥先宗廟弗祀”，遺棄祖父，言其慢之甚也。 \par}

犧牲粢盛，既於兇盜。\footnote{兇人盡盜食之，而紂不罪。粢音諮，黍稷曰粢。盛音成，在器曰盛。}乃曰：‘吾有民有命。’罔懲其侮。\footnote{紂言：“吾所以有兆民，有天命。”故群臣畏罪不爭，無能止其慢心。懲，直承反。爭，爭鬥之爭。}天佑下民，作之君，作之師。\footnote{言天佑助下民，為立君以政之，為立師以教之。為，於偽反。}惟其克相上帝，寵綏四方。\footnote{當能助天寵安天下。相，息亮反。}有罪無罪，予曷敢有越厥志？\footnote{越，遠也。言己志欲為民除惡，是與否，不敢遠其志。否，方有反。}

{\noindent\zhuan\zihao{6}\fzbyks 傳“言天”至“教之”。正義曰：眾民不能自治，立君以治之。立君治民，乃是天意，言天佑助下民為立君也。治民之謂“君”，教民之謂“師”，君既治之，師又教之，故言“作之君，作之師”，“師”謂君與民為師，非謂別置師也。 \par}

{\noindent\zhuan\zihao{6}\fzbyks 傳“當能”至“天下”。正義曰：天愛下民，為立君立師者，當能佑助天意，寵安天下，不奪民之財力，不妄非理刑殺，是助天寵愛民也。 \par}

{\noindent\zhuan\zihao{6}\fzbyks 傳“越遠”至“其志”。正義曰：“越”者，逾越超遠之義，故為遠也。武王伐紂，內實為民除害，外則以臣伐君，故疑其有罪與無罪。“言己志欲為民除害，無問是之與否，不敢遠其志”，言己本志欲伐，何敢遠本志，舍而不伐也？ \par}

{\noindent\shu\zihao{5}\fzkt “天佑”至“厥志”。正義曰：已上數紂之罪，此言伐紂之意。上天佑助下民，不欲使之遭害，故命我為之君上，使臨政之；為之師保，使教誨之。為人君為人師者，天意如此，不可違天。我今惟其當能佑助上，天寵安四方之民，使民免於患難。今紂暴虐，無君師之道，故今我往伐之。不知伐罪之事，為有罪也？為無罪也？不問有罪無罪，志在必伐，我何敢有遠其本志而不伐之？ \par}

“同力度德，同德度義。\footnote{力鈞則有德者勝，德鈞則秉義者強。揆度優劣，勝負可見。度，徒洛反，下注同。}

{\noindent\zhuan\zihao{6}\fzbyks 傳“力鈞”至“可見”。正義曰：“德”者得也，自得於心。“義”者宜也，動合自宜。但德在於身，故言“有德”；義施於行，故言秉執。武王志在養民，動為除害，有君人之明德，執利民之大義，與紂無者為敵，雖未交兵，揆度優劣，勝負可見。示以必勝之道，令士眾勉力而戰也。 \par}

受有臣億萬，惟億萬心。\footnote{人執異心，不和諧。億，十萬曰億。}予有臣三千，惟一心。\footnote{三千一心，言欲同。}商罪貫盈，天命誅之。予弗順天，厥罪惟鈞。\footnote{紂之為惡，一以貫之，惡貫已滿，天畢其命。今不誅紂，則為逆天，與紂同罪。貫，古亂反。}

{\noindent\zhuan\zihao{6}\fzbyks 傳“紂之”至“同罪”。正義曰：紂之為惡，如物在繩索之貫，一以貫之，其惡貫已滿矣。物極則反，天下欲畢其命，故上天命我誅之。今我不誅紂，則是逆天之命，無恤民之心，是我與紂同罪矣。猶如\CJKunderwave{律}“故縱者與同罪”也。 \par}

予小子夙夜祗懼，受命文考，類於上帝，宜於冢土，以爾有眾，厎天之罰。\footnote{祭社曰宜。冢土，社也。言我畏天之威，告文王廟，以事類告天祭社，用汝眾致天罰於紂。類，師祭名。冢，中勇反。厎，之履反。}

{\noindent\zhuan\zihao{6}\fzbyks 傳“祭社”至“於紂”。正義曰：\CJKunderwave{釋天}引\CJKunderwave{詩}云：“乃立冢土，戎醜攸行。”即云：“起大事，動大眾，必先有事乎社而後出，謂之宜。”孫炎曰：“宜，求見福祐也。”是“祭社曰宜”。“冢”訓大也，社是土神，故“冢土，社也”。\CJKunderwave{毛詩傳}云：“冢土，大社也。”“受命文考”是告廟以行，故為“告文王廟”也。\CJKunderwave{毛詩}云：“天子將出，類乎上帝，宜乎社，造乎禰。”此“受命文考”即是“造乎禰”也。\CJKunderwave{王制}以神尊卑為次，故先言“帝”、“社”，後言“禰”，此以廟是己親，若言家內私義，然後告天，故先言“受命文考”，而後言“類於上帝”。\CJKunderwave{舜典}“類於上帝”傳云：“告天及五帝。”此“以事類告天”，亦當如彼也。罰紂是天之意，故“用汝眾致天罰於紂”也。 \par}

天矜於民，民之所欲，天必從之。\footnote{矜，憐也。言天除惡樹善與民同。從，才容反。}爾尚弼予一人，永清四海。\footnote{穢惡除，則四海長清。}時哉弗可失！”\footnote{言今我伐紂，正是天人合同之時，不可違失。}

\section{泰誓中第二【偽】}


惟戊午，王次於河朔。\footnote{次，止也。戊午渡河而誓，既誓而止於河之北。}

{\noindent\zhuan\zihao{6}\fzbyks 傳“次止”至“之北”。正義曰：“次”是止舍之名，\CJKunderwave{穀梁傳}亦云：“次，止也。”序雲“一月戊午，師渡孟津”，則師以戊午日渡也。此戊午日次於河朔，則是師渡之日次止也。上篇是渡河而誓，未及止舍而先誓之，此“次於河朔”者，是“既誓而止於河之北”也。莊三年\CJKunderwave{左傳}例云：“凡師一宿為舍，再宿為信，過宿為次。”此“次”直取止舍之義，非\CJKunderwave{春秋}三日之例也。何則？商郊去河四百餘里，戊午渡河，甲子殺紂，相去才六日耳。是今日次訖又誓，明日誓訖即行，不容三日止於河旁也。 \par}

群后以師畢會，\footnote{諸侯盡會次也。}王乃徇師而誓。曰：“嗚呼！西土有眾，咸聽朕言。\footnote{徇，循也。武王在西，故稱西土。徇，似俊反，\CJKunderwave{字詁}云：“徇，巡也。”}

{\noindent\zhuan\zihao{6}\fzbyks 傳“徇循”至“西土”。正義曰：\CJKunderwave{說文}云：“徇，疾也。循，行也。”“徇”是疾行之意，故以“徇”為循也。下篇“大巡六師”，義亦然也。此誓總戒眾軍，武王國在西偏，此師皆從西而來，故稱“西土”。 \par}

我聞吉人為善，惟日不足。兇人為不善，亦惟日不足。\footnote{言吉人竭日以為善，兇人亦竭日以行惡。竭，苦曷反，又苦蓋反。}今商王受，力行無度，\footnote{行無法度，竭日不足，故曰力行。}播棄黎老，暱比罪人。\footnote{鮐背之耇稱黎老，布棄不禮敬。暱近罪人，謂天下逋逃之小人。黎，力私反，又力兮反。暱，女乙反。比，毗志反。鮐，他來反，又音怡，魚名。逋，布吳反。}

{\noindent\zhuan\zihao{6}\fzbyks 傳“鮐背”至“小人”。正義曰：\CJKunderwave{釋詁}云：“鮐背、耇、老，壽也。”舍人曰：“鮐背，老人氣衰，皮膚消瘠，背若鮐魚也。”孫炎曰：“耇，面凍梨色似浮垢也。”然則老人背皮似鮐，面色似梨，故“鮐背之耇”稱“梨老”。傳以“播”為布。布者，遍也，言遍棄之，不禮敬也。“暱，近”，\CJKunderwave{釋詁}文。孫炎曰：“暱，親近也。”\CJKunderwave{牧誓}數紂之罪云：“四方之多罪逋逃，是崇是長，是信是使。”知紂所親近罪人，“謂天下逋逃之小人”也。 \par}

淫酗肆虐，臣下化之，\footnote{過酗縱虐，以酒成惡，臣下化之。言罪同。酗，況付反。}

{\noindent\zhuan\zihao{6}\fzbyks 傳“過酗”至“罪同”。正義曰：“酗”是酒怒，“淫酗”共文，則“淫”非女色，故以“淫”為過，言飲酒過多也。“肆”是放縱之意，酒過則酗，縱情為虐。以酒成此暴虐之惡，臣下化而為之，由紂惡而臣亦惡，言君臣之罪同也。 \par}

朋家作仇，脅權相滅。無辜籲天，穢德彰聞。\footnote{臣下朋黨，自為仇怨，脅上權命，以相誅滅。籲，呼也。民皆呼天告冤無辜，紂之穢德彰聞天地。言罪惡深。脅，虛業反。籲音喻。穢，於廢反。}

{\noindent\zhuan\zihao{6}\fzbyks 傳“臣下”至“惡深”。正義曰：“脅上”謂紂既昏迷，朝無綱紀，奸宄之臣，脅於在下，假用在上之權命，脅之更相誅滅也。 \par}

{\noindent\shu\zihao{5}\fzkt “朋家”至“彰聞”。正義曰：小人好忿，天性之常，化紂淫酗，怨怒無已。臣下朋黨，共為一家，與前人並作仇敵，脅上權命，以相滅亡。無罪之人，怨嗟呼天，紂之穢惡之德，彰聞天地。言其罪惡深也。 \par}

“惟天惠民，惟闢奉天。\footnote{言君天下者當奉天以愛民。闢，必亦反。}

有夏桀，弗克若天，流毒下國。\footnote{桀不能順天，流毒虐於下國萬民。言兇害。}天乃佑命\CJKunderline{成湯}，降黜夏命。\footnote{言天助湯命，使下退桀命。}惟受罪浮於桀。\footnote{浮，過。}

{\noindent\zhuan\zihao{6}\fzbyks 傳“浮過”。正義曰：物在水上謂之浮，“浮”者高之意，故為過也。桀罪已大，紂又過之，言紂惡之甚，故下句說其過桀之狀。案\CJKunderwave{夏本紀}及\CJKunderwave{帝王世紀}云：“諸侯叛桀，關龍逢引皇圖而諫，桀殺之。\CJKunderline{伊尹}諫桀，桀曰：‘天之有日，如吾之有民，日亡吾乃亡矣。’”是桀亦“賊虐諫輔，謂己有天命”。而云過於桀者，\CJKunderwave{殷本紀}云：“紂剖比干觀其心”，桀殺龍逢，無剖心之事；又桀惟比之於日，紂乃詐命於天；又紂有炮烙之刑，又有刳胎斮脛之事，而桀皆無之，是紂罪過於桀也。 \par}

剝喪元良，賊虐諫輔。\footnote{剝，傷害也。賊，殺也。元，善之長。良，善。以諫輔紂，紂反殺之。喪，息浪反。長，丁丈反。}

{\noindent\zhuan\zihao{6}\fzbyks 傳“剝傷”至“殺之”。正義曰：\CJKunderwave{說文}云：“剝，裂也，一曰剝，割也。”裂與割俱是傷害之義也。殺人謂之“賊”，故“賊”為殺也。“元者，善之長”，\CJKunderwave{易}文言文。“良”之為善，書傳通訓也。“元良”俱善而雙舉之者，言其剝喪善中之善，為害大也。“以諫輔紂，紂反殺之”，即比干是也。上篇言“焚炙忠良”，與此經相類而復言此者，以殺害人為惡之大，故重陳之也。 \par}

謂己有天命，謂敬不足行，謂祭無益，謂暴無傷。\footnote{言紂所以罪過於桀。己音紀。}厥監惟不遠，在彼夏王。\footnote{其視紂罪，與桀同辜。言必誅之。}

{\noindent\zhuan\zihao{6}\fzbyks 傳“其視”至“誅之”。正義曰：紂罪過於桀，而言“與桀同辜”者，罪不過死，合死之罪同，言必誅也。 \par}

天其以予乂民，\footnote{用我治民。當除惡。}朕夢協朕卜，襲於休祥，戎商必克。\footnote{言我夢與卜俱合於美善，以兵誅紂必克之佔。}

{\noindent\zhuan\zihao{6}\fzbyks 傳“言我”至“之佔”。正義曰：夢者事之祥，人之精爽先見者也。吉凶或有其驗，聖王採而用之。我卜伐紂得吉，夢又戰勝。\CJKunderwave{禮記}稱“卜筮不相襲”，“襲”者，重合之義。訓“戎”為兵。夢卜俱合於美，是“以兵誅紂必克之佔”也。聖人逆知來物，不假夢卜，言此以強軍人之意耳。\CJKunderwave{史記·周本紀}云：“武王伐紂，卜，龜兆不吉，群公皆懼，惟太公強之。”太公\CJKunderwave{六韜}云：“卜戰，龜兆焦，筮又不吉，太公曰：‘枯骨朽蓍，不逾人矣。’”彼言“不吉”者，\CJKunderwave{六韜}之書，後人所作，\CJKunderwave{史記}又採用\CJKunderwave{六韜}，好事者妄矜太公，非實事也。 \par}

受有億兆夷人，離心離德。\footnote{平人，凡人也。雖多而執心用德不同。}

{\noindent\zhuan\zihao{6}\fzbyks 傳“平人”至“不同”。正義曰：昭二十四年\CJKunderwave{左傳}此文，服虔、杜預以“夷人”為夷狄之人。即如彼言，惟雲“億兆夷人”，則受率其旅若林，即曾無華夏人矣？故傳訓“夷”為平，平人為凡人，言其智慮齊，識見同。人數雖多，執心用德不同。“心”謂謀慮，“德”謂用行，智識既齊，各欲申意，故“心德不同”也。 \par}

予有亂臣十人，同心同德。\footnote{我治理之臣雖少而心德同。十人，周公旦、召公奭、太公望、畢公、榮公、太顛、閎夭、散宜生、南宮适及文母。治，直吏反。}

{\noindent\zhuan\zihao{6}\fzbyks 傳“我治”至“德同”。正義曰：\CJKunderwave{釋詁}云：“亂，治也。”故謂我治理之臣有十人也。十人皆是上智，咸識周是殷非，故人數雖少而心能同。同佐武王，欲共滅紂也。\CJKunderwave{論語}引此云：“予有亂臣十人。”而\CJKunderline{孔子}論之有一婦人焉，則十人之內其一是婦人，故先儒\CJKunderline{鄭玄}等皆以十人為文母、周公、太公、召公、畢公、榮公、太顛、宏夭、散宜生、南宮括也。 \par}

雖有周親，不如仁人。\footnote{周，至也。言紂至親雖多，不如周家之少仁人。}

{\noindent\zhuan\zihao{6}\fzbyks 傳“周至”至“仁人”。正義曰：\CJKunderwave{詩毛傳}亦以“周”為至，相傳為此訓也。武王三分天下有其二，則紂黨不多於周。但辭有激發，旨有抑揚，欲明多惡不如少善，故言“紂至親雖多，不如周家之少仁人”也。 \par}

“天視自我民視，天聽自我民聽。\footnote{言天因民以視聽，民所惡者天誅之。惡，烏路反，一音如字。}百姓有過，在予一人。\footnote{己能無惡於民，民之有過，在我教不至。}

{\noindent\shu\zihao{5}\fzkt “百姓有過，在予一人”。正義曰：言此者，以上雲民之所惡，天必誅之，己今有善，不為民之所惡，天必佑我。令教化百姓，若不教百姓，使有罪過，實在我一人之身。此“百姓”與下“百姓懍懍”，皆謂天下眾民也。 \par}

今朕必往，我武惟揚，侵於之疆，\footnote{揚，舉也。言我舉武事，侵入紂郊疆伐之。疆，居良反。}取彼兇殘，我伐用張，於湯有光。\footnote{桀流毒天下，湯黜其命。紂行兇殘之德，我以兵取之。伐惡之道張設，比於湯又有光明。}

{\noindent\zhuan\zihao{6}\fzbyks 傳“揚舉”至“伐之”。正義曰：\CJKunderwave{文王世子}論舉賢之法云：“或以事舉，或以言揚。”是“揚”、“舉”義同，故“揚”為舉也。於時猶在河朔，將欲行適商都，言我舉武事，侵入紂之郊疆，往伐之也。\CJKunderwave{春秋}之例有：“鐘鼓曰伐，無曰侵。”此實伐也，言“往侵”者，“侵”是入之意，非如\CJKunderwave{春秋}之例無鐘鼓也。 \par}

{\noindent\shu\zihao{5}\fzkt “今朕”至“有光”。正義曰：既與天下為任，則當為之除害，今我必往伐紂。我之武事惟於此舉之，侵紂之疆境，取彼為兇殘之惡者。若得取而殺之，是我伐兇惡之事用張設矣。湯惟放逐，我能擒取，是比於湯又益有光明。 \par}

勖哉，夫子！罔或無畏，寧執非敵。\footnote{勖，勉也。夫子謂將士。無敢有無畏之心，寧執非敵之志，伐之則克矣。將，子匠反，下篇注同。}

{\noindent\zhuan\zihao{6}\fzbyks 傳“勖勉”至“克矣”。正義曰：“勖，勉”，\CJKunderwave{釋詁}文。呼將士而誓之，知“夫子”是將士也。\CJKunderwave{老子}云：“禍莫大於輕敵。”故今將士“無敢有無畏之心”，令其必以前敵為可畏也。\CJKunderwave{論語}稱：“子路曰：‘子行三軍則誰與？’\CJKunderline{孔子}曰：‘必也臨事而懼。’”令軍士等不欲發意輕前人，寧執非敵之志，恐彼強多，非我能敵，執此志以伐之，則當克矣。 \par}

{\noindent\shu\zihao{5}\fzkt “勖哉”至“非敵”。正義曰：取得紂則功多於湯，宜勉力哉！“夫子”，將士等。呼將士令勉力也。以兵伐人，當臨事而懼，汝將土等無敢有無畏輕敵之心，寧執守似前人之強，非己能敵之志以伐之，如是乃可克矣。 \par}

百姓懍懍，若崩厥角。\footnote{言民畏紂之虐，危懼不安，若崩摧其角，無所容頭。○懍，力甚反。}

{\noindent\zhuan\zihao{6}\fzbyks 傳“言民”至“容頭”。正義曰：“懍懍”是怖懼之意，言民畏紂之虐，危懼不安，其志懍懍然。以畜獸為喻，民之怖懼若似畜獸崩摧其頭角然，無所容頭。顧氏云：“常如人之慾崩其角也，言容頭無地。”隱三年\CJKunderwave{穀梁傳}曰：“高曰崩，頭角之稱崩，體之高也。” \par}

嗚呼！乃一德一心，立定厥功，惟克永世。”\footnote{汝同心立功，則能長世以安民。}

\section{泰誓下第三【偽】}


時厥明，王乃大巡六師，明誓眾士。\footnote{是其戊午明日，師出以律，三申令之，重難之義。眾士，百夫長已上。○令，力政反。重，直用反。長，丁丈反。已音以。上，時掌反。}

{\noindent\zhuan\zihao{6}\fzbyks 傳“是其”至“已上”。正義曰：上篇未次而誓，故略言“大會”。中篇既次乃誓，為文稍詳，故言“以師畢會”。此篇最在其後，為文亦詳，故言“大巡六師”。巡繞周遍大其事，故稱“大”也。“師”者，眾也。天子之行，通以六師為言。於時諸侯盡會，其師不啻六也。“師出以律”，\CJKunderwave{易·師卦}初六爻辭也。“律”，法也。行師以法，即誓敕賞勸是也。禮成於三，故為三篇之誓。三度申重號令，為重慎艱難之義也。\CJKunderwave{孫子兵法}“三令五申之”，此誓三篇，亦為三令之事也。\CJKunderwave{牧誓}王所呼者，從上而下，至“百夫長”而止，知此“眾士”是“百夫長已上”也。 \par}

王曰：“嗚呼！我西土君子，天有顯道，厥類惟彰。\footnote{言天有明道，其義類惟明，言王所宜法則。}

{\noindent\zhuan\zihao{6}\fzbyks 傳“言天”至“法則”。正義曰：\CJKunderwave{孝經}云：“則天之明。”昭二十五年\CJKunderwave{左傳}云：“以象天明。”是治民之事，皆法天之道。天有尊卑之序，人有上下之節，三正五常，皆在於天，有其明道，此天之明道。“其義類惟明”，言明白可效，王者所宜法則之。將言商王不法天道，故先標二句於前。其下乃述商王違天之事，言其罪宜誅也。 \par}

今商王受，狎侮五常，荒怠弗敬。\footnote{輕狎五常之教，侮慢不行，大為怠惰，不敬天地神明。○惰，徒臥反。}

{\noindent\zhuan\zihao{6}\fzbyks 傳“輕狎”至“神明”。正義曰：\CJKunderline{鄭玄}\CJKunderwave{論語注}云：“狎，慣忽之。”言慣見而忽也，意與“侮”同，傳因文重而分之。“五常”即五典，謂父義、母慈、兄友、弟恭、子孝，五者人之常行，法天明道為之。輕狎五常之教，侮慢而不遵行之，是違天顯也。訓“荒”為大，大為怠惰。“不敬”謂“不敬天地神明”也。上篇雲“不事上帝神祇”，知此“不敬天地神明”也。\CJKunderwave{禮}云：“毋不敬。”傳舉“天地”以言，明每事皆不敬也。 \par}

自絕於天，結怨於民。\footnote{不敬天，自絕之。酷虐民，結怨之。}斮朝涉之脛，剖賢人之心，\footnote{冬月見朝涉水者，謂其脛耐寒，斬而視之。比干忠諫，謂其心異於人，剖而觀之。酷虐之甚。○斮，側略反，又士略反。朝，陟遙反。脛，戶定反。剖，普口反。耐，乃代反。}

{\noindent\zhuan\zihao{6}\fzbyks 傳“冬月”至“之甚”。正義曰：\CJKunderwave{釋器}云：“魚曰斮之。”樊光云：“斮，斫也。”\CJKunderwave{說文}云：“斮，斬也。”斬朝涉水之脛，必有所由，知冬月見朝涉水者，謂其脛耐寒，疑其骨髓有異，斬而視之。其事或當有所出也。\CJKunderwave{殷本紀}云：“微子既去，比干曰：‘為人臣者，不得不以死爭。’乃強諫。紂怒曰：‘吾聞聖人心有七竅。’遂剖比干，觀其心。”是紂謂比干心異於人，剖而觀之。言酷虐之甚。 \par}

作威殺戮，毒痡四海。\footnote{痡，病也。言害所及遠。○痡，徐音敷，又普吳反。}

{\noindent\zhuan\zihao{6}\fzbyks 傳“痡病”至“及遠”。正義曰：“痡，病”，\CJKunderwave{釋詁}文。紂之毒害，未必遍及夷狄，而云病四海者，言害所及者遠也。 \par}

崇信奸回，放黜師保，\footnote{回，邪也。奸邪之人，反尊信之。可法以安者，反放退之。○邪，似嗟反。}屏棄典刑，囚奴正士，\footnote{屏棄常法而不顧，箕子正諫而以為囚奴。}郊社不修，宗廟不享，作奇技淫巧以悅婦人。\footnote{言紂廢至尊之敬，營卑褻惡事，作過制技巧，以恣耳目之欲。○技，其綺反。褻，息列反。}

{\noindent\shu\zihao{5}\fzkt “郊社”至“婦人”。正義曰：“不修”謂不掃治也。“不享”謂不祭祀也。與上篇“不祀上帝神祇，遺厥先宗廟不祀”，其事一也，重言之耳。“奇技”謂奇異技能，“淫巧”謂過度工巧，二者本同，但“技”據人身，“巧”指器物為異耳。 \par}

上帝弗順，祝降時喪。\footnote{祝，斷也。天惡紂逆道，斷絕其命，故下是喪亡之誅。○喪，蘇浪反。斷，丁管反。惡，烏路反。}

{\noindent\zhuan\zihao{6}\fzbyks 傳“祝斷”。正義曰：哀十四年\CJKunderwave{公羊傳}云：“子路死，子曰：‘天祝予！’”何休云：“祝，斷也。”是相傳訓也。 \par}

爾其孜孜，奉予一人，恭行天罰。\footnote{孜孜，勸勉不怠。○孜音茲。}“古人有言曰：‘撫我則後，虐我則讎。’\footnote{武王述古言以明義，言非惟今紂惡。}獨夫受，洪惟作威，乃汝世讎。\footnote{言獨夫，失君道也。大作威殺無辜，乃是汝累世之讎。明不可不誅。}樹德務滋，除惡務本，\footnote{立德務滋長，去惡務除本。言紂為天下惡本。}肆予小子，誕以爾眾士殄殲乃讎。\footnote{言欲行除惡之義，絕盡紂。○殄，徒典反。纖，子廉反。}爾眾士其尚迪果毅，以登乃闢。\footnote{迪，進也。殺敵為果，致果為毅。登，成也，成汝君之功。○毅，牛既反。}

{\noindent\zhuan\zihao{6}\fzbyks 傳“迪進”至“之功”。正義曰：“迪，進”、“登，成”皆\CJKunderwave{釋詁}文。“殺敵為果，致果為毅”，宣二年\CJKunderwave{左傳}文。“果”謂果敢，“毅”謂強決。能殺敵人謂之為“果”，言能果敢以除賊。致此果敢是各為“毅”，言能強決以立功。皆言其心不猶豫也。軍法以殺敵為上，故勸令果毅成功也。 \par}

功多有厚賞，不迪有顯戮。\footnote{賞以勸之，戮以威之。}嗚呼！惟我文考，若日月之照臨，光於四方，顯於西土。\footnote{稱父以感眾也。言其明德充塞四方，明著岐周。}惟我有周，誕受多方。\footnote{言文王德大，故受眾方之國，三分天下而有其二。}予克受，非予武，惟朕文考無罪。\footnote{推功於父，言文王無罪於天下，故天佑之，人盡其用。}受克予，非朕文考有罪，惟予小子無良。”\footnote{若紂克我，非我父罪，我之無善之致。}

{\noindent\zhuan\zihao{6}\fzbyks 傳“若紂”至“之致”。正義曰：言克受乃是文王之功，若受克予非是文王之罪。而言“非我父罪，我之無善之致”者，其意言勝非我功，敗非父咎，崇孝罪己，以求眾心耳。 \par}

\section{牧誓第四}


武王戎車三百兩，\footnote{兵車，百夫長所載。車稱兩。一車步卒七十二人，凡二萬一千人，舉全數。○車音居。\CJKunderwave{釋名}云：“古者聲如居，所以居人也。今曰車，聲近舍，車舍也。”韋昭\CJKunderwave{辯釋名}云：“古皆尺遮反，從漢始有音居。”長，丁丈反。卒，子忽反。}虎賁三百人，\footnote{勇士稱也，若虎賁獸，言其猛也。皆百夫長。○賁音奔。稱，尺證反。}與受戰於牧野，作\CJKunderwave{牧誓}。

{\noindent\zhuan\zihao{6}\fzbyks 傳“兵車”至“全數”。正義曰：孔以“虎賁三百人”與戎車數同，王於誓時所呼有“百夫長”，因謂“虎賁”即是百夫之長。一人而乘一車，故云“兵車，百夫長所載”也。數車之法，一車謂之一兩。\CJKunderwave{詩}雲“百兩迓之”，是車稱兩也。\CJKunderwave{風俗通}說車有兩輪，故稱為兩。猶屨有兩隻，亦稱為兩。\CJKunderwave{詩}雲“葛屨五兩”即其類也。“一車步卒七十二人”，\CJKunderwave{司馬法}文也。車有七十二人，三百乘凡二萬一千人。計車有七十二人，三百乘當有二萬一千六百人，孔略六百而不言，故云“舉全數”。顧氏亦同此解。孔既用\CJKunderwave{司馬法}一車七十二人，又云“兵車，百夫長所載”，又下傳以百夫長為“卒帥”，是實領百人，非惟七十二人。依\CJKunderwave{周禮·大司馬法}，天子六軍，出自六鄉，凡起徒役，無過家一人，故一鄉出一軍，鄉為正，遂為副。若鄉遂不足，則徵兵於邦國。則司馬法六十四井為甸，計有五百七十六夫，共出長轂一乘，甲士三人，步卒七十二人。至於臨敵對戰布陳之時，則依六鄉軍法，五人為伍，五伍為兩，四兩為卒，五卒為旅，五旅為師，五師為軍。故\CJKunderwave{左傳}云：“先偏後伍。”又云：“廣有一卒，卒偏之兩。”非直人數如此，車數亦然。故\CJKunderwave{周禮}云：“乃會車之卒伍。”鄭云：“車亦有卒伍。”\CJKunderwave{左傳}“戰於繻葛”，杜注云：“車二十五乘為偏。”是車亦為卒伍之數也。則一車七十二人者，自計元科兵之數。科兵既至，臨時配割，其車雖在，其人分散，前配車之人，臨戰不得還屬本車，當更以虎賁甲士配車而戰。孔舉七十二人元科兵數者，欲總明三百兩人之大數。雲“兵車，百夫長所載者”，欲見臨敵實一車有百人，既“虎賁”與車數相當，又經稱“百夫長”，故孔為此說。 \par}

{\noindent\zhuan\zihao{6}\fzbyks 傳“勇士”至“夫長”。正義曰：\CJKunderwave{周禮}虎賁氏之官，其屬有虎士八百人，是“虎賁”為“勇士稱”也。若虎之賁走逐獸，言其猛也。此“虎賁”必是軍內驍勇選而為之，當時謂之“虎賁”。\CJKunderwave{樂記}雲“虎賁之士說劍”，謂此也。孔意“虎賁”即是經之“百夫長”，故云：“皆百夫長”也。 \par}

{\noindent\shu\zihao{5}\fzkt “武王”至“牧誓”。正義曰：武王以兵戎之車三百兩、虎賁之士三百人與受戰於商郊牧地之野，將戰之時，王設言以誓眾。史敘其事，作\CJKunderwave{牧誓}。 \par}

牧誓\footnote{至牧地而誓眾。○牧如字，徐一音茂，\CJKunderwave{說文}作坶，云：“地名，在朝歌南七十里。”\CJKunderwave{字林}音母。}

時甲子昧爽，\footnote{是克紂之月甲子之日，二月四日。昧，冥；爽，明；早旦。○昧音妹。爽，明也。“昧爽”謂早旦也。馬云：“昧，未旦也。”}

{\noindent\zhuan\zihao{6}\fzbyks 傳“是克”至“早旦”。正義曰：\CJKunderwave{春秋}主書動事，編次為文，於法日月時年皆具，其有不具，史闕耳。\CJKunderwave{尚書}惟記言語，直指設言之日。上篇“戊午，次於河朔”，\CJKunderwave{洛誥}“戊辰，王在新邑”，與此“甲子”，皆言有日無月，史意不為編次，故不具也。是“克紂之月甲子之日，是周之二月四日”，以歷推而知之也。\CJKunderwave{釋言}云：“晦，冥也。”“昧”亦晦義，故為冥也。“冥”是夜，“爽”是明，夜而未明謂早旦之時，蓋雞鳴後也。為下“朝至”發端，“朝”即“昧爽”時也。 \par}

王朝至於商郊牧野，乃誓。\footnote{紂近郊三十里地名牧。癸亥夜陳，甲子朝誓，將與紂戰。○陳，直刃反。}

{\noindent\zhuan\zihao{6}\fzbyks 傳“紂近”至“紂戰”。正義曰：傳言在“紂近郊三十里”，或當有所據也。皇甫謐云：“在朝歌南七十里。”不知出何書也。言“至於商郊牧野”，知“牧”是郊上之地。戰在平野，故言“野”耳。\CJKunderwave{詩}云：“於牧之野。”\CJKunderwave{禮記大傳}云：“牧之野，武王之大事，繼牧言野，明是牧地。”而\CJKunderline{鄭玄}云：“郊外曰野，將戰於郊，故至牧野而誓。”案經“至於商郊牧野乃誓”，豈王行已至於郊，乃後到退適野，誓訖而更進兵乎？何不然之甚也！\CJKunderwave{武成}云：“癸亥夜陳，未畢而雨。”是癸亥夜已布陳，故甲子朝而誓眾，將與紂戰，故戒敕之。 \par}

王左杖黃鉞，右秉白旄以麾，曰：“逖矣，西土之人！”\footnote{鉞，以黃金飾斧。左手杖鉞，示無事於誅。右手把旄，示有事於教。逖，遠也。遠矣，西土之人。勞苦之。○杖，徐直亮反。鉞音越，本又作戌。旄音毛，馬云：“白旄，旄牛尾。”麾，許危反。逖，他歷反。}

{\noindent\zhuan\zihao{6}\fzbyks 傳“鉞以”至“苦之”。正義曰：太公\CJKunderwave{六韜}云：“大柯斧重八斤，一名天鉞。”\CJKunderwave{廣雅}云：“鉞，斧也。”斧稱“黃鉞”，故知“以黃金飾斧”也。鉞以殺戮，殺戮用右手，用左手杖鉞，示無事於誅。右手把旄，示有事於教。其意言惟教軍人，不誅殺也。把旄何以白？旄用白者，取其易見也。“逖，遠”，\CJKunderwave{釋詁}文。 \par}

王曰：“嗟！我友邦冢君，\footnote{同志為友，言志同滅紂。}御事司徒、司馬、司空，\footnote{治事三卿，司徒主民，司馬主兵，司空主土，指誓戰者。}

{\noindent\zhuan\zihao{6}\fzbyks 傳“治事”至“戰者”。正義曰：孔以於時已稱王而有六師，亦應已置六卿。今呼治事惟三卿者，司徒主民，治徒庶之政令；司馬主兵，治軍旅之誓戒；司空主土，治壘壁以營軍；是“指誓戰者”，故不及太宰、大宗、司寇也。其時六卿具否，不可得知，但據此三卿為說耳。此“御事”之文，指三卿而說，是不通於“亞旅”已下。 \par}

亞旅、師氏，\footnote{亞，次。旅，眾也。眾大夫，其位次卿。師氏，大夫，官以兵守門者。}

{\noindent\zhuan\zihao{6}\fzbyks 傳“亞次”至“門者”。正義曰：“亞，次”，\CJKunderwave{釋言}文。“旅，眾”，\CJKunderwave{釋詁}文。此及\CJKunderwave{左傳}皆卿下言“亞旅”，知是“大夫，其位次卿”，而數眾，故以亞次名之，謂諸是四命之大夫，在軍有職事者也。“師氏”亦大夫，其官掌以兵守門，所掌尤重，故別言之。\CJKunderwave{周禮}師氏中大夫，“使其屬帥四夷之隸，各以其兵服守王之門外。朝在野外，則守內列”。\CJKunderline{鄭玄}云：“內列，蕃營之在內者也，守之如守王宮。” \par}

千夫長、百夫長，\footnote{師帥，卒帥。○帥，色類反，下同。}

{\noindent\zhuan\zihao{6}\fzbyks 傳“師帥，卒帥”。正義曰：\CJKunderwave{周禮}二千五百人為師，師帥皆中大夫。百人為卒，卒長皆上士。孔以師雖二千五百人，舉全數亦得為幹夫長，“長”與“帥”其義同，是千夫長亦可以稱“帥”，故以“千夫長”為師帥，“百夫長”為卒帥。王肅雲“師長、卒長”，意與孔同，順經文而稱“長”耳。\CJKunderline{鄭玄}以為“師帥，旅帥也”，與孔不同。 \par}

及庸、蜀、羌、髳、微、盧、彭、濮人。\footnote{八國皆蠻夷戎狄屬文王者國名。羌在西蜀叟，髳、微在巴蜀，盧、彭在西北，庸、濮在江漢之南。}

{\noindent\zhuan\zihao{6}\fzbyks 傳“八國”至“之南”。正義曰：九州之外，四夷大名，則東夷、西戎、南蠻、北狄，其在當方，或南有戎而西有夷。此八國並非華夏，故大判言之，“皆蠻夷戎狄屬文王者國名”也。此八國皆西南夷也，文王國在於西，故西南夷先屬焉。大劉以“蜀”是蜀郡，顯然可知，孔不說。又退“庸”就“濮”解之，故以次先解“羌”。雲“羌在西蜀叟”者，漢世西南之夷，“蜀”名為大，故傳據“蜀”而說。左思\CJKunderwave{蜀都賦}云：“三蜀之豪，時來時往。”是蜀都分為三，羌在其西，故云“西蜀叟”。“叟”者蜀夷之別名，故\CJKunderwave{後漢書}“興平元年，馬騰、劉範謀誅李𠐶，益州牧劉焉遣叟兵五千人助之”，是蜀夷有名“叟”者也。“髳、微在巴蜀”者，巴在蜀之東偏，漢之巴郡所治江州縣也。“盧、彭在西北”者，在東蜀之西北也。文十八年\CJKunderwave{左傳}稱，庸與百濮伐楚，楚遂滅庸。是“庸、濮在江漢之南”。 \par}

稱爾戈，比爾幹，立爾矛，予其誓。”\footnote{稱，舉也。戈，戟。幹，楯也。○比，徐扶志、毗志二反。楯，食準反，又音允。}

{\noindent\zhuan\zihao{6}\fzbyks 傳“稱舉”至“幹楯”。正義曰：“稱，舉”，\CJKunderwave{釋言}文。方言云：“戟，楚謂之孑，吳揚之間謂之戈。”是“戈”即戟也。\CJKunderwave{考工記}云：“戈柲六尺有六寸，車戟常。”鄭云：“八尺曰尋，倍尋曰常。”然則戈戟長短異名，而云“戈”者即戟，戈戟長短雖異，其形制則同，此雲舉戈，宜舉其長者，故以“戈”為戟也。\CJKunderwave{方言}又云：“楯，自關而東或謂之楯，或謂之幹，關西謂之楯。”是“幹”、“楯”為一也。戈短，人執以舉之，故言“稱”。楯則並以捍敵，故言“比”。矛長立之於地，故言“立”也。 \par}

王曰：“古人有言曰：‘牝雞無晨。\footnote{言無晨鳴之道。○牝，類引反，徐扶忍反。}牝雞之晨，惟家之索。’\footnote{索，盡也。喻婦人知外事，雌代雄鳴則家盡，婦奪夫政則國亡。○索，西各反。}

{\noindent\zhuan\zihao{6}\fzbyks 傳“索盡”至“國亡”。正義曰：\CJKunderwave{禮記·檀弓}曰：“吾離群而索居。”則“索居”為散義。\CJKunderline{鄭玄}云：“索，散也。”物散則盡，故“索”為盡也。“牝雞”，雌也。\CJKunderwave{爾雅}飛曰“雌雄”，走曰“牝牡”，而此言“牝雞”者，\CJKunderwave{毛詩}、\CJKunderwave{左傳}稱“雄狐”，是亦飛、走通也。此以牝雞之鳴喻婦人知外事，故重申喻意云：“雌代雄鳴則家盡，婦奪夫政則國亡。”“家”總貴賤為文，言“家”以對“國”耳。將陳紂用婦言，故舉此古人之語。紂直用婦言耳，非能奪其政，舉此言者，專用其言，賞罰由婦，即是奪其政矣。婦人不當知政，是別外內之分，若使賢如文母，可以興助國家，則非牝雞之喻矣。 \par}

今商王受惟婦言是用，\footnote{妲己惑紂，紂信用之。○妲，丹達反；己音紀；紂妻也。}

{\noindent\zhuan\zihao{6}\fzbyks 傳“妲己”至“用之”。正義曰：\CJKunderwave{晉語}云：“殷辛伐有蘇氏，蘇氏以妲己女焉。妲己有寵而亡殷。”\CJKunderwave{殷本紀}云：“紂嬖於婦人，愛妲己，惟妲己之言是從。”\CJKunderwave{列女傳}云：“紂好酒淫樂，不離妲己，妲己所與言者貴之，妲己所憎者誅之。為長夜飲，妲己好之，百姓怨望，而諸侯有叛者。妲已曰：‘罰輕誅薄，威不立耳。’紂乃重刑辟，為炮烙之法，妲已乃笑。武王伐紂，斬妲已頭懸之於小白旗上，以為亡紂者此女也。” \par}

昏棄厥肆祀弗答，\footnote{昏，亂。肆，陳。答，當也。亂棄其所陳祭祀，不復當享鬼神。○復，扶又反。}

{\noindent\zhuan\zihao{6}\fzbyks 傳“昏亂”至“鬼神”。正義曰：昏暗者於事必亂，故“昏”為亂也。\CJKunderwave{詩}云：“肆筵設席。”“肆”者陳設之意，\CJKunderwave{毛傳}亦以“肆”為陳也。對合，相當之事，故“答”為當也。紂身昏亂，棄其宜所陳設祭祀，不復當享鬼神，與上“郊社不修，宗廟不享”亦一也。不事神祗，惡之大者，故\CJKunderwave{泰誓}及此三言之。 \par}

昏棄厥遺王父母弟不迪，\footnote{王父，祖之昆弟。母弟，同母弟。言棄其骨肉，不接之以道。}

{\noindent\zhuan\zihao{6}\fzbyks 傳“王父”至“以道”。正義曰：\CJKunderwave{釋親}雲“父之考為王父”，則“王父”是祖也。紂無親祖可棄，故為“祖之昆弟”。棄其祖之昆弟，則父之昆弟亦棄之矣。\CJKunderwave{春秋}之例，母弟稱“弟”，凡\CJKunderwave{春秋}稱“弟”皆是母弟也。“母弟”謂同母之弟，同母尚棄，別生者必棄矣，舉尊親以見卑疏也。“遺”亦“棄”也，言紂之昏亂，棄其所遺骨肉之親，不接之以道。經先言棄祀、棄親者，\CJKunderline{鄭玄}云：“\CJKunderwave{誓}首言此者，神怒民怨，紂所以亡也。” \par}

乃惟四方之多罪逋逃，是崇是長，\footnote{言紂棄其賢臣，而尊長逃亡罪人，信用之。}是信是使，是以為大夫卿士。\footnote{士，事也。用為卿大夫，典政事。}俾暴虐於百姓，以奸宄於商邑。\footnote{使四方罪人暴虐奸宄于都邑。○俾，必爾反，使也。}

{\noindent\zhuan\zihao{6}\fzbyks 傳“使四”至“都邑”。正義曰：“暴虐”謂殺害，殺害加於人，故言“於百姓”。“奸宄”謂劫奪，劫奪有處，故言“於商邑”。百姓亦是商邑之人，故傳總言“于都邑”也。 \par}

今予發惟恭行天之罰。今日之事，不愆於六步、七步，乃止齊焉。\footnote{今日戰事，就敵不過六步、七步，乃止相齊。言當旅進一心。}

{\noindent\zhuan\zihao{6}\fzbyks 傳“今日”至“一心”。正義曰：戰法布陳然後相向，故設其就敵之限，不過六步、七步，乃止相齊焉。欲其相得力也。\CJKunderwave{樂記}稱“進旅退旅”，是“旅”為眾也，言當眾進一心也。 \par}

夫子勖哉！不愆於四伐、五伐、六伐、七伐，乃止齊焉。\footnote{夫子謂將士，勉勵之。伐謂擊剌，少則四五，多則六七以為例。○勖，許六反。剌，七亦反。}

{\noindent\zhuan\zihao{6}\fzbyks 傳“夫子”至“為例”。正義曰：此及下文三雲“夫子”，此“勖哉”在下，下“勖哉”在上。此先呼其人，然後勉之；此既言然，下先令勉勵，乃呼其人，各與下句為目也。上有“戈”、“矛”，戈謂擊兵，矛謂剌兵，故云“伐謂擊剌”，此“伐”猶伐樹然也。 \par}

勖哉夫子！尚桓桓，\footnote{桓桓，武貌。}

{\noindent\zhuan\zihao{6}\fzbyks 傳“桓桓，武貌”。正義曰：\CJKunderwave{釋訓}云：“桓桓，威也。”\CJKunderwave{詩序}云：“桓，武志也。” \par}

如虎如貔，如熊如羆，於商郊。\footnote{貔，執夷，虎屬也。四獸皆猛健，欲使士眾法之，奮擊於牧野。○貔，彼皮反，\CJKunderwave{爾雅}云：“羆如熊，黃白文。”}

{\noindent\zhuan\zihao{6}\fzbyks 傳“貔，執夷”。正義曰：\CJKunderwave{釋獸}云：“貔，白狐，其子豰。”舍人曰：“貔名白狐，其子名豰。”郭璞曰：“一名執夷，虎豹屬。” \par}

弗迓克奔，以役西土，\footnote{商眾能奔來降者，不迎擊之，如此則所以役我西土之義。○迓，五嫁反，馬作御，禁也。役，馬云：“為也。”為，於偽反。}

{\noindent\zhuan\zihao{6}\fzbyks 傳“商眾”至“之義”。正義曰：“迓”訓迎也，不迎擊商眾能奔來降者，兵法不誅降也。“役”謂使用也，如此不殺降人，則所以使用我西土之義。用義於彼，令彼知我有義也。王肅讀“御”為禦，言“不御能奔走者，如殷民欲奔走來降者，無逆之；奔走去者，可不御止。役，為也，盡力以為我西土”。與孔不同。 \par}

勖哉夫子！爾所弗勖，其於爾躬有戮。”\footnote{臨敵所安，汝不勉，則於汝身有戮矣。}

\section{武成第五【偽】}


武王伐殷，往伐歸獸，\footnote{往誅紂克定，偃武修文，歸馬牛於華山桃林之牧地。○獸，徐始售反；本或作獸，許救反。}識其政事，\footnote{記識殷家政教善事以為法。}作\CJKunderwave{武成}。\footnote{武功成，文事修。}


{\noindent\zhuan\zihao{6}\fzbyks 傳“往誅”至“牧地”。正義曰：此序於經“於徵伐商”，是“往伐”也。“歸馬”、“放牛”是“歸獸”也。故傳引經以解之，\CJKunderwave{爾雅}有\CJKunderwave{釋獸}、\CJKunderwave{釋畜}，畜、獸形相類也，在野自生為獸，人家養之為畜。歸馬放牛，不復乘用，使之自生自死，若野獸然，故謂之“獸”。獸以野澤為家，故言“歸”也。 \par}

{\noindent\zhuan\zihao{6}\fzbyks 傳“記識”至“為法”。正義曰：紂以昏亂而滅，前世政有善者，故訪問殷家政教，記識善事以為治國之法，經雲“列爵惟五，分士惟三”是也。 \par}

{\noindent\shu\zihao{5}\fzkt “武王”至“武成”。正義曰：武王之伐殷也，往則陳兵伐紂，歸放牛馬為獸，記識殷家美政善事而行用之。史敘其事，作\CJKunderwave{武成}。 \par}

武成\footnote{文王受命，有此武功，成於克商。}

{\noindent\zhuan\zihao{6}\fzbyks 傳“文王”至“克商”。正義曰：“文王受命，有此武功”，\CJKunderwave{詩}之文也。彼言“武功”，謂始伐崇耳。殷紂尚在，其功未成，成功在於克商，今武始成矣，故以“武成”名篇，以\CJKunderwave{泰誓}繼文王之年，故本之於文王。鄭云：“著武道至此而成。” \par}

{\noindent\shu\zihao{5}\fzkt “武成”。正義曰：此篇敘事多而王言少，惟辭又首尾不結，體裁異於餘篇。自“惟一月”至“受命於周”，史敘伐殷往反及諸侯大集，為王言發端也。自“王若曰”至“大統未集”，述祖父以來開建王業之事也。自“予小子”至“名山大川”，言己承父祖之意，告神陳紂之罪也。自“曰惟有道”至“無作神羞”，王自陳告神之辭也。“既戊午”已下，又是史敘往伐殺紂，入殷都佈政之事。“無作神羞”以下，惟告神，其辭不結，文義不成，非述作之體。案\CJKunderwave{左傳}荀偃禱河云：“無作神羞，具官臣偃，無敢復濟，惟爾有神裁之。”蒯聵禱祖云：“無作三祖羞，大命不敢請，佩玉不敢愛。”彼二者於“神羞”之下皆更申己意，此經“無作神羞”下更無語，直是與神之言猶尚未訖。且冢君百工初受周命，王當有以戒之，如\CJKunderwave{湯誥}之類。宜應說其除害與民更始，創以為惡之禍，勸以行道之福，不得大聚百官，惟誦禱辭而已。欲徵則殷勤誓眾，既克則空話禱神，聖人有作，理必不爾。竊謂“神羞”之下，更合有言，簡編斷絕，經失其本，所以辭不次耳。或初藏之日，已失其本；或壞壁得之，始有脫漏；故孔稱五十八篇以外，錯亂磨滅，不可復知。明是見在諸篇亦容脫錯，但孔此篇首尾具足，既取其文為之作傳，恥雲有所失落，不復言其事耳。 \par}

惟一月壬辰，旁死魄。\footnote{此本說始伐紂時。一月，周之正月。旁，近也。月二日，近死魄。○旁，步光反。魄，普白反，\CJKunderwave{說文}作霸，匹革反，云：“月始生魄然貌。”近，附近之近。}越翼日癸巳，王朝步自周，於徵伐商。\footnote{翼，明。步，行也。武王以正月三日行自周，往征伐商，二十八日渡孟津。}厥四月，哉生明，王來自商，至於豐。\footnote{其四月。哉，始也。始生明，月三日，與死魄互言。○哉，徐音載。豐，芳弓反，文王所都也。}


{\noindent\zhuan\zihao{6}\fzbyks 傳“此本”至“死魄”。正義曰：將言武成，遠本其始。”此本說始伐紂時。一月，周之正月”，是建子之月，殷十二月也。此月辛卯朔，朔是死魄，故“月二日，近死魄”。“魄”者，形也，謂月之輪郭無光之處名“魄”也。朔後明生而魄死，望後明死而魄生。\CJKunderwave{律曆志}云：“死魄，朔也。生魄，望也。”\CJKunderwave{顧命}云：“惟四月哉生魄。”傳云：“始生魄，月十六日也。”月十六日為始生魄，是一日為始死魄，二日近死魄也。顧氏解“死魄”與小劉同。大劉以三日為始死魄，二日為旁死魄。旁死魄無事而記之者，與下日為發端，猶今之將言日，必先言朔也。 \par}

{\noindent\zhuan\zihao{6}\fzbyks 傳“翼明”至“孟津”。正義曰：“翼，明”，\CJKunderwave{釋言}文。\CJKunderwave{釋宮}云：“堂上謂之行，堂下謂之步。”彼相對為名耳。散則可以通，故“步”為行也。周去孟津千里，以正月三日行自周，二十八日渡孟津，凡二十五日，每日四十許裡，時之宜也。\CJKunderwave{詩}云：“於三十里。”\CJKunderwave{毛傳}云：“師行三十里。”蓋言其大法耳。 \par}

{\noindent\zhuan\zihao{6}\fzbyks 傳“其四”至“互言”。正義曰：“其四月”，此伐商之四月也。“哉，始”，\CJKunderwave{釋詁}文。\CJKunderwave{顧命}傳以“哉生魄”為十六日，則“哉生明”為月初矣。以三日月光見，故傳言“始生明，月三日”也。此經無日，未必非二日也。“生明”、“死魄”俱是月初，上雲“死魄”，此雲“生明”，而魄死明生互言耳。 \par}

乃偃武修文，\footnote{倒載干戈，包以虎皮，示不用。行禮射，設庠序，修文教。}歸馬於華山之陽，放牛於桃林之野，示天下弗服。\footnote{山南曰陽。桃林在華山東。皆非長養牛馬之地，欲使自生自死，示天下不復乘用。○華，胡化、胡瓜二反；華山在恆農。長，丁丈反。復，扶又反。}丁未，祀於周廟，邦甸、侯、衛，駿奔走，執豆籩。\footnote{四月丁未，祭告后稷以下、文考文王以上七世之祖。駿，大也。邦國甸侯、衛服諸侯皆大奔走於廟執事。○駿，荀俊反。豆，本又作梪。籩，音邊。上，時掌反。}越三日庚戌，柴望，大告武成。\footnote{燔柴郊天，望祀山川，先祖後郊，自近始。}

{\noindent\zhuan\zihao{6}\fzbyks 傳“倒載”至“文教”。正義曰：\CJKunderwave{樂記}云，武王克殷，“濟河而西。車甲釁而藏之府庫，倒載干戈，包之以虎皮，天下知武王之不復用兵也。散軍而郊射。左射，\CJKunderwave{貍首}。右射，\CJKunderwave{騶虞}。而貫革之射息也”。是“偃武修文”之事，故傳引之。郊射是禮射也。\CJKunderwave{王制}論四代學名云：“虞謂之庠，夏謂之序。”故言“設庠序，修文教”也。 \par}

{\noindent\zhuan\zihao{6}\fzbyks 傳“山南”至“乘用”。正義曰：\CJKunderwave{釋山}云：“山西曰夕陽，山東曰朝陽。”李巡曰：“山西暮乃見日，故曰夕陽。山東朝乃見日，故云朝陽。”“陽”以見日為名，故知“山南曰陽”。杜預云：“桃林之塞，今宏農華陰縣潼關是也。”是在“華山東”也。指其所往謂之“歸”，據我釋之則雲“放”，“放牛”、“歸馬”互言之耳。華山之旁尤乏水草，非長養牛馬之地，欲使自生自死。此是戰時牛馬，故放之，示天下不復乘用。\CJKunderwave{易·繫辭}云：“服牛乘馬。”“服”、“乘”俱是用義，故以“服”總牛馬。 \par}

{\noindent\zhuan\zihao{6}\fzbyks 傳“四月”至“執事”。正義曰：以“四月”之字，隔文已多，故言“四月丁未”。此以成功設祭，明其遍告群祖，知告“后稷以下”。后稷則始祖以下，容毀廟也。天子七廟，故云“文考文王以上七世之祖”。見是周廟皆祭之，故經總雲“周廟”也。“駿，大”，\CJKunderwave{釋詁}文。\CJKunderwave{周禮}六服侯、甸、男、採、衛、要，此略舉邦國在諸侯服，故云“甸、侯、衛”，其言不次。\CJKunderwave{詩·頌}雲“駿奔走在廟。”故云：“皆大奔走於廟執事”也。“越三日庚戌”。正義曰：\CJKunderwave{召誥}雲“越三日”者，皆從前至今為三日，此從丁未數之，則為四日，蓋史官不同，立文自異。或此“三”當為“四”，由字積與誤。 \par}

{\noindent\shu\zihao{5}\fzkt “惟一”至“武成”。正義曰：此歷敘伐紂往反祖廟告天時日，說武功成之事也。“一月壬辰，旁死魄”，謂伐紂之年周正月辛卯朔，其二日是壬辰也。“翼日癸巳，王朝步自周，於徵伐商”，謂正月三日發鎬京始東行也。其月二十八日戊午渡河。\CJKunderwave{泰誓}序雲“一月戊午，師渡孟津”，\CJKunderwave{泰誓}中篇雲“惟戊午，王次於河朔”是也。二月辛酉朔，甲子殺紂，\CJKunderwave{牧誓}雲“時甲子昧爽，乃誓”是也。其年閏二月庚寅朔，三月庚申朔，四月己丑朔。“厥四月，哉生明，王來自商，至於豐”，謂四月三日，月始生明，其日當是辛卯也。“丁未，祀於周廟”，四月十九日也。“越三日庚戌，柴望”，二十二日也。正月始往伐，四月告成功，史敘其事，見其功成之次也。\CJKunderwave{漢書·律曆志}引\CJKunderwave{武成}篇云：“惟一月壬辰，旁死魄。若翼日癸巳，武王乃朝步自周，於徵伐紂。越若來二月既死魄，越五日甲子，咸劉商王紂。惟四月既旁生魄，越六日庚戌，武王燎於周廟。翼日辛亥，祀於天位。越五日乙卯，乃以庶國祀於周廟。”與此經不同。彼是焚書之後，有人偽為之。漢世謂之“逸書”，其後又亡其篇。\CJKunderline{鄭玄}云：“\CJKunderwave{武成}逸書，建武之際亡。”謂彼偽\CJKunderwave{武成}也。 \par}

既生魄，庶邦冢君暨百工，受命於周。\footnote{魄生明死，十五日之後，諸侯與百官受政命於周。明一統。○暨，其器反。}

{\noindent\zhuan\zihao{6}\fzbyks 傳“魄生”至“一統”。正義曰：月以望虧，望是月半，望在十六日為多，通率在十六日者，四分居三，其一在十五日耳。此言“既生魄”，故言“魄生明死，十五日之後”也。“丁未，祀於周廟”，已是此月十九日矣，此“受命於周”，繼生魄言之，則受命在祀廟之前，故祀廟之時諸侯已奔走執事，豈得未受周命，已助周祭？明其受命在祀廟前矣。史官探其時日，先言告武成既訖，然後卻說受命，故文在下耳。諸侯與百官，舊有未屬周者，今皆受政命於周，於此時始天下一統也。顧氏以既生魄謂庚戌已後，雖十六日始生魄，從十六日至晦皆為生魄，但不知庚戌之後幾日耳。 \par}

王若曰：“嗚呼！群后，\footnote{順其祖業嘆美之，以告諸侯。}惟先王建邦啟土，\footnote{謂后稷也。尊祖，故稱先王。}

{\noindent\zhuan\zihao{6}\fzbyks 傳“謂後”至“先王”。正義曰：此“先王”文在“公劉”之前，知“謂后稷也”。后稷非王，尊其祖，故稱先王。\CJKunderwave{周語}雲“昔我先王后稷”，又曰“我先王不窋”，韋昭云：“王之先祖，故稱王。”\CJKunderwave{商頌}亦以契為“玄王”。文武之功，起於後稷，后稷始封於邰，故言“建邦啟土”。 \par}

公劉克篤前烈，\footnote{后稷曾孫。公，爵。劉，名。能厚先人之業。}

{\noindent\zhuan\zihao{6}\fzbyks 傳“后稷”至“之業”。正義曰：\CJKunderwave{周本紀}云：“后稷卒，子不窋立。卒，子鞠陶立。卒，子公劉立。”是公劉為后稷曾孫也。\CJKunderwave{本紀}云，公劉之後有公非、公祖之類，知“公”是爵。殷時未諱，故稱劉名。先公多矣，獨三人稱“公”，當時之意耳。\CJKunderwave{本紀}云：“公劉復修后稷之業，百姓懷之，多徙而歸保焉。”周道之興，自此之後，是“能厚先人之業”也。 \par}

至於大王，肇基王跡，王季其勤王家。\footnote{大王修德以翦齊商人，始王業之肇跡王季纘統其業，乃勤立王業。○大音太。肇音兆。王跡，於況反，又如字，注“王業”、“王功”同。}

{\noindent\shu\zihao{5}\fzkt 傳“大王”至“王家”。正義曰：\CJKunderwave{詩}云：“后稷之孫，實惟大王。居岐之陽，實始翦商。”是大王翦齊商人，始王業之兆跡也。\CJKunderwave{周本紀}云：“王季修古公之道，諸侯順之。”是能纘統大王之業，勤立王家之基本也。 \par}

我文考文王，克成厥勳，誕膺天命，以撫方夏。\footnote{言我文德之父，能成其王功，大當天命，以撫綏四方中夏。}大邦畏其力，小邦懷其德。\footnote{言天下諸侯，大者畏威，小者懷德，是文王威德之大。}

{\noindent\shu\zihao{5}\fzkt “大邦”至“其德”。正義曰：大邦力足拒敵，故言“畏其力”，小邦必畏矣。小邦或被棄遺，故言“懷其德”，大邦亦懷德矣。量事為文也。 \par}

惟九年，大統未集。\footnote{言諸侯歸之，九年而卒，故大業未就。}

{\noindent\zhuan\zihao{6}\fzbyks 傳“言諸”至“未就”。正義曰：文王斷虞芮之訟，諸侯歸之，改稱元年。至九年而卒，故云“大業未就”也。文王既未稱王，而得輒改元年者，諸侯自於其國各稱元年，是己之所稱，容或中年得改矣。\CJKunderwave{汲冢竹書}魏惠王有後元年，漢初文帝二元，景帝三元，此必有因於古也。伏生、司馬遷、韓嬰之徒不見此書，以為文王受命七年而崩，故\CJKunderline{鄭玄}等皆依用之。 \par}

予小子其承厥志，\footnote{言承文王本意。}厎商之罪，告於皇天后土、所過名山大川，\footnote{致商之罪，謂伐紂之時。后土，社也。名山，華嶽。大川，河。○厎，之履反。}

{\noindent\zhuan\zihao{6}\fzbyks 傳“致商”至“川河”。正義曰：“致商之罪，謂伐紂之時”，欲將伐紂，告天乃發，故文在“所過”之上。\CJKunderwave{禮}天子出征，必類帝宜社。此告皇天后土，即\CJKunderwave{泰誓}上篇“類於上帝，宜於冢土”，故云“后土，社也。”昭二十九年\CJKunderwave{左傳}稱“句龍為后土”，后土為社是也。僖十五年\CJKunderwave{左傳}云，戴皇天而履後上。彼晉大夫要秦伯，故以地神后土而言之，與此異也。自周適商，路過河華，故知所過名山華嶽、大川河也。山川大乃有名，“名”、“大”互言之耳。\CJKunderwave{周禮·大祝}云：“王過大山川，則用事焉。”鄭云：“用事，用祭事告行也。” \par}

曰：‘惟有道曾孫周王發，將有大正於商。\footnote{告天地山川之辭。大正，以兵徵之也。}

{\noindent\shu\zihao{5}\fzkt “曰惟有道曾孫周王發”。正義曰：自稱“有道”者，聖人至公，為民除害，以紂無道，言己有道，所以告神求助，不得飾以謙辭也。稱“曾孫”者，\CJKunderwave{曲禮}說諸侯自稱之辭云：“臨祭祀，內事曰孝子某侯某，外事曰曾孫某侯某。”哀二年\CJKunderwave{左傳}蒯聵禱祖亦自稱曾孫，皆是言己承藉上祖奠享之意。 \par}

今商王受無道，\footnote{無道德。}暴殄天物，害虐烝民，\footnote{暴絕天物，言逆天也。逆天害民，所以為無道。○烝，之承反。}

{\noindent\shu\zihao{5}\fzkt “暴殄”至“烝民”。正義曰：“天物”語闊，人在其間，以人為貴，故別言害民。則“天物”之言，除人外，普謂天下百物、鳥獸草木皆暴絕之。 \par}

為天下逋逃主，萃淵藪。\footnote{逋，亡也。天下罪人逃亡者，而紂為魁主，窟聚淵府藪澤。言大奸。○萃，在醉反。藪，素口反。魁，苦回反。窟，口忽反。}

{\noindent\zhuan\zihao{6}\fzbyks 傳“逋亡”至“大奸”。正義曰：“逋”亦逃也，故以為亡。罪人逃亡，而紂為魁主。“魁”，首也，言受用逃亡者，與之為魁首，為主人。“萃”訓聚也，言若蟲獸入窟，故云“窟聚”。水深謂之“淵”，藏物謂之“府”。史游\CJKunderwave{急就篇}云：“司農少府國之淵。”“淵”、“府”類，故言“淵府”。水鍾謂之“澤”，無水則名“藪”。“藪”、“澤”大同，故言“藪澤”。“萃淵藪”三者各為物室，言紂與亡人為主，亡人歸之若蟲之窟聚，魚歸淵府，獸集藪澤，言紂為大奸也。據傳意，“主”字下讀為便。昭七年\CJKunderwave{左傳}引此文，杜預云：“萃，集也。天下逋逃悉以紂為淵藪，集而歸之。”與孔異也。 \par}

予小子既獲仁人，敢祇承上帝，以遏亂略。\footnote{仁人，謂大公、周、召之徒。略，路也。言誅紂敬承天意以絕亂路。○遏，烏末反。召，上照反，本又作邵。}華夏蠻貊，罔不率俾。恭天成命，\footnote{冕服採章曰華，大國曰夏，及四夷皆相率而使奉天成命。○貊，亡白反。俾，必爾反。}

{\noindent\zhuan\zihao{6}\fzbyks 傳“冕服”至“成命”。正義曰：“冕服採章”對被髮左衽，則為有光華也。\CJKunderwave{釋詁}云：“夏，大也。”故大國曰“夏華”。“夏”謂中國也。言“蠻貊”則戎夷可知也。言華夏及四夷皆相率而充己，使奉天成命，欲其共伐紂也。 \par}

肆予東征，綏厥士女。\footnote{此謂十一年會孟津還時。}惟其士女，篚厥玄黃，昭我周王。\footnote{言東國士女筐篚盛其絲帛，奉迎道次。明我周王為之除害。○篚音匪。為，於偽反。}天休震動，用附我大邑周。\footnote{天之美應，震動民心，故用依附我。○應，應對之應。}惟爾有神，尚克相予，以濟兆民，無作神羞。’”\footnote{神庶幾助我渡民危害，無為神羞辱。○相，息亮反。}既戊午，師逾孟津。癸亥，陳於商郊，俟天休命。\footnote{自河至朝歌，出四百里，五日而至。赴敵宜速，待天休命，謂夜雨止畢陳。○逾，亦作逾。陳，直刃反，計同，徐音塵。}甲子昧爽，受率其旅若林，會於牧野。\footnote{旅，眾也。如林，言盛多。會，逆距戰。}



{\noindent\zhuan\zihao{6}\fzbyks 傳“自河”至“畢陳”。正義曰：“出四百里”，驗地為然。戊午明日猶誓於河朔，癸亥已陳於商郊，凡經五日，日行八十里，所以疾者，“赴敵宜速”也。\CJKunderwave{帝王世紀}云：“王軍至鮪水，紂使膠鬲候周師，見王問曰：‘西伯將焉之？’王曰：‘將攻薛也。’膠鬲曰：‘然，原西伯無我欺。’王曰：‘不子欺也，將之殷。’膠鬲曰：‘何日至？’王曰：‘以甲子日，以是報矣。’膠鬲去而報命於紂。而雨甚，軍卒皆諫王曰：‘卒病，請休之。’王曰：‘吾已令膠鬲以甲子報其主矣。吾雨而行，所以救膠鬲之死也。’遂行，甲子至於商郊。”然則本期甲子，故遠行也。\CJKunderwave{周語}云：“王以二月癸亥夜陳，未畢而雨。”是“雨止畢陳”也。“待天休命”，雨是天之美命也。韋昭云：“雨者，天地神人和同之應也。”天地氣和乃有雨降，是雨為和同之應也。 \par}

{\noindent\zhuan\zihao{6}\fzbyks 傳“旅眾”至“距戰”。正義曰：“旅，眾”，\CJKunderwave{釋詁}文。\CJKunderwave{詩}亦云：“其會如林。”言盛多也。\CJKunderwave{本紀}云：“紂發兵七十萬人以距武王。”紂兵雖則眾多，不得有七十萬人，是史官美其能破強敵，虛言之耳。 \par}

罔有敵於我師，前徒倒戈，攻於後以北，血流漂杵。\footnote{紂眾服周仁政，無有戰心，前徒倒戈，自攻於後以北走，血流漂舂杵。甚之言。○倒，丁老反。漂，四妙反，徐敷妙反，又四消反。杵，昌呂反。}

{\noindent\zhuan\zihao{6}\fzbyks 傳“紂眾”至“之言”。正義曰：“罔有敵於我師”，言紂眾雖多，皆無有敵我之心，故“自攻於後以北走”。自攻其後，必殺人不多，“血流漂舂杵，甚之言”也。\CJKunderwave{孟子}云：“信\CJKunderwave{書}不如無\CJKunderwave{書}，吾於\CJKunderwave{武成}取二三策而已。仁者無敵於天下，以至仁伐不仁，如何其血流漂杵也？”是言不實也。\CJKunderwave{易·繫辭}云：“斷木為杵，掘地為臼。”是“杵”為臼器也。 \par}

{\noindent\shu\zihao{5}\fzkt “既戊午”至“我師”。正義曰：自此以下皆史辭也，其上闕絕，失其本絕，故文無次第。必是王言既終，史乃更敘戰事。於文次當承“自周，於徵伐商”之下，此句次之，故云“既戊午”也。史官敘事，得言“罔有敵於我師”，稱“我”者，猶如自漢至今，文章之士，雖民論國事，莫不稱“我”，皆雲“我大隨”，以心體國，故稱“我”耳，非要王言乃稱“我”也。 \par}

一戎衣，天下大定。\footnote{衣，服也。一著戎服而滅紂，言與眾同心，動有成功。}乃反商政，政由舊。\footnote{反紂惡政，用商先王善政。}釋箕子囚，封比干墓，式商容閭。\footnote{皆武王反紂政。囚，奴，徒隸。封，益其土。商容，賢人，紂所貶退，式其閭巷以禮賢。}

{\noindent\zhuan\zihao{6}\fzbyks 傳“皆武”至“禮賢”。正義曰：紂囚其人而放釋之，紂殺其身而增封其墓，紂退其人而式其門閭，皆是“武王反紂政”也。下句散其財粟，亦是反紂，於此須有所解，因言之耳。上篇雲“囚奴正士”，\CJKunderwave{論語}云：“箕子為之奴”，是紂囚之，又為奴役之。\CJKunderwave{周禮·司厲職}云：“其奴男子入於罪隸。”\CJKunderline{鄭玄}云：“為之奴者，繫於罪隸之官。”是“囚”為奴，以徒隸役之也。“商容”，賢人之姓名，紂所貶退，處於私室。“式”者，車上之橫木，男子立乘，有所敬則俯而憑式，遂以“式”為敬名。\CJKunderwave{說文}云：“閭，族居里門也。”武王過其閭而式之，言此內有賢人，式之禮賢也。\CJKunderwave{帝王世紀}云：“商容及殷民觀周軍之入，見畢公至，殷民曰：‘是吾新君也。’容曰：‘非也，視其為人嚴乎將有急色，故君子臨事而懼。’見太公至，民曰：‘是吾新君也。’容曰：‘非也，視其為人虎據而鷹趾，當敵將眾，威怒自倍，見利即前，不顧其後，故君子臨眾，果於進退。’見周公至，民曰：‘是吾新君也。’容曰：‘非也，視其為人忻忻休休，志在除賊，是非天子，則周之相國也，故聖人臨眾知之。’見武王至，民曰：‘是吾新君也。’容曰：‘然，聖人為海內討惡，見惡不怒，見善不喜，顏色相副，是以知之。’”是說商容之事也。 \par}

散鹿臺之財，發巨橋之粟，\footnote{紂所積之府倉，皆散發以賑貧民。○散，西旦反。}

{\noindent\zhuan\zihao{6}\fzbyks 傳“紂所”至“貧民”。正義曰：藏財為府，藏粟為倉，故言“紂所積之府倉”也。名曰“鹿臺”，“巨橋”則其義未聞。“散”者言其分佈，“發”者言其開出，互相見也。\CJKunderwave{周本紀}云：“命召公釋箕子之囚，命畢公釋百姓之囚，表商容之閭，命閎夭封比干之墓，命南宮括散鹿臺之錢，發巨橋之粟以賑貧弱也。”然則武王親式商容之閭又表之也。\CJKunderwave{新序}雲“鹿臺其大三里，其高千尺”，則容物多矣。此言“鹿臺之財”，則非一物也。\CJKunderwave{史記}作“錢”，後世追論，以錢為主耳。\CJKunderwave{周禮}有泉府之官，\CJKunderwave{周語}稱景王鑄大錢，是周時已名泉為錢也。 \par}

大賚於四海，而萬姓悅服。\footnote{施捨已債，救乏周無，所謂周有大賚，天下皆悅仁服德。○賚，力代反，徐音來。已音以。債，側界反。周音周，本亦作周。}

{\noindent\zhuan\zihao{6}\fzbyks 傳“施捨”至“服德”。正義曰：\CJKunderwave{左傳}成十八年，晉悼公初立，“施捨，已責”。成二年楚將起師，“已責，救乏”。定五年“歸粟於蔡，以周急，矜無資也”。杜預以為“施恩惠，舍勞役”也，已責，“止逋責”也。皆是恤民之事，故傳引之以證“大賚”。“所謂周有大賚”，\CJKunderwave{論語}文。孔安國解\CJKunderwave{堯曰}之篇，有二帝三王之事，“周有大賚”正指此事，故言“所謂”也。“悅”是歡喜，“服”謂聽從，感恩則悅，見義則服，故“天下皆悅仁服德”也。\CJKunderwave{帝王世紀}云：“王命封墓釋囚，又歸施鹿臺之珠玉及傾宮之女於諸侯，殷民咸喜曰：‘王之於仁人也，死者猶封其墓，況生者乎？王之於賢人也，亡者猶表其閭，況存者乎？王之於財也，聚者猶散之，況其復籍之乎？王之於色也，見在者猶歸其父母，況其復徵之乎？’”是悅服之事也。 \par}

列爵惟五，\footnote{即所識政事而法之。爵五等，公侯伯子男。}

分土惟三。\footnote{列地封國，公侯方百里，伯七十里，子男五十里，為三品。}

{\noindent\zhuan\zihao{6}\fzbyks 傳“列地”至“三品”。正義曰：爵五等，地三品，武王於此既從殷法，未知周公制禮亦然以否。\CJKunderwave{孟子}曰：“北宮錡問於孟子曰：‘周之班爵祿如何？’孟子曰：‘其詳不可得聞矣，嘗聞其略。天子之制，地方千里，公侯方百里，伯七十里，子男五十里。’”\CJKunderwave{漢書·地理志}亦云：“周爵五等，其土三等也，公侯百里，伯七十里，子男五十里。”漢世儒者多以為然，包咸注\CJKunderwave{論語}云：“千乘之國，百里之國也，謂大國惟百里耳。”\CJKunderwave{周禮·大司徒}云：“諸公之地，封疆方五百里。侯四百里，伯三百里，子二百里，男一百里。”蓋是周室既衰，諸侯相併，自以國土寬大，皆違禮文，乃除去本經，妄為說耳。\CJKunderline{鄭玄}之徒以為武王時大國百里，周公制禮大國五百里，\CJKunderwave{王制}之注具矣。 \par}

建官惟賢，\footnote{立官以官賢才。}位事惟能。\footnote{居位理事，必任能事。}重民五教，\footnote{所重在民及五常之教。}

{\noindent\shu\zihao{5}\fzkt “重民五教”。正義曰：以“重”總下五事，民與五教，“食喪祭”也。“五教”所以教民，故與“民”同句。下句“食”與“喪”、“祭”三者各為一事，相類而別，故以“惟”目之，言此皆聖王所重也。\CJKunderwave{論語}云：“所重民食喪祭。”以\CJKunderwave{論語}即是此事，而彼無五教，錄\CJKunderwave{論語}者自略之耳。 \par}

惟食喪祭。\footnote{民以食為命，喪禮篤親愛，祭祀崇孝養，皆聖王所重。○養，羊亮反。}惇信明義，\footnote{使天下厚行信，顯忠義。}崇德報功，\footnote{有德尊以爵，有功報以祿。}垂拱而天下治。\footnote{言武王所修皆是，所任得人，故垂拱而天下治。○任，而鴆反。治，直吏反。}

{\noindent\shu\zihao{5}\fzkt “垂拱而天下治”。正義曰：\CJKunderwave{說文}云：“拱，斂手也。”“垂拱而天下治”，謂所任得人，人皆稱職，手無所營，下垂其拱，故美其“垂拱而天下治”也。 \par}

%%% Local Variables:
%%% mode: latex
%%% TeX-engine: xetex
%%% TeX-master: "../Main"
%%% End:

%% -*- coding: utf-8 -*-
%% Time-stamp: <Chen Wang: 2024-04-02 11:42:41>

% {\noindent \zhu \zihao{5} \fzbyks } -> 注 (△ ○)
% {\noindent \shu \zihao{5} \fzkt } -> 疏

\chapter{卷十二}


\section{洪範第六（洪範上、下）}


\CJKunderline{武王}勝殷，殺\CJKunderline{受}，立\CJKunderline{武庚}，\footnote{不放而殺，紂自焚也。\CJKunderline{武庚}，紂子，以為王者後；一名祿父。勝，商證反。父音甫。}以\CJKunderline{箕子}歸，作\CJKunderwave{洪範}。\footnote{歸鎬京，\CJKunderline{箕子}作之。○範音範。鎬，胡老反，本又作鄗，武王所都也。}

{\noindent\zhuan\zihao{6}\fzbyks 傳“不放”至“祿父”。正義曰：放桀也。湯放桀，此不放而殺之者，紂自焚而死也。\CJKunderwave{殷本紀}雲“紂兵敗，紂走入登鹿臺，衣其寶玉衣，赴火而死。武王遂斬紂頭懸之太白旗”是也。\CJKunderwave{泰誓}雲“取彼兇殘”，則志在於殺也。死猶斬之，則生亦不放，傳據實而言之耳。\CJKunderwave{本紀}雲“封紂子\CJKunderline{武庚}祿父以續殷祀”，是以為王者後也。\CJKunderwave{本紀}“\CJKunderline{武庚}祿父”雙言之，伏生\CJKunderwave{尚書}雲“武王勝殷，繼公子祿父”，是一名祿父也。鄭云：“\CJKunderline{武庚}字祿父，春秋之世有齊侯祿父、蔡侯考父、季孫行父，父亦是名，未必為字，故傳言‘一名祿父’。” \par}

{\noindent\zhuan\zihao{6}\fzbyks 傳“歸”至“作之”。正義曰：上篇雲“至於豐”者，文王之廟在豐，至豐先告廟耳。時王都在鎬，知“歸”者，“歸鎬京”也。此經文旨異於餘篇，非直問答而已，不是史官敘述，必是\CJKunderline{箕子}既對武王之問，退而自撰其事，故傳特雲“\CJKunderline{箕子}作之”。\CJKunderwave{書傳}云：“武王釋\CJKunderline{箕子}之囚，\CJKunderline{箕子}不忍周之釋，走之朝鮮。武王聞之，因以朝解封之。\CJKunderline{箕子}既受周之封，不得無臣禮，故於十三祀來朝，武王因其朝而問洪範。”案此序云：“勝殷，以\CJKunderline{箕子}歸。”明既釋其囚，即以歸之，不令其走去而後來朝也。又朝鮮去周，路將萬里，聞其所在，然後封之，受封乃朝，必歷年矣，不得仍在十三祀也。\CJKunderwave{宋世家}云：“既作\CJKunderwave{洪範}，武王乃封\CJKunderline{箕子}於朝鮮。”得其實也。 \par}

{\noindent\shu\zihao{5}\fzkt “武王”至“洪範”。正義曰：武王伐殷，既勝，殺受，立其子\CJKunderline{武庚}為殷後，以\CJKunderline{箕子}歸鎬京，訪以天道，\CJKunderline{箕子}為陳天地之大法，敘述其事，作\CJKunderwave{洪範}。此惟當言“\CJKunderline{箕子}歸”耳，乃言“殺受，立\CJKunderline{武庚}”者，序自相顧為文。上\CJKunderwave{武成}序云：“武王伐紂”，故此言勝之，下\CJKunderwave{微子之命}序雲“黜殷命，殺\CJKunderline{武庚}”，故此言立之，敘言此以順上下也。 \par}

洪範\footnote{洪，大。範，法也。言天地之大法。}


{\noindent\zhuan\zihao{6}\fzbyks 傳“洪大”至“大法”。正義曰：“洪，大”、“範，法”皆\CJKunderwave{釋詁}文。 \par}

{\noindent\shu\zihao{5}\fzkt “洪範”。正義曰：此經開源於首，覆更演說，非復一問一答之勢，必是\CJKunderline{箕子}自為之也。發首二句，自記被問之年，自“王乃言”至“彝倫攸敘”，王問之辭。自“\CJKunderline{箕子}乃言”至“彝倫攸敘”，言\CJKunderline{禹}得九疇之由。自“初一曰”至“威用六極”，言\CJKunderline{禹}第敘九疇之次。自“一五行”已下，\CJKunderline{箕子}更條說九疇之義。此條說者，當時亦以對王，更復退而修撰，定其文辭，使成典教耳。 \par}

惟十有三祀，王訪於\CJKunderline{箕子}。\footnote{商曰祀，\CJKunderline{箕子}稱祀，不忘本。此年四月歸宗，周先告武成，次問天道。}王乃言曰：“嗚呼！\CJKunderline{箕子}，惟天陰騭下民，相協厥居，\footnote{騭，定也。天不言而默定下民，是助合其居，使有常生之資。○陰，默也，馬云：“覆也。”騭，之逸反，馬云：“升也。升猶舉也，舉猶生也。”相，息亮反，助也。}我不知其彝倫攸敘。”\footnote{言我不知天所以定民之常道理次敘。問何由。○彝，以之反。}


{\noindent\zhuan\zihao{6}\fzbyks 傳“商曰”至“天道”。正義曰：“商曰祀，周曰年”，\CJKunderwave{釋天}文。案此\CJKunderwave{周書}也，\CJKunderwave{泰誓}稱“年”，此獨稱“祀”，故解之“\CJKunderline{箕子}稱祀，不忘本”也。此篇\CJKunderline{箕子}所作，\CJKunderline{箕子}商人，故記傳引此篇者，皆雲“\CJKunderwave{商書}曰”，是\CJKunderline{箕子}自作明矣。序言“歸，作\CJKunderwave{洪範}”，似歸即作之，嫌在\CJKunderwave{武成}之前，故云“此年四月歸宗周，先告武成，次問天道”，以次在\CJKunderwave{武成}之後，故知“先告武成”也。 \par}

{\noindent\zhuan\zihao{6}\fzbyks 傳“騭定”至“之資”。正義曰：傳以“騭”即質也，質訓為成，成亦定義，故為定也。言民是上天所生，形神天之所授，故“天不言而默定下民”。群生受氣流形，各有性靈心識，下民不知其然，是天默定也。相，助也。協，和也。“助合其居”者，言民有其心，天佑助之，令其諧合其生。出言是非，立行得失，衣食之用，動止之宜，無不稟諸上，天乃得諧合。失道則死，合道則生，言天非徒賦命於人，授以形體心識，乃復佑助諧合其居業，使有常生之資。九疇施之於民，皆是天助之事也。此問答皆言“乃”者，以天道之大，沈吟乃問，思慮乃答。宣八年\CJKunderwave{公羊傳}曰：“乃，緩辭也。”王肅以“陰騭下民”一句為天事，“相協”以下為民事，注云：“陰，深也。言天深定下民，與之五常之性，王者當助天和合其居所行天之性，我不知常道倫理所以次敘，是問承天順民，何所由。”與孔異也。 \par}

{\noindent\shu\zihao{5}\fzkt “惟十”至“攸敘”。正義曰：此\CJKunderline{箕子}陳王問已之年，被問之事。惟文王受命十有三祀，武王訪問於\CJKunderline{箕子}，即陳其問辭，王乃言曰：“嗚呼！\CJKunderline{箕子}，此上天不言而默定下民，佑助諧合其安居，使有常生之資。我不知此天之定民常道所以次敘。”問天意何由也。 \par}

\CJKunderline{箕子}乃言曰：“我聞在昔，\CJKunderline{鯀}陻洪水，汩陳其五行。\footnote{陻，塞。汩，亂也。治水失道，亂陳其五行。○\CJKunderline{鯀}，工本反。陻音因。汩，工忽反。行，戶更反。}帝乃震怒，不畀洪範九疇，彝倫攸斁。\footnote{畀，與。斁，敗也。天動怒\CJKunderline{鯀}，不與大法九疇。疇，類也。故常道所以敗。○畀，必二反，徐甫至反，注同。斁，多路反，徐同路反，敗也。}


{\noindent\zhuan\zihao{6}\fzbyks 傳“陻塞”至“五行”。正義曰：襄二十五年\CJKunderwave{左傳}說陳之伐鄭雲“井陻木刊”，謂塞其井，斬其木，是“陻”為塞也。“汩”是亂之意，故為亂也。水是五行之一，水性下流，\CJKunderline{鯀}反塞之，失水之性，水失其性，則五行皆失矣。是塞洪水為亂，陳其五行，言五行陳列皆亂也。\CJKunderwave{大禹謨}帝美\CJKunderline{禹}治水之功云：“地平天成。”傳云：“水土治曰平，五行敘曰成。”水既治，五行序，是治水失道，為亂五行也。 \par}

{\noindent\zhuan\zihao{6}\fzbyks 傳“畀與”至“以敗”。正義曰：“畀，與”，\CJKunderwave{釋詁}文。“斁，敗”，相傳訓也。以\CJKunderline{禹}得而\CJKunderline{鯀}不得，故為天動威怒\CJKunderline{鯀}，不與大法九疇。“疇”是輩類之名，故為類也。言其每事自相類者有九，九者各有一章，故\CJKunderwave{漢書}謂之為九章。此謂九類，是天之常道，既不得九類，故常道所以敗也。自古以來得九疇者惟有\CJKunderline{禹}耳，未聞餘人有得之者也。若人皆得之，\CJKunderline{鯀}獨不得，可言天帝怒\CJKunderline{鯀}。餘人皆不得，獨言天怒\CJKunderline{鯀}者，以\CJKunderline{禹}由治水有功，故天賜之，\CJKunderline{鯀}亦治水而天不與，以\CJKunderline{鯀}\CJKunderline{禹}俱是治水，父不得而子得之，所以彰\CJKunderline{禹}之聖當於天心，故舉\CJKunderline{鯀}以彰\CJKunderline{禹}也。 \par}

\CJKunderline{鯀}則殛死，\CJKunderline{禹}乃嗣興，\footnote{放\CJKunderline{鯀}至死不赦。嗣，繼也。廢父興子，堯舜之道。○殛，紀力反，本或作極，音同。}天乃錫\CJKunderline{禹}洪範九疇，彝倫攸敘。\footnote{天與\CJKunderline{禹}洛出書，神龜負文而出，列於背，有數至於九。\CJKunderline{禹}遂因而第之，以成九類，常道所以次敘。○錫，星曆反。}

{\noindent\zhuan\zihao{6}\fzbyks 傳“放\CJKunderline{鯀}”至“之道”。正義曰：傳嫌“殛”謂被誅殺，故辨之雲“放\CJKunderline{鯀}至死不赦”也。“嗣，繼”，\CJKunderwave{釋詁}文。三代以還，父罪子廢，故云“廢父興子，堯舜之道”。賞罰各從其實，為天下之至公也。 \par}

{\noindent\zhuan\zihao{6}\fzbyks 傳“天與”至“次敘”。正義曰：\CJKunderwave{易·繫辭}云：“河出圖，洛出書，聖人則之。”九類各有文字，即是書也。而云“天乃錫\CJKunderline{禹}”，知此天與\CJKunderline{禹}者即是\CJKunderwave{洛書}也。\CJKunderwave{漢書·五行志}：“劉歆以為伏羲系天而王，河出圖，則而畫之，八卦是也。\CJKunderline{禹}治洪水，錫\CJKunderwave{洛書}，法而陳之，\CJKunderwave{洪範}是也。”先達共為此說，龜負\CJKunderwave{洛書}，經無其事，\CJKunderwave{中候}及諸緯多說黃\CJKunderline{帝堯}、舜\CJKunderline{禹}、湯文武受圖書之事，皆雲龍負圖，龜負書。緯候之書，不知誰作，通人討核，謂偽起哀平，雖復前漢之末，始有此書，以前學者必相傳此說，故孔以九類是神龜負文而出，列於背，有數從一而至於九。\CJKunderline{禹}見其文，遂因而第之，以成此九類法也。此九類陳而行之，常道所以得次敘也。言\CJKunderline{禹}第之者，以天神言語，必當簡要，不應曲有次第。丁寧若此，故以為\CJKunderline{禹}次第之。\CJKunderline{禹}既第之，當有成法可傳，應人盡知之，而武王獨問\CJKunderline{箕子}者，\CJKunderwave{五行志}云：“聖人行其道而寶其真，降及於殷，\CJKunderline{箕子}在父師之位而典之。周既克殷，以\CJKunderline{箕子}歸周，武王親虛已而問焉。”言\CJKunderline{箕子}典其事，故武王特問之，其義或當然也。若然，\CJKunderline{大禹}既得九類，常道始有次敘，未有\CJKunderwave{洛書}之前，常道所以不亂者，世有澆淳，教有疏密，三皇已前，無文亦治，何止無\CJKunderwave{洛書}也。但既得九類以後，聖王法而行之，從之則治，違之則亂，故此說常道攸敘攸斁由\CJKunderwave{洛書}耳。 \par}

{\noindent\shu\zihao{5}\fzkt “\CJKunderline{箕子}”至“攸敘”。正義曰：\CJKunderline{箕子}乃言，答王曰：“我聞在昔，\CJKunderline{鯀}障塞洪水，治水失道，是乃亂陳其五行而逆天道也。天帝乃動其威怒，不與\CJKunderline{鯀}大法九類，天之常道所以敗也。\CJKunderline{鯀}則放殛，至死不赦。\CJKunderline{禹}以聖德繼父而興，代治洪水，決道使通，天乃賜\CJKunderline{禹}大法九類，天之常道所以得其次敘。”此說其得九類之由也。 \par}

“初一曰五行，\footnote{九類，類一章，以五行為始。}次二曰敬用五事，\footnote{五事在身，用之必敬乃善。}次三曰農用八政，\footnote{農，厚也，厚用之政乃成。○農，馬云：“食為八政之首，故以農名之。”}次四曰協用五紀，\footnote{協，和也，和天時使得正用五紀。}次五曰建用皇極，\footnote{皇，大。極，中也。凡立事當用大中之道。}次六曰乂用三德，\footnote{治民必用剛柔正直之三德。}次七曰明用稽疑，\footnote{明用卜筮考疑之事。}


{\noindent\zhuan\zihao{6}\fzbyks 傳“農厚”至“乃成”。正義曰：\CJKunderline{鄭玄}云：“農讀為醲。”則“農”是醲意，故為厚也。政施於民，善不厭深，故“厚用之政乃成”也。張晏、王肅皆言“農，食之本也。食為八政之首，故以農言之”。然則農用止為一食，不兼八事，非上下之例，故傳不然。“八政”、“三德”總是治民，但“政”是被物之名，“德”是在己之稱，故分為二疇也。 \par}

{\noindent\zhuan\zihao{6}\fzbyks 傳“協和”至“五紀”。正義曰：“協，和”，\CJKunderwave{釋詁}文。天是積氣，其狀無形，列宿四方，為天之限。天左行，晝夜一周。日月右行，日遲月疾。周天三百六十五度有餘，日則日行一度，月則日行十三度有餘，日月行於星辰，乃為天之歷數。和此天時，令不差錯，使行得正用五紀也。日月逆天道而行，其行又有遲疾，故須調和之。 \par}

{\noindent\zhuan\zihao{6}\fzbyks 傳“皇大”至“之道”。正義曰：“皇，大”，\CJKunderwave{釋詁}文。“極”之為中，常訓也。凡所立事，王者所行皆是，無得過與不及，常用大中之道也。\CJKunderwave{詩}雲“莫匪爾極”，\CJKunderwave{周禮}“以為民極”，\CJKunderwave{論語}“允執其中”，皆謂用大中也。 \par}

次八曰念用庶徵，次九曰鄉用五福，威用六極。\footnote{言天所以鄉勸人用五福，所以威沮人用六極。此已上\CJKunderline{禹}所第敘。○向，許亮反，又許兩反。沮，在汝反。此已上，時掌反。\CJKunderline{禹}所第敘，馬云：“從‘五行’已下至‘六極’，\CJKunderwave{洛書}文也。\CJKunderwave{漢書·五行志}以‘初一’已下皆\CJKunderwave{洛書}文也。”}

{\noindent\zhuan\zihao{6}\fzbyks 傳“言天”至“第敘”。正義曰：“貧”、“弱”等六者，皆謂窮極惡事，故目之“六極”也。“福”者人之所慕，皆鄉望之。“極”者人之所惡，皆畏懼之。“勸”，勉也，勉之為善。“沮”止也，止其為惡。福、極皆上天為之，言天所以鄉望勸勉人用五福，所以畏懼沮止人用六極，自“初一曰”已下至此“六極”已上，皆是\CJKunderline{禹}所次第而敘之。下文更將此九類而演說之，知此九者皆\CJKunderline{禹}所第也。\CJKunderline{禹}為此次者，蓋以五行世所行用，是諸事之本，故“五行”為初也。發見於人則為五事，故“五事”為二也。正身而後及人，施人乃名為政，故“八政”為三也。施人之政，用天之道，故“五紀”為四也。順天佈政，則得大中，故“皇極”為五也。欲求大中，隨德是任，故“三德”為六也。政雖在德，事必有疑，故“稽疑”為七也。行事在於政，得失應於天，故“庶徵”為八也。天監在下，善惡必報，休咎驗於時氣，禍福加於人身，故“五福”、“六極”為九也。“皇極”居中者，總包上下，故“皇極”傳雲“大中之道”。大立其有中，謂行九疇之義是也。“福”、“極”處末者，顧氏云：“前八事俱得，五福歸之。前八事俱失，六極臻之。故福極處末也。”發首言“初一”，其末不言“終九”者，數必以一為始，其九非數之終，故從上言“次”而不言“終”也。五行不言“用”者，五行萬物之本，天地百物莫不用之，不嫌非用也。傳於“五福”、“六極”言天用者，以前並是人君所用，五福六極受之於天，故言天用。傳言此“\CJKunderline{禹}所第敘”，不知\CJKunderwave{洛書}本有幾字。\CJKunderwave{五行志}悉載此一章，乃云：“凡此六十五字，皆\CJKunderwave{洛書}本文。”計天言簡要，必無次第之數。上傳雲“\CJKunderline{禹}因而第之”，則孔以第是\CJKunderline{禹}之所為，“初一曰”等二十七字必是\CJKunderline{禹}加之也。其“敬用”、“農用”等一十八字，大劉及顧氏以為龜背先有總三十八字。小劉以為“敬用”等亦\CJKunderline{禹}所第敘，其龜文惟有二十字。並無明據，未知孰是，故兩存焉。“皇極”不言數者，以總該九疇，理兼萬事，非局數能盡故也。“稽疑”不言數者，以卜五筮二，共成為七，若舉卜不得兼筮，舉筮不得兼卜，且疑事既眾，不可以數總之故也。“庶徵”不言數者，以“庶徵”得為五休，失為五咎，若舉休不兼咎，舉咎不兼休，若休咎並言，便為十事，本是五物，不可言十也。然“五福”、“六極”所以善惡皆言者，以沮勸在下，故丁寧明言善惡也。且“庶徵”雖有休咎，皆以念慮包之。“福”、“極”鄉威相反，不可一言為目，故別為文焉。知“五福”、“六極”非各分為疇，所以其為一者，蓋以龜文“福”、“極”相近一處，故\CJKunderline{禹}第之總為一疇。等行五事，所以福五而極六者，大劉以為“皇極”若得，則分散總為五福，若失則不能為五事之主，與五事並列其咎弱，故為六也。猶\CJKunderwave{詩}平王以後與諸侯並列同為國風焉。咎徵有五而極有六者，\CJKunderwave{五行傳}云：“皇之不極，厥罰常陰。”即與咎徵“常雨”相類，故以“常雨”包之為五也。 \par}

{\noindent\shu\zihao{5}\fzkt “初一”至“六極”。正義曰：天所賜\CJKunderline{禹}大法九類者，初一曰五材氣性流行，次二曰敬用在身五種之行事，次三曰厚用接物八品之政教，次四曰和用天象五物之綱紀，次五曰立治用大為中正之道，次六曰治民用三等之德，次七曰明用小筮以考疑事，次八曰念用天時眾氣之應驗，次九曰鄉勸人用五福，威沮人用六極。此九類之事也。 \par}

“一，五行。一曰水，二曰火，三曰木，四曰金，五曰土。\footnote{皆其生數。}


{\noindent\zhuan\zihao{6}\fzbyks 傳“皆其生數”。正義曰：\CJKunderwave{易·繫辭}曰：“天一，地二，天三，地四，天五，地六，天七，地八，天九，地十。”此即是五行生成之數。天一生水，地二生火，天三生木，地四生金，天五生土，此其生數也。如此則陽無匹，陰無耦，故地六成水，天七成火，地八成木，天九成金，地十成土，於是陰陽各有匹偶，而物得成焉，故謂之成數也。\CJKunderwave{易·繫辭}又曰“天數五，地數五，五位相得而各有合，此所以成變化而行鬼神”，謂此也。又數之所起，起於陰陽。陰陽往來，在於日道。十一月冬至日南極，陽來而陰往。冬，水位也，以一陽生為水數。五月夏至日北極，陰進而陽退。夏，火位也，當以一陰生為火數。但陰不名奇，數必以偶，故以六月二陰生為火數也。是故\CJKunderwave{易說}稱幹貞於十一月子，坤貞於六月未，而皆左行，由此也。冬至以及於夏至，當為陽來。正月為春木位也，三陽已生，故三為木數。夏至以及冬至，當為陰進。八月為秋金位也，四陰已生，故四為金數。三月春之季，四季土位也，五陽已生，故五為土數，此其生數之由也。又萬物之本，有生於無，者生於微，及其成形，亦以微著為漸。五行先後，亦以微著為次。五行之體，水最微，為一。火漸著，為二。木形實，為三。金體固，為四。土質大，為五。亦是次之宜。大劉與顧氏皆以為水火木金，得土數而成，故水成數六，火成數七，木成數八，金成數九，土成數十。義亦然也。 \par}

水曰潤下，火曰炎上，\footnote{言其自然之常性。○炎，榮鉗反。上，時掌反，又如字，下同。}木曰曲直，金曰從革，\footnote{木可以揉曲直，金可以改更。○揉，如酉反。}土爰稼穡。\footnote{種曰稼，斂曰穡。土可以種，可以斂。}


{\noindent\zhuan\zihao{6}\fzbyks 傳“言其自然之常性”。正義曰：\CJKunderwave{易·文言}云：“水流溼，火就燥。”王肅曰：“水之性潤萬物而退下，火之性炎盛而升上。”是“潤下”、“炎上”，言其自然之本性。 \par}

{\noindent\zhuan\zihao{6}\fzbyks 傳“木可”至“改更”。正義曰：此亦言其性也，“揉曲直”者，為器有須曲直也。“可改更”者，可銷鑄以為器也。木可以揉令曲直，金可以從人改更，言其可為人用之意也。由此而觀，水則潤下，可用以灌溉；火則炎上，可用以炊爨，亦可知也。水既純陰，故潤下趣陰。火是純陽，故炎上趣陽。木金陰陽相雜，故可曲直改更也。 \par}

{\noindent\zhuan\zihao{6}\fzbyks 傳“種曰”至“以斂”。正義曰：\CJKunderline{鄭玄}\CJKunderwave{周禮注}云：“種穀曰稼，若嫁女之有所生。”然則“穡”是惜也，言聚畜之可惜也。共為治田之事，分為“種”、“斂”二名耳。土上所為，故為土性。上文“潤下”、“炎上”、“曲直”、“從革”，即是水火木金體有本性。其稼穡以人事為名，非是土之本性，生物是土之本性，其稼穡非土本性也。“爰”亦“曰”也，變“曰”言“爰”，以見此異也。“六府”以“土”、“谷”為二，由其體異故也。 \par}

潤下作咸，\footnote{水滷所生。○咸音咸。滷音魯。}炎上作苦，\footnote{焦氣之味。}曲直作酸，\footnote{木實之性。}從革作辛，\footnote{金之氣味。}稼穡作甘。\footnote{甘味生於百穀。五行以下，\CJKunderline{箕子}所陳。}

{\noindent\zhuan\zihao{6}\fzbyks 傳“水滷所生”。正義曰：水性本甘，久浸其地，變而為滷，滷味乃咸。\CJKunderwave{說文}云：“滷，西方咸地。東方謂之斥，西方謂之滷。”\CJKunderwave{禹貢}云：“海濱廣斥。”是海浸其旁地，使之咸也。\CJKunderwave{月令·冬}雲“其味咸，其臭朽”是也。土言“曰”者，言其本性。此言“作”者，從其發見。指其體則稱“曰”，致其類即言“作”。下“五事”、“庶徵”言“曰”、“作”者，義亦然也。 \par}

{\noindent\zhuan\zihao{6}\fzbyks 傳“焦氣之味”。正義曰：火性炎上，焚物則焦，焦是苦氣。\CJKunderwave{月令·夏}雲“其臭焦，其味苦”，苦為焦味，故云“焦氣之味”也。臭之曰“氣”，在口曰“味”。 \par}

{\noindent\zhuan\zihao{6}\fzbyks 傳“木實之性”。正義曰：木生子實，其味多酸，五果之味雖殊，其為酸一也，是木實之性然也。\CJKunderwave{月令·春}雲“其味酸，其臭羶”是也。 \par}

{\noindent\zhuan\zihao{6}\fzbyks 傳“金之氣味”。正義曰：金之在火，別有腥氣，非苦非酸，其味近辛，故辛為金之氣味。\CJKunderwave{月令·秋}雲“其位辛，其臭腥”是也。 \par}

{\noindent\zhuan\zihao{6}\fzbyks 傳“甘味生於百穀”。正義曰：“甘味生於百穀”，谷是土之所生，故甘為土之味也。\CJKunderwave{月令·中央}雲“其味甘，其臭香”是也。 \par}

{\noindent\shu\zihao{5}\fzkt “一五行”至“作甘”。正義曰：此以下\CJKunderline{箕子}所演陳\CJKunderline{禹}所第疇名於上，條列說以成之。此章所演，文有三重，第一言其名次，第二言其體性，第三言其氣味，言五者性異而味別，各為人之用。\CJKunderwave{書傳}云：“水火者百姓之所飲食也，金木者百姓之所興作也，土者萬物之所資生也。是為人用。”“五行”即五材也，襄二十七年\CJKunderwave{左傳}雲“天生五材，民並用之”，言五者各有材幹也。謂之“行”者，若在天則五氣流行，在地世所行用也。 \par}

“二，五事。一曰貌，\footnote{容儀。○貌，本亦作{\hanaa 䫉}。}二曰言，\footnote{詞章。}三曰視，\footnote{觀正。○視，常止反，徐市止反。}四曰聽，\footnote{察是非。}


{\noindent\zhuan\zihao{6}\fzbyks 傳“察是非”。正義曰：此五事皆有是非，\CJKunderwave{論語}云：“非禮勿視，非禮勿聽，非禮勿言，非禮勿動。”又引\CJKunderwave{詩}云：“思無邪。”故此五事皆有是非也。此經歷言五名，名非善惡之稱，但為之有善有惡，傳皆以是辭釋之。“貌”者言其動有容儀也，“言”道其語有辭章也，“視”者言其觀正不觀邪也，“聽”者受人言察是非也，“思”者心慮所行使行得中也。傳於“聽”雲“察是非”，明五者皆有是非也，所為者為正不為邪也。於“視”不言“視邪正”，於“聽”言“察是非”，亦所以互相明也。 \par}

五曰思，\footnote{心慮所行。○思如字，徐息吏反，下同。}貌曰恭，\footnote{儼恪。○儼，魚檢反。}言曰從，\footnote{是則可從。}視曰明，\footnote{必清審。}聽曰聰，\footnote{必微諦。諦音帝。}思曰睿。\footnote{必通於微。睿，悅歲反，馬云：“通也。”}


{\noindent\zhuan\zihao{6}\fzbyks 傳“必通於微”。正義曰：此一重言敬用之事。貌戒惰容，故“恭”為儼恪。\CJKunderwave{曲禮}曰：“儼若思。”“儼”是嚴正之貌也。“恪”，敬也，貌當嚴正而莊敬也。言非理則人違之，故言是則可從也。視必明於善惡，故必清徹而審察也。聽當別彼是非，必微妙而審諦也。王肅云：“睿，通也。思慮苦其不深，故必深思使通於微也。”此皆敬用使然，故經以善事明之。\CJKunderline{鄭玄}云：“此恭、明、聰、睿行之於我身，其從則是彼人從我，以與上下違者，我是而彼從，亦我所為不乖倒也。”此據人主為文，皆是人主之事，\CJKunderwave{說命}雲“接下思恭，視遠惟明，聽德惟聰”，即此是也。 \par}

恭作肅，\footnote{心敬。}從作乂，\footnote{可以治。}明作晢，\footnote{照了。○晢，之舌反，徐丁列反，又之世反。}聰作謀，\footnote{所謀必成當。○當，丁浪反。}睿作聖。\footnote{於事無不通謂之聖。}

{\noindent\zhuan\zihao{6}\fzbyks 傳“於事”至“之聖”。正義曰：此一重言所致之事也。恭在貌而敬在心，人有心慢而貌恭，必當緣恭以致敬，故貌恭作心敬也。下從上則國治，故人主言必從，其國可以治也。視能清審，則照了物情，故視明致照晢也。聽聰則知其是非，從其是為謀必當，故聽聰致善謀也。睿、聖俱是通名，聖大而睿小，緣其能通微，事事無不通，因睿以作聖也。\CJKunderline{鄭玄}\CJKunderwave{周禮注}云：“聖通而先識也。”是言識事在於眾物之先，無所不通，以是名之為聖。聖是智之上，通之大也。此言人主行其小而致其大，皆是人主之事也。鄭云：“皆謂其政所致也。君貌恭則臣禮肅，君言從則臣職治，君視明則臣照晢，君聽聰則臣進謀，君思睿則臣賢智。”鄭意謂此所致皆是君致臣也。案“庶徵”之意，休徵、咎徵皆肅、乂所致，若肅、乂、明、聰皆是臣事，則休、咎之所致，悉皆不由君矣。又聖大而睿小，若君睿而致臣聖，則臣皆上於君矣，何不然之甚乎！“晢”字王肅及\CJKunderwave{漢書·五行志}皆云：“晢，智也。”定本作“晢”，則讀為哲。 \par}

{\noindent\shu\zihao{5}\fzkt “二五”至“作聖”。正義曰：此章所演亦為三重，第一言其所名，第二言其所用，第三言其所致。“貌”是容儀，舉身之大名也，“言”是口之所出，“視”是目之所見，“聽”是耳之所聞，“思”是心之所慮，一人之上有此五事也。貌必須恭，言必可從，視必當明，聽必當聰，思必當通於微密也。此一重即是敬用之事。貌能恭，則心肅敬也。言可從，則政必治也。視能明，則所見照晢也。聽能聰，則所謀必當也。思通微，則事無不通，乃成聖也。此一重言其所致之事。\CJKunderwave{洪範}本體與人主作法，皆據人主為說。貌總身也，口言之，目視之，耳聽之，心慮之，人主始於敬身，終通萬事，此五事為天下之本也。五事為此次者，鄭云：“此數本諸陰陽，昭明人相見之次也。”\CJKunderwave{五行傳}曰：“貌屬木，言屬金，視屬火，聽屬水，思屬土。”\CJKunderwave{五行傳}伏生之書也。孔於\CJKunderline{太戊}桑谷之下雲“七日大拱，貌不恭之罰”，\CJKunderline{高宗}雊雉之下雲“耳不聰之異”，皆\CJKunderwave{書傳}之文也。孔取\CJKunderwave{書傳}為說，則此次之意亦當如\CJKunderwave{書傳}也。木有華葉之容，故貌屬木。言之決斷若金之斬割，故言屬金。火外光，故視屬火。水內明，故聽屬水。土安靜而萬物生，心思慮而萬事成，故思屬土。又於\CJKunderwave{易}東方震為足，足所以動容貌也。西方兌為口，口出言也。南方離為目，目視物也。北方坎為耳，耳聽聲也。土在內，猶思在心。亦是五屬之義也。 \par}

“三，八政。一曰食，\footnote{勤農業。}二曰貨，\footnote{寶用物。}三曰祀，\footnote{敬鬼神以成教。}四曰司空，\footnote{主空土以居民。}五曰司徒，\footnote{主徒眾，教以禮義。}六曰司寇，\footnote{主奸盜，使無縱。○縱，子用反，或作從，音同。}七曰賓，\footnote{禮賓客，無不敬。}八曰師。\footnote{簡師所任必良，士卒必練。○卒，子忽反。}

{\noindent\zhuan\zihao{6}\fzbyks 傳“寶用物”。正義曰：“貨”者，金玉布帛之總名，皆為人用，故為“用物”。\CJKunderwave{旅獒}雲“不貴異物賤用物”是也。食則勤農以求之，衣則蠶績以求之，但貨非獨衣，不可指言求處，故云得而寶愛之。\CJKunderwave{孝經}云：“謹身節用。”\CJKunderwave{詩序}云：“儉以足用。”是寶物也。 \par}

{\noindent\zhuan\zihao{6}\fzbyks 傳“主空土以居民”。正義曰：\CJKunderwave{周官}篇云：“司空掌邦土，居四民，時地利。司徒掌邦教，敷五典，擾兆民。司寇掌邦禁，詰奸慝，刑暴亂。”\CJKunderwave{周禮}司徒教以禮義，司寇無縱罪人，其文具矣。 \par}

{\noindent\zhuan\zihao{6}\fzbyks 傳“簡師”至“必練”。正義曰：經言“賓”、“師”，當有賓師之法，故傳以“禮賓客，無不敬”，教民待賓客相往來也。“師”者，眾之通名，必當選人為之，故傳言“簡師”，選人為師也。“所任必良”，任良將也。“士卒必練”，“練”謂教習使知義，若練金使精也。\CJKunderwave{論語}：“以不教民戰，是謂棄之。”是士卒必須練也。 \par}

{\noindent\shu\zihao{5}\fzkt “三八政”至“曰師”。正義曰：“八政”者，人主施政教於民有八事也。一曰食，教民使勤農業也。二曰貨，教民使求資用也。三曰祀，教民使敬鬼神也。四曰司空之官，主空土以居民也。五曰司徒之官，教眾民以禮義也。六曰司寇之官，詰治民之奸盜也。七曰賓，教民以禮待賓客，相往來也。八曰師，立師防寇賊，以安保民也。八政如此次者，人不食則死，食於人最急，故食為先也。有食又須衣貨為人之用，故“貨”為二也。所以得食貨，乃是明靈祐之，人當敬事鬼神，故“祀”為三也。足衣食、祭鬼神，必當有所安居，司空主居民，故“司空”為四也。雖有所安居，非禮義不立，司徒教以禮義，故“司徒”為五也。雖有禮義之教，而無刑殺之法，則彊弱相陵，司寇主奸盜，故“司寇”為六也。民不往來，則無相親之好，故“賓”為七也。寇賊為害，則民不安居，故“師”為八也。此用於民緩急而為次也。“食”、“貨”、“祀”、“賓”、“師”指事為之名，三卿舉官為名者，三官所主事多，若以一事為名，則所掌不盡，故舉官名以見義。\CJKunderline{鄭玄}云：“此數本諸其職先後之宜也。食謂掌民食之官，若后稷者也。貨掌金帛之官，若\CJKunderwave{周禮}司貨賄是也。祀掌祭祀之官，若宗伯者也。司空掌居民之官。司徒掌教民之官也。司寇掌詰盜賊之官。賓掌諸侯朝覲之官，\CJKunderwave{周禮}大行人是也。師掌軍旅之官，若司馬也。”王肅云：“賓掌賓客之官也。”即如鄭、王之說，自可皆舉官名，何獨三事舉官也？八政主以教民，非謂公家之事，司貨賄掌公家貨賄，大行人掌王之賓客，若其事如\CJKunderwave{周禮}，皆掌王家之事，非復施民之政，何以謂之“政”乎？且司馬在上，司空在下，今司空在四，司馬在八，非取職之先後也。 \par}

“四，五紀。一曰歲，\footnote{所以紀四時。}二曰月，\footnote{所以紀一月。}三曰日，\footnote{紀一日。}四曰星辰，\footnote{二十八宿迭見以敘氣節，十二辰以紀日月所會。○宿音秀。迭，田節反。見，賢遍反。}五曰歷數。\footnote{歷數節氣之度以為歷，敬授民時。}

{\noindent\zhuan\zihao{6}\fzbyks 傳“二十”至“所會”。正義曰：二十八宿，佈於四方，隨天轉運，昏明迭見。\CJKunderwave{月令}十二月皆紀昏旦所中之星。若\CJKunderwave{月令}孟春昏參中，旦尾中；仲春昏弧中，旦建星中；季春昏七星中，旦牽牛中；孟夏昏翼中，旦婺女中；仲夏昏亢中，旦危中；季夏昏心中，旦奎中；孟秋昏建星中，旦畢中；仲秋昏牽牛中，旦觜觿中；季秋昏虛中，旦柳中；孟冬昏危中，旦七星中；仲冬昏東壁中，旦軫中；季冬昏婁中，旦氐中；皆所以敘氣節也。氣節者，一歲三百六十五日有餘，分為十二月，有二十四氣。一為節氣，謂月初也。一為中氣，謂月半也。以彼迭見之星，敘此月之節氣也。昭七年\CJKunderwave{左傳}晉侯問士文伯曰：“多語寡人辰而莫同，何謂也？”對曰：“日月之會是謂辰。”“會”者，日行遲，月行疾，俱循天度而右行，二十九日過半月行一周天，又前及日而與日會，因謂會處為辰。則\CJKunderwave{月令}孟春日在營室，仲春日在奎，季春日在胃，孟夏日在畢，仲夏日在東井，季夏日在柳，孟秋日在翼，仲秋日在角，季秋日在房，孟冬日在尾，仲冬日在鬥，季冬日在婺女，十二會以為十二辰。“辰”即子醜寅卯之謂也，十二辰所以紀日月之會處也。鄭以為“星，五星也”。然五星所行，下民不以為候，故傳不以“星”為五星也。 \par}

{\noindent\zhuan\zihao{6}\fzbyks 傳“歷數”至“民時”。正義曰：天以積氣無形，二十八宿分之為限，每宿各有度數，合成三百六十五度有餘。日月右行，循此宿度。日行一度，月行十三度有餘，二十九日過半而月一周與日會，每於一會謂之一月，是一歲為十二月，仍有餘十一日。為日行天未周，故置閏以充足。若均分天度以為十二次，則每次三十度有餘。一次之內有節氣、中氣，次之所管，其度多每月之所統，其日入月朔，參差不及，節氣不得在月朔，中氣不得在月半。故聖人歷數此節氣之度，使知氣所在，既得氣在之日，以為一歲之歷，所以敬授民時。王肅院“日月星辰所行，布而數之，所以紀度數”是也。“歲”、“月”、“日”、“星”傳皆言“紀”，“歷數”不言“紀”者，歷數數上四事為紀，所紀非獨一事，故傳不得言“紀”。但成彼四事為紀，故通數以為五耳。 \par}

{\noindent\shu\zihao{5}\fzkt “四五紀”至“歷數”。正義曰：“五紀”者，五事為天時之經紀也。一曰歲，從冬至以及明年冬至為一歲，所以紀四時也。二曰月，從朔至晦，大月三十日，小月二十九日，所以紀一月也。三曰日，從夜半以至明日夜半周十二辰為一日，所以紀一日也。四曰星辰，星謂二十八宿，昏明迭見；辰謂日月別行，會於宿度，從子至於醜為十二辰。星以紀節氣早晚，辰以紀日月所會處也。五曰歷數，算日月行道所歷，計氣朔早晚之數，所以為一歲之歷。凡此五者，皆所以紀天時，故謂之“五紀”也。五紀不言“時”者，以歲月氣節正而四時亦自正，時隨月變，非歷所推，故不言“時”也。五紀為此節者，歲統月，月統日，星辰見於天，其曰“歷數”，總歷四者，故歲為始，歷為終也。 \par}

“五，皇極。皇建其有極，\footnote{大中之道，大立其有中，謂行九疇之義。}斂時五福，用敷錫厥庶民。\footnote{斂是五福之道以為教，用布與眾民使慕之。}


{\noindent\zhuan\zihao{6}\fzbyks 傳“大中”至“之義”。正義曰：此疇以“大中”為名，故演其大中之義。“大中之道，大立其有中”，欲使人主先自立其大中，乃以大中教民也。凡行不迂僻則謂之“中”，\CJKunderwave{中庸}所謂“從容中道”，\CJKunderwave{論語}“允執其中”，皆謂此也。九疇為德，皆求大中，是為善之總，故云“謂行九疇之義”，言九疇之義皆求得中，非獨此疇求大中也。此大中是人君之大行，故特敘以為一疇耳。 \par}

{\noindent\zhuan\zihao{6}\fzbyks 傳“斂是”至“慕之”。正義曰：五福生於五事，五事得中，則福報之。“斂是五福之道”，指其敬用五事也。用五事得中，則各得其福，其福乃散於五處，不相集聚。若能五事皆敬，則五福集來歸之。普敬五事，則是斂聚五福之道。以此敬五事為教，布與眾民，使眾民勸慕為之。福在幽冥，無形可見，敬用五事，則能致之，“斂是五福”，正是敬用五事。不言“敬用五事以教”，而云“斂是五福以為教”者，福是善之見者，故言“福”以勸民，欲其慕而行善也。“汝”者，\CJKunderline{箕子}“汝”王也。 \par}

惟時厥庶民於汝極，錫汝保極。\footnote{君上有五福之教，眾民於君取中，與君以安中之善。言從化。}凡厥庶民，無有淫朋，人無有比德，惟皇作極。\footnote{民有安中之善，則無淫過朋黨之惡、比周之德，為天下皆大為中正。○比，毗志反，注同。}

{\noindent\zhuan\zihao{6}\fzbyks 傳“君上”至“從化”。正義曰：凡人皆有善性，善不能自成，必須人君教之，乃得為善。君上有五福之教，以大中教民，眾民於君取中。“保”訓安也，既學得中，則其心安之。君以大中教民，民以大中向君，是民與君皆以大中之善。君有大中，民亦有大中，言從君化也。 \par}

{\noindent\zhuan\zihao{6}\fzbyks 傳“民有”至“中止”。正義曰：民有安中之善，非中不與為交，安中之人則無淫過朋黨之惡，無有比周之德。“朋黨”、“比周”是不中者。善多惡少，則惡亦化而為善，無復有不中之人，惟天下皆大為中正矣。 \par}

{\noindent\shu\zihao{5}\fzkt “五皇極”至“作極”。正義曰：“皇”，大也。“極”，中也。施政教，治下民，當使大得其中，無有邪僻。故演之云，大中者，人君為民之主，當大自立其有中之道，以施教於民。當先敬用五事，以斂聚五福之道，用此為教，布與眾民，使眾民慕而行之。在上能教如此，惟是其眾民皆效上所為，無不於汝人君取其中道而行。積久漸以成性，乃更與汝人君以安中之道。言皆化也。若能化如是，凡其眾民無有淫過朋黨之行，人無有惡相阿比之德，惟皆大為中正之道。言天下眾民盡得中也。 \par}

凡厥庶民，有猷有為有守，汝則念之。\footnote{民戢有道，有所為，有所執守，汝則念錄敘之。}不協於極，不罹於咎，皇則受之。\footnote{凡民之行，雖不合於中，而不罹於咎惡，皆可進用，大法受之。○罹，馬力馳反，又來多反。行，下孟反。}


{\noindent\zhuan\zihao{6}\fzbyks 傳“民戢”至“敘之”。正義曰：“戢”，斂也，因上“斂是五福”，故傳以“戢”言之。“戢”文兼下三事，民能斂德行智，能使其身有道德，其才能有所施為，用心有所執守。如此人者，汝念錄敘之，宜用之為官也。“有所為”，謂藝能也。“有執守”，謂得善事能守而勿失，言其心正不逆邪也。 \par}

{\noindent\zhuan\zihao{6}\fzbyks 傳“凡民”至“受之”。正義曰：“不合於中，不罹於咎”，謂未為大善，又無惡行，是中人已上，可勸勉有方將者也，故皆可進用，以大法受之。“大法”謂用人之法，取其所長，棄瑕錄用也。上文人君以大中教民，使天下皆為大中，此句印敝痢邦不合於中亦用之者，上文言設教耳。其實天下之大，兆民之眾，不可使皆合大中；且庶官交曠，即須任人，不可待人盡合大中，然後敘用。言各有為，不相妨害。 \par}

而康而色，曰：‘予攸好德。’汝則錫之福，\footnote{汝當安汝顏色，以謙下人。人曰：“我所好者德。”汝則與之爵祿。○好，呼報反。下，遐嫁反。}時人斯其惟皇之極。\footnote{不合於中之人，汝與之福，則是人此其惟大之中。言可勉進。}無虐煢獨而畏高明。\footnote{煢，單，無兄弟也。無子曰獨。單獨者，不侵虐之寵貴者，不枉法畏之。○無虐，馬本作亡侮。煢，岐扃反。畏如字，徐云：“鄭音威。”}

{\noindent\zhuan\zihao{6}\fzbyks 傳“汝當”至“爵祿”。正義曰：安汝顏色，以謙下人，其此不合於中人之，皆人言曰：“我所好者德也。”是有慕善之心，有方將者也。汝則與之爵祿以長進之。上句言“受之”，謂治受以，此言“與爵祿”，謂用為官也。 \par}

{\noindent\zhuan\zihao{6}\fzbyks 傳“不合”至“勉進”。正義曰：“不合於中之人”，初時未閤中也，汝與之爵祿，置之朝廷，見人為善，心必慕之，則是人此其惟大中之道，為大中之人，言可勸勉使進也。\CJKunderwave{荀卿書}曰：“蓬生麻中，不扶自直。白沙在涅，與之俱黑。”斯言信矣。此經或言“時人德”，鄭、王諸本皆無“德”字。此傳不以“德”為義，定本無“德”，疑衍字也。 \par}

{\noindent\zhuan\zihao{6}\fzbyks 傳“煢單”至“畏之”。正義曰：\CJKunderwave{詩}云：“獨行煢煢。”是為單，謂無兄弟也。“無子曰獨”，\CJKunderwave{王制}文。“高明”與“煢獨”相對，非謂才高，知寵貴之人位望高也。“不枉法畏之”，即\CJKunderwave{詩}所謂“不畏強禦”是也。此經皆是據天子，無陵虐煢獨而畏避高明寵貴者。顧氏亦以此經據人君，小劉以為據人臣，謬也。 \par}

{\noindent\shu\zihao{5}\fzkt “凡厥”至“高明”。正義曰：又說用人為官，使之大中。凡其眾民，有道德，有所為，有所執守，汝為人君則當念錄敘之，用之為官。若未能如此，雖不合於中，亦不罹於咎惡，此人可勉進，宜以取人大法則受取之。其受人之大法如何乎？汝當和安汝之顏色，以謙下人。彼欲仕者謂汝曰：“我所好者德也。”汝則與之以福祿，隨其所能，用之為官。是人庶幾必自勉進，此其惟為大中之道。又為君者無侵虐單獨而畏忌高明，高明謂貴寵之人，勿枉法畏之。如是即為大中矣。 \par}

人之有能有為，使羞其行，而邦其昌。\footnote{功能有為之士，使進其所行，汝國其昌盛。○其行，如字，徐下孟反。}凡厥正人，既富方谷，\footnote{凡其正直之人，既當以爵祿富之，又當以善道接之。}汝弗能使有好於而家，時人斯其辜。\footnote{不能使正直之人有好於國家，則是人斯其詐取罪而去。}


{\noindent\zhuan\zihao{6}\fzbyks 傳“功能”至“昌盛”。正義曰：“功能有為之士”，謂其身有才能，所為有成功，此謂已在朝廷任用者也。“使進其行”者，謂人之有善，若上知其有能有為，或以言語勞來之，或以財貨賞賜之，或更任之以大位，如是則其人喜於見知，必當行自進益，人皆漸自修進，汝國其昌盛矣。 \par}

{\noindent\zhuan\zihao{6}\fzbyks 傳“凡其”至“接之”。正義曰：“凡其正直之人”，普謂臣民有正直者。爵祿所設，正直是與。已知彼人正直，必當授之以官。“既當與爵祿富之，又當以善道接之”，言其非徒與官而已，又當數加燕賜，使得其歡心也。 \par}

{\noindent\zhuan\zihao{6}\fzbyks 傳“不能”至“而去”。正義曰：授之以官爵，加之以燕賜，喜於知己，荷君恩德，必進謀樹功，有好善於國家。若雖用為官，心不委任，禮意疏薄，更無恩紀，言不聽，計不用，必將奮衣而去，不肯久留，故言“不能使正直之人有好於國家，則是人斯其詐取罪而去”也。 \par}

於其無好德，汝雖錫之福，其作汝用咎。\footnote{於其無好德之人，汝雖與之爵祿，其為汝用惡道以敗汝善。○其為，於偽反。}

{\noindent\zhuan\zihao{6}\fzbyks 傳“於其”至“汝善”。正義曰：“無好”對“有好”，“有好”謂有善也。“無好德之人”，謂彼性不好德、好惡之人也。\CJKunderwave{論語}曰：“未見好德如好色者。”傳記言好德者多矣，故傳以“好德”言之。定本作“無惡”者，疑誤耳。不好德者性行本惡，君雖與之爵祿，不能感恩行義，其為汝臣，必用惡道以敗汝善也。\CJKunderwave{易·繫辭}云：“無咎者善補過也。”“咎”是過之別名，故為惡耳。 \par}

{\noindent\shu\zihao{5}\fzkt “人之”至“用咎”。正義曰：此又言用臣之法。人之在位者，有才能，有所為，當褒賞之，委任使進其行，汝國其將昌盛也。凡其正直之人，既以爵祿富之，又復以善道接之，使之荷恩盡力。汝若不能使正直之人有好善於汝國家，是人於此其將詐取罪而去矣。於其無好德之人，謂性行惡者，汝雖與之福，賜之爵祿，但本性既惡，必為惡行，其為汝臣，必用惡道以敗汝善。言當任善而去惡。 \par}

無偏無陂，遵王之義。\footnote{偏，不平。陂，不正。言當循先王之正義以治民。○陂音秘，舊本作頗，音普多反。}無有作好，遵王之道。無有作惡，遵王之路。\footnote{言無有亂為私好惡，動必循先王之道路。○好，呼報反。惡，烏路反，注同。}無偏無黨，王道蕩蕩。\footnote{言開闢。○闢，婢必反。}無黨無偏，王道平平。\footnote{言辯治。○平平，婢綿反。治，直吏反。}無反無側，王道正直。\footnote{言所行無反道不正，則王道平直。}會其有極，歸其有極。\footnote{言會其有中而行之，則天下皆歸其有中矣。}


{\noindent\zhuan\zihao{6}\fzbyks 傳“偏下”至“治民”。正義曰：“不平”謂高下，“不正”謂邪僻，與下“好”、“惡”、“反”、“側”其義一也。偏頗阿黨是政之大患，故\CJKunderline{箕子}殷勤言耳。下傳雲“無有亂為私好私惡”者，人有私好惡則亂於正道，故傳以“亂”言之。 \par}

{\noindent\zhuan\zihao{6}\fzbyks 傳“言會”至“中矣”。正義曰：“會”謂集會，言人之將為行也，集會其有中之道而行之，行實得中，則天下皆歸其為有中矣。“天下”者，大言之。\CJKunderwave{論語}云：“一曰克己復禮，天下歸仁焉。”此意與彼同也。 \par}

{\noindent\shu\zihao{5}\fzkt “無偏”至“有極”。正義曰：更言大中之體。為人君者當無偏私，無陂曲，動循先王之正義。無有亂為私好，謬賞惡人，動循先王之正道。無有亂為私惡，濫罰善人，動循先王之正路。無偏私，無阿黨，王家所行之道蕩蕩然開闢矣。無阿黨，無偏私，王者所立之道平平然辯治矣。所行無反道，無偏側，王家之道正直矣。所行得無偏私皆正直者，會集其有中之道而行之。若其行必得中，則天下歸其中矣。言人皆謂此人為大中之人也。 \par}

曰皇極之敷言，是彝是訓，於帝其訓。\footnote{曰者，大其義，言以大中之道布陳言教，不失其常，則人皆是順矣。天且其順，而況於人乎？}凡厥庶民，極之敷言，是訓是行，以近天子之光。\footnote{凡其眾民中心之所陳言，凡順是行之，則可以近益天子之光明。○近，附近之近。}曰天子作民父母，以為天下王。\footnote{言天子佈德惠之教，為兆民之父母，是為天下所歸往，不可不務。}

{\noindent\shu\zihao{5}\fzkt “曰皇”至“下王”。正義曰：既言有中矣，為天下所歸，更美之曰，以大中之道布陳言教，不使失是常道，則民皆於是順矣。天且其順，而況於人乎？以此之故，大中為天下所歸也。又大中之道至矣，何但出於天子為貴？凡其眾民中和之心，所陳之言，謂以善言聞於上者，於是順之，於是行之，悅於民而便於政，則可近益天子之光明矣。又本人君須大中者，更美大之曰，人君於天所子，佈德惠之教，為民之父母，以是之故，為天下所歸往，由大中之道教使然。言人君不可不務大中矣。 \par}

“六，三德。一曰正直，\footnote{能正人之曲直。}二曰剛克，\footnote{剛能立事。○克，馬云：“勝也。”}三曰柔克。\footnote{和柔能治，三者皆德。}平康正直，\footnote{世平安，用正直治之。}彊弗友剛克，\footnote{友，順也。世強御不順，以剛能治之。○御，魚呂反。能治，直吏反。}燮友柔克。\footnote{燮，和也。世和順，以柔能治之。○燮，息協反。}沈潛剛克，\footnote{沈潛謂地，雖柔亦有剛，能出金石。}高明柔克。\footnote{高明謂天，言天為剛德，亦有柔克，不幹四時，喻臣當執剛以正君，君亦當執柔以納臣。}


{\noindent\zhuan\zihao{6}\fzbyks 傳“和柔”至“皆德”。正義曰：剛不恆用，有時施之，故傳言“立事”。柔則常用以治，故傳言“能治”。三德為此次者，正直在剛柔之間，故先言。二者先剛後柔，得其敘矣。王肅意與孔同。\CJKunderline{鄭玄}以為“三德，人各有一德，謂人臣也”。 \par}

{\noindent\zhuan\zihao{6}\fzbyks 傳“友順”至“治之”。正義曰：\CJKunderwave{釋訓}云：“善兄弟為友。”“友”是和順之名，故為順也。傳雲“燮，和也”，\CJKunderwave{釋詁}文。詁此三德是王者一人之德，視世而為之，故傳三者各言“世”。世平安，雖時無逆亂，而民俗未和，其下猶有曲者，須在上以正之，故世平安用正直之德治之。世有強御不順，非剛無以制之，故以剛能治之。世既和順，風俗又安，故以柔能治之。\CJKunderline{鄭玄}以為人臣各有一德，天子擇使之，注云：“安平之國，使中平守一之人治之，使不失舊職而已。國有不順孝敬之行者，則使剛能之人誅治之。其有中和之行者，則使柔能之人治之，差正之。”與孔不同。 \par}

{\noindent\zhuan\zihao{6}\fzbyks 傳“高明”至“納臣”。正義曰：\CJKunderwave{中庸}云：“博厚配地，高明配天。”高而明者惟有天耳，知“高明謂天”也。以此“高明”是天，故上傳“沈潛謂地”也。文五年\CJKunderwave{左傳}云：“天為剛德，猶不幹時。”是言天亦有柔德，不幹四時之序也。地柔而能剛，天剛而能柔，故以“喻臣當執剛以正君，君當執柔以納臣”也。 \par}

惟闢作福，惟闢作威，惟闢玉食。\footnote{言惟君得專威福，為美食。○闢，徐補亦反。玉食，張晏注\CJKunderwave{漢書}云：“玉食，珍食也”。韋昭云：“諸侯備珍異之食。”}臣無有作福作威玉食。臣之有作福作威玉食，其害於而家，兇於而國。人用側頗僻，民用僣忒。\footnote{在位不敦平，則下民僣差。○頗，普多反。僻，匹亦反。僣，子念反。忒，他得反，馬云：“惡也。”}

{\noindent\zhuan\zihao{6}\fzbyks 傳“言惟”至“美食”。正義曰：於“三德”之下說此事者，以德則隨時而用，位則不可假人，故言尊卑之分，君臣之紀，不可使臣專威福，奪君權也。衣亦不得僣君而獨言食者，人之所資，食最為重，故舉言重也。王肅云：“闢，君也。不言王者，關諸侯也，諸侯於國得專賞罰。”其義或當然也。 \par}

{\noindent\zhuan\zihao{6}\fzbyks 傳“在位”至“僣差”。正義曰：此經“福”、“威”與“食”於君每事言“闢”，於臣則並文而略之也。“作福作威”謂秉國之權，勇略震主者也。“人用側頗僻”者，謂在位小臣見彼大臣威福由己，由此之故，小臣皆附下罔上，為此側頗僻也。下民見此在位小臣秉心僻側，用此之故，下民皆不信恆，為此僣差也。言在位由大臣，下民由在位，故皆言“用”也。傳不解“家”，王肅云：“大夫稱家，言秉權之臣必滅家，復害其國也。” \par}

{\noindent\shu\zihao{5}\fzkt “六三得”至“僣忒”。正義曰：此三德者，人君之德，張弛有三也。一曰正直，言能正人之曲使直。二曰剛克，言剛強而能立事。三曰柔克，言和柔而能治。既言人主有三德，又說隨時而用之。平安之世，用正直治之。強御不順之世，用剛能治之。和順之世，用柔能治之。既言三德張弛，隨時而用，又舉天地之德，以喻君臣之交。地之德沉深而柔弱矣，而有剛，能出金石之物也。天之德高明剛強矣，而有柔，能順陰陽之氣也。以喻臣道雖柔，當執剛以正君；君道雖剛，當執柔以納臣也。既言君臣之交，剛柔遞用，更言君臣之分，貴賤有恆。惟君作福，得專賞人也。惟君作威，得專罰人也。惟君玉食，得備珍食也。為臣無得有作福作威玉食，言政當一統，權不可分也。臣之有作福作威玉食者，其必害於汝臣之家，兇於汝君之國，言將得罪喪家，且亂邦也。在位之人，用此大臣專權之故，其行側頗僻。下民用在位頗僻之故，皆言不信，而行差錯。 \par}

“七，稽疑。擇建立卜筮人，\footnote{龜曰卜，蓍曰筮。考正疑事，當選擇知卜筮人而建立之。○蓍音屍。}乃命卜筮。\footnote{建立其人，命以其職。}曰雨，曰霽，\footnote{龜兆形有似雨者，有似雨止者。○霽，子細反。}曰蒙，\footnote{蒙，陰暗。○蒙，武工反，徐亡鉤反。}曰驛，\footnote{氣洛驛不連屬。○驛音亦，注同。屬音燭。}曰克，\footnote{兆相交錯。五者卜兆之常法。}曰貞，曰悔，\footnote{內卦曰貞，外卦曰悔。}凡七。\footnote{卜筮之數。}


{\noindent\zhuan\zihao{6}\fzbyks 傳“龜曰”至“立之”。正義曰：“龜曰卜，蓍曰筮”，\CJKunderwave{曲禮}文也。考正疑事，當選擇知卜筮人而建立之。“建”亦“立”也，復言之耳。鄭、王皆以“建”、“立”為二，言將考疑事，選擇可立者，立為卜人筮人。 \par}

{\noindent\zhuan\zihao{6}\fzbyks 傳“兆相”至“常法”。正義曰：此上五者，灼龜為兆，其璺拆形狀有五種，是“卜兆之常法”也。\CJKunderwave{說文}云：“霽，雨止也。”“霽”似雨止，則“雨”似雨下。\CJKunderline{鄭玄}曰：“霽如雨止者，雲在上也。”“雺”聲近蒙，\CJKunderwave{詩}雲“零雨其蒙”，則蒙是暗之義，故以雺為兆，蒙是陰暗也。“圛”即驛也，故以為兆。“氣落驛不連屬”，“落驛”，希稀之意也。“雨”、“霽”既相對，則“蒙”、“驛”亦相對，故“驛”為落驛氣不連屬，則“雺”為氣連蒙暗也。王肅云：“圛，霍驛消減如雲陰。雺，天氣下地不應，闇冥也。”其意如孔言。\CJKunderline{鄭玄}以“圛”為明，言色澤光明也。“雺”者氣澤鬱郁冥冥也。自以“明”、“暗”相對，異於孔也。“克”謂兆相交錯。王肅云：“兆相侵入，蓋兆為二拆，其拆相交也。”\CJKunderline{鄭玄}云：“克者如雨氣色相侵入。”卜筮之事，體用難明，故先儒各以意說，未知孰得其本。今之用龜，其兆橫者為土，立者為木，斜向徑者為金，背徑者為火，因兆而細曲者曲為水，不知與此五者同異如何。此五兆不言“一曰”、“二曰”者，灼龜所遇，無先後也。 \par}

{\noindent\zhuan\zihao{6}\fzbyks 傳“內卦”至“曰悔”。正義曰：僖十五年\CJKunderwave{左傳}云，秦伯伐晉，卜徒父筮之。其卦遇蠱，蠱卦巽下艮上，說卦云，巽為風，艮為山。其佔云：“蠱之貞，風也；其悔，山也。”是內卦為貞，外卦為悔也。筮法爻從下起，故以下體為內，上體為外。下體為本，因而重之，故以下卦為貞。貞，正也，言下體是其正。\CJKunderline{鄭玄}云：“悔之言晦，晦猶終也。”晦是月之終，故以為終，言上體是其終也。下體言正，以見上體不正；上體言終，以見下體為始；二名互相明也。 \par}

卜五，佔用二，衍忒。立時人作卜筮，三人佔，則從二人之言。\footnote{立是知卜筮人，使為卜筮之事。夏殷周卜筮各異，三法並卜。從二人之言，善鈞從眾。卜筮各三人。○佔用二，馬云：“佔，筮也。”衍，以淺反。}汝則有大疑，謀及乃心，謀及卿士，謀及庶人，謀及卜筮。\footnote{將舉事而汝則有大疑，先盡汝心以謀慮之，次及卿士眾民，然後卜筮以決之。}


{\noindent\zhuan\zihao{6}\fzbyks 傳“立是”至“三人”。正義曰：此經“卜五，佔用二，衍忒”，孔不為傳。\CJKunderline{鄭玄}云：“‘卜五佔用’謂雨、霽、蒙、驛、克也，‘二衍忒’謂貞、悔也。”斷“用”從上句，“二衍忒”者，指謂筮事。王肅云：“‘卜五’者，筮短龜長，故卜多而筮少。‘佔用二’者，以貞、悔佔六爻。‘衍忒’者，當推衍其爻義以極其意。”“卜五，佔二”，其義當如王解，其“衍忒”宜總謂卜筮，皆當衍其義，極其變，非獨筮衍而卜否也。傳言“立是知卜筮人，使為卜筮之事”者，言經之此文覆述上句“立卜筮人”也。言“三人佔”，是佔此卜筮，法當有三人。\CJKunderwave{周禮}：“太卜掌三兆之法，一曰玉兆，二曰瓦兆，三曰原兆。掌三易之法，一曰\CJKunderwave{連山}，二曰\CJKunderwave{歸藏}，三曰\CJKunderwave{周易}。”杜子春以為“玉兆，帝顓頊之兆。瓦兆，\CJKunderline{帝堯}之兆”。又云“\CJKunderwave{連山}，虙犧。\CJKunderwave{歸藏}，黃帝。三兆三易皆非夏殷”。而孔意必以三代夏殷周法者，以\CJKunderwave{周禮}指言“一曰”、“二曰”，不辯時代之名。案\CJKunderwave{考工記}云，夏曰世室，殷曰重屋，周曰明堂。又\CJKunderwave{禮記·郊特牲}雲“夏收，殷冔，周冕”。皆以夏殷周三代相因，明三易亦夏殷周相因之法。子春之言，孔所不取。\CJKunderline{鄭玄}\CJKunderwave{易贊}亦云：“夏曰\CJKunderwave{連山}，殷曰\CJKunderwave{歸藏}。”與孔同也。所言三兆三易，必是三代異法，故傳以為夏殷周卜筮各以三代異法，三法並卜，法有一人，故三人也。“從二人之言”者，二人為善既鈞，故從眾也。若三人之內賢智不等，雖少從賢，不從眾也。“善鈞從眾”，成六年\CJKunderwave{左傳}文。既言“三法並卜”，嫌筮不然，故又云“卜筮各三人”也。經惟言三佔從二，何知不一法而三佔，而知“三法並用”者？\CJKunderwave{金縢}云：“乃卜三龜，一習吉。”\CJKunderwave{儀禮}士喪卜葬，佔者三人，貴賤俱用三龜，知卜筮並用三代法也。 \par}

{\noindent\zhuan\zihao{6}\fzbyks 傳“將舉”至“決之”。正義曰：非有所舉，則自不卜，故云“將舉事”，事有疑，則當卜筮。人君先盡己心以謀慮之，次及卿士眾民，人謀猶不能定，然後問卜筮以決之。故先言“乃心”，後言“卜筮”也。\CJKunderline{鄭玄}云：“卿士，六卿掌事者。”然則“謀及卿士”，以卿為首耳，其大夫及士亦在焉。以下惟言“庶人”，明大夫及士寄“卿”文以見之矣。\CJKunderwave{周禮}：“小司寇掌外朝之政，以致萬民而詢焉。一曰詢國危，二曰詢國遷，三曰詢立君。”是有大疑而詢眾也。又曰“小司寇以敘進而問焉”，是謀及之也。大疑者不要是彼三詢，其謀及則同也。謀及庶人，必是大事，若小事不必詢於萬民，或謀及庶人在官者耳。\CJKunderwave{小司寇}又曰：“以三剌斷庶民獄訟之中，一曰訊群臣，二曰訊群吏，三曰訊萬民。”彼“群臣”、“群吏”分而為二，此惟言“卿士”者，彼將斷獄，令眾議然後行刑，故臣與民為三，其人主待眾議而決之；此則人主自疑，故一人主為一，又總群臣為一也。 \par}

汝則從，龜從，筮從，卿士從，庶民從，是之謂大同。\footnote{人心和順，龜筮從之，是謂大同于吉。}身其康強，子孫其逢吉。\footnote{動不違眾，故後世遇吉。○逢，馬云：“逢，大也。”}汝則從，龜從，筮從，卿士逆，庶民逆，吉。\footnote{三從二逆，中吉，亦可舉事。}卿士從，龜從，筮從，汝則逆，庶民逆，吉。\footnote{君臣不同，決之卜筮，亦中吉。}


{\noindent\zhuan\zihao{6}\fzbyks 傳“人心”至“于吉”。正義曰：人主與卿士庶民皆從，是“人心和順”也。此必臣民皆從，乃問卜筮，而進龜筮於上者，尊神物，故先言之。不在“汝則”之上者，卜當有主，故以人為先，下三事亦然。改“卜”言“龜”者，“卜”是請問之意，吉凶龜佔兆告以人，故改言“龜”也。“筮”則本是蓍名，故不須改也。 \par}

{\noindent\zhuan\zihao{6}\fzbyks 傳“動不”至“遇吉”。正義曰：物貴和同，故大同之吉，延及於後。宣三年\CJKunderwave{左傳}稱“\CJKunderline{成王}定鼎，卜世三十，卜年七百”，是“後世遇吉”。 \par}

{\noindent\zhuan\zihao{6}\fzbyks 傳“三從”至“舉事”。正義曰：此與下二事，皆是三從二逆，除龜、筮以外，有汝與卿士、庶民，分三者各為一從二逆，嫌其貴賤有異，從逆或殊，故三者各以有一從為主，見其為吉同也。方論得吉，以從者為主，故次言“卿士從”，下言“庶民從”也。以從為主，故退“汝則”於下。傳解其意，卿士從者，“君臣不同”也；庶民從者，“民與上異心”也。解臣民與君異心，得其筮之意也。不言四從一逆者，吉可知，不假言之也。四從之內，雖龜筮相違，亦為吉，以其從者多也。若三從之內，龜筮相違，雖不如龜筮俱從，猶勝下龜筮相違，二從三逆。必知然者，以下傳雲“二從三逆，龜筮相違”，既計從之多少，明從多則吉。故杜預云：“龜筮同卿士之數者，是龜筮雖靈，不至越於人也。”上言“庶人”，又言“庶民”者，嫌“庶人”惟指在官者，變“人”言“民”見其同也。民人之賤，得與卿士敵者，貴者雖貴，未必謀慮長，故通以民為一，令與君臣等也。 \par}

庶民從，龜從，筮從，汝則逆，卿士逆，吉。\footnote{民與上異心，亦卜筮以決之。}汝則從，龜從，筮逆，卿士逆，庶民逆，作內吉，作外兇。\footnote{二從三逆，龜筮相違，故可以祭祀冠婚，不可以出師征伐。○冠，官喚反。}龜筮共違於人，\footnote{皆逆。}用靜吉，用作兇。\footnote{安以守常則吉，動則兇。}

{\noindent\zhuan\zihao{6}\fzbyks 傳“民與”至“決之”。正義曰：天子聖人，庶民愚賤，得為識見同者，但聖人生知，不假卜筮，垂教作訓，晦跡同凡，且庶民既眾，以眾情可否，亦得上敵於聖。故\CJKunderwave{老子}雲“聖人無常心，以百姓心為心”是也。 \par}

{\noindent\zhuan\zihao{6}\fzbyks 傳“二從”至“征伐”。正義曰：此“二從三逆”為小吉，故猶可舉事。“內”謂國內，“故可以祭祀冠婚”。“外”謂境外，“故不可以出師征伐”。征伐事大，此非大吉故也。此經“龜從，筮逆”，其筮從、龜逆為吉亦同。故傳言“龜筮相違”，見龜筮之智等也。若龜筮智等，而僖四年\CJKunderwave{左傳}雲“筮短龜長”者，於時晉獻公欲以驪姬為夫人，卜既不吉，而更令筮之，神靈不以實告，筮而得吉，必欲用之，卜人欲令公舍筮從卜，故曰“筮短龜長”，非是龜實長也。\CJKunderwave{易·繫辭}云：“蓍之德圓而神，卦之德方以智。”神以知來，智以藏往，然則知來藏往，是為極妙，雖龜之長，無以加此。聖人演筮為易，所知豈是短乎？明彼長短之說，乃是有為言耳。“此二從三逆”，以汝與龜為二從耳。卿士庶民課有一從，亦是二從，兇吉亦同，故不復設文，同可知也。若然，汝、卿士、庶民皆逆，龜、筮並從，則亦是二從三逆，而經無文者，若君與臣民皆逆，本自不問卜矣，何有龜從筮從之理也？前三從之內，龜筮既從，君與卿士庶民各有一從以配龜筮，凡有三條。若惟君與卿士從配龜為一條，或君與庶民從配龜又為一條，或卿士庶民從配龜又為一條，凡有三條。若筮從龜逆，其事亦然。二從三逆，君配龜從為一條，於經已具。卿士配龜從為二條，庶民配龜從為三條。若筮從龜逆，以人配筮，其事亦同。案\CJKunderwave{周禮·筮人}：“國之大事，先筮而後卜。”\CJKunderline{鄭玄}云：“於筮之兇則止。”何有筮逆龜從及龜筮俱違者？崔靈恩以為筮用三代之佔，若三佔之俱主兇，則止不卜，即鄭注\CJKunderwave{周禮}“筮兇則止”是也。若三佔二逆一從，兇猶不決，雖有筮逆，猶得更卜，故此有筮逆龜從之事。或筮兇則止而不卜，乃是\CJKunderline{鄭玄}之意，非是\CJKunderwave{周禮}經文，未必孔之所取。\CJKunderwave{曲禮}云：“卜筮不相襲。”鄭云：“卜不吉則又筮，筮不吉則又卜，是謂瀆龜筮。”\CJKunderwave{周禮}太卜小事筮，大事卜，應筮而又用卜，應卜而又用筮，及國之大事先筮後卜，不吉之後更作卜筮，如此之等，是為相襲，皆據吉凶分明，不可重為卜筮，若吉凶未決，於事尚疑者，則得更為卜筮。僖二十五年晉侯卜納王，得阪泉之兆，曰：“吾不堪也。”公曰：“筮之。”遇大有之睽；又哀九年晉趙鞅卜救鄭，遇水適火，又筮之，遇泰之需之類是也。\CJKunderwave{周禮}既先筮後卜，而春秋時先卜後筮者，不能依禮故也。 \par}

{\noindent\shu\zihao{5}\fzkt “七稽”至“之言”。正義曰：“稽疑”者，言王者考正疑事。當選擇知卜筮者而建立之，以為卜筮人，謂立為卜人筮人之官也。既立其官，乃命以卜筮之職。雲卜兆有五。曰雨兆，如雨下也。曰霽兆，如雨止也。曰雺兆，氣蒙暗也。曰圛兆，氣落驛不連屬也。曰克兆，相交也。筮卦有二重，二體乃成一卦。曰貞，謂內卦也。曰悔，謂外卦也。卜筮兆卦其法有七事，其卜兆用五，雨、霽、蒙、驛、克也。其筮佔用二，貞與悔也。卜筮皆就此七者推衍其變，立是知卜筮人，使作卜筮之官。其卜筮必用三代之法，三人佔之，若其所佔不同，而其善鈞者，則從二人之言，言以此法考正疑事也。 \par}

“八，庶徵。曰雨，曰暘，曰燠，曰寒，曰風。\footnote{雨以潤物，暘以幹物，暖以長物，寒以成物，風以動物，五者各以其時，所以為眾驗。○暘音陽。幹音幹。暖，乃管反。長，之丈反。}曰時五者來備，各以其敘，庶草蕃廡。\footnote{言五者備至，各以次序，則眾草蕃滋廡豐也。○蕃音煩。廡，無甫反，徐莫柱反。}一極備，兇。一極無，兇。\footnote{一者備極，過甚則兇。一者極無，不至亦兇。謂不時失敘。}

{\noindent\shu\zihao{5}\fzkt 正義曰：“庶”，眾也。“徵”，驗也。王者用九疇，為大中，行“稽疑”以上為善政，則眾驗有美惡，以為人主。自“曰雨”至“一極無兇”，總言五氣之驗，有美有惡。“曰休徵”敘美行之驗，“曰咎徵”敘惡行之驗。自“曰王省”至“家用平康”，言政善致美也。“日月歲時”至“家用不寧”，言政惡致咎也。“庶民惟星”以下，言人君當以常度齊正下民。 \par}

{\noindent\zhuan\zihao{6}\fzbyks 傳“雨以”至“眾驗”。正義曰：\CJKunderwave{易·說卦}云：“風以散之，雨以潤之，日以烜之。”日，暘也；烜，幹也；是“雨以潤物，暘以幹物，風以動物”也。\CJKunderwave{易·繫辭}云：“寒往則暑來，暑往則寒來，寒暑相推而歲成焉。”是言天氣有寒有暑，暑長物而寒成物也。\CJKunderwave{釋言}云：“燠，暖也。”舍人曰：“燠，溫暖也。”是“燠”、“暖”為一，故傳以“暖”言之。不言“暑”而言“燠”者，“燠”是熱之始，“暑”是熱之極，“涼”是冷之始，“寒”是冷之極，長物舉其始，成物舉其極，理宜然也。五者各以其時而至，所以為眾事之驗也。所以言“時”者，謂當至則來，當止則去，無常時也。冬寒夏燠，雖有定時，或夏須漸寒，冬當漸熱，雨足則思暘，暘久則思雨，草木春則待風而長，秋則待風而落，皆是無定時也。不言“一曰”、“二曰”者，為其來無先後也。依五事所致為次，下雲“休徵”、“咎徵”、“雨若”、“風若”，是其致之次也。昭元年\CJKunderwave{左傳}云：“天有六氣，陰、陽、風、雨、晦、明也。”以彼六氣校此五氣，“雨”、“暘”、“風”文與彼同，彼言“晦”、“明”，此言“寒”、“燠”則“晦”是“寒”也，“明”是“燠”也，惟彼“陰”於此無所當耳。\CJKunderwave{五行傳}說五事致此五氣云：“貌之不恭，是謂不肅，厥罰恆雨，惟金沴木。言之不從，是謂不又，厥罰恆暘，惟木沴金。視之不明，是謂不晢，厥罰恆燠，惟水沴火。聽之不聰，是謂不謀，厥罰恆寒，惟火沴水。思之不睿，是謂不聖，厥罰恆風，惟木金水火沴土。”如彼\CJKunderwave{五行傳}言，是雨屬木，暘屬金，燠屬火，寒屬水，風屬土。鄭云：“雨，木氣也，春始施生，故木氣為雨。暘，金氣也，秋物成而堅，故金氣為暘。燠，火氣也。寒，水氣也。風，土氣也。凡氣非風不行，猶金木水火非土不處，故土氣為風。”是用\CJKunderwave{五行傳}為說，孔意亦當然也。六氣有“陰”，五事休咎皆不致陰，\CJKunderwave{五行傳}又曰：“皇之不極，厥罰常陰”，是陰氣不由五事，別自屬皇極也。蓋立用大中，則陰順時為休。大之不中，陰恆若為咎也。 \par}

{\noindent\zhuan\zihao{6}\fzbyks 傳“言五”至“廡豐”。正義曰：五氣所以生成萬物，正可時來時去，不可常無常有，故言“五者備至，各以次序”。須至則來，須止則去，則眾草百物蕃滋廡豐也。\CJKunderwave{釋詁}云：“廡，豐茂也。”“草蕃廡”言草滋多而茂盛也。下言“百穀用成”，此言“眾草蕃廡”者，舉草茂盛則谷成必矣，舉輕以明重也。 \par}

{\noindent\zhuan\zihao{6}\fzbyks 傳“一者”至“失敘”。正義曰：此謂不以時來，其至無次序也。“一者備極，過甚則兇”，謂來而不去也。“一者極無，不至亦兇”，謂去而不來也。即卜雲“恆雨若”、“恆風若”之類是也。有無相刑，去來正反，恆雨則無暘，恆寒則無燠，恆雨亦兇，無暘亦兇，恆寒亦兇，無燠亦兇，謂至不待時，失次序也。如此則草不茂，谷不成也。 \par}

{\noindent\shu\zihao{5}\fzkt “曰雨”至“無兇”。正義曰：將說其驗，先立其名。五者行於天地之間，人物所以得生成也。其名曰雨，所以潤萬物也。曰暘，所以幹萬物也。曰燠，所以長萬物也。曰寒，所以成萬物也。曰風，所以動萬物也。此是五氣之名。“曰時”言五者各以時來，所以為眾事之驗也。更述時與不時之事，五者於是來皆備足，須風則風來，須雨則雨來，其來各以次序，則眾草木蕃滋而豐茂矣。謂來以時也。若不以時，五者之內，一者備極，過甚則兇。一者極無，不至亦兇。雨多則澇，雨少則旱，是備極亦兇，極無亦兇。其餘四者亦然。 \par}

曰休徵：\footnote{敘美行之驗。○行，下孟反。}曰肅，時寒若。\footnote{君行敬，則時雨順之。}曰乂，時暘若。\footnote{君行政治，則時暘順之。○治，直吏反，下“政治”、“治其職”皆同。}曰晢，時燠若。\footnote{君能照晢，則時燠順之。○晢，之設反，徐音制，又音哲。}曰謀，時寒若。\footnote{君能謀，則時寒順之。}曰聖，時風若。\footnote{君能通理，則時風順之。}

{\noindent\shu\zihao{5}\fzkt “曰休徵”至“風若”。正義曰：既言五者次序，覆述次序之事，曰美行致以時之驗，何者是也？曰人君行敬，則雨以時而順之。曰人君政治，則暘以時而順之。曰人君照晢，則燠以時而順之。曰人君謀當，則寒以時而順之。曰人通聖，則風以時而順之。此則致上文“各以其次，敘庶草蕃廡”也。 \par}

曰咎徵：\footnote{敘惡行之驗。○咎，其九反。}曰狂，恆雨若。\footnote{君行狂妄，則常雨順之。}曰僣，恆暘若。\footnote{君行僣差，則常暘順之。}曰豫，恆燠若。\footnote{君行逸豫，則常燠順之。○豫，羊庶反，徐又音舒。}曰急，恆寒若。\footnote{君行急，則常寒順之。}曰蒙，恆風若。\footnote{君行蒙暗，則常風順之。}

{\noindent\zhuan\zihao{6}\fzbyks 傳“君行”至“順之”。正義曰：此“休”、“咎”皆言“若”者，其所致者皆順其所行，故言“若”也。\CJKunderwave{易·文言}云：“雲從龍，風從虎，水流溼，火就燥。”是物各以類相應，故知天氣順人所行以示其驗也。其咎反於休者，人君行不敬則狂妄，故“狂”對“肅”也。政不治則僣差，故“僣”對“乂”也。明不照物則行自逸豫，故“豫”對“晢”也。心無謀慮則行必急躁，故“急”對“謀”也。性不通曉則行必蒙暗，故“蒙”對“聖”也。\CJKunderline{鄭玄}以“狂”為倨慢，以對“不敬”，故為慢也。鄭、王本“豫”作“舒”，鄭雲“舉遲也”，王肅雲“舒，惰也”，以對“昭晢”，故為遲惰。鄭云：“急促，自用也。”以“謀”者用人之言，故“急”為自用己也。鄭云：“蒙，見冒亂也。”王肅云：“蒙，瞽蒙。”以“聖”是通達，故“蒙”為瞽蒙。所見冒亂，言其不曉事，與“聖”反也。與孔各小異耳。 \par}

{\noindent\shu\zihao{5}\fzkt “曰咎之”至“風若”。正義曰：上既言失次序，覆述失次序之事。曰惡行致備極之驗，何者是也？曰君行狂妄，則常雨順之。曰君行僣差，則常暘順之。曰君行逸豫，則常暖順之。曰君行急躁，則常寒順之。曰君行蒙暗，則常風順之。此即致上文“一極備兇，一極無兇”也。 \par}

曰王省惟歲，\footnote{王所省職，兼所總群吏，如歲兼四時。○省，息井反。}卿士惟月，\footnote{卿士各有所掌，如月之有別。○別，彼列反。}師尹惟曰。\footnote{眾正官之吏，分治其職，如日之有歲月。}歲月曰時無易，\footnote{各順常。}百穀用成，乂用明，\footnote{歲月日時無易，則百穀成。君臣無易，則政治明。}俊民用章，家用平康。\footnote{賢臣顯用，國家平寧。}日月歲時既易，\footnote{是三者已易，喻君臣易職。}百穀用不成，乂用昏不明，俊民用微，家用不寧。\footnote{君失其柄，權臣擅命，治暗賢隱，國家亂。}

{\noindent\zhuan\zihao{6}\fzbyks 傳“王所”至“四時”。正義曰：下雲“庶民惟星”以“星”喻民，知此“歲月日”者，皆以喻職事也。於王言“省”，則卿士師尹亦為“省”也。王之所省，職無不兼，所總群吏如歲兼四時。下句惟有“月”、“日”，群臣無喻“時”者，但時以統月，故傳以“四時”言之，言其兼下月日也。 \par}

{\noindent\zhuan\zihao{6}\fzbyks 傳“眾正”至“歲月”。正義曰：師，眾也。尹，正也。“眾正官之吏”謂卿士之下。有正官大夫，與其同類之官為長。\CJKunderwave{周禮}大司樂為樂官之長，太卜為卜官之長，此之類也。此等分治其職，屬王屬卿，“如日之有歲月”，言其有系屬也。\CJKunderwave{詩}稱“赫赫師尹”，乃謂三公之官，此以“師尹”為正官之吏，謂大夫者，以此“師尹”之文在“卿士”之下，卑於卿士，知是大夫。與小官為長，亦是眾官之長，故師尹之名同耳。鄭雲“所以承休徵、咎徵言之者，休咎五事得失之應，其所致尚微，故大陳君臣之象，成皇極之事。其道得，則其美應如此。其道失，則敗德如彼。非徒風雨寒燠而已”是也。 \par}

{\noindent\shu\zihao{5}\fzkt “曰王省”至“不寧”。正義曰：既陳五事之休咎，又言皇極之得失，與上異端，更復言曰，王之省職，兼總群吏，惟如歲也。卿士分居列位，惟如月也。眾正官之長各治其職，惟如日也。此王也，卿士也，師尹也，掌事猶歲月日者，言皆無改易，君秉君道，臣行臣事。則百穀用此而成，歲豐稔也。其治用是而明，世安泰也。俊民用此而章，在官位也。國家用此而平安，風俗和也。若王也，卿士也，師尹也，掌事猶如日月歲者，是已變易，君失其柄權，臣各專恣。百穀用此而不成，歲饑饉也。其治用此昏暗而不明，政事亂也。俊民用此而卑微，皆隱遁也。國家用此而不安泰，時世亂也。此是皇極所致，得中則致善，不中則致惡。歲月日無易，是得中也。既易，是不中也。所致善惡乃大於庶徵，故於此敘之也。 \par}

庶民惟星，星有好風，星有好雨。\footnote{星，民象，故眾民惟若星。箕星好風，畢星好雨，亦民所好。○好，呼報反。}日月之行，則有冬有夏。\footnote{日月之行，冬夏各有常度。君臣政治，小大各有常法。}月之從星，則以風雨。\footnote{月經於箕則多風，離於畢則多雨。政教失常以從民欲，亦所以亂。}

{\noindent\zhuan\zihao{6}\fzbyks 傳“星民”至“所好”。正義曰：星之在天，猶民之在地，星為民象，以其象民，故因以星喻，故“眾民惟若星”也。直言“星有好風”，不知何星，故云“箕星好風”也。畢星好雨亦如民有所好也，不言“畢星好雨”，具於下傳。 \par}

{\noindent\zhuan\zihao{6}\fzbyks 傳“日月”至“常法”。正義曰：日月之行，四時皆有常法，變冬夏為南北之極，故舉以言之。“日月之行，冬夏各有常度”，喻人君為政小大，各有常法。張衡、蔡雍、王蕃等說渾天者皆云，周天三百六十五度四分度之一，天體圓如彈丸，北高南下，北極出地上三十六度，南極入地下三十六度。南極去北極直徑一百二十二度弱，其依天體隆曲。南極去北極一百八十二度強，正當天之中央。南北二極中等之處謂之赤道，去南北極各九十一度。春分日行赤道，從此漸北夏至赤道之北二十四度，去北極六十七度，去南極一百一十五度，日行黑道。從夏至日以後日漸南，至秋分還行赤道，與春分同。冬至行赤道之南二十四度，去南極六十七度，去北極一百一十五度，其日之行處謂之黃道。又有月行之道與日道相近，交絡而過，半在日道之裡，半在日道之表，其當交則兩道相合，交去極遠處，兩道相去六度。此其日月行道之大略也。王肅云：“日月行有常度，君臣禮有常法，以齊其民。” \par}

{\noindent\zhuan\zihao{6}\fzbyks 傳“月經”至“以亂”。正義曰：\CJKunderwave{詩}云：“月離於畢，俾滂沱矣。”是離畢則多雨，其文見於經。經箕則多風，傳記無其事。\CJKunderline{鄭玄}引\CJKunderwave{春秋緯}云：“月離於箕，則風揚沙。”作緯在\CJKunderline{孔君}之後，以前必有此說，孔依用之也。月行雖有常度，時或失道從星，經箕多風，離畢多雨，此天象之自然，以箕為簸揚之器，畢亦捕魚之物故耳。鄭以為“箕星好風者，箕東方木宿，風中央土氣，木克土為妻，從妻所好，故好風也。畢星好雨者，畢西方金宿，雨東方木氣，金克木為妻，從妻所好，故好雨也。推此則南宮好暘，北宮好燠，中宮四季好寒，以各尚妻之所好故也”。未知孔意同否。顧氏所解，亦同於鄭，言“從星者，謂不應從而從，以致此風雨，故喻政教失常以從民欲，亦所以亂也”。上雲“日月之行”，此句惟言“月”者，鄭云：“不言日者，日之從星，不可見故也。” \par}

{\noindent\shu\zihao{5}\fzkt “庶民”至“風雨”。正義曰：既言大中治民，不可改易，又言民各有心，須齊正之。言庶民之性惟若星然。“星有好風，星有好雨”，以喻民有好善，亦有好惡。“日月之行，則有冬有夏”，言日月之行，冬夏各有常道，喻君臣為政小大，各有常法。若日月失其常道，則天氣從而改焉。月之行度失道，從星所好，以致風雨，喻人君政教失常，從民所欲，則致國亂。故常立用大中，以齊正之，不得從民欲也。 \par}

“九，五福：一曰壽，\footnote{百二十年。}二曰富，\footnote{財豐備。}三曰康寧，\footnote{無疾病。}四曰攸好德，\footnote{所好者德福之道。}五曰考終命。\footnote{各成其短長之命以自終，不橫夭。○橫，華孟反，又如字。}


{\noindent\zhuan\zihao{6}\fzbyks 傳“百二十年”。正義曰：人之大期，百年為限，世有長壽雲百二十年者，故傳以最長者言之，未必有正文也。 \par}

{\noindent\zhuan\zihao{6}\fzbyks 傳“所好”至“之道”。正義曰：人所嗜好，稟諸上天，性之所好，不能自已。好善者或當知善是善，好惡者不知惡之為惡，是善，故好之無厭。任其所好，從而觀之，所好者德，是福之道也。好德者天使之然，故為福也。鄭云：“民皆好有德也。”王肅云：“言人君所好者道德為福。”\CJKunderwave{洪範}以人君為主，上之所為，下必從之，人君好德，故民亦好德，事相通也。 \par}

{\noindent\zhuan\zihao{6}\fzbyks 傳“各成”至“橫夭”。正義曰：成十三年\CJKunderwave{左傳}云：“民受天地之中以生，所謂命也。能者養之以福，不能者敗以取禍。”是言命之短長雖有定分，未必能遂其性，不致夭枉，故各成其短長之命以自終，不橫夭者亦為福也。 \par}

六極：一曰兇短折，\footnote{動不遇吉，短未六十，折未三十，言辛苦。○兇，馬云：“終也。”折，時設反，又之舌反。}二曰疾，\footnote{常抱疾苦。}三曰憂，\footnote{多所憂。}四曰貧，\footnote{困於財。}五曰惡，\footnote{醜陋。}六曰弱。”\footnote{尫劣。○尫，鳥黃反。}

{\noindent\zhuan\zihao{6}\fzbyks 傳“動不”至“辛苦”。正義曰：“動不遇吉”者，解“兇”也。傳以“壽”為“百二十年”，“短”者半之，為“未六十”，“折”又半，為“未三十”。“辛苦”者，味也，辛苦之味入口，猶困厄之事在身，故謂殃厄勞役之事為辛苦也。\CJKunderline{鄭玄}以為“兇短折皆是夭枉之名。未齔曰兇，未冠曰短，未婚曰折”。\CJKunderwave{漢書·五行志}云：“傷人曰兇，禽獸曰短，草木曰折。一曰兇，夭是也，兄喪弟曰短，父喪子曰折。”並與孔不同。 \par}

{\noindent\zhuan\zihao{6}\fzbyks 傳“尫劣”。正義曰：“尫”、“劣”並是弱事，為筋力弱，亦為志氣弱。\CJKunderline{鄭玄}云：“愚懦不毅曰弱。”言其志氣弱也。\CJKunderwave{五行傳}有“致極”之文，無致福之事。\CJKunderline{鄭玄}依\CJKunderwave{書傳}云：“兇短折，思不睿之罰。疾，視不明之罰。憂，言不從之罰。貧，聽不聰之罰。惡，貌不恭之罰。弱，皇不極之罰。反此而云，王者思睿則致壽，聽聰則致富，視明則致康寧，言從則致攸好德，貌恭則致考終命。所以然者，不但行運氣性相感，以義言之，以思睿則無擁，神安而保命，故壽。若蒙則不通，殤神夭性，所以短所也。聽聰則謀當，所求而會，故致富。違而失計，故貧也。視明照了，性得而安寧。不明，以擾神而疾也。言從由於德，故好者德也。不從而無德，所以憂耳。貌恭則容儼形美而成性，以終其命。容毀，故致惡也。不能為大中，故所以弱也。”此亦孔所不同焉。此福極之文，雖主於君，亦兼於下，故有“貧”、“富”、“惡”、“弱”之等也。 \par}

{\noindent\shu\zihao{5}\fzkt “九五福”至“曰弱”。正義曰：“五福”者，謂人蒙福祐有五事也。一曰壽，年得長也。二曰富，家豐財貨也。三曰康寧，無疾病也。四曰攸好德，性所好者美德也。五曰考終命，成終長短之命，不橫夭也。“六極”謂窮極惡事有六。一曰兇短折，遇兇而橫夭性命也。二曰疾，常抱疾病。三曰憂，常多憂愁。四曰貧，困之於財。五曰惡，貌狀醜陋。六曰弱，志力尫劣也。“五福”、“六極”，天實得為之，而歷言此者，以人生於世，有此福極，為善致福，為惡致極，勸人君使行善也。“五福”、“六極”如此次者，鄭云：“此數本諸其尤者，福是人之所欲，以尤欲者為先。極是人之所惡，以尤所不欲者為先。以下緣人意輕重為次耳。” \par}

\CJKunderline{武王}既勝殷，邦諸侯，班宗彝，\footnote{賦宗廟彝器酒樽賜諸侯。○班，本又作般，音同。}作\CJKunderwave{分器}。\footnote{言諸侯尊卑，各有分也。亡。○分，扶問反，注同。}

{\noindent\zhuan\zihao{6}\fzbyks 傳“賦宗”至“諸侯”。正義曰：序雲“邦諸侯”者，立邦國，封人為諸侯也。\CJKunderwave{樂記}雲“封有功者為諸侯”，\CJKunderwave{詩·賚}序雲“大封於廟”，謂此時也。\CJKunderwave{釋言}云：“班，賦也。”\CJKunderwave{周禮}有司尊彝之官，鄭云：“彝亦尊也。鬱鬯曰彝。彝，法也，言為尊之法正。”然則盛鬯者為彝，盛酒者為尊，皆祭宗廟之酒器也。分宗廟彝器酒尊以賦諸侯，既封乃賜之也。 \par}

{\noindent\zhuan\zihao{6}\fzbyks 傳“言諸”至“也亡”。正義曰：篇名\CJKunderwave{分器}，知其篇“言諸侯尊卑，各有分也”。昭十二年\CJKunderwave{左傳}楚靈王云：“昔我先王熊繹與\CJKunderline{呂伋}、王孫牟、谿父、禽父並事康王，四國皆有分，我獨無。”十五年傳曰：“諸侯之封也，皆受明器於王室。”杜預云：“謂明德之分器也。”是諸侯各有分也，亡。 \par}

{\noindent\shu\zihao{5}\fzkt “武王”至“分器”。正義曰：武王既已勝殷，制邦國以封有功者為諸侯。既封為國君，乃班賦宗廟彝器以賜之，於時有言誥戒敕。史敘其事，作\CJKunderwave{分器}之篇。 \par}

%%% Local Variables:
%%% mode: latex
%%% TeX-engine: xetex
%%% TeX-master: "../Main"
%%% End:

%% -*- coding: utf-8 -*-
%% Time-stamp: <Chen Wang: 2024-04-02 11:42:41>

% {\noindent \zhu \zihao{5} \fzbyks } -> 注 (△ ○)
% {\noindent \shu \zihao{5} \fzkt } -> 疏

\chapter{卷十三}


\section{旅獒第七【偽】}


西旅獻獒，\footnote{西戎遠國貢大犬。○獒，五羔反。馬云：“作豪，酋豪也。”}太保作\CJKunderwave{旅獒}。\footnote{\CJKunderline{召公}陳戒。○召，時照反，後“\CJKunderline{召公}”皆仿此。}

旅獒\footnote{因獒而陳道義。}


{\noindent\zhuan\zihao{6}\fzbyks 傳“西戎”至“大犬”。正義曰：“西旅”，西方夷名。西方曰“戎”，克商之後乃來，知是“西戎遠國”也。“獒”是犬名，故云“貢大犬”。 \par}

{\noindent\zhuan\zihao{6}\fzbyks 傳“\CJKunderline{召公}陳戒”。正義曰：\CJKunderline{成王}時\CJKunderline{召公}為太保，知此時“太保”亦\CJKunderline{召公}也。\CJKunderwave{釋詁}云：“旅，陳也。”故云“\CJKunderline{召公}陳戒”。上“旅”是國名，此“旅”訓為陳，二“旅”字同而義異。鄭云：“獒讀曰豪，西戎無君名，強大有政者為酋豪。國人遣其遒豪來獻見於周。”良由不見古文，妄為此說。 \par}

{\noindent\shu\zihao{5}\fzkt “西旅”至“旅獒”。正義曰：西方之戎有國名“旅”者，遣獻其大犬，其名曰“獒”，於是太保\CJKunderline{召公}因陳戒。史敘其事，作\CJKunderwave{旅獒}。 \par}

惟克商，遂通道於九夷八蠻。\footnote{四夷慕化，貢其方賄。九、八言非一。皆通道路，無遠不服。○賄，呼罪反。}西旅厎貢厥獒，\footnote{西戎之長，致貢其獒。犬高四尺曰獒，以大為異。○厎，之履反。長，丁丈反。}太保乃作\CJKunderwave{旅獒}，用訓於王。\footnote{陳貢獒之義以訓諫王。}


{\noindent\zhuan\zihao{6}\fzbyks 傳“四夷”至“不服”。正義曰：\CJKunderwave{曲禮}云：“其在東夷、西戎、南蠻、北狄。”經舉“夷”、“蠻”則戎狄可知。“四夷慕化，貢其方賄”，言所貢非獨旅也。四夷各自為國，無大小統領，“九、八言非一也”。釋地云：“九夷、八狄、七戎、六蠻謂之四海。”又云：“八蠻在南方，六戎在西方，五狄在北方。”上下二文三方數目不同。\CJKunderwave{明堂位}稱九夷、八蠻、六戎、五狄，與\CJKunderwave{爾雅}上文不同。\CJKunderwave{周禮}職方氏掌四夷、八蠻、七閩、九貉、五戎、六狄之人。\CJKunderline{鄭玄}云：“四、八、七、九、五、六，周之所服國數也。”遍檢經傳，四夷之數，參差不同，先儒舊解，此\CJKunderwave{爾雅}殷制，\CJKunderwave{明堂位}及\CJKunderwave{職方}並\CJKunderwave{爾雅}下文雲八蠻在南，六戎在西，五狄在北，皆為周制，義或當然。\CJKunderwave{明堂位}言六戎、五狄，\CJKunderwave{職方}言五戎、六狄，趙商以此問鄭，鄭答云：“戎狄但有其國數，其名難得而知。”是鄭亦不能定解。言“克商，遂通道”，是王家遣使通道也。\CJKunderwave{魯語}引此事，韋昭云：“通道，譯使懷柔之。”是王家遣使通彼，彼聞命來獻也。言其通夷蠻而有戎貢，是四夷皆通道路，無所不服。 \par}

{\noindent\zhuan\zihao{6}\fzbyks 傳“西戎”至“為異”。正義曰：“西戎之長”謂旅國之君。“致貢其獒”，或遣使貢之，不必自來也。“犬高四尺曰獒”，\CJKunderwave{釋畜}文。\CJKunderwave{左傳}晉靈公有犬謂之獒。旅國以犬為異，故貢之也。 \par}

{\noindent\shu\zihao{5}\fzkt “惟克”至“於王”。正義曰：惟\CJKunderline{武王}既克商，華夏既定，遂開通道路於九夷八蠻，於是有西戎旅國致貢其大犬名獒。太保\CJKunderline{召公}乃作此篇，陳貢獒之義，用訓諫於王。 \par}

曰：“嗚呼！明王慎德，四夷咸賓。\footnote{言明王慎德以懷遠，故四夷皆賓服。}無有遠邇，畢獻方物，惟服食器用。\footnote{天下萬國無有遠近，盡貢其方土所生之物，惟可以供服食器用者。言不為耳目華侈。○供音恭。為，於偽反。侈，昌氏反，又式氏反。}王乃昭德之致於異姓之邦，無替厥服。\footnote{德之所致，謂遠夷之貢，以分賜異姓諸侯，使無廢其職。}分寶玉於伯叔之國，時庸展親。\footnote{以寶玉分同姓之國，是用誠信其親親之道。}


{\noindent\zhuan\zihao{6}\fzbyks 傳“天下”至“華侈”。正義曰：以言“無有遠近”，是華夷總統之辭。\CJKunderwave{釋詁}云：“畢，盡也。”故云天下萬國無有遠之與近，盡貢其方土所生之物。“惟可以供服食器用”者，玄纁絺紵供服也，橘柚菁茅供食也，羽毛齒革瑤琨篠簜供器用也。下言“不役耳目”，故知言“不為耳目華侈”也。\CJKunderwave{周禮·大行人}云：“九州之外謂之蕃國，世壹見，各以其所貴寶為贄。”\CJKunderline{鄭玄}云：“所貴寶見經傳者，犬戎獻白狼、白鹿是也，餘外則\CJKunderwave{周書·王會}備焉。”案\CJKunderwave{王會}篇諸方致貢，無所不有，此言“惟服食器用”者，遠方所貢雖不充於器用，實亦受之，\CJKunderline{召公}深戒\CJKunderline{武王}，故言此耳。 \par}

{\noindent\zhuan\zihao{6}\fzbyks 傳“德之”至“其職”。正義曰：明王有德，四夷乃貢，是“德之所致，謂遠夷之貢”也。“昭德之致”，正謂賜異姓諸侯，令其見此遠物，服德畏威，無廢其貢獻常職也。\CJKunderwave{魯語}稱，\CJKunderline{武王}時，“肅慎氏來貢楛矢、石砮、長尺有咫。先王欲昭令德之致遠，以示後人，使永監焉，故銘其楛曰‘肅慎氏貢矢’，以分大姬，配虞胡公而封諸陳。古者分異姓以遠方之貢，使無忘服也。故分陳以肅慎氏之矢”。是分異姓之事，禮有異姓庶姓，異姓，王之甥舅；庶姓與王無親。其分庶姓亦當以遠方之貢矣。 \par}

{\noindent\zhuan\zihao{6}\fzbyks 傳“以寶”至“之道”。正義曰：寶玉亦是萬國所貢，但不必是遠方所貢耳。“以寶玉分同姓之國”，示己不愛惜，共諸侯有之，是“用誠信其親親之道”也。言用寶以表誠心，使彼知王親愛之也。定四年\CJKunderwave{左傳}稱分魯公以夏后氏之璜，是“以寶玉分同姓”也。異姓疏，慮其廢職，故賜以遠方之物攝彼心。同姓親，嫌王無恩，賜以寶玉貴物錶王心。此亦互相見也。 \par}

{\noindent\shu\zihao{5}\fzkt “曰嗚呼”至“展親”。正義曰：“嗚呼！”嘆而言也。自古明聖之王，慎其德教以柔遠人，四夷皆來賓服。無有遠之與近，盡貢其方土所生之物。其所獻者惟可以供其服食器用而已，不為耳目華侈供玩好之用也。明王既得所貢，乃明其德之所致，分賜於彼異姓之國，明己德致遠，賜異姓之國，令使無廢其服職事也。分寶玉於同姓伯叔之國，見已無所愛惜，是用誠信其親親之道也。 \par}

人不易物，惟德其物。\footnote{言物貴由人，有德則物貴，無德則物賤，所貴在於德。○易，羊質反。}德盛不狎侮。\footnote{盛德必自敬，何狎易侮慢之有？○易，以豉反。}狎侮君子，罔以盡人心。\footnote{以虛受人，則人盡其心矣。○盡，津忍反，下同。}狎侮小人，罔以盡其力。\footnote{以悅使民，民忘其勞，則力盡矣。}


{\noindent\zhuan\zihao{6}\fzbyks 傳“言物”至“於德”。正義曰：有德不濫賞，賞必加於賢人，得者則以為榮，故“有德則物貴”也。無德則濫賞，賞或加於小人，賢者得之反以為恥，故“無德則物賤”也。所貴不在於物，乃在於德。 \par}

{\noindent\zhuan\zihao{6}\fzbyks 傳“以虛”至“心矣。”。正義曰：“以虛受人”，\CJKunderwave{易·咸卦}象辭也。人主以己為虛，受用人言，執謙以下人，則人皆盡其心矣。 \par}

{\noindent\zhuan\zihao{6}\fzbyks 傳“以悅”至“盡矣。”。正義曰：\CJKunderwave{詩序}云：“悅以使民，民忘其死。”故云“以悅使民，民忘其勞”。在上撫悅之，則人皆盡其力矣。此“君子”謂臣，“小人”謂民，\CJKunderwave{太甲}曰“接下思恭”，不可狎侮臣也。\CJKunderwave{論語}雲“使民如承大祭”，不可狎侮民也。襄九年\CJKunderwave{左傳}雲“君子勞心，小人勞力”，故別言之。 \par}

{\noindent\shu\zihao{5}\fzkt “人不”至“其力”。正義曰：既言分物賜人，因說貴不在物。言有德無德之王，俱是以物賜人，所賜之物一也，不改易其物。惟有德者賜人，其此賜者是物。若無德者賜人，則此物不是物矣。恐人主恃已賜人，不自修德，言此者，戒人主使修德也。又說修德之事，德盛者常自敬身，不為輕狎侮慢之事。狎侮君子，則無以盡人心，君子被君侮慢，不肯盡心矣。狎侮小人，則無以盡其力，小人被君侮慢，不復肯盡力矣。君子不盡心，小人不盡力，則國家之事敗矣。 \par}

不役耳目，百度惟貞。\footnote{言不以聲色自役，則百度正。}玩人喪德，玩物喪志。\footnote{以人為戲弄則喪其德，以器物為戲弄則喪其志。○玩，五貫反。喪，息浪反。}志以道寧，言以道接。\footnote{在心為志，發氣為言，皆以道為本，故君子勤道。}不作無益害有益，功乃成。不貴異物賤用物，民乃足。\footnote{遊觀為無益，奇巧為異物，言明王之道以德義為益，器用為貴，所以化治生民。○觀，官喚反。}


{\noindent\zhuan\zihao{6}\fzbyks 傳“言不”至“度正”。正義曰：昭元年\CJKunderwave{左傳}子產論晉侯之疾云：“茲心不爽，昏亂百度。”杜預云：“百度，百事之節也。”此言志既不營聲色，百事皆自用心，則皆得正也。 \par}

{\noindent\zhuan\zihao{6}\fzbyks 傳“以人”至“其志”。正義曰：“喪德”、“喪志”其義一也。“玩人”為重，以“德”言之；“玩物”為輕，以“志”言之；終是志荒而德喪耳。 \par}

{\noindent\zhuan\zihao{6}\fzbyks 傳“在心”至“勤道”。正義曰：“在心為志”，\CJKunderwave{詩序}文也。“在心為志”謂心動有所向也，“發氣為言”言於志所趣也。志是未發，言是已發，相接而成，本末之異耳。志、言並皆用道，但志未發，故“以道寧”，志不依道，則不得寧耳。言是已發，故“以道接”，言不以道，則不可接物。志、言皆以道為本，故君子須勤道也。 \par}

{\noindent\zhuan\zihao{6}\fzbyks 傳“遊觀”至“生民”。正義曰：遊觀徒費時日，故為“無益”。無益多矣，非徒遊觀而已。奇巧世所希有，故為“異物”。異物多矣，非徒奇巧而已。諸是妄作，皆為無益。諸是世所希，皆為異物。異物、無益不可遍舉，舉此二者以明此類皆是也。“不作”是初造之辭，為作有所害，故以為“無益”。“不貴”是愛好之語，有貴必有賤，故以“異物”對“用物”。雖經言“用物”，傳言“器用”可矣。經言“有益”，“有益”不知所謂，故傳以德義是人之本，故德義為有益。諸是益身之物，皆是有益，亦舉重為言。經之戒人主，人主如此，所以化世俗，生養下民也。此言“生民”，宣十二年\CJKunderwave{左傳}雲“分謗生民”，皆謂生活民也。下雲“生民保厥居”，與\CJKunderwave{孝經}雲“生民之本盡矣”，言民生於世，謂之“生民”，與此傳異也。俗本雲“弗賤”，衍“弗”字也。 \par}

犬馬非其土性不畜，\footnote{非此土所生不畜，以不習其用。○畜，許竹反。}珍禽奇獸不育於國\footnote{皆非所用，有損害故。}。不寶遠物，則遠人格。\footnote{不侵奪其利，則來服矣。}所寶惟賢，則邇人安。\footnote{寶賢任能，則近人安。近人安，則遠人安矣。}

{\noindent\zhuan\zihao{6}\fzbyks 傳“非此”至“其用”。正義曰：此篇為戒，止為此句，以西旅之獒，非中國之大，不用令王愛好之，故言此也。僖十五年\CJKunderwave{左傳}言晉侯乘鄭馬，及戰陷於濘，是非此土所生不習其用也。犬不習用，傳記無文。 \par}

{\noindent\zhuan\zihao{6}\fzbyks 傳“寶賢”至“安矣”。正義曰：\CJKunderwave{詩序}雲“任賢使能，周室中興”，故傳以“任能”配“寶賢”言之。\CJKunderwave{論語}雲“舉直錯諸枉，則民服”，故“寶賢任能，則近人安”。嫌安近不及遠，故云“近人安，則遠人安矣”。\CJKunderwave{楚語}云：“王孫圉聘於晉，定公饗之。趙簡子鳴玉以相，問於王孫圉曰：‘楚之白珩猶在乎？’對曰：‘然。’簡子曰：‘其為寶也幾何矣？’曰：‘未嘗為寶。楚之所寶者，曰觀射父，及左史倚相，此楚國之寶也。若夫白珩，先王之所玩，何寶之焉？’”是謂“寶賢”也。 \par}

{\noindent\shu\zihao{5}\fzkt “不役”至“道接”。正義曰：既言不可狎侮，又言不可縱恣。不以聲色使役耳目，則百事之度惟皆正矣。以聲色自娛，必玩弄人物。既玩弄人者，喪其德也；玩弄物者，喪其志也。人物既不可玩，則當以道自處。志當以道而寧身，言當以道而接物，依道而行，則志自得而言自當。 \par}

嗚呼！夙夜罔或不勤，\footnote{言當早起夜寐，常勤於德。}不矜細行，終累大德。\footnote{輕忽小物，積害毀大，故君子慎其微。○行，下孟反。累，劣偽反。}為山九仞，功虧一簣。\footnote{八尺曰仞，喻向成也。未成一簣，猶不為山，故曰功虧一簣。是以聖人乾乾日昃慎終如始。○仞音刃，字又作刃，七尺曰仞。虧，曲為反。簣，其貴反。向，許亮反。幹，其連反。昃音側。}允迪茲，生民保厥居，惟乃世王。”\footnote{言其能信蹈行此誡，則生人安其居，天子乃世世王天下。\CJKunderline{武王}雖聖，猶設此誡，況非聖人，以無誡乎？其不免於過，則亦宜矣。○世王如字，又於況反，注同。}


{\noindent\zhuan\zihao{6}\fzbyks 傳“輕忽”至“其微”。正義曰：“矜”是憐惜之意，故以不惜細行為“輕忽小物”，謂上狎侮君子小人、愛玩犬馬禽獸之類是小事也。積小害，毀大德，故君子慎其微。\CJKunderwave{易·繫辭}曰：“小人以小善為無益而不為也，以小惡為無傷而不去也，故惡積而不可掩，罪大而不可解。”是故君子當慎微也。 \par}

{\noindent\zhuan\zihao{6}\fzbyks 傳“八尺”至“如始”。正義曰：\CJKunderwave{周禮·匠人}有畎、遂、溝、洫皆廣深等，而澮雲“廣二尋，深二仞”，則澮亦廣深等，仞與尋同，故知“八尺曰仞”。王肅\CJKunderwave{聖證論}及注\CJKunderwave{家語}皆雲“八尺曰仞”，與孔義同。\CJKunderline{鄭玄}雲“七尺曰仞”，與孔意異。\CJKunderwave{論語}云：“譬如為山，未成一簣。”鄭云：“簣，盛土器。”“為山九仞”，欲成山，以喻為善向成也。未成一簣，猶不為山，故曰為山“功虧一簣”。古語云：“行百里者半於九十。”言末路之艱難也。是以聖人乾乾不息，至於日昃，不敢自暇，恐末路之失，同於一簣，故“慎終如始”也。“乾乾”，\CJKunderwave{易·乾卦}文。“日昃”，無逸篇文。 \par}

{\noindent\zhuan\zihao{6}\fzbyks 傳“言其”至“宜矣”。正義曰：此總結上文，“信蹈行此誡”，行此以上言也。言君主於治民，故先雲“生民安其居，天子乃得世世王天下”也。傳以庸君多自用己，不受人言，敘經意而申之云，\CJKunderline{武王}雖聖，\CJKunderline{召公}猶設此誡，況非聖人，可以無誡乎？身既非聖，又無善誡，其不免於過則，亦宜其然矣。 \par}

{\noindent\shu\zihao{5}\fzkt “嗚呼”至“世王”。正義曰：聽戒以終，故嘆以結之。嗚呼！為人君所當早起夜寐，無有不勤於德，言當勤行德也。若不矜惜細行，作隨宜小過，終必損累大德矣。譬如為山，已高九仞，其功虧損在於一簣。惟少一簣而止，猶尚不成山，以喻樹德行政，小有不終，德政則不成矣。必當慎終如始，以成德政。王者信能蹈行此誡，生民皆安其居處，惟天子乃世世王天下也。 \par}

\CJKunderline{巢伯}來朝，\footnote{殷之諸侯。伯，爵也。南方遠國。\CJKunderline{武王}克商，慕義來朝。○巢，仕交反，徐呂交反。}\CJKunderline{芮伯}作\CJKunderwave{旅巢命}。\footnote{\CJKunderline{芮伯}，周同姓，圻內之國，為卿大夫。陳威德以命巢。亡。○芮，如銳反。圻音祁。}


{\noindent\zhuan\zihao{6}\fzbyks 傳“殷之”至“來朝”。正義曰：\CJKunderline{武王}克商，即來受周之王命，知是“殷之諸侯”。“伯”是爵也。\CJKunderwave{仲虺之誥}雲“\CJKunderline{成湯}放桀於南巢”，或此“巢”是也，故先儒相傳皆以為南方之國。今聞\CJKunderline{武王}克商，慕義而來朝也。\CJKunderline{鄭玄}以為“南方世一見者”。孔以夷狄之爵不過子，此君伯爵，夷夏未明，故直言“遠國”也。 \par}

{\noindent\zhuan\zihao{6}\fzbyks 傳“\CJKunderline{芮伯}”至“巢亡”。正義曰：\CJKunderwave{世本}雲“\CJKunderline{芮伯}，姬姓”，是“周同姓”也。杜預云：“芮，馮翊臨晉縣芮鄉是也。”知是“圻內之國”者，\CJKunderline{芮伯}在朝作命，必是王臣。不得其官，故“卿”與“大夫”並言之。“旅”訓為陳，陳王威德以命巢。 \par}

{\noindent\shu\zihao{5}\fzkt “\CJKunderline{巢伯}”至“巢命”。正義曰：“\CJKunderline{巢伯}”，國爵之君，南方遠國也。以\CJKunderline{武王}克商，乃慕義來朝。王之卿大夫有\CJKunderline{芮伯}者，陳王威德以命巢君。史敘其事，作\CJKunderwave{旅巢命}之篇。 \par}

\section{金縢第八}


\CJKunderline{武王}有疾，\CJKunderline{周公}作\CJKunderwave{金縢}。\footnote{為請命之書，藏之於匱，緘之以金，不欲人開之。○\CJKunderline{武王}有疾，馬本作“有疾不豫”。縢，徒登反。緘，工咸反。}

金縢\footnote{遂以所藏為篇名。}

{\noindent\zhuan\zihao{6}\fzbyks 傳“為請”至“開之”。正義曰：經雲“金縢之匱”，則“金縢”是匱之名也。\CJKunderwave{詩}述韔弓之事云：“竹閉緄縢。”\CJKunderwave{毛傳}云：“緄，繩。縢，約也。”此傳言“緘之以金”，則訓“縢”為緘。王、鄭皆云：“縢，束也。”又鄭\CJKunderwave{喪大記}注云：“齊人謂棺束為緘。”\CJKunderwave{家語}稱周廟之內有金人，叄緘其口，則“縢”是束縛之義。“藏之於匱，緘之以金”，若今釘鐷之，不欲人開也。鄭云：“凡藏秘書，藏之於匱，必以金緘其表。”是秘密之書，皆藏於匱，非\CJKunderline{周公}始造此匱，獨藏此書也。“金縢”。正義曰：發首至“\CJKunderline{王季}、\CJKunderline{文王}”，史敘將告神之事也。“史乃策祝”至“屏璧與珪”，告神之辭也。自“乃卜”至“乃瘳”，言卜吉告王差之事也。自“\CJKunderline{武王}既喪”已下，敘\CJKunderline{周公}被流言，東征還反之事也。此篇敘事多而言語少，若使\CJKunderline{周公}不遭流言，則請命之事遂無人知。為\CJKunderline{成王}開書，\CJKunderline{周公}得反，史官美大其事，故敘之以為此篇。 \par}

{\noindent\shu\zihao{5}\fzkt “\CJKunderline{武王}”至“金縢”。正義曰：\CJKunderline{武王}有疾，\CJKunderline{周公}作策書告神，請代\CJKunderline{武王}死。事畢，納書於金縢之匱，遂作\CJKunderwave{金縢}。凡序言“作”者，謂作此篇也。案經\CJKunderline{周公}策命之書，自納金縢之匱，及為流言所謗，\CJKunderline{成王}悟而開之。史敘其事，乃作此篇，非\CJKunderline{周公}作也。序以經具，故略言之。 \par}

既克商二年，王有疾，弗豫。\footnote{伐紂明年，\CJKunderline{武王}有疾，不悅豫。○豫本，又作忬。}二公曰：“我其為王穆卜。”\CJKunderline{周公}曰：“未可以戚我先王。”\footnote{穆，敬。戚，近也。\CJKunderline{召公}、太公言王疾當敬卜吉兇，\CJKunderline{周公}言未可以死近我先王。相順之辭。○為，於偽反。戚，千歷反。}公乃自以為功，\footnote{\CJKunderline{周公}乃自以請命為己事。}為三壇同\xpinyin*{墠}。\footnote{因太王、\CJKunderline{王季}、\CJKunderline{文王}請命於天，故為三壇。壇築土，墠除地，大除地，於中為三壇。○壇，徒丹反，築土也，馬云：“土堂。”墠音善。}

{\noindent\zhuan\zihao{6}\fzbyks 傳“伐紂”至“悅豫”。正義曰：\CJKunderline{武王}以\CJKunderline{文王}受命十三年伐紂，既殺紂，即當稱元年。克紂稱元年，知此“二年”是“伐紂之明年”也。王肅亦云：“克殷明年。”\CJKunderwave{顧命}云：“王有疾，不懌。”“懌，悅也，故不豫為“不悅豫”也。何休因此為例云：“天子曰不豫，諸侯曰負茲，大夫曰犬馬，士曰負薪。” \par}

{\noindent\zhuan\zihao{6}\fzbyks 傳“穆敬”至“之辭”。正義曰：\CJKunderwave{釋訓}云：“穆穆，敬也。”“戚”是親近之義，故為近也。\CJKunderline{武王}時三公惟周、召與太公耳，知“二公”是\CJKunderline{召公}、太公也。言王疾恐死，當敬卜吉兇。\CJKunderline{周公}言\CJKunderline{武王}既定天下，當成就周道，未可以死近我先王。死則神與先王相近，故言近先王。若生則人神道隔，是為遠也。二公恐王死，欲為之卜。\CJKunderline{周公}言王未可以死，是“相順之辭”也。鄭云：“戚，憂也。\CJKunderline{周公}既內知\CJKunderline{武王}有九齡之命，又有\CJKunderline{文王}曰‘吾與爾三之期’，今必瘳，不以此終，故止二公之卜。雲未可以憂怖我先王。”如鄭此言，\CJKunderline{周公}知王不死，先王豈不知乎，而慮先王憂也？ \par}

{\noindent\zhuan\zihao{6}\fzbyks 傳“\CJKunderline{周公}”至“己事”。正義曰：“功”訓事也。\CJKunderline{周公}雖許二公之卜，仍恐王疾不瘳，不復與二公謀之，乃自以請命為己之事，獨請代\CJKunderline{武王}死也。所以\CJKunderline{周公}自請為己事者，\CJKunderline{周公}位居冢宰，地則近親，脫或卜之不善，不可使外人知悉，亦不可苟讓，故自以為功也。 \par}

{\noindent\zhuan\zihao{6}\fzbyks 傳“因大”至“三壇”。正義曰：“請命”請之於天，而告三王者，以三王精神已在天矣，故“因\CJKunderline{大王}、\CJKunderline{王季}、\CJKunderline{文王}以請命於天”。三王每王一罈，故“為三壇”。壇是築土，墠是除地，大除其地，於中為三壇。\CJKunderline{周公}為壇於南方，亦當在此墠內，但其處小別，故下別言之。\CJKunderline{周公}北面，則三壇南面可知，但不知以何方為上耳。\CJKunderline{鄭玄}云：“時為壇墠於豐，壇墠之處猶存焉。” \par}

為壇於南方，北面，\CJKunderline{周公}立焉。\footnote{立壇上，對三王。}植璧秉珪，乃告\CJKunderline{大王}、\CJKunderline{王季}、\CJKunderline{文王}。\footnote{璧以禮神。植，置也，置於三王之坐。\CJKunderline{周公}秉桓珪以為贄。告謂祝辭。○植，時織反，徐音置。贄音至。祝如字，或之疚反，下同。}

{\noindent\zhuan\zihao{6}\fzbyks 傳“立壇”至“三王”。正義曰：\CJKunderwave{禮}“授坐不立”，“授立不坐”，欲其高下均也。神位在壇，故\CJKunderline{周公}“立壇上，對三王”也。 \par}

{\noindent\zhuan\zihao{6}\fzbyks 傳“璧以”至“祝辭”。正義曰：\CJKunderwave{周禮·大宗伯}雲“以蒼璧禮天”，\CJKunderwave{詩}說禱旱雲“圭璧既卒”，是璧以禮神，不知其何色也。鄭云：“植，古置字。”故為置也，言置璧於三王之坐也。\CJKunderwave{周禮}云：“公執桓圭。”知\CJKunderline{周公}秉桓圭，又置以為贄也。“告謂祝辭”，下文是其辭也。 \par}

{\noindent\shu\zihao{5}\fzkt “既克”至“\CJKunderline{文王}”。正義曰：“既克商二年”即伐紂之明年也。王有疾病，不悅豫。\CJKunderline{召公}與太公二公同辭而言曰：“我其為王敬卜吉兇，問王疾病瘳否。”\CJKunderline{周公}曰：“王今有疾，未可以死近我先王，故當須卜也。”\CJKunderline{周公}既為此言，公乃自以請命之事為己事，除地為墠，墠內築壇，為三壇同墠。又為一罈於南方，北面，\CJKunderline{周公}立壇上焉。置璧於三王之坐，公自執珪，乃告\CJKunderline{大王}、\CJKunderline{王季}、\CJKunderline{文王}，告此三王之神也。 \par}

史乃冊，祝曰：“惟爾元孫某，遘厲虐疾。\footnote{史為冊書，祝辭也。元孫，\CJKunderline{武王}。某，名。臣諱君，故曰某。厲，危。虐，暴也。○遘，工豆反，遇也。}若爾三王，是有丕子之責於天，以\CJKunderline{旦}代某之身。\footnote{大子之責，謂疾不可救於天，則當以旦代之。死生有命，不可請代，聖人敘臣子之心，以垂世教。○丕，普悲反，馬同，徐甫眉反，鄭音不。}予仁若考，能多材多藝，能事鬼神。\footnote{我\CJKunderline{周公}仁能順父，又多材多藝，能事鬼神。言可以代\CJKunderline{武王}之意。}乃元孫不若旦多材多藝，不能事鬼神。


{\noindent\zhuan\zihao{6}\fzbyks 傳“史為”至“虐暴”。正義曰：告神之言，書之於策，“祝”是讀書告神之名，故云“史為策書，祝辭”，史讀此策書以祝告神也。\CJKunderline{武王}是\CJKunderline{大王}之曾孫也，尊統於上，繼之於祖，謂“元孫”，是長孫也。“某”者，\CJKunderline{武王}之名，本告神雲“元孫發”，臣諱君，故曰“某”也。\CJKunderwave{易·乾卦}云：“夕惕若厲。”“厲”為危也。“虐”訓為暴。言性命危而疾暴重也。\CJKunderwave{泰誓}、\CJKunderwave{牧誓}皆不諱發而此獨諱之，孔惟言“臣諱君”，不解諱之意。\CJKunderline{鄭玄}云：“諱之者，由\CJKunderline{成王}讀之也。”意雖不明，當謂\CJKunderline{成王}開匱得書，王自讀之，至此字口改為“某”，史官錄為此篇，因遂\CJKunderline{成王}所讀，故諱之。上篇\CJKunderwave{泰誓}、\CJKunderwave{牧誓}王自稱者，令入史製為此典，故不須諱之。 \par}

{\noindent\zhuan\zihao{6}\fzbyks 傳“太子”至“世教”。正義曰：“責”讀如\CJKunderwave{左傳}“施捨已責”之“責”，“責”謂負人物也。“太子之責於天”，言負天一太子，謂必須死，疾不可救於天。必須一子死，則當以旦代之。死生有命，不可請代，今請代者，“聖人敘臣子之心，以垂世教”耳，非謂可代得也。\CJKunderline{鄭玄}弟子趙商問玄曰：“若\CJKunderline{武王}未終，疾固當瘳。信命之終，雖請不得。自古已來，何患不為？”玄答曰：“君父疾病方困，忠臣孝子不忍默爾，視其歔欷，歸其命於天，中心惻然，欲為之請命。\CJKunderline{周公}達於此禮，著在\CJKunderwave{尚書}，若君父之病不為請命，豈忠孝之志也？”然則命有定分，非可代死，\CJKunderline{周公}為此者，自申臣子之心，非謂死實可代。自古不廢，亦有其人，但不見爾，未必\CJKunderline{周公}獨為之。\CJKunderline{鄭玄}云：“丕讀曰不。及子孫曰子。元孫遇疾，若汝不救，是將有不愛子孫之過，為天所責，欲使為之請命也。”與孔讀異。 \par}

{\noindent\zhuan\zihao{6}\fzbyks 傳“我周”至“之意”。正義曰：告神稱“予”，知\CJKunderline{周公}自稱“我”也。“考”是父也，故“仁能順父”。上雲“元孫”，對祖生稱，此言“順父”，從親為始。祖為王考，曾祖為皇考，“考”、“父”，可以通之，傳舉親而言“父”耳。既能順父，又多材多藝，能事鬼神，言己可以代\CJKunderline{武王}之意。上言“丕子之責於天”，則是天欲取\CJKunderline{武王}，非父祖取之，此言己能順父祖，善事鬼神者，假令天意取之，其神必共父祖同處，言己是父祖所欲，欲令請之於天也。 \par}

乃命於帝庭，敷佑四方。\footnote{汝元孫受命於天庭為天子，布其德教，以佑助四方。言不可以死。}用能定爾子孫于下地，四方之民，罔不祗畏。\footnote{言\CJKunderline{武王}用受命帝庭之故，能定先人子孫於天下，四方之民無不敬畏。}嗚呼！無墜天之降寶命，我先王亦永有依歸。\footnote{嘆惜\CJKunderline{武王}，言不救則墜天之寶命，救之則先王長有依歸。}今我即命於元龜，\footnote{就受三王之命於大龜，卜知吉凶。}爾之許我，我其以璧與珪歸俟爾命。\footnote{許謂疾瘳。待命，當以事神。○瘳，敕留反，下同。}爾不許我，我乃屏璧與珪。\footnote{不許謂不愈也。屏，藏也，言不得事神。}

{\noindent\zhuan\zihao{6}\fzbyks 傳“汝元”至“以死”。正義曰：以王者存亡，大運在天，有德於民，天之所與，是“受命天庭”也。以人況天，故言在庭，非王實至天庭受天命也。既受天命以為天子，布其德教以佑助四方之民，當於天心有功，於民言不可以死也。 \par}

{\noindent\shu\zihao{5}\fzkt “史乃”至“與珪”。正義曰：史乃為策書，執以祝之曰，惟爾元孫某，“某”即發也，遇得危暴重疾，今恐其死。若爾三王，是有太子之責於天，謂負天大子責，必須一子死者，請以旦代發之身，令旦死而發生。又告神以代之狀，我仁能順父，又且多材力，多技藝，又能善事鬼神，汝元孫不如旦多材多藝，又不能事鬼神，言取發不如取旦也。然人各有能，發雖不能事鬼，神則有人君之用，乃受命於天帝之庭，能布其德教以佑助四方之民，用能安定汝三王子孫在於下地，四方之民無不敬而畏之。以此之故，不可使死。嗚呼！發之可惜如此，神明當救助之，無得隕墜天之所下寶命。天下寶命謂使為天子，若\CJKunderline{武王}死，是隕墜之也。若不墜命，則我先王亦永有依歸，為宗廟之主，神得歸之。我與三王人神道隔，許我以否不可知，今我就受三王之命於彼大龜，卜其吉凶。吉則許我，兇則為不許我。爾之許我，使卜得吉兆，旦死而發生，我其以璧與珪歸家待汝神命，我死當以珪璧事神。爾不許我，使卜兆不吉，發死而旦生，我乃屏去璧之與珪。言不得事神，當藏珪璧也。 \par}

乃卜三龜，一習吉。\footnote{習，因也。以三王之龜卜，一相因而吉。}啟籥見書，乃並是吉。\footnote{三兆既同吉，開籥見佔兆書，乃亦並是吉。○籥，於若反，徐以略反，馬云：“藏卜兆書管。”並，必政反。}公曰：“體，王其罔害。\footnote{公視兆曰：“如此兆體，王其無害。”言必愈。}予小子新命於三王，惟永終是圖。\footnote{\CJKunderline{周公}言，我小子新受三王之命，\CJKunderline{武王}惟長終是謀周之道。}


{\noindent\zhuan\zihao{6}\fzbyks 傳“習因”至“而吉”。正義曰：“習”則襲也，襲是重衣之名，因前而重之，故以“習”為因也。雖三龜並卜，卜有先後，後者因前，故云“因”也。\CJKunderwave{周禮}：“太卜掌三兆之法，一曰\CJKunderwave{玉兆}，二曰\CJKunderwave{瓦兆}，三曰\CJKunderwave{原兆}”。三兆各別，必三代法也。\CJKunderwave{洪範}卜筮之法，三人佔則從二人之言，是必三代之法並用之矣。故知“三龜”，“三王之龜”。龜形無異代之別，但卜法既別，各用一龜，謂之“三王之龜”耳。每龜一人佔之，其後君與大夫等，部佔三代之龜，定其吉凶。未見佔書知已吉者，卜有大體，見兆之吉凶，粗觀可識，故知吉也。 \par}

{\noindent\zhuan\zihao{6}\fzbyks 傳“三兆”至“是吉”。正義曰：\CJKunderline{鄭玄}云：“籥，開藏之管也。開兆書藏之室以管，乃復見三龜佔書，亦合於是吉。”王肅亦云：“籥，開藏佔兆書管也。”然則佔兆別在於藏。\CJKunderwave{大卜}“三兆”之下云：“其經兆之體，皆百有二十，其頌皆千有二百。”佔兆之書，則彼“頌”是也。略觀三兆，既已同吉，開藏以籥，見彼佔兆之書，乃亦並是吉。言其兆頌符，同為大吉也。 \par}

{\noindent\zhuan\zihao{6}\fzbyks 傳“公視”至“必愈”。正義曰：“如此兆體”，指卜之所得兆也。\CJKunderwave{周禮·占人}云：“凡卜筮，君佔體，大夫佔色，史佔墨，卜人佔坼。”\CJKunderline{鄭玄}云：“體，兆象也。色，兆氣也。墨，兆廣也。坼，兆璺也。尊者視兆象而已，卑者以次詳其餘也。\CJKunderline{周公}卜\CJKunderline{武王}，佔之曰：‘體，王其無害。’”鄭意此言“體”者，即彼“君佔體”也。但\CJKunderline{周公}令卜，汲汲欲王之愈，必當親視灼龜，躬省兆繇，不惟佔體而日。但鄭以“君佔體”與此文同，故引以為證耳。 \par}

茲攸俟，能念予一人。”\footnote{言\CJKunderline{武王}愈，此所以待能念我天子事，成周道。○公歸，乃納冊於金縢之匱中。}王翼日乃瘳。\footnote{從壇歸。翼，明。瘳，差也。○差，初賣反。}

{\noindent\zhuan\zihao{6}\fzbyks 傳“言武”至“周道”。正義曰：此原三王之意也。言\CJKunderline{武王}得愈者，此謂卜吉\CJKunderline{武王}之愈。言天與三王一一須待\CJKunderline{武王}，能念我天子事，成周道。若死，則不復得念天子之事，周道必不成也。\CJKunderwave{禮}天子自稱曰“予一人”，故以“一人”言天子也。 \par}

{\noindent\zhuan\zihao{6}\fzbyks 傳“從壇”至“瘳差”。正義曰：壇所即卜，故“從壇歸”也。“翼，明”，\CJKunderwave{釋言}文。“瘳”訓差，亦為愈，病除之名也。藏此書者，此既告神，即是國家舊事，其書不可捐棄，又不可示諸世人，故藏於金縢之匱耳。 \par}

{\noindent\shu\zihao{5}\fzkt “乃卜”至“乃瘳”。正義曰：祝告已畢，即於壇所乃卜其吉凶。用三王之龜卜，一皆相因而吉。觀兆已知其吉，猶尚未見佔書。佔書在於藏內，啟藏以籥，見其佔書，亦與兆體乃並是吉。公視兆曰，觀此兆體，王身其無患害也。我小子新受命於三王，謂卜得吉也。我\CJKunderline{武王}當惟長終是謀周之道。此卜吉之愈者，上天所以須待\CJKunderline{武王}能念我一人天子之事，成其周道故也。公自壇歸，乃納策於金縢之匱中。王明日乃病瘳。 \par}

\CJKunderline{武王}既喪，\CJKunderline{管叔}及其群弟乃流言於國，\footnote{\CJKunderline{武王}死，\CJKunderline{周公}攝政，其弟\CJKunderline{管叔}及\CJKunderline{蔡叔}、\CJKunderline{霍叔}乃放言於國，以誣\CJKunderline{周公}，以惑\CJKunderline{成王}。○喪，蘇浪反。}曰：“公將不利於孺子。”\footnote{三叔以\CJKunderline{周公}大聖，有次立之勢，遂生流言。孺，稚也。稚子，\CJKunderline{成王}。○孺，如樹反。}\CJKunderline{周公}乃告二公曰：“我之弗闢，我無以告我先王。”\footnote{闢，法也。告\CJKunderline{召公}、太公，言我不以法法三叔，則我無以成周道告我先王。○闢，扶亦反，治也；\CJKunderwave{說文}作壁，雲必亦反，法也；馬、鄭音避，謂避居東都。}\CJKunderline{周公}居東二年，則罪人斯得。\footnote{\CJKunderline{周公}既告二公，遂東征之，二年之中，罪人此得。}


{\noindent\zhuan\zihao{6}\fzbyks 傳“\CJKunderline{武王}死”至“\CJKunderline{成王}”。正義曰：\CJKunderline{武王}既死，\CJKunderline{成王}幼弱，故\CJKunderline{周公}攝政。攝政者，雖以\CJKunderline{成王}為主，政令自公出，不復關\CJKunderline{成王}也。\CJKunderwave{蔡仲之命}云：“群叔流言，乃致闢\CJKunderline{管叔}於商，囚\CJKunderline{蔡叔}於郭鄰，降\CJKunderline{霍叔}於庶人。”則知“群弟”是\CJKunderline{蔡叔}、\CJKunderline{霍叔}也。\CJKunderwave{周語}雲“獸三為群”，則滿三乃稱群。蔡霍二人而言群者，並管故稱群也。傳既言\CJKunderline{周公}攝政，乃雲“其弟\CJKunderline{管叔}”，蓋以\CJKunderline{管叔}為\CJKunderline{周公}之弟。\CJKunderwave{孟子}曰：“\CJKunderline{周公}，弟也。\CJKunderline{管叔}，兄也。”\CJKunderwave{史記}亦以\CJKunderline{管叔}為\CJKunderline{周公}之兄。孔似不用\CJKunderwave{孟子}之說，或可孔以“其弟”謂\CJKunderline{武王}之弟，與\CJKunderwave{史記}亦不違也。“流言”者，宣佈其言，使人聞知，若水流然。“流”即放也，乃放言於國，以誣\CJKunderline{周公}，以惑\CJKunderline{成王}。“王亦未敢誚公”，是王心惑也。\CJKunderline{鄭玄}云：“流公將不利於孺子之言於京師，於時管蔡在東，蓋遣人流傳此言於民間也。” \par}

{\noindent\zhuan\zihao{6}\fzbyks 傳“三叔”至“\CJKunderline{成王}”。正義曰：殷法多兄亡弟立，三叔以\CJKunderline{周公}大聖，又是\CJKunderline{武王}之弟，有次立之勢，今復秉國之權，恐其因即篡奪，遂生流言。不識大聖之度，謂其實有異心，非是故誣之也。但啟商共叛，為罪重耳。 \par}

{\noindent\zhuan\zihao{6}\fzbyks 傳“闢，法也”。正義曰：\CJKunderwave{釋詁}文。 \par}

{\noindent\zhuan\zihao{6}\fzbyks 傳“\CJKunderline{周公}”至“此得”。正義曰：\CJKunderwave{詩·東山}之篇歌此事也，序雲“東征”，知“居東”者，遂東往徵也。雖徵而不戰，故言“居東”也。\CJKunderwave{東山}詩曰：“自我不見，於今三年。”又云“三年而歸”，此言“二年”者，\CJKunderwave{詩}言初去及來，凡經三年；此直數居東之年，除其去年，故二年也。罪人既多，必前後得之，故云“二年之中，罪人此得”。惟言“居東”，不知居在何處。王肅云：“東，洛邑也。管蔡與商奄共叛，故東征鎮撫之。案驗其事，二年之間，罪人皆得。” \par}

於後，公乃為詩以貽王，名之曰\CJKunderwave{鴟鴞}。王亦未敢\xpinyin*{誚}公。\footnote{\CJKunderline{成王}信流言而疑\CJKunderline{周公}，故\CJKunderline{周公}既誅三監，而作詩解所以宜誅之意以遺王，王猶未悟，故欲讓公而未敢。○貽，羊支反。名如字，徐亡政反。鴟，尺夷反。鴞，於嬌反。誚，在笑反。以遺，唯季反。}

{\noindent\zhuan\zihao{6}\fzbyks 傳“\CJKunderline{成王}”至“未敢”。正義曰：\CJKunderline{成王}信流言而疑\CJKunderline{周公}，管蔡既誅，王疑益甚，故\CJKunderline{周公}既誅三監，而作詩解所以宜誅之意。其\CJKunderwave{詩}云：“鴟鴞鴟鴞，既取我子，無毀我室。”\CJKunderwave{毛傳}云：“無能毀我室者，攻堅之故也。寧亡二字，不可以毀我周室。”言宜誅之意也。\CJKunderwave{釋言}云：“貽，道也。”以詩遺王，王猶未悟，故欲讓公而未敢。政在\CJKunderline{周公}，故畏威未敢也。\CJKunderline{鄭玄}以為\CJKunderline{武王}崩，\CJKunderline{周公}為冢宰，三年服終，將欲攝政，管蔡流言，即避居東都。\CJKunderline{成王}多殺公之屬黨，公作\CJKunderwave{鴟鴞}之詩，救其屬臣，請勿奪其官位土地。及遭風雷之異，啟金縢之書，迎公來反，反乃居攝，後方始東征管蔡。解此一篇乃\CJKunderwave{鴟鴞}之詩，皆與孔異。 \par}

{\noindent\shu\zihao{5}\fzkt “\CJKunderline{武王}”至“誚公”。正義曰：\CJKunderline{周公}於\CJKunderline{成王}之世，為管蔡所誣，王開金縢之書，方始明公本意，卒得成就周道，天下太平。史官美大其事，述為此篇，故追言“請命”於前，乃說“流言”於後，自此以下，說\CJKunderline{周公}身事。\CJKunderline{武王}既喪，\CJKunderline{成王}幼弱，\CJKunderline{周公}攝王之政，專決萬機。\CJKunderline{管叔}及其群弟\CJKunderline{蔡叔}、\CJKunderline{霍叔}乃流放其言於國中曰：“公將不利於孺子。”言欲篡王位為不利。\CJKunderline{周公}乃告二公曰：“我之不以法法此三叔，則我無以成就周道，告我先王。”既言此，遂東征之。\CJKunderline{周公}居東二年，則罪人於此皆得，謂獲三叔及諸叛逆者。罪人既得訖，\CJKunderline{成王}猶尚疑公。公於此既得罪人之後，為詩遺王，名之曰\CJKunderwave{鴟鴞}。\CJKunderwave{鴟鴞}言三叔不可不誅之意。王心雖疑，亦未敢責誚公。言王意欲責而未敢也。 \par}

秋，大熟，未獲，天大雷電以風，\footnote{二年秋也。蒙，恆風若，雷以威之，故有風雷之異。○獲，戶郭反。}禾盡偃，大木斯拔，邦人大恐。\footnote{風災所及，邦人皆大恐。○拔，皮八反。}王與大夫盡弁，以啟金縢之書，\footnote{皮弁質服以應天。○弁，皮彥反。徐，扶變反。應，應對之應。}乃得\CJKunderline{周公}所自以為功代\CJKunderline{武王}之說。\footnote{所藏請命冊書本。○說如字，徐始銳反。}


{\noindent\zhuan\zihao{6}\fzbyks 傳“二年”至“之異”。正義曰：上文“居東二年”，未有別年之事，知即是“二年秋”也。嫌別年，故辨之。\CJKunderwave{洪範}“咎徵”云：“蒙，恆風若。”以\CJKunderline{成王}蒙暗，故常風順之。風是暗徵而有雷者，以威怒之故，以示天之威怒有雷風之異。 \par}

{\noindent\zhuan\zihao{6}\fzbyks 傳“風災”至“大恐”。正義曰：言“邦人”，則風災惟在周邦，不及寬遠，故云“風災所及，邦人皆大恐”，言獨畿內恐也。 \par}

{\noindent\zhuan\zihao{6}\fzbyks 傳“皮弁質服以應天”。正義曰：皮弁象古，故為“質服”。祭天尚質，故服以應天也。\CJKunderwave{周禮·司服}云：“王祀昊天上帝，則服大裘而冕。”無旒，乃是冕之質者，是事天宜質服，故服之以應天變也。\CJKunderwave{周禮}：“視朝，則皮弁服。”皮弁是視朝服，每日常服而言“質”者，皮弁白布衣，素積裳，故為質也。\CJKunderline{鄭玄}以為爵弁，“必爵弁者，承天變降服，亦如國家未道焉”。 \par}

二公及王乃問諸史與百執事，\footnote{二公倡王啟之，故先見書。史、百執事皆從\CJKunderline{周公}請命。○倡，昌亮反。從，才用反，又如字。}對曰：“信。噫！公命我勿敢言。”\footnote{史、百執事言信有此事，\CJKunderline{周公}使我勿道，今言之則負\CJKunderline{周公}。噫，恨辭。○噫，於其反，馬本作懿，猶億也。}王執書以泣，曰：“其勿穆卜。\footnote{本欲敬卜吉兇，今天意可知，故止之。}昔公勤勞王家，惟予沖人弗及知。\footnote{言己童幼，不及知\CJKunderline{周公}昔日忠勤。○衝，直忠反。}今天動威，以彰\CJKunderline{周公}之德，\footnote{發雷風之威以明\CJKunderline{周公}之聖德。}惟朕小子其新逆，我國家禮亦宜之。”\footnote{\CJKunderline{周公}以\CJKunderline{成王}未寤，故留東未還，改過自新，遣使者迎之，亦國家禮有德之宜。○新逆，馬本作“親迎”。遣使，所吏反。}

{\noindent\zhuan\zihao{6}\fzbyks 傳“二公”至“請命”。正義曰：二公與王若同而問，當言“王及二公”，今言“二公及王”，則是二公先問，知“二公倡王啟之，故先見書”。鄭云：“開金縢之書者，省察變異所由故事也。”以金縢匱內有先王故事，疑其遭遇災變，必有消伏之術，故倡王啟之。史為公造策書，而百執事給使令，皆從\CJKunderline{周公}請命者。 \par}

{\noindent\zhuan\zihao{6}\fzbyks 傳“史百”至“恨辭”。正義曰：\CJKunderline{周公}使我勿道此事者，公以臣子之情，忠心欲代王死，非是規求名譽，不用使人知之。且\CJKunderline{武王}瘳而\CJKunderline{周公}不死，恐人以公為詐，故令知者勿言。今被問而言之，是違負\CJKunderline{周公}也。“噫”者，心不平之聲，故為“恨辭”。 \par}

{\noindent\zhuan\zihao{6}\fzbyks 傳“\CJKunderline{周公}”至“之宜”。正義曰：公之東征，止為伐罪，罪人既得，公即當還。以\CJKunderline{成王}未寤，恐與公不和，故留東未還，待王之察己也。新迎者，改過自新，遣使者迎之。\CJKunderwave{詩·九罭}之篇是迎之事也。“亦國家禮有德之宜”，言尊崇有德，宜用厚禮。\CJKunderwave{詩}稱“袞衣”、“籩豆”，是國家禮也。 \par}

王出郊，天乃雨，反風，禾則盡起。\footnote{郊以玉幣謝天，天即反風起禾，明郊之是。}二公命邦人，凡大木所偃，盡起而築之。歲則大熟。\footnote{木有偃拔，起而立之，築有其根。桑果無虧，百穀豐熟，\CJKunderline{周公}之德。此已上\CJKunderwave{大誥}後，因\CJKunderline{武王}喪並見之。○築音竹，本亦作築，謂築其根，馬云：“築，拾也。”見，賢遍反。}

{\noindent\zhuan\zihao{6}\fzbyks 傳“郊以”至“之是”。正義曰：祭天於南郊，故謂之“郊”，郊是祭天之處也。“王出郊”者，出城至郊，為壇告天也。\CJKunderwave{周禮·大宗伯}云：“以蒼璧禮天，牲幣如其器之色。”是祭天有玉有幣，今言郊者，以玉幣祭天，告天以謝過也。王謝天，天即反風起禾，明王郊之是也。\CJKunderline{鄭玄}引\CJKunderwave{易傳}云：“陽感天不旋日。陽謂天子也，天子行善以感天，不迴旋經口。”故郊之是得反風也。 \par}

{\noindent\zhuan\zihao{6}\fzbyks 傳“木有”至“見之”。正義曰：上文禾偃木拔，拔必亦偃，故云“木有偃拔，起而立之，築有其根，桑果無虧，百穀豐熟”。鄭、王皆雲“築，拾也。禾為大木所偃者，起其木，拾下禾，無所亡失”。意太曲碎，當非經旨。案序將東征，作\CJKunderwave{大誥}。此上“居東二年”以來，皆是\CJKunderwave{大誥}後事，而編於\CJKunderwave{大誥}之前者，因\CJKunderline{武王}喪並見之。 \par}

{\noindent\shu\zihao{5}\fzkt “秋大”至“大熟”。正義曰：為詩遺王之後，其秋大熟，未及收穫，天大雷電，又隨之以風，禾盡偃仆，大木於此而拔。風災所及，邦人大恐。王見此變，與大夫盡皮弁以開金縢之書，案省故事，求變異所由，乃得\CJKunderline{周公}所自以為功請代\CJKunderline{武王}之說。二公及王問於本從公之人史與百執事，問審然以否。對曰：“信。”言有此事也。乃為不平之聲：“噫！公命我勿敢言。”王執書以泣，曰：“其勿敬卜吉兇。”言天之意已可知也。“昔公勤勞王家，惟我幼童之人不及見知，今天動雷電之威，以彰明\CJKunderline{周公}之德，惟朕小子其改過自新，遣人往迎之。我國家褒崇有德之禮，亦宜行之”。王於是出郊而祭以謝天，天乃雨，反風，禾則盡起。二公命邦人，凡大木所偃仆者，盡扶起而築之。禾木無虧，歲則大熟。言\CJKunderline{周公}之所感致若此也。 \par}

\section{大誥第九}


\CJKunderline{武王}崩，三監及淮夷叛，\footnote{三監，管、蔡、商。淮夷徐奄之屬皆叛周。○監，古懺反，視也。}\CJKunderline{周公}相\CJKunderline{成王}，將黜殷，作\CJKunderwave{大誥}。\footnote{相謂攝政。黜，絕也。將以誅叛者之義大誥天下。○相，息亮反，注同。}

{\noindent\zhuan\zihao{6}\fzbyks 傳“三監”至“叛周”。正義曰：知“三監”是管、蔡、商者，以序上下相顧為文。此言“三監及淮夷叛”，總舉諸叛之人也。下雲“\CJKunderline{成王}既黜殷命，殺\CJKunderline{武庚}，命\CJKunderline{微子}啟代殷後”，又言“\CJKunderline{成王}既伐\CJKunderline{管叔}、\CJKunderline{蔡叔}，以殷餘民封\CJKunderline{康叔}”。此序言三監叛，將徵之，下篇之序歷言伐得三人，足知下文\CJKunderline{管叔}、\CJKunderline{蔡叔}、\CJKunderline{武庚}，即此“三監”之謂，知“三監”是“管、蔡、商”也。\CJKunderwave{漢書·地理志}云：“周既滅殷，分其畿內為三國，\CJKunderwave{詩·風}邶、鄘、衛是也。邶以封紂子\CJKunderline{武庚}，鄘\CJKunderline{管叔}尹之，衛\CJKunderline{蔡叔}尹之，以監殷民，謂之三監。”先儒多同此說，惟\CJKunderline{鄭玄}以三監為管、蔡、霍，獨為異耳。謂之“監”者，當以殷之畿內，被紂化日久，未可以建諸侯，且使三人監此殷民，未是封建之也。三人雖有其分，互相監領，不必獨主一方也。\CJKunderwave{史記·衛世家}云：“\CJKunderline{武王}克殷，封紂子\CJKunderline{武庚}為諸侯，奉其先祀。為\CJKunderline{武庚}未集，恐有賊心，乃令其弟\CJKunderline{管叔}、\CJKunderline{蔡叔}傳相之。”是言輔相\CJKunderline{武庚}，共監殷人，故稱“監”也。序惟言“淮夷叛”，傳言“淮夷徐奄之屬共叛周”者，以下序文雲“\CJKunderline{成王}東伐淮夷，遂踐奄，作\CJKunderwave{成王政}”，又云“\CJKunderline{成王}既黜殷命，滅淮夷，作\CJKunderwave{周官}”，又云“魯侯\CJKunderline{伯禽}宅曲阜，徐夷並興，作\CJKunderwave{費誓}”，彼三序者，一時之事，皆在\CJKunderline{周公}歸政之後也。\CJKunderwave{多方}篇數此諸國之罪雲“至於再，至於三”，得不以\CJKunderline{武王}初崩已叛，\CJKunderline{成王}即政又叛，謂此為再三也。以此知“淮夷叛”者，徐奄之屬皆叛也。 \par}

{\noindent\zhuan\zihao{6}\fzbyks 傳“相謂”至“天下”。正義曰：\CJKunderwave{君奭}序云：“\CJKunderline{召公}為保，\CJKunderline{周公}為師，相\CJKunderline{成王}為左右。”於時\CJKunderline{成王}為天子，自知政事，二公為臣輔助之，此言“相\CJKunderline{成王}”者，有異於彼，故辨之“相謂攝政”。攝政者，教由公出，不復關白\CJKunderline{成王}耳，仍以\CJKunderline{成王}為王，故稱“\CJKunderline{成王}”。\CJKunderline{鄭玄}云：“黜，貶退也。”“黜”實退名，但此“黜”乃殺其身，絕其爵，故以“黜”為絕也。\CJKunderline{周公}此行普伐諸叛，獨言黜殷命者，定四年\CJKunderwave{左傳}云：“管蔡啟商，惎間王室”，則此叛\CJKunderline{武庚}為主，且顧\CJKunderwave{微子}之序，故特言黜殷命也。“以誅叛者之義大誥天下”，經皆是也。 \par}

{\noindent\shu\zihao{5}\fzkt “\CJKunderline{武王}”至“大誥”。正義曰：\CJKunderline{武王}既崩，\CJKunderline{管叔}、\CJKunderline{蔡叔}與紂子\CJKunderline{武庚}三人監殷民者又及淮夷共叛。\CJKunderline{周公}相\CJKunderline{成王}，攝王政，將欲東征，黜退殷君\CJKunderline{武庚}之命，以誅叛之義大誥天下。史敘其事，作\CJKunderwave{大誥}。 \par}

大誥\footnote{陳大道以誥天下，遂以名篇。}

{\noindent\shu\zihao{5}\fzkt “大誥”。正義曰：此陳伐叛之義，以大誥天下，而兵兇戰危，非眾所欲，故言煩重。其自殷勤，多止而更端，故數言“王曰”。大意皆是陳說\CJKunderline{武庚}之罪，自言己之不能，言己當繼父祖之功，須去叛逆之賊，人心既從，卜之又吉，往伐無有不克，勸人勉力用心。此時\CJKunderline{武王}初崩，屬有此亂，\CJKunderline{周公}以臣伐君，天下未察其志，親弟猶尚致惑，何況疏賤者乎？\CJKunderline{周公}慮其有向背之意，故殷勤告之。陳壽云：“\CJKunderline{皋陶}之謨略而雅，\CJKunderline{周公}之誥煩而悉。何則？\CJKunderline{皋陶}與舜\CJKunderline{禹}共談，\CJKunderline{周公}與群下矢誓也。”其意或亦然乎。但\CJKunderwave{君奭}、\CJKunderwave{康誥}乃與\CJKunderline{召公}、\CJKunderline{康叔}語也，其辭亦甚委悉，抑亦當時設言，自好煩複也。管蔡導\CJKunderline{武庚}為亂，此篇略於管蔡者，猶難以伐弟為言，故專說\CJKunderline{武庚}罪耳。 \par}

王若曰：“猷！大誥爾多邦，越爾御事。\footnote{\CJKunderline{周公}稱\CJKunderline{成王}命，順大道以誥天下眾國，及於御治事者盡及之。○猷音由，道也。邦，馬本作“大誥繇爾多邦”。盡，津忍反。}弗弔，天降割於我家，\footnote{言周道不至，故天下兇害於我家不少。謂三監淮夷並作難。○吊音的，又如字。割，馬本作害。不少，馬讀“弗少延”為句。難，乃旦反。}不少延，洪惟我幼沖人，\footnote{兇害延大，惟累我幼童人。\CJKunderline{成王}言其不可不誅之意。○累，劣偽反。}嗣無疆大曆服。弗造哲，迪民康，\footnote{言子孫承繼祖考無窮大數，服行其政，而不能為智道以安人，故使叛。先自責。}矧曰其有能格知天命？\footnote{安人且猶不能，況其有能至知天命者乎？○矧，失忍反。}已！予惟小子，若涉淵水，予惟往求朕攸濟。\footnote{已，發端嘆辭也。我惟小子，承先人之業，若涉淵水，往求我所以濟渡。言祗懼。}


{\noindent\zhuan\zihao{6}\fzbyks 傳“\CJKunderline{周公}”至“及之”。正義曰：序雲“相\CJKunderline{成王}”，則“王若曰”者，稱\CJKunderline{成王}之言，故言“\CJKunderline{周公}稱\CJKunderline{成王}命”。實非王意，\CJKunderline{成王}爾時信流言，疑\CJKunderline{周公}，豈命公伐管蔡乎？“猷”訓道也，故云“順大道以告天下眾國”也。鄭、王本“猷”在“誥”下。\CJKunderwave{漢書}王莽攝位，東郡太守翟義叛莽，莽依此作\CJKunderwave{大誥}，其書亦“道”在“誥”下。此本“猷”在“大”上，言以道誥眾國，於文為便。但此經雲“猷”，\CJKunderwave{大傳}雲“大道”，古人之語多倒，猶\CJKunderwave{詩}稱“中谷”，谷中也。“多邦”之下雲於爾御事，是於諸國治事者盡及之也。\CJKunderline{鄭玄}云：“王，\CJKunderline{周公}也，\CJKunderline{周公}居攝，命大事，則權稱王。”惟名與器不可假人，\CJKunderline{周公}自稱為王，則是不為臣矣，大聖作則，豈為是乎？ \par}

{\noindent\zhuan\zihao{6}\fzbyks 傳“兇害”至“之意”。正義曰：\CJKunderwave{釋詁}云：“延，長也。洪，大也。”此害長大，敗亂國家，經言惟我幼童人，謂損累之，故傳加“累”字，累我童人，言其不可不誅之意。鄭、王皆以“延”上屬為句，言害不少，乃延長之。王肅又以“惟”為念，向下為義，大念我幼童子與繼文武無窮之道。 \par}

{\noindent\zhuan\zihao{6}\fzbyks 傳“言子”至“自責”。正義曰：“嗣”訓繼也。言子孫承繼祖疆，境界則是無窮，大數長遠，“卜世三十，卜年七百”，是長遠也。 \par}

{\noindent\zhuan\zihao{6}\fzbyks 傳“安人”至“者乎”。正義曰：民近而天遠，以易而況難。天子必當至靈，至靈乃知天命，言己猶不能安民，明其不知天命。自責而謙。 \par}

敷賁敷前人受命，茲不忘大功。\footnote{前人，文武也。我求濟渡，在布行大道，在布陳文武受命，在此不忘大功。言任重。○賁，扶雲反，徐音憤。}予不敢閉於天降威用。\footnote{天下威用，謂誅惡也。言我不敢閉絕天所下威用而不行，將欲伐四國。}寧王遺我大寶龜，紹天明，即命。\footnote{安天下之王，謂\CJKunderline{文王}也。遺我大寶龜，疑則卜之，以繼天明，就其命而行之。言卜不可違。○遺，唯季反。}

{\noindent\zhuan\zihao{6}\fzbyks 傳“前人”至“任重”。正義曰：\CJKunderline{成王}前人，故為“文武”也。以涉水為喻，言求濟者，在於布行大道，行天子之政也。文武有大功德，故受天命，又當布陳文武受命所行之事也。陳行天子之政，又陳文武所行之事。在此不忘大功。“大功”，大平之功也。言己所任至重，不得不奉天道行誅伐也。 \par}

{\noindent\zhuan\zihao{6}\fzbyks 傳“天下”至“四國”。正義曰：王者征伐刑獄，象天震曜殺戮，則征伐者，天之所威用，謂誅惡是也。天有此道，王者用之。用之則開，不用則閉，言我不敢閉絕天之所下威用而不行之。既不敢不行，故將伐四國。 \par}

{\noindent\zhuan\zihao{6}\fzbyks 傳“安天”至“可違”。正義曰：紂為昏虐，天下不安，言\CJKunderline{文王}能安之，安天下之王謂\CJKunderline{文王}也。“遺我大寶龜”者，天子寶藏神龜，疑則卜之。繼天明道，就其命而行之，言卜吉則當行，不可違卜也。所以大寶龜皆得繼天明者，以天道玄遠，龜是神靈，能傳天意以示吉凶，故疑則卜之，以繼天明道。\CJKunderline{鄭玄}云：“時既卜，乃後出誥，故先云然。” \par}

{\noindent\shu\zihao{5}\fzkt “王若”至“即命”。正義曰：\CJKunderline{周公}雖攝王政，其號令大事則假\CJKunderline{成王}為辭。言王順大道而為言曰，我今以大道誥汝天下眾國，及於眾治事之臣。以我周道不至，故上天下其兇害於我家不少。言叛逆者多。此害延長寬大，惟累我幼童人。\CJKunderline{成王}自言害及己也。我之致此兇害，以我為子孫，承繼無疆界之大數，服行其政，不能為智道令民安，故使之叛。自責也。安民猶且不能，況曰其能至於知天之大命者乎？言己不能知天意也。復嘆而言，已乎！我惟小子，承先人之業，如涉淵水，惟往求我所以濟渡。言己恐懼之甚。我所求濟者，惟在布行大道，布陳前人\CJKunderline{文王}\CJKunderline{武王}受命之事，在我此身，不忘大功。既不忘大功，當誅叛逆，由此我不敢絕天之所下威用而不行之。言必將伐四國也。寧天下之王，謂\CJKunderline{文王}也。\CJKunderline{文王}遺我大寶龜，疑則就而卜之，以繼天明命，今我就受其命。言己就龜卜其伐之吉凶，已得吉也。 \par}

曰：‘有大艱於西土，西土人亦不靜，越茲蠢。’\footnote{曰，語更端也。四國作大難於京師，西土人亦不安，於此蠢動。○蠢，尺允反。難，乃旦反，下同，又如字。}殷小腆，誕敢紀其敘。\footnote{言殷後小腆腆之祿父，大敢紀其王業，欲復之。○腆，他典反，馬云：“至也。”誕，大旦反。父音甫，後同。}天降威，知我國有疵，\footnote{天下威，謂三叔流言，故祿父知我周國有疵病。○疵，在斯反，馬云：“叚也。”}民不康。


{\noindent\zhuan\zihao{6}\fzbyks 傳“曰語”至“蠢動”。正義曰：\CJKunderline{周公}丁寧其事，止而復言，別加一“曰”，語更端也。下言“王曰”，此不言“王”，史詳略耳。四國作逆於東，京師以為大艱，故言“作大難於京師”。“西土人亦不安”，亦如東方見其亂，不安也。\CJKunderwave{釋詁}云：“蠢，動也。”鄭云：“周民亦不定，其心騷動，言以兵應之。”當時京師無與應者，鄭言妄耳。 \par}

{\noindent\zhuan\zihao{6}\fzbyks 傳“言殷”至“復之”。正義曰：殷本天子之國，\CJKunderline{武庚}比之為小，故言“小腆”，“腆”是小貌也。\CJKunderline{鄭玄}云：“腆謂小國也。”王肅云：“腆，主也，殷小主謂祿父也。”“大敢紀其王業”，經紀王業，望復之也。 \par}

{\noindent\zhuan\zihao{6}\fzbyks 傳“天下”至“疵病”。正義曰：王肅云：“天降威者，謂三叔流言，當誅伐之。”言誅三叔是天下威也。\CJKunderwave{釋詁}云：“疵，病也。”鄭、王皆云：“知我國有疵病之瑕。” \par}

曰：‘予復。’反鄙我周邦。\footnote{祿父言我殷當復，欺惑東國人，令不安，反鄙易我周家。道其罪無狀。○令，力呈反。易，以豉反，下“其易”同。}今蠢，今翼日，民獻有十夫，予翼以於\xpinyin*{敉}\CJKunderline{寧}、\CJKunderline{武}圖功。\footnote{今天下蠢動，今之明日，四國人賢者有十夫來翼佐我周，用撫安武事，謀立其功。言人事先應。○敉，亡婢反。應，應對之應。}我有大事，休？朕卜並吉。\footnote{大事，戎事也。人謀既從，卜又並吉，所以為美。○並，必政反，注及篇末同。}

{\noindent\zhuan\zihao{6}\fzbyks 傳“祿父”至“無狀”。正義曰：祿父以父罪，滅殷身亦當死，幸得繼其先祀，宜荷天恩。反鄙薄輕易我周家，言其不識恩養，道其罪無狀也。漢代止有“無狀”之語，蓋言其罪大無可形狀也。近代已來遭重喪答人書云：“無狀招禍”，是古人之遺語也。 \par}

{\noindent\zhuan\zihao{6}\fzbyks 傳“今天”至“先應”。正義曰：\CJKunderline{武庚}既叛，聞者皆驚，故“今天下蠢動”，謂聞叛之日也。“今之明日”，聞叛之明日。以“獻”為賢，四國民內賢者十夫，來翼佐我周。十人史無姓名，直是在彼逆地，有先見之明，知彼必敗，棄而歸周。\CJKunderline{周公}喜其來降，舉以告眾，謂之為賢，未必是大賢也。“用撫安武事，謀立其功”，用此十夫為之。將欲伐叛，而賢者即來，言人事先應也。 \par}

{\noindent\zhuan\zihao{6}\fzbyks 傳“大事”至“為美”。正義曰：成十三年\CJKunderwave{左傳}云：“國之大事在祀與戎。”今論伐叛，知“大事，戎事也”。十夫來翼，人謀既從，卜又並吉，所以為美，美即經之“休”也。既言其休，乃說我卜，並言以成此休之意。\CJKunderline{鄭玄}云：“卜並吉者，謂三龜皆從也。”王肅云：“何以言美？以三龜一習吉，是言並吉，證其休也。”與孔異矣。 \par}

{\noindent\shu\zihao{5}\fzkt “曰有”至“並吉”。正義曰：上言為害不少，陳欲徵之意，未說\CJKunderline{武庚}之罪。更復發端言之曰，今四國叛逆，有大艱於西土。言作亂於東，與京師為難也。西土之人為此亦不得安靜，於此人情皆蠢蠢然動。殷後小國腆腆然之祿父，大敢紀其王業之次敘，而欲興復之。祿父所以敢然者，上天下威於三叔，以其流言欲下威誅之，祿父知我周國有此疵病，而欺惑東國人，令人不安。祿父謂人曰：“我殷復。”望得更為天子，反鄙易我周國。今天下蠢動，今之明日，四國民之賢者有十夫，不從叛逆，其來為我翼佐我周。於是用撫安武事，謀立其功。明祿父舉事不當，得賢者叛來投我，為我謀用。是人事先應如此，則我有兵戎大事，征伐必休美矣。人謀既從，我卜又並吉，是其休也。言往必克敵安民之意，告眾使知也。 \par}

“肆予告我友邦君，越尹氏、庶士御事，\footnote{以美，故告我友國諸侯，及於正官尹氏卿大夫、眾士御治事者。言謀及之。}曰：‘予得吉卜，予惟以爾庶邦，於伐殷逋播臣。’\footnote{用汝眾國，往伐殷逋亡之臣。謂祿父。○逋，布吾反。}爾庶邦君，越庶士御事，罔不反曰：‘艱大。’\footnote{汝眾國上下無不反曰：“征伐四國為大難。”敘其情以戒之。}


{\noindent\zhuan\zihao{6}\fzbyks 傳“以美”至“及之”。正義曰：“肆”訓故也，承上“休”之下，以其東征必美之故，我告友國君以下共謀之。“尹氏”，即\CJKunderwave{顧命}雲“百尹”氏也。“尹”，正也，諸官之正，謂卿大夫，故傳言“及於正官尹氏卿大夫”。“尹氏”即官也，總呼大夫為官氏也。上文“大誥爾多邦”，越爾御事”，無“尹氏庶士”，下之“爾庶邦君，越庶士御事”亦無“尹氏”，惟此及下文施義二者詳其文，餘略之，從可知也。 \par}

{\noindent\zhuan\zihao{6}\fzbyks 傳“用汝”至“祿父”。正義曰：“逋”，逃也。“播”謂播蕩逃亡之意。祿父殷君，謂之為“殷”。今日叛逆，是背周逃亡，故云用汝眾國，往伐彼殷君於我周家逋逃亡叛之臣。謂祿父也。 \par}

{\noindent\zhuan\zihao{6}\fzbyks 傳“汝眾”至“戒之”。正義曰：王以卜吉之故，將以諸國伐殷，且彼諸國之情，必有不欲伐者，無不反我之意，相與言曰：“征伐四國為大難。”言其情必如此，敘其情以戒之，使勿然也。鄭云：“汝國君及下群臣不與我同志者，無不反我之意，云：‘三監叛，其為難大。’”是言“反”者謂反上意，反是上意，則知“曰”者，相與言也。 \par}

民不靜，亦惟在王宮邦君室。\footnote{言四國不安，亦在天子諸侯教化之過。自責不能綏近以及遠。}越予小子，考翼不可徵，王害不違卜。\footnote{於我小子先卜敬成周道，若謂今四國不可徵，則王室有害，故宜從卜。}

{\noindent\zhuan\zihao{6}\fzbyks 傳“言四”至“及遠”。正義曰：“自責”惟當言天子教化之過，而並言諸侯者，化從天子佈於諸侯，道之不行，亦邦君之咎，見庶邦亦有過，故並言之。教化之過在於君身，而云“王宮邦君室”者，宮室是行化之處，故指以言之。 \par}

{\noindent\zhuan\zihao{6}\fzbyks 傳“於我”至“從卜”。正義曰：“翼”訓敬也，於我小子，先自考卜，欲敬成周道。汝庶邦御事等，若謂今四國不可徵，則周道不成，於王室有害，故宜從卜。“小子先卜”當謂初即位時，卜其欲成周道也。不可違卜，謂上“朕卜並吉”也。言欲徵卜吉，當從卜征之。 \par}

{\noindent\shu\zihao{5}\fzkt “肆予告”至“違卜”。正義曰：以人從卜吉為美之故，故我告汝有邦國之君，及於尹氏卿大夫、眾士治事者曰，我得吉卜，我惟與汝眾國往伐殷逋亡播蕩之臣。謂伐祿父也。汝國君及於眾治事者，無不反我之意，相與言曰：“伐此四國，為難甚大。”言其不欲徵也。汝不欲伐，罪我之由四國之民不安而叛者，亦惟在我天子王宮與邦君之室教化之過使之然。以此令汝難徵，過事在我。雖然，於我小子，先考疑而卜之，欲敬成周道，若謂四國難大不可徵，則於王室有害，不可違卜，宜從卜往徵也。 \par}

肆予沖人永思艱，曰，嗚呼！允蠢鰥寡，哀哉！\footnote{故我童人\CJKunderline{成王}長思此難而嘆曰：“信蠢動天下，使無妻無夫者受其害，可哀哉！”○鰥，故頑反。}予造天役，遺大投艱於朕身。\footnote{我周家為天下役事，遺我甚大，投此艱難於我身。言不得已。○予造，為也。馬云：“遺也。”}越予沖人，不卬自恤。義爾邦君，越爾多士、尹氏御事，\footnote{言徵四國，於我童人不惟自憂而已，乃欲施義於汝胤國君臣上下至御治事者。○卬，五剛反，我也。}綏予曰：‘無毖於恤。不可不成乃\CJKunderline{寧考}圖功。’\footnote{汝眾國君臣，當安勉我曰：“無勞於憂，不可不成汝寧祖聖考文武所謀之功。”責其以善言之助。○毖音秘。}


{\noindent\zhuan\zihao{6}\fzbyks 傳“我周”至“得已”。正義曰：為天子者，當役己以養天下，故“我國家為天下役事”，總言周家當救天下。此事遺我，故為甚大。以大役遺我，以為甚大，而又投擲此艱難之事於我身，謂當已之時有四國叛逆，言已職當靜亂，不得以己也。 \par}

{\noindent\zhuan\zihao{6}\fzbyks 傳“言徵”至“事者”。正義曰：卬，我。恤，憂也。四國叛逆，害及眾國，君得靜亂，則為大美。言徵四國，於我童人，不惟自憂而已。乃欲施義於汝眾國君臣，言難除則義施也。 \par}

{\noindent\zhuan\zihao{6}\fzbyks 傳“汝眾”至“之助”。正義曰：綏，安也。毖，勞也。言我既施義於汝，汝眾國君臣言得我之力，當安慰勉勸我曰：“無勞於憂。”令我無憂四國，眾國自來徵之。經言“寧”即\CJKunderline{文王}，“考”即\CJKunderline{武王}，故言“寧祖聖考”也。王以眾國反己，乃復設為此言，責其無善言助己。 \par}

{\noindent\shu\zihao{5}\fzkt “肆予衝”至“圖功”。正義曰：以汝等有難徵之意，故我童子\CJKunderline{成王}長思此難而嘆曰：“嗚呼！四國今叛，信蠢動天下，使鰥寡受害，尤可哀哉！我周家為天下役事，而遺我甚大，乃投此艱難於我身。此難須平，不可以已。今徵四國，於我童人不惟自憂而已，乃欲施義於汝眾國君，於汝多士尹氏治事之人。如此為汝計，汝君臣當安勉我曰：‘無勞於徵伐之憂，我諸侯當往共徵四國。汝王不可不成汝寧祖聖考所謀之功。’宜出此善言以助我。何謂違我不欲徵也？” \par}

已！予惟小子，不敢替上帝命。\footnote{不敢廢天命，言卜吉當必徵之。}天休於寧王，興我小邦周，寧王惟卜用，克綏受茲命。\footnote{言天美\CJKunderline{文王}興周者，以\CJKunderline{文王}惟卜之用，故能安受此天命。明卜宜用。}今天其相民，矧亦惟卜用。\footnote{人獻十夫，是天助民，況亦用卜乎？吉可知矣。亦亦\CJKunderline{文王}。○相，息亮反。}嗚呼！天明畏，弼我丕丕基。”\footnote{嘆天之明德可畏，輔成我大大之基業。言卜不可違也。○畏如字，徐音威。}王曰：“爾惟舊人，爾丕克遠省，爾知寧王若勤哉！\footnote{特命久老之人，知\CJKunderline{文王}故事者，大能遠省識古事，汝知\CJKunderline{文王}若彼之勤勞哉！目所親見，法之又明。○省，息井反。}天閟毖我成功所，予不敢不極卒寧王圖事。\footnote{閟，慎也。言天慎勞我周家成功所在，我不敢不極盡\CJKunderline{文王}所謀之事。謂致太平。○閟音秘。}肆予大化誘我友邦君，\footnote{我欲極盡\CJKunderline{文王}所謀，故大化天下，道我友國諸侯。}



{\noindent\zhuan\zihao{6}\fzbyks 傳“人獻”至“\CJKunderline{文王}”。正義曰：天之助民，乃是常道，而云“民獻十夫，是天助民”者，下雲“亦惟十人，迪知上帝命”，故以民獻十夫為天助民也。“王曰爾”至“休畢”。正義曰：既述\CJKunderline{文王}之事，王又命於眾曰，汝惟久老之人，汝大能遠省識古事，汝知寧王若此之勤勞哉！以老人目所親見，必知之也。以\CJKunderline{文王}勤勞如此，故天命慎勞來我周家，當至成功所在。天意既然，我不敢不極盡\CJKunderline{文王}所謀之事。\CJKunderline{文王}謀致太平，我欲盡行之。我欲盡\CJKunderline{文王}所謀，故我大為教化，勸誘我所友國君，共伐叛逆。天既輔助我周家有大化誠辭，其必成就我之眾民。天意既如此矣，我何其不於前\CJKunderline{文王}安民之道、謀立其功之處所而終竟之乎？天亦惟勞慎我民，若人有疾病，而欲已去之。天意於民如此之急，我何敢不於前安人\CJKunderline{文王}所受美命終畢之乎？以須終畢之故，故當誅除逆亂，安養下民，使之致太平。 \par}

{\noindent\zhuan\zihao{6}\fzbyks 傳“閟慎”至“太平”。正義曰：“閟，慎”，\CJKunderwave{釋詁}文。“天慎勞我周家”者，美其德當天心，慎惜又勞來勸勉之，使至成功所在，在於致太平也。天意欲使之然，我為\CJKunderline{文王}子孫，敢不極盡\CJKunderline{文王}所謀之事？\CJKunderline{文王}本謀，謂致太平。 \par}

天棐忱辭，其考我民，\footnote{言我周家有大化誠辭，為天所輔，其成我民矣。○棐，徐音匪，又芳鬼反。忱，市林反。}予曷其不於前寧人圖功攸終？\footnote{我何其不於前\CJKunderline{文王}安人之道、謀立其功所終乎？}天亦惟用勤毖我民，若有疾，\footnote{天亦勞慎我民欲安之，如人有疾，欲已去之。}予曷敢不於前寧人攸受休畢？”\footnote{天欲安民，我何敢不於前\CJKunderline{文王}所受美命終畢之？}

{\noindent\zhuan\zihao{6}\fzbyks 傳“言我”至“民矣”。正義曰：\CJKunderwave{釋詁}云：“棐，輔也。忱，誠也。”文承“大化”之下，知輔誠辭者，“言周家有大化誠辭，為天所輔”。“其成我民”，必為民除害，使得成也。 \par}

{\noindent\zhuan\zihao{6}\fzbyks 傳“天亦”至“去之”。正義曰：“亦”者，亦民之義也。君民共為一體，天慎勞使成功，亦當勤勞民使安寧，故言“亦”也。如疾，欲已去之，言天急於民至甚也。 \par}

{\noindent\zhuan\zihao{6}\fzbyks 傳“天欲”至“畢之”。正義曰：上雲“卒寧王圖事”，又云“圖功攸終”，此雲“攸受休畢”，“畢”，終也，三者文辭略同，義不甚異。天意惟言當終\CJKunderline{文王}之業，須徵逆亂之賊，\CJKunderline{周公}重兵慎戰，丁寧以勸民耳。 \par}

{\noindent\shu\zihao{5}\fzkt “已予”至“丕基”。正義曰：既敘眾國之情，告以必徵之意：“已乎！我惟小子，不敢廢上帝之命。卜吉不徵，是廢天命。從卜而興，乃有故事。天休美於安天下之\CJKunderline{文王}興我小國周者，以安民之王，惟卜是用，以此之故，安受此上天之命。明卜宜用之。今天助民矣，十夫佐周，是天助也。人事既驗，況亦如\CJKunderline{文王}惟卜之用，吉可知矣。嗚呼！而嘆天之明德可畏也，輔成我周家大大之基業。卜既得吉，不可違也。” \par}

王曰：“若昔朕其逝，朕言艱日思。\footnote{順古道，我其往東征矣。我所言國家之難備矣，日思念之。○日，人實反。難，乃旦反，下“為難”同。}若考作室，既厎法，厥子乃弗肯堂，矧肯構？\footnote{以作室喻治政也。父已致法，子乃不肯為堂基，況肯構立屋乎？不為其易，則難者可知。○厎，之履反。構，古候反。治，直吏反。}厥父菑，厥子乃弗肯播，矧肯獲？\footnote{又以農喻。其父已菑耕其田，其子乃不肯播種，況肯收穫乎？○菑，側其反，草也，田一歲曰菑。獲，戶郭反。}厥考翼，其肯曰：‘予有後，弗棄基？’\footnote{其父敬事創業，而子不能繼成其功，其肯言我有後，不棄我基業乎？今不正，是棄之。}肆予曷敢不越卬敉寧王大命？\footnote{作室農人，猶惡棄基，故我何敢不於今日撫循\CJKunderline{文王}大命以徵逆乎？○惡，烏路反。}若兄考，乃有友伐厥子，民養其勸弗救。”\footnote{若兄弟父子之家，乃有朋友來伐其子，民養其勸不救者，以子惡故。以此四國將誅而無救者，罪大故。}


{\noindent\zhuan\zihao{6}\fzbyks 傳“又以”至“獲乎”。正義曰：上言作室，此言治田，其取喻一也。上言“若考作室，既厎法”，此類上文，當雲“若父為農，既耕田”，從上省文耳。“菑”謂殺草，故治田一歲曰菑，言其始殺草也。“播”謂布種，后稷播殖百穀是也。定本雲“矧弗肯構”、“矧弗肯獲”，皆有“弗”字，檢孔傳所解，“弗”為衍字。 \par}

{\noindent\zhuan\zihao{6}\fzbyks 傳“其父”至“棄之”。正義曰：治田作室，為喻既同，故以此經結上二事。鄭、王本於“矧肯構”下亦有此一經，然取喻既同，不應重出。蓋先儒見下有而上無，謂其脫而妄增之。 \par}

{\noindent\zhuan\zihao{6}\fzbyks 傳“若兄”至“大故”。正義曰：此經大意，言兄不救弟，父不救子，發首“兄考”備文，“伐厥子”，不言“弟”，互相發見，傳言“兄弟父子之家”以足之。“民養其勸”，“民”謂父兄，為家長者也，養其心不退止也。 \par}

{\noindent\shu\zihao{5}\fzkt “王曰若”至“弗救”。正義曰：子孫成父祖之業，古道當然。王又言曰：“今順古昔之道，我其往東征矣。我所言國家之難備矣，日日思念之。乃以作室為喻，若父作室，營建基趾，既致法矣，其子乃不肯為之堂，況肯構架成之乎？又以治田為喻，其父菑耕其田，殺其草，已堪下種矣，其子乃不肯布種，況肯收穫乎？其此作室治田之父，乃是敬事之人，見其子如此，其肯言曰：‘我有後，不棄我基業乎？’必不肯為此言也。我若不終文武之謀，則文武之神亦如此耳，其肯道我不棄基業乎？作室農人猶惡棄其基業，故我何敢不於我身今日撫循安人之\CJKunderline{文王}大命，以征討叛逆乎？我今東征，無往不克，凡人兄及父與子弟為家長者，乃有朋友來伐其子，則民皆養其勸伐之心不救之。何則？以子惡故也。以逾伐四國，雖親如父兄，亦無救之者，以君惡故也。”言罪大不可不誅，無救所以必克也。顧氏以上“不卬自恤”傳雲“不惟自憂”，遂皆以“卬”為惟。但“卬”之為惟，非是正訓，觀孔意亦以不“卬”為惟義也。 \par}

王曰：“嗚呼！肆哉！爾庶邦君，越爾御事。\footnote{嘆今伐四國必克之故，以告諸侯及臣下御治事者。}爽邦由哲，亦惟十人，迪知上帝命。\footnote{言其故，有明國事、用智道十人蹈知天命。謂人獻十夫來佐周。}越天棐忱，爾時罔敢易法，矧今天降戾於周邦？\footnote{於天輔誠，汝天下是知無敢易天法，況今天下罪於周，使四國叛乎？}惟大艱人，誕鄰胥伐於厥室，爾亦不知天命不易？\footnote{惟大為難之人，謂三叔也。大近相伐於其室家，謂叛逆也。若不早誅汝，天下亦不知天命之不易也。○易，以豉反。}


{\noindent\zhuan\zihao{6}\fzbyks 傳“言其”至“佐周”。正義曰：此其必克之故也。“爽”，明也。“由”，用也。“有明國事、用智道”，言其有賢德也。蹈天者，識天命而履行之。此言“十人”，謂上文民獻十夫來佐周家者。此是賢人，賢人既來，彼無所與，是必克之效也。王肅云：“我未伐而知民弗救者，以民十夫用知天命故也。” \par}

{\noindent\zhuan\zihao{6}\fzbyks 傳“於天”至“叛乎”。正義曰：“於天輔誠”，言天之所輔，必是誠信。汝天下於是觀之，始知無敢變易天法。若易天法，則天不輔之，況今天下罪於周，使四國叛乎？以小況大，易法猶尚不可，況叛逆乎？ \par}

{\noindent\zhuan\zihao{6}\fzbyks 傳“惟大”至“不易”。正義曰：以下句言相伐於其室家，室家自相伐，知“惟大為難之人，謂三叔也”。“大近相伐於其室家”者，三叔為周室至親，而舉兵作亂，是室家自相伐。為叛逆之罪，是變易天法之極，若汝諸國不肯誅之，是汝天下亦不知天命之不可變易也。王肅云：“惟大為難之人，謂管蔡也。大近相伐於其室家，明不可不誅也。管蔡犯天誅而汝不欲伐，則亦不知天命之不易也。” \par}

{\noindent\shu\zihao{5}\fzkt “王曰鳴”至“不易”。正義曰：既言四國無救之者，王曰，又言嘆今伐四國必克之故，告汝眾國君，及於汝治事之臣。所以知必克者，故有明國事、用智道者，亦惟有十人，匆人皆蹈知上天之命。謂民獻十夫來佐周家，此人既來，克之必也。於我天輔誠信之故，汝天下是知無敢變易天法者，若易法無信，則上天不輔，故無敢易法也。況今天下罪於周國，使四國叛逆。惟大為難之人，謂三叔等，大近相伐於其室家，自欲拔本塞源，反害周室，是其為易天法也。彼變易天法，若不早誅之，汝天下亦不知天命之不可變易也。 \par}

予永念曰，天惟喪殷，若穡夫，予曷敢不終朕畝？\footnote{稼穡之夫，除草養苗。我長念天亡殷惡主，亦猶是矣。我何敢不順天，終竟我壟畝乎？言當滅殷。○壟，力勇反。}天亦惟休於前寧人，予曷其極卜，敢弗於從？\footnote{天亦惟美於\CJKunderline{文王}受命，我何其極卜法，敢不於從？言必從也。}率寧人有指疆土，矧今卜並吉？\footnote{循\CJKunderline{文王}所有指意以安疆土則善矣，況今卜並吉乎？言不可不從。}肆朕誕以爾東征。天命不\xpinyin{僣}{jian4}，卜陳惟若茲。”\footnote{以卜吉之故，大以汝眾東征四國。天命不僣差，卜兆陳列惟若此吉，必克之，不可不勉。}


{\noindent\zhuan\zihao{6}\fzbyks 傳“天亦”至“必從”。正義曰：“天亦惟美於\CJKunderline{文王}受命”，言\CJKunderline{文王}德當天心，天每事美之，故得受天命，是\CJKunderline{文王}之德大美也。\CJKunderline{文王}用卜，能受天命，今於我何其窮極\CJKunderline{文王}卜法，敢不從乎？言必從\CJKunderline{文王}卜也。 \par}

{\noindent\zhuan\zihao{6}\fzbyks 傳“循文”至“不從”。正義曰：\CJKunderline{文王}之旨意，欲今天下疆土皆得其宜。有叛逆者，自然須平定之。我直循彼\CJKunderline{文王}所有旨意伐叛，則已善矣，不必須卜筮也，況今卜並吉乎？言不可不從也。王肅云：“順\CJKunderline{文王}安人之道，有旨意盡天下疆土使皆得其所，不必須卜筮也，況今卜三龜皆吉，明不可不從也。” \par}

{\noindent\zhuan\zihao{6}\fzbyks 傳“以卜”至“不勉”。正義曰：“天命不僣”，天意去惡與善，其事必不僣差，言我善而彼惡也。“卜兆陳列惟若此吉”，言往必克之，不可不勉力也。 \par}

{\noindent\shu\zihao{5}\fzkt “予永”至“若茲”。正義曰：所以必當誅四國者，我長思念之曰，天惟喪亡殷國者，若稼穡之夫，務去草也，天意既然，我何敢不終我壟畝也？言穢草盡須除去，殷餘皆當殄滅也。天亦惟美於前寧人\CJKunderline{文王}，我何其極\CJKunderline{文王}卜法，敢不於是從乎？言必從之也。我循彼寧人所有旨意以安疆上，不待卜筮，便即東征，已自善矣，況今卜東征而龜並吉？以吉之故，我大以爾東征四國。天命必不僣差，卜兆陳列惟若此吉，不可不從卜，不可不勉力也。 \par}

\section{微子之命第十【偽】}


\CJKunderline{成王}既黜殷命，殺\CJKunderline{武庚}，\footnote{一名祿父。}命\CJKunderline{微子啟}代殷後，\footnote{啟知紂必亡而奔周，命為宋公，為湯後。}作\CJKunderwave{微子之命}。\footnote{封命之書。}


{\noindent\zhuan\zihao{6}\fzbyks 傳“啟知”至“湯後”。正義曰：啟知紂必亡，告父師少師而遁於荒野，“\CJKunderline{微子}作誥”是其事也。\CJKunderline{武王}既克紂，\CJKunderline{微子}乃歸之，非去紂即奔周也。傳言得封之由，故言其“奔周”耳。僖六年\CJKunderwave{左傳}云，許僖公見楚子，“面縛銜璧，大夫衰絰，士輿櫬。楚子問諸逢伯，對曰：‘昔\CJKunderline{武王}克殷，\CJKunderline{微子}啟如是。\CJKunderline{武王}親釋其縛，受其璧而祓之。焚其櫬，禮而命之，使復其所。’”\CJKunderwave{史記·宋世家}云：“\CJKunderline{武王}克殷，\CJKunderline{微子}啟乃持其祭器造于軍門，肉袒面縛，左牽羊，右把茅，膝行而前以告。\CJKunderline{武王}乃釋\CJKunderline{微子}，復其位如故。”是言\CJKunderline{微子}克殷始歸周也。馬遷之書，辭多錯謬，“面縛”縛手於後，故口銜其璧，又安得“左牽羊，右把茅”也？要言歸周之事是其實耳。\CJKunderwave{樂記}云，\CJKunderline{武王}克殷，既下車，投殷之後於宋。則傳言復其位者，以其自縛為囚，釋之使從本爵，復其卿大夫之位。及下車即封於宋，以其終為殷後，故\CJKunderwave{樂記}雲“投殷之後”，爾時未為殷之後也。\CJKunderline{微子}初封於宋，不知何爵，此時因舊宋命之為公，令為湯後，使祀湯耳，不繼紂也。 \par}

{\noindent\shu\zihao{5}\fzkt “\CJKunderline{成王}”至“之命”。正義曰：\CJKunderline{成王}既黜殷君之命，殺\CJKunderline{武庚}，乃命\CJKunderline{微子}啟代\CJKunderline{武庚}為殷後，為書命之。史敘其事，作\CJKunderwave{微子之命}。“黜殷命”，謂絕其爵也。“殺\CJKunderline{武庚}”，謂誅其身也。 \par}

微子之命\footnote{稱其本爵以名篇。}

{\noindent\shu\zihao{5}\fzkt “\CJKunderline{微子}之命”。正義曰：令寫命書之辭以為此篇，\CJKunderwave{君陳}、\CJKunderwave{君牙}、\CJKunderwave{冏命}皆此類也。 \par}

王若曰：“猷！殷王元子，\footnote{微子，\CJKunderline{帝乙}元子，故順道本而稱之。}惟稽古，崇德象賢。\footnote{惟考古典，有尊德象賢之義。言今法之。}統承先王，修其禮物，\footnote{言二王之後，各修其典禮，正朔服色，與時王並通三統。○正音政。}作賓於王家，與國咸休，永世無窮。\footnote{為時王賓客，與時皆美，長世無竟。}嗚呼！乃祖\CJKunderline{成湯}，克齊聖廣淵，\footnote{言汝祖\CJKunderline{成湯}能齊德聖達廣大深遠，澤流後世。}皇天眷佑，誕受厥命。\footnote{大天眷顧湯，佑助之，大受其命。謂天命。}撫民以寬，除其邪虐，\footnote{撫民以寬政，放桀邪虐。湯之德。}功加於時，德垂後裔。\footnote{言湯立功加於當時，德澤垂及後世。裔，末也。}爾惟踐修厥猷，舊有令聞，\footnote{汝\CJKunderline{微子}，言能踐湯德，久有善譽，昭聞遠近。○令聞，如字，又音問。}恪慎克孝，肅恭神人。予嘉乃德，曰篤不忘。\footnote{言\CJKunderline{微子}敬慎能孝，嚴恭神人，故我善汝德，謂厚不可忘。○篤，本又作竺，東谷反。}


{\noindent\zhuan\zihao{6}\fzbyks 傳“\CJKunderline{微子}”至“稱之”。正義曰：\CJKunderwave{呂氏春秋·仲冬紀}云：“紂之母生\CJKunderline{微子}啟與仲衍，尚為妾，已而為妻後生紂。紂父欲立啟為太子，太史據法而爭之曰：‘有妻之子，又不可立妾之子。’故紂為後。”鄭云：“\CJKunderline{微子}啟，紂同母庶兄也。若，順也。猷，道也。以其本是元子，故順道本而稱之。”\CJKunderwave{釋詁}云：“元、首，始也。”\CJKunderwave{易}曰：“元者，善之長也。” \par}

{\noindent\zhuan\zihao{6}\fzbyks 傳“言二”至“三統”。正義曰：\CJKunderwave{郊特牲}云：“天子存二代之後，猶尊賢也。尊賢不過二代。”\CJKunderwave{書傳}云：“王者存二王之後，與己為三，所以通三統，立三正。周人以日至為正，殷人以日至後三十日為正，夏人以日至後六十日為正。天有三統，土有三王，三王者，所以統天下也。”\CJKunderwave{禮運}云：“杞之郊也，\CJKunderline{禹}也。宋之郊也，契也。”是二王后為郊祭天，以其祖配之。鄭云：“所存二王后者，命使郊天，以天子禮祭其始祖受命之王，自行其正朔服色，此謂通天三統，是立二王后之義也。”此命首言“稽古”，則立先代之後，自古而有此法，不知從何代然也。孔意自夏以上不必改正，縱使正朔不改，典禮服色自當異也。“曰篤不忘”。正義曰：僖十二年\CJKunderwave{左傳}王命管仲之辭曰“謂督不忘”，則“曰”亦謂義。孔訓“篤”為厚，故傳雲“謂厚不可忘”。杜預以“督”為正，可謂正而不可忘也。 \par}

{\noindent\shu\zihao{5}\fzkt “王若曰猷殷王元子”。正義曰：王順道而言曰：“今以大道告汝殷王首子。”告之以下辭也。“曰猷”，如\CJKunderwave{大誥}，言以道誥之。 \par}

上帝時歆，下民祗協，庸建爾於上公，尹茲東夏。\footnote{孝恭之人，祭祀則神歆享，施令則人敬和，用是封立汝於上公之位，正此東方華夏之國。宋在京師東。○歆，許今反。}欽哉！往敷乃訓，慎乃服命，率由典常，以蕃王室。\footnote{敬哉，敬其為君之德。往臨人布汝教訓，慎汝祖服命數，循用舊典，無失其常，以蕃屏周室。戒之。}

{\noindent\shu\zihao{5}\fzkt “慎乃服命”。正義曰：傳言“慎汝祖服命數”，謂祭湯廟得用天子之禮，服其殷之本服，命則上公九命，當慎之，無使乖禮制也。 \par}

弘乃烈祖，律乃有民，永綏厥位，毗予一人。\footnote{大汝烈祖\CJKunderline{成湯}之道，以法度齊汝所有之人，則長安其位，以輔我一人。言上下同榮慶。○毗，房脂反。}世世享德，萬邦作式，\footnote{言\CJKunderline{微子}累世享德，不忝厥祖，雖同公侯，而特為萬國法式。}俾我有周無斁。\footnote{汝世世享德，則使我有周好汝無厭。○俾，必爾反。斁音亦。好，呼報反。厭，於豔反。}嗚呼！往哉惟休，無替朕命。”\footnote{嘆其德，遣往之國。言當惟為美政，無廢我命。}

\CJKunderline{唐叔}得禾，異畝同穎，\footnote{\CJKunderline{唐叔}，\CJKunderline{成王}母弟。食邑內得異禾也。畝，壟。穎，穗也。禾各生一壟而合為一穗。○穎，役領反。穗，似醉反，本亦作遂。}獻諸天子。\footnote{拔而貢之。}王命\CJKunderline{唐叔}歸\CJKunderline{周公}於東，\footnote{異畝同穎，天下和同之象，\CJKunderline{周公}之德所致。\CJKunderline{周公}東征未還，故命\CJKunderline{唐叔}以禾歸\CJKunderline{周公}。\CJKunderline{唐叔}後封晉。}作\CJKunderwave{歸禾}。\footnote{亡。}


{\noindent\zhuan\zihao{6}\fzbyks 傳“\CJKunderline{唐叔}”至“一穗”。正義曰：昭十五年\CJKunderwave{左傳}云：“叔父\CJKunderline{唐叔}，\CJKunderline{成王}之母弟。”指言“\CJKunderline{唐叔}得禾”，知其“所食邑內得異禾”也。\CJKunderline{唐叔}食邑，書傳無文。\CJKunderwave{詩}述后稷種禾，於“實秀”之下乃言“實穎”，\CJKunderwave{毛傳}雲“穎垂”，言穗重而垂，是“穎”為穗也。“禾各生一壟而合為一穗”，言其異也。\CJKunderwave{書傳}云：“\CJKunderline{成王}之時，有三苗貫桑葉而生，同為一穗，其大盈車，長几充箱，民得而上諸\CJKunderline{成王}。”下傳雲“拔而貢之”，若是盈車之穗，不可手拔而貢，孔不用\CJKunderwave{書傳}為說也。 \par}

{\noindent\zhuan\zihao{6}\fzbyks 傳“異畝”至“封晉”。正義曰：禾者，和也，異畝同穎，是天下和同之象，\CJKunderline{成王}以為\CJKunderline{周公}德所感致。於時\CJKunderline{周公}東征未還，故命\CJKunderline{唐叔}以禾歸\CJKunderline{周公}於東也。歸禾年月，史傳無文，不知在啟金縢之先後也。王啟金縢，正當禾熟之月。若是前年得之，於時王疑未解，必不肯歸\CJKunderline{周公}。當是啟金縢之後，喜得東土和平而有此應，故以歸\CJKunderline{周公}也。\CJKunderline{唐叔}後封於晉，經史多矣，傳言此者，欲見此時未封，知在邑內得之。昭元年\CJKunderwave{左傳}稱“\CJKunderline{成王}滅唐，而封太叔焉”，所滅之唐即晉國是也。然則得禾之時，未封於唐，從後稱之為“\CJKunderline{唐叔}”耳。 \par}

{\noindent\shu\zihao{5}\fzkt “\CJKunderline{唐叔}”至“歸禾”。正義曰：\CJKunderline{成王}母弟\CJKunderline{唐叔}，於其食邑之內得禾，下異畝壟，上同穎穗，以其有異，拔而貢於天子，以為\CJKunderline{周公}德所感致。於時\CJKunderline{周公}東征未反，王命\CJKunderline{唐叔}歸\CJKunderline{周公}於東，命有言辭。史敘其事，作\CJKunderwave{歸禾}之篇。 \par}

\CJKunderline{周公}既得命禾，旅天子之命，\footnote{已得\CJKunderline{唐叔}之禾，遂陳\CJKunderline{成王}歸禾之命，而推美\CJKunderline{成王}。善則稱君。}作\CJKunderwave{嘉禾}。\footnote{天下和同，政之善者，故\CJKunderline{周公}作書以“嘉禾”名篇告天下。亡。}

{\noindent\zhuan\zihao{6}\fzbyks 傳“已得”至“稱君”。正義曰：鄭雲“受王歸已禾之命與其禾”，以為既得命，“禾”謂復得禾，義當然矣。\CJKunderline{成王}歸禾之命必歸美\CJKunderline{周公}，\CJKunderline{周公}陳歸禾之命又推美\CJKunderline{成王}，是“善則稱君”之義也。“善則稱君”，\CJKunderwave{坊記}文也。 \par}

{\noindent\zhuan\zihao{6}\fzbyks 傳“天下”至“下亡”。正義曰：“嘉”訓善也，言此禾之善，故以善禾名篇。陳天子之命，故當佈告天下，此以善禾為書之篇名，後世同穎之禾遂名為“嘉禾”，由此也。二篇東征未還時事，\CJKunderline{微子}受命應在此篇後。篇在前者，蓋先封\CJKunderline{微子}，後布此書故也。 \par}

{\noindent\shu\zihao{5}\fzkt “\CJKunderline{周公}”至“嘉禾”。正義曰：\CJKunderline{周公}既得王所命禾，乃陳天子歸禾之命為文辭，稱此禾之善，推美於\CJKunderline{成王}。史敘其事，作\CJKunderwave{嘉禾}之篇。 \par}

%%% Local Variables:
%%% mode: latex
%%% TeX-engine: xetex
%%% TeX-master: "../Main"
%%% End:

%% -*- coding: utf-8 -*-
%% Time-stamp: <Chen Wang: 2024-04-02 11:42:41>

% {\noindent \zhu \zihao{5} \fzbyks } -> 注 (△ ○)
% {\noindent \shu \zihao{5} \fzkt } -> 疏

\chapter{卷十四}


\section{康誥第十一}


\CJKunderline{成王}既伐\CJKunderline{管叔}、\CJKunderline{蔡叔}，\footnote{滅三監。}以殷餘民封\CJKunderline{康叔}，\footnote{以三監之民國\CJKunderline{康叔}為衛侯，\CJKunderline{周公}懲其數叛，故使賢母弟主之。○數叛，上所角反，下亦作畔。}作\CJKunderwave{康誥}、\CJKunderwave{酒誥}、\CJKunderwave{梓材}。

康誥\footnote{命\CJKunderline{康叔}之誥。康，圻內國名。叔，封字。○梓音子。圻，具依反。}


{\noindent\zhuan\zihao{6}\fzbyks 傳“以三”至“主之”。正義曰：此序亦與上相顧為首引。初言“三監叛”，又言“黜殷命”，此雲“既伐\CJKunderline{管叔}、\CJKunderline{蔡叔}”，言“以殷餘民”，圻內之餘民，故云“以三監之民國\CJKunderline{康叔}為衛侯”。然古字“邦”、“封”同，故漢有上邦、下邦縣，“邦”字如“封”字，此亦云“邦\CJKunderline{康叔}”，若\CJKunderwave{分器}序雲“邦諸侯”，故云“國\CJKunderline{康叔}”。並以三監之地封之者，\CJKunderline{周公}懲其數叛，故使賢母弟主之。此始一叛而云“數叛”者，以六州之眾悉來歸周，殷之頑民叛逆天命，至今又叛，據周言之，故云“數叛”。故\CJKunderwave{多方}云：“爾乃不大宅天命，爾乃屑播天命。”以不從天命，故云“叛”也。古者大國不過百里，\CJKunderwave{周禮}上公五百里，侯四百里，孟軻有所不信。\CJKunderwave{費誓}注云，伯禽率七百里之內附庸諸侯，則魯猶非七百里之封。而\CJKunderline{康叔}封千里者，\CJKunderline{康叔}時為方伯，殷之圻內諸侯並屬之，故得總言“三監”，且其實地不方平，計亦不能大於魯也。故\CJKunderwave{左傳}云：“宋、衛，吾匹也。”又曰：“寡君未嘗後衛君。”且言千里，亦大率言之耳。何者？邢在襄國，河內即東圻之限，故以賜諸侯西山，即有黎、潞、河、濟之西，以曹地約有千里也。以此鄭云：“初封於衛，至子孫而並邶、鄘也。”其\CJKunderwave{地理志}邶、鄘之民皆遷，分衛民於邶、鄘，故異國而同風，所以\CJKunderwave{詩}分為三。孔與同否未明也。既三年滅三監，七年始封\CJKunderline{康叔}，則於其間更遣人鎮守，自不知名號耳。 \par}

{\noindent\zhuan\zihao{6}\fzbyks 傳“命\CJKunderline{康叔}”至“封字”。正義曰：以定四年\CJKunderwave{左傳}祝佗雲“命以\CJKunderwave{康誥}”，故以為“命\CJKunderline{康叔}之誥”。知“康，圻內國名”者，以管、蔡、郕、霍皆國名，則康亦國名而在圻內。馬、王亦然，惟\CJKunderline{鄭玄}以康為諡號，以\CJKunderwave{史記·世家}雲“生康伯”故也。則孔以康伯為號諡，而\CJKunderline{康叔}之“康”猶為國，而號諡不見耳。 \par}

{\noindent\shu\zihao{5}\fzkt “\CJKunderline{成王}”至“康誥”。正義曰：既伐叛人三監之\CJKunderline{管叔}、\CJKunderline{蔡叔}等，以殷餘民國\CJKunderline{康叔}為衛侯，\CJKunderline{周公}以王命戒之，作\CJKunderwave{康誥}、\CJKunderwave{酒誥}、\CJKunderwave{梓材}三篇之書也。其\CJKunderwave{酒誥}、\CJKunderwave{梓材}亦戒\CJKunderline{康叔}，但因事而分之。然\CJKunderwave{康誥}戒以德刑，又以化紂嗜酒，故次以\CJKunderwave{酒誥}，卒若梓人之治材為器為善政以結之。 \par}

惟三月哉生魄，\footnote{\CJKunderline{周公}攝政七年三月。始生魄，月十六日，明消而魄生。○魄，字又作𩲸，普白反，馬云：“𩲸，朏也，謂月三日始生兆朏，名曰魄。”}\CJKunderline{周公}初基，作新大邑於東國洛，四方民大和會。\footnote{初造基，建作王城大都邑於東國洛汭，居天下上中，四方之民大和悅而集會。○汭，如銳反。}


{\noindent\zhuan\zihao{6}\fzbyks 傳“\CJKunderline{周公}”至“魄生”。正義曰：知“\CJKunderline{周公}攝政七年之三月”者，以\CJKunderwave{洛誥}即七年反政而言新邑營及獻卜之事，與\CJKunderwave{召誥}參同，俱為七年，此亦言作新邑，又同\CJKunderwave{召誥}，故知七年三月也。若然，\CJKunderwave{書傳}雲四年建衛侯而封\CJKunderline{康叔}，五年營成洛邑，六年制禮作樂。\CJKunderwave{明堂位}雲“昔者\CJKunderline{周公}朝諸侯於明堂之位”，即雲“頒度量，而天下大服”，又云“六年制禮作樂”，是六年已有明堂在洛邑而朝諸侯。言六年已作洛邑而有明堂者，\CJKunderwave{禮記}後儒所錄，\CJKunderwave{書傳}伏生所造，皆孔所不用。始生魄，月十六日戊午，社於新邑之明日。“魄”與“明”反，故云“明消而魄生”。 \par}

{\noindent\zhuan\zihao{6}\fzbyks 傳“初造”至“集會”。正義曰：所以初基東國洛者，以天下土中故也。其\CJKunderwave{召誥}與\CJKunderwave{大司徒}文之所出。\CJKunderwave{釋言}云：“集，會也。”以主治民，故民服悅而見大平也。“初基”者，謂初始營建基址，作此新邑，此史總序言之。鄭以為此時未作新邑，而以“基”為謀，大不辭矣。 \par}

侯、甸、男邦、採、衛，百工播民和，見士於周。\footnote{此五服諸侯，服五百里。侯服去王城千里，甸服千五百里，男服去王城二千里，採服二千五百里，衛服三千里，與\CJKunderwave{禹貢}異制。五服之百官，播率其民和悅，並見即事於周。}\CJKunderline{周公}咸勤，乃洪大誥治。\footnote{\CJKunderline{周公}皆勞勉五服之，人遂乃因大封命，大誥以治道。○陸云：“乃洪治同，一本作‘\CJKunderline{周公}乃洪大誥治’。”勞，力報反。}

{\noindent\zhuan\zihao{6}\fzbyks 傳“此五”至“於周”。正義曰：“男”下獨有“邦”，以五服男居其中，故舉中則五服皆有“邦”可知，言“邦”見其國君焉。以\CJKunderwave{大司馬職}\CJKunderwave{大行人}，故知五服，“服五百里”。\CJKunderwave{禹貢}五服通王畿，此在畿外，去王城五百里，故每畿計之，至衛服三千里，言“與\CJKunderwave{禹貢}異制”也，通王畿與不通為異。以此計畿之均，故須土中。若然，黃帝與帝嚳居偃師，餘非土中者，自出當時之宜。實在土中，因得而美善之也。不見要服者，鄭云：“以遠於役事而恆闕焉。”君行必有臣從，即卿大夫及士，見亦主其勞，故云五服之內，百官播率其民和悅即事。以土功勞事，民之所苦也，而此和悅，見太平也。而\CJKunderwave{書傳}云：“示之以力役，其民猶至，況導之以禮樂乎？”是也。 \par}

{\noindent\zhuan\zihao{6}\fzbyks 傳“\CJKunderline{周公}”至“治道”。正義曰：太保以戊申至，七月庚戌，已雲“庶殷攻位於洛汭”，則庶殷先與之期於前至也。\CJKunderline{周公}以十三日乙卯“朝至於洛，則達觀於新邑營”，此日當勉其民，此因命而並言之。序雲“邦\CJKunderline{康叔}”，“洪”，大也，為大封命大誥\CJKunderline{康叔}以治道也。\CJKunderline{鄭玄}以“洪”為代，言\CJKunderline{周公}代\CJKunderline{成王}誥。何故代誥而反誥王，呼之曰“\CJKunderline{孟侯}”？為不辭矣。 \par}

{\noindent\shu\zihao{5}\fzkt “惟三”至“誥治”。正義曰：言惟以\CJKunderline{周公}攝政七年之三月，始明死而生魄，月十六日己未，於時\CJKunderline{周公}初造基趾，作新大邑於東國洛水之汭，四方之民大和悅而集會，言政治也。此所集之民，即侯、甸、男、採、衛五服。百官播率其民和悅，並見即事於周之東國，而\CJKunderline{周公}皆慰勞勸勉之。乃因大封命以\CJKunderline{康叔}為衛侯，大誥以治道。 \par}

王若曰：“\CJKunderline{孟侯}，朕其弟，小子\CJKunderline{封}。\footnote{\CJKunderline{周公}稱\CJKunderline{成王}命，順\CJKunderline{康叔}之德，命為\CJKunderline{孟侯}。孟，長也。五侯之長，謂方伯，使\CJKunderline{康叔}為之。言王使我命其弟封。封，\CJKunderline{康叔}名。稱小子，明當受教訓。○長，丁丈反，下同。}惟乃丕顯考\CJKunderline{文王}，克明德慎罰，\footnote{惟汝大明父\CJKunderline{文王}，能顯用俊德，慎去刑罰，以為教首。○去，羌呂反，下“欲去”、“去疾”同。}不敢侮鰥寡，庸庸，祇祇，威威，顯民。\footnote{惠恤窮民，不慢鰥夫寡婦，用可用，敬可敬，刑可刑，明此道以示民。}用肇造我區夏，越我一二邦，\footnote{用此明德慎罰之道，始為政於我區域諸夏，故於我一二邦皆以修治。}以修我西土。惟時怙冒，聞於上帝，帝休。\footnote{我西土岐周，惟是怙恃\CJKunderline{文王}之道，故其政教冒被四表，上聞於天，天美其治。○怙音戶。冒，莫報反，覆也。聞如字，徐又音問。}天乃大命\CJKunderline{文王}，殪戎殷，誕受厥命，\footnote{天美\CJKunderline{文王}，乃大命之殺兵殷，大受其王命。謂三分天下有其二，以授\CJKunderline{武王}。○殪，於計反。}越厥邦厥民，惟時敘。\footnote{於其國，於其民，惟是次序，皆\CJKunderline{文王}教。}乃寡兄勖，肆汝小子\CJKunderline{封}，在茲東土。”\footnote{汝寡有之兄\CJKunderline{武王}，勉行\CJKunderline{文王}之道，故汝小子封得在此東土為諸侯。○勖，許玉反。}


{\noindent\zhuan\zihao{6}\fzbyks 傳“\CJKunderline{周公}”至“教訓”。正義曰：以“曰”者為命辭，故曰“\CJKunderline{周公}稱\CJKunderline{成王}命，順\CJKunderline{康叔}之德，命為\CJKunderline{孟侯}。孟，長也。五侯之長，謂方伯”。使\CJKunderline{康叔}為之長者，即州牧也。“五侯之長”，五等諸侯之長也。而\CJKunderwave{左傳}云：“五侯九伯，汝實徵之。”彼謂上公之伯，故徵九伯。而此“五侯”當州牧之“五侯”，與彼不同。\CJKunderwave{王制}有連、屬、率、伯也，孔以五侯亦方伯，則四方者皆可為方伯，而此方伯自是州牧也。\CJKunderline{康叔}以母弟令德受大國封命，固非卒及連、屬也。虞夏及周既有牧，又\CJKunderwave{離騷}雲“伯昌作牧”，殷亦有牧，伯四代皆通也，非如\CJKunderline{鄭玄}雲“殷之州長曰伯”。以“稱小子”為幼弱，故“明當受教訓”，故云“使我命其弟”，為親親而使我用戒故也。此指命\CJKunderline{康叔}為之，而鄭以總告諸侯，依\CJKunderwave{略說}以太子十八為\CJKunderline{孟侯}而呼\CJKunderline{成王}。既禮制無文，義理駢曲，豈\CJKunderline{周公}自許天子，以王為\CJKunderline{孟侯}？皆不可信也。 \par}

{\noindent\zhuan\zihao{6}\fzbyks 傳“惟汝”至“教首”。正義曰：以近而可法，不過子之法父，故舉\CJKunderline{文王}也。法者不過除惡行善，故云“明德慎罰”也。 \par}

{\noindent\zhuan\zihao{6}\fzbyks 傳“惠恤”至“示民”。正義曰：“用可用，敬可敬”，即“明德”也。“用可用”謂小德小官，“敬可敬”謂大德大官，“刑可刑”謂“慎罰”也。 \par}

{\noindent\zhuan\zihao{6}\fzbyks 傳“天美”至“\CJKunderline{武王}”。正義曰：“天美\CJKunderline{文王}，乃天命之殺兵殷”者，“殪”，殺也，“戎”，兵也，用誅殺之道以兵患殷。\CJKunderline{文王}以伐殷事未卒而言“殺兵殷”者，謂三分有二，為滅殷之資也。 \par}

{\noindent\shu\zihao{5}\fzkt “王若”至“東土”。正義曰：言\CJKunderline{周公}稱\CJKunderline{成王}命，順\CJKunderline{康叔}之德而言曰：“命汝為\CJKunderline{孟侯}。王又使我教命其弟小子封。其所教命者，惟汝大明德之父\CJKunderline{文王}，能顯用俊德，慎去刑罰，以為教首。故惠恤窮民，不侮慢鰥夫寡婦，況貴強乎？其明德，用可用，敬可敬，其慎罰，威可威者，顯此道以示民。用此道，故始為政於我區域諸夏，由是於我一二諸國漸以修治也。上政既修，我西土惟是怙恃\CJKunderline{文王}之道，故其政教冒被四表，聞於上天，天美其治道。以此上天乃大命\CJKunderline{文王}以誅殺之道，用兵除害於殷，大受其王命，三分天下而有其二也。其所受二分者，於其國，於其民，惟是皆有次序，以\CJKunderline{文王}之德故也。汝寡有之兄\CJKunderline{武王}，勉行\CJKunderline{文王}之道，故受命克殷，今汝小子封故得在此東土為諸侯。是\CJKunderline{文王}之道，明德慎罰，既用受命，\CJKunderline{武王}無所復加，以為勉行，所以汝必法之。” \par}

王曰：“嗚呼！\CJKunderline{封}，汝念哉！\footnote{念我所以告汝之言。}今民將在祇遹乃文考，紹聞衣德言。\footnote{今治民將在敬循汝文德之父，繼其所聞，服行其德言，以為政教。○遹音聿，又音述，馬紹：“述也。”衣如字，徐於既反。}往敷求於殷先哲王，用保乂民。\footnote{汝往之國，當布求殷先智王之道，用安治民。}汝丕遠惟商耇成人，宅心知訓。\footnote{汝當大遠求商家耇老成人之道，常以居心，則知訓民。○耇音狗。}別求聞由古先哲王，用康保民。\footnote{又當別求所聞父兄用古先智王之道，用其安者以安民。}弘於天，若德裕乃身，不廢在王命。”\footnote{大於天，為順德，則不見廢，常在王命。}


{\noindent\zhuan\zihao{6}\fzbyks 傳“今治”至“政教”。正義曰：“繼其所聞，服行其德言”者，謂\CJKunderline{文王}先有所聞善事，今令\CJKunderline{康叔}繼續其\CJKunderline{文王}所聞善事，被服而施行其德言，以為政教也。 \par}

{\noindent\zhuan\zihao{6}\fzbyks 傳“汝當”至“訓民”。正義曰：上雲“敷求殷先哲王”，謂求殷之賢君，此言“求商家耇老成人”，謂求殷之賢臣。“大遠”者，備遍求之。 \par}

{\noindent\zhuan\zihao{6}\fzbyks 傳“又當”至“安民”。正義曰：以父兄乃所居殷外，故云“別求”。上只言“遹乃文考”，並言“兄”者，以上雲“寡兄勖”，則以文武道同，言文可以兼武，故並言“父兄”也。“古先哲王”，鄭云：“虞夏也。”孔亦當然。以上代與今事遠，不可以同，故言“用其安者”。 \par}

{\noindent\zhuan\zihao{6}\fzbyks 傳“大於”至“王命”。正義曰：以天道人用而光大之，故因雲“大”也。其\CJKunderline{文王}及殷古先哲王，與天其道不異，以前後聖蹟雖殊，同天不二也。以\CJKunderline{康叔}亞聖大賢，治殷餘惡，故使之用天道為順德也。 \par}

{\noindent\shu\zihao{5}\fzkt “王曰嗚呼封汝”至“王命”。正義曰：既言\CJKunderline{文王}“明德慎罰”之訓，\CJKunderline{武王}尚行之，汝既得為君，方別陳明德之事，故稱王命而言曰：“嗚呼！封，汝常念我所以告汝之言哉！今治民所行，將在敬循汝文德之父，繼其所聞者，服行其德言，以為政教。汝往之國，當分佈求於殷先智王之道，用安治民。不但法其先君，汝又當須大遠求商家耇老成人之道，居之於心，即知訓民矣。其外又更當別求所聞父兄用古先智王之道，用其安者以安民。即古虞夏之道也。人事既然，又闡大於天之道而為順德，又加之寬容，則汝身不見廢，常在王命。” \par}

王曰：“嗚呼！小子\CJKunderline{封}，\xpinyin*{恫瘝}乃身，敬哉！\footnote{恫，痛。瘝，病。治民務除惡政，當如痛病在汝身欲去之，敬行我言。○恫音通，又敕動反。敕，古頑反。}天畏棐忱，民情大可見，小人難保。\footnote{天德可畏，以其輔誠。人情大可見，以小人難安。○棐音匪，又芳鬼反。忱，巿林反。}往盡乃心，無康好逸豫，乃其乂民。\footnote{往當盡汝心為政，無自安好逸豫寬身，其乃治民。○盡，徐子忍反。好，呼報反。}我聞曰：‘怨不在大，亦不在小，惠不惠，懋不懋。’\footnote{不在大，起於小；不在小，小至於大。言怨不可為，故當使不順者順，不勉者勉。○懋音茂。}已！汝惟小子，乃服惟弘王，應保殷民，\footnote{已乎！汝惟小子，乃當服行德政，惟弘\CJKunderline{大王}道，上以應天，下以安我所受殷之民眾。○應，應對之應，注同，徐於甑反。}亦惟助王宅天命，作新民。”\footnote{弘王道，安殷民，亦所以惟助王者居順天命，為民日新之教。}


{\noindent\zhuan\zihao{6}\fzbyks 傳“恫痛”至“我言”。正義曰：“恫”聲類於痛，故“恫”為痛也。“瘝，病”，\CJKunderwave{釋詁}文。以痛病在汝身以述治民，故務除惡政如已病也。戒之而言“敬”，故知“敬行我言”也。\CJKunderline{鄭玄}云：“刑罰及已為痛病。”其義不及去惡若已病也。 \par}

{\noindent\zhuan\zihao{6}\fzbyks 傳“天德”至“難安”。正義曰：人情所以大可見者，以小人難安為可見，故須安之。 \par}

{\noindent\zhuan\zihao{6}\fzbyks 傳“不在”至“者勉”。正義曰：以致怨恐謂由大惡，故云“不在大，起於小”，言怨由小事起。“不在小”者，謂為怨不恆在小，言其初小，漸至於大怨，故使不順者順，不勉者勉，其怨自消也。 \par}

{\noindent\zhuan\zihao{6}\fzbyks 傳“弘王”至“之教”。正義曰：“亦所以惟助王”者，言非直\CJKunderline{康叔}身行有益，亦惟助王者居順天命，為民日新之教，謂漸致太平，政教日日益新也。 \par}

{\noindent\shu\zihao{5}\fzkt “王曰嗚呼小”至“新民”。正義曰：所明而云行天人之德者，其要在於治民，故言王曰：“嗚呼！小子封，治民為善而除惡政，當如痛病在汝身欲去之，敬行我言哉！所以去惡政者，以天德可畏者，以其輔誠故也。以民情大率可見，所以可見者，以小人難保也。安之既難，其往治之，當盡汝心為政，無自安好逸豫而寬縱，乃其可以治民。我聞名遺言曰，人之怨不在事大，或由小事而起。雖由小事而起，亦不恆在事小，因小至大。是為民所怨，事不可為。當使施順，令不順者順。勉力勸行，令不勉者勉。則其怨小大都消，令汝消怨者。已乎！汝惟小子，乃當服行政德，惟弘\CJKunderline{大王}道，上以應天，下以安我所受殷民。不但汝身所當行，此亦惟助王者居順天命，為民日新之教。” \par}

王曰：“嗚呼！\CJKunderline{封}，敬明乃罰。\footnote{嘆而敕之，凡行刑罰，汝必敬明之。欲其重慎。}人有小罪，非眚，乃惟終，自作不典，式爾，\footnote{小罪非過失，乃惟終身行之，自為不常，用犯汝。○眚，所領反，本亦作省。}有厥罪小，乃不可不殺。乃有大罪，非終，乃惟眚災，適爾，既道極厥辜，時乃不可殺。”\footnote{汝盡聽訟之理以極其罪，是人所犯，亦不可殺，當以罰宥論之。○宥，於救反。}

{\noindent\shu\zihao{5}\fzkt “王曰嗚呼封敬”至“可殺”。正義曰：以上既言“明德”之理，故此又云“慎罰”之義，而王言曰：“嗚呼！封，又當敬明汝所行刑罰，須明其犯意。人有小罪，非過誤為之，乃惟終身自為不常之行，用犯汝，如此者，有其罪小，乃不可不殺，以故犯而不可赦。若人乃有大罪，非終行之，乃惟過誤為之，以此故，汝當盡斷獄之道以窮極其罪，是人所犯，乃不可以殺，當以罰宥論之，以誤故也。即原心定罪，斷獄之本，所以須敬明之也。” \par}

王曰：“嗚呼！\CJKunderline{封}，有敘，時乃大明服，\footnote{嘆政教有次敘，是乃治理大明，則民服。}惟民其敕懋和。\footnote{民既服化，乃其自敕正勉為和。}若有疾，惟民其畢棄咎。\footnote{化惡為善，如欲去疾，治之以理，則惟民其盡棄惡修善。○咎，其九反。}若保赤子，惟民其康乂。\footnote{愛養紉宴安孩兒赤子，不失其欲，惟民其皆安治。○孩，亥才反。}非汝\CJKunderline{封}刑人殺人，\footnote{言得刑殺罪人。}無或刑人殺人。\footnote{無以得刑殺人，而有妄刑殺非辜者。}非汝封又\CJKunderline{曰}\xpinyin*{劓刵}人，\footnote{劓，截鼻。刵，截耳。刑之輕者，亦言所得行。○劓，魚器反。刵，如志反。}無或劓刵人。”\footnote{所以舉輕以戒，為人輕行之。}


{\noindent\zhuan\zihao{6}\fzbyks 傳“化惡”至“修善”。正義曰：人之有疾，治之以理則疾去。人之有惡，化之以道則惡除。 \par}

{\noindent\zhuan\zihao{6}\fzbyks 傳“愛養”至“安治”。正義曰：既去惡，乃須愛養之為善，人為上養，則化所行，故言其皆安治。子生赤色，故言“赤子”。 \par}

{\noindent\zhuan\zihao{6}\fzbyks 傳“劓截”至“得行”。正義曰：以國君故得專刑殺於國中，而不可濫其刑，即墨、劓、剕、宮也。“劓”在五刑為截鼻，而有“刵”者，\CJKunderwave{周官}五刑所無，而\CJKunderwave{呂刑}亦云“劓刵”，\CJKunderwave{易·噬嗑}上九雲“何校滅耳”。\CJKunderline{鄭玄}以臣從君坐之刑，孔意然否未明，要有刵而不在五刑之類。言“又曰”者，\CJKunderline{周公}述\CJKunderline{康叔}，豈非“汝封”又自言曰得劓刵人？此“又曰”者，述\CJKunderline{康叔}之“又曰”。 \par}

{\noindent\shu\zihao{5}\fzkt “王曰嗚呼封有”至“刵人”。正義曰：以刑者政之助，不得已即用之；非情好殺害，故又本於政不可以濫刑，而王言曰：“嗚呼！封，欲正刑之本，要而汝政教有次序，是乃治理大明則民服。惟民既服從化，其自敕正勉力而平和。然政之化惡為善，若有病而欲去之，治之以理，則惟民其盡棄惡而修善。言愛養人陽母之安赤子，惟民為善，其皆安治。為政保民之如此，不可行以淫刑，豈非汝封得刑人殺人乎？言得刑殺不可以得故，而有濫刑人殺人無辜也。非汝封又曰劓刵人，無以得故，而有所濫劓刵人之無罪者也。” \par}

王曰：“外事，汝陳時臬，司師，茲殷罰有倫。”\footnote{言外土諸侯奉王事，汝當布陳是法，司牧其眾，及此殷家刑罰有倫理者兼用之。○臬，魚列反。}又曰：“要囚，服念五六日，至於旬時，丕蔽要囚。”\footnote{要囚，謂察其要辭以斷獄。既得其辭，服膺思念五六日，至於十日，至於三月，乃大斷之。言必反覆思念，重刑之至也。○要，於宵反。蔽，必世反。斷，丁亂反，下及篇末同。覆，芳服反。}


{\noindent\zhuan\zihao{6}\fzbyks 傳“言外”至“用之”。正義曰：外土以獄事上於州牧之官，為奉土事，汝當用刑書，為布陳是刑法，為司牧其眾，故受而聽之。既衛居殷墟，又周承於殷後，刑書相因，故兼用其有理者。謂當時刑書，或無正條，而殷有故事，可兼用，若今律無條，求故事之比也。“臬”為準限之義，故為法也。 \par}

{\noindent\zhuan\zihao{6}\fzbyks 傳“要囚”至“之至”。正義曰：言“要囚”，明取要辭於囚以思訖事定，故言“乃大斷之”。多至三月，故云“反覆思念，重刑之至”。顧氏云：“‘又曰’者，\CJKunderline{周公}重言之也。” \par}

{\noindent\shu\zihao{5}\fzkt “王曰外事”至“要囚”。正義曰：言不濫刑，不但國內，而王言曰：“若外土諸侯奉王事以至汝，汝當布陳是刑法以司牧其眾，及此殷家刑罰有倫理者兼用之。”\CJKunderline{周公}又重言曰：“既用刑法，要察囚情，得其要辭，以斷其獄。當須服膺思念之，五日六日，次至於十日，遠至於三月，一時乃大斷囚之要辭。”言必反覆重之如此，乃得無濫故耳。 \par}

王曰：“汝陳時臬事，罰蔽殷彝，\footnote{陳是法事，其刑罰斷獄，用殷家常法，謂典刑故事。○彝，以支反。}用其義刑義殺，勿庸以次汝封。\footnote{義，宜也。用舊法典刑，宜於時世者以刑殺，勿庸以就汝封之心所安。}乃汝盡遜曰時敘，惟曰未有遜事。\footnote{乃使汝所行盡順，曰是有次敘，惟當自謂未有順事，君子將興，自以為不足。}已！汝惟小子，未其有若汝\CJKunderline{封}之心，朕心朕德惟乃知。\footnote{已乎！他人未其有若汝封之心。言汝心最善，我心我德惟汝所知。欲其明\CJKunderline{成王}所以命已之款心。○款，苦管反。}


{\noindent\zhuan\zihao{6}\fzbyks 傳“陳是”至“故事”。正義曰：“陳是法事”，即上“汝陳時臬事”，“罰蔽殷彝”即上“殷罰有倫”。上據有初思念得失，此據臨時行事也。 \par}

{\noindent\zhuan\zihao{6}\fzbyks 傳“已乎”至“款心”。正義曰：此言“我”，我王也，以王命，故言王為“我”，以\CJKunderline{康叔}為“己”。若汝不善我王家心德，汝所不知，則我不順命汝款曲之心。只由汝最善，我王心德汝所遍知，故我王命汝以款曲之心。述\CJKunderline{康叔}為言，故云“己”，欲令\CJKunderline{康叔}明識此意也。 \par}

{\noindent\shu\zihao{5}\fzkt “王曰汝”至“乃知”。正義曰：此又申上既要囚思念，定其大斷若為，而王言曰：“汝當陳是刑書之法以行事，其刑法斷獄，用殷家所行常法故事，其陳法殷彝，皆用其合宜者以刑殺，勿用以就汝封意之所安而自行也，以用心不如依法故耳。言汝不但依法，乃使汝所行盡順，曰是有次敘，猶當自惟曰未有順事，其有餘若不足故耳。必期汝於大幸已乎！汝惟小子耳，而他人未其有若汝封之心。言汝心最善，汝心既善，我心我德惟汝所委知也。 \par}

凡民自得罪，寇攘奸宄，殺越人於貨，\footnote{凡民用得罪，為寇盜攘竊奸宄，殺人顛越人，於是以取貨利。○攘，如羊反。宄音軌。}\xpinyin*{暋}不畏死，罔弗憝。”\footnote{暋，強也。自強為惡而不畏死，人無不惡之者，言當消絕之。○暋音敏。憝，徒對反，徐徒猥反。強，其丈反。無不惡，烏路反，下“所大惡”、“疾惡”、“亦惡”並音同。}


{\noindent\zhuan\zihao{6}\fzbyks 傳“凡民”至“貨利”。正義曰：“目”，用也，言所用得罪者，由寇攘也。而為之於外內，既有劫竊，其劫竊皆有殺有傷，“越人”謂不死而傷，皆為之而取貨利故也。 \par}

{\noindent\zhuan\zihao{6}\fzbyks 傳“暋強”至“絕之”。正義曰：“暋，強也”，於\CJKunderwave{盤庚}已訓而此重詳之，以由此得罪，當須絕之。 \par}

{\noindent\shu\zihao{5}\fzkt “凡民”至“弗憝”。正義曰：言人所慎刑者，以凡民所用得罪者，寇盜攘竊於外奸內宄，而殺害及顛越於人以取貨利也。自強為之而不畏死，此為人無不惡之者，以此須刑絕之，故當慎刑罰耳。 \par}

王曰：“\CJKunderline{封}，元惡大憝，矧惟不孝不友。\footnote{大惡之人猶為人所大惡，況不善父母，不友兄弟者乎？言人之罪惡，莫大於不孝不友。}子弗祇服厥父事，大傷厥考心。\footnote{為人子不能敬身服行父道，而怠忽其業，大傷其父心，是不孝。}於父不能字厥子，乃疾厥子。\footnote{於為人父不能字愛其子，乃疾惡其子，是不慈。}於弟弗念天顯，乃弗克恭厥兄。\footnote{於為人弟不念天之明道，乃不能恭事其兄，是不恭。}兄亦不念鞠子哀，大不友于弟。\footnote{為人兄亦不念稚子之可哀，大不篤友于弟，是不友。○鞠，居六反。}


{\noindent\zhuan\zihao{6}\fzbyks 傳“大惡”至“不友”。正義曰：言將有作奸宄大惡，猶為人所大惡，況不孝父母，不善兄弟者乎？\CJKunderwave{孝經}雲“五刑之屬三千，而罪莫大於不孝”是也。\CJKunderwave{釋親}云：“善父母為孝，善兄弟為友。”下文不言“母”，母同於父。父子尊卑而異等，故“孝”名上不通於下。其兄弟雖有長幼而同倫，故共“友”名也。 \par}

{\noindent\zhuan\zihao{6}\fzbyks 傳“為人”至“不孝”。正義曰：“考”亦通生死，即此文及\CJKunderwave{酒誥}是也。下\CJKunderwave{曲禮}雲“死曰考”，是對例耳。人予以述成父事為孝，怠忽其業，即“其肯曰，我有後，不棄基”，故為大傷父心，即是上不孝也。則子不述父事，當輕於盜殺，況以為甚者，此聖人緣心立法，人莫不緣身本於父母也，自親以及物，天然之理，故\CJKunderwave{孝經}曰：“不愛其親而愛他人者，謂之悖德。不敬其親而敬他人者，謂之悖禮。以順則逆，民無則焉。不在於善，而皆在於凶德”是也。以此言賊殺他人，罪小於骨肉相乖阻。但於他人言其極者，於親言其小者，小則有不和詈爭鬥訟相傷者也。於親小則傷心，大乃逆命，毆罵殺害，互相發起而可知也。 \par}

{\noindent\zhuan\zihao{6}\fzbyks 傳“於為”至“不慈”。正義曰：上文不言“不慈”，意以“不孝”為總焉。父當言“義”而云“不慈”者，以父母於子併為慈，因父有愛敬多少而分之。言父義母慈，而由慈以義，故雖義言“不慈”，且見父兼母耳。 \par}

{\noindent\zhuan\zihao{6}\fzbyks 傳“於為”至“不恭”。正義曰：善兄弟曰友，此言“不恭”者，“友”思念之辭，兄弟同倫，故俱言“友”；雖同倫而有長幼，其心友而貌恭，故因兄弟而分“友”文為二而言“恭”也。五教，即\CJKunderwave{左傳}文十八年史克言也。於此言“天之明道”者，父子天性，不嫌非天明，故於兄弟言之。因上先言“不孝”，先言子於父，故此“不友”先言弟於兄，若舉中以見上下，故此言天明，見五教皆是，即\CJKunderwave{孝經}雲“則天之明”，\CJKunderwave{左傳}雲“為父子兄弟姻媾以象天明”，是於天理常然，為天明白之道。 \par}

{\noindent\zhuan\zihao{6}\fzbyks 傳“為人”至“不友”。正義曰：言“亦”者，以兄弟同等而相亦，所謂\CJKunderwave{周書}雲“父子兄弟罪不相及”，即此文也。不孝罪子，非及於父之輩，理所當然。而\CJKunderwave{周官}鄰保以比伍相及，而趙商疑而發問，鄭答云：“\CJKunderwave{周禮}太平制，此為居殷亂而言，斯不然矣。\CJKunderwave{康誥}所云，以骨肉之親，得相容隱，故\CJKunderwave{左傳}云：‘父子兄弟罪不相及。’\CJKunderwave{周禮}所云，據疏人相督率之法，故相連獲罪。故今之律令，大功已上得相容隱，鄰保罪有相及是也。” \par}

惟弔，茲不於我政人得罪，\footnote{惟人至此不孝、不慈、弗友、不恭，不於我執政之人得罪乎？道教不至所致。○吊音的。}天惟與我民彝大湣亂。\footnote{天與我民五常，使父義、母慈、兄友、弟恭、子孝，而廢棄不行，是大滅亂天道。○泯，徐武軫反。}曰，乃其速由\CJKunderline{文王}作罰，刑茲無赦。\footnote{言當速用\CJKunderline{文王}所作違教之罰，刑此亂五常者，無得赦。}

{\noindent\shu\zihao{5}\fzkt “王曰封元”至“無赦”。正義曰：以是所用得其罪，不但寇盜，王命而言曰：“封，非於骨肉之人為大惡，猶尚為人所大惡之，況惟不孝父母，不友兄弟者乎？其罪莫大於不孝也。何者？為人之子不能敬身服行其父事，而怠忽其業，大傷其父心，是不孝也。於為人父不能字愛其子，乃疾惡其子，是不慈也。於為人弟不能念天之明道，故乃不能恭事其兄，是不恭也。為人兄亦不能念稚子之可哀哉，大不友愛於弟，是不友也。惟人所行以至此不孝不友者，豈不由我執政之人道教不至，以得此罪乎？既人罪由教而致，天惟與我民以五常之性，使有恭孝，廢棄不行，是大滅亂天道也。以由我滅亂，曰，乃其疾用\CJKunderline{文王}所作違教之罰，刑此亂五常者，不可赦放也。” \par}

不率大戛，矧惟外庶子、訓人，\footnote{戛，常也。凡民不循大常之教，猶刑之無赦，況在外掌眾子之官主訓民者而親犯乎？○戛，簡八反。}惟厥正人，越小臣諸節。\footnote{惟其正官之人，於小臣諸有符節之吏，及外庶子，其有不循大常者，則亦在無赦之科。}乃別播敷造。民大譽，弗念弗庸，瘝厥君，時乃引惡，惟朕憝。\footnote{汝今往之國，當分別播佈德教，以立民大善之譽。若不念我言、不用我法者，病其君道，是汝長惡，惟我亦惡汝。○別，彼列反，注同。汝長，丁丈反，下同。}


{\noindent\zhuan\zihao{6}\fzbyks 傳“戛常”至“犯乎”。正義曰：“戛”猶楷也，言為楷模之常，故“戛”為常也。述上凡民自得罪，故言“凡民不循大常之教”也。“猶刑之”即上雲“刑茲無赦”故也。亦愚以況智，故言“況在外掌眾子之官主訓民者而親犯乎”，即\CJKunderwave{周官}雲“諸子”，\CJKunderwave{文王世子}雲“庶子”也。以致教諸子，故為“訓人”。\CJKunderwave{周禮}諸子之官亦是王朝之臣，言“在外”者，對父子兄弟為外。惟舉庶子之官者，以其教訓公卿子弟，最為急故也。\CJKunderline{鄭玄}以“訓人”為師長，亦各一家之道也。 \par}

{\noindent\zhuan\zihao{6}\fzbyks 傳“惟其”至“之科”。正義曰：“正官之人”，若\CJKunderwave{周官}三百六十職正官之首。“於小臣諸有符節”者，謂正人之下，非長官之身，下至符吏。“諸有符節”，為教人之故，故言有符節者。非要行道之符節，若為官行文書而有符，今之印者也。以上況之，故言“不循大常，亦在無赦之科”矣。在軍者有旌節，亦得為有符節耳。 \par}

{\noindent\zhuan\zihao{6}\fzbyks 傳“汝今”至“惡汝”。正義曰：言“分別播佈德教”，謂分遣卿大夫為之教民使善。而已有善譽，是“立民以大善之譽”。 \par}

已！汝乃其速由茲義率殺。亦惟君惟長，\footnote{汝乃其速用此典刑宜於時世者，循理以刑殺，則亦惟君長之正道。}不能厥家人，越厥小臣、外正，惟威惟虐，大放王命，乃非德用乂。\footnote{為人君長而不能治其家人之道，則於其小臣外正官之吏，併為威虐，大放棄王命，乃由非德用治之故。}


{\noindent\zhuan\zihao{6}\fzbyks 傳“汝乃”至“正道”。正義曰：此用宜於時以刑殺上不循五常之道者。其“君長”，對則大夫為長，散則人君為長，君而居之，是君亦與長為一。\CJKunderwave{孝經}對例以長為大夫耳。 \par}

{\noindent\zhuan\zihao{6}\fzbyks 傳“為人”至“之故”。正義曰：以五常父母兄弟子即家人之道，\CJKunderwave{易}有家人卦，亦與此同也。不行五教為不能治家人之道，家人不治，則君不明。君既不明，則不察下故則，於其小臣外正官之吏併為威虐，大放棄王命，非德用治，是不明為非德也。 \par}

汝亦罔不克敬典，乃由裕民，惟\CJKunderline{文王}之敬忌。\footnote{常事人之所輕，故戒以無不能敬常。汝用寬民之道，當惟念\CJKunderline{文王}之所敬思而法之。}乃裕民，曰：‘我惟有及。’則予一人以懌。”\footnote{汝行寬民之政，曰：“我惟有及於古。”則我一人以此悅懌汝德。○懌音亦。}

{\noindent\zhuan\zihao{6}\fzbyks 傳“常事”至“法之”。正義曰：“常事”，常所行之事也。人見尋常不為異，故輕之，而以為戒。“\CJKunderline{文王}所敬忌”，即敬德忌刑。鄭云：“‘祇祇、威威’是也。” \par}

{\noindent\zhuan\zihao{6}\fzbyks 傳“汝行”至“汝德”。正義曰：寬則得眾，故五教在寬。上既言“乃由裕民”，此又疊之，汝行寬民之政，曰：“我惟有及於古。”即古賢諸侯。汝惡，我則惡之。汝善，我則愛之。以此，我一人悅懌汝德也。 \par}

{\noindent\shu\zihao{5}\fzkt “不率”至“以懌”。正義曰：言滅五常之害當除，凡民不循大道五常之教，猶刑之，況在外土掌庶子之官，主於訓民，惟其正官之人，及於小臣諸有符節者，併為教首，其心不循大常，豈可赦也？以人之須有五常，汝今往之國，乃當分別播佈德教，以立民大善之譽。若不念我言，不用我法，即病其為君之道，是汝長為惡矣，以此惟我亦惡汝也。已乎！既惡不可為，汝乃其疾用此典刑宜於時世者，循理以刑殺亂常者，則亦惟為人君，惟為人長之正道。既為人君長，不能治其五教，施於家人之道，則於其卑小臣外土正官之吏，惟為威暴，則為酷虐，大放棄王命矣。如是乃由汝非以道德用治之故。由此汝亦無得不能敬其常事，汝用寬民之道，當思惟念用\CJKunderline{文王}之所敬畏而法之。汝以此行寬民之政，曰：“我原惟有及於古。”則我一人天子以此悅懌汝德矣。汝惟宜勤之。 \par}

王曰：“\CJKunderline{封}，爽惟民迪吉康，\footnote{明惟治民之道而善安之。}我時其惟殷先哲王德，用康乂民作求。\footnote{我是其惟殷先智王之德，用安治民，為求等。}矧今民罔迪不適，不迪則罔政在厥邦。”\footnote{治民乃欲求等殷先智王，況今民無道不之。言從教也。不以道訓之，則無善政在其國。}


{\noindent\zhuan\zihao{6}\fzbyks 傳“明惟”至“安之”。正義曰：以慎德刑為明治民之道，教之五常為善，富而不擾為安也。鄭以“迪”為下讀，各為一通也。 \par}

{\noindent\zhuan\zihao{6}\fzbyks 傳“治民”至“其國”。正義曰：以己喻\CJKunderline{康叔}，言我未治之時，乃欲求等殷先智王以致太平者，況今民無道不之。言易從教。不以正道訓民，民不知道，故無善政在其國，為無吉康也。 \par}

{\noindent\shu\zihao{5}\fzkt “王曰封爽”至“厥邦”。正義曰：既言德刑事終而總言之，我所以令汝明德慎罰以施政者，王命所以言曰：“封，為人君，當明惟為治民之道而善安之，故我以是須汝善安民，故我其惟念殷先智聖王之德，用安治民，為求而等之。我於民未治之時，尚求等殷先智王，況今民無道不之而易化，汝若不以道訓之，則無善政在其國，所以須安民以德刑也。” \par}

王曰：“\CJKunderline{封}，予惟不可不監，告汝德之說於罰之行。\footnote{我惟不可不監視古義，告汝施德之說於罰之所行。欲其勤德慎刑。○說如字，徐始銳反。}今惟民不靜，未戾厥心，迪屢未同，\footnote{假令今天下民不安，未定其心，於周教道屢數而未和同。設事之言。○令，力呈反。數，所角反。}爽惟天其罰殛我，我其不怨。\footnote{明惟天其以民不安罰誅我，我其不怨天。汝不治，我罰汝，汝亦不可怨我。○殛，紀力反。}惟厥罪無在大，亦無在多，矧曰其尚顯聞於天？”\footnote{民之不安，雖小邑少民，猶有罰誅，不在多大，況曰不慎罰，明聞於天者乎？言罪大。}


{\noindent\zhuan\zihao{6}\fzbyks 傳“我惟”至“慎刑”。正義曰：以敷求殷先哲王，及別求古先哲王，為己視古義也。德由說而罰須行，故德之言“說”而罰言“行”也。以事終而結上，故云德也。 \par}

{\noindent\zhuan\zihao{6}\fzbyks 傳“假令”至“之言”。正義曰：天下不安，為總說。所以不安，猶“未定其心，於周道屢數而未和同”也。時以大和會，故言“假令”，設不和同事言耳。 \par}

{\noindent\zhuan\zihao{6}\fzbyks 傳“明惟”至“怨我”。正義曰：顧氏云：“明惟天者，言天明察在上，見民不安，乃以刑罰誅戮於我。” \par}

{\noindent\zhuan\zihao{6}\fzbyks 傳“民之”至“罪大”。正義曰：此總德刑而直雲“不慎罰”者，政以德為主，不嫌不明，政失由於濫刑，故舉“罰”以言之。下言“無作怨”，以失罰為罪大。 \par}

{\noindent\shu\zihao{5}\fzkt “王曰封予”至“於天”。正義曰：以汝須善政在國，令我民安，當為政以慎德刑為教，故王又命之曰：“封，我惟不可不視古義，告汝施德之說於罰之所有。”欲其勤德慎刑也。“假令惟天下民不安，未定其心，於周教道屢數而未和同，明惟天其以民不安其罰誅我，我其不怨於天。則汝不治，是其罪，我罰汝，汝亦不可怨我。我以民之不安，惟其罰之，無在大邑，無在多民，以少猶誅罰，況曰為君不慎德刑，其上明聞於天。”是為罪大不可赦。 \par}

王曰：“嗚呼！\CJKunderline{封}，敬哉！無作怨，勿用非謀、非彝，\footnote{言當修己以敬，無為可怨之事，勿用非善謀、非常法。}蔽時忱。丕則敏德，\footnote{斷行是誠道，大法敏德，信則人任焉，敏則有功。}用康乃心，顧乃德，遠乃猷，\footnote{用是誠道安汝心，顧省汝德，無令有非，遠汝謀，思為長久。}裕乃以。民寧，不汝瑕\xpinyin*{殄}。”\footnote{行寬政乃以民安，則我不汝罪過，不絕亡汝。}


{\noindent\zhuan\zihao{6}\fzbyks 傳“斷行”至“有功”。正義曰：以誠在於心，故決斷行之，亦心誠而行敏，為見事之速，事有善而須德法，故云“大法敏德”也。正以此二者，以“信則人任焉，敏則有功”故也。\CJKunderwave{論語}文。 \par}

{\noindent\zhuan\zihao{6}\fzbyks 傳“用是”至“長久”。正義曰：上文有“忱”有“敏”，此惟雲“用是誠道”，不雲“敏”者，“敏”在“誠”下，亦用之可知。 \par}

{\noindent\shu\zihao{5}\fzkt “王曰嗚呼”至“瑕殄”。正義曰：以罰不可失，故王命言曰：“嗚呼！封，當修己以敬哉！無為可怨之事，勿用非善謀、非常法，而以決斷行是誠信之道，大當法為機敏之德。用是信敏安汝心，顧省汝德，廣遠汝謀，能行寬政，乃以民安，則我不於汝罪過而絕亡汝。” \par}

王曰：“嗚呼！肆汝小子\CJKunderline{封}，惟命不於常，\footnote{以民安則不絕亡汝，故當念天命之不於常，汝行善則得之，行惡則失之。}汝念哉！無我享殄，\footnote{無絕棄我言而不念。}明乃服命，\footnote{享有國土，當明汝所服行之命令，使可則。}高乃聽，用康乂民。”\footnote{高汝聽，聽先王道德之言，以安治民。}


{\noindent\zhuan\zihao{6}\fzbyks 傳“享有”至“可則”。正義曰：以“不瑕殄”，即享有國土也。“服行之命”，謂德刑也。 \par}

{\noindent\shu\zihao{5}\fzkt “王曰嗚呼肆”至“乂民”。正義曰：與上相首引。王命言曰：“嗚呼！以民安則不汝絕亡之故，汝小子封，當念天命之不於常也。惟行善則得之，行惡則失之。汝念此無常哉！無絕棄我言而不念。若享有國土，當明汝服行之教令，使可法。高大汝所聽，用先王道德之言以安治民也。” \par}

王若曰：“往哉！\CJKunderline{封}，勿替敬典，\footnote{汝往之國，勿廢所宜敬之常法。}聽朕告汝，乃以殷民世享。”\footnote{順從我所告之言，即汝乃以殷民世世享國，福流後世。}

{\noindent\shu\zihao{5}\fzkt “王若”至“世享”。正義曰：以須高聽治民，故王命順其德而言曰：“汝往之國哉！封乎，勿廢所宜敬之常法，即聽用我誥是也。汝如此，則汝乃得以殷民世世享國。”而言不絕國祚，短長由德也。又言“王若曰”者，一篇終始言之，明於中亦有“若”也。 \par}

\section{酒誥第十二}


酒誥\footnote{\CJKunderline{康叔}監殷民。殷民化紂嗜酒，故以戒酒誥。○嗜，巿志反。}

{\noindent\shu\zihao{5}\fzkt 傳“\CJKunderline{康叔}”至“酒誥”。正義曰：以\CJKunderwave{梓材}雲“若茲監”，故云“\CJKunderline{康叔}監殷民”也。鄭以為“連屬之監，則為牧而言”，然\CJKunderline{康叔}時實為牧，而所戒為居殷墟，化紂餘民，不主於牧；下篇雲“監”，“監”亦指為君言之也，明“監”即國君監一國。故此言“監殷民”，不言“監一州”，若大宰之建牧立監也。 \par}

王若曰：“明大命於妹邦。\footnote{\CJKunderline{周公}以\CJKunderline{成王}命誥\CJKunderline{康叔}，順其事而言之，欲令明施大教命於妹國。妹，地名，紂所都朝歌以北是。○王若，馬本作“\CJKunderline{成王}若曰”，注云：“言\CJKunderline{成王}者，未聞也。俗儒以為\CJKunderline{成王}骨節始成，故曰\CJKunderline{成王}。或曰以\CJKunderline{成王}為少成二聖之功，生號曰\CJKunderline{成王}，沒因為諡。衛、賈以為戒成\CJKunderline{康叔}以慎酒，成就人之道也，故曰成。此三者吾無取焉。吾以為後錄\CJKunderwave{書}者加之，未敢專從，故曰未聞也。”妹邦，馬云：“妹邦，即牧養之地。”欲令，力呈反，下“始令”、“勿令”同。}乃穆考\CJKunderline{文王}，肇國在西土。\footnote{父昭子穆，\CJKunderline{文王}弟稱穆，將言始國在西土。西土，岐周之政。○\CJKunderline{文王}弟稱穆，周自後稷而封，為始祖，后稷生不窋為昭，鞠陶為穆，\CJKunderline{公劉}為昭，慶節為穆，皇僕為昭，羌弗為穆，毀揄為昭，公非為穆，高圉為昭，亞圉為穆，諸盩為昭，\CJKunderline{大王}為穆，\CJKunderline{王季}為昭，\CJKunderline{文王}為穆。故\CJKunderwave{左傳}宮之奇云：“大伯、虞仲，\CJKunderline{大王}之昭也。虢仲、虢叔，\CJKunderline{王季}之穆也。”又富辰云，管蔡已下十六國，文之昭也。昭一音韶。窋音竹律反。揄音投。盩音張流反。大並音太。}


{\noindent\zhuan\zihao{6}\fzbyks 傳“\CJKunderline{周公}”至“北是”。正義曰：此為下之目，故言“明施大教命於妹國”。此“妹”與“沬”一也，故沬為地名，紂所都朝歌以北。但妹為朝歌之所居也，朝歌近妹邑之南，故云“以北是”。\CJKunderwave{詩}又云“沬之東矣”，“沬之鄉矣”，即東與北為鄉也。妹屬鄘，紂所都在妹，又在北與東，是地不方平，偏在鄘多故也。馬、鄭、王本以文涉三家而有“成”字，\CJKunderline{鄭玄}雲“\CJKunderline{成王}所言，成道之王”，三家雲“王年長骨節成立”，皆為妄也。 \par}

{\noindent\zhuan\zihao{6}\fzbyks 傳“父昭”至“之政”。正義曰：以“穆”連“考”，故以昭穆言之。\CJKunderline{文王}廟次為穆，以周自後稷以至\CJKunderline{文王}十五世。案\CJKunderwave{世本}云：“后稷生不窋為昭，不窋生鞠陶為穆，鞠陶生\CJKunderline{公劉}為昭，\CJKunderline{公劉}生慶節為穆，慶節生皇僕為昭，皇僕生羌弗為穆，羌弗生毀榆為昭，毀榆生公飛為穆，公飛生高圉為昭，高圉生亞圉為穆，亞圉生組紺為昭，組紺生\CJKunderline{大王}亶父為穆，亶父生季歷為昭，季歷生\CJKunderline{文王}為穆。”據世次偶為穆也。\CJKunderwave{左傳}曰“大伯、虞仲，\CJKunderline{大王}之昭”，言\CJKunderline{大王}為穆，而子為昭。又曰“虢仲、虢叔，\CJKunderline{王季}之穆”，亦\CJKunderline{王季}為昭而子為穆，與\CJKunderline{文王}同穆也。又管、蔡、郕、霍等十六國亦曰\CJKunderline{文王}之昭，則以\CJKunderline{文王}為穆，其子與\CJKunderline{武王}為昭。又曰“邗晉應韓，武之穆”，以繼\CJKunderline{武王}為昭也。“將言始國在西土。西土，岐周之政”者，據今本先故言“始”，為初始為政，然則居豐前，故云“西土”，欲將言道\CJKunderline{文王}誥毖庶邦以下之政，故先本之雲“肇國在西土”。 \par}

厥誥毖庶邦、庶士越少正、御事，朝夕曰：‘祀茲酒。’\footnote{文王其所告慎眾國眾士於少正官、御治事吏，朝夕敕之：“惟祭祀而用此酒，不常飲。”○毖音秘。少，詩照反。}惟天降命，肇我民，惟元祀。\footnote{惟天下教命，始令我民知作酒者，惟為祭祀。○為，於偽反，下同。}天降威，我民用大亂喪德，亦罔非酒惟行。\footnote{天下威罰，使民亂德，亦無非以酒為行者。言酒本為祭祀，亦為亂行。○惟行，下孟反，注及下注“之行”同。}越小大邦用喪，亦罔非酒惟辜。\footnote{於小大之國所用喪亡，亦無不以酒為罪也。}

{\noindent\zhuan\zihao{6}\fzbyks 傳“\CJKunderline{文王}”至“常飲”。正義曰：告敕使之敬慎，故曰“告慎其眾國”，即眾多國君。“眾士”，朝臣也。既總呼為“士”，則卿大夫俱在內。少正、御治事以其卑賤，更別目之。“朝夕敕之”，丁寧慎之至也。 \par}

{\noindent\zhuan\zihao{6}\fzbyks 傳“惟天”至“祭祀”。正義曰：\CJKunderwave{世本}云，儀狄造酒，夏禹之臣，又云杜康造酒，則人自意所為。言“天下教命”者，以天非人，不因人為者，亦天之所使，故凡造立皆雲本之天。“元祀”者，言酒惟用於大祭祀，見戒酒之深也。顧氏云：“元，大也。\CJKunderwave{洛誥}‘稱秩元祀’，孔以為‘舉秩大祀’。”大劉以“元”為始，誤也。 \par}

{\noindent\zhuan\zihao{6}\fzbyks 傳“天下”至“亂行”。正義曰：民自飲酒致亂，以被威罰，言“天下威”者，亦如上言天之下教命，令民作酒也。為亂而罪，天理當然，故曰“天討有罪，五刑五用哉”。俗本雲“不為亂行”，定本雲“亦為亂行”，俗本誤也。 \par}

{\noindent\zhuan\zihao{6}\fzbyks 傳“於小”至“為罪也”。正義曰：“小大之國”，謂諸侯之國有小大也。上言“民用大亂”，指其身為罪。此言“邦用喪”，言其邦國喪滅。上文總謂貴賤之人，此則專指諸侯之身故也。惟行用酒，惟罪身得罪，亦互相通也。 \par}

{\noindent\shu\zihao{5}\fzkt “王若”至“惟辜”。正義曰：\CJKunderline{周公}以王命誥\CJKunderline{康叔}，順其事而言曰：“汝當明施大教命於妹國而戒之以酒。所以須戒酒者，以汝父於廟以穆考\CJKunderline{文王}，始國在西土岐周為政也。其誥慎所職眾國眾士於少正官、御治事吏，朝夕敕之曰：‘惟祭祀而用此酒，不常為飲也。’所以不常為飲者，以惟天之下教命，始令我民知作酒者，惟為大祭祀，故以酒為祭，不主飲。故天下威罰於我民，用使之大為亂，以喪其德，亦無非以酒為行而用之。故於小大之國，用使之喪亡，亦無非以酒為罪，以此眾事少正，皆須戒酒也。是\CJKunderline{文王}以酒為重戒，汝不可不法也。” \par}

\CJKunderline{文王}誥教小子有正有事，無彝酒。\footnote{小子，民之子孫也。正官治事，謂下群吏。教之皆無常飲酒。}越庶國，飲惟祀，德將無醉。\footnote{於所治眾國，飲酒惟當因祭祀，以德自將，無令至醉。}惟曰我民迪，小子惟土物愛，厥心臧。\footnote{文王化我民，教道子孫，惟土地所生之物皆愛惜之，則其心善。}聰聽祖考之彝訓，越小大德，小子惟一。\footnote{言子孫皆聰聽父祖之常教，於小大之人皆念德，則子孫惟專一。}


{\noindent\zhuan\zihao{6}\fzbyks 傳“小子”至“飲酒”。正義曰：知“小子”謂民之子孫者，以下文二“我民迪小子”，又云“奔走事厥考厥長”，故知“小子”謂民之子孫也。知“有正有事”非士大夫，而云“正官治事，謂下群吏”者，以文與“小子”相連，故知是正官下治事之群吏。 \par}

{\noindent\zhuan\zihao{6}\fzbyks 傳“於所”至“至醉”。正義曰：以述上文內外雙舉，此為小子及民與士大夫可知。其外宜有國君，故下雲指戒\CJKunderline{康叔}為國之事，故總言“眾國”。惟於祭祀得飲酒，猶以德自將，無令至醉。\CJKunderwave{大傳}因此言“宗室將有事，族人皆入侍”，得有醉與不醉而出與不出之事。而以德自將，無令至醉，亦一隅之驗。\CJKunderline{文王}為諸侯而云“眾國”者，\CJKunderline{文王}為西伯，又三分有二諸侯，故得戒眾國也。 \par}

{\noindent\zhuan\zihao{6}\fzbyks 傳“\CJKunderline{文王}”至“心善”。正義曰：以“惟曰”為教辭，故言“\CJKunderline{文王}化我”。民愛惜土物而不損耗，則不嗜酒，故心善。 \par}

{\noindent\shu\zihao{5}\fzkt “\CJKunderline{文王}”至“惟一”。正義曰：前\CJKunderline{文王}戒酒，以為所供當重飲之，則有滅亡之害。此更戒之，令以德自將，不可常飲。故又云，\CJKunderline{文王}誥教其民之小子與正官之下有職事之人。謂群吏。汝等無得常飲酒也。於所治眾國之君臣民眾等，言飲酒惟當因祭祀，以德自將，無令至醉。又自申\CJKunderline{文王}之教小子者，不但身自教之，又化民使自教其子弟。惟教其民曰：“惟我民等，當教道子孫小子，令土地所生之物，皆愛惜之，則其心善矣。”以愛物，則不為酒而損耗故也。既父祖稟\CJKunderline{文王}之教以化其子孫，而子孫能聰審聽用祖考之常訓。言愛物以戒酒也。不但民之小子為然，其於小大德之士大夫等，亦皆能念行\CJKunderline{文王}之德以教其子孫，故子孫亦聰聽之。小子惟皆專一而戒其酒，其民及在位，不問貴賤，子孫皆化，則至成長為德可知也。 \par}

妹土嗣爾股肱純，其藝黍稷，奔走事厥考厥長。\footnote{今往當使妹土之人繼汝股肱之教，為純一之行，其當勤種黍稷，奔走事其父兄。○長，丁丈反，下注“長官諸侯”之長同。}肇牽車牛，遠服賈，用孝養厥父母。\footnote{農功既畢，始牽車牛，載其所有，求易所無，遠行賈賣，用其所得珍異孝養其父母。○賈音古。養，牛亮反。}厥父母慶，自洗腆，致用酒。\footnote{其父母善子之行，子乃自絜厚，致用酒養也。}庶士、有正越庶伯、君子，其爾典聽朕教。\footnote{眾伯君子、長官大夫、統庶士有正者，其汝常聽我教，勿違犯。}


{\noindent\zhuan\zihao{6}\fzbyks 傳“今往”至“父兄”。正義曰：以妹土為所封之都，故言“今往”。“繼汝股肱之教”者，君為元首，臣作股肱，君倡臣行，施由股肱，故言繼其教也。言“奔走”者，顧氏云：“勤種黍稷，奔馳趨走也。” \par}

{\noindent\zhuan\zihao{6}\fzbyks 傳“農功”至“父母”。正義曰：若當農功，則有所廢，故知既畢乃行，故云“始牽車牛”，即牽將大車，載有易無，遠求盈利，所得珍異而本不損，故可孝養其父母，亦愛土物之義也。 \par}

{\noindent\zhuan\zihao{6}\fzbyks 傳“其父”至“酒養”。正義曰：以人父母欲家生之富者，若非盈利，雖得其養，有喪家資，則父母所不善。今勤商得利，富而得養，所以善子之行也。 \par}

{\noindent\zhuan\zihao{6}\fzbyks 傳“眾伯”至“違犯”。正義曰：眾伯君子，統眾士有正者，經雲“庶士有正”者，戒其慎酒，從卑至尊，故先教子孫，乃及庶士眾百君子。 \par}

爾大克羞耇，惟君，爾乃飲食醉飽。\footnote{汝大能進老成人之道，則為君矣。如此汝乃飲食醉飽之道。先戒群吏以聽教，次戒\CJKunderline{康叔}以君義。}丕惟曰，爾克永觀省，作稽中德。\footnote{我大惟教汝曰，汝能長觀省古道，為考中正之德，則君道成矣。○省，悉井反。}爾尚克羞饋祀，爾乃自介用逸。\footnote{能考中德，則汝庶幾能進饋祀於祖考矣。能進饋祀，則汝乃能自大用逸之道。}茲乃允惟王正事之臣，\footnote{汝能以進老成人為醉飽，考中德為用逸，則此乃信任王者正事之大臣。○任音壬。}茲亦惟天若元德，永不忘在王家。”\footnote{言此非但正事之臣，亦惟天順其大德而佑之，長不見忘在王家。}

{\noindent\zhuan\zihao{6}\fzbyks 傳“汝大”至“君義”。正義曰：\CJKunderwave{釋詁}云：“羞，進也。”既以慎酒立教，是大能進行老成人之道，是惟可為人君矣。以人君若治不得，有所民事可憂，雖得酒食，不能醉飽。若能進德，民事可乎，故為飲食可醉飽之道。以群臣言，“聽教”即為臣義，不過慎酒進德，次戒\CJKunderline{康叔}以君義，亦有“聽教”，明為互矣。 \par}

{\noindent\zhuan\zihao{6}\fzbyks 傳“我大”至“成矣”。正義曰：以言“曰”，故以為教辭，即教以“大克羞者”。長省古道，是老成人之德，考其中正，是能大進行，可以惟為君，故云“則君道成矣”。 \par}

{\noindent\zhuan\zihao{6}\fzbyks 傳“能考”至“之道”。正義曰：以聖人為能饗帝，孝子為能饗親，考德為君，則人治之，已成民事，可以祭神，故考中德，能進饋祀於祖考。人愛神助，可以無為，故大用逸之道，即上雲“飲食醉飽之道”也。鄭以為助祭於君，亦非其義勢也。以下然並亦惟天據人事，是惟王正事大臣，本天理，故天順其大德，不見忘在於王家，反覆相成之勢也。 \par}

{\noindent\shu\zihao{5}\fzkt “妹土”至“王家”。正義曰：既上言\CJKunderline{文王}之教，今指戒\CJKunderline{康叔}之身，實如汝當法，\CJKunderline{文王}斷酒之法故今往當使妹土之人繼爾股肱之，教為純一之行。其當勤於耕種黍稷，奔馳趨走供事其父與兄。其農功既畢，始牽車牛遠行賈賣，用其所得珍異孝養其父母，以子如此，善子之行，子乃自洗潔，謹敬厚致用酒以養，此亦小子土物愛也。又謂汝眾士有正之人，及於眾伯君子長官大夫統眾士有正者，其汝亦常聽用我斷酒之教，勿違犯也。汝\CJKunderline{康叔}大能進行老成人之道，則惟可為君矣。如此汝乃為飲食醉飽之道。由須進行老成人，故我大惟教汝曰：“汝能長觀省古道，所為考行中正之德，即是進行老成人，惟堪為君。能考中德，用汝庶幾能進饋祀於祖考矣。以能進饋祀，人神所助，則汝乃能自大用逸之道。如此用逸，則乃信惟王正事之大臣。不但正事大臣，如此亦惟天順其大德而佑助之，長不見遺忘在王家矣。可不務乎？” \par}

王曰：“\CJKunderline{封}，我西土棐徂邦君御事，小子尚克用\CJKunderline{文王}教，不腆於酒。\footnote{我\CJKunderline{文王}在西土，輔訓往日國君及御治事者、下民子孫，皆庶幾能用上教，不厚於酒。言不常飲。}故我至於今，克受殷之命。”\footnote{以不厚於酒，故我周家至於今能受殷王之命。}


{\noindent\zhuan\zihao{6}\fzbyks 傳“我文”至“常飲”。正義曰：“棐”，輔也。“徂”，往也。以事已過，故言“往日”。恐嗜酒不成其德，故以斷酒輔成之。其“御事”謂國君之下眾臣也。“不厚於酒”即“無彝酒”也，故云“不常飲”，總述上也。 \par}

{\noindent\shu\zihao{5}\fzkt “王曰封我西”至“之命”。正義曰：於此乃總言不可不用\CJKunderline{文王}慎酒之教，王命之曰：“封，我\CJKunderline{文王}本在西土，以道輔訓往日國君及治事之臣大夫士與其民之小子，其此等皆庶幾能用\CJKunderline{文王}教，而不厚於酒。故我周家至於今能受殷之王命。以此故，不可不用其教以斷酒。” \par}

王曰：“\CJKunderline{封}，我聞惟曰，在昔殷先哲王，迪畏天顯小民，\footnote{聞之於古。殷先智王，謂湯蹈道畏天，明著小民。}經德秉哲，自\CJKunderline{成湯}咸至於\CJKunderline{帝乙}，成王畏相。\footnote{能常德持智，從湯至\CJKunderline{帝乙}中間之王猶保成其王道，畏敬輔相之臣，不敢為非。○相，息亮反，下同。}惟御事厥棐有恭，不敢自暇自逸，\footnote{惟殷御治事之臣，其輔佐畏相之君，有恭敬之德，不敢自寬暇，自逸豫。○暇，遐嫁反。}矧曰其敢崇飲？\footnote{崇，聚也。自假自逸猶不敢，況敢聚會飲酒乎？明無也。}越在外服，侯、甸、男、衛邦伯，\footnote{於在外國侯服、甸服、男服、衛服國伯諸侯之長。言皆化湯畏相之德。}

{\noindent\zhuan\zihao{6}\fzbyks 傳“聞之”至“小民”。正義曰：言“聞之於古”，是事明眾見也。下言“自\CJKunderline{成湯}”，知此別道湯事也。王者上承天，下恤民，皆由蹈行於為，畏天之罰已故也。又以道教民，故明德著小民。 \par}

{\noindent\zhuan\zihao{6}\fzbyks 傳“能常”至“為非”。正義曰：德在於身，智在於心，故能常德持智，即上迪畏天，顯小民，為自湯後皆爾。 \par}

{\noindent\zhuan\zihao{6}\fzbyks 傳“惟殷”至“逸豫”。正義曰：此事當公卿，故下別雲“越在內服百僚庶尹”也。為君畏相，故輔之。若寬暇與逸豫，則不恭敬，故不敢為也。 \par}

{\noindent\zhuan\zihao{6}\fzbyks 傳“崇聚”至“明無”。正義曰：\CJKunderwave{釋詁}云：“崇，充也。”充實則集聚，故“崇”為聚也。飲必待暇逸，猶尚不敢暇逸，故言“況敢聚集飲酒乎？明無也”。 \par}

{\noindent\zhuan\zihao{6}\fzbyks 傳“於在”至“之德”。正義曰：以公卿與國為體，承君共事，故先言之。然後見廣，故自外及內，舉四者以總六服，又因“衛”為蕃衛，故不言“採”也。“國”謂國君，“伯”言長，連、屬、卒、牧皆是，見遍在外為君，故言“化湯畏相之德”。 \par}

越在內服，百僚、庶尹、惟亞、惟服、宗工，\footnote{於在內服治事百官眾正及次大夫服事尊官，亦不自逸。}越百姓里居，\footnote{於百官族姓及卿大夫致仕居田裡者。}罔敢湎於酒。不惟不敢，亦不暇。\footnote{自外服至里居，皆無敢沈湎於酒。非徒不敢，志在助君敬法，亦不暇飲酒。○湎，面善反。}惟助\CJKunderline{成王}德顯，越尹人祇闢。\footnote{所以不暇飲酒，惟助其君\CJKunderline{成王}道，明其德於正人之道，必正身敬法，其身正，不令而行。闢，扶亦反。}

{\noindent\zhuan\zihao{6}\fzbyks 傳“於在”至“自逸”。正義曰：畿外有服數，畿內無服數，故為“服治事”也。言“百官眾正”，為總之文。但百官眾正除六卿亦有大夫及士，士亦有官首而為政者。“惟亞”，傳雲“次大夫”者，謂雖為大夫不為官首者，亞次官首，故云“亞”。舉大夫尊者為言，其實士亦為亞次之官。必知“惟亞”兼士者，以此經文上下更無別見士之文，故知兼之。“惟服宗工”，總上“百僚庶尹”及“惟亞”，言服治職事尊官之故，亦不自逸。“惟亞”雖不為官首，亦助上服治政事，或可非官首者服事在上之尊官，亦不自逸。 \par}

{\noindent\zhuan\zihao{6}\fzbyks 傳“於百”至“裡者”。正義曰：每言“於”者，繼上君與御事為“於”。此不言“在”，從上“內服”故也。“百官族姓”謂其每官之族姓，而與“里居”為總，故云“卿大夫致仕居田裡者”也。 \par}

{\noindent\zhuan\zihao{6}\fzbyks 傳“自外”至“飲酒”。正義曰：自外服至里居，皆無敢沈湎，亦上御事，雲“亦不暇”，不暇則不逸可知，助君敬法，逆探下經也。 \par}

{\noindent\shu\zihao{5}\fzkt “王曰封我聞”至“祇闢”。正義曰：以周受於殷，\CJKunderline{文王}之前殷代也，今又衛居殷地，故舉殷代以酒興亡得失而為戒。王命之曰：“封，我聞於古，所聞惟曰，殷之先代智道之王\CJKunderline{成湯}，於上蹈道以畏天威，於下明著加於小民，即能常德持智以為政教。自\CJKunderline{成湯}之後皆然，以至於\CJKunderline{帝乙}，猶保成其王道，畏敬輔相之臣。其君既然，惟殷御治事之臣，其輔相於君，有恭敬之德，不敢自寬暇，自逸豫，況曰其敢聚會群飲酒乎？於是在外之服侯、甸、男、衛、國君之長，於是在內之服治事百官眾正惟次大夫惟服事尊官，於百官族姓及致仕在田裡而居者，皆無敢沈湎於酒。不惟不敢，亦自不暇飲。所以不暇者，惟以助其君成其王道，令德顯明；又於正人之道，必正身敬法，正身以化下，不令而行，故不暇飲。是亦可以為法也。” \par}

我聞亦惟曰，在今後嗣王酣身，\footnote{嗣王，紂也。酣樂其身，不憂政事。○酣，戶甘反。樂音洛。}厥命罔顯於民，祇保越怨不易。\footnote{言紂暴虐，施其政令於民，無顯明之德，所敬所安，皆在於怨，不可變易。○易如字，馬以豉反。}誕惟厥縱淫泆於非彝，用燕喪威儀，民罔不\xpinyin*{衋}傷心。\footnote{紂大惟其縱淫泆於非常，用燕安喪其威儀，民無不衋然痛傷其心。○縱，子用反，注同。泆音溢，又作逸，亦作佚。衋，許力反。}惟荒腆於酒，不惟自息乃逸，\footnote{言紂大厚於酒，晝夜不念自息，乃過差。○差，初佳反，又初賣反。}厥心疾很，不克畏死。\footnote{紂疾很其心，不能畏死。言無忌憚。○很，胡懇反。}辜在商邑，越殷國滅無罹。\footnote{紂聚罪人在都邑而任之，於殷國滅亡無憂懼。}弗惟德馨香，祀登聞於天，誕惟民怨。\footnote{紂不念發聞其德，使祀見享，升聞於天，大行淫虐，惟為民所怨咎。}庶群自酒，腥聞在上，故天降喪於殷，罔愛於殷，惟逸。\footnote{紂眾群臣用酒沈荒，腥穢聞在上天，故天下喪亡於殷，無愛於殷，惟以紂奢逸故。○聞音問。}天非虐，惟民自速辜。”\footnote{言凡為天所亡，天非虐民，惟民行惡自召罪。}


{\noindent\zhuan\zihao{6}\fzbyks 傳“言紂”至“變易”。正義曰：“施其政令於民，無顯明之德”，言所施者皆是闇亂之政也。紂意謂之為善，所敬之所安之者，及其施行，皆是害民之事，為民所怨。紂之為惡，執心堅固，不可變易也。 \par}

{\noindent\zhuan\zihao{6}\fzbyks 傳“紂大”至“其心”。正義曰：“誕”訓為大，言紂大惟其縱淫泆於非常之事。 \par}

{\noindent\zhuan\zihao{6}\fzbyks 傳“紂眾”至“逸故”。正義曰：“紂眾群臣用酒沈荒”，“用”者解經之“自”。定本作“自”，俗本多誤為“嗜”。 \par}

{\noindent\zhuan\zihao{6}\fzbyks 傳“言凡”至“召罪”。正義曰：此言“惟人”，謂紂也。今變言“人”者，見雖非紂亦然。 \par}

{\noindent\shu\zihao{5}\fzkt “我聞”至“速辜”。正義曰：既言\CJKunderline{帝乙}以上慎酒以存，故又言紂嗜酒而滅：“我聞亦惟曰，殷之在今\CJKunderline{帝乙}後嗣之謂紂王，酣樂其身，不憂於政事，施其政令，無顯明之德於民，所敬所安，皆在於怨，不可變易。大惟其縱淫泆於非常，用燕安之故，喪其威儀，民見之無不衋然痛傷其心也。皆由惟大愛厚於酒，晝夜不念自止息，乃過逸。其內心疾害很戾，不能畏死。聚罪人在商邑而任之，於殷國滅亡無憂懼也。不念發聞其德令之馨香，使祀見享，升聞於天，大惟行其淫虐，為民下所怨。紂眾群臣集聚用酒荒淫，腥穢聞在上天，故天下喪亡於殷，無愛念於殷，惟以紂奢逸故。非天虐殷以滅之，惟紂為人自召此罪故也。” \par}

王曰：“\CJKunderline{封}，予不惟若茲多誥。\footnote{我不惟若此多誥汝，我親行之。}古人有言曰：‘人無於水監，當於民監。’\footnote{古賢聖有言，人無於水監，當於民監。視水見己形，視民行事見吉凶。○監，工陷反，下及注同。}今惟殷墜厥命，我其可不大監撫於時？\footnote{今惟殷紂無道，墜失天命，我其可不大視此為戒，撫安天下於是？}

{\noindent\shu\zihao{5}\fzkt “王曰封予”至“於時”。正義曰：既陳殷之戒酒與嗜酒以致興亡之異，故誥之，王命言曰：“封，我不惟若此徒多出言以誥汝而已，我自戒酒，己親行之，汝可法之也。所以親行者，古人有言曰：‘人無於水監，當於民監。’以水監但見己形，以民監知成敗故也。以須民監之故，今殷紂無道，墜失其天命，我其可不大視以為戒，撫安天下於今時也？” \par}

予惟曰，汝劼毖殷獻臣，\footnote{劼，固也。我惟告汝曰，汝當固慎殷之善臣信用之。○劼，若八反。}侯、甸、男、衛，矧太史友，內史友、\footnote{侯、甸、男、衛之國當慎接之，況太史、內史掌國典法所賓友乎？}越獻臣、百宗工，矧惟爾事服休、服採？\footnote{於善臣百尊官不可不慎，況汝身事服行美道，服事治民乎？}



{\noindent\zhuan\zihao{6}\fzbyks 傳“劼固”至“用之”。正義曰：“劼，固”，\CJKunderwave{釋詁}文。將欲斷酒為重，故節文以相況。“毖”訓為慎，言誠堅固謹慎，皆敬而釋任之，其文通於下，皆固慎。 \par}

{\noindent\zhuan\zihao{6}\fzbyks 傳“侯甸”至“賓友乎”。正義曰：太史掌國六典，依\CJKunderwave{周禮}，治典、教典、禮典、政典、刑典、事典也。內史掌八柄之法者，爵、祿、廢、置、殺、生、與、奪。此“太史”、“內史”即\CJKunderline{康叔}之國大夫，知者，以下“圻父”、“農父”、“宏父”是諸侯之三卿，明“太史”、“內史”非王朝之官。“所賓友”者，敬也。 \par}

{\noindent\zhuan\zihao{6}\fzbyks 傳“於善”至“民乎”。正義曰：“於善臣”即上經“殷獻臣”也。“百尊官”即上“侯甸男衛”、“太史”、“內史”也。“服行美道，服事治民”即上汝之身事。知“服事”是治民者，民惟邦本，諸侯治民為事故也。\CJKunderline{鄭玄}以“服休”為燕息之近臣，“服採”為朝祭之近臣，非孔意也。 \par}

矧惟若疇圻父，薄違農父？\footnote{圻父，司馬。農父，司徒。身事且宜敬慎，況所順疇諮之司馬乎？況能迫回萬民之司徒乎？言任大。○圻，臣依反。父音甫。薄，蒲各反，徐又扶各反。違如字，徐音回，馬云：“違行也。”}若保宏父，定闢，矧汝剛制於酒？\footnote{宏，大也。宏父，司空。當順安之。司馬、司徒、司空，列國諸侯三卿，慎擇其人而任之，則君道定，況汝剛斷於酒乎？○闢，必亦反。斷，丁亂反。}

{\noindent\zhuan\zihao{6}\fzbyks 傳“圻父”至“任大”。正義曰：司馬主圻封，故云“圻父”。“父”者，尊之辭。以司徒教民五土之藝，故言“農父”也。以司馬征伐在乎閫外所專，故隨順而疇諮之，言君所順疇也。迫近迴繞於萬民，言近民事也。二者皆任大。 \par}

{\noindent\zhuan\zihao{6}\fzbyks 傳“宏大”至“酒乎”。正義曰：“宏，大”，\CJKunderwave{釋詁}文。以司空亦君所順所安和之，故言“當順安之”。諸侯之三卿，以上有司馬、司徒，故知“宏父”是司空。言大父者，以營造為廣大國家之父。因節文而分之，乃總之言“司馬、司徒、司空”。列國三卿，令慎擇其人而任之，則君道定，況剛斷於酒乎？為甚之義也。其“定闢”總上自“劼毖殷獻”已下，獨言三卿者，因文相況而接之，其實總上也。三卿不次者，以司馬征伐為重；次以政教安萬民，司徒為重；司空直指營造，故在下也。司徒言於萬民為迫回者，事務為主故也。司徒不言“若”者，互相明，皆為治民，而君所順也。 \par}

{\noindent\shu\zihao{5}\fzkt “予惟”至“於酒”。正義曰：殷之存亡既可以為監若是，故我惟告汝曰：“汝當堅固愛慎殷之善臣及侯、甸、男、衛之君，則在外尚然，況已下太史所賓友，內史所賓友，於善臣百尊官而不固慎乎？此之卑官猶尚固慎，況惟汝之身事所服行美道，服行美事治民，而可不固慎乎？於己身事猶當固慎，況惟所敬順疇諮之圻父，能迫回萬民之農父，所順所安之宏父？此等大臣能得固慎，則可定其為君之道，固慎大臣，雖非急要，尚能使君道得定，況汝又能剛斷於酒乎？善所莫大，不可加也。” \par}

厥或誥曰：‘群飲。’汝勿佚。\footnote{具有誥汝曰：“民群聚飲酒。”不用上命，則汝收捕之，勿令失也。}盡執拘以歸於周，予其殺。\footnote{盡執拘群飲酒者以歸於京師，我其擇罪重者而殺之。○盡，子忍反。}又惟殷之迪諸臣惟工，乃湎於酒，勿庸殺之，\footnote{又惟殷家蹈惡俗諸臣，惟眾官化紂日久，乃沈湎於酒，勿用法殺之。○惡，烏各反。}姑惟教之，有斯明享。\footnote{以其漸染惡俗，故必三申法令，且惟教之，則汝有此明訓以享國。○三，息暫反，又如字。}乃不用我教辭，惟我一人弗恤，弗蠲乃事，時同於殺。”\footnote{汝若忽怠不用我教辭，惟我一人不憂汝，乃不潔汝政事，是汝同於見殺之罪。}


{\noindent\zhuan\zihao{6}\fzbyks 傳“盡執”至“殺之”。正義曰：言“周”，故為“京師”。但飲有稀數，罪有大小，不可一皆盡殺，故知“擇罪重者殺之”。 \par}

{\noindent\zhuan\zihao{6}\fzbyks 傳“又惟”至“殺之”。正義曰：言“諸臣”，謂尊者，及其下列職眾官，不可用法殺之，明法有張弛。此由殷之諸臣，漸染紂之惡俗日久，故不可即殺。其衛國之民，先非紂之舊臣，乃群聚飲酒，恐增長昏亂，故擇罪重者殺之。據意不同，故殺否有異。 \par}

{\noindent\zhuan\zihao{6}\fzbyks 傳“以其”至“享國”。正義曰：禮成於三，故必三申法令。“有此明訓”，總上之辭，故得享國。 \par}

{\noindent\zhuan\zihao{6}\fzbyks 傳“汝若”至“之罪”。正義曰：汝不用我教辭，則不足憂念，故“惟我一人不憂汝”。“不絜汝之政事”，事惟穢惡，不復教之使絜靜也。 \par}

{\noindent\shu\zihao{5}\fzkt “厥或”至“於殺”。正義曰：以為政莫重於斷酒，故其有人誥汝曰：“民今飲酒，相與群聚。”是不用上命，則汝收捕之，勿令失矣。盡執拘以歸於周之京師，我其擇罪重而殺之也。又惟殷之蹈惡俗諸臣，惟其眾官化紂日久，乃沈湎於酒，勿用法殺之。以漸染惡俗，故三申法令，且惟教之，則汝有此明訓，可以享國。汝若不用我教辭，惟我一人天子不憂汝，不潔汝政事，是汝同於見殺之罪，不可不慎。 \par}

王曰：“\CJKunderline{封}，汝典聽朕毖，\footnote{汝當常聽念我所慎而篤行之。}勿辯乃司民湎於酒。”\footnote{辯，使也。勿使汝主民之吏湎於酒。言當正身以帥民。}

{\noindent\shu\zihao{5}\fzkt “王曰封”至“於酒”。正義曰：以戒酒事終，故結之。王命言曰：“封，汝當常聽念我所使汝慎者，篤而行之。勿使汝主民之吏若宰人者沈湎於酒，當正身以帥民。” \par}

\section{梓材第十三}


梓材\footnote{告\CJKunderline{康叔}以為政之道，亦如梓人治材。○梓音子，本亦作杼，馬云：“古作梓字。治木器曰梓，治土器曰陶，治金器曰冶。”}

{\noindent\shu\zihao{5}\fzkt 傳“告康”至“治材”。正義曰：此取下言“若作梓材，既勤樸斫”，故云“為政之道，如梓人治材”。此古“杍”字，今文作“梓”。“梓”，木名，木之善者，治之宜精，因以為木之工匠之名。下有“稽田”、“作室”，乃言“梓材”，三種獨用“梓材”者，雖三者同喻，田在於外，室總於家，猶非指事之器，故取“梓材”以為功也。因戒德刑與酒事終，言治人似治器而結之故也。 \par}

王曰：“\CJKunderline{封}，以厥庶民暨厥臣達大家，\footnote{言當用其眾人之賢者與其小臣之良者，以通達卿大夫及都家之政於國。○暨，其器反。}以厥臣達王惟邦君。\footnote{汝當信用其臣以通王教於民。言通民事於國，通王教於民，惟乃國君之道。}汝若恆越曰：‘我有師師、\footnote{汝惟君道使順常，於是曰：“我有典常之師可師法。”}司徒、司馬、司空、尹旅曰：‘予罔厲殺人。’\footnote{言國之三卿、正官眾大夫皆順典常，而曰：“我無厲虐殺人之事。”如此則善矣。}


{\noindent\zhuan\zihao{6}\fzbyks 傳“言當”至“於國”。正義曰：“以”，用也。“暨”，與也。言“用”，通“厥臣”可用，明此皆賢與良也。“厥臣”文在“大家”之上，故知“小臣”也。言用之者，既用其言以為政，又用其人以為輔，本之得大家所用統之，即君所遣也。以大夫稱家，對士庶有家而非大，故云“大家”，卿大夫在朝者。“都家”亦卿大夫所得邑也，又公邑而大夫所治亦是也。用此以行政令，上達於國，使人君知之也。即是庶人升為士，又用庶人進在官者，小臣亦得進等而用之。\CJKunderwave{周禮}有都家之官，鄭云：“都謂王子弟所封及公卿所食邑，家謂大夫所食采地。”傳以“大家”言之，總包大臣，故言“卿大夫及都家之政”。卿大夫之政謂在朝所掌者，都家之政謂采邑所有政事，二者並當通達之於國，故連言之。 \par}

{\noindent\zhuan\zihao{6}\fzbyks 傳“汝當”至“之道”。正義曰：言汝當信用臣，即信用卿大夫及都家，自然大家也。傳用小臣與庶人，故得“通王教於民”也。人君上承於王，下治民事，故交通其政，“惟乃國君之道”而已。鄭以“於邑言達大家，於國言達王與邦君，王為二王之後”，即亂名實也。 \par}

{\noindent\zhuan\zihao{6}\fzbyks 傳“汝惟”至“師法”。正義曰：即上民事王教通於國人，是順常也，故總上“惟邦君”，言“汝惟君道使順常”也。“典常可師”即順常也。 \par}

{\noindent\zhuan\zihao{6}\fzbyks 傳“言國”至“善矣”。正義曰：此連上蒙“若恆”之文，故云“國之三卿、正官眾大夫皆順典常”也。不言“士”，從可知也。此曰“予罔厲殺人”，所謂令\CJKunderline{康叔}之語，但在臣下，宜為此也。以上令下行，行之在臣，故云“我無厲虐殺人之事”，互明君及臣皆師法而無虐。 \par}

亦厥君先敬勞，肆徂厥敬勞。\footnote{亦其為君之道，當先敬勞民，故汝往治民，必敬勞來之。○勞，力報反，下同。來，力代反。}肆往，奸宄、殺人、歷人，宥。\footnote{以民當敬勞之故，汝往之國，又當詳察奸宄之人及殺人賊，所過歷之人，有所寬宥，亦所以敬勞之。○宄音軌。}肆亦見厥君事，戕敗人，宥。\footnote{聽訟折獄，當務從寬恕，故往治民，亦當見其為君之事，察民以過誤殘敗人者，當寬宥之。○見如字，徐賢遍反。戕敗，徐在羊反，又七良反，馬云：“殘也。”折，之舌反。}

{\noindent\zhuan\zihao{6}\fzbyks 傳“亦其”至“來之”。正義曰：“亦其為君之道”者，為邦君之道，非直順常，亦須敬勞，故往必敬勞，即\CJKunderwave{論語}雲“先之，勞之”是也。 \par}

{\noindent\zhuan\zihao{6}\fzbyks 傳“以民”至“勞之”。正義曰：上文無罪敬勞，此惟就有罪者原情免宥，亦敬勞也。其實“奸宄”不殺人者，“殺人”亦是奸宄，但重言而別其文。奸宄及殺人，二者並是賊害，自當合罪，不可寬宥。其所過歷之人，情所不知，故詳察寬宥，以為敬勞之。 \par}

{\noindent\zhuan\zihao{6}\fzbyks 傳“聽訟”至“宥之”。正義曰：以君者立於無過之地，使物不失其所，故宥罪原情，當見其為君之事，與上“厥君”始終相承。於“奸”上言“肆往”，此亦以罪事往可知也。言“宥”，明情亦可原，故知“過誤殘敗人”也。 \par}

{\noindent\shu\zihao{5}\fzkt “王曰”至“人宥”。正義曰：王曰：“封，汝為政，當用其眾人之賢者與其小臣之良者，以通達卿大夫及都家等大家之政於國，然後汝當信用其臣以通達王教於民，惟乃可為國君之道。汝為君道，故當使上下順常，於是曰：‘我有典常之師可師法。’是君之順典常也。其下司徒、司馬、司空國之三卿，及正官眾大夫亦皆順典常，而曰：‘我無虐厲殺人之事。’是使臣之順常也。如此君臣皆能順常，則為善矣。為君之道，非但順常，亦須敬勞之。故云亦其為君之道，當先敬心以愛勞民。故汝往治民，必敬勞之。又以民須敬勞之故，汝往之國，詳察其奸宄及殺人之人，二者所過歷之人，原情不知，有所寬宥。以斷獄務從寬，故汝往治亦當見其為君之事，而民有過誤殘敗人者，當寬宥之，此亦為敬勞之也。” \par}

王啟監，厥亂為民。\footnote{言王者開置監官，其治為民，不可不勉。○監，工暫反，劉工銜反，下同。為，於偽反，注同。治，直吏反。}曰：‘無胥戕，無胥虐，至於敬寡，至於屬婦，合由以容。’\footnote{當教民無得相殘傷，相虐殺，至於敬養寡弱，至於存恤妾婦，和合其教，用大道以容之，無令見冤枉。○屬婦，上音蜀，妾之事妻也。令，力呈反，篇末同。冤，紆元反，一本作以冤。}王其效邦君，越御事，厥命曷以？\footnote{王者其效實國君，及於御治事者，知其教命所施何用，不可不勤。}‘引養引恬’，自古王若茲監，罔攸闢。”\footnote{能長養民，長安民，用古王道如此，監無所復罪，當務之。○恬，田廉反。闢，扶亦反。}

{\noindent\zhuan\zihao{6}\fzbyks 傳“當教”至“冤枉”。正義曰：以言“曰”，故知“當教民”也。“殘”謂不死，“虐”，甚則殺，故二文也。經言“屬婦”，傳言“妾婦”者，以妾屬於人，故名“屬婦”。此經“屬婦”與“寡弱”為例，則非關嫡婦也。何者？妻子是家中之貴者，不至冤枉故也。 \par}

{\noindent\zhuan\zihao{6}\fzbyks 傳“王者”至“不勤”。正義曰：以君臣共國事，故並效御治事，而知其所施，則下不得為非，即是王使存省侯伯監治是也，故不可不勤。 \par}

{\noindent\shu\zihao{5}\fzkt “王啟”至“攸闢”。正義曰：\CJKunderline{周公}云：“所以敬勞者，以王者開置監官，其治主為於民故也。以此當教民曰：‘無得相殘傷，無得相虐殺，而為重害也。何但不可為重害，民之相於，當至於敬養寡弱，至於存恤屬婦，合和其教，用大道以相容，無使至冤枉。’所以如此者，以王者其當效實國君，及於御治事者，惟須知其教命所施何用，知其善惡，故不可不勤也。所效實若能長養民，長安民，用古昔明王之，道而治之如此為監，無所復罪，汝當務之。” \par}

“惟曰，若稽田，既勤敷菑，惟其陳修，為厥疆畎。\footnote{言為君監民，惟若農夫之考田，已勞力布發之，惟其陳列修治，為其疆畔畎壟，然後功成。以喻教化。○菑，側其反。畎，工犬反。}若作室家，既勤垣墉，惟其塗塈茨。\footnote{如人為室，家已勤立垣牆，惟其當塗既茨蓋之。○垣音袁。墉音庸，馬云：“卑曰垣，高曰庸。”塈，徐許既反，\CJKunderwave{說文}云：“仰塗也。”\CJKunderwave{廣雅}云：“塗也。”馬云：“堊色。”一音故愛反。茨，徐在私反。}若作梓材，既勤樸斲，惟其塗丹\xpinyin*{雘}。\footnote{為政之術，如梓人治材為器，已勞力樸治斫削，惟其當塗以漆丹以朱而後成。以言教化亦須禮義然後治”。○樸，普角反，馬云：“未成器也。”斫，丁角反。雘，枉略反，徐烏郭反。馬云：“善丹也。”\CJKunderwave{說文}云：“讀與霍同也。”又一郭反，\CJKunderwave{字林}音同。}


{\noindent\zhuan\zihao{6}\fzbyks 傳“為政”至“後治”。正義曰：此三者事別而喻同也。先遠而類疏者，乃漸漸以事近而切者次之。皆言既勤於初，乃言修治於未，明為政孜孜，因前基而修，使善垣墉故也。皆詳而復言之，室器皆雲其事終，而考田止言疆畎，不雲刈獲者，田以一種，但陳修終至收成，故開其初，與下二文互也。二文皆言“斁”，即古“塗”字，明其終而塗飾之。其室言“塗暨”，“暨”亦塗也，總是以物塗之。“茨”謂蓋覆也。器言“塗丹雘”，“塗”、“丹”皆飾物之名，謂塗雘以朱雘。“雘”是彩色之名，有青色者，有硃色者，故\CJKunderline{鄭玄}引\CJKunderwave{山海經}云：“青丘之山，多有青雘。”此經知是“朱”者，與“丹”連文故也。 \par}

{\noindent\shu\zihao{5}\fzkt “惟曰”至“丹雘”。正義曰：既言王者所以效實國君為政之事，故此言國君為政之喻惟為監之事。曰：“若農人之考田也，已勞力遍佈菑而耕發其田，又須為其陳列修治，為疆畔畎壟，以至收穫然後功成。又若人為室家，已勤力立其垣墉，又當惟其塗而暨飾茨蓋之，功乃成也。又若梓人治材為器，已勞力樸治斫削其材，惟其當塗而丹漆以朱雘乃後成。以喻人君為政之道，亦勞心施政，除民之疾，又當惟其飾以禮義，使之行善然後治。” \par}

今王惟曰，先王既勤用明德，懷為夾，\footnote{言文武已勤用明德，懷遠為近，汝治國當法之。○夾音協，近也。}庶邦享作，兄弟方來，，亦既用明德。\footnote{眾國朝享於王，又親仁善鄰為兄弟之國，方方皆來賓服，亦已奉用先王之明德。○朝，直遙反。}後式典集，庶邦丕享。\footnote{君天下能用常法，則和集眾國，大來朝享。}皇天既付中國民，越厥疆土，於先王肆。\footnote{大天已付周家治中國民矣，能遠拓其界壤，則於先王之道遂大。○付如字，馬本作附。拓音託。}王惟德用，和懌先後迷民，用懌先王受命。\footnote{今王惟用德，和悅先後天下迷愚之民。先後謂教訓，所以悅先王受命之義。○懌音亦，字又作斁，下同。先，悉薦反，注同。}已若茲監，惟曰欲至於萬年惟王，\footnote{為監所行已如此所陳法，則我周家惟欲使至於萬年承奉王室。○監，古陷反。為，於威反。}子子孫孫永保民。”\footnote{又欲令其子孫累世長居國以安民。}

{\noindent\zhuan\zihao{6}\fzbyks 傳“言文”至“法之”。正義曰：言“先王”，知謂文武也。“夾”者，是人左右而夾之，故言近也。 \par}

{\noindent\zhuan\zihao{6}\fzbyks 傳“眾國”至“明德”。正義曰：“享”施於王，而“兄弟”為相於之辭，明彼此皆和協。“親仁善鄰”，\CJKunderwave{左傳}文。以先王用明德，於下之所行，今亦奉用，為亦先王耳。 \par}

{\noindent\zhuan\zihao{6}\fzbyks 傳“大天”至“遂大”。正義曰：“肆”，遂也，申遂故為大。“越”，遠也，使天下賓服，故遠柘界壤以益先王，故為“遂大”也。 \par}

{\noindent\zhuan\zihao{6}\fzbyks 傳“今王”至“之義”。正義曰：言“用德”，亦是明德也。“先後”若\CJKunderwave{詩}雲“予曰有先後”，謂於民心先未悟，而啟之已悟，於後化成之，故謂“教訓”也。先王本欲子孫成其事，今化天下使善，是“悅先王受命”。其和悅先王即遠拓疆土，悅其受命即“遂大”也。 \par}

{\noindent\shu\zihao{5}\fzkt “今王”至“保民”。正義曰：此戒\CJKunderline{康叔}已滿三篇，其事將終，須有總結，因其政術言法於明王，上下相承，資以成治，故稱今者王命惟告汝曰：“先王文武在於前世，以自勤用明德，招懷遠人，使來以為親近也。以明德懷柔之故，眾國朝享於王，又相親善為兄弟之國，萬方皆來賓服，亦已化上奉用先王之明德矣。是先王有明德，下亦行明德，以從之而可法也。先王既然，凡為君以君天下者，亦如先王用常法，則和集眾國，使之大來朝享，亦須同先王用明德也。君天下者當如此，今大天已付周家治九州之中國民矣。周家之王，若能為政用明德以懷萬國，遠拓其疆界土壤，則先王之道遂更光大。以此今王須大先王之政，惟明德之大道而用之，以此和悅而先後其天下迷愚之民，使之政治用此，所以悅先王受命使之遂大之義故也。是明德不可不務，故我周王今亦行之。汝為人臣，可以不法乎？當法王家勤用明德治國也。汝若能法我王家而用明德，是為善不可加。”因嘆云：“已乎！如此為監，則我周家惟曰，欲汝至於萬年，惟以承奉王室，今其子子孫孫累世長居國以安民。” \par}

%%% Local Variables:
%%% mode: latex
%%% TeX-engine: xetex
%%% TeX-master: "../Main"
%%% End:

%% -*- coding: utf-8 -*-
%% Time-stamp: <Chen Wang: 2024-04-02 11:42:41>

% {\noindent \zhu \zihao{5} \fzbyks } -> 注 (△ ○)
% {\noindent \shu \zihao{5} \fzkt } -> 疏

\chapter{卷十五}


\section{召誥第十四}


成王在豐，欲宅洛邑，\footnote{武王克商，遷九鼎於洛邑，欲以為都，故成王居焉。}使召公先相宅，\footnote{相所居而卜之，遂以陳戒。○召，詩照反。相，息亮反，下注同。}作\CJKunderwave{召誥}。

召誥\footnote{召公以成王新即政，因相宅以作誥。}


{\noindent\zhuan\zihao{6}\fzbyks 傳“武王”至“居焉”。正義曰：桓二年\CJKunderwave{左傳}云：“昔武王克商，遷九鼎於洛邑。”服虔注云：“今河南有鼎中觀。”雲“九鼎”者，案宣三年\CJKunderwave{左傳}王孫滿云：“昔夏之方有德也，貢金九牧，鑄鼎象物。”然則九牧貢金為鼎，故稱“九鼎”，其實一鼎。案\CJKunderwave{戰國策}顏率說齊王云，昔武王克商，遷九鼎，鼎用九萬人，則以為其鼎有九。但遊說之辭，事多虛誕，不可信用。然鼎之上備載九州山河異物，亦又可疑。未知孰是，故兩解之。 \par}

{\noindent\zhuan\zihao{6}\fzbyks 傳“相所”至“陳戒”。正義曰：孔以序言“相宅”，於經意不盡，故為傳以助成之。召公相所居而卜之，及其經營大作，遂以陳戒，史錄陳戒為篇。其意不在相宅，序以經具，故略之耳。言“先相宅”者，明於時周公攝政，居洛邑是周公之意，周公使召公先行，故言“先”，以見周公自後往也。 \par}

{\noindent\zhuan\zihao{6}\fzbyks 傳“召公”至“作誥”。正義曰：武王既崩，周公即攝王政，至此已積七年，將歸政成王，故經營洛邑，待此邑成，使王即政。召公以成王將新即政，恐王不順周公之意，或將惰於政事，故因相宅以作誥也。作誥之時，王未即政，周公作\CJKunderwave{洛誥}，為反政於成王，召公陳戒，為即政後事，故傳言“新即政”也。 \par}

{\noindent\shu\zihao{5}\fzkt “成王”至“召誥”。正義曰：成王於時在豐，欲居洛邑以為王都，使召公先往相其所居之地，因卜而營之。王與周公從後而往，召公於庶殷大作之時，乃以王命取幣以賜周公，因告王宜以夏殷興亡為戒。史敘其事，作\CJKunderwave{召誥}。 \par}

惟二月既望，\footnote{周公攝政七年二月十五日，日月相望，因紀之。}越六日乙未，王朝步自周，則至於豐。\footnote{於已望後六日，二十一日，成王朝行從鎬京，則至於豐，以遷都之事告文王廟。告文王，則告武王可知，以祖見考。○鎬，胡老反。見，賢遍反，下“不見”同。}惟太保先周公相宅。\footnote{太保，三公官名，召公也。召公於周公前相視洛居，周公後往。○先，息薦反，又如字。}越若來三月，惟丙午。越三日戊申，太保朝至於洛，卜宅。\footnote{朏，明也，月三日明生之名。於順來三月丙午朏。於朏三日，三月五日，召公早朝至於洛邑，相卜所居。○朏，芳尾反，又普沒反，徐又芳憒反。}


{\noindent\zhuan\zihao{6}\fzbyks 傳“周公”至“紀之”。正義曰：\CJKunderwave{洛誥}云：“周公誕保文武受命，惟七年。”\CJKunderwave{洛誥}是攝政七年事也。\CJKunderwave{洛誥}周公云：“予惟乙卯，朝至於洛師。”此篇雲“乙卯，周公朝至於洛”，正是一事，知此“二月”是周公攝政七年之二月也。“望”者，於月之半月，當日衝，日光照月光圓滿，面鄉相當，猶人之相望，故稱“望”也。治歷者必先正朔望，故史官因紀之。將言望後之事，必以望紀之。將言朏後之事，則以朏紀之。猶今人將言日，必先言朔也。望之在月十六日為多，太率十六日者四分之三，十五日者四分之一耳。此年入戊午蔀五十六歲，二月小，乙亥朔。孔雲十五日即為望，是己丑為望，言“已望”者，謂庚寅十六日也。且孔雲“望”與“生魄”、“死魄”皆舉大略而言之，不必恰依歷數。又算術前月大者，後月二日月見，可十五日望也。顧氏亦云：“十五日望，日月正相望也。” \par}

{\noindent\zhuan\zihao{6}\fzbyks 傳“於已”至“見考”。正義曰：“於已望後六日”，是為二十一日也。“步”，行也。此雲“王朝行”，下太保與周公言“朝至”者，君子舉事貴早朝，故皆言“朝”也。宗周者，為天下所宗，止謂王都也。武王已都於鎬，故知宗周是鎬京也。文王居豐，武王未遷之時，於豐立文王之廟，遷都而廟不毀，故成王居鎬京，“則至於豐，以遷都之事告文王廟”也。大事告祖，必告於考，此經不言告武王，以告文王則告武王可知，以告祖見考也。告廟當先祖後考，此必於豐告文王，於鎬京告武王也。 \par}

{\noindent\zhuan\zihao{6}\fzbyks 傳“朏明”至“所居”。正義曰：\CJKunderwave{說文}雲“朏，月未盛之明”，故為“明”也。\CJKunderwave{周書·月令}云：“三日粵朏。”“朏”字從月出，是入月三日明生之名也。“於順來”者，於二月之後依順而來，次三月也。二月乙未而發豐，歷三月丙午朏，又於朏三日，是三月五日，凡發豐至洛為十四日也。“召公早朝至於洛邑，相卜所居”，當以至洛之日即卜也。 \par}

厥既得卜，則經營。\footnote{其已得吉卜，則經營規度城郭郊廟朝市之位處。○度，待洛反。朝，直遙反。處，昌慮反。}越三日庚戌，太保乃以庶殷攻位於洛汭。越五日甲寅，位成。\footnote{於戊申三日庚戌，以眾殷之民治都邑之位於洛水北，今河南城也。於庚戌五日，所治之位皆成。言眾殷，本其所由來。○汭，如銳反。}

{\noindent\zhuan\zihao{6}\fzbyks 傳“其已”至“位處”。正義曰：“經營”，\CJKunderwave{考工記}所云“匠人營國，方九里，左祖右社，面朝後市”是也。下有“丁已郊”，故知“規度城郭郊廟朝市之位處”也。\CJKunderwave{匠人}不言“郊”，以不在國內也。\CJKunderwave{匠人}王城方九里，如\CJKunderwave{典命}文，又以公城方九里，天子城十二里。\CJKunderline{鄭玄}兩說，孔無明解，未知從何文也。“郊”者，\CJKunderwave{司馬法}“百里為郊”，鄭注\CJKunderwave{周禮}雲“近郊五十里”，\CJKunderwave{禮記}祭天於南郊，祭地於北郊，皆謂近郊也。其“廟”，案\CJKunderwave{小宗伯}云：“建國之神位，右社稷，左宗廟。”鄭注\CJKunderwave{朝士職}云，庫門內之左右。其朝者，鄭云，外朝一，在庫門之外，皋門之內，是詢眾庶之朝。內朝二者，其一在路門外，王每日所視，謂之治朝。其一在路門內，路寢之朝，王每日視訖退路寢，謂之燕朝，或與宗人圖私事。顧氏云：“市處王城之北。朝為陽，故在南。市為陰，故處北。”今案\CJKunderwave{周禮·內宰職}“佐後立市”，然則後既主陰，故立市也。 \par}

{\noindent\zhuan\zihao{6}\fzbyks 傳“於戊”至“由來”。正義曰：於戊申後三日庚戌，為三月七日也。水內曰“汭”，蓋以人南面望水，則北為內，故“洛汭”為洛水之北。鄭云：“隈曲中也。”\CJKunderwave{漢書·地理志}河南郡治在洛陽縣，河南城別為河南縣。治都邑之位於洛北，今於漢河南城是也。“所治之位皆成”，佈置處所定也。治位乃是周人，而言“眾殷”者，本其所由來，言本是殷民，今來為我周家役也。莊二十九年\CJKunderwave{左傳}發例云：“凡土功，水昏正而栽，日至而畢。”此以周之三月農時役眾者，彼言尋常土功，此則遷都事大，不可拘以常制也。 \par}

{\noindent\shu\zihao{5}\fzkt “惟二月”至“位成”。正義曰：惟周公攝政七年二月十六日，其日為庚寅，既日月相望矣。於已望後六日乙未，為二月二十一日，王以此日之朝行自周之鎬京，則至於豐，以遷都之事告文王之廟。此日王惟命太保召公先周公往洛水之旁相視所居之處，太保即行。其月小，二十九日癸卯晦。於二月之後順來三月，惟三日丙午朏，而月生明於朏。三日戊申，即三月五日，太保乃以此朝旦至於洛，即卜宅。其已得吉卜，則經營之，規度其城郭郊廟朝市之位處。於戊申三日庚戌，為三月七日，太保乃以眾所受於殷之民，治都邑之位於洛水之汭，謂洛水北也。於庚戌五日，為三月十一日甲寅，而所治之位皆成矣。 \par}

若翼日乙卯，周公朝至於洛，\footnote{周公順位成之明日而朝至於洛汭。}則達觀於新邑營。\footnote{周公通達觀新邑所營。言周遍。}越三日丁巳，用牲於郊，牛二。\footnote{於乙卯三日，用牲告立郊位於天，以後稷配，故二牛。后稷貶於天，有羊豕。羊豕不見，可知。}越翼日戊午，乃社於新邑，牛一，羊一，豕一。\footnote{告立社稷之位，用太牢也。\CJKunderline{共工氏}子曰句龍，能平水土，祀以為社。周祖后稷能殖百穀，祀以為稷。社稷共牢。○共音恭。句，故侯反。}


{\noindent\zhuan\zihao{6}\fzbyks 傳“周公”至“洛汭”。正義曰：周公以順成之明日而朝至，則是三月十二日也。其到洛汭。在召公之後七日。不知初發鎬京以何日也。成王蓋與周公俱來。鄭云：“史不書王往者，王於相宅無事也。” \par}

{\noindent\zhuan\zihao{6}\fzbyks 傳“於乙”至“可知”。正義曰：知此用牲是“告立郊位於天”者，此郊與社，於攻位之時已經營之，今非常祭之月，而特用牲祭天，知是郊位既定，告天使知，而今後常以此處祭天也。\CJKunderwave{禮}郊用特牲，不應用二牛。“以後稷配，故二牛”也。\CJKunderwave{郊特牲}及\CJKunderwave{公羊傳}皆雲養牲必養二，“帝牛不吉，以為稷牛”，言用彼為稷牛者，以之祭帝，其稷牛隨時取用，不在滌養，是帝稷各用一牛，故二牛也。先儒皆雲天神尊，祭天明用犢，貴誠之義。稷是人神，祭用太牢，貶於天神，法有羊豕，因天用牛，遂雲“牛二”，舉其大者，從天言之，羊豕不見，可知也。\CJKunderwave{詩·頌·我將}祀文王於明堂雲“惟羊惟牛”，又\CJKunderwave{月令}雲“以太牢祠於高禖”，皆據配者有羊豕也。 \par}

{\noindent\zhuan\zihao{6}\fzbyks 傳“告立”至“共牢”。正義曰：經有社無稷，稷是社類，知其同告之。告立社稷之位，其祭用太牢，故牛羊豕各一也。句龍能平水土，祀之以為社。后稷能殖百穀，祀以為稷。\CJKunderwave{左傳}、\CJKunderwave{魯語}、\CJKunderwave{祭法}皆有此文。漢世儒者說社稷有二，左氏說社稷惟祭句龍，后稷人神而已，是孔之所用。\CJKunderwave{孝經}說社為土神，稷為穀神，句龍后稷配食者，是鄭之所從。而\CJKunderwave{武成}篇雲“告於皇天后土”，孔以後土為地，言“后土，社也”者，以\CJKunderwave{泰誓}雲“類於上帝，宜於冢土”，故以後土為社也。小劉雲“后土與皇天相對”，以後土為地。若然，\CJKunderwave{左傳}雲“句龍為后土”，豈句龍為地乎？社亦名“后土”，地名“后土”，名同而義異也。“社稷共牢”，經無明說，\CJKunderwave{郊特牲}雲“社稷太牢”，二神共言“太牢”，故傳言“社稷共牢”也。此經上句言“於郊”，此不言“於社”；此言“社於新邑”，上句不言“郊於新邑”；上句言“用牲”，此言牛羊豕，不言“用”；告天不言告地，告社不言告稷；皆互相足，從省文也。\CJKunderwave{洛誥}雲“王在新邑烝祭，王入太室祼”，則洛邑亦立宗廟，此不雲“告廟”，亦從省文也。 \par}

越七日甲子，周公乃朝用書，命庶殷侯、甸、男邦伯。\footnote{於戊午七日甲子，是時諸侯皆會，故周公乃昧爽以賦功屬役書，命眾殷侯、甸、男服之邦伯，使就功。邦伯，方伯，即州牧也。}厥既命殷庶，庶殷丕作。\footnote{其已命殷眾，眾殷之民大作。言勸事。}大保乃以庶邦冢君出取幣，乃復入，\footnote{諸侯公卿並覲於王，王與周公俱至，文不見王，無事。召公與諸侯出取幣，欲因大會顯周公。○復，扶又云。}錫周公，曰：“拜手稽首，旅王若公。\footnote{召公以幣入，稱成王命錫周公，曰：“敢拜手稽首，陳王所宜順周公之事。”}

{\noindent\zhuan\zihao{6}\fzbyks 傳“於戊”至“牧也”。正義曰：\CJKunderwave{康誥}云：“周公初基，作新大邑於東國洛，四方民大和會。侯、甸、男邦、採、衛，百工播民和，見士於周。”與此一事也，故知“是時諸侯皆會，故周公乃昧爽以賦功屬役書，命眾殷在侯、甸、男服之邦伯，使就築作功”也。\CJKunderwave{康誥}五服，此惟三服者，立文有詳略耳。昭三十二年，晉合諸侯城成周，\CJKunderwave{左傳}稱命役書於諸侯，“屬役賦文”，此傳言“賦功屬役”，其意出於彼也。“賦功”謂賦功諸侯之功，科其人夫多少。“屬役”謂付屬役之處，使知得地之尺丈也。“邦伯”，諸國之長，故為方伯州牧。\CJKunderwave{王制}云：“千里之外設方伯。”方伯即州牧也。周公命州牧，使州牧各命其所部。 \par}

{\noindent\zhuan\zihao{6}\fzbyks 傳“諸侯”至“周公”。正義曰：上雲“周公朝用書命庶殷”者，周公自命之，其事不由王也。庶殷既已大作，諸侯公卿乃並覲君王，其時蓋有行宮，王在位而諸侯公卿並覲之。既入見王，乃出取幣。初不言“入”，而經言“出”者，下雲“乃復入”，則上以入可知，從省文也。下賜周公言“旅王若公”，明此出入是覲王之事，而經文不見王至，故傳辯之，王與周公俱至，自此已上於王無事，故不見也。正以經文不見王至，知與周公俱至也。周公居攝功成，將歸政於成王，召公與諸侯出取幣，欲因大會顯周公之功既成。將令王自知政，因賜周公，遂以戒王，故出取幣，復入以待王命。其幣蓋玄纁束帛也。\CJKunderline{鄭玄}云：“所賜之幣，蓋璋以皮，及寶玉大弓，此時所賜。”案鄭注\CJKunderwave{周禮}雲“璋以皮，二王之後享後所用”，寧當以賜臣也？寶玉大弓，魯公之分，伯禽封魯，乃可賜之，不得以此時賜周公也。 \par}

{\noindent\zhuan\zihao{6}\fzbyks 傳“召公”至“之事”。正義曰：太保以庶邦冢君出取幣者，以上太保之意，非王命。幣既入，即雲“賜周公”者，下言召公，不得賜周公，知召公既以幣入，乃稱成王命以賜周公。於時政在周公，成王未得賜周公也。但召公見周公功成作邑，將反王政，欲尊王而顯周公，故稱成王之命以賜周公。\CJKunderline{鄭玄}云：“召公見眾殷之民大作，周公德隆功成，有反政之期，而欲顯之。因大戒天下，故與諸侯出取幣，使戒成王立於位，以其命賜周公。”王肅云：“為戒成王錫周公是也。”曰“拜手稽首”者，召公自言己與冢君等，敢拜手稽首，陳王所宜順周公之事。“宜順之事”，自此以下皆是也。 \par}

{\noindent\shu\zihao{5}\fzkt “若翼”至“若公”。正義曰：順位成之明日乙卯，三月十二日也，周公以此朝旦至於洛，則通達而遍觀於新邑所經營。其位處皆無所改易。於乙卯三日丁巳，三月十四日也，用牲於郊，告立祭天之位，牛二，天與后稷所配各用一牛。於丁巳明日戊午，乃祭社於新邑，用太牢牛一、羊一、豕一。於戊午七日甲子，二十一日也，周公乃以此朝旦用策書，命眾殷在侯、甸、男服之內諸國之長，謂命州牧，使告諸國就功作。其已命殷眾，眾殷皆歡樂歡事而大作矣。太保召公乃以眾國大君諸侯出取幣，乃復入，稱成王命以錫周公，曰：“我敢拜手稽首，以戒王，陳說王所宜順周公之事。” \par}

“誥告庶殷，越自乃御事。\footnote{召公指戒成王，而以眾殷諸侯於自乃御治事為辭，謙也。諸侯在，故託焉。}嗚呼！皇天上帝，改厥元子，茲大國殷之命。\footnote{嘆皇天改其大子，此大國殷之命。言紂雖為天所大子，無道猶改之，言不可不慎。}惟王受命，無疆惟休，亦無疆惟恤。\footnote{所以戒成王，天改殷命，惟王受之，乃無窮惟美，亦無窮惟當憂之。}嗚呼！曷其奈何弗敬？\footnote{何其奈何不憂敬之？欲其行敬。}


{\noindent\zhuan\zihao{6}\fzbyks 傳“嘆皇”至“不慎”。正義曰：\CJKunderwave{釋詁}云：“皇，君也。”天地尊之大，故皇天后土皆以君言之也。“改其大子”，謂改天子之位與他姓，即此大國殷之命，謂紂也。言紂雖為天所大子，無道，尢改之，不可不慎也。以託戒諸侯，故言天子雖大猶改之，況已下乎？\CJKunderwave{釋詁}云：“元，首也。”“首”是體之大，故傳言“大子”。鄭云：“言首子者，凡人皆雲天之子，天子為之首耳。” \par}

{\noindent\shu\zihao{5}\fzkt “誥告”至“弗敬”。正義曰：召公所陳戒王宜順周公之事云：“我為言誥，以告汝庶殷之諸侯，下自汝御事。”欲令君臣皆聽之，其實指以戒王。諸侯皆在，託以為言也。乃曰：“嗚呼！有皇天上帝，改去其大子所受者，即此大國殷之王命也。以其無道，故改命。有德惟王，受得此命，乃無窮惟美，亦無窮惟當憂之。既憂之無窮，嗚呼！何其柰何不敬乎？”欲其長行敬也。“告庶殷”者，告諸侯也。“庶殷”，通尊卑之辭，故民與諸侯同雲“庶殷”，皆謂所受於殷之眾也。 \par}

天既遐終大邦殷之命，茲殷多先哲王在天，\footnote{言天己遠終殷命，此殷多先智王，精神在天不能救者，以紂不行敬故。}越厥後王后民，茲服厥命。\footnote{於其後王后民，謂先智王之後繼世君臣。此服其命，言不忝。}厥終智藏瘝在。\footnote{其終，后王之終，謂紂也。賢智隱藏，瘝病者在位，言無良臣。○瘝，工頑反。}夫知保抱攜持厥婦子，以哀籲天，徂厥亡出執。\footnote{言困於虐政，夫知保抱其子，攜持其妻，以哀號呼天，告冤無辜，往其逃亡，出見執殺，無地自容，所以窮。○夫知，並如字，注同。籲音喻，呼也。號，戶高反。}嗚呼！天亦哀於四方民，其眷命用懋。\footnote{民哀呼天，天亦哀之，其顧視天下有德者，命用勉敬者為民主。}


{\noindent\zhuan\zihao{6}\fzbyks 傳“言天”至“敬故”。正義曰：天既遠終殷命，言其去而不復反也。說天終殷之命，而言智王在天者，言先智王雖精神在天，而不能救紂者，以紂不行敬故也。戒王使行敬。 \par}

{\noindent\zhuan\zihao{6}\fzbyks 傳“於其”至“不忝”。正義曰：“先智王之後繼世君臣”，謂智王之後，紂已前能守位不失者。經言“后王後民”，傳言“君臣”者，見民內有臣。民於此皆服行君之命，言不忝辱父祖也。 \par}

{\noindent\zhuan\zihao{6}\fzbyks 傳“其終”至“良臣”。正義曰：既言“后王”，又復言“其終”，知是“后王之終，謂紂也”。以“瘝”從病類，故言“瘝病”也。鄭、王皆以“瘝”為病，小人在位，殘暴在下，故以病言之。 \par}

{\noindent\zhuan\zihao{6}\fzbyks 傳“言困”至“以窮”。正義曰：言困於虐政，抱子攜妻欲去之。“夫”尢人人，言天下盡然也。“保”訓安也。王肅云：“匹夫知欲安其室，抱其子，攜其妻以悲呼天也。” \par}

{\noindent\shu\zihao{5}\fzkt “天既”至“用懋”。正義曰：更述改殷之事。天既遠終大國殷之王命矣，此殷多有先智之王，精神在天，不能救紂，以紂不行敬故也。於其智王之後人，謂繼世之君及其時之人，皆服行其君之命，由其亦能行敬，故得不忝其先祖。其此後王之終，謂紂之時賢智者隱藏，瘝病者在位，言其時無良臣。多行無禮暴虐，於時之民困於虐政，夫知保抱攜持其婦子，以哀號呼天，告冤枉無辜，往其逃亡，出見執殺，言無地自容以困窮也。天亦哀矜於四方之民，其眷顧天下，選擇賢聖，命用勉力行敬者以為民主，故王今得之也。 \par}

王其疾敬德，相古先民有夏。\footnote{言王當疾行敬德，視古先民有夏之王，以為法戒之。}天迪從子保，面稽天若，今時既墜厥命。\footnote{夏禹能敬德，天道從而子安之。\CJKunderline{禹}亦面考天心而順之，今是桀棄\CJKunderline{禹}之道，天已墜其王命。}今相有殷，\footnote{次復觀有殷。}天迪格保，面稽天若，\footnote{言天道所以至於保安湯者，亦如\CJKunderline{禹}。}今時既墜厥命。\footnote{墜其王命。}今衝子嗣，則無遺壽耇。\footnote{童子，言成王少嗣位治政。無遺棄老成人之言，欲其法之。}曰其稽我古人之德，矧曰其有能稽謀自天？\footnote{衝子成王其考行古人之德則善矣，況曰其有能考謀從天道乎？言至善。}


{\noindent\zhuan\zihao{6}\fzbyks 傳“夏禹”至“王命”。正義曰：勸王疾行敬德，乃言天道安夏，知夏禹能行敬德，天道從而子安之。天既子愛\CJKunderline{禹}，\CJKunderline{禹}亦順天心。鄭雲“面猶迴向也”，則“面”為向義。\CJKunderline{禹}亦志意向天，考天心而順安之，言能同於天心也。\CJKunderline{禹}興夏而桀滅之，知天道子保者是\CJKunderline{禹}也，既墜厥命者是桀也。今桀廢\CJKunderline{禹}之道，已墜失其王命矣。 \par}

{\noindent\zhuan\zihao{6}\fzbyks 傳“言天”至“如\CJKunderline{禹}”。正義曰：此說二代興亡，其意同也。於\CJKunderline{禹}言“從而子安之”，則天於湯亦子安之，故於湯因上略文，直言“格保”。“格”，至也，言至於保安湯者，亦如\CJKunderline{禹}也。 \par}

{\noindent\zhuan\zihao{6}\fzbyks 傳“童子”至“法之”。正義曰：“嗣位治政”，謂周公歸政之後，此時王末蒞政，而言“今衝子嗣”者，召公此戒，戒其即政之後故也。“壽”謂長命，“耇”是老稱，無遺棄長命之老人，欲其取老人之言而法效之，老人之言即下雲“古人之德”也。 \par}

{\noindent\shu\zihao{5}\fzkt “王其”至“自天”。正義曰：既言皇天眷，顧命用勉敬者為人主，故戒王，言其疾行敬德，視古先民有夏之君，取\CJKunderline{大禹}以為法戒。\CJKunderline{禹}以能敬之故，天道從而子安之，\CJKunderline{禹}能面考天心而順以行敬。今是桀棄\CJKunderline{禹}之道，已墜失其王命矣。更復視有殷之君，取\CJKunderline{成湯}以為法戒，湯以能敬之故，天亦從而子。安之天道所以至於保安湯者，亦以湯麵考天心而順以行敬也。今是紂棄湯之道，已墜失其王命矣。夏殷二代，能敬則得之，不敬則失之。今童子為王嗣位治政，則無遺棄壽考成人，宜用老成人之言，法古人為治。曰王其考行古人之德，則已善矣，況曰其有能考行所謀以順從天道乎？若能從順天道，則與\CJKunderline{禹}湯同功，言其善不可加也。 \par}

“嗚呼！有王雖小，元子哉！其丕能諴於小民，今休。\footnote{召公嘆曰：“有成王雖少，而大為天所子，其大能和於小民，成今之美。”勉之。○諴音咸。}王不敢後用，顧畏於民碞。\footnote{王為政當不敢後能用之士，必任之為先。碞，僣也。又當顧畏於下民僣差禮義，能此二者，則德化立，美道成也。○碞，五咸反，徐音吟。}


{\noindent\zhuan\zihao{6}\fzbyks 傳“王為”至“道成”。正義曰：王者為政，任賢使能，有能有用，宜先任之，故“王者為政當不敢後其能用之士，必任之為先”也。“碞”即巖也，參差不齊之意，故為僣也。既任能人，復憂下民，故“又當顧畏於下民僣差禮義”。畏其僣差，當治之使合禮義也。能此二者，則德化立，美道成。“美道成”即“今休”是也。 \par}

{\noindent\shu\zihao{5}\fzkt “嗚呼”至“民碞”。正義曰：召公嘆以戒王：“嗚呼！今所有之王，惟今雖復少小，而大為天所子愛哉！”言任大也。“若其大能和同於天下小民，則成今之美”。以勉之。“故王當不敢後其能用之士，必任以為先。又當顧念畏於下民僣差禮義，能此二者，則德化立，美道成矣”。 \par}

王來紹上帝，自服於土中。\footnote{言王今來居洛邑，繼天為治，躬自服行教化於地勢正中。○治，直吏反，下“為治”、“致治”皆同。}旦曰：‘其作大邑，其自時配皇天。\footnote{稱周公言，其為大邑於土中，其用是大邑，配上天而為治。}毖祀於上下，其自時中乂。\footnote{為治當慎祀於天地，則其用是土中大致治。}王厥有成命，治民今休。’\footnote{用是土中致治，則王其有天之成命，治民今獲太平之美。}


{\noindent\zhuan\zihao{6}\fzbyks 傳“言王”至“正中”。正義曰：傳言“躬自服行”，則不訓自也，鄭、王皆以“自”為用。 \par}

{\noindent\zhuan\zihao{6}\fzbyks 傳“稱周”至“為治”。正義曰：王肅云：“旦，周公名也。禮，君前臣名，故稱周公之言為‘旦曰’。”王者為天所子，代天治民，天有其意，天子繼天使成，謂之“紹上帝”也。天子設法，其理合於天道，是為“配皇天”也。天子將欲配天，必宜治居土中，故稱周公之言，其為大邑於土之中，其當令此成王，用是大邑行化，配上天而為治也。說周公之意然，戒成王使順公也。\CJKunderwave{周禮·大司徒}云：“以土圭之法測土深，正日影，以求地中。日南則影短多暑，日北則影長多寒，日東則影夕多風，日西則影朝多陰。日至之影尺有五寸，謂之地中，天地之所合也，四時之所交也，風雨之所會也，陰陽之所和也。然則百物阜安，乃建王國焉。”馬融云：“王國，東都王城，今河南縣是也。” \par}

{\noindent\zhuan\zihao{6}\fzbyks 傳“為治”至“致治”。正義曰：\CJKunderwave{祭法}云：“有天下者祭百神。”天地為大，“上下”即天地也，故“為治當慎祀於天地”。舉天地則百神之祀皆慎之也。能事神訓民，則其用是土中大致治也。 \par}

{\noindent\zhuan\zihao{6}\fzbyks 傳“用是”至“之美”。正義曰：用是土中致治，當於天心，則王其有天之成命，降福與之，使多歷年歲治民，今獲太平之美。自“旦曰”至此，述周公之意也。 \par}

{\noindent\shu\zihao{5}\fzkt “王來”至“今休”。正義曰：周公之作洛邑，將以反政於王，故召公述其遷洛之意。今王來居洛邑，繼上天為治，躬自服行教化於土地王中之處，故周公旦言曰：“其作大邑於土中，其令成王用是大邑，配大天而為治。為治之道，當事神訓民，謹慎祭祀上下神祇，其用是土中大致治也。既能治，則王其有天之成命，治理下民，今獲太平之美矣。” \par}

王先服殷御事，比介於我有周御事，\footnote{召公既述周公所言，又自陳己意，以終其戒。言當先服治殷家御事之臣，使比近於我有周治事之臣，必和協，乃可一。○比，毗志反，徐扶志反。近，附近之近。}節性，惟日其邁。\footnote{和比殷周之臣，時節其性，令不失中，則道化惟曰其行。○令，力呈反。}王敬作所不可不敬德。\footnote{敬為所不可不敬之德，則下敬奉其命矣。}


{\noindent\zhuan\zihao{6}\fzbyks 傳“召公”至“可一”。正義曰：自“今休”已上，文義相連，知皆是稱周公言也。此一句意異於上，知是“召公自陳己意，以終其戒”。“殷家治事之臣”，謂殷朝舊人，常被殷家任使者也。“周家治事之臣”，謂西土新來翼贊周家初基者也。周臣恃功，或加陵殷士；殷人失勢，或疏忌周臣；新舊不和，政必乖戾。故召公戒王當先治殷臣，使比近周臣，必和協，政乃可一也。不使周臣比殷，而令殷臣比周臣者，周臣奉周之法，當使殷臣從之，故治殷臣使比周臣也。 \par}

{\noindent\zhuan\zihao{6}\fzbyks 傳“和比”至“其行”。正義曰：文承殷周之下，故知“和比殷周之臣”。人各有性，嗜好不同，各恣所欲，必或反道。故以禮義時節其性命，示之限分，令不失中。皆得中道，則各奉王化，故王之道化惟日其行。言日日當行之，日益遠也。顧氏云：“和協殷周新舊之臣，制其性命，勿使怠慢也。” \par}

{\noindent\zhuan\zihao{6}\fzbyks 傳“敬為”至“命矣。”。正義曰：聖王為政，當使易從而難犯，故令行如流水，民從如順風。若使設難從之教，為易犯之令，雖迫以嚴刑，而終不用命。故為其德不可不敬也。王必敬為此不可不敬之德，則下民無不敬奉其命矣民奉其王命，是化行也。 \par}

{\noindent\shu\zihao{5}\fzkt “王先”至“敬德”。正義曰：召公既述周公所言，又自陳己意，戒王今為政，先服治殷家御事之臣，使之比近於我有周治事之臣，令新舊和協，政乃可一。和比殷周之臣，時節其性命，令不失其中，則王之道化惟日其行矣。王當敬為所不可不敬之德，其德為下所敬，則下敬奉其上命，則化必行矣。化在下者，常苦命之不行，故以此為戒。 \par}

“我不可不監於有夏，亦不可不監於有殷。\footnote{言王當視夏殷，法其歷年，戒其不長。}我不敢知曰，有夏服天命，惟有歷年。\footnote{以能敬德，故多歷年數。我不敢獨知，亦王所知。}我不敢知曰，不其延，惟不敬厥德，乃早墜厥命。\footnote{言桀不謀長久，惟以不敬其德，故乃早墜失其王命，亦王所知。}我不敢知曰，有殷受天命，惟有歷年。\footnote{夏言服，殷言受，明受而服行之，互相兼也。殷之賢王，猶夏之賢王，所以歷年，亦王所知。}我不敢知曰，不其延，惟不敬厥德，乃早墜厥命。\footnote{紂早墜其命，猶桀不敬其德，亦王所知。}今王嗣受厥命，我亦惟茲二國命，嗣若功。\footnote{其夏殷也，繼受其王命，亦惟當以此夏殷長短之命為監戒，繼順其功德者而法則之。}


{\noindent\zhuan\zihao{6}\fzbyks 傳“言王”至“不長”。正義曰：“相”、“監”俱訓為視，上言“相有夏”、“相有殷”，今復重言“監有夏”、“監有殷”者，上言順天則興，棄命則滅，此言敬則歷年，不敬則短，故重言視夏殷，欲令王法其歷年，戒其不長故也。 \par}

{\noindent\zhuan\zihao{6}\fzbyks 傳“以能”至“所知”。正義曰：下云：“不敬厥德，乃早墜厥命”，知其“以能敬德者，故多歷年數”也。上言“相夏”、“相殷”皆雲“天迪從子保，面稽天若”，言上天以道安人，人主考天順之，非創業之君不能如是，故傳以\CJKunderline{禹}湯當之。此言“敬德”、“歷年”，則繼體賢君亦能如此，所言“歷年”非獨\CJKunderline{禹}湯而已。下傳雲“殷之賢王，猶夏之賢王”，則此多歷年數者，夏則桀前之賢王，殷則紂前之賢王，不失位者皆是也。召公此誥，指以告王，故知言“我不敢獨知”者，其意言亦是王所知也。王說亦然。 \par}

{\noindent\shu\zihao{5}\fzkt “我不”至“若功”。正義曰：言王所以須慎敬所為不可不敬之德者，以我不可不監視於有夏，亦不可不監視於有殷，皆有歷年，長不與長，由敬與不敬故也，王當法其歷年，戒其不長。更說宜監之意，我不敢獨知，亦王所知，曰有夏之君，服行天命，以敬德之故，惟有多歷年數。謂桀父已前也。其末亦我不敢獨知，亦王所知，曰有夏桀不其長久，惟不敬其德，乃早墜失其王命。是為敬者長，不敬者短，所以我不可不監夏也。我不敢獨知，亦王所知，曰有殷之君受天命，以敬德之故，惟有多歷年數。謂紂父已前也。其末亦我不敢獨知，亦王所知，曰殷紂不其長久，惟不敬其德，乃早墜失其王命。亦是所敬者長，不敬者短，所以我不可不監殷也。夏殷短長既如此矣，今王繼受其命，我亦惟當用此二國夏殷長短之命以為監戒，繼順其功德者而法則之。勸王為敬也。 \par}

王乃初服。嗚呼！若生子，罔不在厥初生，自貽哲命。\footnote{言王新即政，始服行教化，當如子之初生，習為善，則善矣。自遺智命，無不在其初生，為政之道，亦猶是也。○遺，唯季反。}今天其命哲，命吉凶，命歷年。\footnote{今天制此三命，惟人所修。脩敬德則有智，則常吉，則歷年，為不敬德則愚兇不長。雖說之，其實在人。}知今我初服，宅新邑，肆惟王其疾敬德。\footnote{天已知我王今初服政，居新邑洛都，故惟王其當疾行敬德。}王其德之用，祈天永命。\footnote{言王當其德之用，求天長命以歷年。}


{\noindent\zhuan\zihao{6}\fzbyks 傳“言王”至“猶是”。正義曰：以此新即政，始行教化，比子之初生，始欲學習為善，則善矣。若能為善，天必授之以賢智之命，是此賢智之命由己行善而來，是自遺智命矣。初習為惡則惡矣，若其為惡，天必授之以頑愚之命，亦是自遺愚命也。方欲勸王慕善，故惟舉智命而不言愚命者，愚智由學習而至，是無不在其初生。此初生謂年長，以解習學，非初始生也。為政之道亦猶是。為善政得福，為惡政得禍，亦如初生之子習善惡也。 \par}

{\noindent\zhuan\zihao{6}\fzbyks 傳“今天”至“在人”。正義曰：命由天授，遠舉天心，故言“今天制此三命”。有“哲”當有“愚”，有“歷年”當有“不長”，文不備者，以“吉凶”相反，言“命吉凶”，則“哲”對“愚”，“歷年”對“不長”可知矣。天制此三命，善惡由人，惟人所修習也。此篇所云，惟勸脩敬德。故云“脩敬德則有智，則常吉，則多年，惟不敬德則愚兇不長也”。愚智夭壽之外而別言吉凶，於凡人則康強為吉，病患為兇，於王者則太平為吉，禍亂為兇，三者雖以託天說之，其實行之在人。人行之有善惡，天隨以善惡授之耳。此是立教誘人之辭，不可以賢智天枉為難也。 \par}

{\noindent\zhuan\zihao{6}\fzbyks 傳“言王”至“歷年”。正義曰：“其德之用”，言為行當用德，用德與“疾敬德”為一事也。故上傳雲“王其當疾行敬德”，則此文是也。 \par}

其惟王勿以小民淫用非彝，\footnote{勿用小民過用非常。欲其重民秉常。}亦敢殄戮用乂民，\footnote{亦當果敢絕刑戮之道，用治民。戒以慎罰。}若有功，其惟王位在德元。\footnote{順行\CJKunderline{禹}湯所以成功，則其惟王居位在德之首。}小民乃惟刑用於天下，越王顯。\footnote{王在德元，則小民乃惟用法於天下。言治政於王亦有光明。}

{\noindent\zhuan\zihao{6}\fzbyks 傳“勿用”至“秉常”。正義曰：勿用小民非常役，用為非常之義，戒王當使民以時，莫為非常勞役，欲其重民秉常也。 \par}

{\noindent\zhuan\zihao{6}\fzbyks 傳“亦當”至“慎罰”。正義曰：聖人作法，以刑止刑，以殺止殺，若直犯罪之人，亦當果敢致罪之，以此絕刑戮之道，用治民。謂獄事無疑，決斷得理，則果敢為絕刑戮之道。若其獄情疑惑，枉濫者多，是為不能果敢絕刑戮之道也。上戒王以明德，此戒王以慎罰，故言“亦”也。 \par}

{\noindent\zhuan\zihao{6}\fzbyks 傳“順行”至“之首”。正義曰：若有功，必順前世有功者也。上文所云“相夏”、“相殷”，謂\CJKunderline{禹}湯之功，故知此“順行\CJKunderline{禹}湯所有成功”。能順\CJKunderline{禹}湯之功，則惟王居位在德之首。\CJKunderline{禹}湯為有德之首，故王亦為首。 \par}

{\noindent\zhuan\zihao{6}\fzbyks 傳“王在”至“光明”。正義曰：\CJKunderwave{詩}稱“民之秉彝，好是懿德”，故王在德元，則小民乃惟法則於王，行王政於天下。王之為政，民盡行之，是言治政於王道有光明也。 \par}

{\noindent\shu\zihao{5}\fzkt “王乃”至“王顯”。正義曰：既言當法則賢王，又戒王為政之要。王乃初始即政，服行教化。嗚呼！王行教化當如初生之子。子之善惡無不在其初生，若習行善道，此乃自遺智命。“智命”謂身有賢智，命由己來，是自遺也。為政之道亦猶是矣。為政初則能善，天必遺王多福，使王有智則常吉，歷年長久也。今天觀人所為以授之命，其命有智與愚也，其命吉與兇也，其命歷年與不長也。若能敬德，則有智常吉，歷年長久也。若不敬德，則愚兇不長也。天已知我王今初始服政，居此新邑，觀王善惡，欲授之命，故惟王其當疾行敬德。“王其德之用”，言為行當用德，則能求天長命以歷年也。其惟王勿妄役小人過用非常之事，亦當果敢絕刑戮之道，以治下民順行\CJKunderline{禹}湯所有成功，則惟王居天子之位，在德行之首矣。王能如是，小民乃惟法則於王，行用王德於天下，如是則於王道亦有光明也。 \par}

上下勤恤，其曰，我受天命，丕若有夏曆年，式勿替有殷歷年。\footnote{言當君臣勤憂敬德，曰，我受天命，大順有夏之多歷年，勿用廢有殷歷年，庶幾兼之。}欲王以小民受天永命。”\footnote{我欲王用小民受天長命。言常有民。}拜手稽首曰：“予小臣，敢以王之讎民百君子，\footnote{拜手，首至手。稽首，首至地。盡禮致敬，以入其言。言我小臣，謙辭。敢以王之匹民百君子，治民者非一人，言民在下，自上匹之。○讎，字或作酬。}越友民，保受王威命明德。\footnote{言與匹民百君子，於友愛民者共安受王之威命，明德奉行之。}王末有成命，王亦顯。\footnote{臣下安受王命，則王終有天成命，於王亦昭著。}我非敢勤，惟恭奉幣，用供王能祈天永命。”\footnote{言我非敢獨勤而已，惟恭敬奉其幣帛，用供待王，能求天長命。將以慶王多福，必上下勤恤，乃與小民受天永命。○奉如字，又芳孔反。供音恭，徐紀用反，注“供待”同。}


{\noindent\zhuan\zihao{6}\fzbyks 傳“言當”至“兼之”。正義曰：王者不獨治，必當以臣助之。上句惟指勸王，故此又言臣助君。“上下”謂君臣，故言當君臣共勤憂敬德，不獨使王勤也。我周家承夏殷之後受天明命，欲其年過二代，既言大順有夏曆年，又言勿廢有殷歷年，庶幾兼彼二代，歷年長久。勤行敬德，即是大順勿廢也。 \par}

{\noindent\zhuan\zihao{6}\fzbyks 傳“拜手”至“匹之”。正義曰：“拜手”，頭至手。“稽首”，頭至地。謂既為拜，當頭至手，又申頭以至地，故“拜手稽首”重言之。諸言“拜手稽首”者，義皆然也。就此文詳而解之。\CJKunderwave{周禮·太祝}“辨九拜，一曰稽首”，施之於極尊。召公為此拜者，恐王忽而不聽，盡禮致敬以入其言於王。此“拜手稽首”一句，史錄其事，非召公語也。召公設言未盡，為此拜乃更言。鄭云：“拜手稽首者，召公既拜，興曰‘我小臣’以下，言召公拜訖而復言也。”王肅云：“我小臣，召公自謂是小臣，為召公之謙辭。讎訓為匹，敢以王之匹民百君子。百者舉其成數，言治民者非一人。”\CJKunderline{鄭玄}云：“王之諸侯與群吏，是非一人也。”嫌“匹”為齊等，故云“民在下，自上匹之”。 \par}

{\noindent\zhuan\zihao{6}\fzbyks 傳“言我”至“永命”。正義曰：“我非敢勤”，召公自道，言我非敢獨勤而已。“必上下勤恤”，言與眾百君子皆勤也。禮執贄必用幣帛，惟恭敬奉其幣帛，用供待王，能求天長命。將以執贄慶王多福，毛能愛養小民，即是求天長命，待王能愛小民，即欲慶之。 \par}

{\noindent\shu\zihao{5}\fzkt “上下”至“永命”。正義曰：上既勸王敬德，又言臣當助君。言君臣上下勤憂敬德，所以勤者，其言曰：“我周家既受天命，當大順有夏之多歷年歲，用勿廢有殷之多歷年歲。夏殷勤行敬德，故多歷年長久。我君臣亦當行敬德，庶幾兼之。如此者，我欲令王用小民受天長命。”言愛下民，則歷年多也。召公既言此，乃拜手稽首，盡禮致敬，欲王納用其言。既拜而又曰：“我小臣，敢以王之匹配於民眾百君子於友愛民者，共安受王之威命明德，敬奉行之，是上勤恤也。臣下安受王命，則王終有天之成命，於王亦為昭著也。我非敢獨勤而已，眾百君子皆然，言我與眾百君子惟恭敬奉其幣帛，用供待王，能求天長命，將以此慶王受天多福也。” \par}

\section{洛誥第十五}


召公既相宅，周公往營成周，使來告卜，\footnote{召公先相宅卜之，周公自後至，經營作之，遣使以所卜吉兆逆告成王。○相，息亮反，注及下同。使，所吏反，注“遣使”同。}作\CJKunderwave{洛誥}。

洛誥\footnote{既成洛邑，將致政成王，告以居洛之義。}


{\noindent\zhuan\zihao{6}\fzbyks 傳“召公”至“成王”。正義曰：上篇云：“三月戊申，太保朝至於洛，卜宅。厥既得卜，則經營。”是召公先相宅則卜之。又云：“乙卯，周公朝至於洛，則達觀於新邑營。”是周公自後至，經營作之。召公相洛邑，亦相成周，周公營成周，亦營洛邑，各舉其一，互以相明。“卜”者，召公卜也。周公既至洛邑，案行所營之處，遣使以所卜吉兆逆告成王也。案上篇傳雲“王與周公俱至”，何得周公至洛逆告王者？王與周公雖相與俱行，欲至洛之時，必周公先到行處所，故得逆告也。顧氏雲“周公既至洛邑，乃遣以所卜吉兆來告於王”是也。經稱成王言：“公既定宅，伻來，來視予卜休恆吉。”是以得吉兆告成王也。上篇召公以戊申至，周公乙卯至，周公在召公後七日也。至洛較七日，其發鎬京或亦較七日。 \par}

{\noindent\zhuan\zihao{6}\fzbyks 傳“既成”至“之義”。正義曰：周公攝政七年三月經營洛邑，既成洛邑，又歸向西都，其年冬將致政成王，告以居洛之義，故名之曰\CJKunderwave{洛誥}，言以居洛之事告王也。篇末乃雲“戊辰，王在新邑”，明戊辰已上皆是西都時所誥也。 \par}

{\noindent\shu\zihao{5}\fzkt “召公”至“洛誥”。正義曰：序自上下相顧為文，上篇序雲“召公先相宅”，此承其下，故云“召公既相宅”。召公以三月戊申相宅而卜，周公自後而往，以乙卯日至，經營成周之邑。周公即遣使人來告成王以召公所卜之吉兆。及周公將欲歸於成王，乃陳本營洛邑之事，以告成王。王因請教誨之言，周公與王更相報答。史敘其事，作\CJKunderwave{洛誥}。史錄此篇，錄周公與王相對之言，以為後法，非獨相宅告卜而已。但周公因致政本說往前告卜，經文既具，故序略其事，直舉其發言之端耳。 \par}

周公拜手稽首曰：“朕復子明辟。”\footnote{周公盡禮致敬，言我復還明君之政於子。子，成王。年二十成人，故必歸政而退老。○闢，必亦反。}王如弗敢及天基命定命，\footnote{如，往也。言王往日幼少，不敢及知天始命周家安定天下之命，故己攝。}予乃胤保，大相東土，其基作民明闢。\footnote{我乃繼文武安天下之道，大相洛邑，其始為民明君之治。}


{\noindent\zhuan\zihao{6}\fzbyks 傳“周公”至“退老”。正義曰：周公還政而已，明暗在於人君，而云“復還明君之政”者，其意欲令王明，故稱“復子明辟”也。正以此年還政者，以成王年已二十成人，故必歸政而退老也。傳說成王之年，惟此而已。王肅於\CJKunderwave{金縢}篇末云：“武王年九十三而已，冬十一月崩。其明年稱元年，周公攝政，遭流言，作\CJKunderwave{大誥}而東征。二年克殷，殺管叔。三年歸，制禮作樂，出入四年，六年而成。七年營洛邑，作\CJKunderwave{康誥}、\CJKunderwave{召誥}、\CJKunderwave{洛誥}，致政成王。然則武王崩時，成王年已十三矣。周公攝政七年，成王適滿二十。”孔於此言成王年二十，則其義如王肅也。又\CJKunderwave{家語}云：“武王崩時，成王年十三。”是孔之所據也。 \par}

{\noindent\zhuan\zihao{6}\fzbyks 傳“如往”至“己攝”。正義曰：“如，往”，\CJKunderwave{釋詁}文。“及”訓與也，言王往日幼少，志意未成，不敢與知上天始命我周家安定天下之命，故己攝也。天命周家安定天下者，必令天下太平，乃為安定。成王幼少，未能使之安定，故不敢與知之，周公所以攝也。 \par}

{\noindent\zhuan\zihao{6}\fzbyks 傳“我乃”至“之治”。正義曰：“胤”訓繼也，文王受命，武王伐紂，意在安定天下。天下未得安定，故周公言我乃繼續文武安定天下之道，大相洛邑之地，其處可行教化，始營此都，為民明君之政治。言欲為民明君，其意當在此。 \par}

{\noindent\shu\zihao{5}\fzkt “周公”至“民明闢”。正義曰：周公將反歸政，陳成王將居其位，周公拜手稽首，盡禮致敬於王，既拜乃興而言曰，我今復還子明君之政。言王往日幼少，其志意未成，不敢及知天之始命我周家安定天下之命，故我攝王之位，代王為治。我乃繼文王、武王安定天下之道，以此故大視東土洛邑之居，其始欲王居之，為民明君之治。言欲為民明君，必當治土中，故為王營洛邑也。 \par}

予惟乙卯，朝至於洛師。\footnote{致政在冬，本其春來至洛眾，說始卜定都之意。}我卜河朔黎水，我乃卜澗水東、瀍水西，惟洛食。\footnote{我使人卜河北黎水上，不吉。又卜澗瀍之間，南近洛，吉。今河南城也。卜必先墨畫龜，然後灼之，兆順食墨。○河朔，朔北也。瀍，直連反。近，附近之近。}我又卜瀍水東，亦惟洛食。伻來以圖及獻卜。”\footnote{今洛陽也。將定下都，遷殷頑民，故並卜之。遣使以所卜地圖及獻所卜吉兆，來告成王。○伻，普耕反，徐敷耕反，又甫耕反，下同。}


{\noindent\zhuan\zihao{6}\fzbyks 傳“致政”至“之意”。正義曰：下文總結周公攝政之事雲“在十有二月”，是“致政在冬”也。“在冬”，發言嫌此事是冬，故辨之雲“本其春來至洛眾”，追說始卜定都之意也。周公至洛之時，庶殷已集於洛邑，故云“至於洛師”。 \par}

{\noindent\zhuan\zihao{6}\fzbyks 傳“我使”至“食墨”。正義曰：嫌周公自卜，故云“我使人”，謂使召公也。案上篇召公至洛，其日即卜，而得“卜河朔黎水”者，以地合龜，非就地內，此言所卜三處皆一時事也。“黎水”之下不言吉凶者，“我乃”是改卜之辭，明其不吉乃改，故知“卜河北黎水之上，不吉”也。武王定鼎於郟鄏，已有遷都之意，而先卜黎水上者，以帝王所都，不常厥邑，夏殷皆在河北，所以博求吉地，故令先卜河北，不吉乃卜河南也。其“卜澗瀍之間，南近洛，吉。今河南城也”，基趾仍在，可驗而知。所卜黎水之上，其處不可知矣。凡卜之者，必先以墨畫龜，要坼依此墨，然後灼之求其兆，順食此墨畫之處，故云“惟洛食。”顏氏云：“先卜河北黎水者，近於紂都，為其懷土重遷，故先卜近以悅之。”用鄭康成之說，義或然也。 \par}

{\noindent\zhuan\zihao{6}\fzbyks 傳“今洛”至“成王”。正義曰：洛陽即成周，敬王自王城遷而都之。\CJKunderwave{春秋}昭三十二年“城成周”是也。周公慮此頑民未從周化，故既營洛邑，將定下都，以遷殷之頑民，故命召公即並卜之。周公既至，即遣使以所卜地圖及獻所卜吉兆，來告於成王。言己重其事，並獻卜兆者，使王觀兆知其審吉也。 \par}

{\noindent\shu\zihao{5}\fzkt “予惟”至“獻卜”。正義曰：周公追述立東都之事，我惟以七年三月乙卯之日，朝至於洛邑眾作之處，經營此都。其未往之前，我使人卜河北黎水之上，不得吉兆。乃卜澗水東、瀍水西，惟近洛，而其兆得吉，依規食墨。我亦使人卜瀍水東，亦惟近洛，其兆亦吉，依規食墨。我以乙卯至洛，我即使人來以所卜地圖及獻所卜吉兆於王。言卜吉立此都，王宜居之為治也。 \par}

王拜手稽首曰：“公不敢不敬天之休，來相宅，其作周匹休。\footnote{成王尊敬周公，答其拜手稽首而受其言。述而美之，言公不敢不敬天之美，來相宅，其作周以配天之美。}公既定宅，伻來，來視予卜休恆吉。我二人共貞。\footnote{言公前已定宅，遣使來，來視我以所卜之美、常吉之居，我與公共正其美。○貞，正也，馬云：“當也。”}公其以予萬億年敬天之休。”\footnote{公其當用我萬億年敬天之美。十千為萬，十萬為億，言久遠。}拜手稽首誨言。\footnote{成王盡禮致敬於周公，求教誨之言。○盡，子忍反。}


{\noindent\zhuan\zihao{6}\fzbyks 傳“成王”至“之美”。正義曰：拜手稽首，施於極敬。哀十七年\CJKunderwave{左傳}云：“非天子，寡君無所稽首。”諸侯小事大尚不稽首，況於臣乎？成王尊敬周公，答言其拜手稽首而受其言。又述而美之，天命文武使王天下，是天之美事，言公不敢不敬天之美，來相洛邑之宅。 \par}

{\noindent\zhuan\zihao{6}\fzbyks 傳“言公”至“其美”。正義曰：周公追述往前遣使獻卜，故成王複述公言。言公前已定宅，遣使來，來視我所卜之吉兆、常吉之居。自言前已知其卜，既有此美，我當與公二人共正其美意。欲留公輔己，共公正此美事。“來來”重文者，上“來”言使來，下“來”為視我卜也。鄭云：“伻來來者，使二人也。”與孔意異。 \par}

{\noindent\zhuan\zihao{6}\fzbyks 傳“公其”至“久遠”。正義曰：言居洛為治，可以永久，公意其當用我使萬億年敬天之美，言公欲令己祚胤久遠，美公意之深也。\CJKunderwave{王制}云：“方百里者，為方十里者百，為田九十億畝。”方里者萬，則是為田九百萬畝。今\CJKunderwave{記}乃雲“九十億畝”，是名十萬為億也。\CJKunderwave{楚語}雲“百姓、千品、萬官、億醜”，每數相十，是古十萬曰億。今之算術乃萬萬為億也。 \par}

{\noindent\zhuan\zihao{6}\fzbyks 傳“成王”至“之言”。正義曰：此一段史官所錄，非王言也。王求教誨之言，必有求教誨之辭，史略取其意，故直雲“誨言”。為求誨言而拜，故言“成王盡禮致敬於周公，求教誨之言”也。 \par}

{\noindent\shu\zihao{5}\fzkt “王拜手”至“誨言”。正義曰：成王尊敬周公，故亦盡禮致敬，拜手稽首，乃受公之語，述公之美曰：“不敢不敬天之美，來至洛相宅，其意欲作周家配天之美故也。公既定洛邑，即使人來告，亦來視我以所卜之美、常吉之居，我當與公二人共正其美。公定其宅，其當用我萬億年敬天之美故也。”王既言此，又拜手稽首於周公，求教誨之言。 \par}

周公曰：“王肇稱殷禮，祀於新邑，咸秩無文。\footnote{言王當始舉殷家祭祀，以禮典祀於新邑，皆次秩不在禮文者而祀之。}予齊百工，伻從王於周。予惟曰：‘庶有事。’\footnote{我整齊百官，使從王於周，行其禮典。我惟曰：‘庶幾有善政事。’}今王即命曰：‘記功，宗以功，作元祀。’\footnote{今王就行王命於洛邑，曰：“當記人之功，尊人亦當用功大小為序，有大功則列大祀。”謂功施於民者。○曰記，上音越，一音人實反。}


{\noindent\zhuan\zihao{6}\fzbyks 傳“言王”至“祀之”。正義曰：於時制禮已訖，而云“殷禮”者，此“殷禮”即周公所制禮也。雖有損益，以其從殷而來，故稱“殷禮”。猶上篇雲“庶殷”，“本其所由來”，孔於上傳已具，故於此不言。必知殷禮即周禮者，以此雲“祀於新邑”，即下文“烝祭歲”也，既用“騂牛”，明用周禮。雲“始”者，謂於新邑始為此祭。顧氏云：“舉行殷家舊祭祀，用周之常法。”言周禮即殷家之舊禮也。\CJKunderline{鄭玄}云：“王者未制禮樂，恆用先王之禮樂。”是言伐紂以來，皆用殷之禮樂，非始成王用之也。周公制禮樂既成，不使成王即用周禮，仍令用殷禮者，欲待明年即取，告神受職，然後班行周禮。班訖始得用周禮，故告神且印用殷禮也。孔義或然，故復存之。神數多而禮文少，應祭之神名有不在禮文者，故令皆次秩不在禮文而應祀者，皆舉而祀之。 \par}

{\noindent\zhuan\zihao{6}\fzbyks 傳“我整”至“政事”。正義曰：時成王未有留公之意，公以成王初始即政，自慮百官不齊，故雖即致政，猶欲整齊百官，使從王於周，謂從至新邑，行其典禮。周公以成王賢君，今覆成長，故言“我惟曰：‘庶幾有善政事。’”言己私為此言，冀王為政善也。 \par}

{\noindent\zhuan\zihao{6}\fzbyks 傳“今王”至“民者”。正義曰：記臣功者是人主之事，故言“今王就行王命於洛邑”，謂正位為王，臨察臣下，知其有功以否。恐王輕忽此事，故曰“當記人之功”。更言“曰”者，所以致殷勤也。尊人必當用功大小為次序，令功大者居上位，功小者處下位也。“有大功則列為大祀”，謂有殊功，堪載祀典者。\CJKunderwave{祭法}云：“聖王之制祭祀也，法施於民則祀之，以死勤事則祀之，以勞定國則祀之，能御大災則祀之，能捍大患則祀之。”是為大祀“謂功施於民”者也。或時立其祀配享廟廷，亦是也。 \par}

惟命曰：‘汝受命篤，弼丕視功載，乃汝其悉自教工。’\footnote{惟天命我周邦，汝受天命厚矣，當輔大天命，視群臣有功者記載之，乃汝新即政，其當儘自教眾官，躬化之。}孺子其朋，孺子其朋，其往。\footnote{少子慎其朋黨，少子慎朋黨，戒其自今已往。}無若火始焰焰，厥攸灼敘，弗其絕。\footnote{言朋黨敗俗，所宜禁絕。無令若火始然，焰焰尚微，其所及，灼然有次序，不其絕。事從微至著，防之宜以初。○焰音豔。“敘”絕句，馬讀“敘”句字屬下。令，力呈反。}


{\noindent\zhuan\zihao{6}\fzbyks 傳“惟天”至“化之”。正義曰：“惟天命我周邦”，謂天命我文武，故及汝成王復受天命為天子，是天之恩德深厚矣。天以厚德被汝，汝當輔大天命，任賢使能，行合天意，是輔大天也。汝當輔大天命，故宜視群臣有功者記載之，覆上“記功，宗以功”言之也。欲令群臣有功，必須躬自教化之在於初始，故言“乃汝新即政，其當儘自教眾官”。欲令王“躬化之”者，正己之身，使群臣法之，非謂以辭化之也。言“儘自教”者，政有大小，恐王輕大略小，令王儘自親化之。言“惟命曰”，亦是致殷勤。“乃”者，緩辭也。義異上句，故言“乃”耳。王肅云：“此其儘自教百官，謂正身以先之。” \par}

{\noindent\zhuan\zihao{6}\fzbyks 傳“少子”至“已往”。正義曰：鄭云：“孺子，幼少之稱，謂成王也。”此上皆雲成王，此句特言少子者，以明朋黨敗俗，為害尢大，恐年少所忽，故特言“孺子”也。“朋黨”謂臣相朋黨。“慎其朋黨”，令禁絕之。“戒其自今已往”，謂從即政以後，常以此事為戒也。 \par}

{\noindent\zhuan\zihao{6}\fzbyks 傳“言朋”至“以初”。正義曰：“無令若火始然”，以喻無令朋黨始發。若火既然，初雖焰焰尚微，其火所及，灼然有次序，不其復可絕也。以喻朋黨若起，漸漸益大，群黨既成，不可復禁止也。“事從微至著，防之宜以初”，謂朋黨未發之前，防之使不發。 \par}

厥若彝，及撫事如予，惟以在周工。\footnote{其順常道，及撫國事，如我所為，惟用在周之百官。}往新邑，伻鄉即有僚，明作有功，惇大成裕，汝永有辭。”\footnote{往行政化於新邑，當使臣下各鄉就有官，明為有功，厚大成寬裕之德，則汝長有嘆譽之辭於後世。○向，許亮反，注同。惇，都混反。}

{\noindent\zhuan\zihao{6}\fzbyks 傳“其順”至“百官”。正義曰：考古依法，為“順常道”。號令治民，為“撫國事”。周公大聖，動成軌則，“如我所為”，謂如攝政之時事所施為也。惟當用我所為在周之百官，令其行周公之道法於百官也。 \par}

{\noindent\zhuan\zihao{6}\fzbyks 傳“往行”至“後世”。正義曰：此時在西都戒王，故云“往行政化於新邑”。當使臣下各鄉就所有之官，令其各守其職，思不出其位，自當陳力就列，明為有功。在官者當以褊小急躁為累，故令臣下厚大成寬裕之德。臣下既賢，君必明聖，則汝長有嘆譽之辭於後世矣。今\CJKunderwave{周頌}所歌即嘆譽成王之辭也。 \par}

{\noindent\shu\zihao{5}\fzkt “周公”至“有辭”。正義曰：王求教誨之言，公乃誨之。周公曰：“王居此洛邑，當始舉殷家祭祀以為禮典，祀於洛之新邑，皆次秩在禮無文法應祀者，亦次秩而祀之。我雖致政，為王整齊百官，使從王於周，行其禮典。若能如此，我惟曰：‘庶幾有善政事。’今王就行王命於洛邑，曰：‘王當記人之功，尊人亦當用功大小為次序，有大功者則列為大祀。’”又申述所以祀神記臣功者。“政事由臣而立，惟天命我周邦之故，曰：‘汝受天命厚矣，當輔大天命，故須視群臣有功者記載之。君知臣功，則臣皆盡力。欲令群臣盡力，宜於初即教之。乃汝新始即政，其當儘自教誨眾官。’”令王躬自化之，使之立功。又以朋黨害政，尤宜禁絕，故丁寧戒之：“少子慎其朋黨，少子慎其朋黨，戒其自今已往。”令常慎此朋黨之事。“若欲絕止，禁其未犯，無令若火始然。焰焰尚微，火既然焰，其火所及，將灼然有次序矣，不其復可絕也。汝成王其當順此常道，及撫循國事，如我攝政所為。惟當用我此事，在周之百官則當畏服，各立功矣。汝當以此往行政化於新邑，當使臣下百官各向就有官，明為有功，厚大成寬裕之德，則汝長有嘆譽之辭於後世”。此周公誨王之言也。 \par}

公曰：“已！汝惟衝子惟終。\footnote{已乎！汝惟童子，嗣父祖之位，惟當終其美業。}汝其敬識百辟享，亦識其有不享。享多儀，儀不及物，惟曰不享。\footnote{奉上謂之享。言汝為王，其當敬識百君諸侯之奉上者，亦識其有違上者。奉上之道多威儀，威儀不及禮物，惟曰不奉上。}惟不役志於享，凡民惟曰不享，惟事其爽侮。\footnote{言人君惟不役志於奉上，則凡人化之，惟曰不奉上矣。如此則惟政事其差錯侮慢不可治理。}


{\noindent\zhuan\zihao{6}\fzbyks 傳“已乎”至“美業”。正義曰：周公止而復言，故更言“公曰”。“已乎”者，道前言已如是矣，為後言發端也。“童子”者，言其年幼而任重。“嗣父祖之位，當終其美業”，能致太平，是終之也。 \par}

{\noindent\zhuan\zihao{6}\fzbyks 傳“奉上”至“奉上”。正義曰：“享”訓獻也，獻是奉上之辭，故“奉上謂之享”。百官諸侯上事天子，凡所恭承皆是奉上，非獨朝覲貢獻乃為奉上。\CJKunderline{鄭玄}專以朝聘說之，理未盡也。言汝為王當敬識百官諸侯之奉上者，亦識其有違上者，察其恭承王命如法以否，奉上違上皆須記之。奉上者當以禮接之，違上者當以刑威之，所謂賞慶刑威。為君之道，奉上之道，其事非一，故云“多威儀”。威儀既多，皆須合禮，其威儀不及禮物，惟曰不奉上矣。謂旁人觀之，亦言其不奉上也。鄭云：“朝聘之禮至大，其禮之儀不及物，所謂貢篚多而威儀簡也。威儀既簡，亦是不享也。” \par}

{\noindent\shu\zihao{5}\fzkt “公曰”至“爽侮”。正義曰：周公復誨王曰：“嗚呼！前言已如是。”更復教誨：“汝惟童子，嗣父祖之位，惟當終其美業。天子居百官諸侯之上，須知臣下恭之與慢。奉上謂之享，汝為天子，其當恭敬記識百官諸侯奉上者，亦當記識其有不奉上者。奉上之道多威儀，威儀不及禮物，則人惟曰不奉上之道矣。所以須記之者，百官諸侯為下民之君，惟為政教不肯役用其志於此奉上之事，則凡民化之，亦惟曰不奉上矣。百官不承天子，民復不奉百官，上下不相畏敬，惟政事其皆差錯侮慢不可治理矣。故天子須知百官奉上與否也。” \par}

乃惟孺子，頒朕不暇，聽朕教汝於棐民彝。\footnote{我為政常若不暇，汝為小子當分取我之不暇而行之，聽我教汝於輔民之常而用之。○頒音班，徐甫雲反，馬云：“猶也。”棐音匪，又芳鬼反。}汝乃是不蘉，乃時惟不永哉！\footnote{汝乃是不勉為政，汝是惟不可長哉！欲其必勉為可長。○蘉，徐莫剛反，又武剛反，馬云：“勉也。”}

{\noindent\zhuan\zihao{6}\fzbyks 傳“我為”至“用之”。正義曰：“為政常若不暇”，謂居攝時也。聖人為政，務在化人，雖復治致太平，猶恨意之不盡，故謙言己所不暇，若言猶有美事未得施者然。故戒之成王：“汝惟小子，當分取我之不暇而行之。”言己所不暇行者，欲令成王勉行之。\CJKunderline{鄭玄}云：“成王之才，周公倍之，猶未而言分者，誘掖之言也。”生民之為業，雖復志有經營，不能獨自成就，須王者設教以輔助之。“聽我教汝輔民之常法而用之”，謂用善政以安民。\CJKunderwave{說文}云：“頒，分也。” \par}

{\noindent\zhuan\zihao{6}\fzbyks 傳“汝乃”至“可長”。正義曰：成王言“公其以予萬億年”，言欲己長久也，故周公於此戒之：“汝乃於是不勉力為政，汝惟不可長哉！”欲其必勉力勤行政教，為可長久之道，然後可至萬億年耳。“醟”之為勉，相傳訓也。鄭、王皆以為勉。 \par}

篤敘乃正父，罔不若予，不敢廢乃命。\footnote{厚次序汝正父之道而行之，無不順我所為，則天下不敢棄汝命，常奉之。}汝往敬哉！茲予其明農哉！彼裕我民，無遠用戾。”\footnote{汝往居新邑，敬行教化哉！如此我其退老，明教農人以義哉！彼天下被寬裕之政，則我民無遠用來。言皆來。○被，皮寄反，又彼美反。}

{\noindent\zhuan\zihao{6}\fzbyks 傳“厚次”至“奉之”。正義曰：“正父”謂武王，言其德正，故稱“正父”。“厚次序汝正父之道而行之”，令其為武王之政也。武王、周公俱是大聖，無不順我所為。又令法周公之道，既言法武王，又法周公，則天下不敢棄汝命，常奉行之。 \par}

{\noindent\zhuan\zihao{6}\fzbyks 傳“汝往”至“皆來”。正義曰：歸其王政令，汝往居新邑，敬行教化哉！公既歸政，則身當無事，如此我其退老於州里，明教農人以義哉！又令成王行寬裕之政，以治下民。民被寬裕之政，則天下之民無問遠近者，用來歸王，言遠處皆來也。上文使之“惇大成裕”，故此言裕政來民結上事也。伏生\CJKunderwave{書傳}稱，禮致仕之臣，教於州里，大夫為父師，士為少師，朝夕坐於門塾，而教出入之子弟。是“教農人以義”也。 \par}

{\noindent\shu\zihao{5}\fzkt “乃惟”至“用戾”。正義曰：又曰：“己居攝之時，為政常若不暇，汝惟小子，當分取我之不暇而施行之。又聽我教汝於輔民之常而用之。汝乃於是事不勉力為政，則汝是惟不可長久哉！必須勉力為之，乃可長久。此所言皆是汝父所行，汝欲勉之，但厚次序汝正父之道而行之，無不順我所為，則天下不敢棄廢汝命，必常奉而行之。汝往居新邑，敬行教化哉！如此我其退老，明教農人以義哉！汝若能使彼天下之民被寬裕之政，則我天下之民無問遠近者，悉皆用來歸汝矣。” \par}

王若曰：“公，明保予衝子。\footnote{成王順周公意，請留之自輔。言公當明安我童子，不可去之。}公稱丕顯德，以予小子揚文武烈，\footnote{言公當留，舉大明德，用我小子襃揚文武之業而奉順天。○襃，薄謀反，\CJKunderwave{切韻}博毛反。}奉答天命，和恆四方民，居師。\footnote{又當奉當天命，以和常四方之民，居處其眾。}惇宗將禮，稱秩元祀，咸秩無文。\footnote{厚尊大禮，舉秩大祀，皆次秩無禮文而宜在祀典者，凡此待公而行。}


{\noindent\zhuan\zihao{6}\fzbyks 傳“成王”至“去之”。正義曰：成王以周公誨己為善，順周公之意，示己欲行善政，而請留之自輔。王以公若舍我而去，則已政暗而治危，故云“公當明安我童子，不可去”也。 \par}

{\noindent\zhuan\zihao{6}\fzbyks 傳“言公”至“順天”。正義曰：文武受命，功德盛隆，成王自量己身不能繼業，言公當留，舉大明德，以佑助我。“用我小子褒揚文武之業而奉順天”者，下句“奉答天命”是也。孔分經為傳，故探取下句以申之。 \par}

{\noindent\zhuan\zihao{6}\fzbyks 傳“又當”至“其眾”。正義曰：天命周家，欲令民治，故“又當奉當天命，以和常四方之民，居處其眾”也。“奉當”者，尊天意，使允當天心，和協民心，使常行善也。“居處其眾”，使之安土樂業也。 \par}

{\noindent\zhuan\zihao{6}\fzbyks 傳“厚尊”至“而行”。正義曰：\CJKunderwave{釋詁}云：“將，大也。”“厚尊大禮”，謂祭祀之禮。\CJKunderwave{祭統}云：“禮有五經，莫重於祭。”是祭禮最尊大。公誨成王，令“肇稱殷禮，祀於新邑，咸秩無文”，欲答公誨己之事，還述公辭：“舉秩大祀，皆次秩無禮文而宜在祀典者。”其祀事非我所為，凡此皆待公而行者也。言公不可舍我以去也。 \par}

惟公德明，光於上下，勤施於四方。\footnote{言公明德光於天地，勤政施於四海，萬邦四夷服仰公德而化之。}旁作穆穆迓衡，不迷文武勤教，\footnote{四方旁來為敬敬之道，以迎太平之政，不迷惑於文武所勤之教。言化治。○旁，步光反。迓，五嫁反，馬、鄭、王皆音魚據反。}予衝子夙夜毖祀。”\footnote{言政化由公而立，我童子徒早起夜寐，慎其祭祀而已。無所能。}

{\noindent\zhuan\zihao{6}\fzbyks 傳“言公”至“化之”。正義曰：此與下經皆追述居攝時事。\CJKunderwave{堯典}訓“光”為充，此“光”亦為充也。言公之明德充滿天地，即\CJKunderwave{堯典}“格於上下”，勤政施於四方，即\CJKunderwave{堯典}“光被四表”也，意言萬邦四夷皆服仰公德而化之。上言待公乃行之，此言公有是德，言其將來，說其已然，所以深美公也。 \par}

{\noindent\zhuan\zihao{6}\fzbyks 傳“四方”至“化洽”。正義曰：上言施化在公，此言民化公德，四方旁來為敬敬之道，民皆敬鄉公以迎太平之政。言“迎”者，公政從上而下，民皆自下迎之，言其慕化速也。文武勤行教化，欲以教訓利民，民蒙公化，識文武之心，不復迷惑文武所勤之教，言公居攝之時，政化已洽於民也。 \par}

{\noindent\zhuan\zihao{6}\fzbyks 傳“言政”至“所能”。正義曰：此述留公之意，陳自今已後之事。言公若留住，政化由公而立。我童子徒早起夜寐，慎其祭祀而已。於政事無所能，欲惟典祭祀，以政事委公。襄二十六年\CJKunderwave{左傳}云，衛獻公使與寧喜言曰：“苟得反國，政由寧氏，祭則寡人。”亦猶是也。 \par}

{\noindent\shu\zihao{5}\fzkt “王若”至“毖祀”。正義曰：王以周公將退，因誨之而請留公，王順周公之意而言曰：“公當留住而明安我童子，不可去也。所以不可去者，當舉行大明之德，用使我小子褒揚文武之業，而奉當天命，以和常四方之民，居處其眾故也。其厚尊大禮，謂舉秩大祀，皆次秩禮所無文者而皆祀之。凡此皆待公而行，非我能也。”更述居攝時事：“惟公明德光於天地，勤政施於四方，使四方旁來為敬敬之道，以迎太平之政，下民皆不復迷惑於文武所勤之教。”言公化洽使如此也。“今若留輔我童子，惟當早起夜寐，慎其祭祀而已”。言政化由公而立，我無所能也。 \par}

王曰：“公功棐迪篤，罔不若時。”\footnote{公之功輔道我已厚矣，天下無不順而是公之功。}

{\noindent\zhuan\zihao{6}\fzbyks 傳“公之”至“之功”。正義曰：王意言公之居攝，天下若為非，則可舍我而去。公之居攝，天下無不順而是公之功，不可舍我去。 \par}

{\noindent\shu\zihao{5}\fzkt “王曰公功”至“若時”。正義曰：王又重述前言，還說居攝時事也。曰：“公之功輔道我已厚矣，天下無有不順而是公之功者，公所以須留也。” \par}

王曰：“公，予小子其退即闢於周，命公後。\footnote{我小子退坐之後，便就君於周，命正公後，公當留佑我。}四方迪亂，未定於宗禮，亦未克敉公功。\footnote{言四方雖道治，猶未定於尊禮。禮未彰，是亦未能撫順公之大功。明不可以去。○敉，亡婢反。治，直吏反，下同。}迪將其後，監我士師工，\footnote{公留教道，將助我其今已後之政，監篤我政事眾官。委任之言。○監，工銜反，注同。}誕保文武受民亂，為四輔。”\footnote{太安文武所受之民治之，為我四維之輔。明當依倚公。}


{\noindent\zhuan\zihao{6}\fzbyks 傳“我小”至“佑我”。正義曰：“退”者，退朝也。周公於時令成王坐王位而以政歸之，成王順周公言受其政也。言我小子退坐之後，便就君位於周。“周”謂洛邑，許其從公言，適洛邑而行新政也。古者臣有大功，必封為國君，今周公將欲退老，故命立公後，使公子伯禽為國君，公當留佑我。王肅云：“成王前春亦俱至洛邑，是顧無事，既會而還宗周。周公往營成周，還來致政成王也。” \par}

{\noindent\zhuan\zihao{6}\fzbyks 傳“言四”至“以去”。正義曰：王意恐公意以四方既定，不須更留，故謂公云，四方雖已道治，而猶未能定於尊大之禮。言其禮樂未能彰明也。禮既未彰，是天下之民亦未能撫安順行公之大功，公當待其禮法明，公功順乃可去耳。明今不可以去。 \par}

{\noindent\zhuan\zihao{6}\fzbyks 傳“大安”至“倚公”。正義曰：文武受民之於天下，今大安文武所受之民，助我治之，為我四維之輔，明己當依倚公也。“維”者，為之綱紀，猶如用繩維持之。\CJKunderwave{文王世子}雲“設四輔”，謂設眾官為四方輔助。周公一人，事無不統，故一人為四輔。\CJKunderwave{管子}云：“四維不張，國乃滅亡。”傳取\CJKunderwave{管子}之意，故言“四維之輔”也。 \par}

{\noindent\shu\zihao{5}\fzkt “王曰公予”至“四輔”。正義曰：王呼周公曰：“我小子其退此坐，就為君於周。”謂順公之言，行天子之政於洛邑也。“至洛邑當命公後，立公之世子為國君，公當留輔我也。公之攝政，四方雖已道治理，猶自未能定於尊禮，是亦未能撫順公之大功。公當待其定大禮，順公之大功，此時未可去也。公當留教道，將助我其今已後之政，監篤我政事眾官，以此大安文武所受之民而治之，為我四維之輔助”。明己當依倚公也。 \par}

王曰：“公定，予往已。公功肅將祗歡，\footnote{公留以安定我，我從公言，往至洛邑已矣。公功以進大，天下咸敬樂公功。○樂音洛。}公無困哉！我惟無斁其康事，公勿替刑，四方其世享。”\footnote{公必留，無去以困我哉！我惟無厭其安天下事。公勿去以廢法，則四方其世世享公之德。○斁音亦。厭，於豔反。}


{\noindent\zhuan\zihao{6}\fzbyks 傳“公留”至“公功”。正義曰：讀文以“公定”為句，王稱“定”者，言定己也，故傳言“公留以安定我”，“我”字傳加之。“我從公言”，是經之“予”也。“往至洛邑已矣”，言已順從公命，受歸政也。“公功已進大，天下咸敬樂公之功”，亦謂居攝時也。\CJKunderwave{釋詁}云：“肅，進也。” \par}

{\noindent\zhuan\zihao{6}\fzbyks 傳“公必”至“之德”。正義曰：王言己才智淺短，公去則困，故請公“無去以困我哉”。我意欲致太平，惟無厭倦其安天下之事，是以留公，公勿去以廢治國之法，則天下四方之民蒙公之恩，其世世享公之德。“享”謂荷負之。 \par}

{\noindent\shu\zihao{5}\fzkt “王曰公定”至“世享”。正義曰：王又呼公：“公留以安定我，我從公言，往至洛邑已矣。公功已進且大矣，天下皆樂公之功，敬而歡樂。公必留，無去以困我哉！公留助我，我惟無厭其安天下之事。公勿去以廢法，則四方之民其世世享公之德矣。” \par}

周公拜手稽首曰：“王命予來，承保乃文祖受命民，\footnote{拜而後言，許成王留。言王命我來，承安汝文德之祖文王所受命之民，是所以不得去。}越乃光烈考武王，弘朕恭。\footnote{於汝大業之父武王，大使我恭奉其道。敘成王留己意。}孺子來相宅，其大惇典殷獻民，\footnote{少子今所以來相宅於洛邑，其大厚行典常於殷賢人。}亂為四方新闢，作周恭先。\footnote{言當治理天下，新其政化，為四方之新君，為周家見恭敬之王，後世所推先也。}


{\noindent\zhuan\zihao{6}\fzbyks 傳“拜而”至“得去”。正義曰：“拜”是從命之事，故云“拜而後言，許成王留”也。以“退”為去，以“留”為來，故言“王令我來”，來居臣位，為太師也。“承安汝文德之祖文王所受命之民”，天命文王，使為民主，天以民命文王，故民是“文王所受命之民”。“承安”者，承文王之意，安定此民。言王之留己，乃為此事，其事既大，是所以不得去也。 \par}

{\noindent\zhuan\zihao{6}\fzbyks 傳“於汝”至“己意”。正義曰：於汝成王大功業之父武王，王意大使我恭奉其道，敘成王留已之意也。王於文王、武王，皆欲令周公奉其道，安其民，其意一也，周公分言之耳。承安其文王之民，恭奉其武王之道，互相通也。 \par}

{\noindent\zhuan\zihao{6}\fzbyks 傳“少子”至“賢人”。正義曰：“少子”者，呼成王之辭。言我“今所以來相宅於洛邑”者，欲令王居洛，其大厚行典常於殷賢人。而據洛為政，故言“來”。訓“典”為常，故連言“典常”，言其行常道也。周受於殷，故繼之於殷，人有賢性，故稱“賢人”。 \par}

{\noindent\zhuan\zihao{6}\fzbyks 傳“言當”至“推先”。正義曰：\CJKunderwave{易}稱“日新之謂盛德”，雖舊有美政，令王更復新之。言當治理天下，新其政化，為四方之新君，與後人為軌訓，為周家見恭敬之王，後世所推先也。謂周家後世子孫，有德之王被人恭敬推先王。戒成王使為善政，令后王崇重之。 \par}

曰，其自時中乂，萬邦咸休，惟王有成績。\footnote{曰，其當用是土中為治，使萬國皆被美德，如此惟王乃有成功。}予旦以多子越御事，篤前人成烈，答其師，作周孚先。\footnote{我旦以眾卿大夫於御治事之臣，厚率行先王成業，當其眾心，為周家立信者之所推先。}

{\noindent\zhuan\zihao{6}\fzbyks 傳“曰其”至“成功”。正義曰：重以誨王，成其上事，故言“曰”以起之。 \par}

{\noindent\zhuan\zihao{6}\fzbyks 傳“我旦”至“推先”。正義曰：“旦”是周公之名，故自稱“我旦”也。“子”者，有德之稱，大夫皆稱“子”，故以“多子”為眾卿大夫。同欲令成王行善政，為後世賢王所推先。公與群臣盡誠節，為後世賢臣所推先。故欲以眾卿大夫及於御治事之臣，深厚率行先王之業，使當其人眾之心，為周家後世賢臣立信者之所推先也。傳於此不言“後世”，從上省文也。於君言“見恭敬”，於臣言“立信”者，以君尊言人敬，臣卑言自立信，因其所宜以設文也。 \par}

{\noindent\shu\zihao{5}\fzkt “周公”至“孚先”。正義曰：周公拜手稽首，盡禮致敬，許王之留，乃興而為言曰：“王今命我來居臣位，承安汝文德之祖文王所受命之民，令我繼文祖大業，我所以不得去也。又於汝大業父武王，大使我恭奉其道，王意以禮留我，其事甚大，我所以為王留也。”公呼成王云：“小子今所以來相宅於洛邑者，欲其大厚行常道於殷賢人。王當治理天下，新其政化，為四方之新君，為周家後世見恭敬之王所推先也。”重誨王曰：“其當用是土中為治，使萬國皆被美德，如此惟王乃有成功也。”公自稱名曰：“若王居洛邑，則我旦以多眾君子卿大夫等及於御治事之臣，厚率行前人先王成業，使當其眾心，為周家後世人臣立信者之所推先。”言我留輔王，使君臣皆為後世所推先，期於上下俱顯也。 \par}

考朕昭子刑，乃單文祖德，伻來毖殷，乃命寧。\footnote{我所成明子法，乃盡文祖之德，謂典禮也。所以居土中，是文武使己來慎教殷民，乃見命而安之。○單音丹，馬丁但反，信也。}予以秬鬯二卣，曰明禋，拜手稽首，休享。\footnote{周公攝政七年致太平，以黑黍酒二器，明絜致敬，告文武以美享。既告而致政，成王留之。本說之。○秬音巨。鬯，敕亮反，香酒也。卣，由手反，又音由，中樽也。禋音因。}予不敢宿，則禋於文王武王。\footnote{言我見天下太平，則絜告文武，不經宿。}


{\noindent\zhuan\zihao{6}\fzbyks 傳“我所”至“安之”。正義曰：典禮治國，事資聖人，前聖後聖，其終一揆，故言所欲成明子之法，乃盡是汝祖文王之德也。“子”斥成王，下句並告文武，兼用武王可知。又述居洛邑之意，所以居土中者，是文武使己來居此地，周公自言非己意也。文武令我營此洛邑，欲使居土中，慎教殷民，乃是見命於文武而安殷民也。顧氏云：“文武使我來慎教殷民，我今受文武之命以安民也。” \par}

{\noindent\zhuan\zihao{6}\fzbyks 傳“周公”至“說之”。正義曰：\CJKunderwave{康誥}之作，事在七年，雲“四方民大和會”。“和會”即太平之驗，是“周公攝政七年致太平”也。\CJKunderwave{釋草}云：“秬，黑黍。”\CJKunderwave{釋器}云：“卣，中樽也。”以黑黍為酒，煮鬱金之草，築而和之，使芬香調暢，謂之“秬鬯”。鬯酒二器，明絜致敬，告文王武王以美享，謂以太平之美事享祭也。\CJKunderwave{國語}稱：“精意以享謂之禋。”\CJKunderwave{釋詁}云：“禋，敬也。”是明“禋”為“明絜致敬”也。太平是王之美事，故太平告廟是以美享祭也。公既告太平而致政成王，成王留之，故本而說之此事者，欲令成王重其事，厚行之。\CJKunderwave{周禮}鬱鬯之酒實之於彝，此言在卣者，\CJKunderwave{詩·大雅·江漢}及\CJKunderwave{文侯之命}皆言“秬鬯一卣，告於文人”，則未祭實之於卣，祭時實之於彝。彼“一卣”，此“二卣”者，此一告文王，一告武王；彼王賜臣使告其太祖，故惟一卣耳。此經“卣”下言“曰”者，說本盛酒於樽，乃為此辭，故言“曰”也。 \par}

{\noindent\zhuan\zihao{6}\fzbyks 傳“言我”至“經宿”。正義曰：此申述上“明禋”之事，言我見天下太平，則絜告文武，不敢經宿，示虔恭之意也。此三月營洛邑，民已和會，則三月之時已太平矣。既告而致政，則告在歲末，而云“不經宿”者，蓋周公營洛邑，至冬始成，得還鎬京，即告文武，是為“不經宿”也。且太平非一日之事，公雲“不經宿”者，示虔恭之意耳，未必旦見太平，即此日告也。\CJKunderline{鄭玄}以“文祖”為明堂，“曰明禋者，六典成祭於明堂，告五帝太皞之屬也”。既告明堂，則復禋於文武之廟，告成洛邑。 \par}

惠篤敘，無有遘自疾。萬年厭乃德，殷乃引考。\footnote{汝為政當順典常，厚行之使有次序，無有遇用患疾之道者，則天下萬年厭於汝德，殷乃長成為周。○遘，工豆反。厭，於豔反，注同，馬云：“厭，飲也。”徐於廉反。}王伻殷乃承敘萬年，其永觀朕子懷德。”\footnote{王使殷民上下相承有次序，則萬年之道，民其長觀我子孫而歸其德矣勉使終之。}

{\noindent\zhuan\zihao{6}\fzbyks 傳“汝為”至“為周”。正義曰：\CJKunderwave{釋言}云：“惠，順也。”此經述上“惇典”，故言“汝為政當順典常，厚行之使有次序”。\CJKunderwave{釋詁}云：“遘，遇也。”患疾之道謂虐政，使人患疾之。厚行典常，使有次序，則百官諸侯凡為政者皆無有遇用患疾之政以害下民，則經歷萬年厭飽於汝德，則殷國乃長成為周。 \par}

{\noindent\zhuan\zihao{6}\fzbyks 傳“王使”至“終之”。正義曰：上言天下民萬年厭飽王德，此教為王德，使萬年令民厭飽王德也。能使殷民上下有次序，則王德堪至萬年之道。王之子孫當行不怠，則民其長觀我子孫，知其有德，而歸其德矣，此則長成為周。勸勉王使終之。 \par}

{\noindent\shu\zihao{5}\fzkt “考朕”至“懷德”。正義曰：周公又說制禮授王，使王奉之。我所成明子之法，乃盡是汝文祖之德，言用文王之道制禮，其事大不可輕也。又言所以須善治殷獻民者，文武使己來居土中，慎教殷民，乃是見命於文武而安之故也。制典當待太平，我以時既太平，即以秬黍鬯酒，盛於二卣樽內，我言曰：“當以此酒須明絜致敬於文武，我則拜手稽首，告文武以美享。”告云：“今太平，即速告廟，我不敢經宿，則禋告文王武王以致太平之事。”汝王為政，當順典常厚行之，使有次序，則諸為政者無雲有遇用患疾之道苦毒下民，則天下萬年厭飽於汝王之德，殷乃長成為周。王使殷民上下相承有次序，則萬年之道，下民其長觀我子孫而歸其德矣。勸王使終之，皆是誨王之言也。 \par}

戊辰，王在新邑，\footnote{成王既受周公誥，遂就居洛邑，以十二月戊辰晦到。○王在新邑，馬孔絕句，鄭讀“王在新邑烝”。}烝祭歲，文王騂牛一，武王騂牛一。王命作冊，逸祝冊，惟告周公其後。\footnote{明月，夏之仲冬，始於新邑烝祭，故曰“烝祭歲”。古者褒德賞功，必於祭日，示不專也。特加文武各一牛，告曰尊周公，立其後為魯侯。}王賓，殺禋，咸格，王入太室祼。\footnote{王賓異周公，殺牲精意以享文武，皆至其廟親告也。太室，清廟。祼鬯告神。○王賓，絕句。殺禋，絕句，一讀連“咸格”絕句。太室，馬云：“廟中之夾室。”祼，官喚反。}


{\noindent\zhuan\zihao{6}\fzbyks 傳“成王”至“晦到”。正義曰：周公告成王令居洛邑為治，王既受周公之誥，遂東行就居洛邑，以十二月戊辰晦日到洛。指言“戊辰，王在新邑”，知其晦日始到者，此歲入戊午蔀五十六年，三月雲丙午朏，以算術計之，三月甲辰朔大，四月甲戌朔小，五月癸卯朔大，六月癸酉朔小，七月壬寅朔大，八月壬申朔小，九月辛丑朔大，又有閏九月辛未朔小，十月庚子朔大，十一月庚午朔小，十二月己亥朔大，計十二月三十日戊辰晦到洛也。 \par}

{\noindent\zhuan\zihao{6}\fzbyks 傳“明月”至“魯侯”。正義曰：下雲“在十有二月”者，周之十二月，建亥之月也。戊辰是其晦日，故明日即是夏之仲冬建子之月也。言“明月”者，此烝祭非朔日，故言“月”也。自作新邑已來，未嘗於此祭祀，此歲始於新邑烝祭，故曰“烝祭歲”也。\CJKunderwave{周禮·大司馬}仲冬教大閱，遂以享烝是也。王者冬祭，必用仲月，此是周之歲首，故言“歲”耳。王既戊辰晦到，又須戒日致齊，不得以朔日即祭之。\CJKunderwave{祭統}云：“古者明君爵有德而祿有功，必賜爵祿於太廟，示不敢專也。”故云“古者褒德賞功，必於祭日，示不專也”。因封之，特設祭烝之禮。宗廟用大牢，此文武皆言“牛一”，知於太牢之外特加一牛。告白文武之神，言為尊周公，立其後為魯侯。\CJKunderwave{魯頌}所云“王曰叔父，建爾元子，俾侯於魯”，是此時也。“王命作策”者，命有司作策書也。讀策告神謂之“祝”，“逸祝策”者，使史逸讀策書也。\CJKunderline{鄭玄}以“烝祭”上屬。“歲文王騂牛一”者，“歲”是成王元年，正月朔日，特告文武封周公也。案\CJKunderwave{周頌·烈文序}云：“成王即政，諸侯助祭。”鄭箋云：“新王即政，必以朝享之禮祭於祖考，告嗣位也。”則鄭意以朝享之後，特以二牛告文武，封周公之後，與孔義不同。 \par}

{\noindent\zhuan\zihao{6}\fzbyks 傳“王賓”至“告神”。正義曰：“王賓異周公”者，王尊周公為賓，異於其臣。王肅雲“成王尊周公，不敢臣之，以為賓，故封其子”是也。\CJKunderwave{周語}云：“精意以享謂之禋。”既殺二牲，精誠其意以享祭文武。“咸”，皆也。“格”，至也。“皆至其廟”，言王重其事，親告之也。“太室”，室之大者，故為清廟。廟有五室，中央曰“太室”。王肅云：“太室，清廟中央之室。”“清廟”，神之所在，故王入太室祼獻鬯酒以告神也。“祼”者，灌也。王以圭瓚酌鬱鬯之酒以獻屍，屍受祭而灌於地，因奠不飲謂之“祼”。\CJKunderwave{郊特牲}雲“既灌，然後迎牲”，則殺在祼後。此經先言“殺”，後言“祼”者，“殺”者、“咸格”錶王敬公之意，非行事之次也。其“王入太室祼”，乃是祭時行事耳。周人尚臭，祭禮以祼為重，故言王祼。其封伯禽，乃是祭之將末，非祼時也。\CJKunderwave{祭統}賜臣爵祿之法云：“祭之日，一獻，君降立於阼階之南，南向，所命者北面，史由君右，執策命之。”鄭云：“一獻，一酳屍也。”\CJKunderwave{禮}酳屍，屍獻而祭畢，是祭末乃命之，以祼為重，故特言之。 \par}

王命周公後，作冊逸誥。\footnote{王為冊書，使史逸誥伯禽封命之書，皆同在烝祭日，周公拜前，魯公拜後。}在十有二月，惟周公誕保文武受命，惟七年。\footnote{言周公攝政盡此十二月，大安文武受命之事，惟七年，天下太平。自“戊辰”已下，史所終述。○受命，絕句，馬同。惟七年，周公攝政七年，天下太平，馬同；鄭云：“文王武王受命及周公居攝皆七年。”}

{\noindent\zhuan\zihao{6}\fzbyks 傳“王為”至“拜後”。正義曰：“王為策書”，亦命有司為之也。上雲“作策”，作告神之策。此言“作策”，誥伯禽之策。祭於神謂之“祝”，於人謂之“誥”，故云“使史逸誥伯禽封命之書”。封康叔謂之\CJKunderwave{康誥}，此命伯禽，當雲\CJKunderwave{伯禽之誥}。定四年\CJKunderwave{左傳}雲“命以伯禽”，即史逸所讀之策也。上言“逸祝策”，此“誥”下不言“策”者，“祝”是讀書之名，故上雲“祝策”；此“誥”是誥伯禽使知，雖復讀書以誥之，不得言“誥策”也。上告周公，其後已言告神封周公，嫌此“逸誥”以他日告之，故云“皆同在烝祭日”。以\CJKunderwave{祭統}言一獻命之，知此亦祭日也。文十三年\CJKunderwave{公羊傳}曰：“封魯公以為周公也，周公拜乎前，魯公拜乎後。曰生以養周公，死以為周公主。” \par}

{\noindent\zhuan\zihao{6}\fzbyks 傳“言周”至“終述”。正義曰：自“戊辰”已上，周公與成王相對語，未有致政年月，故史於此總結之。自“戊辰”已下非是王與周公之辭，故辨之雲“史所終述”也。 \par}

{\noindent\shu\zihao{5}\fzkt “戊辰”至“七年”。正義曰：自此以下，史終述之。周公歸政，成王既受言誥之，王即東行赴洛邑。其年十二月晦戊辰日，王在新邑。後月是夏之仲冬，為冬節烝祭，其月節是周之歲首，特異常祭，加文王騂牛一，武王騂牛一。王命有司作策書，乃使史官名逸者祝讀此策，惟告文武之神，言周公有功，宜立其後為國君也。其時王尊異周公，以為賓，殺牲享祭文王武王，皆親至其廟，王入廟之大室，行祼鬯之禮。言其尊異周公而禮敬深也。於此祭時，王命周公後，令作策書，使逸讀此策辭以告伯禽，言封之於魯，命為周公後也。又總述之，在十有二月，惟周公大安文武受命之事，於此時惟攝攻七年矣。 \par}

%%% Local Variables:
%%% mode: latex
%%% TeX-engine: xetex
%%% TeX-master: "../Main"
%%% End:

%% -*- coding: utf-8 -*-
%% Time-stamp: <Chen Wang: 2024-04-02 11:42:41>

% {\noindent \zhu \zihao{5} \fzbyks } -> 注 (△ ○)
% {\noindent \shu \zihao{5} \fzkt } -> 疏

\chapter{卷十六}


\section{多士第十六}


成周既成，\footnote{洛陽下都。}遷殷頑民，\footnote{殷大夫士心不則德義之經，故徙近王都教誨之。○不則，如字，或作測，非。近，附近之近。}\CJKunderline{周公}以王命誥，\footnote{稱\CJKunderline{成王}命告令之。}作\CJKunderwave{多士}。

多士\footnote{所告者即眾士，故以名篇。}


{\noindent\zhuan\zihao{6}\fzbyks 傳“洛陽下都”。正義曰：周之成周，於漢為洛陽也。洛邑為王都，故謂此為“下都”。遷殷頑民，以成周道，故名此邑為成周。 \par}

{\noindent\zhuan\zihao{6}\fzbyks 傳“殷大”至“誨之”。正義曰：經雲“商王士”、“殷遺多士”，皆非民，事謂之“頑民”，知是殷之大夫士也。經止雲“士”，而知有大夫者，以經雲“迪簡在王庭，有服在百僚”，其意言將任為王官，以為大臣，不惟告士而已，故知有大夫也。“士”者，在官之總號，故言“士”也。“心不則德義之經”，僖二十四年\CJKunderwave{左傳}文，引之以解釋“頑民”之意。經雲“移爾遐逖，比事臣我宗多遜”，是言“徙近王都教誨之”也。\CJKunderwave{漢書·地理志}及賈逵注\CJKunderwave{左傳}，皆以為遷邶鄘之民於成周，分衛民為三國。計三國俱是從叛，何以獨遷邶鄘？邶鄘在殷畿三分有二，其民眾矣，非一邑能容，民謂之為“士”，其名不類，故孔意不然。 \par}

{\noindent\shu\zihao{5}\fzkt “成周”至“多士”。正義曰：成周之邑既成，乃遷殷之頑民，令居此邑。“頑民”謂殷之大夫士從\CJKunderline{武庚}叛者，以其無知，謂之“頑民”。民性安土重遷，或有怨恨，\CJKunderline{周公}以\CJKunderline{成王}之命誥此眾士，言其須遷之意。史敘其事，作\CJKunderwave{多士}。 \par}

惟三月，\CJKunderline{周公}初於新邑洛，用告商王士。\footnote{\CJKunderline{周公}致政明年三月，始於新邑洛，用王命告商王之眾士。}


{\noindent\zhuan\zihao{6}\fzbyks 傳“\CJKunderline{周公}”至“眾士”。正義曰：以\CJKunderwave{洛誥}之文，成周與洛邑同時成也。王以\CJKunderline{周公}攝政七年十二月來至新邑，明年即政，此篇繼王居洛之後，故知是“致政明年之三月”也。成周南臨洛水，故云“新邑洛”。\CJKunderline{周公}既以致政在王都，故“新邑”，成周。以\CJKunderline{成王}之命告商王之眾士，鄭雲“\CJKunderline{成王}元年三月，\CJKunderline{周公}自王城初往成周之邑，用\CJKunderline{成王}命告殷之眾士以撫安之”是也。 \par}

{\noindent\shu\zihao{5}\fzkt “惟三月”至“王士”。正義曰：惟\CJKunderline{成王}即政之明年三月，\CJKunderline{周公}初始於所造新邑之洛，用\CJKunderline{成王}之命告商王之眾士。言\CJKunderline{周公}親至成周，告新來者。 \par}

王若曰：“爾殷遺多士，\footnote{順其事稱以告殷遺餘眾士。所順在下。}弗吊，旻天大降喪於殷。\footnote{稱天以愍下，言愍道至者，殷道不至，故旻天下喪亡於殷。○吊音的。旻天，上閔巾反，仁覆愍下謂之旻，馬云：“秋日旻天。秋，氣殺也，方言‘降喪’，故稱旻天也。”愍，眉隕反。喪，息浪反。}我有周佑命，將天明威，\footnote{言我有周受天佑助之命，故得奉天明威。}致王罰，敕殷命終於帝。\footnote{天命周致王者之誅罰，王黜殷命，終周於帝王。}肆爾多士，非我小國敢弋殷命。\footnote{天佑我，故汝眾士臣服我。弋，取也。非我敢取殷王命，乃天命。○弋，徐音翼，馬本作翼，義同。}惟天不畀允罔固亂，弼我，我其敢求位？\footnote{惟天不與信無堅固治者，故輔佑我，我其敢求天位乎？○治，直吏反。畀，必利反，下同。}惟帝不畀，惟我下民秉為，惟天明畏。\footnote{惟天不與紂，惟我周家下民秉心為我，皆是天明德可畏之效。○為，於偽反。畏如字，一音威。}

{\noindent\zhuan\zihao{6}\fzbyks 傳“順其”至“在下”。正義曰：順其殷亡之事，稱王命以告之。從紂之臣，或有身已死者，遺餘在者遷於成周，故“告殷遺餘眾士”。“所順在下”，下文皆是順之辭。 \par}

{\noindent\zhuan\zihao{6}\fzbyks 傳“稱天”至“於殷”。正義曰：此經先言“弗吊”，謂殷道不至也。“不至”者，上不至天，事天不以道；下不至民，撫民不以理也。天有多名，獨言“旻天”者，“旻”，愍也，“稱天以愍下”，言天之所愍，愍道至者也。“殷道不至，故旻天下喪亡於殷”，言將覆滅之。 \par}

{\noindent\zhuan\zihao{6}\fzbyks 傳“天命”至“帝王”。正義曰：“天命周致王者之誅罰”，謂奉上天之命，殺無道之王，此乃王者之事，故為“王者之誅罰”。“敕”訓正也，“正黜殷命”謂殺去虐紂，使周受其終事，是“終周於帝王”。“終”猶舜受堯終，言殷祚終而歸於周。 \par}

{\noindent\zhuan\zihao{6}\fzbyks 傳“天祐”至“天命”。正義曰：“肆”訓故也，直雲故爾多士，辭無所結，此經大意，敘其去殷事周，知其故爾眾士言其臣服我。“弋”，射也，射而取之，故“弋”為取也。\CJKunderline{鄭玄}、王肅本“弋”作翼，王亦云：“翼，取也。”鄭云：“翼猶驅也，非我周敢驅取汝殷之王命。”雖訓為驅，亦為取義。周本殷之諸侯，故\CJKunderline{周公}自稱“小國”。 \par}

{\noindent\shu\zihao{5}\fzkt “王若”至“明畏”。正義曰：\CJKunderline{周公}以王命順其事而呼之曰：“汝殷家遺餘之眾士，汝殷家道教不至，旻天以殷道不至之故，天下喪亡於殷，將欲滅殷。我有周受天祐助之命，奉天明白之威，致王者之誅罰，正黜殷命，終我周家於帝王之事，謂使我周家代殷為天子也。天既助我周王，故汝眾士來為我臣。由天助我，我得為之，非我小國敢取殷之王命以為己有。此乃天與我，惟天不與信無堅固於治者，以是故輔弼我。若其不然，我其敢妄求天子之位乎？”言此位天自與我，非我求而得之。“惟天不與紂，故惟我周家下民秉心為我，故我得之，惟天明德可畏之效也。亦既得喪由天，汝等不得不服”。以殷士未服，故以天命喻之。 \par}

我聞曰：‘上帝引逸。’有夏不適逸，則惟帝降格。\footnote{言上天欲民長逸樂，有\CJKunderline{夏桀}為政不之逸樂，故天下至戒以譴告之。○樂音洛，下同。譴，棄戰反。}向於時夏，弗克庸帝，大淫泆有辭。\footnote{天下至戒，是向於時夏，不背棄。桀不能用天戒，大為過逸之行，有惡辭聞於世。○時夏，絕句，馬以“時”字絕句。向，許亮反。泆音逸，又作佾，注同；馬本作屑，云：“過也。”背音佩。行，下孟反。}惟時天罔念聞，厥惟廢元命，降致罰。\footnote{惟是桀惡有辭，故天無所念聞，言不佑，其惟廢其天命，下致天罰。}乃命爾先祖\CJKunderline{成湯}革夏，俊民甸四方。\footnote{天命湯更代夏，用其賢人治四方。○甸，徒遍反。}


{\noindent\zhuan\zihao{6}\fzbyks 傳“言上”至“告之”。正義曰：襄十四年\CJKunderwave{左傳}稱：“天之愛民甚矣。”又曰：“天生民而立之君，使司牧之。”是言上天欲民長得逸樂，故立君養之，使之長逸樂也。\CJKunderline{夏桀}為政，割剝夏邑，使民不得之適逸樂，故上天下此至戒以譴告之。“降”，下。“格”，至也。直言下至，明是“天下至戒”。天所下戒，惟下災異以譴告人主，使之見災而懼，改修德政耳。古書亡失，桀之災異未得盡聞。 \par}

{\noindent\zhuan\zihao{6}\fzbyks 傳“惟是”至“天罰”。正義曰：桀惡流毒於民，乃有惡辭聞於世。惡既有辭，是惡已成矣。惟是桀惡有辭，故天無所念聞，言天不愛念，不聽聞，是其全棄之，不佑助也。棄而不佑，則當更求賢主。“其惟廢大命”，欲奪其王位也。“下致天罰”，欲殺其凶身也。廢大命，知“降致”是下罰也。 \par}

{\noindent\shu\zihao{5}\fzkt “我聞”至“四方”。正義曰：既言天之效驗，法惡與善，更追說往事，比而喻之：“我聞人有言曰：‘上天之情，欲民長得逸樂。’而有夏王桀逆天害民，不得使民之適逸樂。以此則惟上天下災異至戒以譴告之，欲使夏王桀覺悟，改惡為善，是天歸鄉於是夏家，不背棄之。而\CJKunderline{夏桀}不能用天之明戒，改悔己惡，而反大為過逸之行，致有惡辭以聞於世。惟是桀有惡辭，故天無復愛念，無復聽聞，言天不復助桀。其惟廢其大命，欲絕夏祚也。下致天罰，欲誅桀身也。乃命汝先祖\CJKunderline{成湯}，使之改革夏命，用其賢俊之人，以治四方之國。舉桀滅湯興以譬之。 \par}

自\CJKunderline{成湯}至於帝乙，罔不明德恤祀。\footnote{自帝乙以上，無不顯用有德，憂念齊敬，奉其祭祀。言能保宗廟社稷。○上，時掌反。齊，側皆反。}亦惟天丕建保乂有殷，殷王亦罔敢失帝，罔不配天其澤。\footnote{湯既革夏，亦惟天大立安治於殷。殷家諸王皆能憂念祭祀，無敢失天道者，故無不配天布其德澤。}在今後嗣王，誕罔顯於天，矧曰其有聽念於先王勤家？\footnote{後嗣王紂，大無明於天道，行昏虐，天且忽之，況曰其有聽念先祖、勤勞國家之事乎？}誕淫厥泆，罔顧於天顯民祗，\footnote{言紂大過其過，無顧於天，無能明人為敬，暴亂甚。}惟時上帝不保，降若茲大喪。\footnote{惟是紂惡，天不安之，故下若此大喪亡之誅。○喪，息浪反。}惟天不畀不明厥德，凡四方小大邦喪，罔非有辭於罰。”\footnote{惟天不與不明其德者，故凡四方小大國喪滅，無非有辭於天所罰。言皆有闇亂之辭。}


{\noindent\zhuan\zihao{6}\fzbyks 傳“自帝”至“社稷”。正義曰：下篇說中宗、\CJKunderline{高宗}、祖甲三王以外，其後立王，生則逸豫，亦罔或能壽。如彼文，則帝乙以上非無僻王，而此言無不顯用有德，憂念祭祀者，立文之法，辭有抑揚，方說紂之不善，盛言前世皆賢，正以守位不失，故得美而言之。憂念祭祀者，惟有齊肅恭敬，故言“憂念齊敬，奉其祭祀”。言能保宗廟社稷，為天下之主，以見紂不恭敬，故喪亡之。 \par}

{\noindent\zhuan\zihao{6}\fzbyks 傳“湯既”至“德澤”。正義曰：帝乙已上諸王，所以長處天位者，皆由湯之聖德延及後人。“湯既革夏，亦惟天大立安治於殷”者，謂天安治之，故殷家得治理也。殷家諸王自\CJKunderline{成湯}之後，皆能憂念祭祀，無敢失天道者，故得常處王位，無不配天布其德澤於民。為天之子，是“配天”也。號令於民，是“佈德”也。 \par}

{\noindent\zhuan\zihao{6}\fzbyks 傳“言紂”至“亂甚”。正義曰：“淫”、“泆”俱訓為過，“言紂大過其愆過，無顧於天”，言其縱心為惡，不畏天也。“無能明民為敬”，言其多行虐政，不憂民也。不畏於天，不愛於民，言其“暴亂甚”也。此經顧“於天”與“顯民祗”，共蒙上“罔”文，故傳再言“無”也。 \par}

{\noindent\zhuan\zihao{6}\fzbyks 傳“惟天”至“之辭”。正義曰：能明其德，天乃與之，惟天不與不明其德者，紂不明其德，故天喪之。因即廣言天意，凡四方小大邦國，謂諸侯有土之君，其為天所喪滅者，無非皆有惡辭聞於天，乃為上天所罰。言被天罰者，皆有闇亂之辭，上天不罰無辜，紂有闇亂之辭，故天滅之耳。天既滅不明其德，我有明德，為天所立，汝等殷士安得不服我乎？以其心仍不服，故以天道責之。 \par}

{\noindent\shu\zihao{5}\fzkt “自成”至“於罰”。正義曰：既言命湯革夏，又說後世皆賢，至紂始惡，天乃滅之。自\CJKunderline{成湯}至於帝乙，無不顯用有德，憂念祭祀，後世亦賢。非獨\CJKunderline{成湯}以用其行合天意，亦惟天大立安治有殷。殷家諸王皆能明德憂祀，亦無敢失天道者，無不皆配天而布其德澤，以此得天下久為民主。在今後嗣王紂，大無明於天道，敢行昏虐之政於天。天猶且忽之，況曰其有聽念先王父祖、勤勞國家之事乎？乃復大淫過其泆，無所顧於上天，無能明民為敬，以此反於先王，違逆天道。惟是上天不安紂之所為，下若此大喪亡之誅，惟天不與不明其德之人故也。天不與惡，豈獨紂乎？凡四方諸侯小大邦國，其喪滅者，無非皆有惡辭，是以致至於天罰。汝紂以惡而見滅，汝何以不服我也？ \par}

王若曰：“爾殷多士，今惟我周王，丕靈承帝事，\footnote{周王，文武也。大神奉天事，言明德恤祀。}有命曰：‘割殷，告敕於帝。’\footnote{天有命，命周割絕殷命，告正於天。謂既克紂，柴於牧野，告天不頓兵傷士。}惟我事不貳適，惟爾王家我適。\footnote{言天下事已之我周矣，不貳之佗，惟汝殷王家已之我，不復有變。○復，扶又反。}予其曰：‘惟爾洪無度，我不爾動，自乃邑。’\footnote{我其曰，惟汝大無法度，謂紂無道。我不先動誅汝，亂從汝邑起。言自召禍。}予亦念天即於殷大戾，肆不正。”\footnote{我亦念天就於殷大罪而加誅者，故以紂不能正身念法。}


{\noindent\zhuan\zihao{6}\fzbyks 傳“周王”至“恤祀”。正義曰：\CJKunderline{文王}受命，\CJKunderline{武王}伐紂，故知“周王”兼文武也。“大神奉天事”，謂以天為神而勤奉事之，勞身敬神，言亦如湯明德恤祀也。 \par}

{\noindent\zhuan\zihao{6}\fzbyks 傳“天有”至“傷士”。正義曰：以周王奉天之故，故天有命，命我周使割絕殷命，告正於天。謂\CJKunderwave{武成}之篇所云，既克紂，柴於牧野，告天不頓兵傷士是也。前敵即服，故無“頓兵傷士”。師以正行，故為“告正”。\CJKunderwave{武成}正告功成，功成無害，即是不頓傷也。“頓兵”者，昭十五年\CJKunderwave{左傳}文“頓折”也。 \par}

{\noindent\zhuan\zihao{6}\fzbyks 傳“我亦”至“念法”。正義曰：言“我亦念天”者，以紂雖無法度，若使天不命我，我亦不往誅紂。以紂既為大惡，上天命我，我亦念天所遣。我就殷加大罪者何故？以紂不能正身念法也。 \par}

{\noindent\shu\zihao{5}\fzkt “王若”至“不正”。正義曰：\CJKunderline{周公}又稱王順而言曰：“汝殷眾士，今惟我周家文武二王，大神能奉天事，故天有命，命我周王曰：‘當割絕殷命，告正於天。’我受天命，已滅殷告天，惟我天下之事，不有二處之適。”言己之適周，不更適他也。“惟汝殷王家事亦於我之適，不復變改”。又追說初伐紂之事：“我其為汝言曰：‘惟汝殷紂大無法度，故當宜誅絕之。伐紂之時，我不先於汝動，自往誅汝，其亂從汝邑先起，汝紂自召禍耳。’我亦念天所以就於殷致大罪者，故以紂不能正身念法故也。” \par}

王曰：“猷告爾多士，予惟時其遷居西爾。\footnote{以道告汝眾士，我惟汝未達德義，是以徙居西汝於洛邑，教誨汝。}非我一人奉德不康寧，時惟天命。\footnote{我徙汝，非我天子奉德，不能使民安之，是惟天命宜然。}無違，朕不敢有後，無我怨。\footnote{汝無違命，我亦不敢有後誅，汝無怨我。}惟爾知，惟殷先人，有冊有典，殷革夏命。\footnote{言汝所親知，殷先世有冊書典籍。說殷改夏王命之意。}今爾又曰：‘夏迪簡在王庭，有服在百僚。’\footnote{簡，大也。今汝又曰：“夏之眾士蹈道者，大在殷王庭，有服職在百官。”言見任用。}予一人惟聽用德，肆予敢求爾於天邑商。\footnote{言我周亦法殷家，惟聽用有德，故我敢求汝於天邑商，將任用之。}予惟率肆矜爾，非予罪，時惟天命。”\footnote{惟我循殷故事，憐愍汝，故徙教汝，非我罪咎，是惟天命。}

{\noindent\zhuan\zihao{6}\fzbyks 傳“以道”至“誨汝”。正義曰：“猷”訓道也，故云“以道告汝眾士”。上言“惟是”，不言其故，故傳辨之，惟是者，未達德義也。遷使居西，正欲教以德義，是以徙居西汝置於洛邑，近於京師教誨汝也。從殷適洛，南行而西回，故為“居西”也。 \par}

{\noindent\zhuan\zihao{6}\fzbyks 傳“汝無”至“怨我”。正義曰：周既伐紂，又誅\CJKunderline{武庚}，殷士懼更有誅，疑其欲違上命，故設此言以戒之。知“無違朕”者，謂戒之使汝無違命也。汝能用命，我亦不敢有後誅，必無後誅，汝無怨我也。 \par}

{\noindent\zhuan\zihao{6}\fzbyks 傳“言我”至“用之”。正義曰：夏人簡在王庭，為其有德見用。言我亦法殷家，惟聽用有德，汝但有德，我必任用。故我往前敢求汝有德之人於天邑商都，將任用之也。\CJKunderline{鄭玄}云：“言天邑商者，亦本天之所建。”王肅云：“言商今為我之天邑。”二者其言雖異，皆以“天邑商”為殷之舊都。言未遷之時，當求往，遷後有德任用之必矣。 \par}

{\noindent\zhuan\zihao{6}\fzbyks 傳“惟我”至“天命”。正義曰：“循殷故事”，此“故”解經中“肆”字，謂殷用夏人，我亦用殷人。“憐愍汝，故徙之教汝”，此“故”解義之言，非經中“肆”。遷汝來西者，非我罪咎，是惟天命也。 \par}

{\noindent\shu\zihao{5}\fzkt “王曰猷”至“大命”。正義曰：又言曰我以道告汝眾士，我惟是以汝未達德義之故，其今徙居西汝置於洛邑，以教誨汝。我之徙汝，非我一人奉行德義，不能使民安而安之，是惟天命宜然。汝無違我，我亦不敢更有後誅罰，汝等無於我見怨。汝既來遷，當為善事。惟汝所親知，惟汝殷先人往世有策書，有典籍，說殷改夏王命之意，汝當案省知之。汝知先人之故事，今往又有言曰：‘夏之諸臣蹈道者，大在殷王之庭，有服行職事，在於百官。’言其見任用，恐我不任汝。我一人惟聽用有德之者，故我敢求汝有德之人於彼天邑商都，欲取賢而任用之。我惟循殷故事，憐愍汝，故徙教汝。此徙非我有罪，是惟天命當然。”聖人動合天心，故每事惟託天命也。 \par}

王曰：“多士，昔朕來自奄，予大降爾四國民命。\footnote{昔我來從奄，謂先誅三監，後伐奄淮夷。民命謂君也。大下汝民命，謂誅四國君。}我乃明致天罰，移爾遐逖，比事臣我宗，多遜。”\footnote{四國君叛逆，我下其命，乃所以明致天罰。今移徙汝於洛邑，使汝遠於惡俗，比近臣我宗周，多為順道。○逖，他力反。比，毗志反，注同。遠，於萬反。}


{\noindent\zhuan\zihao{6}\fzbyks 傳“昔我”至“國君”。正義曰：\CJKunderwave{金縢}之篇說\CJKunderline{周公}東征，言“居東二年，罪人斯得”，則“昔我來從奄”者，謂攝政三年時也。於時王不親行，而王言“我來自奄”者，\CJKunderline{周公}以王命誅四國，\CJKunderline{周公}師還，亦是王來還也。一舉而誅四國，獨言“來自奄”者，謂先誅三監，後伐奄與淮夷，奄誅在後，誅奄即來，故言“來自奄”也。民以君為命，故“民命謂君也”。大下汝民命，謂誅四國君。王肅云：“君為民命，為君不能順民意，故誅之也。” \par}

{\noindent\zhuan\zihao{6}\fzbyks 傳“四國”至“順道”。正義曰：王之所罰，罰有罪也。四國之君，有叛逆之罪。“我下其命，乃所以明致天罰”，言非苟為之也。“遐”、“逖”俱訓為遠。“今移徙汝於洛邑”，令去本鄉遠也。“使汝遠於惡俗”，令去惡俗遠也。比近京師，臣我周家，使汝從我善化，多為順道，所以救汝之性命也。 \par}

{\noindent\shu\zihao{5}\fzkt “王曰多士”至“多遜”。正義曰：王復言曰：“眾士，昔我來從奄國，大黜下汝管蔡商奄四國民命。民之性命，死生在君，誅殺其君，是下民命。由四國叛逆，我乃明白致行天罰。汝等遣餘，當教之為善，故移徙汝居於遠。令汝遠於惡俗，比近服事臣我宗周，多為順道。翼汝相教為善，永不為惡也。” \par}

王曰：“告爾殷多士，今予惟不爾殺，予惟時命有申。\footnote{所以徙汝，是我不欲殺汝，故惟是教命申戒之。}今朕作大邑於茲洛，予惟四方罔攸賓，\footnote{今我作此洛邑，以待四方，無有遠近，無所賓外。○賓如字，徐音殯，馬云：“卻也。”}亦惟爾多士，攸服奔走臣我，多遜。\footnote{非但待四方，亦惟汝眾士，所當服行奔走臣我，多為順事。}爾乃尚有爾土，爾乃尚寧幹止。\footnote{汝多為順事，乃庶幾還有汝本土，乃庶幾安汝故事止居。以反所生誘之。爾克敬，天惟畀矜爾。○汝能敬行順事，則為天所與，為天所憐。}爾不克敬，爾不啻不有爾土，予亦致天之罰於爾躬。\footnote{汝不能敬順，其罰深重，不但不得還本土而已，我亦致天罰於汝身。言刑殺。○啻，始豉反，徐本作翅，音同，下篇放此。}今爾惟時宅爾邑，繼爾居，爾厥有幹有年於茲洛。\footnote{今汝惟是敬順居汝邑，繼汝所當居為，則汝其有安事，有豐年於此洛。邑言由洛修善，得還本土，有幹有年。}爾小子乃興，從爾遷。”\footnote{汝能敬，則子孫乃起從汝化而遷善。}


{\noindent\zhuan\zihao{6}\fzbyks 傳“今汝”至“有年”。正義曰：殷士遠離本鄉，新來此邑，或當居不安，為棄舊業，故戒之。“今汝惟是敬順，居汝新所受邑，繼汝舊日所當居為”，謂繼其本土之事業也。但能如此，得還本土，其有安事，有豐年也。“有幹有年”謂歸本土。有幹年而言於洛者，言“由在洛修善，得還本土，有幹有年”也。王肅云：“汝其有安事，有長久年於此洛邑。”王解於文甚便，但孔上句為雲“爾乃尚有爾本土”，是誘引之辭，故止為得“還本土，有幹有年”也。 \par}

{\noindent\shu\zihao{5}\fzkt “王曰告”至“爾遷”。正義曰：王又言曰：“告汝殷之多士，所以遠徙汝者，今我惟不欲於汝刑殺，我惟是教命有所申戒由此也。今我作大邑於此洛，非但為我，惟以待四方，無所賓外，亦惟為汝眾士所當服行臣事我宗周，多為順事故也。汝若多為順事，汝乃庶幾還有汝本土，乃庶幾安汝故事止居，可不勉之也？汝能敬行順事，天惟與汝憐汝，況於人乎？汝若不能敬行順事，則汝不啻不得還汝本土，我亦致天之罰於汝身。今汝惟是敬順，居汝所受新邑，繼汝舊日所居為，我當聽汝還歸本鄉，有幹事，有豐年，乃由於此洛邑行善也。汝能敬順，則汝之小子與孫等，乃起從汝化而遷善矣。” \par}

王曰：“又曰時予，乃或言，爾攸居。”\footnote{言汝眾士當是我，勿非我也。我乃有教誨之言，則汝所當居行。}


{\noindent\zhuan\zihao{6}\fzbyks 傳“言汝”至“居行”。正義曰：王以誨之已終，故戒之云：“汝當是我，勿非我。既不非我，我乃有教誨汝之言，則汝所當居行。”令其居於心而行用之。\CJKunderline{鄭玄}\CJKunderwave{論語}注云“或之言有”，此亦“或”為有也。凡言“王曰”，皆是史官錄辭，非王語也。今史錄稱王之言曰，以前事未終，故言“又曰”也。 \par}

{\noindent\shu\zihao{5}\fzkt “王曰又”至“攸居”。正義曰：王之所云，又複稱曰：“汝當是我，勿非我也。我乃有教誨之言，則汝所當居行之。” \par}

\section{無逸第十七}


\CJKunderline{周公}作\CJKunderwave{無逸}。\footnote{中人之性好逸豫，故戒以無逸。○好，呼報反。}

無逸\footnote{\CJKunderline{成王}即政，恐其逸豫，本以所戒名篇。}


{\noindent\zhuan\zihao{6}\fzbyks 傳“\CJKunderline{成王}”至“名篇”。正義曰：篇之次第，以先後為序，\CJKunderwave{多士}、\CJKunderwave{君奭}皆是\CJKunderline{成王}即位之初，知此篇是\CJKunderline{成王}始初即政，\CJKunderline{周公}恐其逸豫，故戒之，使無逸，即以所戒名篇也。 \par}

{\noindent\shu\zihao{5}\fzkt 傳“中人”至“無逸”。正義曰：上智不肯為非，下愚戒之無益，故中人之性，可上可下，不能勉強，多好逸豫，故\CJKunderline{周公}作書以戒之，使無逸。此雖指戒\CJKunderline{成王}，以為人之大法，\CJKunderline{成王}以聖賢輔之，當在中人以上，其實本性亦中人耳。 \par}

\CJKunderline{周公}曰：“嗚呼！君子所其無逸。\footnote{嘆美君子之道，所在唸德，其無逸豫。君子且猶然，況王者乎？}先知稼穡之艱難，乃逸，則知小人之依。\footnote{稼穡農夫之艱難，事先知之，乃謀逸豫，則知小人之所依怙。}相小人，厥父母勤勞稼穡，厥子乃不知稼穡之艱難，\footnote{視小人不孝者，其父母躬勤艱難，而子乃不知其勞。○相，息亮反。}乃逸乃諺。既誕，否則侮厥父母曰：‘昔之人無聞知。’”\footnote{小人之子既不知父母之勞，乃為逸豫遊戲，乃叛諺不恭。已欺誕父母，不欺，則輕侮其父母曰：“古老之人無所聞知。”○諺，魚戰反。}


{\noindent\zhuan\zihao{6}\fzbyks 傳“嘆美”至“者乎”。正義曰：\CJKunderline{周公}意重其事，故嘆而為言。鄭云：“嗚呼者，將戒\CJKunderline{成王}，欲求以深感動之。”是欲深感\CJKunderline{成王}，故“嘆美君子之道”。“君子”者，言其可以君正上位，子愛下民，有德則稱之，不限貴賤。君子之人，念德不怠，故“所在唸德，其無逸豫”也。“君子且猶然，而況王者乎”，言王者日有萬幾，彌復不可逸豫。鄭云：“君子止謂在官長者。所，猶處也。君子處位為政，其無自逸豫也。” \par}

{\noindent\zhuan\zihao{6}\fzbyks 傳“稼穡”至“依怙”。正義曰：民之性命，在於穀食，田作雖苦，不得不為。寒耕熱耘，霑體塗足，是稼穡為農夫艱難之事。在上位者，先知稼穡之艱難，乃可謀其逸豫，使家給人足，乃得思慮不勞，是為“謀逸豫”也。能知稼穡之艱難，則知小人之所依怙，言小人依怙此稼穡之事，不可不勤勞也。上句言君子當無逸，此言“乃謀逸豫”者，君子之事，勞心與形。盤於遊畋，形之逸也；無為而治，心之逸也。君子無形逸而有心逸，既知稼穡之艱難，可以謀心逸也。 \par}

{\noindent\zhuan\zihao{6}\fzbyks 傳“視小人”至“其勞”。正義曰：視小人不孝者，其父母勤苦艱難，勞於稼穡，成於生業，致富以遺之。而其子謂己自然得之，乃不知其父母勤勞。 \par}

{\noindent\zhuan\zihao{6}\fzbyks 傳“小人”至“聞知”。正義曰：上言視小人之身，此言“小人之子”者，“小人”謂無知之人，亦是賤者之稱，躬為稼穡，是賤者之事，故言“小人之子”，謂賤者之子，即上所視之小人也。此子既不知父母之勞，謂己自然得富，恃其家富，乃為逸豫遊戲，乃為叛諺不恭，已是欺誕父母矣。若不欺誕，則輕侮其父母曰：“古老之人無所聞知。”言其罪之深也。\CJKunderwave{論語}曰：“由也諺。”諺則叛諺，欺誕不恭之貌。“昔”訓久也，自今而道遠久，故為“古老之人”。\CJKunderwave{詩}云：“召彼故老。” \par}

{\noindent\shu\zihao{5}\fzkt “\CJKunderline{周公}”至“聞知”。正義曰：\CJKunderline{周公}嘆美君子之道以戒王曰：“嗚呼！君子之人，所在其無逸豫。君子必先知農人稼穡之艱難，然後乃謀為逸豫，如是則知小人之所依怙也。視彼小人不孝者，其父母勤勞稼穡，其子乃不知稼穡之艱難，乃為逸豫遊戲，乃叛諺不恭。既為欺誕父母矣，不欺，則又侮慢其父母曰：‘昔之人無所聞知。’小人與君子如此相反，王宜知其事也。” \par}

\CJKunderline{周公}曰：“嗚呼！我聞曰，昔在殷王中宗，\footnote{\CJKunderline{大戊}也，殷家中世尊其德，故稱宗。}嚴恭寅畏天命，自度，\footnote{言\CJKunderline{太戊}嚴恪恭敬，畏天命，用法度。○嚴如字，又魚檢反，注同，馬作儼。}治民祗懼，不敢荒寧。\footnote{為政敬身畏懼，不敢荒怠自安。○治，直吏反。}肆中宗之享國，七十有五年。\footnote{以敬畏之故，得壽考之福。}


{\noindent\zhuan\zihao{6}\fzbyks 傳“\CJKunderline{太戊}”至“稱宗”。正義曰：“中宗”，廟號。“\CJKunderline{太戊}”，王名。商自\CJKunderline{成湯}已後，政教漸衰，至此王而中興之。王者祖有功，宗有德，殷家中世尊其德，其廟不毀，故稱“中宗”。 \par}

{\noindent\zhuan\zihao{6}\fzbyks 傳“言太”至“法度”。正義曰：\CJKunderwave{祭義}雲“嚴威儼恪”，故引“恪”配“嚴”。\CJKunderline{鄭玄}云：“恭在貌，敬在心。”然則“嚴”是威，“恭”是貌，“敬”是心，三者各異，故累言之。 \par}

{\noindent\shu\zihao{5}\fzkt “\CJKunderline{周公}”至“五年”。正義曰：既言君子不逸，小人反之，更舉前代之王以夭壽為戒。\CJKunderline{周公}曰：“嗚呼！我所聞曰，昔在殷王中宗，威儀嚴恪，貌恭心敬，畏天命，用法度，治民敬身畏懼，不敢荒怠自安，故中宗之享有殷國七十有五年。”言不逸之故，而得歷年長也。 \par}

其在\CJKunderline{高宗}，時舊勞於外，爰暨小人。\footnote{\CJKunderline{武丁}，其父\CJKunderline{小乙}使之久居民間，勞是稼穡，與小人出入同事。}作其即位，乃或亮陰，三年不言。\footnote{\CJKunderline{武丁}起其即王位，則\CJKunderline{小乙}死，乃有信默，三年不言。言孝行著。○行，下孟反。}其惟不言，言乃雍，不敢荒寧。\footnote{在喪則其惟不言，喪畢發言，則天下和。亦法中宗，不敢荒怠自安。}嘉靖殷邦，至於小大，無時或怨。\footnote{善謀殷國，至於小大之政，人無是有怨者。言無非。}肆\CJKunderline{高宗}之享國，五十有九年。\footnote{高宗為政，小大無怨，故亦享國永年。}


{\noindent\zhuan\zihao{6}\fzbyks 傳“\CJKunderline{武丁}其”至“同事”。正義曰：“舊”，久也。在即位之前，而言久勞於外，知是其父\CJKunderline{小乙}使之久居民間，勞是稼穡，與小人出入同為農役，小人之艱難事也。太子使與小人同勞，此乃非常之事，不可以非常怪之。於時蓋未為太子也，殷道雖質，不可既為太子，更得與小人雜居也。 \par}

{\noindent\zhuan\zihao{6}\fzbyks 傳“\CJKunderline{武丁}起”至“行著”。正義曰：以上言久勞於外，為父在時事，故言“起其即王位，則\CJKunderline{小乙}死”也。“亮”，信也。“陰”，默也。三年不言，以舊無功，而今有，故言。乃有說此事者，言其孝行著也。\CJKunderwave{禮記·喪服四制}引\CJKunderwave{書}云：“‘\CJKunderline{高宗}諒暗，三年不言。’善之也。王者莫不行此禮，何以獨善之也？曰，\CJKunderline{高宗}者，\CJKunderline{武丁}。\CJKunderline{武丁}者，殷之賢王也。繼世即位，而慈良於喪。當此之時，殷衰而復興，禮廢而復起，故載之於\CJKunderwave{書}中而高之，故謂之\CJKunderline{高宗}。”三年之喪，君不言也，是說此經“不言”之意也。 \par}

{\noindent\zhuan\zihao{6}\fzbyks 傳“在喪”至“自安”。正義曰：\CJKunderline{鄭玄}云：“其不言之時，時有所言，則群臣皆和諧。”\CJKunderline{鄭玄}意謂此“言乃雍”者，在三年之內，時有所言也。孔意則為出言在三年之外，故云“在喪其惟不言，喪畢發言，則天下大和”。知者，\CJKunderwave{說命}云：“王宅憂，亮陰三祀。既免喪，其惟不言。”除喪猶尚不言，在喪必無言矣，故知喪畢乃發言也。\CJKunderline{高宗}不敢荒寧，與中宗正同，故云“亦法中宗，不敢荒怠自安”。殷家之王，皆是明王，所為善事，計應略同，但古文辭有差異，傳因其文同，故言“法中宗”也。 \par}

{\noindent\zhuan\zihao{6}\fzbyks 傳“善謀”至“無非”。正義曰：\CJKunderwave{釋詁}云：“嘉，善也。靖，謀也。”“善謀殷國”，謀為政教，故至於小大之政，皆允人意。人無是有怨\CJKunderline{高宗}者，言其政無非也。鄭云：“小大謂萬人，上及群臣言。”人臣小大皆無怨王也。 \par}

{\noindent\shu\zihao{5}\fzkt “其在”至“九年”。正義曰：其殷王\CJKunderline{高宗}，父在之時，久勞於外，於時與小人同其事。後為太子，起其即王之位，乃有信默，三年不言。在喪其惟不言，喪畢發言，言得其道，乃天下大和。不敢荒怠自安，善謀殷國，至於小大之政，莫不得所。其時之人，無是有怨恨之者。故\CJKunderline{高宗}之享殷國五十有九年。亦言不逸得長壽也。 \par}

其在祖甲，不義惟王，舊為小人。\footnote{湯孫\CJKunderline{太甲}，為王不義，久為小人之行，\CJKunderline{伊尹}放之桐。}作其即位，爰知小人之依，能保惠於庶民，不敢侮鰥寡。\footnote{在桐三年，思集用光，起就王位，於是知小人之所依。依仁政，故能安順於眾民，不敢侮慢惸獨。○惸，求營反，字又作𦬮。}肆祖甲之享國，三十有三年。\footnote{太甲亦以知小人之依，故得久年。此以德優劣、立年多少為先後，故祖甲在下。殷家亦祖其功，故稱祖。}


{\noindent\zhuan\zihao{6}\fzbyks 傳“湯孫”至“之桐”。正義曰：以文在“\CJKunderline{高宗}”之下，世次顛倒，故特辨之，此祖甲是湯孫\CJKunderline{太甲}也。“為王不義”，謂湯初崩。“久為小人之行，故\CJKunderline{伊尹}放之於桐”，言其廢而復興，為下“作其即位”起本也。王肅亦以祖甲為\CJKunderline{太甲}。\CJKunderline{鄭玄}云：“祖甲，\CJKunderline{武丁}子帝甲也。有兄祖庚賢，\CJKunderline{武丁}欲廢兄立弟，祖甲以此為不義，逃於人間，故云久為小人。”案\CJKunderwave{殷本紀}云：“\CJKunderline{武丁}崩，子祖庚立。祖庚崩，弟祖甲立，是為帝甲，淫亂，殷道復衰。”\CJKunderwave{國語}說殷事云：“帝甲亂之，七代而殞。”則帝甲是淫亂之主，起亡殷之源，寧當與二宗齊名，舉之以戒無逸？\CJKunderline{武丁}賢王，祖庚復賢，以\CJKunderline{武丁}之明，無容廢長立少。祖庚之賢，誰所傳說？\CJKunderline{武丁}廢子，事出何書？妄造此語，是負\CJKunderline{武丁}而誣祖甲也。 \par}

{\noindent\zhuan\zihao{6}\fzbyks 傳“在桐”至“惸獨”。正義曰：“在桐三年”，\CJKunderwave{太甲}序文。“思集用光”，\CJKunderwave{詩·大雅}文。彼“集”作“輯”，“輯”，和也。彼鄭言，\CJKunderline{公劉}之遷豳，“思在和其民人，用光大其道”。此傳之意，蓋言\CJKunderline{太甲}之在桐也，思得安集其身，用光顯王政，故起即王位，於是知小人之依。依於仁政，故能施行政教，安順於眾民，不敢侮慢。惸獨鰥寡之類，尤可憐愍，故特言之。 \par}

{\noindent\zhuan\zihao{6}\fzbyks 傳“\CJKunderline{太甲}”至“稱祖”。正義曰：傳於中宗雲“以敬畏之故，得壽考之福”，“\CJKunderline{高宗}之為政，小大無怨，故亦享國永年”，於此雲\CJKunderline{太甲}，亦以知小人之依，故得久年。各順其文而為之說，其言行善而得長壽，經意三王同也。以其世次顛倒，故解之云，此以德優劣、立年多少為先後，故祖甲在\CJKunderline{太戊}、\CJKunderline{武丁}之下。諸書皆言“\CJKunderline{太甲}”，此言“祖甲”者，殷家亦祖其功，故稱之“祖甲”。與二宗為類，惟見此篇，必言祖其功，亦未知其然。殷之先君有\CJKunderline{祖乙}、祖辛、祖丁，稱祖多矣，或可號之為祖，未必祖其功而存其廟也。 \par}

{\noindent\shu\zihao{5}\fzkt “其在”至“三年”。正義曰：其在殷王祖甲，初遭祖喪，所言行不義。惟亦為王，久為小人之行，\CJKunderline{伊尹}廢諸桐。起其即王之位，於是知小人之所依。依於仁政，乃能安順於眾民，不敢侮鰥寡惸獨，故祖甲之享有殷國三十有三年。亦言不逸得長壽也。 \par}

自時厥後立王，生則逸。\footnote{從是三王，各承其後而立者，生則逸豫無度。}生則逸，不知稼穡之艱難，\footnote{言與小人之子同其敝。}不聞小人之勞，惟耽樂之從。\footnote{過樂謂之耽。惟樂之從，言荒淫。○耽，丁南反，注下同。樂音洛，注下同。}自時厥後，亦罔或克壽。\footnote{以耽樂之故，從是其後，亦無有能壽考。}或十年，或七八年，或五六年，或四三年。”\footnote{高者十年，下者三年，言逸樂之損壽。}

{\noindent\shu\zihao{5}\fzkt “自時”至“三年”。正義曰：從是三王其後所立之王，生則逸豫，不知稼穡之艱難，不聞小人之勞苦，惟耽樂之事則從而為之。故從是其後諸王，無有能壽考者。或十年，或七八年，或五六年，或四三年。言逸樂之損壽，故舉以戒\CJKunderline{成王}也。 \par}

\CJKunderline{周公}曰：“嗚呼！厥亦惟我周太王、\CJKunderline{王季}，克自抑畏。\footnote{太王，\CJKunderline{周公}曾祖。\CJKunderline{王季}即祖。言皆能以義自抑，長敬天命。將說\CJKunderline{文王}，故本其父祖。}\CJKunderline{文王}卑服，即康功田功。\footnote{文王節儉，卑其衣服，以就其安人之功，以就田功，以知稼穡之艱難。○卑如字，馬本作俾，使也。}徽柔懿恭，懷保小民，惠鮮鰥寡。\footnote{以美道和民，故民懷之。以美政恭民，故民安之。又加惠鮮乏鰥寡之人。○鮮，息淺反，注同。}自朝至於日中昃，不遑暇食，用咸和萬民。\footnote{從朝至日昳不暇食，思慮政事，用皆和萬民。○昃音側，本亦作仄。昳，田節反。}


{\noindent\zhuan\zihao{6}\fzbyks 傳“\CJKunderline{大王}”至“父祖”。正義曰：“\CJKunderline{大王}，\CJKunderline{周公}曾祖。\CJKunderline{王季}即祖也”，此乃經傳明文，而須詳言之者，此二王之下辭無所結，陳此不為無逸，\CJKunderline{周公}將說\CJKunderline{文王}，故本其父祖，是以傳詳言也。解其言此之意。“以義自抑”者，言其非無此心，以義自抑而不為耳。 \par}

{\noindent\zhuan\zihao{6}\fzbyks 傳“\CJKunderline{文王}”至“艱難”。正義曰：\CJKunderline{文王}卑其衣服，以就安人之功，言儉於身而厚於人也。立君所以牧人，安人之功，諸有美政皆是也。就安人之內，田功最急，故特雲“田功”，以示知稼穡之艱難也。 \par}

{\noindent\zhuan\zihao{6}\fzbyks 傳“以美”至“之人”。正義曰：“徽”、“懿”皆訓為美，“徽柔懿恭”，此是施人之事，以匆厭恭懷安小民，故傳分而配之。“徽柔”配“懷”，“以美道和民，故民懷之”。“懿恭”配“保”，“以美政恭民，故民安之”，“徽懿”言其美而已，不知何所美也。人君施於民，惟有道與政耳，故傳以“美道”、“美政”言之，政與道亦互相通也。少乏鰥寡尢是可憐，故別言“加惠於鮮乏鰥寡之人”也。 \par}

{\noindent\zhuan\zihao{6}\fzbyks 傳“從朝”至“萬民”。正義曰：昭五年\CJKunderwave{左傳}云：“日上其中，食日為二，旦日為三。”則人之常食在日中之前，謂辰時也。\CJKunderwave{易·豐卦}彖曰：“日中則昃。”謂過中而斜昃也。“昃”亦名“昳”，言日蹉跌而下，謂未時也。故日之十位，食時為辰，日昳為未。言\CJKunderline{文王}勤於政事，從朝不食，或至於日中，或至於日昃，猶不暇食。故經“中”、“昃”並言之。傳舉晚時，故惟言“昳”。“遑”亦“暇”也，重言之者，古人自有復語，猶雲“艱難”也。所以不暇食者，為思慮政事，用皆和萬民。政事雖多，皆是為民，故言“咸”。“咸”訓皆也。 \par}

文王不敢盤於遊田，以庶邦惟正之供。\footnote{文王不敢樂於遊逸田獵，以眾國所取法則，當以正道供待之故。○供音恭。}\CJKunderline{文王}受命惟中身，厥享國五十年。”\footnote{文王九十七而終。中身，即位時年四十七。言中身，舉全數。}

{\noindent\zhuan\zihao{6}\fzbyks 傳“\CJKunderline{文王}”至“之故”。正義曰：\CJKunderwave{釋詁}云：“盤，樂也。”“遊”謂遊逸，“田”謂畋獵，二者不同，故並雲“遊逸田獵”。以眾國皆於\CJKunderline{文王}所取其法則，\CJKunderline{文王}當以正義供待之故也。言\CJKunderline{文王}思為政道以待眾國，故不敢樂於遊田。\CJKunderline{文王}世為西伯，故當為眾國所取法則。禮有田獵而不敢者，順時搜狩，不為取樂，故不敢非時畋獵以為樂耳。 \par}

{\noindent\zhuan\zihao{6}\fzbyks 傳“\CJKunderline{文王}”至“全數”。正義曰：“\CJKunderline{文王}年九十七而終”，\CJKunderwave{禮記·文王世子}文也。於九十七內減享國五十年，是未立之前有四十七。在\CJKunderwave{禮}諸侯逾年即位，此據代父之年，故為“即位時年四十七”也。計九十七年半折以為中身，則四十七時於身非中，言“中身”者，舉全數而稱之也。經言“受命”者，\CJKunderline{鄭玄}雲“受殷王嗣位之命”。然殷之末世，政教已衰，諸侯嗣位何必皆待王命？受先君之命亦可也。王肅云：“\CJKunderline{文王}受命，嗣位為君。”不言受王命也。 \par}

{\noindent\shu\zihao{5}\fzkt “\CJKunderline{周公}”至“十年”。正義曰：殷之三王既如此矣，\CJKunderline{周公}又言曰：“嗚呼！其惟我周家\CJKunderline{大王}、\CJKunderline{王季}，能以義自抑而畏敬天命，故王跡從此起也。\CJKunderline{文王}又卑薄衣服，以就其安人之功與治田之功。以美道柔和其民，以美政恭待其民，以此民歸之。以美政恭民之故，故小民安之，又加恩惠於鮮乏鰥寡之人。其行之也，自朝旦至於日中及昃，尚不遑暇食，用善政以諧和萬民故也。\CJKunderline{文王}專心於政，不敢逸樂於遊戲畋獵，以己為眾國所取法，惟當正身行己以供待之。由是\CJKunderline{文王}受命，嗣位為君，惟於中身受之，其享國五十年，亦以不逸得長壽也。” \par}

\CJKunderline{周公}曰：“嗚呼！繼自今嗣王，\footnote{繼從今已往嗣世之王，皆戒之。}則其無淫於觀、於逸、於遊、于田，以萬民惟正之供。\footnote{所以無敢過於觀遊逸豫田獵者，用萬民當惟正身以供待之故。}無皇曰：‘今日耽樂。’乃非民攸訓，非天攸若，時人丕則有愆。\footnote{無敢自暇曰：“惟今日樂，後日止。”夫耽樂者，乃非所以教民，非所以順天，是人則大有過矣。○愆，起虔反。夫音扶。}無若殷王受之迷亂，酗於酒德哉！”\footnote{以酒為兇謂之酗。言紂心迷政亂，以酗酒為德。戒嗣王無如之。○酗，況付反。}


{\noindent\zhuan\zihao{6}\fzbyks 傳“繼從”至“戒之”。正義曰：先言“繼”者，謂繼此後人，即從今以後嗣世之王也。\CJKunderline{周公}思及長遠后王，盡皆戒之，非獨\CJKunderline{成王}也。 \par}

{\noindent\zhuan\zihao{6}\fzbyks 傳“所以”至“之故”。正義曰：傳意訓“淫”為過，\CJKunderline{鄭玄}云：“淫，放恣也。”“淫”者侵淫不止，其言雖殊，皆是過之義也。言“觀”為非時而行，違禮觀物，如\CJKunderwave{春秋}隱公“如棠觀魚”，莊公“如齊觀社”。\CJKunderwave{穀梁傳}曰：“常事曰視，非常曰觀。”此言“無淫於觀”，禁其非常觀也。“逸”謂逸豫，“遊”謂遊蕩，“田”謂田獵，四者皆異，故每事言“於”。“以”訓用也，用萬民皆聽王命，王者惟當正身待之，故不得淫於觀逸遊田也。 \par}

{\noindent\zhuan\zihao{6}\fzbyks 傳“無敢”至“過矣”。正義曰：“無敢自暇”，謂事不寬不暇，而以為原王之意而為辭，故言曰：“耽以為樂，惟今日樂，而後日止。”惟言“今日樂”，明知“後日止”也。夫“耽樂”者，乃非所以教民，教民當恪勤也；非所以順天，順天當肅恭也。是此耽樂之人，則大有愆過矣。戒王不得如此也。 \par}

{\noindent\zhuan\zihao{6}\fzbyks 傳“以酒”至“如之”。正義曰：“酗”從酉，以兇為聲，是“酗”為兇酒之名，故“以酒為兇謂之酗”。“酗”是飲酒而益兇也。言紂心迷亂，以酗酒為德，飲酒為政，心以兇酒為己德，紂以此亡殷。戒嗣王無如之。 \par}

{\noindent\shu\zihao{5}\fzkt “\CJKunderline{周公}”至“德哉”。正義曰：\CJKunderline{周公}又言而嘆曰：“嗚呼！繼此後世自今以後嗣位之王，則其無得過於觀望，過於逸豫，過於遊戲，過於田獵。所以不得然者，以萬民聽王者之教命，王當正己身以供待萬民，必當早夜恪勤，無敢自閒暇。曰：‘今日且樂，後日乃止。’此為耽樂者，非民之所以教訓也，非天之所以敬順也。若是之人，則有大愆過矣。王當自勤政事，莫如殷王受之述亂國政，酗醟於酒德哉！殷紂藉酒為兇，以酒為德，由是喪亡殷國，王當以紂為戒，無得如之。” \par}

\CJKunderline{周公}曰：“嗚呼！我聞曰，古之人，猶胥訓告，胥保惠，胥教誨，\footnote{嘆古之君臣，雖君明臣良，猶相道告，相安順，相教誨以義方。}民無或胥譸張為幻。\footnote{譸張，誑也。君臣以道相正，故下民無有相欺誑幻惑也。○幻音患。誑，九況反。}此厥不聽，人乃訓之，乃變亂先王之正刑，至於小大。\footnote{此其不聽中正之君，人乃教之以非法，乃變亂先王之正法，至於小大，無不變亂。言己有以致之。}民否則厥心違怨，否則厥口詛祝。”\footnote{以君變亂正法，故民否則其心違怨，否則其口詛祝。言皆患其上。○詛，側助反。祝，之又反。}


{\noindent\zhuan\zihao{6}\fzbyks 傳“嘆古”至“義方”。正義曰：此章二事，善惡相反。下句“不聽人”者，是愚闇之君，知此言“古之人”者，是賢明之君。“相”是兩人相與，故知兼有“臣良”，更相教告。隱三年\CJKunderwave{左傳}石碏曰：“臣聞愛子，教之以義方。”故知相教誨者，使“相教誨以義方”也。則知相訓告者，告之以善道也；相保惠者，相安順以美政也。 \par}

{\noindent\zhuan\zihao{6}\fzbyks 傳“譸張”至“惑也”。正義曰：“譸張，誑也”，\CJKunderwave{釋訓}文。孫炎曰：“眩惑誑欺人也。”民之從上，若影之隨形，君臣以道相正，故下民無有相欺誑幻惑者。“幻”即眩也，惑亂之名，\CJKunderwave{漢書}稱西域有幻人是也。 \par}

{\noindent\zhuan\zihao{6}\fzbyks 傳“此其”至“致之”。正義曰：上言善事，此說惡事。如此其不聽者，是不聽中正之君也。既不聽中正，則好聽邪佞，知此“乃訓之”者，是邪佞之人訓之也。邪佞之人必反正道，故言“人乃教之以非法”。暗君即受用之，變亂先王之正法。“至於小大，無不變亂”，言皆變亂正法盡也。暗君所任同己，由已之暗，致此佞人，言此暗君已身有以致之也。上“君明臣良”，由君明而有良臣，亦是己有致。上之言“胥”，此不言者，君在佞臣，國亡滅矣，不待相教為惡，故不言“胥”也。 \par}

{\noindent\zhuan\zihao{6}\fzbyks 傳“以君”至“其土”。正義曰：君既變亂正法，必將困苦下民。民不堪命，忿恨必起，故民忿君乃有二事，否則心違怨，否則口詛祝，言皆患土而為此也。“違怨”，謂違其命而怨其身。“詛祝”，謂告神明令加殃咎也。以言告神謂之“祝”，請神加殃謂之“詛”。襄十七年\CJKunderwave{左傳}曰：“宋國區區，而有詛有祝。”\CJKunderwave{詩}曰：“侯詛侯祝。”是“詛”、“祝”意小異耳。 \par}

{\noindent\shu\zihao{5}\fzkt “\CJKunderline{周公}”至“詛祝”。正義曰：\CJKunderline{周公}言而嘆曰：“我聞人之言曰，古人之雖君明臣良，猶尚相訓告以善道，相安順以美政，相教誨以義方。君臣相正如此，故於時之民順從上教，無有相誑欺為幻惑者。此其不聽中正之君，人乃教訓之以非法之事，乃從其言，變亂先王之正法，至於小大之事，無不皆變亂之。君既變亂如此，其時之民疾苦，否則其心違上怨上，否則其口詛祝之。”言人患之無已，舉此以戒\CJKunderline{成王}，使之君臣相與養下民也。 \par}

\CJKunderline{周公}曰：“嗚呼！自殷王中宗，及\CJKunderline{高宗}，及祖甲，及我周\CJKunderline{文王}，茲四人迪哲。\footnote{言此四人皆蹈智明德以臨下。}厥或告之曰：‘小人怨汝詈汝。’則皇自敬德，\footnote{其有告之，言小人怨詈汝者，則大自敬德，增修善政。○詈，力智反。}厥愆，曰：‘朕之愆。’允若時不啻不敢含怒。\footnote{其人有禍，則曰：“我過，百姓有過，在予一人。”信如是怨詈，則四王不啻不敢含怒以罪之。言常和悅。}


{\noindent\zhuan\zihao{6}\fzbyks 傳“其有”至“善政”。正義曰：\CJKunderwave{釋詁}云：“皇，大也。”故傳言，“大自敬德者，謂增修善政”也。\CJKunderline{鄭玄}以“皇”為暇，言寬暇自敬。王肅本“皇”作“況”，況滋益用敬德也。 \par}

{\noindent\zhuan\zihao{6}\fzbyks 傳“其人”至“和悅”。正義曰：或告之曰“小人怨汝，詈汝”，其言有虛有實。其言若虛，則民之愆也。民有愆過，則曰“我過”，不責彼為虛言，而引過歸己者，湯所云“百姓有過，在予一人”。故若信有如是怨詈，小人聞之，則含怒以罪彼人。此四王即不啻不敢含怒以罪彼人，乃自原聞其愆言，其顏色常和悅也。\CJKunderline{鄭玄}云：“不但不敢含怒，乃欲屢聞之，以知己政得失之源也。” \par}

{\noindent\shu\zihao{5}\fzkt “\CJKunderline{周公}”至“含怒”。正義曰：既言明君暗君，善惡相反，更述二者之行。\CJKunderline{周公}言而嘆曰：“嗚呼！自殷王中宗，及\CJKunderline{高宗}，及祖甲，及我周\CJKunderline{文王}，此四人者，皆蹈明智之道以臨下民。其有告之曰：‘小人怨恨汝，罵詈汝。’既聞此言，則大自敬德，更增修善政。其民有過，則曰：‘是我之過。’民信有如是怨詈，則不啻不敢含怒以罪彼人，乃欲得數聞此言以自改悔。”言寬弘之若是。 \par}

此厥不聽，人乃或譸張為幻，曰：‘小人怨汝詈汝。’則信之。\footnote{此其不聽中正之君，有人誑惑之，言小人怨憾詛詈汝，則信受之。○憾，胡暗反。}則若時，不永念厥闢，不寬綽厥心，\footnote{則如是信讒者，不長念其為君之道，不寬緩其心。言含怒。}亂罰無罪，殺無辜，怨有同，是叢於厥身。”\footnote{信讒含怒，罰殺無罪，則天下同怨讎之，叢聚於其身。○叢，才公反。}


{\noindent\zhuan\zihao{6}\fzbyks 傳“則如”至“含怒”。正義曰：君人者察獄必審其虛實，然後加罪。“不長念其為君之道”，謂不審察虛實也。“不寬緩其心”，言徑即含怒也。王肅讀“闢”為闢，扶亦反，不長念其刑辟，不當加無罪也。 \par}

{\noindent\shu\zihao{5}\fzkt “此厥”至“厥身”。正義曰：此其不聽中正之人，乃有欺誑為幻惑以告之曰：“小人怨汝詈汝。”不原其本情，則信受之。則知是信讒者，不長念其為君之道，不審虛實，不能寬緩其心，而徑即含怒於人。是亂其正法，罰無罪，殺無辜。罰殺欲以止怨，乃令人怨益甚，天下之民有同怨君，令怨惡聚於其身。言褊急使民之怨若是，教\CJKunderline{成王}勿學此也。 \par}

\CJKunderline{周公}曰：“嗚呼！嗣王其監於茲。”\footnote{視此亂罰之禍以為戒。}

\section{君奭第十八}


\CJKunderline{召公}為保，\CJKunderline{周公}為師，相\CJKunderline{成王}為左右。\footnote{保，太保也。師，太師也。馬云：“保氏、師氏皆大夫官。”相音息亮反。左右，馬云：“分陝為二伯，東為左，西為右。”}\CJKunderline{召公}不說，\CJKunderline{周公}作\CJKunderwave{君奭}。

君奭\footnote{尊之曰君。奭，名，同姓也。陳古以告之，故以名篇。○說音悅。奭，始亦反。}


{\noindent\zhuan\zihao{6}\fzbyks 傳“尊之”至“名篇”。正義曰：\CJKunderline{周公}呼為“君奭”，是\CJKunderline{周公}尊之曰君也。“奭”是其名，“君”非名也。僖二十四年\CJKunderwave{左傳}，富辰言\CJKunderline{文王}之子一十六國，無名“奭”者，則\CJKunderline{召公}必非\CJKunderline{文王}之子。\CJKunderwave{燕世家}云：“\CJKunderline{召公奭}與周同姓姬氏。”譙周曰：“周之支族。”譙周考校古史，不能知其所出。皇甫謐云：“原公名豐，是其一也，是為\CJKunderline{文王}之子一十六國。”然\CJKunderline{文王}之子本無定數，並原、豐為一，當\CJKunderline{召公}於中以為十六，謬矣。此篇多言先世有大臣輔政，是“陳古道以告之”。呼居奭以告之，故以“君奭”名篇。 \par}

{\noindent\shu\zihao{5}\fzkt “\CJKunderline{召公}”至“君奭”。正義曰：\CJKunderline{成王}即政之初，\CJKunderline{召公}為保，\CJKunderline{周公}為師，輔相\CJKunderline{成王}為左右大臣。\CJKunderline{召公}以\CJKunderline{周公}嘗攝王之政，今覆在臣位，其意不說。\CJKunderline{周公}陳己意以告\CJKunderline{召公}，史敘其事，作\CJKunderwave{君奭}之篇也。\CJKunderwave{周官}篇雲“立太師、太傅、太保，茲惟三公”，則此“為保”、“為師”亦為三公官也。此實太師、太保而不言“太”者，意在師法保安王身，言其實為左右爾，不為舉其官名，故不言“太”也。經傳皆言\CJKunderline{武王}之時，太公為太師，此言“\CJKunderline{周公}為師”，蓋太公薨，命\CJKunderline{周公}代之。於時太傅蓋畢公為之，於此無事，不須見也。三公之次，先師後保，此序先言保者，篇之所作，主為\CJKunderline{召公}不說，故先言\CJKunderline{召公}，不以官位為次也。案經\CJKunderline{周公}之言，皆說己留在王朝之意，則\CJKunderline{召公}不說\CJKunderline{周公}之留也。故鄭、王皆云：“\CJKunderline{周公}既攝王政，不宜複列於臣職，故不說。”然則\CJKunderline{召公}大賢，豈不知\CJKunderline{周公}留意而不說者？以\CJKunderline{周公}留在臣職，當時人皆怪之，故欲開道\CJKunderline{周公}之言，以解世人之惑。“\CJKunderline{召公}疑之，作\CJKunderwave{君奭}。”非不知也。\CJKunderwave{史記·燕世家}云：“\CJKunderline{成王}既幻，\CJKunderline{周公}攝政，當國踐阼，\CJKunderline{召公}疑之，作\CJKunderwave{君奭}。”此篇是致政之後言留輔\CJKunderline{成王}之意，其文甚明，馬遷妄為說爾。\CJKunderline{鄭玄}不見\CJKunderwave{周官}之篇，言此師、保為\CJKunderwave{周禮}師氏、保氏大夫之職，言賢聖兼此官，亦謬矣。 \par}

\CJKunderline{周公}若曰：“君奭，\footnote{順古道呼其名而告之。}弗吊，天降喪於殷，殷既墜厥命，我有周既受。\footnote{言殷道不至，故天下喪亡於殷。殷已墜失其王命，我有周道至已受之。○吊音的。}我不敢知曰，厥基永孚於休，若天棐忱。\footnote{廢興之跡，亦君所知，言殷家其始長信於美道，順天輔誠，所以國也。○棐音匪。忱，市林反。}我亦不敢知曰，其終出於不祥。\footnote{言殷紂其終墜厥命，以出於不善之故，亦君所知。}


{\noindent\zhuan\zihao{6}\fzbyks 傳“廢興”至“以國”。正義曰：孔以\CJKunderwave{召誥}雲“我不敢知”者，其意\CJKunderline{召公}言我不敢獨知，亦王所知，則此言“我不敢知”，亦是\CJKunderline{周公}言我不敢獨知，是君奭所知，故以此及下句為說殷之興亡，言與君奭同知。舉其殷興亡為戒，\CJKunderline{鄭玄}亦然也。 \par}

{\noindent\shu\zihao{5}\fzkt “\CJKunderline{周公}”至“不祥”。正義曰：\CJKunderline{周公}留在王朝，\CJKunderline{召公}不說。\CJKunderline{周公}為師，順古道而呼曰：“君奭，殷道以不至之故，故天下喪亡於殷。殷既墜失其王命，我有周已受之矣。今雖受命，貴在能終，若不能終，與殷無異，故視殷以為監戒。我不敢獨知殷家其初始之時，能長信於美道，能安順於上天之，道輔其誠信，所以有國，此亦君之所知。我亦不敢獨知曰，殷紂其終墜失其王命，由出於不善之故，亦君所知也。” \par}

嗚呼！君已！曰，時我，我亦不敢寧於上帝命。\footnote{嘆而言曰：“君已！當是我之留，我亦不敢安於上天之命，故不敢不留。”○已音以。}弗永遠念天威，越我民罔尤違。\footnote{言君不長遠念天之威，而勤化於我民，使無過違之闕。}惟人在我後嗣子孫，大弗克恭上下，遏佚前人光，在家不知。\footnote{惟眾人共存在我後嗣子孫，若大不能恭承天地，絕失先王光大之道，我老在家，則不得知。○遏，於葛反。}天命不易，天難諶，乃其墜命，弗克經歷。\footnote{天命不易，天難信無德者，乃其墜失王命，不能經久歷遠，不可不慎。○易，以豉反，注同。諶，氏壬反。}嗣前人，恭明德，在今予小子旦。\footnote{繼先王之大業，恭奉其明德，正在今我小子旦。言異於餘臣。}非克有正，迪惟前人光，施於我衝子。”\footnote{我留非能有改正，但欲蹈行先王光大之道，施正於我童子。童子，\CJKunderline{成王}。}


{\noindent\zhuan\zihao{6}\fzbyks 傳“嘆而”至“不留”。正義曰：嘆而言曰：“嗚呼！君已！”“已”是引聲之辭，既呼君奭，嘆而引聲，乃復言曰：“君當是我之留。”以其意不說，故令是我而勿非我。“我不敢安於上天之命”，孔意當謂天既命周，我當成就周道，故不敢不留。 \par}

{\noindent\shu\zihao{5}\fzkt “嗚呼”至“衝子”。正義曰：\CJKunderline{周公}又嘆而呼\CJKunderline{召公}曰：“嗚呼！君已！”“已”，辭也。既嘆乃復言曰：“君當是我之留，勿非我也。我亦不敢安於上天之命，故不敢不留。君何不長遠念天之威罰？禍福難量，當勤教於我下民，使無尤過違法之闕。惟今天下眾人，共誠心存在我後嗣子孫。觀其政之善惡，若此嗣王大不能恭承上天下地，絕失先王光大之道，令使眾人失望，我若退老在家，則不能得知，何得不留輔王也？天命不易，言甚難也。天難信，惡則去之，不常在一家，是難信也。天子若不稱天意，乃墜失其王命，不能經久歷遠，其事可不慎乎？繼嗣前人先王之大業，恭奉其明德也，正在今我小子旦。”\CJKunderline{周公}自言已身當恭奉其先王之明德，留輔佐王。“非能有所改正，但欲蹈行先王光大之道，施政於我童子”。童子謂\CJKunderline{成王}，意欲奉行先王之事，以教\CJKunderline{成王}也。 \par}

又曰：“天不可信，我道惟寧王德延，\footnote{無德去之，是天不可信，故我以道惟安寧王之德，謀欲延久。}天不庸釋於\CJKunderline{文王}受命。”\footnote{言天不用令釋廢於\CJKunderline{文王}所受命，故我留佐\CJKunderline{成王}。}


{\noindent\zhuan\zihao{6}\fzbyks 傳“無德”至“延久”。正義曰：此經言“又曰”，傳不明解。鄭雲“人又云”，則\CJKunderline{鄭玄}以此“又曰”為\CJKunderline{周公}稱人之言也。王肅云：“重言天不可信，明己之留蓋畏其天命。”則肅意以\CJKunderline{周公}重言，故稱“又曰”。孔雖不解，當與王肅意同。言“寧王”者，即\CJKunderline{文王}也，鄭、王亦同。 \par}

{\noindent\shu\zihao{5}\fzkt “又曰”至“受命”。正義曰：\CJKunderline{周公}又言曰：“天不可信，無德則去之，是其不可信也。天難信之，故恐其去我周家，故我以道惟安行寧王之德，謀欲延長之。我原上天之意，不用令廢於\CJKunderline{文王}所受命，若嗣王失德，則還廢之，故我當留佐\CJKunderline{成王}也。” \par}

公曰：“君奭，我聞在昔\CJKunderline{成湯}既受命，\footnote{已放桀，受命為天子。}時則有若\CJKunderline{伊尹}，格於皇天。\footnote{尹摯佐湯，功至大天。謂致太平。○摯音至。}在\CJKunderline{太甲}，時則有若保衡。\footnote{太甲繼湯，時則有如此\CJKunderline{伊尹}為保衡，言天下所取安，所取平。}在\CJKunderline{太戊}，\footnote{太甲之孫。}時則有若\CJKunderline{伊陟}、臣扈，格於上帝，巫咸乂王家。\footnote{伊陟、臣扈率\CJKunderline{伊尹}之職，使其君不隕祖業，故至天之功不隕。巫咸治王家，言不及二臣。○隕，于敏反。}


{\noindent\zhuan\zihao{6}\fzbyks 傳“尹摯”至“太平”。正義曰：\CJKunderline{伊尹}名摯，諸子傳記名有其文。“功至大天”猶堯“格於上下”，知其“謂致太平”也。 \par}

{\noindent\zhuan\zihao{6}\fzbyks 傳“\CJKunderline{太甲}”至“取平”。正義曰：據\CJKunderwave{太甲}之篇及諸子傳記，\CJKunderline{太甲}犬臣惟有\CJKunderline{伊尹}，知即保衡也。\CJKunderwave{說命}云：“昔先正保衡，作我先王，佑我烈祖，格於皇天。”\CJKunderwave{商頌·那}祀\CJKunderline{成湯}稱為“烈祖”，“烈祖”，湯之號，言保衡佐湯，明保衡即是\CJKunderline{伊尹}也。\CJKunderwave{詩}稱“實維阿衡，實左右商王”，\CJKunderline{鄭玄}云：“阿，倚。衡，平也。\CJKunderline{伊尹}湯所依倚而取平。”至\CJKunderline{太甲}改曰保衡，保，安也，言天下所取安，所取平。此皆三公之官，當時為之號也。孔以\CJKunderwave{太甲}雲“嗣王不惠於阿衡”，則\CJKunderwave{太甲}亦曰阿衡，與鄭異也。 \par}

{\noindent\zhuan\zihao{6}\fzbyks 傳“\CJKunderline{太甲}之孫”。正義曰：\CJKunderwave{史記·殷本紀}云，\CJKunderline{太甲}崩，子\CJKunderline{沃丁}立。崩，弟太庚立。崩，子小甲立。崩，弟雍已立。崩，弟\CJKunderline{太戊}立。是\CJKunderline{太戊}為\CJKunderline{太甲}之孫，太庚之子。\CJKunderwave{三代表}云，小甲，太庚弟；雍己、\CJKunderline{太戊}又是小甲弟，則\CJKunderline{太戊}亦是\CJKunderline{沃丁}弟，\CJKunderline{太甲}子。\CJKunderwave{本紀}、\CJKunderwave{世表}俱出馬遷，必有一誤。孔於\CJKunderwave{咸乂}序傳雲“\CJKunderline{太戊}，\CJKunderline{沃丁}弟之子”，是\CJKunderline{太戊}為\CJKunderline{太甲}之孫也。 \par}

{\noindent\zhuan\zihao{6}\fzbyks 傳“\CJKunderline{伊陟}”至“二臣”。正義曰：\CJKunderline{伊尹}“格於皇天”，此\CJKunderline{伊陟}、臣扈雲“格於上帝”，其事既同，如此二臣能率循\CJKunderline{伊尹}之職，輔佐其君，使其君不隕祖業，故至天之功亦不隕墜也。\CJKunderwave{夏社}序云：“湯既勝夏，欲遷其社，不可。作\CJKunderwave{夏社}、\CJKunderwave{疑至}、\CJKunderwave{臣扈}。”則湯初有臣扈，已為大臣矣，不得至今仍在，與\CJKunderline{伊尹}之子同時立功。蓋二人名同，或兩字一誤也。案\CJKunderwave{春秋}範武子光輔五君，或臣扈事湯而又事\CJKunderline{太戊}也。“格於上帝”之下乃言“巫咸乂王家”，則巫咸亦是賢臣，俱能紹治王家之事而已，其功不得至天，言不及彼二臣。 \par}

在\CJKunderline{祖乙}，時則有若巫賢。\footnote{祖乙，殷家亦祖其功，時賢臣有如此巫賢。賢，咸子。巫，氏。}在\CJKunderline{武丁}，時則有若\CJKunderline{甘盤}。\footnote{高宗即位，\CJKunderline{甘盤}佐之，後有\CJKunderline{傅說}。○說音悅。}

{\noindent\zhuan\zihao{6}\fzbyks 傳“\CJKunderline{祖乙}”至“巫氏”。正義曰：\CJKunderwave{殷本紀}云，中宗崩，子\CJKunderline{仲丁}立。崩，弟外壬立。崩，弟\CJKunderline{河亶甲}立。崩，子\CJKunderline{祖乙}立。則\CJKunderline{祖乙}是\CJKunderline{太戊}之孫也。孔以其人稱“祖”，故云“殷家亦祖其功”。賢是咸子，相傳云然。父子俱稱為“巫”，知“巫”為氏也。 \par}

{\noindent\zhuan\zihao{6}\fzbyks 傳“\CJKunderline{高宗}”至“\CJKunderline{傅說}”。正義曰：\CJKunderwave{孔命}篇\CJKunderline{高宗}云：“臺小子舊學於\CJKunderline{甘盤}，既乃遯於荒野。”\CJKunderline{高宗}未立之前已有\CJKunderline{甘盤}，免喪不言，乃求\CJKunderline{傅說}，明其即位之初，有\CJKunderline{甘盤}佐之，\CJKunderline{甘盤}卒後有\CJKunderline{傅說}。計\CJKunderline{傅說}當有大功，此惟數六人，不言\CJKunderline{傅說}者，\CJKunderline{周公}意所不言，未知其故。 \par}

{\noindent\shu\zihao{5}\fzkt “公曰君奭”至“\CJKunderline{甘盤}”。正義曰：言“時則有若”者，言當其時有如此人也。指謂如此\CJKunderline{伊尹}、\CJKunderline{甘盤}，非謂別有如此人也。以湯是殷之始王，故言“在昔”、“既受命”，見其為天子也。以下“在\CJKunderline{太甲}”、“在\CJKunderline{武丁}”，亦言其為天子之時，有如此臣也。\CJKunderline{成湯}未為天子，已得\CJKunderline{伊尹}，言“既受命”者，以功格皇天，在受命之後，故言“既受命”也。“皇天”之與“上帝”，俱是天也，變其文爾。其功至於天帝，謂致太平而天下和之也。保衡、\CJKunderline{伊尹}，一人也。異時而別號。“\CJKunderline{伊尹}”之下，已言“格於皇天”，“保衡”之下不言“格於皇天”，從可知也。“\CJKunderline{伊陟}、臣扈”，言“格於上帝”，則其時亦致太平，故與\CJKunderline{伊尹}文異而事同。巫咸、巫賢、\CJKunderline{甘盤}蓋功劣於彼三人，故無格天之言。 \par}

率惟茲有陳，保乂有殷，故殷禮陟配天，多歷年所。\footnote{言\CJKunderline{伊尹}至\CJKunderline{甘盤}六臣佐其君，循惟此道，有陳列之功，以安治有殷，故殷禮能升配天，享國久長，多歷年所。}天惟純佑命，則商實百姓。\footnote{殷禮配天，惟天大佑助其王命，使商家百姓豐實，皆知禮節。}


{\noindent\zhuan\zihao{6}\fzbyks 傳“言伊”至“年所”。正義曰：“率”訓循也。說賢臣佐君雲“循惟此道”，當謂循此為臣之道。盡忠竭力以輔其君，故有陳烈於世，以安治有殷，使殷王得安治民。故殷得此安上治民之禮，能升配上天。天在人上，故謂之“升”。為天之子，是“配天”也。享國久長，多歷年所。 \par}

{\noindent\zhuan\zihao{6}\fzbyks 傳“殷禮”至“禮節”。正義曰：殷能以禮配天，故天降福。天惟大佑助其王命，風雨以時，年穀豐稔，使商家百姓豐實，家給人足。管子曰：“衣食足，知榮辱。倉廩實，知禮節。” \par}

{\noindent\shu\zihao{5}\fzkt “率惟”至“百姓”。正義曰：此\CJKunderline{伊尹}、\CJKunderline{甘盤}六臣等輔佐其君，率循此為臣之道，有陳列之功，以安治有殷，故殷有安上治民之禮，升配上天，享國多歷年之次所。天惟大佑助其為王之命，則使商家富實百姓，為令使商之百姓家給人足，皆知禮節也。 \par}

王人罔不秉德，明恤小臣，屏侯甸。\footnote{自湯至\CJKunderline{武丁}，其王人無不持德立業，明憂其小臣，使得其人，以為蕃屏侯甸之服。小臣且憂得人，則大臣可知。○屏，賓領反。}矧咸奔走，惟茲惟德稱，用又厥闢。\footnote{王猶秉德憂臣，況臣下得不皆奔走？惟王此事，惟有德者舉，用治其君事。○闢，必亦反。}故一人有事於四方，若卜筮，罔不是孚。”\footnote{一人，天子也。君臣務德，故有事於四方，而天下化服。如卜筮，無不是而信之。}


{\noindent\zhuan\zihao{6}\fzbyks 傳“自湯”至“可知”。正義曰：王肅云：“王人猶君人也。”“無不持德立業”，謂持人君之德，立王者之事業。人君之德在官賢人，官得其人，則事業立，故傳以“立業”配“持德”。明憂小臣之不賢，憂欲使得其人，以為蕃屏侯甸之服也。小臣且憂得人，則大臣憂之可知。侯甸尚思得其人，朝廷思之必矣。王肅云：“小臣，巨之微者，舉小以明大也。” \par}

{\noindent\zhuan\zihao{6}\fzbyks 傳“王猶”至“君事”。正義曰：君之所重，莫重於求賢。官之所急，莫急於得人。故此章所陳，惟言君憂得人，臣能舉賢。以王之尊，猶尚秉德憂臣，況其臣下得不皆奔走？惟王此求賢之事，惟有德者必舉之，置於官位用治其君事也。 \par}

{\noindent\zhuan\zihao{6}\fzbyks 傳“一人”至“信之”。正義曰：\CJKunderwave{禮}天子自稱曰予一人，故為天子也。君臣務求有德，眾官得其人，從上至下，遞相師法，職無大小，莫不治理，故天子有事於四方，發號出令而天下化服。譬如卜筮，無不是而信之。事既有驗，言如是則人皆信之。 \par}

{\noindent\shu\zihao{5}\fzkt “王人”至“是孚”。正義曰：“王人”謂與人為王，言此上所說\CJKunderline{成湯}、\CJKunderline{太甲}、\CJKunderline{太戊}、\CJKunderline{祖乙}、\CJKunderline{武丁}，皆王人也。無不持德立業，明憂小臣。雖則小臣，亦憂使得其賢人，以蕃屏侯甸之服。王恐臣之不賢，尚以為憂，況在臣下得不皆勤勞奔走，惟憂王此求賢之事，惟求有德者舉之，用治其君之事乎？君臣共求其有德，所在職事皆治，天子一人有事於四方，天下咸化而服。如有卜筮之驗，無不是而信之。賢臣助君，致使大治，我留不去，亦當如此也。 \par}

公曰：“君奭，天壽平格，保乂有殷，有殷嗣，天滅威。\footnote{言天壽有平至之君，故安治有殷。有殷嗣子紂，不能平至，天滅亡，加之有威。}今汝永念，則有固命，厥亂明我新造邦。”\footnote{今汝長念平至者安治，反是者滅亡。以為法戒，則有堅固王命，其治理足以明我新成國矣”。}


{\noindent\zhuan\zihao{6}\fzbyks 傳“言天”至“以威”。正義曰：“格”訓至也。“平”謂政教均平，“至”謂道有所至也。言“不弔”，謂道有不至者。此言“格”，謂道至者。“天壽有平至之君”，有平至之德，則天與之長壽，則知中宗\CJKunderline{高宗}之屬身是也。由其君有平至之德，故能安治有殷，言有殷國安而民治也。有殷嗣子紂，其德不能平至，國不安，民不治，故天滅亡之而加之以威也。孔傳之意，此經專說君之善惡，其言不及臣也。王肅以為兼言君臣，注云：“殷君臣之有德，故安治有殷。言是者，不可不法殷家有良臣也。”鄭注以為專言臣事，“格”謂至於天也。與孔不同。 \par}

{\noindent\zhuan\zihao{6}\fzbyks 傳“今汝”至“國矣”。正義曰：上句言善者興而惡者亡，此句令其長安治及念明道。念上二者，故言“今汝長念平至者而安治，反是者滅亡”。念此以為法戒，則有堅固王命，王族必不傾壞。若能如此，其治理足以光明我新成國矣。周自\CJKunderline{武王}伐紂，至此年歲末多，對殷而言故為新國。傳意言不及臣，\CJKunderline{周公}說此事者，蓋言興滅由人，我欲輔王，使為平至之君。 \par}

{\noindent\shu\zihao{5}\fzkt “公曰君奭天”至“造邦”。正義曰：\CJKunderline{周公}呼\CJKunderline{召公}曰：“君奭，皇天賦命，壽此有平至之君。”言有德者必壽考也。“殷之先王有平至之德，故能安治有殷”。言故得安治也。“有殷嗣子紂不能平至，故天滅亡而加之以威。今汝奭當長念天道，平至者安治，不平至者滅亡。以此為法戒，則有堅固王命，其治理足以明我新成國矣”。 \par}

公曰：“君奭，在昔上帝，割申勸寧王之德，其集大命於厥躬。\footnote{在昔上天，割制其義，重勸\CJKunderline{文王}之德，故能成其大命於其身。謂勤德以受命。○重，直用反。}惟\CJKunderline{文王}尚克修和我有夏，亦惟有若虢叔，有若閎夭，\footnote{文王庶幾能修政化，以和我所有諸夏，亦惟賢臣之助為治，有如此虢、閎。閎，氏。虢，國；叔，字；\CJKunderline{文王}弟。夭，名。○虢，寡白反，徐公伯反。閎音宏。夭，於表反，徐於驕反。}有若散宜生，有若泰顛，有若南宮括。\footnote{散、泰、南宮皆氏。宜生、顛、括皆名。凡五臣佐\CJKunderline{文王}為胥附、奔走、先後、禦侮之任。}


{\noindent\zhuan\zihao{6}\fzbyks 傳“在昔”至“受命”。正義曰：\CJKunderline{文王}去此未久，但欲遠本天意，故云“在昔上天”，作久遠言之。“割制”謂切割絕斷之意，故云“割制其義”。“重勸\CJKunderline{文王}之德”者，\CJKunderline{文王}既已有德，上天佑助而重勸勉，\CJKunderline{文王}順天之意，故其能成大命於其身。王謂勤行德義，以受天命。 \par}

{\noindent\zhuan\zihao{6}\fzbyks 傳“\CJKunderline{文王}”至“夭名”。正義曰：\CJKunderline{文王}未定天下，庶幾能修政化，以和我所有諸夏，謂三分有二屬己之諸國也。僖五年\CJKunderwave{左傳}雲“虢仲、虢叔，\CJKunderline{王季}之穆也”，是虢叔為\CJKunderline{文王}之弟。虢，國名。叔，字。凡言人之名氏，皆上氏下名，故閎、散、泰、南宮皆氏，夭、宜生、顛、括皆名也。 \par}

{\noindent\zhuan\zihao{6}\fzbyks 傳“散泰”至“之任”。正義曰：\CJKunderwave{詩·綿}之卒章稱\CJKunderline{文王}有疏附、先後、奔奏、禦侮之臣，\CJKunderwave{毛傳}云：“率下親上曰疏附，相通前後曰先後，喻德宣譽曰奔奏，武臣折衝曰禦侮。”鄭箋云：“疏附使疏者親也，奔奏使人歸趨之。”\CJKunderwave{詩}言\CJKunderline{文王}有此四種之臣，經歷言五臣之名，故知五臣佐\CJKunderline{文王}為此任也。此四事者五臣共為此任，非一臣當一事也。鄭云：“不及呂望者，太師致\CJKunderline{文王}以大德，\CJKunderline{周公}謙不可以自比。” \par}

{\noindent\shu\zihao{5}\fzkt “公曰君奭”至“厥躬”。正義曰：公呼\CJKunderline{召公}曰：“君奭，在昔上天斷割其義，重勸\CJKunderline{文王}之德。以\CJKunderline{文王}有德，勸勉使之成功，故\CJKunderline{文王}能成之命於其身。”言\CJKunderline{文王}能順天之意，勤德以受命。 \par}

又曰，無能往來，茲迪彝教\CJKunderline{文王}蔑德，降於國人。\footnote{有五賢臣，猶曰其少，無所能往來。而五人以此道法教\CJKunderline{文王}以精微之德，下政令於國人。言雖聖人，亦須良佐。}亦惟純佑，秉德迪知天威，乃惟時昭\CJKunderline{文王}。\footnote{文王亦如殷家惟天所大佑，\CJKunderline{文王}亦秉德蹈知天威，乃惟是五人明\CJKunderline{文王}之德。}迪見冒聞於上帝，惟時受有殷命哉！\footnote{言能明\CJKunderline{文王}德，蹈行顯見，覆冒下民，彰聞上天，惟是故受有殷之王命。○見，賢遍反，注同。冒，莫報反，下同，馬作勖，勉也。聞音問，或如字。}


{\noindent\zhuan\zihao{6}\fzbyks 傳“有五”至“良佑”。正義曰：“無能往來”一句，\CJKunderline{周公}假為\CJKunderline{文王}之辭。言\CJKunderline{文王}有五賢臣，猶恨其少。又復言曰：“我臣既少，於事無能往來。”謂去還理事，未能周悉，言其好賢之深，不知厭足也。“迪”，道。“彝”，法也。“蔑”，小也，小謂精微也。而五人以此道法教\CJKunderline{文王}以精微之德，用此精微之德下教令於國人。言雖聖人，亦須良佐，以見\CJKunderline{成王}須輔佐之甚也。\CJKunderline{鄭玄}亦云：“蔑，小也。” \par}

{\noindent\shu\zihao{5}\fzkt “又曰”至“命哉”。正義曰：\CJKunderline{文王}既有賢臣五人，又復言曰：“我之賢臣猶少，無所能往來。五人以此道法教\CJKunderline{文王}以微蔑精妙之德，下政令於國人。德政既善，為天所佑。\CJKunderline{文王}亦如殷家，惟為天所大佑。\CJKunderline{文王}亦秉德，蹈知天威。\CJKunderline{文王}德如此者，乃惟是五人明\CJKunderline{文王}之德使然也。五人能明\CJKunderline{文王}德，使蹈行顯見，覆冒下民，聞於上天，惟是之故得受有殷王之命哉！”言\CJKunderline{文王}之聖，猶須良佐，我所以留輔\CJKunderline{成王}。 \par}

\CJKunderline{武王}惟茲四人，尚迪有祿。\footnote{文王沒，\CJKunderline{武王}立，惟此四人，庶幾輔相\CJKunderline{武王}蹈有天祿。虢叔先死，故曰四人。○相，息亮反。}後暨\CJKunderline{武王}，誕將天威，咸劉厥敵。\footnote{言此四人後與\CJKunderline{武王}皆殺其敵。謂誅紂。}惟茲四人，昭\CJKunderline{武王}，惟冒丕單稱德。\footnote{惟此四人，明\CJKunderline{武王}之德，使布冒天下，大盡舉行其德。}


{\noindent\zhuan\zihao{6}\fzbyks 傳“\CJKunderline{文王}”至“四人”。正義曰：\CJKunderline{文王}受命九年而崩，十三年方得殺紂。“\CJKunderline{文王}沒，\CJKunderline{武王}立”，謂\CJKunderline{武王}初立之時，惟此四人而已。“庶幾輔相\CJKunderline{武王}蹈有天祿”，初立則有此志，故下句言後與\CJKunderline{武王}殺紂也。“虢叔先死，故曰四人”，以是\CJKunderline{文王}之弟，其年應長，故言“先死”也。\CJKunderline{鄭玄}疑不知誰死，注云：“至\CJKunderline{武王}時，虢叔等有死者，餘四人也。” \par}

{\noindent\zhuan\zihao{6}\fzbyks 傳“惟此”至“其德”。正義曰：“單”，盡。“稱”，舉也。使\CJKunderline{武王}之德布冒天下，是此四人之力，言此四人大盡舉行\CJKunderline{武王}之德也。 \par}

{\noindent\shu\zihao{5}\fzkt “\CJKunderline{武王}”至“稱德”。正義曰：\CJKunderline{文王}既沒，\CJKunderline{武王}次立，武功初立，惟此四人，庶幾輔相\CJKunderline{武王}蹈有天下之祿。其後四人，與\CJKunderline{武王}大行天之威罰，皆與共殺其強敵，謂其誅紂也。\CJKunderline{武王}之有天下，惟此四人明\CJKunderline{武王}之德，惟\CJKunderline{武王}佈德，覆冒天下，此四人大盡舉行\CJKunderline{武王}之德。言\CJKunderline{武王}亦得良臣之力。 \par}

今在予小子旦，若遊大川，予往暨汝奭其濟小子，同未在位，誕無我責。\footnote{我新還政，今任重在我小子旦，不能同於四方。若遊大川，我往與汝奭其共濟渡\CJKunderline{成王}，同於未在位即政時，汝大無非責我留。}收罔勖不及，耇造德不降，我則鳴鳥不聞，矧曰其有能格？”\footnote{今與汝留輔\CJKunderline{成王}，欲收教無自勉不及道義者，立此化，而老成德不降意為之。我周則鳴鳳不得聞，況曰其有能格於皇天乎？}


{\noindent\zhuan\zihao{6}\fzbyks 傳“我新”至“我留”。正義曰：\CJKunderline{周公}既已還政，則是舍重任矣。而猶言“今任重在我小子旦”者，\CJKunderline{周公}既攝王政，又須傳授得人，若其不能負荷，仍是\CJKunderline{周公}之責，以嗣子劣弱，故言“今任重猶在我小子旦”也。彼四人者能翼贊初基，佑\CJKunderline{成王}業，我不能同於四人，望有大功，惟求救弱而已。\CJKunderwave{詩}雲“泳之遊之”，\CJKunderwave{左傳}稱“閻敖遊湧而逸”，則“遊”者入水浮渡之名。譬若\CJKunderline{成王}在大川，我往與汝奭其同共濟渡\CJKunderline{成王}。若雲從此向川，故言“往”也。 \par}

{\noindent\zhuan\zihao{6}\fzbyks 傳“今與”至“天乎”。正義曰：王朝之臣有不勉力者，今與汝留輔\CJKunderline{成王}者，正欲收斂教誨。無自勉力不及道義者，當教之勉力，使其及道義也。我欲成立此化，而老成德之人不肯降意為之。我周家則鳴鳳尚不聞知，況曰其有能如\CJKunderline{伊尹}之輩，使其功格於皇天乎？言太平不可冀也。經言“耇造德不降”者，\CJKunderline{周公}以己年老應退而留，因即傳言己類，言己若退，則老成德者悉皆退自逸樂，不肯降意為之。政無所成，祥瑞不至，我周家則鳴鳳不得聞。則鳳是難聞之鳥，必為靈瑞之物，故以“鳴鳥”為鳴鳳。\CJKunderline{孔子}稱“鳳鳥不至”，是鳳鳥難聞也。\CJKunderwave{詩·大雅·卷阿}之篇歌\CJKunderline{成王}之德，其九章曰：“鳳皇鳥矣，於彼高岡。”鄭云：“因時鳳皇至，固以喻焉。”則\CJKunderline{成王}之時鳳皇至也。\CJKunderwave{大雅}正經之作，多在\CJKunderline{周公}攝政之後，\CJKunderline{成王}即位之初，則\CJKunderline{周公}言此之時已鳳皇至，見太平矣。而復言此者，恐其不復能然，故戒之。此經之意，言功格上天，難於致鳳，故以鳴鳳況之格天。案\CJKunderwave{禮器}云：“升中於天，而鳳皇降，龜龍假。”“升中”謂功成告天也。如彼\CJKunderwave{記}文，似功至於天，鳳皇乃降，此以鳴鳳易致況格天之難者乎。\CJKunderwave{記}以龍鳳有形，是可見之物，故以鳳降龍至為成功之驗，非言成功告天，然後此物始至也。 \par}

{\noindent\shu\zihao{5}\fzkt “今在”至“能格”。正義曰：\CJKunderline{周公}言：“我新還政\CJKunderline{成王}，今任之重者，其在我小子之身也。我不能同於四人輔文武，使有大功德，但苟求救溺而已。譬如遊於大川，我往與汝奭其共濟渡小子\CJKunderline{成王}，用心輔弼，同於\CJKunderline{成王}未在位之時。恐其未能嗣先人明德，我當與汝輔之，汝大無非責我之留也。我留與汝輔王者，欲收教無自勉力不及道義者。我今欲立此化，而老成德之人不降意為之。我周家則鳴鳳之鳥尚不得聞知，況曰其有能格於皇天者乎？” \par}

公曰：“嗚呼！君，肆其監於茲。我受命無疆惟休，亦大惟艱。\footnote{以朝臣無能立功至天，故其當視於此，我周受命無窮惟美，亦大惟艱難，不可輕忽，謂之易治。}告君乃猷裕，我不以後人迷。”\footnote{告君汝謀寬饒之道，我留與汝輔王，不用後人迷惑，故欲教之。}


{\noindent\zhuan\zihao{6}\fzbyks 傳“告君”至“教之”。正義曰：“猷”訓為謀，告君汝謀寬饒之道，故當以寬饒為法。我留與汝輔王，不用使後人迷惑怪之。無法則迷惑，故欲與汝作法以教之。鄭云：“\CJKunderline{召公}不說似隘急，故令謀於寬裕也。” \par}

{\noindent\shu\zihao{5}\fzkt “公曰嗚呼”至“人迷”。正義曰：\CJKunderline{周公}嘆而呼\CJKunderline{召公}曰：“嗚呼！君，我以朝臣無能立功至天之故，故君其當視於此。”謂視此朝臣無能立功之事。“我周家受天之命，無有境界惟美，亦大惟艱難，不可輕忽，謂之易治。我今告君，汝當謀寬饒之道以治下民，使其事可法，我不用使後世人迷惑，故欲教之也”。 \par}

公曰：“前人敷乃心，乃悉命汝，作汝民極。\footnote{前人文武布其乃心為法度，乃悉以命汝矣，為汝民立中正矣。}曰，汝明勖偶王，在亶乘茲大命。\footnote{汝以前人法度明勉配王，在於成信，行此大命而已。}惟\CJKunderline{文王}德，丕承無疆之恤。”\footnote{惟\CJKunderline{文王}聖德，為之子孫無忝厥祖，大承無窮之憂。}


{\noindent\zhuan\zihao{6}\fzbyks 傳“前人”至“正矣”。正義曰：“乃”，緩辭，不訓為汝。 \par}

{\noindent\zhuan\zihao{6}\fzbyks 傳“汝以”至“而已”。正義曰：“勖”，勉也。“偶”，配也。“亶”，信也。汝當以前人法度明自勉力，配\CJKunderline{成王}，在於誠信行大命而已。言其不復須勞心。傳以“乘”為行，蓋以乘車必行，故訓“乘”為行。 \par}

{\noindent\shu\zihao{5}\fzkt “公曰前”至“之恤”。正義曰：\CJKunderline{周公}又言曰，前人文武布其乃心製法度，乃悉命汝，為民立中正之道矣。治民之法已成就也，戒\CJKunderline{召公}汝當以前人之法度明自勉力，配此\CJKunderline{成王}，在於誠信，行此大命而已。言已有舊法，易可遵行也。惟\CJKunderline{文王}聖德造始周邦，為其子孫欲令無忝厥祖，大承無窮之憂，故我與汝不可不輔。 \par}

公曰：“君，告汝朕允。\footnote{告汝以我之誠信也。}保奭，其汝克敬以予，監於殷喪大否。\footnote{呼其官而名之，敕使能敬以我言，視於殷喪亡大否。言其大不可不戒。}肆念我天威，予不允惟若茲誥，予惟曰：‘襄我二人。’\footnote{以殷喪大故，當念我天德可畏。言命無常，我不信惟若此誥。我惟曰：“當因我文武之道而行之。”}汝有合哉！言曰：‘在時二人，天休滋至，惟時二人弗戡。’\footnote{言汝行事，動當有所合哉！發言常在是文武，則天美周家，日益至矣，惟是文武不勝受。言多福。}其汝克敬德，明我俊民在讓，後人於丕時。\footnote{其汝能敬行德，明我賢人在禮讓，則後代將於此道大且是。}


{\noindent\zhuan\zihao{6}\fzbyks 傳“言汝”至“多福”。正義曰：“動當有所合哉”，舉動皆合文武也。“發言常在是文武”，言非文武道則不言。 \par}

{\noindent\shu\zihao{5}\fzkt “公曰君告”至“丕時”。正義曰：\CJKunderline{周公}呼\CJKunderline{召公}曰：“君，我今告汝以我之誠信。”又呼其官而名之：“太保奭，其汝必須能敬以我之言，視於殷之喪亡。殷之喪亡，其事甚大，不可不戒慎。以殷喪大之故，當念我天德可畏。”言天命無常，無德則去之，甚可畏。“我不信惟若此誥而已。我惟言曰：‘當因我文武二人之道而行之。’汝所行事，舉動必當有所合哉！當與\CJKunderline{文王}\CJKunderline{武王}合也。汝所發言，常在是\CJKunderline{文王}\CJKunderline{武王}二人，則天美我周家，日日滋益至矣。其善既多，惟在是文武二人，不能勝受之矣。其汝能敬行德，明我賢俊之人在於禮讓，則後人於此道大且是也。” \par}

嗚呼！篤棐時二人，我式克至於今日休。\footnote{言我厚輔是文武之道而行之，或用能至於今日其政美。}我咸成\CJKunderline{文王}功於不怠，丕冒海隅出日，罔不率俾。”\footnote{今我周家皆成\CJKunderline{文王}功於不懈怠，則德教大覆冒海隅日所出之地，無不循化而使之。}

{\noindent\shu\zihao{5}\fzkt “嗚呼”至“率俾”。正義曰：\CJKunderline{周公}言而嘆曰：“嗚呼！我厚輔是二人之道而行之，我用能至於今日其政美。”言今日政美，由是文武之道。“我周家若能皆成\CJKunderline{文王}之功，於事常不懈怠，則德教大覆四海之隅，至於日出之處，其民無不循我化，可臣使也”。戒\CJKunderline{召公}與朝臣皆當法\CJKunderline{文王}之功。 \par}

公曰：“君，予不惠若茲多誥，予惟用閔於天越民。”\footnote{我不順若此多誥而已，欲使汝念躬行之閔勉也。我惟用勉於天道加於民。}

{\noindent\shu\zihao{5}\fzkt “公曰君予”至“越民”。正義曰：公呼\CJKunderline{召公}曰：“君，我不徒惟順如此之事多誥而已，欲使汝躬親行之。我惟用勉力自強於天道，行化於民。”顧氏云：“我亦自用勉勸，躬行於天道，加益於民人也。” \par}

公曰：“嗚呼！君，惟乃知民德，亦罔不能厥初，惟其終。\footnote{惟汝所知民德，亦無不能其初，鮮能有終，惟其終則惟君子。戒\CJKunderline{召公}以慎終。○鮮，息淺反。}祗若茲，往敬用治。”\footnote{當敬順我此言，自今以往，敬用治民職事。}


{\noindent\zhuan\zihao{6}\fzbyks 傳“惟汝”至“慎終”。正義曰：\CJKunderwave{詩}云：“靡不有初，鮮克有終。”是凡民之德，無不能其初，少能有終者。凡民皆如是，有終則惟君子。蓋\CJKunderline{召公}至此已說，恐其不能終善，故戒\CJKunderline{召公}汝慎終也。鄭云：“\CJKunderline{召公}是時意說，\CJKunderline{周公}恐其復不說，故依違託言民德以剴切之。” \par}

{\noindent\shu\zihao{5}\fzkt “公曰嗚呼”至“用治”。正義曰：\CJKunderline{周公}嘆而呼\CJKunderline{召公}曰：“嗚呼！君，惟汝知民之德行，亦無有不能其初，惟鮮能其終。”言行之雖易，終之實難，恐\CJKunderline{召公}不能終行善政，故戒之以慎終。“汝當以敬順我此言，自今以往，宜敬用此治民職事”。戒之使行善不懈怠也。 \par}

%%% Local Variables:
%%% mode: latex
%%% TeX-engine: xetex
%%% TeX-master: "../Main"
%%% End:

%% Time-stamp: <Chen Wang: 2024-04-02 11:42:41>
%% -*- coding: utf-8 -*-

% {\noindent \zhu \zihao{5} \fzbyks } -> 注 (△ ○)
% {\noindent \shu \zihao{5} \fzkt } -> 疏

\chapter{卷十七}


\section{蔡仲之命第十九【偽】}


\CJKunderline{蔡叔}既沒，\footnote{以罪放而卒。}王命\CJKunderline{蔡仲}，踐諸侯位，\footnote{\CJKunderline{成王}也。父卒命子，罪不相及。}作\CJKunderwave{蔡仲之命}。\footnote{冊書命之。}

蔡仲之命\footnote{蔡，國名。仲，字。因以名篇。}


{\noindent\zhuan\zihao{6}\fzbyks 傳“\CJKunderline{成王}”至“相及”。正義曰：編書以世先後為次，此篇在\CJKunderline{成王}書內，知“王命\CJKunderline{蔡仲}”是\CJKunderline{成王}命之也。\CJKunderline{蔡叔}之沒，不知何年，其命\CJKunderline{蔡仲}，未必初卒即命，以其繼父命子，故系之\CJKunderline{蔡叔}之後也。\CJKunderline{蔡叔}有罪而命\CJKunderline{蔡仲}者，“父卒命子，罪不相及”也。昭二十年\CJKunderwave{左傳}曰：“父子兄弟，罪不相及。”其言“罪不相及”，謂\CJKunderline{蔡仲}不坐父爾。若父有大罪，罪當絕滅，正可別封他國，不得仍取蔡名，以\CJKunderline{蔡叔}為始祖也。\CJKunderline{蔡叔}身尚不死，明其罪輕。不立\CJKunderline{管叔}之後者，蓋罪重無子，或有而不賢故也。 \par}

{\noindent\shu\zihao{5}\fzkt “\CJKunderline{蔡叔}”至“之命”。正義曰：\CJKunderline{蔡叔}與\CJKunderline{管叔}流言於國，謗毀\CJKunderline{周公}，\CJKunderline{周公}囚之郭鄰，至死不赦。\CJKunderline{蔡叔}既沒，\CJKunderline{成王}命\CJKunderline{蔡叔}之子\CJKunderline{蔡仲}踐諸侯之位，封為國君，以策書命之。史敘其事，故作\CJKunderwave{蔡仲之命}。 \par}

惟\CJKunderline{周公}位冢宰，正百工，\footnote{百官總己以聽冢宰，謂\CJKunderline{武王}崩時。}群叔流言，乃致闢\CJKunderline{管叔}於商；囚\CJKunderline{蔡叔}於郭鄰，以車七乘；\footnote{致法謂誅殺。囚謂制其出入。郭鄰，中國之外地名。從車七乘，言少。管、蔡，國名。○闢，婢亦反，徐扶亦反。乘，繩證反。從，才用反。}降\CJKunderline{霍叔}於庶人，三年不齒。\footnote{罪輕，故退為眾人，三年之後乃齒錄，封為霍侯，子孫為晉所滅。}\CJKunderline{蔡仲}克庸祗德，\CJKunderline{周公}以為卿士。\footnote{蔡仲能用敬德，稱其賢也。明王之法，誅父用子，言至公。\CJKunderline{周公}，圻內諸侯，二卿治事。○圻，巨依反，下同。}叔卒，乃命諸王邦之蔡。\footnote{叔之所封，圻內之蔡。仲之所封，淮汝之間。圻內之蔡名已滅，故取其名以名新國，欲其戒之。}


{\noindent\zhuan\zihao{6}\fzbyks 傳“致法”至“國名”。正義曰：\CJKunderwave{周禮}有掌囚之官，鄭云，囚，拘也，主拘繫當刑殺者。拘繫之是為制其出入，不得輒行。“郭鄰，中國之外地名”，蓋相傳為然，不知在何方。\CJKunderwave{舜典}云“流宥五刑”，謂流之遠地，任其自生，此則徙之郭鄰，而又囚之。\CJKunderwave{管蔡世家}云“封叔鮮于管，封叔度於蔡”，是管、蔡為國名。杜預云：“管在滎陽京縣東北。” \par}

{\noindent\zhuan\zihao{6}\fzbyks 傳“罪輕”至“所滅”。正義曰：言“群叔流言”，則\CJKunderline{霍叔}亦流言也。而知其罪輕者，以其不死不遷，有降黜而已，明其罪輕也。\CJKunderline{霍叔}不監殷民，\CJKunderline{周公}惟伐管蔡，不言伐\CJKunderline{霍叔}，於時\CJKunderline{霍叔}蓋在京邑，聞管蔡之語，流傳其言，謂其實然，不與朝廷同心，故退之。\CJKunderwave{世家}云“\CJKunderline{武王}已克商平天下，封功臣昆弟，封叔處於霍”，則\CJKunderline{武王}已封之矣。後黜為庶人，奪其爵祿，三年之後乃更爵祿，蓋復其舊封，封為霍侯。\CJKunderwave{春秋}閔元年晉侯滅霍，既子孫得為國君，為晉所滅，知三年之後復得封也。\CJKunderwave{世家}惟雲封霍，不雲其爵，傳言“霍侯”，或當有所據而知之。 \par}

{\noindent\zhuan\zihao{6}\fzbyks 傳“\CJKunderline{蔡仲}”至“治事”。正義曰：\CJKunderwave{周禮·冢宰}：“以八則治都鄙。”\CJKunderline{馬融}云：“距王城四百里至五百里謂之都鄙。鄙，邊邑也，以封王之子弟在畿內者。”\CJKunderwave{冢宰}又云：“乃施則于都鄙而建其長，立其兩。”馬、鄭皆云“立卿兩人”，是畿內諸侯立二卿。定四年\CJKunderwave{左傳}說此事云“\CJKunderline{周公}舉之，以為己卿士”，是為\CJKunderline{周公}圻內之卿士也。\CJKunderwave{世家}云“\CJKunderline{周公}舉胡以為魯卿士，魯國治。於是\CJKunderline{周公}言於\CJKunderline{成王}，復封之於蔡。”案\CJKunderwave{魯世家}云，\CJKunderline{成王}封\CJKunderline{周公}於魯，\CJKunderline{周公}不就封，留佐\CJKunderline{成王}。則\CJKunderline{周公}身不就封，安得使胡為卿士？馬遷說之謬爾。 \par}

{\noindent\zhuan\zihao{6}\fzbyks 傳“叔之”至“戒之”。正義曰：“仲之所封，淮汝之間”，\CJKunderwave{左傳}有文。“叔之所封，圻內之蔡”，其事不知所出也。\CJKunderwave{世家}云：“\CJKunderline{蔡叔}居上蔡。”宋仲子云：“胡徙居新蔡。”杜預云：“\CJKunderline{武王}封叔度於汝南上蔡，至平侯徙新蔡，昭侯徙居九江下蔡。”檢其地，上蔡、新蔡皆屬汝南郡，去京師太遠，叔若封於上蔡，不得在圻內也。孔言叔封圻內，或當有以知之。但圻內蔡地，不知所在爾。 \par}

{\noindent\shu\zihao{5}\fzkt “惟周”至“之蔡”。正義曰：惟\CJKunderline{周公}於\CJKunderline{武王}崩後，其位為冢宰之卿，正百官之治，攝王政，治天下。於時管、蔡、霍等群叔流言於國，謗毀\CJKunderline{周公}。\CJKunderline{周公}乃以王命致法，殺\CJKunderline{管叔}於商，就殷都殺之。囚\CJKunderline{蔡叔}，遷之於郭鄰之地，惟與之從車七乘。降黜\CJKunderline{霍叔}於庶人，若今除名為民，三年之內不得與兄弟年齒相次。\CJKunderline{蔡叔}之子\CJKunderline{蔡仲}能用敬德，\CJKunderline{周公}為畿內諸侯，得立二卿，以\CJKunderline{蔡仲}為己之卿士。\CJKunderline{周公}善其為人，及\CJKunderline{蔡叔}既卒，乃將\CJKunderline{蔡仲}命之於王，國之於蔡為諸侯也。 \par}

王若曰：“小子\CJKunderline{胡}，\footnote{言小子，明當受教訓。胡，仲名。順其事而告之。}惟爾率德改行，克慎厥猷，\footnote{言汝循祖之德，改父之行，能慎其道。嘆其賢。}肆予命爾侯於東土。往即乃封，敬哉！\footnote{以汝率德改行之故，故我命汝為諸侯於東土。往就汝所封之國，當修已以敬哉！}爾尚蓋前人之愆，惟忠惟孝，\footnote{汝當庶幾修德，尚蓋前人之過。子能蓋父，所以為惟忠惟孝。}爾乃邁跡自身，克勤無怠，以垂憲乃後。\footnote{汝乃行善跡用汝身，使可蹤跡而法循之，能勤無懈怠，以垂法子孫，世世稱頌，乃當我意。}率乃祖\CJKunderline{文王}之彝訓，無若爾考之違王命。\footnote{言當循文武之常教，以父違命為世戒。}皇天無親，惟德是輔。民心無常，惟惠之懷。\footnote{天之於人，無有親疏，惟有德者則輔佑之。民之於上，無有常主，惟愛己者則歸之。}為善不同，同歸於治。為惡不同，同歸於亂。\footnote{言人為善為惡，各有百端，未必正同。而治亂所歸不殊，宜慎其微。○治，直吏反。}爾其戒哉！慎厥初，惟厥終，終以不困。不惟厥終，終以困窮。\footnote{汝其戒治亂之機哉！作事云為，必慎其初，念其終，則終用不困窮。}懋乃攸績，睦乃四鄰，以蕃王室，以和兄弟。\footnote{勉汝所立之功，親汝四鄰之國，以蕃屏王室，以和協同姓之邦，諸侯之道。○懋音茂。蕃，方元反，注同。}康濟小民，率自中，無作聰明亂舊章。\footnote{汝為政，當安小民之居，成小民之業，循用大中之道，無敢為小聰明，作異辯，以變亂舊典文章。}詳乃視聽，罔以側言改厥度，則予一人汝嘉。”\footnote{詳審汝視聽，非禮義勿視聽，無以邪巧之言易其常度，必斷之以義，則我一人善汝矣。○度如字，注同。斷，丁亂反。}王曰：“嗚呼！小子\CJKunderline{胡}，汝往哉！無荒棄朕命。”\footnote{嘆而敕之，欲其念戒：“小子胡，汝往之國哉！無廢棄我命。”欲其終身奉行，後世遵則。}


{\noindent\zhuan\zihao{6}\fzbyks 傳“汝當”至“惟孝”。正義曰：忠施於君，孝施於父，子能蓋父，惟得為孝，而亦得為忠者，父以不忠獲罪，若能改父之行，蓋父之愆，是為忠臣也。 \par}

{\noindent\shu\zihao{5}\fzkt “侯於東土”。正義曰：此使之為諸侯於東土爾，不知何爵也。\CJKunderwave{世家}云：“\CJKunderline{蔡仲}卒，子蔡伯荒立。卒，子宮侯立。”自此已下遂皆稱侯，則\CJKunderline{蔡仲}初封即為侯也。“蔡伯荒”者，自稱其字，“伯”非爵也。 \par}

\CJKunderline{成王}東伐淮夷，遂踐奄，\footnote{\CJKunderline{成王}即政，淮夷奄國又叛，王親征之，遂滅奄而徙之，以其數反覆。○踐，似淺反，馬同，\CJKunderwave{大傳}云：“藉也。”數，色角反。覆，芳服反。}作\CJKunderwave{成王政}。\footnote{為平淮夷徙奄之政令。亡。○政如字，馬本作徵，雲正。}


{\noindent\zhuan\zihao{6}\fzbyks 傳“\CJKunderline{成王}”至“反覆”。正義曰：\CJKunderwave{洛誥}之篇言\CJKunderline{周公}歸政\CJKunderline{成王}，\CJKunderwave{多士}已下皆是\CJKunderline{成王}即政初事。編篇以先後為次，此篇在\CJKunderline{成王}書內，知是“\CJKunderline{成王}即政”，淮夷奄國又叛，王親征之”。又案\CJKunderwave{洛誥}\CJKunderline{成王}即政，始封\CJKunderline{伯禽}。\CJKunderline{伯禽}既為魯侯，乃居曲阜。\CJKunderwave{費誓}稱“魯侯\CJKunderline{伯禽}宅曲阜”，“淮夷、徐戎並興”，魯侯徵之，作\CJKunderwave{費誓}。彼言淮夷並興，即此“伐淮夷”。王伐淮夷，魯伐徐戎，是同時伐，明是\CJKunderline{成王}即政之年復重叛也。\CJKunderline{鄭玄}謂此伐淮夷與踐奄是攝政三年伐管蔡時事，其編篇於此，即雲未聞。\CJKunderwave{費誓}之篇言淮夷之叛，則是重叛明矣。\CJKunderwave{多方}之篇責殷臣云：“我惟時其戰要囚之，至於再，至於三。”若\CJKunderline{武王}伐紂之後，惟攝政三年之一叛，正可至於再爾，安得至於三乎？故知是\CJKunderline{成王}即政又叛也。\CJKunderline{鄭玄}讀“踐”為翦，翦滅也。孔不破字，蓋以踐其國即是踐滅之事，故孔以“踐”為滅也。下篇序云：“\CJKunderline{成王}既踐奄，將遷其君”，是滅其奄而徙之，以其數反覆故也。 \par}

{\noindent\shu\zihao{5}\fzkt “\CJKunderline{成王}東”至“王政”。正義曰：\CJKunderline{周公}攝政之初，奄與淮夷從管蔡作亂，\CJKunderline{周公}徵而定之。\CJKunderline{成王}即政之初，淮夷與奄又叛，\CJKunderline{成王}親往徵之。\CJKunderline{成王}東伐淮夷，遂踐滅奄國。以其數叛，徙奄民。作誥命之辭，言平淮夷徙奄之政令。史敘其事，作\CJKunderwave{成王政}之篇。“成”訓平也，言平此叛逆之民，以為王者政令，故以“\CJKunderline{成王}政”為篇名。 \par}

\CJKunderline{成王}既踐奄，將遷其君於蒲姑，\footnote{已滅奄，而徙其君及人臣之惡者於蒲姑。蒲姑，齊地，近中國，教化之。○蒲如字，徐又扶各反，馬本作薄。近，附近之近。}\CJKunderline{周公}告\CJKunderline{召公}，作\CJKunderwave{將蒲姑}。\footnote{言將徙奄新立之君於蒲姑，告\CJKunderline{召公}使此冊書告令之。亡。}


{\noindent\zhuan\zihao{6}\fzbyks 傳“已滅”至“化之”。正義曰：昭二十年\CJKunderwave{左傳}晏子云，古人居此地者，有蒲姑氏。杜預云：“樂安博昌縣北有蒲姑城。”是蒲姑為齊地也。\CJKunderline{周公}遷殷頑民於成周，近京師，教化之，知今遷奄君臣於蒲姑，為“近中國，教化之”。必如此言，則奄去中國遠於蒲姑。杜預云：“奄闕，不知所在。”鄭云“奄蓋在淮夷之地”，亦未能詳。\CJKunderline{成王}先伐淮夷，遂滅奄，奄似遠於淮夷也。 \par}

{\noindent\zhuan\zihao{6}\fzbyks 傳“言將”至“之亡”。正義曰：\CJKunderwave{禮}天子不滅國，諸侯有罪，則殺其君而擇立次賢者，故知所徙者言“將徙奄新立之君於蒲姑”也。上言\CJKunderline{周公}告\CJKunderline{召公}，其篇既亡，不知告以何事。孔以意卜之“告\CJKunderline{召公}使為此策書告令之”，不能知其必然否也。 \par}

{\noindent\shu\zihao{5}\fzkt “\CJKunderline{成王}既”至“作蒲姑”。正義曰：\CJKunderline{成王}既踐滅奄國，將遷其君於蒲姑之地，\CJKunderline{周公}告\CJKunderline{召公}，使作冊書，言將遷奄君於蒲姑之地。史敘其事，作\CJKunderwave{將蒲姑}之篇。 \par}

\section{多方第二十}


\CJKunderline{成王}歸自奄，\footnote{伐奄歸。}在宗周，誥庶邦，\footnote{誥以禍福。}作\CJKunderwave{多方}。

多方\footnote{眾方天下諸侯。}

惟五月丁亥，王來自奄，至於宗周。\footnote{\CJKunderline{周公}歸政之明年，淮夷奄又叛。魯徵淮夷，作\CJKunderwave{費誓}。王親征奄，滅其國，五月還至鎬京。○費音秘。}

{\noindent\zhuan\zihao{6}\fzbyks 傳“眾方天下諸侯”。正義曰：自\CJKunderline{武王}伐紂，及\CJKunderline{成王}即政，新封建者甚少。天下諸侯多是殷之舊國，其心未服周家，由是奄君重叛。今因滅奄新歸，故告天下諸侯以興亡之戒，欲令其無二心也。語雖普告天下，意在殷之舊國。篇末亦告殷之多士，獨言“諸侯”者，舉其尊者，以其篇王告殷之諸侯故也。 \par}

{\noindent\zhuan\zihao{6}\fzbyks 傳“\CJKunderline{周公}”至“鎬京”。正義曰：以\CJKunderwave{洛誥}言歸政之事，\CJKunderwave{多士}之篇次之，\CJKunderwave{多士}是歸政明年之事，故知此篇亦歸政明年之事。事猶不明，故取\CJKunderwave{費誓}為證。以\CJKunderwave{成以政}之序言“\CJKunderline{成王}東伐淮夷”，\CJKunderwave{費誓}之篇言“淮夷、徐戎並興”，俱言“淮夷”，明是一事，故言“魯徵淮夷，作\CJKunderwave{費誓}，王親征奄，滅其國”，以明二者為一時之事也。上序言“\CJKunderline{成王}伐淮夷”，而此傳言“魯徵淮夷”者，當時淮夷徐戎並起為亂，魯與二國相近，發意欲並徵二國，故以二國誓眾，但\CJKunderline{成王}恐魯不能獨平二國，故復親往徵之，所以\CJKunderwave{成王政}之序與\CJKunderwave{費誓}之經並言“淮夷”，為此故也。傳言“五月還至鎬京”，明此“宗周”即鎬京也。\CJKunderwave{禮記·祭統}衛孔悝之鼎銘云“即宮於宗周”，彼“宗周”謂洛邑也。是洛邑亦名宗周，知此是鎬京者，\CJKunderline{成王}以\CJKunderline{周公}歸政之時，暫至洛邑，還歸處西都，鎬京是王常居，知“至於宗周”，至鎬京也。且此與\CJKunderwave{周官}同時事也，\CJKunderwave{周官}序云“還歸在豐”，經云“歸於宗周”，豐、鎬相近，即此“宗周”是鎬京也。 \par}

{\noindent\shu\zihao{5}\fzkt “\CJKunderline{成王}”至“多方”。正義曰：\CJKunderline{成王}歸自伐奄，在於宗周鎬京，諸侯以王徵還，皆來朝集，\CJKunderline{周公}稱王命，以禍福咸告天下諸侯國。史敘其事，作\CJKunderwave{多方}。 \par}

\CJKunderline{周公}曰：“王若曰，猷告爾四國多方。\footnote{\CJKunderline{周公}以王命順大道，告四方。稱\CJKunderline{周公}，以別王自告。○別，彼列反。}惟爾殷侯尹民，我惟大降爾命，爾罔不知。\footnote{殷之諸侯王民者，我大降汝命，謂誅紂也。言天下無不知紂暴虐以取亡。}


{\noindent\zhuan\zihao{6}\fzbyks 傳“\CJKunderline{周公}”至“自告”。正義曰：\CJKunderline{成王}新始即政，\CJKunderline{周公}留而輔之。\CJKunderline{周公}以王命告令諸侯，所告實非王言，故加“\CJKunderline{周公}曰”於“王若曰”之上，以明\CJKunderline{周公}宣\CJKunderline{成王}之意也。“猷”，道也，\CJKunderline{周公}以王命順大道告四方也。既言“四國”，又言“多方”，見四方國多也。不直言“王曰”，稱“\CJKunderline{周公}”，以別王自告也。\CJKunderline{王肅}云：“\CJKunderline{周公}攝政，稱\CJKunderline{成王}命以告。及還政，稱‘王曰’嫌自\CJKunderline{成王}辭，故加‘\CJKunderline{周公}’以明之。”然\CJKunderwave{多士}之篇“王若曰”之上不加“\CJKunderline{周公}曰”者，以彼上句云“\CJKunderline{周公}初於新邑洛，用告”，知是\CJKunderline{周公}故也。 \par}

{\noindent\zhuan\zihao{6}\fzbyks 傳“殷之”至“取亡”。正義曰：諸侯為民之主，民所取正，故謂之“正民”。民以君為命，死生在君，天下之命，在於一人紂，言我大黜下汝之民命，正謂\CJKunderline{武王}誅紂也。言天下無不知紂以暴虐取亡，欲使思念之，令其心棄殷而慕周也。 \par}

{\noindent\shu\zihao{5}\fzkt “\CJKunderline{周公}”至“不知”。正義曰：\CJKunderline{周公}以\CJKunderline{成王}之意告眾方之諸侯曰：“我王順大道以告汝四方之國多方諸侯，惟爾殷之諸侯正民者，我\CJKunderline{武王}大下汝天下民命，誅殺虐紂。汝諸侯天下之民，無有不知紂以暴虐取亡。”欲令其思念之。 \par}

洪惟圖天之命，弗永寅念於祀，惟帝降格於夏。\footnote{大惟為王謀天之命，不長敬念於祭祀。謂\CJKunderline{夏桀}。惟天下至戒於夏以譴告之。謂災異。○譴，棄淺反。}有夏誕厥逸，不肯戚言於民，\footnote{有\CJKunderline{夏桀}不畏天戒而大其逸豫，不肯憂言於民。無憂民之言。}乃大淫昏，不克終日勸於帝之迪，\footnote{言桀乃大為過昏之行，不能終日勸於天之道。○迪，徒歷反，馬本作攸，云：“所也。”行，下孟反。}乃爾攸聞。\footnote{言桀之惡乃汝所聞。}


{\noindent\zhuan\zihao{6}\fzbyks 傳“大惟”至“災異”。正義曰：上天之命，去惡與善，凡為民主，皆當謀之。恐天捨己而去，常須敬念祭祀。天所譴告，謂下災異。天不言，故下災異以譴告，責人主，冀自修政也。 \par}

{\noindent\shu\zihao{5}\fzkt “洪惟”至“攸聞”。正義曰：以諸侯心未服周，故舉夏殷為戒。此章皆說桀亡湯興之事，言\CJKunderline{夏桀}大惟居天子之位，謀上天之命，而不能長敬念於祭祀，惟天下至戒於\CJKunderline{夏桀}。謂下災異譴告之，冀其見災而懼，改修政德。而有\CJKunderline{夏桀}不畏天命，乃大其逸豫，不肯憂言於民，惟乃自樂其身，無憂民之言。\CJKunderline{夏桀}乃復大為淫昏之行，不能終竟一日勉於天之道。言不能一日行天道也。桀之此惡，乃是汝之所聞。言不虛也。 \par}

厥圖帝之命，不克開於民之麗。\footnote{桀其謀天之命，不能開於民所施政教。麗，施也。言昏昧。○麗，力馳反。}乃大降罰，崇亂有夏，因甲於內亂。\footnote{桀乃大下罰於民，重亂有夏。言殘虐。外不憂民，內不勤德，因甲於二亂之內。言昏甚。○重，直用反，又直龍反。}不克靈承於旅，罔丕惟進之恭，洪舒於民。\footnote{言桀不能善奉於人眾，無大惟進恭德，而大舒惰於治民。}亦惟有夏之民叨\xpinyin*{懫}，日欽劓割夏邑。\footnote{桀洪舒於民，故亦惟有夏之民貪叨忿懫而逆命，於是桀民尊敬其能劓割夏邑者。謂殘賊臣。○懫，敕二反。劓，魚器反。}


{\noindent\zhuan\zihao{6}\fzbyks 傳“桀乃”至“昏甚”。正義曰：\CJKunderwave{釋詁}云：“崇，重也。”桀既為惡政，無以悛改，乃復大下罪罰於民，重亂有夏之國。言其殘虐大也。“夾”聲近“甲”，古人“甲”與“夾”通用。夾於二事之內，而為亂行，故傳以二事充之。外不憂民，內不勤德，桀身夾於二亂之內，言其昏暗甚也。鄭、王皆以“甲”為狎，王云：“狎習災異於內外為禍亂。”鄭云：“習為鳥獸之行於內為淫亂。”與孔異也。 \par}

{\noindent\zhuan\zihao{6}\fzbyks 傳“言桀”至“治民”。正義曰：民當奉王，而責桀不能善奉於民眾者，君之奉民，謂設美政於民也。以善奉民，當敬以循之，不敢懈惰。桀乃無大惟進於恭德，而大舒緩懈惰於治民，令民益困，而政益亂也。 \par}

{\noindent\zhuan\zihao{6}\fzbyks 傳“桀洪”至“賊臣”。正義曰：\CJKunderwave{禮記}云：“言悖而出，亦悖而入。”桀既不憂於民，故民亦違逆桀命，為貪饕忿懫之行。文十八年\CJKunderwave{左傳}云：“縉雲氏有不才子，貪於飲食，冒於貨賄，天下之民謂之饕餮。”說者皆言貪財為饕，貪食為餮。“饕”即“叨”也，叨餮謂貪財貪食也。“忿懫”言忿怒違理也。民既如此，桀無如之何，惟日日尊敬其能劓割夏邑者，謂性能殘賊者，任用之。 \par}

{\noindent\shu\zihao{5}\fzkt “厥圖”至“夏邑”。正義曰：又言桀惡。桀其謀天之命，不能開發於民之所施政教。正謂不能開發善政，以施於民。桀乃大下罪罰於民，重亂有夏之國。外不憂民，內不勤德，因復甲於二者之內，為亂之行。桀不能以善道奉承於眾民，無大惟進之恭德，而大舒惰於民。言桀不能進行恭德，而舒惰於治民。桀既舒惰於民，故亦惟有夏之民貪饕忿懫而違逆桀命，於是桀日日尊敬殘賊之臣能劓割夏邑者，任用之，使威服下民也。 \par}

天惟時求民主，乃大降顯休命於\CJKunderline{成湯}，\footnote{天惟是桀惡，故更求民主以代之，大下明美之命於\CJKunderline{成湯}，使王天下。}刑殄有夏，惟天不畀純。\footnote{命湯刑絕有夏，惟天不與桀，亦已大。○畀，必二反。}乃惟以爾多方之義民，不克永於多享。\footnote{天所以不與桀，以其乃惟用汝多方之義民為臣，而不能長久多享國故。}惟夏之恭多士，大不克明保享於民，\footnote{惟桀之所謂恭人眾士，大不能明安享於民。言亂主所任，任同己者。}乃胥惟虐於民，至於百為，大不克開。\footnote{桀之眾士，乃相與惟暴虐於民，至於百端所為。言虐非一。大不能開民以善。言與桀合志。}


{\noindent\zhuan\zihao{6}\fzbyks 傳“惟桀”至“己者”。正義曰：惟桀之所謂恭人眾士，實非恭人。亂主所好，好用同己者，以其同己，謂之為恭人，實非善人，故不能明享於民。杜預訓“享”為受，受國者謂受而有之。此言不能安享於民，謂不能安存享受於民眾也。 \par}

{\noindent\shu\zihao{5}\fzkt “天惟”至“克開”。正義曰：天惟桀惡之故，更求民主以代。天乃大下明美之命於\CJKunderline{成湯}，使之代桀王天下。乃命湯施刑罰絕有夏，惟天不與\CJKunderline{夏桀}，亦已大矣。天所不與之者，乃惟此桀用汝多方之義民為臣，而不能長久於多享國故也。義民實賢人也，\CJKunderline{夏桀}不用。惟\CJKunderline{夏桀}之所謂恭人眾士者，大不能用明道安存享於眾民，乃相與惟行暴虐於民，至於百端所為。言虐無所不作。大不能開民以善，其臣與桀同惡，夏家所以滅亡也。 \par}

乃惟\CJKunderline{成湯}，克以爾多方，簡代夏作民主。\footnote{乃惟\CJKunderline{成湯}，能用汝眾方之賢，大代夏政，為天下民主。}慎厥麗乃勸，厥民刑用勸。\footnote{湯慎其施政於民，民乃勸善。其人雖刑，亦用勸善。言政刑清。}以至於\CJKunderline{帝乙}，罔不明德慎罰，亦克用勸。\footnote{言自湯至於\CJKunderline{帝乙}，皆能成其王道，長慎輔相，無不明有德，慎去刑罰，亦能用勸善。○相，息亮反。去，羌呂反。}要囚，殄戮多罪，亦克用勸。開釋無辜，亦克用勸。\footnote{\CJKunderline{帝乙}已上，要察囚情，絕戮眾罪，亦能用勸善。開放無罪之人，必無枉縱，亦能用勸善。○要，一遙反，又一妙反，注同。殄，亭遍反。上，時掌反。}今至於爾闢，弗克以爾多方，享天之命。\footnote{今至於汝君，謂紂，不能用汝眾方，享天之命，故誅滅之。○闢，必亦反。}嗚呼！


{\noindent\zhuan\zihao{6}\fzbyks 傳“乃惟”至“民主”。正義曰：“大代夏”者，言天位之重，湯能代之，謂之“大代夏”也。\CJKunderline{王肅}云：“以大道代夏為民主。” \par}

{\noindent\zhuan\zihao{6}\fzbyks 傳“湯慎”至“刑清”。正義曰：“慎厥麗”者，總謂施政教爾。但下句言“刑用勸”，勸用刑則厥麗之言有賞，賞謂賞用勸也。但所施政教，其事既多，非徒刑賞而已。舉事得中，民皆勸也。政無失，刑無濫，民以是勸善。言政刑清。 \par}

{\noindent\zhuan\zihao{6}\fzbyks 傳“\CJKunderline{帝乙}”至“勸善”。正義曰：將欲斷罪，必受其要辭，察其虛實，故言“要囚”也。“殄戮多罪”，罪者不濫。開釋無罪者，不枉殺人，不縱有罪，亦是政刑清，故能用勸善也。 \par}

{\noindent\shu\zihao{5}\fzkt “乃惟”至“之命”。正義曰：桀殘虐於民，乃惟\CJKunderline{成湯}，能用汝眾方之賢人，大代\CJKunderline{夏桀}，作天下民主。慎其所施政教於民，民乃勸勉為善。其民雖被刑殺，亦用勸勉為善。非徒湯聖，後世亦賢。自湯至於\CJKunderline{帝乙}，皆能成其王道，無不顯用有德，畏慎刑罰，亦能用勸勉為善。要察囚情，絕戮眾罪，亦能用勸勉為善。開放無罪，亦能用勸勉為善。今至於汝君紂，反先王之道，不能用汝多方之民，享有上天之命，由此故被誅滅。汝等宜當知之，不當更令如殷也。 \par}

王若曰：“誥告爾多方，非天庸釋有夏，\footnote{嘆而順其事以告汝眾方，非天用釋棄桀，桀縱惡自棄，故誅放。}非天庸釋有殷，乃惟爾闢，以爾多方，大淫圖天之命，屑有辭。\footnote{非天用棄有殷，乃惟汝君紂，用汝眾方大為過惡者，共謀天之命，惡事盡有辭說，布在天下，故見誅滅也。}

{\noindent\shu\zihao{5}\fzkt “嗚呼”至“有辭”。正義曰：\CJKunderline{周公}先自嘆，而複稱王命云：“王順其事而言曰，以言告人謂之誥，我告汝眾方諸侯，非天用廢有夏，\CJKunderline{夏桀}縱惡自棄也。非天用廢有殷，殷紂縱惡自棄也。”又指說紂惡：“乃惟汝君殷紂，用汝眾方之民大為過惡者，共此惡人，謀天之命。其惡事盡有辭說，布在天下，以此故見誅滅。” \par}

乃惟有夏圖厥政，不集於享，天降時喪，有邦間之。\footnote{更說桀也。言桀謀其政，不成於享，故天下是喪亡以禍之，使天下有國聖人代之。言有國，明皇天無親，佑有德。○間，間廁之間。}

{\noindent\shu\zihao{5}\fzkt “乃惟”至“間之”。正義曰：更說桀亡之由，乃惟有\CJKunderline{夏桀}謀其政，不能成於享國，所謀皆是惡事，故天下是喪亡以禍之，使有國聖人來代之。言皇天無親，惟佑有德，故以聖君代暗主也。湯是夏之諸侯，故云“有國”。 \par}

乃惟爾商后王，逸厥逸，\footnote{后王紂逸豫其過逸。言縱恣無度。}圖厥政，不蠲烝，天惟降時喪。\footnote{紂謀其政，不絜進於善，故天惟下是喪亡。謂誅滅。○蠲，吉玄反，馬云：“明也。”一音圭。烝，絕句，之承反，馬云：“升也。”}惟聖罔念作狂，惟狂克念作聖。\footnote{惟聖人無念於善，則為狂人。惟狂人能念於善，則為聖人。言桀紂非實狂愚，以不念善，故滅亡。}天惟五年，須暇之子孫，誕作民主，罔可念聽。\footnote{天以湯故，五年須暇湯之子孫，冀其改悔。而紂大為民主，肆行無道，事無可念，言無可聽。\CJKunderline{武王}服喪三年，還師二年。}


{\noindent\zhuan\zihao{6}\fzbyks 傳“惟聖”至“滅亡”。正義曰：“聖”者上智之名，“狂”者下愚之稱。\CJKunderline{孔子}曰：“惟上智與下愚不移。”是聖必不可為狂，狂必不能為聖，此事決矣。而此言“惟聖人無念於善，則為狂人。惟狂人能念於善，則為聖人”者，方言天須暇於紂，冀其改悔，說有此理爾，不言此事是實也。謂之為聖，寧肯無念於善？已名為狂，豈能念善？中人念與不念，其實少有所移，欲見念善有益，故舉狂聖極善惡者言之。 \par}

{\noindent\zhuan\zihao{6}\fzbyks 傳“天以”至“二年”。正義曰：湯是創業聖王，理當祚胤長遠。計紂未死五年之前，已合喪滅，但紂是湯之子孫，天以湯聖人之故，故五年須待閒暇湯之子孫，冀其改悔，能念善道。而紂大為民主，肆行無道。所為皆惡事，無可念者；言皆惡言，無可聽者；由是天始滅之。五年者，以\CJKunderline{武王}討紂，初立即應伐之，故從\CJKunderline{武王}初立之年，數至伐紂為五年。\CJKunderline{文王}受命九年而崩，其年\CJKunderline{武王}嗣立。服喪三年，未得征伐。十一年服闋，乃觀兵於孟津，十三年方始殺紂。從九年至十三年，是五年也。然服喪三年，還師二年，乃事理宜然，而云以湯故須暇之者，以殷紂惡盈，久合誅滅，逢\CJKunderline{文王}崩，未暇行師，兼之示弱，凡經五載，聖人因言之以為法教爾。其實非天不知紂狂，望其後改悔，亦非曲念湯德，延此歲年也。 \par}

{\noindent\shu\zihao{5}\fzkt “乃惟”至“念聽”。正義曰：更說紂亡之由。乃惟汝商之後王紂，逸豫其過，縱恣無度。紂謀其為政，不能絜進於善，惟行惡事，天惟下是喪亡以禍之。惟聖人無念於善，則為狂人。惟狂人能念於善，則為聖人。紂雖狂愚，冀其念善也。計紂為惡，早應誅滅，天惟以\CJKunderline{成湯}之故，故積五年須待閒暇湯之子孫。誕緩多年，冀其改悔。而紂大為民主，肆行無道，事無可念，言無可聽，由是天始改意，故誅滅之。 \par}

天惟求爾多方，大動以威，開厥顧天。\footnote{天惟求汝眾方之賢者，大動紂以威，開其能顧天可以代者。}惟爾多方，罔堪顧之。惟我周王，靈承於旅。\footnote{惟汝眾方之中，無堪顧天之道者。惟我周王，善奉於眾。言以仁政得人心。}克堪用德，惟典神天。\footnote{言周文武能堪用德，惟可以主神天之祀，任天王。○任音壬。}天惟式教我用休，簡畀殷命，尹爾多方。\footnote{天以我用德之故，惟用教我用美道代殷，天與我殷之王命，以正汝眾方之諸侯。}


{\noindent\zhuan\zihao{6}\fzbyks 傳“天惟”至“代者”。正義曰：“天惟求汝眾方之賢”，言欲選賢以為天子也。“大動紂以威”，謂誅殺紂也。天意復開其能顧天可以代者，欲使代之。“顧”謂回視，有聖德者，天回視之。\CJKunderwave{詩}所謂“乃眷西顧，此惟與宅”，與彼“顧”同，言天顧\CJKunderline{文王}而與之居，即此意也。但謂天顧人，人亦顧天，此云“開厥顧天”，謂人顧天也。下云“罔堪顧之”，謂天顧人也。言多方人皆無德，不堪使天顧之。傳以顧事通於彼，故皆以天言之。 \par}

{\noindent\zhuan\zihao{6}\fzbyks 傳“天以”至“諸侯”。正義曰：周以能行美道，乃得天顧，復言天用教我美道者，人之美惡，何事非天？由為美道，為天所顧，以美歸功於天，言教我用美道，故得當天意也。 \par}

{\noindent\shu\zihao{5}\fzkt “天惟”至“多方”。正義曰：天以紂惡之故，將選人代之。惟求賢人於汝眾方，大動紂以威。謂誅去紂也。開其有德能顧天之者，欲以伐紂，惟汝眾方之君，悉皆無德，無堪使天顧之。惟我周王善奉於眾，能以仁政得人心，文武能堪用德，惟可以主神天之祀，任作天子也。天惟以我用德之故，故教我使用美道，大與我殷王之命，命我代殷為王，正汝眾方諸侯。言天授我以此世也。 \par}

今我曷敢多誥，我惟大降爾四國民命。\footnote{今我何敢多誥汝而已，我惟大下汝四國民命。謂誅管、蔡、商、奄之君。}爾曷不忱裕之於爾多方？\footnote{汝何不以誠信行寬裕之道於汝眾方？欲其戒四國，崇和協。}爾曷不夾介乂我周王，享天之命？\footnote{夾，近也。汝何不近大見治於我周王，以享天之命，而為不安乎？○夾音協，注同。}今爾尚宅爾宅，畋爾田，爾曷不惠王熙天之命？\footnote{今汝殷之諸侯皆尚得居汝常居，臣民皆尚得畋汝故田，汝何不順從王政，廣天之命，而自懷疑乎？}爾乃迪屢不靜，爾心未愛。\footnote{汝所蹈行，數為不安，汝心未愛我周故。○數，色各反。}爾乃不大宅天命，爾乃屑播天命，\footnote{汝乃不大居安天命，是汝乃盡播棄天命。}爾乃自作不典圖忱於正。\footnote{汝未愛我周，播棄天命，是汝乃自為不常謀信於正道。}


{\noindent\zhuan\zihao{6}\fzbyks 傳“今我”至“之君”。正義曰：今我何敢多為言誥而已，實殺其君，非徒口告。管、蔡、商、奄，皆為叛逆受誅，故今因奄重叛而追說前事，言下四國民命。\CJKunderline{王肅}以“四國”為四方之國，言“從今以後，四方之國苟有此罪，則必誅之”。謂戒其將來之事，與孔不同。 \par}

{\noindent\zhuan\zihao{6}\fzbyks 傳“夾近”至“安乎”。正義曰：夾其旁，旁是近義，故為近也。諸國疏遠周室，不肯以治為功，故責之。顧氏云：“汝眾方諸侯，何不常和協，相親近，大顯見治道於我周王，以享上天之命？而今何以不自安乎？” \par}

{\noindent\zhuan\zihao{6}\fzbyks 傳“今汝”至“疑乎”。正義曰：主遷於上，臣易於下，計汝諸侯之國，應隨殷降黜。今汝殷之諸侯皆尚得居汝常居，臣民畋汝故田。田宅不易，安樂如此，汝何不順從我周王之政，以廣上天之命，使天多佑？汝何故畏我周家，自懷疑乎？諸侯有國，故云“居汝常居”。臣民重田，故云“畋汝故田”。治田謂之“畋”，猶捕魚謂之“漁”，今人以營田求食謂之“畋食”，即此“畋亦田”之義也。 \par}

{\noindent\zhuan\zihao{6}\fzbyks 傳“汝未”至“正道”。正義曰：事君無二臣之道，為人臣者，常宜信之。汝未愛我周家，播棄天命，汝數為叛逆，是汝乃自為此不常謀信於正道。 \par}

我惟時其教告之，我惟時其戰要囚之，\footnote{我惟汝如是不謀信於正道，故其教告之，謂訊以文誥；其戰要囚之，謂討其倡亂，執其朋黨。○要，一遙反。訊音信。倡音唱。}至於再，至於三。\footnote{再，謂三監淮夷叛時。三，謂\CJKunderline{成王}即政又叛。言迪屢不靜之事。}乃有不用我降爾命，我乃其大罰殛之。\footnote{我教告戰要囚汝已至再三，汝其不用我命，我乃大下誅汝君，乃其大罰誅之。○殛，紀力反，本又作極。}非我有周秉德不康寧，乃惟爾自速辜。”\footnote{非我有周執德不安寧，自誅汝，乃惟汝自召罪以取誅。}

{\noindent\zhuan\zihao{6}\fzbyks 傳“我惟”至“朋黨”。正義曰：“教告”與“戰要囚”連文，則告以文辭，是將戰之時，“教告”謂伐紂之事。昭十三年說戰法云：“告之以文辭，董之以武師。”是將戰之時，於法當有文辭告前敵也。我惟汝如是不謀信於正道，故其教告之，謂訊以文辭。“訊”，告也，告以文辭，數其罪也。其“戰要囚之”，謂戰敗其師，執取其人，受其要辭而囚之。謂討其倡亂之人，囚執其朋黨也。此雖總言戰事，但下有至於再三，明此指伐紂也。 \par}

{\noindent\zhuan\zihao{6}\fzbyks 傳“再謂”至“之事”。正義曰：以伐紂為一，故“再”謂攝政之初，三監與淮夷叛時也，“三”謂\CJKunderline{成王}即政又叛也，言上“迪屢不靜”之事。 \par}

{\noindent\shu\zihao{5}\fzkt “今我”至“速辜”。正義曰：今我何敢多以言誥告於汝眾而已，我惟大下黜汝管、蔡、商、奄四國之君也。“民命”，謂民以君為命，謂誅殺四國之君也。我已殺汝四國君矣，汝何不以誠信之心，行寬裕之道於汝眾方諸侯？欲令懲創四國，務崇和協。言汝眾方諸侯何不崇和協，相親近，大顯見治道於我周王，以享愛上天之命，而執心不安乎？今爾殷之諸侯尚得居汝常居，臣民尚得畋汝故田，其安樂如此，汝何得不順從王政，以廣上天之命，而自懷疑乎？汝乃復所蹈行者，數為不安，時或叛逆，是汝心未愛我周家故也。汝乃不大居安天命，是汝乃欲盡播棄天命。汝不愛我周家，播棄天命，是汝乃自為此不常謀信於正道。言其心不常謀正道，故為背違之心。我惟汝如是不謀信於正道之故，其以言辭教告之。我惟汝如是不誠信於正道之故，其用戰伐要察囚繫之。由汝數為不信，故我教告汝，戰伐要囚汝，至於再，至於三。我教告汝，戰伐要囚汝，已至再三，如今而後乃復有不用我命者，我乃其大罰誅之。言我更將殺汝也。非我有周執德不安，數設誅罰，乃惟汝自召罪也。此章反覆殷勤者，恐其更有叛逆，故丁寧戒之。 \par}

王曰：“嗚呼！猷告爾有方多士暨殷多士，\footnote{王嘆而以道告汝眾方與殷多士。}今爾奔走臣我監五祀，\footnote{監謂成周之監，此指謂所遷頑民殷眾士。今汝奔走來徙臣我我監，五年無過，則得還本土。}越惟有胥伯小大多正，爾罔不克\CJKunderline{臬}。\footnote{於惟有相長事小大眾正官之人，汝無不能用法。欲其皆用法。○臬，魚列反，馬作臬刂。長，丁丈反。}自作不和，爾惟和哉！爾室不睦，爾惟和哉！爾邑克明，爾惟克勤乃事。\footnote{小大多正自為不和，汝有方多士，當和之哉！汝親近室家不睦，汝亦當和之哉！汝邑中能明，是汝惟能勤汝職事。}爾尚不忌於凶德，亦則以穆穆在乃位，\footnote{汝庶幾不自忌，入於凶德，亦則用敬敬常在汝位。}克閱於乃邑謀介，爾乃自時洛邑，尚永力畋爾田。\footnote{汝能使我閱具於汝邑，而以汝所謀為大，則汝乃用是洛邑，庶幾長力畋汝田矣。言雖遷徙，而以修善，得反邑里。○閱音悅。}天惟畀矜爾，我有周惟其大介賚爾。\footnote{汝能修善，天惟與汝憐汝，我有周惟其大夫賜汝。言受多福之胙。}迪簡在王庭，尚爾事，有服在大僚。”\footnote{非但受憐賜，又乃蹈大道在王庭，庶幾修汝事，有所服行在大官。}


{\noindent\zhuan\zihao{6}\fzbyks 傳“王嘆”至“多士”。正義曰：言“有方多士與殷多士”，則此二者非一人也。“有方多士”當謂於時所有四方之諸侯也。“與殷多士”當謂遷於成周頑民之眾士也。下雲以“臣我監”者，謂成周之監，明此殷多士也。 \par}

{\noindent\zhuan\zihao{6}\fzbyks 傳“監謂”至“本土”。正義曰：下云“自時洛邑”，此所戒成周之人，故知“監謂成周之監，此指謂所遷頑民殷家眾士”也。五年再閏，天道有成，故期以五年無過，則得還本土。以民性重遷，設期以誘之。 \par}

{\noindent\zhuan\zihao{6}\fzbyks 傳“於惟”至“用法”。正義曰：“胥”，相也。“伯”，長也。顏氏以“相長事”即“小大眾正官之人”也。 \par}

{\noindent\zhuan\zihao{6}\fzbyks 傳“汝庶”至“汝位”。正義曰：和順為善德，怨惡為凶德。“忌”謂自怨忌，上言“自作不和”，是怨忌也。\CJKunderwave{釋訓}云：“穆穆，敬也。”此戒小大正官之人，故云“敬敬常在汝位”。 \par}

{\noindent\zhuan\zihao{6}\fzbyks 傳“汝能”至“邑里”。正義曰：“閱”謂簡閱其事，觀其具足以否，故言“閱具於汝邑”。“介”，大也。以汝所謀為大，善其治理，聽還本國也。是由在洛邑修善，得反其邑里。\CJKunderline{王肅}云：“其無成，雖五年亦不得反也。” \par}

{\noindent\shu\zihao{5}\fzkt “王曰嗚呼猷”至“大僚”。正義曰：王言而嘆曰：“嗚呼！我以道告汝在此所有四方之多士。”謂四方之諸侯及與殷之眾士，謂頑民遷成周者。因告四方諸侯，遂告成周之人，遍使諸侯知之。此章皆告成周之人辭也。“今汝成周之人，奔走勤事，臣我周之監成周者，五年無罪過，則聽汝還本土。於惟有相長事，謂小大眾正官之人，汝無有不能用法。”欲其皆用法也。“小大眾正官之人自為不和，汝眾官等自當和之哉！汝等親近室家不相和親，汝亦當和之哉！汝邑內之內若能明於和睦之道，汝惟能勤於汝之職事”。言是其教之使然。“汝能庶幾不自相怨忌，入於凶德，若能不入於凶德，亦則用敬敬之道，常在汝之職位，不黜退也。汝若能善相教誨，使我簡閱於汝邑，善汝之事，以汝所謀為大，則汝乃用是洛邑，庶幾得反本土，長得勤畋汝故田。汝能修善，天惟與汝憐汝，我有周惟其大大賞賜汝。汝非但受賞而已，其有蹈大道者，得在王庭被任用。庶幾汝事有所服行，在於大官”。恐其心未服，故丁寧勸誘之。 \par}

王曰：“嗚呼！多士，爾不勸忱我命，爾亦則惟不克享，凡民惟曰不享。\footnote{王嘆而言曰：“眾士，汝不能勸信我命，汝亦則惟不能享天祚矣，凡民亦惟曰不享於汝祚矣。”}爾乃惟逸惟頗，大遠王命，則惟爾多方探天之威，我則致天之罰，離逖爾土。”\footnote{若爾乃為逸豫頗僻，大棄王命，則惟汝眾方取天之威，我則致行天罰，離逖汝土，將遠徙之。○頗，破多反。探，吐南反。闢，四亦反。}


{\noindent\zhuan\zihao{6}\fzbyks 傳“王嘆”至“祚矣”。正義曰：“勸信我命”，勸勉而信順之。“凡民亦惟曰不享於汝祚矣”，言民亦不原汝之子孫長久矣。 \par}

{\noindent\zhuan\zihao{6}\fzbyks 傳“若爾”至“徙之”。正義曰：成周一邑之士，不得謂之多方，此蓋意在成周遷者，兼告四方諸國使知。亦如\CJKunderwave{康誥}王告\CJKunderline{康叔}，並使諸侯知之。離遠汝土，更遠徙之。鄭云：“分離奪汝土也。”與孔異也。 \par}

{\noindent\shu\zihao{5}\fzkt “王曰嗚呼”至“爾土”。正義曰：王言而嘆曰：“嗚呼！成周之眾士，汝若不能勸勉信用我之教命，汝則惟不能多受天福祚矣，凡民惟曰不享於汝祚矣。汝乃惟為逸豫，惟為頗僻，大遠棄王命，則惟汝眾方自取天之威刑，我則致天之罰於汝身，將遠徙之，使離遠汝之本土。” \par}

王曰：“我不惟多誥，我惟祗告爾命。”\footnote{我不惟多誥汝而已，我惟敬告汝吉凶之命。}又曰：“時惟爾初，不克敬於和，則無我怨。”\footnote{又誥汝：“是惟汝初不能敬於和道，故誅汝。汝無我怨。”解所以再三加誅之意。}


{\noindent\zhuan\zihao{6}\fzbyks 傳“又誥”至“之意”。正義曰：“又誥”者，更言王意，又謂汝曰也。以上王誥已終，又起別端，故更稱王又復言曰。以序雲\CJKunderline{成王}在豐誥庶邦，則此篇是王親誥之辭，又稱“王曰”者是也。其有\CJKunderline{周公}稱王告者，則上云“\CJKunderline{周公}曰：王若曰”是也，又曰“嗚呼！王若曰”是也。顧氏云：“又曰者，是王又復言曰也。” \par}

{\noindent\shu\zihao{5}\fzkt “王曰我”至“我怨”。正義曰：王曰：“我今告戒汝者，不惟多為言誥汝而已，惟敬告汝吉凶之命。從我則吉，違我則兇，汝命吉凶在此言也。”王又謂：“汝所以再三被誅者，是惟汝初不能敬於和道，故致此爾。汝自取之，則無於我有怨。” \par}

\section{立政第二十一}


\CJKunderline{周公}作\CJKunderwave{立政}。\footnote{\CJKunderline{周公}既致政\CJKunderline{成王}，恐其怠忽，故以君臣立政為戒。}

立政\footnote{言用臣當共立政，故以名篇。}

\CJKunderline{周公}若曰：“拜手稽首，告嗣天子王矣。”\footnote{順古道盡禮致敬，告\CJKunderline{成王}，言：“嗣天子，今已為王矣，不可不慎。”}用咸戒於王曰：“王左右常伯、常任、準人、綴衣、虎賁。”\footnote{\CJKunderline{周公}用王所立政之事皆戒於王曰，常所長事、常所委任，謂三公六卿；準人平法，謂士官；綴衣掌衣服，虎賁以武力事王，皆左右近臣，宜得其人。○任，而鴆反。準，之允反。綴，徐丁衛反，又丁劣反。賁音奔。長，丁丈反，除篇末文注“以長”直良反，餘並同。}\CJKunderline{周公}曰：“嗚呼！休茲，知恤鮮哉！。\footnote{嘆此五者立政之本，知憂得其人者少。○鮮，息淺反。}


{\noindent\zhuan\zihao{6}\fzbyks 傳“順古”至“不慎”。正義曰：\CJKunderline{周公}既拜手稽首，而後發言。還自言“拜手稽首”，示己重其事，欲令受其言，故盡禮致敬以告王也。\CJKunderwave{召誥}云：“拜手稽首，旅王若公。”亦是\CJKunderline{召公}自言已拜手稽首，與此同也。\CJKunderline{成王}嗣世而立，故呼\CJKunderline{成王}為“嗣天子”。\CJKunderline{周公}攝政之時，\CJKunderline{成王}未親王事，此時既已歸政於\CJKunderline{成王}，故言“今已為王矣，不可不慎”也。\CJKunderline{王肅}以為於時\CJKunderline{周公}會群臣共戒\CJKunderline{成王}，其言曰“拜手稽首”者，是\CJKunderline{周公}贊群臣之辭。 \par}

{\noindent\zhuan\zihao{6}\fzbyks 傳“\CJKunderline{周公}”至“其人”。正義曰：此以“立政”名篇，知“用咸戒”者是“\CJKunderline{周公}用王所立政之事皆戒於王”也。三公，臣之尊者，知常所長事謂三公也。六卿分掌國事，王之所任，知常所委任謂六卿也。“準”訓平也，平法之人謂士官也。士，察也，察獄之官用法必當均平，故謂獄官為準人。\CJKunderwave{周禮}司寇之長在常任之內，此士官當謂士師也。衣服必連綴著之，此歷言官人，知“綴衣”是掌衣服者。此言親近大臣，必非造衣裳者。\CJKunderwave{周禮}：“大僕，下大夫。掌正王之服位，出入王之大命。”此掌衣服者，當是大僕之官也。\CJKunderwave{周禮}：“虎賁氏，下大夫。”言其若虎賁獸，是以武力事王者。此皆左右近臣，宜得其人，言其急於餘官。得其人者，文官得其文人，武官得其武人，違才易務，皆為非其人也。 \par}

{\noindent\zhuan\zihao{6}\fzbyks 傳“嘆此”至“者少”。正義曰：此五官皆親近王，故嘆此五者立政之本也。“休”，美也。\CJKunderline{王肅}云“此五官美哉”，是“休茲”為美此五官也。嘆其官之美，美官不可不委賢人用之，故嘆之。“知憂得其人者少”，下句惟言\CJKunderline{禹}、湯、文、武官得其人，是知憂得人者少也。 \par}

{\noindent\shu\zihao{5}\fzkt “\CJKunderline{周公}”至“鮮哉”。正義曰：王之大事在於任賢使能，\CJKunderline{成王}初始即政，猶尚幼少，\CJKunderline{周公}恐其怠忽政事，任非其人，故告以用臣之法。\CJKunderline{周公}順古道而告王曰：“我敢拜手稽首，告嗣世天子\CJKunderline{成王}，今已為王矣。王者當立善政，其事不可不慎。”\CJKunderline{周公}既為此言，乃用王所立政之事皆戒於王曰，王之親近左右，常所長事，謂三公也；常所委任，謂六卿也；平法之人，謂獄官也；綴衣之人，謂掌衣服者也；虎賁，以武力事王者；此等皆近王左右，最須得人。\CJKunderline{周公}既歷言此官，復言而嘆曰：“嗚呼，美哉！此五等之官，立政之本也。知憂此官宜得賢人者少也。” \par}

古之人迪惟有夏，乃有室大競，籥俊尊上帝，\footnote{古之人道惟有夏禹之時，乃有卿大夫室家大強，猶乃招呼賢俊，與共尊事上天。○籥音預。}迪知忱恂於九德之行。\footnote{禹之臣蹈知誠信於九德之行，謂賢智大臣。九德，\CJKunderline{皋陶}所謀。○忱，市林反。恂音荀。行如字，徐下孟反。}乃敢告教厥後曰：‘拜手稽首，後矣！’曰，‘宅乃事，宅乃牧，宅乃準，茲惟後矣。\footnote{知九德之臣乃敢告教其君以立政。君矣，亦猶王矣。宅，居也，居汝事，六卿掌事者。牧，牧民，九州之伯。居內外之官及平法者皆得其人，則此惟君矣。}謀面，用丕訓德，則乃宅人，茲乃三宅無義民。’\footnote{謀所面見之事，無疑則能用大順德，乃能居賢人於眾官。若此則乃能一居無義民。大罪宥之四裔，次九州之外，次中國之外。}\CJKunderline{桀}德惟乃弗作往任，是惟暴德，罔後。\footnote{桀之為德，惟乃不為其先王之法、往所委任，是惟暴德之人，故絕世無後。}


{\noindent\zhuan\zihao{6}\fzbyks 傳“古之”至“上天”。正義曰：經言“古之人迪”，傳言“古之人道”，當說古之求賢人之道也。\CJKunderline{王肅}云：“古之人道，惟有夏之\CJKunderline{大禹}為天子也。”其意言古人之道說有此事，孔意似不然也。孔以“大夫”稱家室，猶家也。“籥”訓呼也，招呼者乃是臣下之事，故以為夏禹之時，乃有卿大夫室家大強，猶乃招呼在外賢俊，與之共立於朝，尊事上天也。言君既求賢臣之助，言天子事天，臣成君事，故言“共尊事上天”。 \par}

{\noindent\zhuan\zihao{6}\fzbyks 傳“\CJKunderline{禹}之”至“所謀”。正義曰：九德之行，非一人能備，言\CJKunderline{禹}之臣蹈知九德之行，極言其賢智大臣也。\CJKunderline{禹}時伯益之輩，乃可以當此。經典之文，更無九德之事，惟有\CJKunderline{皋陶}謀九德，故言九德。\CJKunderline{皋陶}所謀者，即“寬而慄，柔而立，愿而恭，亂而敬，擾而毅，直而溫，簡而廉，剛而塞，強而義”是也。 \par}

{\noindent\zhuan\zihao{6}\fzbyks 傳“知九”至“君矣”。正義曰：進言戒君非大賢不可，故“知九德之臣乃敢告教其君以立政”也。“君矣”亦猶言“王矣”，言已為君矣，不可不慎也。“君”、“王”一也，變文以相避爾。“宅”訓居也，居汝事，須得賢人，六卿各掌其事者也。居汝牧，九州之伯主養民，亦須得賢人養其民也。居汝準，士官主理刑法，亦須賢人平其獄也。六卿掌內，州牧掌外，內外之官及平法三事皆得其人，則此惟為君矣。言群官失職，則不成為君也。上句\CJKunderline{周公}戒王，歷言五官，其內無州牧。此惟言三官，加州牧者，俱是逐急言之，其有詳略爾。\CJKunderwave{曲禮}云：“九州之長曰牧。”\CJKunderwave{王制}云：“千里之外設方伯，八州八伯。”然則“牧”、“伯”一也。“伯”者言一州之長，“牧”者言牧養下民，“牧”、“伯”俱得言之，故孔以“伯”解“牧”。\CJKunderline{鄭玄}云：“殷之州牧曰伯，虞夏及周曰牧。”與孔不同。 \par}

{\noindent\zhuan\zihao{6}\fzbyks 傳“謀所”至“之外”。正義曰：凡人為主，皆欲臣賢，但大佞以忠，賢不可別。欲知其遠，先驗於近，但\CJKunderline{禹}能謀所面見之事。官賢人，既得其官，分別善惡，無所疑惑。仁賢必用，邪佞必退，然後舉直錯諸枉，則為能用大順德，如是乃能居賢人於眾官。賢人既得居官，則能分別善惡，無義之民必獲大罪。量其輕重，斥之遠地，乃能三處居此無義罪人。三居者，“大罪宥之四裔，次九州之外，次中國之外。”“四裔”者，四海之表最遠者也。“次九州之外者”，四海之內，要服之外。“次中國之外”者，謂罪人所居之國外也，猶若衛人居於晉，去本國千里。故孔注\CJKunderwave{舜典}云：“次千里之外是也。”鄭云：“三處者，自九州之外至於四海，三分其地，遠近若周之夷、鎮、蕃也。”與孔不同。 \par}

{\noindent\shu\zihao{5}\fzkt “古之人”至“罔後”。正義曰：既言知憂得人者少，乃遠述上世之事，此言\CJKunderline{禹}與桀也。古之人能用此求賢之道者，惟有夏禹之時。乃有群臣卿大夫皆是賢人，室家大強，猶尚招呼賢俊之人，與共立於朝，尊事上天。\CJKunderline{禹}之臣蹈知誠信於九德之行者，乃敢告教其君曰：“我敢拜手稽首，君今已為君矣，不可不慎也。”戒其君即告曰，居汝掌事之六卿，居汝牧民之州伯，居汝平法之獄官，使此三者皆得其人，則此惟為君矣。言不得賢人，不成為君也。\CJKunderline{禹}能謀所面見之事，無所疑惑，用大明順之德，則乃能居賢人於官。賢人在官，職事修理，乃能三處居無義之民。善人在朝，惡人黜遠，其國乃為治矣。及其末年，桀乃為天子。桀之為德，惟乃不為其先王之法、往所委任，是暴德之人。以此故，絕世無後。得賢人則興，任小人則滅，是須官賢人以立政也。 \par}

亦越\CJKunderline{成湯}陟，丕釐上帝之耿命，\footnote{桀之昏亂，亦於\CJKunderline{成湯}之道得升，大賜上天之光命，王天下。○釐，力之反。耿，工迥反，徐工穎反，又工永反，下同。王，往況反。}乃用三有宅，克即宅，曰三有俊，克即俊。\footnote{湯乃用三有居惡人之法，能使就其居。言服罪。又曰，能用剛柔正直三德之俊，能就其悛事。言明德。}嚴惟丕式，克用三宅三俊。\footnote{言湯所以能嚴威，惟可大法象者，以能用三居三德之法。}其在商邑，用協於厥邑，其在四方，用丕式見德。\footnote{湯在商邑，用三宅三俊之道和其邑。其在四方，用是大法見其聖德。言遠近化。}


{\noindent\zhuan\zihao{6}\fzbyks 傳“桀之”至“天下”。正義曰：“\CJKunderline{成湯}之道得升”，謂從下而升於天，故天賜之以光命，使之得王天下為天子也。“釐，賜”、“耿，光”皆\CJKunderwave{釋詁}文。 \par}

{\noindent\zhuan\zihao{6}\fzbyks 傳“湯乃”至“明德”。正義曰：\CJKunderwave{皋陶謨}“九德”即\CJKunderwave{洪範}之“三德”，細分以為九爾。以此知“三俊”即是\CJKunderwave{洪範}所言“剛克、柔克、正直”三德之俊也。“能就其俊事。言明德”者，用以俊乂居官，顯明其有德也。上句言“則乃宅人，茲乃三宅無義民”，先言用賢，後言去惡。此經先言“三有宅”，後言“曰三有俊”者，用賢去惡，俱是立政之本。上句先說夏禹，言得賢然後去惡，見其須賢之切。及說\CJKunderline{成湯}、文、武，先言去惡，後言用賢，又見惡宜速去。或先或後，所以互相見爾。 \par}

{\noindent\shu\zihao{5}\fzkt “亦越”至“見德”。正義曰：不有所廢，則無以興。桀之滅亡，夏家乃以開道湯德。此言湯之能用人也。桀之昏亂，亦於\CJKunderline{成湯}之道得升聞於天，大賜受上天之光命，得王有天下。湯既為王，乃用三有居惡人之法，能使各就其居處。言皆服其罪也。又曰用三德之俊人，能使各就其俊事。言皆明其德也。湯所以能嚴威，惟可大法象者，以其能用三居三俊之法故也。\CJKunderline{成湯}其在商邑，用此三居三俊之道，和於其邑。其在四方，用是斷罪任賢之大法，見其聖德於民。言遠近皆從化也。 \par}

嗚呼！其在\CJKunderline{受}德暋，惟羞刑暴德之人，同於厥邦。\footnote{受德，紂字。\CJKunderline{帝乙}愛焉，為作善字，而反大惡自強，惟進用刑，與暴德之人同於其國，併為威虐。○受德，紂字，馬云：“受所為德也。”暋，眉謹反，徐亡巾反，一音閔。為，於為反，下“為之”同。強，其丈反。}乃惟庶習逸德之人，同於厥政。\footnote{乃惟眾習為過德之人，同於其政。言不任賢。}帝欽罰之，乃伻我有夏，式商受命，奄甸萬姓。\footnote{天以紂惡，故敬罰之。乃使我周家王有華夏，得用商所受天命，同治萬姓。言皇天無親，佑有德。○伻，普耕反，徐敷耕反，又甫耕反。}


{\noindent\zhuan\zihao{6}\fzbyks 傳“受德”至“威虐”。正義曰：\CJKunderwave{泰誓}三篇，惟單言“受”，而此云“受德”者，則“德”本配“受”，共為一人，故知“受德”是紂字也。既“受”之與“德”共為紂字，而經或言“受”，或言“受德”者，呼之有單復爾。其人實為大惡，“德”字乃為善名，非是時人呼有德。知是\CJKunderline{帝乙}愛焉，為作善字，望其為善，而反為大惡，以其行反其字，明非時人呼也。\CJKunderwave{釋詁}云：“暋，強也。”“暋”即昏也，故訓為強，言紂自強為惡，惟進用刑罰。身既進用刑罰，則愛好暴虐之人，故為與之同於其國，言併為威虐。 \par}

{\noindent\zhuan\zihao{6}\fzbyks 傳“乃惟”至“任賢”。正義曰：“暴德”言以暴虐為德，“逸德”言以過惡為德。習效為之眾者，言其所任多也。紂任眾為過德之人，與之同於其政，言其不任賢也。與暴德同於其國，與惡德同於其政，其事一也，異言之爾。\CJKunderwave{牧誓}所云“四方之多罪逋逃，是信是使，是以為大夫卿士。俾暴虐於百姓，以奸究於商邑”，是其事也。 \par}

{\noindent\zhuan\zihao{6}\fzbyks 傳“天以”至“有德”。正義曰：言天知其惡，熟詳審下罰，故言“敬罰”也。商本受天命，周亦受天命，故言“用商所受天命，同治萬姓”。\CJKunderwave{釋言}云：“弇，同也。”同為天於治萬姓，與商同也。此經之意，言周家有德，皇天親有德也。\CJKunderline{王肅}云：“敬罰者，謂須暇五年。” \par}

{\noindent\shu\zihao{5}\fzkt “嗚呼”至“萬姓”。正義曰：既言湯以用賢而興，又說紂之失人而滅。\CJKunderline{周公}又嘆曰：“嗚呼！其在殷王受德，本性大惡自強，惟進用刑罰，與暴德之人同治其國，併為威虐。乃惟眾習為過德之人，與之同共於其政。由其任同惡之人，故上天敬誅罰之，乃使我周家王有華夏，用商所受天命，同治天下萬姓。”言周能用賢，天親有德，故得為天子。 \par}

“亦越\CJKunderline{文王}、\CJKunderline{武王}，克知三有宅心，灼見三有俊心，\footnote{紂之不善，亦於文武之道大行，以能知三有居惡人之心，灼然見三有賢俊之心。}以敬事上帝，立民長伯。\footnote{言文武知三宅三俊，故能以敬事上天，立民正長。謂郊祀天，建諸侯。}


{\noindent\zhuan\zihao{6}\fzbyks 傳“紂之”至“之心”。正義曰：桀之昏亂開\CJKunderline{成湯}，紂之不善開文武，其事同也。於\CJKunderline{成湯}言能受上天之命，於文武雲能敬事上帝，前聖後聖為行必同，交錯為文，所以互相見爾。\CJKunderline{文王}受命，\CJKunderline{武王}伐紂，二聖共\CJKunderline{成王}道，故文武總言之。猶\CJKunderwave{詩序}云“文武以\CJKunderwave{天保}已上治內，\CJKunderwave{采薇}已下治外”，文武並言，與此同也。\CJKunderline{文王}之時，未定天下，所立之官，亦未具足。下經所言“立政任人”已下，“三亳阪尹”已上，其所舉官屬，多是文武時事，以見二聖同道，父作之，子述之，言其相成爾。故以“能知三有居惡人之心，灼然見三有賢俊之心”，言\CJKunderline{文王}之聖心能揆度知惡人真惡，須屏黜之；知賢人實賢，須舉用之；故去惡進賢，皆得其所。賢人難識，故特言“灼然”，言其知之審也。 \par}

{\noindent\zhuan\zihao{6}\fzbyks 傳“言文”至“諸侯”。正義曰：上天之道，與善去惡，三宅三俊，行合天心。言文武知三宅三俊，故能敬事上帝。“伯”亦長也，故言“立民正長”。天子祭天，知“敬事上帝”，謂“郊祀天”也。天子建國，知“立民長伯”，謂“建諸侯”也。以下句“立政任人”已下，歷言朝廷之臣與蠻夷眾君，知此“立民長伯”主謂諸侯。\CJKunderwave{詩·周頌·維清}述\CJKunderline{文王}之德言“肇禋”，\CJKunderwave{大雅·皇矣}美\CJKunderline{文王}之伐言“是類”，“類”、“禋”皆是祭天之名，是\CJKunderline{文王}已祀天矣。\CJKunderline{文王}未得封建諸侯，其建諸侯，維\CJKunderline{武王}時爾。 \par}

{\noindent\shu\zihao{5}\fzkt “亦越”至“長伯”。正義曰：既言上天去惡與善，滅殷興周，即說\CJKunderline{文王}、\CJKunderline{武王}能用求賢審官之事。桀惡所以興\CJKunderline{成湯}，紂惡所以開文武，言紂之不善，亦於\CJKunderline{文王}、\CJKunderline{武王}使得其道大行。能知居三有惡人之心，居之皆得其所，言服其罪也。灼然見三有賢俊之心，用之皆得其人，言明其德也。文武知此三宅三俊，故能敬事上天，稱天心也。立民正長，合民心也。 \par}

立政，任人、準夫、牧，作三事，\footnote{文武亦法\CJKunderline{禹}湯以立政，常任、準人及牧，治為天地人之三事。}虎賁、綴衣、趣馬小尹，\footnote{趣馬，掌馬之官。言此三者雖小官長，必慎擇其人。○趣，七口反。}左右攜僕、百司庶府，\footnote{雖左右攜持器物之僕，及百官有司主券契藏吏，亦皆擇人。○券音勸。契，苦計反。藏，才浪反。}大都小伯、藝人表臣、百司，\footnote{小臣猶皆慎擇其人，況大都邑之小長，以道藝為表幹之臣及百官有司之職，可以非其任乎？}


{\noindent\zhuan\zihao{6}\fzbyks 傳“文武”至“三事”。正義曰：前聖後聖，其道皆同，未必相放法也。後人法前，自是常事，因其上說\CJKunderline{禹}湯立政，故言“文武亦法\CJKunderline{禹}湯以立政”也。“任人”則前經所云“常任”六卿也，“準夫”則“準人”也，“牧”者前云“宅乃牧”也。前文有“常伯”、“綴衣”、“虎賁”，不言“牧”，此不言“常伯”、“綴衣”、“虎賁”而言“牧”者，以前文先舉朝臣，故不言“牧”，前已備文，故此不言“常伯”。其“綴衣”、“虎賁”而言“牧”者，以下文自詳，故此惟舉內外要官者言之，故內官舉“任人”、“準夫”，外官舉“牧”。故下云：“繼自今，我立政。立事、準人、牧夫，我其克灼知厥若。”又云：“自古商人，亦越我周\CJKunderline{文王}立政、立事、牧夫、準人，則克宅之，克由繹之，茲乃俾乂。”皆據內外要重官以言之。“夫”即人也，立官所以事天地，治人民，為此三事而已，故以“三事”謂天地人也。\CJKunderline{王肅}云：“\CJKunderline{文王}所以立政，任人，常任也；準夫，準人也；牧者，諸侯之長也。”與孔意同。 \par}

{\noindent\zhuan\zihao{6}\fzbyks 傳“趣馬”至“其人”。正義曰：\CJKunderwave{周禮}趣馬為校人屬官，馬一十二匹，立趣馬一人，“掌贊正良馬，而齊其飲食”，是掌馬之小官也。綴衣是大僕也，虎賁、大僕皆下大夫也。此三公六卿，亦為小尹之官。雖文止三官，亦包通在下之屬官。三官之下小官多矣，趣馬即下士，其馬一匹，有圉師一人，是趣馬之下猶有小官也。 \par}

{\noindent\zhuan\zihao{6}\fzbyks 傳“雖左”至“擇人”。正義曰：諸官有所務從業，從王左右攜持器物之僕，謂寺人、內小臣等也。“百司庶府”，謂百官有司之下，主券契府藏之吏，謂其下賤人，非百官有司之身也。言此等亦皆擇人。 \par}

{\noindent\zhuan\zihao{6}\fzbyks 傳“小臣”至“任乎”。正義曰：“小臣猶皆擇人，況大都邑之小長”，謂公卿，都邑之內大夫士及邑宰之屬，以身有道藝、為民之表的楨幹之臣。其都邑之內屬官，謂之“小長”。\CJKunderwave{周禮·太宰職}云：“乃施則于都鄙，而建其長，立其兩，設其伍，陳其殷。”“兩”謂兩卿，“長”謂公卿，“伍”謂大夫，“殷”謂眾士是也。 \par}

太史、尹伯、庶常吉士，\footnote{太史，下大夫，掌邦六典之貳；尹伯，長官大夫；及旅掌常事之善士，皆得其人。}司徒、司馬、司空、亞旅，\footnote{此有三卿及次卿眾大夫，則是文武未伐紂時。舉文武之初以為法則。}夷微、盧烝、三亳、阪尹。\footnote{蠻夷微、盧之眾帥，及亳人之歸\CJKunderline{文王}者，三所為之立監，及阪地之尹長，皆用賢。○阪音反。}

{\noindent\zhuan\zihao{6}\fzbyks 傳“太史”至“其人”。正義曰：\CJKunderwave{周禮}“太史，下大夫二人”，“掌建邦之六典”。又\CJKunderwave{太宰職}亦云“掌建邦之六典”。太史副貳，太宰掌其正，太史掌其貳。“六典”謂治典、教典、禮典、政典、刑典、事典，六卿所掌者也。“掌邦六典之貳”，其所掌事重，故特言之。“尹伯”，長官大夫。\CJKunderwave{周禮}每官各有長，若太史為史官之長，大司樂為樂官之長，如此類皆是也。“及眾掌常事之善士”，謂士為長官者。其大夫及士不為長官者，則前云“百司”也。居官必須善人，此是總舉眾官，故特言“吉士”。 \par}

{\noindent\zhuan\zihao{6}\fzbyks 傳“此有”至“法則”。正義曰：\CJKunderline{周公}攝政之時，制禮作樂，其作\CJKunderwave{立政}之篇，必在制禮之後。\CJKunderwave{周禮}六卿，而“此有三卿及次卿眾大夫”，則是副卿之大夫，有若\CJKunderwave{周禮}小宰之類是也。此文武未伐紂之時也，遠舉文武之初以為法則爾。\CJKunderwave{泰誓}下篇云“王乃大巡六師”，“六師”則六軍也，軍將皆命卿，即伐紂之時已立六卿矣。\CJKunderwave{牧誓}亦云“司徒、司馬、司空”，舉之三卿者，彼傳已解之云“指誓戰者”也。 \par}

{\noindent\zhuan\zihao{6}\fzbyks 傳“蠻夷”至“用賢”。正義曰：\CJKunderwave{牧誓}所云，有“微、盧、彭、濮人”，此舉“夷微、盧”以見彭、濮之等諸夷也。“烝”訓眾也。此篇所言，皆立官之事，此經惟“阪”下言“尹”，則“夷、微”已下以一“尹”總之，故傳言“蠻夷微、盧之眾帥，及亳民之歸\CJKunderline{文王}者，三所為之立監，及阪地之尹長”。故言“帥”，言“監”，亦是言為之立長，義出經文“尹”也。“亳”是湯之舊都，此言“三亳”，必是亳民分為三處。此篇說立官之意，明是分為三亳，必是三所各為立監也。“亳民之歸\CJKunderline{文王}”，經傳未有其事，\CJKunderline{文王}既未伐紂，亳民不應歸之。鄭、王所說皆與孔同。言亳民歸\CJKunderline{文王}者，蓋以此章雜陳\CJKunderline{文王}、\CJKunderline{武王}時事，其言以\CJKunderline{文王}為主，故先儒因言亳民歸\CJKunderline{文王}爾。即如此意，三亳為已歸周，必是\CJKunderline{武王}時也。“及阪地之尹長”，傳言其山阪之地立長爾，不知其指斥何處也。\CJKunderline{鄭玄}以“三亳阪尹”者共為一事，云：“湯舊都之民服\CJKunderline{文王}者，分為三邑，其長居險，故言阪尹。蓋東成皋，南轘轅，西降谷也。”皇甫謐以為“三亳，三處之地，皆名為亳。蒙為北亳，谷熟為南亳，偃師為西亳”。古書亡滅，既無要證，未知誰得旨矣。 \par}

{\noindent\shu\zihao{5}\fzkt “立政”至“阪尹”。正義曰：言文武亦法\CJKunderline{禹}湯，審官以立美政。“任人”謂六卿。“準夫”者，平法之人，謂理獄官也。“牧”者，九州之牧。治為天地人之三事。自“虎賁”已下，歷舉官名，言此官皆須得其人。不以官之尊卑為次，蓋以從近而至遠。虎賁、綴衣、趣馬，三者官雖小，須慎擇其人。乃至左右攜持器物之僕，及百官有司之下至眾府藏之吏，亦須擇其人。既言近王小官，及遠官大者。小官猶須擇人，況乎大都邑之小長，與有道藝之人為表幹之臣，及百官有司之職，可以非其任乎？以近臣況遠臣，以小官況大官。既以近小況遠大，又舉官之次而掌事要者。若太史下大夫、長官大夫及眾掌常事之善士，皆須得其人。更舉官之大者，司徒、司馬、司空之卿，及次卿之眾大夫，皆須得其人。既略言內外之官，又更遠及夷狄蠻夷微、盧之眾帥，與三處亳民之監，及阪地之尹長，皆須用賢人。言文武於此諸官，皆求賢人為之也。 \par}

\CJKunderline{文王}惟克厥宅心，乃克立茲常事司牧人，以克俊有德。\footnote{文王惟其能居心遠惡舉善，乃能立此常事司牧人，用能俊有德者。○遠，於萬反。}\CJKunderline{文王}罔攸兼於庶言，庶獄庶慎，惟有司之牧夫。\footnote{文王無所兼知於譭譽眾言，及眾刑獄，眾當所慎之事，惟慎擇有司牧夫而已。勞於求才，逸於任賢。○譽音餘，又如字。}是訓用違，庶獄庶慎，\CJKunderline{文王}罔敢知於茲。\footnote{是萬民順法，用違法，眾獄眾慎之事，\CJKunderline{文王}一無敢自知於此，委任賢能而已。}


{\noindent\zhuan\zihao{6}\fzbyks 傳“\CJKunderline{文王}”至“德者”。正義曰：上言\CJKunderline{文王}能知三宅三俊，知此言“能居心”者，以遠惡舉善居其心也。既遠惡舉善，乃能立此常事，用賢養民，是人君之常事也。 \par}

{\noindent\zhuan\zihao{6}\fzbyks 傳“\CJKunderline{文王}”至“任賢”。正義曰：下云“是訓用違”，即是在上“庶言”也。“是訓”則稱譽之事，“用違”則毀損之事，但分析言之爾。 \par}

{\noindent\shu\zihao{5}\fzkt “\CJKunderline{文王}”至“於茲”。正義曰：上既總言文武，此又分而說之。\CJKunderline{文王}惟能其居心遠惡舉善，乃能立此常事其主養人之官，用能俊有德者。既任用俊人，每事委之，\CJKunderline{文王}無所兼知於眾人之言，或毀或譽，\CJKunderline{文王}皆不知也。眾獄斷罪得失，\CJKunderline{文王}亦不得知也。眾所當慎之事，\CJKunderline{文王}亦不得知也。惟慎擇在朝有司，在外牧養民之夫。是時萬民或順於法，或用違法，眾刑獄，眾所慎之事，\CJKunderline{文王}一皆無敢自知於此，惟委任賢能而已。 \par}

亦越\CJKunderline{武王}，率惟敉功，不敢替厥義德，\footnote{亦於\CJKunderline{武王}循惟\CJKunderline{文王}無安天下之功，不敢廢其義德，奉遵父道。○敉，亡婢反。}率惟謀從容德，以並受此丕丕基。\footnote{\CJKunderline{武王}循惟謀從\CJKunderline{文王}寬容之德，故君臣並受此大大之基業，傳之子孫。傳，直專反。}


{\noindent\zhuan\zihao{6}\fzbyks 傳“\CJKunderline{武王}”至“子孫”。正義曰：以言“並受”，則非獨王身，故以為“君臣並受此大大之基業”。謀從寬容之德，是與臣謀，及基業成就，則君臣共有，故言“並受”。且王為天子，臣為諸侯，皆受基業，各傳子孫，是亦為“並受”也。 \par}

{\noindent\shu\zihao{5}\fzkt “亦越”至“丕基”。正義曰：亦於\CJKunderline{武王}遵循父道，所循惟\CJKunderline{文王}撫安天下之功，不敢廢其\CJKunderline{文王}義德。言奉行遵父道也。又言\CJKunderline{武王}遵循者，惟謀從\CJKunderline{文王}寬容之德，故\CJKunderline{武王}君臣能並受此大大之基業。謂受命為天子，傳之子孫。 \par}

“嗚呼！孺子王矣。\footnote{嘆稚子今以為王矣，不可不勤法祖考之德。}繼自今，我其立政、立事、準人、牧夫，我其克灼知厥若，丕乃俾亂。\footnote{繼用今已往，我其立政大臣、立事小臣、及準人、牧夫，我其能灼然知其順者，則大乃使治之。言知臣下之勤勞，然後莫不盡其力。○俾，必爾反，下同。治，直吏反，下同。}相我受民，和我庶獄庶慎，時則勿有間之。\footnote{能治我所受天民，和平我眾獄眾慎之事，如是則勿有以代之。言不可復變。○相如字，馬息亮反，下“勱相”同。間，間廁之間。復，扶又反。}自一話一言，我則末惟成德之彥，以乂我受民。\footnote{言政當用一善，善在一言而已。欲其口無擇言。如此我則終惟有成德之美，以治我所受之民。○話，戶怪反。}


{\noindent\zhuan\zihao{6}\fzbyks 傳“繼用”至“心力”。正義曰：自此已下四言“繼自今”者，凡人靡不有初，鮮克有終，恐王不能終之，戒\CJKunderline{成王}使繼續，從今已往常用賢也。“自”訓為從，亦訓為用，此傳言“用今已往”，下傳言“從今已往”，其意同也。“政”、“事”相對，則“政”大“事”小，故以“立政”為大臣，“立事”為小臣。及“準人”、“牧夫”，略舉四者以總諸臣，戒王任此人也。其能灼然知其能順於事者，則大乃使治。顧氏云：“君能知臣下順於事，則臣感君恩，大乃治理，各盡心力也。” \par}

{\noindent\zhuan\zihao{6}\fzbyks 傳“能治”至“復變”。正義曰：“相”訓助也，助君所以治民事，故“相”為治。天命王者，使之治民，則天與王者此民，故言“能治我所受天民”也。能治下民，理眾獄眾慎之事，使得其所，則為政之大要，能如此，則勿有以代之。言此法盡善，不可復變易也。或據臣身既能如此，不可以餘人代之也。 \par}

{\noindent\zhuan\zihao{6}\fzbyks 傳“言政”至“之民”。正義曰：\CJKunderwave{釋詁}云：“自，用也。話，言也。”舍人曰：“話，政之善言也。”孫炎曰：“話善之言也。”然則“話”之與“言”是一物也。“自一話”者，言人君為政，當用純一善言。又云“一言”者，純一善言，在於一言而已。謂發號施令，當須純一，不得差貳，欲令其口無可擇之言也。顧氏云：“人君為政之道，當須用一善而已，為善之法，惟在一言也。‘末’訓為終，‘彥’訓為美，王能出言皆善，口無可擇，如此我王則終惟有成德之美，以治我所受天民矣。”\CJKunderwave{釋訓}云：“美士為彥。”故“彥”為美。 \par}

{\noindent\shu\zihao{5}\fzkt “嗚呼孺子”至“受民”。正義曰：\CJKunderline{周公}既歷說\CJKunderline{禹}湯文武，乃復指戒\CJKunderline{成王}，“嗚呼”而嘆，孺子今已為王矣。既正位為王，事不可不慎。繼續從今已往，我王其與立政，謂大臣也，其與立事，謂小臣也，平法之人及養民之夫，此等諸臣，我王其能察之灼然，知其順於事者，則大乃使之治理。言知其能有勤勞，各盡心力。然後用此賢臣治我所受天民，和平我眾獄訟，及眾當所慎之事，必能如是，則勿復有以代之。言其法不可復變也。政從君出，為人主用是一善之言，善在一言而已，勿以惡言亂之。王能如是，我王則終惟有成德之美，以治我所受天民矣。 \par}

嗚呼！予旦已受人之徽言，咸告孺子王矣。\footnote{嘆所受賢聖說\CJKunderline{禹}湯之美言，皆以告稚子王矣。○稚，直吏反，本亦作稺。}繼自今，文子文孫，其勿誤於庶獄庶慎，惟正是乂之。\footnote{文子文孫，\CJKunderline{文王}之子孫。從今已往，惟以正是之道治眾獄眾慎，其勿誤。}自古商人，亦越我周\CJKunderline{文王}立政、立事、牧夫、準人，則克宅之，克由繹之，茲乃俾乂。\footnote{言用古商湯，亦於我周\CJKunderline{文王}立政立事，用賢人之法，能居之於心，能用陳之，此乃使天下治。○繹音亦。}


{\noindent\zhuan\zihao{6}\fzbyks 傳“言用”至“下治”。正義曰：上陳\CJKunderline{禹}湯文武，此覆上文，惟言湯與\CJKunderline{文王}者，言有詳略，無別意也。“能居之於心”，謂心知其賢也。“能用陳之”，謂陳列於位，用之以為官也。\CJKunderline{王肅}曰：“則能居之在位，能用陳其才力，如此故能使天下治也。” \par}

{\noindent\shu\zihao{5}\fzkt “嗚呼”至“俾乂”。正義曰：“旦”者，\CJKunderline{周公}名也。\CJKunderline{周公}又嘆曰：“嗚呼！我旦已受賢聖人說\CJKunderline{禹}湯之美言，皆以告孺子王矣，王宜依行之。繼續從今以往，\CJKunderline{文王}之子孫，其勿得過誤於眾獄訟眾所慎之事，惟當用是正是之道治之。用古商人\CJKunderline{成湯}，亦於我周家\CJKunderline{文王}其立政、立事、牧夫、準人此等諸官，皆用賢人之法，則能居之於心，能用陳之於位，明識賢人，用之為官，此乃使天下大治。”戒\CJKunderline{成王}使法之。 \par}

國則罔有立政用憸人，不訓於德，是罔顯在厥世。\footnote{商周賢聖之國，則無有立政用憸利之人者。憸人不訓於德，是使其君無顯名在其世。○憸，息廉反，徐七漸反，本又作 ，馬云：“憸利，佞人也。”}繼自今立政，其勿以憸人，其惟吉士，用\xpinyin*{勱}相我國家。\footnote{立政之臣，惟其吉士，用勉治我國家。○勱音邁。}

{\noindent\shu\zihao{5}\fzkt “國則”至“國家”。正義曰：既言湯與\CJKunderline{文王}用賢大治，又言其不宜用小人。商周聖賢之國，無有立政用憸利之人者。此憸利之人不順於德，若其用之，是使其君無顯名在其世也。王當繼續從今已往立其善政，其勿用憸利之人，其惟任用善士，使勉力治我國家。教王使用善士，勿使小人也。 \par}

今文子文孫，孺子王矣。\footnote{告\CJKunderline{文王}之子孫，言稚子已即政為王矣，所以厚戒。}其勿誤於庶獄，惟有司之牧夫。\footnote{獨言眾獄、有司，欲其重刑，慎官人。}其克詰爾戎兵，以陟\CJKunderline{禹}之跡，\footnote{其當能治汝戎服兵器，威懷並設，以升\CJKunderline{禹}治水之舊跡。○詰，起一反，馬云：“賓也。”}方行天下，至於海表，罔有不服。\footnote{方，四方。海表，蠻夷戎狄，無不服化者。}以覲\CJKunderline{文王}之耿光，以揚\CJKunderline{武王}之大烈。\footnote{能使四夷賓服，所以見祖之光明，揚父之大業。}嗚呼！繼自今，后王立政，其惟克用常人。”\footnote{其惟能用賢才為常人，不可以天官有所私。}


{\noindent\zhuan\zihao{6}\fzbyks 傳“獨言”至“官人”。正義曰：上有“庶慎”、“立政”、“立事”、“牧夫”、“準人”，此獨言“庶獄”與“有司之牧夫”者，言“庶獄”欲其重刑，言“有司牧夫”欲其慎官人也。 \par}

{\noindent\zhuan\zihao{6}\fzbyks 傳“其當”至“舊跡”。正義曰：立官所以牧養下民，戒備不虞，故以“詰爾戎兵”為言也。“戎”亦“兵”也，以其並言“戎兵”，故傳以為“戎服兵器，威懷並設，以升\CJKunderline{禹}治水之舊跡”。遠行必登山，故以“陟”言之。如舜之“陟方”，意亦然。 \par}

{\noindent\zhuan\zihao{6}\fzbyks 傳“方四”至“化者”。正義曰：“方行天下”，言無所不至，故以“方”為四方。\CJKunderwave{釋地}云：“九夷、八狄、七戎、六蠻謂之四海。”知“海表”謂“夷狄戎蠻，無有不服化者”。即\CJKunderwave{詩·小雅}云“\CJKunderwave{蓼蕭}，澤及四海”是也。 \par}

{\noindent\zhuan\zihao{6}\fzbyks 傳“其惟”至“所私”。正義曰：官須常得賢人，故惟賢是用。用賢是常，常則非賢不可。人主或知其不賢，以私受用之，代天為官，故言“不可以天官有所私”。 \par}

{\noindent\shu\zihao{5}\fzkt “今文”至“常人”。正義曰：今告汝\CJKunderline{文王}之子，\CJKunderline{文王}之孫，孺子今已即政為王矣，我所以須厚戒之。王其勿誤於眾治獄之官，當須慎刑也。惟有司之牧夫，有司主養民者，宜得賢也。治獄之吏，養民之官，若任得其人，使其能治汝戎服兵器，以此升行\CJKunderline{禹}之舊跡，四方而行，至於天下，至於四海之表，無有不服王之化者，以顯見\CJKunderline{文王}之光明，以播揚\CJKunderline{武王}之大業。言任得賢臣，則光揚父祖。\CJKunderline{周公}又嘆曰：“嗚呼！繼續從今已往，後世之王，立行善政，其惟能用常人，必使常得賢人，不可任非其才，此雖有戒\CJKunderline{成王}，乃是國之常法，因以戒后王，言此法可常行也。 \par}

\CJKunderline{周公}若曰：“太史，\footnote{順其事並告太史。}司寇\CJKunderline{蘇公}，式敬爾由獄，以長我王國。\footnote{\CJKunderline{忿生}為\CJKunderline{武王}司寇，封蘇國，能用法。敬汝所用之獄，以長施行於我王國。言主獄當求\CJKunderline{蘇公}之比。○比，必二反，又如字。}茲式有慎，以列用中罰。”\footnote{此法有所慎行，必以其列用中罰，不輕不重，\CJKunderline{蘇公}所行。太史掌六典，有廢置官人之制，故告之。○行如字。}


{\noindent\zhuan\zihao{6}\fzbyks 傳“\CJKunderline{忿生}”至“之比”。正義曰：成十一年\CJKunderwave{左傳}云：“昔周克商，使諸侯撫封，\CJKunderline{蘇忿生}以溫為司寇。”是“\CJKunderline{忿生}為\CJKunderline{武王}司寇，封蘇國”也。“蘇”是國名，所都之地其邑名“溫”，故傳言“以溫”也。特舉\CJKunderline{蘇公}治獄官以告太史，知其言主獄之官當求\CJKunderline{蘇公}之比類也。 \par}

{\noindent\zhuan\zihao{6}\fzbyks 傳“此法”至“告之”。正義曰：治獄必有定法，此定法有所慎行。\CJKunderwave{周禮·大司寇}云：“刑新國用輕典，刑平國用中典，刑亂國用重典。”輕重各有體式行列，\CJKunderline{周公}言然之時，是法為平國，故必以其列用中罰，使不輕不重。美\CJKunderline{蘇公}治獄，使列用中罰，明中罰不輕不重，是\CJKunderline{蘇公}所行也。\CJKunderwave{周禮}太宰“以八柄詔王，馭群臣”，有爵、祿、廢、置、生、殺、與、奪之法。太史亦掌邦之六典，以副貳太宰，是太史有廢置官人之制，故特呼而告之也。 \par}

{\noindent\shu\zihao{5}\fzkt “\CJKunderline{周公}”至“中罰”。正義曰：\CJKunderline{周公}順其事而言曰：“太史。”以其太史掌廢置官人，故呼而告之：“昔日司寇\CJKunderline{蘇公}既能用法，汝太史當敬汝所用之獄，以長施行於我王國。”欲使太史選主獄之官，當求\CJKunderline{蘇公}之比也。“此刑獄之法，有所慎行，必以其體式，列用中常之罰，不輕不重，當如\CJKunderline{蘇公}所行也”。 \par}

%%% Local Variables:
%%% mode: latex
%%% TeX-engine: xetex
%%% TeX-master: "../Main"
%%% End:

%% -*- coding: utf-8 -*-
%% Time-stamp: <Chen Wang: 2024-04-02 11:42:41>

% {\noindent \zhu \zihao{5} \fzbyks } -> 注 (△ ○)
% {\noindent \shu \zihao{5} \fzkt } -> 疏

\chapter{卷十八}


\section{周官第二十二【偽】}


\CJKunderline{成王}既黜殷命，滅淮夷，\footnote{黜殷在\CJKunderline{周公}東征時，滅淮夷在\CJKunderline{成王}即政後，事相因，故連言之。}還歸在豐，作\CJKunderwave{周官}。\footnote{\CJKunderline{成王}雖作洛邑，猶還西周。}

周官\footnote{言周家設官分職用人之法。}


{\noindent\zhuan\zihao{6}\fzbyks 傳“黜殷”至“言之”。正義曰：據\CJKunderwave{金縢}之經、\CJKunderwave{大誥}之序，知“黜殷命”在\CJKunderline{周公}攝政三年東征之時也。據\CJKunderwave{成王政}之序、\CJKunderwave{費誓}之經，知“滅淮夷”在\CJKunderline{成王}即政之後也。淮夷於攝政之時與\CJKunderline{武庚}同叛，\CJKunderline{成王}既滅淮夷，天下始定。淮夷本因\CJKunderline{武庚}而叛，黜殷命與滅淮夷其事相因，故雖則異年而連言之，以見天下既定，乃作\CJKunderwave{周官}故也。下經言“四徵弗庭”是黜滅之事也，“罔不承德”是安寧之狀也，序顧經文，故追言“黜殷命”，以接“滅淮夷”，見征伐乃安定之意也。 \par}

{\noindent\zhuan\zihao{6}\fzbyks 傳“\CJKunderline{成王}”至“西周”。正義曰：以\CJKunderwave{洛誥}之文言“王在新邑”，今復云“在豐”，故解之也。\CJKunderwave{史記·周本紀}云：“太史公曰：“學者皆稱周伐紂，居洛邑，綜其實不然。\CJKunderline{武王}營之，\CJKunderline{成王}使\CJKunderline{召公}卜，居九鼎焉，而周復都豐、鎬。”是言\CJKunderline{成王}雖作洛邑，猶還西周之事也。\CJKunderwave{多方}云：“王來自奄，至於宗周。”宗周即鎬京也，於彼不解，至此始為傳者，宗周雖是鎬京，文無“豐鎬”之字，故就此解之。\CJKunderline{武王}既以遷鎬京，今王覆在豐者，豐、鎬相近，舊都不毀，豐有\CJKunderline{文王}之廟，故事就豐宣之故也。 \par}

{\noindent\zhuan\zihao{6}\fzbyks 傳“言周”至“之法”。正義曰：\CJKunderwave{周禮}每官言人之員數及職所掌，立其定法，授與\CJKunderline{成王}。\CJKunderline{成王}即政之初，即有淮夷叛逆，未暇得以立官之意號令群臣。今既滅淮夷，天下清泰，故以周家設官分職用人之法以誥群臣，使知立官之大旨也。“設官分職”，\CJKunderwave{周禮}序官之文，言設置群官，分其職掌。經言立三公六卿，是“設官”也。各言所掌，是“分職”也。各舉其官之所掌，示以才堪乃得居之，是說“用人之法”。 \par}

{\noindent\shu\zihao{5}\fzkt “\CJKunderline{成王}”至“周官”。正義曰：\CJKunderline{成王}於\CJKunderline{周公}攝政之時既黜殷命，及其即位之後滅淮夷，於是天下大定。自滅淮夷，還歸在豐，號令群臣，言周家設官分職用人之法。史敘其事，作\CJKunderwave{周官}。 \par}

惟周王撫萬邦，巡侯甸，\footnote{即政撫萬國，巡行天下侯服、甸服。}四徵弗庭，綏厥兆民。\footnote{四面征討諸侯之不直者，所以安其兆民。十億曰兆，言多。}六服群辟，罔不承德。歸於宗周，董正治官。\footnote{六服諸侯，奉承周德。言協服。還歸於豐，督正治理職司之百官。○闢，必亦反。治，直吏反，下至“冢宰”經注同。}


{\noindent\zhuan\zihao{6}\fzbyks 傳“即政”至“甸服”。正義曰：檢\CJKunderwave{成王政}之序與\CJKunderwave{費誓}之經，知\CJKunderline{成王}即政之年，奄與淮夷又叛，叛即往伐，今始還歸。\CJKunderwave{多方}云：“五月丁亥，王來自奄，至於宗周。”與此滅淮夷而還歸在豐為一事也。年初始叛，五月即歸，其間未得巡守於四方也。而此言“撫萬國，巡行天下”，其實止得撫巡向淮夷之道所過諸侯爾，未是用四仲之月大巡守也。以撫諸侯巡守是天子之大事，因即大言之爾。周之法制無萬國也，惟伐淮夷非四徵也，言“萬國”、“四徵”亦是大言之爾。六服而惟言“侯甸”者，二服去圻最近，舉近以言之，言王巡省遍六服也。 \par}

{\noindent\zhuan\zihao{6}\fzbyks 傳“四面”至“言多”。正義曰：“四徵”，從京師而四面徵也。\CJKunderwave{釋詁}云：“庭，直也。綏，安也。”諸侯不直，謂叛逆王命，侵削下民，故“四面征討諸侯之不直者，所以安其兆民”。\CJKunderwave{楚語}云：“十日、百姓、千品、萬官、億醜、兆民。”每數相十，知“十億曰兆”。稱“兆”，言其多也。 \par}

{\noindent\zhuan\zihao{6}\fzbyks 傳“六服”至“百官”。正義曰：\CJKunderwave{周禮}“九服”，此惟言“六”者，夷、鎮、蕃三服在九州之外夷狄之地，王者之於夷狄，羈縻而已，不可同於華夏，故惟舉六服諸侯。奉承周德，言協服也。序云“還歸在豐”，知宗周即豐也。周為天下所宗，王都所在皆得稱之，故豐、鎬與洛邑皆名“宗周”。\CJKunderwave{釋詁}云：“董、督，正也。”是“董”得為督，督正治理職司之百官。下戒敕是“董正”也。 \par}

{\noindent\shu\zihao{5}\fzkt “惟周”至“治官”。正義曰：惟周之王者，佈政教，撫安萬國，巡行天下侯服、甸服，四面征討諸侯之不直者，所以安其海內兆民。六服之內崑崙諸侯之君，無有不奉承周王之德者。自滅淮夷而歸於宗周豐邑，乃督正治理職司之百官。敘王發言之端也。 \par}

王曰：“若昔大猷，制治於未亂，保邦於未危。”\footnote{言當順古大道，制治安國，必於未亂未危之前，思患預防之。}

{\noindent\shu\zihao{5}\fzkt “王曰”至“未危”。正義曰：“治”謂政教，“邦”謂國家。治有失則亂，家不安則危。恐其亂則預為之制，慮其危則謀之使安，制其治於未亂之始，安其國於未危之前。張官設府，使分職明察，任賢委能，令事務順理，如是則政治而國安矣。標此二句於前，以示立官之意。必於未亂未危之前為之者，思患而預防之。“思患而預防之”，\CJKunderwave{易·既濟卦}象辭也。 \par}

曰：“\CJKunderline{唐}、\CJKunderline{虞}稽古，建官惟百。內有百揆四岳，外有州牧侯伯。\footnote{道堯舜考古以建百官，內置百揆四岳，象天之有五行，外置州牧十二及五國之長，上下相維，外內咸治。言有法。○長，丁丈反，下“官長”、“助長”並同。}庶政惟和，萬國咸寧。\footnote{官職有序，故眾政惟和，萬國皆安，所以為正治。}夏商官倍，亦克用乂。\footnote{禹湯建官二百，亦能用治。言不及\CJKunderline{唐}、\CJKunderline{虞}之清要。}明王立政，不惟其官，惟其人。\footnote{言聖帝明王立政修教，不惟多其官，惟在得其人。}


{\noindent\zhuan\zihao{6}\fzbyks 傳“道堯”至“有法”。正義曰：百人無主，不散則亂，有父則有君也。君不獨治，必須輔佐，有君則有臣也。\CJKunderwave{易·序卦}云：“有父子然後有君臣。”則君臣之興，次父子之後，人民之始，則當有之，未知其所由來也。雖遠舉\CJKunderline{唐}、\CJKunderline{虞}，復考古也。\CJKunderwave{說命}曰：“明王奉若天道，建邦設都。”則王者立官，皆象天為之，故“內置百揆四岳，象天之有五行”也。五行佐天，群臣佐主，以此為象天爾。不必其數有五乃象五行，故以“百揆四岳”為五行之象。\CJKunderwave{左傳}少昊立五鳩氏，顓頊已來立五行之官，其數亦有五，故置於五行矣。\CJKunderwave{舜典}云“肇十有二州”，此說虞事，知“置州牧十二”也。“侯伯”謂諸侯之長，\CJKunderwave{益稷}篇\CJKunderline{禹}言治水時事云“外薄四海，咸建五長”，知“侯伯”是五國之長也。\CJKunderline{成王}說此事者，言堯舜所制，上下相維，內外咸治，言有法也。此言“建官惟百”，“夏商官倍”，則\CJKunderline{唐}、\CJKunderline{虞}一百，夏商二百。\CJKunderwave{禮記·明堂位}云“有虞氏官五十，夏后氏官百”者，\CJKunderwave{禮記}是後世之言，不與經典合也。 \par}

{\noindent\shu\zihao{5}\fzkt “曰\CJKunderline{唐}、\CJKunderline{虞}”至“其人”。正義曰：既言須立官之意，乃追述前代之法。止而復言，故更加一“曰”。唐堯\CJKunderline{虞舜}考行古道，立官惟數止一百也。“內有百揆四岳”者，“百揆”，揆度百事，為群官之首，立一人也；“四岳”，內典四時之政，外主太嶽之事，立四人也。“外有州牧侯伯”，牧，一州之長；侯伯，五國之長，各監其所部之國。外內置官，各有所掌，眾政惟以協和，萬邦所以皆安也。夏禹商湯立官倍多於\CJKunderline{唐}、\CJKunderline{虞}，雖不及\CJKunderline{唐}、\CJKunderline{虞}之清簡，亦能用以為治。明王立其政教，不惟多其官，惟在得其人。言自古制法皆明開官司，求賢以處之也。 \par}

今予小子，祗勤於德，夙夜不逮。\footnote{今我小子敬勤於德，雖夙夜匪懈，不能及古人。言自有極。○逮音代，一音大計反。懈，隹賣反。}仰惟前代時若，訓迪厥官。\footnote{言仰惟先代之法是順，訓蹈其所建官而則之，不敢自同堯舜之官，準擬夏殷而蹈之。}立太師、太傅、太保，茲惟三公。論道經邦，燮理陰陽。\footnote{師，天子所師法；傅，傅相天子；保，保安天子於德義者，此惟三公之任。佐王論道，以經緯國事，和理陰陽。言有德乃堪之。○燮，素協反。相，息亮反。}官不必備，惟其人。\footnote{三公之官不必備員，惟其人有德乃處之。○處，昌呂反。}少師、少傅、少保，曰三孤。\footnote{此三官名曰三孤。孤，特也。言卑於公，尊於卿，特置此三者。○少，詩照反，下同。}貳公弘化，寅亮天地，弼予一人。\footnote{副貳三公，弘大道化，敬信天地之教，以輔我一人之治。}冢宰掌邦治，統百官，均四海。\footnote{\CJKunderwave{天官}卿稱太宰，主國政治，統理百官，均平四海之內邦國。言任大。}司徒掌邦教，敷五典，擾兆民。\footnote{\CJKunderwave{地官}卿，司徒主國教化，布五常之教，以安和天下眾民，使小大皆協睦。○擾，而小反，徐音饒。}


{\noindent\zhuan\zihao{6}\fzbyks 傳“師天”至“堪之”。正義曰：三公俱是教道天子，輔相天子，緣其事而為之名。三公皆當運致天子，使歸於德義。傳於“保”下言“保安天子於德義”，總上三者，言皆然也。\CJKunderwave{禮記·文王世子}云：“師也”者，教之以事而喻諸德者也。保也者，慎其身以輔翼之而歸諸道者也。”道、德別掌者，內得於心，出行於道，道德不甚相遠，因其並釋“師”、“保”，故分配之爾。於公云“燮理陰陽”，於孤云“寅亮天地”，“和理”、“敬信”義亦同爾。以孤副貳三公，故其事所掌不異。 \par}

{\noindent\zhuan\zihao{6}\fzbyks 傳“天官”至“任大”。正義曰：此經言六卿所掌之事，撮引\CJKunderwave{周禮}為之總目，或據\CJKunderwave{禮}文，或取\CJKunderwave{禮}意，雖言有小異，義皆不殊。\CJKunderwave{周禮}云：“乃立天官冢宰，使帥其屬而掌邦治。治官之屬，太宰卿一人。”\CJKunderline{馬融}云：“冢，大也。宰，治也。大治者，兼萬事之名也。”\CJKunderline{鄭玄}云：“變冢言大，進退異名也。百官總焉則謂之冢，列職於王則稱大。冢者，大之上也。山頂曰冢。”是解“冢”、“大”異名之意。\CJKunderwave{太宰職}云：“三曰禮典，以統百官。”\CJKunderline{馬融}云：“統，本也。百官是宗伯之事也。”此“統百官”在“冢宰”之下，當以冢尊，故命統治百官為冢宰之事，治官、禮官俱得統之也。\CJKunderwave{禮}云“以佐王均邦國”，此言“均四海”，故傳辨之“均平四海之內邦國”，與孔意不異。 \par}

{\noindent\zhuan\zihao{6}\fzbyks 傳“地官”至“協睦”。正義曰：\CJKunderwave{周禮}云：“乃立地官司徒，使帥其屬而掌邦教，以佐王安擾邦國。”\CJKunderwave{太宰職}云：“二曰教典，以擾萬民。”\CJKunderline{鄭玄}云：“擾亦安也，言饒衍之。”傳亦以“擾”為安。\CJKunderwave{五典}即五教也，布五常之教，以安和天下之人民，使小大協睦也。\CJKunderwave{舜典}云：“契為司徒，敬敷五教。”\CJKunderwave{周禮}：“司徒掌十有二教。一曰以祀禮教敬，則民不苟。二曰以陽禮教讓，則民不爭。三曰以陰禮教親，則民不怨。四曰以樂禮教和，則民不乖。五曰以儀辨等，則民不越。六曰以俗教安，則民不偷。七曰以刑教中，則民不暴。八曰以誓教恤，則民不怠。九曰以度教節，則民知足。十曰以世事教能，則民不失職。十有一曰以賢制爵，則民慎德。十有二曰以庸制祿，則民興功。”\CJKunderline{鄭玄}云：“有虞氏五而周十有二焉，然則十有二細分五教。為之五教，可以常行，謂之五典。五典謂父義、母慈、兄友、弟恭、子孝也。” \par}

宗伯掌邦禮，治神人，和上下。\footnote{\CJKunderwave{春官}卿，宗廟官長，主國禮，治天地神祗人鬼之事，及國之吉、兇、賓、軍、嘉五禮，以和上下尊卑等列。}司馬掌邦政，統六師，平邦國。\footnote{\CJKunderwave{夏官}卿，主戎馬之事，掌國征伐，統正六軍，平治王邦四方國之亂者。}司寇掌邦禁，詰奸慝，刑暴亂。\footnote{\CJKunderwave{秋官}卿，主寇賊法禁，治奸惡，刑強暴作亂者。夏司馬討惡助長物，秋司寇刑奸順時殺。}司空掌邦土，居四民，時地利。\footnote{\CJKunderwave{冬官}卿，主國空土以居民，士農工商四人。使順天時，分地利，授之土。能吐生百穀，故曰土。}六卿分職，各率其屬，以倡九牧，阜成兆民。\footnote{六卿各率其屬官大夫士，治其所分之職，以倡道九州牧伯為政，大成兆民之性命，皆能其官，則政治。○倡，尺亮反，下同。阜音負。治，直吏反。}

{\noindent\zhuan\zihao{6}\fzbyks 傳“春官”至“等列”。正義曰：\CJKunderwave{周禮}云：“乃立春官宗伯，使帥其屬而掌邦禮，以佐王和邦國宗廟也。”“伯”，長也。宗廟官之長，故名其官為“宗伯”。其職云：“掌建邦之天神、人鬼、地祗之禮。”又主吉、兇、賓、軍、嘉之五禮。吉禮之別十有二，凶禮之別有五，賓禮之別有八，軍禮之別有五，嘉禮之別有六，總有三十六禮，皆在宗伯職掌之文，文煩不可具載。\CJKunderwave{太宰職}云：“三曰禮典，以和邦國，以諧萬民。”其職又有“以玉作六瑞，以等邦國。以禽作六贄，以等諸臣”，是“以和上下尊卑等列”也。 \par}

{\noindent\zhuan\zihao{6}\fzbyks 傳“夏官”至“亂者”。正義曰：\CJKunderwave{周禮}云：“乃立夏官司馬，使帥其屬而掌邦政，以佐王平邦國。”其職主戎馬之事，有掌征伐，統正六軍，平治王邦四方國之亂者。天子六軍，軍師之通名也。案其職“掌九伐之法，馮弱犯寡則眚之，賊賢害民則伐之，暴內陵外則壇之，野荒民散則削之，負固不服則侵之，賊殺其親則正之，放弒其君則殘之，犯令陵政則杜之，外內亂、鳥獸行則滅之”。 \par}

{\noindent\zhuan\zihao{6}\fzbyks 傳“秋官”至“時殺”。正義曰：\CJKunderwave{周禮}云：“乃立秋官司寇，使帥其屬而掌邦禁，以佐王刑邦國。”其職云：“刑邦國，詰四方。”\CJKunderline{馬融}云：“詰猶窮也，窮四方之奸也。”孔以“詰”為治，是主寇賊法禁，治奸慝之人，刑殺其強暴作亂者。夏官主征伐，秋官主刑殺，征伐亦殺人而官屬異時者，夏司馬討惡助夏時之長物，秋司寇刑奸順秋時之殺物也。\CJKunderwave{周禮}云“掌邦刑”，此云“掌邦禁”者，避下“刑暴亂”之文，故云“掌邦禁”。 \par}

{\noindent\zhuan\zihao{6}\fzbyks 傳“冬官”至“曰土”。正義曰：\CJKunderwave{周禮·冬官}亡。\CJKunderwave{小宰職}云：“六曰冬官，掌邦事。“又云：“六曰事職，以富國，以養萬民。”\CJKunderline{馬融}云：“事職掌百工、器用、耒耜、弓車之屬。”與此主土居民全不相當。\CJKunderwave{冬官}既亡，不知其本。\CJKunderwave{禮記·王制}記司空之事云：“量地以制邑，度地以居民。”足明\CJKunderwave{冬官}本有主土居民之事也。\CJKunderwave{齊語}云：“管仲製法，令士農工商四民不雜。”即此“居民使順天時，分地利，授之土”也。土則地利為之名，以其吐生百穀，故曰“土”也。\CJKunderwave{周禮}云“事”，此云“土”者，為下有“居四民”，故云“土”以居民，為急故也。 \par}

{\noindent\shu\zihao{5}\fzkt “今予”至“厥官”。正義曰：王言：“今我小子，敬勤於德，雖早夜不懈怠，猶不能及於\CJKunderline{唐}、\CJKunderline{虞}。仰惟先代夏商之法是順，順蹈其前代建官而法則之。”言不敢同堯舜之官，準擬行夏殷之官爾。“若”與“訓”俱訓為順也。 \par}

六年，五服一朝。\footnote{五服，侯、甸、男、採、衛。六年一朝會京師。}又六年，王乃時巡，考制度於四岳。\footnote{周制十二年一巡守，春東、夏南、秋西、冬北，故曰時巡。考正制度禮法於四岳之下，如虞帝巡守然。}諸侯各朝於方岳，大明黜陟。”\footnote{覲四方諸侯，各朝於方岳之下，大明考績黜陟之法。}


{\noindent\zhuan\zihao{6}\fzbyks 傳“周制”至“守然”。正義曰：\CJKunderwave{周禮·大行人}云“十有二歲王巡守殷國”，是“周制十二年一巡守”也。如\CJKunderwave{舜典}所云春東、夏南、秋西、冬北以四時巡行，故曰“時巡”。“考正制度禮法於四岳之下，如虞帝巡守然”，據\CJKunderwave{舜典}，“同律度量衡”已下皆是也。 \par}

{\noindent\shu\zihao{5}\fzkt “六年”至“黜陟”。正義曰：此篇說六卿職掌，皆與\CJKunderwave{周禮}符同，則“六年，五服一朝”亦應是\CJKunderwave{周禮}之法，而\CJKunderwave{周禮}無此法也。\CJKunderwave{周禮·大行人}云：“侯服歲一見，其貢祀物。甸服二歲一見，其貢嬪物。男服三歲一見，其貢器物。採服四歲一見，其貢服物。衛服五歲一見，其貢材物。要服六歲一見，其貢貨物。”先儒說\CJKunderwave{周禮}者，皆云“見”謂來朝也。必如所言，則周之諸侯各以服數來朝，無六年一朝之事。昭十三年\CJKunderwave{左傳}叔向云：“明王之制，使諸侯歲聘以志業，間朝以講禮，再朝而會以示威，再會而盟以顯昭明。自古已來，未之或失也。存亡之道，恆由是興。”說\CJKunderwave{左傳}者以為三年一朝、六年一會、十二年而盟，事與\CJKunderwave{周禮}不同。謂之前代明王之法，先儒未嘗措意，不知異之所由。計彼六年一會，與此“六年，五服一朝”事相當也。再會而盟，與此十二年“王乃時巡”，諸侯各朝於方岳亦相當也。叔向盛陳此法，以懼齊人使盟，若周無此，禮叔向妄說，齊人當以辭拒之，何所畏懼而敬以從命乎？且云“自古以來，未之或失”，則當時猶尚行之，不得為前代之法，脅當時之人明矣。明周有此法，\CJKunderwave{禮}文不具爾。\CJKunderwave{大行人}所云，見者皆言貢物，或可因貢而見，何必見者皆是君自朝乎？遣使貢物亦應可矣。\CJKunderwave{大宗伯}云：“時見曰會，殷見曰同。”“時見”、“殷見”不雲年限，“時見曰會”何必不是“再朝而會”乎？“殷見曰同”何必不是“再會而盟”乎？\CJKunderline{周公}制禮若無此法，豈\CJKunderline{成王}謬言，叔向妄說也？計六年大集，應六服俱來，而此文惟言“五服”。孔以五服為侯、甸、男、採、衛，蓋以要服路遠，外逼四夷，不必常能及期，故寬言之而不數也。 \par}

王曰：“嗚呼！凡我有官君子，欽乃攸司，慎乃出令，令出惟行，弗惟反。\footnote{有官君子，大夫以上。嘆而戒之，使敬汝所司，慎汝出令，從政之本。令出必惟行之，不惟反改。若二三其令，亂之道。}以公滅私，民其允懷。\footnote{從政以公平滅私情，則民其信歸之。}學古入官，議事以制，政乃不迷。\footnote{言當先學古訓，然後入官治政。凡制事必以古義議度終始，政乃不迷錯。○度，待洛反。}其爾典常作之師，無以利口亂厥官。\footnote{其汝為政當以舊典常故事為師法，無以利口辯佞亂其官。}


{\noindent\zhuan\zihao{6}\fzbyks 傳“有官”至“之道”。正義曰：教之出令，使之號令在下，則是尊官，故知“有官君子”是“大夫已上”也，下云“三事暨大夫”是也。安危在於出令，故慎汝出令是從政之本也。令既出口，必須行之，令而不行，是去而更反，故謂之“反”也。“不惟反者”，令其必行之，勿使反也。若前令不行而倒反，別出後令以改前令，二三其政，則在下不知所從，是亂之道也。 \par}

{\noindent\zhuan\zihao{6}\fzbyks 傳“言當”至“迷錯”。正義曰：襄三十一年\CJKunderwave{左傳}子產云：“我聞學而後入政，未聞以政學者也。”言將欲入政，先學古之訓典，觀古之成敗，擇善而從之，然後可以入官治政矣。凡欲制斷當今之事，必以古之義理議論量度其終始，合於古義，然後行之。則其為之政教，乃不迷錯也。 \par}

{\noindent\shu\zihao{5}\fzkt “王曰”至“厥官”。正義曰：王言而嘆曰：“嗚呼！凡我有官君子。”謂大夫已上有職事者。“汝等皆敬汝所主之職事，慎汝所出之號令。令出於口，惟即行之，不惟反之而不用，是去而復反也。為政之法，以公平之心滅己之私慾，則見下民其信汝而歸汝矣。學古之典訓，然後入官治政。論議時事，必以古之制度，如此則政教乃不迷錯矣。其汝為政，當以舊典常故事作師法，無以利口辯佞亂其官”。教之以居官為政之法也。 \par}

蓄疑敗謀，怠忽荒政，不學牆面，蒞事惟煩。\footnote{積疑不決，必敗其謀。怠惰忽略，必亂其政。人而不學，其猶正牆面而立，臨政事必煩。○蓄，敕六反。蒞音利，又音類。}戒爾卿士，功崇惟志，業廣惟勤，惟克果斷，乃罔後艱。\footnote{此戒凡有官位，但言卿士，舉其掌事者。功高由志，業廣由勤，惟能果斷行事，乃無後難。言多疑必致患。○斷，丁亂反，下注同。}

{\noindent\shu\zihao{5}\fzkt “蓄疑”至“後艱”。正義曰：又戒群臣，使彊於割斷，勤於職事。蓄積疑惑，不能彊斷，則必敗其謀慮。怠惰忽略，不能恪勤，則荒廢政事。人而不學，如面向牆，無所睹見，以此臨事，則惟煩亂，不能治理。戒汝卿之有事者，功之高者惟志意強正，業之大者惟勤力在公，惟能果敢決斷，乃無有後日艱難。言多疑必將致後患矣，申說“蓄疑敗謀”也。 \par}

位不期驕，祿不期侈。\footnote{貴不與驕期而驕自至，富不與侈期而侈自來，驕侈以行己，所以速亡。}恭儉惟德，無載爾偽。\footnote{言當恭儉，惟以立德，無行奸偽。}作德，心逸日休。作偽，心勞日拙。\footnote{為德，直道而行，於心逸豫而名日美。為偽，飾巧百端，於心勞苦而事日拙，不可為。}居寵思危，罔不惟畏，弗畏入畏。\footnote{言雖居貴寵，當思危，惟無所不畏。若乃不畏，則入可畏之刑。}推賢讓能，庶官乃和，不和政厖。\footnote{賢能相讓，俊乂在官，所以和諧。厖，亂也。○厖，武江反。}舉能其官，惟爾之能。稱匪其人，惟爾不任。”\footnote{所舉能修其官，惟亦汝之功能。舉非其人，亦惟汝之不勝其任。○勝音升。}王曰：“嗚呼！三事暨大夫，敬爾有官，亂爾有政，\footnote{嘆而敕之，公卿已下，各敬居汝所有之官，治汝所有之職。}以佑乃闢。永康兆民，萬邦惟無斁。”\footnote{言當敬治官政，以助汝君，長安天下兆民，則天下萬國惟乃無厭我周德。○斁音亦。長，直良反。厭，於豔反。}

{\noindent\shu\zihao{5}\fzkt 傳“為德”至“可為”。正義曰：為德者自得於己，直道而行，無所經營，於心逸豫，功成則譽顯，而名益美也。為偽者行違其方，枉道求進，思念欺巧，於心勞苦，詐窮則道屈，而事日益拙也。以此故，偽不可為。申說“無載爾偽”也。 \par}

\CJKunderline{成王}既伐東夷，肅慎來賀。\footnote{海東諸夷駒麗、扶余、馯貌之屬，\CJKunderline{武王}克商，皆通道焉。\CJKunderline{成王}即政而叛，王伐而服之，故肅慎氏來賀。○肅慎，馬本作息慎，云：“北夷也。”駒，俱付反，又如字。麗，力支反。馯，戶旦反，\CJKunderwave{地理志}音寒。貊，孟白反，\CJKunderwave{說文}作“貉，北方豸種。\CJKunderline{孔子}曰，貉之言貊。貊，惡也。”}王俾\CJKunderline{榮伯}作\CJKunderwave{賄肅慎之命}。\footnote{榮，國名。同姓諸侯，為卿大夫。王使之為命書，以幣賄賜肅慎之夷。亡。○俾，必爾反，馬本作辯。}


{\noindent\zhuan\zihao{6}\fzbyks 傳“海東”至“來賀”。正義曰：\CJKunderline{成王}伐淮夷，滅徐奄，指言其國之名。此傳言“東夷”，非徒淮水之上夷也，故以為“海東諸夷駒麗、扶余、馯貊之屬”，此皆於\CJKunderline{孔君}之時有此名也。\CJKunderwave{周禮·職方氏}四夷之名、八蠻、九貉，\CJKunderline{鄭玄}云“北方曰貉”，又云“東北夷也”。\CJKunderwave{漢書}有高駒麗、扶余、韓，無此馯，馯即彼韓也，音同而字異爾。\CJKunderwave{多方}云：“王來自奄。”奄在後滅，言滅奄即來，必非滅奄之後更伐東夷。夷在海東路遠，又不得先伐遠夷，後來滅奄，此云“\CJKunderline{成王}既伐東夷”，不知何時伐之。\CJKunderwave{魯語}云：“\CJKunderline{武王}克商，遂通道於九夷、八蠻，於是肅慎氏來賀，貢楛矢。”則\CJKunderline{武王}之時，東夷服也。\CJKunderline{成王}即政，奄與淮夷近者尚叛，明知遠夷亦叛。蓋\CJKunderline{成王}親伐淮夷而滅之，又使偏師伐東夷而服之。君統臣功，故言王伐，不是\CJKunderline{成王}親自伐也。肅慎之於中國，又遠於所伐諸夷，見諸夷既服，故懼而來賀也。 \par}

{\noindent\zhuan\zihao{6}\fzbyks 傳“榮國”至“夷亡”。正義曰：\CJKunderwave{晉語}云：“\CJKunderline{文王}諏於蔡、原，訪於辛、尹，重之以周、召、畢、榮。”於\CJKunderline{文王}之時，名次\CJKunderline{畢公}之下，則是大臣也。未知此時\CJKunderline{榮伯}是彼榮公以否，或是其子孫也。“同姓諸侯”，相傳為然，注\CJKunderwave{國語}者亦云榮、周同姓。不知時為何官，故並云“卿大夫”。王使\CJKunderline{榮伯}，明使中所作。史錄其篇，名為\CJKunderwave{賄肅慎之命}，明是王使之為命書，以幣賜肅慎氏之夷也。 \par}

{\noindent\shu\zihao{5}\fzkt “\CJKunderline{成王}”至“之命”。正義曰：\CJKunderline{成王}即政之初，東夷背叛。\CJKunderline{成王}既伐而服之，東北遠夷其國有名肅慎氏者，以王戰勝，遠來朝賀。王賜以財賄，使榮國之伯為策書，以命肅慎之夷，嘉其慶賀，慰其勞苦之意。史敘其事，作\CJKunderwave{賄肅慎之命}名篇也。 \par}

\CJKunderline{周公}在豐，\footnote{致政老歸。}將沒，欲葬成周。\footnote{己所營作，示終始念之。}公薨，\CJKunderline{成王}葬於畢，\footnote{不敢臣\CJKunderline{周公}，故使近文武之墓。○近，附近之近。}告\CJKunderline{周公}，作\CJKunderwave{亳姑}。\footnote{\CJKunderline{周公}徙奄君於亳姑，因告柩以葬畢之義，並及奄君已定亳姑，言所遷之功成。亡。○柩，其久反。}


{\noindent\zhuan\zihao{6}\fzbyks 傳“致政老歸”。正義曰：\CJKunderline{周公}既還政\CJKunderline{成王}，\CJKunderline{成王}又留為太師，今言“\CJKunderline{周公}在豐”，則是去離王朝，又致太師之政，告老歸於豐，如\CJKunderline{伊尹}之告歸也。\CJKunderline{成王}封\CJKunderline{伯禽}於魯，以為\CJKunderline{周公}後。公老不歸魯而在豐者，文十三年\CJKunderwave{公羊傳}云：“\CJKunderline{周公}曷為不之魯？欲天下之一乎周也。”何休云：“\CJKunderline{周公}聖人，德至重，功至大，東征則西國怨，西征則東國怨。嫌之魯，恐天下回心趣向之。故封\CJKunderline{伯禽}，命使遙供養，死則奔喪為主，所以一天下之心於周室。”是言\CJKunderline{周公}不歸魯之意也。歸豐者，蓋以先王之都，欲近其宗廟故也。 \par}

{\noindent\zhuan\zihao{6}\fzbyks 傳“\CJKunderline{周公}”至“成亡”。正義曰：序說葬\CJKunderline{周公}之事，其篇乃名\CJKunderwave{亳姑}，篇名與序不相允會。其篇既亡，不知所道，故傳原其意而為之說。上篇將遷亳姑，序言“\CJKunderline{成王}既踐奄，將遷其君於亳姑”者，是\CJKunderline{周公}之意。今告\CJKunderline{周公}之柩以葬畢之義，乃用“亳姑”為篇名，必是告葬之時，並言及奄君已定於亳姑，言\CJKunderline{周公}所遷之功成，故以名篇也。 \par}

{\noindent\shu\zihao{5}\fzkt “\CJKunderline{周公}”至“亳姑”。正義曰：\CJKunderline{周公}既致政於王，歸在豐邑。將沒，遺言欲得葬於成周。以成周是己所營，示己終始念之，故欲葬焉。及公薨，\CJKunderline{成王}葬於畢，以文武之墓在畢，示己不敢臣\CJKunderline{周公}，使近文武之墓。王以葬畢之義告\CJKunderline{周公}之柩，又\CJKunderline{周公}徙奄君於亳姑，因言亳姑功成。史敘其事，作\CJKunderwave{亳姑}之篇。案\CJKunderwave{帝王世紀}云：“文武葬於畢。”畢在杜南，\CJKunderwave{晉書·地道記}亦云畢在杜南，與畢陌別，俱在長安西北。 \par}

\section{君陳第二十三【偽】}


\CJKunderline{周公}既沒，命\CJKunderline{君陳}分正東郊成周，\footnote{\CJKunderline{成王}重\CJKunderline{周公}所營，故命\CJKunderline{君陳}分居，正東郊成周之邑里官司。}作\CJKunderwave{君陳}。\footnote{作書命之。}

君陳\footnote{臣名也，因以名篇。○鄭注\CJKunderwave{禮記}云：“\CJKunderline{周公}之子。”}

王若曰：“\CJKunderline{君陳}，惟爾令德孝恭。\footnote{言其有令德，善事父母，行己以恭。}惟孝，友于兄弟，克施有政。\footnote{言善父母者必友于兄弟，能施有政令。}命汝尹茲東郊，敬哉！\footnote{正此東郊，監殷頑民，教訓之。○監，工銜反。}昔\CJKunderline{周公}師保萬民，民懷其德。往慎乃司，茲率厥常。\footnote{言\CJKunderline{周公}師安天下之民，民歸其德。今往承其業，當慎汝所主，此循其常法而教訓之。}懋昭\CJKunderline{周公}之訓，惟民其乂。\footnote{勉明\CJKunderline{周公}之教，惟民其治。○懋音茂。治，直吏反，下注“政治”同。}


{\noindent\zhuan\zihao{6}\fzbyks 傳“\CJKunderline{成王}”至“官司”。正義曰：“成周”，周之下都。監成周者正是一邑宰爾，而特命\CJKunderline{君陳}大其事者，\CJKunderline{成王}重\CJKunderline{周公}所營，猶恐殷民有不服之者，故命\CJKunderline{君陳}分居，正東郊成周之邑里官司也。以\CJKunderwave{畢命}之序言“分居”，知此“分”亦為分居，分別殷民善惡所居，即\CJKunderwave{畢命}所云“旌別淑慝，表厥宅裡”是也。言“東郊者”，\CJKunderline{鄭玄}云：“天子之國五十里為近郊，今河南洛陽相去則然。”是言成周之邑為周之東郊也。 \par}

{\noindent\zhuan\zihao{6}\fzbyks 傳“臣名”至“名篇”。正義曰：孔直云“臣名”，則非\CJKunderline{周公}子也。\CJKunderline{鄭玄}注\CJKunderwave{中庸}云“\CJKunderline{君陳}，蓋\CJKunderline{周公}子”者，以經云“\CJKunderline{周公}既沒，命\CJKunderline{君陳}”，猶若\CJKunderline{蔡叔}既沒，命\CJKunderline{蔡仲}故也。孔未必然矣。 \par}

{\noindent\zhuan\zihao{6}\fzbyks 傳“言其”至“以恭”。正義曰：“令德”，在身之大名。“孝”是事親之稱，“恭”是身之所行，言其善事父母，行己以恭也。\CJKunderwave{釋訓}云：“善父母為孝，善兄弟為友。” \par}

{\noindent\zhuan\zihao{6}\fzbyks 傳“言善”至“政令”。正義曰：父母尊之極，兄弟親之甚，緣其施孝於極尊，乃能施友于甚親。言善事父母者必友于兄弟，推此親親之心，以至於疏遠，每事以仁恕行之，故能施有政令也。 \par}

{\noindent\shu\zihao{5}\fzkt “\CJKunderline{周公}”至“\CJKunderline{君陳}”。正義曰：\CJKunderline{周公}遷殷頑民於成周，頑民既遷，\CJKunderline{周公}親自監之。\CJKunderline{周公}既沒，\CJKunderline{成王}命其臣名\CJKunderline{君陳}代\CJKunderline{周公}監之，分別居處，正此東郊成周之邑，以策書命之。史錄其事，作策書，為\CJKunderwave{君陳}篇名。 \par}

“我聞曰：‘至治馨香，感於神明。黍稷非馨，明德惟馨。’\footnote{所聞之古聖賢之言，政治之至者，芬芳馨氣動於神明。所謂芬芳非黍稷之氣，乃明德之馨。勵之以德。}爾尚式時\CJKunderline{周公}之猷訓，惟日孜孜，無敢逸豫。\footnote{汝庶幾用是\CJKunderline{周公}之道教殷民，惟當日孜孜勤行之，無敢自寬暇逸豫。○孜音茲。}

{\noindent\shu\zihao{5}\fzkt “我聞”至“逸豫”。正義曰：我聞人之言曰：“有至美治之善者，乃有馨香之氣，感動於神明。所言馨香感神者，黍稷飲食之氣非馨香也，明德之所遠及乃惟為馨香爾。”勉勵\CJKunderline{君陳}使為德也。欲必為明德，惟法\CJKunderline{周公}，汝當庶幾用是\CJKunderline{周公}之道。惟當每日孜孜勤法行之，無敢自寬暇逸豫。教使勤於事也。 \par}

凡人未見聖，若不克見。既見聖，亦不克由聖，\footnote{此言凡人有初無終，未見聖道如不能得見。已見聖道，亦不能用之，所以無成。}爾其戒哉！爾惟風，下民惟草。\footnote{汝戒，勿為凡人之行。民從上教而變，猶草應風而偃，不可不慎。}圖厥政，莫或不艱，有廢有興。出入自爾師虞，庶言同則繹。\footnote{謀其政，無有不先慮其難，有所廢，有所起。出納之事，當用汝總言度之。眾言同，則陳而布之。禁其專。○繹音亦。度，待洛反。}爾有嘉謀嘉猷，則入告爾後於內，爾乃順之於外，\footnote{汝有善謀善道，則入告汝君於內，汝乃順行之於外。}曰：‘斯謀斯猷，惟我後之德。’\footnote{此善謀此善道，惟我君之德。善則稱君，人臣之義。}嗚呼！臣人咸若時，惟良顯哉！”\footnote{嘆而美之曰：“臣於人者皆順此道，是惟良臣，則君顯明於世。”}王曰：“\CJKunderline{君陳}，爾惟弘\CJKunderline{周公}丕訓，無依勢作威，無倚法以削，\footnote{汝為政當闡大\CJKunderline{周公}之大訓，無乘勢位作威人上，無倚法制以行刻削之政。}寬而有制，從容以和。\footnote{寬不失制，動不失和，德教之治。○從，七容反。}殷民在闢，予曰闢，爾惟勿闢。予曰宥，爾惟勿宥。惟厥中。\footnote{殷人有罪在刑法者，我曰：“刑之。”汝勿刑。我曰：“赦宥。”汝勿宥。惟其當以中正平理斷之。○闢，扶亦反，下同。中如字，或丁仲反。斷，丁亂反。}有弗若於汝政，弗化於汝訓，闢以止闢，乃闢。\footnote{有不順於汝政，不變於汝教，刑之而懲止犯刑者，乃刑之。}狃於姦宄，敗常亂俗，三細不宥。\footnote{習於姦宄兇惡，毀敗五常之道，以亂風俗之教，罪雖小，三犯不赦，所以絕惡源。○狃，女九反。}


{\noindent\zhuan\zihao{6}\fzbyks 傳“汝為”至“之政”。正義曰：\CJKunderline{君陳}之智，必不及\CJKunderline{周公}，而令闡大\CJKunderline{周公}訓者，遵行其法，使廣被於民，即是闡揚而大之。非遣\CJKunderline{君陳}為法，使大於\CJKunderline{周公}法也。凡在人上，位貴於人，勢足可畏者，多乘是形勢以作威刑于人，倚附公法以行刻削之政，故禁之也。 \par}

{\noindent\zhuan\zihao{6}\fzbyks 傳“寬不”至“之治”。正義曰：“寬不失制”，則經“寬而有制”。“動不失和”，則經“從容以和”。言“動”，謂“從容”也。 \par}

{\noindent\zhuan\zihao{6}\fzbyks 傳“習於”至“惡源”。正義曰：\CJKunderwave{釋言}云：“狃，復也。”孫炎曰：“狃忕，前復為也。”古言“狃忕”是貫習之義，故以“習”解“狃”。“習於姦宄兇惡”，言為之不知止也。“敗常亂俗”，有大有小，罪雖小者，三犯不赦，恐其滋大，所以絕惡源也。此謂所犯小事，言“三”者，再猶可赦爾。 \par}

{\noindent\shu\zihao{5}\fzkt “王曰”至“不宥”。正義曰：王呼之曰：“\CJKunderline{君陳}，汝今為政，當弘大\CJKunderline{周公}之大訓。\CJKunderline{周公}既有大訓，汝當遵而行之，使其法更寬大。汝奉\CJKunderline{周公}之訓，無得依恃形勢以作威於人，無得倚附法制以行刻削百姓。必當寬容而有法制，使疏而不漏。從容以和協於物，莫為褊急。此成周殷民，有犯事在於刑法未斷決者，我告汝曰：‘刑罰之。’汝惟勿得刑罰之。我告汝曰：‘赦宥之。’汝惟勿得赦宥之。惟其以中正平法斷決之，不得從上意也。其有不順於汝之政令，不化於汝之訓教，其罪既大，當行刑中。刑罰一人可以止息後犯者，故云犯刑者乃刑之。如其罪或輕細，罰不當理，雖刑勿息，故不可輒刑。若有人習於姦宄兇惡，敗五常之道，亂風俗之教，三犯其事者，事雖細小，勿得宥之。以其知而故犯，當殺之以絕惡源也。” \par}

爾無忿疾於頑，無求備於一夫。\footnote{人有頑嚚不喻，汝當訓之，無忿怒疾之。使人當器之，無責備於一夫。}必有忍，其乃有濟。有容，德乃大。\footnote{為人君長，必有所含忍，其乃有所成。有所包容，德乃為大。欲其忍恥藏垢。○長，誅丈反。垢，工口反。}簡厥修，亦簡其或不修。\footnote{簡別其德行修者，亦別其有不修者，善以勸能，惡以沮否。○別，彼列反。沮，在汝反。否，方九反，又音鄙。}進厥良，以率其或不良。\footnote{進顯其賢良者，以率勉其有不良者，使為善。}

{\noindent\shu\zihao{5}\fzkt “爾無”至“不良”。正義曰：民者，冥也，當以漸教之。故戒\CJKunderline{君陳}：“民有不知道者，汝無忿怒疾惡。頑嚚之民，當以漸教訓之。無求備於一人，當取其所能。在為人君，必有所含忍，其事乃有所成。有所寬容，其德乃能大。”欲其寬大不褊隘也。“汝之為政，須知民之善惡，簡別其德行修者，亦簡別其有不修德行者。進顯其賢良，以率勵其不良者”。欲令其化惡，使為善也。 \par}

惟民生厚，因物有遷。\footnote{言人自然之性敦厚，因所見所習之物有遷變之道，故必慎所以示之。}違上所命，從厥攸好。\footnote{人之於上，不從其令，從其所好，故人主不可不慎所好。○好，呼報反。}爾克敬典在德，時乃罔不變。允升於大猷。\footnote{汝治人能敬常在道德，是乃無不變化，其政教則信升於大道。}惟予一人膺受多福，\footnote{汝能升大道，則惟我一人亦當受其多福，無兇危。}其爾之休，終有辭於永世。”\footnote{非但我受多福而已，其汝之美名，亦終見稱誦於長世。言沒而不朽。長如字。朽，許久反。}

{\noindent\shu\zihao{5}\fzkt “惟民”至“永世”。正義曰：惟民初生，自然之性皆敦厚矣。因見所習之物，本性乃有遷變，為惡皆由習效使然。人之情性，好違上所命，命之不必從也，從其君所好。君之所好，民必從之，在上者不可不慎所好也。汝之治民能敬，當從終常在於道德教之。汝以道德教之，是民乃無不變化。民皆變從汝化，則信升於大道矣。汝能如此，惟我一人亦當受其多福，無兇危矣。其汝之美名，亦終有稱誦之美辭於長世矣。 \par}

\section{顧命第二十四（顧命上）}


\CJKunderline{成王}將崩，命\CJKunderline{召公}、\CJKunderline{畢公}。\footnote{二公為二伯，中分天下而治之。}率諸侯相\CJKunderline{康王}，作\CJKunderwave{顧命}。\footnote{臨終之命曰顧命。○相，息亮反。顧，工戶反，馬云：“\CJKunderline{成王}將崩，顧命\CJKunderline{康王}，命\CJKunderline{召公}、\CJKunderline{畢公}率諸侯輔相之。”}


{\noindent\zhuan\zihao{6}\fzbyks 傳“二公”至“治之”。正義曰：\CJKunderwave{禮記·曲禮下}文云，“九州之長曰牧”，“五官之長曰伯，是職方”。\CJKunderline{鄭玄}云：“職，主也。謂為三公者，是伯分主東西者也。”\CJKunderwave{周禮·大宗伯}云：“八命作牧，九命作伯。”鄭云，謂“上公有功德者，加命為二伯”。此\CJKunderwave{禮}文皆伯尊於牧，牧主一州，明伯是中分天下者也。\CJKunderwave{禮}言“職方”，是各主一方也。此二伯即以三公為之。隱五年\CJKunderwave{公羊傳}云：“諸公者何？天子三公。天子三公者何？天子之相也。天子之相何以三？自陝而東者\CJKunderline{周公}主之，自陝而西者\CJKunderline{召公}主之，一相處乎內。”是言三公為二伯也。\CJKunderwave{公羊傳}漢世之書，陝縣者漢之弘農郡所治，其地居二京之中，故以為二伯分掌之界，\CJKunderline{周公}所分亦當然也。\CJKunderwave{公羊傳}所言周、召分主，謂\CJKunderline{成王}即位之初，此時\CJKunderline{周公}已薨，故\CJKunderline{畢公}代之。\CJKunderwave{周官}篇三公之次太師、太傅、太保，太保最在下。此篇以\CJKunderline{召公}為先者，三公命數尊卑同也，王就其中委任賢者，任之重者則在前耳。 \par}

{\noindent\zhuan\zihao{6}\fzbyks 傳“臨終”至“顧命”。正義曰：\CJKunderwave{說文}云：“顧，還視也。”\CJKunderline{鄭玄}云：“回首曰顧，顧是將去之意。”此言“臨終之命曰顧命”，言臨將死去，回顧而為語也。 \par}

{\noindent\shu\zihao{5}\fzkt “\CJKunderline{成王}”至“顧命”。正義曰：\CJKunderline{成王}病困將崩，召集群臣以言，命太保\CJKunderline{召公}、太師\CJKunderline{畢公}，使率領天下諸侯輔相\CJKunderline{康王}。史敘其事。作\CJKunderwave{顧命}。 \par}

顧命\footnote{實命群臣，敘以要言。}

惟四月哉生魄，王不懌。\footnote{\CJKunderline{成王}崩年之四月，始生魄，月十六日，王有疾，故不悅懌。○懌音亦。馬本作“不釋”，云：“不釋，疾不解也。”}甲子，王乃\xpinyin*{洮頮}水。相被冕服，憑玉幾。\footnote{王大發大命，臨群臣，必齋戒沐浴。今疾病，故但洮盥頮面。扶相者，被以冠冕，加朝服，憑玉幾以出命。○洮，他刀反，徐音逃，馬云：“洮，洮發也。”頮音悔，\CJKunderwave{說文}作沬，云：“古文作頮。”馬云：“頮，頮面也。”被，皮義反。徐，扶偽反，注同。憑，皮冰反，下同。\CJKunderwave{說文}作憑，云：“依幾也。”\CJKunderwave{字林}同，父冰反。齊，側皆反。盥音管，又音灌。朝，直遙反。}



{\noindent\zhuan\zihao{6}\fzbyks 傳“實命”至“要言”。正義曰：王之所命，實普命群臣，序以要約為言，直云“命\CJKunderline{召公}、\CJKunderline{畢公}”。傳不於上“\CJKunderline{召公}、\CJKunderline{畢公}”之下而解，於“顧命”之下言之者，以上欲指明二公中分天下之事，非是總語，故“命”不得言之。“顧命”是總命群臣，非但召、畢而已，故於此解也。 \par}

{\noindent\zhuan\zihao{6}\fzbyks 傳“\CJKunderline{成王}”至“悅謂”。正義曰：\CJKunderline{成王}崩年，經典不載，\CJKunderwave{漢書·律曆志}云，\CJKunderline{成王}即位“三十年四月庚戌朔，十五日甲子哉生魄”，即引此\CJKunderwave{顧命}之文。以為\CJKunderline{成王}即位三十年而崩，此是劉歆說也。孔以甲子為十六日，則不得與歆同矣。\CJKunderline{鄭玄}云：“此\CJKunderline{成王}二十八年。”傳惟言“\CJKunderline{成王}崩年”，未知\CJKunderline{成王}即位幾年崩也。\CJKunderwave{志}又云：“死魄，朔也。生魄，望也。”明死魄生從望為始，故始生魄為月十六日，即是望之日也。\CJKunderwave{釋詁}云：“懌，樂也。”有疾，故不悅懌。下云“病日臻，既彌留”，則\CJKunderline{成王}遇病已多日矣。於“哉生魄”下始言“王不懌”者，甲子是發命之日，為“洮頮”張本耳。 \par}

{\noindent\zhuan\zihao{6}\fzbyks 傳“王大發”至“出命”。正義曰：凡有敬事，皆當潔清。王將發大命，臨群臣，必齋戒沐浴。今以病疾之故，不能沐浴，故但洮頮而已。\CJKunderwave{禮}洗手謂之“盥”，洗面謂之“靧”。\CJKunderwave{內則}云，子事父母“面垢，燂潘請靧”。“頮”是洗面，知“洮”為盥手。言“水”謂洮盥俱用水。扶相王者，以冕服加王。\CJKunderline{鄭玄}云：“相者，正王服位之臣，謂太僕。”或當然也。“被以冠冕”，以冕服被王首也。“加朝服”，以服加王身也，謂以袞冕朝諸侯之服加王身也。鄭以為玄冕，知不然者，以顧命群臣，大發大命，以文武之業傳社稷之重，不應惟服玄冕而已。\CJKunderwave{覲禮}王服袞冕而有玉幾。此既“憑玉幾”，明服袞冕也。\CJKunderwave{周禮·司几筵}云，凡大朝覲，王位設黼扆，扆前南向設左右玉幾。是王見群臣當憑玉幾以出命。 \par}

乃同召太保\CJKunderline{奭}、\CJKunderline{芮伯}、\CJKunderline{彤伯}、\CJKunderline{畢公}、\CJKunderline{衛侯}、\CJKunderline{毛公}、\footnote{同召六卿，下至御治事。太、保、畢毛稱公，則三公矣。此先後六卿次第，冢宰第一，\CJKunderline{召公}領之。司徒第二，\CJKunderline{芮伯}為之。宗伯第三，\CJKunderline{彤伯}為之。司馬第四，\CJKunderline{畢公}領之。司寇第五，\CJKunderline{衛侯}為之。司空第六，\CJKunderline{毛公}領之。召、芮、彤、畢、衛、毛皆國名，入為天子公卿。○奭音釋。芮，如銳反。彤，徒冬反。}師氏、虎臣、百尹、御事。\footnote{師氏，大夫官。虎臣，虎賁氏。百尹，百官之長。及諸御治事者。}

{\noindent\zhuan\zihao{6}\fzbyks 傳“同召”至“公卿”。正義曰：下及御事，蒙此同召之，文故云同召六卿，下及御事也。以王病甚，故同時俱召之。“太保”是三公官名，畢、毛又亦稱“公”，知此三人是三公也。三人是三公，而與侯伯相次，知六者是六卿。\CJKunderline{衛侯}為司寇而位第五，知此先後是六卿次第也。以三公尊，故特言“公”，其餘三卿舉其本爵，見其以國君入為卿也。天子三公皆以卿為之，不復別置其人。高官兼攝下司者，漢世以來謂之為“領”，故言“\CJKunderline{召公}領之”，“\CJKunderline{毛公}領之”。定四年\CJKunderwave{左傳}云“\CJKunderline{康叔}為司寇”，知此六人依\CJKunderwave{周禮}次第為六卿也。\CJKunderline{王肅}云：“彤，姒姓之國，其餘五國姬姓。畢、毛，\CJKunderline{文王}庶子。\CJKunderline{衛侯}，\CJKunderline{康叔}所封，\CJKunderline{武王}母弟。”依\CJKunderwave{世本}、\CJKunderwave{史記}為說也。 \par}

{\noindent\zhuan\zihao{6}\fzbyks 傳“師氏”至“事者”。正義曰：\CJKunderwave{周禮}師氏，中大夫，掌以美詔王，居虎門之左，司王朝得失之事，帥其屬守王之門。重其所掌，故於“虎臣”並於“百尹”之上，特言之。“尹”訓正也，故“百尹”為百官之長。“諸御治事”謂諸掌事者，蓋大夫、士皆被召也。\CJKunderline{王肅}云：“治事，蓋群士也。” \par}

{\noindent\shu\zihao{5}\fzkt “顧命”至“御事”。正義曰：發首至“百尹、御事”，敘王以病召臣，為發言之端。自“王曰”至“冒貢於非幾”是顧命之辭也。“茲既受命”至“立於側階”，言命后王崩，欲宣王命，布陳儀衛之事也。自“王麻冕”已下，敘\CJKunderline{康王}受命之事。 \par}

王曰：“嗚呼！疾大漸，惟幾。\footnote{自嘆其疾大進篤，惟危殆。○幾音機，徐音畿，下同。}病日臻，既彌留，恐不獲誓言嗣，茲予審訓命汝。\footnote{病日至，言困甚。已久留，言無瘳。恐不得結信出言嗣續我志，以此故，我詳審教命汝。○瘳，敕留反。}昔君\CJKunderline{文王}、\CJKunderline{武王}宣重光，奠麗陳教則肄。\footnote{言昔先君文武，布其重光，累聖之德，定天命，施陳教，則勤勞。○重光，馬云：“日月星也。太極上元十一月朔旦冬至，日月如疊璧，五星如連珠，故曰重光。”重，直龍反。麗，力馳反。肄，徐以至反，又以制反。}肄不違，用克達殷，集大命。\footnote{文武定命陳教，雖勞而不違道，故能通殷為周，成其大命。}在後之侗，敬迓天威，嗣守\CJKunderline{文}、\CJKunderline{武}大訓，無敢昏逾。\footnote{在文武后之侗稚，\CJKunderline{成王}自斥。敬迎天之威命，言奉順繼守文武大教。無敢昏亂逾越，言戰慄畏懼。○侗徐音同，又敕動反，馬本作詷，云：“共也。”斥，昌亦反。}今天降疾殆，弗興弗悟。爾尚明時朕言，\footnote{今天下疾我身，甚危殆，不起不悟。言必死。汝當庶幾明是我言，勿忽略。}用敬保元子\CJKunderline{釗}，弘濟於艱難，\footnote{用奉我言，敬安太子釗。釗，\CJKunderline{康王}名。大渡於艱難，勤德政。○釗，姜遼反，又音昭，徐之餚反。}柔遠能邇，安勸小大庶邦。\footnote{言當和遠，又能和近，安小大眾國，勸使為善。}思夫人自亂於威儀，爾無以\CJKunderline{釗}冒貢於非幾。”\footnote{群臣皆宜思夫人，夫人自治正於威儀。有威可畏，有儀可象，然後足以率人。汝無以釗冒進於非危之事。}


{\noindent\zhuan\zihao{6}\fzbyks 傳“病日”至“命汝”。正義曰：“病日至”者，言日日益至，遍於身體，困甚也。“已久留”者，言病來多日，無瘳愈也。“恐死不得結信出言嗣續我志”，志欲有言，若不能言，則不得續志。以此，及今能言，故我詳審出言教命汝。言己詳審，欲其敬聽之。 \par}

{\noindent\zhuan\zihao{6}\fzbyks 傳“今天”至“忽略”。正義曰：孔讀“殆”上屬為句，“今天下疾我身，甚危殆”也。“不起”言身不能起，“不悟”言心不能覺悟。病者形弱神亂，“不起不悟”，言必死也。 \par}

{\noindent\shu\zihao{5}\fzkt “王曰”至“非幾”。正義曰：王召群臣既集，乃言而嘆曰，嗚呼！我疾大進益重，惟危殆矣。病日日益至，言病困已甚。病既久留於我身，恐一旦暴死，不得結誓出言語以繼續我志，以此故，我今詳審教訓命誥汝等。昔先君\CJKunderline{文王}、\CJKunderline{武王}，布其重光，累聖之德，安定天命，施陳教誨，則勤勞矣。文武定命陳教，雖勞而不違於道，用能通殷為周，成其大命，代殷為主。至文武后之侗稚，\CJKunderline{成王}自謂己也。言己常敬迎天之威命，終當奉順天道，繼守文武大教，無敢昏逾越。言常戰慄畏懼，恐墜文武之業。今天降疾於我身，甚危殆矣。不能更起，不復覺悟。言己必死。汝等庶幾明是我言，勿忽略之。用我之語，敬安太子釗，大渡於艱難。言當安和遠人，又須能和近人，當為善政，遠近俱安之。又當安勸小大眾國，於彼小大眾國皆安之勸之。安之使國得安存，勸之使相勸為善。汝群臣等思夫人，夫人眾國各自治正於威儀。有威有儀，然後可以率人。無威無儀，則民不從命。戒使慎威儀也。汝無以釗冒進於非事危事。欲令戒其不為惡也。 \par}

茲既受命還，\footnote{此群臣已受顧命，各還本位。}出綴衣於庭。越翼日乙丑，王崩。\footnote{綴衣，幄帳。群臣既退，徹出幄帳於庭。王寢於北墉下，東首，反初生。於其明日，王崩。○出如字，徐尺遂反。綴，丁衛反，下同。王崩，馬本作“\CJKunderline{成王}崩”，注：“安民立政曰成。”幄，於角反，下同。墉音容，本亦作牖。首，手又反。}太保命\CJKunderline{仲桓}、\CJKunderline{南宮毛}，\footnote{冢宰攝政，故命二臣。桓、毛，名。}俾爰齊侯\CJKunderline{呂伋}，以二干戈、虎賁百人，逆子\CJKunderline{釗}於南門之外。\footnote{臣子皆侍左右，將正太子之尊，故出於路寢門外。使桓、毛二臣各執干戈，於齊侯\CJKunderline{呂伋}索虎賁百人，更新逆門外，所以殊之。伋為天子虎賁氏。○俾，必爾反。伋，居及反，齊侯名，太公子。}


{\noindent\zhuan\zihao{6}\fzbyks 傳“此群”至“本位”。正義曰：\CJKunderwave{周禮}：“射人掌國之三公、孤、卿大夫之位。三公北面，孤東面，卿大夫西面。”\CJKunderline{鄭玄}云：“不言士者，此與諸侯之賓射，士不與也。凡朝、燕及射，臣見於君之禮同。”鄭知然者，以\CJKunderwave{周禮}司士掌治朝之位，與\CJKunderwave{射人}同，是天子之澄芻與射禮位同。案\CJKunderwave{燕禮}小臣納卿大夫，卿大夫皆北面，公命爾卿東方西面，爾大夫少進，皆北面。\CJKunderwave{大射禮}其位亦然。是諸侯燕位與射位同，故云“朝、燕及射，臣見於君之禮同”。但天子臣多，故三公北面，孤東面，卿大夫西面。諸侯臣少，故卿西面，大夫北面，其士與天子同，皆門內西方東面。其入門當立定位如是，及王呼與言，必各自前進。已受顧命，退還本位者，謂還本治事之位，故孔下傳云：“朝臣就次，謂退王庭而還治事之處。” \par}

{\noindent\zhuan\zihao{6}\fzbyks 傳“綴衣”至“王崩”。正義曰：“綴衣”者，連綴衣物，出之於庭，則是從內而出。下云“狄設黼扆綴衣”，則綴衣是黼扆之類。黼扆是王坐立之處，知綴衣是施張於王坐之上，故以為“幄帳”也。\CJKunderwave{周禮}：“幕人掌帷幕、幄帟、綬之事。”\CJKunderline{鄭玄}云：“在旁曰帷，在上曰幕。帷幕皆以布為之，四合象宮室曰幄，王所居之帳也。帟，王在幕居幄中，坐上承塵也。幄帟皆以繒為之。”然則幄帳是黼扆之上所張之物。此言“出綴衣於庭”，則亦並出黼扆，故下句雲象王平生之時，更復設之。王發顧命在此黼扆幄帳之坐，命訖，乃復反於寢處。以王病重，不復能臨此坐，故徹出幄帳於庭，將欲為死備也。傳更解徹去幄帳之意，以王病困，寢不在此。\CJKunderwave{喪大記}云：“疾病，君、大夫徹懸，士去琴瑟，寢東首於北墉下，廢床。”\CJKunderline{鄭玄}云：“廢，去也。人始生在地，去床，庶其生氣反。也。”\CJKunderwave{記}言君、大夫、士，則尊卑皆然，故知此時王亦“寢於北墉下，東首，反初生”也。 \par}

{\noindent\zhuan\zihao{6}\fzbyks 傳“臣子”至“賁氏”。正義曰：天子初崩，太子必在其側。解其迎於門外之意，於時臣子皆侍左右，將正太子之尊，故使太子出於路寢門外，更迎入，所以殊之也。經言“以二干戈”，文在“齊侯\CJKunderline{呂伋}”下，似就齊侯取干戈。傳言“使桓、毛二臣各執干戈，於齊侯\CJKunderline{呂伋}索虎賁”，則是執干戈就齊侯。傳似反於經者，於時新遭大禍，內外嚴戒，桓、毛二人必是武臣宿衛，先執干戈，太保就命，使之就干戈以往，傳達其意，故移“干戈”之文於“齊侯”之上，傳言是實也。經言“於齊侯\CJKunderline{呂伋}”，下言“以二干戈、虎賁百人”者，指說迎太子之時有此備衛耳，非言二人干戈亦是齊侯授也。\CJKunderwave{周禮}虎賁氏下大夫，其屬有虎士八百人。知伋為天子虎賁氏，故就伋取虎賁也。 \par}

延入翼室，恤宅宗。\footnote{明室，路寢。延之使居憂，為天下宗主。}丁卯，命作冊度。\footnote{三日，命史為冊書法度，傳顧命於\CJKunderline{康王}。○度，舊音杜洛反，恐誤，注云“作冊書法度”，音宜如字。傳，直專反。}

{\noindent\zhuan\zihao{6}\fzbyks 傳“明室”至“宗主”。正義曰：\CJKunderwave{釋言}云：“翼，明也。”\CJKunderwave{喪大記}云：“君夫人卒於路寢。”以諸侯薨於路寢，知天子亦崩於路寢。今延太子入室，必延入喪所，知“翼室”是明室，謂路寢也。路寢之大者，故以“明”言之。延之使憂居喪主，為天下宗主也。 \par}

{\noindent\zhuan\zihao{6}\fzbyks 傳“三日”至“\CJKunderline{康王}”。正義曰：\CJKunderwave{周禮}內史掌策命，故命內史為策書也。經不言“命史”，史是常職，不假言之。王之將崩，雖口有遺命，未作策書，故以此日作之。既作策書，因作受策法度。下云“曰皇后憑玉幾”，宣\CJKunderline{成王}言，是策書也。將受命時，升階即位，及傳命已後，\CJKunderline{康王}答命，受同祭饗，皆是法度。 \par}

{\noindent\shu\zihao{5}\fzkt “茲既”至“冊度”。正義曰：此群臣既受王命，還複本位，出連綴之衣，王所坐幄帳，置之於庭。於其明日乙丑，王崩矣。太保\CJKunderline{召公}命\CJKunderline{仲桓}、\CJKunderline{南宮毛}，使此二人於齊侯\CJKunderline{呂伋}之所，以二干戈，桓毛各執其一，又取虎賁之士百人，迎太子釗於南門之外。逆此太子，使入於路寢明室。令太子在室當喪憂居，為天下宗主，正其將王之位以系群臣之心也。 \par}

越七日癸酉，伯相命士須材。\footnote{邦伯為相，則\CJKunderline{召公}。於丁卯七日癸酉，\CJKunderline{召公}命士致材木，須待以供喪用。○相，息亮反。供音恭。}狄設黼扆、綴衣。\footnote{狄，下士。扆，屏風，畫為斧文，置戶牖間。復設幄帳，象平生所為。○黼音甫，徐音補。扆，於豈反。屏，步經反。畫，胡卦反。牖音酉。復，扶又反。}


{\noindent\zhuan\zihao{6}\fzbyks 傳“邦伯”至“喪用”。正義曰：\CJKunderline{成王}既崩，事皆聽於冢宰，自非\CJKunderline{召公}無由發命，知“伯相”即\CJKunderline{召公}也。\CJKunderline{王肅}云：“\CJKunderline{召公}為二伯，相王室，故曰伯相。”上言“太保命\CJKunderline{仲桓}”，此改言“伯相”者，於此所命士多，非是國相不得大命諸侯，故改言“伯相”，以見政皆在焉。“於丁卯七日癸酉”，則王乙丑崩，於今已九日矣。於九日始傳顧命，不知其所由也。\CJKunderline{鄭玄}云：“癸酉蓋大斂之明日也。”鄭大夫已上殯斂，皆以死之來日數，天子七日而殯，於死日為八，故以癸酉為殯之明日。孔不為傳，不必如鄭說也。“須”訓待也。今所命者皆為喪事，知“命士須材”者，“\CJKunderline{召公}命士致材木，須待以供喪用”，謂槨與明器，是喪之雜用也。案\CJKunderwave{士喪禮}將葬，筮宅之後，始作槨及明器。此既殯即須材木者，以天子禮大，當須預營之。故\CJKunderwave{禮記}云：“虞人致百祀之木，可為棺槨者斬之。”是與士禮不同。顧氏亦云：“命士供葬槨之材。” \par}

{\noindent\zhuan\zihao{6}\fzbyks 傳“狄下”至“所為”。正義曰：\CJKunderwave{禮記·祭統}云：“狄者樂吏之賤者也。”是賤官有名為狄者，故以狄為下士。\CJKunderwave{喪大記}復魄之禮云“狄人設階”，是喪事使狄與此同也。\CJKunderwave{釋宮}云：“牖戶之間謂之扆。”李巡曰：“謂牖之東戶之西為扆。”郭璞曰：“窗東戶西也。\CJKunderwave{禮}雲‘斧，扆者’，以其所在處名之。”郭璞又云：“\CJKunderwave{禮}有斧扆，形如屏風，畫為斧文，置於扆地，因名為扆。”是先儒相傳，黼扆者，屏風，畫為斧文，在於戶牖之間。\CJKunderwave{考工記}云：“畫繢之事，白與黑謂之黼。”是用白黑畫屏風置之於扆地，故名此物為“黼扆”。上文言“出綴衣於庭”，此復設黼扆帷幄帳者，象王平生時所為也。經於四坐之上言“設黼扆、綴衣”，則四坐皆設之。此經所云“狄設”，亦是伯相命狄使設之。不言“命”者，上云“命士”，此蒙“命”文。設四坐及陳寶玉、兵器與輅車，各有所司，皆是相命，不言所命之人，從上省文也。 \par}

{\noindent\shu\zihao{5}\fzkt “越七日”至“癸酉”。正義曰：自此以下至“立於側階”，惟“命士須材”是擬供喪用，其餘皆是將欲傳命佈設之事。四坐王之所處者，器物國之所寶者，車輅王之所乘者，陳之所以華國，且以示重顧命。其執兵器立於門內堂階者，所以備不虞，亦為國家之威儀也。 \par}

牖間南向，敷重篾席，黼純，華玉仍幾。\footnote{蔑，桃枝竹。白黑雜繒緣之。華，彩色。華玉以飾憑几。仍，因也。因生時，幾不改作。此見群臣、覲諸侯之坐。○向，許亮反。篾，眠結反，馬云：“纖蒻。”純，之允反，又之閏反，下同。緣，悅絹反，本或作純。}西序東向，敷重厎席，綴純，文貝仍幾。\footnote{東西廂謂之序。厎，蒻蘋。綴，雜彩。有文之貝飾幾。此旦夕聽事之坐。○厎，之履反，馬云：“青蒲也。”蒻音弱。蘋音平。}東序西向，敷重豐席，畫純，雕玉仍幾。\footnote{豐，莞。彩色為畫。雕，刻鏤。此養國老饗群臣之坐。○豐，芳弓反。莞音官，又音關。鏤，來豆反。}


{\noindent\zhuan\zihao{6}\fzbyks 傳“篾桃”至“之坐”。正義曰：此篾席與\CJKunderwave{周禮}“次席”一也。鄭注彼云：“次席，桃枝席，有次列成文。”\CJKunderline{鄭玄}不見孔傳，亦言是桃枝席，則此席用桃枝之竹，必相傳有舊說也。鄭注此下則云：“篾，析竹之次青者。”\CJKunderline{王肅}云：“篾席，纖蒻蘋席。”並不知其所據也。\CJKunderwave{考工記}云：“白與黑謂之黼。”\CJKunderwave{釋器}云：“緣謂之純。”知黼純是白黑雜繒緣之，蓋以白繒黑繒錯雜彩以緣之。\CJKunderline{鄭玄}注\CJKunderwave{周禮}云：“斧謂之黼，其繡白黑採也，以絳帛為質。”其意以白黑之線縫剌為黼又以緣席，其事或當然也。“華”是彩之別名，故以為“彩色，用華玉以飾憑几”也。\CJKunderline{鄭玄}云：“華玉，五色玉也。”“仍，因也”，\CJKunderwave{釋詁}文。\CJKunderwave{周禮}云：“凡吉事變幾，凶事仍幾。”禮之於幾有變有仍，故特言“仍幾”，以見“因生時，幾不改作”也。“此見群臣、覲諸侯之坐”，\CJKunderwave{周禮}之文知之。又\CJKunderwave{覲禮}，天子待諸侯，“設斧扆於戶牖之間，左右幾，天子袞冕負斧扆”。彼在朝，此在寢為異，其牖間之坐則同。 \par}

{\noindent\zhuan\zihao{6}\fzbyks 傳“東西”至“之坐”。正義曰：“東西廂謂之序”，\CJKunderwave{釋宮}文。孫炎曰：“堂東西牆，所以別序內外也。”\CJKunderwave{禮}注謂蒲席為蒻蘋，孔以“厎席”為蒻蘋，當謂蒲為蒲蒻之席也。史游\CJKunderwave{急就篇}云“蒲蒻藺席”，“蒲蒻”謂此也。\CJKunderline{王肅}云：“厎席，青蒲席也。”\CJKunderline{鄭玄}云：“厎，致也。篾纖致席也。”鄭謂此“厎席”亦竹蓆也。凡此重席，非有明文可據，各自以意說耳。“綴”者，連綴諸色。席必以彩為緣，故以“綴”為雜彩也。“貝”者水蟲，取其甲以飾器物。\CJKunderwave{釋魚}於貝之下云：“餘蚳，黃白文。餘泉，白黃文。”李巡曰：“貝甲以黃為質，白為文彩，名為餘蚳。貝甲以白為質，黃為文彩，名為餘泉。”“有文之貝飾幾”，謂用此餘蚳、餘泉之貝飾幾也。“此旦夕聽事之坐”，鄭、王亦以為然。牖間是“見群臣、覲諸侯之坐”，見於\CJKunderwave{周禮}。其東序西向，“養國老饗群臣之坐”者，案\CJKunderwave{燕禮}云“坐於阼階上，西向”，則養國老及饗與\CJKunderwave{燕禮}同。其西序之坐在燕饗坐前，以其旦夕聽事，重於燕飲，故西序為“旦夕聽事之坐”。夾室之坐在燕饗坐後，又夾室是隱映之處，又親屬輕於燕饗，故夾室為“親屬私宴之坐”。案朝士職掌治朝之位，王南面，此“西序東向”者，以此諸坐並陳，避牖間南鄉覲諸侯之坐故也。\CJKunderline{王肅}說四坐，皆與孔同。 \par}

{\noindent\zhuan\zihao{6}\fzbyks 傳“豐莞”至“之坐”。正義曰：\CJKunderwave{釋草}云：“莞，苻籬。”郭璞曰：“今之西方人呼蒲為莞，用之為席也。”又云“𦸣，鼠莞”。樊光曰：“\CJKunderwave{詩}云：‘下莞上簟。’”郭璞曰：“似莞而纖細，今蜀中所出莞席是也。”\CJKunderline{王肅}亦云：“豐席，莞。”\CJKunderline{鄭玄}云：“豐席，刮凍竹蓆。”\CJKunderwave{考工記}云：“畫繢之事，雜五色。”是彩色為畫，蓋以五彩色畫帛以為緣。\CJKunderline{鄭玄}云：“似雲氣，畫之為緣。”\CJKunderwave{釋器}云：“玉謂之雕，金謂之鏤，木謂之刻。”是“雕”為刻鏤之類，故以“刻鏤”解“雕”，蓋雜以金玉，刻鏤為飾也。 \par}

西夾南向，敷重筍席，玄紛純，漆仍幾。\footnote{西廂夾室之前。筍，蒻竹。玄紛，黑綬。此親屬私宴之坐，故席幾質飾。○夾，工洽反，徐音頰，注同。筍，息允反，馬云：“箁，箬也。”徐云：“竹子。竹為席。於貧反。”紛，孚雲反。漆音七，徐七利反。綬音受。}越玉五重，陳寶，\footnote{於東西序坐北，列玉五重，又陳先王所寶之器物。○越玉，馬云：“越地所獻玉也。”重，直容反。}赤刀、大訓、弘璧、琬、琰，在西序。\footnote{寶刀，赤刀削。大訓，\CJKunderwave{虞書}典謨。大璧、琬琰之珪為二重。○琬，紆晚反。琰，以冉反。削音笑。}


{\noindent\zhuan\zihao{6}\fzbyks 傳“西廂”至“質飾”。正義曰：下傳云“西房，西夾坐東”，“東房，東廂夾室”，然則“房”與“夾室”實同而異名。天子之室有左右房，房即室也。以其夾中央之大室，故謂之“夾室”。此坐在西廂夾室之前，故系“夾室”言之。\CJKunderwave{釋草}云：“筍，竹萌。”孫炎曰：“竹初萌生謂之筍。”是“筍”為蒻竹，取筍竹之皮以為席也。“紛”則組之小別。\CJKunderline{鄭玄}\CJKunderwave{周禮}注云：“紛如綬，有文而狹者也。”然則紛、綬一物，小大異名，故傳以“玄紛”為黑綬。鄭於此注云：“以玄組為之緣。”\CJKunderwave{周禮·大宗伯}云：“以飲食之禮，親宗族兄弟。”\CJKunderline{鄭玄}云：“親者，使之相親。人君有食宗族飲酒之禮，所以親之也。\CJKunderwave{文王世子}云：‘族食，世降一等。’”是天子有與親屬私宴之事。以骨肉情親，不事華麗，故席幾質飾也。 \par}

{\noindent\zhuan\zihao{6}\fzbyks 傳“於東”至“器物”。正義曰：此經為下總目，下復分別言之越訓於也。“越”訓於也。“於”者，於其處所。上云“西序東向”、“東序西向”，則序旁已有王之坐矣。下句陳玉復云“在西序”、“在東序”者，明於東西序坐北也。“序”者牆之別名，其牆南北長，坐北猶有序牆，故言“在西序”、“在東序”也。西序二重，東序三重，二序共為列玉五重。又陳先王所寶之器物，河圖、大訓、貝、鼓、戈、弓皆是先王之寶器也。 \par}

{\noindent\zhuan\zihao{6}\fzbyks 傳“寶刀”至“二重”。正義曰：上言“陳寶”，非寶則不得陳之，故知“赤刀”為寶刀也。謂之“赤刀”者，其刀必有赤處。刀一名削，故名赤刃削也。\CJKunderwave{禮記·少儀}記執物授人之儀云：“刀授穎，削授拊。”\CJKunderline{鄭玄}云：“避用時也。穎，鐶也。拊謂把也。”然則刀施鐶，削用把。削似小於刀，相對為異，散文則通，故傳以“赤刀”為“赤刀削”。\CJKunderwave{吳錄}稱吳人嚴白虎聚眾反，遣弟興詣孫策，策引白削斫席，興體動，曰：“我見刃為然。”然赤刃為赤削，白刃為白削，是削為刀之別名明矣。\CJKunderwave{周禮·考工記}云：“築氏為削，合六而成規。”鄭注云：“曲刃刀也。”又云：“赤刀者，\CJKunderline{武王}誅紂時刀，赤為飾，周正色。”不知其言何所出也。“大訓，\CJKunderwave{虞書}典謨”\CJKunderline{王肅}亦以為然，鄭云“大訓謂禮法，先王德教”，皆是以意言耳。“弘”訓大也。“大璧、琬琰之圭為二重”，則琬琰共為一重。\CJKunderwave{周禮·典瑞}云“琬圭以治德，琰圭以易行”，則琬琰別玉而共為重者，蓋以其玉形質同，故不別為重也。\CJKunderwave{考工記}琬圭、琰圭皆九寸。\CJKunderline{鄭玄}云：“大璧、琬、琰皆度尺二寸者。”孔既不分為二重，亦不知何所據也。 \par}

大玉、夷玉、天球、河圖，在東序。\footnote{三玉為三重。夷，常也。球，雍州所貢。河圖，八卦。伏犧王天下，龍馬出河，遂則其文以畫八卦，謂之河圖，及典謨皆歷代傳寶之。○夷玉，馬云：“東夷之美玉。”\CJKunderwave{說文}夷玉即珣玗琪。球音求，馬云：“玉磬。”雍，於用反，本亦作邕。}胤之舞衣、大貝、鼖鼓，在西房。\footnote{胤國所為舞者之衣，皆中法。大貝，如車渠。鼖鼓長八尺，商周傳寶之。西房，西夾坐東。○鼖，扶雲反，注同。中，丁仲反。車，尺遮反。車渠，車輈也。}兌之戈、和之弓、垂之竹矢，在東房。\footnote{兌、和，古之巧人。垂，舜\CJKunderline{共工}。所為皆中法，故亦傳寶之。東房，東廂夾室。○兌，徒外反。共音恭。}


{\noindent\zhuan\zihao{6}\fzbyks 傳“三玉”至“寶之”。正義曰：“三玉為三重”，與上共為五重也。“夷，常”，\CJKunderwave{釋詁}文。\CJKunderwave{禹貢}雍州所貢球、琳、琅玕，知球是雍州所貢也。常玉、天球傳不解“常”、“天”之義，未審孔意如何。\CJKunderline{王肅}云：“夷玉，東夷之美玉。天球，玉磬也。”亦不解稱天之意。\CJKunderline{鄭玄}云：“大玉，華山之球也。夷玉，東北之珣玕琪也。天球，雍州所貢之玉，色如天者。皆璞，未見琢治，故不以禮器名之。”\CJKunderwave{釋地}云：“東方之美者，有醫無閭之珣玕琪焉。”東方實有此玉。鄭以夷玉為彼玉，未知經意為然否。“河圖，八卦。是伏羲氏王天下，龍馬出河，遂則其文以畫八卦，謂之河圖”，當孔之時，必有書為此說也。\CJKunderwave{漢書·五行志}：“劉歆以為伏犧氏繼天而王，受河圖，則而畫之，八卦是也。”劉歆亦如孔說，是必有書明矣。\CJKunderwave{易·繫辭}云：“古者包犧氏之王天下也，仰則觀象於天，俯則觀法於地，觀鳥獸之文與地之宜，近取諸身，遠取諸物，於是始作八卦。”都不言法河圖也。而此傳言“河圖”者，蓋\CJKunderwave{易}理寬弘，無所不法，直如\CJKunderwave{繫辭}之言，所法已自多矣，亦何妨更法河圖也。且\CJKunderwave{繫辭}又云：“河出圖，洛出書，聖人則之。”若八卦不則河圖，餘復何所則也？\CJKunderline{王肅}亦云：“河圖，八卦也。”璧，玉人之所貴，是為可寶之物。八卦、典謨非金玉之類，嫌其非寶，故云“河圖及典謨皆歷代傳寶之”。此西序、東序各陳四物，皆是臨時處置，未必別有他義。下二房各有二物，亦應無別意也。 \par}

{\noindent\zhuan\zihao{6}\fzbyks 傳“胤國”至“坐東”。正義曰：以夏有\CJKunderline{胤侯}，知“胤”是國名也。胤是前代之國，舞衣至今猶在，明其所為中法，故常寶之。亦不知舞者之衣是何衣也。大貝必大於餘貝。\CJKunderline{伏生}\CJKunderwave{書傳}云：“\CJKunderline{散宜生}生之江淮，取大貝，如大車之渠。”是言大小如車渠也。\CJKunderwave{考工記}謂車罔為渠。大小如車罔，其貝形曲如車罔，故比之也。\CJKunderwave{考工記}云：“鼓長八尺，謂之鼖鼓。”\CJKunderwave{釋樂}云：“大鼓謂之鼖。”此鼓必有所異。周興至此未久，當是先代之器，故云“商周傳寶之”。西序即是西夾，西夾之前已有南向坐矣，西序亦陳之寶，近在此坐之西，知此“在西房”者，在西夾室東也。 \par}

{\noindent\zhuan\zihao{6}\fzbyks 傳“兌和”至“夾室”。正義曰：戈、弓、竹矢，巧人所作。垂是巧人，知兌、和亦古之巧人也。垂，舜\CJKunderline{共工}，\CJKunderwave{舜典}文。若不中法，即不足可寶，知“所為皆中法，故亦傳寶之”。垂是舜之\CJKunderline{共工}，竹矢蓋舜時之物。其兌、和之所作，則不知寶來幾何世也，故皆言“傳寶之”耳。東夾室無坐，故直言“東廂夾室”，陳於夾室之前也。案鄭注\CJKunderwave{周禮}，宗廟、路寢制如明堂。明堂則五室，此路寢得有東房、西房者，\CJKunderwave{鄭志}張逸以此問，鄭答云：“\CJKunderline{成王}崩在鎬京。鎬京宮室因文武，更不改作，故同諸侯之制，有左右旁也。”孔無明說，或與鄭異，路寢之制不必同明堂也。 \par}

大輅在賓階面，綴輅在阼階面，\footnote{大輅，玉。綴輅，金。面，前。皆南向。○阼，才故反。向，許亮反。}先輅在左塾之前，次輅在右塾之前。\footnote{先輅，象。次輅，木。金、玉、象皆以飾車，木則無飾，皆在路寢門內，左右塾前北面。凡所陳列，皆象\CJKunderline{成王}生時華國之事，所以重顧命。○塾音孰，一音育。重，直用反。}

{\noindent\zhuan\zihao{6}\fzbyks 傳“大輅”至“南向”。正義曰：\CJKunderwave{周禮}巾車“掌王之五輅”，玉輅、金輅、象輅、革輅、木輅，是為五輅也。此經所陳四輅必是\CJKunderwave{周禮}五輅之四。“大輅”，輅之最大，故知大輅玉輅也。“綴輅”，系綴於下，必是玉輅之次，故為金輅也。“面前”者，據人在堂上，面向南方，知面前皆南向，謂轅向南也。地道尊右，故玉輅在西，金輅在東。 \par}

{\noindent\zhuan\zihao{6}\fzbyks 傳“先輅”至“顧命”。正義曰：此經四輅兩兩相配，上言“大輅”、“綴輅”，此言“先輅”、“次輅”，二者各自以前後為文。五輅金即次象，故言“先輅，象”。其木輅在象輅之下，故云“次輅，木”也。又解四輅之名，“金、玉、象皆以飾車”，三者以飾為之名；“木則無飾”，故指木為名耳。\CJKunderline{鄭玄}\CJKunderwave{周禮}注云“革輅，挽之以革而漆之”，“木輅，不挽以革，漆之而已”，以直漆其木，故以“木”為名。木輅之上猶有革輅，不以“次輅”為革輅者，禮五輅而此四輅，於五之內必將少一，蓋以革輅是兵戎之用，於此不必陳之，故不雲革輅而以木輅為次。\CJKunderline{馬融}、\CJKunderline{王肅}皆云：“不陳戎輅者，兵事非常，故不陳之。”孔意或當然也。\CJKunderline{鄭玄}以“綴、次是從後之言，二者皆為副貳之車。先輅是象輅也，綴輅是玉輅之貳，次輅是象輅之貳，不陳金輅、革輅、木輅者，主於朝祀而已”。未知孔、鄭誰得經旨。\CJKunderline{成王}殯在路寢，下云“二人執惠，立於畢門之內”，畢門是路寢之門，知此陳設車輅“皆在路寢門內”也。\CJKunderwave{釋宮}云：“門側之堂謂之塾。”孫炎曰：“夾門堂也。”塾前陳車，必以轅向堂，故知“左右塾前皆北面”也。左塾者謂門內之西，右塾者門內之東，故以北面言之為左右。所陳坐位、器物皆以西為上，由王殯在西序故也。其執兵宿衛之人則先東而後西者，以王在東宿，衛敬新王故也。顧氏云：“先輅在左塾之前，在寢門內之西，北面對玉輅。次輅在右塾之前，在寢門內之東，對金輅也。”凡所陳列自“狄設黼扆”已下至此，皆象\CJKunderline{成王}生時華國之事，所以重顧命也。\CJKunderline{鄭玄}亦云：“陳寶者，方有大事以華國也。”\CJKunderwave{周禮·典路}云：“若有大祭祀，則出路。大喪大賓客亦如之。”是大喪出輅為常禮也。 \par}

{\noindent\shu\zihao{5}\fzkt “牖間”至“漆仍幾”。正義曰：“牖”謂窗也，“間”者窗東戶西戶牖之間也。\CJKunderwave{周禮·司几筵}云：“凡大朝覲、大饗射，凡封國命諸侯，王位設黼扆。扆前南向，設莞筵紛純，加繅席畫純，加次席黼純，左右玉幾。”彼所設者即此坐也。又云：“戶牖之間謂之扆。”彼言“扆前”，此言“牖間”，即一坐也。彼言“次席黼純”，此言“篾席黼純”，亦一物也。\CJKunderwave{周禮}天子之席三重，諸侯之席再重，則此四坐所言敷重席者，其席皆敷三重。舉其上席而言“重”，知其下更有席也。此牖間之坐即是\CJKunderwave{周禮}扆前之坐，篾席之下二重，其次是繅席畫純，其下是莞筵紛純也。此一坐有\CJKunderwave{周禮}可據，知其下二席必然。下文三坐，\CJKunderwave{禮}無其事，以扆前一坐敷三種之席，知下三坐必非一重之席敷三重，但不知其下二重是何席耳。\CJKunderwave{周禮}天子左右幾，諸侯惟右幾，此言“仍幾”，則四坐皆左右幾也。\CJKunderline{鄭玄}云：“左右有幾，優至尊也”。 \par}

二人雀弁，執惠，立於畢門之內。\footnote{士衛殯與在廟同，故雀韋弁。惠，三隅矛。路寢門，一名畢門。○弁，皮彥反，徐扶變反。}四人綦弁，執戈上刃，夾兩階\xpinyin*{戺}。\footnote{綦，文鹿子皮弁。亦士。堂廉曰戺，士所立處。○綦音其，馬本作騏，云：“青黑色。”夾，徐工洽反。戺音俟，徐音士。廉，力佔反，稜也。}


{\noindent\zhuan\zihao{6}\fzbyks 傳“士衛”至“畢門”。正義曰：士入廟助祭，乃服雀弁，於此服雀弁者，士衛王殯，與在廟同，故爵韋弁也。\CJKunderline{鄭玄}云：“赤黑曰雀，言如雀頭色也，雀弁制如冕，黑色，但無藻耳。”然則雀弁所用當與冕同。阮諶\CJKunderwave{三禮圖}云：“雀弁以三十升布為之。”此傳言“雀韋弁”者，蓋以\CJKunderwave{周禮·司服}云“凡兵事，韋弁服”，此人執兵，宜以韋為之，異於祭服，故言“雀韋弁”。下云“綦弁”，孔言鹿子皮為弁，然則下言冕執兵者，不可以韋為冕，未知孔意如何。天子五門，皋、庫、雉、應、路也。下云“王出在應門之內”，出畢門始至應門之內，知畢門即是路寢之門，一名畢門也。此經所陳七種之兵，惟戈經傳多言之，\CJKunderwave{考工記}有其形制，其餘皆無文。傳惟言“惠，三隅矛”，銳亦矛也，“戣、瞿皆戟屬”，不知何所據也。“劉，鉞屬”者，以“劉”與“鉞”相對，故言“屬”以似之，而別又不知何以為異。古今兵器名異體殊，此等形制，皆不可得而知也。\CJKunderline{鄭玄}云：“惠狀蓋斜刃，宜芟刈。戈即今之句孑戟。劉蓋今鑱斧。鉞，大斧。戣、瞿蓋今三鋒矛。銳，矛屬。凡此七兵，或施矜，或著柄。\CJKunderwave{周禮}戈長六尺六寸，其餘未聞長短之數。”\CJKunderline{王肅}惟云：“皆兵器之名也。” \par}

{\noindent\zhuan\zihao{6}\fzbyks 傳“綦文”至“立處”。正義曰：\CJKunderline{鄭玄}云“青黑曰綦”，\CJKunderline{王肅}云“綦，赤黑色”，孔以為“綦，文攏蘗子皮弁”，各以意言，無正文也。大夫則服冕，此服弁，知“亦士”也。“堂廉曰戺”，相傳為然。“廉”者，稜也。所立在堂下，近於堂稜。 \par}

一人冕，執劉，立於東堂。一人冕，執鉞，立於西堂。\footnote{冕，皆大夫也。劉，鉞屬。立於東西廂之前堂。}一人冕，執戣，立於東垂。一人冕，執瞿，立於西垂。\footnote{戣、瞿皆戟屬。立於東西下之階上。○戣音逵。瞿，其俱反，徐音懼。}一人冕，執銳，立於側階。\footnote{銳，矛屬也。側階，北下立階上。○銳，以稅反。}

{\noindent\zhuan\zihao{6}\fzbyks 傳“冕皆”至“前堂”。正義曰：\CJKunderwave{周禮·司服}云：“大夫之服，自玄冕而下。”知服冕者皆大夫也。\CJKunderline{鄭玄}云：“序內半以前曰堂。”謂序內簷下，自室壁至於堂廉，中半以前總名為“堂”。此立於東堂、西堂者，當在東西廂近階而立，以備升階之人也。 \par}

{\noindent\zhuan\zihao{6}\fzbyks 傳“戣瞿”至“階上”。正義曰：\CJKunderwave{釋詁}云：“疆、界、邊、衛、圉、垂也。”則“垂”是遠外之名。此經所言“冕”則在堂上，“弁”則在堂下，此二人服冕，知在堂上也。堂上而言“東垂”、“西垂”，知在堂上之遠地。堂之遠地，當於序外東廂西廂，必有階上堂，知此立於東西堂之階上也。 \par}

{\noindent\zhuan\zihao{6}\fzbyks 傳“銳矛”至“階上”。正義曰：鄭、王皆以側階為東下階也。然立於東垂者已在東下階上，何由此人復共並立？故傳以為“北下階上”，謂堂北階，北階則惟堂北一階而已。“側”猶特也。 \par}

{\noindent\shu\zihao{5}\fzkt “二人”至“側階”。正義曰：\CJKunderwave{禮}大夫服冕，士服弁也。此所執者凡有七兵，立於畢門之內，及夾兩階立堂下者，服爵弁、綦弁者，皆士也。以其去殯遠，故使士為之。其在堂上服冕者，皆大夫也。以其去殯近，皆使大夫為之。先門，次階，次堂，從外向內而敘之也。次東西垂，次側階，又從近向遠而敘之也。在門者兩，守門兩廂各一人，故“二人”。在階者，兩廂各二人，故“四人”。\CJKunderwave{禮記·明堂位}：“三公在中階之前。”\CJKunderwave{考工記}：“夏后氏世室，九階。”\CJKunderline{鄭玄}云：“南面三，三面各二。”\CJKunderline{鄭玄}又云，宗廟及路寢制如明堂，則路寢南面亦當有三階矣。此惟四人夾兩階，不守中階者，路寢制如明堂，惟\CJKunderline{鄭玄}之說耳，路寢三階不書，亦未有明文，縱有中階，中階無人升降，不須以兵衛之。 \par}

王麻冕黼裳，由賓階隮。\footnote{王及群臣皆吉服，用西階升，不敢當主。}卿士、邦君麻冕蟻裳，入即位。\footnote{公卿大夫及諸侯皆同服，亦廟中之禮。蟻，裳名，色玄。○蟻，魚綺反。}太保、太史、太宗皆麻冕彤裳。\footnote{執事各異裳。彤，纁也。太宗，上宗，即宗伯也。}


{\noindent\zhuan\zihao{6}\fzbyks 傳“王及”至“當主”。正義曰：\CJKunderwave{禮}績麻三十升以為冕，故稱“麻冕”。傳嫌麻非吉服，故言“王及群臣皆吉服”也。“王麻冕”者，蓋袞冕也。\CJKunderwave{周禮·司服}：“享先王則袞冕。”此禮授王冊命，進酒祭王，且袞是王之上服，於此正王之尊，明其服必袞冕也。其卿士、邦君當各以命服，服即助祭之冕矣。“袞”，\CJKunderline{鄭玄}\CJKunderwave{周禮注}云：“袞之衣五章，裳四章。”則袞衣之裳，非獨有黼。言“黼裳”者，以裳之章色，黼黻有文，故特取為文。\CJKunderwave{詩·采菽}之篇言王賜諸侯云：“玄袞及黼。”以黼有文，故特言之。\CJKunderline{鄭玄}於此注云：“黼裳者，冕服有文者也。”是言貴文故稱之。\CJKunderwave{禮}“君升阼階”，此用西階升者，以未受顧命，不敢當主也。 \par}

{\noindent\zhuan\zihao{6}\fzbyks 傳“公卿”至“色玄”。正義曰：“卿士”，卿之有事者，公則卿兼之。此行大禮，大夫亦與焉。略舉“卿士”為文，公與大夫必在，故傳言“公卿大夫及諸侯皆同服”，言同服吉服，此“亦廟中之禮”也。言其如助祭，各服其冕服也。\CJKunderwave{禮}無“蟻裳”，今云“蟻”者，裳之名也。“蟻”者，蚍蜉蟲也，此蟲色黑，知蟻裳色玄，以色玄如蟻，故以蟻名之。\CJKunderwave{禮}祭服皆玄衣纁裳，此獨雲玄裳者，卿士、邦君於此無事，不可全與祭同，改其裳以示變於常也。太保、太史有所主者，則純如祭服，暫從吉也。“入即位”者，\CJKunderline{鄭玄}云：“卿西面，諸侯北面。”\CJKunderline{鄭玄}惟據經“卿士、邦君”言之，其公亦北面，孤東面也。 \par}

{\noindent\zhuan\zihao{6}\fzbyks 傳“執事”至“宗伯”。正義曰：此三官者皆執事，俱“彤裳”，而言“各異裳”者，各自異於卿士、邦君也。“彤”，赤也。\CJKunderwave{禮}祭服纁裳。纁是赤色之淺者，故以“彤”為纁，言是常祭服也。“太宗”與下文“上宗”一人，即宗伯之卿也。 \par}

太保承介圭，上宗奉同、瑁，由阼階隮。\footnote{大圭尺二寸，天子守之，故奉以奠\CJKunderline{康王}所位。同，爵名。瑁，所以冒諸侯圭，以齊瑞信，方四寸，邪刻之。用阼階升，由便不嫌。○冒，莫報反。}太史秉書，由賓階隮，御王冊命。\footnote{太史持冊書顧命進\CJKunderline{康王}，故同階。}

{\noindent\zhuan\zihao{6}\fzbyks 傳“大圭”至“不嫌”。正義曰：\CJKunderwave{考工記·玉人}云：“鎮圭尺有二寸，天子守之。”鎮圭，圭之大者。“介”訓大也，故知是彼鎮圭。天子之所守，故奉之以奠\CJKunderline{康王}所位，以明正位為天子也。\CJKunderwave{禮}又有“大圭長三尺”，知“介圭”非彼三尺圭者，\CJKunderwave{典瑞}云：“王搢大圭，執鎮圭以朝日。”\CJKunderwave{玉人}云：“大圭長三尺，天子服之。”彼搢於紳帶是天子之笏，不是天子所守，故知非幣妖尺之大圭也。“上宗奉同、瑁”，則下文雲天子“受同、瑁”。大保必奠於位，其奉介圭，下文不言“受介圭”者，以同、瑁並在手中，故不得執之，太保必奠於其位，但文不見耳。\CJKunderwave{禮}於奠爵無名“同”者，但下文祭酢皆用同奉酒，知“同”是酒爵之名也。\CJKunderwave{玉人}云：“天子執冒四寸以朝諸侯。”\CJKunderline{鄭玄}注云：“名玉曰冒者，言德能覆蓋天下也。四寸者，方。以尊接卑，以小為貴。”\CJKunderwave{禮}天子所以執瑁者，諸侯即位，天子賜之以命圭，圭頭邪銳，其瑁當下邪刻之，其刻闊狹長短如圭頭。諸侯來朝，執圭以授天子，天子以冒之刻處冒彼圭頭，若大小相當，則是本所賜；其或不同，則圭是偽作，知諸侯信與不信。故天子執瑁，所以冒諸侯之圭以齊瑞信，猶今之合符。然經傳惟言圭之長短，不言闊狹。瑁方四寸，容彼圭頭，則圭頭之闊無四寸也。天子以一瑁冒天下之圭，則公侯伯之圭闊狹等也。此瑁惟冒圭耳，不得冒璧。璧亦稱瑞，不知所以齊信，未得而聞之也。“阼階”者，東階也。謂之“阼”者，\CJKunderline{鄭玄}\CJKunderwave{士冠禮}注云“阼猶酢也。東階所以答酢賓客”，是其義也。\CJKunderwave{禮}凶事設洗於西階西南，吉事設洗於東階東南。此太保、上宗皆行吉事，盥洗在東，故用阼階升，由便，以卑不嫌為主人也。\CJKunderline{鄭玄}云：“上宗猶太宗，變其文者，宗伯之長，大宗伯一人，與小宗伯二人，凡三人，使其上二人也，一人奉同，一人奉瑁。”傳無明解，當同於鄭也。 \par}

{\noindent\zhuan\zihao{6}\fzbyks 傳“太史”至“同階”。正義曰：訓“御”為進。太史持策書顧命欲以進王，故與王同升西階。\CJKunderline{鄭玄}云：“御猶向也。王此時正立賓階上少東，太史東面於殯西南而讀策書，以命王嗣位之事。”孔雖以“御”為進，其意當如鄭言。不言王面北，可知也。篇以“顧命”為名，指上文為言。顧命策書，稟王之意為言，亦是顧命之事，故傳言“策書顧命”。 \par}

{\noindent\shu\zihao{5}\fzkt “王麻”至“冊命”。正義曰：此將傳顧命，佈設位次，即上所作法度也。凡諸行禮，皆賤者先至，此必卿、下士、邦君即位既定，然後王始升階。但以君臣之序，先言王服，因服之下即言升階，從省文。卿士、邦君無所執事，故直言“即位”而已。太保、太史、太宗皆執事之人，故別言衣服。各自所職，不得即言升階，故別言所執，各從升階為文次也。卿士王臣，故先於邦君。太史乃是大宗之屬，而先於太宗者，太史之職掌冊書，此禮主以為冊命，太史所掌事重，故先言之。 \par}

曰：“皇后憑玉幾，道揚末命，命汝嗣訓，\footnote{冊命之辭。大君，\CJKunderline{成王}。言憑玉幾所道，稱揚終命，所以感動\CJKunderline{康王}。命汝繼嗣其道，言任重，因以託戒。○憑，皮冰反。}臨君周邦，率循大卞，\footnote{用是道臨君周國，率群臣循大法。卞，皮彥反，徐扶變反。}燮和天下，用答揚\CJKunderline{文}、\CJKunderline{武}之光訓。”\footnote{言用和道和天下，用對揚聖祖文武之大教。敘\CJKunderline{成王}意。}王再拜，興，答曰：“眇眇予末小子，其能而亂四方，以敬忌天威？”\footnote{言微微我淺末小子，其能如父祖治四方，以敬忌天威德乎？謙辭，託不能。○眇，彌小反。}


{\noindent\zhuan\zihao{6}\fzbyks 傳“冊命”至“託戒”。正義曰：言“憑玉幾所道”，以示不憑玉幾則不能言，所以感動\CJKunderline{康王}，令其哀而聽之，不敢忽也。以“訓”為道，命汝繼嗣其道，繼父道為天下之主。言所任者重，因以託戒也。 \par}

{\noindent\zhuan\zihao{6}\fzbyks 傳“用是”至“大法”。正義曰：“卞”之為法，無正訓也。告以為法之道，令率群臣循之，明所循者法也，故以“大卞”為大法。\CJKunderline{王肅}亦同也。 \par}

{\noindent\shu\zihao{5}\fzkt “曰皇”至“光訓”。正義曰：此即丁卯命作之冊書也。誥\CJKunderline{康王}曰：“大君\CJKunderline{成王}病困之時，憑玉幾所道，稱揚將終之教命。命汝繼嗣其道，代為民主。用是道以臨君周邦，率群臣循大法，用和道和天下，用對揚聖祖文武之大教。”敘\CJKunderline{成王}之意，言\CJKunderline{成王}命汝如此也。 \par}

乃受同、瑁，王三宿，三祭，三吒。\footnote{王受瑁為主，受同以祭。禮成於三，故酌者實三爵於王，王三進爵，三祭酒，三奠爵，告已受群臣所傳顧命。○吒，陟嫁反，字亦作宅，又音妒，徐又音託，又豬夜反。\CJKunderwave{說文}作詫，丁故反，奠爵也。馬作詫，與\CJKunderwave{說文}音義同。}上宗曰：“饗！”\footnote{祭必受福，贊王曰：“饗福酒。”}太保受同，降，\footnote{受王所饗同，下堂反於篚。}盥，以異同，秉璋以酢。\footnote{太保以盥手洗異同，實酒秉璋以酢祭。半圭曰璋，臣所奉。王已祭，太保又祭。報祭曰酢。○酢，才各反。}


{\noindent\zhuan\zihao{6}\fzbyks 傳“王受”至“顧命”。正義曰：天子執瑁，故“受瑁為主”。“同”是酒器，故“受同以祭”。\CJKunderline{鄭玄}云：“王既對神，則一手受同，一手受瑁。”然既受之後，王受同而祭，則瑁以授人。“禮成於三，酌者實三爵於王”，當是實三爵而續送，三祭各用一同，非一同而三反也。\CJKunderwave{釋詁}云：“肅，進也。”“宿”即肅也，故以宿爵而續送。祭各用一同為一進，“三宿”謂三進爵，從立處而三進至神所也。三祭酒，三酹酒於神坐也。每一酹酒，則一奠爵，三奠爵於地也。為此祭者，告神言已已受群臣所傳顧命，白神使知也。經典無此“吒”字，“吒”為奠爵，傳記無文。正以既祭必當奠爵，既言“三祭”，知“三吒”為三奠爵也。\CJKunderline{王肅}亦以“吒”為奠爵。\CJKunderline{鄭玄}云：“徐行前曰肅，卻行曰吒。王徐行前三祭，又三卻複本位。”與孔異也。 \par}

{\noindent\zhuan\zihao{6}\fzbyks 傳“祭必”至“福酒”。正義曰：\CJKunderwave{禮}於祭末必飲神之酒，受神之福，其大祭則有受嘏之福。\CJKunderwave{禮·特牲}、\CJKunderwave{少牢}主人受嘏福，是受神之福也。其告祭小祀，則不得備儀，直飲酒而已。此非大祭，故於王三奠爵訖，上宗以同酌酒進王，贊王曰：“饗福酒也。”王取同嚌之，乃以同授太保也。 \par}

{\noindent\zhuan\zihao{6}\fzbyks 傳“受王”至“於篚”。正義曰：上宗贊王以饗福酒也，即云“太保受同”，明是“受王所饗同”也。祭祀飲酒之禮，爵未用皆實於篚，既飲皆反於篚，知此“下堂反於篚”也。 \par}

{\noindent\zhuan\zihao{6}\fzbyks 傳“太保”至“曰酢”。正義曰：祭祀以變為敬，不可即用王同，故太保以盥手更洗異同，實酒於同中，乃秉璋以酢祭。於上祭後更復報祭，猶如正祭大禮之亞獻也。\CJKunderwave{周禮·典瑞}云：“四圭有邸，以祀天。兩圭有邸，以祀地。圭璧以祀日月，璋邸射以祀山川。”從上而下，遞減其半，知“半圭曰璋”。\CJKunderwave{祭統}云：“君執圭瓚，大宗執璋瓚。”謂亞獻用璋瓚，此非正祭，亦是亞獻之類，故亦執璋。若助祭，公侯伯子男自得執圭璧也。秉璋以酢，是報祭之事。王已祭，太保又報祭也。“酢”訓報也，故“報祭曰酢”。飲酒之禮稱“獻酢”者，亦是報之義也。 \par}

授宗人同，拜，王答拜。\footnote{宗人，小宗伯，佐大宗伯。太宗供王，宗人供太保。拜白已傳顧命，故授宗人同。拜，王答拜，尊所受命。}太保受同，祭，嚌，\footnote{太宗既拜而祭，既祭，受福。嚌至齒，則王亦至齒。王言饗，太保言嚌，互相備。○嚌，才細反。互音戶。}宅，授宗人同，拜，王答拜。\footnote{太保居其所，授宗人同，拜白\CJKunderline{成王}以事畢。王答拜，敬所白。○宅如字，馬同，徐殆故反。}太保降，收。\footnote{太保下堂，則王下可知。有司於此盡收徹。○徹，醜列反，又徐直列反。}

{\noindent\zhuan\zihao{6}\fzbyks 傳“宗人”至“受命”。正義曰：“上宗”為太宗伯，知“宗人”為小宗伯也。太保所以拜者，白\CJKunderline{成王}言已已傳顧命訖也。將欲拜，故先授宗人同。拜者，自為拜神，不拜\CJKunderline{康王}，但白神言已傳顧命之事，先告王已受顧命。王答拜者，尊所受之命，亦告神使知，故答拜也。王既祭，則奠同於地。太保不敢奠於地，故以同授宗人，然後拜也。太保既酢祭而拜，則王之奠爵，每奠必拜。於王祭不言“拜”者，祭酒必拜，乃是常禮。於王不言“拜”，於太保言“拜”者，足以見王拜也。 \par}

{\noindent\zhuan\zihao{6}\fzbyks 傳“太保”至“相備”。正義曰：“太保受同”者，謂太保既拜之後，於宗人邊受前所授之同，而進以祭神。既祭神之後，遂更受福酒，嚌以至齒。禮之通例，啐入口，是嚌至於齒，示飲而實不飲也。太保報王之祭，事與王祭禮同，而史錄其事，二文不等，故傳辨其意，於太保言“嚌至齒”，則王饗福酒，亦嚌至齒也；於王言“上宗曰饗”，則太保亦應有宗人曰饗，二文不同，互見以相備。 \par}

{\noindent\zhuan\zihao{6}\fzbyks 傳“太保”至“所白”。正義曰：“宅”訓居也。太保居其所，於受福酒之處，足不移，為將拜，故授宗人同。祭祀既畢而更拜者，白\CJKunderline{成王}以事畢也。既拜白\CJKunderline{成王}以傳顧命事畢，則王受顧命亦畢，王答拜，敬所白也。 \par}

{\noindent\shu\zihao{5}\fzkt “乃受”至“降收”。正義曰：王受冊命之時，立於西階上少東，北面。太史於柩西南，東面讀策書。讀冊既訖，王再拜。上宗於王西南，北面奉同、瑁以授王。王一手受同，一手受瑁。王又以瑁授宗人。王乃執同，就樽於兩楹之間，酌酒，乃於殯東西面立，三進於神坐前，祭神如前祭。凡前祭酒酹地而奠爵訖，復位。再拜，王又於樽所別以同酌酒，祭神如前。復三祭，故云“三宿、三祭、三吒”。然後酌福酒以授王，上宗贊王曰：“饗福酒。”王再拜，受酒，跪而祭。先嚌至齒，興，再拜。太保受同，降自東階，反於篚。又盥以異同，執璋升自東階，適樽所酌酒，至殯東西面報祭之。欲祭之時，授宗人同，拜白王柩云：“已傳顧命訖。”王則答拜。拜柩，尊所受命。太保乃於宗人處受同，祭柩如王禮，但一祭而已。祭訖，乃受福。祝酌同以授太保，宗人贊太保曰：“饗福酒。”太保再拜，受同，亦祭，先而嚌至齒，與，再拜訖，於所居位授宗人同。太保更拜白柩以事畢。王又答拜，拜柩，敬所白。王與太保降階而下堂，有司於是收徹器物。 \par}

諸侯出廟門俟。\footnote{言諸侯，則卿士已下亦可知。殯之所處，故曰廟。待王后命。○處，昌呂反。}

{\noindent\shu\zihao{5}\fzkt “諸侯出廟門俟”。正義曰：“廟門”謂路寢門也。出門待王后命，即作後篇。後篇云，二伯率諸侯入應門，則諸侯之出應門之外，非出廟門而已。以其在廟行事，事畢出於廟門，不言出廟門即止也。 \par}

%%% Local Variables:
%%% mode: latex
%%% TeX-engine: xetex
%%% TeX-master: "../Main"
%%% End:

%% -*- coding: utf-8 -*-
%% Time-stamp: <Chen Wang: 2024-04-02 11:42:41>

% {\noindent \zhu \zihao{5} \fzbyks } -> 注 (△ ○)
% {\noindent \shu \zihao{5} \fzkt } -> 疏

\chapter{卷十九}


\section{康王之誥第二十五（顧命下）}


康王既屍天子，\footnote{屍，主也，主天子之正號。○馬本此句上更有“\CJKunderline{成王}崩”三字。}遂誥諸侯，作\CJKunderwave{康王之誥}。\footnote{既受顧命，群臣陳戒，遂報誥之。因事曰遂。}

康王之誥\footnote{求諸侯之見匡弼。}

{\noindent\shu\zihao{5}\fzkt “康王既”至“之誥”。正義曰：康王既受顧命，主天子之位，群臣進戒於王，王遂報誥諸侯。史敘其事，作\CJKunderwave{康王之誥}。伏生以此篇合於\CJKunderwave{顧命}，共為一篇，後人知其不可，分而為二。馬、鄭、王本此篇自“高祖寡命”已上內於\CJKunderwave{顧命}之篇，“王若曰”已下始為\CJKunderwave{康王之誥}，諸侯告王，王報誥諸侯，而使告報異篇，失其義也。 \par}

王出，在應門之內，\footnote{出畢門，立應門內之中庭，南面。}太保率西方諸侯，入應門左，\CJKunderline{畢公}率東方諸侯，入應門右，\footnote{二公為二伯，各率其所掌諸侯，隨其方為位，皆北面。}皆布乘黃朱。\footnote{諸侯皆陳四黃馬朱鬛以為庭實。○乘音繩證反。鬛，力輒反。}


{\noindent\zhuan\zihao{6}\fzbyks 傳“出畢”至“南面”。正義曰：出在門內，不言“王坐”，諸侯既拜，王即答拜，復不言“興”，知立庭中南面也。 \par}

{\noindent\zhuan\zihao{6}\fzbyks 傳“二公”至“北面”。正義曰：二公率領諸侯，知其“為二伯，各率其所掌諸侯”。\CJKunderwave{曲禮}所謂“職方”者，此之義也。王肅云：“\CJKunderline{畢公}代\CJKunderline{周公}為東伯，故率東方諸侯。”然則\CJKunderline{畢公}是太師也。當太師之名，在太保之上，此先言太保者，於時太保領冢宰，相王室，任重，故先言西方。若使東伯任重，亦當先言東方。北面，以東為右，西為左，入左入右隨其方為位。嫌東西相向，故云“皆北面”，將拜王，明北面也。 \par}

{\noindent\zhuan\zihao{6}\fzbyks 傳“諸侯”至“庭實”。正義曰：諸侯朝見天子，必獻國之所有，以表忠敬之心，故“諸侯皆陳四黃馬朱鬛以為庭實”，言實之於王庭也。四馬曰“乘”，言“乘黃”正是馬色黃矣。“黃”下言“朱”，“朱”非馬色。定十年\CJKunderwave{左傳}云：“宋公子地有白馬四，公嬖向魋，魋欲之。公取而朱其尾、鬛以與之。”是古人貴朱鬛。知“朱”者，朱其尾、鬛也。於時諸侯其數必眾，眾國皆陳四馬，則非王庭所容。諸侯各有所獻，必當少陳之也。案\CJKunderwave{周禮·小行人}云：“合六幣，圭以馬，璋以皮，璧以帛，琮以錦，琥以繡，璜以黼。此六物者，以和諸侯之好。”\CJKunderline{鄭玄}云：“六幣所以享也，五等諸侯享天子用壁，享後用琮。用圭璋者，二王之後也。”如鄭彼言，則諸侯之享天子，惟二王之後用馬。此雲皆陳馬者，下雲“奉圭兼幣”，幣即馬是也，圭是文馬之物。鄭云：“此幣圭以馬，蓋舉王者之後以言耳。諸侯當璧以帛，亦有庭實。”然則此陳馬者是二王之後享王物也。獨取此物以總表諸侯之意，故云諸侯皆陳馬也。圭亦享王之物，下言“奉圭”，此不陳圭者，圭奉以文命，不陳之也。案\CJKunderwave{覲禮}諸侯享天子“馬卓上，九馬隨之”。此用乘黃者，因喪禮而行朝，故略之。 \par}

賓稱奉圭兼幣，曰：“一二臣衛，敢執壤奠。”\footnote{賓，諸侯也。舉奉圭兼幣之辭，言“一二”，見非一也。為蕃衛，故曰“臣衛”。來朝而遇國喪，遂因見新王，敢執壤地所出而奠贄也。○壤，如丈反。見，賢遍反，下同。蕃，方袁反。朝，直遙反。喪，息浪反。贄音至。}皆再拜稽首。王義嗣德，答拜。\footnote{諸侯拜送幣而首至地，盡禮也。康王以義繼先人明德，答其拜，受其幣。○盡，子忍反。}

{\noindent\zhuan\zihao{6}\fzbyks 傳“賓諸”至“奠贄”。正義曰：天子於諸侯有不純臣之義，故以諸侯為賓。“稱”訓舉也。“舉奉圭兼幣之辭”，以圭幣奉王而為之作辭。辭出一人之口而言“一二”者，見諸侯同為此意，意非一人也。\CJKunderline{鄭玄}雲“釋辭者一人，其餘奠幣，拜者稽首而已”是也。言“衛”者，諸侯之在四方，皆為天子蕃衛，故曰“臣衛”。此時\CJKunderline{成王}始崩，即得有諸侯在京師者，來朝而遇國喪，遂因見新王也。諸侯享天子，其物甚眾，非徒圭馬而已。皆是土地所有，故云“敢執壤地所出而奠贄”也。然“舉奉圭兼幣”，乃是享禮。凡享禮，則每一國事畢，乃更餘國復入，其朝則侯氏總入。故\CJKunderline{鄭玄}注\CJKunderwave{曲禮}云：“春“受贄於朝，受享於廟”，是朝與享別。此既諸侯總入而得有庭實享禮者，以新朝嗣王，因行享禮，故鄭注云：“朝兼享禮也，與常禮不同。” \par}

{\noindent\zhuan\zihao{6}\fzbyks 傳“諸侯”至“其幣”。正義曰：\CJKunderwave{周禮·太祝}“辨九拜，一曰稽首”，施之於極尊，故為“盡禮”也。“義嗣德”三字，史言王答拜之意也。康王先是太子，以義繼先人明德，今為天子，無所嫌，故答其拜，受其幣，自許與諸侯為主也。 \par}

{\noindent\shu\zihao{5}\fzkt “王出”至“答拜”。正義曰：此敘諸侯見新王之事。王出畢門，在應門之內，立於中庭。太保\CJKunderline{召公}為西伯，率西方諸侯，入應門左，立於門內之西廂也。太師\CJKunderline{畢公}為東伯，率東方諸侯，入應門右，立於門內之東廂也。諸侯皆布陳一乘四匹之黃馬朱鬛，以為見新王之庭實。諸侯為王之賓，實共使一人少前進，舉奉圭兼幣之辭，言曰：“一二天子之臣，在外為蕃衛者，敢執土壤所有奠之於庭。”既為此言，乃皆再拜稽首，用盡禮致敬，以正王為天子也。康王先為太子，以義嗣先人明德，不以在喪為嫌，答諸侯之拜，以示受其圭幣，與之為主也。 \par}

太保暨\CJKunderline{芮伯}咸進，相揖，皆再拜稽首，\footnote{冢宰與司徒皆共群臣諸侯並進陳戒。不言諸侯，以內見外。}曰：“敢敬告天子，皇天改大邦殷之命，\footnote{大天改大國殷之王命，謂誅紂也。}惟周文武，誕受羑若，克恤西土。\footnote{言文武大受天道而順之，能憂我西土之民。本其所起。○羑，羊久反，馬云：“道也。”}惟新陟王，畢協賞罰，戡定厥功，用敷遺後人休。\footnote{惟周家新升王位，當盡和天下賞罰，能定其功，用布遺後人之美。言施及子孫無窮。○戡音堪。遺，唯季反，注及下同。施，以豉反。}今王敬之哉！\footnote{敬天道，務崇先人之美。}張皇六師，無壞我高祖寡命。”\footnote{言當張大六師之眾，無壞我高德之祖寡有之教命。○壞音怪。}


{\noindent\zhuan\zihao{6}\fzbyks 傳“冢宰”至“見外”。正義曰：\CJKunderline{召公}為冢宰，\CJKunderline{芮伯}為司徒，司徒位次冢宰，故言“太保與\CJKunderline{芮伯}咸進”。\CJKunderline{芮伯}已下，共告群臣諸侯並皆進也。“相揖”者，揖之使俱進也。大保揖群臣，群臣又報揖太保，故言“相揖”。動足然後相揖，故“相揖”之文在“咸進”之下。 \par}

{\noindent\zhuan\zihao{6}\fzbyks 傳“言文”至“所起”。正義曰：“羑”聲近猷，故訓之為道。王肅云：“羑，道也。”文武所憂，非憂西土而已，特言“能憂西土之民”，本其初起於西土故也。 \par}

{\noindent\zhuan\zihao{6}\fzbyks 傳“言當”至“教命”。正義曰：“皇”訓大也。國之大事，在於強兵，故令張大六師之眾。“高德之祖”，謂\CJKunderline{文王}也。王肅云：“美\CJKunderline{文王}少有及之，故曰‘寡有’也。” \par}

{\noindent\shu\zihao{5}\fzkt “太保”至“寡命”。正義曰：太保\CJKunderline{召公}與司徒\CJKunderline{芮伯}皆共諸侯並進，相顧而揖，乃並再拜稽首，起而言曰：“敢告天子，大天改大國殷之王命，誅殺殷紂。惟周家\CJKunderline{文王}、\CJKunderline{武王}大受天道而順之，能憂我西土之民，以此王有天下。惟我周家新升王位，當盡和天下賞罰，戡定其為王之功用，布遺後人之美，將使施及子孫，無有窮盡之期。今王新即王位，其敬之哉！當張大我之六師，令國常強盛，無令傾壞我高祖寡有之命。”戒王使繼先王之業也。 \par}

王若曰：“庶邦侯、甸、男、衛，\footnote{順其戒而告之，不言群臣，以外見內。○馬本從此以下為\CJKunderwave{康王之誥}，又云：“與\CJKunderwave{顧命}差異。敘歐陽、大小夏侯同為\CJKunderwave{顧命}。”}惟予一人釗報誥。\footnote{報其戒。}昔君文武丕平富，不務咎，\footnote{言先君文武道大，政化平美，不務咎惡。}厎至齊\footnote{馬讀絕句。}信，用昭明於天下。\footnote{致行至中信之道，用顯明於天下。言聖德治。○厎，之履反。}則亦有熊羆之士，不二心之臣，保乂王家。\footnote{言文武既聖，則亦有勇猛如熊羆之士，忠一不二心之臣，共安治王家。○熊音雄。羆，彼皮反。}用端命於上帝，皇天用訓厥道，付畀四方。\footnote{君聖臣良，用受端直之命於上天。大天用順其道，付與四方之國，王天下。○畀，必利反，徐甫至反。王，於況反。}乃命建侯樹屏，在我後之人。\footnote{言文武乃施政令，立諸侯，樹以為蕃屏，傳王業在我後之人。謂子孫。}今予一二伯父，尚胥暨顧，綏爾先公之臣，服於先王。\footnote{天子稱同姓諸侯曰伯父。言今我一二伯父，庶幾相與顧念文武之道，安汝先公之臣，服於先王而法循之。}雖爾身在外，乃心罔不在王室，\footnote{言雖汝身在外士之為諸侯，汝心常當忠篤，無不在王室。“熊羆之士”，勵朝臣，此督諸侯。○督，丁木反。}用奉恤厥若，無遺鞠子羞。”\footnote{當各用心奉憂其所行順道，無自荒怠，遺我稚子之羞辱。稚子，康王自謂也。○鞠，居六反。}


{\noindent\zhuan\zihao{6}\fzbyks 傳“順其”至“見內”。正義曰：群臣戒王使勤，王又戒之使輔己，是順其事而告之也。上文太保、\CJKunderline{芮伯}進言，不言諸侯，以內見外。此王告庶邦，不言朝臣，以外見內，欲令互相備也。周制六服，此惟四服，不言採、要者，略舉其事。猶\CJKunderwave{武成}雲“甸、侯、衛，駿奔走”，亦略舉之矣。“予一人釗”。正義曰：\CJKunderwave{禮}天子自稱予一人，不言名。此王自稱名者，新即王位，謙也。 \par}

{\noindent\zhuan\zihao{6}\fzbyks 傳“言先”至“咎惡”。正義曰：孔以“富”為美，故云“政化平美”。不務咎惡於人，言哀矜下民，不用刑罰之。王肅雲“文武道大，天下以平，萬民以富”是也。 \par}

{\noindent\zhuan\zihao{6}\fzbyks 傳“致行”至“德洽”。正義曰：孔以“齊”為中，致行中正誠信之道。王肅云：“立大中之道也。” \par}

{\noindent\zhuan\zihao{6}\fzbyks 傳“天子”至“循之”。正義曰：\CJKunderwave{覲禮}言天子呼諸侯之禮云：“同姓大國則曰伯父，其異姓則曰伯舅，同姓小邦則曰叔父，其異姓則曰叔舅。”計此時諸侯多矣，獨雲“伯父”，舉同姓大國言之也。諸侯先公以臣道服於先王，其事有法，故令安汝先公之用臣，服於先王，以臣之道而法循之。 \par}

{\noindent\zhuan\zihao{6}\fzbyks 傳“言雖”至“諸侯”。正義曰：王之此誥，並誥群臣諸侯，但互相發見，其言不備。言先王有熊羆之士，勵朝臣使用力如先世之臣也。此言汝身在外土，心念王室，督諸侯使然。 \par}

{\noindent\shu\zihao{5}\fzkt “王若”至“子羞”。正義曰：群臣諸侯既進戒王，王順其戒呼而告之曰：“眾邦在侯、甸、男、衛諸服內之國君，惟我一人釗報誥卿士群公。昔先君\CJKunderline{文王}、\CJKunderline{武王}其道甚大，政化平美，專以美道教化，不務咎惡於人，致行至美中正誠信之道，用是顯明於天下”。言聖道博洽也。“文武既聖，時臣亦賢，則亦有如熊如羆之勇士，不二心之忠臣，共安治王家。以君聖臣良之故，用能受端直之命於上天。大天用順其道，付與四方之國，使文武受此諸國，王有天下。”言文武得賢臣之力也。“文武以得臣力之故，乃施政令，封立賢臣為諸侯者，樹之以為蕃屏，令屏衛在我後之人”。先王所立諸侯，即今諸侯之祖，故舉先世之事以告今之諸侯。“今我一二伯父，庶幾相與顧念文武之道，安汝先公之用臣，服於先王之道而法循之，亦當以忠誠輔我天子。雖汝身在外土為國君，汝心常當無有不在王室，當各用心奉憂其所行順道，無自荒怠，以遺我稚子之羞辱”。“稚子”，康王自謂。戒令匡弼己也。 \par}

群公既皆聽命，相揖趨出。\footnote{已聽誥命，趨出罷退，諸侯歸國，朝臣就次。}王釋冕，反喪服。\footnote{脫去黼冕，反服喪服，居倚盧。○去，羌呂反。}

{\noindent\shu\zihao{5}\fzkt “群公”至“喪服”。正義曰：“群公”總謂朝臣與諸侯也。\CJKunderline{鄭玄}云：“群公主為諸侯與王之三公，諸臣亦在焉。王釋冕，反喪服，朝臣諸侯亦反喪服。\CJKunderwave{禮·喪服}篇臣為君，諸侯為天子，皆斬衰。” \par}

\section{畢命第二十六【偽】}


康王命作冊畢，\footnote{命為冊書，以命\CJKunderline{畢公}。}分居里，成周郊，\footnote{分別民之居里，異其善惡。成定東周郊境，使有保護。○別，彼列反。}作\CJKunderwave{畢命}。

畢命\footnote{言\CJKunderline{畢公}見命之書。}


{\noindent\zhuan\zihao{6}\fzbyks 傳“命為”至“\CJKunderline{畢公}”。正義曰：\CJKunderwave{周禮·內史}云：“凡命諸侯及孤卿大夫，則策命之。”此雲“命作冊”者，命內史為冊書以命\CJKunderline{畢公}，故云以冊命\CJKunderline{畢公}。 \par}

{\noindent\zhuan\zihao{6}\fzbyks 傳“分別”至“保護”。正義曰：殷之頑民，遷居此邑，歷世化之，已得純善，恐其變改，故更命\CJKunderline{畢公}分別民之居里，異其善惡。即經所云“旌別淑慝，表厥宅裡，彰善癉惡，樹之風聲”，“殊厥井疆，俾克畏慕”皆是也。“分”者令其善惡分別，使惡者慕善，非分別其處，使之異居也。此邑本名成周，欲以成就周道。民不純善，則是未成，故命\CJKunderline{畢公}教之。“成定東周郊境”，即經“申畫郊圻，慎固封守”，是其使有保護。 \par}

{\noindent\shu\zihao{5}\fzkt “康王”至“畢命”。正義曰：康王命史官作冊書命\CJKunderline{畢公}，使\CJKunderline{畢公}分別民之居里，令善惡有異。於成周之邑，成定東周之郊境。史敘其事，作\CJKunderwave{畢命}。 \par}

惟十有二年，六月庚午朏，\footnote{康王即位十二年六月三日庚午。○朏，普忽反，徐芳尾反，又芳憒反。}越三日壬申，王朝步自宗周，至於豐。\footnote{於朏三日壬申，王朝行自宗周，至於豐。宗周，鎬京。豐，\CJKunderline{文王}所都。○朝，直遙反。鎬，戶老反。}以成周之眾，命\CJKunderline{畢公}保釐東郊。\footnote{用成周之民眾，命\CJKunderline{畢公}使安理治正成周東郊，令得所。○釐，力之反。治，直吏反，一本作“治政”，則依字讀。令，力呈反。}


{\noindent\zhuan\zihao{6}\fzbyks 傳“康王”至“庚午”。正義曰：漢初不得此篇，有偽作其書以代之者。\CJKunderwave{漢書·律曆志}云：“康王十二年六月戊辰朔，三日庚午，故\CJKunderwave{畢命豐刑}曰‘惟十有二年六月庚午朏，王命作策書\CJKunderwave{豐刑}’。”此偽作者傳聞舊語，得其年月，不得以下之辭，妄言作\CJKunderwave{豐刑}耳，亦不知\CJKunderwave{豐刑}之言何所道也。\CJKunderline{鄭玄}云：“今其逸篇有冊命霍侯之事，不同與此序相應，非也。”\CJKunderline{鄭玄}所見又似異於\CJKunderwave{豐刑}，皆妄作也。\CJKunderwave{說文}云：“朏，月未盛之明也。”此日未有事而記此“庚午朏”者，為下言壬申張本，猶如記朔望與生魄死魄然也。 \par}

{\noindent\shu\zihao{5}\fzkt “惟十”至“東郊”。正義曰：惟康王即位十有二年六月三日庚午，月光朏然而明也。於朏後三日壬申，王早朝行從宗周鎬京，至於豐邑，就\CJKunderline{文王}之廟。以成周之民眾命太師\CJKunderline{畢公}，使安理東郊之民，令得其所。 \par}

王若曰：“嗚呼！父師，惟\CJKunderline{文王}、\CJKunderline{武王}，敷大德於天下，用克受殷命。\footnote{王順其事嘆告\CJKunderline{畢公}代，\CJKunderline{周公}為大師，為東伯，命之代\CJKunderline{君陳}。言文武布大德於天下，故天佑之，用能受殷之王命。○大音泰。}惟\CJKunderline{周公}左右先王，綏定厥家。\footnote{言\CJKunderline{周公}助先王安定其家。}毖殷頑民，遷於洛邑，密邇王室，式化厥訓。\footnote{慎殷頑民，恐其叛亂，故徙於洛邑，密近王室，用化其教。○毖音秘。近如字，又附近之近。}既歷三紀，世變風移，四方無虞，予一人以寧。\footnote{言殷民遷周已經三紀，世代民易，頑者漸化，四方無可度之事，我天子用安矣。十二年曰紀。父子曰世。○度，待洛反，舊作待路反。}道有升降，政由俗革，不臧厥臧，民罔攸勸。\footnote{天道有上下交接之義，政教有用俗改更之理。民之俗善，以善養之。俗有不善，以法御之。若乃不善其善，則民無所勸慕。○上，時掌反。更，古衡反。}


{\noindent\zhuan\zihao{6}\fzbyks 傳“王順”至“王命”。正義曰：\CJKunderline{畢公}代\CJKunderline{周公}為太師，故王呼為“父師”，率東方諸侯，是為“東伯”也。蓋\CJKunderline{君陳}卒，命之使代\CJKunderline{君陳}也。 \par}

{\noindent\zhuan\zihao{6}\fzbyks 傳“言周”至“其家”。正義曰：\CJKunderwave{釋詁}云：“左、右，助也。”言\CJKunderline{周公}助先王安定其家。伐殷之時，\CJKunderline{周公}已有其功，復能遷殷頑民，言其功之多也。 \par}

{\noindent\zhuan\zihao{6}\fzbyks 傳“言殷”至“曰世”。正義曰：\CJKunderline{周公}以攝政七年營成周，\CJKunderline{成王}元年遷殷頑民，\CJKunderline{成王}在位之年雖未知，其實當在三十左右，至今應三十六年，是殷民遷周已歷三紀。十二年者，天之大數。歲星、太歲皆十二年而一周天，故“十二年曰紀”。父子易人為世。\CJKunderwave{大禹謨}云：“賞延於世。”謂緣父及子也。 \par}

{\noindent\zhuan\zihao{6}\fzbyks 傳“天道”至“勸慕”。正義曰：天氣下降，地氣上騰，而有寒暑生焉。刑新國用輕典，刑亂國用重典，輕重隨俗而有寬猛異焉。天道有上下交接之義，故寒暑易節。政教有用俗改更之理，故寬猛相濟。天道有寒暑遞來，政教以寬猛相濟。民之風俗，善惡無常，或善變為惡，或惡變為善，不可以其既善，謂善必不變。民之俗善，須以善養之，令善遂不變。人之俗有不善，當以善法御之，使變而為善。若乃不善其善，則下民無所勸慕。民無所慕，則變為惡矣。殷民今雖已善，更當以善教之。欲以屈\CJKunderline{畢公}之意。 \par}

惟公懋德，克勤小物，弼亮四世，正色率下，罔不祗師言。\footnote{言公勉行德，能勤小物，輔佐文、武、成、康，四世為公卿，正色率下，下人無不敬仰師法。○懋音茂。}嘉績多於先王，予小子垂拱仰成。”\footnote{公之善功多大先人之美。我小子為王，垂拱仰公成理。言其上顯父兄，下施子孫。○拱，九勇反。仰如字，徐五亮反。}

{\noindent\zhuan\zihao{6}\fzbyks 傳“言公”至“師法”。正義曰：“小物”猶小事也，能勤小事，則大事必能勤矣，故舉能勤小事以為\CJKunderline{畢公}之善。\CJKunderwave{釋詁}云：“亮，佐也。”\CJKunderwave{晉語}說\CJKunderline{文王}之事云，“詢於八虞，訪於辛、尹，重之以周、召、畢、榮”，則\CJKunderline{畢公}於\CJKunderline{文王}之世已為大臣，是“輔佐文、武、成、康四世為公卿”也。“正色”謂嚴其顏色，不惰慢，不阿諂。以此率下，下民無不敬仰師法之。 \par}

{\noindent\zhuan\zihao{6}\fzbyks 傳“公之”至“子孫”。正義曰：先王之功，無由可及。言公之善功多大先人之美，方欲委之以事，盛言之，重其功美矣。 \par}

{\noindent\shu\zihao{5}\fzkt “王若”至“仰成”。正義曰：康王順其事嘆而呼\CJKunderline{畢公}曰：“嗚呼！父師，惟\CJKunderline{文王}、\CJKunderline{武王}布大德於天下，用此能受殷之王命，代殷為天子。惟\CJKunderline{周公}佐助先王，安定其家。慎彼殷之頑民，恐其或有叛逆，故遷於洛邑，令之北近王室，用使化其教訓。自爾已來，既歷三紀，人世既變，風俗亦移，四方無可度之事，我天子一人用是而得安寧。但天道有上下交接之義，政教有用俗改更之理。今日雖善，或變為惡，若不善其善，則民無所勸慕。更須選賢教之，舉善勸之，宜此任者，莫先於公。惟公勉力行德，能勤小事，輔佐四世，正色率下，無有不敬仰師法公言者。公之善功多於先王，我小子垂衣拱手，仰公成理。”將欲任之，故盛稱其德也。 \par}

王曰：“嗚呼！父師，今予祗命公以\CJKunderline{周公}之事，往哉！\footnote{今我敬命公以\CJKunderline{周公}所為之事，往為之哉！言非\CJKunderline{周公}所為，不敢枉公往治。○治，直吏反。}旌別淑慝，表厥宅裡，彰善癉惡，樹之風聲。\footnote{言當識別頑民之善惡，表異其居里，明其為善，病其為惡，立其善風，揚其善聲。○別音彼列反。癉音丁但反。}弗率訓典，殊厥井疆，俾克畏慕。\footnote{其不循教道之常，則殊其井居田界，使能畏為惡之禍，慕為善之福，所以沮勸。○俾，必爾反。沮，辭汝反，又慈呂反。}申畫郊圻，慎固封守，以康四海。\footnote{郊圻雖舊所規畫，當重分明之。又當謹慎堅固封疆之守備，以安四海。京圻安，則四海安矣。○守，徐始救反。重，直用反。}政貴有恆，辭尚體要，不惟好異。\footnote{政以仁義為常，辭以理實為要，故貴尚之。若異於先王，君子所不好。○好，呼報反。}商俗靡靡，利口惟賢，餘風未殄，公其念哉！\footnote{紂以靡靡利口為賢，覆亡國家。今殷民利口餘風未絕，公其念絕之。○覆，芳服反。}


{\noindent\zhuan\zihao{6}\fzbyks 傳“言當”至“善聲”。正義曰：旌旗所以表識貴賤，故傳以“旌”為識。“淑”，善也。“慝”，惡也。言當識別頑民之善，惡知其善者，表異其所居之裡，若今孝子、順孫、義夫、節婦，表其門閭者也。表其善者，則惡者自見。明其為善，當褒賞之。病其為惡，當罪罰之。其有善人，立其善風，令邑里使放效之；揚其善聲，告之疏遠，使聞知之。 \par}

{\noindent\zhuan\zihao{6}\fzbyks 傳“其不”至“沮勸”。正義曰：\CJKunderwave{孟子}云，“方里為井，井九百畝”，使民“死徙無出鄉，鄉田同井，出入相友，守望相助，疾病相扶持，則百姓親睦”。然則先王制之為井田也，欲使民相親愛，生相佐助，死相殯葬。不循道教之常者，其人不可親近，與善民雜居，或染善為惡，故殊其井田居界，令民不與來往。猶今下民有大罪過不肯服者，則擯出族黨之外，吉凶不與交通，此之義也。亦既殊其井田，必當思自改悔，使其能畏為惡之禍，慕為善之福，所以沮止為惡者，勸勉為善者。 \par}

{\noindent\zhuan\zihao{6}\fzbyks 傳“郊圻”至“安矣”。正義曰：“郊圻”謂邑之境界。境界雖舊有規畫，而年世久遠，或相侵奪，當重分明畫之，以防後相侵犯。雖舉邑之郊境為言，其民田疆畔亦令更重畫之，不然何以得“殊其井疆”也？王城之立，四郊以為京師屏障，預備不虞，又當謹慎牢固封疆之守備，以安四海之內。此是王之近郊，牢設守備，惟可以安京師耳。而云“安四海”者，京師安，則四海安矣。 \par}

{\noindent\zhuan\zihao{6}\fzbyks 傳“紂以”至“絕之”。正義曰：韓宣子稱，紂使師延作靡靡之樂。“靡靡”者，相隨順之意。紂之為人，拒諫飾非，惡聞其短，惟以靡靡相隨順、利口捷給、能隨從上意者以之為賢。商人效之，遂成風俗，由此所以覆亡國家。殷民利口餘風，至今不絕，公其念絕之。欲令其變惡俗也。 \par}

{\noindent\shu\zihao{5}\fzkt “王曰”至“念哉”。正義曰：王更嘆而呼\CJKunderline{畢公}曰：“嗚呼！父師，今日我敬命公以\CJKunderline{周公}所為之事，公其往為之哉！公往至彼，當識別善之與惡，表異其善者所居之裡，彰明其為善，病其為惡。其為善之人，當立其善風，揚其善聲。其有不循道教之常者，則殊其井田疆界，使之能畏為惡之禍，慕為善之福。更重畫郊圻境界，謹慎牢固其封疆守備，以安彼四海之內。為政貴在有常，言辭尚其禮實要約，當不惟好其奇異。商之舊俗，靡靡然好相隨順，利口辯捷、阿諛順旨者惟以為賢。餘風至今未絕，公其念絕之哉！”戒\CJKunderline{畢公}以治殷民之法。 \par}

我聞曰：‘世祿之家，鮮克由禮，以蕩陵德，實悖天道。\footnote{特言我聞自古有之，世有祿位而無禮教，少不以放蕩陵邈有德者，如此實亂天道。○鮮，息淺反。悖，布內反。}敝化奢麗，萬世同流。’\footnote{言敝俗相化，車服奢麗，雖相去萬世，若同一流。○敝，步寐反。}茲殷庶士，席寵惟舊，怙侈滅義，服美於人。\footnote{此殷眾士，居寵日久，怙恃奢侈，以滅德義。服飾過制，美於其民。言僣上。○怙音戶。}驕淫矜侉，將由惡終。雖收放心，閒之惟艱。\footnote{言殷眾士驕恣過制，矜其所能，以自侉大，如此不變，將用惡自終。雖今順從周制，心未厭服，以禮閒御，其心惟難。○侉音苦瓜反。壓，於葉反，又於甲反，又於豔反。}資富能訓，惟以永年。惟德惟義，時乃大訓。不由古訓，於何其訓？”\footnote{以富資而能順義，則惟可以長年命矣。惟有德義，是乃大順。若不用古訓典籍，於何其能順乎？}王曰：“嗚呼！父師，邦之安危，惟茲殷士，不剛不柔，厥德允修。\footnote{言邦國所以安危，惟在和此殷士而已。治之不剛不柔，寬猛相濟，則其德政信修立。}惟\CJKunderline{周公}克慎厥始，惟\CJKunderline{君陳}克和厥中，惟公克成厥終。\footnote{\CJKunderline{周公}遷殷頑民以消亂階，能慎其始。\CJKunderline{君陳}弘\CJKunderline{周公}之訓，能和其中。\CJKunderline{畢公}闡二公之烈，能成其終。}三後協心，同底於道，道洽政治，澤潤生民。\footnote{三君合心為一，終始相成，同致於道。道至普洽，政化治理，其德澤惠施，乃浸潤生民。言三君之功，不可不尚。○治，直吏反。施，始鼓反。浸，子鴆反。}四夷左衽，罔不咸賴，予小子永膺多福。\footnote{言東夷、西戎、南蠻、北狄被髮左衽之人，無不皆恃賴三君之德，我小子亦長受其多福。○衽，而甚反，又而鴆反。}公其惟時成周，建無窮之基，亦有無窮之聞。\footnote{公其惟以是成周之治，為周家立無窮之基業，於公亦有無窮之名，以聞於後世。○為，於偽反。}子孫訓其成式，惟乂。\footnote{言後世子孫順公之成法，惟以治。}嗚呼！罔曰弗克，惟既厥心。\footnote{人之為政，無曰不能，惟在盡其心而已。}罔曰民寡，惟慎厥事。\footnote{無曰人少不足治也，惟在慎其政事，無敢輕之。○少，詩照反。}欽若先王成烈，以休於前政。”\footnote{敬順文武成業，以美於前人之政。所以勉\CJKunderline{畢公}。}


{\noindent\zhuan\zihao{6}\fzbyks 傳“特言”至“天道”。正義曰：凡以善言教化，無非古之訓典，於此特言“我聞”者，言此事自古有之，所以尢須嚴禁故也。世有祿位，財多勢重，縱恣其心而無禮教，如此之人，少能不以放蕩之心陵邈有德者。天道以上臨下，以善率惡，今乃以下慢上，以惡陵善，如此者實亂天道也。 \par}

{\noindent\zhuan\zihao{6}\fzbyks 傳“此殷”至“僣上”。正義曰：“席”者人之所處，故為居之義。“舊”，久也。殷士多是世貴之家，故為“居寵日久”。怙恃己之奢侈，自謂奢侈為賢，德義廢而不行，故為“以滅德義”。又以人輕位卑，美服盛飾，是“服飾過制度”。“美於其人”，言僣上服，服勝人也。 \par}

{\noindent\zhuan\zihao{6}\fzbyks 傳“言殷”至“惟難”。正義曰：“淫”訓過也，故為“過制”。強梁者不得其死，好勝者必遇其敵，故矜侉“不變，將用惡自終”。言“雖收放心”，則已收之矣。雖令順從周制，思威自止，故怨猶在，心未壓服，故“以禮閒御其心，惟難”也。“閒”謂防閒御止也。 \par}

{\noindent\zhuan\zihao{6}\fzbyks 傳“敬順”至“\CJKunderline{畢公}”。正義曰：“美於前人之政”，謂光前人之政。所以勉勵\CJKunderline{畢公}。 \par}

{\noindent\shu\zihao{5}\fzkt “我聞”至“其訓”。正義曰：我聞古人言曰：“世有祿位之家，恃富驕恣，少能用禮，以放蕩之心陵邈有德之士，如此者實悖亂天道。敝俗相化，奢侈華麗，雖相去萬世，而共同一流。”此殷之眾士，皆是富貴之家，居處寵勢，惟已久矣。怙恃奢侈，以滅德義。身卑而僣上，飾其服，美於其人。驕恣過制，矜能自侉，行如此不變，將用惡自終。令以法約之，雖收斂其放佚之心，恆防閒之，惟大艱難。資財富足，能順道義，則惟可以長年命矣。惟能用德，惟能行義，是乃為大順德也。若不用古之訓典，則於何其能順乎？欲令\CJKunderline{畢公}以古之訓典教殷民也。 \par}

\section{君牙第二十七【偽】}


穆王命君牙為周大司徒，\footnote{穆王，康王孫，昭王子。○穆王，名滿。君牙，或作君雅。}作\CJKunderwave{君牙}。\footnote{君牙，臣名。}

君牙\footnote{命以其名，遂以名篇。}

王若曰：“嗚呼！君牙，\footnote{順其事而嘆，稱其名而命之。}惟乃祖乃父，世篤忠貞，服勞王家，厥有成績，紀於太常。\footnote{言汝父祖，世厚忠貞，服事勤勞王家，其有成功，見紀錄書於王之太常，以表顯之。王之旌旗畫日月曰太常。○畫，胡卦反。}惟予小子，嗣守文、武、成、康遺緒，亦惟先正之臣，克左右亂四方。\footnote{惟我小子，繼守先王遺業，亦惟父祖之臣，能佐助我治四方。言己無所能。}心之憂危，若蹈虎尾，涉於春冰。\footnote{言祖業之大，己才之弱，故心懷危懼。虎尾畏噬，春冰畏陷，危懼之甚。○蹈，徒報反。噬，市制反。陷，陷沒之陷。}


{\noindent\zhuan\zihao{6}\fzbyks 傳“言汝”至“太常”。正義曰：\CJKunderwave{周禮·司勳}云：“凡有功者，銘書於王之太常，祭於大烝。”\CJKunderline{鄭玄}云：“銘之言名也。生則書於王旌，以識其人與其功也。死則於烝先王祭之。”是有功者書於王之太常，以表顯之也。\CJKunderwave{周禮·司常}云：“日月為常。”王建太常，是王之旌旗晝日月名之曰太常也。 \par}

{\noindent\shu\zihao{5}\fzkt “穆王”至“君牙”。正義曰：穆王命其臣名君牙者為周大司徒之卿，以策書命之。史錄其策書，作\CJKunderwave{君牙}。 \par}

今命爾予翼，作股肱心膂。\footnote{今命汝為我輔翼股肱心體之臣。言委任。○膂音旅。}纘乃舊服，無忝祖考，弘敷五典，式和民則。\footnote{繼汝先祖故所服，忠勤無辱累祖考之道，大布五常之教，用和民令有法則。○累，劣偽反。令，力呈反。}爾身克正，罔敢弗正，民心罔中，惟爾之中。\footnote{言汝身能正，則下無敢不正。民心無中，從汝取中。必當正身示民以中正。}夏暑雨，小民惟曰怨諮。\footnote{夏月暑雨，天之常道，小人惟曰怨嘆諮嗟。言心無中也。}冬祁寒，小民亦惟曰怨諮。\footnote{冬大寒，亦天之常道，民猶怨諮。}厥惟艱哉！思其艱以圖其易，民乃寧。\footnote{天不可怨，民猶怨嗟，治民其惟難哉！當思慮其難以謀其易，民乃安。○易，以豉反。}


{\noindent\zhuan\zihao{6}\fzbyks 傳“今命”至“委任”。正義曰：“股”，足也。“肱”，臂也。“膂”，背也。汝為我輔翼，當如我之身，故舉四支以喻為股肱心體之臣，言委任如身也。傳以“膂”為體，以見四者皆體，非獨“膂”為體也。\CJKunderwave{禮記·緇衣}云：“民以君為心，君以民為體。”此舉四體，今以臣為心者，君臣合體，則亦同心。\CJKunderwave{詩}云：“赳赳武夫，公侯腹心”，是臣亦為君心也。 \par}

{\noindent\zhuan\zihao{6}\fzbyks 傳“冬大”至“怨嗟”。正義曰：傳以“祁”為大，故云“冬大寒”。寒言大，則夏暑雨是大雨，於此言“祁”以見之。上言“暑雨”，此不言“寒雪”者，於上言“雨”以見之，互相備也。 \par}

{\noindent\shu\zihao{5}\fzkt “今命”至“乃寧”。正義曰：王言我以危懼之故，今命汝為大司徒，汝當作我股肱心膂。言將任之如已身也。繼汝先世舊所服行，亦如父祖忠勤，無為不忠，辱累汝祖考。當須大布五常之教，用和天下兆氏，令有法則。凡欲率下，當先正身，汝身能正，則下無敢不正。民心無能中正，惟取汝之中正，汝當正身心以率之。夏月大暑大雨，天之常也，小民惟曰怨恨而諮嗟。冬月大寒，亦天之常也，小民亦惟曰怨恨而諮嗟。天不可怨，民尚怨之，治民欲使無怨，其惟難哉！思慮其難，以謀其易，為治不違道，不逆民，民乃安矣。 \par}

嗚呼！丕顯哉，\CJKunderline{文王}謨！\footnote{嘆\CJKunderline{文王}所謀大顯明。}丕承哉，\CJKunderline{武王}烈！\footnote{言\CJKunderline{武王}業美，大可承奉。}啟佑我後人，咸以正罔缺。\footnote{文武之謀業，大明可承奉，開助我後嗣，皆以正道無邪缺。○缺，若穴反。}爾惟敬明乃訓，用奉若於先王，\footnote{汝惟當敬明汝五教，用奉順於先王之道。}對揚文武之光命，追配於前人。”\footnote{言當答揚文武光明之命，君臣各追配於前令名之人。}


{\noindent\zhuan\zihao{6}\fzbyks 傳“言武”至“承奉”。正義曰：\CJKunderline{文王}未克殷，始謀造周，故美其謀。\CJKunderline{武王}以殺紂功成業就，故美其業。謀則明白可遵，業則功成可奉，故謀言“顯”，烈言“承”。\CJKunderwave{詩·周頌·武}篇曰，“於皇\CJKunderline{武王}，無競維烈”，亦美\CJKunderline{武王}業之大也。 \par}

{\noindent\zhuan\zihao{6}\fzbyks 傳“文武”至“邪缺”。正義曰：文始謀之，武卒成之。文謀大明，武業可奉。言先王以此成功開道佑助我之後人，使我得安其事而奉行之。以正道見其無邪，罔缺失見其周備，故傳言“無邪缺”。 \par}

{\noindent\shu\zihao{5}\fzkt “嗚呼”至“前人”。正義曰：王又嘆言：“嗚呼！大是顯明哉，\CJKunderline{文王}之謀也！大可承奉哉，\CJKunderline{武王}之業也！\CJKunderline{文王}之謀、\CJKunderline{武王}之業，開道佑助我在後之人，皆以正道無邪缺。”言先王之道易可遵也。“汝惟敬明汝之五教，用奉順於先王之道。汝當答揚文武光明之命，追配於前世令名之人”。令其順先王之道，同古之大賢也。 \par}

王若曰：“君牙，乃惟由先正舊典時式，民之治亂在茲。\footnote{汝惟當奉用先正之臣所行故事、舊典、文籍是法，民之治亂在此而已，用之則民治，廢之則民亂。○治，直吏反，下注同。}率乃祖考之攸行，昭乃闢之有乂。”\footnote{言當循汝父祖之所行，明汝君之有治功。○闢，必亦反。}

{\noindent\shu\zihao{5}\fzkt “王若”至“有乂”。正義曰：王順而呼之曰：“君牙，汝為大司徒，惟當奉用先世正官之法，諸臣所行故事舊典，於是法則之，民之治亂在此而已。汝必奉而用之，循汝祖考之所行，明汝君之有治功。”“汝君”，王自謂也。 \par}

\section{冏命第二十八【偽】}


穆王命伯冏為周太僕正，\footnote{伯冏，臣名也。太僕長，太御中大夫。○冏，九永反，字亦作煛。長，誅丈反。}作\CJKunderwave{冏命}。

冏命\footnote{以冏見命名篇。}


{\noindent\zhuan\zihao{6}\fzbyks 傳“伯冏”至“大夫”。正義曰：“正”訓長也。\CJKunderwave{周禮}“太御中大夫”，“太僕下大夫”，孔以此言“太僕正”，則官高於太僕，故以為\CJKunderwave{周禮}太御者，知非\CJKunderwave{周禮}太僕。若是\CJKunderwave{周禮}太僕，則此雲太僕足矣，何須雲“正”乎？且此經雲“命汝作大正，正於群僕”，案\CJKunderwave{周禮}“太馭中大夫”而下，有戎僕、齊僕、道僕、田僕，太御最為長，既稱正於群僕，故以為太御中大夫。且與君同車，最為親近，故\CJKunderwave{春秋}隨侯寵少師以為車右，\CJKunderwave{漢書}文帝愛趙同命之為御。凡御者最為密暱，故此經雲“汝無暱於憸人，充耳目之官”。故以為太御中大夫，掌御玉輅之官。戎僕雖中大夫，以戎事為重，敘在太御之下，故以太御為長。太僕雖掌燕朝，非親近之任，又是下大夫，不得為長。 \par}

{\noindent\shu\zihao{5}\fzkt “穆王”至“冏命”。正義曰：穆王命其臣名伯冏者為周太僕正之官，以策書命之。史錄其策書，作\CJKunderwave{冏命}。 \par}

王若曰：“伯冏，惟予弗克於德，嗣先人宅丕後，\footnote{順其事以命伯冏，言我不能於道德，繼先人居大君之位，人輕任重。}怵惕惟厲，中夜以興，思免厥愆。\footnote{言常悚懼惟危，夜半以起，思所以免其過悔。○怵，敕律反。惕，他歷反。}昔在文武，聰明齊聖，小大之臣，咸懷忠良，\footnote{聰明，視聽遠。齊通，無滯礙。臣雖官有尊卑，無不忠良。○礙，五代反。}其侍御僕從，罔匪正人。\footnote{雖給侍、進御、僕役從官，官雖微，無不用中正之人。○御如字，一音禦。從，才用反，注及下注“侍從”同。}以旦夕承弼厥闢，出入起居，罔有不欽，\footnote{小臣皆良，僕役皆正，以旦夕承輔其君，故君出入起居，無有不敬。}發號施令，罔有不臧。下民祇若，萬邦咸休。\footnote{言文武發號施令，無有不善。下民敬順其命，萬國皆美其化。}


{\noindent\zhuan\zihao{6}\fzbyks 傳“言常”至“過悔”。正義曰：\CJKunderwave{禮記·祭義}云：“春雨露既濡，君子履之，必有怵惕之心。”“怵惕”是心動之名，多憂懼之意也。“厲”訓危也，言常悚懼，惟恐傾危。\CJKunderwave{易}稱“夕惕若厲”，即此義也。 \par}

{\noindent\zhuan\zihao{6}\fzbyks 傳“聰明”至“忠良”。正義曰：聰發於耳，明發於目，故為“視聽遠”也。“齊”訓中也，“聖”訓通也，動必得中，通而先識，是“無滯礙”也。 \par}

{\noindent\shu\zihao{5}\fzkt “王若”至“咸休”。正義曰：王順其事而呼之曰：“伯冏，惟我不能於道德，而繼嗣先人居大君之位。人輕任重，終常悚懼。心內怵惕，惟恐傾危，中夜以起，思望免其愆過。昔在\CJKunderline{文王}、\CJKunderline{武王}，聰無所不聞，明無所不見。齊，中也，每事得中。聖，通也，通知諸事。其身明聖如此，又小大之臣無不皆思忠良，其左右侍御僕從無非中正之人。以旦夕承輔其君，故其君出入起居無有不敬，文武發號施令無有不善。以此之故，下民敬順其命，萬邦皆美其化。”由臣善故也。 \par}

“惟予一人無良，實賴左右前後有位之士，匡其不及，\footnote{惟我一人無善，實恃左右前後有職位之士，匡正其不及。言此責群臣正己。}繩愆糾謬，格其非心，俾克紹先烈。\footnote{言恃左右之臣彈正過誤，檢其非妄之心，使能繼先王之功業。○繩，市陵反。俾，必爾反。}


{\noindent\zhuan\zihao{6}\fzbyks 傳“言恃”至“功業”。正義曰：木不正者，以繩正之。“繩”謂彈正，“糾”謂發舉，有愆過則彈正之，有錯謬則發舉之。“格”謂檢括，其有非理枉妄之心，檢括使妄心不作。臣當如此匡君，使能繼先王之功業。言己無能，責臣使如此也。 \par}

{\noindent\shu\zihao{5}\fzkt “惟予”至“先烈”。正義曰：王言：“惟我一人無善，亦既無知，實恃賴左右前後有職位之臣，匡正其智所不及者。”責群臣使正己也，即言正己之事。“繩其愆過，糾其錯謬，格其非妄之心，心有妄作，則格正之，使能繼先王之功業”。言得臣匡輔，乃可繼世也。 \par}

今予命汝作大正，正於群僕侍御之臣，\footnote{欲其教正群僕，無敢佞偽。}懋乃後德，交修不逮。\footnote{言侍御之臣，無小大親疏，皆當勉汝君為德，更代修進其所不及。○更，古衡反。}慎簡乃僚，無以巧言令色、便辟側媚，其惟吉士。\footnote{當謹慎簡選汝僚屬侍臣，無得用巧言無實、令色無質、便辟足恭、側媚諂諛之人，其惟皆吉良正士。○便，婢綿反。闢，匹亦反，徐扶亦反。足，將住反。諛，徐以朱反。}


{\noindent\zhuan\zihao{6}\fzbyks 傳“欲其”至“佞偽”。正義曰：“作大正”，“正”，長也，作僕官之長。“正於群僕”，令教正之。二“正”義不同也。群僕雖官有小大，皆近天子。近人主者多以諂佞自容，今大僕教正群僕，明使教之無敢佞偽也。案\CJKunderwave{周禮}太馭中大夫掌御玉輅，戎僕中大夫掌御戎車，齊僕下大夫掌馭金輅，道僕上士掌馭象輅，田僕上士掌馭田輅。“群僕”謂此也。 \par}

{\noindent\zhuan\zihao{6}\fzbyks 傳“當謹”至“正士”。正義曰：府史已下，官長所自闢除命，士以上皆應人主自選。此令太僕正謹慎簡選僚屬者，人主所用皆由臣下，臣下銓擬，可者然後用之，故令太僕正慎簡僚屬也。\CJKunderwave{論語}稱：“巧言、令色、足恭，左聰明恥之。”“便辟”是巧言令色之類，知是彼“足恭”也。“巧言”者，巧為言語以順從上意，無情實也。“令色”者，善為顏色以媚說人主，無本質也。“便僻”者，前卻俯仰，以是為恭。“側媚”者，為僻側之事以求媚於君。此等皆是諂諛之人，不可用為近官也。“媚”，愛也。襄三十一年\CJKunderwave{左傳}云，鄭子產謂子皮曰：“誰敢求愛於子？”知此為“側媚”者，為側行以求愛，非是愛前人也。若能愛在上，則忠臣也，不當禁其無用。 \par}

{\noindent\shu\zihao{5}\fzkt “今予”至“吉士”。正義曰：今我命汝作太僕官大正，汝當教正於群僕侍御之臣，勸勉汝君為德，汝與同僚交更修進汝君智所不及之事。汝為僕官之長，當慎簡汝之僚屬，必使皆得正人，無得用巧言令色、便辟側媚之人，其惟皆當用吉良善士。令選其在下屬官，小臣僕隸之等，皆用善人。 \par}

僕臣正，厥後克正。僕臣諛，厥後自聖。\footnote{言僕臣皆正，則其君乃能正。僕臣諂諛，則其君乃自謂聖。}後德惟臣，不德惟臣。\footnote{君之有德，惟臣成之。君之無德，惟臣誤之。言君所行善惡，專在左右。}爾無暱於憸人，充耳目之官，迪上以非先王之典。\footnote{汝無親近於憸利小子之人，充備侍從在視聽之官，道君上以非先王之法。○暱，女乙反。憸，息廉反，徐七漸反，利口也，本亦作𢘼。近附近之近。道，導也。}非人其吉，惟貨其吉。\footnote{若非人其實吉良，惟以貨財配其吉良，以求入於僕侍之臣，汝當清審。}若時瘝厥官，\footnote{若用是行貨之人，則病其官職。○瘝，故頑反。}惟爾大弗克祇厥闢，惟予汝辜。”\footnote{用行貨之人，則惟汝大不能敬其君，惟我則亦以此罪汝。言不忠也。}王曰：“嗚呼，欽哉！永弼乃後於彝憲。”\footnote{嘆而敕之，使敬用所言，當長輔汝君於常法。此穆王庶幾欲蹈行常法。}

\section{呂刑第二十九（呂刑上、下）}


呂命，\footnote{呂侯見命為天子司寇。}穆王訓夏贖刑，\footnote{呂侯以穆王命作書，訓暢夏禹贖刑之法，更從輕以佈告天下。○贖音蜀，注及下同。}作\CJKunderwave{呂刑}。

呂刑\footnote{後為甫侯，故或稱\CJKunderwave{甫刑}。}


{\noindent\zhuan\zihao{6}\fzbyks 傳“呂侯”至“司寇”。正義曰：呂侯得王命，必命為王官。\CJKunderwave{周禮}司寇掌刑，知呂侯見命為天子司寇。\CJKunderline{鄭玄}云：“呂侯受王命，入為三公。”引\CJKunderwave{書說}云：“周穆王以呂侯為相。”\CJKunderwave{書說}謂\CJKunderwave{書緯·刑得放}之篇有此言也。以其言“相”，知為三公。即如鄭言，當以三公領司寇，不然，何以得專王刑也。 \par}

{\noindent\zhuan\zihao{6}\fzbyks 傳“呂侯”至“天下”。正義曰：名篇謂之\CJKunderwave{呂刑}，其經皆言“王曰”，知“呂侯以穆王命作書”也。經言陳罰贖之事，不言何代之禮，故序言“訓夏”，以明經是夏法。王者代相革易，刑罰世輕世重，殷以變夏，周又改殷。夏法行於前代，廢已久矣。今複訓暢夏禹贖刑之法，以周法傷重，更從輕以佈告天下。以其事合於當時，故\CJKunderline{孔子}錄之以為法。經多說治獄之事，是訓釋申暢之也。金作贖刑，\CJKunderline{唐}、\CJKunderline{虞}之法。\CJKunderwave{周禮}職金“掌受士之金罰、貨罰，入於司兵”，則周亦有贖刑。而遠訓夏之贖刑者，\CJKunderwave{周禮}惟言“士之金罰”，人似不得贖罪。縱使亦得贖罪，贖必異於夏法。以夏刑為輕，故祖而用之。罪實則刑之，罪疑則贖之，故當並言贖刑，非是惟訓贖罰也。\CJKunderwave{周禮}“司刑掌五刑之法，以麗萬民之罪。墨罪五百，劓罪五百，宮罪五百，刖罪五百，殺罪五百”。五刑惟有二千五百。此經“五刑之屬三千”，案刑數乃多於\CJKunderwave{周禮}，而言變從輕者，\CJKunderwave{周禮}五刑皆有五百，此則輕刑少而重刑多；此經墨、劓皆千，刖刑五百，宮刑三百，大辟二百，輕刑多而重刑少，變周用夏，是改重從輕也。然則\CJKunderline{周公}聖人，相時製法而使刑罰太重，令穆王改易之者，穆王遠取夏法，殷刑必重於夏。夏承堯舜之後，民淳易治，故制刑近輕。輕則民慢，故殷刑稍重。自湯已後，世漸苛酷，紂作炮烙之刑，明知刑罰益重。周承暴虐之後，不可頓使太輕。雖減之輕，猶重於夏法。成康之間，刑措不用，下及穆王，民猶易治。故呂侯度時制宜，勸王改從夏法。聖人之法非不善也，而不以經遠。呂侯之智非能高也，而法可以適時。苟適於時，事即可為善，亦不言呂侯才高於\CJKunderline{周公}，法勝於前代。所謂觀民設教，遭時制宜，刑罰所以世輕世重，為此故也。 \par}

{\noindent\zhuan\zihao{6}\fzbyks 傳“後為”至“甫刑”。正義曰：\CJKunderwave{禮記}書傳引此篇之言多稱為“\CJKunderwave{甫刑}曰”，故傳解之“後為甫侯，故或稱\CJKunderwave{甫刑}”。知“後為甫侯”者，以\CJKunderwave{詩·大雅·崧高}之篇宣王之詩，雲“生甫及申”；\CJKunderwave{揚之水}為平王之詩，雲“不與我戍甫”，明子孫改封為甫侯。不知因呂國改作甫名？不知別封餘國而為甫號？然子孫封甫，穆王時未有甫名而稱為\CJKunderwave{甫刑}者，後人以子孫之國號名之也。猶若叔虞初封於唐，子孫封晉，而\CJKunderwave{史記}稱\CJKunderwave{晉世家}。然宣王以後，改呂為甫，\CJKunderwave{鄭語}史伯之言幽王之時也，乃雲“申呂雖衰，齊許猶在”，仍得有呂者，以彼史伯論四嶽治水，其齊、許、申、呂是其後也。因上“申呂”之文而云“申呂雖衰”，呂即甫也。 \par}

{\noindent\shu\zihao{5}\fzkt “呂命”至“呂刑”。正義曰：呂侯得穆王之命為天子司寇之卿，穆王於是用呂侯之言，訓暢夏禹贖刑之法。呂侯稱王之命而佈告天下。史錄其事，作\CJKunderwave{呂刑}。 \par}

惟呂命，王享國百年，耄荒，\footnote{言呂侯見命為卿，時穆王以享國百年，耄亂荒忽。穆王即位過四十矣，言百年，大其雖老而能用賢以揚名。○耄，本亦作𦿗，毛報反，\CJKunderwave{切韻}莫報反。}度作刑以詰四方。\footnote{度時世所宜，訓作贖刑，以治天下四方之民。○度，待洛反。注同，馬如字，云：“法度也。”詰，起一反。}


{\noindent\zhuan\zihao{6}\fzbyks 傳“言呂”至“揚名”。正義曰：史述呂侯見命而記王年，知其得命之時王已享國百年也。\CJKunderwave{曲禮}云：“八十九十曰耄。”是“耄荒”為年老精神耄亂荒忽也。穆王即位之時，已年過四十矣，比至命呂侯之年，未必已有百年。言“百年”者，美大其事，雖則年老而能用賢以揚名，故記其百年之耄荒也。\CJKunderwave{周本紀}云：“甫侯言於王，作修刑辟。”是修刑法者皆呂侯之意，美王能用之。穆王即位過四十者，不知出何書也。\CJKunderwave{周本紀}云：“穆王即位，春秋已五十矣”，“立五十五年崩”。司馬遷若在孔後，或當各有所據。\CJKunderwave{無逸}篇言殷之三王及\CJKunderline{文王}享國若干年者，皆謂在位之年。此言“享國百年”，乃從生年而數，意在美王年老能用賢，而言其長壽，故舉從生之年，以“耄荒”接之，美其老之意也。文不害意，不與彼同。 \par}

{\noindent\shu\zihao{5}\fzkt “惟呂”至“四方”。正義曰：惟呂侯見命為卿，於時穆王享有周國已積百年，王精神耄亂而荒忽矣。王雖老耄，猶能用賢，取呂侯之言，度時世所宜，作夏贖刑以治天下四方之民也。 \par}

王曰：“若古有訓，蚩尢惟始作亂，延及於平民，\footnote{順古有遺訓，言蚩尤造始作亂，惡化相易，延及於平善之人。九黎之君號曰蚩尢。○蚩，尺之反；尤，有牛反；馬云：“少昊之末九黎君名。”}罔不寇賊鴟義，奸宄奪攘矯虔。\footnote{平民化之，無不相寇賊，為鴟梟之義。以相奪攘，矯稱上命，若固有之。亂之甚。○鴟，尺之反；鴟梟，惡鳥；馬云：“鴟，輕也。”義，本亦作誼。宄音軌。攘，如羊反。矯，居表反。虔，其然反。}


{\noindent\zhuan\zihao{6}\fzbyks 傳“順古”至“蚩尢”。正義曰：古有遺訓，順而言之，故為“順古有遺訓”也。“蚩尢造始作亂”，其事往前未有，蚩尢今始造之，必是亂民之事，不知造何事也。下說三苗之主習蚩尢之惡，作五虐之刑，此章主說虐刑之事，蚩尢所作，必亦造虐刑也。以峻法治民，民不堪命，故惡化轉相染易，延及於平善之民，亦化為惡也。“九黎之君號曰蚩尢”，當有舊說云然，不知出何書也。\CJKunderwave{史記·五帝本紀}云：“神農氏世衰，諸侯相侵伐，蚩尢最為暴虐，莫能伐之。黃帝乃徵師諸侯，與蚩尢戰於涿鹿之野，遂擒殺蚩尢，而諸侯咸尊軒轅為天子。”如\CJKunderwave{本紀}之言，蚩尢是炎帝之末諸侯名也。應劭云：“蚩尢，古天子。”鄭云：“蚩尢霸天下，黃帝所伐者。”\CJKunderwave{漢書音義}有臣瓚者，引\CJKunderwave{\CJKunderline{孔子}三朝記}云：“蚩尢，庶人之貪者。”諸說不同，未知蚩尢是何人也。\CJKunderwave{楚語}曰：“少昊氏之衰也，九黎亂德，顓頊受之，使復舊常。”則九黎在少昊之末，非蚩尢也。韋昭云：“九黎氏九人，蚩尢之徒也。”韋昭雖以九黎為蚩尢，要\CJKunderwave{史記}蚩尢在炎帝之末，\CJKunderwave{國語}九黎在少昊之末，二者不得同也。“九黎”之文惟出楚語，孔以蚩尢為九黎。下傳又云：“蚩尢黃帝所滅”，言“黃帝所滅”，則與\CJKunderwave{史記}同矣。孔非不見\CJKunderwave{楚語}，而為此說，蓋以蚩尢是九黎之君，黃帝雖滅蚩尢，猶有種類尚在，故下至少昊之末，更復作亂。若其不然，孔意不可知也。\CJKunderline{鄭玄}云：“學蚩尢為亂者，九黎之君，在少昊之代也。”其意以蚩尢當炎帝之末，九黎當少昊之末，九黎學蚩尢，九黎非蚩尢也。 \par}

{\noindent\zhuan\zihao{6}\fzbyks 傳“平民”至“之甚”。正義曰：蚩尢作亂，當是作重刑以亂民。以峻法酷刑，民無所措手足，困於苛虐所酷，人皆苛且，故平民化之，無有不相寇賊。群行攻劫曰“寇”，殺人曰“賊”，言攻殺人以求財也。“鴟梟”，貪殘之鳥。\CJKunderwave{詩}云：“為梟為鴟。”梟是鴟類。\CJKunderline{鄭玄}云：“盜賊狀如鴟梟，鈔掠良善，劫奪人物。”傳言“鴟梟之義”，如鄭說也。\CJKunderwave{釋詁}云：“虔，固也”。“若固有之”，言取得人物，如己自有也。 \par}

苗民弗用靈，制以刑，惟作五虐之刑曰法。\footnote{三苗之君習蚩尢之惡，不用善化民，而制以重刑。惟為五虐之刑，自謂得法。蚩尢黃帝所滅，三苗\CJKunderline{帝堯}所誅，言異世而同惡。}殺戮無辜，爰始淫為劓、刵、椓、黥。\footnote{三苗之主，頑兇若民，敢行虐刑，以殺戮無罪，於是始大為截人耳鼻，椓陰，黥面，以加無辜，故曰“五虐”。○劓，魚器反。刵，徐如志反。椓，丁角反。黥，其京反。}越茲麗刑並制，罔差有辭。\footnote{苗民於此施刑，並制無罪，無差有直辭者。言淫濫。○麗，力馳反。並，必政反。}民興胥漸，泯泯棼棼，罔中於信，以覆詛盟。\footnote{三苗之民瀆於亂政，起相漸化，泯泯為亂，棼棼同惡，皆無中於信義，以反背詛盟之約。○泯，面忍反，徐音民。棼，芳雲反，徐扶雲反。覆，芳服反，徐敷目反。詛，側助反。背音佩。約如字，又於妙反。}


{\noindent\zhuan\zihao{6}\fzbyks 傳“三苗”至“同惡”。正義曰：上說蚩尢之惡，即以“苗民”繼之，知經意言“三苗之君習蚩尢之惡”。“靈”，善也。不用善化民，而制以重刑。學蚩尢制之，用五刑而虐為之，故為“五虐之刑”，不必\CJKunderline{皋陶}五刑之外，別有五也。“曰法”者，述苗民之語，自謂所作得法，欲民行而畏之。如\CJKunderwave{史記}之文，蚩尢黃帝所滅，下句所說“三苗\CJKunderline{帝堯}所誅”，\CJKunderwave{楚語}雲“三苗復九黎之惡”，是“異世而同惡”也。\CJKunderline{鄭玄}以為“苗民即九黎之後。顓頊誅九黎，至其子孫為三國。高辛之衰，又復九黎之惡。堯興，又誅之。堯末，又在朝，舜臣堯又竄之。後\CJKunderline{禹}攝位，又在洞庭逆命，\CJKunderline{禹}又誅之。穆王深惡此族三生凶德，故著其惡而謂之民”。孔惟言“異世同惡”，不言三苗是蚩尢之子孫。韋昭云：“三苗，炎帝之後諸侯\CJKunderline{共工}也。” \par}

{\noindent\zhuan\zihao{6}\fzbyks 傳“三苗”至“五虐”。正義曰：三苗之主，實國君也。頑兇若民，故謂之“苗民”。不於上經為傳者，就此惡行解之，以其頑兇，敢行虐刑，以殺戮無罪。\CJKunderwave{釋詁}云：“淫，大也。”“於是大為截人耳鼻，椓陰，黥面”，苗民為此刑也。“椓陰”即宮刑也。“黥面”即墨刑也。\CJKunderwave{康誥}\CJKunderline{周公}戒\CJKunderline{康叔}雲“無或劓刵人”，即周世有劓刵之刑，非苗民別造此刑也。以加無辜，故曰“五虐”。\CJKunderline{鄭玄}云：“刵，斷耳。劓，截鼻。椓謂椓破陰，黥為羈黥人面。苗民大為此四刑者，言其特深刻，異於\CJKunderline{皋陶}之為。”鄭意蓋謂截耳截鼻多截之，椓陰苦於去勢，黥面甚於墨頟，孔意或亦然也。 \par}

{\noindent\zhuan\zihao{6}\fzbyks 傳“三苗”至“之約”。正義曰：“三苗之民”，謂三苗國內之民也。“瀆”謂慣瀆，苗君久行虐刑，民慣見亂政，習以為常，起相漸化。“泯泯”，相似之意。“棼棼”，擾攘之狀。“泯泯為亂”，習為亂也。“棼棼同惡”，共為惡也。“中”猶當也，“皆無中於信義”，言為行無與信義合者。\CJKunderwave{詩}云：“君子屢盟，亂是用長。”亂世之民，多相盟詛，既無信義，必皆違之，以此無中於信，反背詛盟之約也。 \par}

虐威庶戮，方告無辜於上。上帝監民，罔有馨香，德刑發聞惟腥。\footnote{三苗虐政作威，眾被戮者方方各告無罪於天，天視苗民無有馨香之行，其所以為德刑，發聞惟乃腥臭。○聞音問，又如字，注同。腥音星。行，下孟反。}皇帝哀矜庶戮之不辜，報虐以威，遏絕苗民，無世在下。\footnote{皇帝，\CJKunderline{帝堯}也。哀矜眾被戮者之不辜，乃報為虐者以威，誅遏絕苗民，使無世位在下國也。○君帝，君宜作皇字，\CJKunderline{帝堯}也。遏，於葛反。}

{\noindent\zhuan\zihao{6}\fzbyks 傳“三苗”至“腥臭”。正義曰：“方方各告無罪於上天”，言其處處告也。天矜於下，俯視苗民，無有馨香之行。“馨香”以喻善也。“其所以為德刑”，苗民自謂是德刑者，發聞於外，惟乃皆是腥臭。“腥臭”喻惡也。 \par}

{\noindent\zhuan\zihao{6}\fzbyks 傳“君帝”至“下國”。正義曰：\CJKunderwave{釋詁}云：“皇，君也。”此言“遏絕苗民”，下句即雲“乃命重黎”，重黎是\CJKunderline{帝堯}之事，知此滅苗民亦\CJKunderline{帝堯}也。此滅苗民在堯之初興，使無世位在於下國，而堯之末年，又有竄三苗者，禮天子不滅國，擇立其次賢者。此為五虐之君，自無世位在下，其改立者復得在朝。但此族數生凶德，故歷代每被誅耳。 \par}

{\noindent\shu\zihao{5}\fzkt “王曰”至“在下”。正義曰：呂侯進言於王，使用輕刑。又稱王之言以告天下，說重刑害民之義。王曰：“順古道有遺餘典訓，記法古人之事。昔炎帝之末，有九黎之國君號蚩尢者，惟造始作亂，惡化遞相染易，延及末平善之民。平民化之，亦變為惡，無有不相寇盜，相賊害，為鴟梟之義。鈔掠良善，外奸內宄，劫奪人物，攘竊人財，矯稱上命，以取人財，若己固自有之。然蚩尢之惡已如此矣，至於高辛氏之末，又有三苗之國君，習蚩尢之惡，不肯用善化民，而更制重法。惟作五虐之刑，乃言曰此得法也。殺戮無罪之人，於是始大為四種之刑。刵，截人耳。劓，截人鼻。劅，椓人陰。黥，割人面。苗民於此施刑之時，並制無罪之人。對獄有罪者無辭，無罪者有辭，苗民斷獄，並皆罪之，無差簡有直辭者。言濫及無罪者也。三苗之民，慣瀆亂政，起相漸染，皆化為惡。泯泯為亂，棼棼同惡，小大為惡。民皆巧詐，無有中於信義。以此無中於信，反背詛盟之約，雖有要約，皆違背之。三苗虐政作威，眾被戮者方方各告無罪於上天。上天下視苗民，無有馨香之行。其所以為德刑者，發聞於外，惟乃皆腥臭，無馨香也。君帝\CJKunderline{帝堯}哀矜眾被殺戮者，不以其罪，乃報為暴虐者以威，止絕苗民，使無世位在於下國。”言以刑虐，故滅之也。 \par}

乃命重黎，絕地天通，罔有降格。\footnote{重即羲，黎即和。堯命義和世掌天地四時之官，使人神不擾，各得其序，是謂絕地天通。言天神無有降地，地民不至於天，明不相干。○重，直龍反。黎，力兮反。}群后之逮在下，明明棐常，鰥寡無蓋。\footnote{群后諸侯之逮在下國，皆以明明大道輔行常法，故使鰥寡得所，無有掩蓋。○棐音匪，又芳鬼反。鰥，居頑反。}皇帝清問下民，鰥寡有辭於苗。\footnote{\CJKunderline{帝堯}詳問民患，皆有辭怨於苗民。○清問，馬云：“清，訊。”}德威惟畏，德明惟明。\footnote{言堯監苗民之見怨，則又增修其德，行威則民畏服，明賢則德明人，所以無能名焉。}


{\noindent\zhuan\zihao{6}\fzbyks 傳“重即”至“相干”。正義曰：\CJKunderwave{楚語}云：“宅王問於觀射父曰：‘\CJKunderwave{周書}所謂重黎實使天地不通者，何也？若無然，民將能登天乎？’對曰：‘非此之謂也。古者民神不雜。少昊氏之衰也，九黎亂德，家為巫史，民神同位，禍災荐臻。顓頊受之，乃命南正重司天以屬神，命火正黎司地以屬民，使復舊常，無相侵瀆，是謂絕地天通。其後三苗復九黎之德，堯覆育重黎之後，不忘舊者，使復典之。’”彼言主說此事，而\CJKunderwave{堯典}雲“乃命羲和，欽若昊天”，即所謂育重黎之後，使典之也。以此知“重即羲”也，“黎即和”也。言羲是重之子孫，和是黎之子孫，能不忘祖之舊業，故以“重黎”言之。傳言“堯乃命羲和掌天地四時之官”，\CJKunderwave{堯典}文也。“民神不擾，是謂絕地天通”，\CJKunderwave{楚語}文也。孔惟加“各得其序”一句耳。\CJKunderwave{楚語}又云，司天屬神，司地屬民。令神與天在上，民與地在下，定上下之分，使民神不雜，則祭享有度，災厲不生。經言民神分別之意，故言“罔有降格”。言天神無有降至於地者，謂神不幹民。孔因互文雲地民不有上至於天者，言民不幹神也。乃總之雲“明不相干”，即是民神不雜也。“地民”或作“地祇”。學者多聞神祇，又“民”字似“祇”，因妄改使謬耳。如\CJKunderwave{楚語}雲“乃命重黎”，是顓頊命之。\CJKunderline{鄭玄}以“‘皇帝哀矜庶戮之不辜’至‘罔有降格’，皆說顓頊之事。乃命重黎即是命重黎之身，非羲和也。‘皇帝清問’以下乃說堯事。顓頊與堯再誅苗民，故上言‘遏絕苗民’，下雲‘有辭於苗’，異代別時，非一事也”。案\CJKunderwave{楚語}雲“少昊氏之衰也，九黎亂德”，又云“其後三苗復九黎之德”，則“九黎”、“三苗”非一物也。顓頊誅九黎謂之“遏絕苗民”，於鄭義為不愜。\CJKunderwave{楚語}言顓頊命重黎，解為\CJKunderline{帝堯}命羲和，於孔說又未允，不知二者誰得經意也。 \par}

{\noindent\zhuan\zihao{6}\fzbyks 傳“言堯”至“名焉”。正義曰：此經二句說\CJKunderline{帝堯}之德事也，而其言不順。文在“苗民”之下，故傳以為“堯監苗民之見怨，則又增修其德”，敦德以臨之，以德行其威罰，則民畏之而不敢為非。“明賢則德明人”者，若凡人雖欲以德明賢者，不能照察。今堯德明賢者，則能以德明識賢人，故皆勸慕為善。明與上句相互，則“德威”者，凡人雖欲以德行威，不能威肅。今堯行威罰，則能以德威罰罪人，故人皆畏威服德也。 \par}

{\noindent\shu\zihao{5}\fzkt “乃命”至“推明”。正義曰：三苗亂德，民神雜擾。\CJKunderline{帝堯}既誅苗民，乃命重黎二氏，使絕天地相通，令民神不雜。於是天神無有下至地，地民無有上至天，言天神地民不相雜也。群后諸侯相與在下國，群臣皆以明明大道輔行常法，鰥寡皆得其所，無有掩蓋之者。君帝\CJKunderline{帝堯}清審詳問下民所患，鰥寡皆有辭怨於苗民。言誅之合民意。堯視苗民見怨，則又增修其德。以德行威，則民畏之，不敢為非。以德明人，人皆勉力自修，使德明。言堯所行賞罰得其所也。 \par}

乃命三後，恤功於民。\CJKunderline{\CJKunderline{伯夷}}降典，折民惟刑。\CJKunderline{禹}平水土，主名山川。稷降播種，農殖嘉穀。\footnote{\CJKunderline{\CJKunderline{伯夷}}下典禮教民而斷以法。\CJKunderline{禹}治洪水，山川無名者主名之。后稷下教民播種，農畝生善谷。所謂堯命三君，憂功於民。○折，之設反，下同；馬、鄭、王皆音悊，馬云：“智也。”種音章用反。殖，承力反。斷，丁亂反，下同。}三後成功，惟殷於民。\footnote{各成其功，惟所以殷盛於民。言禮教備，衣食足。}士制百姓於刑之中，以教祇德。\footnote{言\CJKunderline{\CJKunderline{伯夷}}道民典禮，斷之以法。\CJKunderline{皋陶}作士，制百官於刑之中，助成道化，以教民為敬德。○祇，止而反。}


{\noindent\zhuan\zihao{6}\fzbyks 傳“\CJKunderline{\CJKunderline{伯夷}}”至“於民”。正義曰：\CJKunderline{\CJKunderline{伯夷}}與稷言“降”，\CJKunderline{禹}不言“降”，“降”可知降下也，從上而下於民也。\CJKunderwave{舜典}\CJKunderline{\CJKunderline{伯夷}}主禮典，“教民而斷以法”，即\CJKunderwave{論語}所謂“齊之以禮”也。山川與天地並生，民應先與作名。但\CJKunderline{禹}治水，萬事改新，古老既死，其名或滅。故當時無名者，\CJKunderline{禹}皆主名之。言此者，以見\CJKunderline{禹}治山川，為民於此耕稼故也。此三事者皆是為民，故傳既解三事，乃結上句，此即“所謂堯命三君，憂功於民”，憂欲與民施功也。此三事之次，當\CJKunderline{禹}功在先。先治水土，乃得種穀。民得穀食，乃能行禮。\CJKunderwave{管子}云：“衣食足，知榮辱。倉稟實，知禮節。”是言足食足衣然後行禮也。此經先言“\CJKunderline{\CJKunderline{伯夷}}”者，以民為國之本，禮是民之所急，將言制刑，先言用禮，刑禮相須，重禮，故先言之也。 \par}

{\noindent\zhuan\zihao{6}\fzbyks 傳“言伯”至“敬德”。正義曰：此經大意，言\CJKunderline{禹}、稷教民，使衣食充足。\CJKunderline{\CJKunderline{伯夷}}道民，使知禮節，有不從教者，乃以刑威之。故先言三君之功，乃說用刑之事。言\CJKunderline{禹}、稷教民稼穡，衣食既已充足。\CJKunderline{\CJKunderline{伯夷}}道民典禮，又能折之以法。禮法既行，乃使\CJKunderline{皋陶}作士，制百官於刑之中。令百官用刑，皆得中正，使不僣不濫，不輕不重，助成道化，以教民為敬德。言從\CJKunderline{\CJKunderline{伯夷}}之法，敬德行禮也。 \par}

{\noindent\shu\zihao{5}\fzkt “乃命”至“祇德”。正義曰：堯既誅苗民，乃命三君伯、夷、\CJKunderline{禹}、稷憂施功於民。使\CJKunderline{\CJKunderline{伯夷}}下禮典教民，折斷下民，惟以典法。\CJKunderline{伯禹}身平治水土，主名天下山川，其無名者皆與作名。后稷下教民布種，在於農畝種殖嘉。谷三君者各成其功，惟以殷盛於民，使民衣食充足。乃使士官制御百官之姓於刑之中正，以教民為敬德。言先以禮法化民，民既富而後教之，非苟欲刑殺也。 \par}

穆穆在上，明明在下，灼於四方，罔不惟德之勤。\footnote{堯躬行敬敬在上，三後之徒秉明德明君道於下，灼然彰著四方，故天下之士無不惟德之勤。}故乃明於刑之中，率乂於民棐彝。\footnote{天下皆勤立德，故乃能明於用刑之中正，循道以治於民，輔成常教。○治，直吏反。}


{\noindent\zhuan\zihao{6}\fzbyks 傳“堯躬”至“之勤”。正義曰：\CJKunderwave{釋訓}云：“穆穆，敬也。”“明明”重明，則“穆穆”重敬，當敬天敬民，在於上位也。“明明在下”，則是臣事，知是“三後之徒秉明德明君道於下”也。彰著於四方，四方皆法效之，故天下之士無不惟德之勤。 \par}

{\noindent\zhuan\zihao{6}\fzbyks 傳“天下”至“常教”。正義曰：刑者所以助教而不可專用，非是身有明德，則不能用刑。以天下之大，萬方之眾，必當盡能用刑，天下乃治。此美堯能使“天下皆勤立德，故乃能明於用刑之中正”，言天下皆能用刑，盡得中正，循治民之道以治於民，輔成常教。\CJKunderline{\CJKunderline{伯夷}}所典之禮，是常行之教也。 \par}

{\noindent\shu\zihao{5}\fzkt “穆穆”至“棐彝”。正義曰：言堯躬行敬敬之道在於上位，三後之徒躬秉明德明君道在於下，君臣敬明與德，灼然著於四方，故天下之事無不惟德之勤，悉皆勤行德矣。天下之士皆勤立德，故乃能明於用刑之中正，循大道以治於民，輔成常教。美堯君臣明德，能用刑得中以輔禮教。 \par}

典獄，非訖於威，惟訖於富。\footnote{言堯時主獄，有威有德有恕，非絕於威，惟絕於富。世治，貨賂不行。○賂，來故反。}敬忌，罔有擇言在身。\footnote{堯時典獄皆能敬其職，忌其過，故無有可擇之言在其身。}惟克天德，自作元命，配享在下。”\footnote{凡明於刑之中，無擇言在身，必是惟能天德，自為大命，配享天意，在於天下。}


{\noindent\zhuan\zihao{6}\fzbyks 傳“言堯”至“不行”。正義曰：堯時主獄之官，有威嚴，有德行，有恕心。有犯罪必罪之，是“有威”也。無罪則赦之，是“有德”也。有威有德有恕心，行之不受貨賂，是恕心也。“訖”是盡也，故傳以“訖”為絕。不可能使民不犯，非絕於威。能使不受貨賂，惟絕於富。言以恕心行之，世治則貨賂不行，故獄官無得富者。得凡明”至“天下”。正義曰：“惟克天德”，言能效天為德，當謂天德平均，獄官效天為平均。凡能明於刑之中正矣，又能使無可擇之言在身者，此人必是惟能為天平均之德，斷獄必平矣。“皇天無親，惟德是輔”，若能斷獄平均者，必壽長久大命。大命由己而來，是“自為大命”。“享”訓當也，是此人能配當天命，在於天之下。鄭云：“大命謂延期長久也。” \par}

{\noindent\shu\zihao{5}\fzkt “典獄”至“在下”。正義曰：堯時典獄之官，非能止絕於威，有犯必當行威，威刑不可止也。惟能止絕於富，受貨然後得富，無貨富自絕矣。言於時世治，貨賂不行。堯時典獄之官皆能敬其職事，忌其過失，無有可釋之言在於其身。天德平均，惟能為天之德。志性平均，自為長久大命。配當天意，在於天下。言堯德化之深，於時典獄之官皆能賢也。 \par}

王曰：“嗟！四方司政典獄，非爾惟作天牧？\footnote{主政典獄，謂諸侯也。非汝惟為天牧民乎？言任重是汝。○為，於偽反。任，而鴆反。重，輕重之重。}今爾何監，非時\CJKunderline{\CJKunderline{伯夷}}播刑之迪？\footnote{言當視是\CJKunderline{\CJKunderline{伯夷}}布刑之道而法之。}其今爾何懲？惟時苗民匪察於獄之麗。\footnote{其今汝何懲戒乎？所懲戒惟是苗民非察於獄之施刑，以取滅亡。○麗，力馳反。}罔擇吉人，觀於五刑之中，惟時庶威奪貨，\footnote{言苗民無肯選擇善人，使觀視五刑之中正，惟是眾為威虐者任之，以奪取人貨，所以為亂。}斷制五刑，以亂無辜。上帝不蠲，降咎於苗。\footnote{苗民任奪貨奸人，斷制五刑，以亂加無罪。天不潔其所為，故下咎罪。謂誅之。○蠲，吉緣反。咎，其九反。}苗民無辭於罰，乃絕厥世。”\footnote{言罪重，無以辭於天罰，故堯絕其世。申言之為至戒。}


{\noindent\zhuan\zihao{6}\fzbyks 傳“言當”至“法之”。正義曰：\CJKunderline{\CJKunderline{伯夷}}典禮，\CJKunderline{皋陶}主刑，刑禮相成以為治。不使視\CJKunderline{皋陶}而令視\CJKunderline{\CJKunderline{伯夷}}者，欲其先禮而後刑。道之以禮，禮不從乃刑之，則刑亦\CJKunderline{\CJKunderline{伯夷}}之所布，故令視\CJKunderline{\CJKunderline{伯夷}}布刑之道而法之。王肅云：“\CJKunderline{\CJKunderline{伯夷}}道之以禮，齊之以刑。” \par}

{\noindent\zhuan\zihao{6}\fzbyks 傳“其今”至“滅亡”。正義曰：上言“非時”，此言“惟時”，文異者，“非時”者言豈非是事也，“惟時”者言惟當是事也，雖文異而意同。“惟是苗民非察於獄之施刑以取滅亡”也，言其正謂察於獄之施刑不當於罪以取滅亡。 \par}

{\noindent\zhuan\zihao{6}\fzbyks 傳“苗民”至“誅之”。正義曰：“以亂加無罪”者，正謂以罪加無罪，是亂也。“蠲”訓絜也。“天不絜其所為”者，\CJKunderline{鄭玄}云：“天以苗民所行腥臊不絜，故下禍誅之。” \par}

{\noindent\shu\zihao{5}\fzkt “王曰”至“厥世”。正義曰：王呼諸侯戒之曰：“諮嗟！汝四方主政事典獄訟者諸侯之君等，非汝惟為天牧養民乎？”言汝等皆為天養民，言任重也。“受任既重，當觀古成敗，今汝何所監視乎？其所視者，非是\CJKunderline{\CJKunderline{伯夷}}布刑之道也”。言當效\CJKunderline{\CJKunderline{伯夷}}善布刑法，受令名也。“其今汝何所懲創乎？其所創者惟是苗民非察於獄之施刑乎？”言當創苗民施刑不當取滅亡也。“彼苗民之為政也，無肯選擇善人，使觀視於五刑之中正，惟是眾為威虐者任之，以奪取人之貨賂，任用此人，使斷制五刑，以亂加無罪之人。上天不絜其所為，故下咎惡於苗民。苗民無以辭於天罰，堯乃絕滅其世。汝等安得不懲創乎！” \par}

王曰：“嗚呼！念之哉！\footnote{念以\CJKunderline{\CJKunderline{伯夷}}為法，苗民為戒。}伯父、伯兄、仲叔、季弟、幼子、童孫，皆聽朕言，庶有格命。\footnote{皆石同姓，有父兄弟子孫。列者伯仲叔季，順少長也。舉同姓包異姓，言不殊也。聽從我言，庶幾有至命。○聽如字，又他經反。少，詩照反。長，丁丈反。}今爾罔不由慰日勤，爾罔或戒不勤。\footnote{今汝無不用安自居，日當勤之。汝無有徒念戒而不勤。○日，人實反，一音曰。}天齊於民，俾我一日，非終惟終在人。\footnote{天整齊於下民，使我為之，一日所行，非為天所終，惟為天所終，在人所行。○天齊於民，絕句。馬云：“齊，中也。”俾我，絕句。俾，必爾反，馬本作矜；矜，哀也。}爾尚敬逆天命，以奉我一人。雖畏勿畏，雖休勿休，\footnote{汝當庶幾敬逆天命，以奉我一人之戒。行事雖見畏，勿自謂可敬畏。雖見美，勿自謂有德美。}惟敬五刑，以成三德。一人有慶，兆民賴之，其寧惟永。”\footnote{先戒以勞謙之德，次教以惟敬五刑，所以成剛柔正直之三德也。天子有善，則兆民賴之，其乃安寧長久之道。}


{\noindent\zhuan\zihao{6}\fzbyks 傳“皆王”至“至命”。正義曰：此總告諸侯，不獨告同姓，知“舉同姓包異姓”也。“格”訓至也，言庶幾有至命。“至命”當謂至善之命，不知是何命也。\CJKunderline{鄭玄}云：“格，登也。登命謂壽考者。”傳雲“至命”亦謂壽考。 \par}

{\noindent\zhuan\zihao{6}\fzbyks 傳“今汝”至“不勤”。正義曰：“由”，用也。“慰”，安也。人之行事多有始無終，從而不改。王既殷勤教誨，恐其知而不行，或當曰欲勤行而中道倦怠，故以此言戒之。今汝等諸侯無不用安道以自居，言曰我當勤之。“安道”者，謂勤其職，是安之道。若不勤其職，是危之道也。 \par}

{\noindent\zhuan\zihao{6}\fzbyks 傳“天整”至“所行”。正義曰：“天整齊於下民”者，欲使之順道依理，以性命自終也。以民不能自治，故使我為之，使我為天子。我既受天委付，務欲稱天之心。墜失天命，是不為天所終。保全祿位，是為天所終。我一日所行善之與惡，非為天所終，惟為天所終，皆在人所行。王言已冀欲使為行稱天意也。 \par}

{\noindent\zhuan\zihao{6}\fzbyks 傳“汝當”至“德美”。正義曰：“逆”，迎也。上天授人為主，是下天命也。諸侯上輔天子，是逆天命也，言與天意相迎逆也。“汝當庶幾敬逆天命，以奉我一人之戒”，欲使之順天意而用己命。凡人被人畏，必當自謂己有可畏敬；被人譽，必自謂已實有德美。故戒之，汝等所行事，雖見畏，勿自謂可敬畏；雖見美，勿自謂有德美。教之令謙而不自恃也。 \par}

{\noindent\zhuan\zihao{6}\fzbyks 傳“先戒”至“之道”。正義曰：上句“雖畏勿畏，雖休勿休”，是“先戒以勞謙之德”也。“勞謙”，\CJKunderwave{易·謙卦}九三爻辭。謙則心勞，故云“勞謙”。天子有善，以善事教天下，則兆民蒙賴之。 \par}

{\noindent\shu\zihao{5}\fzkt “王曰”至“惟永”。正義曰：王言而嘆曰：“嗚呼！汝等諸侯其當念之哉！”念以\CJKunderline{\CJKunderline{伯夷}}為法，苗民為戒。既令念此法戒，又呼同姓諸侯曰：“伯父、伯兄、仲叔、季弟、幼子、童孫等，汝皆聽從我言，依行用之，庶幾有至善之命，命必長壽也。今汝等諸侯無不用安道以自居，曰我當勤之哉。汝已許自勤，即當必勤，汝無有徒念我戒，許欲自勤而身竟不勤。”戒使必自勤也。“上天欲整齊於下民，使我為之令，我為天子整齊下民也。我一日所行失其道，非為天所終。一日所行得其理，惟為天所終。此事皆在人所行”。言已當慎行以順天也。“我已冀欲順天，汝等當庶幾敬逆天命，以奉用我一人之戒。汝所行事，雖見畏，勿自謂可敬畏。雖見美，勿自謂有德美”。欲令其謙而勿自取也。“汝等惟當敬慎用此五刑，以成剛柔正直之三德，以輔我天子。我天子一人有善事，則億兆之民蒙賴之。若能如此，其乃安寧，惟久長之道也”。 \par}

王曰：“吁！來，有邦有土，告爾祥刑。\footnote{籲，嘆也。有國土諸侯，告汝以善用刑之道。○籲，況於反，馬作於；於，於也。}在今爾安百姓，何擇非人？何敬非刑？何度非及？\footnote{在今爾安百姓兆民之道，當何所擇？非惟吉人乎？當何所敬？非惟五刑乎？當何所度？非惟及世輕重所宜乎？○度，待洛反，注同，馬云：“造謀也。”}兩造具備，師聽五辭。\footnote{兩謂囚、證。造，至也。兩至具備，則眾獄官共聽其入五刑之辭。○造，七報反，注同。}五辭簡孚，正於五刑。\footnote{五辭簡核，信有罪驗，則正之於五刑。○核，幸革反。}


{\noindent\zhuan\zihao{6}\fzbyks 傳“在今”至“宜乎”。正義曰：“何度非及”，其言不明。以論刑事，而言度所及，知所度者，度及世之用刑輕重所宜。王肅云：“度，謀也。非當與主獄者謀慮刑事，度世輕重所宜也。” \par}

{\noindent\zhuan\zihao{6}\fzbyks 傳“兩謂”至“之辭”。正義曰：“兩”謂兩人，謂囚與證也。凡競獄必有兩人為敵，各言有辭理。或時兩皆須證，則囚之與證非徒兩人而已。兩人謂囚與證，不為兩敵至者，將斷其罪，必須得證，兩敵同時在官，不須待至；且兩人競理，或並皆為囚，各自須證，故以“兩”為囚與證也。兩至具備，謂囚證具足。各得其辭，乃據辭定罪。與眾獄官共聽其辭，觀其犯狀，斟酌入罪，或入墨劓，或入宮剕，故云“聽其入五刑之辭”也。 \par}

{\noindent\zhuan\zihao{6}\fzbyks 傳“五辭”至“五刑”。正義曰：既得囚證將入五刑之辭，更復簡練核實，知其信有罪狀，與刑書正同，則依刑書斷之，應墨者墨之，應殺者殺之。 \par}

五刑不簡，正於五罰。\footnote{不簡核，謂不應五刑。當正五罰，出金贖罪。○應，應對之應，下同。}五罰不服，正於五過。\footnote{不服，不應罰也。正於五過，從赦免。}五過之疵：惟官，惟反，惟內，惟貨，惟來。\footnote{五過之所病，或嘗同官位，或詐反囚辭，或內親用事，或行貨枉法，或舊相往來，皆病所在。○疵，才斯反。來，馬本作求，云：“有求，請賕也。”}


{\noindent\zhuan\zihao{6}\fzbyks 傳“不簡”至“贖罪”。正義曰：“不簡核”者謂覆審囚證之辭，不如簡核之狀。既囚與證辭不相符合，則是犯狀不定，謂“不應五刑”。不與五刑書同，獄官疑不能決，則當正之於五罰，令其出金贖罪。依準五刑，疑則從罰，故為“五罰”，即下文是也。今律：“疑罪各依所犯以贖。”論虛實之證，等是非之理，均或事涉疑似，旁無證見，或雖有證見，事非疑似，如此者皆為疑罪。 \par}

{\noindent\zhuan\zihao{6}\fzbyks 傳“不服”至“赦免”。正義曰：“不服，不應罰”者，欲令贖罪，而其人不服，獄官重加簡核，無復疑似之狀，本情非罪，不可強遣出金，如是者則正之於五過。雖事涉疑似有罪，乃是過失，過則可原，故從赦免。下文惟有“五刑”、“五罰”而無“五過”，亦稱“五”者，緣五罰為過，故謂之“五過”。五者之過，皆可原也。 \par}

{\noindent\zhuan\zihao{6}\fzbyks 傳“五過”至“所在”。正義曰：\CJKunderwave{釋詁}云：“疵，病也。”此五過之所病，皆謂獄吏故出入人罪，應刑不刑，應罰不罰，致之五過而赦免之，故指言“五過之疵”。於五刑五罰，不赦其罪，未有此病，故不言“五刑之疵”、“五罰之疵”。應刑而罰，亦是其病，於赦免言病，則赦刑從罰亦是病可知。損害王道，於政為病，故謂之“病”。“惟官”謂嘗同官位，與吏舊同僚也。“或詐反囚辭”，拒諱實情，不承服也。“或內親用事”，囚有親戒在官吏，或望其意而曲筆也。或行貨於吏，吏受財枉法也。或囚與吏舊相往來。此五事皆是病之所在。五事皆是枉法，但枉法多是為貨，故於“貨”言“枉”，餘皆枉可知。 \par}

其罪惟均，其審克之。\footnote{以病所在，出入人罪，使在五過，罪與犯法者同。其當清察，能使之不行。}五刑之疑有赦，五罰之疑有赦，其審克之。\footnote{刑疑赦從罰，罰疑赦從免。其當清察，能得其理。}簡孚有眾，惟貌有稽。\footnote{簡核誠信，有合眾心。惟察其貌，有所考合，重刑之至。}


{\noindent\zhuan\zihao{6}\fzbyks 傳“以病”至“不行”。正義曰：以五病所在，出入人罪，不罰不刑使得在於五過，妄赦免之，此獄吏之罪與犯法者同。諸侯國君清證審察，能使之不行，乃為善也。此以病所在，惟出人罪耳，而傳並言“入”者，有罪而妄出與無罪而妄入，獄吏之罪等，故以“出入”言之。今律：“故出入者與同罪。”而此是也。 \par}

{\noindent\zhuan\zihao{6}\fzbyks 傳“刑疑”至“其理”。正義曰：刑疑有赦，赦從罰也。罰疑有赦，赦從免也。上雲“五罰不服，正於五過”，即是免之也。不言五過之疑有赦者，知過則赦之，不得疑也。“其當清察，能得其理”，不使應刑妄得罰，應罰妄得免也。\CJKunderwave{舜典}雲“眚災肆赦”，\CJKunderwave{大禹謨}雲“宥過無大”，\CJKunderwave{易·解卦}象雲“君子以赦過宥罪”，\CJKunderwave{論語}雲“赦小過”，是過失之罪，皆當赦放，故知過即是赦之。\CJKunderline{鄭玄}云：“不言五過之疑有赦者，過不赦也。\CJKunderwave{禮記}云：‘凡執禁以齊眾者，不赦過。’”如鄭此言，五罰不服正於五過者，五過皆當罪之也。五刑之疑赦刑取贖，五罰疑者反使服刑，是刑疑而輸贖，罰疑而受刑，不疑而更輕，可疑而益重，事之顛倒一至此乎？謂之“祥刑”，豈當若是？然則“不赦過”者，復何所謂“執禁以齊眾”非謂之平常之過失也。人君故設禁約，將以齊整大眾，小事易犯，人必輕之，過犯悉皆赦之，眾人不可復禁，是故不赦小過，所以齊整眾人，令其不敢犯也。今律：“和合御藥誤不如本方，御幸舟舡誤不牢固，罪皆死。乏軍興者斬。”故失等皆是不赦過也。 \par}

{\noindent\zhuan\zihao{6}\fzbyks 傳“簡核”至“之”至“”。正義曰：“簡核誠信，有合眾心”，或皆以為可刑，或可以為赦，未得即斷之，惟當察其囚貌，更有所考合，考合復同，乃從眾議斷之，重刑之至也。“察其貌”者，即\CJKunderwave{周禮}五聽，辭聽、色聽、氣聽、耳聽、目聽也。\CJKunderline{鄭玄}以為辭聽“觀其出言，不直則煩”；色聽“觀其顏色，不直則赧然”；氣聽“觀其氣息，不直則喘”；耳聽“觀其聽聆，不直則惑”；目聽“觀其眸子，視不直則眊然”。是“察其貌，有所考合”也。 \par}

無簡不聽，具嚴天威。\footnote{無簡核誠信，不聽理具獄，皆當嚴敬天威，無輕用刑。}墨闢疑赦，其罰百鍰，閱實其罪。\footnote{刻其顙而涅之曰墨刑，疑則赦從罰。六兩曰鍰。鍰，黃鐵也。閱實其罪，使與罰名相當。○闢，婢亦反。鍰，徐戶關反，六兩也。鄭及\CJKunderwave{爾雅}同。\CJKunderwave{說文}云：“六鋝也。”“鋝，十一銖二十五分述之十三也。”馬同，又云：“賈逵說俗儒以鋝重六兩，\CJKunderwave{周官}劍重九鋝，俗儒近是。”閱音悅。顙，素黨反。涅，乃結反。}劓闢疑赦，其罰惟倍，閱實其罪。\footnote{截鼻曰劓。刑倍百為二百鍰。}剕闢疑赦，其罰倍差，閱實其罪。\footnote{刖足曰剕。倍差謂倍之又半，為五百鍰。}


{\noindent\zhuan\zihao{6}\fzbyks 傳“無簡”至“用刑”。正義曰：“無簡核誠信”者，謂簡核之，於罪無誠信效驗可簡核，即是無罪之人，當赦之。 \par}

{\noindent\zhuan\zihao{6}\fzbyks 傳“刻其”至“相當”。正義曰：五刑之名，見於經傳，\CJKunderline{唐}、\CJKunderline{虞}已來皆有之矣，未知上古起在何時也。漢文帝始除肉刑，其刻顙、截鼻、刖足、割勢皆法傳於先代，\CJKunderline{孔君}親見之。\CJKunderwave{說文}云：“顙，頟也。”“墨”一名黥。\CJKunderline{鄭玄}\CJKunderwave{周禮}注云：“墨，黥也。先刻其面，以墨窒之。”言刻頟為瘡，以墨塞瘡孔，令變色也。“六兩曰鍰”，蓋古語，存於當時，未必有明文也。\CJKunderwave{考工記}云，戈矛重三鋝。馬融云：“鋝，量名。當與\CJKunderwave{呂刑}鍰同。俗儒雲鋝六兩為一川，不知所出耳。”\CJKunderline{鄭玄}云：“鍰，稱輕重之名。今代、東萊稱，或以太半兩為鈞，十鈞為鍰，鍰重六兩太半兩。鍰、鋝似同也。或有存行之者，十鈞為鍰，二鍰四鈞而當一斤，然則鍰重六兩三分兩之二。\CJKunderwave{周禮}謂鍰為鋝。”如\CJKunderline{鄭玄}之言，一鍰之重六兩，多於孔、王所說，惟校十六銖爾。\CJKunderwave{舜典}云：“金作贖刑。”傳以金為黃金，此言“黃鐵”者，古者金銀銅鐵總號為“金”，今別之以為四名，此傳言“黃鐵”，\CJKunderwave{舜典}傳言“黃金”，皆是今之銅也。古人贖罪悉皆用銅，而傳或稱“黃金”，或言“黃鐵”，謂銅為金為鐵爾。“閱實其罪”，撿閱核實其所犯之罪，使與罰名相當，然後收取其贖。此既罪疑而取贖，疑罪不定，恐受贖參差，故五罰之下皆言“閱實其罪”，慮其不相當故也。 \par}

{\noindent\zhuan\zihao{6}\fzbyks 傳“刖足”至“百鍰”。正義曰：\CJKunderwave{釋詁}云：“剕，刖也。”李巡云：“斷足曰刖。”\CJKunderwave{說文}云：“刖，絕也。”是“刖”者斷絕之名，故“刖足曰剕”。贖劓倍墨，剕應倍劓，而云“倍差”，倍之又有差，則不啻一倍也。下句贖宮六百鍰，知倍之又半之為五百鍰也。截鼻重於黥頟，相校猶少。刖足重於截鼻，所校則多。刖足之罪，近於宮刑，故使贖剕不啻倍劓，而多少近於贖宮也。 \par}

宮闢疑赦，其罰六百鍰，閱實其罪。\footnote{宮，淫刑也。男子割勢，婦人幽閉，次死之刑。序五刑，先輕轉至重者，事之宜。}大辟疑赦，其罰千鍰，閱實其罪。\footnote{死刑也。五刑疑各入罰，不降相因，古之制也。}墨罰之屬千，劓罰之屬千，剕罰之屬五百，宮罰之屬三百，大辟之罰其屬二百。五刑之屬三千。\footnote{別言罰屬，合言刑屬，明刑罰同屬，互見其義以相備。○見，賢遍反。}

{\noindent\zhuan\zihao{6}\fzbyks 傳“宮淫”至“之宜”。正義曰：伏生\CJKunderwave{書傳}云：“男女不以義交者，其刑宮。”是宮刑為淫刑也。男子之陰名為勢，割去其勢，與椓去其陰，事亦同也。“婦人幽閉”，閉於宮使不得出也。本制宮刑，主為淫者，後人被此罪者，未必盡皆為淫。昭五年\CJKunderwave{左傳}楚子“以羊舌肸為司宮”，非坐淫也。漢除肉刑，除墨、劓、剕耳，宮刑猶在。近代反逆緣坐，男子十五已下不應死者皆宮之。大隋開皇之初，始除男子宮刑，婦人猶閉於宮。宮是次死之刑，宮於四刑為最重也。人犯輕刑者多，犯重刑者少，又以鍰數以倍相加，序五刑先輕後重，取事之宜。 \par}

{\noindent\zhuan\zihao{6}\fzbyks 傳“死刑”至“制也”。正義曰：\CJKunderwave{釋詁}云：“闢，罪也。”死是罪之大者，故謂死刑為“大辟”。經歷陳罰之鍰數，五刑之疑各自入罰。“不降相因”，不合死疑入宮，宮疑入剕者，是古之制也。所以然者，以其所犯疑不能決，故使贖之。次刑非其所犯，故不得降相因。 \par}

{\noindent\zhuan\zihao{6}\fzbyks 傳“別言”至“相備”。正義曰：此經歷言“二百”、“三百”、“五百”者，各是刑之條也。每於其條有犯者，實則刑之，疑則罰之，刑屬罰屬其數同也。別言罰屬，五者各言其數，合言刑屬，但總雲“三千”，明刑罰同其屬數，互見其義以相備也。經雲“大辟之罰，其屬二百”，文異於上四罰者，以“大辟”二字不可雲“大辟罰之屬”，故分為二句，以其二字足使成文。 \par}

{\noindent\shu\zihao{5}\fzkt “王曰”至“天威”。正義曰：凡與人言，必呼使來前。“籲”，嘆聲也。王嘆而呼諸侯曰：“吁！來，有邦國、有土地諸侯國君等，告汝以善用刑之道。在於今日，汝安百姓兆民之道，何所選擇？非惟選擇善人乎？何所敬慎？非惟敬慎五刑乎？何所謀度？非惟度及世之用刑輕重所宜乎”即教諸侯以斷獄之法。“凡斷獄者，必令囚之與證兩皆來至。囚證具備，取其言語，乃與眾獄官共聽其入五刑之辭。其五刑之辭簡核，信實有罪，則正之於五刑，以五刑之罪罪其身也。五刑之辭不如眾所簡核，不合入五刑，則正之於五罰。罰謂其取贖也。於五罰論之，又有辭不服，則正之於五過，過失可宥，則教宥之。從刑入罰，從罰入過。此五過之所病者，惟嘗同官位，惟詐反囚辭，惟內親用事，惟行貨枉法，惟舊相往來。以此五病出入人罪，其罪與犯法者均。其當清證審察，能使五者不行，乃為能耳。五刑之疑有赦，赦從罰也。五罰之疑有赦，赦從過也，過則赦之矣。其當清證審察使能之，勿使妄入人罪，妄得赦免。既得囚辭，簡核誠信，有合眾心。或記可刑，或皆可放，雖雲合罪，惟更審察其貌，有所考合”。謂貌又當罪，乃決斷之。“無簡不聽”者，謂雖似罪狀，無可簡核誠信合罪者，則不聽理其獄，當放赦之。皆當嚴敬天威，勿輕聽用刑也。 \par}

上下比罪，無僣亂辭，勿用不行，\footnote{上下比方其罪，無聽僣亂之辭以自疑，勿用折獄，不可行。○僣，子念反。}惟察惟法，其審克之。\footnote{惟當清察罪人之辭，附以法理，其當詳審能之。}上刑適輕，下服。\footnote{重刑有可以虧減則之輕，服下罪。}下刑適重，上服。輕重諸罰有權。\footnote{一人有二罪，則之重而輕並數。輕重諸刑罰各有權宜。○並，必政反。數，色住反。}刑罰世輕世重，惟齊非齊，有倫有要。\footnote{言刑罰隨世輕重也。刑新國用輕典，刑亂國用重典，刑平國用中典。凡刑所以齊非齊，各有倫理，有要善。}


{\noindent\zhuan\zihao{6}\fzbyks 傳“上下”至“可行”。正義曰：罪條雖有多數，犯者未必當條，當取故事並之，上下比方其罪之輕重。上比重罪，下比輕罪，觀其所犯當與誰同。獄官不可盡賢，其間或有阿曲，宜預防之。“僣”，不信也。獄官與囚等或作不信之辭，以惑亂在上，人君無得聽此僣亂之辭以自疑惑，勿即用此僣亂之辭以之斷獄，此僣亂之言不可行用也。 \par}

{\noindent\zhuan\zihao{6}\fzbyks 傳“一人”至“權宜”。正義曰：“一人有二罪，則之重而輕並數”者，謂若一人有二罪，則應兩罪俱治，今惟斷獄以重條，而輕者不更別數，與重並數為一。劉君以為“上刑適輕、下刑適重皆以為一人有二罪。上刑適輕者，若今律重罪應贖，輕罪應居作官當者，以居作官當為重，是為上刑適輕。下刑適重者，謂若二者俱是贓罪，罪從重科，輕贓亦備，是為而輕並數也”。知不然者，案經既言“下刑適重，上服”，則是重上服而已，何得為輕贓亦備？又今律雲“重罪應贖，輕重應居作官當者，以居作官當為重”者，此即是下刑適重之條，而以為上刑適輕之例，實為未允。且孔傳下經始雲“一人有二罪”，則上經所云非一人有二罪者也。劉君妄為其說，故今不從。 \par}

{\noindent\zhuan\zihao{6}\fzbyks 傳“言刑”至“要善”。正義曰：“刑罰隨世輕重”，言觀世而制刑也。“刑新國用輕典，刑亂國用重典，刑平國用中典”，\CJKunderwave{周禮·大司寇}文也。\CJKunderline{鄭玄}云：“新國者，新闢地立君之國。用輕法者，為其民未習於教也。平國，承平守成之國。用中典者，常行之法也。亂國，篡弒叛逆之國。用重典者，以其化惡，伐滅之也。” \par}

{\noindent\shu\zihao{5}\fzkt “上下”至“有要”。正義曰：此又述斷獄之法。將斷獄訟，當上下比方其罪之輕重，乃與獄官眾議斷之。其囚有僣亂之虛辭者，無得聽之，勿用此辭斷獄，此僣亂之辭，言不可行也。惟當清察罪人之辭，惟當附以法理，其當詳審使能之，勿使僣失為不能也。“上刑適輕”者，謂一人雖犯一罪，狀當輕重兩條，據重條之上有可以虧減者，則之輕條，服下罪也。“下刑適重”者，謂一人之身輕重二罪俱發，則以重罪而從上服，令之服上罪。或輕或重，諸所罪罰，皆有權宜，當臨時斟酌其狀，不得雷同加罪。刑罰有世輕世重，當視世所宜，權而行之。行罰者所以齊非齊者，有倫理，有要善。戒令審量之。 \par}

罰懲非死，人極於病。\footnote{刑罰所以懲過，非殺人，欲使惡人極於病苦，莫敢犯者。}非佞折獄，惟良折獄，罔非在中。\footnote{非口才可以斷獄，惟平良可以斷獄，無不在中正。}察辭於差，非從惟從。\footnote{察囚辭其難在於差錯，非從其偽辭，惟從其本情。}哀敬折獄，明啟刑書，胥佔，咸庶中正。\footnote{當憐下人之犯法，敬斷獄之害人，明開刑書，相與佔之，使刑當其罪，皆庶幾必得中正之道。○當，丁浪反。}其刑其罰，其審克之。\footnote{其所刑，其所罰，其當詳審能之，無失中正。}



{\noindent\zhuan\zihao{6}\fzbyks 傳“當憐”至“之道”。正義曰：\CJKunderwave{論語}云，陽膚為士師，曾子戒之云：“如得其情，則哀矜而勿喜。”是斷獄者於斷之時，當憐下民之犯法也。死者不可復生，斷者不可復續，當須敬慎斷獄之害人，勿得輕耳即決之。五刑之屬三千，皆著在刑書，使斷獄者依案用之，宜令斷獄諸官明開刑書，相與佔之，使刑書當其罪。令人之所犯，不必當條，須探測刑書之意，比附以斷其罪，若卜筮之佔然，故稱“佔”也。“皆庶幾必得中正之道”，令獄官同心思使中也。此言“明啟刑書”，而\CJKunderwave{左傳}雲“昔先王議事以制，不為刑辟”者，彼鑄刑書以宣示百姓，故云臨事制宜，不預明刑辟。人有犯罪，原其情之善惡，斷定其輕重，乃於刑書比附而罪之。故彼此各據其一，義不相違也。 \par}

獄成而孚，輸而孚。\footnote{斷獄成辭而信，當輸汝信於王。謂上其鞫劾文辭。○上，時掌反，下注同。鞫，九六反。劾，亥代反，\CJKunderwave{玉篇}胡得反。}其刑上備，有並兩刑。”\footnote{其斷刑文書上王府皆當備具，有並兩刑，亦具上之。}

{\noindent\zhuan\zihao{6}\fzbyks 傳“斷獄”至“文辭”。正義曰：“孚”，信也。“輸”，寫也。下“而”為汝也。斷獄成辭而得信實，當輸寫汝之信實以告於王，勿藏隱其情不告王也。曲必隱情，直則無隱，令其不隱情者，欲使之無阿曲也。漢世問罪謂之“鞫”，斷獄謂之“劾”，謂上其鞫劾文辭也。 \par}

{\noindent\zhuan\zihao{6}\fzbyks 傳“其斷”至“上之”。正義曰：“其斷刑文書上王府皆當備具”，若今曹司寫案申尚書省也。“有並兩刑”，謂人犯兩事，刑有上下，雖罪從重斷，有兩刑者，亦並具上之，使王知其事。王或時以下刑為重，改下為上，故並亦上之。 \par}

{\noindent\shu\zihao{5}\fzkt “罰懲”至“兩刑”。正義曰：言聖人之制刑罰，所以懲創罪過，非要使人死也，欲使惡人極於病苦，莫敢犯之而已。非口才辯佞之人可以斷獄，惟良善之人乃可以斷獄。言斷獄無非在其中正，佞人即不能然也。察囚之辭其難在於言辭差錯，斷獄者非從其偽辭，惟從其本情。斷獄之時，當哀憐之下民之犯法，敬慎斷獄之害人，勿得輕耳斷之，必令典獄諸官明開刑書，相與佔之，皆無幾得中正之道，其所刑罰，其當詳審能之，勿使失中。其斷獄成辭，得其信實，又當輸汝信實之狀而告於王。其斷刑文書上於王府，皆使備具，勿有疏漏。其囚若犯二事，罪雖從重，有並兩刑上之者，言有兩刑，亦具上之。恐獄官有所隱沒，故戒之。 \par}

王曰：“嗚呼！敬之哉！官伯、族姓，朕言多懼。\footnote{敬之哉，告使敬刑。官長，諸侯。族，同族。姓，異姓也。我言多可戒懼，以儆之。○儆音景。}朕敬於刑，有德惟刑。\footnote{我敬於刑，當使有德者惟典刑。}今天相民，作配在下，明清於單辭。\footnote{今天治民，人君為配天在下，當承天意，聽訟當清審單辭。單辭特難聽，故言之。○相如字，馬息亮反，助也。}


{\noindent\zhuan\zihao{6}\fzbyks 傳“敬之”至“儆之”。正義曰：此篇主多戒諸侯百官之長，故知“官長”即諸侯也。襄十二年\CJKunderwave{左傳}哭諸侯之例云：“異姓臨於外，同族於禰廟。”是相對則“族”為同姓，“姓”為異姓也。告之以“我言多可戒懼”者，以儆戒之也。下言民無善政，則天罰人主，是儆戒諸侯也。 \par}

{\noindent\zhuan\zihao{6}\fzbyks 傳“我敬”至“典刑”。正義曰：“當使有德者惟典刑”，言將選有德之人使為刑官，刑官不用無德之人也。 \par}

{\noindent\zhuan\zihao{6}\fzbyks 傳“今天”至“言之”。正義曰：傳以“相”為治，“今天治民”者，天有意治民，而天不自治，使人治之。人君為配天在下，當承天意治民，治之當使稱天心也。欲稱天心，聽獄當清審單辭。“單辭”謂一人獨言，未有與對之人。訟者多直已以曲彼，構辭以誣人，單辭特難聽，故言之也。\CJKunderline{孔子}美子路云：“片言可以折獄者，其由也與。”“片言”即“單辭”也。子路行直聞於天下，不肯自道已長，妄稱彼短，得其單辭即可以斷獄者，惟子路爾。凡人少能然，故難聽也。 \par}

民之亂，罔不中聽獄之兩辭，\footnote{民之所以治，由典獄之無不以中正聽獄之兩辭，兩辭棄虛從實，刑獄清則民治。○治，直吏反。}無或私家於獄之兩辭。\footnote{典獄無敢有受貨聽詐，成私家於獄之兩辭。}獄貨非寶，惟府辜功，報以庶尤。\footnote{受獄貨非家寶也，惟聚罪之事，其報則以眾人見罪。}


{\noindent\zhuan\zihao{6}\fzbyks 傳“民之”至“民治”。正義曰：“獄之兩辭”，謂兩人競理，一虛一實，實者枉屈，虛者得理，則此民之所以不得治也。民之所以得治者，由典獄之官其無不以有中正之心聽獄之兩辭，棄虛從實，實者得理，虛者受刑，虛者不敢更訟，則刑獄清而民治矣。\CJKunderline{孔子}稱“必也使無訟乎”，謂此也。 \par}

{\noindent\zhuan\zihao{6}\fzbyks 傳“典獄”至“兩辭”。正義曰：典獄知其虛，受其貨，而聽其詐。詐者虛而得理，獄官致富成私家，此民之所以亂也。故戒諸侯無使獄官成私家於獄之兩辭。 \par}

{\noindent\zhuan\zihao{6}\fzbyks 傳“受獄”至“見罪”。正義曰：“府”，聚也。“功”，事也。受獄貨非是家之寶也，惟最聚近罪之事爾。罪多必有惡報，其報則以眾人見罪也。眾人見罪者多，天必報以禍罰，故下句戒令畏天罰也。 \par}

永畏惟罰，非天不中，惟人在命。\footnote{當長畏懼惟為天所罰，非天道不中，惟人在教命使不中，不中則天罰之。}天罰不極，庶民罔有令政在於天下。”\footnote{天道罰不中，令眾民無有善政在於天下，由人主不中，將亦罰之。○令，力呈反。}

{\noindent\zhuan\zihao{6}\fzbyks 傳“當長”至“罰之”。正義曰：眾人見罪者多，天必報以禍罰，汝諸侯等當長畏懼為天所罰。天之罰人，非天道不得其中，惟人在其教命自使不申，教命不中，則天罰之。諸侯一，施教命於民者也，故戒以施教命中否也。 \par}

{\noindent\zhuan\zihao{6}\fzbyks 傳“天道”至“罰之”。正義曰：天道下罰，罰不中者，令使眾民無有善政在於天下，由人主不中。為人主不中，故無善政，天將亦罰人主。“人主”謂諸侯，此言戒諸侯也。 \par}

{\noindent\shu\zihao{5}\fzkt “王曰”至“天下”。正義曰：王嘆而呼諸侯曰：“嗚呼！刑罰事重，汝當敬之哉！謂諸侯官之長，此同族異姓等，我言多可戒懼。我敬於刑，當刑命有德者惟典刑事。今上天治民，命人君為天子，配天在於下，承天之意，為事甚重。其聽獄訟，當明白清審於獄之單辭。民之所以治者，由獄官無有不用中正聽訟之兩辭。由以中正之故，下民得治。汝獄官無有敢受貨賂，成私家於獄之兩辭。勿於獄之兩家受貨致富，治獄受貨非家寶也，惟是聚罪之事。”言汝身多違則不達，虛言戒行急惡，疏非虛論矣。“多聚罪則天報汝，以眾人見被尤怨而罰責之。汝當長畏惟天所罰，天罰汝者非是天道不中，惟人在於自作教命，使不中爾。教命不中，則天罰汝。天道罰不中也，若令眾民無有善政在於天下，則是人主不中，天亦將罰人主”。諸侯為民之主，故以天罰懼之。 \par}

王曰：“嗚呼！嗣孫，今往何監？非德於民之中？尚明聽之哉！\footnote{嗣孫，諸侯嗣世子孫，非一世。自今已往，當何監視？非當立德於民，為之中正乎？庶幾明聽我言而行之哉！}哲人惟刑，無疆之辭，屬於五極，咸中有慶。\footnote{言智人惟用刑，乃有無窮之善，辭名聞於後世。以其折獄屬五常之中正，皆中有善，所以然也。○屬音燭。}受王嘉師，監於茲祥刑。”\footnote{有邦有土受王之善眾而治之者，視於此善刑。欲其勤而法之，為無疆之辭。}


{\noindent\zhuan\zihao{6}\fzbyks 傳“言智”至“然也”。正義曰：“屬”謂屬著也。“極”，中也。“慶”，善也。“五常”謂仁義禮智信，人所常行之道也。言得有善辭，名聞於後世者，以其斷獄能屬著於五常之中正，皆得其理而法之有善，所以得然也。知“五”是五常者，以人所常行惟有五事，知是五常也。 \par}

{\noindent\shu\zihao{5}\fzkt “王曰”至“祥刑”。正義曰：戒之既終，王又言而嘆曰：“嗚呼！汝諸侯嗣世子孫等，從自今已往，當何所監視？非當視立德於民而為之中正乎？”言諸侯並嗣世惟當視此立德於民為之中正之事。“汝必視此，庶幾明聽我言而行之哉！有智之人惟能用刑，乃有無疆境之善辭。得有無疆善辭者，以其折獄能屬於五常之中正，皆中其理而法有善政故也。汝有邦有土之君，受王之善眾而治之，當視於此善刑。”從上已來舉善刑以告之，欲其勤而法之，使有無窮之美譽。 \par}

%%% Local Variables:
%%% mode: latex
%%% TeX-engine: xetex
%%% TeX-master: "../Main"
%%% End:

%% -*- coding: utf-8 -*-
%% Time-stamp: <Chen Wang: 2024-04-02 11:42:41>

% {\noindent \zhu \zihao{5} \fzbyks } -> 注 (△ ○)
% {\noindent \shu \zihao{5} \fzkt } -> 疏

\chapter{卷二十}


\section{文侯之命第三十}


平王錫晉文侯秬鬯圭瓚，\footnote{以圭為杓柄謂之圭瓚。○平王，馬無平字。錫，星曆反，馬本作賜。秬音巨，鬯，敕亮反。瓚，才但反。杓，上灼反。柄，彼病反。}作\CJKunderwave{文侯之命}。\footnote{所以名篇。幽王為犬戎所殺，平王立而東遷洛邑，晉文侯迎送安定之，故錫命焉。}

文侯之命\footnote{平王命為侯伯。}

{\noindent\zhuan\zihao{6}\fzbyks 傳“以圭”至“圭瓚”。正義曰：祭之初，酌鬱鬯之酒以灌屍。“圭瓚”者，酌鬱鬯之杓，杓下有槃，“瓚”即槃之名也；是以圭為杓之柄，故謂之“圭瓚”。\CJKunderwave{周禮·典瑞}云：“祼圭有瓚，以肆先王，以祼賓客。”鄭司農云：“於圭頭為器，可以挹鬯祼祭謂之瓚。以肆先王，灌先王祭也。”\CJKunderline{鄭玄}云：“肆，解牲體以祭，因以為名。”爵行曰“祼”。漢禮瓚槃大五升，口徑八寸，下有槃，口徑一尺。\CJKunderwave{詩}云：“瑟彼玉瓚，黃流在中。”\CJKunderwave{毛傳}云：“玉瓚，圭瓚也，黃金所以飾流鬯也。”鄭云：“黃流，秬鬯也。圭瓚之狀，以圭為柄，黃金為勺，青金為外，朱中央。”是說圭瓚之形狀也。\CJKunderwave{禮}無明文，而知其然者，\CJKunderwave{祭統}云：“君執圭瓚祼屍，大宗執璋瓚亞祼。”鄭云：“圭瓚、璋瓚，祼器也。以圭璋為柄，酌鬱鬯曰祼。”然則圭瓚、璋瓚惟柄以圭、璋為異，其瓚形則同。\CJKunderwave{考工記·玉人}云：“祼圭尺有二寸，有瓚，以祀廟。大璋、中璋九寸，邊璋七寸，厚寸，黃金勺，青金外，朱中，鼻寸。”鄭云：“鼻，勺流也，凡流皆為龍口也。三璋之勺，形如圭瓚。”是鄭以璋形如此，知圭瓚亦然。\CJKunderwave{毛傳}又云“九命然後錫以秬鬯圭瓚”，則晉文侯於時九命為東西大伯，故得受此賜也。“秬鬯”從經為傳，故此惟解“圭瓚”。 \par}

{\noindent\zhuan\zihao{6}\fzbyks 傳“所以”至“命焉”。正義曰：\CJKunderwave{周本紀}云，幽王嬖褒姒，褒姒生子伯服。幽王廢申後，並去太子，用褒姒為後，伯服為太子。申侯怒，乃與西夷犬戎共攻殺幽王。於是諸侯乃與申侯共立太子宜臼，是為平王。東徙於洛邑，避戎寇。隱六年\CJKunderwave{左傳}：“周桓公言於王曰：‘我周之東遷，晉鄭焉依。’”鄭語云：“晉文侯於是乎定天子。”是迎送安定之，故平王錫命焉。 \par}

{\noindent\zhuan\zihao{6}\fzbyks 傳“平王命為侯伯”。正義曰：“伯”，長也，諸侯之長謂之伯也。僖元年\CJKunderwave{左傳}云：“凡侯伯，救患、分災、討罪，禮也。”是與諸侯之長為“侯伯”。王肅云：“幽王既滅，平王東遷，晉文侯、鄭武公夾輔王室，晉為大國，功重，故平王命為侯伯。” \par}

{\noindent\shu\zihao{5}\fzkt “平王”至“之命”。正義曰：幽王嬖褒姒，廢申後，逐太子宜臼。宜臼奔申。申侯與犬戎既殺幽王，晉文侯與鄭武公迎宜臼立之，是為平王，遷於東都。平王乃以文侯為方伯，賜其秬鬯之酒，以圭瓚副焉，作策書命之。史錄其策書，作\CJKunderwave{文侯之命}。 \par}

王若曰：“父義和，\footnote{順其功而命之。文侯同姓，故稱曰父。義和，字也。稱父者非一人，故以字別之。○義和，馬云：“能以義和諸侯。”義本作誼。別，彼列反。}丕顯文武，克慎明德，\footnote{大明乎！\CJKunderline{文王}、\CJKunderline{武王}之道，能詳慎顯用有德。}昭升於上，敷聞在下，惟時上帝集厥命於\CJKunderline{文王}。\footnote{更述\CJKunderline{文王}所以王也。言\CJKunderline{文王}聖德明升於天，而布聞在下居。惟以是，故上天集成其王命，德流子孫。○聞音問。王，於況反。}亦惟先正，克左右昭事厥闢，\footnote{言君既聖明，亦惟先正官賢臣能，左右明事其君，所以然。○闢，必亦反。}越小大謀猷，罔不率從，肆先祖懷在位。\footnote{文王君聖臣良，於小大所謀道德，天下無不循從其化，故我後世先祖歸在王位。}


{\noindent\zhuan\zihao{6}\fzbyks 傳“順其”至“別之”。正義曰：\CJKunderwave{覲禮}說天子呼諸侯之義曰：“同姓大國則曰伯父，其異姓則曰伯舅，同姓小國則曰叔父，其異姓則曰叔舅。”\CJKunderline{鄭玄}\CJKunderwave{禮}注云：“稱之以父與舅，親親之辭。”晉文侯唐叔之後，與王同姓，故稱曰“父”。\CJKunderwave{曲禮}天子謂二伯為伯父伯舅。計文侯為侯伯，天子當呼為“伯父”，此不雲“伯”而直稱“父”者，尢親之也。\CJKunderwave{左傳}以文侯名仇，今呼曰“義和”，知是字也。天子於同姓諸侯皆呼為“父”，稱“父”者非一人，若不稱其字，無以知是文侯，故以字別之。\CJKunderline{鄭玄}讀“義”為“儀”，儀、仇皆訓匹也，故名仇，字儀。古人名字不可皆令相配，不必然也。 \par}

{\noindent\zhuan\zihao{6}\fzbyks 傳“\CJKunderline{文王}”至“王位”。正義曰：“後世先祖”謂文武之後，在今王之先祖，成康以至宣幽皆是也。“懷”，歸也。“歸在王位”者，王位是其所有也，若歸向家然，故稱“歸”也。 \par}

{\noindent\shu\zihao{5}\fzkt “王若”至“在位”。正義曰：平王順文侯之功，親之，敬而呼其字曰：“父義和。”既呼其字，乃告以上世之事：“大明乎！\CJKunderline{文王}、\CJKunderline{武王}之道，能詳順顯用有德之人以為大臣。\CJKunderline{文王}之為王也，聖德明升於天。”言其道至天也。“又布聞於在下”。言其德被民也。“惟以是，故上天成其大命於\CJKunderline{文王}，使之身為天子，澤流後世。文武聖明如此，亦惟先世長官之臣，能左右明事其君，君聖臣賢之故。於小大所謀道德，天下無有不循從其化，故我之先祖文武之後諸王，皆得歸在王位”。言先世聖王得賢臣之力，將說己無賢臣，故言此也。 \par}

嗚呼！閔予小子嗣，造天丕愆。\footnote{嘆而自痛傷也。言我小子而遭天大罪過，父死國敗，祖業隤隕。○予如字，又音與。愆，去虔反。隤，杜回反。隕，于敏反。}殄資澤於下民，侵戎我國家純。\footnote{言周邦喪亂，絕其資用惠澤於下民，侵兵傷我國及卿大夫之家，禍甚大。○殄，大見反。}即我御事，罔或耆壽俊在厥服，予則罔克。\footnote{所以遇禍，即我治事之臣，無有耆宿壽考俊德在其服位，我則材劣無能之致。}


{\noindent\zhuan\zihao{6}\fzbyks 傳“言周”至“甚大”。正義曰：此經所言，追敘幽王滅事。民不自治，立君以養之。民之資用，是王者佑助以得之。言周邦喪亂，不能撫佑下民，絕其資用惠澤於下民也。幽王之滅，由夷狄交侵，兵傷我國及卿大夫之家，其禍甚大。諸言“國家”者，皆謂國為“國家”，傳意欲見君臣俱被其害，故以“家”為卿大夫之家。王肅云：“遭天之大愆，謂幽王為犬戎所殺，殄絕其先祖之澤於下民。侵犯兵寇，傷我國家甚大，謂犬戎也。” \par}

{\noindent\zhuan\zihao{6}\fzbyks 傳“所以”至“之致”。正義曰：此經亦是追敘往事，言幽王所以遇禍者，即我周家治事之臣，無有耆宿壽考俊德之人在其服位，致使有犬戎之禍，亦是我材劣無能之致。幽王之時，平王被逐在外，國之興亡，非平王所知，言我無能之致者，引過歸己，自懼將來複然，故下句思得賢臣。 \par}

{\noindent\shu\zihao{5}\fzkt “嗚呼！”至“罔克”。正義曰：王又嘆而自傷：“嗚呼！疲病者，是我小子繼嗣先王之位，遭天大罪過於我周家，父死國敗，傾覆祖業，致使周邦喪亂，絕其資用惠澤於下民。”言下民資用盡，致使而王澤竭也。“西夷犬戎，侵兵傷我國及卿大夫之家，其禍亦甚大也。所以遇此禍者，即我治事之臣，無有耆宿壽考俊德之人在其服位，我則材弱無能之致”。自恨已弱不能致得賢臣，恐又不能自立也。 \par}

曰惟祖惟父，其伊恤朕躬。嗚呼！有績，予一人永綏在位。\footnote{王曰：“同姓諸侯在我惟祖惟父列者，其惟當憂念我身。嗚呼！能有成功，則我一人長安在王位。”言恃諸侯。}父義和，汝克昭乃顯祖，\footnote{重稱字，親之。不稱名，尊之。言汝能明汝顯祖唐叔之道，獎之。}汝肇刑文武，用會紹乃闢，追孝於前文人。\footnote{言汝今始法文武之道矣。當用是道合會繼汝君以善，使追孝於前文德之人。汝君，平王自謂也。繼先祖之志為孝。○闢，扶亦反。}汝多修，扞我於艱，若汝，予嘉。”\footnote{戰功曰多，言汝之功多，甚修矣。乃扞我於艱難，謂救周，誅犬戎，汝功我所善之。○扞，下旦反，注同。}


{\noindent\zhuan\zihao{6}\fzbyks 傳“王曰”至“諸侯”。正義曰：文侯是同姓諸侯，王言已未得文侯之時，常望同姓助己。王私為言曰：“同姓諸侯在我惟祖惟父列者，惟當憂念我身。”“伊”訓惟也。望得同姓之間有憂己者。以思惟未得，更嘆而為言：“嗚呼！同姓諸侯若有能助我有功，則我一人長得安在王位。”言己恃賴諸侯，思得其人，在後果得文侯。告文侯以此言，言己思文侯之功。 \par}

{\noindent\zhuan\zihao{6}\fzbyks 傳“重稱”至“獎之”。正義曰：天子之於諸侯，當稱“父舅”而已。既呼其“父”，又稱其字，所以別他人也。初則別於他人，重則可以已矣。重稱其字者，親之也。\CJKunderwave{禮}君父之前曰名，朋友之交曰字。是名重於字也。輕前人則斥其名，尊前人則避其重。故不稱其名，尊之也。不於上文作傳，於此言“尊之”者，就此“親之”，並解之也。“昭乃顯祖”，不知所斥，以晉之上世有功名者惟有唐叔耳，故知“明汝顯祖唐叔之道”。所以勸獎之，令其繼唐叔之業也。 \par}

{\noindent\zhuan\zihao{6}\fzbyks 傳“言汝”至“為孝”。正義曰：以其初有大功，終當不殞其業，故言“始法文武之道”。“當用是文武之道合會繼汝君以善”，令以功德佐汝君，使汝君繼前世，追行孝道於前世文德之人。“汝君”者，平王自謂也。先祖之志，在於平定天下，故子孫繼父祖之志為孝也。 \par}

{\noindent\zhuan\zihao{6}\fzbyks 傳“戰功”至“所善”。正義曰：“戰功曰多”者，\CJKunderwave{周禮·司勳}文。又云：“王功曰勳，國功曰功，民功曰庸，事功曰勞，治功曰力，戰功曰多。”彼有此六功也。言功多殊於他人，故云“汝之功多，甚修矣”。言其功修整，美其功之善也。文侯之功，在於誅犬戎，立平王，言“乃扞蔽我於艱難”，知“謂救周，誅犬戎”也。“若”訓如也，如汝之功，我所嘉也。王肅亦云：“如汝之功，我所善也。” \par}

{\noindent\shu\zihao{5}\fzkt “曰惟”至“予嘉”。正義曰：“王又言我以無能之致，私為言曰，同姓諸侯，惟我祖之列者，惟我父之列者，其惟當憂念我身。”又自傷嘆：“嗚呼！此諸侯等若有能助我有功，則我一人長安在王位。”言己無能，惟將賴諸侯也。又呼文侯字曰：“父義和，汝能明汝顯祖唐叔之道，汝始法文武之道，用是道合會繼汝君以善，追孝於前世文德之人。救周之國，汝功為多，甚修矣。乃能扞蔽我於艱難。”謂救周，誅犬戎也。“如汝之功，是我所善”。陳其前功，以勸勉之。 \par}

王曰：“父義和，其歸視爾師，寧爾邦。\footnote{遣令還晉國，其歸視汝眾，安汝國內上下。○令，力呈反。}用賚爾秬鬯一卣，\footnote{黑黍曰秬，釀以鬯草。不言圭瓚，可知。卣，中樽也。當以錫命告其始祖，故賜鬯。○賚，力代反。卣音酉，又音由。釀，女亮反。}彤弓一，彤矢百，盧弓一，盧矢百，\footnote{彤，赤。盧，黑也。諸侯有大功，賜弓矢，然後專征伐。彤弓以講德習射，藏示子孫。○彤，徒冬反。}馬四匹。\footnote{馬供武用。四匹曰乘。侯伯之賜無常，以功大小為度。○供音恭。}


{\noindent\zhuan\zihao{6}\fzbyks 傳“黑黍”至“賜鬯”。正義曰：\CJKunderwave{釋草}云：“秬，黑黍。”李巡曰：“黑黍一名秬。”\CJKunderwave{周禮}：“鬱人掌和鬱鬯，以實彝而陳之。”鄭云：“鬱，鬱金香草也。築鬱金煮之以和鬯酒。”鄭眾云：“鬱為草若蘭。”又有“鬯人掌供秬鬯”。\CJKunderline{鄭玄}云：“鬯釀秬為酒，芬香調暢於上下也。”如彼鄭說，釀黑黍之米為酒，築鬱金之草煮以和之。此傳言“釀以鬯草”，似用鬯草合釀。不同者終是以鬯和黍米之酒，或先或後言之耳。\CJKunderwave{詩}美宣王賜召穆公云：“釐爾圭瓚，秬鬯一卣，告於文人。”知賜秬鬯者必以圭瓚副焉。此不言“圭瓚”，明並賜之，可知也。“卣，中尊也”，\CJKunderwave{釋器}文。孫炎云：“樽，彝為上，罍為下，卣居中。”郭璞曰：“在罍彝之間。”即犧象壺著大山，等六尊是也。\CJKunderwave{周禮·司尊彝}云：“春祠夏禴，祼用雞彝鳥彝。秋嘗冬烝，祼用斝彝黃彝。”則祭時實鬯酒於彝。此用卣者，未祭則盛於卣，及祭則實於彝，此初賜未祭，故盛以卣也。\CJKunderwave{詩}稱“告於文人”，\CJKunderwave{毛傳}云：“文人，文德之人也。”\CJKunderline{鄭玄}云：“王賜召虎以鬯酒一尊，使以祭其宗廟，告其先祖諸有德美見記也。”然則得秬鬯之賜，當遍告宗廟，此傳惟言告始祖者，舉祖之尊者言之耳。 \par}

{\noindent\zhuan\zihao{6}\fzbyks 傳“彤赤”至“子孫”。正義曰：“彤”字從丹，“玈”字從玄，故“彤，赤。玈，黑”也。是“諸侯有大弓，賜弓矢，然後專征伐”，\CJKunderwave{禮記·王制}文也。\CJKunderwave{周禮}“司弓矢掌六弓”，其名王、弧、夾、庾、唐、大。\CJKunderline{鄭玄}云：“六者弓異體之名也。往體寡，來體多，曰王、弧。往體多，來體寡，曰夾、庾。往體來體若一，曰唐、大。”經又云：“唐弓大弓，以授學射者、使者、勞者。”鄭云：“學射者弓用中，後習強，弱則易也。使者、勞者，弓亦用中，遠近可也。勞者勤勞王事，若晉文侯受弓矢之賜者。”\CJKunderline{鄭玄}以此“彤弓”、“玈弓”為\CJKunderwave{周禮}“唐弓”、“大弓”。“唐”、“大”是弓強弱之名，“彤”、“玈”是弓赤黑之色，孔意亦當然也。此傳及\CJKunderwave{毛傳}皆雲“彤弓以講德習射”，用\CJKunderwave{周禮}為說也。唐弓大弓以授學射者，是習射也；授使者、勞者，是講德也。講論知其有德，乃賜之耳。襄八年\CJKunderwave{左傳}云，晉範宣子來聘，季武子賦\CJKunderwave{彤弓}。宣子曰：“城濮之役，我先君文公受彤弓於襄王，以為子孫藏。”杜預云：“藏之以示子孫。” \par}

{\noindent\zhuan\zihao{6}\fzbyks 傳“馬供”至“為度”。正義曰：六畜特以馬賜之者，為“馬供武用”故也。\CJKunderwave{周禮·校人}云：“乘馬一師四圉。”圉養一馬，是四匹曰乘，乘車必駕四馬故也。\CJKunderwave{司勳}云：“凡賞無常，輕重視功。”是“侯伯之賜無常，以功大小為度”。 \par}

父往哉！柔遠能邇，惠康小民，無荒寧。\footnote{父往歸國哉！懷柔遠人，必以文德。能柔遠者必能柔近，然後國安。安小人之道必以順，無荒廢人事而自安。}簡恤爾都，用成爾顯德。”\footnote{當簡閡汝所任，憂治汝都鄙之人，人和政治，則汝顯用有德之功成矣。不言鄙，由近以及遠。○核，戶革反。治，直吏反。}

{\noindent\zhuan\zihao{6}\fzbyks 傳“父往”至“自安”。正義曰：\CJKunderwave{論語}云：“遠人不服，則修文德以來之。”是“懷柔遠人，必以文德”也。能柔遠者必能柔近，遠近俱安，然後國安。“惠”，順也。“康”，安也。言順安小民者，安小民之道，必以順道安之，故言順安也。“順”者，順小民之心為其政也。\CJKunderwave{論語}雲“因民之所利而利之”，是順安也。 \par}

{\noindent\zhuan\zihao{6}\fzbyks 傳“當簡”至“及遠”。正義曰：“簡恤”者，共有“爾都”之文，當簡閡殃都內善人而任之，令以德憂治汝都鄙之人。人和政治，則汝顯用有德之功成矣。言用賢之名既成，國君之治亦成也。鄭云：“都，國都也。”“鄙”，邊邑也。言“都”不言“鄙”，由近以及遠也。 \par}

{\noindent\shu\zihao{5}\fzkt “王曰”至“顯德”。正義曰：王既陳其功，乃賚賜之。王曰：“父義和，其當歸汝晉國，視汝眾民，安汝國內上下。用賜汝秬鬯之酒一卣樽，歸以告祭汝之始祖。又賜汝彤弓一，彤矢百，玈弓一，玈矢百，馬四匹。父往歸國哉！必以文德安彼遠人。欲安遠，必能安近，是遠近乃得安耳。當以順道安汝之小民，無得荒廢人事以自安逸。簡閡汝所任之臣，憂治汝都鄙之人民，用成汝顯明之德。”戒使歸國善治民也。 \par}

\section{費誓第三十一}


魯侯伯禽宅曲阜，\footnote{治封之國居曲阜。}\footnote{伯禽，魯侯名。}徐、夷並興，東郊不開，\footnote{徐戎、淮夷並起，為寇於魯，故東郊不開。○開，舊讀皆作開，馬本作闢。}作\CJKunderwave{費誓}。\footnote{魯侯徵之於費地而誓眾也。諸侯之事而連帝王，\CJKunderline{孔子}序\CJKunderwave{書}以魯有治戎征討之備，秦有悔過自誓之戒，足為世法，故錄以備王事，猶\CJKunderwave{詩}錄商魯之\CJKunderwave{頌}。○費音秘。}

費誓\footnote{費，魯東郊之地名。}


{\noindent\zhuan\zihao{6}\fzbyks 傳“徐戎”至“不開”。正義曰：經稱“淮夷、徐戎”，序言“徐、夷”，略之也。此戎夷在魯之東，諸侯之制，於郊有門，恐其侵逼魯境，故東郊之門不開。 \par}

{\noindent\zhuan\zihao{6}\fzbyks 傳“費魯”至“地名”。正義曰：\CJKunderwave{甘誓}、\CJKunderwave{牧誓}皆至戰地而誓，知“費”非戰地者，“東郊不開”，則戎、夷去魯近矣。此誓令其治兵器，具糗糧，則是未出魯境，故知“費”是魯東郊地名，非戰處也。 \par}

{\noindent\shu\zihao{5}\fzkt “魯侯”至“費誓”。正義曰：魯侯伯禽於成王即政元年始就封於魯，居曲阜之地。於時徐州之戎、淮浦之夷並起，為寇於魯，東郊之門不敢開闢。魯侯時為方伯，率諸侯徵之，至費地而誓戒士眾。史錄其誓辭，作\CJKunderwave{費誓}。 \par}

公曰：“嗟！人無譁，聽命。\footnote{伯禽為方伯，監七百里內之諸侯，帥之以徵。嘆而敕之，使無喧譁，欲其靜聽誓命。○譁，戶瓜反。監，工銜反。}徂茲淮夷、徐戎並興。\footnote{今往徵此淮浦之夷、徐州之戎，並起為寇。此戎夷帝王所羈縻統敘，故錯居九州之內，秦始皇逐出之。}


{\noindent\zhuan\zihao{6}\fzbyks 傳“伯禽”至“誓命”。正義曰：\CJKunderwave{禮}諸侯不得專征伐，惟州牧於當州之內有不順者，得專征之。於時伯禽為方伯，監七百里內之諸侯，故得帥之以徵戎夷。\CJKunderwave{王制}云：“千里之外設方伯。”以八州八伯，是州別立一賢侯以為方伯，即\CJKunderwave{周禮·大宗伯}雲“八命作牧”是也。\CJKunderwave{禮記·明堂位}云：“封周公於曲阜，地方七百里。”孔意以周之大國不過百里，\CJKunderwave{禮記}雲“七百里”者，監此七百里內之諸侯，非以七百里地並封伯禽也。下雲“魯人三郊三遂”，指言“魯人”，明於時軍內更有諸侯之人，故知帥七百里內諸侯之人，以之共徵也。鄭云：“人謂軍之士眾及費地之民。”案下句令填塞坑阱，必使軍旁之民塞之，或當如鄭言也。 \par}

{\noindent\zhuan\zihao{6}\fzbyks 傳“今往”至“出之”。正義曰：\CJKunderwave{詩}美宣王命程伯休父，“率彼淮浦，省此徐土”。知“淮夷”是淮浦之夷，“徐戎”是徐州之戎也。四海之名，東方曰夷，西方曰戎，謂在九州之外。此徐州、淮浦，中夏之地而得有戎夷者，此戎夷帝王之所羈縻而統敘之，不以中國之法齊其風俗，故得雜錯居九州之內。此伯禽之時有淮浦者，淮浦之夷並起，\CJKunderwave{詩}美宣王命召穆公平淮夷，則戎夷之處中國久矣。漢時內地無戎夷者，秦始皇逐出之。始皇之崩至孔之初，惟可三四十年，古老尢在，及見其事，故孔得親知之也。王肅云：“皆紂時錯居中國。”經傳不說其事，無以知紂時來也。 \par}

善敹乃甲冑，敿乃幹，無敢不弔。\footnote{言當善簡汝甲鎧冑兜鍪，施汝楯紛，無敢不令至攻堅使可用。○敹，了雕反。敿，居表反。吊音的。鎧，苦代反。兜，丁侯反。鍪音矛。楯，常準反，又音允。紛，芳雲反。令，力呈反。}備乃弓矢，鍛乃戈矛，礪乃鋒刃，無敢不善。\footnote{備汝弓矢，弓調矢利。鍛練戈矛，磨礪鋒刃。皆使無敢不功善。○鍛，丁亂反。礪，力世反。練，來見反。}

{\noindent\zhuan\zihao{6}\fzbyks 傳“言當”至“可用”。正義曰：\CJKunderwave{世本}云：“杼作甲。”宋仲子云：“少康子杼也。”\CJKunderwave{說文}云：“胄兜鍪也。”“兜鍪”，首鎧也，經典皆言“甲冑”。秦世已來始有“鎧”、“兜鍪”之文。古之作甲用皮，秦漢已來用鐵，“鎧”、“鍪”二字皆從金，蓋用鐵為之，而因以作名也。甲冑為有善有惡，故令 簡，取其善者。鄭云：“ 謂穿徹之。”謂甲繩有斷絕，當使 理穿治之。“幹”是楯也，“敿乃幹”，必施功於楯，但楯無施功之處，惟系紛於楯，故以為“施汝楯紛”。紛如綬而小，繫於楯以持之，其以為飾。鄭雲“敿尢系也”，王肅雲“敿楯當有紛系持之”，是相傳為此說也。“吊”訓至也，無敢不令至極，攻堅使可用。鄭云：“至，尢善也。” \par}

{\noindent\zhuan\zihao{6}\fzbyks 傳“備汝”至“功善”。正義曰：“備”訓具也。每弓百矢，弓十矢千，使其數備足，令弓調矢利。案\CJKunderwave{毛傳}雲“五十矢為束”，或臨戰用五十矢為束。凡金為兵器，皆須鍛礪，有刃之兵，非獨戈矛而已。雲“鍛練戈矛，磨礪鋒刃，令其文互相通稱。諸侯兵器，皆使無敢不功善，令皆利快也”。 \par}

{\noindent\shu\zihao{5}\fzkt “公曰”至“不善”。正義曰：魯侯將徵徐戎，召集士眾，嘆而敕之。公曰：“嗟！在軍之人，無得喧譁，皆靜而聽我誓命。在往徵此淮浦之夷、徐州之戎，以其並起為寇故也。汝等善簡擇汝之甲冑，施汝楯紛，無敢不令至攻極堅。備汝弓矢，一弓百矢，令弓調矢利。鍛練汝之戈矛，磨礪汝之鋒刃，無敢不使皆善。”戒之使善，言不善將得罪也。 \par}

“今惟淫舍牿牛馬，\footnote{今軍人惟大放舍牿牢之牛馬，言軍所在必放牧也。○牿，工毒反。}杜乃擭，敜乃阱，無敢傷牿。牿之傷，汝則有常刑。\footnote{擭，捕獸機檻，當杜塞之。阱，穿地陷獸，當以土窒敜之。無敢令傷所以牿牢之牛馬。牛馬之傷，汝則有殘人畜之常刑。○杜，本又作度攵。擭，華化反，徐戶覆反。敜，徐乃協反，又乃結反。阱，在性反。檻，戶減反。窒，珍慄反。畜，許六反，又醜六反。}


{\noindent\zhuan\zihao{6}\fzbyks 傳“今軍”至“放牧”。正義曰：“淫”訓大也。\CJKunderwave{周禮}：“充人掌系祭祀之牲牷。祀五帝，則繫於牢，芻之三月。”\CJKunderline{鄭玄}云：“牢，閒也。”“校人掌王馬之政，天子十有二閒，馬六種。”然則養牛馬之處謂之牢閒，牢閒是周衛之名也。此言大舍牿牛馬，則是出之牢閒，牧於野澤，令其逐草而牧之。故謂此牢閒之牛馬為“牿牛馬”，而知“牿”即閒牢之謂也。故言“大放舍牿牢之牛馬”，言軍人所在，必須放牧此告軍旁之民也。既言牛馬在牿，遂以“牿”為牛馬之名，下雲“無敢傷牿”，謂傷牛馬，牿之傷謂牛馬傷也。\CJKunderline{鄭玄}以“牿為桎梏之梏，施梏於牛馬之腳，使不得走失”。 \par}

{\noindent\zhuan\zihao{6}\fzbyks 傳“擭捕”至“常刑”。正義曰：\CJKunderwave{周禮}：冥氏掌“為阱擭以攻猛獸”。知“阱”“擭”皆是捕獸之器也。擭以捕虎豹，穿地為深坑，又設機於上，防其躍而出也。阱以捕小獸，穿地為深坑，入必不能出，其上不設機也。阱以穿地為名，擭以得獸為名，擭亦設於阱中，但阱不設機為異耳。“杜”，塞之；“窒”，敜之；皆閉塞之義。使之填坑廢機，無敢令傷所放牿牢之牛馬。牛馬之傷，汝則有殘人畜之常刑。今律文：“施機槍作坑阱者，杖一百。梢閹之畜產者，償所減價。”王肅云：“杜，閉也。擭，所以捕禽獸機檻之屬。敜，塞也。阱，穿地為之，所以陷墮之。恐害牧牛馬，故使閉塞之。”\CJKunderline{鄭玄}云：“山林之田，春始穿地為阱，或設擭其中，以遮獸。擭，作㓵也。 \par}

{\noindent\shu\zihao{5}\fzkt “今惟”至“常刑”。正義曰：此戒軍旁之民也。今軍人惟欲大放舍牿牢之牛馬，令牧於野澤杜。汝捕獸之擭，塞汝陷獸之阱，無敢令傷所放牿牢之牛馬。牛馬之傷，汝則有殘害人畜之常刑。 \par}

馬牛其風，臣妾逋逃，勿敢越逐，\footnote{馬牛其有風佚，臣妾逋亡，勿敢棄越壘伍而求逐之。役人賤者男曰臣，女曰妾。○逋，布吳反。佚音逸。}祗復之，我商賚爾。\footnote{眾人其有得佚馬牛、逃臣妾，皆敬還復之，我則商度汝功，賜與汝。○商如字，徐音章。賚，力代反，徐音來。度，待洛反。}乃越逐，不復，汝則有常刑。\footnote{越逐為失伍，不還為攘盜，汝則有此常刑。○攘，如羊反。}無敢寇攘，逾垣牆，\footnote{軍人無敢暴劫人，逾越人垣牆，物有自來者，無敢取之。○垣音袁。}竊馬牛，誘臣妾，汝則有常刑。\footnote{軍人盜竊馬牛，誘偷奴婢，汝則有犯軍令之常刑。}甲戌，我惟徵徐戎。\footnote{誓後甲戌之日，我惟徵之。}峙乃糗糧，無敢不逮，汝則有大刑。\footnote{皆當儲峙汝糗糒之糧，使足食，無敢不相逮及，汝則有乏軍興之死刑。○峙，直裡反。\CJKunderwave{爾雅}云：“具也。”糗，去九反，一音昌紹反。糧音良。糒音備。}“魯人三郊三遂，峙乃楨幹。甲戌，我惟築，\footnote{總諸國之兵，而但稱魯人。峙具楨幹，道近也。題曰楨，旁曰幹。言“三郊三遂”，明東郊距守不峙，甲戌日當築攻敵壘距堙之屬。○楨，徐音貞。幹，工翰反。築，陟六反。守，手又反。堙音因。}


{\noindent\zhuan\zihao{6}\fzbyks 傳“馬牛”至“曰妾”。正義曰：僖四年\CJKunderwave{左傳}云：“唯是風馬牛不相及也。賈逵云：“風，放也。牝牡相誘謂之風。”然則馬牛風佚，因牝牡相逐而遂至放佚遠去也。“逋”亦逃也。軍士在軍，當各守部署，止則有壘壁，行則有隊伍，勿敢棄越壘伍而遠求逐之。\CJKunderwave{周禮}太宰“以九職任萬民”，“八曰臣妾，聚斂疏材”。僖十七年\CJKunderwave{左傳}云，晉惠公之妻“梁嬴孕，過期。卜招父與其子卜之。其子曰：‘將生一男一女。’招曰：‘然。男為人臣，女為人妾。’”是“役人賤者男曰臣，女曰妾”也。古人或以婦女從軍，故云“臣妾逋逃”也。 \par}

{\noindent\zhuan\zihao{6}\fzbyks 傳“皆當”至“死刑”。正義曰：“峙”，具也。預貯米粟謂之“儲峙”。鄭眾云：“糗，熬大豆及米也。”\CJKunderwave{說文}云：“糗，熬米麥也。”\CJKunderline{鄭玄}云：“糗，搗熬谷也。”謂熬米麥使熟，又搗之以為粉也。“糒”，乾飯也。“糗糒”是行軍之糧。皆當儲峙汝糗糒之糧，使在軍足食。“無敢不相逮及”，謂儲糧少，不及眾人，汝則有乏軍興之死刑。興軍征伐而有乏少，謂之“乏軍興”。今律：“乏軍興者斬。” \par}

{\noindent\zhuan\zihao{6}\fzbyks 傳“總諸”至“之屬”。正義曰：指言“魯人”，明更有他國之人。總諸國之兵，而但謂魯人。峙具楨幹，為道近故也。峙具楨幹以擬築之用。“題曰楨”，謂當牆兩端者也。“旁曰幹”，謂在牆兩邊者也。\CJKunderwave{釋詁}云：“楨，幹也。”舍人曰：“楨，正也，築牆所立兩木也。幹所以當牆兩邊障土者。”“三郊三遂”謂魯人三軍。\CJKunderwave{周禮·司徒}萬二千五百家為鄉。\CJKunderwave{司馬法}：“萬二千五百人為軍。”\CJKunderwave{小司徒}云：“凡起徒役，無過家一人。”是家出一人，一鄉為一軍。天子六軍，出自六鄉，則諸侯大國三軍，亦當出自三鄉也。\CJKunderwave{周禮}又云，萬二千五百家為遂。\CJKunderwave{遂人職}云：“以歲時稽其人民，簡其兵器，以起徵役。”則六遂亦當出六軍，鄉為正，遂為副耳。鄭眾云：“六遂之地在王國百里之外。”然則王國百里為郊，鄉在郊內，遂在郊外。\CJKunderwave{釋地}云：“邑外謂之郊。”孫炎曰：“邑，國都也。設百里之國，去國十里為郊。”則諸侯之制，亦當鄉在郊內，遂在郊外。此言“三郊三遂”者，“三郊”謂三鄉也。蓋使三鄉之民，分在四郊之內，三遂之民，分在四郊之外，鄉近於郊，故以郊言之。鄉遂之民，分在國之四面，當有四郊四遂，惟言“三郊三遂”者，明東郊令留守，不令峙楨幹也。上雲“甲戌，我惟徵徐戎”，此雲“甲戌，我惟築”，期以至日即築，當築攻敵之壘距堙之屬。\CJKunderwave{兵法}：“攻城築土為山，以闚望城內，謂之距堙。”襄六年\CJKunderwave{左傳}云：“晏弱城東陽，而遂圍萊。甲寅，堙之環城，傅於堞。”杜預云：“堞，女牆也。堙，土山也。周城為土山及女牆。”宣十五年\CJKunderwave{公羊傳}楚子圍宋，“使司馬子反乘堙而窺宋城，宋華元亦乘堙而出見之”。何休云：“堙，距堙，上城具也。”是攻敵城壘必有距堙，知“築”者，築距堙之屬也。 \par}

無敢不供，汝則有無餘刑，非殺。\footnote{峙具楨幹，無敢不供。不供，汝則有無餘之刑。刑者非一也，然亦非殺汝。○供音恭。}魯人三郊三遂，峙乃芻茭，無敢不多，汝則有大刑。”\footnote{郊遂多積芻茭，供軍牛馬。不多，汝則亦有乏軍興之大刑。○芻，初俱反，茭音交。}

{\noindent\zhuan\zihao{6}\fzbyks 傳“峙具”至“殺汝”。正義曰：上雲“無敢不逮”，此雲“無敢不供”，下雲“無敢不多”，文異者，糗糧難備，不得偏少，故云“無敢不逮”；楨幹易得，惟恐闕事，故云“無敢不供”；芻茭賤物，惟多為善，故云“無敢不多”，量事而為文也。“不供，汝則有無餘之刑”者，言刑者非一，謂閤家盡刑之。王肅云：“汝則有無餘刑，父母妻子同產皆坐之，無遺免之者，故謂無餘之刑。然入於罪隸，亦不殺之。”\CJKunderline{鄭玄}云：“無餘刑非殺者，謂盡奴其妻子，不遺其種類，在軍使給廝役，反則入於罪隸、舂槁，不殺之。”\CJKunderwave{周禮·司厲}云：“其奴，男子入於罪隸，女子入於舂槁。”\CJKunderline{鄭玄}云：“奴，從坐而沒入縣官者，男女同名。”鄭眾云：“輸於罪隸、舂人、槁人之官也。”然不供楨幹，雖是大罪，未應緣坐盡及家人，蓋亦權以脅之，使勿犯耳。○“芻茭”。正義曰：鄭云：“茭，乾芻也。” \par}

{\noindent\shu\zihao{5}\fzkt “馬牛”至“常刑”。正義曰：馬牛其有放佚，臣妾其有逋逃，汝無敢棄越壘伍而遠求逐之。其有得逸馬牛、逃臣妾，皆敬還復之，歸於本主，我則商度汝功，賞賜汝。汝若棄越壘伍，遠求逐馬牛臣妾，及有得馬牛臣妾不肯敬還復歸本主者，汝則有常刑。” \par}

\section{秦誓第三十二}


秦穆公伐鄭，\footnote{遣三帥帥師往伐之。○事見魯僖公三十三年。三帥，謂孟明視、西乞術、白乙丙。帥，色類反，下注同。}晉襄公帥師敗諸崤，\footnote{崤，晉要塞也。以其不假道，伐而敗之，囚其三帥。○崤，戶交反。塞，悉代反。假，工下反。}還歸，作\CJKunderwave{秦誓}。\footnote{晉舍三帥，還歸秦，穆公悔過作誓。}

秦誓\footnote{貪鄭取敗，悔而自誓。}


{\noindent\zhuan\zihao{6}\fzbyks 傳“遣三”至“伐之”。正義曰：\CJKunderwave{左傳}僖三十年，晉文公與秦穆公圍鄭，鄭使燭之武說秦伯，秦伯竊與鄭人盟，使杞子、逢孫、揚孫戍之，乃還。三十二年，杞子自鄭使告於秦曰：“鄭人使我掌其北門之管，若潛師以，來國可得也。”穆公訪諸蹇叔，蹇叔曰：“不可。”公辭焉。召孟明、西乞、白乙，使出師伐鄭。是“遣三帥帥師往伐之”事也。序言“穆公伐鄭”，嫌似穆公親行，故辨之耳。 \par}

{\noindent\zhuan\zihao{6}\fzbyks 傳“崤晉”至“三帥”。正義曰：杜預云：“殽在弘農澠池縣西。”築城守道謂之“塞”，言其要塞盜賊之路也。崤山險阨，是晉之要道關塞也。從秦向鄭，路經晉之南境，於南河之南崤關而東適鄭。\CJKunderwave{禮}征伐朝聘，過人之國，必遣使假道。晉以秦不假道，故伐之。\CJKunderwave{左傳}僖三十二年，晉文公卒。三十三年，秦師及滑，鄭商人弦高將市於周，遇之，矯鄭伯之命以牛十二犒師。孟明曰：“鄭有備矣，不可冀也。攻之不克，圍之不繼，吾其還也。”滅滑而還。晉先軫請伐秦師。襄公在喪，墨縗絰。夏四月，敗秦師於殽，獲百里孟明視、西乞術、白乙丙以歸。是襄公親自帥師伐而敗之，囚其三帥也。\CJKunderwave{春秋}之例，君將不言“帥師”，舉其重者。此言“襄公帥師”，依實為文，非彼例也。又\CJKunderwave{春秋}經書此事云：“晉人及羑戎敗秦師於殽。”實是晉侯而書“晉人”者，杜預云：“晉侯諱背喪用兵，通以賤者告也。”是言晉人告魯，不言晉侯親行，而云大夫將兵。大夫賤，不合書名氏，故稱“人”也。直言敗秦師於殽，不言秦之將帥之名，亦諱背喪用兵，故言辭略也。 \par}

{\noindent\zhuan\zihao{6}\fzbyks 傳“晉舍”至“作誓”。正義曰：\CJKunderwave{左傳}又稱，晉文公之夫人文嬴，秦女也，請三帥曰：“彼實構吾二君，寡君若得而食之，不厭，君何辱討焉？使歸就戮於秦，以逞寡君之志，若何？”公許之。秦伯素服郊次，向師而哭曰：“孤違蹇叔，以辱二三子，孤之罪也。不替孟明，孤之過也。”是晉舍三帥而得還，秦穆公於是悔過作誓。序言“還歸”，謂三帥還也，嫌穆公身還，故辨之。\CJKunderwave{公羊傳}說此事云：“四馬只輪無反者。”\CJKunderwave{左傳}稱秦伯“向師而哭”，則師亦少有還者， \par}

{\noindent\shu\zihao{5}\fzkt “秦穆”至“秦誓”。正義曰：秦穆公使孟明視、西乞術、白乙丙三帥帥師伐鄭，未至鄭而還。晉襄公帥師敗之於崤山，囚其三帥。後晉舍三帥，得還歸於秦。秦穆公自悔己過，誓戒群臣。史錄其誓辭，作\CJKunderwave{秦誓}。 \par}

公曰：“嗟！我士，聽無譁。\footnote{誓其群臣，通稱士也。}予誓告汝總言之首。\footnote{總言之本要。}古人有言曰：‘民訖自若，是多盤。’\footnote{言民之行己，盡用順道，是多樂。稱古人言，悔前不順忠臣。○樂音洛。}責人斯無難，惟受責俾如流，是惟艱哉！\footnote{人之有非，以義責之，此無難也。若己有非，惟受人責，即改之如水流下，是惟艱哉。○俾，必爾反，下同。}我心之憂，日月逾邁，若弗雲來。\footnote{言我心之憂，欲改過自新，如日月並行過，如不復雲來，雖欲改悔，恐死及之，無所益。○復，扶又反。}


{\noindent\zhuan\zihao{6}\fzbyks 傳“誓其”至“稱士”。正義曰：“士”者，男子之大號，故群臣通稱之。鄭云：“誓其群臣，下及萬民，獨雲士者，舉中言之。” \par}

{\noindent\zhuan\zihao{6}\fzbyks 傳“言民”至“忠臣”。正義曰：“訖”，盡也。“自”，用。“若”，順。“盤”，樂也。盡用順道則有福，有福則身樂，故云“是多樂”也。“稱古人言”者，悔前不用古人之言，不順忠臣之謀故也。昔漢明帝問東平王劉蒼云：“在家何者為樂？”對曰：“為善最樂。”是其用順道則多樂。 \par}

{\noindent\zhuan\zihao{6}\fzbyks 傳“言我”至“所益”。正義曰：“逾”，益。“邁”，行也。“員”即“雲”也。言日月益為疾行，並皆過去，如似不復雲來。畏其去而不復來，夜而不復明，言己年老，前途稍近，雖欲改悔，恐死及之，不得修改，身無所益也。王肅云：“年已衰老，恐命將終，日月遂往，若不雲來，將不復見日月，雖欲改過，無所及益。自恨改過遲晚，深自咎責之辭。” \par}

{\noindent\shu\zihao{5}\fzkt “公曰”至“雲來”。正義曰：穆公自悔伐鄭，召集群臣而告之。公曰：“諮嗟！我之朝廷之士，聽我告於汝，無得喧譁。我誓告汝眾言之首，誥汝以言中之最要者。古人有言曰：‘民之行己，盡用順道。是多樂。’言順善事，則身大樂也。見他有非理，以義責之，此無難也。惟己有非理，受人之責，即能改之，使如水之流下，此事是惟難哉！”言己已往之前不受人言，故自悔也。“今我心憂，欲自改過自新，但日月益為疾行，如似不復雲來，恐己老死不得改悔也”。 \par}

惟古之謀人，則曰未就予忌。\footnote{惟為我執古義之謀人，謂忠賢蹇叔等也，則曰未成我所欲，反忌之耳。○為，於偽反，下“為我謀”同。}惟今之謀人，姑將以為親。\footnote{惟指今事為我所謀之人，我且將以為親而用之。悔前違古從今，以取破敗。}

{\noindent\shu\zihao{5}\fzkt “惟古”至“為親”。正義曰：此穆公自說己之前過。我欲伐鄭之時，群臣共為謀計，惟為我執古義之謀人，我則曰未成我之所欲，反猜忌之。惟指今事為我所謀之人，我且將以為親己而用之。悔前違古從今，自取破敗也。其“古之謀人”，當謂忠賢之臣若蹇叔之等。“今之謀人”，勸穆公使伐鄭者，蓋謂杞子之類，國內亦當有此人。 \par}

“雖則云然，尚猷詢茲黃髮，則罔所愆。\footnote{言前雖則有云然之過，今我庶幾以道謀此黃髮賢老，則行事無所過矣。}番番良士，旅力既愆，我尚有之。\footnote{勇武番番之良士，雖眾力已過老，我今庶幾欲有此人而用之。○番音波。}仡仡勇夫，射御不違，我尚不欲。\footnote{仡仡壯勇之夫，雖射御不違，我庶幾不欲用。自悔之至。○仡，許乞反。}


{\noindent\shu\zihao{5}\fzkt “雖則”至“不欲”。正義曰：言我前事雖則有云然之過，我今庶幾以道謀此黃髮賢老，受用其言，則行事無所過也。番番然勇武之善士，雖眾力既過老，而謀計深長，我庶幾欲有此人而用之。仡仡然壯勇之夫，雖射御不有違失，而智慮淺近，我庶幾不欲用之。自悔往前用壯勇之計失也。 \par}

惟截截善諞言，俾君子易辭，我皇多有之，昧昧我思之。\footnote{惟察察便巧善為辨佞之言，使君子迴心易辭，我前多有之，以我昧昧思之不明故也。○截，才節反。馬云：“辭語截削省要也。”諞音辨，徐敷連反，又甫淺反，馬本作偏，云：“少也，辭約損明，大辨佞之人。”易，羊石反。昧音妹。}如有一介臣，斷斷猗，無他伎，其心休休焉，其如有容。\footnote{如有束脩一介臣，斷斷猗然專一之臣，雖無他伎藝，其心休休焉樂善，其如是，則能有所容。言將任之。○介音界，馬本作界，云：“一介，耿介，一心端愨者。”字又作個，音工佐反。斷，丁亂反，又音短。猗，於綺反，又於宜反。技，其綺反，本亦作伎。他，本亦作它，吐何反。樂音洛。}


{\noindent\zhuan\zihao{6}\fzbyks 傳“惟察”至“故也”。正義曰：“截截”猶“察察”，明辯便巧之意。“諞”猶辯也，由其便巧善為辯佞之言，使君子聽之迴心易辭。“皇”訓大也，我前大多有之，謂杞子之等，及在國從己之人。以我昧昧而暗，思之不明，故有此輩在我側也。 \par}

{\noindent\zhuan\zihao{6}\fzbyks 傳“如有”至“任之”。正義曰：孔注\CJKunderwave{論語}，以“束脩”為“束帶修飾”，此亦當然。“一介”謂一心耿介。“斷斷”，守善之貌。“休休”，好善之意。如有束帶修飾，一心耿介，斷斷然守善猗然專一之臣，雖復無他技藝，休休焉好樂善道，其心行如是，則能有所含容。言得此人將任用之。“猗”者，足句之辭，不為義也。\CJKunderwave{禮記·太學}引此作“斷斷兮”，“猗”是“兮”之類，\CJKunderwave{詩}雲“河水清且漣漪”是也。王肅云：“一介，耿介，一心端愨，斷斷守善之貌。無他技能，徒守善而已。休休，好善之貌。其如是，人能有所容忍小過，寬則得眾。穆公疾技巧多端，故思斷斷無他技者。” \par}

{\noindent\shu\zihao{5}\fzkt “惟截截”至“有容”。正義曰：惟察察然便巧善為辯佞之言，能使君子迴心易辭。我前大多有之，昧昧然我思之不明故也。如有一心耿介之臣，斷斷守善猗然，雖無他技藝，而其心樂善休休焉，其如是，則能有所含容。如此者，我將任用之。悔前用巧佞之人，今將任寬容善士也。 \par}

人之有技，若己有之。人之彥聖，其心好之，不啻若自其口出，是能容之。\footnote{人之有技，若己有之樂，善之至也。人之美聖，其心好之，不啻如自其口出，心好之至也。是人必能容之。○好，呼報反。啻，失豉反。}以保我子孫黎民，亦職有利哉！\footnote{用此好技聖之人，安我子孫眾人，亦主有利哉！言能興國。}

{\noindent\shu\zihao{5}\fzkt “人之”至“利哉”。正義曰：此說大賢之行也。大賢之人，見人之有技，如似己自有之。見人之有美善通聖者，其心愛好之，不啻如自其口出。愛彼美聖，口必稱揚而薦達之，其心愛之，又甚於口，言其愛之至也。是人於民必能含容之。用此愛好技聖之人，安我子孫眾民，則我子孫眾民亦主有利益哉！言其能興邦也。 \par}

人之有技，冒疾以惡之。人之彥聖，而違之，俾不達。\footnote{見人之有技藝，蔽冒疾害以惡之。人之美聖，而違背壅塞之，使不得上通。○冒，莫報反。惡，烏路反。背音佩。壅，於勇反。塞，先得反。}是不能容，以不能保我子孫黎民，亦曰殆哉！\footnote{冒疾之人，是不能容人用之，不能安我子孫眾人，亦曰危殆哉！○殆，唐在反。}


{\noindent\zhuan\zihao{6}\fzbyks 傳“見人”至“上通”。正義曰：傳以“冒”為覆冒之“冒”，謂蔽障掩蓋之也。“疾”謂疾惡之，謂憎疾患害之也。見人之美善通聖而違背之，不從其言；壅塞之，使不得上通，皆是佞人害賢之行也。 \par}

{\noindent\shu\zihao{5}\fzkt “人之”至“殆哉”。正義曰：此說大佞之行也。大佞之人，見人之有技，蔽冒疾害以惡之。見人之有美善通聖者，而違背壅塞之，使不達於在上。是人之不能含容人也。用此疾惡技聖之人，不能安我子孫眾民，則我子孫眾民亦曰危殆哉！言其必亂邦也。 \par}

邦之杌隉，曰由一人。\footnote{杌隉，不安，言危也。一人所任用，國之傾危，曰由所任不用賢。○杌，五骨反。隉，五結反，徐語折反。}邦之榮懷，亦尚一人之慶。”\footnote{國之光榮，為民所歸，亦庶幾其所任用賢之善也。穆公陳戒，背賢則危，用賢則榮，自誓改前過之意。}

{\noindent\shu\zihao{5}\fzkt “邦之”至“之慶”。正義曰：既言賢佞行異，又言用之安否。邦之杌隉，危而不安，曰由所任一人之不賢也；邦之光榮，為民所歸，亦庶幾所任一人之有慶也。言國家用賢則榮，背賢則危，穆公自誓將改前過，用賢人者也。人者也。 \par}

%%% Local Variables:
%%% mode: latex
%%% TeX-engine: xetex
%%% TeX-master: "../Main"
%%% End:



% Draft only
% \input{all_new}

%% 索隐
% \printindex
% \addcontentsline{toc}{part}{索引}

\end{document}

%%% Local Variables:
%%% mode: latex
%%% TeX-engine: xetex
%%% TeX-master: t
%%% End:
